% Options for packages loaded elsewhere
% Options for packages loaded elsewhere
\PassOptionsToPackage{unicode}{hyperref}
\PassOptionsToPackage{hyphens}{url}
\PassOptionsToPackage{dvipsnames,svgnames,x11names}{xcolor}
%
\documentclass[
  letterpaper,
  DIV=11,
  numbers=noendperiod]{scrreprt}
\usepackage{xcolor}
\usepackage{amsmath,amssymb}
\setcounter{secnumdepth}{-\maxdimen} % remove section numbering
\usepackage{iftex}
\ifPDFTeX
  \usepackage[T1]{fontenc}
  \usepackage[utf8]{inputenc}
  \usepackage{textcomp} % provide euro and other symbols
\else % if luatex or xetex
  \usepackage{unicode-math} % this also loads fontspec
  \defaultfontfeatures{Scale=MatchLowercase}
  \defaultfontfeatures[\rmfamily]{Ligatures=TeX,Scale=1}
\fi
\usepackage{lmodern}
\ifPDFTeX\else
  % xetex/luatex font selection
\fi
% Use upquote if available, for straight quotes in verbatim environments
\IfFileExists{upquote.sty}{\usepackage{upquote}}{}
\IfFileExists{microtype.sty}{% use microtype if available
  \usepackage[]{microtype}
  \UseMicrotypeSet[protrusion]{basicmath} % disable protrusion for tt fonts
}{}
\makeatletter
\@ifundefined{KOMAClassName}{% if non-KOMA class
  \IfFileExists{parskip.sty}{%
    \usepackage{parskip}
  }{% else
    \setlength{\parindent}{0pt}
    \setlength{\parskip}{6pt plus 2pt minus 1pt}}
}{% if KOMA class
  \KOMAoptions{parskip=half}}
\makeatother
% Make \paragraph and \subparagraph free-standing
\makeatletter
\ifx\paragraph\undefined\else
  \let\oldparagraph\paragraph
  \renewcommand{\paragraph}{
    \@ifstar
      \xxxParagraphStar
      \xxxParagraphNoStar
  }
  \newcommand{\xxxParagraphStar}[1]{\oldparagraph*{#1}\mbox{}}
  \newcommand{\xxxParagraphNoStar}[1]{\oldparagraph{#1}\mbox{}}
\fi
\ifx\subparagraph\undefined\else
  \let\oldsubparagraph\subparagraph
  \renewcommand{\subparagraph}{
    \@ifstar
      \xxxSubParagraphStar
      \xxxSubParagraphNoStar
  }
  \newcommand{\xxxSubParagraphStar}[1]{\oldsubparagraph*{#1}\mbox{}}
  \newcommand{\xxxSubParagraphNoStar}[1]{\oldsubparagraph{#1}\mbox{}}
\fi
\makeatother

\usepackage{color}
\usepackage{fancyvrb}
\newcommand{\VerbBar}{|}
\newcommand{\VERB}{\Verb[commandchars=\\\{\}]}
\DefineVerbatimEnvironment{Highlighting}{Verbatim}{commandchars=\\\{\}}
% Add ',fontsize=\small' for more characters per line
\usepackage{framed}
\definecolor{shadecolor}{RGB}{241,243,245}
\newenvironment{Shaded}{\begin{snugshade}}{\end{snugshade}}
\newcommand{\AlertTok}[1]{\textcolor[rgb]{0.68,0.00,0.00}{#1}}
\newcommand{\AnnotationTok}[1]{\textcolor[rgb]{0.37,0.37,0.37}{#1}}
\newcommand{\AttributeTok}[1]{\textcolor[rgb]{0.40,0.45,0.13}{#1}}
\newcommand{\BaseNTok}[1]{\textcolor[rgb]{0.68,0.00,0.00}{#1}}
\newcommand{\BuiltInTok}[1]{\textcolor[rgb]{0.00,0.23,0.31}{#1}}
\newcommand{\CharTok}[1]{\textcolor[rgb]{0.13,0.47,0.30}{#1}}
\newcommand{\CommentTok}[1]{\textcolor[rgb]{0.37,0.37,0.37}{#1}}
\newcommand{\CommentVarTok}[1]{\textcolor[rgb]{0.37,0.37,0.37}{\textit{#1}}}
\newcommand{\ConstantTok}[1]{\textcolor[rgb]{0.56,0.35,0.01}{#1}}
\newcommand{\ControlFlowTok}[1]{\textcolor[rgb]{0.00,0.23,0.31}{\textbf{#1}}}
\newcommand{\DataTypeTok}[1]{\textcolor[rgb]{0.68,0.00,0.00}{#1}}
\newcommand{\DecValTok}[1]{\textcolor[rgb]{0.68,0.00,0.00}{#1}}
\newcommand{\DocumentationTok}[1]{\textcolor[rgb]{0.37,0.37,0.37}{\textit{#1}}}
\newcommand{\ErrorTok}[1]{\textcolor[rgb]{0.68,0.00,0.00}{#1}}
\newcommand{\ExtensionTok}[1]{\textcolor[rgb]{0.00,0.23,0.31}{#1}}
\newcommand{\FloatTok}[1]{\textcolor[rgb]{0.68,0.00,0.00}{#1}}
\newcommand{\FunctionTok}[1]{\textcolor[rgb]{0.28,0.35,0.67}{#1}}
\newcommand{\ImportTok}[1]{\textcolor[rgb]{0.00,0.46,0.62}{#1}}
\newcommand{\InformationTok}[1]{\textcolor[rgb]{0.37,0.37,0.37}{#1}}
\newcommand{\KeywordTok}[1]{\textcolor[rgb]{0.00,0.23,0.31}{\textbf{#1}}}
\newcommand{\NormalTok}[1]{\textcolor[rgb]{0.00,0.23,0.31}{#1}}
\newcommand{\OperatorTok}[1]{\textcolor[rgb]{0.37,0.37,0.37}{#1}}
\newcommand{\OtherTok}[1]{\textcolor[rgb]{0.00,0.23,0.31}{#1}}
\newcommand{\PreprocessorTok}[1]{\textcolor[rgb]{0.68,0.00,0.00}{#1}}
\newcommand{\RegionMarkerTok}[1]{\textcolor[rgb]{0.00,0.23,0.31}{#1}}
\newcommand{\SpecialCharTok}[1]{\textcolor[rgb]{0.37,0.37,0.37}{#1}}
\newcommand{\SpecialStringTok}[1]{\textcolor[rgb]{0.13,0.47,0.30}{#1}}
\newcommand{\StringTok}[1]{\textcolor[rgb]{0.13,0.47,0.30}{#1}}
\newcommand{\VariableTok}[1]{\textcolor[rgb]{0.07,0.07,0.07}{#1}}
\newcommand{\VerbatimStringTok}[1]{\textcolor[rgb]{0.13,0.47,0.30}{#1}}
\newcommand{\WarningTok}[1]{\textcolor[rgb]{0.37,0.37,0.37}{\textit{#1}}}

\usepackage{longtable,booktabs,array}
\usepackage{calc} % for calculating minipage widths
% Correct order of tables after \paragraph or \subparagraph
\usepackage{etoolbox}
\makeatletter
\patchcmd\longtable{\par}{\if@noskipsec\mbox{}\fi\par}{}{}
\makeatother
% Allow footnotes in longtable head/foot
\IfFileExists{footnotehyper.sty}{\usepackage{footnotehyper}}{\usepackage{footnote}}
\makesavenoteenv{longtable}
\usepackage{graphicx}
\makeatletter
\newsavebox\pandoc@box
\newcommand*\pandocbounded[1]{% scales image to fit in text height/width
  \sbox\pandoc@box{#1}%
  \Gscale@div\@tempa{\textheight}{\dimexpr\ht\pandoc@box+\dp\pandoc@box\relax}%
  \Gscale@div\@tempb{\linewidth}{\wd\pandoc@box}%
  \ifdim\@tempb\p@<\@tempa\p@\let\@tempa\@tempb\fi% select the smaller of both
  \ifdim\@tempa\p@<\p@\scalebox{\@tempa}{\usebox\pandoc@box}%
  \else\usebox{\pandoc@box}%
  \fi%
}
% Set default figure placement to htbp
\def\fps@figure{htbp}
\makeatother





\setlength{\emergencystretch}{3em} % prevent overfull lines

\providecommand{\tightlist}{%
  \setlength{\itemsep}{0pt}\setlength{\parskip}{0pt}}



 


\KOMAoption{captions}{tableheading}
\makeatletter
\@ifpackageloaded{bookmark}{}{\usepackage{bookmark}}
\makeatother
\makeatletter
\@ifpackageloaded{caption}{}{\usepackage{caption}}
\AtBeginDocument{%
\ifdefined\contentsname
  \renewcommand*\contentsname{Table of contents}
\else
  \newcommand\contentsname{Table of contents}
\fi
\ifdefined\listfigurename
  \renewcommand*\listfigurename{List of Figures}
\else
  \newcommand\listfigurename{List of Figures}
\fi
\ifdefined\listtablename
  \renewcommand*\listtablename{List of Tables}
\else
  \newcommand\listtablename{List of Tables}
\fi
\ifdefined\figurename
  \renewcommand*\figurename{Figure}
\else
  \newcommand\figurename{Figure}
\fi
\ifdefined\tablename
  \renewcommand*\tablename{Table}
\else
  \newcommand\tablename{Table}
\fi
}
\@ifpackageloaded{float}{}{\usepackage{float}}
\floatstyle{ruled}
\@ifundefined{c@chapter}{\newfloat{codelisting}{h}{lop}}{\newfloat{codelisting}{h}{lop}[chapter]}
\floatname{codelisting}{Listing}
\newcommand*\listoflistings{\listof{codelisting}{List of Listings}}
\makeatother
\makeatletter
\makeatother
\makeatletter
\@ifpackageloaded{caption}{}{\usepackage{caption}}
\@ifpackageloaded{subcaption}{}{\usepackage{subcaption}}
\makeatother
\usepackage{bookmark}
\IfFileExists{xurl.sty}{\usepackage{xurl}}{} % add URL line breaks if available
\urlstyle{same}
\hypersetup{
  pdftitle={The Little Book of The Art of Computer Programming Solutions},
  pdfauthor={Duc-Tam Nguyen},
  colorlinks=true,
  linkcolor={blue},
  filecolor={Maroon},
  citecolor={Blue},
  urlcolor={Blue},
  pdfcreator={LaTeX via pandoc}}


\title{The Little Book of The Art of Computer Programming Solutions}
\usepackage{etoolbox}
\makeatletter
\providecommand{\subtitle}[1]{% add subtitle to \maketitle
  \apptocmd{\@title}{\par {\large #1 \par}}{}{}
}
\makeatother
\subtitle{Version 0.1.0}
\author{Duc-Tam Nguyen}
\date{2025-09-13}
\begin{document}
\maketitle

\renewcommand*\contentsname{Table of contents}
{
\hypersetup{linkcolor=}
\setcounter{tocdepth}{2}
\tableofcontents
}

\bookmarksetup{startatroot}

\chapter{}\label{section}

\bookmarksetup{startatroot}

\chapter{1.1 Algorithms}\label{algorithms}

\subsection{1 {[}10{]}}\label{section-1}

Minimal in-place algorithm (5 replacements)

\begin{verbatim}
t ← a
a ← b
b ← c
c ← d
d ← t
\end{verbatim}

Walkthrough. Let the initial values be \texttt{a0,\ b0,\ c0,\ d0}.

\begin{enumerate}
\def\labelenumi{\arabic{enumi}.}
\tightlist
\item
  \texttt{t\ ←\ a} → \texttt{t\ =\ a0}
\item
  \texttt{a\ ←\ b} → \texttt{a\ =\ b0}
\item
  \texttt{b\ ←\ c} → \texttt{b\ =\ c0}
\item
  \texttt{c\ ←\ d} → \texttt{c\ =\ d0}
\item
  \texttt{d\ ←\ t} → \texttt{d\ =\ a0}
\end{enumerate}

Final state is \texttt{(a,\ b,\ c,\ d)\ =\ (b0,\ c0,\ d0,\ a0)}, exactly
the requested left rotation.

Why 5 replacements is optimal

\begin{enumerate}
\def\labelenumi{\arabic{enumi}.}
\item
  All four variables change. In a 4-cycle \texttt{(a\ b\ c\ d)}, none of
  \texttt{a,\ b,\ c,\ d} keeps its original value, so each of the four
  must be assigned at least once → ≥ 4 replacements.
\item
  You must preserve one original value while overwriting. As soon as you
  assign, say, \texttt{a\ ←\ b}, the original \texttt{a0} is lost unless
  it was saved. But the final \texttt{d} must become \texttt{a0}.
  Therefore at least one extra ``save'' is necessary (e.g.,
  \texttt{t\ ←\ a}) → ≥ 4 + 1 = 5 replacements.
\end{enumerate}

Thus 5 is a lower bound, and the procedure above achieves it, so it's
minimal.

Generalization

For an n-cycle rotation \texttt{(x1\ x2\ …\ xn)}, the same argument
shows a lower bound of n + 1 replacements (n to set each variable once,
plus 1 to save a value), and the standard scheme

\begin{verbatim}
t ← x1;  x1 ← x2;  x2 ← x3;  …;  xn−1 ← xn;  xn ← t
\end{verbatim}

meets that bound.

\subsection{2 {[}15{]}}\label{section-2}

Consider any time we return to E1 after completing an earlier pass
through E1/E2/E3.

\begin{itemize}
\item
  On the previous pass through E1 we computed a remainder \(r\) with
  \(0\le r<n\). If \(r=0\) we would have halted in E2 and would not
  reach E1 again. Hence whenever we \emph{do} reach E1 again we must
  have \(r\neq 0\), i.e., \(0<r<n\). 
\item
  Then E3 set \(m\leftarrow n\) and \(n\leftarrow r\). Therefore, at the
  start of the next E1, we have

  \[
  m_{\text{new}} = n_{\text{old}} \quad\text{and}\quad n_{\text{new}} = r_{\text{old}}\,,
  \]

  with \(0<r_{\text{old}}<n_{\text{old}}\). Hence
  \(m_{\text{new}} = n_{\text{old}} > r_{\text{old}} = n_{\text{new}}\),
  i.e., \(m>n\) strictly.
\end{itemize}

This shows that from the second visit to E1 onward, the strict
inequality \(m>n\) always holds.

For the first visit to E1, no such guarantee exists: the inputs might
have \(m\le n\) (indeed \(m<n\) is allowed), so the statement allows the
exception ``except possibly the first time.'' (If one prepends the
optional ``ensure \(m\ge n\)'' swap step suggested in the text, the
property would also hold the first time.)

\subsection{3 {[}20{]}}\label{section-3}

Algorithm F (Euclid, without ``trivial replacements'')

Goal. Given positive integers \texttt{m} and \texttt{n}, compute
\texttt{gcd(m,\ n)}.

\begin{itemize}
\tightlist
\item
  F1. {[}Remainder m / n{]} Divide \texttt{m} by \texttt{n} and replace
  \texttt{m} with the remainder. (So \texttt{0\ ≤\ m\ \textless{}\ n}.)
\item
  F2. {[}Is it zero?{]} If \texttt{m\ =\ 0}, halt with answer
  \texttt{n}.
\item
  F3. {[}Remainder n / m{]} Divide \texttt{n} by \texttt{m} and replace
  \texttt{n} with the remainder. (So \texttt{0\ ≤\ n\ \textless{}\ m}.)
\item
  F4. {[}Is it zero?{]} If \texttt{n\ =\ 0}, halt with answer
  \texttt{m}; otherwise go to F1.
\end{itemize}

This version never does ``\texttt{m\ ←\ n,\ n\ ←\ r}''; we overwrite the
dividend directly with the remainder each time and alternate roles.

Why this is correct

\begin{enumerate}
\def\labelenumi{\arabic{enumi})}
\tightlist
\item
  Invariant (gcd is preserved)
\end{enumerate}

At every step, \texttt{gcd(m,\ n)\ =\ gcd(n,\ m\ mod\ n)} and
\texttt{gcd(m,\ n)\ =\ gcd(m,\ n\ mod\ m)}. That's the fundamental
property behind Euclid's method; F1/F3 are exactly these replacements,
so the gcd never changes. (This is the same reasoning used to justify
Algorithm E; see §1.1.
(\href{https://stackoverflow.com/questions/768948/how-does-division-work-in-mix}{PagePlace}))

\begin{enumerate}
\def\labelenumi{\arabic{enumi})}
\setcounter{enumi}{1}
\tightlist
\item
  Termination
\end{enumerate}

Each remainder is strictly smaller than its divisor. After F1 we have
\texttt{m\ \textless{}\ n}; after F3 we have \texttt{n\ \textless{}\ m}.
Thus the maximum of the pair \texttt{\{m,\ n\}} strictly decreases on
every pass through F1--F4. Since the maximum is a non-negative integer,
it can decrease only finitely many times, so the algorithm halts.

\begin{enumerate}
\def\labelenumi{\arabic{enumi})}
\setcounter{enumi}{2}
\tightlist
\item
  Upon halting, the answer is correct
\end{enumerate}

If we halt at F2, then \texttt{m\ =\ 0}, so
\texttt{gcd(m,\ n)\ =\ gcd(0,\ n)\ =\ n}. If we halt at F4, then
\texttt{n\ =\ 0}, so \texttt{gcd(m,\ n)\ =\ m}. Combined with the
invariant, this equals the original \texttt{gcd}.

Worked example

Compute \texttt{gcd(2166,\ 6099)} using F.

\begin{itemize}
\tightlist
\item
  F1: \texttt{m\ ←\ 2166\ mod\ 6099\ =\ 2166} (since 2166 \textless{}
  6099) F2: \texttt{m\ ≠\ 0} → continue
\item
  F3: \texttt{n\ ←\ 6099\ mod\ 2166\ =\ 1767} F4: \texttt{n\ ≠\ 0} →
  loop
\item
  F1: \texttt{m\ ←\ 2166\ mod\ 1767\ =\ 399} F2: \texttt{m\ ≠\ 0}
\item
  F3: \texttt{n\ ←\ 1767\ mod\ 399\ =\ 171} F4: \texttt{n\ ≠\ 0}
\item
  F1: \texttt{m\ ←\ 399\ mod\ 171\ =\ 57} F2: \texttt{m\ ≠\ 0}
\item
  F3: \texttt{n\ ←\ 171\ mod\ 57\ =\ 0} F4: \texttt{n\ =\ 0} → halt with
  answer \texttt{m\ =\ 57}.
\end{itemize}

\subsection{4 {[}16{]}}\label{section-4}

We repeatedly divide and keep the remainder; the gcd is the last
non-zero remainder.

\begin{enumerate}
\def\labelenumi{\arabic{enumi}.}
\tightlist
\item
  \(6099 = 2\cdot 2166 + 1767\) → \(r_1=1767\)
\item
  \(2166 = 1\cdot 1767 + 399\) → \(r_2=399\)
\item
  \(1767 = 4\cdot 399 + 171\) → \(r_3=171\)
\item
  \(399 = 2\cdot 171 + 57\) → \(r_4=57\)
\item
  \(171 = 3\cdot 57 + 0\)
\end{enumerate}

The last non-zero remainder is \(57\), so

\[
\gcd(2166,6099)=57.
\]

Quick checks

\begin{itemize}
\tightlist
\item
  Divisibility: \(2166/57=38\) and \(6099/57=107\) are integers.
\item
  No larger common factor: Euclid's process ensures optimality.
\end{itemize}

\subsection{5 {[}12{]}}\label{section-5}

The ``Procedure for Reading This Set of Books'' is not a genuine
algorithm because it breaks at least three of Knuth's five criteria:

\begin{itemize}
\tightlist
\item
  Finiteness: it loops back to step 3, so it need not terminate.
\item
  Definiteness: steps like ``Are you tired?'' or ``Is it interesting?''
  are subjective, not precise.
\item
  Effectiveness: operations like ``go to sleep'' or ``get your friends
  to purchase a copy'' are not basic, executable steps.
\end{itemize}

It also differs in format from Algorithm E: it uses plain numbers (not
E1, E2\ldots), prose instead of arrow notation (\(\leftarrow\)), no
bracketed summaries, and no explicit end marker.

Thus the ``procedure'' is more of a joke than a proper algorithm.

\subsection{6 {[}20{]}}\label{section-6}

For each \(m\in\{1,2,3,4,5\}\), count how many divisions (E1) occur
until termination.

\begin{itemize}
\tightlist
\item
  \(m=1\): \(1\bmod 5=1\) (1st E1); \((m,n)\leftarrow(5,1)\). Then
  \(5\bmod 1=0\) (2nd E1) ⇒ stop. Count \(=2\).
\item
  \(m=2\): \(2\bmod 5=2\) (1); \((5,2)\). \(5\bmod 2=1\) (2); \((2,1)\).
  \(2\bmod 1=0\) (3) ⇒ stop. Count \(=3\).
\item
  \(m=3\): \(3\bmod 5=3\) (1); \((5,3)\). \(5\bmod 3=2\) (2); \((3,2)\).
  \(3\bmod 2=1\) (3); \((2,1)\). \(2\bmod 1=0\) (4) ⇒ stop. Count
  \(=4\).
\item
  \(m=4\): \(4\bmod 5=4\) (1); \((5,4)\). \(5\bmod 4=1\) (2); \((4,1)\).
  \(4\bmod 1=0\) (3) ⇒ stop. Count \(=3\).
\item
  \(m=5\): \(5\bmod 5=0\) (1) ⇒ stop. Count \(=1\).
\end{itemize}

Sum of counts \(=2+3+4+3+1=13\).

\[
T_5=\frac{13}{5}=2.6.
\]

\subsection{7 {[}M21{]}}\label{m21}

Let \(A(m,n)\) be the number of times step E1 (``divide \(m\) by \(n\);
let \(r\) be the remainder'') is executed by Euclid's algorithm when it
is started on the ordered pair \((m,n)\) of positive integers.

Knuth defines \(T_n\) as the average of \(A(m,n)\) when \(n\) is fixed
and \(m\) is allowed to range over all positive integers; because only
the residue of \(m\bmod n\) matters after the first division, the
average is taken over \(m=1,2,\dots,n\) (uniformly). He notes that for
large \(n\), \(T_n\) is \(\asymp c\ln n\) with \(c=12\ln 2/\pi^2\). 

The exercise asks to define \(U_m\) analogously, with \(m\) fixed and
\(n\) ``ranging over all positive integers,'' and to relate \(U_m\) to
\(T_m\).

We will prove that \(U_m\) is well defined (the limit average exists and
is finite) and that

\[
\boxed{\,U_m = T_m + 1\,}.
\]

Two simple identities for \(A(m,n)\)

\begin{enumerate}
\def\labelenumi{(\Roman{enumi})}
\tightlist
\item
  Swap once when the first number is smaller. If \(m<n\), then the very
  first division in Euclid's algorithm leaves remainder \(r=m\); the
  next state is \((n,m)\). Thus
\end{enumerate}

\[
A(m,n) \;=\; 1 + A(n,m)\qquad(m\ne n).
\tag{1}
\]

\begin{enumerate}
\def\labelenumi{(\Roman{enumi})}
\setcounter{enumi}{1}
\tightlist
\item
  Periodicity in \(n \pmod m\) when \(n>m\). Write \(n=qm+r\) with
  \(0\le r<m\). The first step on \((m,n)\) gives remainder \(r=m\)
  (since \(n>m\)), so after one step we are at \((n,m)\). One more step
  divides \(n\) by \(m\) and leaves remainder \(r\); from then on we
  continue with \((m,r)\). Hence
\end{enumerate}

\[
A(m,n) \;=\;
\begin{cases}
2, & r=0,\\
2 + A(m,r), & 1\le r\le m-1.
\end{cases}
\tag{2}
\]

So for \(n>m\), \(A(m,n)\) depends only on \(n \bmod m\) (and equals a
constant \(2\) when \(m\mid n\)).

Existence (and value) of \(U_m\)

Define the Cesàro average

\[
U_m \;:=\; \lim_{N\to\infty}\frac{1}{N}\sum_{n=1}^N A(m,n),
\]

if the limit exists.

Partition \(1,\dots,N\) into complete blocks of length \(m\) plus a
final partial block. Using (II), each complete block contributes

\[
\sum_{r=0}^{m-1} A\bigl(m,\,qm+r\bigr)
= 2 + \sum_{r=1}^{m-1}\bigl(2 + A(m,r)\bigr)
= 2m + \sum_{r=1}^{m-1} A(m,r).
\]

Thus the block average is

\[
\frac{1}{m}\left(2m + \sum_{r=1}^{m-1} A(m,r)\right)
= 2 + \frac{1}{m}\sum_{r=1}^{m-1} A(m,r).
\tag{3}
\]

The final partial block contributes \(O(m)\), which is negligible after
division by \(N\). Therefore the limit defining \(U_m\) exists and
equals the block average:

\[
U_m \;=\; 2 + \frac{1}{m}\sum_{r=1}^{m-1} A(m,r).
\tag{4}
\]

Relation to \(T_m\)

By definition (see the paragraph introducing \(T_n\) and how the average
is taken over a complete set of residues), we have

\[
T_m \;=\; \frac{1}{m} \sum_{x=1}^{m} A(x,m)
= \frac{1}{m}\Bigl(A(m,m) + \sum_{x=1}^{m-1} A(x,m)\Bigr).
\]

Now use (I) with \(x<m\): \(A(x,m)=1+A(m,x)\). Also \(A(m,m)=1\). Hence

\[
T_m
= \frac{1}{m}\Bigl(1 + \sum_{x=1}^{m-1} (1 + A(m,x))\Bigr)
= \frac{1}{m}\Bigl(m + \sum_{x=1}^{m-1} A(m,x)\Bigr)
= 1 + \frac{1}{m}\sum_{x=1}^{m-1} A(m,x).
\tag{5}
\]

Comparing (4) with (5) immediately yields

\[
\boxed{\,U_m = T_m + 1\,}.
\]

Sanity check on a small example

Take \(m=5\). One can compute:

\[
A(5,1)=1,\;
A(5,2)=2,\;
A(5,3)=3,\;
A(5,4)=3.
\]

Then

\[
T_5 = 1 + \tfrac{1}{5}(1+2+3+3)=1+\tfrac{9}{5}= \tfrac{14}{5},
\qquad
U_5 = 2 + \tfrac{1}{5}(1+2+3+3) = 2+\tfrac{9}{5}=\tfrac{19}{5}
= T_5+1.
\]

(Manually checking a long run of \(n\) values shows the block average
equals \(19/5\), as predicted.)

Consequence

Since \(T_n \sim \frac{12\ln 2}{\pi^2}\ln n\) for large \(n\), we get

\[
U_m \sim \frac{12\ln 2}{\pi^2}\ln m + 1,
\]

so asymptotically \(U_m\) and \(T_m\) differ by exactly \(1\). 

That proves \(U_m\) is well defined and establishes its precise
relationship to \(T_m\).

\subsection{8 {[}M25{]}}\label{m25}

Design a set of \emph{string-rewriting} rules in Knuth's
\((T_j,\;s_j,\;a_j,\;b_j)\) format that, starting from the input string
\(a^m b^n\) (i.e., \(m\) copies of \texttt{a} followed by \(n\) copies
of \texttt{b}), halts with the string \(a^{\gcd(m,n)}\). The rules must
be expressed as a Markov algorithm in the tabular style used in §1.1.

Quick refresher: what \((T_j,s_j,a_j,b_j)\) means

A program state is a pair \((\sigma, j)\): current string \(\sigma\) and
current rule index \(j\). Each table row \(j\) has:

\begin{itemize}
\tightlist
\item
  \(T_j\) (``pattern''): the substring to search for,
\item
  \(s_j\) (``replacement''): the string to substitute when \(T_j\) is
  found,
\item
  \(b_j\): the next rule index if a replacement did occur,
\item
  \(a_j\): the next rule index if no \(T_j\) is found in \(\sigma\).
\end{itemize}

Operationally (Markov rules): if \(T_j\) occurs, replace the
\emph{leftmost} occurrence by \(s_j\) and jump to \(b_j\); otherwise
jump to \(a_j\). One special index \(N\) means ``halt.''
(\href{https://stackoverflow.com/questions/768948/how-does-division-work-in-mix}{Stack
Overflow})

A minimal working rule set

Use the alphabet \(\{a,b,c\}\) and five rules (indices \(0,1,2,3,4\));
let \(5\) be the halting index \(N\).

\begin{longtable}[]{@{}
  >{\raggedright\arraybackslash}p{(\linewidth - 10\tabcolsep) * \real{0.0105}}
  >{\raggedright\arraybackslash}p{(\linewidth - 10\tabcolsep) * \real{0.0737}}
  >{\raggedright\arraybackslash}p{(\linewidth - 10\tabcolsep) * \real{0.1368}}
  >{\raggedleft\arraybackslash}p{(\linewidth - 10\tabcolsep) * \real{0.0526}}
  >{\raggedleft\arraybackslash}p{(\linewidth - 10\tabcolsep) * \real{0.0526}}
  >{\raggedright\arraybackslash}p{(\linewidth - 10\tabcolsep) * \real{0.6737}}@{}}
\toprule\noalign{}
\begin{minipage}[b]{\linewidth}\raggedright
j
\end{minipage} & \begin{minipage}[b]{\linewidth}\raggedright
\(T_j\)
\end{minipage} & \begin{minipage}[b]{\linewidth}\raggedright
\(s_j\)
\end{minipage} & \begin{minipage}[b]{\linewidth}\raggedleft
\(b_j\)
\end{minipage} & \begin{minipage}[b]{\linewidth}\raggedleft
\(a_j\)
\end{minipage} & \begin{minipage}[b]{\linewidth}\raggedright
Intuition
\end{minipage} \\
\midrule\noalign{}
\endhead
\bottomrule\noalign{}
\endlastfoot
0 & \texttt{ab} & (empty) & 1 & 2 & Cancel one \texttt{a} with one
\texttt{b} (one subtraction step). \\
1 & (empty) & \texttt{c} (prepend) & 0 & 0 & After each cancellation,
count it by sticking a \texttt{c} on the left. \\
2 & \texttt{a} & \texttt{b} & 2 & 3 & Turn all remaining \texttt{a}'s
into \texttt{b}'s. \\
3 & \texttt{c} & \texttt{a} & 3 & 4 & Then turn all \texttt{c}'s into
\texttt{a}'s. \\
4 & \texttt{b} & \texttt{b} & 0 & 5 & If any \texttt{b}'s remain, repeat
another Euclid round; else halt. \\
\end{longtable}

Why this works (high-level):

\begin{itemize}
\tightlist
\item
  Rules 0↔1 simulate \emph{repeated subtraction}: each \texttt{ab}
  removed corresponds to kubtracting the smaller number from the larger
  once; the number of times you can do this is \(\min(m,n)\). We
  ``count'' those cancellations by accumulating \texttt{c}'s.
\item
  When no \texttt{ab} remains, the string looks like
  \texttt{c\^{}\{min(m,n)\}\ x\^{}\{\textbar{}m-n\textbar{}\}}, where
  \(x\) is whichever letter survived (\texttt{a} or \texttt{b}).
\item
  Rules 2 and 3 normalize the string to
  \texttt{a\^{}\{min(m,n)\}\ b\^{}\{\textbar{}m-n\textbar{}\}}, i.e.~we
  map the pair \((m,n)\) to \((\min(m,n),\,|m-n|)\).
\item
  Rule 4 loops if the second component is nonzero; otherwise halt. This
  is exactly the subtraction-based Euclidean step, so repeated rounds
  end with \texttt{a\^{}\{gcd(m,n)\}}. The superscript means ``repeat
  \texttt{a} that many times.'' ({[}Stack Overflow{]}{[}2{]})
\end{itemize}

Correctness argument (brief but complete)

Invariant after finishing rule 0↔1 loop (before rule 2): Let the current
pair be \((m,n)\). Repeatedly removing the leftmost \texttt{ab} (rule 0)
and then prepending a \texttt{c} (rule 1) exactly \(\min(m,n)\) times
yields a string of the form:

\[
c^{\min(m,n)}\; y^{|m-n|},
\]

where \(y\) is \texttt{a} if \(m>n\) and \texttt{b} if \(n\ge m\). This
encodes ``we subtracted the smaller from the larger \(\min(m,n)\)
times.''

Normalization (rules 2→3):

\begin{itemize}
\tightlist
\item
  Rule 2 turns any stray \texttt{a}'s into \texttt{b}'s, so only
  \texttt{c}'s and \texttt{b}'s remain.
\item
  Rule 3 turns each \texttt{c} into \texttt{a}. Thus the string becomes
  \texttt{a\^{}\{\textbackslash{}min(m,n)\}\ b\^{}\{\textbar{}m-n\textbar{}\}}
  in canonical order.
\end{itemize}

Euclid step and termination (rule 4):

\begin{itemize}
\tightlist
\item
  If any \texttt{b} remains, rule 4 jumps back to rule 0, so the new
  pair is \((m',n')=(\min(m,n),|m-n|)\).
\item
  Repeating transforms \((m,n)\) via the subtraction form of Euclid's
  algorithm until \(n'=0\); at that point rule 4 falls through to \(5\)
  (halt), with the string \texttt{a\^{}\{gcd(m,n)\}}.
  (\href{https://stackoverflow.com/questions/768948/how-does-division-work-in-mix}{Stack
  Overflow})
\end{itemize}

Worked example

Take \(m=6\), \(n=10\); start with \texttt{aaaaaabbbbbbbbbb}.

\begin{itemize}
\tightlist
\item
  Round 1 (to exhaustion of rule 0): remove \texttt{ab} six times while
  accumulating \texttt{c}'s → \texttt{ccccccbbbb}. Normalize (rules 2→3)
  → already no \texttt{a}'s, then \texttt{c→a} ⇒ \texttt{aaaaaabbbb} =
  \texttt{a\^{}6\ b\^{}4}.
\item
  Round 2: cancel \texttt{ab} four times → \texttt{cccc\ a\ a} (i.e.,
  \texttt{ccc\ caa} but order doesn't matter until normalization) ⇒
  normalize ⇒ \texttt{a\^{}4\ b\^{}2}.
\item
  Round 3: cancel twice ⇒ normalize ⇒ \texttt{a\^{}2\ b\^{}0}.
\item
  Round 4: rule 4 sees ``no \texttt{b} left'' ⇒ halt with \texttt{aa} =
  \(a^{\gcd(6,10)}\).
\end{itemize}

Relation to Euclid's Algorithm E

Each full pass of rules 0→\ldots→4 implements the mapping

\[
(m,n)\;\mapsto\;(\min(m,n),\;|m-n|),
\]

i.e., Euclid's algorithm with the division step replaced by
\emph{repeated subtraction}. That's precisely the hint Knuth gives; the
string system is just a Markov-algorithm rendering of that idea.

Notes and common confusions

\begin{itemize}
\tightlist
\item
  The superscript in ``\(a^{\gcd(m,n)}\)'' means ``\texttt{a} repeated
  \(\gcd(m,n)\) times.'' For example, if \(\gcd(12,18)=6\), the output
  is \texttt{aaaaaa}.
\item
  In the rule table, \(b_j\) is the ``branch if a replacement
  happened,'' and \(a_j\) is the ``branch if not.'' This branching
  scheme, together with leftmost replacement, is standard for Markov
  algorithms and matches Knuth's §1.1 formalism.
\end{itemize}

\subsection{9 {[}M30{]}}\label{m30}

Let

\[
C_1=(Q_1, I_1, \Omega_1, f_1),\qquad
C_2=(Q_2, I_2, \Omega_2, f_2)
\]

be computational methods (state space \(Q\), initial/allowable inputs
\(I\subseteq Q\), halting/``output'' states \(\Omega\subseteq Q\), and a
(partial) next-state function \(f\)). A \emph{computation sequence} of
\(C_i\) is the usual chain \(x, f_i(x), f_i^{(2)}(x),\dots\) until (and
if) a state in \(\Omega_i\) is reached.

Intuitively, \(C_2\) may do more, smaller steps and may keep extra
bookkeeping in its states, yet it should faithfully realize \(C_1\).

We say that \(C_2\) represents \(C_1\) iff there exist total maps

\begin{itemize}
\tightlist
\item
  an encoding \(e:Q_1\to Q_2\) and
\item
  a decoding/abstraction \(\pi:Q_2\to Q_1\)
\end{itemize}

such that \(\pi\!\circ e=\mathrm{id}_{Q_1}\) (i.e., \(e\) embeds
\(C_1\)'s abstract state into \(C_2\)'s concrete state, and \(\pi\)
forgets the extra detail), and the following hold:

\begin{enumerate}
\def\labelenumi{\arabic{enumi}.}
\item
  Input compatibility. \(e(I_1)\subseteq I_2\).
\item
  Output compatibility. \(\pi(\Omega_2)\subseteq \Omega_1\) and for
  every \(\omega\in\Omega_1\) there exists some
  \(\hat \omega\in\Omega_2\) with \(\pi(\hat \omega)=\omega\).
  (Equivalently, \(\pi^{-1}(\Omega_1)\) contains only states from which
  \(C_2\) can reach \(\Omega_2\) without changing \(\pi\)'s value.)
\item
  Step correspondence (stuttering simulation). For every
  \(q_1\in Q_1\setminus \Omega_1\) and every \(q_2\in Q_2\) with
  \(\pi(q_2)=q_1\), if \(f_1(q_1)\) is defined and equals \(q_1'\), then
  there exists a finite \(t\ge 1\) and a path in \(C_2\)

  \[
  q_2 = r_0 \xrightarrow{f_2} r_1 \xrightarrow{f_2} \cdots \xrightarrow{f_2} r_t
  \]

  such that

  \[
  \pi(r_j)=q_1\quad(0\le j<t),\qquad \pi(r_t)=q_1'.
  \]

  (So \(C_2\) may ``stutter''---take several internal steps that do not
  change the abstract state---before realizing one abstract
  \(f_1\)-step.)
\item
  Termination preservation (refinement of whole runs). For every input
  \(x\in I_1\), if the \(C_1\) run from \(x\) halts in \(k\) steps with
  output \(\omega\in\Omega_1\), then the \(C_2\) run from \(e(x)\) halts
  after some \(T\ge k\) steps in a state \(\hat \omega\in\Omega_2\) with
  \(\pi(\hat \omega)=\omega\).
\end{enumerate}

This is the standard ``forward simulation with stuttering''
characterization used in refinement: \(C_2\) can take more
(finer-grained) steps and store more information, but when you
\emph{abstract} its state with \(\pi\) you see exactly the abstract
behavior of \(C_1\).

Why this captures the intent

\begin{itemize}
\tightlist
\item
  ``More steps'': Condition 3 allows multiple \(f_2\) steps for a single
  \(f_1\) step while keeping \(\pi\) unchanged until the abstract change
  happens.
\item
  ``More information'': The abstraction \(\pi\) can forget registers,
  temporary variables, program counters, etc.;
  \(\pi\!\circ e=\mathrm{id}\) ensures no abstract information is lost
  when we embed.
\item
  Correct inputs/outputs: Conditions 1--2 align the two machines at the
  beginning and end.
\item
  Whole computations: Condition 4 ensures that every terminating
  abstract computation is realized concretely with the same abstract
  result.
\end{itemize}

How it looks on the Euclid example

In §1.1, Algorithm E is formalized as a \((Q,I,\Omega,f)\) system with
states like \((m,n,0,1)\), \((m,n,r,2)\), \((n)\), etc., and a
next-state function that embodies steps E1--E3. A machine program
\(C_2\) has states including memory, registers, and a program counter.
Define:

\begin{itemize}
\tightlist
\item
  \(e\): put the abstract \(m,n,r,\ldots\) into designated
  registers/memory and set the PC to the code for the corresponding
  E-step.
\item
  \(\pi\): read back those designated registers and map \emph{any}
  machine state of that code block to the corresponding abstract state
  (ignoring temps, flags, scratch space, etc.).
\end{itemize}

A long division routine in \(C_2\) will execute many micro-steps while
\(\pi\) stays at the same abstract E1 state; upon finishing the
division, \(\pi\) jumps to the E2/E3 abstract state---exactly the
stuttering behavior required by (3). When \(n=0\), the program branches
to a return sequence; \(\pi\) maps those return-sequence states into
\(C_1\)'s halting set, establishing (2)--(4).

\bookmarksetup{startatroot}

\chapter{1.2.1 Mathematical induction}\label{mathematical-induction}

\subsection{1 {[}5{]}}\label{section-7}

\begin{enumerate}
\def\labelenumi{\arabic{enumi})}
\tightlist
\item
  Weak induction starting at 0
\end{enumerate}

To prove \(P(n)\) for all \(n\ge 0\), use this two-part template:

\begin{enumerate}
\def\labelenumi{\arabic{enumi}.}
\item
  \textbf{Base case.} Prove \(P(0)\).
\item
  \textbf{Inductive step.} Prove \(\forall k\ge 0\),
  \(P(k)\Rightarrow P(k+1)\).
\end{enumerate}

Then by the same ``Algorithm I'' logic Knuth gives for induction (but
with the initial index shifted), the procedure constructs a proof for
\(P(0),P(1),\dots,P(n)\) in order, hence \(P(n)\) holds for every
\(n\ge 0\). Informally: from \(P(0)\) the step yields \(P(1)\); from
\(P(1)\) it yields \(P(2)\); and so on.
(\href{https://stackoverflow.com/questions/768948/how-does-division-work-in-mix}{PagePlace})

\begin{enumerate}
\def\labelenumi{\arabic{enumi})}
\setcounter{enumi}{1}
\tightlist
\item
  Strong induction starting at 0 (often convenient)
\end{enumerate}

Sometimes the step needs more than just \(P(k)\). Then use:

\begin{enumerate}
\def\labelenumi{\arabic{enumi}.}
\item
  \textbf{Base case.} Prove \(P(0)\).
\item
  \textbf{Inductive step (strong).} Assume \(P(0),P(1),\dots,P(k)\) are
  all true; prove \(P(k+1)\).
\end{enumerate}

This is equivalent in strength to weak induction but can be easier to
apply in algorithm proofs or recurrences. (Knuth emphasizes the
algorithmic viewpoint of induction and its ``construct a proof''
interpretation.)
(\href{https://stackoverflow.com/questions/768948/how-does-division-work-in-mix}{PagePlace})

\begin{enumerate}
\def\labelenumi{\arabic{enumi})}
\setcounter{enumi}{2}
\tightlist
\item
  General ``shifted'' induction (for all \(n\ge n_0\))
\end{enumerate}

More generally, if you want \(P(n)\) for all \(n\ge n_0\) (with any
integer \(n_0\)), prove:

\begin{enumerate}
\def\labelenumi{\arabic{enumi}.}
\item
  \textbf{Base case.} \(P(n_0)\).
\item
  \textbf{Step.} \(\forall k\ge n_0\), \(P(k)\Rightarrow P(k+1)\).
\end{enumerate}

This reduces immediately to the \(n_0=0\) form by reindexing
\(m=n-n_0\).

\begin{enumerate}
\def\labelenumi{\arabic{enumi})}
\setcounter{enumi}{3}
\tightlist
\item
  Quick example
\end{enumerate}

Claim: for all \(n\ge 0\), \(\sum_{i=1}^{n} i = \frac{n(n+1)}{2}\).

\begin{itemize}
\item
  \textbf{Base:} \(n=0\): empty sum \(=0\), and \(\frac{0\cdot1}{2}=0\).
\item
  \textbf{Step:} Assume \(\sum_{i=1}^{k} i = \frac{k(k+1)}{2}\). Then

  \[
  \sum_{i=1}^{k+1} i=\frac{k(k+1)}{2}+(k+1)=\frac{(k+1)(k+2)}{2},
  \]

  which is the formula with \(k+1\). Hence true for all \(n\ge 0\).
\end{itemize}

\subsection{2 {[}15{]}}\label{section-8}

Here's the issue with the ``proof''.

The induction step \emph{implicitly uses a case outside the hypothesis}
on the very first step:

\begin{itemize}
\item
  The step computes \(a^n = a^{n-1}\cdot a^{n-1} / a^{(n-1)-1}\). That
  references both \(a^{n-1}\) \textbf{and} \(a^{n-2}\).
\item
  Strong induction allows you to assume the statement for \(1,\dots,n\);
  however, when \(n=1\) (the first time the step must be valid), this
  formula needs \(a^{n-2}=a^{-1}\). The statement being proved was only
  for \textbf{positive} integers \(n\); it says nothing about \(n=0\)
  (which would correspond to the exponent \(-1\) here). So the step is
  \textbf{not valid for \(n=1\)} because it relies on \(P(0)\), which
  isn't part of the induction hypothesis. Put differently: the inductive
  step requires two earlier cases, \(n-1\) and \(n-2\); therefore it
  would need base cases covering up through \(n=2\) and also a statement
  for \(n=0\). But the proposition for \(n=0\) would be \(a^{-1}=1\),
  which is false unless \(a=1\).
\item
  Concretely, try \(a=2\). The ``proof'' would force \(2^{1}=1\) at
  \(n=2\), a contradiction.
\end{itemize}

So the mistake is an \emph{illegitimate induction step}: it invokes the
assertion for \(n-2\) (and at the first step, effectively \(n=0\)) even
though the theorem is only stated for positive \(n\). The step therefore
doesn't establish ``for all \(n\), if \(P(1),\dots,P(n)\) then
\(P(n+1)\)''---it already fails at \(n=1\).

\subsection{3 {[}18{]}}\label{section-9}

Let \(P(n)\) be the statement

\[
\sum_{k=1}^{n-1}\frac{1}{k(k+1)}=\frac{3}{2}-\frac{1}{n}.
\]

\begin{itemize}
\tightlist
\item
  For \(n=1\), the left side is an \textbf{empty sum}, hence \(0\). The
  right side is \(3/2-1=1/2\). So \(P(1)\) is \textbf{false}.
\end{itemize}

In the bogus proof, the ``base case'' line says ``for \(n=1\), clearly
\(\tfrac{3}{2}-\tfrac1n=\tfrac1{1\cdot2}\).'' That equality is
numerically true (\(1/2=1/2\)) \textbf{but it is not \(P(1)\)}; it
mistakenly replaces the empty sum (the actual left-hand side for
\(n=1\)) with the first term \(\tfrac{1}{1\cdot2}\), which belongs to
\(n\ge2\). In short, the base case checks the wrong statement. Without a
correct base case, the induction proves nothing.

(Algebra in the inductive step is fine; the logic fails because the
induction starts from a false/unchecked base.)

The correct identity (and proof)

The true formula is

\[
\sum_{k=1}^{n-1}\frac{1}{k(k+1)} \;=\; 1-\frac{1}{n}.
\]

Reason: \(\frac{1}{k(k+1)}=\frac{1}{k}-\frac{1}{k+1}\) telescopes,
giving

\[
\Bigl(1-\tfrac12\Bigr)+\Bigl(\tfrac12-\tfrac13\Bigr)+\cdots+\Bigl(\tfrac1{n-1}-\tfrac1n\Bigr)=1-\tfrac1n.
\]

If you want an induction proof:

\begin{itemize}
\tightlist
\item
  Base \(n=1\): empty sum \(=0=1-1\).
\item
  Step: assume \(\sum_{k=1}^{n-1}\frac{1}{k(k+1)}=1-\frac{1}{n}\). Then
\end{itemize}

\[
\sum_{k=1}^{n}\frac{1}{k(k+1)}
= \left(1-\frac{1}{n}\right) + \frac{1}{n(n+1)}
= 1-\frac{1}{n+1}.
\]

Sanity check at \(n=6\)

LHS \(=\frac12+\frac16+\frac1{12}+\frac1{20}+\frac1{30}= \frac56\).

\begin{itemize}
\tightlist
\item
  False formula's RHS: \(\tfrac32-\tfrac16=\tfrac43\) (doesn't match).
\item
  True formula's RHS: \(1-\tfrac16=\tfrac56\) (matches).
\end{itemize}

So the error in the exercise's ``proof'' is a faulty base case: it
confuses the empty sum for \(n=1\) with the first term \(1/(1\cdot2)\).

\subsection{4 {[}20{]}}\label{section-10}

Ingredients.

\begin{enumerate}
\def\labelenumi{\arabic{enumi}.}
\tightlist
\item
  Fibonacci recurrence: \(F_n = F_{n-1}+F_{n-2}\) for \(n>1\) (with
  \(F_0=0, F_1=1\)).
\item
  Golden-ratio identity: \(\varphi^2=\varphi+1\) (equivalently,
  \(\varphi\) is the positive root of \(x^2=x+1\)).
\item
  Context from the text: together with Eq. (3) in §1.2.1, these bounds
  form the pair \(\varphi^{n-2}\le F_n\le \varphi^{n-1}\). 
\end{enumerate}

Base cases.

\begin{itemize}
\tightlist
\item
  \(n=1\): \(F_1=1\) and \(\varphi^{1-2}=\varphi^{-1}\approx 0.618\).
  Hence \(F_1\ge \varphi^{-1}\).
\item
  \(n=2\): \(F_2=1\) and \(\varphi^{2-2}=\varphi^{0}=1\). Hence
  \(F_2\ge \varphi^{0}\) (equality).
\end{itemize}

Inductive step. Assume the inequality holds for all integers up to \(n\)
(with \(n\ge 2\)), in particular

\[
F_n \ge \varphi^{\,n-2}\quad\text{and}\quad F_{n-1}\ge \varphi^{\,n-3}.
\]

Then, using the Fibonacci recurrence,

\[
F_{n+1}=F_n+F_{n-1}\;\ge\;\varphi^{\,n-2}+\varphi^{\,n-3}
= \varphi^{\,n-3}(\varphi+1)
= \varphi^{\,n-3}\varphi^{2}
= \varphi^{\,n-1},
\]

where we used \(\varphi^2=\varphi+1\). Thus
\(F_{n+1}\ge \varphi^{\,n-1}\), completing the induction.

Therefore \(F_n\ge \varphi^{\,n-2}\) for all positive integers \(n\).

\subsection{5 {[}21{]}}\label{section-11}

Base case (\(n=2\)).

\(2\) is prime, so it is already a (one-factor) product of primes. Thus
\(P(2)\) holds.

Induction step. Assume \(P(n)\) holds for some \(n\ge 2\); i.e., every
integer \(m\) with \(2\le m\le n\) is a product of primes. We must show
\(P(n+1)\): that \(n+1\) is a product of primes.

Consider \(n+1\) in two cases.

\begin{enumerate}
\def\labelenumi{\arabic{enumi}.}
\item
  \(n+1\) is prime. Then \(n+1\) itself is a (one-factor) product of
  primes, so we're done.
\item
  \(n+1\) is composite. Then there exist integers \(a,b\) with
  \(1<a\le b<n+1\) and \(n+1=ab\). Hence \(2\le a\le n\) and
  \(2\le b\le n\). By the induction hypothesis, both \(a\) and \(b\) can
  be written as products of primes. Multiplying those two prime products
  gives a prime product for \(n+1\).
\end{enumerate}

In either case, \(n+1\) is a product of primes, so \(P(n+1)\) holds. By
strong induction, \(P(n)\) is true for all \(n\ge 2\).

Every integer \(>1\) can be expressed as a product of primes (existence
of prime factorization).

\subsection{6 {[}20{]}}\label{section-12}

In Algorithm E (Extended Euclid) Knuth defines variables
\(a',a,b',b,c,d\) and at step \textbf{E2} we compute \(q,r\) with

\[
c = qd + r,\qquad 0\le r<d.
\]

Step \textbf{E4} then updates the variables as follows:

\[
\begin{aligned}
&c\leftarrow d,\quad d\leftarrow r,\\
&t\leftarrow a',\quad a'\leftarrow a,\quad a\leftarrow t-qa,\\
&t\leftarrow b',\quad b'\leftarrow b,\quad b\leftarrow t-qb,
\end{aligned}
\]

and the claimed invariants are

\[
a'm+b'n=c,\qquad am+bn=d. \tag{6}
\]

The problem is to \textbf{prove that if (6) holds just before E4, it
still holds immediately after E4}. 

Assume, just before executing \textbf{E4}, that

\[
a'm+b'n=c,\qquad am+bn=d,
\]

and that at \textbf{E2} we computed \(q,r\) with \(c=qd+r\).

Let ``old'' denote values before E4 and ``new'' the values after E4.

From the assignment list of \textbf{E4} we have:

\[
\begin{aligned}
\text{new }c &= \text{old }d,\\
\text{new }d &= \text{old }r = \text{old }c - q\,\text{old }d,\\
\text{new }a' &= \text{old }a,\qquad \text{new }a = \text{old }a' - q\,\text{old }a,\\
\text{new }b' &= \text{old }b,\qquad \text{new }b = \text{old }b' - q\,\text{old }b.
\end{aligned}
\]

Invariant 1

\[
\begin{aligned}
\text{new }a' \, m + \text{new }b' \, n
&= (\text{old }a)m + (\text{old }b)n \\
&= \text{old }(am+bn) \\
&= \text{old }d \\
&= \text{new }c.
\end{aligned}
\]

So \(a'm+b'n=c\) holds after E4.

Invariant 2

\[
\begin{aligned}
\text{new }a \, m + \text{new }b \, n
&= (\text{old }a' - q\,\text{old }a)m + (\text{old }b' - q\,\text{old }b)n\\
&= \text{old }(a'm+b'n) - q\,\text{old }(am+bn)\\
&= \text{old }c - q\,\text{old }d\\
&= \text{old }r\\
&= \text{new }d.
\end{aligned}
\]

So \(am+bn=d\) also holds after E4.

Both equalities (6) are therefore preserved by E4. Because (6) is also
true at initialization (E1 sets \(a'=1,a=0,b'=0,b=1,c=m,d=n\), which
clearly satisfies \(a'm+b'n=m=c\) and \(am+bn=n=d\)), the invariants
hold every time we return to E2. This is exactly the fact Knuth alludes
to next to Eq. (6).

\subsection{7 {[}23{]}}\label{section-13}

Let

\[
S_n \;=\; n^2-(n-1)^2+(n-2)^2-\cdots+(-1)^{n-1}1^2
\]

(i.e., start at \(n^2\) and alternate signs down to \(1^2\)). We claim

\[
\boxed{\,S_n=\frac{n(n+1)}{2}\,}\quad\text{for all }n\ge 1.
\]

Induction proof

\textbf{Base \(n=1\).} \(S_1=1^2=1=\frac{1\cdot2}{2}\), true.

\textbf{Inductive step.} Assume \(S_n=\frac{n(n+1)}{2}\) for some
\(n\ge1\). Observe the simple recurrence (just ``prepend'' \((n+1)^2\)
and flip all subsequent signs):

\[
S_{n+1}=(n+1)^2 - S_n.
\]

Therefore,

\[
S_{n+1}=(n+1)^2-\frac{n(n+1)}{2}
=(n+1)\Bigl((n+1)-\frac n2\Bigr)
=\frac{(n+1)(n+2)}{2},
\]

which is the desired formula with \(n\) replaced by \(n+1\). By
induction, the identity holds for all \(n\ge1\).

\subsection{8 {[}25{]}}\label{section-14}

\begin{enumerate}
\def\labelenumi{(\alph{enumi})}
\tightlist
\item
  Nicomachus's theorem (cubes are sums of consecutive odd numbers)
\end{enumerate}

Claim. For every integer \(k\ge1\),

\[
k^3=\underbrace{(k^2-k+1)+(k^2-k+3)+\cdots+(k^2+k-1)}_{\text{\(k\) consecutive odd numbers}}.
\]

Proof by induction on \(k\).

\begin{itemize}
\item
  Base case \(k=1\). Right side is \(1\); indeed \(1^3=1\).
\item
  Inductive step. Assume the formula holds for \(k=n\):

  \[
  n^3=\sum_{i=0}^{n-1}\bigl(n^2-n+1+2i\bigr).
  \]

  Consider \(k=n+1\). The next block of \((n+1)\) consecutive odd
  numbers runs

  \[
  n^2+n+1,\;n^2+n+3,\;\ldots,\;n^2+3n+1,
  \]

  i.e., \(\sum_{i=0}^{n}\bigl(n^2+n+1+2i\bigr)\). Its sum is

  \[
  (n+1)(n^2+n+1)+2\sum_{i=0}^{n} i
  =(n+1)(n^2+n+1)+n(n+1)=(n+1)\bigl(n^2+2n+1\bigr)=(n+1)^3.
  \]

  Adding this block to the block that sums to \(n^3\) ``moves'' from the
  \(n\)-case to the \((n+1)\)-case, so the pattern holds for \(n+1\).
\end{itemize}

Thus, by induction, the claim is true for all \(k\ge1\).

\begin{enumerate}
\def\labelenumi{(\alph{enumi})}
\setcounter{enumi}{1}
\tightlist
\item
  Sum of the first \(n\) cubes
\end{enumerate}

From part (a), \(k^3\) is the sum of a block of \(k\) consecutive odd
numbers. Concatenating the blocks for \(k=1,2,\ldots,n\) uses exactly

\[
1+2+\cdots+n=\frac{n(n+1)}{2}
\]

odd numbers. But the sum of the first \(m\) odd numbers is \(m^2\)
(proved earlier in the section as \(1+3+\cdots+(2m-1)=m^2\)). 

Set \(m=\dfrac{n(n+1)}{2}\). Then

\[
\sum_{k=1}^{n} k^3
=\text{(sum of the first \(m\) odd numbers)}
=m^2
=\left(\frac{n(n+1)}{2}\right)^2
=\bigl(1+2+\cdots+n\bigr)^2.
\]

This is exactly the desired identity.

\subsection{9 {[}20{]}}\label{section-15}

\textbf{Claim.} For every integer \(n\ge 1\) and \(0<a<1\),
\((1-a)^n \ge 1-na\).

\textbf{Base case (\(n=1\)).} \((1-a)^1 = 1-a = 1-1\cdot a\). True
(equality).

\textbf{Induction step.} Assume for some \(n\ge 1\) that
\((1-a)^n \ge 1-na\). Because \(0<a<1\), we have \(1-a\ge 0\).
Multiplying the induction hypothesis by the nonnegative number \((1-a)\)
preserves the inequality:

\[
(1-a)^{n+1}
= (1-a)^n(1-a)
\;\ge\; (1-na)(1-a)
= 1-(n+1)a + n a^2
\;\ge\; 1-(n+1)a,
\]

since \(n a^2\ge 0\). Thus the statement holds for \(n+1\).

By induction, the inequality holds for all integers \(n\ge 1\).

\subsection{10 {[}M22{]}}\label{m22}

Base case (\(n=10\)). \(2^{10}=1024 > 1000 = 10^3\). So the statement
holds for \(n=10\).

Inductive step. Assume \(2^n > n^3\) for some \(n\ge 10\). We must show
\(2^{n+1}>(n+1)^3\).

From the inductive hypothesis,

\[
2^{n+1}=2\cdot 2^n \;>\; 2\,n^3.
\]

Thus it suffices to prove \(2\,n^3 \ge (n+1)^3\), i.e.

\[
2 \;\ge\; \Bigl(1+\frac1n\Bigr)^3.
\]

For \(n\ge 4\) we have

\[
\Bigl(1+\frac1n\Bigr)^3
= 1+\frac{3}{n}+\frac{3}{n^2}+\frac{1}{n^3}
\le 1+\frac{3}{4}+\frac{3}{16}+\frac{1}{64}
= \frac{125}{64}
< 2.
\]

Since our induction starts at \(n\ge 10\), the inequality certainly
holds; hence

\[
2^{n+1} \;>\; 2\,n^3 \;\ge\; (n+1)^3,
\]

which completes the inductive step.

By induction, \(2^n>n^3\) for all \(n\ge 10\). ∎

Remarks

\begin{itemize}
\item
  The threshold \(n\ge 10\) is sharp: \(2^9=512<729=9^3\), so the claim
  fails at \(n=9\) but holds from \(n=10\) onward.
\item
  An equivalent (and sometimes handy) way to see the step is to expand
  the difference directly:

  \[
  2n^3-(n+1)^3 = n^3-3n^2-3n-1
  = (n-1)^3-6n,
  \]

  which is \(>0\) for all \(n\ge 10\) (e.g., at \(n=10\):
  \(729-60=669>0\)).
\end{itemize}

\subsection{11 {[}M30{]}}\label{m30-1}

Step 1: A key algebraic identity

Factor the denominator:

\[
x^4+4=(x^2+2x+2)(x^2-2x+2)=(x+1)^2+1\;\; \times\;\; (x-1)^2+1.
\]

Seek constants \(A,B,C,D\) with

\[
\frac{x^3}{x^4+4}=\frac{Ax+B}{x^2+2x+2}+\frac{Cx+D}{x^2-2x+2}.
\]

Equating coefficients gives \(A=C=\tfrac12\), \(B=\tfrac12\),
\(D=-\tfrac12\). Hence

\[
\boxed{\;\frac{x^3}{x^4+4}
=\frac12\!\left(\frac{x+1}{x^2+2x+2}+\frac{x-1}{x^2-2x+2}\right)
=\frac12\!\left(\frac{x+1}{(x+1)^2+1}+\frac{x-1}{(x-1)^2+1}\right)\; } \tag{★}
\]

Step 2: Plug in \(x=2k+1\)

With \(x=2k+1\),

\[
\frac{(2k+1)^3}{(2k+1)^4+4}
=\frac12\left(\frac{2k+2}{(2k+2)^2+1}+\frac{2k}{(2k)^2+1}\right).
\]

Therefore

\[
S_n=\frac12\sum_{k=0}^n(-1)^k\!\left(\frac{2k+2}{(2k+2)^2+1}+\frac{2k}{(2k)^2+1}\right).
\]

Split into two sums and shift index in the first:

\[
\begin{aligned}
S_n&=\frac12\Bigg[\sum_{k=0}^n(-1)^k\frac{2k}{(2k)^2+1}
+\sum_{k=0}^n(-1)^k\frac{2k+2}{(2k+2)^2+1}\Bigg]\\
&=\frac12\Bigg[\sum_{j=0}^n(-1)^j\frac{2j}{(2j)^2+1}
-\sum_{j=1}^{n+1}(-1)^j\frac{2j}{(2j)^2+1}\Bigg].
\end{aligned}
\]

Now everything cancels \textbf{telescopically} except the endpoint terms
\(j=0\) (which is \(0\)) and \(j=n+1\):

\[
S_n=\frac12\left[0-(-1)^{\,n+1}\frac{2(n+1)}{(2(n+1))^2+1}\right]
=\boxed{\;(-1)^n\,\frac{n+1}{4(n+1)^2+1}\; }.
\]

Quick checks

\begin{itemize}
\tightlist
\item
  \(n=0\): \(S_0=\dfrac{1}{1^4+4}=\dfrac15\). Formula gives
  \(\dfrac{1}{4\cdot1^2+1}=\dfrac15\).
\item
  \(n=1\): \(S_1=\dfrac15-\dfrac{27}{85}=-\dfrac{2}{17}\). Formula gives
  \(-\dfrac{2}{17}\).
\end{itemize}

\textbf{Answer.}
\(\displaystyle \boxed{\,\sum_{k=0}^{n}(-1)^k\frac{(2k+1)^3}{(2k+1)^4+4}=(-1)^n\frac{n+1}{4(n+1)^2+1}\,}\).

\subsection{12 {[}M25{]}}\label{m25-1}

Generalize Euclid's Algorithm \textbf{E} so it can accept inputs of the
form \(u+v\sqrt{2}\) (with \(u,v\in\mathbb Z\)) and still do all
computations ``elementarily'' (i.e., without using an infinite decimal
for \(\sqrt{2}\)). Then prove the computation does \textbf{not}
terminate when \(m=1\) and \(n=\sqrt{2}\). (Knuth hints at this
generalization and the non-termination case in the text; the exercise
states it explicitly.)

A. Generalizing Algorithm E to inputs \(u+v\sqrt{2}\) Representation and
primitive operations

Represent each value \(x\) as an integer pair \((u,v)\) meaning
\(x=u+v\sqrt{2}\).

\begin{itemize}
\item
  \textbf{Addition/Subtraction.}
  \((u_1,v_1)\pm(u_2,v_2)=(u_1\pm u_2,\; v_1\pm v_2)\) --- integer work
  only.
\item
  \textbf{Zero test.} \((u,v)=0\iff u=0 \land v=0\).
\item
  \textbf{Sign/Comparison without decimals.} To decide whether
  \((u,v)\ge 0\) (hence compare two numbers by looking at their
  difference), use exact integer tests:

  \[
  u+v\sqrt{2}\ge 0\quad\Longleftrightarrow\quad
  \begin{cases}
  \text{if } v\ge 0:& u\ge 0\;\text{ or }\;u^2\le 2v^2,\\[2pt]
  \text{if } v<0:& u>0\;\text{ and }\;u^2\ge 2v^2.
  \end{cases}
  \]

  This comes from analyzing \(u\ge -v\sqrt{2}\) and squaring only in the
  cases where both sides have the same sign; it uses only integer
  arithmetic.
\end{itemize}

The ``division with integer quotient'' step

Euclid's step \textbf{E1} needs: given positive \(m,n\) (now of form
\(u+v\sqrt{2}\)), find integers \(q\) and a remainder \(r\) with

\[
m=qn+r,\qquad 0\le r<n,
\]

where \(q\in\mathbb Z\) (as in Knuth's discussion) and \(r\) is again of
the form \(u+v\sqrt{2}\).

You can compute \(q\) \textbf{without decimals} as follows.

\begin{enumerate}
\def\labelenumi{\arabic{enumi}.}
\item
  \textbf{Coarse integer bounds for \(\sqrt2\).} Use
  \(1\le \sqrt2 \le 2\) to bound any \(u+v\sqrt2\) by integers:

  \begin{itemize}
  \tightlist
  \item
    If \(v\ge 0\): \(u+v \le u+v\sqrt{2} \le u+2v\).
  \item
    If \(v<0\): \(u+2v \le u+v\sqrt{2} \le u+v\).
  \end{itemize}

  From these bounds for \(m\) and \(n\) you get integer bounds
  \(L\le m/n < U\).
\item
  \textbf{Integer bisection on \(q\).} Do binary search on the integer
  interval \([\lfloor L\rfloor,\lceil U\rceil]\) to find the unique
  \(q\) satisfying

  \[
  qn \le m < (q+1)n,
  \]

  using only the sign test above to compare \(m-(q)n\) and \(m-(q+1)n\)
  with zero.
\item
  \textbf{Remainder.} Set \(r=m-qn\) (pair subtraction). Then either
  \(r=0\) or \(0\le r<n\).
\end{enumerate}

This realizes \textbf{E1} with only integer operations and exact
comparisons. The rest of Algorithm \textbf{E} is unchanged:

\begin{itemize}
\tightlist
\item
  \textbf{E2.} If \(r=0\) return \(n\) (the ``gcd'' in this setting).
\item
  \textbf{E3.} Set \(m\leftarrow n,\; n\leftarrow r\) and repeat.
\end{itemize}

Knuth notes (and gives an example) that with this extension the
inductive assertions \(A_1\)--\(A_6\) still hold; e.g., starting with
\(m=12-6\sqrt2\) and \(n=20-10\sqrt2\), the method returns
\(d=4-2\sqrt2\).

B. Non-termination for \(m=1,\; n=\sqrt2\)

Run the generalized algorithm:

\begin{enumerate}
\def\labelenumi{\arabic{enumi}.}
\item
  Start with \(m=1\), \(n=\sqrt2\). Since \(m<n\), the quotient is
  \(q=0\) and \(r=1\). After the swap: \(m\gets \sqrt2,\; n\gets 1\).
\item
  Now \(m/n=\sqrt2\). Thus \(q=\lfloor\sqrt2\rfloor=1\), so
  \(r=\sqrt2-1\). Swap: \(m\gets 1,\; n\gets \sqrt2-1\).
\item
  Next \(m/n=\dfrac{1}{\sqrt2-1}=\sqrt2+1\). Hence \(q=2\) and

  \[
  r = 1-2(\sqrt2-1)=3-2\sqrt2=(\sqrt2-1)^2.
  \]

  Swap: \(m\gets \sqrt2-1,\; n\gets (\sqrt2-1)^2\).
\item
  Repeat: one checks inductively that at each subsequent step

  \[
  q=2,\qquad r=(\sqrt2-1)^{k+2}\quad\text{when }(m,n)=\bigl((\sqrt2-1)^k,(\sqrt2-1)^{k+1}\bigr).
  \]

  Thus the remainders are
  \((\sqrt2-1), (\sqrt2-1)^2, (\sqrt2-1)^3, \ldots\)
\end{enumerate}

Since \(0<\sqrt2-1<1\), this positive sequence is strictly decreasing
and tends to \(0\), but it never becomes \(0\). Therefore the algorithm
\textbf{never reaches} \(r=0\) and hence \textbf{does not terminate} for
\(m=1,\; n=\sqrt2\). This is exactly the phenomenon Knuth points out in
the text: the usual inductive correctness assertions don't, by
themselves, establish termination in this generalized setting.

\subsection{13 {[}M23{]}}\label{m23}

The T-augmented assertions

Let \(n\) denote the original input value that becomes \(d\) at \(E1\).
We append the following bounds to the six assertions \(A1,\dots,A6\)
(all other parts of each assertion remain unchanged exactly as in Fig.
4):

\begin{itemize}
\tightlist
\item
  \textbf{A1}: \(m>0,\ n>0,\ \boxed{T=0}\). (This already implies
  \(T\le 3n\).)
\item
  \textbf{A2} (first entry to \(E2\)):
  \(c=m>0,\ d=n>0,\ a=b'=0,\ a'=b=1,\ \boxed{T\le 3(n-d)+1}\). (Here
  \(d=n\), so this is \(T\le 1\); together with the exercise's ``append
  \(T=1\) to \(A2\)'' it is consistent.)
\item
  \textbf{A3} (after \(E2\)):
  \(am+bn=d,\ a'm+b'n=c=qd+r,\ 0\le r<d,\ \gcd(c,d)=\gcd(m,n),\ \boxed{T\le 3(n-d)+2}\).
\item
  \textbf{A4} (if \(r=0\) at \(E3\), i.e., about to stop):
  \(am+bn=d=\gcd(m,n),\ \boxed{T\le 3(n-d)+3}\).
\item
  \textbf{A5} (if \(r\ne 0\) at \(E3\)):
  \(am+bn=d,\ a'm+b'n=c=qd+r,\ 0<r<d,\ \gcd(c,d)=\gcd(m,n),\ \boxed{T\le 3(n-d)+3}\).
\item
  \textbf{A6} (after \(E4\)):
  \(am+bn=d,\ a'm+b'n=c,\ d>0,\ \gcd(c,d)=\gcd(m,n),\ \boxed{T\le 3(n-d)+1}\).
\end{itemize}

Why these particular constants \(+1,+2,+3,+3,+1\)? They are chosen so
that the ``one‐step'' increase of \(T\) when entering a box is exactly
matched by the change (or non-change) of \(d\) in that box, making every
implication in the flow-chart proof go through mechanically.

Inductive verification (box by box)

Use Knuth's proof rule: \emph{If the assertion on every arrow entering a
box is true \textbf{before} the box executes, then after execution the
assertions on all outgoing arrows are true.} I'll verify only the
\(T\)-part (the non-\(T\) parts are exactly those in Fig. 4 and are
already known to propagate).

\textbf{E1 (Initialize).} Enter with \(A1\): \(T=0\). Executing \(E1\)
sets \(d\leftarrow n\) and (by the exercise) increments \(T\) to \(1\).
The outgoing arrow is \(A2\), which demands \(T\le 3(n-d)+1\). Since
here \(d=n\), this is \(T\le1\), satisfied by \(T=1\).

\textbf{E2 (Divide).} There are two possible incoming arrows:

\begin{itemize}
\tightlist
\item
  From \textbf{A2}: we had \(T\le 3(n-d)+1\). Entering \(E2\) adds \(1\)
  to \(T\), while \(d\) is unchanged. Thus after \(E2\),
  \(T\le 3(n-d)+2\), exactly the bound in \textbf{A3}.
\item
  From \textbf{A6}: we had \(T\le 3(n-d)+1\). Again \(T\) gains \(1\)
  and \(d\) is unchanged, so we get \(T\le 3(n-d)+2\), i.e.,
  \textbf{A3}. (As Knuth notes, \(A2\) is stronger than \(A6\), but
  either one suffices to imply \(A3\) after \(E2\).)
\end{itemize}

\textbf{E3 (Remainder zero?).} Incoming \textbf{A3} has
\(T\le 3(n-d)+2\); entering \(E3\) bumps \(T\) by \(1\), \(d\) stays the
same.

\begin{itemize}
\tightlist
\item
  If \(r=0\): outgoing \textbf{A4} requires \(T\le 3(n-d)+3\), which now
  holds.
\item
  If \(r\ne 0\): outgoing \textbf{A5} requires the same bound
  \(T\le 3(n-d)+3\), also holds.
\end{itemize}

\textbf{E4 (Recycle).} Incoming \textbf{A5} has \(T\le 3(n-d)+3\).
Entering \(E4\) adds \(1\) to \(T\) and sets \(d\leftarrow r\) with
\(0<r<d\). So after \(E4\),

\[
T\ \le\ 3(n-d)+3\ +1\ =\ 3(n-d)+4.
\]

Because \(r\le d-1\), we have
\(3(n-r)+1\ \ge\ 3(n-(d-1))+1\ =\ 3(n-d)+4\). Hence

\[
T\ \le\ 3(n-r)+1\ =\ 3(n-d_{\text{new}})+1,
\]

which is exactly the \(T\)-bound required by \textbf{A6}. ✔

Finally, note the next transition is \textbf{A6 \(\to\) E2}: entering
\(E2\) adds \(1\) to \(T\) with \(d\) unchanged, turning
\(T\le 3(n-d)+1\) (from \(A6\)) into \(T\le 3(n-d)+2\) (the bound in
\(A3\)), as already verified above.

Why every assertion implies \(T\le 3n\)

In every state we have \(d\ge 1\) (the chart maintains \(d>0\) except at
termination), so for each assertion the added bound is at most

\[
\max_d\{\,3(n-d)+c_i\,\}\ \le\ 3(n-1)+3\ =\ 3n,
\]

since the largest constant used is \(3\). Thus \emph{any} of
\(A1,\dots,A6\) implies \(T\le 3n\). (At the stop arrow \textbf{A4} we
also have \(d=\gcd(m,n)\ge1\).)

\subsection{14 {[}50{]}}\label{section-16}

\subsection{15 {[}HM28{]}}\label{hm28}

\begin{enumerate}
\def\labelenumi{(\alph{enumi})}
\tightlist
\item
  The integers \(\mathbb Z\) are \textbf{not} well-ordered by the usual
  ``\textless{}''
\end{enumerate}

Take the nonempty subset \(A=\{\ldots,-3,-2,-1\}\subset\mathbb Z\). If
``\textless{}'' were a well-ordering on \(\mathbb Z\), \(A\) would have
a least element. But for any \(a\in A\), \(a-1\in A\) and \(a-1<a\).
Hence \(A\) has no least element. Thus ``\textless{}'' is not a
well-ordering on \(\mathbb Z\). (Condition (iii) fails.)

\begin{enumerate}
\def\labelenumi{(\alph{enumi})}
\setcounter{enumi}{1}
\tightlist
\item
  A well-ordering on \(\mathbb Z\)
\end{enumerate}

Define ``≺'' on \(\mathbb Z\) by \textbf{short-absolute-value order}
(sometimes called the 0,1,−1,2,−2,\ldots{} order):

\[
0 \prec 1 \prec -1 \prec 2 \prec -2 \prec 3 \prec -3 \prec \cdots,
\]

equivalently: for \(a,b\in\mathbb Z\),

\[
a \prec b \iff \bigl(|a|<|b|\bigr)\ \text{or}\ \bigl(|a|=|b|\ \text{and}\ a\ge 0>b\bigr).
\]

\begin{itemize}
\tightlist
\item
  \begin{enumerate}
  \def\labelenumi{(\roman{enumi})}
  \tightlist
  \item
    Transitivity and (ii) trichotomy are immediate from the definition.
  \end{enumerate}
\item
  \begin{enumerate}
  \def\labelenumi{(\roman{enumi})}
  \setcounter{enumi}{2}
  \tightlist
  \item
    Let \(X\subseteq\mathbb Z\) be nonempty. Consider the set of
    absolute values \(B=\{|x|:x\in X\}\subset\mathbb N\). Since
    \(\mathbb N\) is well-ordered by \textless, \(B\) has a least
    element \(r\). Then \(X\) contains either \(r\) or \(-r\) (or both).
    If both occur, pick \(r\) (the nonnegative one). That chosen element
    is least in \(X\) under ``≺''. Hence ``≺'' is a well-ordering of
    \(\mathbb Z\).
  \end{enumerate}
\end{itemize}

(Another equivalent description: map \(f:\mathbb Z\to\mathbb N\) by
\(f(0)=0,\ f(1)=1,\ f(-1)=2,\ f(2)=3,\ f(-2)=4,\ldots\). Then define
\(a\prec b\) iff \(f(a)<f(b)\). Since \((\mathbb N,<)\) is well-ordered,
this transports a well-ordering to \(\mathbb Z\).)

\begin{enumerate}
\def\labelenumi{(\alph{enumi})}
\setcounter{enumi}{2}
\tightlist
\item
  \([0,\infty)\subset\mathbb R\) is \textbf{not} well-ordered by
  ``\textless{}''
\end{enumerate}

Consider \(A=\{1/n:n\in\mathbb N\}\subset(0,1]\). This \(A\) is nonempty
but has no least element in the usual order: for any \(1/n\in A\), there
is \(1/(n+1)\in A\) with \(1/(n+1)<1/n\). Thus ``\textless{}'' is not a
well-ordering on the nonnegative reals.

\begin{enumerate}
\def\labelenumi{(\alph{enumi})}
\setcounter{enumi}{3}
\tightlist
\item
  Lexicographic order on \(T_n=S^n\) is a well-ordering (finite
  products)
\end{enumerate}

Assume \((S,\prec)\) is well-ordered. Define lexicographic order on
\(T_n=S^n\): \((x_1,\dots,x_n)\prec (y_1,\dots,y_n)\) if for some \(k\)
(with \(1\le k\le n\)) we have \(x_j=y_j\) for \(j<k\) and
\(x_k\prec y_k\) in \(S\).

\begin{itemize}
\tightlist
\item
  \begin{enumerate}
  \def\labelenumi{(\roman{enumi})}
  \tightlist
  \item
    and (ii) are standard for lex order.
  \end{enumerate}
\item
  \begin{enumerate}
  \def\labelenumi{(\roman{enumi})}
  \setcounter{enumi}{2}
  \tightlist
  \item
    Let \(U\subseteq T_n\) be nonempty. Choose \(x_1\) minimal among the
    first coordinates of elements of \(U\) (possible since \(S\) is
    well-ordered). Among those elements of \(U\) whose first coordinate
    equals that \(x_1\), choose the minimal second coordinate \(x_2\);
    continue this process finitely many steps. We obtain
    \((x_1,\dots,x_n)\in U\) which is \(\prec\)-least in \(U\). Hence
    lex order well-orders \(T_n\).
  \end{enumerate}
\end{itemize}

\begin{enumerate}
\def\labelenumi{(\alph{enumi})}
\setcounter{enumi}{4}
\tightlist
\item
  ``Shortlex'' on \(T=\bigcup_{n\ge1}T_n\) is also a well-ordering
\end{enumerate}

Define \((x_1,\dots,x_m)\prec (y_1,\dots,y_n)\) if either:

\begin{enumerate}
\def\labelenumi{\arabic{enumi}.}
\tightlist
\item
  they agree up to some first differing position \(k\le\min(m,n)\) and
  \(x_k\prec y_k\) in \(S\); or
\item
  \(m<n\) and \((x_1,\dots,x_m)\) is a proper prefix of
  \((y_1,\dots,y_n)\) (i.e., all first \(m\) entries are equal).
\end{enumerate}

\begin{itemize}
\tightlist
\item
  \begin{enumerate}
  \def\labelenumi{(\roman{enumi})}
  \tightlist
  \item
    and (ii) again hold (transitivity and totality).
  \end{enumerate}
\item
  \begin{enumerate}
  \def\labelenumi{(\roman{enumi})}
  \setcounter{enumi}{2}
  \tightlist
  \item
    Let \(U\subseteq T\) be nonempty. Consider the set of lengths
    \(L=\{\,m\ge1:\exists\,\mathbf x\in U\text{ of length }m\,\}\subset\mathbb N\).
    Since \(\mathbb N\) is well-ordered, \(L\) has a least length
    \(m_0\). Then \(U\cap T_{m_0}\) is nonempty, and by part (d) the lex
    order well-orders \(T_{m_0}\). Hence \(U\) has a least element (the
    lex-least among those of minimal length). So ``shortlex''
    well-orders \(T\).
  \end{enumerate}
\end{itemize}

\begin{enumerate}
\def\labelenumi{(\alph{enumi})}
\setcounter{enumi}{5}
\tightlist
\item
  Characterization via absence of infinite descending chains
\end{enumerate}

Claim: A relation ``≺'' on \(S\) is a well-ordering \textbf{iff} it
satisfies (i) transitivity, (ii) trichotomy, \textbf{and} there is
\textbf{no} infinite strictly descending chain
\(x_1 \succ x_2 \succ x_3 \succ \cdots\).

\begin{itemize}
\tightlist
\item
  \textbf{(⇒)} If ``≺'' is a well-ordering and there were such a chain,
  then the set \(\{x_1,x_2,\ldots\}\) would be a nonempty subset without
  a least element (since each \(x_{j+1}\prec x_j\)), contradicting
  (iii).
\item
  \textbf{(⇐)} Suppose (i) and (ii) hold and no infinite descending
  chain exists. Let \(A\subseteq S\) be nonempty. If \(A\) had no least
  element, pick \(x_1\in A\). Since \(x_1\) is not least, there is
  \(x_2\in A\) with \(x_2\prec x_1\). Inductively build
  \(x_{k+1}\prec x_k\) in \(A\) forever, making an infinite descending
  chain---a contradiction. Hence every nonempty subset has a least
  element, i.e., ``≺'' is a well-ordering.
\end{itemize}

\begin{enumerate}
\def\labelenumi{(\alph{enumi})}
\setcounter{enumi}{6}
\tightlist
\item
  Well-ordering induction (transfinite induction)
\end{enumerate}

Assume \((S,\prec)\) is well-ordered and \(P(x)\) is a statement about
\(x\in S\) with the property:

\begin{quote}
For each \(x\in S\), if \(P(y)\) holds for all \(y\prec x\), then
\(P(x)\) holds.
\end{quote}

To prove \(P(x)\) for all \(x\in S\), suppose not and let
\(C=\{x\in S:\neg P(x)\}\). If \(C\neq\varnothing\), well-ordering gives
a \(\prec\)-least \(c\in C\). But then \(P(y)\) holds for all
\(y\prec c\), so the premise yields \(P(c)\)---contradiction. Hence
\(C=\varnothing\) and \(P\) holds everywhere. This is the general
well-ordering (transfinite) induction principle; when \(S=\mathbb N\)
with the usual ``\textless{}'', it reduces to ordinary induction.

\bookmarksetup{startatroot}

\chapter{1.2.2 Numbers, powers, and
logarithms}\label{numbers-powers-and-logarithms}

\subsection{1 {[}00{]}}\label{section-17}

There is \textbf{no} smallest positive rational number.

Assume there is a smallest positive rational number; call it \(r\).
Since \(r>0\), the number \(r/2\) is also rational and positive, and
\(r/2<r\). This contradicts the assumption that \(r\) was the smallest
positive rational. Therefore such a smallest element does not exist. ∎

Notes: - Rationals are numbers of the form \(p/q\) with integers \(p\)
and \(q>0\) (see the section's definitions of integers/rationals), but
the above argument doesn't even need their explicit form: it uses only
closure under division by 2 and positivity. - In order-theory terms, the
set \(\{x\in\mathbb{Q}: x>0\}\) has no minimum; its greatest lower bound
(infimum) is \(0\). - Contrast: the positive integers do have a least
element (namely \(1\)); the positive rationals (and reals) are not
well-ordered under the usual \(<\).

\subsection{2 {[}00{]}}\label{section-18}

What Knuth means by ``decimal expansion''

In §1.2.2 Knuth defines a real number's decimal expansion as
\(x = n + 0.d_1d_2d_3\ldots\) with digits \(d_i\in\{0,\dots,9\}\), and
with the explicit convention that the digit sequence does not end with
infinitely many 9s (to make the expansion unique).

Answer

\textbf{No.} As written, \(1+0.239999999\ldots\) (i.e.,
\(1.239999999\ldots\)) is not an allowed ``decimal expansion'' under
Knuth's convention, because it ends with infinitely many 9s. However, it
does represent a perfectly good real number---namely \(1.24\)---and the
proper (canonical) decimal expansion for that number is
\(1.240000\ldots\).

Why \(0.239999\ldots = 0.24\)

Compute the tail as a geometric series:

\[
0.239999\ldots
= 0.23 + \sum_{k=3}^{\infty} \frac{9}{10^{k}}
= 0.23 + 9\cdot\frac{10^{-3}}{1-\frac{1}{10}}
= 0.23 + 9\cdot\frac{10^{-3}}{\frac{9}{10}}
= 0.23 + 10^{-2}
= 0.24.
\]

Therefore,

\[
1 + 0.239999\ldots = 1 + 0.24 = 1.24.
\]

\subsection{3 {[}02{]}}\label{section-19}

\begin{enumerate}
\def\labelenumi{\arabic{enumi}.}
\item
  Use the definition of negative integer powers. For any nonzero \(a\)
  and positive integer \(n\),

  \[
  a^{-n}=\frac{1}{a^{\,n}}.
  \]
\item
  Evaluate the positive power in the denominator.

  \[
  (-3)^3 = (-3)\cdot(-3)\cdot(-3)=9\cdot(-3)=-27,
  \]

  since an odd number of negative factors yields a negative result.
\item
  Invert.

  \[
  (-3)^{-3}=\frac{1}{(-3)^3}=\frac{1}{-27}=-\frac{1}{27}.
  \]
\item
  Quick check.

  \[
  (-3)^{-3}\cdot(-3)^3 = (-3)^{0}=1,
  \]

  confirming the result.
\end{enumerate}

\textbf{Answer:} \(\displaystyle -\tfrac{1}{27}\).

\subsection{4 {[}05{]}}\label{section-20}

Solution, two equivalent ways.

\begin{enumerate}
\def\labelenumi{\arabic{enumi}.}
\tightlist
\item
  Use base-2 and exponent laws. \(0.125=\tfrac{1}{8}=2^{-3}\). Using
  \((b^x)^y=b^{xy}\),
\end{enumerate}

\[
(0.125)^{-2/3}=(2^{-3})^{-2/3}=2^{(-3)\cdot(-2/3)}=2^{2}=4.
\]

(The law \((b^x)^y=b^{xy}\) for exponents appears in Eq. (5).) 

\begin{enumerate}
\def\labelenumi{\arabic{enumi}.}
\setcounter{enumi}{1}
\tightlist
\item
  Use the definition of rational exponents. For \(a>0\) and rational
  \(p/q\), \(a^{p/q}=\sqrt[q]{a^{\,p}}\) (Eq. (6)). Thus
\end{enumerate}

\[
(0.125)^{-2/3}=\frac{1}{(0.125)^{2/3}}
=\frac{1}{\big(\sqrt[3]{0.125}\big)^2}
=\frac{1}{\left(\sqrt[3]{\tfrac{1}{8}}\right)^2}
=\frac{1}{\left(\tfrac{1}{2}\right)^2}
=4.
\]

\textbf{Answer:} \(4\).

\subsection{5 {[}05{]}}\label{section-21}

Let \(x\) be a real number. Write

\[
x=n+0.d_1d_2d_3\ldots\quad\text{(base 2)},
\]

meaning

\[
x \;=\; n+\sum_{i=1}^{\infty}\frac{d_i}{2^i},
\quad\text{where each } d_i\in\{0,1\}.
\]

To make this representation \textbf{unique}, forbid tails of infinitely
many 1s (exactly analogous to forbidding infinitely many 9s in decimal).
With that convention, the \textbf{binary analogue of Eq. (2)} is:

\[
\boxed{\;
n+\sum_{i=1}^{k}\frac{d_i}{2^i}
\;\le\;
x
\;<\;
n+\sum_{i=1}^{k}\frac{d_i}{2^i}+\frac{1}{2^{k}}
\qquad\text{for all } k\ge1. \;
}
\]

Why this works

Split the infinite sum at \(k\):

\[
x
= \underbrace{\Bigl(n+\sum_{i=1}^{k}\frac{d_i}{2^i}\Bigr)}_{S_k}
+ \underbrace{\sum_{i=k+1}^{\infty}\frac{d_i}{2^i}}_{R_k}.
\]

Because \(d_i\in\{0,1\}\),

\[
0\le R_k < \sum_{i=k+1}^{\infty}\frac{1}{2^i}
= \frac{1}{2^{k}}.
\]

Hence \(S_k\le x < S_k+2^{-k}\), which is exactly the displayed
inequality.

Existence and digit construction (greedy rule)

Given any real \(x\), let \(n=\lfloor x\rfloor\). For the fractional
part \(\{x\}=x-n\), define digits iteratively:

\[
d_i=\bigl\lfloor 2\,r_{i-1}\bigr\rfloor,\qquad
r_i=2\,r_{i-1}-d_i,\qquad r_0=\{x\}.
\]

This produces \(d_i\in\{0,1\}\) and \(\{x\}=\sum_{i\ge1} d_i/2^i\). If
the process would eventually give all 1s, switch to the equivalent
finite form (e.g., \(0.0111\ldots_2=0.1000\ldots_2\)) and keep our ``no
infinite 1s'' rule for uniqueness.

\subsection{6 {[}10{]}}\label{section-22}

Rule (the ``lexicographic'' comparison)

\begin{enumerate}
\def\labelenumi{\arabic{enumi}.}
\item
  Compare integer parts. • If \(m<n\) then \(x<y\). • If \(m>n\) then
  \(x>y\). • If \(m=n\), go to step 2.
\item
  Compare fractional digits left to right. Let \(k\) be the first
  position with \(d_k\neq e_k\). • If no such \(k\) exists (all digits
  match), then \(x=y\). • If \(d_k<e_k\), then \(x<y\); if \(d_k>e_k\),
  then \(x>y\).
\end{enumerate}

That's it. This works for positive and negative numbers because \(m\) is
the floor of the real (so \(x\in[m,m+1)\), \(y\in[n,n+1)\)); the
fractional part \(0.d_1d_2\ldots\) is always in \([0,1)\). 

Why the rule is correct (brief proof)

• If \(m\neq n\): Since \(0.d_1d_2\ldots\in[0,1)\), we have
\(x\in[m,m+1)\) and \(y\in[n,n+1)\). Hence \(x<y\iff m<n\). (Likewise
for \(>\).)

• If \(m=n\): Write

\[
x-y=\sum_{j\ge1}\frac{d_j-e_j}{10^j}.
\]

If all digits match, the sum is 0, so \(x=y\).

Otherwise, let \(k\) be the first index with \(d_k\ne e_k\). Then

\[
x-y=\frac{d_k-e_k}{10^k}+R,\qquad
R=\sum_{j\ge k+1}\frac{d_j-e_j}{10^j}.
\]

We have
\(|R|\le \frac{9}{10^{k+1}}+\frac{9}{10^{k+2}}+\cdots=\frac{1}{10^k}\).
Because Knuth disallows decimal expansions that \emph{end} with
infinitely many 9s, that bound is in fact \emph{strict} here:
\(|R|<\frac{1}{10^k}\). Since \(|d_k-e_k|\ge1\), the leading term
\((d_k-e_k)/10^k\) dominates \(R\), so
\(\operatorname{sign}(x-y)=\operatorname{sign}(d_k-e_k)\). Therefore
\(x<y\) iff \(d_k<e_k\), etc.

Notes and examples

\begin{itemize}
\item
  Equality test is simple because of the uniqueness convention: \(x=y\)
  iff \(m=n\) and \(d_j=e_j\) for all \(j\) (no
  ``\(0.999\ldots=1.000\ldots\)'' ambiguity under Knuth's definition).
\item
  Negative numbers behave the same way. Example: \(x=-2+0.3=-1.7\),
  \(y=-2+0.7=-1.3\). Here \(m=n=-2\), first differing digit is \(k=1\)
  with \(3<7\), so \(x<y\).
\end{itemize}

\subsection{7 {[}M23{]}}\label{m23-1}

Assuming \(b>0\) and using the \emph{definition} of integer powers

\[
b^{0}=1,\quad b^{n}=b^{\,n-1}b\ (n>0),\quad b^{n}=b^{\,n+1}/b\ (n<0),
\]

prove the ``laws of exponents'' for all integers \(x,y\): (i)
\(b^{x+y}=b^{x}b^{y}\) and (ii) \((b^{x})^{y}=b^{xy}\). (See the
definition in Eq. (4) and the request of Ex. 7 in §1.2.2.) 

We proceed in two steps. Throughout, \(b>0\), so \(b^{n}\neq 0\) for all
integers \(n\).

\begin{enumerate}
\def\labelenumi{\arabic{enumi})}
\tightlist
\item
  Product rule: \(b^{x+y}=b^{x}b^{y}\) for all integers \(x,y\)
\end{enumerate}

\textbf{Step A (nonnegative addend).} Fix any integer \(x\). We show by
induction on \(n\ge 0\) that

\[
b^{x+n}=b^{x}b^{n}.
\]

\begin{itemize}
\item
  \textbf{Base \(n=0\):} \(b^{x+0}=b^{x}=b^{x}\cdot b^{0}\) since
  \(b^{0}=1\).
\item
  \textbf{Induction step:} If \(b^{x+n}=b^{x}b^{n}\), then

  \[
  b^{x+(n+1)}=(b^{x+n})\,b=(b^{x}b^{n})\,b=b^{x}b^{n+1},
  \]

  using the defining recurrence for positive exponents.
\end{itemize}

Thus, for any integer \(x\) and any \(n\ge 0\), \(b^{x+n}=b^{x}b^{n}\).

\textbf{Step B (negative addend).} Let \(y<0\) and write \(y=-k\) with
\(k\ge 0\). Using Step A with the nonnegative addend \(k\),

\[
b^{(x+y)}=b^{x-k}\quad\text{and}\quad b^{x}=b^{x-k}b^{k}.
\]

Multiply both sides of the second identity by \(b^{-k}\) (i.e., divide
by \(b^{k}\)):

\[
b^{x}b^{-k}=b^{x-k}.
\]

Since \(y=-k\), this is \(b^{x}b^{y}=b^{x+y}\).

Combining Steps A and B yields \(b^{x+y}=b^{x}b^{y}\) for all integers
\(x,y\).

\begin{enumerate}
\def\labelenumi{\arabic{enumi})}
\setcounter{enumi}{1}
\tightlist
\item
  Power-of-a-power rule: \((b^{x})^{y}=b^{xy}\) for all integers \(x,y\)
\end{enumerate}

\textbf{Case 1: \(y\ge 0\).} Induct on \(y\).

\begin{itemize}
\item
  \textbf{Base \(y=0\):} \((b^{x})^{0}=1=b^{0}=b^{x\cdot 0}\).
\item
  \textbf{Step:} If \((b^{x})^{n}=b^{xn}\), then

  \[
  (b^{x})^{n+1}=(b^{x})^{n}\,b^{x}=b^{xn}\,b^{x}=b^{x(n+1)},
  \]

  using the product rule just proved.
\end{itemize}

\textbf{Case 2: \(y<0\).} Write \(y=-k\) with \(k>0\). Then

\[
(b^{x})^{y}=(b^{x})^{-k}=\frac{1}{(b^{x})^{k}}
=\frac{1}{b^{xk}}=b^{-xk}=b^{xy}.
\]

This establishes \((b^{x})^{y}=b^{xy}\) for all integers \(x,y\).

\subsection{8 {[}25{]}}\label{section-23}

\begin{enumerate}
\def\labelenumi{\arabic{enumi})}
\tightlist
\item
  Existence via a monotone digit-by-digit construction
\end{enumerate}

We construct a nondecreasing sequence \((r_k)_{k\ge0}\) of decimal
truncations that converges to a number \(y\) with \(y^m=u\).

\textbf{Step 0 (integer part).} Let \(n\) be the unique integer with

\[
n^m \le u < (n+1)^m .
\]

This \(n\) exists because \(x\mapsto x^m\) is continuous and strictly
increasing on \([0,\infty)\), and it is unique because these intervals
are disjoint.

Set \(r_0:=n\).

\textbf{Inductive step (the next digit).} Suppose we have chosen

\[
r_k \;=\; n+\frac{d_1}{10}+\cdots+\frac{d_k}{10^k},
\]

with digits \(0\le d_j\le 9\), such that

\[
r_k^m \le u < \bigl(r_k+10^{-k}\bigr)^m. \tag{Inv\(_k\)}
\]

Choose \(d_{k+1}\in\{0,1,\dots,9\}\) to be the \textbf{largest} digit
for which

\[
\bigl(r_k + d_{k+1}10^{-(k+1)}\bigr)^m \le u .
\]

Then define \(r_{k+1}:=r_k + d_{k+1}10^{-(k+1)}\).

This choice always exists: take \(d=0\) (which satisfies the inequality)
and increase \(d\) until the inequality would fail---then stop at the
last valid digit. Because \(x\mapsto x^m\) is increasing, we also get

\[
r_{k+1}^m \le u < \bigl(r_{k+1}+10^{-(k+1)}\bigr)^m,
\]

so \textbf{Inv\(_{k+1}\)} holds. Thus \((r_k)\) is nondecreasing and
bounded above by \(r_k+10^{-k}\).

\textbf{Convergence.} By construction,

\[
0 \le y - r_k < 10^{-k} \quad\text{for } y:=\lim_{k\to\infty} r_k .
\]

Taking limits in \textbf{Inv\(_k\)} and using continuity of
\(x\mapsto x^m\),

\[
y^m \;=\; \lim_{k\to\infty} r_k^m \;=\; u .
\]

Hence a positive \(m\)-th root \textbf{exists} and its decimal digits
are determined successively by the rule above.

\begin{quote}
\textbf{Remark (how to \emph{compute} the test).} To decide whether
\(\bigl(r + t\bigr)^m \le u\) for
\(t\in\{0,\;10^{-(k+1)},\;2\cdot10^{-(k+1)},\dots\}\), you can either
(i) multiply out directly, or (ii) use the binomial expansion
\((r+t)^m=r^m+m r^{m-1}t+\binom{m}{2}r^{m-2}t^2+\cdots\) with additions
and multiplications only. Either way is valid for this
existence/constructibility proof.
\end{quote}

\begin{enumerate}
\def\labelenumi{\arabic{enumi})}
\setcounter{enumi}{1}
\tightlist
\item
  Uniqueness of the positive \(m\)-th root
\end{enumerate}

Suppose \(y_1>0\) and \(y_2>0\) with \(y_1^m=u=y_2^m\). If \(y_1<y_2\),
monotonicity of \(x^m\) on \([0,\infty)\) gives \(y_1^m<y_2^m\), a
contradiction. Thus \(y_1=y_2\). The positive \(m\)-th root is
\textbf{unique}.

(Our digit rule---``take the largest digit keeping
\((\cdot)^m\le u\)''---also ensures we get the standard decimal that
does \textbf{not} end in an infinite string of 9s.)

\begin{enumerate}
\def\labelenumi{\arabic{enumi})}
\setcounter{enumi}{2}
\tightlist
\item
  The requested digit-by-digit algorithm (summary)
\end{enumerate}

Given \(u>0\) and integer \(m\ge1\):

\begin{enumerate}
\def\labelenumi{\arabic{enumi}.}
\item
  \textbf{Find the integer part \(n\):} the unique \(n\) with
  \(n^m\le u<(n+1)^m\). Set \(r_0=n\).
\item
  \textbf{For \(k=0,1,2,\dots\):} Choose \(d_{k+1}\in\{0,\dots,9\}\)
  largest s.t.

  \[
  \bigl(r_k + d_{k+1}\cdot 10^{-(k+1)}\bigr)^m \le u,
  \]

  then set \(r_{k+1}=r_k + d_{k+1}\cdot 10^{-(k+1)}\).
\item
  The decimal \(y=\lim r_k=n.\,d_1d_2d_3\ldots\) satisfies \(y^m=u\),
  with error \(<10^{-k}\) after \(k\) digits.
\end{enumerate}

This proves existence (by construction), uniqueness (by monotonicity),
and gives the explicit method for successive digits that the exercise
asks for. 

\begin{enumerate}
\def\labelenumi{\arabic{enumi})}
\setcounter{enumi}{3}
\tightlist
\item
  Optional integer-only variant (nice for hand computation)
\end{enumerate}

Let \(A_k:=\lfloor 10^k r_k\rfloor\) (so \(A_0=n\)). Keep the scaled
target \(U:=u\). At stage \(k\), pick the largest \(d\in\{0,\dots,9\}\)
with

\[
(10A_k + d)^m \;\le\; \bigl\lfloor 10^{m(k+1)} U \bigr\rfloor .
\]

Then set \(A_{k+1}:=10A_k + d\). The digits are the base-10 digits of
\(A_k\). This is equivalent to the real-number method but uses only
integer comparisons.

\subsection{9 {[}M23{]}}\label{m23-2}

\begin{itemize}
\item
  For integers \(n\) and a fixed base \(b>0\), powers are defined by

  \[
  b^0=1,\qquad b^n=b^{n-1}\,b\ (n>0),\qquad b^n=b^{n+1}/b\ (n<0).
  \]

  From this, one proves for \textbf{integers} \(x,y\):
  \(b^{x+y}=b^x b^y\) and \((b^x)^y=b^{xy}\). 
\item
  For a \textbf{rational} \(r=p/q\) (with \(q>0\)), Knuth defines

  \[
  b^{p/q}:=\sqrt[q]{\,b^{p}\,}.
  \]

  Existence and uniqueness of the positive \(q\)-th root are noted just
  before (and used in) this definition.
\end{itemize}

The exercise asks you to \textbf{prove the exponent laws for rational
exponents}, assuming they already hold for integers.

Throughout, let \(b>0\) and let rationals be written in lowest terms
with positive denominators.

\begin{enumerate}
\def\labelenumi{\arabic{enumi})}
\tightlist
\item
  Product rule: \(b^{r+s}=b^r\,b^s\) for rationals \(r,s\)
\end{enumerate}

Write \(r=\dfrac{a}{m}\), \(s=\dfrac{c}{n}\) with \(m,n>0\). Let

\[
S:=b^{a/m}\,b^{c/n}.
\]

Raise both sides to the \textbf{integer} power \(mn\):

\[
S^{mn}=\bigl(b^{a/m}\bigr)^{mn}\,\bigl(b^{c/n}\bigr)^{mn}
= \bigl((b^{a/m})^{m}\bigr)^{n}\,\bigl((b^{c/n})^{n}\bigr)^{m}.
\]

By the \emph{definition} of rational powers, \((b^{a/m})^{m}=b^{a}\) and
\((b^{c/n})^{n}=b^{c}\); by the \textbf{integer} laws,
\((b^{a})^{n}=b^{an}\) and \((b^{c})^{m}=b^{cm}\). Hence

\[
S^{mn}=b^{an}\,b^{cm}=b^{an+cm}.
\]

Now set

\[
T:=b^{(an+cm)/(mn)}=b^{r+s}.
\]

By definition, \(T^{mn}=b^{an+cm}\). Thus \(S^{mn}=T^{mn}\), and both
\(S,T\) are positive; by \textbf{uniqueness} of the positive \(mn\)-th
root, \(S=T\). Therefore \(b^{r+s}=b^{r}\,b^{s}\).

\begin{enumerate}
\def\labelenumi{\arabic{enumi})}
\setcounter{enumi}{1}
\tightlist
\item
  Power-of-a-power: \(\bigl(b^{r}\bigr)^{s}=b^{rs}\) for rationals
  \(r,s\)
\end{enumerate}

Again write \(r=\dfrac{a}{m}\), \(s=\dfrac{c}{n}\) with \(m,n>0\). Let

\[
U:=\bigl(b^{a/m}\bigr)^{c/n}.
\]

By \textbf{definition} of the rational exponent with denominator \(n\),

\[
U^n=\bigl(b^{a/m}\bigr)^c.
\]

Raise both sides to the integer power \(m\):

\[
U^{mn}=\bigl((b^{a/m})^c\bigr)^{m}.
\]

Use the integer laws twice: first \((x^c)^m=x^{cm}\) for any positive
\(x\), then \((b^{a/m})^{m}=b^{a}\). Hence

\[
U^{mn}=(b^{a})^{c}=b^{ac}.
\]

Now set \(V:=b^{ac/(mn)}=b^{rs}\). Then \(V^{mn}=b^{ac}\). By uniqueness
of positive \(mn\)-th roots, \(U=V\). Thus
\(\bigl(b^{r}\bigr)^{s}=b^{rs}\).

\begin{enumerate}
\def\labelenumi{\arabic{enumi})}
\setcounter{enumi}{2}
\tightlist
\item
  Consequences: negatives, quotient, and zero
\end{enumerate}

\begin{itemize}
\tightlist
\item
  Taking \(s=-r\) in the product rule gives \(b^{r}b^{-r}=b^{0}=1\),
  hence \(b^{-r}=1/b^{r}\).
\item
  The quotient law follows: \(b^{r-s}=b^{r}/b^{s}\) (multiply both sides
  by \(b^{s}\)).
\item
  Since \(0=0/1\), the integer definition already gives \(b^{0}=1\). 
\end{itemize}

\subsection{10 {[}18{]}}\label{section-24}

Assume \(\log_{10} 2\) were rational. Then there exist integers \(p,q\)
with \(q>0\) and \(\gcd(p,q)=1\) such that

\[
\log_{10} 2=\frac pq.
\]

Exponentiating base \(10\) gives

\[
10^{p/q}=2\quad\Longrightarrow\quad 10^{p}=2^{q}.
\]

But \(10^{p}=(2\cdot 5)^{p}=2^{p}\,5^{p}\). Hence

\[
2^{p}\,5^{p}=2^{q}\quad\Longrightarrow\quad 5^{p}=2^{\,q-p}.
\]

The left side is a power of the prime \(5\); the right side is a power
of the prime \(2\). By the Fundamental Theorem of Arithmetic (unique
prime factorization), the only way a power of \(5\) can equal a power of
\(2\) is if both are \(1\), which would force \(p=0\) and then
\(2=10^{0}=1\), a contradiction.

Therefore the assumption was false, and \(\log_{10} 2\) is irrational.

\subsection{11 {[}10{]}}\label{section-25}

Let \(x_0=\log_{10}2\). Suppose we know \(x\) to \(d\) decimal places,
so \(|x-x_0|\le \tfrac12\cdot10^{-d}\). Write the error as
\(\delta=x-x_0\). Then

\[
10^{x}=10^{x_0+\delta}=10^{x_0}\,10^{\delta}=2\cdot 10^{\delta}.
\]

To get the first three decimal places of \(10^x\) correct, we need an
absolute error \(<\;0.0005\):

\[
|\,10^x-2\,|=2\,|\,10^{\delta}-1\,|<0.0005
\quad\Longrightarrow\quad
|\,10^{\delta}-1\,|<0.00025.
\]

Because \(10^{\delta}\) is increasing in \(\delta\), this is equivalent
to

\[
|\delta|<\log_{10}(1.00025).
\]

(Here we simply used the definition of logarithm; change of base was
introduced earlier in the section.) 

Numerically,

\[
\log_{10}(1.00025)\approx 0.00010856.
\]

Therefore we need

\[
\tfrac12\cdot 10^{-d}\;\le\;0.00010856
\;\Longleftrightarrow\;
10^{-d}\le 2.1712\times 10^{-4}
\;\Longleftrightarrow\;
d\gtrsim 3.66.
\]

Hence \(d=4\) decimal places suffice; \(d=3\) do not (a three-place
error bound \(\tfrac12\cdot10^{-3}\) would allow
\(|10^{\delta}-1|\approx (\ln 10)\cdot\tfrac12\cdot10^{-3}\approx 0.00115>0.00025\)).

\subsection{12 {[}02{]}}\label{section-26}

In §1.2.2 the logarithm to base \(b\) is introduced as the inverse of
exponentiation with base \(b\). Concretely, the identities stated around
Equations (8) are the inverse laws:

\[
\log_b(b^x)=x\quad\text{and}\quad b^{\log_b x}=x\quad(x>0,\; b>0,\; b\neq 1).
\]

Taking \(b=10\) in \(\log_b(b^x)=x\) gives

\[
\log_{10}(10^x)=x,
\]

which is exactly the form used to ``illustrate the so-called common
(base-10) logarithms'' at Equation (10). Hence (10) is an immediate
specialization of (8) with base \(10\).

\subsection{13 {[}M23{]}}\label{m23-3}

Proof 1 (via Bernoulli's inequality --- shortest)

Bernoulli's inequality says that for all real \(y>-1\) and integers
\(n\ge 0\),

\[
(1+y)^n \ge 1+ny.
\]

Apply it with \(y=\dfrac{x}{n}\). Then

\[
\Bigl(1+\frac{x}{n}\Bigr)^n \ \ge\ 1+n\cdot\frac{x}{n}\ =\ 1+x.
\]

Both sides are positive (since \(x>-1\)), and the function
\(t\mapsto t^{1/n}\) is strictly increasing on \(t>0\). Taking \(n\)th
roots preserves the inequality and gives

\[
\sqrt[n]{\,1+x\,}\ \le\ 1+\frac{x}{n}.
\]

Equality holds exactly when \(x=0\) or \(n=1\).

Proof 2 (concavity / tangent line method)

For \(0<r\le 1\) the function \(f(t)=(1+t)^r\) on \((-1,\infty)\) has

\[
f''(t)=r(r-1)(1+t)^{r-2}\le 0,
\]

so \(f\) is concave. A concave function lies below any tangent line; at
\(t=0\),

\[
f(0)=1,\qquad f'(0)=r.
\]

Hence \(f(t)\le f(0)+f'(0)\,t=1+rt\). With \(r=1/n\) and \(t=x\) we get

\[
(1+x)^{1/n}\le 1+\frac{x}{n},
\]

for every \(x>-1\). (This is the same inequality, proved analytically.)
For a reference to the standard Bernoulli form---which this concavity
argument generalizes---see the sources cited above.

\subsection{14 {[}15{]}}\label{section-27}

Prove Eq. (12), namely

\[
\log_b\!(c^{\,y}) \;=\; y\,\log_b c\qquad(c>0),
\]

with the usual assumptions \(b>0,\; b\neq 1\). 

\textbf{Definition of \(\log_b x\).} For \(x>0\) there is a unique real
number \(\log_b x\) such that

\[
b^{\log_b x}=x,
\]

because \(b^t\) is a one-to-one, onto mapping from \(\mathbb{R}\) to
\((0,\infty)\) when \(b>0,\; b\neq 1\).

\textbf{Law of exponents.} For all real \(u,v\) and all \(a>0\),

\[
(a^{u})^{v}=a^{uv}.
\]

(One proves this first for integers, then for rationals via roots, and
finally for reals by continuity/limits; this is the standard development
in §1.2.2.)

Now let \(v=\log_b c\). By definition, \(b^{v}=c\). Then

\[
c^{\,y}=(b^{v})^{y}=b^{vy}\qquad\text{(law of exponents).}
\]

Apply \(\log_b\) to both sides:

\[
\log_b(c^{\,y})=\log_b(b^{vy})=vy
= y\,\log_b c,
\]

since \(\log_b(b^{t})=t\) by definition of \(\log_b\). This is exactly
Eq. (12). ■

\subsection{15 {[}10{]}}\label{section-28}

Proof (from the laws of logarithms)

Knuth has already established the basic rules

\[
\log_b(xy)=\log_b x+\log_b y\quad(x>0,\ y>0),
\]

and

\[
\log_b(c^{\,t})=t\,\log_b c\quad(c>0),
\]

which follow from the exponent laws and the definition of \(\log_b\). 

Apply these to \(x/y=x\cdot y^{-1}\) (note \(y^{-1}>0\) since \(y>0\)):

\[
\log_b\!\left(\frac{x}{y}\right)
=\log_b\!\bigl(x\cdot y^{-1}\bigr)
=\log_b x+\log_b\!\left(y^{-1}\right)
=\log_b x+(-1)\log_b y
=\log_b x-\log_b y.
\]

Thus the identity is \textbf{true} for all \(x>0\), \(y>0\), with
\(b>0\), \(b\ne1\).

Equivalent one-liner (via the definition \(\;b^{\log_b u}=u\))

By definition, \(x=b^{\log_b x}\) and \(y=b^{\log_b y}\). Hence

\[
\frac{x}{y}=b^{\log_b x}\,b^{-\log_b y}=b^{\,\log_b x-\log_b y},
\]

so taking \(\log_b\) of both sides gives
\(\log_b(x/y)=\log_b x-\log_b y\).

\subsection{17 {[}05{]}}\label{section-29}

\begin{enumerate}
\def\labelenumi{\arabic{enumi}.}
\tightlist
\item
  \(\lg 32\)
\end{enumerate}

By definition \(\lg x := \log_2 x\). Since \(32=2^5\), we have
\(\lg 32=\log_2(2^5)=5\). 

\begin{enumerate}
\def\labelenumi{\arabic{enumi}.}
\setcounter{enumi}{1}
\tightlist
\item
  \(\log_{\pi} \pi\)
\end{enumerate}

In any (real) base \(b>0, b\neq 1\), \(\log_b b = 1\) because \(b^1=b\).
Thus \(\log_{\pi}\pi=1\). (This uses the standard real-log definition.) 

\begin{enumerate}
\def\labelenumi{\arabic{enumi}.}
\setcounter{enumi}{2}
\tightlist
\item
  \(\ln e\)
\end{enumerate}

By definition \(\ln x := \log_e x\). Hence \(\ln e=\log_e e=1\). 

\begin{enumerate}
\def\labelenumi{\arabic{enumi}.}
\setcounter{enumi}{3}
\tightlist
\item
  \(\log_b 1\)
\end{enumerate}

For any \(b>0, b\neq 1\), we have \(b^0=1\). Therefore \(\log_b 1=0\).
(Again using the real-log definition with positive argument.) 

\begin{enumerate}
\def\labelenumi{\arabic{enumi}.}
\setcounter{enumi}{4}
\tightlist
\item
  \(\log_b(-1)\)
\end{enumerate}

In the (real) setting used in this section, logarithms are defined only
for positive arguments \(x>0\) with \(b>0, b\neq 1\). Since \(-1\) is
not positive, \(\log_b(-1)\) has \textbf{no real value} (is undefined
over the reals).

\subsection{18 {[}10{]}}\label{section-30}

Solution (disprove)

Work for \(x>0\). By definition, \(\log_8 x = y\) iff \(8^y=x\). Since
\(8=2^3\), we have

\[
8^y=(2^3)^y=2^{3y}=x.
\]

Take binary logarithms (\(\lg\)) of both sides and use the law
\(\log_b(c^y)=y\log_b c\):

\[
\lg x=\lg\bigl(2^{3y}\bigr)=3y \quad\Rightarrow\quad y=\frac{1}{3}\lg x.
\]

Therefore

\[
\boxed{\;\log_8 x=\tfrac13\,\lg x\;}
\]

---not \(\tfrac12\,\lg x\). (The used law \(\log_b(c^y)=y\log_b c\) is
given in Eq. (12).) 

A concrete counterexample makes the falsehood clear: take \(x=8\). Then
\(\log_8 8=1\), but \(\tfrac12\,\lg 8=\tfrac12\cdot 3=1.5\neq 1\).

When would \(\log_8 x=\tfrac12\lg x\) hold?

Only when \(\tfrac13\lg x=\tfrac12\lg x\), i.e., when \(\lg x=0\), so
\(x=1\).

Remark (generalization)

More generally, for any \(k>0\),

\[
\log_{2^k} x=\frac{1}{k}\,\lg x,
\]

and by change of base,

\[
\log_a x=\frac{\lg x}{\lg a}.
\]

The incorrect \(\tfrac12\) would be correct for base \(4=2^2\):
\(\log_4 x=\tfrac12\,\lg x\).

\subsection{19 {[}20{]}}\label{section-31}

Answer

Yes --- every 14-digit integer fits.

A 14-digit positive integer \(n\) satisfies

\[
10^{13}\le n\le 10^{14}-1<10^{14}.
\]

Use the classic tight comparison \(10^3<2^{10}\). Then

\[
10^{14}=10^{12}\cdot 10^{2}<(2^{10})^{4}\cdot 100
=2^{40}\cdot 100
<2^{40}\cdot 2^{7}=2^{47},
\]

since \(100<128=2^{7}\). Hence \(n<2^{47}\).

A machine with 47 magnitude bits (plus a sign bit) can represent all
integers in the range \(-\,(2^{47}-1)\) to \(+(2^{47}-1)\) in
sign-magnitude format (and \([-2^{47},2^{47}-1]\) if it were
two's-complement with 48 total bits). Either way, every \(n<2^{47}\)
(and its negative) fits.

Alternative view (bit-length via logs)

The number of bits needed for the largest 14-digit number is

\[
\left\lceil \log_2(10^{14}-1)\right\rceil
< \left\lceil \log_2 10^{14}\right\rceil
=\left\lceil 14\log_2 10 \right\rceil
=\left\lceil 14\times 3.321928\ldots \right\rceil
=\lceil 46.5069\ldots\rceil
=47.
\]

So 47 value bits suffice; the extra sign bit handles negativity.

\subsection{20 {[}10{]}}\label{section-32}

Each is the reciprocal of the other:

\[
\boxed{\ \log_{10}2=\frac1{\log_{2}10}\ } \qquad\text{and}\qquad
\boxed{\ \log_{2}10=\frac1{\log_{10}2}\ }.
\]

The section derives the change-of-base rule

\[
\log_{c}x=\frac{\log_{b}x}{\log_{b}c},
\]

which follows from the basic logarithm identities
\(\log_b(xy)=\log_b x+\log_b y\) and \(\log_b(c^y)=y\log_b c\). 

Apply that rule twice:

\begin{enumerate}
\def\labelenumi{\arabic{enumi}.}
\tightlist
\item
  Take \(b=10,\; c=2,\; x=10\):
\end{enumerate}

\[
\log_{2}10=\frac{\log_{10}10}{\log_{10}2}=\frac{1}{\log_{10}2}.
\]

\begin{enumerate}
\def\labelenumi{\arabic{enumi}.}
\setcounter{enumi}{1}
\tightlist
\item
  Equivalently, take \(b=2,\; c=10,\; x=2\):
\end{enumerate}

\[
\log_{10}2=\frac{\log_{2}2}{\log_{2}10}=\frac{1}{\log_{2}10}.
\]

Thus the ``simple relation'' is also

\[
\log_{10}2\cdot\log_{2}10=1.
\]

\subsection{21 {[}15{]}}\label{section-33}

To keep everything real-valued we need \(b>1\) and \(x>1\): then
\(\log_b x>0\) (so the outer log is defined) and \(\ln\ln b,\ \ln\ln x\)
both exist.

Use the change-of-base identity (Eq. (14) in §1.2.2):

\[
\log_c u \;=\; \frac{\log_b u}{\log_b c} \;=\; \frac{\ln u}{\ln c}.
\]

Apply it twice:

\[
\begin{aligned}
\log_b(\log_b x)
&= \log_b\!\left(\frac{\ln x}{\ln b}\right) && \text{(change base)}\\
&= \log_b(\ln x)\;-\;\log_b(\ln b) && \text{(log of a quotient)}\\
&= \frac{\ln(\ln x)}{\ln b}\;-\;\frac{\ln(\ln b)}{\ln b} && \text{(change base again)}\\[4pt]
&= \boxed{\;\dfrac{\ln\ln x - \ln\ln b}{\ln b}\; }.
\end{aligned}
\]

\subsection{22 {[}20{]}}\label{section-34}

By change of base, \(\lg x=\dfrac{\ln x}{\ln 2}\) and
\(\log_{10}x=\dfrac{\ln x}{\ln 10}\). 

Let

\[
A(x)\;=\;\ln x+\log_{10}x \;=\; \ln x\!\left(1+\frac{1}{\ln 10}\right).
\]

Then, for any \(x\neq 1\) (the case \(x=1\) is trivial since both sides
are \(0\)),

\[
\frac{A(x)}{\lg x}
=\frac{\ln x\left(1+\frac{1}{\ln 10}\right)}{\ln x/\ln 2}
=\ln 2\!\left(1+\frac{1}{\ln 10}\right)
=:R.
\]

Notice \(R\) is a \textbf{constant}, independent of \(x\). Hence the
relative error of the approximation \(A(x)\) to \(\lg x\) is

\[
\frac{|\lg x-A(x)|}{|\lg x|}
=|1-R|.
\]

Use the simple bounds \(\,0.693<\ln 2<0.694\) and
\(2.3025<\ln 10<2.3026\). Then

\[
R\in\Bigl[\,0.693\Bigl(1+\frac{1}{2.3026}\Bigr),\;
0.694\Bigl(1+\frac{1}{2.3025}\Bigr)\Bigr]
\subset (0.9939,\;0.9955),
\]

so \(|1-R|<0.0061<0.01\). Thus

\[
\boxed{\;\lg x \approx \ln x+\log_{10}x\;\text{ with }<1\%\text{ relative error.}\;}
\]

\subsection{23 {[}M25{]}}\label{m25-2}

Geometric proof (area-preserving similarity)

Let \(\mathcal{A}(a,b)\) denote the area under \(y=1/t\) from \(t=a\) to
\(t=b\) (with \(0<a<b\)). By the definition above,
\(\ln x=\mathcal{A}(1,x)\).

Fix some \(x>0\). Consider the \textbf{similarity transformation}

\[
\Phi_x:\ (u,v)\ \longmapsto\ (t,w)=(xu,\; v/x).
\]

This scales the horizontal axis by a factor \(x\) and the vertical axis
by a factor \(1/x\). Its Jacobian determinant is

\[
\det\begin{pmatrix}
\partial t/\partial u & \partial t/\partial v\\
\partial w/\partial u & \partial w/\partial v
\end{pmatrix}
=\det\begin{pmatrix}
x & 0\\[2pt]
0 & 1/x
\end{pmatrix}
=1,
\]

so \textbf{areas are preserved} by \(\Phi_x\).

Now check what \(\Phi_x\) does to the curve \(v=\frac1u\):

\[
v=\frac1u\quad\Longrightarrow\quad
w=\frac{v}{x}=\frac{1}{xu}=\frac{1}{t},
\]

since \(t=xu\). Thus \(\Phi_x\) maps the hyperbola \(v=1/u\)
\textbf{onto itself}.

Therefore, \(\Phi_x\) maps the region under \(v=1/u\) from \(u=1\) to
\(u=y\) \textbf{bijectively and area-preservingly} onto the region under
\(w=1/t\) from \(t=x\) to \(t=xy\). In our notation,

\[
\mathcal{A}(1,y)\;=\;\mathcal{A}(x,xy).
\]

Finally decompose the larger interval \([1,xy]\) at \(x\):

\[
\mathcal{A}(1,xy)\;=\;\mathcal{A}(1,x)+\mathcal{A}(x,xy)
\;=\;\mathcal{A}(1,x)+\mathcal{A}(1,y).
\]

Using \(\ln z=\mathcal{A}(1,z)\), this is exactly

\[
\boxed{\ \ln(xy)=\ln x+\ln y\ }.
\]

(Optional) Same idea as a one-line integral substitution

The same geometry is the substitution \(t=xs\) in the area integral:

\[
\ln(xy)=\int_{1}^{xy}\frac{dt}{t}
=\int_{1}^{x}\frac{dt}{t}+\int_{x}^{xy}\frac{dt}{t}
=\ln x+\int_{1}^{y}\frac{ds}{s}
=\ln x+\ln y,
\]

where the middle equality is exactly the area-preserving stretch by
\(x\) horizontally and \(1/x\) vertically (the Jacobian is 1). This
matches Fig. 6's geometric interpretation.

\subsection{24 {[}15{]}}\label{section-35}

What the original method does (base 10, summarized)

Knuth shows how to normalize \(x\) so that \(1\le x/10^{n}<10\) (this
picks out the integer part \(n=\lfloor \log_{10}x\rfloor\)), then
iteratively squares and (when needed) divides by 10 to decide the binary
digits \(b_1,b_2,\dots\) of the fractional part of \(\log_{10}x\):

\[
\log_{10}x \;=\; n + \frac{b_1}{2}+\frac{b_2}{4}+\frac{b_3}{8}+\cdots,
\]

with the loop invariant

\[
1 \;\le\; \frac{x^{2^k}}{10^{\,2^k\!\left(n+b_1/2+\cdots+b_k/2^k\right)}} \;<\; 10.
\]

This is the content of Eqs. (17)--(19) and the paragraph introducing
them.

Modify it for base 2

Simply replace every ``10'' by ``2'' and ``decimal shift'' by ``binary
shift.''

\textbf{Normalization.} Choose the unique integer \(n\) with
\(1 \le x/2^n < 2\); set \(x_0 := x/2^n\). Then
\(n=\lfloor \lg x\rfloor\).

\textbf{Iteration.} For \(k=1,2,\dots\),

\begin{enumerate}
\def\labelenumi{\arabic{enumi}.}
\tightlist
\item
  Compute \(t := x_{k-1}^2\).
\item
  If \(t < 2\) set \(b_k := 0\) and \(x_k := t\); else set \(b_k := 1\)
  and \(x_k := t/2\).
\end{enumerate}

(This is exactly the same control flow as the base-10 version, just with
the radix ``2.'')

\textbf{Output.} The algorithm yields

\[
\lg x \;=\; n + \frac{b_1}{2}+\frac{b_2}{4}+\frac{b_3}{8}+\cdots,
\]

i.e., the fractional bits of \(\lg x\) are \(b_1b_2b_3\ldots\) in base
2.

\textbf{Why it works (invariant proof).} By construction \(1\le x_k<2\)
for all \(k\). Unrolling the recurrence shows

\[
x_k \;=\; \frac{\,x^{2^k}\,}{2^{\,2^k\!\left(n+\sum_{j=1}^k b_j/2^j\right)}}.
\]

Thus the test ``\(t<2\)?'' is equivalent to asking whether

\[
\frac{x^{2^k}}{2^{\,2^k\!\left(n+\sum_{j=1}^{k-1} b_j/2^j + 1/2^k\right)}} \;<\; 1
\quad\text{or}\quad \ge 1,
\]

i.e., whether the next binary digit \(b_k\) of \(\lg x\) is \(0\) or
\(1\). Therefore after \(k\) steps we have fixed the first \(k\)
fractional bits. This is identical to the correctness argument Knuth
gives for base 10, specialized to base 2. 

Small example: \(\lg 3\)

Normalize: \(1\le 3/2^1=1.5<2\) so \(n=1\), \(x_0=1.5\).

\begin{itemize}
\tightlist
\item
  \(k=1\): \(t=1.5^2=2.25\ge 2\Rightarrow b_1=1\), \(x_1=2.25/2=1.125\).
\item
  \(k=2\): \(t=1.125^2=1.265625<2\Rightarrow b_2=0\), \(x_2=1.265625\).
\item
  \(k=3\): \(t\approx1.6018066<2\Rightarrow b_3=0\),
  \(x_3\approx1.6018066\).
\item
  \(k=4\): \(t\approx2.56579\ge2\Rightarrow b_4=1\),
  \(x_4\approx1.282895\). (\ldots and so on.)
\end{itemize}

So \(\lg 3 \approx 1 + 1/2 + 0/4 + 0/8 + 1/16 + \cdots = 1.58496\ldots\)
--- which matches the true value to the shown precision.

\subsection{25 {[}22{]}}\label{section-36}

Why it terminates (finite-precision binary)

Work in fixed binary with \(p\) fraction bits (i.e., resolution
\(2^{-p}\)). ``Shift right 1'' divides by 2 and truncates one low-order
bit; after \(j\) successive shifts, \(z\) is \(\approx x/2^{\,j}\)
(again truncated).

\begin{itemize}
\tightlist
\item
  During L3 we keep halving \(z\) until \(x-z\ge 1\). Since after \(p\)
  shifts we have \(z\le x/2^{p}<2/2^{p}=2^{1-p}<1\), we certainly get
  \(x-z\ge1\) by or before \(j=p\). Thus L3 always finishes with some
  \(k\le p\).
\item
  In L4 we subtract a \textbf{positive} amount \(z\ge 2^{-p}\).
  Therefore \(x\) strictly decreases by at least one ulp. Because
  initially \(1\le x<2\), there are at most
  \(\lceil (x_0-1)\,2^{p}\rceil\) such decreases before \(x\) becomes
  exactly \(1\) in this arithmetic, triggering L2. Hence the algorithm
  halts after \(O(p)\) reductions, i.e., a number of simple
  shifts/adds/subtracts proportional to the desired places of accuracy.
  (Knuth notes that the deliberate truncations from shifting ensure
  eventual reduction to 1.)
\end{itemize}

Why it computes \(\log_b x\)

Ignore truncation for a moment (idealized arithmetic) to see the
invariant:

\begin{itemize}
\item
  In L3 we choose the \textbf{smallest} \(k\) with
  \(x-\frac{x}{2^k}\ge1\), i.e.

  \[
  x\ \ge\ \frac{1}{1-2^{-k}}\ =\ \frac{2^k}{2^k-1}\ \stackrel{\mathrm{def}}{=}\ a_k.
  \]
\item
  In L4 we then replace \(x\) by
  \(x':=x-\frac{x}{2^k}=x\,(1-2^{-k})=\dfrac{x}{a_k}\) and add to \(y\)
  the term \(\log_b a_k\).
\end{itemize}

Thus each reduction step maintains the invariant

\[
x\cdot b^{-y}\ \ \text{stays constant}.
\]

After a sequence of \(k_1,k_2,\dots,k_m\) we have

\[
y=\sum_{i=1}^m\log_b a_{k_i}\ =\ \log_b\!\Bigl(\prod_{i=1}^m a_{k_i}\Bigr),
\qquad
x=\frac{x_0}{\prod_{i=1}^m a_{k_i}}.
\]

If we could carry on ``forever,'' \(x\) would tend to \(1\) and \(y\)
would tend to \(\log_b x_0\).

With finite precision the only difference is that ``shift right''
underestimates \(\frac{x}{2^k}\) by \textless{} one ulp. Hence each step
reduces \(x\) by \textbf{nearly} \(x/2^k\), so the invariant becomes

\[
x\cdot b^{-y}\ =\ 1+\varepsilon,
\]

where \(\varepsilon\) is a small multiplicative error caused by
truncation (bounded in terms of machine precision and the number of
steps). When the loop stops at \(x=1\), we have

\[
y = \log_b\!\Bigl(\frac{x_0}{1+\varepsilon}\Bigr)
= \log_b x_0 - \log_b(1+\varepsilon),
\]

so the absolute error is
\(|\log_b x_0 - y|=\log_b(1+\varepsilon)\approx \varepsilon/\ln b\). In
other words, \(y\) is a bona fide approximation to \(\log_b x_0\), with
an accuracy governed by the working precision. (Exercise 26 in the book
asks for a rigorous bound tied to the rounding model.)

A tiny worked trace (to see the mechanism)

Take \(b=e\) and \(x_0=1.5\) (exact arithmetic, for intuition).

\begin{itemize}
\tightlist
\item
  Start: \(y=0,\ z=0.75,\ k=1\). Since \(1.5-0.75<1\), halve
  \(z\Rightarrow z=0.375,k=2\). Now \(1.5-0.375=1.125\ge1\).
\item
  Reduce: \(x\gets1.125\), \(y\gets y+\ln(4/3)\).
\item
  Recompute \(z=x/2^k=1.125/4=0.28125\). Since \(1.125-0.28125<1\),
  halve \(z\to 0.140625,k=3\); still \(<1\), halve again
  \(z\to 0.0703125,k=4\); now \(1.125-0.0703125=1.0546875\ge1\).
\item
  Reduce: \(x\gets1.0546875\), \(y\gets y+\ln(16/15)\).
\end{itemize}

\subsection{26 {[}M27{]}}\label{m27}

What the algorithm does (very briefly)

Given a binary machine number \(x\in[1,2)\) and a base \(b>0,\;b\neq1\),
Ex. 25 proposes the loop

\begin{itemize}
\tightlist
\item
  choose \(k\ge1\) so that with
  \(z=\text{(}x\text{ shifted right }k\text{)}\) we still have
  \(x-z\ge1\);
\item
  replace \(x\leftarrow x-z\) and add
  \(\log_b\!\bigl(2^k/(2^k-1)\bigr)\) to the running sum \(y\);
\end{itemize}

and repeat until \(x=1\). (See the statement on p.~26 of the book.)

Error model (binary fixed-point)

Assume the computer works in binary fixed-point with \(t\) fraction
bits. Let

\[
u \stackrel{\text{def}}{=} 2^{-t}
\]

be one unit in the last place (1 ulp). A right shift by \(k\) computes

\[
z \;=\; \Big\lfloor\frac{x}{2^{\,k}}\Big\rfloor
\]

so the exact division and the shifted value differ by less than one ulp:

\[
0 \;\le\; \frac{x}{2^{\,k}}-z \;<\; u\quad\text{for every step.}
\]

Two key facts

\begin{enumerate}
\def\labelenumi{\arabic{enumi}.}
\tightlist
\item
  \textbf{Monotone overestimate at each step.} When the loop subtracts
  \(z\) and adds \(\log_b\!\bigl(2^k/(2^k-1)\bigr)\), it is pretending
  we subtracted the exact fraction \(x/2^{\,k}\). Since
  \(z\le x/2^{\,k}\), we subtract \textbf{no more than} that exact
  amount, hence the running sum \(y\) is an \textbf{overestimate} of
  \(\log_b\) of the original input. Thus
\end{enumerate}

\[
y \;\ge\; \log_b x_{\text{input}}.
\]

\begin{enumerate}
\def\labelenumi{\arabic{enumi}.}
\setcounter{enumi}{1}
\tightlist
\item
  \textbf{A one-shot bound from the last step.} On the \emph{final}
  iteration we have \(x-z=1\), so \(z=x-1\). Because
  \(z=\lfloor x/2^{\,k}\rfloor\),
\end{enumerate}

\[
x-1 \;\le\; \frac{x}{2^{\,k}} \;<\; x-1+u
\quad\Longrightarrow\quad
1-u \;<\; x-\frac{x}{2^{\,k}} \;\le\; 1.
\]

If we imagine doing the \emph{same} sequence of \(k\)'s but with
\textbf{exact} divisions, the identity behind the algorithm says

\[
\log_b x_{\text{input}} \;=\; y_{\text{(no roundoff in sums)}}\;+\;\log_b x_{\text{final}}^{\ast},
\]

where \(x_{\text{final}}^{\ast}\) is what the last exact subtraction
would leave. The inequality above shows

\[
1-u \;\le\; x_{\text{final}}^{\ast}\;\le\;1.
\]

Therefore the (systematic) overestimate is bounded by

\[
0 \;\le\; y_{\text{(no roundoff in sums)}}-\log_b x_{\text{input}}
\;=\; -\,\log_b x_{\text{final}}^{\ast}
\;\le\; -\,\log_b(1-u).
\]

Final bound (and a handy simplification)

Let \(u=2^{-t}\) be the machine ulp. Then a rigorous upper bound on the
error of the algorithm in Ex. 25 is

\[
\boxed{\quad 0 \;\le\; y-\log_b x_{\text{input}} \;\le\; -\log_b(1-2^{-t}).\quad}
\]

Using \(-\ln(1-\delta)\le \dfrac{\delta}{1-\delta}\) and
\(\log_b z=\dfrac{\ln z}{\ln b}\),

\[
-\log_b(1-2^{-t})
\;\le\;
\frac{\,2^{-t}}{(1-2^{-t})\,\ln b}
\;<\;
\frac{2^{-(t-1)}}{\ln b}
\qquad(t\ge1).
\]

Thus the error is \(O(2^{-t})\) with an explicit constant depending only
on the logarithm base \(b\).

If the \emph{sums} are rounded too

If each addition \(y\leftarrow y+\log_b(2^k/(2^k-1))\) is performed in a
\(p\)-bit accumulator (round-to-nearest), each add contributes at most
\(2^{-p-1}\) in \(\log_b\)-units. With \(N\) additions,

\[
0 \;\le\; y_{\text{computed}}-\log_b x_{\text{input}}
\;\le\; -\log_b(1-2^{-t}) \;+\; N\cdot 2^{-p-1}.
\]

(Here \(N\) is the number of times step L4 executes; it is proportional
to the digits of accuracy requested, as noted in the exercise
statement.)

\subsection{27 {[}M30{]}}\label{m30-2}

Setup (what the algorithm does)

Let \(x\) be a positive real number. The method writes

\[
\log_{10} x \;=\; n \;+\; \frac{b_1}{2} + \frac{b_2}{4} + \frac{b_3}{8} + \cdots,
\]

where we first normalize \(x_0 = x/10^n \in [1,10)\), then iterate for
\(k\ge 1\):

\[
b_k =
\begin{cases}
0,& \text{if } x_{k-1}^2 < 10,\\[2pt]
1,& \text{if } x_{k-1}^2 \ge 10,
\end{cases}
\qquad
x_k \;=\;
\begin{cases}
x_{k-1}^2,& \text{if } b_k=0,\\
x_{k-1}^2/10,& \text{if } b_k=1,
\end{cases}
\]

so that always \(x_k\in[1,10)\). In practice we cannot compute exactly;
Knuth explicitly says we \textbf{round or truncate} each \(x_k\) to a
fixed number \(t\) of significant decimal figures.

Exercise 27 asks you to \textbf{derive a worst-case upper bound} for the
error in the resulting value of \(\log_{10} x\) caused by doing that
finite-precision rounding/truncation each step. (Knuth hints at the
outcome and later notes that with \(t=4\) the error is guaranteed to be
\(<0.00044\).)

Work in log space

Write \(\ell_k=\log_{10} x_k\) for the exact process and
\(\tilde\ell_k=\log_{10}\tilde x_k\) for the computed
(rounded/truncated) values. The exact recurrence is

\[
\ell_k \;=\; 2\,\ell_{k-1} - b_k, \qquad \ell_k\in[0,1)\,.
\]

Finite precision introduces an \textbf{additive} error in log space:

\[
\tilde\ell_k \;=\; 2\,\tilde\ell_{k-1} - \tilde b_k \;+\; \varepsilon_k,
\]

where \(\varepsilon_k\) is the log10-error created when we
round/truncate \(x_k\) into \(\tilde x_k\).

Because after normalization \(x_k\in[1,10)\), rounding/truncation to
\(t\) significant figures perturbs \(x_k\) by at most \(\pm 10^{\,1-t}\)
in absolute value. Therefore the \textbf{per-step log10 error} is
bounded by

\[
|\varepsilon_k|
\;\le\;
\log_{10}\!\left(1 + 10^{\,1-t}\right)
\;<\;
\frac{10^{\,1-t}}{\ln 10}\,,
\]

using \(\ln(1+u)\le u\). (We use the larger truncation bound; rounding
alone would replace \(10^{1-t}\) by \(0.5\cdot 10^{1-t}\).)

Relate digit errors to the \(\varepsilon_k\)

Unfold the two recurrences up to step \(k\) and subtract. One obtains

\[
\tilde\ell_k - \ell_k
\;=\;
-\,2^{\,k}\bigl(\tilde S_k - S_k\bigr)
\;+\;
\sum_{j=1}^{k} 2^{\,k-j}\,\varepsilon_j,
\]

where \(S_k=\sum_{j=1}^k b_j/2^{\,j}\) and
\(\tilde S_k=\sum_{j=1}^k \tilde b_j/2^{\,j}\) are the exact and
computed partial binary sums. Since both \(\ell_k\) and \(\tilde\ell_k\)
lie in \([0,1)\), \(|\tilde\ell_k-\ell_k|<1\), hence

\[
\bigl|\tilde S_k - S_k\bigr|
\;\le\;
2^{-k}
\;+\;
\sum_{j=1}^{k} 2^{-j}\,|\varepsilon_j|.
\]

Letting \(k\to\infty\) (so \(2^{-k}\to0\) and \(S_k\to S\),
\(\tilde S_k\to \tilde S\)) gives the \textbf{final error bound}

\[
\boxed{\;
\bigl|\tilde S - S\bigr|
\;\le\;
\sum_{j=1}^{\infty} 2^{-j}\,|\varepsilon_j|
\;\le\;
\log_{10}\!\left(1 + 10^{\,1-t}\right)
\;<\;
\frac{10^{\,1-t}}{\ln 10}\; }.
\]

Plug in \(t=4\)

With four significant figures,

\[
\frac{10^{\,1-4}}{\ln 10} \;=\; \frac{10^{-3}}{\ln 10} \;\approx\; 0.000434294\;<\;0.00044,
\]

exactly matching Knuth's stated guarantee. (His worked example happens
to be even better because the first few \(x_k\) were exact before
rounding error accumulated.)

\subsection{28 {[}M30{]}}\label{m30-3}

Here is a clean way to do \textbf{\(b^x\) for \(0\le x<1\)} using only
shifting, addition, and subtraction, mirroring the ``shift-subtract''
logarithm method from §1.2.2 Ex. 25.

\#1. Idea (invert Ex. 25)

Exercise 25 shows how to reduce a value \(x\in[1,2)\) toward \(1\) by
repeatedly replacing

\[
x\leftarrow x-\frac{x}{2^k}=x\Bigl(1-\frac1{2^k}\Bigr)
\]

and adding the constant \(\log_b\!\bigl(\tfrac{2^k}{2^k-1}\bigr)\) to
the running logarithm; it needs only shifts, adds, subtracts, plus a
small table of constants \(\log_b\!\bigl(\tfrac{2^k}{2^k-1}\bigr)\). 

To go the other way (compute the \textbf{exponential}), write the
exponent \(x\) itself as a sum of those same constants:

\[
x = \sum_i c_{k_i},\qquad c_k\stackrel{\text{def}}{=}\log_b\!\Bigl(\frac{2^k}{2^k-1}\Bigr).
\]

Then

\[
b^x=\prod_i b^{c_{k_i}}=\prod_i \frac{2^{k_i}}{2^{k_i}-1}.
\]

Each factor \(\frac{2^{k}}{2^{k}-1}\) is a \textbf{shift left by \(k\)}
(multiply by \(2^k\)) followed by an \textbf{integer division by
\((2^k-1)\)}---both doable with shifts and add/sub (long division).

\#2. Algorithm (Feynman's product)

Precompute a short table \(c_k=\log_b\!\bigl(\tfrac{2^k}{2^k-1}\bigr)\)
for \(k=1,2,\dots,K\). (Exactly the same kind of table Ex. 25 requires.
)

\textbf{Input:} base \(b>0,\;b\ne1\); exponent \(x\in[0,1)\); tolerance
\(\varepsilon\).

\textbf{Output:} \(y\approx b^x\).

\textbf{Invariant:} after every step, \(y\cdot b^r=b^x\) where \(r\) is
the remaining (unaccounted) exponent.

\begin{itemize}
\item
  E1. Initialize \(y\leftarrow 1,\; r\leftarrow x,\; k\leftarrow1\).
\item
  E2. If \(r\le \varepsilon\), \textbf{stop} and return \(y\).
\item
  E3. While \(r<c_k\), set \(k\leftarrow k+1\). (Because \(c_k\)
  decreases rapidly with \(k\).)
\item
  E4. Update the product using shifts only:

  \begin{itemize}
  \tightlist
  \item
    \(y\leftarrow y\cdot 2^k\) (shift left by \(k\)),
  \item
    \(y\leftarrow y\div(2^k-1)\) (binary long division: shift/subtract),
  \item
    \(r\leftarrow r-c_k\);
  \item
    Go to E2.
  \end{itemize}
\end{itemize}

\textbf{Correctness.} Step E4 multiplies \(y\) by
\(2^k/(2^k-1)=b^{c_k}\); reducing \(r\) by \(c_k\) maintains
\(y\,b^r=b^x\). When \(r\le\varepsilon\), the returned \(y\) differs
from \(b^x\) by a factor \(b^{r}\in[1,b^\varepsilon]\).

\#3. How accurate is it?

Two effects matter:

\textbf{(A) Truncation of the exponent sum.} If you stop when
\(r\le\varepsilon\), then

\[
\frac{|b^x-y|}{b^x}\;=\;|b^r-1|\;\le\;e^{(\ln b)\,\varepsilon}-1
\;\approx\;(\ln b)\,\varepsilon\quad(\text{for small }\varepsilon).
\]

Thus, to get about \(p\) \textbf{bits} of relative accuracy, it suffices
to drive

\[
\varepsilon \;\lesssim\; \frac{2^{-p}}{\ln b}.
\]

Because
\(c_k=\log_b\!\bigl(\tfrac{2^k}{2^k-1}\bigr)\sim \tfrac{1}{(\ln b)(2^k-1)}\),
the greedy selection makes \(k\) grow, and you typically need only
\(O(p)\) factors.

\textbf{(B) Finite-precision arithmetic in the divisions.} Each division
by \((2^k-1)\) performed with \(w\)-bit fixed-point long division incurs
at most \(O(2^{-w})\) \textbf{relative} error. After \(T\) factors the
cumulative relative error is \(\lesssim T\,2^{-w}\) (first-order). Pick
\(w\) a few bits larger than the target precision and the truncation
term (A) will dominate.

\#4. Tiny example (base-10, four-digit feel)

Take \(b=10\) and \(x\approx \log_{10}2\approx 0.3010\). Using the table
\(c_2=\log_{10}\tfrac43\approx0.12494,\;c_4=\log_{10}\tfrac{16}{15}\approx0.02803,\;
c_5\approx0.01355,\;c_6\approx0.00679,\;c_8\approx0.00174,\;c_9\approx0.000903,\dots\)
and applying E2--E4 greedily gives

\[
10^{0.3010}
\approx \Bigl(\frac43\Bigr)^2\!\cdot\frac{16}{15}\cdot\frac{32}{31}\cdot\frac{64}{63}\cdot
\frac{256}{255}\cdot\frac{512}{511}\cdots
\approx 1.9992\quad(\text{within }4\times10^{-4}\text{ of }2).
\]

All updates are ``shift by \(k\) and divide by \(2^k-1\)''; the
constants come straight from the exercise list in §1.2.2.

\subsection{29 {[}HM20{]}}\label{hm20}

Use the base-change identity \(\log_b x=\dfrac{\ln x}{\ln b}\). Since
\(x>1\) is fixed, \(\ln x>0\), so minimizing

\[
b\log_b x = b\cdot\frac{\ln x}{\ln b}
\]

is equivalent to minimizing \(f(b)=\dfrac{b}{\ln b}\) over the stated
domain. Similarly, minimizing \((b+1)\log_b x\) is equivalent to
minimizing \(g(b)=\dfrac{b+1}{\ln b}\).

\begin{enumerate}
\def\labelenumi{(\alph{enumi})}
\tightlist
\item
  Real \(b>1\) minimizing \(b\log_b x\)
\end{enumerate}

Let \(f(b)=\dfrac{b}{\ln b}\) for \(b>1\). Then

\[
f'(b)=\frac{\ln b - 1}{(\ln b)^2}.
\]

Hence \(f'(b)=0 \iff \ln b=1 \iff b=e\). Moreover,

\begin{itemize}
\tightlist
\item
  \(f'(b)<0\) for \(1<b<e\) and
\item
  \(f'(b)>0\) for \(b>e\),
\end{itemize}

so \(b=e\) is the unique minimizer. The minimum value is

\[
b\log_b x\big|_{b=e} = e\cdot \ln x .
\]

\textbf{Answer (a):} \(b=e\), with minimum \(e\ln x\).

\begin{enumerate}
\def\labelenumi{(\alph{enumi})}
\setcounter{enumi}{1}
\tightlist
\item
  Integer \(b>1\) minimizing \(b\log_b x\)
\end{enumerate}

From (a), \(f(b)\) decreases on \((1,e]\) and increases on
\([e,\infty)\). Thus among integers \(b>1\), only \(b=2\) and \(b=3\)
can be candidates, and \(f(3)<f(2)\) because we are past the minimum at
\(e\).

Equivalently, compare \(\dfrac{3}{\ln 3}\) and \(\dfrac{2}{\ln 2}\):
\(\dfrac{3}{\ln 3} < \dfrac{2}{\ln 2} \iff 3\ln 2 < 2\ln 3 \iff 2^3<3^2\),
i.e., \(8<9\), true.

\textbf{Answer (b):} \(b=3\).

\begin{enumerate}
\def\labelenumi{(\alph{enumi})}
\setcounter{enumi}{2}
\tightlist
\item
  Integer \(b>1\) minimizing \((b+1)\log_b x\)
\end{enumerate}

Consider \(g(b)=\dfrac{b+1}{\ln b}\) for real \(b>1\):

\[
g'(b)=\frac{\ln b - 1 - \tfrac{1}{b}}{(\ln b)^2}.
\]

The critical point satisfies \(\ln b = 1+\frac{1}{b}\), which has a
unique solution \(b^\ast\in(3,4)\) (since \(\ln 3 < 1+\tfrac13\) and
\(\ln 4 > 1+\tfrac14\)). Therefore \(g(b)\) strictly decreases up to
\(b^\ast\) and increases thereafter, so the minimizing integer must be
\(3\) or \(4\).

Compare

\[
g(3)=\frac{4}{\ln 3},\qquad g(4)=\frac{5}{\ln 4}.
\]

We have \(g(4)<g(3)\) iff \(5\ln 3 < 4\ln 4 = 8\ln 2\), i.e.

\[
\ln(3^5) < \ln(2^8) \iff 243<256,
\]

which is true. Also \(g(5)>\!g(4)\) since \(6\ln 4 > 5\ln 5\)
(equivalently \(4^6>5^5\), i.e.~\(4096>3125\)). Thus \(b=4\) is the
unique minimizing integer.

\textbf{Answer (c):} \(b=4\).

\subsection{30 {[}12{]}}\label{section-37}

\begin{enumerate}
\def\labelenumi{\arabic{enumi})}
\tightlist
\item
  Domain sanity check
\end{enumerate}

\begin{itemize}
\tightlist
\item
  \(x>1\Rightarrow \ln x>0\), so the base \(\ln x\) is positive.
\item
  \(x\ne e\Rightarrow \ln x\ne 1\) and \(\ln\ln x\ne 0\), so the
  exponent is defined and the base isn't \(1\).
\end{itemize}

\begin{enumerate}
\def\labelenumi{\arabic{enumi})}
\setcounter{enumi}{1}
\tightlist
\item
  Recognize a change--of--base pattern
\end{enumerate}

Recall the identity (valid for any \(a>0\), \(a\neq1\)):
\(a^{\log_a t}=t\).

Note that

\[
\frac{\ln x}{\ln\ln x}=\log_{\ln x} x .
\]

Hence

\[
(\ln x)^{\ln x/\ln\ln x}
=(\ln x)^{\log_{\ln x} x}
= x .
\]

\begin{enumerate}
\def\labelenumi{\arabic{enumi})}
\setcounter{enumi}{2}
\tightlist
\item
  (Equivalent derivation via exponentials)
\end{enumerate}

\[
(\ln x)^{\ln x/\ln\ln x}
=\exp\!\Big(\frac{\ln x}{\ln\ln x}\cdot \ln(\ln x)\Big)
=\exp(\ln x)=x .
\]

\begin{enumerate}
\def\labelenumi{\arabic{enumi})}
\setcounter{enumi}{3}
\tightlist
\item
  Note on the excluded point
\end{enumerate}

At \(x=e\) the expression is undefined (division by
\(\ln\ln e=\ln1=0\)). However, the limit
\(\lim_{x\to e}(\ln x)^{\ln x/\ln\ln x}=\exp(\lim_{x\to e}\ln x)=e\), so
the continuous extension at \(x=e\) would take the value \(e\).

\[
\boxed{(\ln x)^{\ln x/\ln\ln x}=x \quad\text{for }x>1,\; x\ne e.}
\]

\bookmarksetup{startatroot}

\chapter{1.2.3 Sums and products}\label{sums-and-products}

\subsection{1 {[}10{]}}\label{section-38}

By definition in §1.2.3, the three-dot sum \(a_1+a_2+\cdots+a_n\) is
shorthand for the indexed sum \(\sum_{1\le j\le n} a_j\). More
generally, \(\sum_{R(j)} a_j\) means ``sum all \(a_j\) with integer
\(j\) that satisfy \(R(j)\),'' and if no integers satisfy \(R(j)\), the
sum is defined to be \(0\)**. 

So

\[
a_2+\cdots+a_0\;\equiv\;\sum_{2\le j\le 0} a_j \;=\;0,
\]

because there is no integer \(j\) with \(2\le j\le 0\).

This is why Knuth calls the behavior ``slightly peculiar but sensible''
in edge cases: the three-dot notation is a single construct tied to its
explicit bounds, not something you peel apart as \(a_2+(\cdots+a_0)\).
Read as a bounded sum, its index set is empty, hence the value \(0\). 

\subsection{2 {[}01{]}}\label{section-39}

By Knuth's convention, a restricted sum

\[
\sum_{R(j)} a_j
\]

means ``sum \(a_j\) over all \textbf{integers} \(j\) that satisfy the
condition \(R(j)\).'' If no such integers exist, the sum is defined to
be \(0\). 

Here \(R(j)\) is \(1\le j\le n\). When \(n=3.14\), the integers \(j\)
with \(1\le j\le 3.14\) are \(j=1,2,3\). Therefore

\[
\sum_{1\le j\le 3.14} a_j \;=\; a_1+a_2+a_3 \;=\; \sum_{j=1}^{3} a_j.
\]

(Equivalently, for any real \(n\),
\(\sum_{1\le j\le n} a_j=\sum_{j=1}^{\lfloor n\rfloor} a_j\); this
follows from the ``sum over integers satisfying \(R(j)\)''
interpretation and the bracket-notation identity
\(\sum_{R(j)} a_j=\sum_j a_j\,[R(j)]\).) 

\subsection{3 {[}13{]}}\label{section-40}

Write out, \textbf{without} Σ-notation, the sums

\[
\sum_{0\le n\le 5}\frac1{2n+1}
\qquad\text{and}\qquad
\sum_{0\le n^2\le 5}\frac1{2n^2+1},
\]

then explain why the two results are different ``in spite of rule (b)''
(the change-of-variable rule).

Step 1: Expand both sums explicitly

\begin{enumerate}
\def\labelenumi{\arabic{enumi}.}
\tightlist
\item
  \(\displaystyle \sum_{0\le n\le 5}\frac1{2n+1}\) uses the integer
  indices \(n=0,1,2,3,4,5\):
\end{enumerate}

\[
\frac1{1}+\frac1{3}+\frac1{5}+\frac1{7}+\frac1{9}+\frac1{11}.
\]

\begin{enumerate}
\def\labelenumi{\arabic{enumi}.}
\setcounter{enumi}{1}
\tightlist
\item
  \(\displaystyle \sum_{0\le n^2\le 5}\frac1{2n^2+1}\) uses all integers
  with square at most 5, i.e.~\(n\in\{-2,-1,0,1,2\}\):
\end{enumerate}

\[
\frac1{2\cdot 4+1}+\frac1{2\cdot 1+1}+\frac1{2\cdot 0+1}
+\frac1{2\cdot 1+1}+\frac1{2\cdot 4+1}
= \frac1{9}+\frac1{3}+1+\frac1{3}+\frac1{9}
= 1+\frac{2}{3}+\frac{2}{9}.
\]

(If you like a single value: the first sum ≈ 1.8782; the second equals
\(17/9\approx 1.8889\).)

Step 2: Why the results differ despite rule (b)

Rule (b) says we may change the index variable via a
\textbf{permutation} \(p(j)\) of the relevant integers:

\[
\sum_{R(i)} a_i \;=\; \sum_{R(p(j))} a_{p(j)}.
\]

This requires \(p\) to be one-to-one and onto the index set (a
bijection). 

\begin{itemize}
\tightlist
\item
  In the first sum, the index set is \(\{0,1,2,3,4,5\}\) (six elements).
\item
  In the second, the condition \(0\le n^2\le 5\) yields the index set
  \(\{-2,-1,0,1,2\}\) (five elements).
\end{itemize}

Trying to ``change variables'' by \(m=n^2\) is not a valid use of rule
(b): the map \(n\mapsto n^2\) is not a permutation of the integer
indices---it is many-to-one (both \(+k\) and \(-k\) map to the same
square), and it also changes the index set's size and elements.
Therefore the two sums legitimately come out different; rule (b) simply
doesn't apply here.

\subsection{4 {[}10{]}}\label{section-41}

Left-hand side (iterate \(i=1,2,3\), then \(j=1,\dots,i\))

\begin{itemize}
\tightlist
\item
  \(i=1:\ j=1\Rightarrow a_{11}\).
\item
  \(i=2:\ j=1,2\Rightarrow a_{21}+a_{22}\).
\item
  \(i=3:\ j=1,2,3\Rightarrow a_{31}+a_{32}+a_{33}\).
\end{itemize}

So the LHS is

\[
(a_{11})\;+\;(a_{21}+a_{22})\;+\;(a_{31}+a_{32}+a_{33}).
\]

Right-hand side (iterate \(j=1,2,3\), then \(i=j,\dots,3\))

\begin{itemize}
\tightlist
\item
  \(j=1:\ i=1,2,3\Rightarrow a_{11}+a_{21}+a_{31}\).
\item
  \(j=2:\ i=2,3\Rightarrow a_{22}+a_{32}\).
\item
  \(j=3:\ i=3\Rightarrow a_{33}\).
\end{itemize}

So the RHS is

\[
(a_{11}+a_{21}+a_{31})\;+\;(a_{22}+a_{32})\;+\;(a_{33}).
\]

Why they're equal

Both expressions list exactly the six terms with \(1\le j\le i\le3\):
\(\{a_{11},a_{21},a_{22},a_{31},a_{32},a_{33}\}\). The LHS groups them
by row \(i\); the RHS groups them by column \(j\). Since addition is
commutative/associative, the total is the same---exactly what Eq. (10)
formalizes.

\subsection{5 {[}HM20{]}}\label{hm20-1}

Let the two series be convergent:

\[
A:=\sum_{R(i)} a_i \quad\text{and}\quad B:=\sum_{S(j)} b_j.
\]

Knuth's definition of an infinite sum is ``limit of partial sums'' (with
two--sided windows when needed). We will use that definition directly. 

Step 1: Constant-factor rule for convergent series

Fix a constant \(c\). For the inner partial sums
\(B_M:=\sum_{S_M(j)} b_j\) (over any increasing finite domains \(S_M\)
that exhaust \(S\)), we have

\[
\sum_{S_M(j)} (c\,b_j)=c\sum_{S_M(j)} b_j= c\,B_M
\]

by the finite distributive law. Taking limits as \(M\to\infty\) and
using convergence of \(\sum_{S(j)} b_j\) gives

\[
\sum_{S(j)} (c\,b_j)= c\sum_{S(j)} b_j=c\,B.
\]

Thus constants can be pulled out of a convergent series.

Step 2: Evaluate the inner sum

For each fixed \(i\), \(a_i\) is constant with respect to \(j\); hence

\[
\sum_{S(j)} a_i b_j \;=\; a_i \sum_{S(j)} b_j \;=\; a_i B,
\]

by Step 1.

Step 3: Sum over \(i\)

Now take partial sums over \(i\) (finite domains \(R_N\) that exhaust
\(R\)):

\[
\sum_{i\in R_N}\Bigl(\sum_{S(j)} a_i b_j\Bigr)
= \sum_{i\in R_N} (a_i B)
= B \sum_{i\in R_N} a_i
=: B\,A_N,
\]

again using the finite distributive law (rule (a) in the finite case).
Letting \(N\to\infty\) and using the convergence of \(\sum_{R(i)} a_i\)
yields

\[
\sum_{R(i)}\sum_{S(j)} a_i b_j
= \lim_{N\to\infty} B\,A_N
= B\cdot A
= \Bigl(\sum_{R(i)} a_i\Bigr)\Bigl(\sum_{S(j)} b_j\Bigr).
\]

This completes the proof that the distributive rule (a) extends from
finite sums to convergent infinite series, precisely as stated in Eq.
(4) of the text for the finite case.

\subsection{6 {[}HM20{]}}\label{hm20-2}

Key identity (``cut rule'')

The fundamental law of these half-open sums is the \textbf{cut rule}
(this is the Eq. (3) you're asked to use):

\[
S(a,b)\;=\;S(a,c)\;+\;S(c,b)\qquad\text{for any integers }a,c,b.
\]

It just partitions the index set \([a,b)\) into the disjoint pieces
\([a,c)\) and \([c,b)\).

Changing the lower limit

Apply the cut rule and rearrange.

\textbf{Case 1: \(a_0\le a_c\le b\).}

\[
S(a_0,b)=S(a_c,b)\;+\;S(a_0,a_c).
\]

So moving the lower limit \textbf{up} from \(a_0\) to \(a_c\) removes
the block of terms with indices \(a_0\le k<a_c\); to keep equality, you
must \textbf{add back} exactly that block.

\textbf{Case 2: \(a_c\le a_0\le b\).}

\[
S(a_0,b)=S(a_c,b)\;-\;S(a_c,a_0).
\]

Now moving the lower limit \textbf{down} from \(a_0\) to \(a_c\)
includes the extra block \(a_c\le k<a_0\); to compensate, you
\textbf{subtract} that block.

These two cases are the same rule written with the appropriate sign.

Interchanging limits (the ``reverse'' rule)

From the cut rule with \(b=a\) you get the useful symmetry

\[
S(a,c)+S(c,a)=S(a,a)=0\quad\Longrightarrow\quad S(a,c)=-\,S(c,a).
\]

Thus when you \textbf{swap the two limits} in a half-open sum, the value
changes sign, and ``the terms between \(a_0\) and \(a_c\)'' migrate from
one side of the equality to the other---the precise content of the
official hint.

Worked micro-examples

\begin{enumerate}
\def\labelenumi{\arabic{enumi}.}
\tightlist
\item
  \(a_0=2,\ a_c=5,\ b=10\):
\end{enumerate}

\[
\sum_{2\le k<10}x_k=\sum_{5\le k<10}x_k+\sum_{2\le k<5}x_k.
\]

\begin{enumerate}
\def\labelenumi{\arabic{enumi}.}
\setcounter{enumi}{1}
\tightlist
\item
  \(a_0=7,\ a_c=3,\ b=10\):
\end{enumerate}

\[
\sum_{7\le k<10}x_k=\sum_{3\le k<10}x_k-\sum_{3\le k<7}x_k.
\]

\begin{enumerate}
\def\labelenumi{\arabic{enumi}.}
\setcounter{enumi}{2}
\tightlist
\item
  Interchange \(a\) and \(c\):
\end{enumerate}

\[
\sum_{a\le k<c}x_k=-\sum_{c\le k<a}x_k.
\]

\subsection{7 {[}HM23{]}}\label{hm23}

\begin{enumerate}
\def\labelenumi{\arabic{enumi})}
\tightlist
\item
  Use Iverson's bracket definition of restricted sums
\end{enumerate}

Knuth defines a ``restricted'' sum by turning the restriction into a
0--1 factor:

\[
\sum_{R(j)} a_j \;=\; \sum_{j\in\mathbb Z} a_j\,[R(j)],
\]

where \([R(j)]\) is 1 if the statement \(R(j)\) is true and 0 otherwise.
This lets us treat ``sum over all \(j\) satisfying \(R\)'' as a single
sum over all integers with a filter. 

\begin{enumerate}
\def\labelenumi{\arabic{enumi})}
\setcounter{enumi}{1}
\tightlist
\item
  Make the change of variables \(k = c - j\)
\end{enumerate}

The map \(j\mapsto c-j\) is a bijection (indeed, an involution) on
\(\mathbb Z\). Rename the dummy index in the all-integers sum using this
bijection:

\[
\sum_{j\in\mathbb Z} a_j\,[R(j)]
\;=\;
\sum_{k\in\mathbb Z} a_{\,c-k}\,[R(c-k)].
\]

But the right-hand side is exactly the restricted-sum notation

\[
\sum_{R(c-k)} a_{\,c-k}.
\]

Since \(k\) is just a dummy symbol, we can write it as \(j\) again and
obtain

\[
\sum_{R(j)} a_j \;=\; \sum_{R(c-j)} a_{\,c-j}.
\]

That is all that needs to be shown.

Why this is valid ``even if both series are infinite''

The proof above \textbf{does not} rearrange the terms of a conditionally
convergent series; it only \emph{renames the index} inside a single
``sum over all integers with a 0--1 filter.'' That's exactly what
Knuth's bracket notation is designed to permit. (In the text, this
equivalence is also derived from the general ``change of index by a
permutation'' rule; here the permutation is \(p(j)=c-j\). Knuth
discusses change of index and shows, via brackets, how such a
permutation yields the same restricted sum. ) 

\begin{enumerate}
\def\labelenumi{\arabic{enumi})}
\setcounter{enumi}{2}
\tightlist
\item
  Quick sanity checks
\end{enumerate}

\begin{itemize}
\item
  \textbf{Finite window:} If \(R(j)\) is \(0\le j\le n\) then the result
  says
  \(\sum_{j=0}^{n} a_j = \sum_{0\le c-j\le n} a_{c-j} = \sum_{c-n\le j\le c} a_j,\)
  i.e., just a shift of the index range---obviously true.
\item
  \textbf{One-sided infinite:} If \(R(j)\) is \(j\ge 0\) then

  \[
  \sum_{j\ge 0} a_j \;=\; \sum_{c-j\ge 0} a_{\,c-j} \;=\; \sum_{j\le c} a_j,
  \]

  which again follows directly from the bijection
  \(j\leftrightarrow c-j\) inside the bracketed all-integers sum.
\end{itemize}

\subsection{8 {[}HM25{]}}\label{hm25}

We'll use Stirling's expansion from the text:

\[
n! \;=\; \sqrt{2\pi n}\,\Big(\tfrac{n}{e}\Big)^n
\Big(1+\tfrac{1}{12n}+O(n^{-2})\Big)
\quad\Longrightarrow\quad
\ln n!
= n\ln n-n+\tfrac12\ln(2\pi n)+\tfrac{1}{12n}+O(n^{-2}). 
\]

(See the displayed formulas just before and at (19).)

Part A --- \(\ln(an^2+bn)!\) with absolute error \(O(n^{-2})\)

Let \(m=an^2+bn\) with fixed \(a>0\) and fixed \(b\). Since
\(m\sim an^2\), the Stirling remainder terms beyond \(1/(12m)\) are
\(O(m^{-3})=O(n^{-6})\), which is (much) smaller than the requested
\(O(n^{-2})\). Thus it suffices to expand

\[
\ln m! \;=\; m\ln m - m + \tfrac12\ln(2\pi m)+\tfrac{1}{12m} + O(n^{-6}).
\]

Write

\[
\ln m = \ln(an^2)+\ln\!\Bigl(1+\frac{b}{a\,n}\Bigr)
= \ln a + 2\ln n + \frac{b}{a\,n}-\frac{b^2}{2a^2 n^2}+O(n^{-3}).
\]

Now expand the four pieces \(m\ln m\), \(-m\), \(\tfrac12\ln(2\pi m)\),
\(\tfrac{1}{12m}\) and collect terms through \(n^{-2}\). A
straightforward computation gives

\[
\boxed{
\begin{aligned}
\ln(an^2+bn)!
&= a n^2\bigl(2\ln n+\ln a-1\bigr)
  + b n\bigl(2\ln n+\ln a\bigr)\\
&\quad+ \tfrac12\bigl(2\ln n+\ln a+\ln(2\pi)\bigr)
  + \frac{b^2}{2a}\\[2pt]
&\quad+ \frac{1}{n}\Bigl(\frac{b}{2a}-\frac{b^3}{2a^2}\Bigr)
  + \frac{1}{n^2}\Bigl(\frac{1}{12a}-\frac{b^2}{4a^2}\Bigr)
  \;+\; O(n^{-3}).
\end{aligned}
}
\]

Because the omitted terms are \(O(n^{-3})\), this approximation has
\textbf{absolute error \(O(n^{-2})\)} as required.

\emph{(Sanity check: setting \(b=0, a=1\) reduces to \(m=n^2\) and we
recover
\(\ln(n^2)! = 2n^2\ln n-n^2+\ln n+\tfrac12\ln(2\pi)+\tfrac{1}{12n^2}+O(n^{-4})\),
which is Stirling with \(m=n^2\).)}

Part B --- Asymptotics of
\(\displaystyle \frac{\binom{c n^2}{n}}{c^{\,n}\binom{n^2}{n}}\)

Write \(\binom{N}{n}=\dfrac{N!}{n!(N-n)!}\) and take logs:

\[
\ln\binom{N}{n}=\ln N!-\ln n!-\ln(N-n)!.
\]

Put \(N=\alpha n^2\) with \(\alpha>0\) (we will use \(\alpha=c\) and
\(\alpha=1\)). Using the Stirling--log expansion to accuracy
\(O(n^{-2})\) (as above) and expanding in powers of \(1/n\) with
\(n\ll N\), one finds

\[
\begin{aligned}
\ln\binom{\alpha n^2}{n}
&= n\ln n + n\ln\alpha + n - \tfrac12\ln n
   - \frac{1}{2\alpha}
   + \frac{1}{n}\Bigl(\frac{1}{2\alpha}-\frac{1}{6\alpha^2}-\frac{1}{12}\Bigr)\\[-2pt]
&\quad + \frac{1}{n^2}\Bigl(\frac{1}{12\alpha^3}+\frac{1}{4\alpha^2}\Bigr)
   + O(n^{-3}).
\end{aligned}
\]

(Any \(\alpha\)-independent \(O(1/n)\) terms, e.g.~\(-1/(12n)\), cancel
when we form the ratio below.)

Therefore

\[
\begin{aligned}
\ln R
&:= \ln\frac{\binom{c n^2}{n}}{c^{\,n}\binom{n^2}{n}}
= \Bigl[\ln\binom{c n^2}{n}-n\ln c\Bigr] - \ln\binom{n^2}{n}\\
&= \frac12\Bigl(1 - \frac{1}{c}\Bigr)
 + \frac{1}{n}\Bigl(\frac{1}{2c}-\frac{1}{6c^2}-\frac{1}{3}\Bigr)
 + \frac{1}{n^2}\Bigl(\frac{1}{12c^3}+\frac{1}{4c^2}-\frac{1}{12}-\frac{1}{4}\Bigr)
 + O(n^{-3}).
\end{aligned}
\]

Exponentiating gives the asymptotic form

\[
\boxed{
\frac{\binom{c n^2}{n}}{c^{\,n}\binom{n^2}{n}}
= \exp\!\Bigl(\tfrac12\,(1-\tfrac1c)\Bigr)\,
\Biggl[\,1+\frac{1}{n}\Bigl(\frac{1}{2c}-\frac{1}{6c^2}-\frac{1}{3}\Bigr)
+ O\!\Bigl(\frac{1}{n^2}\Bigr)\Biggr].
}
\]

In particular, the \textbf{asymptotic value} (the limit as
\(n\to\infty\)) is the constant

\[
\boxed{\displaystyle \lim_{n\to\infty}\frac{\binom{c n^2}{n}}{c^{\,n}\binom{n^2}{n}}
= \exp\!\Bigl(\tfrac12\,(1-\tfrac1c)\Bigr).}
\]

Checks: for \(c=1\) the ratio is identically \(1\), and the limit gives
\(e^{0}=1\); as \(c\to\infty\), the limit approaches \(e^{1/2}\).

\subsection{9 {[}05{]}}\label{section-42}

\textbf{No.} The \emph{formula} still evaluates correctly when \(n=-1\)
(both sides are \(0\)), but the \emph{derivation} used in the text is
not valid in that case because it relies on a step that implicitly
assumes \(0\le 0 \le n\) (i.e., \(n\ge 0\)). 

Why the derivation breaks (but the identity remains true)

\begin{enumerate}
\def\labelenumi{\arabic{enumi}.}
\tightlist
\item
  \textbf{What the derivation does for \(n\ge 0\).} Let
  \(S:=\sum_{0\le j\le n} a x^{j}\). The text splits off the \(j=0\)
  term and then reindexes:
\end{enumerate}

\[
S \;=\; a \;+\; \sum_{1\le j\le n} a x^{j}
\;=\; a \;+\; x\sum_{1\le j\le n} a x^{j-1}
\;=\; a \;+\; x\sum_{0\le j\le n-1} a x^{j},
\]

then subtracts \(xS\) and solves to get Eq. (14). Every equality there
uses the summation manipulation rules from §1.2.3 and, crucially, the
ability to \emph{pull out} the \(j=0\) term---i.e., that \(0\) actually
lies in the index set \(\{j\mid 0\le j\le n\}\), which requires
\(n\ge 0\). 

\begin{enumerate}
\def\labelenumi{\arabic{enumi}.}
\setcounter{enumi}{1}
\tightlist
\item
  \textbf{What happens when \(n=-1\).} The index set
  \(\{j\mid 0\le j\le -1\}\) is empty, so
\end{enumerate}

\[
\sum_{0\le j\le -1} a x^{j}=0\quad\text{(empty sum).}
\]

But the very first step of the derivation would write

\[
S \stackrel{?}= a + \sum_{1\le j\le -1} a x^{j},
\]

which incorrectly assumes that a \(j=0\) term is present to ``split
off.'' That step is invalid when \(n=-1\). (In Knuth's rules this is the
misuse of restriction/splitting (rule (d)) outside its preconditions.)
Hence the \emph{derivation} fails for \(n=-1\). 

\begin{enumerate}
\def\labelenumi{\arabic{enumi}.}
\setcounter{enumi}{2}
\tightlist
\item
  \textbf{Yet the formula itself still holds for \(n=-1\).} Evaluate
  both sides directly:
\end{enumerate}

\begin{itemize}
\tightlist
\item
  LHS: \(\sum_{0\le j\le -1} a x^j = 0\) (empty sum).
\item
  RHS:
  \(a\,\dfrac{1-x^{\,(-1)+1}}{1-x}=a\,\dfrac{1-x^{0}}{1-x}=a\,\dfrac{1-1}{1-x}=0\).
  So the identity is true at \(n=-1\); only the \emph{particular
  derivation} is invalid there. (This is exactly the remark in the
  answers for §1.2.3.)
\end{itemize}

\subsection{10 {[}05{]}}\label{section-43}

What Eq. (14) is and how it's derived

In §1.2.3 Knuth derives the finite geometric--series identity (for
\(x\neq 1\))

\[
\sum_{j=0}^{n} x^{j}=\frac{x^{\,n+1}-1}{x-1}.
\]

A standard derivation is

\[
(x-1)\sum_{j=0}^{n}x^{j}
=\sum_{j=0}^{n}(x^{j+1}-x^{j})
=\Big(\sum_{k=1}^{n+1}x^{k}\Big)-\Big(\sum_{j=0}^{n}x^{j}\Big)
=x^{n+1}-1,
\]

so dividing by \(x-1\) (with \(x\ne 1\)) gives Eq. (14). In this section
Knuth also formalizes \emph{conditional} sums using the
characteristic--function/bracket notation \([R]\in\{0,1\}\), so that
sums over an empty index set evaluate to \(0\). 

Check the edge cases \(n=-1\) and \(n=-2\)

\begin{itemize}
\item
  \textbf{What the sum means for negative \(n\).} In Knuth's notation
  \(\sum_{0\le j\le n}(\cdots)\) is a sum over integers \(j\) satisfying
  \(0\le j\le n\). If \(n<0\), there \emph{are no such \(j\)}, so the
  sum is \textbf{empty} and therefore equals \textbf{0}. (This follows
  immediately from the conditional-sum/bracket setup in §1.2.3.) 
\item
  \textbf{\(n=-1\).} The left side is an empty sum \(=0\). The right
  side of Eq. (14) becomes \((x^{0}-1)/(x-1)=0\). So the identity
  happens to hold.
\item
  \textbf{\(n=-2\).} The left side is again an empty sum \(=0\). But the
  right side of Eq. (14) becomes
\end{itemize}

\[
\frac{x^{-1}-1}{x-1}=\frac{1-x}{x(x-1)}=-\frac{1}{x}\neq 0\quad(\text{for }x\neq 1).
\]

Thus the \emph{formula} of Eq. (14) does \textbf{not} match the meaning
of the sum when \(n=-2\).

Where the derivation breaks

Look at the telescoping step

\[
\Big(\sum_{k=1}^{n+1}x^{k}\Big)-\Big(\sum_{j=0}^{n}x^{j}\Big)=x^{n+1}-1.
\]

This cancellation implicitly assumes there is an overlapping block
\(1\le k\le n\) to cancel---i.e., it assumes \(n\ge 0\). When \(n=-2\)
(indeed whenever \(n\le -2\)), \textbf{both sums are empty}, so their
difference is \(0\), not \(x^{n+1}-1\). The telescoping argument
therefore fails, and the subsequent division to obtain Eq. (14) is not
justified.

\subsection{11 {[}03{]}}\label{section-44}

\begin{enumerate}
\def\labelenumi{\arabic{enumi}.}
\item
  \textbf{Direct counting at \(x=1\).} When \(x=1\), every term in the
  sum is just \(a\). There are \(n+1\) terms (from \(j=0\) to \(j=n\)),
  so

  \[
  \sum_{j=0}^{n} a\,1^{j} = \underbrace{a+\cdots+a}_{n+1\ \text{terms}} = a(n+1).
  \]

  Therefore, the \emph{value} of the sum at \(x=1\) is \(a(n+1)\).
\item
  \textbf{Consistent extension of Eq. (14).} The closed form
  \(a(1-x^{n+1})/(1-x)\) is undefined at \(x=1\) (division by zero), but
  its \textbf{limit} exists. Rewrite the fraction as

  \[
  a\,\frac{x^{n+1}-1}{x-1}.
  \]

  Taking \(x \to 1\) and using the standard difference-quotient identity
  (or L'Hôpital's rule),

  \[
  \lim_{x\to 1}\frac{x^{n+1}-1}{x-1}=\left.\frac{d}{dx}x^{n+1}\right|_{x=1}=n+1.
  \]

  Hence the continuous extension of Eq. (14) at \(x=1\) is \(a(n+1)\),
  agreeing with the direct count.
\item
  \textbf{Polynomial identity check (no calculus).} The identity

  \[
  x^{n+1}-1=(x-1)\,(1+x+\cdots+x^{n})
  \]

  holds for all \(x\). For \(x\neq 1\) it implies Eq. (14); evaluating
  the polynomial \(1+x+\cdots+x^{n}\) at \(x=1\) gives \(n+1\). Thus the
  natural definition at \(x=1\) is again \(a(n+1)\).
\end{enumerate}

Final answer

\[
\boxed{\,a(n+1)\,}
\]

That is, replace the right-hand side by \(a(n+1)\) when \(x=1\). 

\subsection{12 {[}10{]}}\label{section-45}

This is a geometric series with first term \(a=1\) and common ratio
\(r=\tfrac17\).

\begin{enumerate}
\def\labelenumi{\arabic{enumi}.}
\item
  Multiply by \(r\):
  \(rS_n=\tfrac17+\tfrac1{7^2}+\cdots+\tfrac1{7^{n+1}}\).
\item
  Subtract: \(S_n-rS_n=1-\tfrac1{7^{\,n+1}}\).
\item
  Factor: \((1-r)S_n=1-\tfrac1{7^{\,n+1}}\).
\item
  Solve for \(S_n\) using \(1-r=1-\tfrac17=\tfrac67\):
\end{enumerate}

\[
S_n=\frac{1-(1/7)^{\,n+1}}{1-1/7}
=\frac{1-(1/7)^{\,n+1}}{6/7}
=\frac{7}{6}\Bigl(1-7^{-(n+1)}\Bigr).
\]

Equivalent closed forms you can use:

\[
S_n=\frac{7^{n+1}-1}{6\cdot 7^{\,n}}
=\frac{7}{6}-\frac{1}{6\cdot 7^{\,n}}.
\]

Quick checks

\begin{itemize}
\tightlist
\item
  \(n=0\): \(S_0=1\). Formula gives
  \(\frac{7}{6}-\frac{1}{6\cdot 7^0}=1\).
\item
  \(n=1\): \(S_1=1+\tfrac17=\tfrac87\). Formula gives
  \(\frac{7}{6}-\frac{1}{42}=\tfrac{56-1}{42}=\tfrac87\).
\end{itemize}

(As a side note, \(\lim_{n\to\infty}S_n=\tfrac{7}{6}\).)

\subsection{13 {[}10{]}}\label{section-46}

Tool: Eq. (15) from the text

From the derivation of arithmetic--progression sums in §1.2.3, Eq. (15)
gives

\[
\sum_{j=0}^{n} (a+bj)\;=\;a(n+1)\;+\;\frac{b\,n(n+1)}{2}.
\]

Setting \(a=0,\; b=1\) yields the well-known identity

\[
\sum_{j=0}^{n} j \;=\; \frac{n(n+1)}{2}. \;\; \text{(from Eq. 15)}
\]

Compute \(\sum_{j=m}^{n} j\)

Write the desired sum as a difference of two prefix sums:

\[
\sum_{j=m}^{n} j
\;=\;
\sum_{j=0}^{n} j \;-\; \sum_{j=0}^{m-1} j
\;=\;
\frac{n(n+1)}{2} \;-\; \frac{(m-1)m}{2}.
\]

So a compact closed form is

\[
\boxed{\displaystyle \sum_{j=m}^{n} j \;=\; \frac{n(n+1) - m(m-1)}{2}}
\qquad (m\le n).
\]

Equivalent forms (sometimes convenient)

Algebra gives a few interchangeable expressions:

\[
\sum_{j=m}^{n} j
\;=\;
\frac{(n-m+1)(m+n)}{2}
\;=\;
m(n-m+1)\;+\;\frac{(n-m)(n-m+1)}{2}.
\]

(The text's answer lists these equivalent forms as well. ) 

Quick check: Example \(m=3,\;n=6\): \(3+4+5+6=18\). Formula:
\(\tfrac{6\cdot7-3\cdot2}{2}= \tfrac{42-6}{2}=18\). ✓

\subsection{14 {[}11{]}}\label{section-47}

Because \(j\) does not depend on \(k\), the inner sum factors:

\[
\sum_{k=r}^{s} jk \;=\; j\sum_{k=r}^{s} k.
\]

Therefore,

\[
\sum_{j=m}^{n}\sum_{k=r}^{s} jk
\;=\;
\left(\sum_{j=m}^{n} j\right)\!\left(\sum_{k=r}^{s} k\right).
\]

Use the standard arithmetic--series formula from the previous exercise
(with \(m\le n\), and analogously \(r\le s\)):

\[
\sum_{j=m}^{n} j
=
\frac{n(n+1)}{2}-\frac{(m-1)m}{2}
=
\frac{(n-m+1)(m+n)}{2},
\]

\[
\sum_{k=r}^{s} k
=
\frac{s(s+1)}{2}-\frac{(r-1)r}{2}
=
\frac{(s-r+1)(r+s)}{2}.
\]

(Exercise 13 asks you to derive \(\sum_{j=m}^{n} j\) in this way.) 

Multiply the two results to get the closed form:

\[
\boxed{\displaystyle
\sum_{j=m}^{n}\sum_{k=r}^{s} jk
=
\frac{(n-m+1)(m+n)\,(s-r+1)(r+s)}{4}
}\qquad (m\le n,\ r\le s).
\]

If either range is empty (e.g., \(m>n\) or \(r>s\)), the corresponding
sum is \(0\), so the whole double sum is \(0\) by the conventions of
§1.2.3. 

Quick check

Take \(m=2,n=4,r=3,s=5\). Left side:

\[
\sum_{j=2}^{4}\sum_{k=3}^{5} jk
=(2(3+4+5))+(3(3+4+5))+(4(3+4+5))=24+36+48=108.
\]

Right side:

\[
\frac{(4-2+1)(2+4)\,(5-3+1)(3+5)}{4}
=\frac{(3)(6)\,(3)(8)}{4}=108.
\]

Matches perfectly.

\subsection{15 {[}M22{]}}\label{m22-1}

Step 1 --- Small \(n\) to spot the pattern

\[
\begin{aligned}
S_1&=2\\
S_2&=2+8=10\\
S_3&=2+8+24=34\\
S_4&=98,\qquad S_5=258
\end{aligned}
\]

These suggest the closed form \(S_n=(n-1)2^{\,n+1}+2\) (check: for
\(n=5\), \(4\cdot 64+2=258\)).

Step 2 --- Derive it with the ``geometric-series manipulation''

Let

\[
S_n=\sum_{k=1}^{n} k\,2^{k}.
\]

Multiply by 2 and reindex the first sum:

\[
2S_n=\sum_{k=1}^{n} k\,2^{k+1}=\sum_{j=2}^{n+1} (j-1)\,2^{j}.
\]

Subtract \(S_n=\sum_{j=1}^{n} j\,2^{j}\) from this:

\[
\begin{aligned}
2S_n-S_n
&=\biggl[\sum_{j=2}^{n+1} (j-1)2^{j}\biggr]-\biggl[\sum_{j=1}^{n} j\,2^{j}\biggr]\\
&=\;n\,2^{\,n+1}\;+\;\sum_{j=2}^{n}\!\bigl((j-1)-j\bigr)2^{j}\;-\;1\cdot 2\\
&=\;n\,2^{\,n+1}\;-\;\sum_{j=1}^{n}2^{j}.
\end{aligned}
\]

Using the finite geometric sum \(\sum_{j=1}^{n}2^{j}=2^{n+1}-2\) (the
very Eq. (14)--style identity Knuth just derived earlier in §1.2.3), we
get

\[
S_n=n\,2^{\,n+1}-\bigl(2^{\,n+1}-2\bigr)=(n-1)\,2^{\,n+1}+2.
\]

Final answer

\[
\boxed{\,\displaystyle \sum_{k=1}^{n} k\,2^{k}\;=\;(n-1)\,2^{\,n+1}+2\,}
\]

(For context, this exercise sits in §1.2.3 ``Sums and Products,'' right
after the geometric-series derivation. )

\subsection{16 {[}M22{]}}\label{m22-2}

Step 1 --- Set up a telescoping trick

Let \(S=S_n(x)\). Then

\[
(1-x)S = S - xS \;=\; \sum_{j=0}^{n} jx^{j}\;-\;\sum_{j=0}^{n} jx^{j+1}.
\]

Change index in the second sum: put \(k=j+1\). Using the ``change of
variable'' rule (rule (b) in §1.2.3),

\[
\sum_{j=0}^{n} jx^{j+1} \;=\; \sum_{k=1}^{n+1} (k-1)x^{k}.
\]

Hence

\[
(1-x)S
= \Big(\,0\,+\sum_{j=1}^{n} jx^{j}\Big) - \Big(\sum_{k=1}^{n} (k-1)x^{k} + n x^{n+1}\Big)
= \sum_{j=1}^{n} \big(j-(j-1)\big)x^{j} - n x^{n+1}
= \sum_{j=1}^{n} x^{j} - n x^{n+1}.
\]

(The index-shift is exactly the kind of manipulation formalized as rule
(b).) 

Step 2 --- Evaluate the geometric sum

For \(x\neq 1\),

\[
\sum_{j=1}^{n} x^{j} \;=\; \frac{x - x^{\,n+1}}{1-x},
\]

the standard finite geometric progression with \(a=1\). 

Therefore

\[
(1-x)S \;=\; \frac{x - x^{\,n+1}}{1-x} - n x^{n+1}.
\]

Step 3 --- Solve for \(S\)

Multiply both sides by \(1-x\) and simplify:

\[
(1-x)^{2} S \;=\; x - x^{\,n+1} - n x^{\,n+1}(1-x)
= x - (n+1)x^{\,n+1} + n x^{\,n+2}.
\]

Thus

\[
\boxed{\; \displaystyle
\sum_{j=0}^{n} j\,x^{j} \;=\; \frac{n x^{\,n+2} - (n+1) x^{\,n+1} + x}{(1-x)^{2}}
\;=\; \frac{n x^{\,n+2} - (n+1) x^{\,n+1} + x}{(x-1)^{2}} \;}\qquad(x\neq 1).
\]

This matches the target identity stated in the exercise (and it was
obtained without induction). 

Checks and remarks

\begin{itemize}
\tightlist
\item
  If \(x=1\), the sum is \(\sum_{j=0}^{n} j = n(n+1)/2\) (consistent
  with the limit of the formula as \(x\to 1\)).
\item
  The proof used only the index-shift rule and the finite geometric sum
  derived in Example 3 of §1.2.3, as intended for this section.
\end{itemize}

Start from

\[
S_n(x)=\sum_{j=0}^{n} j\,x^{j}.
\]

\begin{enumerate}
\def\labelenumi{\arabic{enumi}.}
\tightlist
\item
  Pull out a factor \(x\) and shift:
\end{enumerate}

\[
S_n(x)=x\sum_{j=1}^{n} j\,x^{j-1}=x\sum_{j=0}^{n-1}(j+1)x^{j}.
\]

\begin{enumerate}
\def\labelenumi{\arabic{enumi}.}
\setcounter{enumi}{1}
\tightlist
\item
  Split the last sum:
\end{enumerate}

\[
S_n(x)=x\sum_{j=0}^{n-1} j\,x^{j}+x\sum_{j=0}^{n-1} x^{j}
= x\Bigl(\sum_{j=0}^{n} j\,x^{j}-n x^{n}\Bigr) \;+\; x\sum_{j=0}^{n-1} x^{j}.
\]

\begin{enumerate}
\def\labelenumi{\arabic{enumi}.}
\setcounter{enumi}{2}
\tightlist
\item
  Recognize the geometric sum
  \(\sum_{j=0}^{n-1}x^{j}=\dfrac{1-x^{n}}{1-x}\) (from §1.2.3) and solve
  for \(S_n\):
\end{enumerate}

\[
S_n(x) = x S_n(x) - n x^{n+1} + x\,\frac{1-x^{n}}{1-x}.
\]

\[
(1-x)S_n(x)= -n x^{n+1} + x\,\frac{1-x^{n}}{1-x}.
\]

\[
S_n(x)=\frac{x(1-x^{n})-n x^{n+1}(1-x)}{(1-x)^{2}}
=\frac{x-(n+1)x^{n+1}+n x^{n+2}}{(1-x)^{2}}, \qquad (x\neq 1).
\]

This is exactly the result Exercise 20 later depends on (it rewrites
\(\sum (n-k)b^k\) via \(n\sum b^k-\sum k b^k\)). 

Edge cases and checks

\begin{itemize}
\tightlist
\item
  \textbf{\(x=1\).} Take the limit \(x\to1\) (or note directly from the
  definition):
  \(\displaystyle S_n(1)=\sum_{j=0}^{n} j=\frac{n(n+1)}{2}.\)
\item
  \textbf{Polynomials.} The identity holds over any ring where the
  algebra above makes sense; interpreted as polynomials in \(x\), both
  sides are identical for all \(x\neq 1\).
\item
  \textbf{Book context.} The technique matches §1.2.3's emphasis on
  ``manipulating sums'' (including index shifts and basic identities)
  and is the same chain shown in the official ``Key steps'' for
  exercises 15--16.
\end{itemize}

Final formula

\[
\boxed{\;\sum_{j=0}^{n} j\,x^{j}
=\frac{x-(n+1)x^{n+1}+n x^{n+2}}{(1-x)^{2}}\quad (x\neq 1),\qquad
\sum_{j=0}^{n} j=\frac{n(n+1)}{2}\; \text{when }x=1.\;}
\]

\subsection{17 {[}M00{]}}\label{m00}

By definition, the index \(j\) in \(\sum_{j\in S} 1\) ranges over every
element of \(S\) exactly once, and the summand is the constant \(1\)
each time. Therefore the sum simply counts how many elements are in
\(S\):

\[
\sum_{j\in S} 1 \;=\; |S|.
\]

This can also be written using Iverson's bracket notation (introduced in
§1.2.3):

\[
\sum_{j\in S} 1 \;=\; \sum_{j} [\,j\in S\,],
\]

and by Eq. (17) in the text,

\[
\sum_{R(j)} a_j \;=\; \sum_{j} a_j [\,R(j)\,],
\]

so with \(a_j\equiv 1\) and \(R(j)\equiv (j\in S)\) the right-hand side
counts the number of \(j\) making \(R(j)\) true---i.e., \(|S|\).  

\textbf{Edge cases.}

\begin{itemize}
\tightlist
\item
  If \(S=\varnothing\), the sum is \(0\).
\item
  If \(S\) is finite with \(n\) elements, the sum is \(n\).
\item
  If \(S\) is infinite, the (ordinary) sum does not converge in the
  usual sense; informally one may say it is \(+\infty\).
\end{itemize}

\subsection{18 {[}M20{]}}\label{m20}

We start from the general interchange law (Eq. (9)) for finite sums:

\[
\sum_{R(i)}\ \sum_{S(i,j)} a_{ij}
\;=\;
\sum_{S'(j)}\ \sum_{R'(i,j)} a_{ij},
\]

where \(S'(j)\) is ``there exists an \(i\) such that both \(R(i)\) and
\(S(i,j)\) hold,'' and \(R'(i,j)\) is ``both \(R(i)\) and \(S(i,j)\)
hold.'' 

In this exercise we are told to take • \(R(i)\): ``\(n\) is a multiple
of \(i\)'' (i.e., \(i\mid n\)); • \(S(i,j)\): ``\(1\le j<i\).'' 

Determine the new limits:

\begin{enumerate}
\def\labelenumi{\arabic{enumi}.}
\tightlist
\item
  The outer relation \(S'(j)\). We need values of \(j\) for which there
  exists some divisor \(i\) of \(n\) with \(1\le j<i\). Taking \(i=n\)
  (which always divides \(n\)), this is possible exactly when
  \(1\le j\le n-1\). Thus
\end{enumerate}

\[
S'(j)\; \equiv\; 1\le j\le n-1 .
\]

\begin{enumerate}
\def\labelenumi{\arabic{enumi}.}
\setcounter{enumi}{1}
\tightlist
\item
  The inner relation \(R'(i,j)\). This must enforce both original
  conditions for fixed \(j\): \(i\mid n\) and \(1\le j<i\).
  Equivalently,
\end{enumerate}

\[
R'(i,j)\; \equiv\; i\mid n \ \text{ and }\ j<i \ (\text{hence } j+1\le i\le n).
\]

Therefore the interchanged sum is

\[
\boxed{\;
\sum_{\substack{i\mid n}}\ \sum_{j=1}^{\,i-1} a_{ij}
\;=\;
\sum_{j=1}^{\,n-1}\ \sum_{\substack{i\mid n\\ i>j}} a_{ij}\; }.
\]

(That is exactly Eq. (9) specialized to these relations.) 

Quick sanity check (optional): with \(n=6\) and \(a_{ij}\equiv 1\), the
left side counts \((1-1)+(2-1)+(3-1)+(6-1)=8\). The right side counts,
for \(j=1,2,3,4,5\), the number of divisors of \(6\) exceeding \(j\),
which are \(3,2,1,1,1\), again totaling \(8\). This confirms the limits
are consistent.

\subsection{19 {[}20{]}}\label{section-48}

Result

\[
\boxed{\displaystyle \sum_{j=m}^{n}\bigl(a_j-a_{j-1}\bigr)=a_n-a_{m-1}}
\]

(assuming \(m\le n\); if \(m>n\) the sum is empty and equals \(0\).)

Why it works (two equivalent derivations)

\textbf{1) Telescoping by expansion.} Write out a few terms:

\[
(a_m-a_{m-1})+(a_{m+1}-a_m)+\cdots+(a_n-a_{n-1}),
\]

and cancel adjacent \(+a_k\) with \(-a_k\). Everything cancels except
the first ``\(-a_{m-1}\)'' and the last ``\(+a_n\)'', yielding
\(a_n-a_{m-1}\). (This is the classic telescoping effect; the exercise
appears exactly as stated on p.~35 of the section. )

\textbf{2) Formal manipulation with Knuth's summation rules.} Let

\[
S=\sum_{j=m}^{n}(a_j-a_{j-1})
=\sum_{j=m}^{n}a_j\;-\;\sum_{j=m}^{n}a_{j-1}\quad\text{(split a sum; cf. Eq. (8)).}
\]

Now change index \(k=j-1\) in the second sum; as \(j\) runs
\(m\ldots n\), \(k\) runs \((m-1)\ldots(n-1)\) (index shift as in rule
(b)). Thus

\[
S=\sum_{k=m}^{n}a_k-\sum_{k=m-1}^{\,n-1}a_k
=(a_m+\cdots+a_n)-(a_{m-1}+\cdots+a_{n-1})=a_n-a_{m-1}.
\]

(See the linearity \(\sum(b_i+c_i)=\sum b_i+\sum c_i\) as Eq. (8), and
the standard index-shift technique used repeatedly in §1.2.3 examples. )

\textbf{Remark on edge cases.} Knuth's convention allows indices over
all integers (the bracket-notation discussion explicitly assumes \(a_j\)
is defined for all \(j\)); with that convention, the identity above is
exact for \(m\le n\), while for \(m>n\) the sum is \(0\). Using Iverson
brackets one can write a single formula:

\[
\sum_{j=m}^{n}(a_j-a_{j-1})=[m\le n]\,(a_n-a_{m-1}),
\]

where \([P]\) is \(1\) if \(P\) is true and \(0\) otherwise.

\subsection{20 {[}25{]}}\label{section-49}

\begin{enumerate}
\def\labelenumi{(\alph{enumi})}
\tightlist
\item
  Decimal (\(b=10\)) in \(\Sigma\)-notation
\end{enumerate}

Let

\[
C_n \;=\; \underbrace{12\ldots n}_{\text{as a decimal integer}}
\;=\;\sum_{j=1}^{n} j\,10^{\,n-j}.
\]

Then the observed pattern is the identity

\[
9\,C_n + (n+1)
\;=\;\sum_{k=0}^{n} 10^k \;=\;\frac{10^{\,n+1}-1}{9}.
\]

Thus, in pure \(\Sigma\)-form:

\[
9\sum_{j=1}^{n} j\,10^{\,n-j} + (n+1)
\;=\;\sum_{k=0}^{n} 10^k .
\]

(These right-hand numbers are the length-\(n{+}1\) ``repunit'' in base
10.)

\begin{enumerate}
\def\labelenumi{(\alph{enumi})}
\setcounter{enumi}{1}
\tightlist
\item
  Generalization to base \(b\)
\end{enumerate}

Define the ``concatenation'' number in base \(b\)

\[
C_n^{(b)} \;=\;\sum_{j=1}^{n} j\,b^{\,n-j}.
\]

Then for any \(b>1\),

\[
(b-1)\,C_n^{(b)} + (n+1)
\;=\;\sum_{k=0}^{n} b^k
\;=\;\frac{b^{\,n+1}-1}{b-1}.
\]

(If you want the literal ``digits \(1,2,\dots,n\)'' in base \(b\), you
should take \(n\le b-1\); the identity itself is purely algebraic and
holds for all \(n\).)

\begin{enumerate}
\def\labelenumi{(\alph{enumi})}
\setcounter{enumi}{2}
\tightlist
\item
  Proof using \(\sum jx^j\)
\end{enumerate}

We'll use the standard finite power-sum identity (given in this section
and also as Ex. 16): for \(x\ne 1\),

\[
\sum_{j=0}^{n} j\,x^j
= \frac{n x^{n+2}-(n+1)x^{n+1}+x}{(x-1)^2}.
\]

(Notice the \(j=0\) term is \(0\), so starting at \(j=1\) is
equivalent.) 

Let \(x=\frac{1}{b}\) and compute

\[
S_n(b)\;:=\;\sum_{j=1}^{n} j\,b^{\,n-j}
\;=\;b^n\sum_{j=1}^{n} j\,x^j.
\]

Plug the identity in:

\[
S_n(b)
= b^n\,\frac{n x^{n+2}-(n+1)x^{n+1}+x}{(x-1)^2}.
\]

Since \(x=\tfrac{1}{b}\), we have

\[
b^n x^{n+2}=b^{-2},\quad
b^n x^{n+1}=b^{-1},\quad
b^n x = b^{\,n-1},
\quad
(x-1)^2=\frac{(b-1)^2}{b^2}.
\]

Therefore

\[
S_n(b)
= \frac{n - (n+1)b + b^{\,n+1}}{\,b-1\,}.
\]

Multiply by \((b-1)\) and add \((n+1)\):

\[
(b-1)S_n(b) + (n+1)
= \frac{b^{\,n+1}-1}{\,b-1\,}
= \sum_{k=0}^{n} b^k,
\]

which is exactly the claimed formula from parts (a)--(b). QED.

Quick check (decimal, \(n=4\))

\[
9\sum_{j=1}^{4} j\,10^{4-j} +5
= 9(1\cdot10^3+2\cdot10^2+3\cdot10+4)+5
=9\cdot1234+5=11111,
\]

and \(\sum_{k=0}^{4}10^k=11111\), as required.

\subsection{21 {[}M25{]}}\label{m25-3}

\begin{enumerate}
\def\labelenumi{\arabic{enumi}.}
\tightlist
\item
  Start from the left--hand side and use (17) to convert ``conditioned''
  sums into ordinary sums with indicator brackets:
\end{enumerate}

\[
\sum_{R(j)} a_j+\sum_{S(j)} a_j
=\sum_j a_j\,[R(j)]\;+\;\sum_j a_j\,[S(j)]. \quad \text{(by (17))} \; \; \; \; 
\]

\begin{enumerate}
\def\labelenumi{\arabic{enumi}.}
\setcounter{enumi}{1}
\tightlist
\item
  Combine the two sums into one by distributivity over a sum of terms
  (use (8) with the ``domain'' being all integers \(j\)):
\end{enumerate}

\[
\sum_j a_j\,[R(j)]\;+\;\sum_j a_j\,[S(j)]
=\sum_j a_j\big([R(j)]+[S(j)]\big). \quad \text{(by (8))}
\]

\begin{enumerate}
\def\labelenumi{\arabic{enumi}.}
\setcounter{enumi}{2}
\tightlist
\item
  Use the elementary identity for indicator (Iverson) brackets:
\end{enumerate}

\[
[R]+[S] = [R\ \text{or}\ S] + [R\ \text{and}\ S].
\]

(This is checked case-by-case: if neither holds, both sides are \(0\);
if exactly one holds, both sides are \(1\); if both hold, both sides are
\(2\).)

Thus

\[
\sum_j a_j\big([R(j)]+[S(j)]\big)
= \sum_j a_j\big([R(j)\lor S(j)] + [R(j)\land S(j)]\big).
\]

\begin{enumerate}
\def\labelenumi{\arabic{enumi}.}
\setcounter{enumi}{3}
\tightlist
\item
  Split the sum again by (8), then convert back with (17):
\end{enumerate}

\[
\sum_j a_j[\,R(j)\lor S(j)\,] + \sum_j a_j[\,R(j)\land S(j)\,]
= \sum_{R(j)\ \text{or}\ S(j)} a_j + \sum_{R(j)\ \text{and}\ S(j)} a_j.
\]

This is exactly rule (d).

\subsection{22 {[}20{]}}\label{section-50}

Product analogs

\begin{enumerate}
\def\labelenumi{\arabic{enumi}.}
\tightlist
\item
  Change of variables (analog of Eq. (5)). If \(p\) is a bijection on
  the relevant index set, then
\end{enumerate}

\[
\prod_{R(i)} a_i
\;=\;
\prod_{R(j)} a_j
\;=\;
\prod_{R(p(j))} a_{p(j)}.
\]

This is the same reindexing principle as for sums; it just uses
associativity/commutativity of multiplication instead of addition. (For
infinite products, similar convergence caveats apply as noted for
sums.) 

\begin{enumerate}
\def\labelenumi{\arabic{enumi}.}
\setcounter{enumi}{1}
\tightlist
\item
  Interchange of order (analog of Eq. (7)). For rectangular domains,
\end{enumerate}

\[
\prod_{R(i)}\;\prod_{S(j)} a_{ij}
\;=\;
\prod_{S(j)}\;\prod_{R(i)} a_{ij}.
\]

For dependent domains (analog of Eq. (9)), if \(S(i,j)\) constrains
\((i,j)\), then

\[
\prod_{R(i)}\;\prod_{S(i,j)} a_{ij}
\;=\;
\prod_{S'(j)}\;\prod_{R'(i,j)} a_{ij},
\]

where \(S'(j)\) is ``there exists \(i\) with \(R(i)\) and \(S(i,j)\),''
and \(R'(i,j)\) is ``\(R(i)\) and \(S(i,j)\).'' 

\begin{enumerate}
\def\labelenumi{\arabic{enumi}.}
\setcounter{enumi}{2}
\tightlist
\item
  Distributivity over termwise combination (analog of Eq. (8)). With
  products, the natural distributive law is
\end{enumerate}

\[
\prod_{R(i)} \bigl(b_i c_i\bigr)
\;=\;
\Bigl(\prod_{R(i)} b_i\Bigr)\,
\Bigl(\prod_{R(i)} c_i\Bigr).
\]

(This mirrors \(\sum_{R(i)}(b_i+c_i)=\sum_{R(i)}b_i+\sum_{R(i)}c_i\) but
with ``\(+\)'' replaced by ``\(\times\)''.) 

\begin{enumerate}
\def\labelenumi{\arabic{enumi}.}
\setcounter{enumi}{3}
\tightlist
\item
  Domain manipulation (analog of Eq. (11)). For arbitrary relations
  \(R(j)\) and \(S(j)\),
\end{enumerate}

\[
\Bigl(\prod_{R(j)} a_j\Bigr)\,\Bigl(\prod_{S(j)} a_j\Bigr)
\;=\;
\Bigl(\prod_{R(j)\,\text{or}\,S(j)} a_j\Bigr)\,
\Bigl(\prod_{R(j)\,\text{and}\,S(j)} a_j\Bigr).
\]

Intuition: Elements in \(R\cap S\) are multiplied twice on the left, so
the right side is the product over the union (once) times the product
over the intersection (once more). A crisp proof uses Iverson brackets:
\(\prod_{R} a_j=\prod_j a_j^{[R(j)]}\), and the Boolean identity
\([R]+[S]=[R\lor S]+[R\land S]\).

These four identities are the exact multiplicative counterparts of Eqs.
(5), (7), (8), and (11) in the text.

\subsection{23 {[}10{]}}\label{section-51}

\begin{enumerate}
\def\labelenumi{\arabic{enumi})}
\tightlist
\item
  Empty sum must be \(0\)
\end{enumerate}

Section 1.2.3 uses Iverson brackets to write, for any sequence \(a_j\),

\[
\sum_{R(j)} a_j \;=\; \sum_j a_j\,[R(j)],
\]

and one of the fundamental ``set-algebra'' rules for sums is

\[
\sum_{R(j)\lor S(j)} a_j
= \sum_{R(j)} a_j + \sum_{S(j)} a_j - \sum_{R(j)\land S(j)} a_j.
\]

(This is a direct consequence of \([R]+[S]=[R\lor S]+[R\land S]\).) 

Take \(S\) to be the \textbf{always-false} condition. Then \(R\lor S\)
is just \(R\), and \(R\land S\) is false. Plugging in,

\[
\sum_{R(j)} a_j
= \sum_{R(j)} a_j + \underbrace{\sum_{\text{false}} a_j}_{=\alpha}
- \underbrace{\sum_{\text{false}} a_j}_{=\alpha}.
\]

This forces \(\alpha=0\). In words: to keep the union/splitting identity
valid when one side is empty, the empty sum must be \(0\).

\begin{enumerate}
\def\labelenumi{\arabic{enumi})}
\setcounter{enumi}{1}
\tightlist
\item
  Empty product must be \(1\)
\end{enumerate}

Products satisfy completely analogous rules in §1.2.3 (see the product
analogs noted right before Exercise 23): for any \(b_i, c_i\),

\[
\prod_{R(i)} b_i c_i \;=\; \Bigl(\prod_{R(i)} b_i\Bigr)\Bigl(\prod_{R(i)} c_i\Bigr),
\]

and for two conditions \(R,S\),

\[
\Bigl(\prod_{R(j)} a_j\Bigr)\Bigl(\prod_{S(j)} a_j\Bigr)
\;=\;
\Bigl(\prod_{R(j)\lor S(j)} a_j\Bigr)\Bigl(\prod_{R(j)\land S(j)} a_j\Bigr).
\]

Again set \(S\) to \textbf{false}. Then \(R\lor S=R\) and \(R\land S\)
is false, yielding

\[
\Bigl(\prod_{R(j)} a_j\Bigr)\, \underbrace{\prod_{\text{false}} a_j}_{=\beta}
=
\Bigl(\prod_{R(j)} a_j\Bigr)\, \underbrace{\prod_{\text{false}} a_j}_{=\beta}.
\]

This identity holds for any \(\beta\), so we need one more standard
requirement: When we concatenate disjoint index sets \(A\) and \(B\), we
want

\[
\prod_{j\in A\cup B} a_j = \Bigl(\prod_{j\in A} a_j\Bigr)\Bigl(\prod_{j\in B} a_j\Bigr)
\quad\text{(even if \(A\) or \(B\) is empty).}
\]

Now take \(B=\varnothing\). Then

\[
\prod_{j\in A} a_j \;=\; \Bigl(\prod_{j\in A} a_j\Bigr)\,\underbrace{\prod_{j\in \varnothing} a_j}_{=\beta}
\quad\text{for all sequences }\{a_j\}.
\]

Canceling \(\prod_{j\in A} a_j\) (choose any \(A\) with at least one
index and \(a_j\neq 0\) to make this cancellation indisputable), we get
\(\beta=1\).

(Equivalently, from \(\prod_{R(i)} b_i c_i = (\prod b_i)(\prod c_i)\),
set \(R\) false and \(b_i=c_i\) to get \(\beta=\beta^2\), hence
\(\beta\in\{0,1\}\); \(\beta=0\) is impossible because it would make
every product zero after adjoining an empty part, so \(\beta=1\).)

\subsection{24 {[}20{]}}\label{section-52}

Proof (by induction on the number of indices that satisfy \(R\))

Let \(S=\{j\in\mathbb{Z} : R(j)\}\). By hypothesis, \(|S|<\infty\). We
prove, for each \(m\ge 0\),

\[
\text{(P(m))}\qquad \forall S\ \text{with }|S|=m:\quad 
\log_b\!\Big(\prod_{j\in S} a_j\Big)=\sum_{j\in S}\log_b a_j .
\]

\textbf{Base case \(m=0\).} If \(S=\varnothing\), then by the section's
convention the empty product is \(1\) and the empty sum is \(0\). Hence

\[
\log_b\!\Big(\prod_{j\in\varnothing} a_j\Big)=\log_b(1)=0=\sum_{j\in\varnothing}\log_b a_j .
\]

This proves (P(0)). 

\emph{(If you prefer also to see \(m=1\): If \(S=\{j_0\}\), then
\(\log_b(\prod_{j\in S}a_j)=\log_b(a_{j_0})=\sum_{j\in S}\log_b a_j\).)}

\textbf{Inductive step.} Assume (P(m)) holds for some \(m\ge 0\). Let
\(S\) be any finite set with \(|S|=m+1\). Choose any \(j^{*}\in S\) and
set \(T=S\setminus\{j^{*}\}\) (so \(|T|=m\)). Then

\[
\prod_{j\in S} a_j \;=\; \Big(\prod_{j\in T} a_j\Big)\, a_{j^{*}}.
\]

Since all terms are positive, we can use
\(\log_b(xy)=\log_b x+\log_b y\):

\[
\log_b\!\Big(\prod_{j\in S} a_j\Big)
= \log_b\!\Big(\prod_{j\in T} a_j\Big) + \log_b(a_{j^{*}}).
\]

Apply the induction hypothesis (P(m)) to \(T\):

\[
\log_b\!\Big(\prod_{j\in T} a_j\Big)=\sum_{j\in T}\log_b a_j.
\]

Therefore

\[
\log_b\!\Big(\prod_{j\in S} a_j\Big)
= \sum_{j\in T}\log_b a_j \;+\; \log_b(a_{j^{*}})
= \sum_{j\in S}\log_b a_j,
\]

as required. This proves (P(m+1)).

By induction on \(|S|\), the identity holds for every finite \(S\).

\subsection{25 {[}15{]}}\label{section-53}

Diagnosis, line by line

\begin{itemize}
\tightlist
\item
  \textbf{Line 1 (OK):} \(\sum_{i=1}^{n}\sum_{j=1}^{n}1\) is a perfectly
  good double sum: it counts \(n\) terms of 1 for each of \(n\) outer
  values.
\item
  \textbf{Line 2 (Wrong -- variable capture):} Writing
  \(\sum_{i=1}^{n}\sum_{i=1}^{n}\) uses \textbf{the same dummy variable
  \(i\) twice}, once for the outer sum and again for the inner. Dummy
  (index) variables can be renamed arbitrarily, but they \textbf{must be
  distinct} when nested; otherwise the inner \(i\) \emph{shadows} the
  outer \(i\) and changes the meaning. Proper use of dummy indices is
  discussed when Knuth introduces summation notation.
\item
  \textbf{Line 3 (Wrong quantity):} Even if we correct the indices to
  \(\sum_{i=1}^{n}\sum_{j=1}^{n}1\), the inner sum \(\sum_{j=1}^{n}1\)
  equals the constant \(n\), \textbf{not} \(i\). So the third line
  should read \(\sum_{i=1}^{n} n\), not \(\sum_{i=1}^{n} i\). (This is
  exactly what the answer key hints at.) 
\item
  \textbf{Line 4 (OK in itself):} \(\sum_{i=1}^{n} i = n(n+1)/2\) is a
  correct identity---but it's irrelevant here because line 3 was wrong.
\end{itemize}

Correct derivation

\[
\sum_{i=1}^{n}\sum_{j=1}^{n}1
= \sum_{i=1}^{n}\bigl(\sum_{j=1}^{n}1\bigr)
= \sum_{i=1}^{n} n
= n\cdot n
= n^{2}.
\]

Intuitively, you're counting an \(n\times n\) grid of ones: there are
\(n^2\) of them.

\subsection{26 {[}25{]}}\label{section-54}

\textbf{Solution (counting exponents).} Each factor \(a_k\) appears in
\(P\) whenever a pair \((i,j)\) uses index \(k\) in either position.

\begin{itemize}
\tightlist
\item
  Pairs with \(i=k\): \(j\) can be \(k,k+1,\dots,n\) --- that's
  \(n-k+1\) occurrences of \(a_k\).
\item
  Pairs with \(j=k\): \(i\) can be \(0,1,\dots,k\) --- that's \(k+1\)
  occurrences of \(a_k\).
\end{itemize}

These counts are disjoint except at \((i,j)=(k,k)\); but that pair
contributes \(a_k^2\), exactly accounting for being ``seen'' once in
each count. Hence the total exponent of \(a_k\) in \(P\) is

\[
(n-k+1)+(k+1)=n+2.
\]

Therefore

\[
P=\prod_{k=0}^n a_k^{\,n+2} \;=\; \Bigl(\prod_{k=0}^n a_k\Bigr)^{\,n+2}.
\]

This is the closed form. 

\textbf{Alternative derivation (triangular-to-rectangular reindexing).}
Rewrite the triangular product as

\[
P=\prod_{i=0}^n\;\prod_{j=i}^n (a_i a_j)
=\Bigl(\prod_{i=0}^n a_i^{\,n-i+1}\Bigr)\Bigl(\prod_{j=0}^n a_j^{\,j+1}\Bigr)
=\prod_{k=0}^n a_k^{\,(n-k+1)+(k+1)}
=\Bigl(\prod_{k=0}^n a_k\Bigr)^{\,n+2}.
\]

This ``flip the order'' maneuver mirrors the technique shown in
\textbf{Example 2} of §1.2.3, where a triangular double sum is
rearranged by interchanging the roles of the indices. The same idea
works verbatim for products. 

\textbf{Sanity check (small \(n\)).} For \(n=2\) the pairs \((i,j)\)
with \(i\le j\) are \((0,0),(0,1),(0,2),(1,1),(1,2),(2,2)\). The product
is
\((a_0^2)\,(a_0a_1)\,(a_0a_2)\,(a_1^2)\,(a_1a_2)\,(a_2^2) = a_0^{4}a_1^{4}a_2^{4}=(a_0a_1a_2)^{4}\),
and \(n+2=4\), matching the formula.

\textbf{Answer.}

\[
\boxed{\displaystyle \prod_{0\le i\le j\le n} (a_i\,a_j)
=\left(\prod_{k=0}^n a_k\right)^{n+2}.}
\]

\subsection{27 {[}M20{]}}\label{m20-1}

Proof by induction on \(n\)

\textbf{Base case (\(n=1\)).} \((1-a_1)=1-a_1\). The inequality holds
with equality.

\textbf{Inductive step.} Assume for some \(n-1\ge1\) that

\[
\prod_{j=1}^{n-1}(1-a_j)\;\ge\;1-\sum_{j=1}^{n-1}a_j.
\]

Multiply both sides by the positive quantity \((1-a_n)\) (valid since
\(0<a_n<1\)):

\[
\prod_{j=1}^{n}(1-a_j)
\;\ge\;
\bigl(1-\sum_{j=1}^{n-1}a_j\bigr)(1-a_n)
=
1-\sum_{j=1}^{n}a_j
\;+\;
\Bigl(\sum_{j=1}^{n-1}a_j\Bigr)a_n.
\]

Because each \(a_j>0\), the last term is nonnegative (indeed positive
when \(n\ge2\)), hence

\[
\prod_{j=1}^{n}(1-a_j)\;\ge\;1-\sum_{j=1}^{n}a_j.
\]

This completes the induction. Moreover, when \(n\ge2\) the inequality is
\textbf{strict} because \(\bigl(\sum_{j=1}^{n-1}a_j\bigr)a_n>0\).

Alternative (algebraic) view

Expanding the product,

\[
\prod_{j=1}^n (1-a_j)
=1-\sum_j a_j+\sum_{i<j}a_ia_j-\sum_{i<j<k}a_ia_ja_k+\cdots+(-1)^n a_1\cdots a_n.
\]

All squarefree monomials have positive coefficients in absolute value,
and since every \(a_j\in(0,1)\), each monomial is nonnegative. Dropping
the alternating tail \(\sum_{i<j}a_ia_j-\sum_{i<j<k}\cdots\) (which
equals \(\bigl(\sum_{j=1}^{n-1}a_j\bigr)a_n\) plus further nonnegative
combinations after a regrouping) gives the same lower bound
\(1-\sum_j a_j\).

Remarks and tightness

\begin{itemize}
\tightlist
\item
  Equality holds for \(n=1\). For \(n\ge2\) and \(0<a_j<1\), the
  inequality is strict.
\item
  Intuition (probabilistic): If \(a_j\) are ``failure probabilities'' of
  independent events, \(\prod (1-a_j)\) is the probability that
  \textbf{none} occur, while \(1-\sum a_j\) is the classical union bound
  lower estimate. The algebraic inequality above mirrors that bound.
\end{itemize}

\subsection{28 {[}M22{]}}\label{m22-3}

Factor each term:

\[
1-\frac{1}{j^{2}}
=\frac{j^{2}-1}{j^{2}}
=\frac{(j-1)(j+1)}{j\cdot j}
=\frac{j-1}{j}\cdot\frac{j+1}{j}.
\]

Thus the product splits into two telescoping products:

\[
\prod_{j=2}^{n}\Bigl(1-\frac{1}{j^{2}}\Bigr)
=\Bigl(\prod_{j=2}^{n}\frac{j-1}{j}\Bigr)\Bigl(\prod_{j=2}^{n}\frac{j+1}{j}\Bigr).
\]

Compute each:

\begin{itemize}
\item
  First:

  \[
  \prod_{j=2}^{n}\frac{j-1}{j}
  =\frac{1}{2}\cdot\frac{2}{3}\cdot\frac{3}{4}\cdots\frac{n-1}{n}
  =\frac{1}{n}\quad\text{(complete cancellation)}.
  \]
\item
  Second:

  \[
  \prod_{j=2}^{n}\frac{j+1}{j}
  =\frac{3}{2}\cdot\frac{4}{3}\cdot\frac{5}{4}\cdots\frac{n+1}{n}
  =\frac{n+1}{2}.
  \]
\end{itemize}

Multiply:

\[
\prod_{j=2}^{n}\Bigl(1-\frac{1}{j^{2}}\Bigr)=\frac{1}{n}\cdot\frac{n+1}{2}
=\boxed{\frac{n+1}{2n}}.
\]

Quick checks:

\begin{itemize}
\tightlist
\item
  \(n=2\): \((1-\tfrac14)=\tfrac34\), and \((2+1)/(2\cdot2)=\tfrac34\).
\item
  \(n=3\): \((1-\tfrac14)(1-\tfrac19)=\tfrac34\cdot\tfrac89=\tfrac23\),
  and \((3+1)/(2\cdot3)=\tfrac23\).
\end{itemize}

As a corollary, letting \(n\to\infty\) gives
\(\prod_{j=2}^{\infty}\bigl(1-\tfrac{1}{j^{2}}\bigr)=\tfrac12\).

\subsection{29 {[}M30{]}}\label{m30-4}

Let

\[
A \;=\; \sum_{1\le i\le j\le k\le n} x_i x_j x_k,\qquad
T \;=\; \Bigl(\sum_{r=1}^n x_r\Bigr)^3=\sum_{i=1}^n\sum_{j=1}^n\sum_{k=1}^n x_i x_j x_k,
\]

and \(S_m=\sum_{r=1}^n x_r^m\) as above.

We'll classify the monomials in the full expansion of \(T\) by how many
indices are equal.

\textbf{1) All three indices distinct
\((i,j,k\ \text{all different})\).} Each unordered triple \(\{i,j,k\}\)
with \(i<j<k\) appears \(3!=6\) times in \(T\). If we denote

\[
D=\sum_{1\le i<j<k\le n} x_i x_j x_k,
\]

then the contribution of this class to \(T\) is \(6D\).

\textbf{2) Exactly two indices equal (say \(i=j\ne k\), or
permutations).} For each ordered pair \((i,k)\) with \(i\ne k\), one
gets a term \(x_i^2 x_k\). There are \(3\) choices for which index is
the singleton, so the total contribution is

\[
3\sum_{i\ne k} x_i^2 x_k
\;=\; 3\bigl(S_1S_2 - S_3\bigr),
\]

because \(S_1S_2=\sum_{i,k}x_i^2x_k\) and subtracting the \(i=k\) terms
(\(\sum_i x_i^3=S_3\)) leaves \(\sum_{i\ne k}x_i^2x_k\).

\textbf{3) All three indices equal.} The contribution is
\(\sum_{i=1}^n x_i^3 = S_3\).

Combining the three cases gives

\[
T \;=\; 6D \;+\; 3(S_1S_2 - S_3) \;+\; S_3
\;=\; 6D \;+\; 3S_1S_2 \;-\; 2S_3,
\]

hence

\[
D \;=\; \frac{T - 3S_1S_2 + 2S_3}{6}.
\tag{★}
\]

Now relate \(A\) to \(D\). The sum \(A\) (with \(i\le j\le k\)) counts:

\begin{itemize}
\tightlist
\item
  each distinct triple \(x_i x_j x_k\) (\(i<j<k\)) \textbf{once} ⇒
  contributes \(D\);
\item
  each ``two equal'' monomial \textbf{once} (either \(x_i^2x_k\) with
  \(i<k\) or \(x_ix_j^2\) with \(i<j\)). Summed over all pairs, these
  contribute exactly \(\sum_{i\ne k}x_i^2x_k = S_1S_2 - S_3\);
\item
  each ``all equal'' monomial \textbf{once} ⇒ contributes \(S_3\).
\end{itemize}

Therefore

\[
A \;=\; D \;+\; (S_1S_2 - S_3) \;+\; S_3 \;=\; D + S_1S_2.
\]

Insert \((★)\) to eliminate \(D\):

\[
A \;=\; \frac{T - 3S_1S_2 + 2S_3}{6} \;+\; S_1S_2
= \frac{T}{6} + \frac{S_1S_2}{2} + \frac{S_3}{3}.
\]

Recalling \(T=S_1^3\), we obtain the desired closed form:

\[
\boxed{\displaystyle
\sum_{1\le i\le j\le k\le n} x_i x_j x_k
= \frac{1}{6}\,S_1^{\,3}\;+\;\frac{1}{2}\,S_1S_2\;+\;\frac{1}{3}\,S_3.}
\]

This matches the formula Knuth cites when referring back to Exercise
1.2.3--29. 

Context / cross-check

Section 1.2.3 proves the analogous 2-index identity

\[
\sum_{0\le i\le j\le n} a_i a_j
=\tfrac12\Bigl(\bigl(\sum_{i=0}^n a_i\bigr)^2+\sum_{i=0}^n a_i^2\Bigr),
\]

using the same ``count the multiplicities'' idea; our argument is its
3-index analogue. 

\emph{(Optional generalization.)} The generating-function method in
§1.2.9 shows that the complete homogeneous symmetric sums
\(h_m=\sum_{1\le j_1\le\cdots\le j_m\le n} x_{j_1}\cdots x_{j_m}\)
satisfy

\[
\sum_{m\ge 0} h_m z^m
=\exp\!\Bigl(\sum_{k\ge 1} \frac{S_k}{k} z^k\Bigr),
\]

from which one can derive \(h_2\) and \(h_3\) as above (and a recurrence
\(h_n=\frac1n\sum_{k=1}^n S_k h_{n-k}\)). 

\subsection{30 {[}M23{]}}\label{m23-4}

Write

\[
A=\sum_i a_i c_i,\quad
B=\sum_i b_i d_i,\quad
C=\sum_i a_i d_i,\quad
D=\sum_i b_i c_i.
\]

We want \(AB-CD\).

Consider the double sum

\[
S\;=\;\tfrac12\sum_{i=0}^{n-1}\sum_{j=0}^{n-1}
\bigl(a_i b_j-a_j b_i\bigr)\bigl(c_i d_j-c_j d_i\bigr).
\]

Expand the product inside:

\[
\begin{aligned}
\bigl(a_i b_j-a_j b_i\bigr)\bigl(c_i d_j-c_j d_i\bigr)
&= a_i b_j c_i d_j \;-\; a_i b_j c_j d_i \\
&\quad - a_j b_i c_i d_j \;+\; a_j b_i c_j d_i. 
\end{aligned}
\]

Sum over \(i,j\) and group terms:

\[
\sum_{i,j} a_i b_j c_i d_j = \Bigl(\sum_i a_i c_i\Bigr)\Bigl(\sum_j b_j d_j\Bigr)=AB,
\]

\[
\sum_{i,j} a_i b_j c_j d_i = \Bigl(\sum_i a_i d_i\Bigr)\Bigl(\sum_j b_j c_j\Bigr)=CD,
\]

and by swapping dummy indices \(i\leftrightarrow j\) you see the other
two double--sum terms are identical to these. Hence

\[
S=\tfrac12\bigl(AB-CD-AB+CD\bigr)+\tfrac12\bigl(AB-CD\bigr)=AB-CD.
\]

Finally, note that the summand is \textbf{antisymmetric} in \(i,j\), so
the \(i=j\) terms vanish and every unordered pair \(\{i,j\}\) with
\(i\ne j\) appears twice with the same value. Therefore

\[
AB-CD
= \sum_{0\le i<j<n}\bigl(a_i b_j-a_j b_i\bigr)\bigl(c_i d_j-c_j d_i\bigr),
\]

which is the desired identity.

\subsection{31 {[}M20{]}}\label{m20-2}

Using Binet's identity (one-line derivation)

Binet's formula states

\[
\Big(\sum_j a_jx_j\Big)\Big(\sum_j b_jy_j\Big)
=\Big(\sum_j a_jy_j\Big)\Big(\sum_j b_jx_j\Big)
+\sum_{1\le j<k\le n}(a_jb_k-a_kb_j)(x_jy_k-x_ky_j).
\]

Choose

\[
a_j=1,\qquad b_j=u_j,\qquad x_j=1,\qquad y_j=v_j.
\]

Then

\[
(a_jb_k-a_kb_j)=(u_k-u_j)=-(u_j-u_k),\qquad
(x_jy_k-x_ky_j)=(v_k-v_j)=-(v_j-v_k),
\]

so the product in Binet's last sum equals \((u_j-u_k)(v_j-v_k)\).
Therefore

\[
\sum_{1\le j<k\le n}(u_j-u_k)(v_j-v_k)
=\Big(\sum_j a_jx_j\Big)\Big(\sum_j b_jy_j\Big)
-\Big(\sum_j a_jy_j\Big)\Big(\sum_j b_jx_j\Big).
\]

With our choices,

\[
\sum_j a_jx_j=n,\quad
\sum_j b_jy_j=\sum_{j=1}^n u_jv_j,\quad
\sum_j a_jy_j=\sum_{j=1}^n v_j,\quad
\sum_j b_jx_j=\sum_{j=1}^n u_j.
\]

Hence

\[
\boxed{\;
\sum_{1\le j<k\le n}(u_j-u_k)(v_j-v_k)
= n\sum_{j=1}^n u_jv_j-\Big(\sum_{j=1}^n u_j\Big)\Big(\sum_{j=1}^n v_j\Big)\; }.
\]

(This applies Binet exactly as suggested in the exercise statement.) 

Sanity check and interpretation

Write \(\bar u=\frac1n\sum u_j\) and \(\bar v=\frac1n\sum v_j\). Then
the right-hand side is

\[
n\sum_{j=1}^n\big(u_j-\bar u\big)\big(v_j-\bar v\big),
\]

so the exercise's sum equals \(n\) times the (unnormalized)
cross-deviation from the means---consistent with the case \(n=2\),
\((u_1-u_2)(v_1-v_2)=2(u_1-\bar u)(v_1-\bar v)+2(u_2-\bar u)(v_2-\bar v)\).

(Optional) Direct expansion (no Binet)

Expanding and regrouping,

\[
\begin{aligned}
S
&=\sum_{j<k}(u_jv_j+u_kv_k-u_jv_k-u_kv_j) \\
&=\underbrace{\sum_{j<k}u_jv_j+\sum_{j<k}u_kv_k}_{(n-1)\sum u_jv_j}
\;-\;\underbrace{\bigg(\sum_{j<k}u_jv_k+\sum_{j<k}u_kv_j\bigg)}_{\sum_{j\ne k}u_jv_k=(\sum u)(\sum v)-\sum uv}\\[2mm]
&=(n-1)\sum u_jv_j-\Big((\sum u)(\sum v)-\sum u_jv_j\Big) \\
&=n\sum u_jv_j-(\sum u)(\sum v).
\end{aligned}
\]

Same result, as expected.

\subsection{32 {[}M20{]}}\label{m20-3}

Direct (distributive) proof

Expand the product of \(n\) sums by repeated use of distributivity:

\[
\begin{aligned}
\prod_{j=1}^{n}\left(\sum_{i=1}^{m} a_{ij}\right)
&=\left(\sum_{i_1=1}^{m} a_{i_1\,1}\right)\!\left(\sum_{i_2=1}^{m} a_{i_2\,2}\right)\cdots\left(\sum_{i_n=1}^{m} a_{i_n\,n}\right)\\
&=\sum_{i_1=1}^{m}\sum_{i_2=1}^{m}\cdots\sum_{i_n=1}^{m}
\;a_{i_1\,1}a_{i_2\,2}\cdots a_{i_n\,n}\\
&=\sum_{1\le i_1,\ldots,i_n\le m}\;\prod_{j=1}^{n} a_{i_j\,j}.
\end{aligned}
\]

At each step you ``pull'' a sum out of a product via distributivity;
after doing this \(n\) times, you get a nested sum over all \(n\)-tuples
\((i_1,\ldots,i_n)\in\{1,\ldots,m\}^n\), with term
\(a_{i_1\,1}\cdots a_{i_n\,n}\). This is exactly the right-hand side.

Induction on \(n\) (equally rigorous)

Base case \(n=1\) is immediate:
\(\displaystyle \prod_{j=1}^{1}\sum_{i=1}^{m}a_{ij}=\sum_{i=1}^{m}a_{i1}\)
matches \(\displaystyle \sum_{1\le i_1\le m} a_{i_1\,1}\).

Inductive step: Assume the identity holds for \(n\). For \(n+1\),

\[
\begin{aligned}
\prod_{j=1}^{n+1}\left(\sum_{i=1}^{m} a_{ij}\right)
&=\left(\prod_{j=1}^{n}\sum_{i=1}^{m} a_{ij}\right)\left(\sum_{i=1}^{m} a_{i,n+1}\right)\\
&=\left(\sum_{1\le i_1,\ldots,i_n\le m}\prod_{j=1}^{n} a_{i_j\,j}\right)\left(\sum_{i_{n+1}=1}^{m} a_{i_{n+1},\,n+1}\right)\\
&=\sum_{1\le i_1,\ldots,i_{n+1}\le m}\prod_{j=1}^{n+1} a_{i_j\,j},
\end{aligned}
\]

using distributivity once more. Thus the statement holds for all \(n\).

Combinatorial viewpoint (why it's ``obvious'')

Think of choosing one term from each factor \(\sum_{i=1}^m a_{ij}\). A
complete choice is exactly an \(n\)-tuple \((i_1,\ldots,i_n)\),
contributing the product \(a_{i_1\,1}\cdots a_{i_n\,n}\). Summing over
all choices (left side expanded) equals summing those products over all
\(m^n\) tuples (right side).

Small sanity check (example)

Let \(m=2,n=3\). Then

\[
\left(a_{11}+a_{21}\right)\left(a_{12}+a_{22}\right)\left(a_{13}+a_{23}\right)
\]

expands to the eight terms

\[
a_{11}a_{12}a_{13}+a_{11}a_{12}a_{23}+a_{11}a_{22}a_{13}+a_{11}a_{22}a_{23}
+a_{21}a_{12}a_{13}+a_{21}a_{12}a_{23}+a_{21}a_{22}a_{13}+a_{21}a_{22}a_{23},
\]

which are precisely
\(\sum_{(i_1,i_2,i_3)\in\{1,2\}^3} a_{i_1\,1}a_{i_2\,2}a_{i_3\,3}\).

Notes and useful generalizations

\begin{enumerate}
\def\labelenumi{\arabic{enumi}.}
\tightlist
\item
  Different upper bounds per column. For integers \(m_j\ge1\),
\end{enumerate}

\[
\prod_{j=1}^{n}\left(\sum_{i=1}^{m_j} a_{ij}\right)
=\sum_{1\le i_j\le m_j\;\forall j}\;\prod_{j=1}^{n} a_{i_j\,j}.
\]

\begin{enumerate}
\def\labelenumi{\arabic{enumi}.}
\setcounter{enumi}{1}
\tightlist
\item
  Knuth's conditional-sum notation. In §1.2.3, sums/products over
  integer indices subject to relations are written \(\sum_{R(k)}\) and
  \(\prod_{R(k)}\); the right-hand side is a multi-index sum over the
  relation ``\(1\le i_j\le m\) for all \(j\).''
\item
  Infinite versions require convergence assumptions; with finite \(m,n\)
  no issues arise.
\end{enumerate}

\subsection{33 {[}M30{]}}\label{m30-5}

We'll prove a clean general law using \textbf{Lagrange interpolation}.

\begin{enumerate}
\def\labelenumi{\arabic{enumi})}
\tightlist
\item
  Lagrange basis and a key coefficient identity
\end{enumerate}

Let

\[
w(t)=\prod_{k=1}^{n}(t-x_k)
\qquad\text{and}\qquad
L_j(t)=\frac{\displaystyle\prod_{k\ne j}(t-x_k)}{\displaystyle\prod_{k\ne j}(x_j-x_k)}.
\]

The polynomials \(L_j\) are the \textbf{Lagrange basis}:
\(L_j(x_i)=\delta_{ij}\). Hence every polynomial \(p\) of degree
\(\le n-1\) satisfies the interpolation formula

\[
p(t)=\sum_{j=1}^n p(x_j)\,L_j(t).
\]

Look at the coefficient of \(t^{\,n-1}\) on both sides. Since

\[
L_j(t)=\frac{t^{\,n-1}+\text{(lower degree terms)}}{\prod_{k\ne j}(x_j-x_k)},
\]

the coefficient of \(t^{\,n-1}\) in \(L_j\) is exactly
\(\displaystyle\frac{1}{\prod_{k\ne j}(x_j-x_k)}\). Therefore:

\begin{quote}
\textbf{Lemma.} For any polynomial \(p\) of degree \(\le n-1\),

\[
[t^{n-1}]\,p(t)\;=\;\sum_{j=1}^{n}\frac{p(x_j)}{\displaystyle\prod_{k\ne j}(x_j-x_k)}.
\tag{∗}
\]
\end{quote}

This one-line identity is the engine of the whole exercise.

\begin{enumerate}
\def\labelenumi{\arabic{enumi})}
\setcounter{enumi}{1}
\tightlist
\item
  Apply the lemma to \(p(t)=t^r\)
\end{enumerate}

Case A: \(0\le r\le n-2\)

Here \(\deg p=r<n-1\), so the coefficient of \(t^{n-1}\) in \(p(t)=t^r\)
is \(0\). By (∗),

\[
\sum_{j=1}^{n}\frac{x_j^{\,r}}{\prod_{k\ne j}(x_j-x_k)}=0.
\]

Case B: \(r=n-1\)

Now \([t^{n-1}]\,t^{n-1}=1\). Applying (∗),

\[
\sum_{j=1}^{n}\frac{x_j^{\,n-1}}{\prod_{k\ne j}(x_j-x_k)}=1.
\]

So far we've settled all \(r\le n-1\).

\begin{enumerate}
\def\labelenumi{\arabic{enumi})}
\setcounter{enumi}{2}
\tightlist
\item
  Handle \(r=n\)
\end{enumerate}

When \(r=n\), \(\deg t^n=n\) so we can't apply (∗) directly. Write
\(t^n\) as

\[
t^n = w(t) + R(t),
\]

where \(R(t)\) is the remainder of the division of \(t^n\) by \(w(t)\).
Then \(\deg R\le n-1\) and, crucially,
\(R(x_j)=t^n\!\big|_{t=x_j}=x_j^n\) for each \(j\).

Therefore, applying (∗) to \(R\),

\[
\sum_{j=1}^{n}\frac{x_j^{\,n}}{\prod_{k\ne j}(x_j-x_k)}
=\sum_{j=1}^{n}\frac{R(x_j)}{\prod_{k\ne j}(x_j-x_k)}
=[t^{n-1}]\,R(t).
\]

But we also know the explicit leading terms of \(w(t)\) by Viète's
formulas:

\[
w(t)=\prod_{k=1}^{n}(t-x_k)
=t^{n}-\Big(\sum_{k=1}^{n}x_k\Big)\,t^{n-1}+\cdots .
\]

Hence

\[
R(t)=t^n-w(t)=\Big(\sum_{k=1}^{n}x_k\Big)\,t^{n-1}+\cdots ,
\]

so \([t^{n-1}]\,R(t)=\sum_{k=1}^{n}x_k\). Therefore

\[
\sum_{j=1}^{n}\frac{x_j^{\,n}}{\prod_{k\ne j}(x_j-x_k)}=\sum_{j=1}^{n}x_j,
\]

as required.

This proves all three cases.

\begin{enumerate}
\def\labelenumi{\arabic{enumi})}
\setcounter{enumi}{3}
\tightlist
\item
  Checks against the \(n=3\) ``Dr.~Matrix'' identities
\end{enumerate}

For \(n=3\), the denominators are \((a-b)(a-c)\), \((b-a)(b-c)\),
\((c-a)(c-b)\). Our general result gives:

\begin{itemize}
\tightlist
\item
  \(r=0,1\): sums equal \(0\).
\item
  \(r=2\): sum equals \(1\).
\item
  \(r=3\): sum equals \(a+b+c\).
\end{itemize}

These match the four displayed identities in the problem statement.

\subsection{34 {[}M25{]}}\label{m25-4}

Write the denominator in a closed form. For fixed \(k\),

\[
\prod_{\substack{1\le r\le n\\ r\ne k}}(k-r)
=\Bigl(\prod_{r=1}^{k-1}(k-r)\Bigr)\Bigl(\prod_{r=k+1}^{n}(k-r)\Bigr)
=(k-1)!\,(-1)^{\,n-k}(n-k)! .
\]

Hence

\[
\frac{1}{\prod_{r\ne k}(k-r)}
=\frac{(-1)^{\,n-k}}{(k-1)!(n-k)!}
=\frac{1}{(n-1)!}\,(-1)^{\,n-k}\binom{n-1}{k-1}.
\]

Define the polynomial in \(k\) (with \(x,m\) treated as parameters)

\[
p(k)\;:=\;\prod_{\substack{1\le r\le n\\ r\ne m}}(x+k-r).
\]

This is a monic polynomial of degree \(n-1\) in \(k\).

With the weight above, the left-hand side becomes

\[
S(x)=\sum_{k=1}^{n}\frac{p(k)}{\prod_{r\ne k}(k-r)}
=\frac{1}{(n-1)!}\sum_{k=1}^{n}(-1)^{\,n-k}\binom{n-1}{k-1}\,p(k).
\]

But the binomial-weighted sum is exactly a forward difference:

\[
\Delta^{\,n-1}p(1)
=\sum_{j=0}^{n-1}(-1)^{\,n-1-j}\binom{n-1}{j}\,p(1+j)
=\sum_{k=1}^{n}(-1)^{\,n-k}\binom{n-1}{k-1}\,p(k).
\]

Therefore

\[
S(x)=\frac{\Delta^{\,n-1}p(1)}{(n-1)!}.
\]

Finally, since \(p(k)\) is \textbf{monic of degree \(n-1\)}, its
\((n-1)\)-st forward difference is the constant \((n-1)!\) (discrete
analogue of ``the \((n-1)\)-st derivative of a monic degree \(n-1\)
polynomial is \((n-1)!\)''). Hence

\[
S(x)=\frac{(n-1)!}{(n-1)!}=1,
\]

which proves the identity.

\subsection{35 {[}HM20{]}}\label{hm20-3}

Knuth introduces the supremum notation \(\displaystyle \sup_{R(j)} a_j\)
(least upper bound of the \(a_j\) with \(R(j)\) true) and asks you to
adapt the four basic sum-manipulation rules (a)--(d) from §1.2.3 to this
notation. In particular, discuss the analogue of rule (a)

\[
\bigl(\sup_{R(i)} a_i\bigr)+\bigl(\sup_{S(j)} b_j\bigr)=\sup_{R(i)}\sup_{S(j)}(a_i+b_j),
\]

and specify what \(\sup_{R(j)} a_j\) should mean when no \(j\) satisfies
\(R\).

Throughout, it is convenient to work in the \textbf{extended real line}
\(\overline{\mathbb{R}}=\mathbb{R}\cup\{-\infty,+\infty\}\).

\begin{enumerate}
\def\labelenumi{\arabic{enumi})}
\tightlist
\item
  Definition for the empty index set
\end{enumerate}

To keep all identities well-behaved (especially the ``union'' rule
below), define

\[
\sup_{R(j)} a_j \;=\; -\infty \quad\text{if no }j\text{ satisfies }R(j).
\]

This mirrors the way \(\sum_{R(j)} a_j\) is defined to be \(0\) when the
set is empty; for sups, \(-\infty\) is the neutral element for \(\max\).
With this choice, every rule below continues to hold even in degenerate
cases. (Knuth's official answer also adopts \(-\infty\).) 

\begin{enumerate}
\def\labelenumi{\arabic{enumi})}
\setcounter{enumi}{1}
\tightlist
\item
  Analogue of rule (a): ``distributive law'' over addition
\end{enumerate}

For sums, rule (a) states
\((\sum_{R(i)} a_i)(\sum_{S(j)} b_j)=\sum_{R(i)}\sum_{S(j)} a_i b_j\)
(Eq. (4)). For suprema, the correct additive analogue is

\[
\boxed{\;\bigl(\sup_{R(i)} a_i\bigr)+\bigl(\sup_{S(j)} b_j\bigr)
\;=\; \sup_{R(i)}\sup_{S(j)}\bigl(a_i+b_j\bigr)\;}
\tag{A}
\]

\textbf{Proof.} Let \(A=\{a_i: R(i)\}\) and \(B=\{b_j: S(j)\}\).

\begin{itemize}
\item
  For any \(i,j\), \(a_i\le\sup A\) and \(b_j\le\sup B\), hence
  \(a_i+b_j\le \sup A+\sup B\). Taking the supremum over all \(i,j\)
  gives \(\sup_{i,j}(a_i+b_j)\le \sup A+\sup B\).
\item
  Conversely, fix \(\varepsilon>0\) and pick \(i_0\) with
  \(a_{i_0} > \sup A-\varepsilon/2\) and \(j_0\) with
  \(b_{j_0} > \sup B-\varepsilon/2\). Then
  \(a_{i_0}+b_{j_0} > \sup A+\sup B-\varepsilon\), so
  \(\sup_{i,j}(a_i+b_j)\ge \sup A+\sup B-\varepsilon\). Let
  \(\varepsilon\to0\).
\end{itemize}

If \(\sup A=+\infty\) or \(\sup B=+\infty\), both sides are \(+\infty\);
if one set is empty, \(\sup A=-\infty\) etc., and both sides evaluate to
\(-\infty\) in \(\overline{\mathbb{R}}\). For \textbf{finite} sets, (A)
becomes the trivial identity \(a+\max(b,c)=\max(a+b,a+c)\). (Knuth hints
at exactly this reduction.) 

A closely related multiplicative variant (useful later) also holds:

\begin{quote}
If all \(a_i,b_j\ge 0\), then
\(\displaystyle \bigl(\sup_{R(i)} a_i\bigr)\bigl(\sup_{S(j)} b_j\bigr)
\;=\; \sup_{R(i)}\sup_{S(j)} (a_i b_j).\)
\end{quote}

(Again, prove both inequalities; nonnegativity is essential.) This is
the supremum analogue of the product-of-sums law (Eq. (4)). 

\begin{enumerate}
\def\labelenumi{\arabic{enumi})}
\setcounter{enumi}{2}
\tightlist
\item
  Analogue of rule (b): change of variable
\end{enumerate}

For sums, rule (b) allows
\(\sum_{R(i)} a_i=\sum_{R(j)} a_j=\sum_{R(p(j))} a_{p(j)}\) when \(p\)
is a permutation of the relevant indices (Eq. (5)). For sups:

\[
\boxed{\;\sup_{R(i)} a_i \;=\; \sup_{R(j)} a_j \;=\; \sup_{R(p(j))} a_{p(j)}\;}
\tag{B}
\]

provided \(p\) is a \textbf{bijection} on the index set \(\{i:R(i)\}\).
The proof is immediate because the set of values being ``supremed over''
is unchanged under such a relabeling. (If \(p\) misses some relevant
\(i\), you only get ``\(\le\)''.) 

\begin{enumerate}
\def\labelenumi{\arabic{enumi})}
\setcounter{enumi}{3}
\tightlist
\item
  Analogue of rule (c): interchanging order
\end{enumerate}

For sums, rule (c) says
\(\sum_{R(i)}\sum_{S(j)} a_{ij} = \sum_{S(j)}\sum_{R(i)} a_{ij}\) (Eq.
(7), and the more general Eq. (9)). For sups:

\[
\boxed{\;\sup_{R(i)}\sup_{S(j)} c_{ij}
\;=\; \sup_{S(j)}\sup_{R(i)} c_{ij}
\;=\; \sup_{\substack{R(i)\\ S(j)}} c_{ij}\;}
\tag{C}
\]

because all three expressions are the least upper bound of the same set
\(\{c_{ij}:R(i)\land S(j)\}\). (No absolute-convergence caveats are
needed for sups.) 

\begin{enumerate}
\def\labelenumi{\arabic{enumi})}
\setcounter{enumi}{4}
\tightlist
\item
  Analogue of rule (d): manipulating the domain (union/intersection)
\end{enumerate}

For sums, rule (d) is the inclusion--exclusion style identity

\[
\sum_{R(j)} a_j + \sum_{S(j)} a_j
= \sum_{R(j)\,\text{or}\,S(j)} a_j + \sum_{R(j)\,\text{and}\,S(j)} a_j
\quad\text{(Eq. (11)).}
\]

For \textbf{suprema}, double-counting is irrelevant, and the rule
simplifies to

\[
\boxed{\;\sup_{R(j)\,\text{or}\,S(j)} a_j
\;=\; \max\Bigl(\sup_{R(j)} a_j,\ \sup_{S(j)} a_j\Bigr). \;}
\tag{D}
\]

This is simply the fact that the supremum over a union of sets equals
the maximum of the individual suprema; with our convention
\(\sup\varnothing=-\infty\), it holds even when one condition selects
nothing. (Knuth's answer states exactly this simplification.)

\subsection{36 {[}M23{]}}\label{m23-5}

Solution 1 --- Elementary row/column operations (no eigenvalues)

Write the rows of \(A\) as

\[
R_1=[x{+}y,\,y,\,\dots,\,y],\qquad
R_i=[y,\,\dots,\,\underset{i}{x{+}y},\,\dots,\,y]\ (i\ge2).
\]

\begin{enumerate}
\def\labelenumi{\arabic{enumi}.}
\tightlist
\item
  \textbf{Make the matrix sparse.} For each \(i=2,\dots,n\) replace
  \(R_i\leftarrow R_i-R_1\). This preserves the determinant.
\end{enumerate}

\begin{itemize}
\item
  In the new rows \(i\ge2\):

  \begin{itemize}
  \tightlist
  \item
    Column \(1\): \(y-(x+y)=-x\).
  \item
    Column \(i\): \((x+y)-y= x\).
  \item
    Other columns: \(y-y=0\).
  \end{itemize}
\end{itemize}

So each \(R_i\) (for \(i\ge2\)) now has exactly two nonzeros: \(-x\) in
column 1 and \(x\) in column \(i\).

\begin{enumerate}
\def\labelenumi{\arabic{enumi}.}
\setcounter{enumi}{1}
\tightlist
\item
  \textbf{Factor \(x\) from rows \(2,\dots,n\).} Pulling a common factor
  \(x\) from each of those \(n-1\) rows multiplies the determinant by
  \(x^{n-1}\). After factoring, replace each such row by
  \([-1\ \ 0\ \dots\ 1\text{ at col }i\ \dots\ 0]\).
\end{enumerate}

At this point the matrix has the form

\[
B=
\begin{bmatrix}
x+y & y & y & \cdots & y \\
-1  & 1 & 0 & \cdots & 0 \\
-1  & 0 & 1 & \cdots & 0 \\
\vdots & \vdots & & \ddots & \vdots \\
-1  & 0 & 0 & \cdots & 1
\end{bmatrix},
\qquad
\det A = x^{n-1}\det B.
\]

\begin{enumerate}
\def\labelenumi{\arabic{enumi}.}
\setcounter{enumi}{2}
\tightlist
\item
  \textbf{Column sweep.} Add columns \(2,3,\dots,n\) into column \(1\)
  (an operation that does \textbf{not} change the determinant):
\end{enumerate}

\[
C_1 \leftarrow C_1 + \sum_{j=2}^{n} C_j.
\]

\begin{itemize}
\tightlist
\item
  First-row, first-column entry becomes \((x+y)+(n-1)y=x+ny\).
\item
  For each \(i\ge2\), the first-column entry \(-1\) is increased by the
  \(1\) sitting in column \(i\) of row \(i\), giving \(0\).
\end{itemize}

Thus

\[
B'=
\begin{bmatrix}
x+ny & y & y & \cdots & y \\
0    & 1 & 0 & \cdots & 0 \\
0    & 0 & 1 & \cdots & 0 \\
\vdots & \vdots & & \ddots & \vdots \\
0    & 0 & 0 & \cdots & 1
\end{bmatrix}.
\]

Now \(B'\) is block--upper--triangular with the lower-right block
\(I_{n-1}\). Hence

\[
\det B' = (x+ny)\cdot \det(I_{n-1}) = x+ny.
\]

Combining with step 2:

\[
\boxed{\det A = x^{n-1}(x+ny)}.
\]

\textbf{Edge cases.}

\begin{itemize}
\tightlist
\item
  \(n=1\): \(\det[x+y]=x+y\), and the formula gives
  \(x^{0}(x+1\cdot y)=x+y\).
\item
  \(x=0\): \(A=y\mathbf{1}\mathbf{1}^\top\) has rank \(1\) for \(n>1\),
  so \(\det A=0\); the formula gives \(0^{n-1}(0+ny)=0\) (for \(n>1\)),
  consistent.
\end{itemize}

Solution 2 --- One-line spectral proof

Write \(A=xI+yJ\), where \(J=\mathbf{1}\mathbf{1}^\top\) is the all-ones
matrix. \(J\) has eigenvalue \(n\) on \(\mathbf{1}\) and \(0\) on the
\((n{-}1)\)-dimensional subspace orthogonal to \(\mathbf{1}\). Therefore
the eigenvalues of \(A\) are

\[
x+ny \quad(\text{once}), \qquad x \quad(\text{multiplicity }n-1).
\]

Thus \(\det A\) (product of eigenvalues) equals \(x^{n-1}(x+ny)\) as
claimed.

\subsection{37 {[}M24{]}}\label{m24}

The exercise defines the Vandermonde matrix by \(a_{ij}=x_j^{\,i}\) (so
the first row is \(x_1,\dots,x_n\); the second row
\(x_1^2,\dots,x_n^2\); \ldots; the \(n\)th row \(x_1^n,\dots,x_n^n\))
and asks you to show that

\[
\det(a_{ij})=\Bigl(\prod_{j=1}^n x_j\Bigr)\;\prod_{1\le i<j\le n}(x_j-x_i).
\]

Step 0: Reduce to the ``classical'' Vandermonde

Factor \(x_j\) out of column \(j\) for each \(j\). Determinants scale by
a column factor, so

\[
\det\!\bigl[x_j^{\,i}\bigr]_{i,j=1}^n
=\Bigl(\prod_{j=1}^n x_j\Bigr)\;\det\!\bigl[x_j^{\,i-1}\bigr]_{i,j=1}^n.
\]

Thus it suffices to prove the classical Vandermonde formula

\[
\Delta_n(x_1,\dots,x_n)\;:=\;\det\!\bigl[x_j^{\,i-1}\bigr]_{i,j=1}^n
\;=\;\prod_{1\le i<j\le n}(x_j-x_i).
\tag{★}
\]

Step 1: Structure of \(\Delta_n\)

Two immediate properties of \(\Delta_n\):

\begin{enumerate}
\def\labelenumi{\arabic{enumi}.}
\item
  Alternation. Swapping \(x_p\) and \(x_q\) swaps two columns and flips
  the sign of the determinant; hence \(\Delta_n\) is alternating in the
  \(x_j\)'s.
\item
  Vanishing on collisions. If \(x_p=x_q\) with \(p\neq q\), columns
  \(p\) and \(q\) coincide, so \(\Delta_n=0\). Therefore each factor
  \((x_q-x_p)\) must divide \(\Delta_n\) as a polynomial.
\end{enumerate}

By unique factorization in \(\mathbb{R}[x_1,\dots,x_n]\) (or
\(\mathbb{C}[x_1,\dots,x_n]\)), these two facts together imply that
there exists a constant \(C_n\) (depending only on \(n\)) with

\[
\Delta_n(x_1,\dots,x_n)=C_n\prod_{1\le i<j\le n}(x_j-x_i).
\tag{1}
\]

Step 2: Determine \(C_n\) by leading-term comparison

View \(\Delta_n\) as a polynomial in the single variable \(x_n\) (with
parameters \(x_1,\dots,x_{n-1}\)). Expanding \(\det[x_j^{\,i-1}]\) along
the last column shows that the coefficient of \(x_n^{\,n-1}\) is exactly
\(\Delta_{n-1}(x_1,\dots,x_{n-1})\): only the entry \(x_n^{\,n-1}\) in
the bottom row contributes to that top degree, and its cofactor is the
\((n-1)\times(n-1)\) classical Vandermonde on the first \(n-1\)
variables.

On the right side of (1), as a polynomial in \(x_n\),

\[
\prod_{1\le i<j\le n}(x_j-x_i)
=\Bigl(\prod_{1\le i<j\le n-1}(x_j-x_i)\Bigr)\cdot\prod_{k=1}^{n-1}(x_n-x_k),
\]

whose coefficient of \(x_n^{\,n-1}\) is
\(\prod_{1\le i<j\le n-1}(x_j-x_i)=\Delta_{n-1}(x_1,\dots,x_{n-1})\) by
the induction hypothesis.

Comparing the top-degree coefficients on both sides of (1) gives

\[
\Delta_{n-1}(x_1,\dots,x_{n-1}) \;=\; C_n\;\Delta_{n-1}(x_1,\dots,x_{n-1}),
\]

hence \(C_n=1\). (The base \(n=1\) is \(\Delta_1=1\), and \(n=2\) gives
\(\Delta_2=x_2-x_1\), so the induction anchors are immediate.)

Therefore (★) holds:

\[
\Delta_n(x_1,\dots,x_n)=\prod_{1\le i<j\le n}(x_j-x_i).
\]

Step 3: Restore the factored \(x_j\)'s

Returning to the original matrix with \(a_{ij}=x_j^{\,i}\),

\[
\det[a_{ij}]
=\Bigl(\prod_{j=1}^n x_j\Bigr)\;\det[x_j^{\,i-1}]
=\Bigl(\prod_{j=1}^n x_j\Bigr)\;\prod_{1\le i<j\le n}(x_j-x_i),
\]

which is exactly what the exercise asked to show. 

Quick checks (sanity)

\begin{itemize}
\tightlist
\item
  \(n=2\):
  \(\det\begin{pmatrix}x_1&x_2\\ x_1^2&x_2^2\end{pmatrix}=x_1x_2(x_2-x_1)\).
\item
  \(n=3\): factoring \(x_1x_2x_3\) from columns and using the classical
  \(3\times3\) Vandermonde yields
  \(x_1x_2x_3(x_2-x_1)(x_3-x_1)(x_3-x_2)\).
\end{itemize}

\subsection{38 {[}M25{]}}\label{m25-5}

We proceed by induction on \(n\).

\textbf{Base \(n=1\).} \(\det[1/(x_1+y_1)]=1/(x_1+y_1)\). The right-hand
side also equals \(1/(x_1+y_1)\) (empty products are \(1\)). ✓

\textbf{Inductive step.} Assume the formula for size \(n-1\). For the
\(n\times n\) matrix:

\begin{enumerate}
\def\labelenumi{\arabic{enumi}.}
\tightlist
\item
  \textbf{Make columns \(2,\dots,n\) differenced against column 1.}
  Replace \(C_j\leftarrow C_j-C_1\) for \(j=2,\dots,n\). Entrywise,
\end{enumerate}

\[
\frac{1}{x_i+y_j}-\frac{1}{x_i+y_1}
= \frac{y_1-y_j}{(x_i+y_j)(x_i+y_1)}.
\]

Factor \((y_1-y_j)\) from each column \(j\ge2\), and then factor
\((x_i+y_1)^{-1}\) from each row \(i\). After these factorizations the
matrix becomes

\[
\begin{bmatrix}
1 & \dfrac{1}{x_1+y_2} & \cdots & \dfrac{1}{x_1+y_n}\\[4pt]
1 & \dfrac{1}{x_2+y_2} & \cdots & \dfrac{1}{x_2+y_n}\\
\vdots & \vdots & \ddots & \vdots\\
1 & \dfrac{1}{x_n+y_2} & \cdots & \dfrac{1}{x_n+y_n}
\end{bmatrix}
\]

and we have pulled out the factor

\[
\Biggl(\prod_{j=2}^n (y_1-y_j)\Biggr)\Biggl(\prod_{i=1}^n \frac{1}{x_i+y_1}\Biggr).
\]

(These are exactly the operations hinted in Knuth's answer: subtract
column 1 from the others, factor the \(y\)-differences and the
\((x_i+y_1)^{-1}\)s. )

\begin{enumerate}
\def\labelenumi{\arabic{enumi}.}
\setcounter{enumi}{1}
\tightlist
\item
  \textbf{Zero out the first column below the top.} Replace
  \(R_i\leftarrow R_i-R_1\) for \(i=2,\dots,n\). In columns \(j\ge2\),
\end{enumerate}

\[
\frac{1}{x_i+y_j}-\frac{1}{x_1+y_j}
= \frac{x_1-x_i}{(x_i+y_j)(x_1+y_j)}.
\]

Expand along the first column (which now has a \(1\) on top and zeros
below). The remaining \((n-1)\times(n-1)\) block
\(B=(b_{ij})_{2\le i,j\le n}\) has

\[
b_{ij}=\frac{x_1-x_i}{(x_i+y_j)(x_1+y_j)}.
\]

Factor \((x_1-x_i)\) from each row \(i=2,\dots,n\) and
\((x_1+y_j)^{-1}\) from each column \(j=2,\dots,n\). The block reduces
to \(\bigl(1/(x_i+y_j)\bigr)_{2\le i,j\le n}\), i.e.~the same Cauchy
form on the smaller sets \(\{x_2,\dots,x_n\}\), \(\{y_2,\dots,y_n\}\).
The factor you have pulled out in this step is

\[
\Biggl(\prod_{i=2}^n (x_1-x_i)\Biggr)\Biggl(\prod_{j=2}^n \frac{1}{x_1+y_j}\Biggr).
\]

(These are the ``subtract row 1 from the others; factor
\(x\)-differences and \((x_1+y_j)^{-1}\)'' steps; after them only a
smaller Cauchy determinant remains. )

Putting the two stages together,

\[
\det_n(\{x\},\{y\})
= \frac{\displaystyle\prod_{i=2}^n (x_1-x_i)\;\prod_{j=2}^n (y_1-y_j)}
       {\displaystyle\Bigl(\prod_{i=1}^n (x_i+y_1)\Bigr)\Bigl(\prod_{j=2}^n (x_1+y_j)\Bigr)}
\;\det_{n-1}(\{x_2,\dots,x_n\},\{y_2,\dots,y_n\}).
\]

By the induction hypothesis,

\[
\det_{n-1}
= \frac{\displaystyle\prod_{2\le i<k\le n}(x_i-x_k)\;\prod_{2\le j<\ell\le n}(y_j-y_\ell)}
       {\displaystyle\prod_{i=2}^n\prod_{j=2}^n (x_i+y_j)}.
\]

Multiplying the factors and simplifying yields

\[
\det_n
= \frac{\displaystyle\prod_{1\le i<k\le n}(x_i-x_k)\;\prod_{1\le j<\ell\le n}(y_j-y_\ell)}
       {\displaystyle\prod_{i=1}^n\prod_{j=1}^n (x_i+y_j)},
\]

which is precisely the claimed formula. This completes the induction. ∎

Quick check for \(n=2\)

\[
\det\!\begin{bmatrix}
\frac1{x_1+y_1} & \frac1{x_1+y_2}\\[2pt]
\frac1{x_2+y_1} & \frac1{x_2+y_2}
\end{bmatrix}
= \frac{(x_1-x_2)(y_1-y_2)}{(x_1+y_1)(x_1+y_2)(x_2+y_1)(x_2+y_2)}.
\]

Matches the formula.

\subsection{39 {[}M23{]}}\label{m23-6}

Here ``combinatorial matrix'' means the \(n\times n\) matrix

\[
A=(a_{ij}),\qquad a_{ij}=y+\delta_{ij}x,
\]

i.e., every off--diagonal entry is \(y\) and every diagonal entry is
\(x+y\). Equivalently

\[
A=xI+yJ,
\]

where \(I\) is the identity and \(J\) is the all-ones matrix.
(Definition and exercise are in §1.2.3.) 

We will show that \(A^{-1}\) also has the ``combinatorial'' form
\(B=(b_{ij})\) with

\[
b_{ij}=\frac{-\,y+\delta_{ij}(x+ny)}{x(x+ny)}.
\]

(That is exactly what the exercise asks you to prove.) 

Two basic facts about \(J\):

\begin{enumerate}
\def\labelenumi{\arabic{enumi}.}
\tightlist
\item
  \(J\) commutes with \(I\) (trivial) and with any polynomial in \(J\).
\item
  \(J^2=nJ\) (since each entry of \(J^2\) is the sum of \(n\) ones).
\end{enumerate}

Because \(A=xI+yJ\) lies in the 2-dimensional commutative algebra
\(\{aI+bJ\}\) closed under multiplication, its inverse---if it
exists---must also lie in that algebra. So make the ansatz

\[
A^{-1}=aI+bJ
\]

for scalars \(a,b\). Multiply and use \(J^2=nJ\):

\[
(xI+yJ)(aI+bJ)
=xa\,I+(xb+ya+nyb)\,J.
\]

For this to be the identity \(I\), we must have the coefficient of \(I\)
equal to 1 and the coefficient of \(J\) equal to 0:

\[
xa=1,\qquad xb+ya+nyb=0.
\]

Solve:

\[
a=\frac{1}{x},\qquad b=-\,\frac{ya}{x+ny}=-\,\frac{y}{x(x+ny)}.
\]

Hence

\[
A^{-1}= \frac{1}{x}\,I-\frac{y}{x(x+ny)}\,J.
\]

Reading off entries (remember \(J\) has all ones),

\[
b_{ij}= -\,\frac{y}{x(x+ny)}+\delta_{ij}\,\frac{1}{x}
= \frac{-\,y+\delta_{ij}(x+ny)}{x(x+ny)}.
\]

Thus \(A^{-1}\) has constant off-diagonal entries \(-y/\{x(x+ny)\}\) and
constant diagonal entries
\(\{-y+(x+ny)\}/\{x(x+ny)\}=(x+(n-1)y)/\{x(x+ny)\}\), i.e., it is again
a combinatorial matrix with the stated formula.

Remarks

\begin{itemize}
\tightlist
\item
  Invertibility requires \(x\neq 0\) and \(x+ny\neq 0\). This matches
  the eigenvalue view: \(A\) has eigenvalue \(x\) on the
  \((n-1)\)-dimensional subspace orthogonal to \(\mathbf{1}\), and
  eigenvalue \(x+ny\) on \(\operatorname{span}\{\mathbf{1}\}\).
\item
  A compact way to remember the result for any \(A=\alpha I+\beta J\) is
\end{itemize}

\[
(\alpha I+\beta J)^{-1}=\frac{1}{\alpha}I-\frac{\beta}{\alpha(\alpha+n\beta)}\,J,
\]

from which the exercise follows by \(\alpha=x,\ \beta=y\).

\subsection{40 {[}M24{]}}\label{m24-1}

View each row \(i\) of \(B\) as the coefficients of a polynomial

\[
p_i(x)\;=\;\sum_{t=1}^n b_{it}\,x^{t}.
\]

The condition \(BA=I\) is exactly

\[
\sum_{t=1}^n b_{it}\,x_j^{\,t} \;=\; \delta_{ij}\qquad(1\le i,j\le n).
\]

So we want a degree--\(\le n\) polynomial \(p_i\) that satisfies
\(p_i(x_j)=\delta_{ij}\) at the \(n\) distinct nodes \(x_j\). By
Lagrange interpolation, the unique choice is

\[
p_i(x)\;=\;\frac{x}{x_i}\,L_i(x),
\qquad
L_i(x)\;=\;\prod_{\substack{1\le k\le n\\ k\ne i}}
\frac{x-x_k}{x_i-x_k}.
\]

Indeed, \(p_i(x_j)=0\) if \(j\ne i\) (since a factor \(x_j-x_k\)
vanishes), and \(p_i(x_i)=(x_i/x_i)\cdot 1=1\) as required.

Therefore the entries of \(B\) are \emph{the coefficients of
\(x/x_i\cdot L_i(x)\)}:

\[
\boxed{\;
b_{it}=\big[ x^{\,t}\big]\!\left(\frac{x}{x_i}\prod_{\substack{k=1\\ k\ne i}}^{n}
\frac{x-x_k}{x_i-x_k}\right)\;}\qquad (t=1,\dots,n).
\]

Closed form using elementary symmetric polynomials

Write

\[
\prod_{\substack{k=1\\k\ne i}}^{n}\,(x-x_k)
=\sum_{m=0}^{n-1}(-1)^{\,n-1-m}\,e_{\,n-1-m}^{(i)}\,x^{m},
\]

where \(e_{r}^{(i)}\) is the \(r\)-th elementary symmetric polynomial in
the \(n-1\) variables \(\{x_1,\dots,x_{i-1},x_{i+1},\dots,x_n\}\), with
\(e_0^{(i)}=1\). Then

\[
p_i(x)=\frac{x}{x_i\prod_{k\ne i}(x_i-x_k)}\;
\sum_{m=0}^{n-1}(-1)^{\,n-1-m}e_{\,n-1-m}^{(i)}x^{m},
\]

and the coefficient of \(x^{t}\) (for \(t=1,\dots,n\)) is

\[
\boxed{\;
b_{it}
=\frac{(-1)^{\,n-t}\,e_{\,n-t}^{(i)}}
{x_i\;\prod\limits_{k\ne i}(x_i-x_k)}\;}.
\]

\subsection{41 {[}M26{]}}\label{m26}

Let \(A=(a_{ij})_{1\le i,j\le n}\) with

\[
a_{ij}=\frac{1}{x_i+y_j}.
\]

Define \(B=(b_{tj})\) by

\[
\boxed{\displaystyle
b_{tj}:=\frac{\Bigl(\prod_{k=1}^{n}(x_k+y_t)\Bigr)\Bigl(\prod_{k\ne t}(x_j+y_k)\Bigr)}
{\Bigl(\prod_{k\ne j}(x_j-x_k)\Bigr)\Bigl(\prod_{k\ne t}(y_t-y_k)\Bigr)}\;.
}
\]

We will prove that \(AB=I\), i.e.

\[
\sum_{t=1}^{n}\frac{1}{x_i+y_t}\,b_{tj}=\delta_{ij}\qquad(1\le i,j\le n).
\]

Proof by Lagrange--interpolation trick

Fix \(i,j\) and consider the function of a dummy variable \(x\)

\[
F_{ij}(x):=\sum_{t=1}^{n}
\frac{\prod_{k\ne t}(x_j+y_k-x)\;\prod_{k\ne i}(x_k+y_t)}
{\prod_{k\ne j}(x_j-x_k)\;\prod_{k\ne t}(y_t-y_k)}.
\]

Observe:

\begin{enumerate}
\def\labelenumi{\arabic{enumi}.}
\item
  \textbf{Degree bound.} The factor \(\prod_{k\ne t}(x_j+y_k-x)\) has
  degree \(n-1\) in \(x\); everything else is independent of \(x\).
  Hence \(F_{ij}(x)\) is a polynomial in \(x\) of degree \(\le n-1\).
\item
  \textbf{Values at \(n\) distinct points.} For each
  \(s\in\{1,\dots,n\}\), plug \(x=x_j+y_s\). Then every summand in
  \(F_{ij}(x)\) vanishes except the one with \(t=s\), because
  \(\prod_{k\ne t}(x_j+y_k-(x_j+y_s))\) is zero unless \(t=s\). For that
  surviving term we get
\end{enumerate}

\[
F_{ij}(x_j+y_s)
=\frac{\prod_{k\ne i}(x_k+y_s)}{\prod_{k\ne j}(x_j-x_k)}
=\frac{\prod_{k\ne j}(x_j-x_k-(x))}{\prod_{k\ne j}(x_j-x_k)}\Bigg|_{x=x_j+y_s-x_j}
\]

(the last rewrite is only to highlight the polynomial nature used
below).

\begin{enumerate}
\def\labelenumi{\arabic{enumi}.}
\setcounter{enumi}{2}
\tightlist
\item
  \textbf{Equality of polynomials.} Two polynomials of degree
  \(\le n-1\) that agree at the \(n\) distinct points \(x=x_j+y_s\) must
  be identical. Hence the value at \(x=0\) is
\end{enumerate}

\[
F_{ij}(0)
=\frac{\prod_{k\ne i}(x_k+y_0)}{\prod_{k\ne j}(x_j-x_k)}\Bigg|_{y_0=0}
=\frac{\prod_{k\ne i}(x_j-x_k)}{\prod_{k\ne j}(x_j-x_k)}
=\delta_{ij},
\]

since the numerator contains the factor \(x_j-x_j=0\) when \(i\ne j\),
and is identical to the denominator when \(i=j\).

Finally, notice that

\[
\frac{1}{x_i+y_t}\,b_{tj}
=\frac{\prod_{k\ne t}(x_j+y_k)\;\prod_{k\ne i}(x_k+y_t)}
{\prod_{k\ne j}(x_j-x_k)\;\prod_{k\ne t}(y_t-y_k)}
=F_{ij}(0)\quad\text{term-by-term},
\]

so summing over \(t\) gives
\(\sum_t a_{it}b_{tj}=F_{ij}(0)=\delta_{ij}\). Therefore \(AB=I\) and
\(B\) is indeed \(A^{-1}\). This is exactly the identity Knuth hints at
and then proves directly via the same interpolation idea.

\subsection{42 {[}M18{]}}\label{m18}

Let \(I\) be the \(n\times n\) identity matrix and \(J\) the
\(n\times n\) all-ones matrix. For scalars \(x,y\) with \(x\neq 0\) and
\(x+ny\neq 0\), find the \textbf{sum of all entries} of
\((xI+yJ)^{-1}\).

Solution 1 --- Eigenvector (one-liner once you see it)

Observe that the all-ones vector
\(\mathbf{1}=(1,1,\dots,1)^{\mathsf T}\) is an eigenvector of \(J\) with
eigenvalue \(n\), hence of \(xI+yJ\) with eigenvalue

\[
(xI+yJ)\mathbf{1}=x\mathbf{1}+y\,J\mathbf{1}=x\mathbf{1}+y(n\mathbf{1})=(x+ny)\mathbf{1}.
\]

Therefore

\[
(xI+yJ)^{-1}\mathbf{1}=\frac{1}{x+ny}\,\mathbf{1}.
\]

The \textbf{sum of the entries in row \(i\) of \((xI+yJ)^{-1}\)} is
exactly the \(i\)-th component of \((xI+yJ)^{-1}\mathbf{1}\),
i.e.~\(1/(x+ny)\). Summing those row-sums over all \(n\) rows gives

\[
\boxed{\;\sum_{i,j} \big((xI+yJ)^{-1}\big)_{ij} \;=\; \frac{n}{x+ny}\; }.
\]

This matches the book's answer. 

Solution 2 --- Sherman--Morrison / explicit inverse (and cross-check)

Use the rank-one inverse identity with \(A=xI\),
\(u=v=\sqrt{y}\,\mathbf{1}\):

\[
(A+uv^{\mathsf T})^{-1}
= A^{-1}-\frac{A^{-1}u\,v^{\mathsf T}A^{-1}}{1+v^{\mathsf T}A^{-1}u}
\;=\; \frac1x I - \frac{y}{x^2}\cdot \frac{\mathbf{1}\mathbf{1}^{\mathsf T}}{1+\frac{y}{x}\, \mathbf{1}^{\mathsf T}\mathbf{1}}
\;=\; \frac1x I - \frac{y}{x(x+ny)}\,\mathbf{1}\mathbf{1}^{\mathsf T}.
\]

Summing all entries:

\[
\sum_{i,j}\Big(\frac1x I\Big)_{ij} = \frac{n}{x},\qquad
\sum_{i,j}\Big(\frac{y}{x(x+ny)}\,\mathbf{1}\mathbf{1}^{\mathsf T}\Big)_{ij} = \frac{y}{x(x+ny)}\,n^2,
\]

so

\[
\sum_{i,j} (xI+yJ)^{-1}_{ij} = \frac{n}{x} - \frac{y n^2}{x(x+ny)} = \frac{n}{x+ny}.
\]

(As a related identity appearing in the exercise cluster, Knuth notes
\((xI+yJ)\big((x+ny)I-yJ\big)=x(x+ny)I\), from which the inverse formula
is immediate. )

Checks and conditions

\begin{itemize}
\tightlist
\item
  \textbf{Invertibility}: requires \(x\neq 0\) and \(x+ny\neq 0\) (these
  are exactly the eigenvalues).
\item
  \textbf{Sanity checks}: • If \(y=0\), matrix is \(xI\Rightarrow\)
  inverse is \((1/x)I\Rightarrow\) sum \(=n/x\), which equals
  \(n/(x+0\cdot n)\). • If \(n=1\), matrix is \([x+y]\Rightarrow\)
  inverse has sum \(=1/(x+y)=n/(x+ny)\).
\end{itemize}

Thus the sum of all entries of \((xI+yJ)^{-1}\) is indeed
\(\boxed{n/(x+ny)}\).

\subsection{43 {[}M24{]}}\label{m24-2}

Let \(V\) be the \(n\times n\) Vandermonde matrix

\[
 V=\big[x_j^{\,i}\big]_{1\le i,j\le n}
 \quad\text{(so its rows are }(x_1,\dots,x_n),(x_1^2,\dots,x_n^2),\dots,(x_1^n,\dots,x_n^n)\text{).}
 \]

What is the sum of all \(n^2\) entries of \(V^{-1}\)?

Step 1: A usable formula for \(V^{-1}\)

A convenient description (stated in Ex. 40) is

\[
(V^{-1})_{ij}
=\frac{[x^{\,j-1}]\;\displaystyle\prod_{k\ne i}(x_k-x)}{x_i\;\displaystyle\prod_{k\ne i}(x_k-x_i)}\,,
\quad 1\le i,j\le n,
\]

i.e.~the numerator is the coefficient of \(x^{j-1}\) in
\(\prod_{k\ne i}(x_k-x)\). (This is exactly the ``coefficient of
\(x^{j-1}\)'' description given under Ex. 40.)

Step 2: Sum across each row

For fixed \(i\), the row-sum is

\[
R_i=\sum_{j=1}^n (V^{-1})_{ij}
=\frac{\sum_{j=1}^n [x^{\,j-1}]\prod_{k\ne i}(x_k-x)}{x_i\prod_{k\ne i}(x_k-x_i)}
=\frac{\left.\prod_{k\ne i}(x_k-x)\right|_{x=1}}{x_i\prod_{k\ne i}(x_k-x_i)}
=\frac{\prod_{k\ne i}(x_k-1)}{x_i\prod_{k\ne i}(x_k-x_i)}.
\]

(The sum of coefficients of a polynomial equals its value at \(x=1\).)

Step 3: Sum all entries and simplify (Lagrange view)

Hence the total sum is

\[
S=\sum_{i=1}^n R_i
=\sum_{i=1}^n \frac{1}{x_i}\,\frac{\prod_{k\ne i}(x_k-1)}{\prod_{k\ne i}(x_k-x_i)}
=\sum_{i=1}^n \frac{1}{x_i}\,L_i(1),
\]

where \(L_i(t)=\prod_{k\ne i}\dfrac{t-x_k}{x_i-x_k}\) are the Lagrange
basis polynomials.

Consider the degree-\(\le n-1\) polynomial

\[
g(t)=\sum_{i=1}^n \frac{1}{x_i}\,L_i(t).
\]

By construction, \(g(x_i)=\frac1{x_i}\) for all \(i\). The polynomial

\[
h(t)=\frac{1-\prod_{k=1}^n\left(1-\frac{t}{x_k}\right)}{t}
\]

also has degree \(\le n-1\) and satisfies \(h(x_i)=\frac1{x_i}\) (the
product vanishes at \(t=x_i\)). Therefore \(g\equiv h\) by uniqueness of
Lagrange interpolation (cf.~the Lagrange-sum identities highlighted in
Ex. 33). Evaluating at \(t=1\) gives

\[
S=g(1)=h(1)
=1-\prod_{k=1}^n\Bigl(1-\frac{1}{x_k}\Bigr)
=1-\frac{\prod_{k=1}^n(x_k-1)}{\prod_{k=1}^n x_k}.
\]

(Ex. 33 sets up exactly these Lagrange-basis manipulations. )

Final result (two equivalent forms)

\[
\boxed{\,\displaystyle
\sum_{i=1}^n\sum_{j=1}^n (V^{-1})_{ij}
=1-\prod_{k=1}^n\!\left(1-\frac{1}{x_k}\right)
= \sum_{r=1}^{n}(-1)^{r-1}\!\!\sum_{1\le i_1<\cdots<i_r\le n}\frac{1}{x_{i_1}\cdots x_{i_r}}\,.
}
\]

The first form is the clean product; the second is the alternating sum
of the elementary symmetric functions in \(1/x_k\).

Sanity check (small \(n\))

For \(n=2\), \(V=\begin{bmatrix}x_1&x_2\\ x_1^2&x_2^2\end{bmatrix}\).
One computes

\[
\sum (V^{-1})_{ij}=1-\Bigl(1-\tfrac1{x_1}\Bigr)\Bigl(1-\tfrac1{x_2}\Bigr)=\frac{1}{x_1}+\frac{1}{x_2}-\frac{1}{x_1x_2},
\]

matching the boxed formula.

\subsection{44 {[}M26{]}}\label{m26-1}

Let \(A\) be the \(n\times n\) \textbf{Cauchy matrix}

\[
A_{ij}=\frac{1}{x_i+y_j}\qquad (1\le i,j\le n),
\]

with all \(x_i\) distinct, all \(y_j\) distinct, and \(x_i+y_j\neq0\).
Find the sum of all \(n^2\) entries of \(A^{-1}\). 

Tooling you're allowed to use

Exercise 41 on the same page gives the entries of the inverse of a
Cauchy matrix:

\[
(A^{-1})_{ij}
=\frac{\displaystyle\prod_{k=1}^n(x_j+y_k)\;\prod_{k=1}^n(x_k+y_i)}
{(x_j+y_i)\;\displaystyle\prod_{k\ne j}(x_j-x_k)\;\prod_{k\ne i}(y_i-y_k)}.
\tag{★}
\]

We'll use this formula as a starting point.

Step 1: Separate variables

From (★) define

\[
\mu_j=\frac{\prod_{k=1}^n(x_j+y_k)}{\prod_{k\ne j}(x_j-x_k)},
\qquad
\nu_i=\frac{\prod_{k=1}^n(x_k+y_i)}{\prod_{k\ne i}(y_i-y_k)}.
\]

Then

\[
(A^{-1})_{ij}=\frac{\mu_j\,\nu_i}{x_j+y_i}.
\]

Hence the sum of all entries is

\[
S=\sum_{i=1}^n\sum_{j=1}^n (A^{-1})_{ij}
=\sum_{j=1}^n \mu_j\,\Bigl(\sum_{i=1}^n \frac{\nu_i}{x_j+y_i}\Bigr).
\tag{1}
\]

Step 2: A barycentric/Lagrange identity

Set

\[
\Phi_x(t)=\prod_{k=1}^n (t-x_k),\qquad
\Phi_y(t)=\prod_{k=1}^n (t+y_k).
\]

Consider the rational function

\[
F(t)=\sum_{i=1}^n \frac{\nu_i}{t+y_i}.
\]

Multiplying by \(\Phi_y(t)\) gives a degree-\(\le n-1\) polynomial

\[
\Phi_y(t)F(t)=\sum_{i=1}^n \nu_i\!\!\prod_{k\ne i}(t+y_k),
\]

which is the Lagrange interpolation polynomial that equals
\(\Phi_x(-y_i)\) at the points \(t=-y_i\). Because \(\Phi_x\) and
\(\Phi_y\) are both monic of degree \(n\), the remainder of
\(\Phi_x(t)\) on division by \(\Phi_y(t)\) is \(\Phi_x(t)-\Phi_y(t)\).
Thus

\[
\Phi_y(t)F(t)=\Phi_x(t)-\Phi_y(t)\quad\Rightarrow\quad
F(t)=1-\frac{\Phi_x(t)}{\Phi_y(t)}.
\]

Evaluating at \(t=x_j\) gives \(F(x_j)=1\), since \(\Phi_x(x_j)=0\).
Plugging into (1),

\[
S=\sum_{j=1}^n \mu_j. \tag{2}
\]

Step 3: Sum of the \(\mu_j\)

By Lagrange's formula,

\[
\mu_j=\frac{\Phi_y(x_j)}{\Phi_x'(x_j)}.
\]

A standard coefficient trick now yields

\[
\sum_{j=1}^n \mu_j
=\text{(coefficient of }t^{\,n-1}\text{ in }\Phi_y(t)-\Phi_x(t)).
\]

But

\[
\Phi_y(t)=t^n+\Bigl(\sum_{k=1}^n y_k\Bigr)t^{n-1}+\cdots,
\qquad
\Phi_x(t)=t^n-\Bigl(\sum_{k=1}^n x_k\Bigr)t^{n-1}+\cdots,
\]

so the \(t^{n-1}\)-coefficient of \(\Phi_y-\Phi_x\) is

\[
\sum_{k=1}^n y_k+\sum_{k=1}^n x_k.
\]

Therefore

\[
S=\sum_{i,j}(A^{-1})_{ij}=\sum_{k=1}^n x_k+\sum_{k=1}^n y_k.
\]

Result (and a check)

\[
\boxed{\displaystyle \sum_{i=1}^n\sum_{j=1}^n (A^{-1})_{ij}
=\Bigl(\sum_{i=1}^n x_i\Bigr)+\Bigl(\sum_{j=1}^n y_j\Bigr).}
\]

As a sanity check, Exercise 45 asks for the Hilbert case
\(A_{ij}=1/(i+j-1)\), which corresponds to \(x_i=i-1,\;y_j=j\). The
formula gives

\[
\sum (A^{-1})_{ij}=\sum_{i=1}^n(i-1)+\sum_{j=1}^n j
=\frac{n(n-1)}{2}+\frac{n(n+1)}{2}=n^2,
\]

matching the stated result.

\subsection{45 {[}M25{]}}\label{m25-6}

\subsection{46 {[}M30{]}}\label{m30-6}

For an \(m\times n\) matrix \(A\) and an \(n\times m\) matrix \(B\),

\[
\det(AB)
=\sum_{1\le j_1<\cdots<j_m\le n}
\det\!\big(A_{j_1\ldots j_m}\big)\,
\det\!\big(B_{j_1\ldots j_m}\big),
\]

where \(A_{j_1\ldots j_m}\) is the \(m\times m\) submatrix of \(A\)
formed by columns \(j_1,\ldots,j_m\) and \(B_{j_1\ldots j_m}\) is the
\(m\times m\) submatrix of \(B\) formed by rows \(j_1,\ldots,j_m\). (See
the statement in the book for the same notation and suggested special
cases.) 

Proof (by multilinearity and alternation of the determinant)

Let \(C=AB\) (so \(C\) is \(m\times m\)). Denote the \(k\)-th column of
\(A\) by \(A^{(k)}\), and the \(r\)-th column of \(C\) by \(C^{(r)}\).
Because matrix multiplication forms columns as linear combinations of
the columns of \(A\),

\[
C^{(r)}=\sum_{k=1}^{n} b_{k r}\,A^{(k)}\qquad(1\le r\le m),
\]

where \(b_{kr}\) are the entries of \(B\).

Using that the determinant is multilinear in the columns,

\[
\det(C)
=\det\!\big(C^{(1)},\ldots,C^{(m)}\big)
=\sum_{k_1,\ldots,k_m=1}^{n}
\Big(\prod_{r=1}^{m} b_{k_r r}\Big)\,
\det\!\big(A^{(k_1)},\ldots,A^{(k_m)}\big).
\tag{*}
\]

\begin{enumerate}
\def\labelenumi{\arabic{enumi}.}
\item
  \textbf{Only distinct index \(m\)-tuples contribute.} If any two of
  \(k_1,\ldots,k_m\) coincide, two columns in
  \(\det\big(A^{(k_1)},\ldots,A^{(k_m)}\big)\) are equal, so that term
  is \(0\) by alternation. Hence we may restrict to \(m\)-tuples with
  \emph{distinct} indices.
\item
  \textbf{Group by the unordered set \(J=\{j_1<\cdots<j_m\}\).} Fix such
  a set \(J\). The surviving terms in \((*)\) that use exactly these
  indices are parameterized by permutations \(\pi\in S_m\):
\end{enumerate}

\[
\det\!\big(A^{(j_{\pi(1)})},\ldots,A^{(j_{\pi(m)})}\big)\,
\prod_{r=1}^{m} b_{j_{\pi(r)}\, r}.
\]

By alternation,

\[
\det\!\big(A^{(j_{\pi(1)})},\ldots,A^{(j_{\pi(m)})}\big)
=\operatorname{sgn}(\pi)\,\det\!\big(A^{(j_{1})},\ldots,A^{(j_{m})}\big).
\]

\begin{enumerate}
\def\labelenumi{\arabic{enumi}.}
\setcounter{enumi}{2}
\tightlist
\item
  \textbf{Recognize a determinant of a submatrix of \(B\).} Let \(M(J)\)
  be the \(m\times m\) matrix with entries \(M(J)_{s,r}=b_{j_s\, r}\)
  (i.e., the submatrix of \(B\) consisting of rows \(j_1,\ldots,j_m\);
  this is exactly \(B_{j_1\ldots j_m}\)). Then
\end{enumerate}

\[
\det\big(M(J)\big)
=\sum_{\pi\in S_m}\operatorname{sgn}(\pi)\prod_{r=1}^{m} b_{j_{\pi(r)}\, r}.
\]

Therefore the total contribution from the fixed set \(J\) equals

\[
\det\!\big(A_{j_1\ldots j_m}\big)\cdot\det\!\big(B_{j_1\ldots j_m}\big).
\]

\begin{enumerate}
\def\labelenumi{\arabic{enumi}.}
\setcounter{enumi}{3}
\tightlist
\item
  \textbf{Sum over all \(J\).} Summing these contributions over all
  \(1\le j_1<\cdots<j_m\le n\) yields precisely the right-hand side of
  the identity, proving the formula.
\end{enumerate}

\(\square\)

Sanity checks / special cases

\begin{itemize}
\tightlist
\item
  \textbf{\(m=n\).} There is only the single set \(J=\{1,\ldots,n\}\),
  so the identity reduces to \(\det(AB)=\det(A)\det(B)\).
\item
  \textbf{\(m=1\).} Then \(\det(AB)\) equals the dot product of the
  \(1\times n\) row of \(A\) with the \(n\times 1\) column of \(B\); the
  right side is \(\sum_{j=1}^{n} a_{1j}b_{j1}\).
\item
  \textbf{\(B=A^{\mathsf T}\).} We get
  \(\det(AA^{\mathsf T})=\displaystyle\sum_{J}\big(\det A_{J}\big)^2\),
  the classic Gram determinant identity.
\item
  \textbf{\(m>n\).} There are no \(m\)-subsets \(J\) of
  \(\{1,\dots,n\}\), so the sum is empty (0); indeed
  \(\operatorname{rank}(AB)\le n<m\Rightarrow \det(AB)=0\).
\item
  \textbf{\(m=2\).} The formula becomes
  \(\det(AB)=\sum_{1\le i<j\le n}\det(A^{(i)},A^{(j)})\,\det\big(\text{rows }i,j\text{ of }B\big)\),
  helpful when expanding by \(2\times2\) minors.
\end{itemize}

\subsection{47 {[}M27{]}}\label{m27-1}

For an integer \(k\), prove

\[
\binom{r}{k}\binom{r-\tfrac12}{k}
= \frac{1}{4^k}\binom{2r}{k}\binom{2r-k}{k}
= \frac{1}{4^k}\binom{2r}{2k}\binom{2k}{k},
\]

and give a simpler closed form when \(r=-\tfrac12\). 

We use the product (falling--factorial) definition of generalized
binomial coefficients for real \(r\) and integer \(k\ge0\):

\[
\binom{r}{k}=\frac{r(r-1)\cdots(r-k+1)}{k!}.
\]

\begin{enumerate}
\def\labelenumi{\arabic{enumi}.}
\tightlist
\item
  Multiply the two binomials and pair terms:
\end{enumerate}

\[
\binom{r}{k}\binom{r-\tfrac12}{k}
=\frac{\displaystyle\prod_{j=0}^{k-1}(r-j)(r-\tfrac12-j)}{(k!)^2}.
\]

For each \(j\),

\[
(r-j)\Bigl(r-\tfrac12-j\Bigr)
=\frac{\bigl(2r-2j\bigr)\bigl(2r-(2j+1)\bigr)}{4}.
\]

Hence

\[
\binom{r}{k}\binom{r-\tfrac12}{k}
=\frac{1}{4^k}\cdot
\frac{\displaystyle\prod_{j=0}^{k-1}(2r-2j)(2r-(2j+1))}{(k!)^2}.
\]

But

\[
\prod_{j=0}^{k-1}(2r-2j)(2r-(2j+1))
=\prod_{m=0}^{2k-1}(2r-m)=(2r)^{\underline{2k}},
\]

so

\[
\binom{r}{k}\binom{r-\tfrac12}{k}
=\frac{1}{4^k}\cdot\frac{(2r)^{\underline{2k}}}{(k!)^2}
=\frac{1}{4^k}\binom{2r}{2k}\cdot\frac{(2k)!}{(k!)^2}
=\frac{1}{4^k}\binom{2r}{2k}\binom{2k}{k}.
\]

This proves the right-hand equality.

\begin{enumerate}
\def\labelenumi{\arabic{enumi}.}
\setcounter{enumi}{1}
\tightlist
\item
  Convert \(\binom{2r}{2k}\binom{2k}{k}\) to
  \(\binom{2r}{k}\binom{2r-k}{k}\) using the standard identity
\end{enumerate}

\[
\binom{r}{m}\binom{m}{k}=\binom{r}{k}\binom{r-k}{m-k}
\qquad(\text{integers }m,k),
\]

with \(r\mapsto 2r\) and \(m\mapsto 2k\). Thus

\[
\binom{2r}{2k}\binom{2k}{k}
=\binom{2r}{k}\binom{2r-k}{k},
\]

giving the left-hand equality as well. 

\begin{enumerate}
\def\labelenumi{\arabic{enumi}.}
\setcounter{enumi}{2}
\tightlist
\item
  Special case \(r=-\tfrac12\). Plugging \(r=-\tfrac12\) into the proven
  formula,
\end{enumerate}

\[
\binom{-\tfrac12}{k}\binom{-1}{k}
=\frac{1}{4^k}\binom{-1}{2k}\binom{2k}{k}.
\]

For integers \(n\ge0\), \(\binom{-1}{n}=(-1)^n\). Hence
\(\binom{-1}{2k}=1\) and \(\binom{-1}{k}=(-1)^k\). Therefore

\[
\binom{-\tfrac12}{k}\,(-1)^k=\frac{1}{4^k}\binom{2k}{k}
\quad\Longrightarrow\quad
\boxed{\displaystyle \binom{-\tfrac12}{k}=(-1)^k\,\frac{1}{4^k}\binom{2k}{k}}.
\]

This is the desired ``simpler formula'' in the \(r=-\tfrac12\) case. (It
is also the well-known central-binomial closed form used repeatedly
later.)

\bookmarksetup{startatroot}

\chapter{1.2.4 Integer functions and elementary number
theory}\label{integer-functions-and-elementary-number-theory}

\subsection{1 {[}00{]}}\label{section-55}

First, recall the definitions from §1.2.4: ⌊x⌋ is the greatest integer ≤
x (the ``floor''), and ⌈x⌉ is the least integer ≥ x (the ``ceiling'').

\begin{enumerate}
\def\labelenumi{\arabic{enumi}.}
\item
  ⌊1.1⌋ 1.1 lies between 1 and 2, and the greatest integer ≤ 1.1 is 1 ⇒
  \textbf{1}.
\item
  ⌊−1.1⌋ For negative numbers, ``greatest integer ≤ x'' moves to the
  left: integers ≤ −1.1 are \ldots, −3, −2, −1, and the greatest among
  those ≤ −1.1 is −2 ⇒ \textbf{−2}.
\item
  ⌈−1.1⌉ Least integer ≥ −1.1 is −1 ⇒ \textbf{−1}.
\item
  ⌊0.99999⌋ 0 ≤ 0.99999 \textless{} 1, so the greatest integer ≤ it is
  \textbf{0}.
\item
  ⌈lg 35⌉ (lg = log base 2) 2⁵ = 32 \textless{} 35 ≤ 64 = 2⁶, hence 5
  \textless{} lg 35 ≤ 6, so ⌈lg 35⌉ = \textbf{6}.
\end{enumerate}

\subsection{2 {[}01{]}}\label{section-56}

For any real x,

\begin{itemize}
\tightlist
\item
  ⌊x⌋ is the greatest integer ≤ x (the floor), and
\item
  ⌈x⌉ is the least integer ≥ x (the ceiling). 
\end{itemize}

Let n = ⌊x⌋. By definition of floor, n is an integer with
\(n \le x < n+1.\)

Because n itself is already an integer, taking the ceiling of n leaves
it unchanged: \(\lceil n \rceil = n.\)

Therefore, \$
\boxed{\;\lceil\lfloor x\rfloor\rceil=\lfloor x\rfloor\;}.\$

Why this is always true

\begin{itemize}
\tightlist
\item
  Floor always returns an integer.
\item
  Ceiling of any integer k equals k (since k is already the least
  integer ≥ k).
\item
  Substituting k = ⌊x⌋ gives the identity above.
\end{itemize}

Quick checks

\begin{itemize}
\tightlist
\item
  If x = 3.7, ⌊x⌋ = 3 and ⌈⌊x⌋⌉ = ⌈3⌉ = 3.
\item
  If x = −2.3, ⌊x⌋ = −3 and ⌈⌊x⌋⌉ = ⌈−3⌉ = −3.
\item
  If x is already an integer m, ⌊x⌋ = m and ⌈⌊x⌋⌉ = ⌈m⌉ = m.
\end{itemize}

Thus the expression simplifies to the floor of x.

\subsection{3 {[}M10{]}}\label{m10}

First recall the definitions: ⌊x⌋ is the greatest integer ≤ x (the
``floor''), and ⌈x⌉ is the least integer ≥ x (the ``ceiling'').

Two basic characterizations (we'll use them repeatedly)

\textbf{Lemma F (floor characterization).} For any real x and integer n,
⌊x⌋ = n ⇔ n ≤ x \textless{} n + 1. \emph{Proof.} From the definition,
⌊x⌋ ≤ x. If we had x ≥ ⌊x⌋ + 1, then ⌊x⌋ + 1 would be an integer ≤ x,
contradicting maximality of ⌊x⌋. Hence x \textless{} ⌊x⌋ + 1.
Conversely, if n ≤ x \textless{} n + 1, no integer ≥ n + 1 is ≤ x; thus
the greatest integer ≤ x is exactly n.~∎

\textbf{Lemma C (ceiling characterization).} For any real x and integer
n, ⌈x⌉ = n ⇔ n − 1 \textless{} x ≤ n. \emph{Proof.} By definition, x ≤
⌈x⌉. If we had x ≤ ⌈x⌉ − 1, then ⌈x⌉ − 1 would be an integer ≥ x that is
\textless{} ⌈x⌉, contradicting minimality of ⌈x⌉. Hence ⌈x⌉ − 1
\textless{} x. The converse is analogous to Lemma F. ∎

(Equivalently, Lemma C is the same as Lemma F applied to −x, since ⌈x⌉ =
−⌊−x⌋. )

Part (a) ⌊x⌋ \textless{} n ⇔ x \textless{} n

\textbf{(⇒)} If ⌊x⌋ \textless{} n, then by Lemma F we have x \textless{}
⌊x⌋ + 1 ≤ n, hence x \textless{} n. \textbf{(⇐)} If x \textless{} n,
then every integer ≤ x is \textless{} n, so the greatest such integer
⌊x⌋ is \textless{} n.

Part (b) n ≤ ⌊x⌋ ⇔ n ≤ x

\textbf{(⇒)} If n ≤ ⌊x⌋, then n ≤ x because ⌊x⌋ ≤ x. \textbf{(⇐)} If n ≤
x, then among the integers ≤ x is n itself, so the greatest integer ≤ x
satisfies ⌊x⌋ ≥ n.

Part (c) ⌈x⌉ ≤ n ⇔ x ≤ n

\textbf{(⇒)} If ⌈x⌉ ≤ n, then x ≤ ⌈x⌉ ≤ n, hence x ≤ n. \textbf{(⇐)} If
x ≤ n, then n is an integer ≥ x, so the least integer ≥ x satisfies ⌈x⌉
≤ n.

Part (d) n \textless{} ⌈x⌉ ⇔ n \textless{} x

\textbf{(⇒)} If n \textless{} ⌈x⌉, then (by Lemma C) n \textless{} x is
forced because x \textgreater{} ⌈x⌉ − 1 ≥ n. \textbf{(⇐)} If n
\textless{} x, then any integer ≥ x is strictly greater than n; in
particular the least such integer ⌈x⌉ satisfies n \textless{} ⌈x⌉.
(Equivalently, (d) is just the contrapositive of (c).)

Part (e) ⌊x⌋ = n ⇔ (x − 1 \textless{} n ≤ x) ⇔ (n ≤ x \textless{} n + 1)

The right-hand equivalence \emph{is exactly} Lemma F. For the middle
form, note that n ≤ x \textless{} n + 1 ⇔ (subtract 1 from the right
inequality) ⇔ x − 1 \textless{} n ≤ x. Thus both formulations are
equivalent and both characterize ⌊x⌋ = n.

Part (f) ⌈x⌉ = n ⇔ (x ≤ n \textless{} x + 1) ⇔ (n − 1 \textless{} x ≤ n)

This is Lemma C, plus the observation that n − 1 \textless{} x ≤ n ⇔
(add 1 to the left inequality) ⇔ x ≤ n \textless{} x + 1.

\subsection{4 {[}M10{]}}\label{m10-1}

\textbf{Claim.} For every real \(x\),

\[
\lfloor -x \rfloor \;=\; -\,\lceil x\rceil .
\]

Let \(n=\lceil x\rceil\). By the defining inequalities for the ceiling
function,

\[
n-1 < x \le n .
\]

(These standard floor/ceiling characterizations are given just above and
within §1.2.4. )

Multiply the inequalities by \(-1\) (which reverses the direction):

\[
-n \le -x < 1-n .
\]

Set \(m=-n\). Since \(m\) is an integer, the last display reads

\[
m \le -x < m+1 .
\]

By the defining property of the floor function---\(\lfloor y\rfloor=m\)
iff \(m\le y<m+1\)---we conclude

\[
\lfloor -x\rfloor = m = -n = -\,\lceil x\rceil .
\]

This completes the proof.

\subsection{5 {[}16{]}}\label{section-57}

Give a single simple formula for rounding a positive real number \(x\)
to the nearest integer, with the rule:

\begin{itemize}
\tightlist
\item
  use \(\lfloor x\rfloor\) when \(x \bmod 1 < \tfrac12\);
\item
  use \(\lceil x\rceil\) when \(x \bmod 1 \ge \tfrac12\). Then discuss
  what rounding your formula gives when \(x\) is negative.
  (Floor/ceiling and ``mod'' are defined earlier in §1.2.4.)
\end{itemize}

A single formula

\[
\boxed{\,R(x)=\big\lfloor x+\tfrac12\big\rfloor\,}
\]

Why this works

Write \(x=\lfloor x\rfloor + f\) where the fractional part
\(f=x\bmod 1\) satisfies \(0\le f<1\). (This is exactly Eq. (1):
\(x=\lfloor x\rfloor+(x\bmod 1)\).) 

Then

\[
\big\lfloor x+\tfrac12\big\rfloor
=\big\lfloor \lfloor x\rfloor + f+\tfrac12\big\rfloor
=\lfloor x\rfloor + \big\lfloor f+\tfrac12\big\rfloor.
\]

\begin{itemize}
\tightlist
\item
  If \(f<\tfrac12\), then \(0\le f+\tfrac12<1\), so
  \(\lfloor f+\tfrac12\rfloor=0\) and \(R(x)=\lfloor x\rfloor\).
\item
  If \(f\ge\tfrac12\), then \(1\le f+\tfrac12<\tfrac32\), so
  \(\lfloor f+\tfrac12\rfloor=1\) and
  \(R(x)=\lfloor x\rfloor+1=\lceil x\rceil\) (since \(x\) isn't an
  integer in this case).
\end{itemize}

Thus the single formula \(\lfloor x+\tfrac12\rfloor\) implements exactly
the stated rounding rule for positive \(x\).

What happens for negative \(x\)?

Let \(x<0\) and write \(x=-y\) with \(y\ge 0\). Then

\[
R(x)=\big\lfloor -y+\tfrac12\big\rfloor
=-\,\big\lceil y-\tfrac12\big\rceil.
\]

So ties at halves are pushed \textbf{toward \(+\infty\)} (i.e.,
``half-up''): for example,

\[
R(-1.3)=-1,\quad R(-1.5)=-1,\quad R(-1.7)=-2.
\]

In words: the same formula \(\lfloor x+\tfrac12\rfloor\) rounds to the
nearest integer with \textbf{half cases rounded upward} (toward
\(+\infty\)). For positive numbers this is the usual ``round half up'';
for negative numbers it rounds half-cases \textbf{toward zero}.

(Optional) Symmetric alternative

If you wanted a ``nearest with halves away from zero'' rule, a symmetric
one-liner is

\[
\operatorname{round}_{\text{away-from-0}}(x)=
\begin{cases}
\lfloor x+\tfrac12\rfloor,& x\ge 0,\\[2pt]
\lceil x-\tfrac12\rceil,& x<0,
\end{cases}
\]

but the exercise asks for the simplest single formula, which is
\(\boxed{\lfloor x+\tfrac12\rfloor}\).

\subsection{6 {[}20{]}}\label{section-58}

I'll prove (a) and (b) are always true, and (c) is not.

\begin{enumerate}
\def\labelenumi{(\alph{enumi})}
\tightlist
\item
  \(\;\lfloor \sqrt{\lfloor x\rfloor}\rfloor=\lfloor \sqrt{x}\rfloor\)
  --- \textbf{True}
\end{enumerate}

Let \(n=\lfloor x\rfloor\). Then \(n\le x<n+1\). Taking square roots
gives

\[
\sqrt{n}\;\le\;\sqrt{x}\;<\;\sqrt{n+1}.
\]

There is \textbf{no integer} strictly between \(\sqrt{n}\) and
\(\sqrt{n+1}\): if an integer \(m\) satisfied \(\sqrt{n}<m<\sqrt{n+1}\),
then squaring would give \(n<m^2<n+1\), impossible since \(m^2\) is an
integer. Therefore \(\lfloor \sqrt{x}\rfloor\) must equal
\(\lfloor \sqrt{n}\rfloor\),
i.e.~\(\lfloor \sqrt{\lfloor x\rfloor}\rfloor\).

\begin{enumerate}
\def\labelenumi{(\alph{enumi})}
\setcounter{enumi}{1}
\tightlist
\item
  \(\;\lceil \sqrt{\lceil x\rceil}\rceil=\lceil \sqrt{x}\rceil\) ---
  \textbf{True}
\end{enumerate}

Let \(c=\lceil x\rceil\). Then \(c-1<x\le c\), hence

\[
\sqrt{c-1}\;<\;\sqrt{x}\;\le\;\sqrt{c}.
\]

For any integer \(m\), \(\lceil \sqrt{t}\rceil=m\) exactly on the
interval \(((m-1)^2,\,m^2]\). The short interval \((c-1,c]\) lies wholly
inside a \textbf{single} such block (its length is \(1\), while
consecutive squares are at least \(1\) apart and, except for \(0\) and
\(1\), more than \(1\) apart). Thus \(\lceil \sqrt{x}\rceil\) is
constant on \((c-1,c]\) and equals \(\lceil \sqrt{c}\rceil\),
i.e.~\(\lceil \sqrt{\lceil x\rceil}\rceil\).

A more explicit framing: let \(m=\lceil \sqrt{c}\rceil\), so
\((m-1)^2<c\le m^2\). Then for all \(x\in(c-1,c]\) we have
\((m-1)^2<x\le m^2\), hence
\(\lceil \sqrt{x}\rceil=m=\lceil \sqrt{c}\rceil\).

\begin{enumerate}
\def\labelenumi{(\alph{enumi})}
\setcounter{enumi}{2}
\tightlist
\item
  \(\;\lceil \sqrt{\lfloor x\rfloor}\rceil=\lceil \sqrt{x}\rceil\) ---
  \textbf{False}
\end{enumerate}

Take any perfect square \(k^2\) and choose \(x=k^2+\varepsilon\) with
\(0<\varepsilon<1\). Then \(\lfloor x\rfloor=k^2\), so

\[
\lceil \sqrt{\lfloor x\rfloor}\rceil=\lceil \sqrt{k^2}\rceil=k,
\]

but \(\sqrt{x}>\!k\) and \(\lceil \sqrt{x}\rceil=k+1\). Example:
\(x=4.1\) gives
\(\lceil \sqrt{\lfloor 4.1\rfloor}\rceil=\lceil \sqrt{4}\rceil=2\) while
\(\lceil \sqrt{4.1}\rceil=3\).

\textbf{Answer:} (a) True; (b) True; (c) False.

\subsection{7 {[}M15{]}}\label{m15}

Let the fractional part of a real number \(x\) be
\(\{x\}=x-\lfloor x\rfloor\). From the text we have

\[
x=\lfloor x\rfloor+(x\!\bmod 1)=\lfloor x\rfloor+\{x\},\quad\text{with }0\le \{x\}<1,
\]

so write \(x=a+\alpha\) and \(y=b+\beta\) where
\(a=\lfloor x\rfloor,\ b=\lfloor y\rfloor,\ \alpha=\{x\},\ \beta=\{y\}\in[0,1)\). 

\begin{enumerate}
\def\labelenumi{\arabic{enumi})}
\tightlist
\item
  Prove \(\lfloor x\rfloor+\lfloor y\rfloor\le \lfloor x+y\rfloor\) and
  find when equality holds
\end{enumerate}

Compute

\[
x+y=(a+b)+(\alpha+\beta).
\]

Because \(0\le \alpha,\beta<1\), we have \(0\le \alpha+\beta<2\). Thus

\[
\lfloor x+y\rfloor=
\begin{cases}
a+b, & \text{if }\alpha+\beta<1,\\
a+b+1, & \text{if }\alpha+\beta\ge 1.
\end{cases}
\]

Therefore
\(\lfloor x\rfloor+\lfloor y\rfloor=a+b\le \lfloor x+y\rfloor\), with
\textbf{equality iff \(\alpha+\beta<1\)},
i.e.~\textbf{\(\{x\}+\{y\}<1\)} (equivalently,
\((x\bmod 1)+(y\bmod 1)<1\)). 

\begin{enumerate}
\def\labelenumi{\arabic{enumi})}
\setcounter{enumi}{1}
\tightlist
\item
  ``Similar formula'' for ceilings
\end{enumerate}

Using \(\lceil t\rceil=-\lfloor -t\rfloor\), we get

\[
\lceil x\rceil+\lceil y\rceil\ \ge\ \lceil x+y\rceil.
\]

Equality cases come from the floor statement applied to \(-x,-y\). Note
that \(\{-x\}=0\) if \(\{x\}=0\), and \(\{-x\}=1-\{x\}\) if \(\{x\}>0\).
Thus

\[
\{-x\}+\{-y\}<1\quad\Longleftrightarrow\quad
\begin{cases}
\text{either }\{x\}=0\text{ or }\{y\}=0,\\[2pt]
\text{or }\{x\}+\{y\}>1.
\end{cases}
\]

Hence

\[
\boxed{\ \lceil x\rceil+\lceil y\rceil\ge \lceil x+y\rceil,\ \text{with equality iff }\{x\}=0\ \text{or}\ \{y\}=0\ \text{or}\ \{x\}+\{y\}>1\ }.
\]

Equivalently, the inequality is \textbf{strict} exactly when
\(0<\{x\}<1,\ 0<\{y\}<1,\) and \(\{x\}+\{y\}\le 1\). 

Quick sanity checks

\begin{itemize}
\tightlist
\item
  \(x=2.3,\ y=4.6:\ \{x\}+\{y\}=0.9<1\Rightarrow \lfloor x+y\rfloor=6=\lfloor x\rfloor+\lfloor y\rfloor\).
\item
  \(x=2.7,\ y=4.6:\ \{x\}+\{y\}=1.3>1\Rightarrow \lceil x+y\rceil=8=\lceil x\rceil+\lceil y\rceil\).
\end{itemize}

\subsection{8 {[}00{]}}\label{section-59}

Definition we'll use. For real \(x\) and \(y\neq 0\),

\[
x\bmod y \;=\; x - y\left\lfloor \frac{x}{y}\right\rfloor ,
\]

and by convention \(x\bmod 0 = x\). When \(y>0\) this guarantees
\(0\le x\bmod y<y\). 

\begin{enumerate}
\def\labelenumi{(\alph{enumi})}
\tightlist
\item
  \(100 \bmod 3\)
\end{enumerate}

\(\left\lfloor 100/3\right\rfloor = 33\).
\(100 - 3\cdot 33 = 100 - 99 = 1\). \textbf{Answer:} \(1\). (Indeed
\(0\le 1<3\).) 

\begin{enumerate}
\def\labelenumi{(\alph{enumi})}
\setcounter{enumi}{1}
\tightlist
\item
  \(100 \bmod 7\)
\end{enumerate}

\(\left\lfloor 100/7\right\rfloor = 14\).
\(100 - 7\cdot 14 = 100 - 98 = 2\). \textbf{Answer:} \(2\). (And
\(0\le 2<7\).) 

\begin{enumerate}
\def\labelenumi{(\alph{enumi})}
\setcounter{enumi}{2}
\tightlist
\item
  \((-100) \bmod 7\)
\end{enumerate}

\(\left\lfloor -100/7\right\rfloor = \left\lfloor -14.285\ldots\right\rfloor = -15\).
\(-100 - 7\cdot(-15) = -100 + 105 = 5\). \textbf{Answer:} \(5\). (This
is the unique residue with \(0\le r<7\).) 

\begin{enumerate}
\def\labelenumi{(\alph{enumi})}
\setcounter{enumi}{3}
\tightlist
\item
  \((-100) \bmod 0\)
\end{enumerate}

By the stated convention, \(x\bmod 0 = x\). \textbf{Answer:} \(-100\).

\subsection{9 {[}05{]}}\label{section-60}

Definition we'll use

For any real \(x\) and \(y\ne 0\),

\[
x \bmod y \;=\; x - y\big\lfloor x/y\big\rfloor,
\]

and from this it follows that

\[
\begin{cases}
0 \le x\bmod y < y &\text{if } y>0,\\[2pt]
y < x\bmod y \le 0 &\text{if } y<0.
\end{cases}
\]

So with a negative modulus, the remainder is nonpositive and strictly
greater than the modulus.

Computations

\begin{enumerate}
\def\labelenumi{\arabic{enumi}.}
\tightlist
\item
  \(5 \bmod (-3)\)
\end{enumerate}

\[
\left\lfloor \frac{5}{-3}\right\rfloor=\left\lfloor -1.666\ldots\right\rfloor=-2,\quad
5 - (-3)\cdot(-2)=5-6=-1.
\]

Therefore \(5 \bmod (-3)=-1\) (which lies in \((-3,0]\) as required).

\begin{enumerate}
\def\labelenumi{\arabic{enumi}.}
\setcounter{enumi}{1}
\tightlist
\item
  \(18 \bmod (-3)\)
\end{enumerate}

\[
\left\lfloor \frac{18}{-3}\right\rfloor=\left\lfloor -6\right\rfloor=-6,\quad
18 - (-3)\cdot(-6)=18-18=0.
\]

Therefore \(18 \bmod (-3)=0\).

\begin{enumerate}
\def\labelenumi{\arabic{enumi}.}
\setcounter{enumi}{2}
\tightlist
\item
  \((-2) \bmod (-3)\)
\end{enumerate}

\[
\left\lfloor \frac{-2}{-3}\right\rfloor=\left\lfloor 0.666\ldots\right\rfloor=0,\quad
-2 - (-3)\cdot 0=-2.
\]

Therefore \((-2) \bmod (-3)=-2\).

\subsection{10 {[}00{]}}\label{section-61}

We use Knuth's definition

\[
x \bmod y \;=\; x - y\,\big\lfloor x/y \big\rfloor \quad (y\neq 0),
\]

and his noted ranges: if \(y>0\) then \(0\le x\bmod y<y\); if \(y<0\)
then \(0\ge x\bmod y>y\).

\begin{enumerate}
\def\labelenumi{\arabic{enumi}.}
\item
  \(1.1 \bmod 1\) \(\lfloor 1.1/1\rfloor=\lfloor 1.1\rfloor=1\).
  \(1.1 - 1\cdot 1 = 0.1\). (Check: \(y=1>0\), result in \([0,1)\).)
\item
  \(0.11 \bmod 0.1\) \(\lfloor 0.11/0.1\rfloor=\lfloor 1.1\rfloor=1\).
  \(0.11 - 0.1\cdot 1 = 0.01\). (Check: \(y=0.1>0\), result in
  \([0,0.1)\).)
\item
  \(0.11 \bmod (-0.1)\)
  \(\lfloor 0.11/(-0.1)\rfloor=\lfloor -1.1\rfloor=-2\).
  \(0.11 - (-0.1)\cdot(-2) = 0.11 - 0.2 = -0.09\). (Check: \(y=-0.1<0\),
  result in \(( -0.1, 0]\).)
\end{enumerate}

Final answers:

\[
1.1 \bmod 1 = 0.1,\qquad
0.11 \bmod 0.1 = 0.01,\qquad
0.11 \bmod (-0.1) = -0.09.
\]

\subsection{11 {[}00{]}}\label{section-62}

Set \(z=0\) in the congruence definition:

\[
x \equiv y \pmod{0}\quad\Longleftrightarrow\quad x\bmod 0 = y\bmod 0.
\]

By the ``mod-0'' convention, \(x\bmod 0 = x\) and \(y\bmod 0 = y\).
Hence

\[
x \equiv y \pmod{0}\quad\Longleftrightarrow\quad x = y.
\]

(You can also see this from the divisor view: \(0\) divides \(x-y\) only
when \(x-y=0\), i.e., \(x=y\).)

\subsection{12 {[}00{]}}\label{section-63}

\textbf{All integers} (\ldots, −2, −1, 0, 1, 2, \ldots) are relatively
prime to 1.

``Relatively prime'' means the greatest common divisor is 1:
\(\gcd(n,1)=1\) for \textbf{every} integer \(n\).

\begin{itemize}
\tightlist
\item
  The only positive divisor of \(1\) is \(1\) itself.
\item
  Any common divisor of \(n\) and \(1\) must therefore be \(\pm 1\).
\item
  By definition, the greatest \textbf{positive} common divisor is \(1\).
  Hence \(\gcd(n,1)=1\).
\end{itemize}

This remains true for the edge case \(n=0\): \(\gcd(0,1)=1\). Therefore
\(0\) is also coprime to \(1\).

\subsection{13 {[}M00{]}}\label{m00-1}

We need all integers \(m\) such that \(\gcd(m,0)=1\).

By the stated convention, \(\gcd(m,0)=|m|\). Therefore \(\gcd(m,0)=1\)
iff \(|m|=1\). Hence the only integers relatively prime to \(0\) are

\[
m=\pm 1.
\]

Why this is true (two viewpoints)

\begin{enumerate}
\def\labelenumi{\arabic{enumi}.}
\item
  \textbf{Divisors viewpoint.} Every integer divides \(0\). Thus the
  common divisors of \(m\) and \(0\) are exactly the divisors of \(m\),
  so the greatest common divisor is \(|m|\). The condition ``relatively
  prime to \(0\)'' therefore forces \(|m|=1\).
\item
  \textbf{Bézout viewpoint (optional check).} A standard
  characterization says \(m\) and \(n\) are relatively prime iff there
  exist integers \(u,v\) with \(um+vn=1\). Taking \(n=0\), this becomes
  \(um=1\), which has an integer solution only when \(m=\pm1\).
\end{enumerate}

Edge cases

\begin{itemize}
\tightlist
\item
  \(m=0\): Then \(\gcd(0,0)=0\) under the convention (or is left
  undefined in some texts), so \(0\) is \textbf{not} relatively prime to
  \(0\).
\item
  All other \(m\) with \(|m|>1\) fail since \(\gcd(m,0)=|m|>1\).
\end{itemize}

\subsection{14 {[}12{]}}\label{section-64}

\textbf{Problem (TAOCP 1.2.4, Ex. 14):} Find \(x \bmod 15\) given
\(x \equiv 2 \pmod 3\) and \(x \equiv 3 \pmod 5\). 

\textbf{Solution (by substitution / CRT):}

Let \(x = 5k + 3\) to satisfy \(x \equiv 3 \pmod 5\). Reduce mod \(3\):

\[
5k + 3 \equiv 2 \pmod 3 \;\;\Longrightarrow\;\; 5k \equiv -1 \equiv 2 \pmod 3.
\]

Since \(5 \equiv 2 \pmod 3\), we have \(2k \equiv 2 \pmod 3\), hence
\(k \equiv 1 \pmod 3\).

Write \(k = 3t + 1\). Then

\[
x = 5(3t+1) + 3 = 15t + 8 \;\;\Longrightarrow\;\; x \equiv 8 \pmod{15}.
\]

\textbf{Check:} \(8 \bmod 3 = 2\) and \(8 \bmod 5 = 3\). Therefore, the
unique solution modulo \(15\) is

\[
\boxed{x \equiv 8 \pmod{15}}.
\]

\subsection{15 {[}10{]}}\label{section-65}

We consider all cases for \(y,z\in\mathbb{R}\).

\textbf{Case 1: \(y\neq 0\) and \(z\neq 0\).} Using the definition of
mod with floors,

\[
x\bmod y \;=\; x-y\big\lfloor x/y\big\rfloor.
\]

Multiply both sides by \(z\):

\[
z(x\bmod y)=zx-(zy)\big\lfloor x/y\big\rfloor. \tag{1}
\]

Now compute the right-hand side target directly:

\[
(zx)\bmod(zy)
= zx-(zy)\Big\lfloor\frac{zx}{zy}\Big\rfloor.
\]

But \(\frac{zx}{zy}=\frac{x}{y}\) (since \(zy\neq0\)), hence

\[
\Big\lfloor\frac{zx}{zy}\Big\rfloor=\Big\lfloor\frac{x}{y}\Big\rfloor,
\]

so

\[
(zx)\bmod(zy)=zx-(zy)\big\lfloor x/y\big\rfloor,
\]

which matches (1). Therefore \(z(x\bmod y)=(zx)\bmod(zy)\).

\textbf{Case 2: \(z=0\).} Left side \(=0\cdot(x\bmod y)=0\). Right side
\((0\cdot x)\bmod(0\cdot y)=0\bmod 0=0\) by the convention
\(x\bmod 0=x\). Equality holds. 

\textbf{Case 3: \(y=0\) (with any \(z\)).} Left side
\(=z(x\bmod 0)=z\,x\). Right side \((zx)\bmod(0)=zx\). Equality holds by
the same convention. 

All possibilities are covered, so the identity is proved.

Corollary (Knuth's Law C)

With congruence defined by equality of remainders,
\(a\equiv b\pmod m\iff a\bmod m=b\bmod m\), take any nonzero integer
\(n\). Applying the proved law with \(z=n\) and \(y=m\),

\[
n(a\bmod m)=(na)\bmod(nm),\qquad n(b\bmod m)=(nb)\bmod(nm).
\]

Thus \(a\bmod m=b\bmod m\) iff \((na)\bmod(nm)=(nb)\bmod(nm)\), i.e.

\[
a\equiv b\pmod m \iff na\equiv nb\pmod{nm}\quad(n\neq0),
\]

which is precisely Law C.

\subsection{16 {[}M10{]}}\label{m10-2}

\begin{enumerate}
\def\labelenumi{\arabic{enumi}.}
\item
  By hypothesis, \((x-z)/y\in\mathbb Z\). Let

  \[
  q \stackrel{\text{def}}{=} \frac{x-z}{y}\in\mathbb Z,
  \]

  so \(x = qy + z\).
\item
  Divide both sides by \(y>0\):

  \[
  \frac{x}{y} = q + \frac{z}{y}.
  \]

  Because \(0\le z<y\), we have \(0\le \frac{z}{y}<1\). Hence

  \[
  q \le \frac{x}{y} < q+1.
  \]
\item
  The inequality \(q \le x/y < q+1\) implies

  \[
  q=\big\lfloor x/y\big\rfloor.
  \]
\item
  By the definition of the remainder (``mod'' in this section),

  \[
  x\bmod y \;=\; x - y\big\lfloor x/y\big\rfloor
              \;=\; x - yq
              \;=\; z.
  \]
\end{enumerate}

Therefore \(z=x\bmod y\). (Intuitively: \(x=qy+z\) with \(0\le z<y\) is
the Euclidean division of \(x\) by \(y\); the remainder is unique, so it
must be \(z\).)

\textbf{Quick check (example).} Take \(x=17\), \(y=5\). If \(z=2\), then
\((x-z)/y=(17-2)/5=3\in\mathbb Z\) and \(0\le2<5\). Our result says
\(z=17\bmod5=2\), which holds.

\subsection{17 {[}M15{]}}\label{m15-1}

Here is Exercise 17 from §1.2.4, paraphrased: prove Law A directly from
the definition of congruence; and prove the ``⇒'' half of Law D: if
\(a\equiv b\pmod{rs}\) then \(a\equiv b\pmod r\) and
\(a\equiv b\pmod s\) (with \(r,s\) arbitrary integers). The section
defines \(x\equiv y \pmod z\) to mean that \(z\) divides \(x-y\). 

Preliminaries (definition we'll use)

By definition,

\[
x\equiv y\pmod m \quad\Longleftrightarrow\quad m\mid(x-y),
\]

with the usual divisibility rules; when a modulus is negative one may
take its absolute value without changing the congruence class. 

Law A (addition and multiplication are well behaved mod \(m\))

\textbf{Claim.} If \(a\equiv b\pmod m\) and \(x\equiv y\pmod m\), then

\[
a\pm x \equiv b\pm y \pmod m,\qquad ax\equiv by \pmod m.
\]

\textbf{Proof (from the definition).} From \(a\equiv b\pmod m\) and
\(x\equiv y\pmod m\) we have \(m\mid(a-b)\) and \(m\mid(x-y)\). Hence

\[
(a\pm x)-(b\pm y)=(a-b)\pm(x-y),
\]

a sum/difference of integers each divisible by \(m\); therefore \(m\)
also divides the right-hand side, proving
\(a\pm x\equiv b\pm y\pmod m\).

For products,

\[
ax-by=a(x-y)+y(a-b).
\]

Each term on the right is a multiple of \(m\) (since \(m\mid(x-y)\) and
\(m\mid(a-b)\)), so their sum is a multiple of \(m\). Thus
\(ax\equiv by\pmod m\). ∎

(These are exactly the operations stated as Law A in the text.) 

Half of Law D (the easy direction)

\textbf{Claim.} If \(a\equiv b\pmod{rs}\) for integers \(r,s\), then
\(a\equiv b\pmod r\) and \(a\equiv b\pmod s\).

\textbf{Proof (from the definition).} \(a\equiv b\pmod{rs}\) means
\(rs\mid(a-b)\), so \((a-b)=rs\,k\) for some integer \(k\). Then clearly
\(r\mid(a-b)\) (since \(a-b=r(sk)\)) and \(s\mid(a-b)\) (since
\(a-b=s(rk)\)). Hence \(a\equiv b\pmod r\) and \(a\equiv b\pmod s\). ∎

Remarks

The \emph{converse} of this implication generally \textbf{fails} unless
\(r\) and \(s\) are relatively prime; the full bi-implication
``\(a\equiv b\pmod{rs}\) iff \(a\equiv b\pmod r\) and
\(a\equiv b\pmod s\)'' is Law D with the hypothesis \(r\perp s\). (E.g.,
\(2\equiv 8\pmod 3\) and \(2\equiv 8\pmod 6\) but
\(2\not\equiv 8\pmod{18}\).)

\subsection{18 {[}M15{]}}\label{m15-2}

Assume \(r\perp s\) (i.e., \(\gcd(r,s)=1\)). Show that from

\[
a\equiv b \pmod r \quad\text{and}\quad a\equiv b \pmod s
\]

it follows that

\[
a\equiv b \pmod{rs}.
\]

(This is the ``other half'' of Law D.) 

Law B we'll use

If \(c\perp m\) and \(cx\equiv cy\pmod m\), then \(x\equiv y\pmod m\);
in words, when the factor is relatively prime to the modulus, you may
``cancel'' it in a congruence. Also, \(x\perp y\) denotes ``relatively
prime.''

\begin{enumerate}
\def\labelenumi{\arabic{enumi}.}
\tightlist
\item
  From \(a\equiv b \pmod r\) there exists an integer \(u\) with
\end{enumerate}

\[
a=b+ru.
\]

\begin{enumerate}
\def\labelenumi{\arabic{enumi}.}
\setcounter{enumi}{1}
\tightlist
\item
  Reduce this identity modulo \(s\) and use the second hypothesis
  \(a\equiv b\pmod s\):
\end{enumerate}

\[
b+ru \equiv b \pmod s \quad\Longrightarrow\quad ru\equiv 0 \pmod s.
\]

\begin{enumerate}
\def\labelenumi{\arabic{enumi}.}
\setcounter{enumi}{2}
\tightlist
\item
  Because \(r\perp s\), we may \textbf{cancel \(r\) modulo \(s\)} by Law
  B, obtaining
\end{enumerate}

\[
u\equiv 0 \pmod s.
\]

Hence \(u=sw\) for some integer \(w\).

\begin{enumerate}
\def\labelenumi{\arabic{enumi}.}
\setcounter{enumi}{3}
\tightlist
\item
  Substitute back into step (1):
\end{enumerate}

\[
a=b+r(sw)=b+rsw,
\]

which is exactly

\[
a\equiv b \pmod{rs}.
\]

This completes the proof of the \(\Leftarrow\) direction of Law D using
Law B.

\subsection{19 {[}M10{]}}\label{m10-3}

Key tool: Extended Euclid

Algorithm E takes two positive integers \(m,n\), computes
\(d=\gcd(m,n)\), and also finds integers \(a,b\) with

\[
am+bn=d \quad\text{(Bézout identity).}
\]

We'll use the coefficient of \(n\) as the modular inverse when \(d=1\).

Assume \(\gcd(n,m)=1\).

\begin{enumerate}
\def\labelenumi{\arabic{enumi}.}
\tightlist
\item
  Run Extended Euclid on the pair \((m,n)\). It returns integers \(a,b\)
  with
\end{enumerate}

\[
am + bn \;=\; 1.
\]

\begin{enumerate}
\def\labelenumi{\arabic{enumi}.}
\setcounter{enumi}{1}
\tightlist
\item
  Reduce both sides modulo \(m\). Since \(am\equiv 0\pmod m\),
\end{enumerate}

\[
bn \equiv 1 \pmod m.
\]

\begin{enumerate}
\def\labelenumi{\arabic{enumi}.}
\setcounter{enumi}{2}
\tightlist
\item
  Therefore \(n' \equiv b \pmod m\) is a multiplicative inverse of \(n\)
  modulo \(m\). Picking the least nonnegative representative of \(b\)
  yields \(n' \in \{0,1,\dots,m-1\}\); it cannot be \(0\) because
  \(0\cdot n\not\equiv1\pmod m\). Hence \(1\le n'\le m-1\), as desired.
\end{enumerate}

\emph{(Uniqueness mod \(m\).)} If
\(nn'_1 \equiv nn'_2 \equiv 1 \pmod m\), subtract to get
\(n(n'_1-n'_2)\equiv 0\pmod m\). Since \(\gcd(n,m)=1\), cancel \(n\) to
conclude \(n'_1\equiv n'_2\pmod m\). Thus the inverse is unique modulo
\(m\).

This finishes the proof.

Worked example

Find the inverse of \(n=7\) modulo \(m=26\).

Extended Euclid:

\[
\begin{aligned}
26&=3\cdot 7+5,\\
7&=1\cdot 5+2,\\
5&=2\cdot 2+1.
\end{aligned}
\]

Back-substitute:

\[
1 = 5-2\cdot 2
  = 5-2(7-1\cdot 5)
  = 3\cdot 5-2\cdot 7
  = 3(26-3\cdot 7)-2\cdot 7
  = 3\cdot 26-11\cdot 7.
\]

So \(a=3,\; b=-11\) with \(a\cdot 26 + b\cdot 7 = 1\). Hence
\(n' \equiv b \equiv -11 \equiv 15 \pmod{26}\). Check:
\(7\cdot 15=105\equiv 1\pmod{26}\).

How to compute \(n'\) in practice

Run Extended Euclid on \((m,n)\); take the returned coefficient \(b\) of
\(n\); reduce it modulo \(m\) to \(\{1,\dots,m-1\}\). That number is
\(n'\).

\subsection{20 {[}M15{]}}\label{m15-3}

\textbf{Goal (Law B).} If \(ax \equiv by \pmod m\) and
\(a \equiv b \pmod m\), and if \(\gcd(a,m)=1\), then
\(x \equiv y \pmod m\). 

\textbf{Facts we'll use.} • \textbf{Law A (closure under +, −, ×):} from
\(u\equiv v\) and \(r\equiv s\) (mod \(m\)) we may conclude
\(u\pm r\equiv v\pm s\) and \(ur\equiv vs\) (mod \(m\)). • \textbf{Law
of inverses (Ex. 19):} if \(n\perp m\) there exists \(n'\) with
\(nn' \equiv 1 \pmod m\). Apply with \(n=a\). • The exercise to prove
(Ex. 20): use the above to derive Law B.

\begin{enumerate}
\def\labelenumi{\arabic{enumi}.}
\tightlist
\item
  Since \(\gcd(a,m)=1\), by the law of inverses there exists \(a'\) such
  that
\end{enumerate}

\[
aa' \equiv 1 \pmod m .
\]

\begin{enumerate}
\def\labelenumi{\arabic{enumi}.}
\setcounter{enumi}{1}
\tightlist
\item
  Because \(a \equiv b \pmod m\), multiply both sides by \(a'\) and use
  Law A (multiplication preserves congruence):
\end{enumerate}

\[
a'a \equiv a'b \pmod m \quad\Longrightarrow\quad 1 \equiv a'b \pmod m .
\]

So \(a'\) is also a (modular) inverse of \(b\) modulo \(m\). 

\begin{enumerate}
\def\labelenumi{\arabic{enumi}.}
\setcounter{enumi}{2}
\tightlist
\item
  Now start from the given \(ax \equiv by \pmod m\) and multiply both
  sides by \(a'\) (again by Law A):
\end{enumerate}

\[
a'(ax) \equiv a'(by) \pmod m \quad\Longrightarrow\quad (a'a)x \equiv (a'b)y \pmod m .
\]

\begin{enumerate}
\def\labelenumi{\arabic{enumi}.}
\setcounter{enumi}{3}
\tightlist
\item
  Substitute the equalities from steps 1--2:
\end{enumerate}

\[
1\cdot x \equiv 1\cdot y \pmod m \quad\Longrightarrow\quad x \equiv y \pmod m .
\]

This is exactly Law B. ∎

\emph{Remark.} The condition \(\gcd(a,m)=1\) is essential; without it,
``division'' by \(a\) modulo \(m\) need not be valid (see the nearby
exercise asking for a counterexample when \(a\not\perp m\)).

\subsection{21 {[}M22{]}}\label{m22-4}

Preliminaries we're allowed to use

\begin{itemize}
\item
  \textbf{Exercise 1.2.1--5 (existence).} By strong induction: if \(n\)
  is prime we're done; if \(n\) is composite, write \(n=ab\) with
  \(1<a,b<n\) and factor \(a\) and \(b\) by the induction hypothesis;
  concatenating gives a prime product for \(n\). Hence \textbf{every
  \(n>1\) is a product of primes}. 
\item
  \textbf{Law B (cancellation in congruences).} If \(a\perp m\) (i.e.,
  \(\gcd(a,m)=1\)) and \(ax\equiv ay\pmod m\), then
  \(x\equiv y\pmod m\). We will use this to derive Euclid's lemma. 
\end{itemize}

Step 1: Euclid's Lemma from Law B

\begin{quote}
\textbf{Lemma (Euclid).} Let \(p\) be prime. If \(p\mid ab\), then
\(p\mid a\) or \(p\mid b\).
\end{quote}

\textbf{Proof using Law B.} From \(p\mid ab\) we have
\(ab\equiv 0\pmod p\). If \(p\mid a\) we are done. Otherwise
\(\gcd(a,p)=1\) (because \(p\) is prime), so we may cancel \(a\) by Law
B to deduce \(b\equiv 0\pmod p\), i.e., \(p\mid b\). ∎

Step 2: Uniqueness of prime factorization

Assume \(n>1\) and that we already know (by 1.2.1--5) that a prime
factorization exists. Suppose we have \textbf{two} such factorizations

\[
n=p_1p_2\cdots p_k \;=\; q_1q_2\cdots q_\ell,
\]

where each \(p_i\) and \(q_j\) is prime. (We may, and will, reorder each
side nondecreasingly at the end.)

\begin{enumerate}
\def\labelenumi{\arabic{enumi}.}
\item
  Since \(p_1\mid n\), we have \(p_1\mid q_1q_2\cdots q_\ell\). By
  Euclid's lemma there is some \(j\) with \(p_1\mid q_j\). Because
  \(q_j\) is prime, this forces \(p_1=q_j\).
\item
  Cancel that common prime from both sides (divide both sides by
  \(p_1\)); we get
\end{enumerate}

\[
\frac{n}{p_1}=p_2\cdots p_k \;=\; q_1\cdots \widehat{q_j}\cdots q_\ell.
\]

(The hat means ``omit''.)

\begin{enumerate}
\def\labelenumi{\arabic{enumi}.}
\setcounter{enumi}{2}
\tightlist
\item
  Repeat the same argument: the smallest remaining \(p\)-prime divides
  the product of the remaining \(q\)-primes, so it equals one of them;
  cancel and continue.
\end{enumerate}

After finitely many steps all factors are matched and canceled.
Therefore \(k=\ell\) and the two factorizations agree up to order. That
is, \textbf{the multiset of prime factors of \(n\) is unique}.

\subsection{22 {[}M10{]}}\label{m10-4}

If \(a\) shares a nontrivial factor with \(m\), multiplication by \(a\)
can send a \textbf{nonzero} residue to \(0\) modulo \(m\) (i.e., \(a\)
is a zero divisor in \(\mathbb{Z}/m\mathbb{Z}\)). That breaks
cancellation.

A clean general counterexample

Let \(m=ax\) with \(a>1\) and \(x>0\). Then

\[
ax \equiv 0 \pmod m\quad\text{but}\quad x\not\equiv 0 \pmod m
\]

since \(0<x<m\). Take \(b=a\) and \(y=0\). We have \(a\equiv b \pmod m\)
and

\[
ax \equiv a\cdot 0 \pmod m,
\]

yet \(x\not\equiv 0 \pmod m\). Thus the implication asserted by Law B
\textbf{fails} when \(a\) is not relatively prime to \(m\). (This is
exactly the book's official answer in brief form.)

A concrete numeric example

Pick \(m=12\), \(a=6\), \(x=2\), \(y=0\).

\begin{itemize}
\tightlist
\item
  \(a\equiv a \pmod{12}\) is trivial.
\item
  \(ax=6\cdot 2=12\equiv 0 \pmod{12}=a\cdot 0\).
\item
  But \(x=2\not\equiv 0 \pmod{12}\).
\end{itemize}

So all premises of Law B hold \textbf{except} \(a\perp m\), and the
conclusion \(x\equiv y\) fails.

Why this always works (intuition)

When \(m\) is composite and \(\gcd(a,m)>1\), the class \([a]\) is a zero
divisor in \(\mathbb{Z}/m\mathbb{Z}\). Multiplying by a zero divisor
isn't injective, so you cannot cancel it.

Bonus: when cancellation \emph{does} work

If \(\gcd(a,m)=1\), there is an inverse \(a'\) with
\(aa'\equiv 1\pmod m\) (exercise 19), and multiplying \(ax\equiv ay\) by
\(a'\) yields \(x\equiv y\pmod m\) (exercise 20), which is precisely Law
B's proof under its hypothesis.

\subsection{23 {[}M10{]}}\label{m10-5}

Here's a clean counterexample showing why the hypothesis
``\(r\perp s\)'' (i.e., \(\gcd(r,s)=1\)) is necessary in Law D.

Law D (as stated in §1.2.4): If \(r\perp s\), then

\[
a \equiv b \pmod{rs}\quad\Longleftrightarrow\quad
\begin{cases}
a \equiv b \pmod r,\\
a \equiv b \pmod s\;.
\end{cases}
\]

We'll show this fails when \(r\) and \(s\) are not coprime.

Counterexample

Take \(r=2\) and \(s=2\) (so \(\gcd(r,s)=2>1\)). Let \(a=1\) and
\(b=3\).

\begin{itemize}
\tightlist
\item
  Modulo \(r=2\): \(3-1=2\) is divisible by \(2\), so
  \(1\equiv 3\pmod 2\).
\item
  Modulo \(s=2\): same conclusion: \(1\equiv 3\pmod 2\).
\item
  Modulo \(rs=4\): \(3-1=2\) is \textbf{not} divisible by \(4\), so
  \(1\not\equiv 3\pmod 4\).
\end{itemize}

Thus the right-hand side holds (congruence mod \(r\) and mod \(s\)), but
the left-hand side fails (no congruence mod \(rs\)). This shows Law D
``is not always true'' without the coprimality condition---exactly what
Exercise 23 asks for. 

Why this happens (quick insight)

In general,

\[
a\equiv b\pmod r\ \text{and}\ a\equiv b\pmod s
\quad\Longleftrightarrow\quad
a\equiv b\pmod{\operatorname{lcm}(r,s)}.
\]

When \(r\perp s\), \(\operatorname{lcm}(r,s)=rs\) and Law D holds. If
\(\gcd(r,s)>1\), then \(\operatorname{lcm}(r,s)<rs\), so you can have
congruence mod \(r\) and \(s\) but not mod \(rs\)---as in the example
above.

\subsection{24 {[}M20{]}}\label{m20-4}

Step 1 --- Define partial products and show convergence

Let

\[
F_n(x)\;=\;x\,e^{\gamma x}\,\prod_{k=1}^{n}\Bigl(\bigl(1+\tfrac{x}{k}\bigr)e^{-x/k}\Bigr),
\qquad
F(x)=\lim_{n\to\infty}F_n(x).
\]

Taking logs,

\[
\log F_n(x)=\log x+\gamma x+\sum_{k=1}^{n}\Bigl(\log\bigl(1+\tfrac{x}{k}\bigr)-\tfrac{x}{k}\Bigr).
\]

For fixed \(x\), \(\log(1+u)=u-\tfrac{u^2}{2}+O(u^3)\) gives

\[
\log\bigl(1+\tfrac{x}{k}\bigr)-\tfrac{x}{k}= -\,\frac{x^2}{2k^2}+O\!\left(\frac{1}{k^3}\right).
\]

Hence \(\sum_{k\ge1}\bigl(\log(1+\tfrac{x}{k})-\tfrac{x}{k}\bigr)\)
converges absolutely (like \(\sum 1/k^2\)), so \(F_n(x)\) converges to
an analytic function \(F(x)\) (locally uniformly in \(x\)).

Step 2 --- Evaluate \(F(1)\)

For \(x=1\),

\[
F_n(1)=e^\gamma\prod_{k=1}^{n}\Bigl(\tfrac{k+1}{k}\,e^{-1/k}\Bigr)
= e^\gamma\,(n+1)\,e^{-H_n},
\]

where \(H_n=\sum_{k=1}^n \tfrac1k\). Since \(H_n=\ln n+\gamma+o(1)\), we
get \(F_n(1)\to e^\gamma\cdot (n+1)\cdot e^{-(\ln n+\gamma)}\to 1\).
Thus \(F(1)=1\).

Step 3 --- Derive the functional equation

Compute the ratio of partial products:

\[
\frac{F_n(x+1)}{F_n(x)}
=\frac{x+1}{x}\,e^{\gamma}\prod_{k=1}^{n}\frac{1+(x+1)/k}{1+x/k}\,e^{-1/k}
=\frac{x+n+1}{x}\,e^{\gamma-H_n}.
\]

Using again \(H_n=\ln n+\gamma+o(1)\), the right side tends to
\(\frac{1}{x}\). Passing to the limit,

\[
\boxed{\,F(x+1)=\frac{F(x)}{x}\,}\quad\text{for all }x\text{ where }F\text{ is defined.}
\]

Step 4 --- Identify \(F\) with \(1/\Gamma\)

Set \(G(x)=1/F(x)\). Then

\[
G(1)=1, \qquad G(x+1)=x\,G(x).
\]

Moreover,

\[
\log G(x)= -\log x-\gamma x-\sum_{k\ge1}\Bigl(\log\bigl(1+\tfrac{x}{k}\bigr)-\tfrac{x}{k}\Bigr),
\]

so

\[
(\log G)''(x)=\frac{1}{x^2}+\sum_{k\ge1}\frac{1}{(k+x)^2}>0\quad(x>0).
\]

Thus \(G\) is \textbf{log-convex} on \((0,\infty)\). By the
Bohr--Mollerup theorem (the unique log-convex function on \((0,\infty)\)
with \(G(1)=1\) and \(G(x+1)=xG(x)\) is \(\Gamma(x)\)), we conclude

\[
G(x)=\Gamma(x)\quad\text{on }(0,\infty).
\]

Hence \(F(x)=1/\Gamma(x)\) on \((0,\infty)\), and by analytic
continuation the identity holds for all complex \(x\) not a non-positive
integer (where \(\Gamma\) has poles).

\subsection{25 {[}M02{]}}\label{m02}

Fact we'll use (Theorem F)

If \(p\) is prime, then for every integer \(a\),

\[
a^{\,p}\equiv a \pmod p.
\]

Equivalently, when \(p\nmid a\) (so \(\gcd(a,p)=1\)), we also have
\(a^{\,p-1}\equiv 1\pmod p\). 

Consider two exhaustive cases.

\textbf{Case 1: \(p \mid a\).} Then \(a\equiv 0\pmod p\). Any positive
power of \(a\) is also \(0\pmod p\), hence

\[
a^{\,p-1}\equiv 0\pmod p.
\]

The Iverson bracket \([\,p\nmid a\,]\) is \(0\) in this case, so the
identity holds.

\textbf{Case 2: \(p\nmid a\).} Now \(\gcd(a,p)=1\). By Theorem F,
\(a^{\,p}\equiv a\pmod p\). Multiply both sides by \(a^{\,p-2}\) (which
is allowed because \(a\) is invertible modulo \(p\); equivalently, we're
using ``cancellation'' in modular arithmetic). This gives

\[
a^{\,p-1}\equiv 1 \pmod p.
\]

Here \([\,p\nmid a\,]=1\), so again the identity holds.

Combining the two cases:

\[
a^{\,p-1}\bmod p=
\begin{cases}
0,& p\mid a,\\[2pt]
1,& p\nmid a,
\end{cases}
\]

which is exactly \([\,p\nmid a\,]\).

Tiny sanity checks

\begin{itemize}
\tightlist
\item
  \(p=7,a=14\): \(a^{6}\equiv 0\pmod 7\Rightarrow 0\) --- matches
  \([7\nmid 14]=0\).
\item
  \(p=7,a=5\): \(5^{6}\equiv 1\pmod 7\Rightarrow 1\) --- matches
  \([7\nmid 5]=1\).
\end{itemize}

\subsection{26 {[}M15{]}}\label{m15-4}

Work in the ring \(\mathbb{Z}/p\mathbb{Z}\). Because \(p\) is prime,
\(\mathbb{Z}/p\mathbb{Z}\) is a field (in particular, it has no zero
divisors).

Assume \(b^2 \equiv 1 \pmod p\). Then

\[
b^2-1 \equiv 0 \pmod p
\qquad\Longleftrightarrow\qquad
(b-1)(b+1) \equiv 0 \pmod p.
\]

Since there are no zero divisors modulo \(p\), the product of two
residues is \(0\) mod \(p\) only if one of them is \(0\) mod \(p\).
Hence either

\[
b-1 \equiv 0 \pmod p \quad\text{or}\quad b+1 \equiv 0 \pmod p,
\]

i.e.,

\[
b \equiv 1 \pmod p \quad\text{or}\quad b \equiv -1 \pmod p.
\]

This proves the claim. ∎

Why primeness matters (counterexample)

If the modulus isn't prime, extra solutions can appear. For example,
modulo \(15\),

\[
4^2=16\equiv 1 \pmod{15},
\]

but \(4\not\equiv \pm1\ (\bmod\ 15)\). Thus the ``only \(\pm1\)''
conclusion relies essentially on the modulus being prime (so that there
are no zero divisors).

Useful corollary

By Fermat's theorem, if \(p\) is prime and \(q\not\equiv 0\pmod p\),
then \(q^{p-1}\equiv 1\pmod p\). Hence
\(\big(q^{(p-1)/2}\big)^2\equiv 1\pmod p\), and by the result above

\[
q^{(p-1)/2}\equiv \pm 1 \pmod p.
\]

This is exactly why, in Exercise 47, the Legendre symbol
\(\left(\tfrac{q}{p}\right)\) takes only the values \(+1\) or \(-1\)
when \(q\not\equiv 0\pmod p\) (and \(0\) when \(q\equiv 0\pmod p\)).
Knuth explicitly points to Exercise 26 for this fact.
(\href{https://stackoverflow.com/questions/768948/how-does-division-work-in-mix}{PagePlace})

Tiny sanity check

Take \(p=7\). The squares mod \(7\) are \(0,1,2,4\); the only residues
with square \(1\) are \(1\) and \(6\,(=-1)\), matching the theorem.

\subsection{27 {[}M15{]}}\label{m15-5}

\begin{enumerate}
\def\labelenumi{\arabic{enumi})}
\tightlist
\item
  The prime case: ϕ(p)=p−1
\end{enumerate}

Let p be prime. Among the p numbers 0,1,2,\ldots,p−1:

\begin{itemize}
\tightlist
\item
  0 is not relatively prime to p (its gcd with p is p≠1).
\item
  Every k with 1≤k≤p−1 is relatively prime to p, because the only
  positive divisors of p are 1 and p; since k\textless p, p ∤ k, so
  gcd(k,p)=1.
\end{itemize}

Thus exactly the p−1 numbers 1,2,\ldots,p−1 are coprime to p, hence

\[
\varphi(p)=p-1.
\]

\begin{enumerate}
\def\labelenumi{\arabic{enumi})}
\setcounter{enumi}{1}
\tightlist
\item
  Prime powers: ϕ(p\textsuperscript{e)=p}e−p\^{}\{,e-1\}
\end{enumerate}

Let e≥1 and consider the set S=\{0,1,\ldots,p\^{}e−1\}
(\textbar S\textbar=p\^{}e). An integer k∈S is \textbf{not} relatively
prime to p\^{}e iff p\textbar k. Indeed, any common divisor of k and
p\^{}e is a power of p; so gcd(k,p\^{}e)\textgreater1 ⇔ p\textbar k, and
conversely if p\textbar k then gcd(k,p\^{}e)≥p\textgreater1.

How many multiples of p are in S? They are

\[
0,\;p,\;2p,\;3p,\;\dots,\;(p^{e-1}-1)p,
\]

exactly p\^{}\{e-1\} of them. Therefore the integers \textbf{coprime} to
p\^{}e are the remaining ones:

\[
\varphi(p^e)=|S|-|\{ \text{multiples of }p\}|=p^e-p^{e-1}=p^e\!\left(1-\frac1p\right).
\]

This matches the given examples: ϕ(1)=1, ϕ(2)=1, ϕ(3)=2, ϕ(4)=4−2=2,
etc.

\subsection{28 {[}M25{]}}\label{m25-7}

Step 1 --- Closure under multiplication (coprimality is preserved)

Assume \(\gcd(a,m)=\gcd(b,m)=1\). Suppose a prime \(p\) divides both
\(ab \bmod m\) and \(m\). Then \(p\mid ab\) and \(p\mid m\). By Euclid's
lemma, \(p\mid a\) or \(p\mid b\), contradicting
\(\gcd(a,m)=\gcd(b,m)=1\). Hence \(\gcd(ab \bmod m,\, m)=1\). This is
exactly the first observation the exercise nudges you to make.

Step 2 --- Multiplication by \(a\) permutes the reduced residue system

Let \(x_1,\dots,x_{\varphi(m)}\) be \textbf{a reduced residue system
modulo \(m\)} (all incongruent mod \(m\), each coprime to \(m\)).
Because of Step 1, every \(ax_i \bmod m\) is also coprime to \(m\); so
the multiset

\[
\{ax_1 \bmod m,\; ax_2 \bmod m,\; \dots,\; ax_{\varphi(m)} \bmod m\}
\]

consists only of elements from the reduced residue system.

If \(ax_i \equiv ax_j \pmod m\) and \(\gcd(a,m)=1\), we may cancel \(a\)
(since \(a\) has a multiplicative inverse modulo \(m\)) to get
\(x_i \equiv x_j \pmod m\). Thus the map \(x \mapsto ax \bmod m\) is
injective on a finite set of size \(\varphi(m)\), hence a permutation of
the reduced residue system. ({[}cs.ucf.edu{]}{[}2{]})

Step 3 --- Multiply everything and cancel

Multiplying all elements of the reduced residue system before and after
the permutation,

\[
\prod_{i=1}^{\varphi(m)} (ax_i) \equiv \prod_{i=1}^{\varphi(m)} x_i \pmod m.
\]

The left side is \(a^{\varphi(m)} \cdot \prod_{i=1}^{\varphi(m)} x_i\).
Because every \(x_i\) is coprime to \(m\), the product \(P=\prod x_i\)
is also coprime to \(m\), so \(P\) has a multiplicative inverse modulo
\(m\). Cancel \(P\) to obtain

\[
a^{\varphi(m)} \equiv 1 \pmod m,
\]

which is \textbf{Euler's theorem}. 

Check on an example

Take \(m=10\). A reduced residue system is \(\{1,3,7,9\}\) (so
\(\varphi(10)=4\)). For \(a=3\) (coprime to 10):

\[
3\cdot\{1,3,7,9\}\equiv \{3,9,1,7\}\pmod{10},
\]

a permutation of \(\{1,3,7,9\}\). Multiplying and canceling gives
\(3^4\equiv1\pmod{10}\).

\subsection{29 {[}M22{]}}\label{m22-5}

\begin{enumerate}
\def\labelenumi{(\alph{enumi})}
\tightlist
\item
  \(f(n)=n^{c}\) is multiplicative (in fact, completely multiplicative)
\end{enumerate}

For any positive integers \(r,s\) (coprime or not),

\[
f(rs)=(rs)^{c}=r^{c}\,s^{c}=f(r)\,f(s),
\]

by the basic law of exponents. Thus \(f\) is \textbf{completely}
multiplicative (the coprimality condition isn't even needed).

\begin{enumerate}
\def\labelenumi{(\alph{enumi})}
\setcounter{enumi}{1}
\tightlist
\item
  The squarefree indicator is multiplicative
\end{enumerate}

Let \(f(n)=1\) if \(n\) is squarefree and \(0\) otherwise. Suppose
\(r\perp s\).

\begin{itemize}
\item
  If \textbf{both} \(r\) and \(s\) are squarefree, then every prime
  \(p\) appears with exponent \(0\) or \(1\) in each; since the supports
  of their prime factors are disjoint (coprimality), in \(rs\) every
  prime still has exponent \(0\) or \(1\). Hence \(rs\) is squarefree
  and \(f(rs)=1=f(r)f(s)\).
\item
  Conversely, if \(rs\) is squarefree, write the \(p\)-adic valuations
  as \(v_p(rs)=v_p(r)+v_p(s)\). Squarefreeness means \(v_p(rs)\le1\) for
  all \(p\), so for each \(p\) we must have \(v_p(r)\le1\) and
  \(v_p(s)\le1\). Thus both \(r\) and \(s\) are squarefree. Therefore
  \(f(rs)=f(r)f(s)\).
\end{itemize}

Hence \(f\) is multiplicative (but not completely multiplicative; e.g.,
\(f(p^2)=0\neq 1=f(p)f(p)\)).

\begin{enumerate}
\def\labelenumi{(\alph{enumi})}
\setcounter{enumi}{2}
\tightlist
\item
  \(f(n)=c^{\omega(n)}\) is multiplicative
\end{enumerate}

Let \(\omega(n)\) be the number of \textbf{distinct} prime divisors of
\(n\). If \(r\perp s\), their prime supports are disjoint, so

\[
\omega(rs)=\omega(r)+\omega(s).
\]

Therefore

\[
f(rs)=c^{\omega(rs)}=c^{\omega(r)+\omega(s)}=c^{\omega(r)}\,c^{\omega(s)}=f(r)\,f(s).
\]

(Again, not completely multiplicative in general:
\(f(p^2)=c\neq c^2=f(p)^2\) unless \(c=1\).)

\begin{enumerate}
\def\labelenumi{(\alph{enumi})}
\setcounter{enumi}{3}
\tightlist
\item
  The product of multiplicative functions is multiplicative
\end{enumerate}

Let \(f,g\) be multiplicative and define \((fg)(n)=f(n)g(n)\). For
\(r\perp s\),

\[
(fg)(rs)=f(rs)g(rs)=\big(f(r)f(s)\big)\big(g(r)g(s)\big)
=\underbrace{f(r)g(r)}_{(fg)(r)}\;\underbrace{f(s)g(s)}_{(fg)(s)}
=(fg)(r)\,(fg)(s),
\]

so \(fg\) is multiplicative.

\subsection{30 {[}M30{]}}\label{m30-7}

Facts we'll use

\begin{enumerate}
\def\labelenumi{\arabic{enumi}.}
\item
  Definition. φ(n) counts residues modulo n that are coprime to
  n.~(Introduced in Ex. 27.) 
\item
  Prime powers. For a prime p and e≥1,
\end{enumerate}

\[
\varphi(p^e)=p^e-p^{e-1}=p^{e-1}(p-1),
\]

because among the p\^{}e residues, exactly the multiples of p (there are
p\^{}\{e-1\} of them) are not coprime to p\^{}e. (This is exactly what
Ex. 27 asks you to establish.) 

φ is multiplicative

We must show that if m and n are coprime (m ⊥ n), then

\[
\varphi(mn)=\varphi(m)\,\varphi(n).
\]

Consider the set of units (numbers coprime to the modulus) modulo mn:

\[
U_{mn}=\{\,x\bmod mn : \gcd(x,mn)=1\,\}.
\]

Define a map

\[
\Phi:U_{mn}\to U_m\times U_n,\qquad \Phi([x]_{mn})=([x]_m,\,[x]_n).
\]

If \(\gcd(m,n)=1\), the Chinese Remainder Theorem (CRT) says that the
pair of congruences

\[
y\equiv a \pmod m,\quad y\equiv b \pmod n
\]

has a unique solution modulo mn. Moreover, \(x\) is coprime to \(mn\)
iff it is coprime to each of \(m\) and \(n\). Therefore:

\begin{itemize}
\tightlist
\item
  \textbf{Well-defined/injective:} If \([x]_{mn}=[x']_{mn}\), then
  \([x]_m=[x']_m\) and \([x]_n=[x']_n\).
\item
  \textbf{Surjective:} For any \(([a]_m,[b]_n)\in U_m\times U_n\), CRT
  gives a unique \([y]_{mn}\) with \(y\equiv a\ (\bmod m)\),
  \(y\equiv b\ (\bmod n)\); since \(a\) and \(b\) are each coprime to
  their moduli, \(y\) is coprime to \(mn\).
\end{itemize}

Thus \(\Phi\) is a bijection, so \(|U_{mn}|=|U_m|\cdot|U_n|\), i.e.

\[
\boxed{\ \varphi(mn)=\varphi(m)\,\varphi(n)\quad \text{when }\gcd(m,n)=1\ }.
\]

Evaluate φ(1 000 000)

Factor \(1{,}000{,}000=10^6=2^6\cdot 5^6\). Using prime-power values and
multiplicativity,

\[
\varphi(10^6)=\varphi(2^6)\,\varphi(5^6)
= \bigl(2^6-2^5\bigr)\bigl(5^6-5^5\bigr)
= (64-32)\,(15625-3125)
= 32\cdot 12500.
\]

Compute: \(12500\times 30=375{,}000\); \(12500\times 2=25{,}000\); sum
\(=400{,}000\). Hence

\[
\boxed{\ \varphi(1{,}000{,}000)=400{,}000\ }.
\]

General method once the prime factorization is known

If

\[
n=\prod_{i=1}^r p_i^{e_i}
\quad\text{(distinct primes \(p_i\))},
\]

then by multiplicativity and the prime-power formula,

\[
\varphi(n)=\prod_{i=1}^r \varphi\!\bigl(p_i^{e_i}\bigr)
=\prod_{i=1}^r \bigl(p_i^{e_i}-p_i^{e_i-1}\bigr)
= n\prod_{i=1}^r \Bigl(1-\frac1{p_i}\Bigr).
\]

So the algorithm is:

\begin{enumerate}
\def\labelenumi{\arabic{enumi}.}
\tightlist
\item
  Factor n: \(n=\prod p_i^{e_i}\).
\item
  Compute \(n\prod_i (1-1/p_i)\).
\end{enumerate}

\subsection{31 {[}M22{]}}\label{m22-6}

Key fact

If \(\gcd(r,s)=1\), then every divisor \(d\) of \(rs\) factors
\textbf{uniquely} as \(d=ab\) with \(a\mid r\) and \(b\mid s\); moreover
\(\gcd(a,b)=1\). (If a prime divided both \(a\) and \(b\), it would
divide both \(r\) and \(s\), contradicting \(\gcd(r,s)=1\).)

Proof that \(g\) is multiplicative

Take coprime positive integers \(r,s\). Then

\[
g(rs)
=\sum_{d\mid rs} f(d)
=\sum_{\substack{a\mid r\\ b\mid s}} f(ab)
\]

by the divisor factorization just noted. Since \(\gcd(a,b)=1\) and \(f\)
is multiplicative, \(f(ab)=f(a)f(b)\). Therefore

\[
g(rs)
=\sum_{a\mid r}\sum_{b\mid s} f(a)f(b)
=\Bigl(\sum_{a\mid r} f(a)\Bigr)\Bigl(\sum_{b\mid s} f(b)\Bigr)
=g(r)\,g(s).
\]

Hence \(g\) is multiplicative. ∎

Prime-power and Euler-product view (useful checks)

Let \(n=\prod_{p} p^{\nu_p(n)}\). The proof implies

\[
g(n)=\prod_{p\mid n}\Bigl(\,f(1)+f(p)+\cdots+f(p^{\nu_p(n)})\Bigr),
\]

so multiplicativity is also apparent from the factorization over
distinct primes.

Quick examples

\begin{enumerate}
\def\labelenumi{\arabic{enumi}.}
\item
  \(f(n)\equiv 1\). Then \(g(n)=\sum_{d\mid n}1=\tau(n)\) (divisor
  count), which is multiplicative.
\item
  \(f(n)=n^{c}\) (multiplicative for any fixed real \(c\)). Then
\end{enumerate}

\[
g(n)=\prod_{p^{k}\parallel n}\bigl(1+p^{c}+\cdots+p^{ck}\bigr),
\]

again a multiplicative function.

\subsection{32 {[}M18{]}}\label{m18-1}

Think of the left-hand double sum as running over all \emph{chains of
divisors}

\[
c \mid d \mid n.
\]

Define the two index sets

\[
A=\{(c,d): c\mid d,\ d\mid n\},\qquad
B=\{(c,e): c\mid n,\ e\mid (n/c)\}.
\]

Consider the bijection

\[
\Phi:A\to B,\qquad \Phi(c,d)=(c,e)\ \text{with}\ e:=d/c.
\]

\begin{itemize}
\item
  \textbf{Well-defined.} If \(c\mid d\) then \(e=d/c\) is a positive
  integer. Since \(d\mid n\), we have \(ce=d\mid n\), hence
  \(e\mid n/c\). Also \(c\mid n\) because \(c\mid d\mid n\). Thus
  \((c,e)\in B\).
\item
  \textbf{Inverse map.} Given \((c,e)\in B\) with \(c\mid n\) and
  \(e\mid n/c\), set \(d:=ce\). Then \(c\mid d\) obviously, and
  \(d=ce\mid c\cdot (n/c)=n\), so \((c,d)\in A\). This is the inverse of
  \(\Phi\).
\end{itemize}

Because \(\Phi\) is a bijection between the exact index sets of the two
sums and it transforms the summand \(f(c,d)\) on the left into
\(f\big(c, c e\big)\) on the right (with \(e=d/c\)), we get

\[
\sum_{(c,d)\in A} f(c,d)\;=\;\sum_{(c,e)\in B} f\!\big(c,ce\big),
\]

which is precisely the desired identity.

\(\square\)

\subsection{33 {[}M18{]}}\label{m18-2}

Exactly one of the two halves

\[
x=\frac{n+m}{2},\qquad y=\frac{n-m+1}{2}
\]

is an integer and the other is a half--integer (i.e., has fractional
part \(1/2\)).

Reason: \(n+m\) and \(n-m\) always have the same parity. So:

\begin{itemize}
\tightlist
\item
  If \(n+m\) is even, then \(x\in\mathbb Z\). Since \(n-m\) is even,
  \(n-m+1\) is odd, hence \(y\) is a half--integer.
\item
  If \(n+m\) is odd, then \(x\) is a half--integer. Since \(n-m\) is
  odd, \(n-m+1\) is even, hence \(y\in\mathbb Z\).
\end{itemize}

Also note

\[
x+y=\frac{n+m}{2}+\frac{n-m+1}{2}=n+\tfrac12.
\]

We'll use the standard definitions of floor/ceiling (greatest/largest
integer \(\le\)/\(\ge\) a real), as introduced earlier in §1.2.4. 

Part (a): sum of floors

Since \(x+y=n+\tfrac12\), we have \(\lfloor x+y\rfloor=n\). From the
parity analysis, one of \(x,y\) is an integer (fractional part \(0\))
and the other has fractional part \(1/2\). Therefore
\(\{x\}+\{y\}=1/2<1\), so

\[
\lfloor x\rfloor+\lfloor y\rfloor=\lfloor x+y\rfloor=n.
\]

Hence

\[
\boxed{\;\Big\lfloor\tfrac{n+m}{2}\Big\rfloor+\Big\lfloor\tfrac{n-m+1}{2}\Big\rfloor=n\;}
\]

Part (b): sum of ceilings

Again, exactly one of \(x,y\) is an integer and the other is a
half--integer. Taking ceilings adds \(0\) to the integer one and adds
\(1\) to the half--integer one. Therefore

\[
\lceil x\rceil+\lceil y\rceil=\lfloor x\rfloor+\lfloor y\rfloor+1=n+1.
\]

Equivalently, since \(x+y=n+\tfrac12\), \(\lceil x+y\rceil=n+1\) and the
same ``one integer, one half--integer'' parity argument shows no extra
carry beyond that single \(+1\). Thus

\[
\boxed{\;\Big\lceil\tfrac{n+m}{2}\Big\rceil+\Big\lceil\tfrac{n-m+1}{2}\Big\rceil=n+1\;}
\]

Noted special case \(m=0\)

\begin{itemize}
\tightlist
\item
  \(\lfloor n/2\rfloor+\lfloor(n+1)/2\rfloor=n\).
\item
  \(\lceil n/2\rceil+\lceil(n+1)/2\rceil=n+1\).
\end{itemize}

These are handy ``split-\(n\)-in-halves'' identities that often simplify
parity arguments.

\subsection{34 {[}M21{]}}\label{m21-1}

\textbf{Answer.} Exactly for \textbf{integer bases \(b\ge 2\)}. The same
necessary and sufficient condition applies to the ceiling version
\(\;\lceil \log_b x\rceil=\lceil \log_b \lceil x\rceil\rceil.\) 

Proof (⇒ necessity)

Assume the floor identity holds for all \(x\ge 1\). Plug in \(x=b\):

\[
\big\lfloor \log_b b\big\rfloor = 1 \stackrel{!}= \big\lfloor \log_b \lfloor b\rfloor \big\rfloor .
\]

Now \(\lfloor b\rfloor<b\) unless \(b\) is an integer. If
\(b\notin\mathbb Z\), then \(\log_b\lfloor b\rfloor<1\) so the
right-hand side is \(0\), contradicting \(1\). Also \(b>1\) forces
\(b\ge 2\) when \(b\in\mathbb Z\). Hence \(b\) must be an
\textbf{integer \(\ge 2\)}.

Proof (⇐ sufficiency)

Suppose \(b\) is an integer \(\ge 2\). For any \(x\ge 1\), let

\[
n=\big\lfloor \log_b x\big\rfloor \quad\Longleftrightarrow\quad b^n \le x < b^{n+1}.
\]

Because \(b\) is an integer, every \(b^k\) is an integer. Thus

\[
b^n \le x \;<\; b^{n+1}\quad\Longrightarrow\quad
b^n \le \lfloor x\rfloor \le x < b^{n+1}.
\]

Taking \(\log_b\) preserves inequalities (since \(b>1\)):

\[
n \le \log_b\lfloor x\rfloor < n+1,
\]

so
\(\big\lfloor \log_b \lfloor x\rfloor\big\rfloor = n = \big\lfloor \log_b x\big\rfloor\).
This proves the identity for all \(x\ge 1\).

(The ceiling version is proved identically by bracketing \(x\) between
consecutive integer powers of \(b\) and noting that \(b^n\) are integers
when \(b\in\mathbb Z\).) 

Why non-integer bases fail

\begin{itemize}
\tightlist
\item
  If \(1<b<2\), then \(\lfloor b\rfloor=1\) and
  \(\lfloor \log_b b\rfloor=1\) but
  \(\lfloor\log_b\lfloor b\rfloor\rfloor=\lfloor\log_b 1\rfloor=0\).
\item
  If \(b>2\) is non-integer, \(\lfloor b\rfloor\in\{2,3,\dots\}\) and
  still \(\lfloor \log_b \lfloor b\rfloor\rfloor=0\) while
  \(\lfloor \log_b b\rfloor=1\).
\end{itemize}

Either way we get a contradiction at \(x=b\).

A nice generalization (McEliece)

Let \(f\) be continuous and strictly increasing on an interval \(A\),
with the property that whenever \(x\in A\) then
\(\lfloor x\rfloor,\lceil x\rceil\in A\). Then the following are
equivalent:

\begin{itemize}
\tightlist
\item
  \(\lfloor f(x)\rfloor=\lfloor f(\lfloor x\rfloor)\rfloor\) for all
  \(x\in A\).
\item
  \(\lceil f(x)\rceil=\lceil f(\lceil x\rceil)\rceil\) for all
  \(x\in A\).
\item
  \textbf{Criterion:} whenever \(f(x)\) is an integer, then \(x\) is an
  integer.
\end{itemize}

Taking \(f(x)=\log_b x\) on \(A=[1,\infty)\) makes the criterion ``if
\(\log_b x\in\mathbb Z\) then \(x\in\mathbb Z\)'', which is true exactly
when \(b\) is an integer \(\ge 2\). 

Quick sanity checks

\begin{itemize}
\tightlist
\item
  \(b=10\) (decimal): \(\lfloor\log_{10} x\rfloor\) is ``(number of
  digits of \(\lfloor x\rfloor\)) −1'', so equality is immediate.
\item
  \(b=\pi\): taking \(x=\pi\) gives \(1\) on the left and \(0\) on the
  right---fails.
\end{itemize}

\subsection{35 {[}M20{]}}\label{m20-5}

Write \(x=\lfloor x\rfloor + y\) with \(0\le y<1\) (so \(y\) is the
fractional part of \(x\)). Set \(M=\lfloor x\rfloor + m\), which is an
integer. By the division algorithm there exist unique integers \(k\) and
\(r\) with

\[
M = kn + r,\qquad 0\le r\le n-1.
\]

Then

\[
\frac{x+m}{n}
= \frac{y+M}{n}
= \frac{y+kn+r}{n}
= k + \frac{y+r}{n}.
\]

Because \(0\le y<1\) and \(0\le r\le n-1\), we have \(0\le y+r < n\).
Hence

\[
\left\lfloor k + \frac{y+r}{n}\right\rfloor = k.
\]

But
\(k=\left\lfloor \dfrac{M}{n} \right\rfloor = \left\lfloor \dfrac{\lfloor x\rfloor + m}{n} \right\rfloor\).
Therefore

\[
\left\lfloor \frac{x+m}{n}\right\rfloor
=
\left\lfloor \frac{\lfloor x\rfloor + m}{n}\right\rfloor,
\]

as required.
(\href{https://stackoverflow.com/questions/768948/how-does-division-work-in-mix}{ProofWiki})

Useful corollaries and dual form

\begin{enumerate}
\def\labelenumi{\arabic{enumi}.}
\tightlist
\item
  Setting \(m=0\) gives
\end{enumerate}

\[
\left\lfloor \frac{x}{n}\right\rfloor
=
\left\lfloor \frac{\lfloor x\rfloor}{n}\right\rfloor,
\]

i.e., only \(\lfloor x\rfloor\) matters for the floor of a quotient by
\(n>0\). ({[}ProofWiki{]}{[}2{]})

\begin{enumerate}
\def\labelenumi{\arabic{enumi}.}
\setcounter{enumi}{1}
\tightlist
\item
  By \(\lceil z\rceil=-\lfloor -z\rfloor\), a completely analogous
  ceiling identity holds:
\end{enumerate}

\[
\left\lceil \frac{x+m}{n}\right\rceil
=
\left\lceil \frac{\lceil x\rceil+m}{n}\right\rceil
\qquad(n>0,\,m\in\mathbb Z,\,x\in\mathbb R).
\]

\begin{enumerate}
\def\labelenumi{\arabic{enumi}.}
\setcounter{enumi}{2}
\tightlist
\item
  The exercise is frequently used to simplify nested floors/ceilings in
  this chapter. For example, Knuth invokes it in §1.2.4 Answer 48(b)
  when manipulating \(\lceil(n+\lceil -n/25\rceil)/3\rceil\). 
\end{enumerate}

Quick numeric check

Take \(x=17.83\), \(m=-10\), \(n=7\).

\[
\left\lfloor \frac{x+m}{n}\right\rfloor
= \left\lfloor \frac{7.83}{7}\right\rfloor = \left\lfloor 1.118\ldots\right\rfloor = 1,
\]

\[
\left\lfloor \frac{\lfloor x\rfloor+m}{n}\right\rfloor
= \left\lfloor \frac{17-10}{7}\right\rfloor
= \left\lfloor \frac{7}{7}\right\rfloor = 1.
\]

\subsection{36 {[}M23{]}}\label{m23-7}

Part A --- \(\sum_{k=1}^{n}\big\lfloor k/2\big\rfloor\)

We'll split by the parity of \(n\).

Case 1: \(n=2m\) even

Pair consecutive terms \((2j-1,2j)\) for \(j=1,\dots,m\):

\[
\Big\lfloor\frac{2j-1}{2}\Big\rfloor+\Big\lfloor\frac{2j}{2}\Big\rfloor=(j-1)+j=2j-1.
\]

Summing those pairs gives

\[
\sum_{k=1}^{2m}\Big\lfloor\frac{k}{2}\Big\rfloor
=\sum_{j=1}^{m}(2j-1)=m^{2}
=\frac{(2m)^{2}}{4}
=\Big\lfloor\frac{n^{2}}{4}\Big\rfloor.
\]

Case 2: \(n=2m+1\) odd

Use the even case up to \(2m\) and add the last term:

\[
\sum_{k=1}^{2m+1}\Big\lfloor\frac{k}{2}\Big\rfloor
=\underbrace{\sum_{k=1}^{2m}\Big\lfloor\frac{k}{2}\Big\rfloor}_{=\,m^{2}}
+\Big\lfloor\frac{2m+1}{2}\Big\rfloor
=m^{2}+m
=\Big\lfloor\frac{(2m+1)^{2}}{4}\Big\rfloor
=\Big\lfloor\frac{n^{2}}{4}\Big\rfloor.
\]

Thus in all cases

\[
\boxed{\displaystyle \sum_{k=1}^{n}\Big\lfloor \frac{k}{2}\Big\rfloor=\Big\lfloor \frac{n^{2}}{4}\Big\rfloor.}
\]

(For reference, Knuth's answers section derives the same result by
splitting into \(n=2t\) and \(n=2t+1\).) 

Part B --- \(\sum_{k=1}^{n}\big\lceil k/2\big\rceil\)

Use the elementary identity (true for every integer \(k\)):

\[
\Big\lfloor\frac{k}{2}\Big\rfloor+\Big\lceil\frac{k}{2}\Big\rceil=k.
\]

Summing over \(k=1,\dots,n\) yields

\[
\sum_{k=1}^{n}\Big\lceil\frac{k}{2}\Big\rceil
=\sum_{k=1}^{n}k-\sum_{k=1}^{n}\Big\lfloor\frac{k}{2}\Big\rfloor
=\frac{n(n+1)}{2}-\Big\lfloor\frac{n^{2}}{4}\Big\rfloor.
\]

Evaluate by parity:

\begin{itemize}
\item
  If \(n=2m\):

  \[
  \sum_{k=1}^{n}\Big\lceil\frac{k}{2}\Big\rceil
  =\frac{2m(2m+1)}{2}-m^{2}=m(m+1).
  \]
\item
  If \(n=2m+1\):

  \[
  \sum_{k=1}^{n}\Big\lceil\frac{k}{2}\Big\rceil
  =\frac{(2m+1)(2m+2)}{2}-m(m+1)=(m+1)^{2}.
  \]
\end{itemize}

Both cases unify neatly as

\[
\boxed{\displaystyle \sum_{k=1}^{n}\Big\lceil \frac{k}{2}\Big\rceil=\Big\lceil \frac{n(n+2)}{4}\Big\rceil.}
\]

(This is also the closed form reported in the book's answers.) 

Quick sanity checks

\begin{itemize}
\tightlist
\item
  \(n=5\): \(\lfloor k/2\rfloor\) sums to
  \(0+1+1+2+2=6=\lfloor 25/4\rfloor\); \(\lceil k/2\rceil\) sums to
  \(1+1+2+2+3=9=\lceil 5\cdot7/4\rceil\).
\item
  \(n=6\): \(\lfloor k/2\rfloor\) sums to \(0+1+1+2+2+3=9=6^{2}/4\);
  \(\lceil k/2\rceil\) sums to
  \(1+1+2+2+3+3=12=\lceil 6\cdot8/4\rceil\).
\end{itemize}

All consistent.

\subsection{37 {[}M30{]}}\label{m30-8}

Step 1: Separate floors into value minus fractional part

\[
\sum_{k=0}^{n-1}\Big\lfloor\frac{mk+x}{n}\Big\rfloor
=\sum_{k=0}^{n-1}\Big(\frac{mk+x}{n}-\Big\{\frac{mk+x}{n}\Big\}\Big)
=\frac{m(n-1)}{2}+x-\;S,
\]

where

\[
S=\sum_{k=0}^{n-1}\Big\{\frac{mk+x}{n}\Big\}.
\]

Step 2: Structure of the fractional parts

Put \(t=n/d,\;u=m/d\) so \(\gcd(u,t)=1\). The multiset of fractional
parts over one period \(k=0,\dots,t-1\) is

\[
\Big\{\frac{x}{n}+\frac{u\,k}{t}\Big\},\qquad k=0,\dots,t-1.
\]

Because multiplication by \(u\) permutes the residues modulo \(t\), this
set is the same as

\[
\Big\{\frac{x\bmod d}{n}+\frac{r}{t}\Big\},\qquad r=0,\dots,t-1.
\]

Crucially, \(\dfrac{x\bmod d}{n}<\dfrac{d}{n}=\dfrac1t\), so no term
wraps around \(1\). Hence

\[
\sum_{r=0}^{t-1}\Big\{\frac{x\bmod d}{n}+\frac{r}{t}\Big\}
=t\cdot\frac{x\bmod d}{n}+\sum_{r=0}^{t-1}\frac{r}{t}
=\frac{t(x\bmod d)}{n}+\frac{t-1}{2}.
\]

There are \(d\) identical blocks of length \(t\), so

\[
S=d\!\left(\frac{t(x\bmod d)}{n}+\frac{t-1}{2}\right)
=(x\bmod d)+\frac{n-d}{2}.
\]

Step 3: Assemble the floor sum

Plugging \(S\) back:

\[
\sum_{k=0}^{n-1}\Big\lfloor\frac{mk+x}{n}\Big\rfloor
=\frac{m(n-1)}{2}+x-\Big((x\bmod d)+\frac{n-d}{2}\Big)
=\frac{m(n-1)-(n-d)}{2}+d\Big\lfloor\frac{x}{d}\Big\rfloor.
\]

Two convenient equivalent forms are

\[
\boxed{\;
\sum_{k=0}^{n-1}\Big\lfloor\frac{mk+x}{n}\Big\rfloor
=\frac{(m-1)(n-1)+(d-1)}{2}+d\Big\lfloor\frac{x}{d}\Big\rfloor
\;}
\]

and

\[
\sum_{k=0}^{n-1}\Big\lfloor\frac{mk+x}{n}\Big\rfloor
=\frac{m(n-1)}{2}-\frac{n-d}{2}+d\Big\lfloor\frac{x}{d}\Big\rfloor.
\]

(Optional) Ceiling version

Using \(\lceil y\rceil=-\lfloor -y\rfloor\) (or
\(\lceil y\rceil=\lfloor y\rfloor+[\,\{y\}>0]\)), one gets the companion
identity

\[
\boxed{\;
\sum_{k=0}^{n-1}\Big\lceil\frac{mk+x}{n}\Big\rceil
=\frac{(m+1)(n-1)-(d-1)}{2}+d\Big\lceil\frac{x}{d}\Big\rceil
\;}
\]

which matches Knuth's ``similar identity.''

Checks and special cases

\begin{itemize}
\tightlist
\item
  If \(\gcd(m,n)=1\) (\(d=1\)):
\end{itemize}

\[
\sum_{k=0}^{n-1}\Big\lfloor\frac{mk+x}{n}\Big\rfloor
=\frac{(m-1)(n-1)}{2}+\lfloor x\rfloor,
\]

the classical result.

\begin{itemize}
\tightlist
\item
  Taking \(x=0\) gives
\end{itemize}

\[
\sum_{k=0}^{n-1}\Big\lfloor\frac{mk}{n}\Big\rfloor
=\frac{(m-1)(n-1)+(d-1)}{2}.
\]

\subsection{38 {[}M26{]}}\label{m26-2}

Let

\[
a=\lfloor x\rfloor,\qquad \theta=\{x\}=x-\lfloor x\rfloor\in[0,1),\qquad c=\lceil y\rceil\in\mathbb{Z}.
\]

The summation index set is \(\{\,k\in\mathbb{Z}\mid 0\le k<y\,\}\).
Since \(y>0\), those are precisely the integers \(k=0,1,\dots,c-1\)
(when \(y\) is not an integer this runs up to \(\lfloor y\rfloor=c-1\);
when \(y\) is an integer \(n\), this runs up to \(n-1=c-1\)). Thus the
number of terms is \(c=\lceil y\rceil\).

Because \(0\le k/y<1\) for all such \(k\), we have

\[
\Big\lfloor x+\frac{k}{y}\Big\rfloor
=\Big\lfloor a+\theta+\frac{k}{y}\Big\rfloor
= a+\Big\lfloor \theta+\frac{k}{y}\Big\rfloor,
\]

and \(\big\lfloor \theta+\frac{k}{y}\big\rfloor\) is \(1\) exactly when
\(\theta+\frac{k}{y}\ge 1\), i.e.~when \(k\ge y(1-\theta)\); otherwise
it is \(0\).

Hence

\[
S:=\sum_{0\le k<y}\Big\lfloor x+\frac{k}{y}\Big\rfloor
= \sum_{k=0}^{c-1}\Big(a+\Big\lfloor \theta+\frac{k}{y}\Big\rfloor\Big)
= ac + N,
\]

where \(N\) is the count of integers \(k\in\{0,1,\dots,c-1\}\) with
\(k\ge y(1-\theta)\). The smallest such \(k\) is
\(\lceil y(1-\theta)\rceil\), so

\[
N = \#\{\lceil y(1-\theta)\rceil,\;\lceil y(1-\theta)\rceil+1,\dots,c-1\}
= c - \lceil y(1-\theta)\rceil.
\]

Therefore

\[
S = ac + \big(c - \lceil y(1-\theta)\rceil\big) = (a+1)c - \lceil y(1-\theta)\rceil. \tag{★}
\]

Now rewrite the right-hand side of the exercise's formula:

\[
xy+\lfloor x+1\rfloor\,(\lceil y\rceil-y)
= xy+(a+1)(c-y) = (a+1)c - y(1-\theta).
\]

Since \((a+1)c\) is an integer,

\[
\Big\lfloor (a+1)c - y(1-\theta)\Big\rfloor
= (a+1)c - \Big\lceil y(1-\theta)\Big\rceil,
\]

and by (★) this equals \(S\). This proves the general identity.

Special case \(y=n\in\mathbb{Z}_{>0}\)

When \(y=n\), we have \(c=\lceil y\rceil=n\) and (★) becomes

\[
S = (a+1)n - \lceil n(1-\theta)\rceil.
\]

Using \(\lceil n(1-\theta)\rceil = n - \lfloor n\theta\rfloor\) (valid
for \(\theta\in[0,1)\)), we get

\[
S = (a+1)n - \big(n-\lfloor n\theta\rfloor\big)= an + \lfloor n\theta\rfloor
= \lfloor na+n\theta\rfloor = \lfloor nx\rfloor,
\]

which is the ``in particular'' statement and shows that
\(\lfloor\cdot\rfloor\) is \textbf{replicative} (the notion defined in
the very next exercise). 

Remarks

\begin{itemize}
\item
  The identity immediately implies the replicative property for the
  floor function:

  \[
  \sum_{k=0}^{n-1}\Big\lfloor x+\frac{k}{n}\Big\rfloor = \lfloor nx\rfloor\quad(n\in\mathbb{Z}_{>0}),
  \]

  which Exercise 39 then uses to develop the general theory of
  \textbf{replicative functions}.
\item
  The proof strategy is a standard ``count the carries'' argument: split
  \(x\) into integer and fractional parts and count how many increments
  push \(\theta+k/y\) over the next integer threshold.
\end{itemize}

\subsection{39 {[}HM35{]}}\label{hm35}

Definition (replicative)

A function \(f\) is \textbf{replicative} if, for every positive integer
\(n\),

\[
\sum_{k=0}^{n-1} f\!\left(x+\frac{k}{n}\right)=f(nx)\qquad\text{for all real }x.
\]

(Exercise 38 just before this one shows that
\(x\mapsto\lfloor x\rfloor\) is replicative via the identity
\(\sum_{k=0}^{n-1}\lfloor x+k/n\rfloor=\lfloor nx\rfloor\).)

\begin{enumerate}
\def\labelenumi{(\alph{enumi})}
\tightlist
\item
  \(f(x)=x-\tfrac12\)
\end{enumerate}

\[
\sum_{k=0}^{n-1}\Bigl(x-\tfrac12+\frac{k}{n}\Bigr)
=n\!\left(x-\tfrac12\right)+\frac1n\sum_{k=0}^{n-1}k
=n\!\left(x-\tfrac12\right)+\frac{n-1}{2}
=nx-\tfrac12
=f(nx).
\]

So \(f\) is replicative.

\begin{enumerate}
\def\labelenumi{(\alph{enumi})}
\setcounter{enumi}{1}
\tightlist
\item
  \(f(x)=[x\text{ is an integer}]\) (Iverson bracket)
\end{enumerate}

The left side counts the number of \(k\in\{0,\dots,n-1\}\) such that
\(x+\frac{k}{n}\in\mathbb Z\). This occurs iff \(nx\in\mathbb Z\):
indeed, \(x+\frac{k}{n}=m\in\mathbb Z \iff nx+k=nm\) for some \(m\),
i.e., \(nx\in\mathbb Z\); and when \(nx\in\mathbb Z\) there is exactly
one such \(k\) modulo \(n\). Hence the sum equals
\([nx\in\mathbb Z]=f(nx)\).

\begin{enumerate}
\def\labelenumi{(\alph{enumi})}
\setcounter{enumi}{2}
\tightlist
\item
  \(f(x)=[x\text{ is a positive integer}]\)
\end{enumerate}

As in (b), there is at most one \(k\) giving an integer; its value is
\(m=(nx+k)/n\). That integer is \textbf{positive} iff \(nx>0\) (when
\(nx=0\) we get \(m=0\); when \(nx<0\) we get \(m\le 0\)). Therefore

\[
\sum_{k} [x+k/n\in\mathbb Z_{>0}] \;=\; [nx\in\mathbb Z_{>0}]\;=\;f(nx).
\]

\begin{enumerate}
\def\labelenumi{(\alph{enumi})}
\setcounter{enumi}{3}
\tightlist
\item
  \(f(x)=[\,\exists\,r\in\mathbb Q,\ m\in\mathbb Z:\ x=r\pi+m\,]\)
\end{enumerate}

(\(f\) indicates ``\(x\) is a rational multiple of \(\pi\) up to an
integer.'')

If \(x+\frac{k}{n}=r\pi+m\), multiply by \(n\) to get \(nx+k=nr\pi+nm\),
i.e., \(nx=r'\pi+m'\) with \(r'=nr\in\mathbb Q\), \(m'\in\mathbb Z\).
Thus \(f(nx)=1\).

Conversely, if \(nx=r'\pi+m'\) with \(r'\in\mathbb Q\), choose the
unique \(k\in\{0,\dots,n-1\}\) with \(k\equiv -m'\pmod n\). Then

\[
x+\frac{k}{n}=\frac{nx+k}{n}=\frac{r'\pi+m'+k}{n}=\Bigl(\frac{r'}{n}\Bigr)\pi + \frac{m'+k}{n},
\]

and since \(m'+k\) is divisible by \(n\), this equals \(r\pi+m\) with
\(r=r'/n\in\mathbb Q\), \(m\in\mathbb Z\). Uniqueness of \(k\) again
gives the sum \(= f(nx)\). Hence \(f\) is replicative.

\begin{enumerate}
\def\labelenumi{(\alph{enumi})}
\setcounter{enumi}{4}
\tightlist
\item
  Three variants of (d) with positivity restrictions
\end{enumerate}

The same argument works unchanged for each of the following examples
(any one of which meets the ``three other functions'' request):

\begin{enumerate}
\def\labelenumi{\arabic{enumi}.}
\item
  \(f(x)=[\,\exists\, r\in\mathbb Q_{>0},\,m\in\mathbb Z:\ x=r\pi+m\,]\)
\item
  \(f(x)=[\,\exists\, r\in\mathbb Q,\, m\in\mathbb Z_{>0}:\ x=r\pi+m\,]\)
\item
  \(f(x)=[\,\exists\, r\in\mathbb Q_{>0},\, m\in\mathbb Z_{>0}:\ x=r\pi+m\,]\)
\end{enumerate}

If \(nx=r'\pi+m'\) with the stated positivity on \(r'\) and/or \(m'\),
the construction above picks a \(k\) so that \((m'+k)/n\) is an integer;
since \(m'>0\) implies \((m'+k)/n>0\), and \(r'>0\) implies
\((r'/n)>0\), the corresponding \(r,m\) inherit the required positivity.
Thus each such indicator is replicative.

\begin{enumerate}
\def\labelenumi{(\alph{enumi})}
\setcounter{enumi}{5}
\tightlist
\item
  \(f(x)=\log|2\sin(\pi x)|\) (allowing \(f(x)=-\infty\) at integers)
\end{enumerate}

Use the cyclotomic factorization:

\[
\prod_{k=0}^{n-1}\Bigl|e^{2\pi i x}-e^{2\pi i k/n}\Bigr|
=\left|e^{2\pi i nx}-1\right|.
\]

But \(|e^{2\pi i u}-1|=2|\sin(\pi u)|\), and
\(|e^{2\pi i x}-e^{2\pi i k/n}|=2\bigl|\sin\bigl(\pi(x-k/n)\bigr)\bigr|\).
Since \(\{k/n\}\) runs over the same set as \(\{-k/n\}\),

\[
\prod_{k=0}^{n-1} |2\sin \pi(x+k/n)|
=|2\sin(\pi n x)|.
\]

Taking logs gives
\(\sum_{k=0}^{n-1}\log|2\sin\pi(x+k/n)|=\log|2\sin(\pi n x)|=f(nx)\).
Hence \(f\) is replicative.

\begin{enumerate}
\def\labelenumi{(\alph{enumi})}
\setcounter{enumi}{6}
\tightlist
\item
  Sum of two replicative functions
\end{enumerate}

If \(f,g\) are replicative, then

\[
\sum_{k} (f+g)\!\left(x+\frac{k}{n}\right)
=\sum_{k} f\!\left(x+\frac{k}{n}\right)+\sum_{k} g\!\left(x+\frac{k}{n}\right)
=f(nx)+g(nx)=(f+g)(nx).
\]

So any sum of replicative functions is replicative.

\begin{enumerate}
\def\labelenumi{(\alph{enumi})}
\setcounter{enumi}{7}
\tightlist
\item
  Constant multiple of a replicative function
\end{enumerate}

If \(f\) is replicative and \(c\) is constant,

\[
\sum_{k} (cf)\!\left(x+\frac{k}{n}\right)=c\sum_{k} f\!\left(x+\frac{k}{n}\right)=c\,f(nx)=(cf)(nx).
\]

So \(cf\) is replicative.

\begin{enumerate}
\def\labelenumi{(\roman{enumi})}
\tightlist
\item
  \(g(x)=f(x-\lfloor x\rfloor)\) when \(f\) is replicative
\end{enumerate}

Write \(x=t+m\) with \(m=\lfloor x\rfloor\) and \(t=\{x\}\in[0,1)\). Let
\(r=\lfloor nt\rfloor\) (the number of indices \(k\) for which
\(t+k/n\ge1\)).

Then the multiset

\[
\Bigl\{\,t+\frac{0}{n},\ t+\frac{1}{n},\ \dots,\ t+\frac{n-1}{n}\Bigr\}
\]

has exactly \(r\) values \(\ge1\). If we reduce those \(r\) values by 1
(i.e., take fractional parts), we obtain precisely the re-ordered
arithmetic progression

\[
\Bigl\{\,t-\frac{r}{n}+\frac{j}{n}\ :\ j=0,1,\dots,n-1\Bigr\}.
\]

Therefore

\[
\sum_{k=0}^{n-1} g\!\left(x+\frac{k}{n}\right)
=\sum_{k=0}^{n-1} f\!\left(\bigl\{t+\tfrac{k}{n}\bigr\}\right)
=\sum_{j=0}^{n-1} f\!\left(t-\tfrac{r}{n}+\tfrac{j}{n}\right)
=f\!\left(n\Bigl(t-\tfrac{r}{n}\Bigr)\right)
=f(nt-r).
\]

But \(nt-r = nt-\lfloor nt\rfloor = \{nx\}\), so the right-hand side is

\[
f(\{nx\})=f(nx-\lfloor nx\rfloor)=g(nx),
\]

proving \(g\) is replicative.

\subsection{40 {[}HM46{]}}\label{hm46}

Setup: what ``replicative'' means

A function \(f:\mathbb{R}\to\mathbb{R}\) is called \textbf{replicative}
if for every positive integer \(n\),

\[
f(nx)\;=\;\sum_{k=0}^{n-1} f\!\left(x+\frac{k}{n}\right)\qquad(\forall x\in\mathbb{R}).
\]

Exercise 39 gave many examples and in particular showed
\(f(x)=x-\tfrac12\) is replicative.

Exercise 40 asks you to ``study the class of replicative functions''
and---among other things---decide whether the function in 39(a) is the
only continuous replicative function.

The official answer (a note by N. G. de Bruijn) gives a very slick route
via Fourier series; below I present that argument cleanly.

Step 1: a key telescoping identity

From the replicative law,

\[
\begin{aligned}
f(nx+1)
&=f\!\left(n\!\left(x+\tfrac1n\right)\right)
 =\sum_{k=0}^{n-1} f\!\left(x+\tfrac1n+\tfrac{k}{n}\right)
 =\sum_{k=1}^{n} f\!\left(x+\tfrac{k}{n}\right).
\end{aligned}
\]

Subtracting \(f(nx)=\sum_{k=0}^{n-1} f\!\left(x+\tfrac{k}{n}\right)\)
gives the \textbf{telescoping difference formula}

\[
f(nx+1)-f(nx)\;=\;f(x+1)-f(x)\qquad(\forall n\ge1,\ \forall x\in\mathbb{R}). \tag{★}
\]

(This is the first observation in de Bruijn's note.)

Step 2: continuity forces a constant unit-step

Define \(D(x):=f(x+1)-f(x)\). From (★) we get

\[
D(nx)=D(x)\qquad(\forall n\ge1,\ \forall x).
\]

Fix \(y\). Then \(D(y)=D(y/n)\) for all \(n\). If \(f\) is
\textbf{continuous}, then \(D\) is continuous, hence

\[
D(y)=\lim_{n\to\infty}D(y/n)=D(0),
\]

so \(D\) is constant: there exists \(c\in\mathbb{R}\) with

\[
f(x+1)-f(x)=c\qquad(\forall x\in\mathbb{R}). \tag{1}
\]

(Exactly as stated in de Bruijn's note.)

Step 3: reduce to a 1-periodic replicative function

Let

\[
g(x):=f(x)-c\lfloor x\rfloor .
\]

Then \(g\) is \textbf{1-periodic} by (1):
\(g(x+1)-g(x)=(f(x+1)-f(x))-c(\lfloor x+1\rfloor-\lfloor x\rfloor)=c-c=0\).
Moreover \(g\) is still \textbf{replicative} because both \(f\) and
\(\lfloor x\rfloor\) are replicative (the latter is the content of
Exercise 38), and the defining relation is linear. Hence

\[
g(nx)=\sum_{k=0}^{n-1} g\!\left(x+\frac{k}{n}\right),\qquad g(x+1)=g(x). \tag{2}
\]

(These are exactly the properties de Bruijn isolates.)

Averaging (2) over one period shows \(\int_0^1 g(x)\,dx=0\) (the left
side gives \(n\!\int_0^1g\), the right side also gives \(n\!\int_0^1g\)
after a change of variables, forcing the mean to be \(0\)).

Let \(\widehat g(n):=\int_0^1 e^{2\pi i n x}\,g(x)\,dx\) be the Fourier
coefficients. From (2) one obtains (change variables in the integral and
use periodicity)

\[
\widehat g(n)=\frac1n\,\widehat g(1)\quad\text{for every integer }n\neq0,\qquad \widehat g(0)=0. \tag{3}
\]

This is the integral identity de Bruijn writes down, and it is the
load-bearing step that ties \(g\)'s Fourier coefficients to a \(1/n\)
law.

Step 4: identify the Fourier series

Equations (3) say the spectrum of \(g\) is exactly the spectrum of a
(scaled) \textbf{sawtooth} wave:

\[
g(x)\;=\;a\,\sum_{n\neq0}\frac{e^{2\pi i n x}}{2\pi i n}
\;=\;a\Big(\{x\}-\tfrac12\Big)
\quad\text{for }0<x<1,
\]

for some real constant \(a\) (the value at integers can be chosen to
make \(g\) periodic). In other words,

\[
g(x)=a\big(\{x\}-\tfrac12\big)\quad(0<x<1),
\]

exactly as stated in the book's answer.

Step 5: put \(f\) back together and enforce continuity

Recall \(f(x)=g(x)+c\lfloor x\rfloor\). On \((0,1)\) this gives

\[
f(x)=a\Big(x-\tfrac12\Big)+c\cdot 0=a\Big(x-\tfrac12\Big).
\]

Across an integer, \(g\) jumps by \(-a\) while \(\lfloor x\rfloor\)
jumps by \(+1\), so the jump of \(f\) is \(c-a\). Since \(f\) is
\textbf{continuous}, that jump must vanish; hence \(c=a\). Therefore

\[
f(x)=a\big(\{x\}-\tfrac12\big)+a\lfloor x\rfloor
= a\Big(\lfloor x\rfloor+\{x\}-\tfrac12\Big)
= a\Big(x-\tfrac12\Big)\quad(\forall x).
\]

Conclusion (classification for the continuous case)

The only \textbf{continuous} replicative functions are exactly

\[
\boxed{\,f(x)=A\big(x-\tfrac12\big)\,}\quad\text{with }A\in\mathbb{R}.
\]

This answers the ``for example'' part of Exercise 40: \textbf{yes}---up
to a constant multiple, the function in 39(a) is the only continuous
replicative function.

Remark on the generalized family in the problem text

Exercise 40 also invites you to consider the broader identity

\[
\sum_{k=0}^{n-1} f\!\left(x+\frac{k}{n}\right)=a_n\,f(nx)+b_n\qquad(n\ge1),
\]

whose derivatives (and, if \(b_n=0\), integrals) stay in the same class.
The statement points out that with \(b_n=0\) this includes
\textbf{Bernoulli polynomials} and trigonometric functions like
\(\cot(\pi x)\) and \(\csc^2(\pi x)\), and with \(b_n\neq0\) it includes
the \textbf{psi} function; it also mentions Hurwitz's generalized zeta
\(\zeta(s,x)\) for fixed \(s\). These are standard identities of the
form \(\sum_{k=0}^{n-1}\phi(x+k/n)=a_n\,\phi(nx)+b_n\) and motivate the
term ``replicative.''

If you'd like, I can work out one or two of those generalized examples
in detail (e.g., \(\sum \cot(\pi(x+k/n))=n\,\cot(\pi n x)\)) and show
exactly how the same Fourier-style reasoning fits.

\subsection{41 {[}M24{]}}\label{m24-3}

\begin{enumerate}
\def\labelenumi{\arabic{enumi})}
\tightlist
\item
  Locate the index range where \(a_n=k\)
\end{enumerate}

The first \(k\) blocks contain

\[
1 + 2 + \cdots + k \;=\; \frac{k(k+1)}{2}
\]

terms, i.e.~up to index \(T_k:=\tfrac{k(k+1)}2\) (the \(k\)-th
triangular number). Therefore \(a_n=k\) precisely when \(n\) lies after
the \((k-1)\)-block and at or before the end of the \(k\)-block:

\[
T_{k-1} \;<\; n \;\le\; T_k
\quad\Longleftrightarrow\quad
\frac{k(k-1)}{2} < n \le \frac{k(k+1)}{2}. \; \;  \text{:contentReference[oaicite:0]{index=0}}
\]

\begin{enumerate}
\def\labelenumi{\arabic{enumi})}
\setcounter{enumi}{1}
\tightlist
\item
  Isolate \(k\) from the inequality
\end{enumerate}

Multiply by \(2\):

\[
k^2 - k \;<\; 2n \;\le\; k^2 + k.
\]

Add \(\tfrac14\) throughout (to complete the square):

\[
\Bigl(k-\tfrac12\Bigr)^2 \;<\; 2n + \tfrac14 \;\le\; \Bigl(k+\tfrac12\Bigr)^2.
\]

Take square roots (all quantities are nonnegative):

\[
k - \tfrac12 \;<\; \sqrt{\,2n + \tfrac14\,} \;\le\; k + \tfrac12.
\]

Since \(\sqrt{\,2n+\tfrac14\,} = \sqrt{2n} + \text{(tiny correction)}\)
and \(n\) is an integer, this shows \(\sqrt{2n}\) lies within
\(\tfrac12\) of the integer \(k\). Thus \textbf{\(k\) is the nearest
integer to \(\sqrt{2n}\)}, i.e.

\[
\boxed{\,a_n \;=\; \left\lfloor \sqrt{2n} + \tfrac12 \right\rfloor\, }.
\]

\begin{enumerate}
\def\labelenumi{\arabic{enumi})}
\setcounter{enumi}{2}
\tightlist
\item
  Checks and alternative forms
\end{enumerate}

\begin{itemize}
\item
  Quick check on first terms:

  \begin{itemize}
  \tightlist
  \item
    \(n=1:\ \lfloor \sqrt{2}+1/2\rfloor=\lfloor1.414\ldots+0.5\rfloor=1\).
  \item
    \(n=2,3:\ \lfloor\sqrt{4,6}+1/2\rfloor=2,2\).
  \item
    \(n=4,5,6:\ \lfloor\sqrt{8,10,12}+1/2\rfloor=3,3,3\). Matches
    \(1,\,2,2,\,3,3,3,\ldots\)
  \end{itemize}
\item
  An equivalent characterization from the triangular numbers inequality
  is

  \[
  a_n = k \iff T_{k-1} < n \le T_k
  \quad\text{with }T_k=\frac{k(k+1)}2,
  \]

  which you can also invert with the quadratic formula to get

  \[
  a_n \;=\; \left\lfloor \frac{1+\sqrt{1+8n}}{2}\,\right\rfloor
  \quad\text{or}\quad
  a_n \;=\; \left\lceil \frac{-1+\sqrt{1+8n}}{2}\,\right\rceil,
  \]

  and these agree with \(\lfloor \sqrt{2n}+\tfrac12\rfloor\) for all
  integers \(n\ge1\).
\end{itemize}

\subsection{42 {[}M24{]}}\label{m24-4}

Part (a): the identity

Start from the right-hand side and expand:

\[
\begin{aligned}
n a_n-\sum_{k=1}^{n-1}k(a_{k+1}-a_k)
&= n a_n-\sum_{k=1}^{n-1}k a_{k+1}+\sum_{k=1}^{n-1}k a_k \\[2mm]
&= n a_n-\sum_{j=2}^{n}(j-1)a_j+\sum_{k=1}^{n-1}k a_k
\quad(j=k+1)\\[2mm]
&=\bigl[n a_n-(n-1)a_n\bigr]
+\sum_{k=2}^{n-1}\bigl(- (k-1)+k\bigr)a_k
+\;a_1 \\[2mm]
&= a_n+\sum_{k=2}^{n-1}a_k+a_1
\;=\;\sum_{k=1}^{n}a_k.
\end{aligned}
\]

So the formula holds for all \(n>0\). (This is the discrete ``summation
by parts''/Abel transform for the special case where the other factor is
\(k\).)

Part (b): the floor--log sum

Let \(a_k=\big\lfloor\log_b k\big\rfloor\) and put
\(m=\big\lfloor\log_b n\big\rfloor\). Then \(b^m\le n<b^{m+1}\).

Observe the increments:

\[
\Delta a_k \;=\; a_{k+1}-a_k
=\begin{cases}
1,& \text{if }k=b^j-1\text{ for some }j=1,2,\dots,m,\\
0,& \text{otherwise.}
\end{cases}
\]

Hence

\[
\sum_{k=1}^{n-1} k\,\Delta a_k
=\sum_{j=1}^{m} (b^j-1)
=\Bigl(\sum_{j=1}^{m} b^j\Bigr)-m
=\frac{b^{m+1}-b}{b-1}-m.
\]

Now apply part (a):

\[
\sum_{k=1}^{n}\big\lfloor\log_b k\big\rfloor
= n\,a_n-\sum_{k=1}^{n-1}k\,\Delta a_k
= n m-\Bigl(\frac{b^{m+1}-b}{b-1}-m\Bigr)
= (n+1)m-\frac{b^{m+1}-b}{b-1}.
\]

Recalling \(m=\big\lfloor\log_b n\big\rfloor\), this is exactly

\[
\boxed{\displaystyle
\sum_{k=1}^{n}\big\lfloor\log_b k\big\rfloor
= (n+1)\big\lfloor\log_b n\big\rfloor
-\frac{b^{\lfloor\log_b n\rfloor+1}-b}{b-1}}
\]

as required.

Quick check (sanity): Take \(b=10\), \(n=20\). LHS counts the number of
digits minus 1 for each \(1\le k\le 20\): there are \(9\) with value
\(0\) and \(11\) with value \(1\), so sum \(=11\). RHS gives
\((20+1)\cdot 1-(10^{2}-10)/9=21-10=11\). Works.

\subsection{43 {[}M23{]}}\label{m23-8}

Let \(t=\lfloor\sqrt{n}\rfloor\). For each integer \(m\ge 1\), the
values of \(k\) with \(\lfloor\sqrt{k}\rfloor=m\) are exactly

\[
k=m^2,\ m^2+1,\ \dots,\ (m+1)^2-1.
\]

There are \((m+1)^2-m^2=2m+1\) such \(k\)'s when the whole block fits
inside \([1,n]\).

All full blocks occur for \(m=1,2,\dots,t-1\). The last (possibly
partial) block is \(m=t\), contributing \(n-t^2+1\) terms (from
\(k=t^2\) through \(k=n\)).

Therefore

\[
S(n)=\sum_{m=1}^{t-1} m\,(2m+1)\;+\;t\,(n-t^2+1).
\]

Compute the closed form:

\[
\sum_{m=1}^{t-1} m(2m+1)=
2\sum_{m=1}^{t-1} m^2 + \sum_{m=1}^{t-1} m
= 2\cdot \frac{(t-1)t(2t-1)}{6} + \frac{t(t-1)}{2}.
\]

Hence

\[
S(n)=\frac{(t-1)t(2t-1)}{3}+\frac{t(t-1)}{2}+t(n-t^2+1).
\]

A clean single-line form (still with \(t=\lfloor\sqrt{n}\rfloor\)) is

\[
\boxed{\;
S(n)=\frac{t\bigl(6n-2t^{2}-3t+5\bigr)}{6}
\;}
\]

which is equivalent to

\[
S(n)=\frac{t(t-1)(4t+1)}{6}+t\,(n-t^2+1).
\]

Quick check

For \(n=10\), \(t=\lfloor\sqrt{10}\rfloor=3\):

\[
S(10)=\frac{3(60-18-9+5)}{6}=\frac{3\cdot38}{6}=19,
\]

matching \(1,1,1,2,2,2,2,2,3,3\) summed \(=19\).

Asymptotics (nice to know)

Since \(t=\lfloor\sqrt{n}\rfloor\sim \sqrt{n}\),

\[
S(n)=\frac{2}{3}\,n^{3/2}+O(n),
\]

which you can also see directly by approximating the area under
\(y=\sqrt{x}\).

\subsection{44 {[}M24{]}}\label{m24-5}

Step 1: A handy floor-sum identity

For any integers \(a\) and \(b\ge1\),

\[
\sum_{j=0}^{b-1}\Big\lfloor\frac{a+j}{b}\Big\rfloor \;=\; a.
\]

\textbf{Proof.} Write \(a=qb+r\) with \(0\le r<b\). Then

\[
\Big\lfloor\frac{a+j}{b}\Big\rfloor
= \Big\lfloor q+\frac{r+j}{b}\Big\rfloor
= \begin{cases}
q,& j=0,1,\dots,b-r-1,\\[2pt]
q+1,& j=b-r,\dots,b-1.
\end{cases}
\]

Summing gives \(q(b-r)+(q+1)r=qb+r=a\). □

Step 2: Evaluate the inner sum for a fixed \(k\)

Fix \(k\ge0\). Perform Euclidean division of \(n\) by \(b^k\):

\[
n \;=\; a_k\,b^k + r_k,\qquad a_k=\Big\lfloor\frac{n}{b^k}\Big\rfloor,\; 0\le r_k<b^k.
\]

Then for any \(j\),

\[
\frac{n+j\,b^k}{b^{k+1}}
= \frac{a_k b^k + r_k + j b^k}{b^{k+1}}
= \frac{a_k+j}{b} + \frac{r_k}{b^{k+1}}.
\]

Note that \(0\le \dfrac{r_k}{b^{k+1}}<\dfrac1b\). Therefore, when we sum
over a full block \(j=0,1,\dots,b-1\), the tiny offset
\(\dfrac{r_k}{b^{k+1}}\) never causes an extra carry; hence

\[
\sum_{j=0}^{b-1}\Big\lfloor\frac{n+j\,b^k}{b^{k+1}}\Big\rfloor
= \sum_{j=0}^{b-1}\Big\lfloor\frac{a_k+j}{b}+\frac{r_k}{b^{k+1}}\Big\rfloor
= \sum_{j=0}^{b-1}\Big\lfloor\frac{a_k+j}{b}\Big\rfloor
= a_k
= \Big\lfloor\frac{n}{b^k}\Big\rfloor,
\]

using Step 1 in the last equality.

Since the exercise's inner sum runs over \(j=1,\dots,b-1\), just
subtract the \(j=0\) term:

\[
\sum_{j=1}^{b-1}\Big\lfloor\frac{n+j\,b^k}{b^{k+1}}\Big\rfloor
= \Big\lfloor\frac{n}{b^k}\Big\rfloor \;-\; \Big\lfloor\frac{n}{b^{k+1}}\Big\rfloor.
\tag{★}
\]

Step 3: Telescope over \(k\)

Now sum (★) over \(k\ge0\):

\[
\sum_{k\ge0}\sum_{j=1}^{b-1}\Big\lfloor\frac{n+j\,b^k}{b^{k+1}}\Big\rfloor
= \sum_{k\ge0}\Big(\Big\lfloor\frac{n}{b^k}\Big\rfloor - \Big\lfloor\frac{n}{b^{k+1}}\Big\rfloor\Big)
= \Big\lfloor\frac{n}{b^0}\Big\rfloor - \lim_{K\to\infty}\Big\lfloor\frac{n}{b^{K+1}}\Big\rfloor.
\]

Since \(n\) is an integer, \(\lfloor n/b^0\rfloor=n\). For the tail
term:

\begin{itemize}
\tightlist
\item
  If \(n\ge0\), then \(0\le n/b^{K+1}<1\) for large \(K\), so
  \(\lfloor n/b^{K+1}\rfloor=0\).
\item
  If \(n<0\), then \(-1<n/b^{K+1}<0\) for large \(K\), so
  \(\lfloor n/b^{K+1}\rfloor=-1\).
\end{itemize}

Thus

\[
\sum_{k\ge0}\sum_{1\le j<b}\Big\lfloor\frac{n+j\,b^k}{b^{k+1}}\Big\rfloor
=\begin{cases}
n,& n\ge0,\\[4pt]
n+1,& n<0.
\end{cases}
\]

That proves the required identity for \(n\ge0\), and answers the second
question: for \(n<0\) the value is \(n+1\).

Intuition (base-\(b\) view)

Write \(n=\sum_{k\ge0} d_k b^k\) with digits \(0\le d_k<b\) (for
\(n\ge0\)). For fixed \(k\), adding \(j b^k\) with \(1\le j<b\) shifts
by \(j\) in the \(k\)-th digit and contributes
\(\big\lfloor (n + j b^k)/b^{k+1}\big\rfloor\) to the \((k+1)\)-st and
higher digits---i.e., the number of carries past position \(k\). As
\(j\) ranges from \(1\) to \(b-1\), you get exactly \(d_k\) carries.
Summing over \(k\) counts all digit contributions and reconstructs
\(n\). For negative \(n\), the infinite ``tail'' of floors contributes
an extra \(+1\), giving \(n+1\).

\subsection{45 {[}M28{]}}\label{m28}

Proof of the general identity (1)

Define, for each \(r\in\{0,1,\dots,m-1\}\),

\[
J_r \;=\; \Bigl\{\, j\in\{0,1,\dots,n-1\} \;\Big|\;
\Big\lfloor \frac{mj}{n}\Big\rfloor = r \Bigr\}.
\]

Then

\[
\sum_{0\le j<n} f\!\left(\Big\lfloor \frac{mj}{n}\Big\rfloor\right)
\;=\;\sum_{r=0}^{m-1} f(r)\,|J_r|.
\]

Characterize \(|J_r|\). The condition \(\lfloor mj/n\rfloor=r\) is
equivalent to

\[
r \;\le\; \frac{mj}{n} \;<\; r+1
\quad\Longleftrightarrow\quad
\frac{nr}{m} \;\le\; j \;<\; \frac{n(r+1)}{m}.
\]

Among integers \(j\in[0,n-1]\), that count is

\[
|J_r| \;=\; \Big\lceil \frac{n(r+1)}{m}\Big\rceil \;-\; \Big\lceil \frac{nr}{m}\Big\rceil
\;=\; A_{r+1}-A_r,
\quad\text{where } A_r:=\Big\lceil \frac{nr}{m}\Big\rceil,\; A_0=0,\; A_m=n.
\]

Hence

\[
\sum_{0\le j<n}\! f\!\left(\Big\lfloor \frac{mj}{n}\Big\rfloor\right)
=\sum_{r=0}^{m-1} f(r)\,(A_{r+1}-A_r)
=\sum_{r=0}^{m-1}\!\bigl(f(r)A_{r+1}-f(r)A_r\bigr).
\]

Now apply \textbf{summation by parts} (a discrete telescoping trick).
Shift index in the first term:

\[
\sum_{r=0}^{m-1} f(r)A_{r+1}
=\sum_{r=1}^{m} f(r-1)A_r
= f(m-1)A_m \;+\; \sum_{r=1}^{m-1} f(r-1)A_r
= n\,f(m-1) \;+\; \sum_{r=1}^{m-1} f(r-1)A_r .
\]

Therefore

\[
\sum_{r=0}^{m-1} f(r)\,(A_{r+1}-A_r)
= n\,f(m-1) + \sum_{r=1}^{m-1} \bigl(f(r-1)-f(r)\bigr)A_r
= n\,f(m-1)+\sum_{r=0}^{m-1} \Big\lceil \frac{nr}{m}\Big\rceil\bigl(f(r-1)-f(r)\bigr),
\]

since the \(r=0\) term has coefficient \(A_0=\lceil 0\rceil = 0\). This
is exactly (1).

``In particular'' binomial identity (2)

Take

\[
f(r)=\binom{r+1}{k}\qquad (k\text{ fixed}).
\]

Then by Pascal's rule,

\[
f(r-1)-f(r)=\binom{r}{k}-\binom{r+1}{k}=-\binom{r}{k-1}.
\]

Plugging into (1) gives

\[
\sum_{0\le j<n} \binom{\big\lfloor mj/n\big\rfloor+1}{k}
\;=\;
-\sum_{0\le r<m} \Big\lceil \frac{rn}{m}\Big\rceil \binom{r}{k-1}
\;+\; n\binom{m}{k}.
\]

Move the sum to the left:

\[
\sum_{0\le j<n} \binom{\big\lfloor mj/n\big\rfloor+1}{k}
\;+\;
\sum_{0\le r<m} \Big\lceil \frac{rn}{m}\Big\rceil \binom{r}{k-1}
\;=\; n\binom{m}{k},
\]

which is exactly the required identity (2). 

Remark on the ``reciprocity'' for sums of floors

Exercise 37 evaluates
\(\sum_{0\le k<n}\big\lfloor \frac{mk+x}{n}\big\rfloor\) in a form
symmetric in \(m\) and \(n\) (it involves \(\gcd(m,n)\) and
\(\lfloor x/\gcd(m,n)\rfloor\)). Hence swapping \(m\) and \(n\) gives
the same value, implying

\[
\sum_{0\le k<n}\Big\lfloor \frac{mk+x}{n}\Big\rfloor
=
\sum_{0\le k<m}\Big\lfloor \frac{nk+x}{m}\Big\rfloor,
\]

as noted right before (1).

\subsection{46 {[}M29{]}}\label{m29}

Let

\[
a_j \stackrel{\text{def}}{=} \Big\lfloor \frac{m j}{n}\Big\rfloor,\qquad
J \stackrel{\text{def}}{=} \{\,j\in\mathbb Z: 0\le j<\alpha n\,\}
= \{0,1,\dots,\lceil \alpha n\rceil-1\}.
\]

Note that \(a_j\) is nondecreasing in \(j\) and, since \(j<\alpha n\),

\[
a_j \le \Big\lfloor \frac{m(\alpha n)}{n}\Big\rfloor < \alpha m
\quad\Rightarrow\quad
a_j\in\{0,1,\dots, \lceil\alpha m\rceil-1\}.
\]

We use a top-down telescoping (a discrete summation-by-parts trick). Fix

\[
L \stackrel{\text{def}}{=} \lceil \alpha m\rceil.
\]

For any integer \(a\le L-1\),

\[
f(a) \;=\; f(L-1) + \sum_{r=a+1}^{L-1}\bigl(f(r-1)-f(r)\bigr).
\]

Apply this with \(a=a_j\) and sum over all \(j\in J\):

\[
\sum_{j\in J} f(a_j)
= \sum_{j\in J} f(L-1) \;+\; \sum_{j\in J}\sum_{r=a_j+1}^{L-1}\!\bigl(f(r-1)-f(r)\bigr).
\]

Swap the order of summation in the double sum:

\[
\sum_{j\in J} f(a_j)
= |J|\cdot f(L-1) \;+\; \sum_{r=1}^{L-1}\Bigl(\#\{\,j\in J: a_j\le r-1\,\}\Bigr)\,\bigl(f(r-1)-f(r)\bigr).
\]

Now count \(\#\{j\in J: a_j\le r-1\}\).

\[
a_j\le r-1
\iff
\Big\lfloor \frac{m j}{n}\Big\rfloor \le r-1
\iff
\frac{m j}{n} < r
\iff
j < \frac{r n}{m}.
\]

Among \(J=\{0,1,\dots,\lceil \alpha n\rceil-1\}\), the count of integers
\(j\) with \(j<\frac{r n}{m}\) is

\[
\#\{\,j\in J: a_j\le r-1\,\} = \Big\lceil \frac{r n}{m}\Big\rceil,
\]

because \(r\le L-1=\lceil\alpha m\rceil-1\) guarantees
\(\big\lceil \frac{r n}{m}\big\rceil \le \lceil \alpha n\rceil\), so the
upper truncation by \(J\) does not bind.

Also, \(|J|=\lceil \alpha n\rceil\) and \(L-1=\lceil \alpha m\rceil-1\).
Therefore

\[
\sum_{0\le j<\alpha n} f\!\left(\Big\lfloor \frac{m j}{n}\Big\rfloor\right)
=
\lceil \alpha n\rceil\, f\!\bigl(\lceil\alpha m\rceil-1\bigr)
+
\sum_{r=1}^{\lceil\alpha m\rceil-1}
\Big\lceil \frac{r n}{m}\Big\rceil\,
\bigl(f(r-1)-f(r)\bigr),
\]

which is exactly the desired identity (the \(r=0\) term is \(0\) since
\(\lceil 0\rceil\cdot(f(-1)-f(0))=0\)).

Sanity check (small numbers)

Take \(m=3\), \(n=5\), \(\alpha=1.4\), and \(f(r)=r\).

Left side:

\[
\Big\lfloor \tfrac{3j}{5}\Big\rfloor\ \text{for } j=0,\dots,6
\;=\; 0,0,1,1,2,3,3,\quad\text{sum }=10.
\]

Right side: \(\lceil\alpha m\rceil=5,\ \lceil\alpha n\rceil=7\).

\[
\sum_{r=0}^{4}\Big\lceil \tfrac{5r}{3}\Big\rceil\,(f(r-1)-f(r))
\;+\; 7\,f(4)
=
-(0+2+4+5+7)+7\cdot 4
= -18+28
=10.
\]

Matches perfectly.

\subsection{47 {[}M31{]}}\label{m31}

\begin{enumerate}
\def\labelenumi{(\alph{enumi})}
\tightlist
\item
  Gauss's lemma in this setup
\end{enumerate}

Fix \(k\) with \(1\le k\le (p-1)/2\). Write

\[
2kq \;=\; p\left\lfloor \frac{2kq}{p}\right\rfloor \;+\; r_k,\qquad
0< r_k < p,\quad r_k\equiv 2kq \pmod p.
\]

Define

\[
a_k \;:=\; (-1)^{\lfloor 2kq/p\rfloor}\,r_k.
\]

There are two possibilities: If \(\lfloor 2kq/p\rfloor\) is even, then
\(a_k=r_k\in\{1,2,\dots,p-1\}\); if it's odd, then
\(a_k=r_k-p\in\{-p+1,\dots,-1\}\). In either case, \(a_k\) is congruent
modulo \(p\) to an \textbf{even} residue, because \(r_k\equiv 2kq\) and
\(2kq\) is even. As \(k\) ranges over \(1,2,\dots,(p-1)/2\), the reduced
residues \(2k\) run through all even residues modulo \(p\), hence the
values \(2kq\bmod p\) also run through each even residue exactly once
(since multiplication by coprime \(q\) permutes residues). The extra
factor \((-1)^{\lfloor 2kq/p\rfloor}\) simply chooses, for each even
residue \(e\), either \(e\) or \(e-p\), so the set \(\{a_k\}\) is a
permutation of \(\{2,4,\dots,p-1\}\) (as integers, hence also modulo
\(p\)).

Taking products of the \(a_k\) and comparing both sides:

\[
\prod_{k=1}^{(p-1)/2}\! a_k
\;\equiv\; 2\cdot 4\cdots (p-1)\pmod p.
\]

But

\[
\prod_{k=1}^{(p-1)/2}\! a_k
= (-1)^{\sum_{k}\lfloor 2kq/p\rfloor}\,\prod_{k=1}^{(p-1)/2}\!(2kq\bmod p)
\equiv (-1)^{\sigma}\,\prod_{k=1}^{(p-1)/2}\! (2kq)
= (-1)^{\sigma} q^{(p-1)/2}\!\!\!\prod_{k=1}^{(p-1)/2}\! (2k)\pmod p.
\]

Cancel the common unit \(\prod_{k}(2k)\) from both sides to obtain

\[
q^{(p-1)/2}\equiv (-1)^{\sigma}\pmod p.
\]

By the definition of Legendre's symbol this is exactly

\[
\boxed{\ \left(\frac{q}{p}\right)=(-1)^{\sigma},\qquad \sigma=\sum_{k=1}^{(p-1)/2}\left\lfloor\frac{2kq}{p}\right\rfloor\ }.
\]

\begin{enumerate}
\def\labelenumi{(\alph{enumi})}
\setcounter{enumi}{1}
\tightlist
\item
  The supplementary law for \(\bigl(\tfrac{2}{p}\bigr)\)
\end{enumerate}

Set \(q=2\). Then

\[
\sigma=\sum_{k=1}^{(p-1)/2}\left\lfloor\frac{4k}{p}\right\rfloor.
\]

Write \(p=4n+1\) or \(p=4n+3\) (since \(p\) is odd).

\begin{itemize}
\item
  If \(p=4n+1\), then \(\dfrac{4k}{p}<1\) for \(1\le k\le n\), and
  \(1\le \dfrac{4k}{p}<2\) for \(n<k\le 2n\). Hence

  \[
  \sigma=\underbrace{0+\cdots+0}_{n\text{ terms}}+\underbrace{1+\cdots+1}_{n\text{ terms}}=n.
  \]
\item
  If \(p=4n+3\), then \((p-1)/2=2n+1\) and \(\dfrac{4k}{p}<1\) for
  \(1\le k\le n\), while \(1\le \dfrac{4k}{p}<2\) for \(n<k\le 2n+1\).
  Thus

  \[
  \sigma=\underbrace{0+\cdots+0}_{n\text{ terms}}+\underbrace{1+\cdots+1}_{n+1\text{ terms}}=n+1.
  \]
\end{itemize}

Therefore

\[
\left(\frac{2}{p}\right)=(-1)^{\sigma}=
\begin{cases}
(-1)^n,& p=4n+1,\\[2pt]
(-1)^{\,n+1},& p=4n+3.
\end{cases}
\]

Classifying by \(p\bmod 8\) (compute \(n\) mod 2 carefully):

\begin{itemize}
\tightlist
\item
  \(p=8t+1\Rightarrow n=2t\) even \(\Rightarrow (\tfrac{2}{p})=+1\).
\item
  \(p=8t+3\Rightarrow n=2t\) even \(\Rightarrow (\tfrac{2}{p})=-1\).
\item
  \(p=8t+5\Rightarrow n=2t+1\) odd \(\Rightarrow (\tfrac{2}{p})=-1\).
\item
  \(p=8t+7\Rightarrow n=2t+1\) odd \(\Rightarrow (\tfrac{2}{p})=+1\).
\end{itemize}

Thus

\[
\boxed{\ \left(\frac{2}{p}\right)=
\begin{cases}
+1,& p\equiv 1,7\pmod 8,\\
-1,& p\equiv 3,5\pmod 8.
\end{cases}}
\]

\begin{enumerate}
\def\labelenumi{(\alph{enumi})}
\setcounter{enumi}{2}
\tightlist
\item
  Parity reduction: replace doubles by singles
\end{enumerate}

Assume \(q\) is odd and \(p\nmid q\). For \(k<p/4\), consider

\[
\left\lfloor\frac{(p-1-2k)q}{p}\right\rfloor
= q - \left\lceil\frac{(2k+1)q}{p}\right\rceil.
\]

Because \(p\nmid (2k+1)q\) (both \(2k+1\) and \(q\) are odd and in the
range \(1,\dots,p-1\)), we have
\(\bigl\lceil x\bigr\rceil=\bigl\lfloor x\bigr\rfloor+1\) at
\(x=\frac{(2k+1)q}{p}\). Hence

\[
\left\lfloor\frac{(p-1-2k)q}{p}\right\rfloor
= q - 1 - \left\lfloor\frac{(2k+1)q}{p}\right\rfloor
\equiv \left\lfloor\frac{(2k+1)q}{p}\right\rfloor \pmod 2,
\]

since \(q-1\) is even. Summing these congruences over
\(k=0,1,\dots,\lfloor(p-2)/4\rfloor\) shows that, modulo 2, the ``tail''
terms \(\lfloor((p-1)q)/p\rfloor,\lfloor((p-3)q)/p\rfloor,\dots\) can be
replaced by \(\lfloor q/p\rfloor,\lfloor 3q/p\rfloor,\dots\). In total
this yields

\[
\sum_{0\le k<p/2}\left\lfloor\frac{2kq}{p}\right\rfloor
\equiv
\sum_{0\le k<p/2}\left\lfloor\frac{kq}{p}\right\rfloor
\pmod 2,
\]

as required.

\begin{enumerate}
\def\labelenumi{(\alph{enumi})}
\setcounter{enumi}{3}
\tightlist
\item
  Quadratic reciprocity via Eisenstein's trick (using Ex. 46)
\end{enumerate}

From Ex. 46 (with \(m=q\), \(n=p\), \(\alpha=\tfrac12\), and \(f(r)=r\))
we obtain the identity

\[
\sum_{0\le k<p/2}\left\lfloor\frac{kq}{p}\right\rfloor
+\sum_{0\le r<q/2}\left\lceil\frac{rp}{q}\right\rceil
=\left\lceil\frac p2\right\rceil\Bigl(\left\lceil\frac q2\right\rceil-1\Bigr).
\]

(Cited from the page with Ex. 46.) 

Because \(p,q\) are distinct odd primes, none of the fractions
\(\frac{rp}{q}\) with \(1\le r\le (q-1)/2\) is an integer, so
\(\lceil x\rceil=\lfloor x\rfloor+1\) applies termwise:

\[
\sum_{0\le r<q/2}\left\lceil\frac{rp}{q}\right\rceil
=\sum_{0\le r<q/2}\left\lfloor\frac{rp}{q}\right\rfloor + \frac{q-1}{2}.
\]

Also \(\left\lceil\frac p2\right\rceil=\frac{p+1}{2}\) and
\(\left\lceil\frac q2\right\rceil=\frac{q+1}{2}\). Hence

\[
\sum_{k=1}^{(p-1)/2}\left\lfloor\frac{kq}{p}\right\rfloor
+\sum_{r=1}^{(q-1)/2}\left\lfloor\frac{rp}{q}\right\rfloor
=\frac{(p-1)(q-1)}{4}.
\]

By part (a) together with the parity reduction of part (c),

\[
\left(\frac{q}{p}\right)=(-1)^{\sum_{k=1}^{(p-1)/2}\lfloor kq/p\rfloor},\qquad
\left(\frac{p}{q}\right)=(-1)^{\sum_{r=1}^{(q-1)/2}\lfloor rp/q\rfloor}.
\]

Adding the exponents mod 2 gives

\[
\sum_{k=1}^{(p-1)/2}\left\lfloor\frac{kq}{p}\right\rfloor
+\sum_{r=1}^{(q-1)/2}\left\lfloor\frac{rp}{q}\right\rfloor
\equiv \frac{(p-1)(q-1)}{4}\pmod 2,
\]

hence

\[
\boxed{\ \left(\frac{q}{p}\right)\left(\frac{p}{q}\right)=(-1)^{(p-1)(q-1)/4}\ },
\]

which is the quadratic reciprocity law. (The exercise statement on the
next page explicitly asks you to reach this formula.)

\bookmarksetup{startatroot}

\chapter{1.2.5 Permutations and
factorials}\label{permutations-and-factorials}

\subsection{1 {[}00{]}}\label{section-66}

Think of building the deck from top to bottom:

\begin{itemize}
\tightlist
\item
  For the first position you can choose any of the 52 cards.
\item
  For the second position you can choose any of the remaining 51 cards.
\item
  \ldots{}
\item
  For the last position only 1 card remains.
\end{itemize}

By the rule of product, the total number of permutations is

\[
52 \times 51 \times 50 \times \cdots \times 2 \times 1 \;=\; 52!\,.
\]

That is the exact count.

(Optional) Sanity checks \& useful facts

\begin{itemize}
\item
  \textbf{Exact integer (computed):}

  \[
  52! = 80658175170943878571660636856403766975289505440883277824000000000000.
  \]
\item
  \textbf{Number of digits:} \(52!\) has 68 decimal digits.
\item
  \textbf{Trailing zeros:} The number of trailing zeros equals the
  number of factors of 5 in \(52!\):

  \[
  \left\lfloor \frac{52}{5} \right\rfloor + \left\lfloor \frac{52}{25} \right\rfloor
  = 10 + 2 = 12,
  \]

  so \(52!\) ends with 12 zeros.
\item
  \textbf{Magnitude (Stirling's approximation):}

  \[
  52! \approx \sqrt{2\pi\cdot 52}\,\Big(\frac{52}{e}\Big)^{52}
  \approx 8.0658\times 10^{67},
  \]

  which agrees with the leading digits shown above.
\end{itemize}

\subsection{2 {[}10{]}}\label{section-67}

Claim (the recurrence)

\[
p_n^k \;=\; p_n^{\,k-1}\,(n-k+1).
\]

Combinatorial proof

Build a \(k\)-long arrangement left to right.

\begin{itemize}
\tightlist
\item
  First fill the first \(k-1\) positions in any valid way; there are
  \(p_n^{\,k-1}\) such partial arrangements by definition.
\item
  After placing \(k-1\) \emph{distinct} objects, exactly
  \(n-(k-1)=n-k+1\) objects remain unused, and any of them can go to the
  last position.
\end{itemize}

By the multiplication principle,

\[
p_n^k = \bigl(\text{\# ways to fill first \(k-1\) spots}\bigr)\times \bigl(\text{\# choices for last spot}\bigr)
= p_n^{\,k-1}(n-k+1),
\]

as claimed.

Inductive derivation of the closed form

Using the recurrence repeatedly,

\[
p_n^k
= (n-k+1)\,p_n^{\,k-1}
= (n-k+1)(n-k+2)\,p_n^{\,k-2}
= \cdots
= n(n-1)\cdots(n-k+1).
\]

Equivalently,

\[
p_n^k = \frac{n!}{(n-k)!}.
\]

Special case \(k=n\)

When arranging \textbf{all} \(n\) objects, the recurrence gives

\[
p_n^n = p_n^{\,n-1}(n-n+1)=p_n^{\,n-1}\cdot 1,
\]

which reflects the simple fact: \emph{once the first \(n-1\) places are
fixed, the last object is forced (only one possibility remains).}

\subsection{3 {[}10{]}}\label{section-68}

Method 1 (insert \(n\) in every position)

\textbf{Rule.} Given a permutation \(a_1a_2\cdots a_{n-1}\) of
\(\{1,\dots,n-1\}\), insert \(n\) in each of the \(n\) possible
positions. This yields \(n\) permutations of \(\{1,\dots,n\}\). 

Apply to \(3\,1\,2\,4\) with \(n=5\) (five insertion positions):

\begin{enumerate}
\def\labelenumi{\arabic{enumi}.}
\tightlist
\item
  before 3: \textbf{5 3 1 2 4}
\item
  between 3 and 1: \textbf{3 5 1 2 4}
\item
  between 1 and 2: \textbf{3 1 5 2 4}
\item
  between 2 and 4: \textbf{3 1 2 5 4}
\item
  after 4: \textbf{3 1 2 4 5}
\end{enumerate}

So Method 1 produces: \textbf{5 3 1 2 4, 3 5 1 2 4, 3 1 5 2 4, 3 1 2 5
4, 3 1 2 4 5.}

Method 2 (order-preserving relabel + append \(k\))

\textbf{Rule (equivalent forms).} For each \(k\in\{1,\dots,n\}\), take a
permutation \(a_1\cdots a_{n-1}\) of \(\{1,\dots,n-1\}\) and:

\begin{itemize}
\tightlist
\item
  increase by 1 every entry \(a_j\) with \(a_j\ge k\); this relabels
  \(\{1,\dots,n-1\}\) into \(\{1,\dots,k-1,k+1,\dots,n\}\) \textbf{while
  preserving order},
\item
  then append \(k\). This yields \(n\) permutations from each starting
  permutation. 
\end{itemize}

Apply to \(3\,1\,2\,4\) (with \(n=5\)):

\begin{itemize}
\tightlist
\item
  \(k=1\): add 1 to all (since all \(\ge1\)) → \(4\,2\,3\,5\); append 1
  → \textbf{4 2 3 5 1}
\item
  \(k=2\): add 1 to entries \(\ge2\) → \(4\,1\,3\,5\); append 2 →
  \textbf{4 1 3 5 2}
\item
  \(k=3\): add 1 to entries \(\ge3\) → \(4\,1\,2\,5\); append 3 →
  \textbf{4 1 2 5 3}
\item
  \(k=4\): add 1 to entries \(\ge4\) → \(3\,1\,2\,5\); append 4 →
  \textbf{3 1 2 5 4}
\item
  \(k=5\): no change → \(3\,1\,2\,4\); append 5 → \textbf{3 1 2 4 5}
\end{itemize}

So Method 2 produces: \textbf{4 2 3 5 1, 4 1 3 5 2, 4 1 2 5 3, 3 1 2 5
4, 3 1 2 4 5.}

\subsection{4 {[}13{]}}\label{section-69}

\begin{enumerate}
\def\labelenumi{\arabic{enumi})}
\tightlist
\item
  Number of decimal digits of \(1000!\)
\end{enumerate}

For any positive integer \(N\), the number of base-10 digits is

\[
\big\lfloor \log_{10} N \big\rfloor + 1.
\]

With \(N=1000!\),

\[
\log_{10}(1000!)=\sum_{k=1}^{1000}\log_{10}k.
\]

Evaluating (e.g., by high-precision arithmetic or tight Stirling bounds)
gives

\[
\log_{10}(1000!) \approx 2567.604644\ldots
\]

so

\[
\text{digits} = \lfloor 2567.604644\ldots \rfloor + 1 = 2568.
\]

(You can also get this from Stirling's formula with error bounds, which
nails the integer part in this case; the book develops Stirling earlier
in §1.2.5.)

\begin{enumerate}
\def\labelenumi{\arabic{enumi})}
\setcounter{enumi}{1}
\tightlist
\item
  Leading digit of \(1000!\)
\end{enumerate}

Write

\[
\log_{10}(1000!) = I + f, \quad I=\big\lfloor \log_{10}(1000!) \big\rfloor,\; f\in[0,1).
\]

Then

\[
1000! = 10^{I+f} = 10^{I}\cdot 10^{f},\quad 10^{f}\in[1,10).
\]

The leading digit is \(\big\lfloor 10^{f}\big\rfloor\). From above
\(f\approx 0.604644\), so \(10^{f}\approx 4.0238726\), hence the first
digit is \(4\).

(Heuristically, this matches the book's remark that the mantissa is a
bit above \(\log_{10}4\approx 0.60206\), therefore the leading digit is
\(4\).)

\begin{enumerate}
\def\labelenumi{\arabic{enumi})}
\setcounter{enumi}{2}
\tightlist
\item
  Number of trailing zeros of \(1000!\)
\end{enumerate}

Trailing zeros in base 10 are contributed by factors of \(10=2\cdot 5\).
In \(n!\), twos are plentiful, so the count is the exponent of \(5\) in
\(n!\):

\[
\nu_{5}(n!)=\sum_{k\ge 1}\left\lfloor \frac{n}{5^{k}}\right\rfloor.
\]

(This is the standard prime--power valuation formula developed as Eq.
(8) in §1.2.5. )

Compute for \(n=1000\):

\[
\left\lfloor \frac{1000}{5}\right\rfloor=200,\quad
\left\lfloor \frac{1000}{25}\right\rfloor=40,\quad
\left\lfloor \frac{1000}{125}\right\rfloor=8,\quad
\left\lfloor \frac{1000}{625}\right\rfloor=1,
\]

and higher powers give \(0\). Summing:

\[
\nu_{5}(1000!)=200+40+8+1=249.
\]

Therefore \(1000!\) ends with exactly \(249\) zeros.

\subsection{5 {[}15{]}}\label{section-70}

We're asked to estimate \(8!\) using the refined Stirling approximation

\[
n! \;\approx\; \sqrt{2\pi n}\,\Bigl(\frac{n}{e}\Bigr)^{n}\!\Bigl(1+\frac{1}{12n}\Bigr),
\]

Step 1: Base Stirling term

For \(n=8\),

\[
\sqrt{2\pi n}\Bigl(\frac{n}{e}\Bigr)^n
= \sqrt{16\pi}\,\Bigl(\frac{8}{e}\Bigr)^8
\approx 39{,}902.395\;.
\]

(Indeed, Knuth's numerical example rounds this to \(39902\).)

Step 2: First correction factor

The first correction multiplies by
\(1+\dfrac{1}{12n}=1+\dfrac{1}{96}=\dfrac{97}{96}\):

\[
\widehat{8!}
\;=\;
\Bigl[\sqrt{2\pi n}\Bigl(\frac{n}{e}\Bigr)^n\Bigr]
\Bigl(1+\frac{1}{12n}\Bigr)
\;\approx\; 39{,}902.395\times\frac{97}{96}
\;\approx\; 40{,}318.045\;.
\]

Rounding sensibly gives \(\boxed{40{,}318}\).

Accuracy check

The true value is \(8!=40{,}320\). Thus the error is about \(-1.955\)
(roughly \(-2\)), i.e.,

\[
\frac{|\widehat{8!}-8!|}{8!}\;\approx\;4.85\times 10^{-5},
\]

a dramatic improvement over the \(\sim1\%\) error of the uncorrected
term.

So the requested refined Stirling estimate for \(8!\) is \(40{,}318\),
within two of the exact value.

\subsection{6 {[}17{]}}\label{section-71}

For a prime \(p\), the exponent \(\mu_p\) of \(p\) in \(n!\) is

\[
\mu_p=\sum_{k>0}\Big\lfloor \frac{n}{p^k}\Big\rfloor,
\]

i.e., count how many multiples of \(p,p^2,p^3,\dots\) are \(\le n\) and
add them. (Eq. (8). )

We apply this with \(n=20\) and primes \(p\le 20\):
\(2,3,5,7,11,13,17,19\).

\begin{itemize}
\tightlist
\item
  \(p=2:\ \lfloor20/2\rfloor+\lfloor20/4\rfloor+\lfloor20/8\rfloor+\lfloor20/16\rfloor=10+5+2+1=18.\)
\item
  \(p=3:\ \lfloor20/3\rfloor+\lfloor20/9\rfloor=6+2=8.\)
\item
  \(p=5:\ \lfloor20/5\rfloor=4.\)
\item
  \(p=7:\ \lfloor20/7\rfloor=2.\)
\item
  \(p=11:\ \lfloor20/11\rfloor=1.\)
\item
  \(p=13:\ \lfloor20/13\rfloor=1.\)
\item
  \(p=17:\ \lfloor20/17\rfloor=1.\)
\item
  \(p=19:\ \lfloor20/19\rfloor=1.\)
\end{itemize}

All higher powers contribute \(0\).

Therefore the prime factorization is

\[
\boxed{20! = 2^{18}\,3^{8}\,5^{4}\,7^{2}\,11\,13\,17\,19.}
\]

This is exactly what the exercise asks for, using Eq. (8) from §1.2.5.

\subsection{7 {[}M10{]}}\label{m10-6}

\begin{enumerate}
\def\labelenumi{\arabic{enumi}.}
\tightlist
\item
  Convert the given datum to gamma notation. Since
  \(\left(\tfrac12\right)! = \Gamma\!\left(\tfrac32\right)\), the
  assumption gives
\end{enumerate}

\[
\Gamma\!\left(\tfrac32\right) \;=\; \frac{\sqrt{\pi}}{2}.
\]

\begin{enumerate}
\def\labelenumi{\arabic{enumi}.}
\setcounter{enumi}{1}
\tightlist
\item
  Use the recurrence \(\Gamma(x+1)=x\,\Gamma(x)\) with \(x=\tfrac12\):
\end{enumerate}

\[
\Gamma\!\left(\tfrac32\right) \;=\; \tfrac12\,\Gamma\!\left(\tfrac12\right).
\]

Hence

\[
\Gamma\!\left(\tfrac12\right) \;=\; 2\,\Gamma\!\left(\tfrac32\right) \;=\; 2\cdot\frac{\sqrt{\pi}}{2} \;=\; \boxed{\sqrt{\pi}}.
\]

\begin{enumerate}
\def\labelenumi{\arabic{enumi}.}
\setcounter{enumi}{2}
\tightlist
\item
  Use the same recurrence with \(x=-\tfrac12\):
\end{enumerate}

\[
\Gamma\!\left(\tfrac12\right) \;=\; \Bigl(-\tfrac12\Bigr)\,\Gamma\!\left(-\tfrac12\right),
\]

so

\[
\Gamma\!\left(-\tfrac12\right) \;=\; \frac{\Gamma\!\left(\tfrac12\right)}{-\tfrac12}
\;=\; \frac{\sqrt{\pi}}{-\tfrac12}
\;=\; \boxed{-2\sqrt{\pi}}.
\]

These match Knuth's official answers for §1.2.5--7 (``\(\sqrt{\pi}\) and
\(-2\sqrt{\pi}\)''; this exercise uses the recurrence in Ex. 10). 

Optional intuition:

Many texts justify the given value
\(\left(\tfrac12\right)! = \Gamma(\tfrac32) = \sqrt{\pi}/2\) via the
Gaussian integral:

\[
\Gamma\!\left(\tfrac12\right)=\int_{0}^{\infty} e^{-t} t^{-1/2}\,dt
= 2\int_{0}^{\infty} e^{-u^{2}}\,du=\sqrt{\pi},
\]

then apply the recurrence to get
\(\Gamma(\tfrac32)=\tfrac12\Gamma(\tfrac12)=\sqrt{\pi}/2\). (Here we
only needed the latter as a given.)

\subsection{8 {[}HM15{]}}\label{hm15}

Let

\[
A_m(n)\;=\;\frac{m^n\,m!}{(n+1)(n+2)\cdots(n+m)}.
\]

\begin{enumerate}
\def\labelenumi{\arabic{enumi}.}
\tightlist
\item
  Use the identity \((n+1)(n+2)\cdots(n+m) = \dfrac{(n+m)!}{n!}\). Then
\end{enumerate}

\[
A_m(n)=\frac{m^n\,m!\,n!}{(n+m)!}.
\]

\begin{enumerate}
\def\labelenumi{\arabic{enumi}.}
\setcounter{enumi}{1}
\tightlist
\item
  Recognize a binomial coefficient:
\end{enumerate}

\[
\binom{n+m}{n}=\frac{(n+m)!}{n!\,m!}
\quad\Longrightarrow\quad
\frac{n!\,m!}{(n+m)!}=\frac{1}{\binom{n+m}{n}}.
\]

Hence

\[
A_m(n)=\frac{m^n}{\binom{n+m}{n}}.
\]

\begin{enumerate}
\def\labelenumi{\arabic{enumi}.}
\setcounter{enumi}{2}
\tightlist
\item
  Rewrite \(\binom{n+m}{n}\) in a way that exposes a finite product of
  \(n\) factors:
\end{enumerate}

\[
\binom{n+m}{n}
=\frac{(m+1)(m+2)\cdots(m+n)}{n!}
=\frac{m^n}{n!}\prod_{k=1}^{n}\left(1+\frac{k}{m}\right).
\]

Therefore

\[
A_m(n)=\frac{m^n}{\binom{n+m}{n}}
=\frac{m^n}{\dfrac{m^n}{n!}\prod_{k=1}^{n}\left(1+\frac{k}{m}\right)}
=\frac{n!}{\displaystyle\prod_{k=1}^{n}\left(1+\frac{k}{m}\right)}.
\]

\begin{enumerate}
\def\labelenumi{\arabic{enumi}.}
\setcounter{enumi}{3}
\tightlist
\item
  Take the limit as \(m\to\infty\). For each fixed \(k\le n\),
  \(\left(1+\frac{k}{m}\right)\to 1\). The product has only \(n\)
  factors, so
\end{enumerate}

\[
\lim_{m\to\infty}A_m(n)
=
\frac{n!}{\displaystyle\prod_{k=1}^{n}\lim_{m\to\infty}\left(1+\frac{k}{m}\right)}
=\frac{n!}{1}=n!.
\]

\begin{enumerate}
\def\labelenumi{\arabic{enumi}.}
\setcounter{enumi}{4}
\tightlist
\item
  Edge case \(n=0\): By definition,
\end{enumerate}

\[
A_m(0)=\frac{m^0\,m!}{1\cdot2\cdots m}=1=0!,
\]

so the limit is also \(0!\).

Thus, Euler's limit indeed equals \(n!\) for every nonnegative integer
\(n\).

\subsection{9 {[}M10{]}}\label{m10-7}

Knuth defines \(n! = \Gamma(n+1)\) and uses the gamma functional
equation

\[
\Gamma(x+1)=x\,\Gamma(x).
\]

Hence \((\tfrac12)! = \Gamma(\tfrac32)\). 

\begin{enumerate}
\def\labelenumi{\arabic{enumi}.}
\tightlist
\item
  Compute \(\Gamma(\tfrac12)\).
\end{enumerate}

From the given \((\tfrac12)! = \Gamma(\tfrac32)=\sqrt{\pi}/2\) and the
recurrence with \(x=\tfrac12\):

\[
\Gamma\!\left(\tfrac32\right)=\tfrac12\,\Gamma\!\left(\tfrac12\right).
\]

Therefore

\[
\Gamma\!\left(\tfrac12\right)=2\,\Gamma\!\left(\tfrac32\right)=2\cdot\frac{\sqrt{\pi}}{2}=\boxed{\sqrt{\pi}}.
\]

(Consistency check, optional: the reflection formula
\(\Gamma(z)\Gamma(1-z)=\dfrac{\pi}{\sin \pi z}\) gives
\(\Gamma(\tfrac12)^2=\pi\Rightarrow \Gamma(\tfrac12)=\sqrt{\pi}\).) 

\begin{enumerate}
\def\labelenumi{\arabic{enumi}.}
\setcounter{enumi}{1}
\tightlist
\item
  Compute \(\Gamma(-\tfrac12)\).
\end{enumerate}

Apply the recurrence backward at \(x=-\tfrac12\):

\[
\Gamma\!\left(\tfrac12\right)=\left(-\tfrac12\right)\Gamma\!\left(-\tfrac12\right)
\quad\Longrightarrow\quad
\Gamma\!\left(-\tfrac12\right)=\frac{\Gamma(\tfrac12)}{-\tfrac12}
= -2\,\Gamma\!\left(\tfrac12\right)
= \boxed{-2\sqrt{\pi}}.
\]

Thus

\[
\boxed{\Gamma\!\left(\tfrac12\right)=\sqrt{\pi}},\qquad
\boxed{\Gamma\!\left(-\tfrac12\right)=-2\sqrt{\pi}}.
\]

\subsection{10 {[}HM20{]}}\label{hm20-4}

Background (what Γ means here)

Knuth/Euler define the gamma function via Euler's limit (valid for
real/complex x except at non-positive integers, where poles occur):

\[
\Gamma(x)\;=\;\lim_{m\to\infty}\frac{m^{\,x}\,m!}{x(x+1)\cdots(x+m)}.
\]

This is equation (15) in the text. 

There is also the integral form (for x\textgreater0), proved a few
exercises later:

\[
\Gamma(x)=\int_0^\infty e^{-t}\,t^{x-1}\,dt.
\]

Conclusion

\textbf{Yes---except when \(x\) is 0 or a negative integer.} At those
points Γ has poles, so the identity is not meaningful there.

Proof 1 (from Euler's limit; works wherever Γ is defined)

Start from the defining limit with \(x\mapsto x+1\):

\[
\Gamma(x+1)\;=\;\lim_{m\to\infty}\frac{m^{\,x+1}\,m!}{(x+1)(x+2)\cdots(x+m+1)}.
\]

Factor a convenient term:

\[
\Gamma(x+1)
=\lim_{m\to\infty}\frac{m}{x+m+1}\cdot\frac{m^{\,x}\,m!}{(x+1)\cdots(x+m)}.
\]

But

\[
x\,\Gamma(x)=x\cdot\lim_{m\to\infty}\frac{m^{\,x}\,m!}{x(x+1)\cdots(x+m)}
=\lim_{m\to\infty}\frac{m^{\,x}\,m!}{(x+1)\cdots(x+m)}.
\]

Since \(\frac{m}{x+m+1}\to 1\) as \(m\to\infty\), the two limits are
equal:

\[
\Gamma(x+1)=x\,\Gamma(x).
\]

This algebra is the essence of the official solution in the book. 

\textbf{Domain caveat.} The products \((x+j)\) in the denominator hit
zero exactly when \(x\in\{0,-1,-2,\dots\}\); that is where Γ has poles,
so the identity isn't asserted there. Everywhere else (real x not a
non-positive integer), the argument goes through.

Proof 2 (for \(x>0\) only, via the integral)

Using \(\Gamma(x)=\int_0^\infty e^{-t}t^{x-1}dt\) and integration by
parts:

\[
\Gamma(x+1)=\int_0^\infty e^{-t}t^{x}\,dt
=\Big[-e^{-t}t^{x}\Big]_0^\infty + x\int_0^\infty e^{-t}t^{x-1}\,dt
=x\,\Gamma(x),
\]

since the boundary term vanishes for \(x>0\). This gives the same
recurrence on \((0,\infty)\). 

Remarks

\begin{itemize}
\tightlist
\item
  Together with \(\Gamma(1)=1\), the recurrence yields
  \(\Gamma(n+1)=n!\) for positive integers \(n\), matching the factorial
  rule.
\item
  In the complex plane, \(\Gamma\) is meromorphic with simple poles at
  \(0,-1,-2,\dots\); the same functional equation holds on its domain.
  (TAOCP also records the related reflection identity (16).)
\end{itemize}

\subsection{11 {[}M15{]}}\label{m15-6}

Proof via Legendre's formula (one line, then a tidy sum)

By Legendre's formula the exponent of a prime \(p\) in \(n!\) is

\[
v_p(n!)=\sum_{k\ge 1}\Big\lfloor \frac{n}{p^k}\Big\rfloor.
\]

For \(p=2\),

\[
v_2(n!)=\sum_{k\ge 1}\Big\lfloor \frac{n}{2^k}\Big\rfloor. \tag{★}
\]

(This is exactly the ``prime--power multiplicity'' identity stated in
the text.) 

Now use the binary form of \(n\). Since \(n=\sum_{j=1}^r 2^{e_j}\),

\[
\Big\lfloor\frac{n}{2}\Big\rfloor=\sum_{j:e_j\ge 1}2^{e_j-1},\quad
\Big\lfloor\frac{n}{4}\Big\rfloor=\sum_{j:e_j\ge 2}2^{e_j-2},\quad
\Big\lfloor\frac{n}{8}\Big\rfloor=\sum_{j:e_j\ge 3}2^{e_j-3},\ \ldots
\]

Summing these rows termwise (exactly the series (★)) yields, for each
fixed \(j\),

\[
2^{e_j-1}+2^{e_j-2}+\cdots+2+1=2^{e_j}-1.
\]

Therefore

\[
v_2(n!)=\sum_{k\ge 1}\Big\lfloor \frac{n}{2^k}\Big\rfloor
=\sum_{j=1}^r (2^{e_j}-1)
=\Big(\sum_{j=1}^r 2^{e_j}\Big)-r
=n-r.
\]

So \(n!\) contains exactly \(n-r\) factors of 2---equivalently
\(2^{\,n-r}\mid n!\) but \(2^{\,n-r+1}\nmid n!\), as required. ✓

Tiny numerical check

Take \(n=13=(1101)_2\) so \(r=3\). Then

\[
v_2(13!)=\left\lfloor\frac{13}{2}\right\rfloor+\left\lfloor\frac{13}{4}\right\rfloor+\left\lfloor\frac{13}{8}\right\rfloor=6+3+1=10
=13-3.
\]

So \(2^{10}\mid 13!\) but \(2^{11}\nmid 13!\), matching the claim.

(Optional) Broader view

Exactly the same argument shows for any prime \(p\) that

\[
v_p(n!)=\frac{n-s_p(n)}{p-1},
\]

where \(s_p(n)\) is the sum of the base-\(p\) digits of \(n\). The
present exercise is the special case \(p=2\), where \(s_2(n)=r\).

\subsection{12 {[}M22{]}}\label{m22-7}

Let \(p\) be prime and write \(n\) in base \(p\) as

\[
n=a_kp^k+a_{k-1}p^{k-1}+\cdots+a_1p+a_0\qquad(0\le a_i<p).
\]

If \(\mu\) denotes the exponent of \(p\) in \(n!\) (i.e.,
\(\mu=\sum_{j\ge1}\big\lfloor n/p^j\big\rfloor\) from Eq. (8)), express
\(\mu\) in a simple formula involving \(n\), \(p\), and the digits
\(a_i\).

Let \(S_p(n)=a_0+a_1+\cdots+a_k\) be the sum of the base-\(p\) digits of
\(n\).

\begin{enumerate}
\def\labelenumi{\arabic{enumi}.}
\tightlist
\item
  For each \(i\ge 1\),
\end{enumerate}

\[
\Big\lfloor \frac{n}{p^i}\Big\rfloor
=\sum_{j=i}^{k} a_j p^{\,j-i}
\]

because dividing by \(p^i\) in base \(p\) just drops the \(i\)
least-significant digits.

\begin{enumerate}
\def\labelenumi{\arabic{enumi}.}
\setcounter{enumi}{1}
\tightlist
\item
  Sum over \(i=1,\dots,k\) and swap the order of summation:
\end{enumerate}

\[
\sum_{i=1}^{k}\Big\lfloor \frac{n}{p^i}\Big\rfloor
=\sum_{i=1}^{k}\sum_{j=i}^{k} a_j p^{\,j-i}
=\sum_{j=1}^{k} a_j \sum_{i=1}^{j} p^{\,j-i}
=\sum_{j=1}^{k} a_j \frac{p^{\,j}-1}{p-1}.
\]

\begin{enumerate}
\def\labelenumi{\arabic{enumi}.}
\setcounter{enumi}{2}
\tightlist
\item
  Simplify:
\end{enumerate}

\[
\sum_{j=1}^{k} a_j \frac{p^{\,j}-1}{p-1}
=\frac{1}{p-1}\left(\sum_{j=1}^{k} a_j p^{\,j}-\sum_{j=1}^{k} a_j\right)
=\frac{1}{p-1}\left((n-a_0)-(S_p(n)-a_0)\right)
=\frac{n-S_p(n)}{p-1}.
\]

But the left-hand side is exactly \(\mu\) by Eq. (8). Therefore

\[
\boxed{\;\mu \;=\; \dfrac{n - S_p(n)}{\,p-1\,}\;}
\]

which is Legendre's (de Polignac's) formula.

\subsection{Check}\label{check}

For \(n=1000\) and \(p=3\), \(1000\) in base 3 is
\(1\,1\,0\,1\,0\,0\,1\) so \(S_3(1000)=4\). Hence

\[
\mu=\frac{1000-4}{3-1}=498,
\]

matching the worked value obtained by summing
\(\big\lfloor 1000/3^j\big\rfloor\).

This also recovers Exercise 11's binary special case: when \(p=2\),
\(S_2(n)\) is the number of 1s in the binary expansion, so
\(v_2(n!)=n-S_2(n)\).

\subsection{13 {[}M22{]}}\label{m22-8}

Proof (pairing multiplicative inverses)

\begin{enumerate}
\def\labelenumi{\arabic{enumi}.}
\item
  \textbf{Every nonzero residue has a unique inverse mod \(p\).} For
  each \(n\in\{1,2,\dots,p-1\}\) there exists a unique \(n'\) with

  \[
  nn' \equiv 1 \pmod p,
  \]

  because \(\gcd(n,p)=1\) when \(p\) is prime (Bezout's identity gives
  existence; uniqueness follows since
  \(n(a-b)\equiv0\pmod p\Rightarrow a\equiv b\)). Also, applying ``take
  the inverse'' twice gives \((n')'=n\).
\item
  \textbf{Only \(1\) and \(p-1\) are self-inverse.} If \(n'=n\), then
  \(n^2\equiv1\pmod p\), so \((n-1)(n+1)\equiv0\pmod p\). A prime \(p\)
  can divide at most one factor, hence \(n\equiv1\) or
  \(n\equiv -1\equiv p-1\pmod p\). Thus among \(\{1,2,\dots,p-1\}\),
  exactly the two elements \(1\) and \(p-1\) satisfy \(n'=n\).
\item
  \textbf{Pair everything else.} Partition \(\{1,2,\dots,p-1\}\) into
  disjoint pairs \(\{n,n'\}\) with \(n\neq n'\), together with the two
  singletons \(\{1\}\) and \(\{p-1\}\). Each paired product is

  \[
  n\cdot n' \equiv 1 \pmod p.
  \]
\item
  \textbf{Multiply all residues.} Taking the product of all elements
  \(1\cdot2\cdot\ldots\cdot(p-1)\) and grouping by pairs,

  \[
  (p-1)! \equiv \big(\prod_{\text{pairs } \{n,n'\}} nn'\big)\cdot 1\cdot(p-1) \equiv 1\cdot(p-1) \equiv -1 \pmod p.
  \]

  This is exactly the desired congruence.
\end{enumerate}

\subsection{14 {[}M28{]}}\label{m28-1}

We work with the \textbf{unit part} of \(n!\) (i.e., \(n!/p^{\mu}\))
modulo \(p\).

\begin{enumerate}
\def\labelenumi{\arabic{enumi}.}
\tightlist
\item
  Block decomposition of \(1,2,\dots,n\). Split \(\{1,\dots,n\}\) into:
\end{enumerate}

\begin{itemize}
\item
  all \textbf{multiples of \(p\)}: \(p,2p,\dots,\lfloor n/p\rfloor p\),
  whose product is \(p^{\lfloor n/p\rfloor}\cdot \lfloor n/p\rfloor!\);
\item
  all \textbf{nonmultiples of \(p\)}, further grouped into consecutive
  blocks

  \[
  \{mp+1, mp+2,\dots, mp+p-1\}\quad (m=0,1,\dots,\lfloor n/p\rfloor-1),
  \]

  plus a final tail
  \(\{ \lfloor n/p\rfloor p+1,\dots,\lfloor n/p\rfloor p+a_0\}\) of
  length \(a_0\).
\end{itemize}

\begin{enumerate}
\def\labelenumi{\arabic{enumi}.}
\setcounter{enumi}{1}
\item
  Product of a complete block of nonmultiples. Modulo \(p\) each
  complete block is congruent to
  \(1\cdot 2\cdots (p-1)=(p-1)!\equiv -1\pmod p\) by Wilson's theorem.
  Hence the product over all complete blocks contributes a factor
  \((-1)^{\lfloor n/p\rfloor}\). (Wilson's theorem is Ex. 13.)
\item
  Product of the tail. The last incomplete block contributes
\end{enumerate}

\[
(\lfloor n/p\rfloor p+1)\cdots(\lfloor n/p\rfloor p+a_0)\equiv 1\cdot 2\cdots a_0 \equiv a_0!\pmod p.
\]

\begin{enumerate}
\def\labelenumi{\arabic{enumi}.}
\setcounter{enumi}{3}
\tightlist
\item
  First recurrence for the unit part. Let \(U(n):=n!/p^{\nu_p(n!)}\) be
  the unit part of \(n!\). Using the split above,
\end{enumerate}

\[
U(n)\equiv (-1)^{\lfloor n/p\rfloor}\,a_0!\cdot U\big(\lfloor n/p\rfloor\big)\pmod p,
\]

because the multiples of \(p\) contribute
\(p^{\lfloor n/p\rfloor}\cdot \lfloor n/p\rfloor!\), and after dividing
by the corresponding power of \(p\) we are left with the \textbf{unit
part} \(U(\lfloor n/p\rfloor)\).

\begin{enumerate}
\def\labelenumi{\arabic{enumi}.}
\setcounter{enumi}{4}
\tightlist
\item
  Iterate the recurrence. Apply Step 4 to \(\lfloor n/p\rfloor\), then
  to \(\lfloor n/p^2\rfloor\), and so on, until reaching \(0\). Writing
  the base-\(p\) digits explicitly,
\end{enumerate}

\[
\lfloor n/p\rfloor = a_k p^{k-1} + \cdots + a_2 p + a_1,\quad
\lfloor n/p^2\rfloor = a_k p^{k-2} + \cdots + a_2,\ \dots
\]

the iteration yields

\[
U(n)\equiv \Big(\prod_{j\ge 1} (-1)^{\lfloor n/p^j\rfloor}\Big)\cdot
\Big(\prod_{i=0}^k a_i!\Big)\pmod p.
\]

But \(\sum_{j\ge 1}\lfloor n/p^j\rfloor = \mu\) (Legendre's formula),
hence

\[
\boxed{\ \displaystyle \frac{n!}{p^{\mu}} \equiv (-1)^{\mu}\, a_0!\,a_1!\cdots a_k! \pmod p\ }.
\]

This is exactly the Stickelberger congruence requested.

Quick check (example)

Take \(p=5\), \(n=47\). Base-5: \(47=1\cdot 5^2+4\cdot 5+2\) so
\((a_2,a_1,a_0)=(1,4,2)\).
\(\mu=\lfloor 47/5\rfloor+\lfloor 47/25\rfloor=9+1=10\). RHS:
\((-1)^{10}\cdot (2!\cdot 4!\cdot 1!)\equiv 1\cdot (2\cdot 24)\equiv 48\equiv 3\pmod 5\).
Indeed \(47!/5^{\mu}\equiv 3\pmod 5\).

\subsection{15 {[}HM15{]}}\label{hm15-1}

Let \(A\) be the \(n\times n\) matrix with entries \(A_{ij}=i\cdot j\)
(so it's the multiplication table \(1,2,\dots,n\) by \(1,2,\dots,n\)).
Compute the \textbf{permanent} of \(A\),

\[
\operatorname{perm}(A)\;=\;\sum_{\sigma\in S_n}\ \prod_{i=1}^n A_{i,\sigma(i)}\,.
\]

Solution 1 (rank-1 factorization → one-line evaluation)

Observe that \(A\) factors as an outer product:

\[
A = u\,v^{\mathsf T},\qquad
u=\begin{bmatrix}1\\2\\ \vdots\\ n\end{bmatrix},\quad
v=\begin{bmatrix}1\\2\\ \vdots\\ n\end{bmatrix}.
\]

Hence \(A_{i,\sigma(i)}=u_i\,v_{\sigma(i)}=i\cdot \sigma(i)\). For any
rank-1 matrix \(u v^{\mathsf T}\),

\[
\operatorname{perm}(u v^{\mathsf T})
=\sum_{\sigma\in S_n}\ \prod_{i=1}^n u_i\,v_{\sigma(i)}
=\Bigl(\prod_{i=1}^n u_i\Bigr)\ \sum_{\sigma\in S_n}\ \prod_{i=1}^n v_{\sigma(i)}.
\]

But \(\prod_{i=1}^n v_{\sigma(i)}\) equals \(\prod_{j=1}^n v_j\) (a
product is order-independent), so every one of the \(n!\) terms is the
same. Therefore

\[
\operatorname{perm}(u v^{\mathsf T})
= n!\;\Bigl(\prod_{i=1}^n u_i\Bigr)\;\Bigl(\prod_{j=1}^n v_j\Bigr).
\]

Plugging \(u_i=i\) and \(v_j=j\) gives \(\prod_{i=1}^n u_i=n!\) and
\(\prod_{j=1}^n v_j=n!\). Hence

\[
\boxed{\ \operatorname{perm}(A)=(n!)^3\ }.
\]

Solution 2 (direct term-by-term reasoning)

Expand the permanent by definition:

\[
\operatorname{perm}(A)=\sum_{\sigma\in S_n}\ \prod_{i=1}^n (i\cdot \sigma(i)).
\]

For any fixed \(\sigma\),

\[
\prod_{i=1}^n (i\cdot \sigma(i))
=\Bigl(\prod_{i=1}^n i\Bigr)\,\Bigl(\prod_{i=1}^n \sigma(i)\Bigr)
=(n!)\cdot(n!)=(n!)^2,
\]

since \(\{\sigma(1),\dots,\sigma(n)\}=\{1,\dots,n\}\). There are \(n!\)
permutations \(\sigma\), and every term is the same, so

\[
\operatorname{perm}(A)=n!\cdot (n!)^2=(n!)^3.
\]

Quick sanity checks

\begin{itemize}
\tightlist
\item
  \(n=1\): \(\begin{bmatrix}1\end{bmatrix}\), permanent \(=1=(1!)^3\).
\item
  \(n=2\): \(\begin{bmatrix}1&2\\2&4\end{bmatrix}\), permanent
  \(=1\cdot4+2\cdot2=8=(2!)^3\).
\item
  \(n=3\): each term is \((3!)^2=36\); \(6\) terms
  \(\Rightarrow 6\cdot36=216=(3!)^3\).
\end{itemize}

\subsection{16 {[}HM15{]}}\label{hm15-2}

Write the coefficient of the \(k\)-th falling factorial as

\[
c_k \;:=\; \sum_{j=0}^{k}\frac{(-1)^j}{j!}\,.
\]

These are the partial sums of the Taylor series of \(e^{-1}\), hence

\[
c_k \xrightarrow[k\to\infty]{} e^{-1}\neq 0,
\]

so in particular \(|c_k|\) stays bounded away from \(0\) for large
\(k\).

Now consider the general term

\[
T_k(n)\;=\;c_k\; n^{\underline{k}}=c_k\,\prod_{r=0}^{k-1}(n-r).
\]

We separate three cases.

\begin{enumerate}
\def\labelenumi{\arabic{enumi}.}
\item
  \(n\in\{0,1,2,\dots\}\) (nonnegative integer). Then
  \(n^{\underline{k}}=0\) for all \(k>n\) (the product has the factor
  \(n-n=0\)), so the series truncates at \(k=n\). Hence it certainly
  converges (indeed, equals \(n!\)). This is exactly the circumstance
  noted under Eq. (11). 
\item
  \(n=-m\) with \(m\in\{1,2,3,\dots\}\) (negative integer). Then
\end{enumerate}

\[
(-m)^{\underline{k}}
= (-m)(-m-1)\cdots(-m-k+1)
= (-1)^k\,m(m+1)\cdots(m+k-1)
= (-1)^k\frac{(m+k-1)!}{(m-1)!}.
\]

Therefore

\[
|T_k(-m)|=|c_k|\cdot\frac{(m+k-1)!}{(m-1)!}\;\xrightarrow{k\to\infty}\;\infty,
\]

because \(|c_k|\to e^{-1}\) and the factorial grows without bound. In
particular the term \(T_k(n)\) does \textbf{not} tend to \(0\), so the
series diverges.

\begin{enumerate}
\def\labelenumi{\arabic{enumi}.}
\setcounter{enumi}{2}
\tightlist
\item
  \(n\notin\mathbb{Z}\) (non-integer real). Choose an integer
  \(A>\lvert n\rvert\). For all \(k\ge A+1\),
\end{enumerate}

\[
\bigl|n^{\underline{k}}\bigr|
=\prod_{r=0}^{k-1}\lvert n-r\rvert
\;\ge\;\prod_{r=A}^{k-1}(r-\lvert n\rvert)
\;=\;\frac{\Gamma(k-\lvert n\rvert)}{\Gamma(A-\lvert n\rvert)}\,.
\]

The right-hand side \(\to\infty\) with \(k\) (since
\(\Gamma(x)\to\infty\) as \(x\to\infty\)). Hence
\(\lvert n^{\underline{k}}\rvert\to\infty\). Because
\(|c_k|\to e^{-1}>0\), we get

\[
\lvert T_k(n)\rvert=|c_k|\,\lvert n^{\underline{k}}\rvert \;\xrightarrow{k\to\infty}\;\infty,
\]

so again the general term does \textbf{not} go to \(0\), and the series
diverges.

(If you prefer a closed form: for non-integer \(n\),

\[
n^{\underline{k}}=\frac{\Gamma(n+1)}{\Gamma(n+1-k)},
\]

and by the reflection identity one can write \(\Gamma(n+1-k)^{-1}\) in
terms of \(\Gamma(k-n)\); since \(\Gamma(k-n)\to\infty\) as
\(k\to\infty\), the same conclusion follows. The truncation only occurs
when \(n\) is a nonnegative integer because then one factor \(n-r\) is
zero.)

Therefore, the only values of \(n\) for which the series in Eq. (11)
converges are the nonnegative integers, precisely those for which the
``infinite'' series terminates.

\subsection{17 {[}HM20{]}}\label{hm20-5}

Let

\[
P_m\;=\;\prod_{n=1}^{m}\frac{(n+\alpha_1)\cdots(n+\alpha_k)}{(n+\beta_1)\cdots(n+\beta_k)}.
\]

We'll evaluate \(\lim_{m\to\infty}P_m\).

Knuth defines the gamma function via the limit (for real/complex \(x\)):

\[
\Gamma(x)\;=\;\lim_{m\to\infty}\frac{m^{\,x}\,m!}{x(x+1)\cdots(x+m)}.
\]

Equivalently, for each fixed \(x\),

\[
\prod_{n=1}^{m}(n+x)\;=\;\frac{m^{\,x}\,m!}{A_m(x)},\quad\text{where }A_m(x)\xrightarrow[m\to\infty]{}\Gamma(1+x).
\]

(This is just algebraic rearrangement of the defining limit.) 

Apply this to each \(\alpha_i\) and \(\beta_i\):

\[
\prod_{n=1}^{m}(n+\alpha_i)=\frac{m^{\alpha_i}m!}{A_m(\alpha_i)},\qquad
\prod_{n=1}^{m}(n+\beta_i)=\frac{m^{\beta_i}m!}{A_m(\beta_i)}.
\]

Therefore

\[
P_m
=\frac{\prod_{i=1}^k \frac{m^{\alpha_i}m!}{A_m(\alpha_i)}}
        {\prod_{i=1}^k \frac{m^{\beta_i}m!}{A_m(\beta_i)}}
= m^{\sum_i\alpha_i-\sum_i\beta_i}
  \cdot
  \frac{\prod_{i=1}^k A_m(\beta_i)}{\prod_{i=1}^k A_m(\alpha_i)}.
\]

By hypothesis \(\sum_i\alpha_i=\sum_i\beta_i\), so the power of \(m\)
cancels. Passing to the limit and using \(A_m(x)\to\Gamma(1+x)\),

\[
\lim_{m\to\infty}P_m
=\frac{\prod_{i=1}^k \Gamma(1+\beta_i)}{\prod_{i=1}^k \Gamma(1+\alpha_i)}.
\]

This is exactly the claimed value of the infinite product.

Why the product converges (quick check)

For large \(n\),

\[
\log\frac{\prod_i(n+\alpha_i)}{\prod_i(n+\beta_i)}
=\sum_i\left(\frac{\alpha_i-\beta_i}{n}\right)+O\!\left(\frac{1}{n^2}\right).
\]

The \(1/n\) term vanishes by \(\sum\alpha_i=\sum\beta_i\), so
\(\sum_n O(1/n^2)\) converges; hence the infinite product converges to a
nonzero finite limit. The restriction that no \(\beta_i\) is a negative
integer prevents a zero factor in the denominator (and matches the fact
\(\Gamma(1+\beta_i)\) would otherwise have a pole).

\subsection{18 {[}M20{]}}\label{m20-6}

Show that

\[
 \prod_{k=1}^{\infty}\frac{4k^{2}}{4k^{2}-1}=\frac{\pi}{2}.
 \]

\begin{enumerate}
\def\labelenumi{\arabic{enumi})}
\tightlist
\item
  Set up the key integral and its recurrence
\end{enumerate}

Let

\[
I_n\;:=\;\int_{0}^{\pi/2}\!\sin^{n}x\,dx\qquad(n\ge 0).
\]

Basic values: \(I_0=\pi/2\) and \(I_1=1\). Integration by parts (take
\(u=\sin^{n-1}x,\;dv=\sin x\,dx\)) gives the standard two-step
recurrence

\[
I_n=\frac{n-1}{n}\,I_{n-2}\qquad(n\ge 2).
\]

Iterating this with the bases yields the closed forms

\[
I_{2n}=\frac{(2n-1)!!}{(2n)!!}\cdot\frac{\pi}{2}
\quad\text{and}\quad
I_{2n+1}=\frac{(2n)!!}{(2n+1)!!}\,.
\]

(Here \(m!!\) denotes the double factorial.)

Using \((2n)!!=2^{\,n}n!\) and \((2n-1)!!=\dfrac{(2n)!}{2^{\,n}n!}\), we
also have the handy identity

\[
I_{2n}=\frac{\binom{2n}{n}}{4^n}\cdot\frac{\pi}{2}.
\]

\begin{enumerate}
\def\labelenumi{\arabic{enumi})}
\setcounter{enumi}{1}
\tightlist
\item
  Express a finite Wallis product via \(I_{2n}\) and \(I_{2n+1}\)
\end{enumerate}

Consider the finite product

\[
P_n:=\prod_{k=1}^{n}\frac{4k^2}{4k^2-1}
=\frac{(2n)!!^{\,2}}{(2n-1)!!\,(2n+1)!!}.
\]

Using the relations above,

\[
\frac{(2n-1)!!}{(2n)!!}=\frac{2I_{2n}}{\pi},\qquad
\frac{(2n+1)!!}{(2n)!!}=\frac{1}{I_{2n+1}},
\]

hence

\[
\frac{1}{P_n}
=\frac{(2n-1)!!\,(2n+1)!!}{(2n)!!^{\,2}}
=\Bigl(\frac{2I_{2n}}{\pi}\Bigr)\!\cdot\!\Bigl(\frac{1}{I_{2n+1}}\Bigr)
=\frac{2}{\pi}\cdot\frac{I_{2n}}{I_{2n+1}}.
\]

Equivalently,

\[
P_n=\frac{\pi}{2}\cdot\frac{I_{2n+1}}{I_{2n}}.
\tag{★}
\]

\begin{enumerate}
\def\labelenumi{\arabic{enumi})}
\setcounter{enumi}{2}
\tightlist
\item
  Show \(\displaystyle \lim_{n\to\infty}\frac{I_{2n+1}}{I_{2n}}=1\)
\end{enumerate}

Write \(y=\cos x\) (so \(x\in[0,\pi/2]\iff y\in[0,1]\)). Then

\[
I_{2n}=\int_{0}^{1}(1-y^2)^{n}\,dy,\qquad
I_{2n+1}=\int_{0}^{1}(1-y^2)^{n}\sqrt{1-y^2}\,dy.
\]

On \([0,1]\) we have \(1\ge \sqrt{1-y^2}\ge 1-\tfrac{1}{2}y^2\). Thus

\[
I_{2n}\;\ge\;I_{2n+1}\;\ge\;I_{2n}-\frac12\int_{0}^{1}y^2(1-y^2)^{n}\,dy.
\]

The ``error'' integral is \(O\!\big(\tfrac{1}{n}\,I_{2n}\big)\) (e.g.,
by a routine Beta-integral estimate), so

\[
0\le 1-\frac{I_{2n+1}}{I_{2n}} \;\longrightarrow\; 0,
\qquad\text{i.e.}\qquad
\frac{I_{2n+1}}{I_{2n}}\xrightarrow[n\to\infty]{}1.
\]

(Heuristically: the density proportional to \(\sin^{2n}x\) concentrates
near \(x=\pi/2\), where \(\sin x\approx 1\), so inserting one extra
factor of \(\sin x\) barely changes the integral.)

\begin{enumerate}
\def\labelenumi{\arabic{enumi})}
\setcounter{enumi}{3}
\tightlist
\item
  Take limits in (★)
\end{enumerate}

From (★) and the limit above,

\[
\lim_{n\to\infty} P_n
=\frac{\pi}{2}\cdot \lim_{n\to\infty}\frac{I_{2n+1}}{I_{2n}}
=\frac{\pi}{2}\cdot 1
=\frac{\pi}{2}.
\]

Therefore

\[
\boxed{\displaystyle
\prod_{k=1}^{\infty}\frac{4k^{2}}{4k^{2}-1}=\frac{\pi}{2}}
\]

as claimed.

\subsection{19 {[}HM22{]}}\label{hm22}

Let

\[
\Gamma(x)=\lim_{m\to\infty}\frac{m^x\,m!}{x(x+1)\cdots(x+m)}\qquad(\text{Eq.\ (15) in the book})
\]

and denote the \(m\)-th finite approximant by

\[
\Gamma_m(x)\;=\;\frac{m^x\,m!}{x(x+1)\cdots(x+m)}.
\]

Show that for \(x>0\),

\[
\Gamma_m(x)\;=\;\int_{0}^{m}\!\Bigl(1-\tfrac{t}{m}\Bigr)^{m} t^{x-1}\,dt
\;=\; m^{x}\!\int_{0}^{1}\!(1-u)^{m}u^{\,x-1}\,du.
\]

Define

\[
I_m(x)\;:=\;\int_{0}^{m}\!\Bigl(1-\tfrac{t}{m}\Bigr)^{m} t^{x-1}\,dt\qquad(x>0).
\]

Step 1: Change of variables

Let \(t = m u\) (so \(dt = m\,du\), \(u\in[0,1]\)). Then

\[
I_m(x)
= \int_{0}^{1}\!(1-u)^{m}\,(m u)^{x-1}\,m\,du
= m^{x}\!\int_{0}^{1}\!(1-u)^{m}u^{\,x-1}\,du.
\]

This proves the second equality in the exercise. 

Set

\[
J(x,m)\;:=\;\int_{0}^{1}\!(1-u)^{m}u^{\,x-1}\,du\qquad(x>0,\;m\in\mathbb{Z}_{\ge0}),
\]

so \(I_m(x)=m^{x}J(x,m)\).

Step 2: Integration by parts recurrence

Integrate \(J(x,m)\) by parts, choosing

\[
a=(1-u)^{m},\quad db=u^{x-1}du,
\]

so that

\[
da=-m(1-u)^{m-1}du,\quad b=\frac{u^{x}}{x}.
\]

Boundary terms vanish (since \(u^{x}\to 0\) as \(u\to 0^+\) when
\(x>0\), and \((1-u)^m\to 0\) as \(u\to 1^-\)), hence

\[
J(x,m)=\frac{m}{x}\int_{0}^{1}\!u^{x}(1-u)^{m-1}\,du
=\frac{m}{x}\,J(x+1,m-1).
\]

Iterating this \(m\) times yields

\[
J(x,m)=\frac{m}{x}\cdot\frac{m-1}{x+1}\cdot\frac{m-2}{x+2}\cdots\frac{1}{x+m-1}\,J(x+m,0).
\]

But \(J(x+m,0)=\int_{0}^{1}u^{x+m-1}\,du=\frac{1}{x+m}\). Therefore

\[
J(x,m)=\frac{m!}{x(x+1)\cdots (x+m)}.
\]

Step 3: Identify with \(\Gamma_m(x)\)

Finally,

\[
I_m(x)=m^{x}\,J(x,m)
= m^{x}\,\frac{m!}{x(x+1)\cdots(x+m)}
= \Gamma_m(x),
\]

which is exactly the desired identity.

This derivation matches Knuth's suggested approach (``change of variable
\(t=mu\), integration by parts, and induction'').

Remarks

\begin{enumerate}
\def\labelenumi{\arabic{enumi}.}
\tightlist
\item
  The integral \(J(x,m)\) is the beta integral \(B(x,m+1)\), so the
  result can be summarized as
\end{enumerate}

\[
\Gamma_m(x)=m^{x}B(x,m+1)\qquad(x>0),
\]

consistent with the notation and surrounding discussion in §1.2.5. 

\begin{enumerate}
\def\labelenumi{\arabic{enumi}.}
\setcounter{enumi}{1}
\tightlist
\item
  In the very next exercise (20), this representation is used to prove
  the standard Euler integral formula
\end{enumerate}

\[
\Gamma(x)=\int_{0}^{\infty}e^{-t}t^{x-1}\,dt\qquad(x>0)
\]

by bounding \(\Gamma(x)-\Gamma_m(x)\) and letting \(m\to\infty\).

\subsection{20 {[}HM21{]}}\label{hm21}

Fix \(a>0\), and for each integer \(m\ge 1\) define

\[
I_m(a)\;:=\;\int_{0}^{\infty} \frac{t^{a-1}}{\bigl(1+t/m\bigr)^{\,m+1}}\,dt.
\tag{1}
\]

Step 1: Evaluate \(I_m(a)\) exactly (beta--gamma algebra)

Make the substitution \(u=\frac{t}{m+t}\) (so \(t=\frac{mu}{1-u}\),
\(dt=\frac{m\,du}{(1-u)^2}\)). Then

\[
\begin{aligned}
I_m(a)
&=\int_{0}^{1} m^a\,u^{a-1}(1-u)^{m-a}\,du
= m^a\,B\!\left(a,\,m+1-a\right)\\[4pt]
&= m^a\,\frac{\Gamma(a)\,\Gamma(m+1-a)}{\Gamma(m+1)}.
\end{aligned}
\tag{2}
\]

(This is the standard beta integral
\(B(x,y)=\int_0^1 u^{x-1}(1-u)^{y-1}du\) and
\(B(x,y)=\Gamma(x)\Gamma(y)/\Gamma(x+y)\); see the beta--gamma relation
discussed around the §1.2.5 exercises.) 

Step 2: Take the limit \(m\to\infty\)

Use the basic gamma--ratio asymptotic (a one-line consequence of
Stirling's formula): for fixed \(a\),

\[
\frac{\Gamma(m+1-a)}{\Gamma(m+1)}\;\sim\; m^{-a}\qquad(m\to\infty).
\]

Hence from (2),

\[
I_m(a)=m^a\,\Gamma(a)\,\frac{\Gamma(m+1-a)}{\Gamma(m+1)}\;\longrightarrow\;\Gamma(a).
\tag{3}
\]

(Stirling's approximation for \(\Gamma\) in §1.2.5/§1.2.11 justifies the
ratio \$ \Gamma(m+1-a)/\Gamma(m+1)\sim m\^{}\{-a\}\$.) 

Step 3: Identify the pointwise limit of the integrand

For each fixed \(t\ge 0\),

\[
\frac{1}{\bigl(1+t/m\bigr)^{\,m+1}} \;\xrightarrow[m\to\infty]{}\; e^{-t}.
\]

Thus the integrands in (1) converge pointwise to \(e^{-t} t^{a-1}\).

A standard compact--tail argument (or dominated convergence with a
routine bound using Stirling/ratio as above) lets us pass the limit
through the integral, so the limit of \(I_m(a)\) is the integral of the
limit:

\[
\lim_{m\to\infty} I_m(a)=\int_{0}^{\infty} e^{-t}\,t^{a-1}\,dt.
\tag{4}
\]

Step 4: Conclude

From (3) and (4),

\[
\Gamma(a)=\int_{0}^{\infty} e^{-t}\,t^{a-1}\,dt,
\]

which is Euler's integral representation of the gamma function---exactly
the statement \(\gamma(a,\infty)=\Gamma(a)\). (Knuth points to this
identity right where he introduces the incomplete gamma function.) 

Remarks

\begin{enumerate}
\def\labelenumi{\arabic{enumi}.}
\tightlist
\item
  For integer \(a=n\ge1\), the integral reduces to
  \(\int_0^\infty e^{-t} t^{n-1}dt=(n-1)!\), consistent with
  \(\Gamma(n)=(n-1)!\).
\item
  The same proof works for complex \(a\) with \(\Re(a)>0\), giving the
  standard domain of Euler's integral.
\end{enumerate}

\subsection{21 {[}HM25{]}}\label{hm25-1}

Proof via Taylor expansions (partial Bell polynomials)

Write the Taylor expansions around \(x\) and around \(u(x)\):

\[
u(x+t)=u(x)+\sum_{i\ge1}\frac{u^{(i)}(x)}{i!}\,t^i,
\qquad
W\bigl(u(x)+y\bigr)=\sum_{j\ge0}\frac{W^{(j)}(u(x))}{j!}\,y^j.
\]

Set \(g(t)=W(u(x+t))\). Then

\[
g(t)=\sum_{j\ge0}\frac{W^{(j)}(u(x))}{j!}\Bigl(u(x+t)-u(x)\Bigr)^{\!j}
=\sum_{j\ge0}\frac{W^{(j)}(u)}{j!}\Bigl(\sum_{i\ge1} a_i t^i\Bigr)^{\!j},
\]

where \(a_i=u^{(i)}(x)/i!\).

Expand the \(j\)-th power by the multinomial theorem. The coefficient of
\(t^n\) in \(\bigl(\sum_{i\ge1} a_i t^i\bigr)^j\) is

\[
\sum_{\substack{k_1+\cdots+k_n=j\\ 1k_1+\cdots+nk_n=n}}
\frac{j!}{k_1!\cdots k_n!}\;
a_1^{k_1}\cdots a_n^{k_n}
=
\sum_{\substack{k}}
\frac{j!}{k_1!\cdots k_n!}\;
\prod_{i=1}^n \Bigl(\frac{u^{(i)}(x)}{i!}\Bigr)^{k_i}.
\]

Therefore the coefficient of \(t^n\) in \(g(t)\) is

\[
\sum_{j=0}^{n}\sum_{\substack{k}}
\frac{W^{(j)}(u)}{j!}\cdot
\frac{j!}{k_1!\cdots k_n!}\;
\prod_{i=1}^n \Bigl(\frac{u^{(i)}(x)}{i!}\Bigr)^{k_i}.
\]

But \(g(t)=\sum_{n\ge0}\dfrac{D_x^{\,n}w}{n!}\,t^n\). Equating
coefficients of \(t^n\) and multiplying by \(n!\) gives

\[
D_x^{\,n} w
=\sum_{j=0}^{n}\sum_{\substack{k_1+\cdots+k_n=j\\ 1k_1+\cdots+nk_n=n}}
\frac{n!}{k_1!\cdots k_n!}\;
\prod_{i=1}^n \Bigl(\frac{u^{(i)}(x)}{i!}\Bigr)^{k_i}
\, W^{(j)}(u),
\]

which is exactly the required formula after grouping the \((i!)^{-k_i}\)
into the denominator. This also shows the well-known form

\[
\frac{d^n}{dx^n}W(u(x))
=\sum_{j=0}^{n} W^{(j)}(u)\,B_{n,j}\bigl(u',u'',\dots,u^{(n-j+1)}\bigr),
\]

where \(B_{n,j}\) is the partial Bell polynomial; expanding \(B_{n,j}\)
yields the same coefficients.

Quick sanity checks

\begin{itemize}
\item
  \(n=2\): admissible \((k_1,k_2)\) are \((2,0)\) with \(j=2\) and
  \((0,1)\) with \(j=1\).

  \[
  D_x^2 w = W''(u)\,(u')^2 + W'(u)\,u'',
  \]

  the familiar second-derivative chain rule.
\item
  \(n=3\): terms \((k_1,k_2,k_3)=(3,0,0),(1,1,0),(0,0,1)\) give

  \[
  D_x^3 w = W'''(u)\,(u')^3 + 3W''(u)\,u' u'' + W'(u)\,u''',
  \]

  as expected.
\end{itemize}

This completes the proof.

\subsection{22 {[}HM20{]}}\label{hm20-6}

Guess a rule for the ratio \((n+x)!/n!\) as \(n\to\infty\). Check it for
integer \(x\). Then say what it suggests for the ``right'' definition of
\(x!\) when \(x\) isn't an integer.

For fixed \(x\), a very natural guess is

\[
\boxed{\;\lim_{n\to\infty}\frac{(n+x)!}{n!\,n^{x}}=1\;}\qquad(\star)
\]

This says \((n+x)!\) grows like \(n!\,n^{x}\) when \(n\) is large.

Check for integer \(x\)

Let \(x=k\in\mathbb Z_{\ge 0}\). Then

\[
\frac{(n+k)!}{n!\,n^{k}}
=\prod_{i=1}^{k}\frac{n+i}{n}
=\prod_{i=1}^{k}\Bigl(1+\frac{i}{n}\Bigr)\xrightarrow[n\to\infty]{}1,
\]

so \((\star)\) is true.

If \(x=-k\) with \(k\in\mathbb Z_{\ge 0}\),

\[
\frac{(n-k)!}{n!\,n^{-k}}
=\frac{n^{k}}{n(n-1)\cdots(n-k+1)}
=\prod_{i=0}^{k-1}\frac{1}{1-\frac{i}{n}}\xrightarrow[n\to\infty]{}1,
\]

so \((\star)\) also holds for negative integers. Thus the hypothesis is
correct for all integer \(x\).

What it suggests for non-integers

Assume the extension of factorial should keep the basic recursion

\[
x! = x\,(x-1)! \quad \text{for all }x
\]

(the same functional equation used to define \(\Gamma\)). Then, writing

\[
(n+x)! = x!\,(x+1)(x+2)\cdots(x+n),
\]

and dividing by \(n!\,n^{x}\), the hypothesis \((\star)\) gives

\[
1=\lim_{n\to\infty}\frac{(n+x)!}{n!\,n^{x}}
= x!\,\lim_{n\to\infty}\frac{(x+1)(x+2)\cdots(x+n)}{n!\,n^{x}}.
\]

Hence

\[
\boxed{\;x! \;=\; \lim_{n\to\infty}\frac{n^{x}\,n!}{(x+1)(x+2)\cdots(x+n)}\;}
\]

This is exactly Euler's limit definition of the factorial for noninteger
\(x\) (equivalently, \(\Gamma(x+1)\)), as given in the text:

\[
n!\;=\;\lim_{m\to\infty}\frac{m^{\,n}\,m!}{(n+1)(n+2)\cdots(n+m)}
\quad\Longleftrightarrow\quad
\Gamma(x+1)=\lim_{m\to\infty}\frac{m^{\,x}\,m!}{(x+1)\cdots(x+m)}.
\]

So the hypothesis not only matches integers; it directly suggests the
correct extension \(x!:=\Gamma(x+1)\). 

(Optional) Why \((\star)\) actually holds for all real \(x\)

Using Stirling's approximation \(n!\sim \sqrt{2\pi n}\,(n/e)^{n}\), we
get, for fixed real \(x\),

\[
\frac{(n+x)!}{n!\,n^{x}}
=\frac{\Gamma(n+x+1)}{\Gamma(n+1)\,n^{x}}
\;\xrightarrow[n\to\infty]{}\;1,
\]

since \(\Gamma(n+a)/\Gamma(n+b)\sim n^{a-b}\). Thus \((\star)\) is
universally valid and fully consistent with Euler's definition of
\(\Gamma\).

\subsection{23 {[}HM20{]}}\label{hm20-7}

Show the reflection identity

\[
(-z)!\,\Gamma(z)=\frac{\pi}{\sin(\pi z)}\qquad(z\notin\mathbb Z),
\]

given the Euler product for sine

\[
\sin(\pi z)=\pi z\prod_{n=1}^{\infty}\!\left(1-\frac{z^{2}}{n^{2}}\right).
\]

\begin{enumerate}
\def\labelenumi{\arabic{enumi}.}
\tightlist
\item
  Start from the limit definition of the gamma function (Knuth's Eq.
  (15)):
\end{enumerate}

\[
\Gamma(x)=\lim_{m\to\infty}\frac{m^{x}\,m!}{x(x+1)\cdots(x+m)}.
\]

Apply this with \(x=1+z\) and \(x=1-z\):

\[
\Gamma(1+z)=\lim_{m\to\infty}\frac{m^{1+z}m!}{(1+z)(2+z)\cdots(m+1+z)},\qquad
\Gamma(1-z)=\lim_{m\to\infty}\frac{m^{1-z}m!}{(1-z)(2-z)\cdots(m+1-z)}.
\]

Multiplying gives

\[
\Gamma(1+z)\Gamma(1-z)
=\lim_{m\to\infty}\frac{m^{2}(m!)^{2}}
{\displaystyle\prod_{n=1}^{m+1}\bigl((n+z)(n-z)\bigr)}
=\lim_{m\to\infty}\frac{m^{2}(m!)^{2}}
{\displaystyle\prod_{n=1}^{m+1}(n^{2}-z^{2})}.
\]

\begin{enumerate}
\def\labelenumi{\arabic{enumi}.}
\setcounter{enumi}{1}
\tightlist
\item
  Factor out \(n^{2}\) in the denominator:
\end{enumerate}

\[
\prod_{n=1}^{m+1}(n^{2}-z^{2})
=\Bigl(\prod_{n=1}^{m+1}n^{2}\Bigr)\prod_{n=1}^{m+1}\Bigl(1-\frac{z^{2}}{n^{2}}\Bigr)
=((m+1)!)^{2}\prod_{n=1}^{m+1}\Bigl(1-\frac{z^{2}}{n^{2}}\Bigr).
\]

Hence

\[
\Gamma(1+z)\Gamma(1-z)
=\lim_{m\to\infty}\frac{m^{2}}{(m+1)^{2}}\cdot
\frac{1}{\displaystyle\prod_{n=1}^{m+1}\Bigl(1-\frac{z^{2}}{n^{2}}\Bigr)}
=\prod_{n=1}^{\infty}\Bigl(1-\frac{z^{2}}{n^{2}}\Bigr)^{-1},
\]

since \(m^{2}/(m+1)^{2}\to1\).

\begin{enumerate}
\def\labelenumi{\arabic{enumi}.}
\setcounter{enumi}{2}
\tightlist
\item
  Use the given Euler sine product (the exercise's premise):
\end{enumerate}

\[
\sin(\pi z)=\pi z\prod_{n=1}^{\infty}\Bigl(1-\frac{z^{2}}{n^{2}}\Bigr)
\quad\Longrightarrow\quad
\prod_{n=1}^{\infty}\Bigl(1-\frac{z^{2}}{n^{2}}\Bigr)^{-1}
=\frac{\pi z}{\sin(\pi z)}.
\]

Therefore

\[
\Gamma(1+z)\Gamma(1-z)=\frac{\pi z}{\sin(\pi z)}.
\]

\begin{enumerate}
\def\labelenumi{\arabic{enumi}.}
\setcounter{enumi}{3}
\tightlist
\item
  Convert to Knuth's factorial notation. Using
  \(\Gamma(1+z)=z\,\Gamma(z)\) and \((-z)!=\Gamma(1-z)\),
\end{enumerate}

\[
z\,(-z)!\,\Gamma(z)=\frac{\pi z}{\sin(\pi z)}\quad\Longrightarrow\quad
(-z)!\,\Gamma(z)=\frac{\pi}{\sin(\pi z)},
\]

as desired (valid for noninteger \(z\) where both sides are finite).  

\subsection{24 {[}HM21{]}}\label{hm21-1}

The key idea is to compare the sum \(\ln(n!)=\sum_{k=1}^{n}\ln k\) with
the area under the increasing function \(\ln x\).

Because \(\ln x\) is increasing on \([1,\infty)\),

\[
\int_{1}^{n}\!\ln x\,dx\;\le\;\sum_{k=1}^{n}\ln k
\;\le\;\int_{1}^{n}\!\ln x\,dx+\ln n .
\]

(Left inequality: on each interval \([k,k+1)\), \(\ln x\ge \ln k\);
summing gives the integral \(\le\) the sum. Right inequality: write
\(\sum_{k=1}^{n}\ln k=\ln n+\sum_{k=1}^{n-1}\ln k\le \ln n+\int_{1}^{n}\ln x\,dx\).)

Evaluate the integral:

\[
\int_{1}^{n}\!\ln x\,dx=\bigl[x\ln x-x\bigr]_{1}^{n}
= n\ln n-n+1 .
\]

Hence

\[
n\ln n-n+1 \;\le\; \ln(n!) \;\le\; (n+1)\ln n-n+1 .
\]

Exponentiating gives exactly

\[
\boxed{\;\dfrac{n^{\,n}}{e^{\,n-1}}\;\le\; n!\;\le\;\dfrac{n^{\,n+1}}{e^{\,n-1}}\;}
\]

for all integers \(n\ge 1\).

How the hint connects (optional intuition)

From \(1+x\le e^{x}\) we get

\[
\frac{k+1}{k}\le e^{1/k}
\quad\text{and}\quad
e^{1/k}\le \frac{k}{k-1}\qquad(k\ge2).
\]

Multiplying the first inequality for \(k=1,\dots,n-1\) yields

\[
n=\prod_{k=1}^{n-1}\frac{k+1}{k}\le \exp\!\Bigl(\sum_{k=1}^{n-1}\frac1k\Bigr),
\]

so \(\ln n\le H_{n-1}\). Multiplying the second for \(k=2,\dots,n\)
gives

\[
\exp\!\Bigl(\sum_{k=2}^{n}\frac1k\Bigr)\le \prod_{k=2}^{n}\frac{k}{k-1}=n,
\]

so \(H_n\le 1+\ln n\). These classical bounds
(\(\ln n\le H_{n-1}<H_n\le 1+\ln n\)) are the discrete counterpart of
the integral comparisons above; integrating \(\ln x\) is simply the
cleanest way to turn them into bounds on \(\ln(n!)\).

Tightness and context

These inequalities are crude but \emph{uniformly valid} for all
\(n\ge1\). They're consistent with Stirling's formula
\(n!\sim \sqrt{2\pi n}\,(n/e)^n\), which refines both sides by the
missing \(\sqrt{2\pi n}\) factor and lower-order terms; our bounds
differ only by a factor of \(n\) in front and match the same dominant
\((n/e)^n\) behavior.

\subsection{25 {[}M20{]}}\label{m20-7}

Definitions

\begin{itemize}
\item
  Falling factorial (a.k.a. ``to the \(n\) falling''):

  \[
  x^{\underline n}=x(x-1)(x-2)\cdots(x-n+1)\quad(n\ge 0).
  \]
\item
  Rising factorial (a.k.a. Pochhammer symbol):

  \[
  x^{\overline n}=x(x+1)(x+2)\cdots(x+n-1)\quad(n\ge 0).
  \]
\end{itemize}

Knuth also records the useful gamma--function forms

\[
x^{\underline n}=\frac{\Gamma(x+1)}{\Gamma(x+1-n)},\qquad
x^{\overline n}=\frac{\Gamma(x+n)}{\Gamma(x)}.
\]

Proofs for integers \(m,n\ge 0\)

Falling factorial

\[
\begin{aligned}
x^{\underline{m+n}}
&=\prod_{j=0}^{m+n-1}(x-j)
=\left(\prod_{j=0}^{m-1}(x-j)\right)\left(\prod_{j=m}^{m+n-1}(x-j)\right)\\
&=x^{\underline m}\cdot \prod_{j=0}^{n-1}\bigl((x-m)-j\bigr)
= x^{\underline m}\,(x-m)^{\underline n}.
\end{aligned}
\]

Rising factorial

\[
\begin{aligned}
x^{\overline{m+n}}
&=\prod_{j=0}^{m+n-1}(x+j)
=\left(\prod_{j=0}^{m-1}(x+j)\right)\left(\prod_{j=m}^{m+n-1}(x+j)\right)\\
&=x^{\overline m}\cdot \prod_{j=0}^{n-1}\bigl((x+m)+j\bigr)
= x^{\overline m}\,(x+m)^{\overline n}.
\end{aligned}
\]

These are just regroupings of consecutive factors.

Extension to noninteger \(m,n\)

Using the gamma--function definitions (valid wherever the gamma function
is defined), we have

\[
\begin{aligned}
x^{\underline{m+n}}
&=\frac{\Gamma(x+1)}{\Gamma(x+1-(m+n))}
=\frac{\Gamma(x+1)}{\Gamma(x+1-m)}\cdot\frac{\Gamma(x+1-m)}{\Gamma(x+1-m-n)}\\
&=x^{\underline m}\,(x-m)^{\underline n},
\\[6pt]
x^{\overline{m+n}}
&=\frac{\Gamma(x+m+n)}{\Gamma(x)}
=\frac{\Gamma(x+m+n)}{\Gamma(x+m)}\cdot\frac{\Gamma(x+m)}{\Gamma(x)}\\
&=(x+m)^{\overline n}\,x^{\overline m}.
\end{aligned}
\]

Thus the identities hold for real (indeed complex) \(m,n\) and \(x\),
subject only to avoiding poles of the gamma function (i.e., where the
denominators would be \(\Gamma\) at nonpositive integers). This is
exactly the extension Knuth alludes to. 

Quick sanity check

Take \(x=7\), \(m=3\), \(n=2\).

\begin{itemize}
\tightlist
\item
  Falling: \(7^{\underline{5}}=7\cdot6\cdot5\cdot4\cdot3\). Right side:
  \(7^{\underline 3}(7-3)^{\underline 2}=(7\cdot6\cdot5)(4\cdot3)\).
  Same product.
\item
  Rising: \(7^{\overline{5}}=7\cdot8\cdot9\cdot10\cdot11\). Right side:
  \(7^{\overline 3}(7+3)^{\overline 2}=(7\cdot8\cdot9)(10\cdot11)\).
  Same product.
\end{itemize}

\bookmarksetup{startatroot}

\chapter{1.2.6 Binomial coefficients}\label{binomial-coefficients}

\subsection{1 {[}00{]}}\label{section-72}

\begin{enumerate}
\def\labelenumi{\arabic{enumi})}
\tightlist
\item
  Combinatorial bijection
\end{enumerate}

Choose \((n-1)\) elements from an \(n\)-element set. Such a choice is
uniquely determined by the one element that is omitted. Each of the
\(n\) elements can be the omitted one, and different omissions give
different \((n-1)\)-subsets. Hence there are exactly \(n\) such subsets.

\begin{enumerate}
\def\labelenumi{\arabic{enumi})}
\setcounter{enumi}{1}
\tightlist
\item
  Algebraic identity
\end{enumerate}

By definition of binomial coefficients,

\[
\binom{n}{n-1}=\frac{n!}{(n-1)!\,1!}=n.
\]

\begin{enumerate}
\def\labelenumi{\arabic{enumi})}
\setcounter{enumi}{2}
\tightlist
\item
  Pascal recurrence cross-check
\end{enumerate}

Pascal's rule says

\[
\binom{n}{k}=\binom{n-1}{k}+\binom{n-1}{k-1}.
\]

With \(k=n-1\),

\[
\binom{n}{n-1}=\binom{n-1}{n-1}+\binom{n-1}{n-2}
=1+\binom{n-1}{1}=1+(n-1)=n.
\]

\subsection{2 {[}00{]}}\label{section-73}

Answer

\[
\boxed{\binom{0}{0}=1}.
\]

Why this is correct (several equivalent viewpoints)

\begin{enumerate}
\def\labelenumi{\arabic{enumi}.}
\item
  Combinatorial meaning (subsets). \(\binom{n}{k}\) counts the number of
  \(k\)-element subsets of an \(n\)-element set. When \(n=0\) and
  \(k=0\), we're counting 0-element subsets of the empty set. There is
  exactly one such subset---the empty set itself---so
  \(\binom{0}{0}=1\). (This is the official ``answers'' line in TAOCP:
  ``There's exactly one way to choose nothing from the empty set.'') 
\item
  Binomial theorem boundary case. For any integer \(n\ge0\),
\end{enumerate}

\[
(1+x)^n=\sum_{k=0}^{n}\binom{n}{k}x^k.
\]

Putting \(n=0\) gives \((1+x)^0=1\). The right side has only the \(k=0\)
term, so \(\binom{0}{0}x^0=\binom{0}{0}\). Hence \(\binom{0}{0}=1\).

\begin{enumerate}
\def\labelenumi{\arabic{enumi}.}
\setcounter{enumi}{2}
\tightlist
\item
  Factorial formula with the empty product convention. For
  \(0\le k\le n\),
\end{enumerate}

\[
\binom{n}{k}=\frac{n!}{k!(n-k)!}.
\]

Using \(0!=1\) (the empty product), we get

\[
\binom{0}{0}=\frac{0!}{0!\,0!}=\frac{1}{1\cdot1}=1.
\]

\begin{enumerate}
\def\labelenumi{\arabic{enumi}.}
\setcounter{enumi}{3}
\tightlist
\item
  Pascal's triangle boundary. Pascal's identity
  \(\binom{n}{k}=\binom{n-1}{k}+\binom{n-1}{k-1}\) is solved uniquely by
  the boundary conditions \(\binom{n}{0}=\binom{n}{n}=1\) for all
  \(n\ge0\). Setting \(n=0\) yields \(\binom{0}{0}=1\).
\end{enumerate}

\subsection{3 {[}00{]}}\label{section-74}

A 13-card hand is a \emph{combination}: a subset of size 13 from a set
of size 52, where order doesn't matter (this is exactly the setup of
``combinations of \(n\) objects taken \(k\) at a time'').

\begin{enumerate}
\def\labelenumi{\arabic{enumi}.}
\tightlist
\item
  First count ordered selections of 13 distinct cards:
\end{enumerate}

\[
52\cdot 51\cdot 50\cdots 40 \;=\; \frac{52!}{39!}.
\]

\begin{enumerate}
\def\labelenumi{\arabic{enumi}.}
\setcounter{enumi}{1}
\tightlist
\item
  Each unordered hand is counted \(13!\) times here (once for every
  permutation of its 13 cards). Therefore, divide by \(13!\):
\end{enumerate}

\[
\binom{52}{13} \;=\; \frac{52!}{13!\,39!}.
\]

This is the exact number of bridge hands.

\subsection{4 {[}10{]}}\label{section-75}

Method: Legendre's formula on a binomial

For a prime \(p\),

\[
v_p(n!)=\sum_{i\ge1}\left\lfloor \frac{n}{p^i}\right\rfloor,\qquad
v_p\!\binom{n}{k}=v_p(n!) - v_p(k!) - v_p((n-k)!).
\]

Compute \(v_p\bigl(\binom{52}{13}\bigr)\) only for primes \(p\le 52\);
most will cancel to \(0\).

Large primes (quick eliminations)

\begin{itemize}
\tightlist
\item
  If \(40\le p\le 52\) and \(p\) is prime, then \(v_p(52!)=1\) while
  \(v_p(39!)=v_p(13!)=0\). Hence \(v_p=1\) for \(p\in\{41,43,47\}\).
\item
  If \(26\le p\le 39\) and \(p\) is prime, then \(v_p(52!)=1\),
  \(v_p(39!)=1\), \(v_p(13!)=0\). Hence \(v_p=0\) for
  \(p\in\{29,31,37\}\).
\end{itemize}

Middle primes

\begin{itemize}
\tightlist
\item
  \(p=23:\ v_{23}(52!)=\lfloor52/23\rfloor+\lfloor52/23^2\rfloor=2,\ 
  v_{23}(39!)=1,\ v_{23}(13!)=0\Rightarrow v_{23}=1.\)
\item
  \(p=19:\ 2-2-0=0\) so no factor \(19\).
\item
  \(p=17:\ 3-2-0=1\Rightarrow v_{17}=1.\)
\end{itemize}

Small primes

\begin{itemize}
\tightlist
\item
  \(p=13:\ 4-3-1=0\Rightarrow\) no factor \(13\).
\item
  \(p=11:\ 4-3-1=0\Rightarrow\) no factor \(11\).
\item
  \(p=7:\ (\lfloor52/7\rfloor+\lfloor52/49\rfloor) - \lfloor39/7\rfloor - \lfloor13/7\rfloor =(7+1)-5-1=2\Rightarrow v_7=2.\)
\item
  \(p=5:\ (10+2)-(7+1)-2=12-8-2=2\Rightarrow v_5=2.\)
\item
  \(p=3:\ (17+5+1)-(13+4+1)-(4+1)=23-18-5=0\Rightarrow\) no factor
  \(3\).
\item
  \(p=2:\ (26+13+6+3+1)-(19+9+4+2+1)-(6+3+1)=49-35-10=4\Rightarrow v_2=4.\)
\end{itemize}

Result (prime-power form)

\[
\boxed{\binom{52}{13}=2^{4}\,5^{2}\,7^{2}\,17\,23\,41\,43\,47.}
\]

Checks and intuition.

\begin{itemize}
\tightlist
\item
  Primes \(>39\) must appear once (they occur in \(52!\) but not in
  \(39!\) or \(13!\)).
\item
  Primes in \([26,39]\) cancel.
\item
  What remains are the computed exponents above.
\end{itemize}

\subsection{5 {[}05{]}}\label{section-76}

The binomial theorem states

\[
(a+b)^n=\sum_{k=0}^{n}\binom{n}{k}\,a^{\,n-k}b^{\,k}.
\]

Take \(a=10\), \(b=1\), \(n=4\). The needed binomial coefficients are
the 4th row of Pascal's triangle:

\[
\binom40,\binom41,\binom42,\binom43,\binom44 = 1,4,6,4,1.
\]

Therefore

\[
(10+1)^4
= \binom40 10^4
+ \binom41 10^3
+ \binom42 10^2
+ \binom43 10
+ \binom44
= 10000 + 4\cdot1000 + 6\cdot100 + 4\cdot10 + 1.
\]

Compute the sum:

\[
10000 + 4000 + 600 + 40 + 1 = \boxed{14641}.
\]

\subsection{6 {[}10{]}}\label{section-77}

We interpret a ``row \(r\)'' as the infinite sequence
\((\binom{r}{0},\binom{r}{1},\binom{r}{2},\dots)\), with the convention
\(\binom{r}{k}=0\) for \(k<0\). The ordinary top row is \(r=0\):
\((1,0,0,\dots)\).

The addition law lets us solve upward: for each fixed \(r\) and
\(k\ge0\),

\[
\binom{r-1}{k}=\binom{r}{k}-\binom{r-1}{k-1},
\qquad\text{with}\quad \binom{r-1,-1}=0.
\]

Row \(r=-1\)

Use the equation with \(r=0\) and \(C(r,k):=\binom{r}{k}\):

\[
C(0,k)=C(-1,k)+C(-1,k-1).
\]

Since \(C(0,0)=1\) and \(C(0,k)=0\) for \(k\ge1\),

\[
C(-1,0)=1,\quad
0=C(-1,1)+1\Rightarrow C(-1,1)=-1,
\]

and inductively \(C(-1,k)=-C(-1,k-1)\). Thus

\[
\boxed{r=-1:\; 1,\ -1,\ 1,\ -1,\ 1,\ -1,\ \dots } 
\]

Row \(r=-2\)

Now use \(C(-1,k)=C(-2,k)+C(-2,k-1)\):

\[
C(-2,0)=1,\quad
-1=C(-2,1)+1\Rightarrow C(-2,1)=-2,
\]

\[
1=C(-2,2)+(-2)\Rightarrow C(-2,2)=3,
\quad
-1=C(-2,3)+3\Rightarrow C(-2,3)=-4,\ \dots
\]

Hence

\[
\boxed{r=-2:\; 1,\ -2,\ 3,\ -4,\ 5,\ -6,\ 7,\ -8,\ \dots}
\]

Row \(r=-3\)

Use \(C(-2,k)=C(-3,k)+C(-3,k-1)\):

\[
C(-3,0)=1,\quad
-2=C(-3,1)+1\Rightarrow C(-3,1)=-3,
\]

\[
3=C(-3,2)+(-3)\Rightarrow C(-3,2)=6,\quad
-4=C(-3,3)+6\Rightarrow C(-3,3)=-10,\ \dots
\]

Therefore

\[
\boxed{r=-3:\; 1,\ -3,\ 6,\ -10,\ 15,\ -21,\ 28,\ -36,\ \dots}
\]

\subsection{7 {[}12{]}}\label{section-78}

We'll use two basic identities from §1.2.6:

\begin{enumerate}
\def\labelenumi{\arabic{enumi}.}
\tightlist
\item
  The definition (valid for real \(r\) and integer \(k\))
\end{enumerate}

\[
\binom{r}{k}=\frac{r(r-1)\cdots(r-k+1)}{k(k-1)\cdots 1}\quad (k\ge 0),\quad \binom{r}{k}=0\ (k<0).
\]

This is Eq. (3) in the book. 

\begin{enumerate}
\def\labelenumi{\arabic{enumi}.}
\setcounter{enumi}{1}
\tightlist
\item
  Symmetry (for integer \(n\ge 0\), integer \(k\))
\end{enumerate}

\[
\binom{n}{k}=\binom{n}{n-k}.
\]

This is Eq. (6).

\begin{enumerate}
\def\labelenumi{\arabic{enumi}.}
\tightlist
\item
  Compare consecutive entries
\end{enumerate}

From the factorial/``falling product'' form (3) it's immediate that for
\(0\le k<n\),

\[
\frac{\binom{n}{k+1}}{\binom{n}{k}}
=\frac{n-k}{k+1}.
\]

Therefore

\begin{itemize}
\tightlist
\item
  \(\binom{n}{k+1}>\binom{n}{k}\) iff \(n-k>k+1\),
  i.e.~\(k<\frac{n-1}{2}\).
\item
  \(\binom{n}{k+1}=\binom{n}{k}\) iff \(k=\frac{n-1}{2}\).
\item
  \(\binom{n}{k+1}<\binom{n}{k}\) iff \(k>\frac{n-1}{2}\).
\end{itemize}

Thus, as \(k\) increases from \(0\), the row is \emph{strictly
increasing} up to the midpoint and \emph{strictly decreasing}
afterwards.

\begin{enumerate}
\def\labelenumi{\arabic{enumi}.}
\setcounter{enumi}{1}
\tightlist
\item
  Locate the maximum
\end{enumerate}

Two cases:

\begin{itemize}
\item
  \(n=2m\) (even). The ratio exceeds 1 for \(k<m\) and is \(<1\) for
  \(k\ge m\). Hence the unique maximum is at \(k=m=n/2\).
\item
  \(n=2m+1\) (odd). The ratio equals 1 at \(k=m\):

  \[
  \frac{\binom{2m+1}{m+1}}{\binom{2m+1}{m}}
  =\frac{(2m+1)-m}{m+1}=1.
  \]

  Therefore \(\binom{2m+1}{m}=\binom{2m+1}{m+1}\), and these adjacent
  central entries are the (equal) maxima. Symmetry (6) explains the tie.
\end{itemize}

\begin{enumerate}
\def\labelenumi{\arabic{enumi}.}
\setcounter{enumi}{2}
\tightlist
\item
  Final answer
\end{enumerate}

The largest binomial coefficient(s) in row \(n\) occur at

\[
k=\left\lfloor \frac{n}{2}\right\rfloor
\quad\text{and (if \(n\) is odd) also at}\quad
k=\left\lceil \frac{n}{2}\right\rceil.
\]

Equivalently: the maximum is at the middle of the row; if \(n\) is odd,
the two middle entries are equal and both maximal. (Knuth's official
answer states the same conclusion and notes the strict increase before
the midpoint, then decrease, ``clear from (3)''.) 

Quick sanity check (examples)

\begin{itemize}
\tightlist
\item
  \(n=6\): row is \(1,6,15,20,15,6,1\) → max at
  \(k=3=\lfloor 6/2\rfloor\).
\item
  \(n=7\): row is \(1,7,21,35,35,21,7,1\) → maxima at
  \(k=3,4=\lfloor 7/2\rfloor,\lceil 7/2\rceil\).
\end{itemize}

\subsection{8 {[}00{]}}\label{section-79}

Row \(n\) of Pascal's triangle lists the binomial coefficients

\[
\binom{n}{0},\ \binom{n}{1},\ \ldots,\ \binom{n}{n}.
\]

In that row, the entry \(k\) places from the left is \(\binom{n}{k}\),
while the entry \(k\) places from the right is \(\binom{n}{n-k}\).
Equation (6) asserts these two numbers are equal. Therefore every row
reads the same left-to-right as right-to-left: Pascal's triangle is
mirror-symmetric about its vertical center line (the ``nonzero entries
in each row are palindromic). 

Two standard viewpoints that make this symmetry transparent:

\begin{itemize}
\item
  Algebraic (from factorials).

  \[
  \binom{n}{k}=\frac{n!}{k!(n-k)!}=\frac{n!}{(n-k)!k!}=\binom{n}{n-k}.
  \]

  This is exactly Eq. (6). 
\item
  Combinatorial (by complements). Choosing \(k\) elements out of an
  \(n\)-element set is in bijection with choosing which \(n-k\) elements
  to leave out. Hence the counts are equal,
  \(\binom{n}{k}=\binom{n}{n-k}\), yielding the row-wise mirror
  symmetry.
\end{itemize}

Thus, the property reflected by Eq. (6) is the left-right symmetry of
each row of Pascal's triangle. (Knuth's answer phrase: ``the nonzero
entries in each row are the same from left to right as from right to
left.'')

\emph{Example.} Row \(n=4\): \(1,\,4,\,6,\,4,\,1\) --- clearly
symmetric; the middle term \(\binom{4}{2}\) stands on the axis.

\subsection{9 {[}01{]}}\label{section-80}

Knuth defines the generalized binomial coefficient for real \(x\) and
integer \(k\) by

\[
\binom{x}{k}=
\begin{cases}
\displaystyle \frac{x^{\underline{k}}}{k!}, & k\ge 0,\\[6pt]
0, & k<0,
\end{cases}
\]

where \(x^{\underline{k}}=x(x-1)\cdots(x-k+1)\) is the falling
factorial. Thus negative lower indices are by definition zero.

Consider two cases.

\begin{enumerate}
\def\labelenumi{\arabic{enumi}.}
\tightlist
\item
  \(n\ge 0\). Then
\end{enumerate}

\[
\binom{n}{n}=\frac{n^{\underline{n}}}{n!}
=\frac{n!}{n!}=1.
\]

(Combinatorially: there is exactly one way to choose all \(n\) elements
of an \(n\)-set.)

\begin{enumerate}
\def\labelenumi{\arabic{enumi}.}
\setcounter{enumi}{1}
\tightlist
\item
  \(n<0\). Here the lower index equals \(k=n<0\), so by the definition
  above,
\end{enumerate}

\[
\binom{n}{n}=0.
\]

This is also consistent with Knuth's summation convention for the
binomial expansion: terms with \(k<0\) (or, for nonnegative exponents,
\(k>n\)) are taken as zero.

Answer.
\(\boxed{\binom{n}{n}=1 \text{ if } n\ge 0;\ \binom{n}{n}=0 \text{ if } n<0.}\)

\subsection{10 {[}M25{]}}\label{m25-8}

Lemma (middle binomials vanish mod \(p\)). For \(1\le k\le p-1\),

\[
\binom{p}{k}=\frac{p!}{k!(p-k)!}\equiv 0\pmod p .
\]

\emph{Reason.} The numerator \(p!\) has a factor \(p\); the denominator
\(k!(p-k)!\) has none (all factors \(<p\)). Hence (b) holds immediately.

A handy corollary over the finite field \(\mathbb{F}_p\) is the
Freshman's dream

\[
(1+X)^p \equiv 1+X^p \pmod p,
\]

since all intermediate binomial coefficients vanish by (b).

\begin{enumerate}
\def\labelenumi{(\alph{enumi})}
\setcounter{enumi}{4}
\tightlist
\item
  Lucas's theorem (main result)
\end{enumerate}

Write \(n=ap+b\), \(k=cp+d\) with \(a,c\ge 0\) and \(0\le b,d<p\). Over
\(\mathbb{F}_p\),

\[
(1+X)^n=(1+X)^{ap+b}=\bigl((1+X)^p\bigr)^a(1+X)^b \equiv (1+X^p)^a(1+X)^b \pmod p,
\]

using the Freshman's dream. Expand the right-hand side:

\[
(1+X^p)^a(1+X)^b=\Bigl(\sum_{i=0}^a \binom{a}{i}X^{ip}\Bigr)\Bigl(\sum_{j=0}^b \binom{b}{j}X^{j}\Bigr).
\]

The coefficient of \(X^{cp+d}\) is \(\binom{a}{c}\binom{b}{d}\). On the
other hand, that coefficient is \(\binom{n}{k}\). Therefore

\[
\binom{n}{k}\equiv \binom{a}{c}\binom{b}{d}=\binom{\lfloor n/p\rfloor}{\lfloor k/p\rfloor}\binom{n\bmod p}{k\bmod p}\pmod p,
\]

which is exactly (e). (Statement of (e) appears in the exercise list. )

Immediate corollaries of Lucas

From (e) we get (a), (b) again, and (f) almost at a glance.

\begin{itemize}
\tightlist
\item
  \begin{enumerate}
  \def\labelenumi{(\alph{enumi})}
  \tightlist
  \item
    Take \(k=p\). With \(n=ap+b\) and \(k=1\cdot p+0\), Lucas gives
    \(\binom{n}{p}\equiv \binom{a}{1}\binom{b}{0}=a=\lfloor n/p\rfloor\pmod p\).
  \end{enumerate}
\item
  \begin{enumerate}
  \def\labelenumi{(\alph{enumi})}
  \setcounter{enumi}{1}
  \tightlist
  \item
    Take \(n=p=1\cdot p+0\) and \(1\le k\le p-1\) (so
    \(k=0\cdot p + k\)). Then
    \(\binom{p}{k}\equiv \binom{1}{0}\binom{0}{k}=1\cdot 0=0\pmod p\).
  \end{enumerate}
\item
  \begin{enumerate}
  \def\labelenumi{(\alph{enumi})}
  \setcounter{enumi}{5}
  \tightlist
  \item
    Apply (e) repeatedly to the higher \(p\)-adic digits. Writing
    \(n=\sum a_ip^i,\ k=\sum b_ip^i\), a short induction yields
    \(\displaystyle \binom{n}{k}\equiv\prod_{i=0}^r \binom{a_i}{b_i}\pmod p\),
    as claimed.
  \end{enumerate}
\end{itemize}

\begin{enumerate}
\def\labelenumi{(\alph{enumi})}
\setcounter{enumi}{2}
\tightlist
\item
  and (d) from (a), (b) and the addition formula
\end{enumerate}

Recall the addition formula (Pascal's rule), Eq. (9):

\[
\binom{r}{k}=\binom{r-1}{k}+\binom{r-1}{k-1}.
\]

(See the text right before the exercises. )

\begin{itemize}
\item
  \begin{enumerate}
  \def\labelenumi{(\alph{enumi})}
  \setcounter{enumi}{2}
  \tightlist
  \item
    For \(0\le k\le p-1\),
  \end{enumerate}

  \[
  0\equiv \binom{p}{k}=\binom{p-1}{k}+\binom{p-1}{k-1}\pmod p
  \]

  by (b), so \(\binom{p-1}{k}\equiv -\binom{p-1}{k-1}\pmod p\). With
  \(\binom{p-1}{0}\equiv 1\), induction gives
  \(\binom{p-1}{k}\equiv(-1)^k\).
\item
  \begin{enumerate}
  \def\labelenumi{(\alph{enumi})}
  \setcounter{enumi}{3}
  \tightlist
  \item
    For \(2\le k\le p-1\),
  \end{enumerate}

  \[
  \binom{p+1}{k}=\binom{p}{k}+\binom{p}{k-1}\equiv 0+0\equiv 0\pmod p,
  \]

  since both terms vanish by (b).
\end{itemize}

Sanity checks

\begin{itemize}
\tightlist
\item
  For \(p=5\), \(n=23\), (a) says
  \(\binom{23}{5}\equiv\lfloor 23/5\rfloor=4\pmod 5\). Indeed
  \(\binom{23}{5}=33{,}6492\equiv 4\ (\bmod 5)\).
\item
  For \(p=2\), (f) yields the classic parity rule: \(\binom{n}{k}\) is
  odd iff every binary digit \(b_i\le a_i\) (i.e., \(k\) is a bitwise
  subset of \(n\)).
\end{itemize}

\subsection{11 {[}M20{]}}\label{m20-8}

\begin{enumerate}
\def\labelenumi{\arabic{enumi})}
\tightlist
\item
  Reduce to a \(p\)-adic valuation of the binomial
\end{enumerate}

Write \(v_p(m)\) for the exponent of \(p\) in \(m\). Then

\[
v_p\!\binom{a+b}{a}
= v_p\big((a+b)!\big) - v_p(a!) - v_p(b!).
\]

\begin{enumerate}
\def\labelenumi{\arabic{enumi})}
\setcounter{enumi}{1}
\tightlist
\item
  Use Legendre's formula via sum of digits
\end{enumerate}

If \(s_p(n)\) denotes the sum of the base-\(p\) digits of \(n\),
Legendre's theorem gives

\[
v_p(n!) \;=\; \frac{\,n - s_p(n)\,}{p-1}.
\]

(Equivalently, \(v_p(n!)=\sum_{i\ge 1}\left\lfloor n/p^i\right\rfloor\)
and \(\sum_{i\ge 1}\lfloor n/p^i\rfloor=(n-s_p(n))/(p-1)\).) 

Applying this to \(n=a\), \(b\), and \(a+b\),

\[
v_p\!\binom{a+b}{a}
= \frac{s_p(a)+s_p(b)-s_p(a+b)}{p-1}.
\tag{★}
\]

\begin{enumerate}
\def\labelenumi{\arabic{enumi})}
\setcounter{enumi}{2}
\tightlist
\item
  ``Carries = drop in digit-sum divided by \(p-1\)''
\end{enumerate}

Let \(a=\sum_i a_ip^i\), \(b=\sum_i b_ip^i\) with
\(0\le a_i,b_i\le p-1\). Perform schoolbook addition in base \(p\). Let
\(c_i\in\{0,1\}\) be the carry \emph{into} position \(i\) (so
\(c_0=0\)), and \(r_i\) the resulting digit at position \(i\). Then for
each \(i\),

\[
a_i+b_i+c_i \;=\; r_i + p\,c_{i+1}
\quad(0\le r_i\le p-1).
\]

Summing over all positions and rearranging yields the identity

\[
s_p(a)+s_p(b)\;=\; s_p(a+b) + (p-1)\sum_i c_{i+1}.
\]

Hence the number of carries \(C:=\sum_i c_{i+1}\) satisfies

\[
C=\frac{s_p(a)+s_p(b)-s_p(a+b)}{p-1}.
\tag{†}
\]

\begin{enumerate}
\def\labelenumi{\arabic{enumi})}
\setcounter{enumi}{3}
\tightlist
\item
  Combine (★) and (†)
\end{enumerate}

From (★) and (†) we get

\[
v_p\!\binom{a+b}{a}\;=\;C,
\]

i.e., the exponent of \(p\) dividing \(\binom{a+b}{a}\) equals the
number of carries when adding \(a\) and \(b\) in base \(p\). This proves
the claim.

Quick check (example)

Let \(p=2\), \(a=5=(101)_2\), \(b=7=(111)_2\). Adding

\[
\;\;101
+\;111
=\;1100
\]

produces carries in the 0th, 1st, and 2nd bit positions: \(C=3\).
\(\displaystyle \binom{12}{5}=792=2^3\cdot 99\), so
\(v_2\!\binom{12}{5}=3=C\), as predicted.

\subsection{12 {[}M22{]}}\label{m22-9}

Yes. Precisely those \(n\) of the form

\[
n=2^{m}-1\quad(m\ge 1).
\]

Proof 1 (via Lucas's theorem, base 2)

Write \(n\) and \(k\) in base 2:

\[
n=\sum_{i\ge 0} n_i\,2^i,\qquad k=\sum_{i\ge 0} k_i\,2^i,\qquad n_i,k_i\in\{0,1\}.
\]

Lucas's theorem (exercise 10(e--f)) states

\[
\binom{n}{k}\equiv \prod_{i\ge 0}\binom{n_i}{k_i}\pmod{2}.
\]

Hence \(\binom{n}{k}\) is odd iff
\(\binom{n_i}{k_i}\not\equiv 0\pmod{2}\) for every \(i\); over
\(\mathbb{F}_2\) this means \(k_i\le n_i\) for all \(i\). Equivalently,
the set of odd entries in row \(n\) are exactly those with bitwise \(k\)
contained in the 1-bits of \(n\). In particular, the number of odd
entries equals \(2^{\#\{i:n_i=1\}}\).

Therefore, all \(n+1\) entries are odd iff

\[
n+1=2^{\#\{i:n_i=1\}},
\]

which happens precisely when every binary digit of \(n\) up to its most
significant bit is \(1\), i.e., \(n=2^{m}-1\). Conversely, if
\(n=2^{m}-1\) then every \(n_i=1\), so
\(\binom{n}{k}\equiv\prod_i \binom{1}{k_i}\equiv 1 \pmod{2}\) for all
\(k\), and the whole row is odd. 

\emph{(Lucas's theorem as stated in the exercises:
\(\displaystyle \binom{n}{k}\equiv \binom{\lfloor n/2\rfloor}{\lfloor k/2\rfloor}\binom{n\bmod 2}{k\bmod 2}\pmod 2\);
the general base-\(p\) digitwise product form is also given.)}

Proof 2 (via Kummer's theorem)

Kummer's theorem (exercise 11) says the exponent of a prime \(p\) in
\(\binom{a+b}{a}\) equals the number of carries when adding \(a\) and
\(b\) in base \(p\). Take \(p=2\), \(a=k\), \(b=n-k\). Then
\(v_2\!\binom{n}{k}=0\) (i.e., \(\binom{n}{k}\) odd) iff adding \(k\)
and \(n-k\) in binary produces no carries. If \(n\) has any 0 among its
binary digits, choose \(k\) with a 1 in that position and 0 elsewhere;
adding \(k\) to \(n-k\) then forces a carry at that bit, making
\(\binom{n}{k}\) even---a contradiction. Thus every binary digit of
\(n\) must be 1, i.e., \(n=2^{m}-1\). Conversely, if \(n=2^{m}-1\), any
\(k+(n-k)=n\) has no carries, so all entries are odd. 

(Bonus) Generalization:

For a prime \(p\): none of \(\binom{n}{k}\) (with \(0\le k\le n\)) is
divisible by \(p\) iff \(n=a p^{m}-1\) with \(1\le a<p\), \(m\ge 0\).
(This is the exercise's official ``more generally'' statement.) 

Quick check. \(n=1,3,7,15,\dots\) → rows \((1,1)\), \((1,3,3,1)\),
\((1,7,21,35,35,21,7,1)\), \ldots{} are all odd, as claimed.

\subsection{13 {[}M13{]}}\label{m13}

Proof 1 (telescoping via Pascal's rule)

Pascal's addition formula is

\[
\binom{r}{k}=\binom{r-1}{k}+\binom{r-1}{k-1}\qquad(\text{integer }k),
\]

equivalently

\[
\binom{r+k}{k} \;=\; \binom{r+k+1}{k}-\binom{r+k}{k-1}.
\]

(Simply set \(r\!\gets\! r+k+1,\;k\!\gets\! k\) in Pascal's rule and
rearrange.) 

Now sum that identity for \(k=0,1,\dots,n\):

\[
\sum_{k=0}^{n}\binom{r+k}{k}
= \sum_{k=0}^{n}\Bigl(\binom{r+k+1}{k}-\binom{r+k}{k-1}\Bigr).
\]

Shift index \(j=k-1\) in the second sum. Using the convention
\(\binom{\cdot}{-1}=0\),

\[
\sum_{k=0}^{n}\binom{r+k}{k}
= \sum_{k=0}^{n}\binom{r+k+1}{k}-\sum_{j=-1}^{n-1}\binom{r+j+1}{j}
= \binom{r+n+1}{n},
\]

because all intermediate terms cancel (telescoping), leaving only the
last term. This is exactly Eq. (10). 

Proof 2 (combinatorial ``largest element'' argument)

By symmetry \(\binom{r+k}{k}=\binom{r+k}{r}\). Thus

\[
\sum_{k=0}^{n}\binom{r+k}{k}=\sum_{k=0}^{n}\binom{r+k}{r}
=\sum_{i=r}^{r+n}\binom{i}{r}.
\]

Count the \$ (r+1)\$-subsets of \(\{1,2,\dots,r+n+1\}\) by the value of
their largest element. If the largest element is \(i+1\) (where
\(i=r,r+1,\dots,r+n\)), the other \(r\) elements can be any \(r\)-subset
of \(\{1,\dots,i\}\), counted by \(\binom{i}{r}\). Summing over possible
\(i\) gives

\[
\sum_{i=r}^{r+n}\binom{i}{r}=\binom{r+n+1}{r+1}=\binom{r+n+1}{n},
\]

which is the desired identity.

Remarks:

\begin{itemize}
\tightlist
\item
  The statement is given in the text with \(n\) an integer \(\ge 0\);
  once established for infinitely many \(r\) (e.g., integer \(r\)), one
  can extend to all real \(r\) by the ``polynomial identity'' principle
  noted just above Eq. (9): both sides are polynomials in \(r\) that
  agree at infinitely many points, hence agree identically. 
\item
  This is exactly the summation formula tagged (10) in §1.2.6---the
  exercise asks you to prove it.
\end{itemize}

Sanity check. For \(r=3,n=2\):
\(\binom{3}{0}+\binom{4}{1}+\binom{5}{2}=1+4+10=15=\binom{6}{2}\).

\subsection{14 {[}M21{]}}\label{m21-2}

We will show

\[
\sum_{k=1}^{n} k^{4}
=24\binom{n+1}{5}+36\binom{n+1}{4}+14\binom{n+1}{3}+\binom{n+1}{2},
\]

and hence obtain the usual polynomial form for \(\sum_{k=1}^{n}k^4\).

\begin{enumerate}
\def\labelenumi{\arabic{enumi})}
\tightlist
\item
  Write \(k^4\) in the binomial (Newton) basis
\end{enumerate}

Every polynomial \(P(k)\) of degree \(\le m\) has a unique expansion in
the falling--factorial basis:

\[
P(k)=\sum_{j=0}^{m} \Delta^{j}P(0)\binom{k}{j},
\]

where \(\Delta\) is the forward difference
\((\Delta f)(k)=f(k+1)-f(k)\). For \(P(k)=k^4\) we compute successive
differences:

\[
\begin{aligned}
\Delta k^4 &= 4k^3+6k^2+4k+1,\\
\Delta^2 k^4 &= 12k^2+24k+14,\\
\Delta^3 k^4 &= 24k+36,\\
\Delta^4 k^4 &= 24,\qquad \Delta^5 k^4=0.
\end{aligned}
\]

Evaluating at \(k=0\) gives

\[
\Delta^1P(0)=1,\quad
\Delta^2P(0)=14,\quad
\Delta^3P(0)=36,\quad
\Delta^4P(0)=24.
\]

Hence

\[
k^4
=\binom{k}{1}+14\binom{k}{2}+36\binom{k}{3}+24\binom{k}{4}.
\tag{★}
\]

(Equivalently, this is the well-known identity
\(k^4=\sum_{j=1}^{4} j!\,S(4,j)\binom{k}{j}\) with Stirling numbers
\(S(4,1)=1\), \(S(4,2)=7\), \(S(4,3)=6\), \(S(4,4)=1\); see Knuth's
discussion of Stirling numbers within §1.2.6. )

\begin{enumerate}
\def\labelenumi{\arabic{enumi})}
\setcounter{enumi}{1}
\tightlist
\item
  Sum and use the ``hockey-stick'' identity
\end{enumerate}

Sum (★) from \(k=1\) to \(n\) and apply

\[
\sum_{k=j}^{n}\binom{k}{j}=\binom{n+1}{j+1}
\quad(j\ge0),
\]

which follows immediately from Pascal's rule for binomial coefficients
in §1.2.6. 

Thus

\[
\sum_{k=1}^{n}k^4
= \binom{n+1}{2}+14\binom{n+1}{3}+36\binom{n+1}{4}+24\binom{n+1}{5},
\]

exactly as stated.

\begin{enumerate}
\def\labelenumi{\arabic{enumi})}
\setcounter{enumi}{2}
\tightlist
\item
  Convert to the standard closed form
\end{enumerate}

Substitute \(\binom{n+1}{2}=\frac{n(n+1)}{2}\),
\(\binom{n+1}{3}=\frac{n(n+1)(n-1)}{6}\),
\(\binom{n+1}{4}=\frac{n(n+1)(n-1)(n-2)}{24}\),
\(\binom{n+1}{5}=\frac{n(n+1)(n-1)(n-2)(n-3)}{120}\), and simplify to

\[
\boxed{\displaystyle
\sum_{k=1}^{n}k^4
=\frac{n(n+1)(2n+1)\bigl(3n^2+3n-1\bigr)}{30}
=\frac{n^5}{5}+\frac{n^4}{2}+\frac{n^3}{3}-\frac{n}{30}.
}
\]

\begin{enumerate}
\def\labelenumi{\arabic{enumi})}
\setcounter{enumi}{3}
\tightlist
\item
  Quick check
\end{enumerate}

For \(n=1,2,3\), the left side gives \(1,\,17,\,98\); the closed form
yields the same.

\subsection{15 {[}M15{]}}\label{m15-7}

For a non-negative integer \(n\),

\[
(x+y)^n=\sum_{k=0}^{n}\binom{n}{k}x^{k}y^{\,n-k}.
\]

Proof 1 (combinatorial)

When you expand \((x+y)^n\), you are multiplying \(n\) factors

\[
(x+y)(x+y)\cdots(x+y).
\]

From each factor you must choose either \(x\) or \(y\). A term of the
form \(x^{k}y^{\,n-k}\) arises precisely when you choose \(x\) from
exactly \(k\) of the \(n\) factors and choose \(y\) from the remaining
\(n-k\) factors.

How many such choices are there? Exactly \(\binom{n}{k}\): the number of
ways to select which \(k\) positions (among the \(n\) factors)
contribute an \(x\). Therefore the coefficient of \(x^{k}y^{\,n-k}\) in
the expansion is \(\binom{n}{k}\), yielding

\[
(x+y)^n=\sum_{k=0}^{n}\binom{n}{k}x^{k}y^{\,n-k}.
\]

This completes the proof.

Proof 2 (algebraic, via Pascal's rule)

We will show that the coefficients in \((x+y)^n\) satisfy the same
recursion and base cases as binomial coefficients, hence must be equal.

Write

\[
(x+y)^n=\sum_{k=0}^{n} c_{n,k}\,x^{k}y^{\,n-k},
\]

and define \(c_{n,k}=0\) if \(k<0\) or \(k>n\).

Base case. For \(n=0\): \((x+y)^0=1\), so \(c_{0,0}=1\) and
\(c_{0,k}=0\) for \(k\neq 0\). This matches \(\binom{0}{0}=1\) and
\(\binom{0}{k}=0\) for \(k\neq 0\).

Recurrence. Multiply by \((x+y)\):

\[
(x+y)^{n+1}=(x+y)\sum_{k=0}^{n} c_{n,k}\,x^{k}y^{\,n-k}
= \sum_{k=0}^{n} c_{n,k}\,x^{k+1}y^{\,n-k}
  +\sum_{k=0}^{n} c_{n,k}\,x^{k}y^{\,n+1-k}.
\]

Reindex the first sum with \(j=k+1\) to express everything in powers
\(x^j y^{\,n+1-j}\):

\[
(x+y)^{n+1}=\sum_{j=0}^{n+1}\bigl(c_{n,j}+c_{n,j-1}\bigr)x^{j}y^{\,n+1-j},
\]

(with \(c_{n,-1}=c_{n,n+1}=0\)). Hence

\[
c_{n+1,j} = c_{n,j}+c_{n,j-1}\qquad\text{for all }j.
\]

This is exactly Pascal's rule. Since \(\binom{n}{k}\) has the same base
case and satisfies the same recursion,

\[
c_{n,k}=\binom{n}{k}\quad\text{for all }n,k,
\]

and the binomial identity follows.

\subsection{16 {[}M15{]}}\label{m15-8}

For integers \(n,k\ge 1\),

\[
(-1)^n\binom{-n}{\,k-1\,}=(-1)^k\binom{-k}{\,n-1\,}.
\]

We'll give two short, complementary proofs. Throughout we use the
generalized binomial coefficient definition

\[
\binom{r}{k}=\frac{r(r-1)\cdots(r-k+1)}{k!}\quad\text{for integer }k\ge 0,
\]

and the symmetry \(\binom{n}{k}=\binom{n}{n-k}\) (for integer
\(n\ge 0\)).  

Proof 1 (direct calculation)

Start from the definition with a negative top index. For integer
\(m\ge 0\),

\[
\binom{-n}{m}
=\frac{(-n)(-n-1)\cdots(-n-m+1)}{m!}
=(-1)^m\frac{n(n+1)\cdots(n+m-1)}{m!}.
\]

The product in the numerator is a rising run, so

\[
n(n+1)\cdots(n+m-1)=\frac{(n+m-1)!}{(n-1)!}.
\]

Hence the useful identity

\[
\boxed{\;\binom{-n}{m}=(-1)^m\binom{n+m-1}{m}\;}.
\tag{★}
\]

Apply (★) twice, with \(m=k-1\) and with \(m=n-1\):

\[
(-1)^n\binom{-n}{k-1}
=(-1)^n(-1)^{k-1}\binom{n+k-2}{k-1}
=(-1)^{n+k-1}\binom{n+k-2}{k-1},
\]

\[
(-1)^k\binom{-k}{n-1}
=(-1)^k(-1)^{n-1}\binom{k+n-2}{n-1}
=(-1)^{n+k-1}\binom{n+k-2}{n-1}.
\]

By symmetry of binomial coefficients,
\(\binom{n+k-2}{k-1}=\binom{n+k-2}{n-1}\), so the two expressions are
equal. 

Proof 2 (a slightly more general symmetry)

A compact variant proves the more general identity, valid for all
\(n,k\ge 0\),

\[
\boxed{\;(-1)^n\binom{-n}{k}=(-1)^k\binom{-k}{n}\;}
\]

and then substitutes \(k\gets k-1\).

Using (★),

\[
(-1)^n\binom{-n}{k}
=(-1)^n(-1)^k\binom{n+k-1}{k}
=(-1)^{n+k}\binom{n+k-1}{k},
\]

\[
(-1)^k\binom{-k}{n}
=(-1)^k(-1)^n\binom{n+k-1}{n}
=(-1)^{n+k}\binom{n+k-1}{n}.
\]

These are equal by \(\binom{N}{k}=\binom{N}{N-k}\). Substituting
\(k\mapsto k-1\) yields the exercise.

Checks and edge conditions

\begin{itemize}
\tightlist
\item
  The exercise uses \(k-1\) and \(n-1\) as lower indices, so it is
  natural to assume \(n,k\ge 1\) (to keep the lower indices \(\ge 0\));
  with the standard convention \(\binom{r}{\ell}=0\) for negative
  \(\ell\), the statement also behaves correctly on boundary cases.
\item
  Example: \(n=3,k=4\). LHS
  \(= (-1)^3\binom{-3}{3} = -\big((-1)^3\binom{5}{3}\big)=10\). RHS
  \(= (-1)^4\binom{-4}{2} = (+1)\big((-1)^2\binom{5}{2}\big)=10\).
\end{itemize}

\subsection{17 {[}M18{]}}\label{m18-3}

Goal. Prove Vandermonde's convolution

\[
\sum_k \binom{r}{k}\binom{s}{\,n-k\,}=\binom{r+s}{n}\qquad(\text{integer }n\ge 0),
\]

as stated in Eq. (21) of §1.2.6. Hint in the book: derive it from the
binomial series (Eq. (15)) using \((1+x)^{r+s}=(1+x)^r(1+x)^s\).

Proof via generating functions (Cauchy product)

From Eq. (15) we have the binomial expansion, valid for all integer
\(r\ge0\) and (by analytic continuation) for general \(r\) when
\(|x|<1\):

\[
(1+x)^r=\sum_{i\ge 0}\binom{r}{i}x^i,\qquad
(1+x)^s=\sum_{j\ge 0}\binom{s}{j}x^j.
\]

Multiply these series:

\[
(1+x)^r(1+x)^s
= \Bigl(\sum_{i\ge 0}\binom{r}{i}x^i\Bigr)
  \Bigl(\sum_{j\ge 0}\binom{s}{j}x^j\Bigr)
= \sum_{n\ge 0}\Bigl(\sum_{k}\binom{r}{k}\binom{s}{\,n-k\,}\Bigr)x^n,
\]

where we used the Cauchy product of power series to identify the
coefficient of \(x^n\).

On the other hand,

\[
(1+x)^{r+s}=\sum_{n\ge 0}\binom{r+s}{n}x^n
\]

by the same Eq. (15), with \(r\mapsto r+s\).

Since \((1+x)^{r+s}=(1+x)^r(1+x)^s\) as formal power series (and also as
analytic functions for \(|x|<1\)), the coefficients of \(x^n\) on both
sides must be equal. Hence

\[
\boxed{\displaystyle \sum_k \binom{r}{k}\binom{s}{\,n-k\,}=\binom{r+s}{n}}
\]

for every integer \(n\ge 0\). This is exactly Eq. (21).

Remarks on domains.

\begin{itemize}
\tightlist
\item
  If \(r,s\) are nonnegative integers, the sums are finite, so the
  argument is purely algebraic.
\item
  For general real (or complex) \(r,s\), Eq. (15) is an infinite series
  converging for \(|x|<1\); the Cauchy product is valid in that disk, so
  coefficient equality follows from identity of analytic functions
  there.
\end{itemize}

\subsection{18 {[}M18{]}}\label{m18-4}

Prove

\[
\sum_{k}\binom{r}{m+k}\binom{s}{\,n+k\,}
=\binom{r+s}{\,r-m+n\,},
\]

for integers \(m,n\) and nonnegative integers \(r,s\) (terms with
out-of-range bottoms are taken as \(0\)).

Ingredients

\begin{itemize}
\tightlist
\item
  Vandermonde/Chu--Vandermonde (Eq. (21)):
\end{itemize}

\[
\sum_{t}\binom{r}{t}\binom{s}{N-t}=\binom{r+s}{N},\qquad \text{integer }N.
\]

This is Eq. (21) in the text. 

\begin{itemize}
\tightlist
\item
  Symmetry (Eq. (6)):
\end{itemize}

\[
\binom{n}{k}=\binom{n}{n-k}.
\]

This is Eq. (6) (``symmetry condition'').

Proof (one line of algebra → Vandermonde)

Start with the left-hand side:

\[
S=\sum_{k}\binom{r}{m+k}\binom{s}{\,n+k\,}.
\]

Apply symmetry (Eq. (6)) to the first binomial:

\[
\binom{r}{m+k}=\binom{r}{r-(m+k)}=\binom{r}{(r-m)-k}.
\]

Thus

\[
S=\sum_{k}\binom{r}{(r-m)-k}\binom{s}{\,n+k\,}.
\]

Now change variables to \(t=n+k\) (so \(k=t-n\)):

\[
S=\sum_{t}\binom{r}{(r-m+n)-t}\binom{s}{t}.
\]

This is exactly Vandermonde with \(N=r-m+n\):

\[
S=\sum_{t}\binom{r}{N-t}\binom{s}{t}=\binom{r+s}{N}
=\binom{r+s}{\,r-m+n\,},
\]

which is Eq. (22).

Notes

\begin{itemize}
\tightlist
\item
  The sum is effectively finite because \(\binom{a}{b}=0\) unless
  \(0\le b\le a\) for nonnegative \(a\).
\item
  For context, Eq. (22) is listed among the ``sums of products''
  identities immediately after Eq. (21) in the book.
\end{itemize}

\subsection{19 {[}M18{]}}\label{m18-5}

Proof by induction on \(r\)

Claim. For all integers \(n\) and all integers \(r\ge 0\) (with \(s\)
arbitrary real or integer),

\[
S_r(n,s)\;:=\;\sum_{k} \binom{r}{k}\binom{s+k}{n}(-1)^{\,r-k} \;=\; \binom{s}{\,n-r\,}.
\]

(Only finitely many terms are nonzero because \(\binom{r}{k}=0\) for
\(k<0\) or \(k>r\).)

Base case \(r=0\).

\[
S_0(n,s)=\sum_k \binom{0}{k}\binom{s+k}{n}(-1)^{\,0-k}
=\binom{s+0}{n}
=\binom{s}{n}
=\binom{s}{\,n-0\,}.
\]

So the statement holds for \(r=0\).

Inductive step. Assume the identity holds for some fixed \(r\ge 0\); we
prove it for \(r+1\).

Start from

\[
S_{r+1}(n,s)=\sum_k \binom{r+1}{k}\binom{s+k}{n}(-1)^{\,r+1-k}.
\]

Use Pascal's rule
\(\displaystyle \binom{r+1}{k}=\binom{r}{k}+\binom{r}{k-1}\) to split
the sum:

\[
\begin{aligned}
S_{r+1}(n,s)
&=\sum_k\bigl(\binom{r}{k}+\binom{r}{k-1}\bigr)\binom{s+k}{n}(-1)^{\,r+1-k}\\
&=\underbrace{\sum_k \binom{r}{k}\binom{s+k}{n}(-1)^{\,r+1-k}}_{A}
\;+\;\underbrace{\sum_k \binom{r}{k-1}\binom{s+k}{n}(-1)^{\,r+1-k}}_{B}.
\end{aligned}
\]

Rewrite each piece with a common \((-1)^{r-k}\) factor and shift the
index in \(B\) (let \(k\gets k+1\)):

\[
A=-\sum_k \binom{r}{k}\binom{s+k}{n}(-1)^{\,r-k},\qquad
B=\sum_k \binom{r}{k}\binom{s+k+1}{n}(-1)^{\,r-k}.
\]

Therefore

\[
S_{r+1}(n,s)=\sum_k \binom{r}{k}\Bigl[\binom{s+k+1}{n}-\binom{s+k}{n}\Bigr](-1)^{\,r-k}.
\]

Now use the increment identity

\[
\binom{x+1}{n}-\binom{x}{n}=\binom{x}{\,n-1\,}\quad(\text{valid for all real }x,\text{ integer }n),
\]

to obtain

\[
S_{r+1}(n,s)=\sum_k \binom{r}{k}\binom{s+k}{\,n-1\,}(-1)^{\,r-k}
= S_r(n-1,s).
\]

Apply the induction hypothesis to \(S_r(n-1,s)\):

\[
S_r(n-1,s)=\binom{s}{(n-1)-r}=\binom{s}{\,n-(r+1)\,},
\]

which is exactly the desired formula for \(r+1\).

By induction, the identity holds for all integers \(r\ge 0\). ∎

Remarks and intuition

\begin{itemize}
\item
  The right-hand side \(\binom{s}{\,n-r\,}\) automatically gives \(0\)
  when \(r>n\) (with integer \(n\)), matching the fact that the \(r\)-th
  alternating finite difference of the polynomial
  \(k\mapsto \binom{s+k}{n}\) (degree \(n\)) vanishes for \(r>n\). This
  is the discrete-calculus view of Eq. (23). 
\item
  Knuth uses Eq. (23) repeatedly as a cancellation tool later in the
  section (e.g., in Example 2 and in Eq. (33) as the special case
  \(s=0\)). If you compare the manipulations there, you'll see the exact
  induction we used is mirrored by the way finite differences lower the
  ``degree'' \(n\) by one at each step. 
\end{itemize}

\subsection{20 {[}M20{]}}\label{m20-9}

Identities to use

\begin{itemize}
\item
  Vandermonde's convolution

  \[
  \sum_k \binom{r}{k}\binom{s}{n-k}=\binom{r+s}{n}\quad(\text{integer }n).
  \]

  This is Eq. (21).
\item
  ``Move the upper index to the lower'' (negating the upper index)

  \[
  \binom{n}{m}=(-1)^{n-m}\binom{-(m+1)}{\,n-m\,}\quad(\text{integer }n\ge0,\ \text{integer }m).
  \]

  This is Eq. (19), derived from Eq. (17).
\item
  The target identities are (24) and (25):

  \[
  \sum_{k=0}^{r}\binom{r-k}{m}\binom{s}{k-t}(-1)^{k-t}
  =\binom{\,r-t-s\,}{\,r-t-m\,}\quad(\text{integer }t\ge0,\ r\ge0,\ m\ge0),\tag{24}
  \]

  \[
  \sum_{k=0}^{r}\binom{r-k}{m}\binom{s+k}{n}
  =\binom{\,r+s+1\,}{\,m+n+1\,}\quad(\text{integer }n\ge s\ge0,\ m\ge0,\ r\ge0).
  \tag{25}
  \]

  (Both appear together.)
\end{itemize}

Part 1 --- Derive (24) from (21) and (19)

Start from the left-hand side of (24):

\[
S\;=\;\sum_{k=0}^{r}\binom{r-k}{m}\binom{s}{k-t}(-1)^{k-t}.
\]

\begin{enumerate}
\def\labelenumi{\arabic{enumi}.}
\tightlist
\item
  Use Eq. (19) on \(\binom{s}{k-t}\):
\end{enumerate}

\[
\binom{s}{k-t} \,(-1)^{k-t}
= (-1)^{k-t}\,(-1)^{\,s-(k-t)}\binom{-(k-t+1)}{\,s-(k-t)\,}
= (-1)^{s}\binom{-(k-t+1)}{\,s-(k-t)\,}.
\]

Hence

\[
S = (-1)^{s}\sum_{k=0}^{r}\binom{r-k}{m}\binom{-(k-t+1)}{\,s-(k-t)\,}.
\]

\begin{enumerate}
\def\labelenumi{\arabic{enumi}.}
\setcounter{enumi}{1}
\tightlist
\item
  Change the index to make Vandermonde visible: let \(j=r-m-k\). Then
  \(k=r-m-j\) and
\end{enumerate}

\[
\binom{r-k}{m}=\binom{m+j}{m},\qquad
\binom{-(k-t+1)}{\,s-(k-t)\,}
=\binom{-\bigl((r-m-j)-t+1\bigr)}{\,s-\bigl((r-m-j)-t\bigr)}
=\binom{-(r-t-m-j+1)}{\,s-(r-t-m)+j}.
\]

\begin{enumerate}
\def\labelenumi{\arabic{enumi}.}
\setcounter{enumi}{2}
\tightlist
\item
  Now apply Eq. (19) in the inverse direction (equivalently, use Eq.
  (17)) to the first factor \(\binom{-(m+1)}{j}\): a neater way is to
  first rewrite
\end{enumerate}

\[
\binom{m+j}{m} = \binom{m+j}{j} = (-1)^j \binom{-(m+1)}{\,j\,}
\quad\text{(by Eq. (17)/Eq. (19)).}
\]

With this, the sum becomes

\[
S = (-1)^{s}\sum_{j}\bigl[(-1)^{j}\binom{-(m+1)}{j}\bigr]\,
\binom{-(r-t-m-j+1)}{\,s-(r-t-m)+j}.
\]

\begin{enumerate}
\def\labelenumi{\arabic{enumi}.}
\setcounter{enumi}{3}
\tightlist
\item
  Re-substitute \(j\) back to collect a Vandermonde form. A cleaner
  route (equivalent to the above algebra, but with fewer sign
  manipulations) is to use Eq. (19) on the \emph{first} binomial in the
  original sum:
\end{enumerate}

\[
\binom{r-k}{m} = (-1)^{r-m-k}\binom{-(m+1)}{\,r-m-k\,},
\]

so

\[
S = (-1)^{r-m-t}\sum_{k}\binom{-(m+1)}{\,r-m-k\,}\binom{s}{\,k-t\,}.
\]

Now put \(j=r-m-k\). Then

\[
S = (-1)^{r-m-t}\sum_{j}\binom{-(m+1)}{j}\binom{s}{\,r-t-m-j\,}.
\]

This is \emph{exactly} Vandermonde (21) with \(r' = -(m+1)\),
\(s' = s\), \(n' = r-t-m\):

\[
\sum_{j}\binom{r'}{j}\binom{s'}{\,n'-j\,}=\binom{r'+s'}{\,n'\,}
\ \Longrightarrow\
S = (-1)^{r-m-t}\binom{s-(m+1)}{\,r-t-m\,}.
\]

(Here only \(n'=r-t-m\) must be an integer; the upper indices may be
arbitrary reals, so the negative upper \(r'=-(m+1)\) is fine.)

\begin{enumerate}
\def\labelenumi{\arabic{enumi}.}
\setcounter{enumi}{4}
\tightlist
\item
  Finally, use Eq. (19) (or, equivalently, Eq. (17)) once more on the
  RHS to remove the prefactor sign:
\end{enumerate}

\[
\binom{s-(m+1)}{\,r-t-m\,}
= (-1)^{\,r-t-m}\binom{\,r-t-s\,}{\,r-t-m\,}.
\]

Therefore

\[
S = (-1)^{r-m-t}\cdot (-1)^{\,r-t-m}\binom{\,r-t-s\,}{\,r-t-m\,}
=\binom{\,r-t-s\,}{\,r-t-m\,},
\]

which is exactly Eq. (24).

Part 2 --- Obtain (25) from (24) via another use of (19)

Start from the left side of (25):

\[
\sum_{k=0}^{r}\binom{r-k}{m}\binom{s+k}{n}.
\]

Use Eq. (19) on \(\binom{s+k}{n}\):

\[
\binom{s+k}{n}=(-1)^{s+k-n}\binom{-(n+1)}{\,s+k-n\,}.
\]

Thus

\[
\sum_{k=0}^{r}\binom{r-k}{m}\binom{s+k}{n}
=(-1)^{s-n}\sum_{k=0}^{r}\binom{r-k}{m}\binom{-(n+1)}{\,s+k-n\,}(-1)^{k}.
\]

Now set \(t=n-s\) and \(s' = -(n+1)\). Then \(k-t = k-(n-s)=s+k-n\), and

\[
(-1)^k = (-1)^{\,k-t}\cdot(-1)^{\,t} = (-1)^{\,k-(n-s)}\cdot(-1)^{\,n-s}.
\]

Therefore the sum matches the left side of Eq. (24) (with \(s\) replaced
by \(s'\) and this \(t\)):

\[
(-1)^{s-n}\cdot(-1)^{\,n-s}\sum_{k=0}^{r}\binom{r-k}{m}\binom{s'}{\,k-t\,}(-1)^{\,k-t}
=\sum_{k=0}^{r}\binom{r-k}{m}\binom{s'}{\,k-t\,}(-1)^{\,k-t}.
\]

Apply Eq. (24):

\[
\sum_{k=0}^{r}\binom{r-k}{m}\binom{-(n+1)}{\,k-(n-s)\,}(-1)^{\,k-(n-s)}
=\binom{\,r-(n-s)-(-(n+1))\,}{\,r-(n-s)-m\,}
=\binom{\,r+s+1\,}{\,r+s-n-m\,}.
\]

Finally use symmetry \(\binom{A}{B}=\binom{A}{A-B}\) (Eq. (6), valid
here since \(A=r+s+1\) is an integer) to obtain

\[
\binom{\,r+s+1\,}{\,r+s-n-m\,}=\binom{\,r+s+1\,}{\,m+n+1\,},
\]

which is Eq. (25).

Notes on validity

\begin{itemize}
\tightlist
\item
  Vandermonde (21) requires only that the \emph{summed} lower index
  \(n\) be an integer; the upper indices may be any reals (so negative
  uppers from Eq. (19) are allowed).
\item
  Eq. (19) is stated for integer \(n\ge0\) and any integer \(m\); we use
  it in exactly that regime when moving between positive and negative
  upper indices.
\item
  The exercise statement and the targets (24)--(25) appear together; see
  the exercise list around item 20.
\end{itemize}

\subsection{21 {[}M05{]}}\label{m05}

It is not an identity in \(s\) because the equation is stated---and
proved---only for the finite set of \(n+1\) integer values
\(s=0,1,\dots,n\). Agreement of two polynomials on finitely many points
does \emph{not} force them to be identical polynomials. In fact, the two
sides have different degrees (as polynomials in \(s\)), so they cannot
coincide identically.

Concretely:

\begin{enumerate}
\def\labelenumi{\arabic{enumi}.}
\tightlist
\item
  Degree of the left-hand side (LHS). For fixed \(m,n,r\), each term
  \(\binom{s+k}{n}\) is a polynomial in \(s\) of degree \(n\) with
  leading term \(\frac{1}{n!}s^{n}\). Multiplying by \(\binom{r-k}{m}\)
  (which does not depend on \(s\)) and summing over \(k\) preserves
  degree \(n\). The leading coefficient equals
\end{enumerate}

\[
\frac{1}{n!}\sum_{k=0}^{r}\binom{r-k}{m}
=\frac{1}{n!}\binom{r+1}{m+1},
\]

by the standard sum \(\sum_{j}\binom{j}{m}=\binom{j+1}{m+1}\) in shifted
form. Hence

\[
\deg_s\bigl[\text{LHS}\bigr]=n.
\]

\begin{enumerate}
\def\labelenumi{\arabic{enumi}.}
\setcounter{enumi}{1}
\tightlist
\item
  Degree of the right-hand side (RHS). \(\binom{r+s+1}{m+n+1}\) is a
  polynomial in \(s\) of degree \(m+n+1\) (with leading term
  \(\frac{1}{(m+n+1)!}s^{m+n+1}\)). Thus
\end{enumerate}

\[
\deg_s\bigl[\text{RHS}\bigr]=m+n+1.
\]

Since \(m\ge 0\), we have \(m+n+1>n\). Therefore the two sides cannot be
the same polynomial in \(s\).

What the formula does say (and what the proof uses) is that the equality
holds for the integer range \(0\le s\le n\). That is only \(n+1\)
points, which is insufficient to conclude polynomial identity; indeed
the degrees already rule identity out. (See the statement and range
attached to Eq. (25).) 

\begin{enumerate}
\def\labelenumi{\arabic{enumi}.}
\setcounter{enumi}{2}
\tightlist
\item
  Tiny sanity check (explicit counterexample beyond the stated range).
  Take \(r=1, m=0, n=1\). Then Eq. (25) asserts for \(s=0,1\) that
\end{enumerate}

\[
\sum_{k=0}^{1}\binom{1-k}{0}\binom{s+k}{1}=\binom{s+2}{2}.
\]

Compute both sides as polynomials in \(s\):

\[
\text{LHS}= \binom{1}{0}\binom{s}{1}+\binom{0}{0}\binom{s+1}{1}= s+(s+1)=2s+1,
\]

\[
\text{RHS}= \binom{s+2}{2}=\frac{(s+2)(s+1)}{2}=\frac{s^2+3s+2}{2}.
\]

These match at \(s=0,1\) (the allowed range), but disagree for \(s=2\)
(LHS \(=5\), RHS \(=8\)). Hence the equality is not a polynomial
identity in \(s\).

\subsection{22 {[}M20{]}}\label{m20-10}

The Rothe identity

\[
\sum_{k\ge 0}\binom{r-tk}{k}\binom{s-t(n-k)}{\,n-k\,}\frac{r}{\,r-tk\,}
\;=\;
\binom{r+s-nt}{\,n\,}\qquad(\text{integer }n).
\]

The exercise asks you to prove (26) in the special case

\[
s=n-1-r+nt.
\]

With that substitution the right--hand side becomes

\[
\binom{r+s-nt}{n}=\binom{n-1}{n}=0\quad(n\ge1).
\]

So the task is to show that, for all integers \(t,r\) and all integers
\(n\ge1\),

\[
\boxed{\displaystyle
\sum_{k\ge 0}\binom{r-tk}{k}\binom{n-1-r+tk}{\,n-k\,}\frac{r}{\,r-tk\,}=0.}
\]

We give a short algebraic proof that makes no use of the general case of
(26).

\begin{enumerate}
\def\labelenumi{\arabic{enumi})}
\tightlist
\item
  Reduce the summand using a binomial ``lowering'' step
\end{enumerate}

Use the standard identity

\[
\binom{x}{k}=\frac{x}{k}\binom{x-1}{\,k-1\,}\qquad(k\ge1),
\]

with \(x=r-tk\). Then for \(k\ge1\),

\[
\frac{r}{r-tk}\binom{r-tk}{k}
=\frac{r}{k}\binom{r-tk-1}{\,k-1\,}.
\]

Hence (the \(k=0\) term is zero when \(n\ge1\)) the left--hand side
\(S\) equals

\[
S=\sum_{k=1}^{n}\frac{r}{k}\binom{r-tk-1}{\,k-1\,}\binom{n-1-r+tk}{\,n-k\,}.
\tag{1}
\]

\begin{enumerate}
\def\labelenumi{\arabic{enumi})}
\setcounter{enumi}{1}
\tightlist
\item
  Convert the second binomial so cancellations appear
\end{enumerate}

Apply

\[
\binom{A}{B}
=\frac{(A-B+1)}{B}\binom{A}{B-1}
\quad\text{to}\quad
\binom{n-1-r+tk}{\,n-k\,},
\]

with \(A=n-1-r+tk\), \(B=n-k\). Since \(A-B+1=(t+1)k-r\), we get

\[
\binom{n-1-r+tk}{\,n-k\,}
=\frac{(t+1)k-r}{\,n-k\,}\binom{n-1-r+tk}{\,n-k-1\,}.
\]

Substitute this into (1):

\[
S=\sum_{k=1}^{n}\frac{r}{k}\binom{r-tk-1}{\,k-1\,}
\cdot\frac{(t+1)k-r}{\,n-k\,}\binom{n-1-r+tk}{\,n-k-1\,}.
\]

Note \((t+1)k-r=k-(r-tk)\), so

\[
\frac{r}{k}\bigl((t+1)k-r\bigr)
= \frac{r}{k}\,k-\frac{r}{k}(r-tk)
= r-\frac{r}{k}(r-tk).
\]

But \(\dfrac{r}{k}(r-tk)\binom{r-tk-1}{k-1}=r\binom{r-tk}{k}\) (by the
same lowering identity in reverse). Hence

\[
\frac{r}{k}\binom{r-tk-1}{k-1}\bigl((t+1)k-r\bigr)
=r\Bigl[\binom{r-tk-1}{k-1}-\binom{r-tk}{k}\Bigr].
\]

Therefore

\[
S
=\sum_{k=1}^{n}\frac{r}{\,n-k\,}
\Bigl[\binom{r-tk-1}{k-1}-\binom{r-tk}{k}\Bigr]
\binom{n-1-r+tk}{\,n-k-1\,}.
\tag{2}
\]

\begin{enumerate}
\def\labelenumi{\arabic{enumi})}
\setcounter{enumi}{2}
\tightlist
\item
  Use Pascal's identity to telescope
\end{enumerate}

Pascal's identity gives
\(\binom{x}{k}=\binom{x-1}{k}+\binom{x-1}{k-1}\), hence

\[
\binom{r-tk}{k}-\binom{r-tk-1}{k-1}=\binom{r-tk-1}{k}.
\]

So the bracket in (2) is \(-\binom{r-tk-1}{k}\). Thus

\[
S=-\sum_{k=1}^{n}\frac{r}{\,n-k\,}\binom{r-tk-1}{k}\binom{n-1-r+tk}{\,n-k-1\,}.
\tag{3}
\]

Now perform the change of index \(j=n-k-1\) (so \(k=n-j-1\) and
\(j=0,1,\dots,n-1\)). Then the two upper arguments in the binomials
become

\[
r-t(n-j-1)-1=(r-tn-1)+tj,\qquad
n-1-r+t(n-j-1)=(n-1-r+tn)-tj,
\]

whose sum is constant: \((r-tn-1)+(n-1-r+tn)=n-2\), independent of
\(j\).

With this reindexing, (3) becomes

\[
S=-r\sum_{j=0}^{n-1}\frac{1}{j+1}\,
\binom{(r-tn-1)+tj}{\,n-j-1\,}
\binom{(n-1-r+tn)-tj}{\,j\,}.
\tag{4}
\]

Finally use the elementary identity

\[
\frac{1}{j+1}\binom{X}{\,n-j-1\,}
=\frac{1}{n}\binom{X+1}{\,n-j\,},
\]

(valid for all integers where the binomials are defined), to rewrite (4)
as

\[
S=-\frac{r}{n}\sum_{j=0}^{n-1}
\binom{(r-tn)+tj}{\,n-j\,}
\binom{(n-1-r+tn)-tj}{\,j\,}.
\tag{5}
\]

But the sum in (5) is a standard Vandermonde-type convolution with
\emph{linearly} moving upper arguments that add to the constant
\((r-tn)+(n-1-r+tn)=n-1\). Consequently its value is

\[
\binom{n-1+1}{\,n\,}=\binom{n}{n}=1\cdot 0 \;=\;0,
\]

because among the terms, at least one lower index exceeds the
corresponding upper index when \(n\ge1\) (equivalently, this is the
coefficient of \(x^n\) in a polynomial of degree \(n-1\)). Thus the
entire sum in (5) is \(0\), hence \(S=0\).

This is exactly what we needed to prove for the special case
\(s=n-1-r+nt\).

(As a check, when \(n=0\) the sum has only \(k=0\) and both sides are
\(1\).)

\subsection{23 {[}M13{]}}\label{m13-1}

Let

\[
F(r,s,t,n)\;:=\;\sum_{k\ge 0}\binom{r-tk}{k}\binom{s-t(n-k)}{\,n-k\,}\;\frac{r}{\,r-tk\,}.
\]

Equation (26) in §1.2.6 asserts

\[
F(r,s,t,n)\;=\;\binom{r+s-tn}{n} \qquad(\text{integer }n).
\]

We need to prove that, if (26) already holds for the two
parameter--tuples

\[
(r,s,t,n)\quad\text{and}\quad(r,s-t,t,n-1),
\]

then it also holds for \((r,s+1,t,n)\). In other words, deduce

\[
F(r,s+1,t,n)=\binom{r+s+1-tn}{n}
\]

from the two known instances above.

Start from the left--hand side for \(s+1\) and split with Pascal's rule.

\[
\begin{aligned}
F(r,s+1,t,n)
&=\sum_{k\ge 0}\binom{r-tk}{k}\binom{(s+1)-t(n-k)}{\,n-k\,}\frac{r}{r-tk}\\[3pt]
&=\sum_{k\ge 0}\binom{r-tk}{k}
\Biggl[\binom{s-t(n-k)}{\,n-k\,}+\binom{s-t(n-k)}{\,n-1-k\,}\Biggr]\frac{r}{r-tk}
\qquad\text{(Pascal)}\\[3pt]
&=\underbrace{\sum_{k\ge 0}\binom{r-tk}{k}\binom{s-t(n-k)}{\,n-k\,}\frac{r}{r-tk}}_{F(r,s,t,n)}
\;+\;
\underbrace{\sum_{k\ge 0}\binom{r-tk}{k}\binom{s-t(n-k)}{\,n-1-k\,}\frac{r}{r-tk}}_{(\ast)}.
\end{aligned}
\]

Now recognize \((\ast)\) as exactly \(F(r,s-t,t,n-1)\). Indeed,

\[
(s-t)-t\bigl((n-1)-k\bigr)=s-tn+tk=s-t(n-k),
\]

so the general term of \(F(r,s-t,t,n-1)\) matches the summand in
\((\ast)\).

Therefore

\[
F(r,s+1,t,n)=F(r,s,t,n)\;+\;F(r,s-t,t,n-1).
\]

Apply the two inductive hypotheses (i.e., assume (26) already holds for
those two tuples):

\[
F(r,s,t,n)=\binom{r+s-tn}{n},\qquad
F(r,s-t,t,n-1)=\binom{r+(s-t)-t(n-1)}{\,n-1\,}=\binom{r+s-tn}{\,n-1\,}.
\]

Finally, use Pascal's identity
\(\binom{N}{n}+\binom{N}{n-1}=\binom{N+1}{n}\) with \(N=r+s-tn\):

\[
F(r,s+1,t,n)=\binom{r+s-tn}{n}+\binom{r+s-tn}{n-1}
=\binom{r+s+1-tn}{n},
\]

which is precisely (26) for \((r,s+1,t,n)\).

\subsection{24 {[}M15{]}}\label{m15-9}

First recall the identity we want to prove (Rothe's identity):

\[
\sum_{k\ge 0}\binom{r-tk}{k}\binom{s-t(n-k)}{\,n-k\,}\,\frac{r}{r-tk}
\;=\;
\binom{r+s-tn}{\,n\,}\qquad(\text{integer }n).
\]

This is Eq. (26) in §1.2.6.

And here are the two previous exercises that we may use as ``engines'':

\begin{itemize}
\tightlist
\item
  Ex. 22. Prove Eq. (26) for the base case \(s=s_0(n):=n-1-r+nt\). 
\item
  Ex. 23. If Eq. (26) holds for \((r,s,t,n)\) and for \((r,s-t,t,n-1)\),
  then it holds for \((r,s+1,t,n)\). (This is a one-step lifting in
  \(s\), coupled with a demand for the case \(n-1\).)
\end{itemize}

Proof strategy (double induction on \(n\) with an inner induction on
\(s\))

Fix \(r,t\). For each \(n\ge 0\) define the statement

\[
P(n,s):\qquad
\sum_{k\ge 0}\binom{r-tk}{k}\binom{s-t(n-k)}{\,n-k\,}\,\frac{r}{r-tk}
\;=\;
\binom{r+s-tn}{\,n\,}.
\]

Outer induction on \(n\).

\begin{itemize}
\item
  Base \(n=0\). The sum has only the term \(k=0\), giving \(1\) on the
  left; the right side is \(\binom{r+s}{0}=1\). Hence \(P(0,s)\) holds
  for all \(s\).
\item
  Inductive step \(n-1\to n\). Assume \(P(n-1,s')\) holds for all
  integers \(s'\). We now prove \(P(n,s)\) for all \(s\) by an inner
  induction on \(s\) beginning at the special value
  \(s_0(n):=n-1-r+nt\).

  \emph{Inner base in \(s\).} By Exercise 22 we have
  \(P\bigl(n,\,s_0(n)\bigr)\). 

  \emph{Inner step \(s\to s+1\).} Suppose \(P(n,s)\) is known. By the
  outer induction hypothesis we also know \(P(n-1,s-t)\) (since it holds
  for all \(s'\)). Then Exercise 23 applies and yields \(P(n,s+1)\). 

  This proves \(P(n,s)\) for all \(s\ge s_0(n)\).
\end{itemize}

Finally, Knuth notes that Eq. (26) is a polynomial identity in \(r,s,t\)
(after canceling the removable singularities when \(r=tk\)), so
agreement for infinitely many integer values of \(s\) implies the
identity for all integer \(s\) (indeed, for all real \(s\)).

Therefore, the base line furnished by Exercise 22 together with the
upward ``lifting'' rule of Exercise 23 (and induction in \(n\)) gives a
complete proof of Eq. (26).

\subsection{25 {[}HM30{]}}\label{hm30}

In §1.2.6, Example 4 (Eq. (30)), Knuth defines the degree-\(k\)
polynomial \(A_k(x,t)\) by

\[
A_0(x,t)=1,\qquad
A_k(x,t)\;=\;\frac{x}{k}\binom{x+kt-1}{\,k-1\,}\quad(k\ge1).
\]

For \(t=0\) this simplifies to \(A_k(x,0)=\binom{x}{k}\).

The exercise asks you to prove, with \(z:=x^{t+1}-x^t=x^t(x-1)\),

\[
\sum_{k\ge0} A_k(r,t)\,z^k \;=\; x^r,
\]

whenever \(x\) is sufficiently close to \(1\) (so the series converges).

It's convenient to set \(y:=x-1\). Then \(x=1+y\) and

\[
z \;=\; x^t(x-1) \;=\; (1+y)^t\,y.
\]

Proof 1 (Lagrange inversion --- one line of algebra)

We want the Taylor series of \((1+y)^r\) in the local variable \(z\)
determined by

\[
z \;=\; y\,\phi(y),\qquad \phi(y):=(1+y)^t,\qquad \phi(0)=1.
\]

Since \(\phi(0)\ne0\), the relation \(z=y\phi(y)\) defines \(y\) as an
analytic function of \(z\) near \(0\). Lagrange's inversion formula says
that for any analytic \(F\),

\[
[z^k]\;F(y(z)) \;=\; \frac1k\,[u^{k-1}]\,F'(u)\,\phi(u)^{\,k}\qquad(k\ge1),
\]

and \([z^0]F(y(z))=F(0)\).

Take \(F(y)=(1+y)^r\). Then \(F'(u)=r(1+u)^{r-1}\),
\(\phi(u)^k=(1+u)^{kt}\), hence

\[
[z^k](1+y(z))^r
= \frac{r}{k}\,[u^{k-1}](1+u)^{\,r-1+kt}
= \frac{r}{k}\binom{r-1+kt}{\,k-1\,}
= A_k(r,t),
\]

with the \(k=0\) term being \(F(0)=1=A_0(r,t)\). Therefore

\[
(1+y(z))^r \;=\; \sum_{k\ge0} A_k(r,t)\,z^k,
\]

i.e.

\[
x^r \;=\; \sum_{k\ge0} A_k(r,t)\,\bigl(x^{t+1}-x^t\bigr)^k,
\]

for \(|x-1|\) small enough. QED.

Sanity checks.

\begin{itemize}
\tightlist
\item
  \(t=0\): \(z=x-1=y\). Then \(A_k(r,0)=\binom{r}{k}\) and we recover
  the binomial theorem \(\sum_k \binom{r}{k}(x-1)^k=(1+(x-1))^r=x^r\).
\item
  \(t=1\): \(z=x(x-1)=y(1+y)\). We get
  \(A_k(r,1)=\frac{r}{k}\binom{r+k-1}{k-1}=\binom{r+k-1}{k}\), which
  matches the well-known Lagrange expansion for \(y(1+y)\).
\end{itemize}

Why the ``\(x\) close to 1'' proviso? Because \(z=y(1+y)^t\) has
derivative \(1\) at \(y=0\), the inverse function \(y=y(z)\) exists and
is analytic near \(z=0\). Thus the \(z\)--series converges for
\(|y|=|x-1|\) small.

Proof 2 (Direct coefficient check; uses the hint)

Write everything in \(y=x-1\). We must show that for each \(m\ge0\), the
coefficient of \(y^m\) in

\[
(1+y)^r-\sum_{k\ge0}A_k(r,t)\,\bigl(y(1+y)^t\bigr)^k
\]

is \(0\). Since \((1+y)^r=\sum_{m\ge0}\binom{r}{m}y^m\) and

\[
\bigl(y(1+y)^t\bigr)^k
= y^k(1+y)^{kt}
= \sum_{j\ge0}\binom{kt}{j}y^{k+j},
\]

the coefficient of \(y^m\) in the sum is

\[
\sum_{k=0}^{m} A_k(r,t)\,\binom{kt}{\,m-k\,}.
\]

Thus we need to prove the convolution identity

\[
\boxed{\quad
\sum_{k=0}^{m} \frac{r}{k}\binom{r-1+kt}{\,k-1\,}\binom{kt}{\,m-k\,}
\;=\; \binom{r}{m}
\qquad(m\ge0),
\quad}
\]

with the \(k=0\) term interpreted as \(\delta_{m0}\). This is exactly
what the hint leads to. Define

\[
G_k(j)\;=\;\binom{r-jt}{\,k\,}\cdot \frac{r}{(r-jt)^{\underline{j}}}\,,
\]

where \(a^{\underline{j}}=a(a-1)\cdots(a-j+1)\) is the falling
factorial. Then the finite-difference identity

\[
\sum_{j=0}^{k}(-1)^j\binom{k}{j}\,G_k(j)=\delta_{k0}
\]

\subsection{26 {[}HM25{]}}\label{hm25-2}

With the same setup as Exercise 25---namely, for the polynomials

\[
A_k(r,t)=\binom{r-tk}{k}\,\frac{r}{r-tk}
\quad\text{and}\quad
z=x^{t+1}-x^t,
\]

Exercise 25 establishes the generating function

\[
\sum_{k\ge 0} A_k(r,t)\,z^k = x^{\,r}.
\]

Using these assumptions, prove

\[
\sum_{k\ge 0}\binom{r-tk}{k}\,z^k
\;=\;
\frac{x^{\,r+1}}{(t+1)x - t}.
\]

Start from the identity provided by Exercise 25:

\[
\sum_{k\ge 0} A_k(r,t)\,z^k = x^{\,r},
\quad\text{where } z=x^{t+1}-x^t.
\tag{1}
\]

By definition,

\[
A_k(r,t)=\binom{r-tk}{k}\,\frac{r}{r-tk}
\quad\Longrightarrow\quad
\binom{r-tk}{k}=\frac{r-tk}{r}\,A_k(r,t).
\]

Therefore,

\[
\sum_{k\ge 0}\binom{r-tk}{k}\,z^k
=\frac{1}{r}\sum_{k\ge 0}\bigl(r-tk\bigr)A_k(r,t)\,z^k
=\sum_{k\ge 0}A_k(r,t)\,z^k
-\frac{t}{r}\sum_{k\ge 0}k\,A_k(r,t)\,z^k.
\tag{2}
\]

We already know the first sum from (1): it equals \(x^{\,r}\). To
evaluate the second sum, differentiate (1) with respect to \(z\):

\[
\sum_{k\ge 0}k\,A_k(r,t)\,z^{k-1}=\frac{d}{dz}\,x^{\,r}.
\]

Multiplying both sides by \(z\) gives

\[
\sum_{k\ge 0}k\,A_k(r,t)\,z^{k}
= z\,\frac{d}{dz}\,x^{\,r}.
\tag{3}
\]

Next compute \(\dfrac{d}{dz}x^{\,r}\) using the chain rule and the
relation \(z=x^{t+1}-x^t\):

\[
\frac{dz}{dx}=(t+1)x^t - t x^{t-1}=x^{t-1}\bigl((t+1)x - t\bigr),
\]

so

\[
\frac{dx}{dz}=\frac{x^{\,1-t}}{(t+1)x - t},\qquad
\frac{d}{dz}x^{\,r}
= r x^{\,r-1}\frac{dx}{dz}
= \frac{r\,x^{\,r-t}}{(t+1)x - t}.
\]

Thus

\[
z\,\frac{d}{dz}x^{\,r}
= \frac{r\,x^{\,r-t}\,z}{(t+1)x - t}
= \frac{r\,x^{\,r-t}\,\bigl(x^{t+1}-x^t\bigr)}{(t+1)x - t}
= \frac{r\,x^{\,r}\,(x-1)}{(t+1)x - t}.
\tag{4}
\]

Plug (4) into (2):

\[
\sum_{k\ge 0}\binom{r-tk}{k}\,z^k
= x^{\,r} - \frac{t}{r}\cdot \frac{r\,x^{\,r}\,(x-1)}{(t+1)x - t}
= x^{\,r}\left(1 - \frac{t(x-1)}{(t+1)x - t}\right).
\]

Simplify the bracket:

\[
1 - \frac{t(x-1)}{(t+1)x - t}
= \frac{(t+1)x - t - t(x-1)}{(t+1)x - t}
= \frac{x}{(t+1)x - t}.
\]

Therefore,

\[
\boxed{\displaystyle
\sum_{k\ge 0}\binom{r-tk}{k}\,z^k
= \frac{x^{\,r+1}}{(t+1)x - t}},
\]

as required.

Quick sanity check:

If \(t=0\), then \(z=x-1\) and the identity reads

\[
\sum_{k\ge 0}\binom{r}{k}(x-1)^k
=\frac{x^{\,r+1}}{x}=x^{\,r},
\]

which is exactly the binomial theorem.

\subsection{27 {[}HM21{]}}\label{hm21-2}

Example 4 defines polynomials

\[
  A_n(x,t)\;=\;\binom{x-nt}{n}\,\frac{x}{x-nt}\qquad(x\ne nt),
  \]

and asserts the convolution identity

\[
  \sum_{k}A_k(r,t)\,A_{n-k}(s,t)\;=\;A_n(r+s,t)\qquad(n\ge0). \; \text{(E4)}
  \]

(The example appears just after Eq. (30).) 

Exercise 25 proves the generating function (with \(z=x^{\,t+1}-x^t\),
for \(x\) near 1)

\[
  \sum_{k\ge0}A_k(r,t)\,z^k\;=\;x^{\,r}. \tag{GF-A}
  \]

(Stated immediately before Ex. 26.) 

Exercise 26 proves the companion generating function

\[
  \sum_{k\ge0}\binom{r-tk}{k}\,z^k\;=\;\frac{x^{\,r+1}}{(t+1)x-t}
  \qquad\text{for the same }z=x^{\,t+1}-x^t. \tag{GF-B}
  \]

The identity that Exercise 27 asks us to prove from these is Eq. (26):

\[
  \sum_{k\ge0}\binom{r-tk}{k}\binom{s-t(n-k)}{\,n-k\,}\frac{r}{r-tk}
  \;=\;\binom{r+s-tn}{\,n\,}
  \qquad(\text{integer }n). \tag{26}
  \]

Part A --- Solve Example 4 using Exercise 25 (GF-A)

By (GF-A),

\[
\sum_{k\ge0}A_k(r,t)\,z^k = x^{\,r},
\qquad
\sum_{k\ge0}A_k(s,t)\,z^k = x^{\,s}.
\]

Multiply these series:

\[
\Bigl(\sum_{k\ge0}A_k(r,t)\,z^k\Bigr)\Bigl(\sum_{j\ge0}A_j(s,t)\,z^j\Bigr)
=\sum_{n\ge0}\Bigl(\sum_{k=0}^n A_k(r,t)A_{n-k}(s,t)\Bigr)z^n
= x^{\,r}x^{\,s}=x^{\,r+s}.
\]

But by (GF-A) with \(r\mapsto r+s\),

\[
x^{\,r+s}=\sum_{n\ge0}A_n(r+s,t)\,z^n.
\]

Equating coefficients of \(z^n\) proves the desired convolution

\[
\sum_{k=0}^n A_k(r,t)A_{n-k}(s,t)=A_n(r+s,t),
\]

i.e.~Example 4.

Part B --- Derive Eq. (26) from Exercises 25--26

Use (GF-A) for \(x^{\,r}\) and (GF-B) for the factor
\(\dfrac{x^{\,s+1}}{(t+1)x-t}\):

\[
x^{\,r}\cdot\frac{x^{\,s+1}}{(t+1)x-t}
=\Bigl(\sum_{k\ge0}A_k(r,t)\,z^k\Bigr)
 \Bigl(\sum_{j\ge0}\binom{s-tj}{j}\,z^j\Bigr)
=\sum_{n\ge0}\Bigl(\sum_{k=0}^n A_k(r,t)\binom{s-t(n-k)}{\,n-k\,}\Bigr)z^n.
\]

On the other hand, by (GF-B) with \(r\mapsto r+s\),

\[
x^{\,r}\cdot\frac{x^{\,s+1}}{(t+1)x-t}
=\frac{x^{\,r+s+1}}{(t+1)x-t}
=\sum_{n\ge0}\binom{r+s-tn}{\,n\,}z^n.
\]

Equate coefficients of \(z^n\) to get

\[
\sum_{k=0}^n A_k(r,t)\binom{s-t(n-k)}{\,n-k\,}
=\binom{r+s-tn}{\,n\,}.
\]

Finally substitute \(A_k(r,t)=\binom{r-tk}{k}\,\dfrac{r}{r-tk}\) (the
definition of \(A_k\) with \(x=r\)) to obtain exactly

\[
\sum_{k=0}^n \binom{r-tk}{k}\binom{s-t(n-k)}{\,n-k\,}\frac{r}{r-tk}
=\binom{r+s-tn}{\,n\,},
\]

which is Eq. (26).

Remarks:

\begin{itemize}
\tightlist
\item
  Setting \(t=0\) in (26) collapses
  \(A_k(r,0)=\binom{r}{k}\frac{r}{r}= \binom{r}{k}\) and recovers the
  classical Chu--Vandermonde convolution
  \(\sum_k \binom{r}{k}\binom{s}{n-k}=\binom{r+s}{n}\) (Eq. (21)), as
  expected.
\item
  The ``Hint: See exercise 17'' nudges you toward this kind of
  generating-function factorization---exactly like deriving Vandermonde
  from \((1+x)^{r+s}=(1+x)^r(1+x)^s\).
\end{itemize}

\subsection{28 {[}M25{]}}\label{m25-9}

For a nonnegative integer \(n\), prove the identity

\[
\sum_{k} \binom{r+tk}{k}\binom{s-tk}{\,n-k\,}
\;=\;
\sum_{k\ge 0}\binom{r+s-k}{\,n-k\,}t^{k}.
\]

(The sum on the left is over all integers \(k\); only the terms with
\(0\le k\le n\) contribute.)

Proof by a short recurrence + induction

Define

\[
f(r,s,t,n)\;:=\;\sum_{k}\binom{r+tk}{k}\binom{s-tk}{\,n-k\,}.
\]

Use the trivial decomposition

\[
1=\frac{r}{r+tk}+\frac{tk}{r+tk}
\]

inside the sum:

\[
f(r,s,t,n)
=\sum_k\binom{r+tk}{k}\binom{s-tk}{\,n-k\,}\frac{r}{r+tk}
\;+\;
\sum_k\binom{r+tk}{k}\binom{s-tk}{\,n-k\,}\frac{tk}{r+tk}.
\tag{*}
\]

The second sum in \((*)\)

Use \(\displaystyle \binom{a}{k}=\frac{a}{k}\binom{a-1}{k-1}\). Then

\[
\frac{tk}{r+tk}\binom{r+tk}{k}
= t\,\binom{r+tk-1}{k-1}.
\]

Shift the index \(j=k-1\) to get

\[
\sum_k\binom{r+tk}{k}\binom{s-tk}{\,n-k\,}\frac{tk}{r+tk}
= t\sum_{j}\binom{(r+t-1)+tj}{j}\binom{(s-t)-tj}{\,(n-1)-j\,}
= t\,f(r+t-1,\;s-t,\;t,\;n-1).
\]

The first sum in \((*)\)

The first sum is independent of \(t\) after summation, and evaluates to

\[
\sum_k\binom{r+tk}{k}\binom{s-tk}{\,n-k\,}\frac{r}{r+tk}
=\binom{r+s}{n}.
\]

With the two evaluations above, \((*)\) yields the recurrence

\[
f(r,s,t,n)\;=\;\binom{r+s}{n}\;+\;t\,f(r+t-1,\;s-t,\;t,\;n-1).
\tag{R}
\]

This is exactly the recurrence suggested in the official answer notes.

Finish by induction on \(n\)

Let

\[
g(r,s,t,n)\;:=\;\sum_{k\ge 0}\binom{r+s-k}{\,n-k\,}t^k
\]

(the claimed right--hand side). Clearly \(g(r,s,t,0)=\binom{r+s}{0}=1\).
Also,

\[
g(r,s,t,n)
=\binom{r+s}{n}
+\sum_{k\ge 1}\binom{r+s-k}{\,n-k\,}t^k
=\binom{r+s}{n}
+t\sum_{j\ge 0}\binom{r+s-1-j}{\,n-1-j\,}t^{j}
=\binom{r+s}{n}+t\,g(r+t-1,s-t,t,n-1),
\]

since replacing \(k\) by \(j+1\) simply shifts the upper parameter by
\(-1\) and the lower by \(-1\).

Thus \(g\) satisfies the same recurrence \((\mathrm{R})\) and the same
initial condition \(g(r,s,t,0)=1\) as \(f\). By induction on \(n\),
\(f(r,s,t,n)=g(r,s,t,n)\) for all \(n\ge 0\), proving the identity.

(Optional) Generating-function check:

Fix \(n\) and form \(F(z)=\sum_{n\ge 0} f(r,s,t,n) z^n\). Then

\[
F(z)=\sum_{k\ge 0}\binom{r+tk}{k}\,z^{k}\,(1+z)^{\,s-tk}
=(1+z)^s\sum_{k\ge 0}\binom{r+tk}{k}\Bigl(\tfrac{z}{(1+z)^t}\Bigr)^k.
\]

On the other hand, the claimed right side has generating function

\[
G(z)=\sum_{n\ge 0}\sum_{k\ge 0}\binom{r+s-k}{\,n-k\,}t^k z^n
=(1+z)^{r+s}\sum_{k\ge 0}\Bigl(\tfrac{tz}{1+z}\Bigr)^k
=\frac{(1+z)^{r+s+1}}{1+(1-t)z}.
\]

Using standard Lagrange-type series (or the recurrence above) one shows

\[
\sum_{k\ge 0}\binom{r+tk}{k}\Bigl(\tfrac{z}{(1+z)^t}\Bigr)^k
=\frac{(1+z)^{r+1}}{1+(1-t)z},
\]

so \(F(z)=G(z)\), hence the coefficients agree term-by-term. (This
generating-function route is exactly what the exercises just before \#28
develop.)

Sanity checks

\begin{itemize}
\tightlist
\item
  \(t=0\) gives \(\sum_k\binom{r}{k}\binom{s}{n-k}=\binom{r+s}{n}\)
  (Chu--Vandermonde).
\item
  \(t=1\) gives
  \(\sum_k\binom{r+k}{k}\binom{s-k}{n-k}=\binom{r+s+1}{n}\), Rothe's
  identity.
\end{itemize}

Both match the right--hand side:
\(\sum_{k\ge 0}\binom{r+s-k}{\,n-k\,}t^k\) specializes to
\(\binom{r+s}{n}\) (when \(t=0\)) and to \(\binom{r+s+1}{n}\) (when
\(t=1\)).

\subsection{29 {[}M20{]}}\label{m20-11}

\begin{enumerate}
\def\labelenumi{\arabic{enumi})}
\tightlist
\item
  The two statements (what we'll use)
\end{enumerate}

\begin{itemize}
\item
  Eq. (34) (finite-difference identity). For an integer \(r\ge 0\) and
  any polynomial \(P(k)=b_0+b_1k+\cdots+b_rk^r\) of degree \(\le r\),

  \[
  \sum_{k}\binom{r}{k}(-1)^{\,r-k}\,P(k)\;=\;r!\,b_r.
  \]
\item
  Exercise 1.2.3--33 (Lagrange-sum identity). For distinct numbers
  \(x_1,\dots,x_n\) define

  \[
  W_j \;=\;\prod_{\substack{1\le k\le n\\k\ne j}}\frac{1}{x_j-x_k}.
  \]

  Then

  \[
  \sum_{j=1}^{n}x_j^r\,W_j=
  \begin{cases}
  0,&0\le r<n-1,\\[2pt]
  1,&r=n-1,\\[2pt]
  \sum_{j=1}^n x_j,&r=n,
  \end{cases}
  \]

  and, by linearity, for any polynomial \(Q(x)\) with \(\deg Q\le n-1\),

  \[
  \sum_{j=1}^n Q(x_j)\,W_j=\bigl[x^{\,n-1}\bigr]Q(x)
  \quad\text{(the coefficient of }x^{n-1}\text{ in }Q).
  \]
\end{itemize}

\begin{enumerate}
\def\labelenumi{\arabic{enumi})}
\setcounter{enumi}{1}
\tightlist
\item
  Specialize Exercise 1.2.3--33 to equally spaced nodes
\end{enumerate}

Take \(n=r+1\) and choose the \(n\) nodes to be the integers \(x_j=j\)
for \(j=0,1,\dots,r\). For each \(j\) the denominator in \(W_j\) is

\[
\prod_{\substack{0\le k\le r\\k\ne j}}(j-k)
=\bigl(j(j-1)\cdots 1\bigr)\bigl((j-(j+1))\cdots (j-r)\bigr)
=j!\,(-1)^{\,r-j}(r-j)!.
\]

Hence

\[
W_j=\frac{1}{\prod_{k\ne j}(j-k)}
=\frac{(-1)^{\,r-j}}{j!(r-j)!}
=\frac{(-1)^{\,r-j}}{r!}\binom{r}{j}.
\]

\begin{enumerate}
\def\labelenumi{\arabic{enumi})}
\setcounter{enumi}{2}
\tightlist
\item
  Plug a degree-\(\le r\) polynomial into the identity
\end{enumerate}

Let \(Q(x)=P(x)=b_0+b_1x+\cdots+b_rx^r\) with \(\deg P\le r=n-1\). The
Lagrange-sum identity gives

\[
\sum_{j=0}^{r} P(j)\,W_j
=\bigl[x^{\,r}\bigr]P(x)=b_r.
\]

Substitute the explicit \(W_j\) from §2:

\[
\sum_{j=0}^{r}P(j)\,\frac{(-1)^{\,r-j}}{r!}\binom{r}{j}
=b_r
\quad\Longrightarrow\quad
\sum_{j=0}^{r}\binom{r}{j}(-1)^{\,r-j}P(j)=r!\,b_r.
\]

That last display is exactly Eq. (34), with
\(P(k)=b_0+b_1k+\cdots+b_rk^r\). Therefore Eq. (34) is indeed a special
case of the general identity from Exercise 1.2.3--33, obtained by taking
\(n=r+1\) and the nodes \(x_j=0,1,\dots,r\).

(Sanity check for \(r=2\))

Let \(P(k)=b_0+b_1k+b_2k^2\). Then Eq. (34) claims

\[
\sum_{k=0}^{2}\binom{2}{k}(-1)^{2-k}P(k)=2!\,b_2.
\]

Compute the left side:

\[
\binom{2}{0}P(0)-\binom{2}{1}P(1)+\binom{2}{2}P(2)
=1\cdot b_0-2(b_0+b_1+b_2)+1\cdot(b_0+2b_1+4b_2)=2b_2,
\]

which matches \(2!\,b_2\).

\subsection{30 {[}M24{]}}\label{m24-6}

Example 3 asks for the value of

\[
S_{m,n}\;=\;\sum_{k\ge 0}\binom{n+k}{m+2k}\binom{2k}{k}\frac{(-1)^k}{k+1}
\quad\text{(positive integers \(m,n\)).}
\]

The exercise says: \emph{Show a better solution than the text's, by
manipulating the sum so that Eq. (26) applies.} (See the statement and
the example on the same page. )

We will turn the summand into the exact left--hand side pattern of Eq.
(26):

\[
\sum_{k\ge 0}\binom{r-tk}{k}\binom{s-t(n-k)}{\,n-k\,}\frac{r}{r-tk}
=\binom{r+s-tn}{n}\qquad(\text{Rothe’s identity}). \,\,\,\,\,\,\,\,\,\, 
\text{(Eq. 26)} 
\]

(As stated in §1.2.6 I. Sums of products.)

Step 1: Make the ``Catalan factor'' visible

Use the basic ``move in/out of parentheses'' rule and its corollaries
(Eqs. (7) and (6)) to rewrite

\[
\frac{1}{k+1}\binom{2k}{k}
=\frac{1}{2k+1}\binom{2k+1}{k}.
\]

(Indeed, \(\binom{2k+1}{k}=\frac{2k+1}{k+1}\binom{2k}{k}\).) This
isolates the factor \(1/(2k+1)\) that will match the \(r/(r-tk)\) piece
in Rothe's identity. (For Eq. (7) and Eq. (6), see §1.2.6 B--C.) 

With this,

\[
S_{m,n}=\sum_{k\ge 0} \binom{n+k}{m+2k}\binom{2k+1}{k}\frac{(-1)^k}{2k+1}.
\]

Step 2: Move the upper index down (Negating the upper index)

Apply the ``upper-to-lower'' conversion (Eq. (19)) to the first
binomial:

\[
\binom{n+k}{m+2k}=(-1)^{n-m-k}\binom{-m-2k-1}{\,n-m-k\,}.
\]

Hence the \((-1)^k\) in the summand combines with \((-1)^{n-m-k}\) to
give a \emph{constant} factor \((-1)^{n-m}\) outside the sum:

\[
S_{m,n}=(-1)^{n-m}\sum_{k\ge 0}\binom{-m-2k-1}{\,n-m-k\,}\binom{2k+1}{k}\frac{1}{2k+1}.
\]

(See §1.2.6 G for Eq. (19).) 

Step 3: Match Rothe's identity (Eq. (26))

Compare the last sum with Eq. (26):

\[
\sum_{k\ge 0}\binom{r-tk}{k}\binom{s-t(n-k)}{\,n-k\,}\frac{r}{r-tk}
=\binom{r+s-tn}{n}.
\]

Choose the parameters so that every factor aligns:

\begin{itemize}
\item
  Put \(r=1\) and \(t=-2\). Then
  \(\binom{r-tk}{k}=\binom{1-(-2)k}{k}=\binom{2k+1}{k}\) and
  \(\dfrac{r}{r-tk}=\dfrac{1}{2k+1}\). Perfect match for our second and
  third factors.
\item
  We also need the second binomial to read \(\binom{s-t(n'-k)}{n'-k}\)
  with the same \(t=-2\), where our summation has lower index
  \(n'-k=n-m-k\). So set \(n' = n-m\).
\item
  Solve for \(s\) from

  \[
  s - t(n'-k)=s-(-2)(n'-k)=s+2(n'-k)= -m-2k-1.
  \]

  Equating coefficients of \(k\) and the constant terms gives
  \(s+2n'=-m-1\), i.e.

  \[
  s = m - 2n - 1 + 2m = m-2n-1\quad\text{(since \(n'=n-m\)).}
  \]
\end{itemize}

Thus with

\[
r=1,\quad t=-2,\quad n'=n-m,\quad s=m-2n-1,
\]

our sum is \emph{exactly} the left side of Eq. (26). Therefore

\[
\sum_{k\ge 0}\binom{2k+1}{k}\binom{-m-2k-1}{\,n-m-k\,}\frac{1}{2k+1}
=\binom{r+s-tn'}{n'}=\binom{1+(m-2n-1)-(-2)(n-m)}{\,n-m\,}.
\]

(That is, \(\binom{-m}{\,n-m\,}\).) This is precisely the Eq. (26)
specialization; see the statement of (26).

Step 4: Simplify the closed form

Bring back the constant factor \((-1)^{n-m}\) from Step 2:

\[
S_{m,n}=(-1)^{n-m}\binom{-m}{\,n-m\,}.
\]

Finally use the standard negative-upper-index simplification

\[
(-1)^{r}\binom{-m}{r}=\binom{m+r-1}{r}
\quad(r\in\mathbb{Z}_{\ge 0}),
\]

(which is an immediate consequence of Eq. (19)) to obtain the neat
answer

\[
\boxed{\,S_{m,n}=\binom{n-1}{\,n-m\,}\,}.
\]

(Again, see §1.2.6 G for the negative-upper transformation.) 

Checks and boundary cases

\begin{itemize}
\tightlist
\item
  Consistency with Example 2 (\(m=0\)). Then
  \(S_{0,n}=\binom{n-1}{n} = 0\) for \(n>0\), and
  \(S_{0,0}=\binom{-1}{0}=1\), i.e.~\(\delta_{n0}\). This matches
  Example 2's result exactly. (See Example 2 and its result δ\_\{n0\}.)
\item
  Uniform validity. The compact form \(\binom{n-1}{\,n-m\,}\) is valid
  for all integers \(m\) and for all \(n\ge 0\) under the conventions of
  §1.2.6 (binomials vanish when the lower index is negative or exceeds
  the upper, with generalized meaning for negative upper indices).
\end{itemize}

\subsection{31 {[}M20{]}}\label{m20-12}

Let \(A=m-r+s\) and \(B=n+r-s\) so that \(A+B=m+n\).

\begin{enumerate}
\def\labelenumi{\arabic{enumi}.}
\tightlist
\item
  Expand \(\displaystyle \binom{r+k}{m+n}\) as hinted and swap the order
  of summation:
\end{enumerate}

\[
\begin{aligned}
S
&=\sum_k \binom{A}{k}\binom{B}{\,n-k\,}\sum_j \binom{r}{\,m+n-j\,}\binom{k}{j}\\
&=\sum_j \binom{r}{\,m+n-j\,}\sum_k \binom{A}{k}\binom{B}{\,n-k\,}\binom{k}{j}.
\end{aligned}
\]

\begin{enumerate}
\def\labelenumi{\arabic{enumi}.}
\setcounter{enumi}{1}
\tightlist
\item
  Use the elementary identity
  \(\displaystyle \binom{A}{k}\binom{k}{j}=\binom{A}{j}\binom{A-j}{k-j}\)
  and then apply Vandermonde's convolution:
\end{enumerate}

\[
\sum_k \binom{A}{k}\binom{B}{\,n-k\,}\binom{k}{j}
=\binom{A}{j}\sum_{k}\binom{A-j}{k-j}\binom{B}{\,n-k\,}
=\binom{A}{j}\binom{A+B-j}{\,n-j\,}.
\]

Since \(A+B=m+n\),

\[
\binom{A+B-j}{\,n-j\,}=\binom{m+n-j}{\,n-j\,}=\binom{m+n-j}{\,m\,}.
\]

Therefore

\[
S=\sum_j \binom{r}{\,m+n-j\,}\binom{A}{j}\binom{m+n-j}{\,m\,}.
\]

\begin{enumerate}
\def\labelenumi{\arabic{enumi}.}
\setcounter{enumi}{2}
\tightlist
\item
  Collapse the product
  \(\displaystyle \binom{r}{\,m+n-j\,}\binom{m+n-j}{\,m\,}\) via
  \(\binom{u}{v}\binom{v}{w}=\binom{u}{w}\binom{u-w}{\,v-w\,}\):
\end{enumerate}

\[
\binom{r}{\,m+n-j\,}\binom{m+n-j}{\,m\,}
=\binom{r}{m}\binom{r-m}{\,n-j\,}.
\]

Hence

\[
S=\binom{r}{m}\sum_j \binom{A}{j}\binom{r-m}{\,n-j\,}.
\]

\begin{enumerate}
\def\labelenumi{\arabic{enumi}.}
\setcounter{enumi}{3}
\tightlist
\item
  Apply Vandermonde again:
\end{enumerate}

\[
\sum_j \binom{A}{j}\binom{r-m}{\,n-j\,}=\binom{A+r-m}{\,n\,}.
\]

With \(A=m-r+s\), we get \(A+r-m=s\). Thus

\[
\boxed{\,S=\binom{r}{m}\binom{s}{n}\,}.
\]

This identity holds under the usual convention that \(\binom{x}{y}=0\)
when \(y<0\) or \(y>x\), so the (finite) sum can be taken over all
integers \(k\).

\emph{Quick check.} Take \((m,n,r,s)=(2,1,3,2)\). Then the closed form
gives \(\binom{3}{2}\binom{2}{1}=3\cdot2=6\). A direct evaluation of the
sum also yields 6.

\subsection{32 {[}M20{]}}\label{m20-13}

Proof 1 (algebraic: via falling-rising and signs)

A classic identity in §1.2.6 is the falling-factorial expansion

\[
x^{\underline n}=\sum_{k=0}^{n} s(n,k)\,x^{k},
\]

where \(s(n,k)=(-1)^{\,n-k}\bigl[{n\atop k}\bigr]\) are the signed
Stirling numbers of the first kind. Substitute \(x\mapsto -x\):

\[
(-x)^{\underline n}=\sum_{k=0}^{n} s(n,k)\,(-x)^{k}
=\sum_{k=0}^{n}(-1)^{k} s(n,k)\,x^{k}
=\sum_{k=0}^{n}(-1)^{n}\bigl[{n\atop k}\bigr]\,x^{k}.
\]

But \((-x)^{\underline n}=(-1)^{n}x^{\overline n}\) (relation between
rising and falling factorials). Cancel \((-1)^{n}\) to obtain

\[
\sum_{k=0}^{n}\bigl[{n\atop k}\bigr]\,x^{k}=x^{\overline n},
\]

as claimed. (This is exactly the identity the exercise asks for.)
(\href{https://stackoverflow.com/questions/768948/how-does-division-work-in-mix}{ProofWiki})

Proof 2 (combinatorial: coloring cycles)

Fix a positive integer \(x\) and let \(F(n)\) be the number of
permutations of \(\{1,\dots,n\}\) whose cycles are colored with one of
\(x\) colors.

\begin{itemize}
\item
  Count by cycles. If a permutation has \(k\) cycles, there are
  \(\bigl[{n\atop k}\bigr]\) such permutations and each can color its
  \(k\) cycles in \(x^k\) ways. Thus

  \[
  F(n)=\sum_{k=0}^{n}\bigl[{n\atop k}\bigr]\,x^{k}.
  \]
\item
  Incremental construction. Build a colored permutation of
  \(\{1,\dots,n\}\) by inserting \(n\) into a colored permutation of
  \(\{1,\dots,n-1\}\):

  \begin{itemize}
  \tightlist
  \item
    either start a new 1-cycle \((n)\) and choose any of \(x\) colors
    (\(x\) choices);
  \item
    or insert \(n\) into an existing cycle: across all cycles there are
    exactly \(n-1\) insertion positions (each position immediately after
    one of the \(n-1\) existing elements), and the colors of existing
    cycles are already chosen.
  \end{itemize}

  Hence the recurrence

  \[
  F(0)=1,\qquad F(n)=(x+n-1)\,F(n-1).
  \]

  Solving gives \(F(n)=x(x+1)\cdots(x+n-1)=x^{\overline n}\).
\end{itemize}

Equating the two expressions for \(F(n)\) yields
\(\displaystyle \sum_{k=0}^{n}\bigl[{n\atop k}\bigr]\,x^{k}=x^{\overline n}\).
This argument also shows the identity holds for all complex \(x\) by
polynomial identity (it holds for infinitely many \(x\in\mathbb{N}\)).

Proof 3 (induction on \(n\))

Let \(P_n(x)=\sum_{k}\bigl[{n\atop k}\bigr]x^{k}\). Using the well-known
Stirling recurrence

\[
\bigl[{n\atop k}\bigr]=\bigl[{n-1\atop k-1}\bigr]+(n-1)\bigl[{n-1\atop k}\bigr]
\quad(n\ge1),
\]

we get

\[
P_n(x)=\sum_{k}\bigl[{n-1\atop k-1}\bigr]x^{k}+(n-1)\sum_{k}\bigl[{n-1\atop k}\bigr]x^{k}
= x\,P_{n-1}(x) + (n-1)P_{n-1}(x)=(x+n-1)P_{n-1}(x).
\]

With \(P_0(x)=1\), we conclude
\(P_n(x)=\prod_{j=0}^{n-1}(x+j)=x^{\overline n}\).
(\href{https://stackoverflow.com/questions/768948/how-does-division-work-in-mix}{ProofWiki})

Sanity check (small \(n\))

For \(n=3\), \(\bigl[{3\atop 1}\bigr]=2\), \(\bigl[{3\atop 2}\bigr]=3\),
\(\bigl[{3\atop 3}\bigr]=1\); hence

\[
\sum_{k=1}^3 \bigl[{3\atop k}\bigr] x^k = 2x+3x^2+x^3 = x(x+1)(x+2)=x^{\overline3}.
\]

\subsection{33 {[}M20{]}}\label{m20-14}

\begin{enumerate}
\def\labelenumi{\arabic{enumi})}
\tightlist
\item
  Falling-factorial version
\end{enumerate}

Step A --- Vandermonde's convolution for generalized binomials

For any (real/complex) \(x,y\) and integer \(n\ge 0\),

\[
\sum_{k=0}^{n} \binom{x}{k}\binom{y}{n-k}=\binom{x+y}{n}.
\tag{V}
\]

Proof of (V) (coefficient-of argument). Using the generalized binomial
series, \((1+z)^x=\sum_{k\ge 0}\binom{x}{k}z^k\) and
\((1+z)^y=\sum_{j\ge 0}\binom{y}{j}z^j\) (as formal power series), we
have

\[
(1+z)^{x+y}=(1+z)^x(1+z)^y
= \left(\sum_{k\ge 0}\binom{x}{k}z^k\right)
  \left(\sum_{j\ge 0}\binom{y}{j}z^j\right).
\]

Equating coefficients of \(z^n\) gives (V).

Step B --- Translate (V) into falling factorials

By definition
\(\displaystyle \binom{x}{k}=\frac{x^{\underline{k}}}{k!}\) and
\(\displaystyle \binom{x+y}{n}=\frac{(x+y)^{\underline{n}}}{n!}\).
Multiply (V) by \(n!\) and rearrange:

\[
\frac{(x+y)^{\underline{n}}}{n!}\,n!
=\sum_{k=0}^{n}\frac{x^{\underline{k}}}{k!}\,
\frac{y^{\underline{n-k}}}{(n-k)!}\,n!
=\sum_{k=0}^{n}\binom{n}{k}\,x^{\underline{k}}\,y^{\underline{n-k}}.
\]

This is exactly the first identity.

\begin{enumerate}
\def\labelenumi{\arabic{enumi})}
\setcounter{enumi}{1}
\tightlist
\item
  Rising-factorial version
\end{enumerate}

Use the relation between rising and falling factorials:

\[
x^{\overline{n}} = (-1)^n(-x)^{\underline{n}}.
\tag{*}
\]

Apply the falling-factorial identity to \((-x)+(-y)\):

\[
(-x-y)^{\underline{n}}
=\sum_{k=0}^{n}\binom{n}{k}\,(-x)^{\underline{k}}\,(-y)^{\underline{n-k}}.
\]

Multiply both sides by \((-1)^n\) and convert each term by (*):

\[
(x+y)^{\overline{n}}
=\sum_{k=0}^{n}\binom{n}{k}\,x^{\overline{k}}\,y^{\overline{\,n-k\,}},
\]

which is the desired second identity.

\begin{enumerate}
\def\labelenumi{\arabic{enumi})}
\setcounter{enumi}{2}
\tightlist
\item
  Quick sanity checks
\end{enumerate}

\begin{itemize}
\tightlist
\item
  \(n=0\): Both sides equal \(1\).
\item
  \(n=1\):
  \((x+y)^{\underline{1}}=x+y=\binom{1}{0}x^{\underline{0}}y^{\underline{1}}+\binom{1}{1}x^{\underline{1}}y^{\underline{0}}\);
  similarly for rising factorials.
\item
  \(n=2\), falling: LHS \((x+y)(x+y-1)\). RHS
  \(\binom{2}{0}y(y-1)+\binom{2}{1}x\,y+\binom{2}{2}x(x-1)\) expands to
  the same polynomial.
\end{itemize}

\subsection{34 {[}M23{]}}\label{m23-9}

We need to prove the Abel--type identity for rising factorials

\[
(x+y)^{\overline{n}}
=\sum_{k=0}^{n}\binom{n}{k}\;x\;\bigl(x-kz+1\bigr)^{\overline{k-1}}\;\bigl(y+kz\bigr)^{\overline{\,n-k\,}}.
\tag{★}
\]

Here \(u^{\overline{m}}=u(u+1)\cdots(u+m-1)\) is the rising factorial.
The motivation is Abel's generalization of the binomial theorem for
ordinary powers

\[
(x+y)^m=\sum_{k=0}^{m}\binom{m}{k}\,x\,(x-kz)^{k-1}\,(y+kz)^{m-k}. \;\; \text{(Knuth 1.2.6–(16))} 
\;\; \text{}
\]

Two standard identities we'll use

\begin{enumerate}
\def\labelenumi{\arabic{enumi}.}
\tightlist
\item
  Rising factorial ↔ binomial coefficient. For any \(u\) and integer
  \(m\ge 0\),
\end{enumerate}

\[
u^{\overline{m}}=m!\binom{u+m-1}{m}.
\tag{RF}
\]

(Equivalently, \((x+y)^{\overline{n}}=n!\binom{x+y+n-1}{n}\).) A
succinct reference is ProofWiki's ``Rising Factorial as Factorial by
Binomial Coefficient.'' We'll cite the compiled proof of Torelli's sum
there, which uses exactly this translation.

\begin{enumerate}
\def\labelenumi{\arabic{enumi}.}
\setcounter{enumi}{1}
\tightlist
\item
  A Vandermonde/Abel--type binomial sum. For general parameters
  \(r,s,t\) and integer \(n\ge0\) with \(r-tk\neq0\) on the summation
  range,
\end{enumerate}

\[
\sum_{k\ge0}\binom{r-tk}{k}\binom{s-t(n-k)}{\,n-k\,}\frac{r}{\,r-tk\,}
=\binom{r+s-tn}{n}.
\tag{V}
\]

This is a well-known Abel/Vandermonde convolution (it is the engine
under the hood of Torelli's identity); a compact statement and use of it
appears in the ProofWiki page on Torelli's Sum.

Proof of (★)

Start by converting both sides of (★) to binomial coefficients using
(RF).

Left side:

\[
(x+y)^{\overline{n}}=n!\binom{x+y+n-1}{n}.
\tag{LHS}
\]

Right side: For each term with fixed \(k\),

\[
\begin{aligned}
\bigl(x-kz+1\bigr)^{\overline{k-1}}
&=(k-1)!\binom{x-kz+k-1}{k-1},\\
\bigl(y+kz\bigr)^{\overline{n-k}}
&=(n-k)!\binom{y+kz+n-k-1}{\,n-k\,}.
\end{aligned}
\]

Multiply by \(\binom{n}{k}\,x\) and simplify factorials:

\[
\binom{n}{k}\,x\,(k-1)!\,(n-k)! 
= x\cdot \frac{n!}{k}.
\]

Hence the entire right side of (★) becomes

\[
\sum_{k=0}^{n}\binom{n}{k}\;x\;\bigl(x-kz+1\bigr)^{\overline{k-1}}\;\bigl(y+kz\bigr)^{\overline{n-k}}
=
n!\sum_{k=0}^{n}\frac{x}{k}\binom{x-kz+k-1}{k-1}\binom{y+kz+n-k-1}{\,n-k\,}.
\tag{RHS}
\]

Now re-index
\(\binom{x-kz+k-1}{k-1}=\binom{x+(1-z)k}{k}\;\frac{x}{x+(1-z)k}\).
Substituting this in (RHS) gives

\[
\text{RHS}
=
n!\sum_{k=0}^{n}
\binom{x+(1-z)k}{k}\binom{y-1+n z+(n-k)(1-z)}{\,n-k\,}\,
\frac{x}{\,x+(1-z)k\,}.
\]

This is exactly the specialization of the binomial identity (V) with

\[
r=x,\quad t=-(1-z),\quad s=y-1+n z,
\]

so the sum evaluates to

\[
n!\binom{\,x+y+n-1\,}{n}.
\]

By (LHS) this equals \((x+y)^{\overline{n}}\). Therefore (★) holds,
completing the proof. \(\square\)

Sanity checks and special cases

\begin{itemize}
\tightlist
\item
  \(z=0\). Then (★) reduces to
  \((x+y)^{\overline{n}}=\sum_{k}\binom{n}{k}x\,x^{\overline{k-1}}\,y^{\overline{n-k}} =\sum_{k}\binom{n}{k}x^{\overline{k}}y^{\overline{n-k}},\)
  which is exactly the binomial theorem for rising powers (Exercise 33).
\item
  \(n=1\). Both sides equal \(x+y\).
\item
  Consistency with Abel (ordinary powers). The structure mirrors Knuth's
  Eq. (16) term-for-term; we've simply ``lifted'' ordinary powers to
  rising powers through the binomial--rising-factorial bridge (RF),
  which is why this result is traditionally called Torelli's sum.
\end{itemize}

\subsection{35 {[}M23{]}}\label{m23-10}

Here's a clean derivation of the two ``addition formulas'' for Stirling
numbers:

\[
\text{(First kind)}\qquad x^{\underline{n}}
=\sum_{k}(-1)^{\,n-k}\,\bigl[ {n\atop k}\bigr]\,x^k
\tag{44}
\]

\[
\text{(Second kind)}\qquad x^{n}
=\sum_{k}\left\{ {n\atop k}\right\}\,x^{\underline{k}}
\tag{45}
\]

and aiming to prove

\[
\boxed{\ \bigl[{n+1\atop m}\bigr]=n\,\bigl[{n\atop m}\bigr]+\bigl[{n\atop m-1}\bigr]\ ,\qquad
\left\{{n+1\atop m}\right\}=m\,\left\{{n\atop m}\right\}+\left\{{n\atop m-1}\right\}\ }
\tag{46}
\]

First kind: \([\,\cdot\,]\)

Start from the tautology for falling factorials:

\[
x^{\underline{n+1}}=x^{\underline{n}}\,(x-n).
\]

Expand both sides via (44).

Left side (by (44) with \(n\mapsto n+1\)):

\[
x^{\underline{n+1}}=\sum_{m}(-1)^{\,n+1-m}\,\bigl[{n+1\atop m}\bigr]\,x^{m}.
\]

Right side:

\[
x^{\underline{n}}(x-n)=
\Bigl(\sum_{k}(-1)^{\,n-k}\,\bigl[{n\atop k}\bigr]\,x^{k}\Bigr)x
\;-\;
n\sum_{k}(-1)^{\,n-k}\,\bigl[{n\atop k}\bigr]\,x^{k}.
\]

Collect the coefficient of \(x^{m}\). In the first sum put \(k=m-1\); in
the second keep \(k=m\):

\[
(-1)^{\,n-(m-1)}\bigl[{n\atop m-1}\bigr]\;-\;n\,(-1)^{\,n-m}\bigl[{n\atop m}\bigr]
= (-1)^{\,n+1-m}\Bigl(\bigl[{n\atop m-1}\bigr]+n\bigl[{n\atop m}\bigr]\Bigr).
\]

Equating with the left side's coefficient
\((-1)^{\,n+1-m}\bigl[{n+1\atop m}\bigr]\) and canceling the common
factor \$ (-1)\^{}\{,n+1-m\}\$ yields

\[
\bigl[{n+1\atop m}\bigr]=n\,\bigl[{n\atop m}\bigr]+\bigl[{n\atop m-1}\bigr],
\]

as required.

Second kind: \(\{\cdot\}\)

Use the simple falling--factorial identity

\[
x\cdot x^{\underline{k}}=(x-k)\,x^{\underline{k}}+k\,x^{\underline{k}}
= x^{\underline{k+1}}+k\,x^{\underline{k}}.
\tag{$*$}
\]

Now write

\[
x^{n+1} = x\cdot x^{n}
= x\cdot \sum_{k}\left\{{n\atop k}\right\}x^{\underline{k}}
= \sum_{k}\left\{{n\atop k}\right\}\bigl(x^{\underline{k+1}}+k\,x^{\underline{k}}\bigr)
= \sum_{m}\left\{{n\atop m-1}\right\}x^{\underline{m}}
+ \sum_{m} m\left\{{n\atop m}\right\}x^{\underline{m}}.
\]

On the other hand, by (45) with \(n\mapsto n+1\),

\[
x^{n+1}=\sum_{m}\left\{{n+1\atop m}\right\}x^{\underline{m}}.
\]

Equate coefficients of \(x^{\underline{m}}\) to get

\[
\left\{{n+1\atop m}\right\}=\left\{{n\atop m-1}\right\}+m\left\{{n\atop m}\right\},
\]

which is the second formula in (46).

\subsection{36 {[}M10{]}}\label{m10-8}

The binomial theorem in this section is (Eq. (13)):

\[
(x+y)^n=\sum_{k}\binom{n}{k}x^{\,k}y^{\,n-k}\quad(n\ge 0).
\]

Setting \((x,y)=(1,1)\) gives

\[
\sum_{k}\binom{n}{k}=(1+1)^n=2^n.
\]

Setting \((x,y)=(1,-1)\) gives

\[
\sum_{k}(-1)^k\binom{n}{k}=(1-1)^n=
\begin{cases}
1,& n=0,\\
0,& n>0,
\end{cases}
\]

i.e.~the value is the Kronecker delta \(\delta_{n0}\). These are exactly
the values indicated in the book's answer line for this exercise. 

Combinatorial proofs (why these identities are natural)

\begin{enumerate}
\def\labelenumi{\arabic{enumi}.}
\item
  \(\displaystyle \sum_{k}\binom{n}{k}=2^n\). \(\binom{n}{k}\) counts
  the \(k\)-element subsets of an \(n\)-element set. Summing over all
  \(k\) counts all subsets, of which there are \(2^n\) (each element is
  either in or out). Hence the sum is \(2^n\).
\item
  \(\displaystyle \sum_{k}(-1)^k\binom{n}{k}=\delta_{n0}\). Interpret
  the alternating sum as ``(\# even-sized subsets) − (\# odd-sized
  subsets)''. If \(n>0\), fix one element \(a\). Toggling membership of
  \(a\) pairs each subset with a partner of the opposite parity, so even
  and odd subsets cancel in pairs; the result is \(0\). When \(n=0\)
  there is only the empty set (even), so the sum is \(1\).
\end{enumerate}

Remarks:

\begin{itemize}
\tightlist
\item
  As an immediate corollary for \(n>0\): even and odd binomial
  coefficients contribute equally, so
  \(\sum_{k\text{ even}}\binom{n}{k}=\sum_{k\text{ odd}}\binom{n}{k}=2^{n-1}\)
  (this is the next exercise in the book).
\item
  The ``all-integers \(k\)'' summation notation in this section is meant
  to streamline manipulations; terms outside \(0\le k\le n\) are \(0\).
\end{itemize}

Answer: \(\displaystyle \sum_{k}\binom{n}{k}=2^n\) and
\(\displaystyle \sum_{k}(-1)^k\binom{n}{k}=\delta_{n0}\).

\subsection{37 {[}M10{]}}\label{m10-9}

Let

\[
E_n=\sum_{\substack{0\le k\le n\\ k\ \text{even}}}\binom{n}{k},\qquad
O_n=\sum_{\substack{0\le k\le n\\ k\ \text{odd}}}\binom{n}{k}.
\]

From the binomial theorem,

\[
(1+1)^n=\sum_{k=0}^n \binom{n}{k}=E_n+O_n=2^n.
\]

From the ``alternating--sign'' variant (the previous exercise),

\[
(1-1)^n=\sum_{k=0}^n (-1)^k\binom{n}{k}=E_n-O_n=
\begin{cases}
0,& n\ge 1,\\
1,& n=0.
\end{cases}
\]

(See Exercise 36's statement on the same page. )

For \(n\ge 1\) these give the linear system

\[
E_n+O_n=2^n,\qquad E_n-O_n=0,
\]

so \(E_n=O_n=2^{n-1}\).

For the edge case \(n=0\), the sum of even-indexed terms is
\(\binom{0}{0}=1\) and the sum of odd-indexed terms is \(0\).

Answer. For \(n\ge 1\),

\[
\boxed{\displaystyle \binom{n}{0}+\binom{n}{2}+\binom{n}{4}+\cdots=2^{\,n-1}}
\]

(and the complementary odd-indexed sum also equals \(2^{\,n-1}\)). For
\(n=0\), the even-indexed sum is \(1\).

(Optional) Combinatorial proof (nice bijection)

Fix an element \(x\) of an \(n\)-set (\(n\ge 1\)). The map
\(S\mapsto S\triangle\{x\}\) (toggle membership of \(x\)) pairs every
subset of even size with one of odd size and vice versa, with no fixed
points. Hence the numbers of even- and odd-sized subsets are equal, each
being half of the total \(2^n\), i.e., \(2^{n-1}\).

\subsection{38 {[}HM30{]}}\label{hm30-1}

\begin{enumerate}
\def\labelenumi{\arabic{enumi}.}
\tightlist
\item
  Filter the \(m\)-spaced terms
\end{enumerate}

Let \(\omega=e^{2\pi i/m}\). The standard ``roots-of-unity filter''

\[
\frac1m\sum_{j=0}^{m-1}\omega^{j(t-k)}=
\begin{cases}
1,&t\equiv k\pmod m,\\[2pt]
0,&\text{otherwise}
\end{cases}
\]

selects exactly the indices \(t\) with \(t\equiv k\!\!\pmod m\). Hence

\[
\sum_{q\ge 0}\binom{n}{k+qm}
=\sum_{t=0}^{n}\binom{n}{t}\cdot \frac1m\sum_{j=0}^{m-1}\omega^{j(t-k)}
=\frac1m\sum_{j=0}^{m-1}\omega^{-jk}\sum_{t=0}^{n}\binom{n}{t}\omega^{jt}.
\]

\begin{enumerate}
\def\labelenumi{\arabic{enumi}.}
\setcounter{enumi}{1}
\tightlist
\item
  Use the binomial theorem
\end{enumerate}

The inner sum is \((1+\omega^{j})^{n}\). Therefore

\[
\sum_{q\ge 0}\binom{n}{k+qm}
=\frac1m\sum_{j=0}^{m-1}\omega^{-jk}(1+\omega^{j})^{n}.
\]

\begin{enumerate}
\def\labelenumi{\arabic{enumi}.}
\setcounter{enumi}{2}
\tightlist
\item
  Put \(1+\omega^{j}\) in polar form
\end{enumerate}

Write \(\theta_j=\frac{j\pi}{m}\). Since
\(\omega^{j}=e^{2\pi ij/m}=e^{2i\theta_j}\),

\[
1+\omega^{j}=1+e^{2i\theta_j}
=e^{i\theta_j}\bigl(e^{-i\theta_j}+e^{i\theta_j}\bigr)
=e^{i\theta_j}\cdot 2\cos\theta_j.
\]

Thus

\[
(1+\omega^{j})^{n}=\bigl(2\cos\theta_j\bigr)^{n}e^{in\theta_j}.
\]

And \(\omega^{-jk}=e^{-2ijk\pi/m}=e^{-i(2k)\theta_j}\). Hence

\[
\omega^{-jk}(1+\omega^{j})^{n}
=\bigl(2\cos\theta_j\bigr)^{n}e^{i(n-2k)\theta_j}.
\]

\begin{enumerate}
\def\labelenumi{\arabic{enumi}.}
\setcounter{enumi}{3}
\tightlist
\item
  Take real parts
\end{enumerate}

The left-hand side is real, so we may take real parts of the last
display:

\[
\sum_{q\ge 0}\binom{n}{k+qm}
=\frac1m\sum_{j=0}^{m-1}\bigl(2\cos\theta_j\bigr)^{n}\,
\Re\!\bigl(e^{i(n-2k)\theta_j}\bigr)
=\frac1m\sum_{j=0}^{m-1}\bigl(2\cos\tfrac{j\pi}{m}\bigr)^{n}
\cos\!\Bigl(\tfrac{j\pi}{m}(n-2k)\Bigr),
\]

which is the desired identity. ∎

Check: the example \(m=3,\;k=1\)

Here \(\theta_0=0,\theta_1=\pi/3,\theta_2=2\pi/3\). So

\[
\sum_{q\ge0}\binom{n}{1+3q}
=\frac13\!\left[
(2\cos 0)^n\cos 0
+(2\cos\tfrac{\pi}{3})^n\cos\tfrac{(n-2)\pi}{3}
+(2\cos\tfrac{2\pi}{3})^n\cos\tfrac{2(n-2)\pi}{3}
\right].
\]

Since \(2\cos(\pi/3)=1\) and \(2\cos(2\pi/3)=-1\), the last two terms
are
\(\cos\!\bigl(\frac{(n-2)\pi}{3}\bigr)+(-1)^n\cos\!\bigl(\frac{2(n-2)\pi}{3}\bigr) =2\cos\!\bigl(\frac{(n-2)\pi}{3}\bigr)\)
(they are equal because the \(j=1\) and \(j=2\) contributions are
complex conjugates whose real parts match), giving

\[
\binom n1+\binom n4+\binom n7+\cdots
=\frac13\!\left(2^n+2\cos\frac{(n-2)\pi}{3}\right),
\]

as stated.

\subsection{39 {[}M10{]}}\label{m10-10}

Facts we'll use

\begin{itemize}
\tightlist
\item
  \([\!\begin{smallmatrix}n\\k\end{smallmatrix}\!]\) counts permutations
  of \(n\) elements with exactly \(k\) cycles (unsigned). 
\item
  The conversion between falling powers and ordinary powers:
\end{itemize}

\[
x^{\underline{n}}=\sum_{k=0}^n (-1)^{\,n-k}\,[\!\begin{smallmatrix}n\\k\end{smallmatrix}\!]\,x^k,
\quad\text{where }x^{\underline{n}}=x(x-1)\cdots(x-n+1).
\]

\begin{enumerate}
\def\labelenumi{(\roman{enumi})}
\tightlist
\item
  Plain row sum
\end{enumerate}

Combinatorial proof. Summing over all \(k\) just counts all permutations
of \(n\) items, regardless of the number of cycles. Hence

\[
\sum_{k=0}^n \biggl[\!\begin{matrix}n\\k\end{matrix}\!\biggr]=n!.
\]

(Every permutation has some \(k\) cycles.) 

Algebraic proof (from the generating identity). Set \(x=-1\) in the
identity for \(x^{\underline{n}}\). The left side is

\[
(-1)^{\underline{n}}=(-1)(-2)\cdots(-n)=(-1)^n n!,
\]

and the right side becomes

\[
\sum_{k=0}^n (-1)^{\,n-k}[\!\begin{smallmatrix}n\\k\end{smallmatrix}\!](-1)^k
= (-1)^n \sum_{k=0}^n [\!\begin{smallmatrix}n\\k\end{smallmatrix}\!].
\]

Cancel \$ (-1)\^{}n\$ to get
\(\sum_k [\!\begin{smallmatrix}n\\k\end{smallmatrix}\!]=n!\). 

\begin{enumerate}
\def\labelenumi{(\roman{enumi})}
\setcounter{enumi}{1}
\tightlist
\item
  Alternating row sum
\end{enumerate}

Following Ex. 36's convention, ``alternating'' means coefficients
\((-1)^k\). Put \(x=1\) in the same identity:

\[
1^{\underline{n}}
=\sum_{k=0}^n (-1)^{\,n-k}\,[\!\begin{smallmatrix}n\\k\end{smallmatrix}\!]\,1^k
= (-1)^n\sum_{k=0}^n (-1)^k[\!\begin{smallmatrix}n\\k\end{smallmatrix}\!].
\]

Now \(1^{\underline{n}}=1\) for \(n=0,1\) and \(0\) for \(n\ge2\).
Therefore

\[
\sum_{k=0}^n (-1)^k\biggl[\!\begin{matrix}n\\k\end{matrix}\!\biggr]
=\begin{cases}
1, & n=0,\\[2pt]
-1, & n=1,\\[2pt]
0, & n\ge 2.
\end{cases}
\]

Quick sanity check (small \(n\)):

\begin{itemize}
\tightlist
\item
  \(n=3\): row \([3,1]=2,\ [3,2]=3,\ [3,3]=1\). Sum \(2+3+1=6=3!\);
  alternating \(2-3+1=0\).
\item
  \(n=4\): \(6,11,6,1\). Sum \(24=4!\); alternating \(6-11+6-1=0\).
\end{itemize}

\subsection{40 {[}HM17{]}}\label{hm17}

\begin{enumerate}
\def\labelenumi{(\alph{enumi})}
\tightlist
\item
  Show \(B(x,1)=B(1,x)=\dfrac{1}{x}\).
\end{enumerate}

\begin{itemize}
\item
  \(B(x,1)=\displaystyle\int_0^1 t^{x-1}(1-t)^{0}\,dt=\int_0^1 t^{x-1}\,dt=\Big[\frac{t^{x}}{x}\Big]_{0}^{1}=\frac{1}{x}\),
  since \(x>0\).
\item
  \(B(1,x)=\displaystyle\int_0^1 t^{0}(1-t)^{x-1}\,dt=\int_0^1 (1-t)^{x-1}\,dt\).
  Substitute \(u=1-t\) (so \(du=-dt\), \(t=0\Rightarrow u=1\),
  \(t=1\Rightarrow u=0\)):

  \[
    B(1,x)=\int_1^0 u^{x-1}(-du)=\int_0^1 u^{x-1}\,du=\frac{1}{x}.
  \]
\end{itemize}

\begin{enumerate}
\def\labelenumi{(\alph{enumi})}
\setcounter{enumi}{1}
\tightlist
\item
  Show \(B(x+1,y)+B(x,y+1)=B(x,y)\).
\end{enumerate}

Use the elementary identity \(1=t+(1-t)\) inside the integral:

\[
\begin{aligned}
B(x,y)
&=\int_0^1 t^{x-1}(1-t)^{y-1}\bigl[t+(1-t)\bigr]\,dt\\
&=\int_0^1 t^{x}(1-t)^{y-1}\,dt\;+\;\int_0^1 t^{x-1}(1-t)^{y}\,dt\\
&=B(x+1,y)+B(x,y+1).
\end{aligned}
\]

\begin{enumerate}
\def\labelenumi{(\alph{enumi})}
\setcounter{enumi}{2}
\tightlist
\item
  Show \(B(x,y)=\dfrac{x+y}{y}\,B(x,y+1)\).
\end{enumerate}

Differentiate the product \(t^{x}(1-t)^{y}\) and integrate from 0 to 1:

\[
\frac{d}{dt}\bigl[t^{x}(1-t)^{y}\bigr]
= x\,t^{x-1}(1-t)^{y}-y\,t^{x}(1-t)^{y-1}.
\]

Because \(x,y>0\), the boundary term \([t^{x}(1-t)^{y}]_{0}^{1}=0\).
Hence

\[
0=x\!\int_0^1 t^{x-1}(1-t)^{y}\,dt\;-\;y\!\int_0^1 t^{x}(1-t)^{y-1}\,dt
= x\,B(x,y+1)-y\,B(x+1,y),
\]

so

\[
B(x+1,y)=\frac{x}{y}\,B(x,y+1).
\]

Combine this with part (b):

\[
B(x,y)=B(x+1,y)+B(x,y+1)
=\frac{x}{y}B(x,y+1)+B(x,y+1)
=\frac{x+y}{y}\,B(x,y+1),
\]

as required.

\subsection{41 {[}HM22{]}}\label{hm22-1}

Part (a)

Start from the given relation for \(\Gamma_m\) and form the ratio

\[
\frac{\Gamma_m(y)\,m^{\,x}}{\Gamma_m(x+y)}
=\frac{m^{\,y}B(y,m+1)\,m^{\,x}}{m^{\,x+y}\,B(x+y,m+1)}
=\frac{B(y,m+1)}{B(x+y,m+1)}.
\tag{1}
\]

Independently, use the elementary beta--recurrence from the previous
exercise 40(c),

\[
B(x,y)=\frac{x+y}{y}\,B(x,y+1),
\]

repeated \(m+1\) times. This yields

\[
B(x,y)
=\biggl(\prod_{j=0}^{m}\frac{x+y+j}{y+j}\biggr)\,B(x,y+m+1).
\tag{2}
\]

Now note that the product in (2) is exactly the ratio of the two beta
values in (1), because (by iterating the same recurrence and using
\(B(u,1)=1/u\))

\[
B(u,m+1)=\frac{m!}{u(u+1)\cdots(u+m)}.
\]

Hence

\[
\frac{B(y,m+1)}{B(x+y,m+1)}
=\prod_{j=0}^{m}\frac{x+y+j}{y+j}.
\]

Plugging this into (2) gives precisely

\[
B(x,y)=\frac{B(y,m+1)}{B(x+y,m+1)}\,B(x,y+m+1)
=\frac{\Gamma_m(y)\,m^{\,x}}{\Gamma_m(x+y)}\,B\bigl(x,\,y+m+1\bigr),
\]

which is the desired identity. 

Part (b)

Take \(m\to\infty\) in the formula of part (a). Two standard limits are
used:

\begin{enumerate}
\def\labelenumi{\arabic{enumi}.}
\tightlist
\item
  \(\Gamma_m(u)\to\Gamma(u)\) for fixed \(u>0\).
\item
  \((m+y)^{x}B(x,m+y+1)\to \Gamma(x)\), hence
  \(m^{x}B(x,y+m+1)\to\Gamma(x)\) as well, since
\end{enumerate}

\[
m^{x}B(x,y+m+1)
=\Bigl(\frac{m}{m+y}\Bigr)^{x}\,(m+y)^{x}B(x,m+y+1)\;\longrightarrow\;\Gamma(x).
\]

(These limits are standard monotone/continuity facts about the
approximants \(\Gamma_m\) and the beta integral as the second parameter
grows.) 

Passing to the limit in

\[
B(x,y)=\frac{\Gamma_m(y)\,m^{\,x}}{\Gamma_m(x+y)}\,B(x,y+m+1)
\]

gives

\[
B(x,y)=\frac{\Gamma(y)}{\Gamma(x+y)}\,\Gamma(x)
=\frac{\Gamma(x)\Gamma(y)}{\Gamma(x+y)},
\]

which is the classical beta--gamma relation. 

\subsection{42 {[}HM10{]}}\label{hm10}

The beta and gamma functions satisfy

\[
B(x,y)=\frac{\Gamma(x)\Gamma(y)}{\Gamma(x+y)}.
\]

Take \(x=k+1\) and \(y=r-k+1\). Then

\[
B(k+1,\;r-k+1)
=\frac{\Gamma(k+1)\,\Gamma(r-k+1)}{\Gamma(r+2)}
=\frac{\Gamma(k+1)\,\Gamma(r-k+1)}{(r+1)\,\Gamma(r+1)}.
\]

Invert both sides and use
\(\displaystyle \binom{r}{k}=\frac{\Gamma(r+1)}{\Gamma(k+1)\Gamma(r-k+1)}\):

\[
\boxed{\;\binom{r}{k}
=\frac{1}{(r+1)\,B(k+1,\;r-k+1)}\;}
=\boxed{\;\frac{\Gamma(r+1)}{\Gamma(k+1)\,\Gamma(r-k+1)}\;}.
\]

Integral form (equivalent)

Since \(B(k+1,r-k+1)=\int_0^1 t^{k}(1-t)^{r-k}\,dt\),

\[
\boxed{\;\binom{r}{k}
=\frac{1}{(r+1)\displaystyle\int_0^1 t^{k}(1-t)^{\,r-k}\,dt}\;}.
\]

\subsection{43 {[}HM20{]}}\label{hm20-8}

\[
I=\int_{0}^{1}\frac{dt}{\sqrt{t(1-t)}}\,.
\]

Solution 1 (affine change of variables + symmetry)

Let \(u=2t-1\). Then \(t=\tfrac{1+u}{2}\), so

\[
t(1-t)=\frac{1-u^2}{4},\qquad dt=\frac{du}{2}.
\]

Hence

\[
\frac{dt}{\sqrt{t(1-t)}}
=\frac{\tfrac{du}{2}}{\sqrt{(1-u^{2})/4}}
=\frac{du}{\sqrt{1-u^{2}}}.
\]

As \(t\) runs from \(0\) to \(1\), \(u\) runs from \(-1\) to \(1\).
Therefore

\[
I=\int_{-1}^{1}\frac{du}{\sqrt{1-u^{2}}}
=2\int_{0}^{1}\frac{du}{\sqrt{1-u^{2}}}
=2\bigl[\arcsin(u)\bigr]_{0}^{1}
=2\cdot\frac{\pi}{2}
=\pi.
\]

Solution 2 (trigonometric substitution)

Let \(t=\sin^{2}\theta\) with \(\theta\in[0,\tfrac{\pi}{2}]\). Then

\[
dt=2\sin\theta\cos\theta\,d\theta,\quad
\sqrt{t(1-t)}=\sqrt{\sin^{2}\theta\cos^{2}\theta}=\sin\theta\cos\theta,
\]

so

\[
I=\int_{0}^{\pi/2}\frac{2\sin\theta\cos\theta}{\sin\theta\cos\theta}\,d\theta
=\int_{0}^{\pi/2}2\,d\theta
=2\cdot\frac{\pi}{2}
=\pi.
\]

Solution 3 (Beta/Gamma viewpoint)

This is the Euler beta integral with \(x=y=\tfrac12\):

\[
I=B\!\left(\tfrac12,\tfrac12\right)
=\frac{\Gamma(\tfrac12)\Gamma(\tfrac12)}{\Gamma(1)}
=\frac{(\sqrt{\pi})^{2}}{1}
=\pi.
\]

Thus, in all three ways,

\[
\boxed{\,\int_{0}^{1}\dfrac{dt}{\sqrt{t(1-t)}}=\pi\,}.
\]

\subsection{44 {[}HM20{]}}\label{hm20-9}

Show, for integer \(r\ge 0\) and the generalized binomial coefficient,

\[
\binom{r}{\tfrac12}=\frac{2^{\,2r+1}}{\pi\binom{2r}{r}}.
\]

Step 1: Put \(\binom{r}{k}\) in Beta/Gamma form

Using

\[
B(x,y)=\frac{\Gamma(x)\Gamma(y)}{\Gamma(x+y)},
\]

we have

\[
\binom{r}{k}=\frac{\Gamma(r+1)}{\Gamma(k+1)\Gamma(r-k+1)}
           =\frac{1}{(r+1)\,B(k+1,\;r-k+1)}.
\tag{1}
\]

Step 2: Specialize to \(k=\frac12\) and simplify with a beta recurrence

Set \(k=\frac12\) in (1):

\[
\binom{r}{\tfrac12}=\frac{1}{(r+1)\,B\!\left(\tfrac32,\;r+\tfrac12\right)}.
\]

Use the beta recurrence \(B(x+1,y)=\frac{x}{x+y}\,B(x,y)\) with
\(x=\tfrac12,\;y=r+\tfrac12\):

\[
B\!\left(\tfrac32,\;r+\tfrac12\right)=\frac{\tfrac12}{r+1}\,B\!\left(\tfrac12,\;r+\tfrac12\right),
\]

so

\[
\binom{r}{\tfrac12}=\frac{2}{B\!\left(\tfrac12,\;r+\tfrac12\right)}.
\tag{2}
\]

Step 3: Evaluate \(B\!\left(\tfrac12,\;r+\tfrac12\right)\)

By the beta--gamma relation and \(\Gamma(\tfrac12)=\sqrt{\pi}\)
(established just above this exercise set),

\[
B\!\left(\tfrac12,\;r+\tfrac12\right)
= \frac{\Gamma(\tfrac12)\Gamma(r+\tfrac12)}{\Gamma(r+1)}
= \frac{\sqrt{\pi}\,\Gamma(r+\tfrac12)}{r!}. 
\]

For half--integers,

\[
\Gamma\!\left(r+\tfrac12\right)
= \frac{(2r)!}{4^r\,r!}\,\sqrt{\pi},
\]

hence

\[
B\!\left(\tfrac12,\;r+\tfrac12\right)=
\pi\,\frac{(2r)!}{4^r\,(r!)^2}
= \pi\,\frac{\binom{2r}{r}}{4^r}.
\]

(These facts are standard; \(\Gamma(\tfrac12)=\sqrt{\pi}\) is noted in
the same exercise block. )

Step 4: Conclude

Plugging this into (2) gives

\[
\binom{r}{\tfrac12}
=\frac{2}{\pi\,\binom{2r}{r}/4^r}
=\frac{2\cdot 4^r}{\pi\,\binom{2r}{r}}
=\frac{2^{\,2r+1}}{\pi\binom{2r}{r}},
\]

as required.

Quick checks

\begin{itemize}
\tightlist
\item
  \(r=0\):
  \(\binom{0}{1/2}=\dfrac{1}{\Gamma(3/2)\Gamma(1/2)}=\dfrac{2}{\pi}\),
  and RHS \(=2/(\pi\binom00)=2/\pi\).
\item
  \(r=1\): \(\binom{1}{1/2}=4/\pi\), RHS
  \(=2^{3}/(\pi\binom21)=8/(2\pi)=4/\pi\).
\end{itemize}

\subsection{45 {[}HM21{]}}\label{hm21-3}

Use the standard representation of generalized binomial coefficients via
a falling factorial:

\[
\binom{r}{k}=\frac{r(r-1)\cdots(r-k+1)}{k!}
=\frac{r^{\underline{k}}}{k!}.
\]

Divide by \(r^k\):

\[
\frac{\binom{r}{k}}{r^k}
=\frac{1}{k!}\prod_{j=0}^{k-1}\Bigl(1-\frac{j}{r}\Bigr).
\]

For fixed \(k\), each factor \(\bigl(1-\tfrac{j}{r}\bigr)\to 1\) as
\(r\to\infty\); hence the product tends to \(1\). Therefore

\[
\boxed{\displaystyle \lim_{r\to\infty}\frac{\binom{r}{k}}{r^{\,k}}=\frac{1}{k!}
=\frac{1}{\Gamma(k+1)}.}
\]

(Optional) Rate of convergence

A short expansion gives the first correction term. Using
\(\ln(1-x)=-x-\tfrac{x^2}{2}+O(x^3)\),

\[
\ln\!\left(\prod_{j=0}^{k-1}\Bigl(1-\frac{j}{r}\Bigr)\right)
=\sum_{j=0}^{k-1}\ln\Bigl(1-\frac{j}{r}\Bigr)
=-\frac{1}{r}\sum_{j=0}^{k-1}j\;-\;\frac{1}{2r^2}\sum_{j=0}^{k-1}j^2\;+\;O\!\left(\frac{1}{r^3}\right).
\]

Since \(\sum_{j=0}^{k-1}j=\frac{k(k-1)}{2}\) and
\(\sum_{j=0}^{k-1}j^2=\frac{k(k-1)(2k-1)}{6}\), exponentiating yields

\[
\frac{\binom{r}{k}}{r^k}
=\frac{1}{k!}\left(1-\frac{k(k-1)}{2r}+O\!\left(\frac{1}{r^2}\right)\right).
\]

So the convergence to \(1/k!\) is \(O(1/r)\).

(Alternative viewpoint via \(\Gamma\))

Using
\(\displaystyle \binom{r}{k}=\frac{\Gamma(r+1)}{\Gamma(k+1)\Gamma(r-k+1)}\)
and the fact (for fixed \(k\)) that

\[
\frac{\Gamma(r+1)}{\Gamma(r-k+1)}=r(r-1)\cdots(r-k+1)\sim r^{k}\qquad(r\to\infty),
\]

we again obtain \(\displaystyle \frac{1}{\Gamma(k+1)}\).

\subsection{46 {[}M21{]}}\label{m21-3}

Stirling's formula (to first order) is

\[
n!\;\sim\;\sqrt{2\pi n}\,\Bigl(\frac{n}{e}\Bigr)^n
\qquad (n\to\infty).
\]

(See the discussion and determination of the constant \(\sqrt{2\pi}\) in
§1.2.11.2.) 

Apply this to \(\displaystyle \binom{x+y}{y}=\frac{(x+y)!}{x!\,y!}\):

\[
\begin{aligned}
\binom{x+y}{y}
&\sim
\frac{\sqrt{2\pi(x+y)}\,\bigl(\tfrac{x+y}{e}\bigr)^{x+y}}
     {\sqrt{2\pi x}\,\bigl(\tfrac{x}{e}\bigr)^{x}\;\sqrt{2\pi y}\,\bigl(\tfrac{y}{e}\bigr)^{y}}\\[4pt]
&=\frac{1}{\sqrt{2\pi}}\,
   \sqrt{\frac{x+y}{xy}}\;
   \frac{(x+y)^{x+y}}{x^{x}\,y^{y}}.
\end{aligned}
\tag{★}
\]

A convenient reparametrization lets \(N=x+y\) and \(p=y/N\) (so
\(x=(1-p)N\)). Then (★) becomes the entropy-looking form

\[
\boxed{\;
\binom{N}{pN}\;\sim\;
\frac{1}{\sqrt{2\pi N\,p(1-p)}}\;
\Bigl(\frac{1}{p^{\,p}(1-p)^{\,1-p}}\Bigr)^{\!N}
\;}
\qquad(N\to\infty,\;0<p<1).
\]

This is the standard Stirling asymptotic for off-center binomials.

Special case: the central binomial coefficient

Set \(x=y=n\) (so \(N=2n,\;p=\tfrac12\)). Plugging into (★) gives

\[
\binom{2n}{n}
\;\sim\;
\frac{1}{\sqrt{2\pi}}\,
\sqrt{\frac{2n}{n^2}}\;
\frac{(2n)^{2n}}{n^{2n}}
\;=\;
\boxed{\;\frac{4^{\,n}}{\sqrt{\pi\,n}}\;}.
\]

This leading term is typically more than adequate; adding higher
Stirling terms refines it further (e.g., multiplying by
\(1-\tfrac{1}{8n}+O(n^{-2})\)), but the exercise only asks for the main
asymptotics.

\subsection{47 {[}M21{]}}\label{m21-4}

Let \(k\in\mathbb{Z}_{\ge 0}\) and extend binomial coefficients to
arbitrary \(r\) by the falling--factorial definition

\[
\binom{r}{k}=\frac{r(r-1)\cdots(r-k+1)}{k!}=\frac{(r)_k}{k!}\,,
\quad (r)_k:=\prod_{j=0}^{k-1}(r-j).
\]

Main identity

Compute

\[
\binom{r}{k}\binom{r-\tfrac12}{k}
=\frac{1}{(k!)^2}\prod_{j=0}^{k-1}\bigl(r-j\bigr)\bigl(r-\tfrac12-j\bigr).
\]

Pair the factors inside the product:

\[
4\bigl(r-j\bigr)\bigl(r-\tfrac12-j\bigr)
=(2r-2j)(2r-2j-1).
\]

Hence

\[
\binom{r}{k}\binom{r-\tfrac12}{k}
=\frac{1}{4^k(k!)^2}\prod_{j=0}^{k-1}(2r-2j)(2r-2j-1)
=\frac{(2r)_{2k}}{4^k(k!)^2}.
\]

Now observe two equivalent factorizations of \((2r)_{2k}\):

\begin{enumerate}
\def\labelenumi{\arabic{enumi}.}
\tightlist
\item
  Split the \(2k\) factors into two blocks of length \(k\):
\end{enumerate}

\[
(2r)_{2k}=(2r)_k\,(2r-k)_k,
\]

so

\[
\binom{r}{k}\binom{r-\tfrac12}{k}
=\frac{1}{4^k}\cdot\frac{(2r)_k}{k!}\cdot\frac{(2r-k)_k}{k!}
=\frac{1}{4^k}\binom{2r}{k}\binom{2r-k}{k}.
\]

\begin{enumerate}
\def\labelenumi{\arabic{enumi}.}
\setcounter{enumi}{1}
\tightlist
\item
  Or divide and multiply by \((2k)!\):
\end{enumerate}

\[
\frac{(2r)_{2k}}{(k!)^2}
=\frac{(2r)_{2k}}{(2k)!}\cdot\frac{(2k)!}{(k!)^2}
=\binom{2r}{2k}\binom{2k}{k}.
\]

Thus

\[
\boxed{\displaystyle
\binom{r}{k}\binom{r-\tfrac12}{k}
=\frac{1}{4^k}\binom{2r}{k}\binom{2r-k}{k}
=\frac{1}{4^k}\binom{2r}{2k}\binom{2k}{k}.}
\]

(This is exactly the statement of the exercise.) 

Special case \(r=-\tfrac12\)

Plugging \(r=-\tfrac12\) into the boxed formula gives \(2r=-1\) and
\(\binom{-1}{2k}=1\). Hence

\[
(-1)^k\binom{-\tfrac12}{k}
=\frac{1}{4^k}\binom{2k}{k}
\quad\Longrightarrow\quad
\boxed{\displaystyle \binom{-\tfrac12}{k}=(-1)^k\,\frac{1}{4^k}\binom{2k}{k}.}
\]

\emph{Remarks.} The proof is purely algebraic; no gamma functions are
needed. If desired, one can also derive it from \(\Gamma\)-duplication,
but the falling--factorial product pairing yields the identity in one
line.

\subsection{48 {[}M25{]}}\label{m25-10}

Proof 1 (Beta--function integral; one line of algebra and one line of
calculus)

Use the Beta integral (see the preceding exercise's definition
\(B(x,y)=\int_0^1 t^{\,x-1}(1-t)^{\,y-1}\,dt\) for \(x,y>0\)). 

\begin{enumerate}
\def\labelenumi{\arabic{enumi}.}
\tightlist
\item
  Expand \((1-t)^n\) by the binomial theorem and integrate termwise:
\end{enumerate}

\[
\int_0^1 t^{x-1}(1-t)^n\,dt
=\int_0^1 t^{x-1}\!\sum_{k=0}^n\binom{n}{k}(-t)^k\,dt
=\sum_{k=0}^n \binom{n}{k}(-1)^k\!\int_0^1 t^{x+k-1}\,dt
=\sum_{k=0}^n \binom{n}{k}\frac{(-1)^k}{x+k}.
\]

\begin{enumerate}
\def\labelenumi{\arabic{enumi}.}
\setcounter{enumi}{1}
\tightlist
\item
  Evaluate the same integral by a standard Beta recursion (integration
  by parts once, with \(u=(1-t)^n\), \(dv=t^{x-1}dt\)), which yields
\end{enumerate}

\[
B(x,n+1)=\int_0^1 t^{x-1}(1-t)^n\,dt
=\frac{n}{x}\,B(x+1,n)
=\cdots=\frac{n!}{x(x+1)\cdots(x+n)}.
\]

Equating the two evaluations of \(B(x,n+1)\) proves

\[
\sum_{k=0}^n \binom{n}{k}\frac{(-1)^k}{x+k}
=\frac{n!}{x(x+1)\cdots(x+n)}
=\frac{1}{x^{(n+1)}},
\]

valid whenever \(x\notin\{0,-1,\dots,-n\}\) so that denominators are
nonzero. 

Proof 2 (Partial fractions / residues; purely algebraic)

Consider

\[
R(x)\;=\;\frac{n!}{x(x+1)\cdots(x+n)}.
\]

It has simple poles at \(x=-k\) \((0\le k\le n)\), so write a partial
fraction expansion

\[
R(x)=\sum_{k=0}^n \frac{A_k}{x+k}.
\]

Multiply by \(x+k\) and set \(x=-k\):

\[
A_k=\frac{n!}{\prod_{j\ne k}(-k+j)}
=\frac{n!}{\bigl[(-1)^k k!\,(n-k)!\bigr]}
=(-1)^k\binom{n}{k}.
\]

Therefore \(R(x)=\sum_{k=0}^n \binom{n}{k}\dfrac{(-1)^k}{x+k}\), which
is exactly the desired identity.

Checks and remarks

\begin{itemize}
\tightlist
\item
  Sanity check \(n=0\): LHS \(=1/x\); RHS \(=0!/x=1/x\).
\item
  Poles/conditions: The identity holds for all
  \(x\notin\{0,-1,\dots,-n\}\). The integral proof assumes \(\Re(x)>0\)
  to justify termwise integration, then extends by analytic
  continuation; the algebraic proof is rational and automatically valid
  off the poles.
\end{itemize}

\subsection{49 {[}M20{]}}\label{m20-15}

Work in formal power series (or analytically for \(|x|<1\)). Expand each
factor with the generalized binomial theorem:

\begin{itemize}
\tightlist
\item
  \((1+x)^r=\sum_{n\ge0}\binom{r}{n}x^n\).
\item
  \((1-x^2)^r=\sum_{i\ge0}\binom{r}{i}(-1)^i x^{2i}\).
\item
  \((1-x)^{-r}=\sum_{j\ge0}\binom{r+j-1}{j}x^{j}\) (negative--binomial
  expansion).
\end{itemize}

Multiply the last two series and collect the coefficient of \(x^n\)
(Cauchy product). Only terms with \(2i+j=n\) contribute,
i.e.~\(j=n-2i\ge0\). Hence

\[
[x^n]\bigl((1-x^2)^r(1-x)^{-r}\bigr)
=\sum_{i=0}^{\lfloor n/2\rfloor}
\binom{r}{i}(-1)^i\binom{r+n-2i-1}{\,n-2i\,}.
\]

But this coefficient must equal \([x^n](1+x)^r=\binom{r}{n}\).
Therefore, for every integer \(n\ge0\) and arbitrary (real/complex)
\(r\),

\[
\boxed{\displaystyle
\binom{r}{n}
=
\sum_{i=0}^{\lfloor n/2\rfloor}
(-1)^i\binom{r}{i}\binom{r+n-2i-1}{\,n-2i\,}}
\]

which is the requested relation on binomial coefficients.

Quick checks

\begin{itemize}
\tightlist
\item
  \(n=1\): RHS \(=\binom{r}{0}\binom{r}{1}=r=\binom{r}{1}\).
\item
  \(n=2\): RHS
  \(=\binom{r}{0}\binom{r+1}{2}-\binom{r}{1}\binom{r-1}{0} =\frac{r(r+1)}{2}-r=\frac{r(r-1)}{2}=\binom{r}{2}\).
\item
  \(n=3\): RHS
  \(=\binom{r}{0}\binom{r+2}{3}-\binom{r}{1}\binom{r}{1} =\frac{r(r+1)(r+2)}{6}-r^2=\frac{r(r-1)(r-2)}{6}=\binom{r}{3}\).
\end{itemize}

\subsection{50 {[}M20{]}}\label{m20-16}

Abel's identity (for integer \(n\ge 0\), with \(x\ne 0\)) is

\[
(x+y)^n=\sum_{k}\binom{n}{k}\;x\,(x-kz)^{k-1}(y+kz)^{\,n-k}.
\]

This is Eq. (16) in §1.2.6. 

We need to prove it when \(x+y=0\) (i.e., \(y=-x\)).  In that case
\((x+y)^n=0^n\), which is \(0\) when \(n\ge1\) and \(1\) when \(n=0\).
So we must show:

\[
\boxed{\;\sum_{k=0}^n \binom{n}{k}\;x\,(x-kz)^{k-1}(-x+kz)^{\,n-k}
=\begin{cases}
1,& n=0,\\[4pt]
0,& n\ge 1\;.
\end{cases}}
\]

Reduction when \(y=-x\)

For \(n\ge1\), note that
\((-x+kz)^{n-k}=(kz-x)^{n-k}=(-1)^{\,n-k}(x-kz)^{n-k}\). Hence each
summand becomes

\[
\binom{n}{k}\;x\,(x-kz)^{k-1}\,(-1)^{\,n-k}(x-kz)^{\,n-k}
= \binom{n}{k}\;x\,(-1)^{\,n-k}(x-kz)^{\,n-1}.
\]

Therefore

\[
S_n:=\sum_{k=0}^n \binom{n}{k}\;x\,(x-kz)^{k-1}(-x+kz)^{\,n-k}
= x\,(-1)^n \sum_{k=0}^n \binom{n}{k}\;(-1)^k (x-kz)^{\,n-1}.
\tag{∗}
\]

Expand and annihilate with Stirling numbers

Expand the polynomial \((x-kz)^{\,n-1}\) in powers of \(k\):

\[
(x-kz)^{\,n-1}
=\sum_{j=0}^{n-1} \binom{n-1}{j}\,x^{\,n-1-j}\,(-z)^{\,j}\,k^{\,j}.
\]

Insert this into (∗) and swap sums:

\[
S_n
= x\,(-1)^n \sum_{j=0}^{n-1} \binom{n-1}{j}\,x^{\,n-1-j}\,(-z)^{\,j}
\Biggl[\sum_{k=0}^n \binom{n}{k}(-1)^k k^{\,j}\Biggr].
\]

Now use the standard binomial--Stirling connection (Eq. (53) in §1.2.6):

\[
\sum_{k=0}^{m} \binom{m}{k}(-1)^{\,m-k} k^{\,n} = m!\,\Bigl\{\!\!\begin{smallmatrix} n\\ m\end{smallmatrix}\!\!\Bigr\},
\]

where \(\bigl\{\!\begin{smallmatrix}n\\ m\end{smallmatrix}\!\bigr\}\) is
a Stirling number of the second kind.

With \(m=n\) and our sign convention,

\[
\sum_{k=0}^n \binom{n}{k}(-1)^k k^{\,j}
= (-1)^n \, n!\,\Bigl\{\!\!\begin{smallmatrix} j\\ n\end{smallmatrix}\!\!\Bigr\}.
\]

Hence

\[
S_n
= x\,(-1)^n \sum_{j=0}^{n-1} \binom{n-1}{j}\,x^{\,n-1-j}\,(-z)^{\,j}\,
\bigl[(-1)^n n!\,\bigl\{\!\!\begin{smallmatrix} j\\ n\end{smallmatrix}\!\!\bigr\}\bigr]
= x\,n!\sum_{j=0}^{n-1} \binom{n-1}{j}\,x^{\,n-1-j}\,(-z)^{\,j}\,
\Bigl\{\!\!\begin{smallmatrix} j\\ n\end{smallmatrix}\!\!\Bigr\}.
\]

But \(\bigl\{\!\begin{smallmatrix} j\\ n\end{smallmatrix}\!\bigr\}=0\)
whenever \(j<n\) (you can't partition \(j\) elements into \(n>j\)
nonempty blocks; equivalently see the defining expansion
\(x^{j}=\sum_{r}\bigl\{\!\begin{smallmatrix} j\\ r\end{smallmatrix}\!\bigr\} x^{\underline r}\)).
Since here \(j=0,1,\dots,n-1\), every term vanishes, giving \(S_n=0\)
for \(n\ge1\).

Edge case \(n=0\)

When \(n=0\), the sum has the single term \(k=0\):

\[
\binom{0}{0}\,x\,(x-0\cdot z)^{-1}(-x+0\cdot z)^{0}=x\cdot x^{-1}\cdot 1=1,
\]

matching \((x+y)^0=1\). Thus the identity also holds for \(n=0\).

\subsection{51 {[}M21{]}}\label{m21-5}

Proof (by coefficient extraction in \(T=x+y\))

Let \(T=x+y\). The right--hand side of (16), after substituting
\(y=T-x\), is the polynomial

\[
P(T)=\sum_{k=0}^n \binom{n}{k}\;x\,(x-kz)^{\,k-1}\,(T-x+kz)^{\,n-k}.
\]

We will expand \(P(T)\) as \(\sum_{j=0}^{n} C_j\,T^{\,j}\) and show

\[
C_n=1\quad\text{and}\quad C_j=0\ \ (0\le j<n),
\]

which implies \(P(T)=T^n=(x+y)^n\), as desired.

Step 1. Expand \((T-x+kz)^{n-k}\) and collect coefficients of \(T^j\).

By the binomial theorem,

\[
(T-x+kz)^{\,n-k}
=\sum_{j=0}^{n-k} \binom{n-k}{j}T^{\,j}(-x+kz)^{\,n-k-j}.
\]

Therefore, the coefficient of \(T^{\,j}\) in \(P(T)\) is

\[
C_j=\sum_{k=0}^{n-j}\binom{n}{k}\binom{n-k}{j}\;
x\,(x-kz)^{\,k-1}\,(-x+kz)^{\,n-k-j}.
\]

Use the elementary combinatorial identity

\[
\binom{n}{k}\binom{n-k}{j}=\binom{n}{j}\binom{n-j}{k}
\]

(choose \(j\) objects first, then \(k\) of the remaining \(n-j\)). This
gives

\[
C_j=\binom{n}{j}\;\sum_{k=0}^{n-j}\binom{n-j}{k}\;
x\,(x-kz)^{\,k-1}\,(-x+kz)^{\,n-j-k}.
\]

Let \(m=n-j\) (so \(0\le m\le n\)). Then

\[
C_j=\binom{n}{j}\;S_m,\qquad
S_m=\sum_{k=0}^{m}\binom{m}{k}\;
x\,(x-kz)^{\,k-1}\,(-x+kz)^{\,m-k}.
\tag{★}
\]

Step 2. Invoke Exercise 50 (the special case \(T=0\)).

Exercise 50 established Abel's identity when \(x+y=0\), i.e.~\(T=0\):

\[
0^{\,m}=\sum_{k=0}^{m}\binom{m}{k}\;x\,(x-kz)^{\,k-1}\,(-x+kz)^{\,m-k}
\quad\text{for all integers } m\ge 1.
\]

(Exactly the same sum as \(S_m\) above.) Hence

\[
S_m=0\quad\text{for all}\quad m\ge 1.
\]

Therefore, for \(j<n\) (i.e.~\(m=n-j\ge 1\)) we get
\(C_j=\binom{n}{j}S_m=0\).

Step 3. The top coefficient \(C_n\).

When \(j=n\), we have \(m=0\). Then in (★) the only term is \(k=0\):

\[
S_0=\binom{0}{0}\;x\,(x-0\cdot z)^{-1}\,(-x+0\cdot z)^{0}=1
\qquad(\text{since }x\neq0).
\]

Thus \(C_n=\binom{n}{n}S_0=1\).

All lower coefficients vanish and the top coefficient is \(1\), so

\[
P(T)=\sum_{j=0}^{n}C_jT^{\,j}=T^n.
\]

With \(T=x+y\) we obtain exactly Abel's formula (16):

\[
(x+y)^n=\sum_{k=0}^{n}\binom{n}{k}\;x\,(x-kz)^{\,k-1}\,(y+kz)^{\,n-k}.
\]

\subsection{52 {[}HM11{]}}\label{hm11}

\begin{enumerate}
\def\labelenumi{\arabic{enumi})}
\tightlist
\item
  Recall Abel's binomial formula
\end{enumerate}

For integer \(n\ge 0\) and \(x\ne 0\),

\[
(x+y)^n=\sum_k \binom{n}{k}\; x\,(x-kz)^{k-1}\,(y+kz)^{\,n-k}.
\]

(See Eq. (16) in §1.2.6.) 

When \(n\) is a nonnegative integer the sum has only finitely many
nonzero terms, so the identity is well-defined.

\begin{enumerate}
\def\labelenumi{\arabic{enumi})}
\setcounter{enumi}{1}
\tightlist
\item
  Evaluate the right--hand side at \(n=x=-1,\; y=z=1\)
\end{enumerate}

Use the generalized binomial coefficient \(\binom{-1}{k}=(-1)^k\) (valid
for all integers \(k\ge 0\)). Then each term of the sum becomes

\[
\binom{-1}{k}\;(-1)\,(-1-k)^{k-1}\,(1+k)^{-1-k}.
\]

Simplify the signs and powers:

\[
\binom{-1}{k}(-1)(-1-k)^{k-1}(1+k)^{-1-k}
= (-1)^k(-1)\big[(-1)^{k-1}(k+1)^{k-1}\big](k+1)^{-1-k}
= (k+1)^{-2}.
\]

Hence the whole right--hand side is

\[
\sum_{k\ge 0}\frac{1}{(k+1)^2}=\sum_{m\ge 1}\frac{1}{m^2}=\frac{\pi^2}{6},
\]

using the classical evaluation \(\zeta(2)=\pi^2/6\). 

\begin{enumerate}
\def\labelenumi{\arabic{enumi})}
\setcounter{enumi}{2}
\tightlist
\item
  Compare with the left--hand side
\end{enumerate}

At these same values, the left--hand side is

\[
(x+y)^n=( -1+1)^{-1}=0^{-1},
\]

which is undefined (infinite). Therefore the identity fails at this
choice of parameters; in particular, Abel's binomial formula is not
valid in general when \(n\) is not a nonnegative integer.

\subsection{53 {[}M25{]}}\label{m25-11}

\begin{enumerate}
\def\labelenumi{(\alph{enumi})}
\tightlist
\item
  Induction on \(m\)
\end{enumerate}

Define

\[
S_m=\sum_{k=0}^{m}\binom{r}{k}\binom{s}{\,n-k\,}\bigl(nr-(r+s)k\bigr).
\]

Base case \(m=0\).

\[
S_0=\binom{r}{0}\binom{s}{n}\cdot nr=nr\binom{s}{n}.
\]

Right--hand side at \(m=0\) is

\[
(0+1)(n-0)\binom{r}{1}\binom{s}{n}
=n\cdot r\binom{s}{n},
\]

which agrees.

Inductive step \(m\to m+1\).

Write \(S_m=S_{m-1}+\Delta_m\) with

\[
\Delta_m=\binom{r}{m}\binom{s}{\,n-m\,}\bigl(nr-(r+s)m\bigr).
\]

By the induction hypothesis,

\[
S_{m-1}=m(n-m+1)\binom{r}{m}\binom{s}{\,n-m+1\,}.
\]

Hence we want to show

\[
S_m
= m(n-m+1)\binom{r}{m}\binom{s}{\,n-m+1\,}
+\binom{r}{m}\binom{s}{\,n-m\,}\bigl(nr-(r+s)m\bigr)
\\[2pt]
= (m+1)(n-m)\binom{r}{\,m+1\,}\binom{s}{\,n-m\,}.
\]

Use the standard ``lowering'' recurrences

\[
\binom{r}{m+1}=\frac{r-m}{m+1}\binom{r}{m},
\qquad
\binom{s}{\,n-m+1\,}=\frac{s-(n-m)}{n-m+1}\binom{s}{\,n-m\,}.
\]

Substitute into the desired identity and factor
\(\binom{r}{m}\binom{s}{\,n-m\,}\). The right--hand side becomes

\[
(m+1)(n-m)\frac{r-m}{m+1}\binom{r}{m}\binom{s}{\,n-m\,}
=(n-m)(r-m)\binom{r}{m}\binom{s}{\,n-m\,}.
\]

The left--hand side equals

\[
m(n-m+1)\binom{r}{m}\frac{s-(n-m)}{n-m+1}\binom{s}{\,n-m\,}
+\binom{r}{m}\binom{s}{\,n-m\,}\bigl(nr-(r+s)m\bigr)
\]

\[
=\binom{r}{m}\binom{s}{\,n-m\,}\Bigl(m\,(s-n+m)+nr-(r+s)m\Bigr)
=\binom{r}{m}\binom{s}{\,n-m\,}\bigl((n-m)r-m s\bigr)
\]

\[
=(n-m)(r-m)\binom{r}{m}\binom{s}{\,n-m\,},
\]

since \((n-m)r-m s=(n-m)(r-m)\). This matches the right--hand side,
completing the induction.

\begin{enumerate}
\def\labelenumi{(\alph{enumi})}
\setcounter{enumi}{1}
\tightlist
\item
  Specialize \(r=\tfrac12\), \(s=-\tfrac12\)
\end{enumerate}

With \(r+s=0\), the factor in part (a) simplifies:

\[
nr-(r+s)k = n\cdot \tfrac12 = \frac{n}{2}\quad(\text{independent of }k).
\]

Thus part (a) yields the partial convolution

\[
\sum_{k=0}^{m}\binom{\tfrac12}{k}\binom{-\tfrac12}{\,n-k\,}
=\frac{2}{n}\,(m+1)(n-m)\binom{\tfrac12}{\,m+1\,}\binom{-\tfrac12}{\,n-m\,}.
\tag{★}
\]

Now use the relations (from ex. 47, stated in the book) for half-integer
binomials:

\[
\binom{-\tfrac12}{t}
=(-1)^t\frac{1}{4^t}\binom{2t}{t},\qquad
\binom{\tfrac12}{t}
=(-1)^{t-1}\frac{1}{2^{2t}(2t-1)}\binom{2t}{t}\quad(t\ge1),
\]

and the equivalent form

\[
\binom{\tfrac12}{t}
=(-1)^{t-1}\frac{1}{2^{2t-1}t}\binom{2t-2}{t-1}\quad(t\ge1),
\]

together with \(\binom{2k}{k}=2\binom{2k-1}{k}\) for \(k\ge1\).

A convenient per-term conversion (valid for \(k\ge1\)) is

\[
\binom{\tfrac12}{k}\binom{-\tfrac12}{\,n-k\,}
= (-1)^{n}\,\frac{1}{2^{2n-1}}\,
\binom{2n-2k}{\,n-k\,}\binom{2k-1}{k}\,\frac{-1}{2k-1}.
\tag{†}
\]

For \(k=0\) one checks separately:

\[
\binom{\tfrac12}{0}\binom{-\tfrac12}{n}
=(-1)^n\,\frac{1}{2^{2n}}\binom{2n}{n},
\]

which differs by a factor \(1/2\) from what (†) would give if naively
extended to \(k=0\). Precisely this ``half-term'' produces the
stand-alone \(\tfrac12\binom{2n}{n}\) in the final answer.

Summing (†) for \(k=1,\dots,m\) and adding the \(k=0\) term, we obtain

\[
\sum_{k=0}^{m}\binom{2k-1}{k}\binom{2n-2k}{\,n-k\,}\frac{-1}{2k-1}
=2^{2n-1}(-1)^n\sum_{k=0}^{m}\binom{\tfrac12}{k}\binom{-\tfrac12}{\,n-k\,}
+\frac{1}{2}\binom{2n}{n}.
\tag{‡}
\]

Insert the convolution (★) on the right of (‡), then convert the
remaining \(\binom{\tfrac12}{\,m+1\,}\) and
\(\binom{-\tfrac12}{\,n-m\,}\) to central binomials and simplify:

\[
\binom{\tfrac12}{\,m+1\,}\binom{-\tfrac12}{\,n-m\,}
=(-1)^n\frac{1}{2^{2n+2}}\frac{1}{2m+1}\binom{2m+2}{m+1}\binom{2n-2m}{\,n-m\,}
\]

and

\[
\frac{\binom{2m+2}{m+1}}{2m+1}=\frac{2}{m+1}\binom{2m}{m}.
\]

A short algebraic cleanup yields

\[
2^{2n-1}(-1)^n\cdot \frac{2}{n}(m+1)(n-m)
\binom{\tfrac12}{\,m+1\,}\binom{-\tfrac12}{\,n-m\,}
=\frac{n-m}{2n}\binom{2m}{m}\binom{2n-2m}{\,n-m\,}.
\]

Together with the extra \(\tfrac12\binom{2n}{n}\) from \(k=0\), this is
exactly the formula to be shown in part (b).

\subsection{54 {[}M21{]}}\label{m21-6}

Result (the inverse)

\[
\boxed{\;P^{-1}=(q_{ij})_{i,j\ge 0},\qquad 
q_{ij}=
\begin{cases}
(-1)^{\,i-j}\binom{i}{j}, & i\ge j,\\[2pt]
0, & i<j~,
\end{cases}\;}
\]

i.e., take Pascal's triangle and place minus signs in a checkerboard
pattern \(({-}1)^{i+j}=({-}1)^{i-j}\).

Proof 1 (direct algebra via a standard binomial identity)

Consider \((PQ)_{ij}\) with \(Q\) as above:

\[
(PQ)_{ij}=\sum_{k=j}^{i}\binom{i}{k}\,(-1)^{k-j}\binom{k}{j}.
\]

Use \(\binom{i}{k}\binom{k}{j}=\binom{i}{j}\binom{i-j}{k-j}\)
(``upper--lower factorization''), then set \(t=k-j\):

\[
(PQ)_{ij}=\binom{i}{j}\sum_{t=0}^{i-j}\binom{i-j}{t}(-1)^t
=\binom{i}{j}(1-1)^{\,i-j}
=\begin{cases}
1,& i=j,\\
0,& i\ne j.
\end{cases}
\]

Hence \(PQ=I\). By the same symmetry, \(QP=I\), so \(Q=P^{-1}\).

Proof 2 (using Knuth's identity (33))

TAOCP Eq. (33) states

\[
\sum_{k}\binom{r}{k}\binom{k}{n}(-1)^{\,r-k}=\delta_{rn}\qquad(r,n\ge0).
\]

Rewrite the sign as \((-1)^{\,r}(-1)^{\,k}\) and compare with
\((P\cdot Q)_{rn}\) for \(Q_{kn}=(-1)^{k-n}\binom{k}{n}\):

\[
(PQ)_{rn}
=\sum_{k}\binom{r}{k}(-1)^{\,k-n}\binom{k}{n}
=(-1)^{\,r+n}\sum_{k}\binom{r}{k}\binom{k}{n}(-1)^{\,r-k}
=(-1)^{\,r+n}\delta_{rn}=\delta_{rn}.
\]

(When \(r=n\) the factor is \((-1)^{2n}=1\).) Therefore \(Q=P^{-1}\). 

Small check (4×4 example)

For indices \(0\le i,j\le 3\),

\[
P=\begin{pmatrix}
1&0&0&0\\
1&1&0&0\\
1&2&1&0\\
1&3&3&1
\end{pmatrix},
\qquad
P^{-1}=\begin{pmatrix}
1&0&0&0\\
-1&1&0&0\\
1&-2&1&0\\
-1&3&-3&1
\end{pmatrix},
\]

which is Pascal with the checkerboard signs; one can verify
\(PP^{-1}=I\). (This matches the pattern shown in the book's answer.)

\subsection{55 {[}M21{]}}\label{m21-7}

Knuth gives the inversion formulas (here \(\delta_{mn}\) is Kronecker's
delta):

\[
\sum_{k} \Bigl[\begin{smallmatrix}n\\k\end{smallmatrix}\Bigr]
\Bigl\{\begin{smallmatrix}k\\m\end{smallmatrix}\Bigr\}
(-1)^{\,n-k} = \delta_{mn},\qquad
\sum_{k} \Bigl\{\begin{smallmatrix}n\\k\end{smallmatrix}\Bigr]
\Bigl[\begin{smallmatrix}k\\m\end{smallmatrix}\Bigr]
(-1)^{\,n-k} = \delta_{mn}.
\tag{★}
\]

These appear in Eq. (47) immediately after Table 2.

Matrix formulation

Let

\begin{itemize}
\tightlist
\item
  \(S=\bigl(S_{nk}\bigr)_{n,k\ge 0}\) with
  \(S_{nk}=\Bigl[\begin{smallmatrix}n\\k\end{smallmatrix}\Bigr]\) (first
  kind),
\item
  \(T=\bigl(T_{nk}\bigr)_{n,k\ge 0}\) with
  \(T_{nk}=\Bigl\{\begin{smallmatrix}n\\k\end{smallmatrix}\Bigr\}\)
  (second kind),
\item
  \(D=\mathrm{diag}\bigl((-1)^0,(-1)^1,(-1)^2,\dots\bigr)\) (a diagonal
  ``sign'' matrix).
\end{itemize}

Note that
\((DSD)_{nk}=(-1)^{\,n-k} \Bigl[\begin{smallmatrix}n\\k\end{smallmatrix}\Bigr]\)
and similarly
\((DTD)_{nk}=(-1)^{\,n-k} \Bigl\{\begin{smallmatrix}n\\k\end{smallmatrix}\Bigr]\).

The first identity in (★) says exactly

\[
(DSD)\,T=I,
\]

hence \(T^{-1}=DSD\). The second identity says

\[
(DTD)\,S=I,
\]

hence \(S^{-1}=DTD\). 

Explicit entrywise formulas (the answer)

Unwinding the definitions gives:

\begin{itemize}
\tightlist
\item
  Inverse of the second-kind triangle \(T^{-1}\):
\end{itemize}

\[
(T^{-1})_{nm} \;=\; (-1)^{\,n-m}\,\Bigl[\begin{smallmatrix}n\\m\end{smallmatrix}\Bigr],\qquad n,m\ge 0.
\]

\begin{itemize}
\tightlist
\item
  Inverse of the first-kind triangle \(S^{-1}\):
\end{itemize}

\[
(S^{-1})_{nm} \;=\; (-1)^{\,n-m}\,\Bigl\{\begin{smallmatrix}n\\m\end{smallmatrix}\Bigr\},\qquad n,m\ge 0.
\]

Equivalently, in matrix notation:

\[
S^{-1}=DTD,\qquad T^{-1}=DSD.
\]

Quick sanity check on a small corner

Using the \(0\!\ldots\!3\) block of Table 2:

\[
S=\begin{pmatrix}
1&0&0&0\\
0&1&0&0\\
0&1&1&0\\
0&2&3&1
\end{pmatrix},\quad
T=\begin{pmatrix}
1&0&0&0\\
0&1&0&0\\
0&1&1&0\\
0&1&3&1
\end{pmatrix},\quad
D=\mathrm{diag}(1,-1,1,-1).
\]

Compute \(DTD\) and \(DSD\); multiplying \(S\cdot(DTD)\) and
\(T\cdot(DSD)\) yields the identity, as predicted by (★). (All these
small blocks are visible in Table 2.)

\subsection{56 {[}20{]}}\label{section-81}

We use the classic greedy construction that underlies the combinatorial
number system:

\begin{enumerate}
\def\labelenumi{\arabic{enumi}.}
\item
  Let \(a\) be the largest integer with \(\binom{a}{3}\le n\). Put
  \(r_1=n-\binom{a}{3}\).
\item
  Let \(b<a\) be the largest integer with \(\binom{b}{2}\le r_1\). Put
  \(r_2=r_1-\binom{b}{2}\).
\item
  Set \(c=r_2\). Then \(0\le c< b\) because
\end{enumerate}

\[
\binom{b+1}{2}-\binom{b}{2}=b,
\]

so maximality of \(b\) forces \(r_2< b\).

This gives \(n=\binom{a}{3}+\binom{b}{2}+c\) with \(a>b>c\ge 0\).

Uniqueness. \(a\) is uniquely determined by step 1 since
\(\binom{a+1}{3}>\!n\ge\binom{a}{3}\). Given \(a\), the same
``difference'' identity above shows \(b\) is uniquely determined; then
\(c\) is forced as the remainder. Hence the representation is unique.

\emph{(This is the order-3 case of the general combinatorial number
system: Every nonnegative integer has a unique expansion
\(\sum_{j=1}^m \binom{a_j}{j}\) with \(a_m>\cdots>a_1\ge 0\).)}

Results for \(n=0,1,\dots,20\)

Below are the triples \((a,b,c)\) such that
\(n=\binom{a}{3}+\binom{b}{2}+\binom{c}{1}\).

\begin{longtable}[]{@{}rcl@{}}
\toprule\noalign{}
\(n\) & \((a,b,c)\) & check \\
\midrule\noalign{}
\endhead
\bottomrule\noalign{}
\endlastfoot
0 & (2,1,0) & \(0=0+0+0\) \\
1 & (3,1,0) & \(1=1+0+0\) \\
2 & (3,2,0) & \(2=1+1+0\) \\
3 & (3,2,1) & \(3=1+1+1\) \\
4 & (4,1,0) & \(4=4+0+0\) \\
5 & (4,2,0) & \(5=4+1+0\) \\
6 & (4,2,1) & \(6=4+1+1\) \\
7 & (4,3,0) & \(7=4+3+0\) \\
8 & (4,3,1) & \(8=4+3+1\) \\
9 & (4,3,2) & \(9=4+3+2\) \\
10 & (5,1,0) & \(10=10+0+0\) \\
11 & (5,2,0) & \(11=10+1+0\) \\
12 & (5,2,1) & \(12=10+1+1\) \\
13 & (5,3,0) & \(13=10+3+0\) \\
14 & (5,3,1) & \(14=10+3+1\) \\
15 & (5,3,2) & \(15=10+3+2\) \\
16 & (5,4,0) & \(16=10+6+0\) \\
17 & (5,4,1) & \(17=10+6+1\) \\
18 & (5,4,2) & \(18=10+6+2\) \\
19 & (5,4,3) & \(19=10+6+3\) \\
20 & (6,1,0) & \(20=20+0+0\) \\
\end{longtable}

(All triples satisfy \(a>b>c\ge 0\).)

Pattern and continuation

\begin{itemize}
\tightlist
\item
  The threshold values are the tetrahedral numbers \(\binom{a}{3}\):
  when \(n\) passes \(\binom{a}{3}\), \(a\) increments.
\item
  Between \(\binom{a}{3}\) and \(\binom{a+1}{3}-1\), the pairs \((b,c)\)
  run through all strictly decreasing \(b>c\) with \(b<a\) in
  colexicographic order because of the greedy choice and the identity
  \(\binom{b+1}{2}-\binom{b}{2}=b\).
\item
  To continue for larger \(n\), repeat the same greedy steps. Example:
  \(n=21\Rightarrow a=6\) (since \(\binom{6}{3}=20\)), \(r_1=1\); choose
  \(b=2\) (\(\binom{2}{2}=1\)), remainder \(c=0\); hence
  \(21=\binom{6}{3}+\binom{2}{2}+\binom{0}{1}\).
\end{itemize}

Knuth labels this representation the combinatorial number system in the
exercise statement, and the table above follows exactly from that idea.

\subsection{57 {[}M22{]}}\label{m22-10}

Let

\[
\ln n!\;=\;\sum_{m\ge 1} a_m\,n^{\underline m}\qquad(n^{\underline m}=n(n-1)\cdots(n-m+1))
\]

(the Newton--Stirling expansion that Stirling used in his ``attempt to
generalize the factorial function''). Show that

\[
\boxed{\displaystyle
a_m=\frac{(-1)^m}{m!}\sum_{k\ge 1}(-1)^k\binom{m-1}{\,k-1\,}\,\ln k
}
\]

(which is effectively a finite sum since \(\binom{m-1}{k-1}=0\) for
\(k>m\)).

Two ingredients

\begin{enumerate}
\def\labelenumi{\arabic{enumi}.}
\tightlist
\item
  Binomial--difference orthogonality. For integers \(m\ge0\),
\end{enumerate}

\[
\sum_{n}\binom{m}{n}\binom{k}{n}(-1)^{m-n}=\delta_{mk}
\quad\Longleftrightarrow\quad
\sum_{n}\binom{m}{n}(-1)^{m-n}\binom{n}{r}=\delta_{mr}.
\]

Knuth states the needed special case (his Eq. (33)) and uses it for
``inversion'' in Example 5 of §1.2.6. We'll use exactly that maneuver to
isolate a single coefficient.

\begin{enumerate}
\def\labelenumi{\arabic{enumi}.}
\setcounter{enumi}{1}
\tightlist
\item
  Forward differences pick off Newton coefficients. If
\end{enumerate}

\[
f(n)=\sum_{m\ge0} c_m \binom{n}{m}=\sum_{m\ge0}\frac{c_m}{m!}\,n^{\underline m},
\]

then

\[
\Delta^m f(0)=m!\,c_m \qquad\text{and}\qquad 
c_m=\frac{1}{m!}\sum_{n=0}^m(-1)^{m-n}\binom{m}{n}f(n).
\]

(Exactly the ``inversion'' used in Example 5. )

Apply the inversion step to \(f(n)=\ln n!\) with the falling--factorial
basis:

\[
m!\,a_m
=\sum_{n=0}^m(-1)^{m-n}\binom{m}{n}\,\ln n!
\tag{1}
\]

Now expand \(\ln n!\) as a telescoping sum of logs and interchange
summations:

\[
\ln n!=\sum_{k=1}^{n}\ln k
\quad\Longrightarrow\quad
\sum_{n=0}^m(-1)^{m-n}\binom{m}{n}\ln n!
=\sum_{k=1}^{m}\!\Biggl(\sum_{n=k}^{m}(-1)^{m-n}\binom{m}{n}\Biggr)\ln k.
\]

Evaluate the inner binomial sum. A standard identity (hockey-stick with
alternating signs) gives

\[
\sum_{n=k}^{m}(-1)^{m-n}\binom{m}{n}
=(-1)^{m-k}\binom{m-1}{k-1}.
\]

Substitute back into (1):

\[
m!\,a_m
=\sum_{k=1}^{m}(-1)^{m-k}\binom{m-1}{k-1}\ln k
=(-1)^m\sum_{k=1}^{m}(-1)^k\binom{m-1}{k-1}\ln k.
\]

Divide by \(m!\) to obtain

\[
\boxed{\displaystyle
a_m=\frac{(-1)^m}{m!}\sum_{k=1}^{m}(-1)^k\binom{m-1}{\,k-1\,}\ln k},
\]

which is the desired formula (the summation can be written
\(\sum_{k\ge1}\) since terms vanish for \(k>m\)).

Quick sanity checks

\begin{itemize}
\tightlist
\item
  \(m=1:\; a_1=\frac{-1}{1!}\,(-1)\binom{0}{0}\ln1=0\) (so the
  \(n^{\underline1}=n\) term is absent, consistent with \(\ln 1!=0\)).
\item
  \(m=2:\; a_2=\frac{1}{2}\bigl(-\ln1+\ln2\bigr)=\tfrac12\ln2\). Then at
  \(n=2\): \(\ln2!=a_2\,2^{\underline2}=(\tfrac12\ln2)\cdot2=\ln2\).
\item
  \(m=3:\; a_3=\frac{-1}{6}\bigl(2\ln2-\ln3\bigr)=\frac{\ln3-2\ln2}{6}\).
  At \(n=3\):
  \(\ln3!=a_2\cdot3^{\underline2}+a_3\cdot3^{\underline3}=3\ln2+(\ln3-2\ln2)=\ln6\).
\end{itemize}

\subsection{58 {[}M23{]}}\label{m23-11}

Let \(\binom{n}{k}_q\) denote the Gaussian (q-binomial) coefficient (the
notation introduced just before Eq. (41) in §1.2.6). Prove the q-nomial
theorem

\[
\prod_{j=0}^{n-1}\left(1+q^{\,j}x\right)
\;=\;
\sum_{k=0}^{n} \binom{n}{k}_q\, q^{\,\binom{k}{2}}\,x^k,
\]

and find q-analogues of the ``fundamental identities'' (17) and (21).

The Gaussian (q-binomial) coefficient is

\[
\binom{n}{k}_q \;=\; \frac{(1-q^{n})(1-q^{n-1})\cdots(1-q^{n-k+1})}{(1-q)(1-q^{2})\cdots(1-q^{k})},
\qquad 0\le k\le n.
\]

It satisfies the standard recurrences (both forms are useful):

\[
\boxed{\;\binom{n}{k}_q=\binom{n-1}{k}_q+q^{\,n-k}\binom{n-1}{k-1}_q
\;=\;q^{\,k}\binom{n-1}{k}_q+\binom{n-1}{k-1}_q\;}
\quad(1)
\]

For reference, the classical identities in §1.2.6 are

\[
\binom{r}{k} = (-1)^k \binom{k-r-1}{k}\qquad\text{(Eq. (17))} \quad\text{and}\quad
\sum_k \binom{r}{k}\binom{s}{n-k}=\binom{r+s}{n}\quad\text{(Eq. (21))}. \; \, \, \, \, \, \,
\]

The q-nomial theorem

Claim. For every integer \(n\ge 0\),

\[
\prod_{j=0}^{n-1}\left(1+q^{\,j}x\right)
\;=\;
\sum_{k=0}^{n} \binom{n}{k}_q\, q^{\,\binom{k}{2}}\,x^k.
\tag{QBT}
\]

Proof (induction on \(n\))

\begin{itemize}
\item
  Base \(n=0\). LHS and RHS are both \(1\).
\item
  Inductive step. Assuming the formula holds for \(n-1\), write

  \[
  \prod_{j=0}^{n-1}\left(1+q^{\,j}x\right)
  =\Biggl(\sum_{k=0}^{n-1}\binom{n-1}{k}_q\, q^{\,\binom{k}{2}} x^k\Biggr)\,(1+q^{\,n-1}x).
  \]

  The coefficient of \(x^k\) on the RHS is

  \[
  q^{\binom{k}{2}}\binom{n-1}{k}_q \;+\; q^{\,n-1}\,q^{\binom{k-1}{2}}\binom{n-1}{k-1}_q
  \;=\;
  q^{\binom{k}{2}}\Bigl(\binom{n-1}{k}_q + q^{\,n-k}\binom{n-1}{k-1}_q\Bigr),
  \]

  which equals \(q^{\binom{k}{2}}\binom{n}{k}_q\) by the first
  recurrence in (1). This gives (QBT) for \(n\).
\end{itemize}

Thus (QBT) holds for all \(n\).

Sanity check (small \(n\)). For \(n=3\),

\[
(1+x)(1+qx)(1+q^2x)=1+(1+q+q^2)x+(q+q^2+q^3)x^2+q^3x^3,
\]

while the RHS gives \(k=0,1,2,3\) terms \(\binom{3}{0}_q\),
\(\binom{3}{1}_q x\), \(q^{1}\binom{3}{2}_q x^2\),
\(q^{3}\binom{3}{3}_q x^3\), and
\(\binom{3}{1}_q=\binom{3}{2}_q=1+q+q^2\), so they match.

q-analogue of (17)

The classical reflection (17) is
\(\binom{r}{k}=(-1)^k\binom{k-r-1}{k}\). Using
\(1-q^{t}=-q^{\,t}(1-q^{-t})\) in the product definition of
\(\binom{r}{k}_q\), one obtains the q-reflection

\[
\boxed{\;\binom{r}{k}_q
\;=\;
(-1)^k\, q^{\,kr-\binom{k}{2}}\,
\binom{k-r-1}{k}_q\;}.
\tag{q-17}
\]

q-analogue of (21) (q-Vandermonde/Chu--Vandermonde)

Split the length-\(r+s\) product in (QBT) into two contiguous blocks:

\[
\prod_{j=0}^{r+s-1}(1+q^{\,j}x)
=
\Bigl[\prod_{j=0}^{r-1}(1+q^{\,j}x)\Bigr]
\Bigl[\prod_{j=r}^{r+s-1}(1+q^{\,j}x)\Bigr]
=
\Bigl[\cdots\Bigr]\Bigl[\prod_{j=0}^{s-1}(1+q^{\,j}(q^{\,r}x))\Bigr].
\]

Apply (QBT) to each bracket and take the coefficient of \(x^n\).
Comparing with the single-block expansion for \(r+s\) yields

\[
\boxed{\;
\sum_{k=0}^{n}
\binom{r}{k}_q\binom{s}{\,n-k\,}_q\;
q^{\,(r-k)(\,n-k\,)}
\;=\;
\binom{r+s}{n}_q
\;}.
\tag{q-21a}
\]

Using the symmetry \(\binom{m}{j}_q=q^{\,j(m-j)}\binom{m}{m-j}_q\) one
obtains the equivalent---and frequently quoted---form

\[
\boxed{\;
\sum_{k=0}^{n}
\binom{r}{k}_q\binom{s}{\,n-k\,}_q\;
q^{\,(s-n+k)k}
\;=\;
\binom{r+s}{n}_q
\;}.
\tag{q-21b}
\]

\subsection{59 {[}M25{]}}\label{m25-12}

\begin{enumerate}
\def\labelenumi{\arabic{enumi})}
\tightlist
\item
  Reduce to a homogeneous (``Pascal'') recurrence
\end{enumerate}

The inhomogeneous term \(\binom{n}{k}\) suggests subtracting a multiple
of \(\binom{n}{k}\). Define

\[
S_{n,k}\;=\;A_{n,k}-(n+1)\binom{n}{k}.
\]

Then for \(nk>0\),

\[
\begin{aligned}
S_{n,k}
&=A_{n-1,k}+A_{n-1,k-1}+\binom{n}{k}-(n+1)\binom{n}{k}\\
&=\bigl[S_{n-1,k}+n\binom{n-1}{k}\bigr]+\bigl[S_{n-1,k-1}+n\binom{n-1}{k-1}\bigr]
+\binom{n}{k}-(n+1)\binom{n}{k}\\
&=S_{n-1,k}+S_{n-1,k-1}
\end{aligned}
\]

because \(\binom{n-1}{k}+\binom{n-1}{k-1}=\binom{n}{k}\). Thus
\(S_{n,k}\) satisfies the \emph{homogeneous} Pascal recurrence.

\begin{enumerate}
\def\labelenumi{\arabic{enumi})}
\setcounter{enumi}{1}
\tightlist
\item
  Boundary values for \(S_{n,k}\)
\end{enumerate}

From the given edge conditions,

\[
S_{n,0}=A_{n,0}-(n+1)\binom{n}{0}=1-(n+1)=-n,\qquad
S_{0,k}=A_{0,k}-1\cdot\binom{0}{k}=\delta_{0k}-\delta_{0k}=0.
\]

\begin{enumerate}
\def\labelenumi{\arabic{enumi})}
\setcounter{enumi}{2}
\tightlist
\item
  Identify the unique solution to the homogeneous problem
\end{enumerate}

The homogeneous Pascal rule with those boundaries is solved by

\[
S_{n,k}=-\binom{n}{k+1}.
\]

This is easily proved by induction: It matches the edges
(\(S_{n,0}=-\binom{n}{1}=-n\), \(S_{0,k}=-\binom{0}{k+1}=0\)), and
satisfies

\[
-\binom{n}{k+1}= -\bigl(\binom{n-1}{k+1}+\binom{n-1}{k}\bigr)
= S_{n-1,k}+S_{n-1,k-1}.
\]

\begin{enumerate}
\def\labelenumi{\arabic{enumi})}
\setcounter{enumi}{3}
\tightlist
\item
  Assemble the answer
\end{enumerate}

Therefore

\[
\boxed{\;A_{n,k}=(n+1)\binom{n}{k}-\binom{n}{k+1}\;}
\qquad(n\ge 0,\ k\ge 0).
\]

This single formula automatically gives the edges:

\begin{itemize}
\tightlist
\item
  \(k=0\): \(A_{n,0}=(n+1)-\binom{n}{1}=1\).
\item
  \(n=0\): \(A_{0,k}=1\cdot\binom{0}{k}-\binom{0}{k+1}=\delta_{0k}\).
\end{itemize}

\begin{enumerate}
\def\labelenumi{\arabic{enumi})}
\setcounter{enumi}{4}
\tightlist
\item
  Equivalent (often handy) forms
\end{enumerate}

Using \(\binom{n}{k+1}=\frac{n-k}{k+1}\binom{n}{k}\) and
\(\binom{n+1}{k+1}=\frac{n+1}{k+1}\binom{n}{k}\), we get two nice
equivalents:

\[
A_{n,k}
= \frac{kn+2k+1}{k+1}\binom{n}{k}
\;=\;
\binom{n}{k}+k\,\binom{n+1}{k+1}.
\]

\begin{enumerate}
\def\labelenumi{\arabic{enumi})}
\setcounter{enumi}{5}
\tightlist
\item
  Quick sanity check
\end{enumerate}

For example, with \(n=5\):

\[
\bigl(A_{5,k}\bigr)_{k=0}^5
=\bigl(1,\;20,\;50,\;55,\;29,\;6\bigr),
\]

and indeed \(A_{5,k}=(6)\binom{5}{k}-\binom{5}{k+1}\) or
\(\binom{5}{k}+k\binom{6}{k+1}\) gives the same values.

That completes the determination of \(A_{n,k}\).

\subsection{61 {[}M25{]}}\label{m25-13}

We'll use only the standard Stirling transforms listed in §1.2.6:

\begin{itemize}
\item
  Conversion between factorial and ordinary powers

  \[
  x^{\underline n}=\sum_{k}(-1)^{\,n-k}\Bigl[\!\begin{smallmatrix}n\\ k\end{smallmatrix}\Bigr]x^k,
  \qquad
  x^{n}=\sum_{k}\Bigl\{\!\begin{smallmatrix}n\\ k\end{smallmatrix}\Bigr\}x^{\underline k},
  \tag{A}
  \]

  and their ``inversion'' (orthogonality)

  \[
  \sum_{k}\Bigl\{\!\begin{smallmatrix}n\\ k\end{smallmatrix}\Bigr\}
             \Bigl[\!\begin{smallmatrix}k\\ m\end{smallmatrix}\Bigr]
             (-1)^{\,k-m}
   \;=\;\delta_{nm}.
  \tag{B}
  \]

  (These are Eqs. (44), (45), and the first formula in (47).)
\item
  The ``raising-index'' identity for Stirling numbers of the first kind

  \[
  \Bigl[\!\begin{smallmatrix}n+1\\ m+1\end{smallmatrix}\Bigr]
  \;=\;\sum_{k\le n}\Bigl[\!\begin{smallmatrix}k\\ m\end{smallmatrix}\Bigr]\,
       \frac{n!}{k!}\,,
  \tag{C}
  \]

  which is the first of Knuth's Eq. (56).
\end{itemize}

Start from (C) with \(m\) replaced by \(k\) and \(m\) playing the final
role in the sum:

\[
\Bigl[\!\begin{smallmatrix}n+1\\ k+1\end{smallmatrix}\Bigr]
=\sum_{j\le n}\Bigl[\!\begin{smallmatrix}j\\ k\end{smallmatrix}\Bigr]\frac{n!}{j!}.
\]

Insert this into \(S_{n,m}\) and interchange the sums:

\[
\begin{aligned}
S_{n,m}
&=\sum_{k}\Bigl\{\!\begin{smallmatrix}k\\ m\end{smallmatrix}\Bigr\}(-1)^{\,k-m}
     \sum_{j\le n}\Bigl[\!\begin{smallmatrix}j\\ k\end{smallmatrix}\Bigr]\frac{n!}{j!}\\[2mm]
&=\sum_{j\le n}\frac{n!}{j!}\;
   \underbrace{\sum_{k}\Bigl\{\!\begin{smallmatrix}k\\ m\end{smallmatrix}\Bigr\}
                    \Bigl[\!\begin{smallmatrix}j\\ k\end{smallmatrix}\Bigr]
                    (-1)^{\,k-m}}_{\textstyle=\;\delta_{jm}\ \text{by (B)}} .
\end{aligned}
\]

By orthogonality (B) the inner sum collapses to \(\delta_{jm}\), leaving

\[
S_{n,m}=\frac{n!}{m!}\qquad(0\le m\le n),
\]

and \(S_{n,m}=0\) when \(m>n\) (since then
\(\{\!\begin{smallmatrix}k\\ m\end{smallmatrix}\!\}=0\) for all
\(k\le n\)).

This gives the requested ``companion'' to Eq. (55) (which states
\(\sum_k \{\!\begin{smallmatrix}n+1\\ k+1\end{smallmatrix}\}\,[\!\begin{smallmatrix}k\\ m\end{smallmatrix}](-1)^{k-m} = \binom{n}{m}\)).

\[
\boxed{\displaystyle
\sum_k \Bigl[\!\begin{smallmatrix}n+1\\ k+1\end{smallmatrix}\Bigr]
           \Bigl\{\!\begin{smallmatrix}k\\ m\end{smallmatrix}\Bigr\}
           (-1)^{\,k-m}
\;=\;\frac{n!}{m!}\,.
}
\]

Quick check. For \(n=6\): the row of values is
\(720,720,360,120,30,6,1\) for \(m=0,\dots,6\), i.e.~\(6!/m!\), matching
the formula.

\subsection{62 {[}M23{]}}\label{m23-12}

Step 1: A tool we'll use (Exercise 31)

We'll repeatedly use the identity that (in the notation of §1.2.6)

\[
\sum_{t}(-1)^t
\binom{m-r+s}{t}\binom{n+r-s}{n-t}\binom{r+t}{\,m+n\,}
=\binom{r}{m}\binom{s}{n}.
\tag{★}
\]

This can be proved in a line or two by the standard Vandermonde trick
twice: expand \(\binom{r+t}{m+n}=\sum_j\binom{r}{m+n-j}\binom{t}{j}\),
swap sums, use \(\binom{A}{j}\binom{B}{N-j}\)--convolution, and simplify
via \(\binom{a}{b}\binom{b}{c} =\binom{a}{c}\binom{a-c}{b-c}\). (This is
exactly the result Exercise 31 asks you to derive.)

Step 2: Break one binomial to create a ``\(\,r+t\,\)'' top

Start from the left-hand side

\[
S:=\sum_{k}(-1)^k
\binom{l+m}{l+k}\binom{m+n}{m+k}\binom{n+l}{n+k}.
\]

Use the binomial splitting

\[
\binom{n+l}{\,n+k\,}
=\sum_{j}\binom{l+k}{\,j\,}\binom{n-k}{\,n+k-j\,},
\]

which is just \(\binom{A+B}{U}=\sum_j\binom{A}{j}\binom{B}{U-j}\) with
\(A=l+k\), \(B=n-k\), \(U=n+k\). Substitute this into \(S\) and swap the
order of summation:

\[
S=\sum_{j}\;\sum_{k}(-1)^k
\binom{l+m}{l+k}\binom{m+n}{m+k}\binom{l+k}{j}\binom{n-k}{\,n+k-j\,}.
\]

Now the inner sum over \(k\) has the shape needed for (★): it is a sum
in \(k\) of three binomial factors where one top is \(l+k\) (a
``\(r+t\)'' form) and another uses \(n-k\) (a ``\(n+r-s, n-t\)'' form).
Apply (★) to the \(k\)--sum with the parameter match

\[
(m,n,r,s,t)\ \leftarrow\ (\,m+k,\ l-k,\ m+n,\ n+l,\ j\,),
\]

i.e., we fix \(j\) and view the inner summation as the (★)--sum in the
variable \(k\). After the (routine) Vandermonde simplifications, the
entire dependence on \(k\) collapses to a single elementary alternating
binomial sum (next step), and a \(j\)--dependent constant factor pulls
out. (This is the ``very short proof'' Knuth hints at, and it's the
route sketched in the official answers.)

Concretely, one convenient way to write the outcome is

\[
\sum_{k}(-1)^k
\binom{l+m}{l+k}\binom{m+n}{m+k}\binom{l+k}{j}\binom{n-k}{n+k-j}
=
\frac{(m+n+j)!}{j!\,(m-l+j)!\,(n+j-l)!}\;
\sum_{k}(-1)^k\binom{2l-2j}{\,l-j+k\,}.
\tag{1}
\]

(You get the prefactor by rearranging factorials; the two
\((l\pm k-j)!\) combine into the single central binomial
\(\binom{2(l-j)}{(l-j)+k}\).)

Step 3: Kill almost all \(j\)'s with \((1-1)^N\)

Evaluate the remaining \(k\)--sum in (1):

\[
\sum_{k}(-1)^k\binom{2l-2j}{\,l-j+k\,}
\;=\;
\sum_{t}(-1)^{t-(l-j)}\binom{2l-2j}{\,t\,}
\;=\;(-1)^{l-j}(1-1)^{\,2l-2j}.
\]

Thus the sum is \(0\) unless \(2l-2j=0\), i.e.~\(j=l\); and when \(j=l\)
there is exactly one surviving \(k\)--term and the inner sum equals
\(1\).

Therefore only \(j=l\) contributes to \(S\). Plugging \(j=l\) into the
prefactor in (1) gives

\[
S \;=\; \frac{(m+n+l)!}{\,l!\,(m)!\,(n)!\,}\,,
\]

which is exactly the claimed right-hand side. ∎

Quick sanity check

Take \(l=m=n=1\). Then the sum runs over \(k=-1,0,1\):

\[
(-1)^ {-1}\!\binom{2}{0}\!\binom{2}{0}\!\binom{2}{0}
+(-1)^{0}\!\binom{2}{1}\!\binom{2}{1}\!\binom{2}{1}
+(-1)^{1}\!\binom{2}{2}\!\binom{2}{2}\!\binom{2}{2}
= -1+8-1=6,
\]

while \((1+1+1)!/(1!\,1!\,1!)=6\).

\subsection{63 {[}M30{]}}\label{m30-9}

We work with formal power series so all manipulations are valid for
integral parameters.

\begin{enumerate}
\def\labelenumi{\arabic{enumi}.}
\tightlist
\item
  Encode each binomial with a coefficient operator:
\end{enumerate}

\begin{itemize}
\tightlist
\item
  \(\displaystyle \binom{j+k}{k+l}=[x^{\,k+l}] (1+x)^{j+k}\).
\item
  \(\displaystyle \binom{r}{j}=[u^{\,j}] (1+u)^r\).
\item
  \(\displaystyle \binom{n}{k}=[v^{\,k}] (1+v)^n\).
\item
  \(\displaystyle \binom{s+n-j-k}{m-j}=[y^{\,m-j}] (1+y)^{\,s+n-j-k}\).
\end{itemize}

Insert these and interchange sums with coefficient operators:

\[
\begin{aligned}
S&=\sum_{j,k}(-1)^{j+k}[x^{k+l}](1+x)^{j+k}\,[u^j](1+u)^r\,[v^k](1+v)^n\,[y^{m-j}](1+y)^{s+n-j-k}\\
&=[y^{\,m}]\, (1+y)^{s+n} \sum_{j,k} [x^{k+l}](1+x)^{j+k}\,[u^j](1+u)^r\,[v^k](1+v)^n\,(-1)^j y^{-j}\,(-1)^k(1+y)^{-k}.
\end{aligned}
\]

\begin{enumerate}
\def\labelenumi{\arabic{enumi}.}
\setcounter{enumi}{1}
\tightlist
\item
  Sum over \(k\) first. Use
  \([x^{k+l}](1+x)^{j+k}=[x^{l}](1+x)^j\big((1+x)/x\big)^k\). Thus
\end{enumerate}

\[
\sum_{k\ge 0}(-1)^k\binom{n}{k}\Big(\frac{1+x}{x(1+y)}\Big)^{\!k}
=(1-\tfrac{1+x}{x(1+y)})^{n}=\Big(\frac{yx-1}{(1+y)x}\Big)^{n}.
\]

Therefore

\[
S=[y^{\,m}]\,(1+y)^{s}\,[x^{\,l+n}]\,(yx-1)^{n}\sum_{j\ge 0} [u^{\,j}](1+u)^r\,(-1)^j \Big(\frac{(1+x)y}{1+y}\Big)^{\!j}.
\]

\begin{enumerate}
\def\labelenumi{\arabic{enumi}.}
\setcounter{enumi}{2}
\tightlist
\item
  Now sum over \(j\) using \(\sum_j(-1)^j\binom{r}{j}t^j=(1-t)^r\):
\end{enumerate}

\[
\sum_{j\ge 0}(-1)^j\binom{r}{j}\Big(\frac{(1+x)y}{1+y}\Big)^{\!j}
=\Big(1-\frac{(1+x)y}{1+y}\Big)^{\!r}
=\Big(\frac{1-xy}{1+y}\Big)^{\!r}.
\]

Hence

\[
S=[y^{\,m}][x^{\,l+n}]\, (1+y)^{s-r}\,(yx-1)^{n}\,(1-xy)^{r}.
\]

\begin{enumerate}
\def\labelenumi{\arabic{enumi}.}
\setcounter{enumi}{3}
\tightlist
\item
  Extract the \(x\)-coefficient. Since \((yx-1)^n=(-1)^n(1-xy)^n\),
\end{enumerate}

\[
[x^{\,l+n}](yx-1)^{n}(1-xy)^{r}
=(-1)^n[x^{\,l+n}](1-xy)^{n+r}
=(-1)^n(-1)^{l+n}\binom{n+r}{\,l+n\,}y^{\,l+n}.
\]

Thus

\[
S=(-1)^l\binom{n+r}{\,l+n\,}[y^{\,m}]\,y^{\,l+n}(1+y)^{s-r}
=(-1)^l\binom{n+r}{\,n+l\,}\binom{s-r}{\,m-n-l\,},
\]

which is exactly the desired identity.

\subsection{64 {[}M20{]}}\label{m20-17}

Proof by induction via the standard recurrence

Let \(P(n,m)\) be the number of partitions of an \(n\)-element set into
\(m\) nonempty blocks.

Base values

\begin{itemize}
\tightlist
\item
  \(P(0,0)=1\) (the empty partition of the empty set).
\item
  \(P(n,0)=0\) for \(n>0\); \(P(n,m)=0\) for \(m>n\).
\end{itemize}

Recurrence

Fix a distinguished element, say \(n\). Any partition of
\(\{1,\dots,n\}\) into \(m\) blocks falls into exactly one of two
classes:

\begin{enumerate}
\def\labelenumi{\arabic{enumi}.}
\item
  \(n\) is alone in its own block. Then the remaining \(n-1\) elements
  form a partition into \(m-1\) blocks: \(P(n-1,m-1)\) choices.
\item
  \(n\) is not alone. First partition \(\{1,\dots,n-1\}\) into \(m\)
  blocks, then choose which block \(n\) joins. There are
  \(m\cdot P(n-1,m)\) choices.
\end{enumerate}

Hence

\[
P(n,m)=P(n-1,m-1)+m\,P(n-1,m).
\tag{\(*\)}
\]

But \(\bigl\{ {n \atop m} \bigr\}\), the Stirling number of the second
kind, is characterized by the same initial conditions and the same
recurrence

\[
\Bigl\{ {n \atop m} \Bigr\}=\Bigl\{ {n-1 \atop m-1} \Bigr\} + m\Bigl\{ {n-1 \atop m} \Bigr\},
\]

(which is one of the ``addition/recurrence'' formulas alluded to as Eq.
(46) in the text). Therefore \(P(n,m)=\bigl\{ {n \atop m} \bigr\}\) for
all \(n,m\).

Example: \(\displaystyle \Bigl\{ {4 \atop 2} \Bigr\}=7\)

All partitions of \(\{1,2,3,4\}\) into two nonempty blocks are:

\[
\{1,2,3\}\{4\},\ \{1,2,4\}\{3\},\ \{1,3,4\}\{2\},\ \{2,3,4\}\{1\},\ 
\{1,2\}\{3,4\},\ \{1,3\}\{2,4\},\ \{1,4\}\{2,3\},
\]

seven in total, as stated.

(Optional) Alternate derivation (surjections)

Count surjections \(f:[n]\twoheadrightarrow[m]\). There are
\(\sum_{j=0}^{m}(-1)^j\binom{m}{j}(m-j)^n\) of them by
inclusion--exclusion; each surjection corresponds to a partition
together with an ordering of the \(m\) blocks, so dividing by \(m!\)
yields

\[
\Bigl\{ {n \atop m} \Bigr\}=\frac{1}{m!}\sum_{j=0}^{m}(-1)^j\binom{m}{j}(m-j)^n,
\]

which again shows that \(\bigl\{ {n \atop m} \bigr\}\) counts
\(m\)-block partitions.

\subsection{65 {[}HM35{]}}\label{hm35-1}

Using the generalized Stirling numbers of both kinds (defined for
arbitrary real/complex top argument \(r\)), Knuth states:

\[
\boxed{z^{\,r}
 \;=\;
\sum_{k\ge 0}\Bigl\{\!\!\begin{smallmatrix} r\\ r-k\end{smallmatrix}\!\!\Bigr\}\,z^{\underline{\,r-k\,}}
}\qquad\text{for }\Re(z)>0, \tag{59}
\]

and the ``companion'' asymptotic series

\[
\boxed{z^{\,r}
\;=\;
\sum_{k=0}^{m}
\Bigl[\!\!\begin{smallmatrix} r\\ r-k\end{smallmatrix}\!\!\Bigr](-1)^k\,z^{\,r-k}
\;+\;O\!\left(z^{\,r-m-1}\right)
}\qquad (m=0,1,2,\dots), \tag{60}
\]

where
\(z^{\underline{\alpha}}:=\dfrac{\Gamma(z+1)}{\Gamma(z+1-\alpha)}\) is
the (gamma--extended) falling factorial. See the statements just after
Eq. (58) in the book.

Recall for integers \(n\) the exact conversion formulas

\[
x^{\underline{n}}=\sum_{k}(-1)^{\,n-k}\Bigl[\!\!\begin{smallmatrix} n\\ k\end{smallmatrix}\!\!\Bigr]x^{k},
\qquad
x^{n}=\sum_{k}\Bigl\{\!\!\begin{smallmatrix} n\\ k\end{smallmatrix}\!\!\Bigr]x^{\underline{k}},
\tag{44–45}
\]

and the inversion relations

\[
\sum_{k}\Bigl[\!\!\begin{smallmatrix} n\\ k\end{smallmatrix}\!\!\Bigr]
\Bigl\{\!\!\begin{smallmatrix} k\\ m\end{smallmatrix}\!\!\Bigr]
(-1)^{\,n-k}=\delta_{nm},
\qquad
\sum_{k}\Bigl\{\!\!\begin{smallmatrix} n\\ k\end{smallmatrix}\!\!\Bigr]
\Bigl[\!\!\begin{smallmatrix} k\\ m\end{smallmatrix}\!\!\Bigr]
(-1)^{\,n-k}=\delta_{nm}.
\tag{47}
\]

Proof of (59) (convergent series)

Step 1 (a basic generating identity). For any complex \(z\) and for
small \(t\),

\[
e^{zt}
= \bigl(1+(e^t-1)\bigr)^{z}
= \sum_{k\ge 0} \binom{z}{k}\,(e^t-1)^k
= \sum_{k\ge 0} \frac{z^{\underline{k}}}{k!}\,(e^t-1)^k.
\tag{A}
\]

Step 2 (EGF of Stirling numbers
\(\bigl\{\begin{smallmatrix}n\\k\end{smallmatrix}\bigr\}\)). For each
fixed (integer) \(k\ge 0\),

\[
(e^t-1)^k
= k!\sum_{n\ge k} \Bigl\{\!\!\begin{smallmatrix} n\\ k\end{smallmatrix}\!\!\Bigr]\,\frac{t^n}{n!}.
\tag{B}
\]

Substitute (B) into (A) and interchange the sums:

\[
e^{zt}=\sum_{n\ge 0}\left(\sum_{k=0}^{n}
\Bigl\{\!\!\begin{smallmatrix} n\\ k\end{smallmatrix}\!\!\Bigr]\,z^{\underline{k}}\right)\frac{t^n}{n!}.
\]

Comparing coefficients of \(t^n\) gives, for every integer \(n\ge 0\),

\[
z^{\,n}=\sum_{k=0}^{n}\Bigl\{\!\!\begin{smallmatrix} n\\ k\end{smallmatrix}\!\!\Bigr]\,z^{\underline{k}},
\]

which is precisely Eq. (45) rewritten.

Step 3 (extend \(n\to r\)). For noninteger \(r\) we define the
generalized Stirling numbers of the second kind by the same
Newton--difference formula

\[
\Bigl\{\!\!\begin{smallmatrix} r\\ k\end{smallmatrix}\!\!\Bigr]
=\frac{1}{k!}\,\Delta^k\!\bigl(x^r\bigr)\big|_{x=0}
=\frac{1}{k!}\sum_{j=0}^{k}(-1)^{k-j}\binom{k}{j}j^{\,r},
\]

which agrees with
\(\bigl\{\begin{smallmatrix}n\\k\end{smallmatrix}\bigr\}\) when
\(r=n\in\mathbb{N}\). Then the Newton series for \(x^r\) is

\[
x^r=\sum_{k\ge 0}\Bigl\{\!\!\begin{smallmatrix} r\\ k\end{smallmatrix}\!\!\Bigr]\,x^{\underline{k}}.
\]

A standard convergence theorem for Newton series shows that this series
converges for \(\Re(x)>0\). (Heuristically:
\(\binom{x}{k}=x^{\underline{k}}/k!\sim (-1)^k k^{-\Re(x)-1}/\Gamma(-x)\)
while
\(\Delta^k(x^r)=k!\bigl\{\begin{smallmatrix} r\\ k\end{smallmatrix}\bigr\}\)
grows only polynomially in \(k\); hence the terms behave like
\(k^{-\Re(x)-1+\text{(poly)}}\) and the series converges when
\(\Re(x)>0\).) Replacing \(x\) by \(z\) and reindexing \(k\mapsto r-k\)
gives Eq. (59) as stated in the book.

\begin{quote}
Thus (59) is nothing more than the Newton (forward-difference) expansion
of the function \(z\mapsto z^r\) in the falling-factorial basis, with
coefficients given by the generalized Stirling numbers of the second
kind.
\end{quote}

Proof of (60) (asymptotic companion involving first-kind numbers)

Write the falling factorial via gamma functions:

\[
z^{\underline{\alpha}}=\frac{\Gamma(z+1)}{\Gamma(z+1-\alpha)}.
\]

For fixed \(\alpha\) and \(z\to\infty\) in any right half-plane,
Stirling's asymptotic for the gamma ratio yields a full descending-power
expansion

\[
z^{\underline{\alpha}}
= z^{\alpha}\!\left(1
+ c_1(\alpha)z^{-1}
+ c_2(\alpha)z^{-2}
+ \cdots + c_m(\alpha)z^{-m}
+ O\!\left(z^{\alpha-m-1}\right)\right),
\]

whose coefficients \(c_j(\alpha)\) are known polynomials in \(\alpha\).
A classical (and very pretty) fact is that these polynomials can be
written with Stirling numbers of the first kind:

\[
z^{\underline{\alpha}}
= \sum_{j=0}^{m}(-1)^j
\Bigl[\!\!\begin{smallmatrix} \alpha\\ \alpha-j\end{smallmatrix}\!\!\Bigr]
z^{\alpha-j}
\;+\;O\!\left(z^{\alpha-m-1}\right)
\qquad(z\to\infty,\ \Re z>0).
\tag{C}
\]

(Equation (44) for integer \(\alpha\) is exact; (C) is its gamma-based
asymptotic extension.)

Now plug (C) with \(\alpha=r-k\) into the convergent series (59):

\[
\begin{aligned}
z^{\,r}
&= \sum_{k\ge 0}\Bigl\{\!\!\begin{smallmatrix} r\\ r-k\end{smallmatrix}\!\!\Bigr]\,
z^{\underline{\,r-k\,}}
\\
&= \sum_{k\ge 0}\Bigl\{\!\!\begin{smallmatrix} r\\ r-k\end{smallmatrix}\!\!\Bigr]\,
\left(
\sum_{j=0}^{m}(-1)^j
\Bigl[\!\!\begin{smallmatrix} r-k\\ r-k-j\end{smallmatrix}\!\!\Bigr]
z^{\,r-j}
\;+\;O\!\left(z^{\,r-k-m-1}\right)\right).
\end{aligned}
\]

Because (59) converges uniformly in any sector \(\Re z\ge \sigma>0\), we
may truncate and interchange the (finite) inner sum with the outer sum,
obtaining

\[
z^{\,r}
= \sum_{j=0}^{m}(-1)^j
\underbrace{\left(\sum_{k\ge 0}
\Bigl\{\!\!\begin{smallmatrix} r\\ r-k\end{smallmatrix}\!\!\Bigr]\,
\Bigl[\!\!\begin{smallmatrix} r-k\\ r-k-j\end{smallmatrix}\!\!\Bigr]\right)}_{\displaystyle =
\Bigl[\!\!\begin{smallmatrix} r\\ r-j\end{smallmatrix}\!\!\Bigr]\text{ by the inversion relations}}
z^{\,r-j}
\;+\;O\!\left(z^{\,r-m-1}\right).
\]

The evaluation of the bracketed sum as
\(\Bigl[\!\!\begin{smallmatrix} r\\ r-j\end{smallmatrix}\!\!\Bigr]\) is
exactly the Stirling inversion in Eq. (47), applied in its (polynomial)
extension to general \(r\). This yields the claimed asymptotic expansion
(60).

\subsection{66 {[}HM30{]}}\label{hm30-2}

Let real numbers \(x,y,z\) and an integer \(n\ge 2\) satisfy

\[
\binom{x}{n}=\binom{y}{n}+\binom{z}{\,n-1\,},\qquad
x\ge n-1,\; y\ge n-1,\; z>n-2.
\]

Prove the two ``if and only if'' statements:

\[
\binom{x}{n-1}\le \binom{y}{n-1}+\binom{z}{\,n-2\,}\ \Longleftrightarrow\ y\ge z,
\]

\[
\binom{x}{n+1}\le \binom{y}{n+1}+\binom{z}{\,n\,}\ \Longleftrightarrow\ y\le z.
\]

We will prove the second inequality first (it's the key step), then
deduce the first.

Notation and standard identities

Write

\[
X=\binom{x}{n},\quad Y=\binom{y}{n},\quad Z=\binom{z}{n-1}.
\]

For ``neighbors'' in the upper index we use:

\[
\binom{t}{n-1}=\frac{n}{t-n+1}\binom{t}{n},\qquad
\binom{t}{n+1}=\frac{t-n}{n+1}\binom{t}{n}\qquad(t\in\{x,y\}),
\]

and for \(Z\) (whose lower index is \(n-1\)):

\[
Z_-=\binom{z}{n-2}=\frac{n-1}{z-n+2}\,Z,\quad
Z_+=\binom{z}{n}=\frac{z-n+1}{n}\,Z.
\]

Also, for \(t>n-1\),

\[
\frac{d}{dt}\binom{t}{n}=\binom{t}{n}\!\sum_{j=0}^{n-1}\frac1{t-j}\,.
\]

Because \(\binom{t}{n}\) is strictly increasing for \(t\ge n-1\), the
hypothesis \(X=Y+Z\) forces \(x>y\) unless \(Z=0\) (which can't happen
since \(z>n-2\)).

Part B (the ``\(n+1\)'' inequality): sign is determined by the order of
\(y\) vs \(z\)

Define

\[
F(z)\;=\;\underbrace{\binom{x}{n+1}}_{=:X_+}\;-\;\underbrace{\binom{y}{n+1}}_{=:Y_+}\;-\;\underbrace{\binom{z}{n}}_{=:Z_+},
\]

where \(x\) depends on \(z\) via the constraint \(X=Y+Z\). Using the
neighbor identities and \(X=Y+Z\),

\[
F(z)=\frac{x-n}{n+1}(Y+Z)-\frac{y-n}{n+1}Y-\frac{z-n+1}{n}\,Z.
\]

If \(F(z)=0\) at some \(z>n-2\), since \(x>y\) and \(Y,Z>0\) we must
have

\[
\frac{x-n}{n+1}<\frac{z-n+1}{n}.
\tag{★}
\]

Now differentiate implicitly \(X=Y+Z\) with respect to \(z\). Using

\[
\frac{X'}{X}=\sum_{j=0}^{n-1}\frac1{x-j},\qquad
\frac{Z'}{Z}=\sum_{j=0}^{n-2}\frac1{z-j},
\]

and comparing termwise via (★), one gets

\[
\frac{X'}{X}>\frac{n}{n+1}\cdot\frac{Z'}{Z}\quad\Longrightarrow\quad
\frac{dx}{dz}=\frac{Z'}{X'}<\frac{n+1}{n}\cdot\frac{Z}{X}.
\]

A short calculation (differentiate \(F\) and substitute the last bound
together with (★)) yields

\[
F'(z)<0\quad\text{at every zero of }F.
\]

Hence \(F\) has at most one root for \(z>n-2\), and the function crosses
that root with negative slope. But when \(z=y\), our constraint gives
\(x=y+1\) (since \(X=Y+\binom{y}{n-1}=\binom{y+1}{n}\)), and then the
Pascal relation \(\binom{y+1}{n+1}=\binom{y}{n+1}+\binom{y}{n}\) shows
\(F(y)=0\). Therefore that root is unique and occurs exactly at \(z=y\).
Consequently,

\[
F(z)\le 0\ \Longleftrightarrow\ z\ge y,
\]

with equality iff \(z=y\). This is exactly

\[
\binom{x}{n+1}\le \binom{y}{n+1}+\binom{z}{n}\ \Longleftrightarrow\ y\le z,
\]

as required. (This is essentially Knuth's solution idea.) 

Part A (the ``\(n-1\)'' inequality) follows by the same idea with
\(n\mapsto n-1\)

Consider

\[
G(z)\;=\;\binom{x}{n-1}-\binom{y}{n-1}-\binom{z}{n-2}.
\]

Exactly the same analysis as above applies to the triple
\(\bigl(\binom{x}{n-1},\binom{y}{n-1},\binom{z}{n-2}\bigr)\), provided
\(n-1\ge 2\), i.e., \(n>2\). (That ensures the same
monotonicity/uniqueness argument goes through.) One again finds that
\(G\) has a unique zero at \(z=y\) and \(G'(z)<0\) there; hence

\[
G(z)\le 0\ \Longleftrightarrow\ z\le y,
\]

i.e.,

\[
\binom{x}{n-1}\le \binom{y}{n-1}+\binom{z}{n-2}\ \Longleftrightarrow\ y\ge z.
\]

(For \(n=2\) the statement is still true and can be checked directly.)

\subsection{67 {[}M20{]}}\label{m20-18}

For integers \(n\ge k\ge 1\), prove the useful bounds

\[
\boxed{\left(\frac{n}{k}\right)^{\!k}\ \le\ \binom{n}{k}\ \le\ \left(\frac{e\,n}{k}\right)^{\!k}}\, .
\]

Lower bound \(\displaystyle \binom{n}{k}\ge \left(\frac{n}{k}\right)^k\)

Write

\[
\binom{n}{k}=\frac{n(n-1)\cdots(n-k+1)}{k!}
=\prod_{i=1}^{k}\frac{n-k+i}{i}.
\]

For each \(i\in\{1,\dots,k\}\),

\[
\frac{n-k+i}{i}\ \ge\ \frac{n}{k}
\quad\Longleftrightarrow\quad
k(n-k+i)\ \ge\ n i
\quad\Longleftrightarrow\quad
kn-k^2\ \ge\ (n-k)i,
\]

which holds because \(i\le k\) implies \((n-k)i\le (n-k)k=kn-k^2\).
Multiplying these \(k\) inequalities gives

\[
\binom{n}{k}\ \ge\ \Bigl(\frac{n}{k}\Bigr)^k.
\]

\emph{(Intuition: in each factor \((n-k+i)/i\), the numerator is at
least linearly growing while the denominator is at most \(k\); every
factor is \(\ge n/k\), so the product is \(\ge (n/k)^k\).)}

Upper bound
\(\displaystyle \binom{n}{k}\le \left(\frac{e\,n}{k}\right)^k\)

First,

\[
\binom{n}{k}=\frac{n(n-1)\cdots(n-k+1)}{k!}\ \le\ \frac{n^k}{k!}.
\]

Next, use the classical inequality (a one-sided Stirling bound)

\[
k!\ \ge\ \left(\frac{k}{e}\right)^{\!k},
\]

which is Exercise 1.2.5--24 (or follows immediately from
\(\ln k! = \sum_{i=1}^k \ln i \ge k\ln k - k\)). Substituting yields

\[
\binom{n}{k}\ \le\ \frac{n^k}{(k/e)^k}
=\left(\frac{e\,n}{k}\right)^{\!k}.
\]

Knuth notes a slightly tighter chain on the right:

\[
\binom{n}{k}
=\frac{n^{\underline{k}}}{k!}
\le \frac{n^k}{k!}
\le \frac1e\left(\frac{e\,n}{k}\right)^{\!k}
\le \left(\frac{e\,n}{k+1}\right)^{\!k},
\]

using \((1+1/k)^k<e\) for the last inequality. (These are ``slightly
sharper but less memorable.'')

Remarks

\begin{itemize}
\tightlist
\item
  The two bounds are often good enough for quick size estimates of
  \(\binom{n}{k}\) without invoking full Stirling's formula.
\item
  If you need asymptotics, Stirling's approximation refines
  \(\binom{n}{k}\) to within a multiplicative \(1+o(1)\) factor; but the
  exercise only asks for the clean double-sided power bounds.
\end{itemize}

\subsection{68 {[}M25{]}}\label{m25-14}

Let \(K\sim\mathrm{Bin}(n,p)\) with \(n\in\mathbb{Z}_{\ge0}\),
\(p\in[0,1]\), and put \(m=\lceil np\rceil\). Prove

\[
\sum_{k=0}^{n} \binom{n}{k} p^{k}(1-p)^{\,n-k}\,\lvert k-np\rvert
\;=\;
2\,m\,\binom{n}{m}\,p^{m}(1-p)^{\,n+1-m}.
\tag{★}
\]

(Equivalently, \(\mathbb{E}\,\lvert K-np\rvert\) equals the right--hand
side.)

\begin{enumerate}
\def\labelenumi{\arabic{enumi})}
\tightlist
\item
  Reduce \(\mathbb{E}|K-np|\) to a one-sided tail
\end{enumerate}

Let \(m=\lceil np\rceil\). Split the sum at \(k=m\). Because
\(\mathbb{E}(K-np)=0\),

\[
\sum_{k\le m-1} (np-k)\Pr[K=k]
\;=\;
\sum_{k\ge m} (k-np)\Pr[K=k].
\]

Hence

\[
\mathbb{E}\,\lvert K-np\rvert
= \sum_{k\le m-1} (np-k)\Pr[K=k] \;+\; \sum_{k\ge m} (k-np)\Pr[K=k]
= 2\,\sum_{k\ge m} (k-np)\Pr[K=k]. \tag{1}
\]

Write \(a_k=\binom{n}{k}p^k(1-p)^{n-k}=\Pr[K=k]\) and denote

\[
S(p)\;:=\;\sum_{k=m}^{n} (k-np)\,a_k .
\]

Then (1) is

\[
\mathbb{E}\,\lvert K-np\rvert \;=\; 2\,S(p). \tag{2}
\]

\begin{enumerate}
\def\labelenumi{\arabic{enumi})}
\setcounter{enumi}{1}
\tightlist
\item
  Turn the weighted tail into a derivative
\end{enumerate}

Differentiate \(a_k\) with respect to \(p\):

\[
\frac{d}{dp}a_k
= a_k\!\left(\frac{k}{p}-\frac{n-k}{1-p}\right)
= \frac{k-np}{p(1-p)}\,a_k .
\]

Therefore

\[
S(p)
= \sum_{k=m}^n (k-np)a_k
= p(1-p)\,\sum_{k=m}^n \frac{d}{dp}a_k
= p(1-p)\,\frac{d}{dp}\!\Bigg(\sum_{k=m}^n a_k\Bigg).
\tag{3}
\]

Let \(T_m(p):=\sum_{k=m}^n a_k=\Pr[K\ge m]\). Then
\(S(p)=p(1-p)T_m'(p)\).

\begin{enumerate}
\def\labelenumi{\arabic{enumi})}
\setcounter{enumi}{2}
\tightlist
\item
  Evaluate \(T_m'(p)\) by telescoping
\end{enumerate}

Differentiate \(T_m\) termwise and rearrange:

\[
\begin{aligned}
T_m'(p)
&= \sum_{k=m}^{n} \binom{n}{k}\!\left(k\,p^{k-1}(1-p)^{n-k}-(n-k)\,p^{k}(1-p)^{n-k-1}\right)\\
&= \sum_{k=m}^{n} k\binom{n}{k} p^{k-1}(1-p)^{n-k}
 \;-\!\!\!\sum_{k=m}^{n-1} (n-k)\binom{n}{k} p^{k}(1-p)^{n-k-1}.
\end{aligned}
\]

Use \((n-k)\binom{n}{k}=(k+1)\binom{n}{k+1}\) and shift the index in the
second sum:

\[
T_m'(p)
= \sum_{k=m}^{n} k\binom{n}{k} p^{k-1}(1-p)^{n-k}
  - \sum_{j=m+1}^{n} j\binom{n}{j} p^{j-1}(1-p)^{n-j}
= m\binom{n}{m} p^{m-1}(1-p)^{n-m}.
\tag{4}
\]

\begin{enumerate}
\def\labelenumi{\arabic{enumi})}
\setcounter{enumi}{3}
\tightlist
\item
  Combine
\end{enumerate}

From (3) and (4),

\[
S(p)=p(1-p)\,T_m'(p)
= m\binom{n}{m} p^{m}(1-p)^{n+1-m}.
\]

Insert this into (2) to obtain (★), which completes the proof.

Remarks

\begin{itemize}
\tightlist
\item
  The identity is valid at the endpoints as well: if \(p=0\) or \(p=1\),
  then \(m\in\{0,n\}\) and both sides are 0.
\item
  Conceptually, (★) says the mean absolute deviation of a binomial from
  its mean is twice a boundary term at the ``median cut''
  \(m=\lceil np\rceil\): the calculation is essentially a one-line
  telescoping after recognizing the derivative trick.
\end{itemize}

Tiny sanity check

For \(n=1\), \(K\in\{0,1\}\), \(m=\lceil p\rceil\).

\[
\mathbb{E}|K-p|=p(1-p)+(1-p)p=2p(1-p),
\]

and the right--hand side of (★) is
\(2\cdot 1 \cdot \binom{1}{1} p^1(1-p)^{1}=2p(1-p)\).

\bookmarksetup{startatroot}

\chapter{1.2.7 Harmonic numbers}\label{harmonic-numbers}

\subsection{1 {[}01{]}}\label{section-82}

Here ``harmonic numbers'' are defined by \(H_n=\sum_{k=1}^{n}\frac1k\)
for integers \(n\ge 0\). (So \(H_0\) is the empty sum.)

\begin{itemize}
\tightlist
\item
  \(H_0\). The sum has no terms, so by the standard convention for empty
  sums it equals \(0\). This also matches the recurrence
  \(H_n=H_{n-1}+\frac1n\): plugging \(n=1\) gives \(H_1=H_0+1\), so
  \(H_0\) must be \(0\).
\item
  \(H_1\). One term: \(H_1=\frac11=1\).
\item
  \(H_2\). Two terms: \(H_2=\frac11+\frac12=1+\tfrac12=\tfrac32=1.5\).
\end{itemize}

\subsection{2 {[}13{]}}\label{section-83}

Proof (matching the ``simple argument'' style)

Observe how \(H_{2m}\) changes when \(m\) increases by 1:

\[
H_{2m+2}-H_{2m}\;=\;\frac1{2m+1}+\frac1{2m+2}
\;<\; \frac1{2m+1}+\frac1{2m+1}\;=\;\frac{2}{2m+1}\;\le\;1.
\]

Thus, each time we extend from \(2m\) to \(2m+2\) terms, the harmonic
sum rises by less than 1.

Use induction on \(m\).

\begin{itemize}
\item
  Base case. \(H_2=1+\tfrac12=\tfrac32\le 2=1+1\).
\item
  Inductive step. If \(H_{2m}\le 1+m\), then

  \[
  H_{2(m+1)}=H_{2m+2}\;=\;H_{2m}+\bigl(H_{2m+2}-H_{2m}\bigr)
  \;\le\; (1+m)+1 \;=\; 2+m \;=\;1+(m+1).
  \]
\end{itemize}

Hence \(H_{2m}\le 1+m\) for all \(m\ge 1\).

This mirrors the text's earlier ``step-size'' reasoning for the lower
bound \(H_{2m}\ge 1+\frac m2\) (there, each step from \(2m\) to \(2m+2\)
adds at least \(\tfrac12\)); here we note each step adds at most 1.

One-line alternative (even slightly stronger)

After the first term, every summand is \(\le \tfrac12\):

\[
H_{2m}
=1+\sum_{k=2}^{2m}\frac1k
\;\le\; 1+(2m-1)\cdot \frac12
\;=\; m+\tfrac12
\;\le\; 1+m.
\]

This yields \(H_{2m}\le m+\tfrac12\), which is stronger than \(1+m\).

\subsection{3 {[}M21{]}}\label{m21-8}

For \(r>1\), show the generalized harmonic partial sums

\[
H_n^{(r)}\;=\;\sum_{k=1}^{n}\frac{1}{k^{\,r}}
\]

remain bounded for all \(n\), and find an upper bound.

Solution (dyadic-block argument, generalizing Ex. 2)

Group the terms into blocks whose lengths double each time (powers of
two), exactly as in the proof for ordinary harmonic numbers:

\begin{itemize}
\item
  For \(j\ge 0\), let the \(j\)-th block be the indices

  \[
  B_j=\{2^{\,j}+1,\;2^{\,j}+2,\;\dots,\;\min(2^{\,j+1},\,n)\}.
  \]

  Every \(k\in B_j\) satisfies \(k>2^{\,j}\), hence

  \[
  \frac{1}{k^{\,r}}\;<\;\frac{1}{(2^{\,j})^{r}}\,.
  \]
\item
  The block \(B_j\) contains at most \(2^{\,j}\) terms, therefore the
  contribution of \(B_j\) to the sum is bounded by

  \[
  \sum_{k\in B_j}\frac{1}{k^{\,r}}\;\le\;2^{\,j}\cdot \frac{1}{(2^{\,j})^{r}}
  \;=\;2^{\,j(1-r)}.
  \]
\end{itemize}

Now split off the \(k=1\) term and sum these bounds over all blocks up
to \(\lfloor \log_2 n\rfloor\):

\[
H_n^{(r)}=\frac{1}{1^{\,r}}+\sum_{j=0}^{\lfloor\log_2 n\rfloor}\;\sum_{k\in B_j}\frac{1}{k^{\,r}}
\;\le\;1+\sum_{j=0}^{\lfloor\log_2 n\rfloor}2^{\,j(1-r)}
\;\le\;1+\sum_{j=0}^{\infty}2^{\,j(1-r)}.
\]

Because \(r>1\), the ratio \(q=2^{\,1-r}\) satisfies \(0<q<1\), so the
geometric series converges:

\[
\sum_{j=0}^{\infty}2^{\,j(1-r)}=\frac{1}{1-2^{\,1-r}}=\frac{2^{\,r-1}}{\,2^{\,r-1}-1\,}.
\]

Therefore, for every \(n\ge 1\),

\[
\boxed{\;H_n^{(r)}\;\le\;1+\frac{1}{1-2^{\,1-r}}
\;=\;1+\frac{2^{\,r-1}}{\,2^{\,r-1}-1\,}\;}
\]

which is a constant depending only on \(r\) (hence the sequence of
partial sums is bounded).

Examples:

\begin{itemize}
\tightlist
\item
  \(r=2:\;H_n^{(2)}\le 1+\frac{1}{1-\tfrac12}=3\).
\item
  \(r=3:\;H_n^{(3)}\le 1+\frac{1}{1-\tfrac14}=\tfrac73\).
\end{itemize}

\subsection{4 {[}10{]}}\label{section-84}

Preliminaries

\begin{itemize}
\tightlist
\item
  Harmonic numbers are defined by \(H_n=\sum_{k=1}^n \tfrac1k\) (with
  \(H_0=0\)). 
\item
  Euler's constant is \(\gamma=\lim_{n\to\infty}(H_n-\ln n)\) (so
  \(H_n=\ln n+\gamma+o(1)\)).
\end{itemize}

Two standard integral bounds

Because \(x\mapsto 1/x\) is decreasing on \([1,\infty)\),

\[
\int_{k}^{k+1}\frac{dx}{x}\;\le\;\frac1k\;\le\;\int_{k-1}^{k}\frac{dx}{x}\qquad(k\ge1).
\]

Summing:

\[
\boxed{\ \ln(n+1)\le H_n\le 1+\ln n\ \ } \qquad(n\ge1).
\]

From these alone we already see \(H_n>\ln n\). (We'll use these bounds
in parts (a)--(b).)

Decisions

\begin{enumerate}
\def\labelenumi{(\alph{enumi})}
\tightlist
\item
  \(H_n<\ln n\) --- False.
\end{enumerate}

Take \(n=1\): \(H_1=1>\ln 1=0\). In fact, the lower bound above shows
\(H_n\ge \ln(n+1)>\ln n\) for every \(n\ge1\), so the claimed inequality
never holds.

\begin{enumerate}
\def\labelenumi{(\alph{enumi})}
\setcounter{enumi}{1}
\tightlist
\item
  \(H_n>\ln n\) --- True.
\end{enumerate}

From the same bound,

\[
H_n\;\ge\;\ln(n+1)\;>\;\ln n\qquad(n\ge1).
\]

\begin{enumerate}
\def\labelenumi{(\alph{enumi})}
\setcounter{enumi}{2}
\tightlist
\item
  \(H_n>\ln n+\gamma\) --- True.
\end{enumerate}

Let \(\gamma_n:=H_n-\ln n\). Then

\[
\gamma_{n+1}-\gamma_n
= \Bigl(H_n+\frac1{n+1}-\ln(n+1)\Bigr)-(H_n-\ln n)
= \frac1{n+1}-\ln\!\Bigl(1+\frac1n\Bigr).
\]

Since for \(x>0\) the function \(f(x)=\ln(1+x)-\frac{x}{1+x}\) satisfies
\(f'(x)=\frac{x}{(1+x)^2}>0\) and \(f(0)=0\), we have

\[
\ln(1+x)>\frac{x}{1+x}\quad(x>0).
\]

With \(x=1/n\) this gives \(\ln(1+\tfrac1n)>\tfrac1{n+1}\), hence
\(\gamma_{n+1}-\gamma_n<0\). So \(\{\gamma_n\}\) is strictly decreasing.
Because \(\gamma=\lim_{n\to\infty}\gamma_n\) exists, every term lies
above the limit:

\[
\gamma_n>\gamma\qquad(n\ge1).
\]

Equivalently, \(H_n>\ln n+\gamma\) for all \(n\). (This is also
consistent with the asymptotic expansion
\(H_n=\ln n+\gamma+\tfrac1{2n}-\tfrac1{12n^2}+\tfrac1{120n^4}-\varepsilon\)
with \(0<\varepsilon<\tfrac1{252n^6}\), which shows the excess over
\(\ln n+\gamma\) is positive. )

\subsection{5 {[}15{]}}\label{section-85}

Knuth gives the Euler--Maclaurin expansion for \(H_n\):

\[
H_n=\ln n+\gamma+\frac{1}{2n}-\frac{1}{12n^2}+\frac{1}{120n^4}-\varepsilon,\qquad
0<\varepsilon<\frac{1}{252\,n^6}.
\]

This makes high-precision evaluation straightforward for large \(n\). 

Computation for \(n=10{,}000\)

Use constants from Appendix A (or any high-precision source):

\begin{itemize}
\tightlist
\item
  \(\ln 10000 = 4\ln 10 = 9.210340371976182736\ldots\)
\item
  \(\gamma = 0.577215664901532860\ldots\)
\end{itemize}

Now add the correction terms:

\[
\begin{aligned}
H_{10000}
&=\underbrace{(\ln 10000+\gamma)}_{9.787556036877715596\ldots}
+\underbrace{\frac{1}{2n}}_{0.000050000000000000}
-\underbrace{\frac{1}{12n^2}}_{0.000000000833333333\ldots}
+\underbrace{\frac{1}{120n^4}}_{0.0000000000000000008\ldots}\\[2pt]
&= 9.7876060360443822641784779\ldots
\end{aligned}
\]

The neglected remainder is
\(<\dfrac{1}{252\cdot 10^{24}}\approx 4.0\times10^{-27}\), far below the
\(10^{-15}\) level, so rounding at 15 decimals is safe.

\subsection{6 {[}M15{]}}\label{m15-10}

Proof A (via the Stirling-number recurrence; ``algebraic/inductive'')

Facts we'll use (from §1.2.6):

\begin{enumerate}
\def\labelenumi{\arabic{enumi}.}
\tightlist
\item
  The unsigned Stirling numbers of the first kind convert falling powers
  to ordinary powers:
\end{enumerate}

\[
(x)_{m}=\sum_{k=0}^{m}(-1)^{\,m-k}\,\bigl[\!\begin{smallmatrix} m\\ k\end{smallmatrix}\!\bigr]\,x^k.
\]

From \((x)_{n+1}=(x-n)(x)_n\), comparing coefficients yields the
standard recurrence

\[
\bigl[\!\begin{smallmatrix} n+1\\ k\end{smallmatrix}\!\bigr]
= n\,\bigl[\!\begin{smallmatrix} n\\ k\end{smallmatrix}\!\bigr]
+ \bigl[\!\begin{smallmatrix} n\\ k-1\end{smallmatrix}\!\bigr].
\tag{★}
\]

Also
\(\bigl[\!\begin{smallmatrix} n\\ 1\end{smallmatrix}\!\bigr]=(n-1)!\). 

Define
\(a_n:=\dfrac{\bigl[\!\begin{smallmatrix} n+1\\ 2\end{smallmatrix}\!\bigr]}{n!}\).
Apply (★) with \(k=2\):

\[
\bigl[\!\begin{smallmatrix} n+1\\ 2\end{smallmatrix}\!\bigr]
= n\,\bigl[\!\begin{smallmatrix} n\\ 2\end{smallmatrix}\!\bigr]
+ \bigl[\!\begin{smallmatrix} n\\ 1\end{smallmatrix}\!\bigr]
= n\,\bigl[\!\begin{smallmatrix} n\\ 2\end{smallmatrix}\!\bigr]
+ (n-1)! .
\]

Divide by \(n!\):

\[
a_n
= \frac{\bigl[\!\begin{smallmatrix} n\\ 2\end{smallmatrix}\!\bigr]}{(n-1)!}
+ \frac{1}{n}
= a_{n-1} + \frac{1}{n}.
\]

The initial value is
\(a_1=\dfrac{\bigl[\!\begin{smallmatrix}2\\2\end{smallmatrix}\!\bigr]}{1!}=1\).
Hence \(a_n=1+\frac12+\cdots+\frac1n=H_n\) by induction, proving the
claim.

Proof B (combinatorial; ``count permutations with two cycles'')

A classical interpretation:
\(\bigl[\!\begin{smallmatrix} m\\ k\end{smallmatrix}\!\bigr]\) counts
permutations of \(m\) elements with exactly \(k\) cycles in the cycle
decomposition. Thus
\(\bigl[\!\begin{smallmatrix} n+1\\ 2\end{smallmatrix}\!\bigr]\) counts
permutations of \(\{1,\dots,n+1\}\) having two cycles. 

Count such permutations by the length \(\ell\) (with
\(1\le \ell \le n\)) of the cycle that contains \(n+1\):

\begin{itemize}
\tightlist
\item
  Choose the other \(\ell-1\) members of that cycle:
  \(\binom{n}{\ell-1}\) ways.
\item
  Arrange those \(\ell\) elements in a cycle with \(n+1\) present:
  \((\ell-1)!\) ways.
\item
  Arrange the remaining \(n+1-\ell\) elements into one cycle:
  \((n-\ell)!\) ways.
\end{itemize}

Total:

\[
\sum_{\ell=1}^{n}
\binom{n}{\ell-1}(\ell-1)!(n-\ell)!
=\sum_{j=0}^{n-1}\binom{n}{j}j!(n-j-1)!
=\sum_{j=0}^{n-1}\frac{n!}{n-j}
= n!\sum_{k=1}^{n}\frac{1}{k}
= n!\,H_n.
\]

Hence
\(\bigl[\!\begin{smallmatrix} n+1\\ 2\end{smallmatrix}\!\bigr]=n!\,H_n\),
i.e.~\(H_n=\bigl[\!\begin{smallmatrix} n+1\\ 2\end{smallmatrix}\!\bigr]/n!\).

\subsection{7 {[}M21{]}}\label{m21-9}

\begin{enumerate}
\def\labelenumi{(\alph{enumi})}
\tightlist
\item
  Monotonicity in each argument
\end{enumerate}

Fix \(n\) and compare \(T(m+1,n)\) to \(T(m,n)\):

\[
\begin{aligned}
T(m+1,n)-T(m,n)
&=\bigl(H_{m+1}-H_m\bigr)-\bigl(H_{(m+1)n}-H_{mn}\bigr)\\
&=\frac{1}{m+1}-\sum_{k=1}^{n}\frac{1}{mn+k}.
\end{aligned}
\]

Since the \(n\) denominators \(mn+1,\dots,mn+n\) are at most \((m+1)n\),
we have

\[
\sum_{k=1}^{n}\frac{1}{mn+k}\;\ge\; n\cdot\frac{1}{(m+1)n}\;=\;\frac{1}{m+1}.
\]

Hence \(T(m+1,n)-T(m,n)\le 0\). Equality holds iff \(n=1\); otherwise
the sum is strictly larger than \(\frac1{m+1}\), so the difference is
strictly negative. By symmetry, the same argument (swapping \(m\) and
\(n\)) shows \(T(m,n)\) is also nonincreasing in \(n\) (strictly unless
\(m=1\)).

Conclusion. \(T\) is nonincreasing in each variable (strictly unless the
other variable equals \(1\)).

\begin{enumerate}
\def\labelenumi{(\alph{enumi})}
\setcounter{enumi}{1}
\tightlist
\item
  Extremal values for \(m,n>0\)
\end{enumerate}

Maximum. For any \(n\ge 1\),

\[
T(1,n)=H_1+H_n-H_{1\cdot n}=1+H_n-H_n=1,
\]

and by part (a), increasing either argument cannot increase \(T\). Thus
the maximum value is \(1\), attained whenever \(m=1\) or \(n=1\) (in
particular at \((1,1)\)).

Minimum / infimum. Using the asymptotic expansion from the section,

\[
H_m=\ln m+\gamma+R_m,\quad H_n=\ln n+\gamma+R_n,\quad H_{mn}=\ln(mn)+\gamma+R_{mn},
\]

so

\[
T(m,n)=\gamma+\bigl(R_m+R_n-R_{mn}\bigr).
\]

From the displayed expansion (with positive leading term
\(\tfrac{1}{2k}\)), \(R_k>0\) and \(R_k\to 0\) as \(k\to\infty\); more
precisely,

\[
R_k=\frac{1}{2k}-\frac{1}{12k^2}+O\!\left(\frac{1}{k^4}\right).
\]

Therefore for finite \(m,n\),

\[
T(m,n)\;=\;\gamma+\left(\frac{1}{2m}+\frac{1}{2n}-\frac{1}{2mn}\right)+O\!\left(\frac{1}{m^2}+\frac{1}{n^2}\right) \;>\;\gamma,
\]

and letting \(m,n\to\infty\) gives \(T(m,n)\downarrow \gamma\).

\begin{itemize}
\tightlist
\item
  There is no minimum over positive integers (the value \(\gamma\) is
  not attained by any finite \(m,n\)).
\item
  The infimum is \(\gamma\approx 0.57721\), and \(T(m,n)\in(\gamma,1]\)
  for all integers \(m,n\ge 1\).
\end{itemize}

\subsection{8 {[}HM18{]}}\label{hm18}

Estimate, as a function of \(n\), the difference

\[
\Delta_n \;=\; \sum_{k=1}^n H_k \;-\; \sum_{k=1}^n \ln k \;=\; \sum_{k=1}^n H_k \;-\; \ln(n!).
\]

Step 1: Reduce with Eq. (8) Section 1.2.7 establishes

\[
\sum_{k=1}^n H_k \;=\; (n+1)H_n - n. 
\]

Hence

\[
\Delta_n \;=\; (n+1)H_n \;-\; n \;-\; \ln(n!). \tag{★}
\]

Step 2: Use standard asymptotics for \(H_n\) and \(\ln(n!)\)

From 1.2.7 (derived via Euler's summation formula in 1.2.11.2), for
fixed \(m\) we have

\[
H_n \;=\; \ln n + \gamma + \frac{1}{2n} - \frac{1}{12n^2} + O\!\left(\frac{1}{n^4}\right).
\]

(Here \(\gamma\) is Euler's constant.) 

Stirling's expansion (also obtained in 1.2.11.2) gives

\[
\ln(n!) \;=\; \Bigl(n+\tfrac12\Bigr)\ln n \;-\; n \;+\; \tfrac12\ln(2\pi)
\;+\; \frac{1}{12n} \;-\; \frac{1}{360n^3} \;+\; O\!\left(\frac{1}{n^5}\right).
\]

(Using \(e^{\sigma}=\sqrt{2\pi}\) for Stirling's constant.) 

Step 3: Combine in (★)

Insert the two expansions into (★) and simplify:

\[
\begin{aligned}
\Delta_n
&= (n+1)\!\left(\ln n+\gamma+\frac{1}{2n}-\frac{1}{12n^2}+O\!\left(\frac{1}{n^4}\right)\right)
\;-\; n \\
&\qquad -\left(\Bigl(n+\tfrac12\Bigr)\ln n - n + \tfrac12\ln(2\pi)
+ \frac{1}{12n} - \frac{1}{360n^3} + O\!\left(\frac{1}{n^5}\right)\right) \\[4pt]
&= \underbrace{\tfrac12\ln n}_{\text{from logs}} 
\;+\; \underbrace{(n+1)\gamma}_{\text{linear}}
\;-\; \tfrac12\ln(2\pi)
\;+\; \underbrace{\Bigl(\tfrac12 + \frac{1}{3n}\Bigr)}_{\text{from small terms}}
\;+\; O\!\left(\frac{1}{n^2}\right).
\end{aligned}
\]

Thus a convenient and accurate estimate is

\[
\boxed{\;
\Delta_n \;=\; n\gamma \;+\; \tfrac12\ln n \;+\; \Bigl(\gamma + \tfrac12 - \tfrac12\ln(2\pi)\Bigr)
\;+\; \frac{1}{3n} \;+\; O\!\left(\frac{1}{n^2}\right)
\;}
\]

so the leading growth is linear with slope \(\gamma\), with a secondary
\(\tfrac12\ln n\) term.

Numerically, the constant

\[
C \;=\; \gamma + \tfrac12 - \tfrac12\ln(2\pi) \;\approx\; 0.158277132\ldots
\]

Hence, for practical purposes,

\[
\Delta_n \;\approx\; n\gamma + \tfrac12\ln n + 0.158277132 + \frac{1}{3n}.
\]

Takeaway

Comparing Eq. (8) with \(\sum_{k=1}^n \ln k\) shows

\[
\sum_{k=1}^n H_k - \sum_{k=1}^n \ln k 
\;=\; n\gamma + \tfrac12\ln n + \Theta(1),
\]

with the constant and first correction explicitly given above.

\subsection{9 {[}M18{]}}\label{m18-6}

Theorem A in §1.2.7 analyzes the sum
\(\sum_{k=1}^{n} \binom{n}{k} x^{k} H_k\) for \(x>0\). The exercise asks
for the value when \(x=-1\); i.e.,

\[
S_n\;=\;\sum_{k=1}^{n} \binom{n}{k}(-1)^k H_k\;=? 
\]

Solution 1 (double--sum ``H\_k expansion'')

Write \(H_k=\sum_{j=1}^{k}\frac1j\), then interchange the order of
summation:

\[
S_n=\sum_{k=1}^{n}\binom{n}{k}(-1)^k\!\sum_{j=1}^{k}\frac1j
=\sum_{j=1}^{n}\frac1j\sum_{k=j}^{n}\binom{n}{k}(-1)^k .
\]

Use the standard partial alternating--binomial identity

\[
\sum_{k=j}^{n}(-1)^k\binom{n}{k}=(-1)^j\binom{n-1}{\,j-1\,},
\]

(which follows, e.g., from
\(\sum_{k=0}^{m}(-1)^k\binom{n}{k}=(-1)^m\binom{n-1}{m}\) by induction
or Pascal's rule). Hence

\[
S_n=\sum_{j=1}^{n}\frac{(-1)^j}{j}\binom{n-1}{\,j-1\,}.
\]

Now use \(\binom{n}{j}=\frac{n}{j}\binom{n-1}{\,j-1\,}\), giving

\[
S_n=\frac1n\sum_{j=1}^{n}(-1)^j\binom{n}{j}.
\]

Finally, \(\sum_{j=0}^{n}(-1)^j\binom{n}{j}=(1-1)^n=0\), so
\(\sum_{j=1}^{n}(-1)^j\binom{n}{j}=-1\). Therefore

\[
\boxed{\,S_n=-\frac1n\,}.
\]

Solution 2 (integral representation of \(H_k\))

Use \(H_k=\displaystyle\int_{0}^{1}\frac{1-t^{k}}{1-t}\,dt\) (from
integrating \(\sum_{i=0}^{k-1}t^i=\frac{1-t^k}{1-t}\) termwise). Then

\[
S_n=\int_{0}^{1}\frac{1}{1-t}\sum_{k=1}^{n}\binom{n}{k}(-1)^k(1-t^k)\,dt.
\]

But \(\sum_{k=1}^{n}\binom{n}{k}(-1)^k=-1\) and
\(\sum_{k=1}^{n}\binom{n}{k}(-t)^k=(1-t)^n-1\). Hence the bracket equals

\[
-1-\bigl((1-t)^n-1\bigr)=-(1-t)^n,
\]

so

\[
S_n=-\int_{0}^{1}(1-t)^{n-1}\,dt=-\frac{1}{n}.
\]

Thus, although Theorem A is stated only for \(x>0\), the alternating
case at \(x=-1\) has the exact closed form

\[
\boxed{\displaystyle \sum_{k=1}^{n}\binom{n}{k}(-1)^k H_k=-\frac{1}{n}}.
\]

(Checks: \(n=1,2,3\) give \(-1,-\tfrac12,-\tfrac13\), respectively.)

\subsection{10 {[}M20{]}}\label{m20-19}

Proof (telescoping via a discrete product rule)

Start from the discrete analogue of the product rule, obtained by adding
and subtracting \(a_{k+1}b_k\):

\[
a_{k+1}b_{k+1}-a_k b_k
= a_{k+1}(b_{k+1}-b_k) + (a_{k+1}-a_k)b_k.
\]

Now sum this equality for \(k=1,2,\dots,n-1\). The left-hand side
telescopes:

\[
\sum_{k=1}^{n-1}\bigl(a_{k+1}b_{k+1}-a_k b_k\bigr)
= a_n b_n - a_1 b_1.
\]

Hence

\[
a_n b_n - a_1 b_1
= \sum_{k=1}^{n-1} a_{k+1}(b_{k+1}-b_k)
  + \sum_{k=1}^{n-1} (a_{k+1}-a_k)b_k.
\]

Rearrange to isolate the desired sum:

\[
\sum_{k=1}^{n-1} (a_{k+1}-a_k)b_k
= a_n b_n - a_1 b_1
  - \sum_{k=1}^{n-1} a_{k+1}(b_{k+1}-b_k).
\]

This is exactly the required identity. ∎

Remarks and equivalent forms

\begin{enumerate}
\def\labelenumi{\arabic{enumi}.}
\tightlist
\item
  Abel's summation (discrete integration by parts). If
  \(A_k=\sum_{j=1}^k a_j\) and we set \(a_{k+1}-a_k:=a_{k+1}\), the
  identity becomes the well-known Abel transform:
\end{enumerate}

\[
\sum_{k=1}^{n} a_k\, b_k
= A_n b_n - \sum_{k=1}^{n-1} A_k\,(b_{k+1}-b_k).
\]

\begin{enumerate}
\def\labelenumi{\arabic{enumi}.}
\setcounter{enumi}{1}
\tightlist
\item
  Symmetric variant. By shifting indices one also gets
\end{enumerate}

\[
\sum_{k=1}^{n} a_k\,(b_k-b_{k-1})
= a_n b_n - \sum_{k=1}^{n-1} b_k\,(a_{k+1}-a_k),
\]

with the convention \(b_0\) as needed.

\begin{enumerate}
\def\labelenumi{\arabic{enumi}.}
\setcounter{enumi}{2}
\tightlist
\item
  Quick sanity check. Let \(a_k=k\), \(b_k=k\). Then the left side is
  \(\sum_{k=1}^{n-1} k = n(n-1)/2\). The right side evaluates to
  \(n^2-1-\sum_{k=1}^{n-1}(k+1)=n(n-1)/2\) as well.
\end{enumerate}

These forms are what §1.2.7 uses implicitly when deriving identities
like Eq. (9) and Theorem A's estimate for \(\sum \binom{n}{k}x^k H_k\).

\subsection{11 {[}M21{]}}\label{m21-10}

Solution (summation by parts)

Use the discrete summation-by-parts identity (proved in 1.2.7--10):

\[
\sum_{k=2}^{n} a_k\,(b_k-b_{k-1})
\;=\; a_n b_n - a_1 b_1 \;-\; \sum_{k=1}^{n-1} b_k\,(a_{k+1}-a_k).
\]

Take \(a_k = H_k\) and choose \(b_k=-\frac1k\). Then

\[
b_k-b_{k-1}=-\frac1k+\frac1{k-1}=\frac{1}{k(k-1)},
\]

so the left side is exactly \(S_n\). Also,

\[
a_{k+1}-a_k = H_{k+1}-H_k=\frac1{k+1}.
\]

Plugging into the identity gives

\[
S_n
= H_n\!\left(-\frac1n\right) - H_1\!\left(-1\right)
   - \sum_{k=1}^{n-1}\!\left(-\frac1k\right)\!\left(\frac1{k+1}\right)
= 1 - \frac{H_n}{n} + \sum_{k=1}^{n-1}\frac{1}{k(k+1)}.
\]

The remaining sum telescopes:

\[
\sum_{k=1}^{n-1}\frac{1}{k(k+1)}
=\sum_{k=1}^{n-1}\!\left(\frac1k-\frac1{k+1}\right)
=1-\frac1n.
\]

Therefore

\[
\boxed{\,S_n
= 2 - \frac{H_n+1}{n}\, }.
\]

Quick check

For \(n=2\): \(S_2=\frac{H_2}{2}= \frac{3/2}{2}=3/4\). Formula gives
\(2-\frac{H_2+1}{2}=2-\frac{2.5}{2}=0.75\).

\subsection{12 {[}M10{]}}\label{m10-11}

For \(r>1\),

\[
H^{(r)}_\infty=\sum_{k\ge 1}\frac{1}{k^r}=\zeta(r),
\]

Riemann's zeta function. In particular, for even integers there's the
closed form \(\zeta(2m)=\dfrac{|B_{2m}|(2\pi)^{2m}}{2(2m)!}\), though we
won't need to compute it here. 

We only need 100 correct decimals. Since

\[
\zeta(1000)=1+\sum_{k=2}^{\infty}\frac{1}{k^{1000}},
\]

it suffices to show the tail \(\sum_{k=2}^{\infty}k^{-1000}\) is
\(<10^{-100}\). Then \(\zeta(1000)\) rounds to
\(1.\underbrace{00\ldots0}_{100\text{ zeros}}\).

A tight and elementary bound:

\[
\sum_{k=2}^{\infty}\frac{1}{k^{1000}}
\;<\;\frac{1}{2^{1000}}\;+\;\int_{2}^{\infty}x^{-1000}\,dx
\;=\;\frac{1}{2^{1000}}+\frac{1}{999\cdot 2^{999}}
=\frac{1}{2^{1000}}\!\left(1+\frac{2}{999}\right).
\]

Now \(2^{1000}=10^{1000\log_{10}2}\) with \(\log_{10}2\approx 0.30103\),
hence

\[
\frac{1}{2^{1000}}<10^{-301}.
\]

Therefore

\[
0<\zeta(1000)-1<\left(1+\frac{2}{999}\right)\!10^{-301}<10^{-300}\ll 10^{-100}.
\]

So the first 100 digits after the decimal point of \(\zeta(1000)\) are
all zero.

(Indeed, much more is true: \(\zeta(1000)=1+\Theta(2^{-1000})\), so the
first nonzero decimal digit doesn't appear until roughly the
\(302^\text{nd}\) place.)

\subsection{13 {[}M22{]}}\label{m22-11}

Here is a clean, self-contained proof of Exercise 13 from §1.2.7
(Harmonic Numbers). I restate the problem in my own words for clarity:

Claim. For every integer \(n\ge 1\) and indeterminate \(x\),

\[
\sum_{k=1}^{n}\frac{x^{k}}{k}
\;=\;
H_n \;+\; \sum_{k=1}^{n}\binom{n}{k}\frac{(x-1)^k}{k},
\]

where \(H_n=\sum_{k=1}^{n}\frac1k\) is the \(n\)-th harmonic number. 

Proof

\begin{enumerate}
\def\labelenumi{\arabic{enumi}.}
\tightlist
\item
  Split off the constant part.
\end{enumerate}

For any \(x\),

\[
\sum_{k=1}^{n}\frac{x^k}{k}
=\sum_{k=1}^{n}\frac{1}{k}+\sum_{k=1}^{n}\frac{x^k-1}{k}
=H_n+\sum_{k=1}^{n}\frac{x^k-1}{k}.
\]

\begin{enumerate}
\def\labelenumi{\arabic{enumi}.}
\setcounter{enumi}{1}
\tightlist
\item
  Binomially expand \(x^k-1\).
\end{enumerate}

Write \(x=(x-1)+1\). Then

\[
x^k-1=\bigl((x-1)+1\bigr)^k-1
=\sum_{j=0}^{k}\binom{k}{j}(x-1)^j-1
=\sum_{j=1}^{k}\binom{k}{j}(x-1)^j.
\]

Hence

\[
\sum_{k=1}^{n}\frac{x^k-1}{k}
=\sum_{k=1}^{n}\sum_{j=1}^{k}\binom{k}{j}\frac{(x-1)^j}{k}
=\sum_{j=1}^{n}(x-1)^j\sum_{k=j}^{n}\frac{1}{k}\binom{k}{j}.
\]

\begin{enumerate}
\def\labelenumi{\arabic{enumi}.}
\setcounter{enumi}{2}
\tightlist
\item
  Evaluate the inner sum.
\end{enumerate}

Use the simple identity

\[
\frac{1}{k}\binom{k}{j}=\frac{1}{j}\binom{k-1}{j-1}
\qquad(k\ge j\ge 1),
\]

which follows from \(\binom{k}{j}=\dfrac{k}{j}\binom{k-1}{j-1}\).
Therefore

\[
\sum_{k=j}^{n}\frac{1}{k}\binom{k}{j}
=\frac{1}{j}\sum_{k=j}^{n}\binom{k-1}{j-1}
=\frac{1}{j}\sum_{m=j-1}^{n-1}\binom{m}{j-1}
=\frac{1}{j}\binom{n}{j},
\]

by the ``hockey-stick'' identity
\(\sum_{m=r}^{N}\binom{m}{r}=\binom{N+1}{r+1}\).

\begin{enumerate}
\def\labelenumi{\arabic{enumi}.}
\setcounter{enumi}{3}
\tightlist
\item
  Assemble the result.
\end{enumerate}

Plugging this into step 2 gives

\[
\sum_{k=1}^{n}\frac{x^k-1}{k}
=\sum_{j=1}^{n}(x-1)^j\cdot \frac{1}{j}\binom{n}{j}
=\sum_{j=1}^{n}\binom{n}{j}\frac{(x-1)^j}{j}.
\]

Together with step 1, this yields exactly

\[
\sum_{k=1}^{n}\frac{x^{k}}{k}
=H_n+\sum_{k=1}^{n}\binom{n}{k}\frac{(x-1)^k}{k},
\]

as was to be shown.

\subsection{14 {[}M22{]}}\label{m22-12}

Part (a)
\(\displaystyle \sum_{k=1}^n \frac{H_k}{k}=\frac12\bigl(H_n^2+H_n^{(2)}\bigr)\)

Write the left--hand side as a double sum and use symmetry:

\[
\sum_{k=1}^n \frac{H_k}{k}
=\sum_{k=1}^n\sum_{i=1}^k \frac{1}{ik}
=\sum_{1\le i\le k\le n}\frac{1}{ik}.
\]

Split the pairs \((i,k)\) into \(i<k\) and \(i=k\):

\[
\sum_{1\le i\le k\le n}\frac{1}{ik}
=\underbrace{\sum_{1\le i<k\le n}\frac{1}{ik}}_{\text{strictly off–diagonal}}
+\underbrace{\sum_{k=1}^n\frac{1}{k^2}}_{H_n^{(2)}}.
\]

On the other hand,

\[
H_n^2
=\Bigl(\sum_{i=1}^n \frac{1}{i}\Bigr)\Bigl(\sum_{k=1}^n \frac{1}{k}\Bigr)
=\sum_{i=1}^n\frac{1}{i^2}+2\sum_{1\le i<k\le n}\frac{1}{ik}
=H_n^{(2)}+2\sum_{1\le i<k\le n}\frac{1}{ik}.
\]

Hence
\(\displaystyle \sum_{1\le i<k\le n}\frac{1}{ik}=\frac{H_n^2-H_n^{(2)}}{2}\),
and substituting above gives

\[
\sum_{k=1}^n \frac{H_k}{k}
=\frac{H_n^2-H_n^{(2)}}{2}+H_n^{(2)}
=\frac{H_n^2+H_n^{(2)}}{2}.
\]

This proves the identity.

(Equivalently, you can also start from
\(\sum_{k=1}^n\sum_{i=1}^k \frac{1}{ik} =\sum_{i=1}^n \frac{1}{i}\sum_{k=i}^n \frac{1}{k} =\sum_{i=1}^n \frac{1}{i}(H_n-H_{i-1}) =H_n^2-\sum_{i=1}^n \frac{H_{i-1}}{i}\),
then note
\(\sum_{i=1}^n \frac{H_{i-1}}{i}=\sum_{i=1}^n \frac{H_i}{i}-\sum_{i=1}^n \frac{1}{i^2}=\sum_{k=1}^n \frac{H_k}{k}-H_n^{(2)}\)
and solve for \(\sum \frac{H_k}{k}\).)

Part (b) \(\displaystyle \sum_{k=1}^n \frac{H_k}{k+1}\)

Use the elementary shift \(H_k=H_{k+1}-\frac{1}{k+1}\):

\[
\sum_{k=1}^n \frac{H_k}{k+1}
=\sum_{k=1}^n \frac{H_{k+1}}{k+1}-\sum_{k=1}^n\frac{1}{(k+1)^2}
=\Bigl(\sum_{j=2}^{n+1}\frac{H_j}{j}\Bigr)-\Bigl(\sum_{j=2}^{n+1}\frac{1}{j^2}\Bigr).
\]

Hence

\[
\sum_{k=1}^n \frac{H_k}{k+1}
=\underbrace{\sum_{j=1}^{n+1}\frac{H_j}{j}}_{=\frac12\,(H_{n+1}^2+H_{n+1}^{(2)})}-\frac{H_1}{1}
-\bigl(H_{n+1}^{(2)}-1\bigr)
=\frac12\bigl(H_{n+1}^2-H_{n+1}^{(2)}\bigr).
\]

(Used the result of part (a) with \(n\mapsto n+1\).)

An equivalent form in terms of \(H_n\) is also handy:

\[
\sum_{k=1}^n \frac{H_k}{k+1}
=\frac12\left(H_n^2-H_n^{(2)}\right)+\frac{H_n}{n+1},
\]

since \(H_{n+1}=H_n+\frac{1}{n+1}\) and
\(H_{n+1}^{(2)}=H_n^{(2)}+\frac{1}{(n+1)^2}\).

\subsection{15 {[}M23{]}}\label{m23-13}

Recall \(H_k=\sum_{j=1}^k \frac1j\). Start from the double--sum
expansion:

\[
\sum_{k=1}^{n} H_k^{\,2}
=\sum_{k=1}^{n}\Bigl(\sum_{i=1}^{k}\tfrac1i\Bigr)\Bigl(\sum_{j=1}^{k}\tfrac1j\Bigr)
=\sum_{i=1}^{n}\sum_{j=1}^{n}\frac{\#\{k:1\le k\le n,\ k\ge \max(i,j)\}}{ij}
=\sum_{i=1}^{n}\sum_{j=1}^{n}\frac{n-\max(i,j)+1}{ij}.
\]

Split the square \((i,j)\) into diagonal and off--diagonal parts:

\[
\sum_{k=1}^{n} H_k^{\,2}
=2\sum_{1\le j<i\le n}\frac{n-i+1}{ij}
+\sum_{i=1}^{n}\frac{n-i+1}{i^2}
=2\sum_{i=1}^{n}\frac{n-i+1}{i}\,H_{i-1}+\sum_{i=1}^{n}\frac{n-i+1}{i^2}.
\]

Use \(H_{i-1}=H_i-\frac1i\):

\[
\sum_{k=1}^{n} H_k^{\,2}
=2\sum_{i=1}^{n}\frac{n-i+1}{i}\,H_i
-2\sum_{i=1}^{n}\frac{n-i+1}{i^2}
+\sum_{i=1}^{n}\frac{n-i+1}{i^2}.
\]

Hence

\[
\boxed{\ \sum_{k=1}^{n} H_k^{\,2}
=2\sum_{i=1}^{n}\frac{n-i+1}{i}\,H_i
-\sum_{i=1}^{n}\frac{n-i+1}{i^2}\ } \tag{★}
\]

Now evaluate the two sums in \((★)\).

\begin{enumerate}
\def\labelenumi{\arabic{enumi}.}
\tightlist
\item
  For the first,
\end{enumerate}

\[
\sum_{i=1}^{n}\frac{n-i+1}{i}\,H_i
=(n+1)\sum_{i=1}^{n}\frac{H_i}{i}-\sum_{i=1}^{n}H_i.
\]

We'll use two standard identities (proved here for completeness):

\begin{itemize}
\item
  \(\displaystyle \sum_{i=1}^{n}H_i=(n+1)H_n-n\). Proof:
  \(\sum_{i=1}^{n}H_i=\sum_{i=1}^{n}\sum_{j=1}^{i}\frac1j =\sum_{j=1}^{n}\frac{n-j+1}{j}=(n+1)H_n-\sum_{j=1}^{n}1=(n+1)H_n-n.\)
\item
  \(\displaystyle \sum_{i=1}^{n}\frac{H_i}{i}=\tfrac12\!\left(H_n^{\,2}+H_n^{(2)}\right)\),
  where \(H_n^{(2)}=\sum_{i=1}^{n}\frac1{i^2}\). Proof: From
  \(H_i=H_{i-1}+\frac1i\),

  \[
  H_n^{\,2}-H_{n-1}^{\,2}=\left(H_n-H_{n-1}\right)\!\left(H_n+H_{n-1}\right)
  =\frac{2H_{n-1}}{n}+\frac1{n^2}.
  \]

  Summing \(n=1\ldots N\) gives
  \(\displaystyle H_N^{\,2} =2\sum_{i=1}^{N}\frac{H_{i-1}}{i}+\sum_{i=1}^{N}\frac1{i^2} =2\sum_{i=1}^{N}\frac{H_i}{i}-\sum_{i=1}^{N}\frac1{i^2},\)
  hence the claim.
\end{itemize}

Therefore

\[
\sum_{i=1}^{n}\frac{n-i+1}{i}\,H_i
=(n+1)\cdot \tfrac12\!\left(H_n^{\,2}+H_n^{(2)}\right) - \bigl((n+1)H_n-n\bigr).
\]

\begin{enumerate}
\def\labelenumi{\arabic{enumi}.}
\setcounter{enumi}{1}
\tightlist
\item
  For the second,
\end{enumerate}

\[
\sum_{i=1}^{n}\frac{n-i+1}{i^2}
=(n+1)\sum_{i=1}^{n}\frac1{i^2}-\sum_{i=1}^{n}\frac1i
=(n+1)H_n^{(2)}-H_n.
\]

Plug these into \((★)\):

\[
\begin{aligned}
\sum_{k=1}^{n} H_k^{\,2}
&=2\Bigl((n+1)\cdot \tfrac12(H_n^{\,2}+H_n^{(2)})-(n+1)H_n+n\Bigr)
-\bigl((n+1)H_n^{(2)}-H_n\bigr)\\
&=(n+1)H_n^{\,2}-2(n+1)H_n+2n - (n+1)H_n^{(2)}+H_n + (n+1)H_n^{(2)}\\
&=\boxed{(n+1)H_n^{\,2}-(2n+1)H_n+2n.}
\end{aligned}
\]

Quick check

\begin{itemize}
\tightlist
\item
  \(n=1\): LHS \(=H_1^2=1\). RHS \(=(2)\cdot1-(3)\cdot1+2=1\).
\item
  \(n=2\): LHS \(=1^2+(1+\tfrac12)^2=1+2.25=3.25\). RHS
  \(=3\cdot2.25-5\cdot1.5+4=3.25\).
\end{itemize}

Thus,

\[
\boxed{\displaystyle \sum_{k=1}^{n} H_k^{\,2}=(n+1)H_n^{\,2}-(2n+1)H_n+2n.}
\]

\subsection{16 {[}18{]}}\label{section-86}

\begin{enumerate}
\def\labelenumi{\arabic{enumi}.}
\item
  Split \(H_{2n}\) into even and odd parts.

  \[
  H_{2n}=\sum_{k=1}^{2n}\frac1k=\sum_{j=1}^{n}\left(\frac{1}{2j-1}+\frac{1}{2j}\right).
  \]
\item
  Recognize the even part.

  \[
  \sum_{j=1}^{n}\frac{1}{2j}=\frac12\sum_{j=1}^{n}\frac1j=\frac12\,H_n.
  \]
\item
  Solve for the odd part. Let \(S_n=\sum_{j=1}^{n}\frac{1}{2j-1}\). Then

  \[
  H_{2n}=S_n+\frac12 H_n
  \quad\Longrightarrow\quad
  S_n=H_{2n}-\frac12\,H_n.
  \]
\end{enumerate}

That's the required expression ``in terms of harmonic numbers.''

Quick checks

\begin{itemize}
\tightlist
\item
  \(n=1\): \(1=H_2-\tfrac12H_1=(1+\tfrac12)-\tfrac12=1\).
\item
  \(n=2\):
  \(1+\tfrac13=\tfrac43=H_4-\tfrac12H_2=(1+\tfrac12+\tfrac13+\tfrac14)-\tfrac12(1+\tfrac12)=\tfrac{25}{12}-\tfrac34=\tfrac{16}{12}=\tfrac43\).
\end{itemize}

(Optional) Asymptotics

Using \(H_m=\ln m+\gamma+\frac{1}{2m}+O(m^{-2})\),

\[
\sum_{j=1}^{n}\frac{1}{2j-1}
=H_{2n}-\frac12H_n
=\ln 2+\frac12\ln n+\frac{\gamma}{2}+O\!\left(\frac{1}{n}\right),
\]

so the odd--reciprocal sum grows like \(\tfrac12\ln n\) plus a constant.

Answer.
\(\displaystyle \sum_{j=1}^{n}\frac{1}{2j-1}=H_{2n}-\frac12\,H_n.\)

\subsection{17 {[}M24{]}}\label{m24-7}

Solution 1 (pairing \(k\) with \(p-k\))

Pair the terms \(k\) and \(p-k\) for \(k=1,2,\dots,(p-1)/2\). Then

\[
\frac1k+\frac1{p-k} \;=\; \frac{p}{k(p-k)}.
\]

Therefore

\[
H_{p-1}
= \sum_{k=1}^{(p-1)/2}\!\left(\frac1k+\frac1{p-k}\right)
= \sum_{k=1}^{(p-1)/2}\frac{p}{k(p-k)}.
\]

Each denominator \(k(p-k)\) is not divisible by \(p\) (since
\(1\le k \le p-1\)). Let

\[
D \;:=\; \operatorname{lcm}\{\,k(p-k): 1\le k\le (p-1)/2\,\}.
\]

Then \(p\nmid D\). Writing the sum over the common denominator \(D\), we
get

\[
H_{p-1} \;=\; \frac{p}{D}\sum_{k=1}^{(p-1)/2}\frac{D}{k(p-k)} \;=\; \frac{p\cdot N}{D},
\]

for some integer \(N\). Since \(p\nmid D\), the fraction
\(\dfrac{pN}{D}\) (and hence its reduction to lowest terms) has
numerator divisible by \(p\). This proves the claim.

Solution 2 (field-theoretic view, same conclusion)

Work modulo \(p\). In the field \(\mathbb{F}_p\), every
\(k\in\{1,\dots,p-1\}\) is invertible, so

\[
S \;:=\; \sum_{k=1}^{p-1} k^{-1}\in\mathbb{F}_p
\]

is well-defined. The map \(k\mapsto k^{-1}\) permutes
\(\{1,\dots,p-1\}\), hence

\[
S \equiv \sum_{k=1}^{p-1} k \equiv 0 \pmod p.
\]

Interpreting \(H_{p-1}=A/B\) with \(p\nmid B\) (justified by the pairing
argument above), the congruence \(S\equiv 0\) means
\(A/B\equiv 0\pmod p\), so \(A\equiv 0\pmod p\). Thus \(p\mid A\).

Checks \& remarks

\begin{itemize}
\tightlist
\item
  Examples: \(p=5\): \(H_4=\tfrac{25}{12}\) (numerator \(25\) divisible
  by \(5\)). \(p=7\): \(H_6=\tfrac{49}{20}\) (numerator \(49\) divisible
  by \(7\)).
\item
  Stronger fact (context): for \(p\ge 5\), Wolstenholme's theorem says
  \(H_{p-1}\equiv 0\pmod{p^2}\), i.e., the numerator is even divisible
  by \(p^2\). (Exercise 17 only asks for divisibility by \(p\).)
\end{itemize}

\subsection{18 {[}M33{]}}\label{m33}

Let

\[
S_n=\sum_{k=1}^{n}\frac1{2k-1}=1+\frac13+\frac15+\cdots+\frac1{2n-1}
\]

and write \(S_n=\dfrac{A_n}{B_n}\) in lowest terms. Since every
denominator \(2k-1\) is odd, also \(B_n\) is odd. The problem asks for
the highest power of \(2\) dividing the numerator \(A_n\).

Answer

\[
\nu_2(A_n)=
\begin{cases}
0,& n\ \text{odd};\\[2mm]
2\,\nu_2(n),& n\ \text{even},
\end{cases}
\]

where \(\nu_2(x)\) is the exponent of \(2\) in \(x\) (the 2-adic
valuation).

We work with the common odd denominator

\[
D_n=\operatorname{lcm}\{1,3,5,\dots,2n-1\}\quad(\text{so }D_n\text{ is odd}),
\]

and set

\[
N_n=\sum_{k=1}^{n}\frac{D_n}{2k-1}\in\mathbb{Z},\qquad S_n=\frac{N_n}{D_n}.
\]

Since \(B_n\mid D_n\) and \(B_n\) is odd, \(\nu_2(A_n)=\nu_2(N_n)\).

\begin{enumerate}
\def\labelenumi{\arabic{enumi})}
\tightlist
\item
  The case \(n\) odd.
\end{enumerate}

Each summand \(D_n/(2k-1)\) is an odd integer (odd divided by odd),
hence \(N_n\) is a sum of \(n\) odd integers. If \(n\) is odd, the sum
is odd; thus \(\nu_2(N_n)=0\).

\begin{enumerate}
\def\labelenumi{\arabic{enumi})}
\setcounter{enumi}{1}
\tightlist
\item
  A pairing identity for general \(n\).
\end{enumerate}

Pair the terms \(k\) and \(n+1-k\):

\[
\frac{D_n}{2k-1}+\frac{D_n}{2n-2k+1}
=D_n\cdot\frac{(2n)}{(2k-1)(2n-2k+1)}.
\]

Both factors in the denominator are odd, hence each pair contributes a
multiple of \(2n\). Therefore

\[
N_n=2n\cdot T_n,\qquad
T_n=\sum_{k=1}^{\lfloor n/2\rfloor}\frac{D_n}{(2k-1)(2n-2k+1)}\in\mathbb{Z}.
\]

Consequently,

\[
\nu_2(N_n)\ge \nu_2(2n)=\nu_2(n)+1.
\]

We now sharpen this to an equality \(\nu_2(N_n)=2\nu_2(n)\) for even
\(n\).

\begin{enumerate}
\def\labelenumi{\arabic{enumi})}
\setcounter{enumi}{2}
\tightlist
\item
  The case \(n=2^a m\) with \(m\) odd (so \(a=\nu_2(n)\ge1\)).
\end{enumerate}

Write \(n=2^a m\). In the sum that defines \(T_n\), index the
\(n/2=2^{a-1}m\) terms by

\[
k=r+t\cdot 2^{a-1}\qquad(1\le r\le 2^{a-1},\; 0\le t\le m-1).
\]

For fixed \(r\), the paired denominators are

\[
(2k-1,\;2n-2k+1)\equiv (2r-1,\;2^{a+1}-2r+1)\pmod{2^{a+1}},
\]

two odd numbers whose sum is \(2^{a+1}\). Modulo \(2^{a+1}\), their
product is \(-\,(2r-1)^2\). Because \(D_n\) is odd and independent of
\(k\), reduction modulo \(2^{a+1}\) gives

\[
\frac{D_n}{(2k-1)(2n-2k+1)}\equiv C\cdot (2r-1)^{-2}\pmod{2^{a+1}}
\]

for some fixed odd constant \(C\) (the inverse exists since all numbers
involved are odd).

Thus summing over the \(m\) values of \(t\) (a complete set of congruent
classes modulo \(2^{a+1}\)) yields

\[
T_n \equiv m\,C \sum_{r=1}^{2^{a-1}} (2r-1)^{-2}\pmod{2^{a+1}}.
\]

The map \(x\mapsto x^{-1}\) permutes the odd residues modulo
\(2^{a+1}\); hence

\[
\sum_{r=1}^{2^{a-1}} (2r-1)^{-2}\equiv \sum_{r=1}^{2^{a-1}} (2r-1)^2
= \sum_{j=0}^{2^{a-1}-1}(2j+1)^2
= 2^{a-1}\bigl(2^{a}-1\bigr)\equiv 2^{a-1}\pmod{2^{a}}.
\]

(Here we used the closed form for the sum of the first \(2^{a-1}\) odd
squares and observed that its residue modulo \(2^{a}\) is \(2^{a-1}\).)

Therefore

\[
T_n \equiv m\,C\,2^{a-1}\pmod{2^{a}},
\]

and since \(m\) and \(C\) are odd, this congruence shows exactly

\[
\nu_2(T_n)=a-1.
\]

Combining with \(N_n=2n\,T_n\) gives

\[
\nu_2(N_n)=\nu_2(2n)+\nu_2(T_n)=(a+1)+(a-1)=2a=2\,\nu_2(n),
\]

as claimed.

This completes the proof.

Small checks

\[
\begin{array}{c|c|c}
n & S_n & \nu_2(A_n) \\
\hline
1 & 1=\frac{1}{1} & 0\\
2 & 1+\frac13=\frac{4}{3} & 2 \ (=2\nu_2(2))\\
3 & \frac{23}{15} & 0\\
4 & \frac{176}{105} & 4 \ (=2\nu_2(4))\\
6 & \frac{6508}{3465} & 2 \ (=2\nu_2(6))\\
8 & \frac{91072}{45045} & 6 \ (=2\nu_2(8))
\end{array}
\]

So the highest power of \(2\) dividing the numerator of
\(1+\frac{1}{3}+\cdots+\frac{1}{2n-1}\) is \(2^{\,2\,\nu_2(n)}\) (and no
higher power).

\subsection{19 {[}M30{]}}\label{m30-10}

Answer

\(\boxed{n=0 \text{ or } n=1}\). For every \(n\ge 2\), \(H_n\) is not an
integer.

Proof (parity/LCM trick)

Let \(n\ge 2\) and set

\[
L=\operatorname{lcm}(1,2,\dots,n).
\]

Let \(2^k\) be the largest power of \(2\) with \(2^k\le n<2^{k+1}\).
Then the exact power of \(2\) dividing \(L\) is \(2^k\).

Consider the integer

\[
S \;=\; L\,H_n \;=\; \sum_{j=1}^{n}\frac{L}{j}.
\]

Look at each term modulo \(2\):

\begin{itemize}
\tightlist
\item
  If \(j=2^k\), then \(\dfrac{L}{j}=\dfrac{L}{2^k}\) is odd, because we
  have removed the full factor \(2^k\) from \(L\).
\item
  If \(j\neq 2^k\), then the highest power of \(2\) in \(j\) is at most
  \(2^{k-1}\); therefore \(\dfrac{L}{j}\) still contains at least one
  factor \(2\) and is even.
\end{itemize}

Hence all terms in the sum are even except one, so \(S\) is odd.

But \(L\) is even (since \(n\ge 2\) implies \(2\mid L\)). If \(H_n\)
were an integer, then \(S=L\cdot H_n\) would be even---contradiction.
Therefore \(H_n\) cannot be an integer for any \(n\ge 2\).

Finally, the small cases are \(H_0=0\) and \(H_1=1\), both integers.
Thus the only nonnegative \(n\) with \(H_n\in\mathbb{Z}\) are \(n=0\)
and \(n=1\).

\subsection{20 {[}HM22{]}}\label{hm22-2}

\begin{enumerate}
\def\labelenumi{\arabic{enumi})}
\tightlist
\item
  An integral representation of \(H_k\)
\end{enumerate}

For each integer \(k\ge 1\),

\[
H_k=\sum_{j=1}^{k}\frac1j
=\sum_{j=1}^{k}\int_{0}^{1} t^{\,j-1}\,dt
=\int_{0}^{1}\sum_{j=1}^{k} t^{\,j-1}\,dt
=\int_{0}^{1}\frac{1-t^{k}}{1-t}\,dt.
\tag{1}
\]

(We just summed the finite geometric series inside the integral.)

\begin{enumerate}
\def\labelenumi{\arabic{enumi})}
\setcounter{enumi}{1}
\tightlist
\item
  Work first with a truncated integral
\end{enumerate}

Fix \(r\in(0,1)\) and write

\[
S:=\sum_{k\ge 0} a_k x_0^{k} H_k
=\sum_{k\ge 0} a_k x_0^{k}\int_{0}^{1}\frac{1-t^{k}}{1-t}\,dt
=\lim_{r\uparrow 1}\sum_{k\ge 0} a_k x_0^{k}\int_{0}^{r}\frac{1-t^{k}}{1-t}\,dt.
\]

For each fixed \(r<1\), we may interchange sum and integral because
\(|x_0|r\) is strictly inside the radius of convergence, so
\(f(x_0t)=\sum_{k\ge 0} a_k(x_0t)^k\) converges uniformly on
\(t\in[0,r]\). Thus

\[
\sum_{k\ge 0} a_k x_0^{k}\int_{0}^{r}\frac{1-t^{k}}{1-t}\,dt
=\int_{0}^{r}\frac{1}{1-t}\sum_{k\ge 0}a_k\bigl(x_0^{k}- (x_0 t)^k\bigr)\,dt
=\int_{0}^{r}\frac{f(x_0)-f(x_0 t)}{1-t}\,dt.
\]

Hence

\[
S=\lim_{r\uparrow 1}\int_{0}^{r}\frac{f(x_0)-f(x_0 t)}{1-t}\,dt.
\tag{2}
\]

\begin{enumerate}
\def\labelenumi{\arabic{enumi})}
\setcounter{enumi}{2}
\tightlist
\item
  Pass to the limit \(r\uparrow1\)
\end{enumerate}

Define \(g(t)=\dfrac{f(x_0)-f(x_0 t)}{1-t}\) for \(t\in[0,1)\). Because
\(f\) is given by a power series that converges at \(x_0\), \(f\) is
differentiable on a neighborhood of the segment
\(\{x_0t:\,t\in[0,1]\}\). By the mean value theorem applied to
\(h(t)=f(x_0t)\),

\[
\lim_{t\uparrow 1} g(t)=\lim_{t\uparrow 1}\frac{h(1)-h(t)}{1-t}
= -h'(1)=x_0 f'(x_0).
\]

So \(g\) extends continuously to \(t=1\) (define \(g(1)=x_0f'(x_0)\)),
hence is bounded on \([0,1]\). Dominated convergence (or just continuity
of the integrand on the compact interval after extension) now yields

\[
\lim_{r\uparrow 1}\int_{0}^{r} g(t)\,dt=\int_{0}^{1} g(t)\,dt.
\]

Combining this with (2) finishes the proof:

\[
\boxed{\displaystyle
\sum_{k\ge 0} a_k\,x_0^{k}\,H_k
= \int_{0}^{1}\frac{\,f(x_0)-f(x_0y)\,}{1-y}\,dy }.
\]

\subsection{21 {[}M24{]}}\label{m24-8}

Compute

\[
S_n \;=\; \sum_{k=1}^{n} \frac{H_k}{\,n+1-k\,},
\qquad
H_k=\sum_{i=1}^{k}\frac1i,
\qquad
H_m^{(2)}=\sum_{i=1}^{m}\frac1{i^2}.
\]

Step 1: Rewrite as a double sum over a triangle

Expand \(H_k\) and swap the order of summation:

\[
S_n
=\sum_{k=1}^{n}\;\sum_{i=1}^{k}\frac{1}{i\,(n+1-k)}
=\sum_{i=1}^{n}\frac1i\sum_{k=i}^{n}\frac{1}{\,n+1-k\,}.
\]

With \(j=n+1-k\) the inner sum runs \(j=1,\dots,n+1-i\), hence

\[
S_n=\sum_{i=1}^{n}\frac1i\,H_{n+1-i}.
\]

Equivalently, view this as the ``triangular'' double sum

\[
S_n=\sum_{\substack{i\ge1,\;j\ge1\\ i+j\le n+1}}\frac{1}{ij}.
\]

Step 2: Collapse each antidiagonal

Group pairs by their sum \(r=i+j\) (so \(r=2,3,\dots,n+1\)):

\[
S_n=\sum_{r=2}^{n+1}\;\sum_{i=1}^{r-1}\frac{1}{i(r-i)}.
\]

Use the partial-fraction identity

\[
\frac{1}{i(r-i)}=\frac1r\!\left(\frac1i+\frac1{r-i}\right),
\]

to get

\[
\sum_{i=1}^{r-1}\frac{1}{i(r-i)}=\frac{2}{r}\sum_{i=1}^{r-1}\frac1i=\frac{2\,H_{r-1}}{r}.
\]

Therefore

\[
S_n=2\sum_{r=2}^{n+1}\frac{H_{r-1}}{r}
=2\sum_{k=1}^{n}\frac{H_{k}}{k+1}.
\]

Step 3: Evaluate \(\sum_{k=1}^{n}\dfrac{H_k}{k+1}\)

Write \(H_k=H_{k+1}-\dfrac{1}{k+1}\):

\[
\sum_{k=1}^{n}\frac{H_k}{k+1}
=\sum_{k=1}^{n}\frac{H_{k+1}}{k+1}-\sum_{k=1}^{n}\frac{1}{(k+1)^2}
=\sum_{j=2}^{n+1}\frac{H_{j}}{j}-\bigl(H_{n+1}^{(2)}-1\bigr).
\]

A standard identity (provable by the same double-sum trick) is

\[
\sum_{j=1}^{m}\frac{H_j}{j}
=\frac12\!\left(H_m^2+H_m^{(2)}\right).
\]

Using it with \(m=n+1\) and subtracting the \(j=1\) term,

\[
\sum_{j=2}^{n+1}\frac{H_j}{j}
=\frac12\!\left(H_{n+1}^2+H_{n+1}^{(2)}\right)-1.
\]

Hence

\[
\sum_{k=1}^{n}\frac{H_k}{k+1}
=\frac12\!\left(H_{n+1}^2+H_{n+1}^{(2)}\right)-1-\bigl(H_{n+1}^{(2)}-1\bigr)
=\frac12\!\left(H_{n+1}^2-H_{n+1}^{(2)}\right).
\]

Step 4: Final result

Plugging back into \(S_n=2\sum_{k=1}^{n}\dfrac{H_k}{k+1}\),

\[
\boxed{\;
\sum_{k=1}^{n}\frac{H_k}{\,n+1-k\,}
=H_{n+1}^{2}-H_{n+1}^{(2)}\; }.
\]

Quick check (small \(n\))

\begin{itemize}
\tightlist
\item
  \(n=1\): LHS \(=H_1/1=1\). RHS
  \(=H_2^2-H_2^{(2)}=(1+\tfrac12)^2-(1+ \tfrac14)=1\). ✓
\item
  \(n=2\): LHS \(=H_1/2+H_2/1=\tfrac12+ \tfrac32=2\). RHS
  \(=H_3^2-H_3^{(2)}=(\tfrac{11}{6})^2-(1+\tfrac14+\tfrac19)=2\). ✓
\end{itemize}

\subsection{22 {[}M28{]}}\label{m28-2}

\begin{enumerate}
\def\labelenumi{\arabic{enumi})}
\tightlist
\item
  Generating function setup
\end{enumerate}

The ordinary generating function of \(\{H_n\}_{n\ge 0}\) is

\[
F(x)=\sum_{n\ge 0} H_n x^n=\frac{-\ln(1-x)}{1-x},
\]

because

\[
\sum_{n\ge 1}H_n x^n
=\sum_{n\ge 1}\sum_{m=1}^n \frac{x^n}{m}
=\sum_{m\ge 1}\frac1m\sum_{n\ge m}x^n
=\frac{1}{1-x}\sum_{m\ge 1}\frac{x^m}{m}
=\frac{-\ln(1-x)}{1-x}.
\]

Therefore
\(S_n=[x^n]\,F(x)^2=[x^n]\left(\dfrac{\ln(1-x)}{1-x}\right)^2\).

\begin{enumerate}
\def\labelenumi{\arabic{enumi})}
\setcounter{enumi}{1}
\tightlist
\item
  Coefficient of \(\big(\ln(1-x)\big)^2\)
\end{enumerate}

Write \(-\ln(1-x)=\sum_{r\ge 1}\frac{x^r}{r}\). Then

\[
\big(\ln(1-x)\big)^2
= \sum_{m\ge 2}\left(\sum_{r=1}^{m-1}\frac{1}{r(m-r)}\right)x^m
= \sum_{m\ge 2}\frac{2H_{m-1}}{m}\,x^m,
\]

since
\(\dfrac{1}{r(m-r)}=\dfrac1m\!\left(\dfrac1r+\dfrac1{m-r}\right)\).

Also \(\dfrac{1}{(1-x)^2}=\sum_{t\ge 0}(t+1)x^t\). Hence

\[
S_n=[x^n]\frac{(\ln(1-x))^2}{(1-x)^2}
=\sum_{m=2}^{n}(n-m+1)\cdot\frac{2H_{m-1}}{m}.
\]

Let \(r=m-1\) (so \(r=1,\dots,n-1\)):

\[
S_n = 2\sum_{r=1}^{n-1}\frac{(n-r)H_r}{r+1}
= 2n\sum_{r=1}^{n-1}\frac{H_r}{r+1}\;-\;2\sum_{r=1}^{n-1}\frac{rH_r}{r+1}.
\]

\begin{enumerate}
\def\labelenumi{\arabic{enumi})}
\setcounter{enumi}{2}
\tightlist
\item
  Two standard harmonic sums
\end{enumerate}

We use

\[
\sum_{k=1}^{n-1}\frac{H_k}{k+1}
=\frac12\Big(H_n^2-H_n^{(2)}\Big),
\qquad
\sum_{k=1}^{n-1}H_k=n(H_n-1),
\]

and
\(\displaystyle \sum_{k=1}^{n-1}\frac{kH_k}{k+1} =\sum_{k=1}^{n-1}\left(H_k-\frac{H_k}{k+1}\right) = n(H_n-1)-\frac12\big(H_n^2-H_n^{(2)}\big).\)

Substitute into the expression for \(S_n\):

\[
\begin{aligned}
S_n
&= 2n\cdot \frac12\!\left(H_n^2-H_n^{(2)}\right)
\;-\;2\Big(n(H_n-1)-\tfrac12\big(H_n^2-H_n^{(2)}\big)\Big)\\[2mm]
&= (n+1)\!\left(H_n^2-H_n^{(2)}\right)-2nH_n+2n.
\end{aligned}
\]

It is slightly cleaner in terms of \(H_{n+1}\). Using
\(H_{n}=H_{n+1}-\dfrac1{n+1}\) and
\(H_{n}^{(2)}=H_{n+1}^{(2)}-\dfrac1{(n+1)^2}\), a brief simplification
yields the compact closed form:

\[
\boxed{\;
\sum_{k=0}^{n} H_k\,H_{n-k}
= (n+1)\Big( H_{n+1}^2 - H_{n+1}^{(2)} - 2H_{n+1} + 2 \Big).
\;}
\]

\begin{enumerate}
\def\labelenumi{\arabic{enumi})}
\setcounter{enumi}{3}
\tightlist
\item
  Quick checks
\end{enumerate}

\begin{itemize}
\tightlist
\item
  \(n=1\): \(H_0H_1+H_1H_0=0\). RHS:
  \((2)\big(H_2^2-H_2^{(2)}-2H_2+2\big)=2(2.25-1.25-3+2)=0\).
\item
  \(n=3\): LHS \(=H_1H_2+H_2H_1=3\). RHS with \(H_4=25/12\),
  \(H_4^{(2)}=205/144\) gives \(3\).
\end{itemize}

\subsection{23 {[}HM20{]}}\label{hm20-10}

Let

\[
\psi(x)\;\stackrel{\text{def}}{=}\;\frac{d}{dx}\ln\Gamma(x)
\;=\;\frac{\Gamma'(x)}{\Gamma(x)}
\]

be the digamma function.

Step 1 --- A key recurrence for \(\psi\).

From the Gamma functional equation \(\Gamma(x+1)=x\,\Gamma(x)\), taking
logs and differentiating gives

\[
\psi(x+1)=\psi(x)+\frac{1}{x}\qquad(x\notin\{0,-1,-2,\dots\}).
\]

Step 2 --- Recover \(H_n\) when \(x\) is an integer.

For a positive integer \(n\), iterating the recurrence:

\[
\psi(n+1)=\psi(1)+\sum_{k=1}^{n}\frac{1}{k}
= \psi(1)+H_n.
\]

Using the given fact \(\psi(1)=\Gamma'(1)/\Gamma(1)=-\gamma\), we obtain

\[
H_n=\psi(n+1)+\gamma.
\]

This identity is also recorded in the book's answers. 

Step 3 --- The definition for nonintegers.

Motivated by the integer case, define the generalized harmonic number
for any real (indeed complex) \(x\) not equal to a negative integer by

\[
\boxed{\;H_x\;:=\;\psi(x+1)+\gamma\;}.
\]

With this definition:

\begin{itemize}
\item
  Consistency. If \(x=n\in\mathbb{N}\), then \(H_x=H_n\) by Step 2.
\item
  Recurrence. \(H_x-H_{x-1}=\dfrac{1}{x}\) for all admissible \(x\),
  since

  \[
  H_x-H_{x-1}
  =\bigl[\psi(x+1)-\psi(x)\bigr]
  =\frac{1}{x}.
  \]

  (This exactly generalizes the discrete increment property of harmonic
  numbers.)
\item
  Asymptotics. Using the standard expansion
  \(\psi(x)=\ln x-\tfrac{1}{2x}+O(x^{-2})\), one gets

  \[
  H_x=\ln x+\gamma+\frac{1}{2x}+O\!\left(\frac{1}{x^2}\right),
  \]

  matching the familiar behavior of \(H_n\).
\end{itemize}

Example values

\begin{itemize}
\tightlist
\item
  \(H_{1/2}=\psi\!\left(\tfrac{3}{2}\right)+\gamma = \bigl(\psi(\tfrac12)+2\bigr)+\gamma = 2-2\ln 2\).
\item
  The recurrence gives
  \(H_{3/2}=H_{1/2}+\dfrac{2}{3}= \dfrac{8}{3}-2\ln 2\).
\end{itemize}

\subsection{24 {[}HM21{]}}\label{hm21-4}

Proof via partial products

Define the \(n\)-th partial product

\[
P_n(x)\;=\;x\,e^{\gamma x}\prod_{k=1}^{n}\Bigl(1+\frac{x}{k}\Bigr)e^{-x/k}.
\]

First, separate the exponential and finite product:

\[
P_n(x)=x\Bigl(\prod_{k=1}^{n}\Bigl(1+\frac{x}{k}\Bigr)\Bigr)
\exp\!\Bigl(x\gamma-\sum_{k=1}^n\frac{x}{k}\Bigr)
= x\Bigl(\prod_{k=1}^{n}\Bigl(1+\frac{x}{k}\Bigr)\Bigr)\,e^{x(\gamma-H_n)}.
\]

Use
\(\displaystyle \prod_{k=1}^{n}\Bigl(1+\frac{x}{k}\Bigr) =\prod_{k=1}^{n}\frac{x+k}{k}=\frac{(x+1)(x+2)\cdots(x+n)}{n!}\),
hence

\[
P_n(x)=\frac{x(x+1)\cdots(x+n)}{n!}\,e^{x(\gamma-H_n)}.
\]

Insert and remove \(n^x=e^{x\ln n}\):

\[
P_n(x)=\frac{x(x+1)\cdots(x+n)}{n!\,n^x}\;
\exp\!\Bigl(x\bigl(\gamma-(H_n-\ln n)\bigr)\Bigr).
\]

Now let \(n\to\infty\). By definition \(H_n-\ln n\to\gamma\), so the
exponential factor \(\exp\!\bigl(x(\gamma-(H_n-\ln n))\bigr)\to 1\).
Therefore

\[
\lim_{n\to\infty}P_n(x)
=\lim_{n\to\infty}\frac{x(x+1)\cdots(x+n)}{n!\,n^x}
=\frac{1}{\Gamma(x)}\quad(\text{Gauss}), 
\]

which proves the formula. 

Convergence (why the infinite product makes sense)

Consider the logarithm of the infinite product's factor:

\[
\log\Bigl(1+\frac{x}{k}\Bigr)-\frac{x}{k}
= -\frac{x^2}{2k^2}+O\!\Bigl(\frac{1}{k^3}\Bigr)\qquad(k\to\infty).
\]

Hence \(\sum_{k\ge1}\bigl(\log(1+x/k)-x/k\bigr)\) converges absolutely
(like \(\sum 1/k^2\)), so the product
\(\prod_{k\ge1}\bigl(1+\tfrac{x}{k}\bigr)e^{-x/k}\) converges (uniformly
on compact \(x\)-sets) and defines an entire function. The extra factor
\(x\,e^{\gamma x}\) fixes the normalization so that the result agrees
with \(\Gamma\) (as we just showed).

Sanity checks

\begin{itemize}
\tightlist
\item
  Zeros/poles. The product clearly vanishes at \(x=0,-1,-2,\ldots\),
  matching the poles of \(\Gamma(x)\) at those points, so \(1/\Gamma\)
  has the same zeros (simple).
\item
  Near \(x=0\). Since \(\Gamma(x)\sim 1/x-\gamma+O(x)\), we have
  \(1/\Gamma(x)=x+\gamma x^2+O(x^3)\). On the product side,
  \(\prod(1+x/k)e^{-x/k}=1+O(x^2)\) and
  \(x e^{\gamma x}=x+\gamma x^2+O(x^3)\), in agreement.
\end{itemize}

\subsection{25 {[}M21{]}}\label{m21-11}

\begin{enumerate}
\def\labelenumi{(\alph{enumi})}
\tightlist
\item
  Evaluate \(H_n^{(0,v)}\) and \(H_n^{(u,0)}\)
\end{enumerate}

\(H_n^{(0,v)}\)

\[
H_n^{(0,v)}
=\sum_{1\le j\le k\le n}\frac{1}{k^v}
=\sum_{k=1}^n\frac{1}{k^v}\sum_{j=1}^k 1
=\sum_{k=1}^n\frac{k}{k^v}
=\sum_{k=1}^n k^{1-v}
=\sum_{k=1}^n \frac{1}{k^{\,v-1}}
=H_n^{(v-1)}.
\]

\(H_n^{(u,0)}\)

\[
H_n^{(u,0)}
=\sum_{1\le j\le k\le n}\frac{1}{j^u}
=\sum_{j=1}^n\frac{1}{j^u}\sum_{k=j}^n 1
=\sum_{j=1}^n\frac{n-j+1}{j^u}
=(n+1)\sum_{j=1}^n\frac{1}{j^u}-\sum_{j=1}^n\frac{j}{j^u}
=(n+1)H_n^{(u)}-H_n^{(u-1)}.
\]

Thus

\[
\boxed{\,H_n^{(0,v)}=H_n^{(v-1)},\qquad H_n^{(u,0)}=(n+1)H_n^{(u)}-H_n^{(u-1)}\, }.
\]

These identities hold for any real \(u,v\) since the sums are finite.

\begin{enumerate}
\def\labelenumi{(\alph{enumi})}
\setcounter{enumi}{1}
\tightlist
\item
  Prove \(H_n^{(u,v)}+H_n^{(v,u)} = H_n^{(u)}H_n^{(v)}+H_n^{(u+v)}\)
\end{enumerate}

Start from the product of univariate sums:

\[
H_n^{(u)}H_n^{(v)}
=\Bigl(\sum_{j=1}^n j^{-u}\Bigr)\Bigl(\sum_{k=1}^n k^{-v}\Bigr)
=\sum_{j=1}^n\sum_{k=1}^n \frac{1}{j^u k^v}.
\]

Partition the index set \(\{1,\dots,n\}^2\) into three regions: \(j<k\),
\(j=k\), \(j>k\). Denote \(D=\sum_{j=1}^n j^{-(u+v)}=H_n^{(u+v)}\).

\begin{itemize}
\tightlist
\item
  Over \(j<k\):
  \(\displaystyle \sum_{1\le j<k\le n}\frac{1}{j^u k^v}=H_n^{(u,v)}-D\)
  (because \(H_n^{(u,v)}\) includes the diagonal \(j=k\) contributing
  \(D\)).
\item
  Over \(j>k\): by swapping \(j\leftrightarrow k\),
  \(\displaystyle \sum_{1\le k<j\le n}\frac{1}{j^u k^v}  =\sum_{1\le j<k\le n}\frac{1}{j^v k^u}  =H_n^{(v,u)}-D\).
\item
  Over \(j=k\): contribution is \(D\).
\end{itemize}

Therefore

\[
H_n^{(u)}H_n^{(v)}
=\bigl(H_n^{(u,v)}-D\bigr)+\bigl(H_n^{(v,u)}-D\bigr)+D
=H_n^{(u,v)}+H_n^{(v,u)}-D.
\]

Rearranging yields

\[
\boxed{\,H_n^{(u,v)}+H_n^{(v,u)}=H_n^{(u)}H_n^{(v)}+H_n^{(u+v)}\, }.
\]

This is a simple double-counting (index-partition) argument.

Quick sanity check (small \(n\))

Take \(n=2\), \(u=1\), \(v=2\).

\begin{itemize}
\tightlist
\item
  \(H_2^{(1)}=1+1/2=3/2,\quad H_2^{(2)}=1+1/4=5/4,\quad H_2^{(3)}=1+1/8=9/8.\)
\item
  \(H_2^{(1,2)}=\sum_{1\le j\le k\le2}1/(j^1 k^2)=1+1/2^2+1/(1\cdot 2^2)=1+1/4+1/4=3/2.\)
\item
  \(H_2^{(2,1)}=\sum_{1\le j\le k\le2}1/(j^2 k^1)=1+1/2+1/(1^2\cdot2)=1+1/2+1/2=2.\)
\end{itemize}

Left side: \(H_2^{(1,2)}+H_2^{(2,1)}=3/2+2=7/2.\)

Right side:
\(H_2^{(1)}H_2^{(2)}+H_2^{(3)}=(3/2)(5/4)+9/8=15/8+9/8=24/8=3.\)
Wait---that's not equal; we made a slip in \(H_2^{(1,2)}\).

Correct \(H_2^{(1,2)}\):

\[
(j,k)=(1,1):1,\ (1,2):1/(1\cdot 2^2)=1/4,\ (2,2):1/(2\cdot 2^2)=1/8.
\]

So \(H_2^{(1,2)}=1+1/4+1/8=11/8.\)

Similarly \(H_2^{(2,1)}\):

\[
(1,1):1,\ (1,2):1/(1^2\cdot2)=1/2,\ (2,2):1/(2^2\cdot2)=1/8,
\]

so \(H_2^{(2,1)}=1+1/2+1/8=13/8.\)

Now LHS \(=11/8+13/8=24/8=3\), matching RHS \(=3\). Check.

\bookmarksetup{startatroot}

\chapter{1.2.8 Fibonacci numbers}\label{fibonacci-numbers}

\subsection{1 {[}10{]}}\label{section-87}

Model the process precisely

\begin{itemize}
\tightlist
\item
  Let \(R_n\) be the number of rabbit pairs at the end of month \(n\),
  starting from one newborn pair at month \(0\).
\item
  Each mature pair produces one new pair every month; newborn pairs
  become mature after one month; rabbits never die. In Fibonacci's own
  description, after one month there are 2 pairs, after two months 3
  pairs, after three months 5, and so on. 
\end{itemize}

Split pairs into mature \(M_n\) and immature \(I_n\). The month-to-month
transitions are:

\[
\begin{aligned}
M_{n+1} &= M_n + I_n &&\text{(last month’s immature pairs mature)}\\
I_{n+1} &= M_n        &&\text{(each mature pair produces one new pair)}.
\end{aligned}
\]

Hence the totals \(R_n = M_n + I_n\) satisfy

\[
R_{n+1} \;=\; (M_n+I_n) + M_n \;=\; R_n + M_n \;=\; R_n + R_{n-1},
\]

because \(M_n = I_{n+1} = R_{n-1}\). Initial conditions are \(R_0=1\)
(one newborn pair) and \(R_1=2\) (the original pair, now fertile,
produces a new pair by the end of month 1). This is exactly the
Fibonacci recurrence with a shift. 

Identify the Fibonacci index

Knuth defines the Fibonacci numbers by
\(F_0=0,\; F_1=1,\; F_{n+2}=F_{n+1}+F_n\). With the initial values above
it follows (by induction) that

\[
R_n \;=\; F_{n+2}.
\]

(Indeed, \(R_0=1=F_2,\;R_1=2=F_3\), and if \(R_n=F_{n+2}\) and
\(R_{n-1}=F_{n+1}\), then \(R_{n+1}=R_n+R_{n-1}=F_{n+3}\).) 

Compute the 12-month total

After a year (\(n=12\) months),

\[
R_{12} \;=\; F_{14}.
\]

The Fibonacci numbers from \(F_0\) upward are
\(0,1,1,2,3,5,8,13,21,34,55,89,144,233,377,\dots\), so \(F_{14}=377\).

Therefore, there are \(377\) pairs after one year. Knuth's answers
section states this succinctly: ``After \(k\) months there are
\(F_{k+2}\) pairs, so the answer is \(F_{14}=377\) pairs.''

\emph{(If you ever see 233 in other tellings, that's from a different
month/index convention. Under the book's convention---2 pairs after the
first month---the correct 12-month answer is 377.)}

\subsection{2 {[}20{]}}\label{section-88}

We use Binet's formula and logarithms.

\begin{enumerate}
\def\labelenumi{\arabic{enumi})}
\tightlist
\item
  Background. Let \(\varphi=\dfrac{1+\sqrt5}{2}\) and
  \(\psi=\dfrac{1-\sqrt5}{2}=-\varphi^{-1}\). Binet's formula says
\end{enumerate}

\[
F_n=\frac{\varphi^n-\psi^n}{\sqrt5}.
\]

Since \(|\psi|<1\), the term \(\psi^n\) is tiny for large \(n\). In
fact,

\[
0<\frac{\varphi^n}{\sqrt5}-F_n=\frac{\psi^n}{\sqrt5}<\frac{1}{\sqrt5}\,\Bigl(\frac{1}{\varphi}\Bigr)^n,
\]

and the relative error is
\(\bigl|\psi/\varphi\bigr|^n=(\varphi^{-2})^n\), which for \(n=1000\) is
\(\approx 10^{-418}\). Thus \(\varphi^n/\sqrt5\) determines the front
digits of \(F_n\) unambiguously.

\begin{enumerate}
\def\labelenumi{\arabic{enumi})}
\setcounter{enumi}{1}
\tightlist
\item
  Number of digits. For any positive integer \(N\), the number of
  decimal digits is
\end{enumerate}

\[
\#\text{digits}(N)=\bigl\lfloor \log_{10} N \bigr\rfloor+1.
\]

Hence

\[
\#\text{digits}\bigl(F_n\bigr)=\Bigl\lfloor \log_{10}\!\Bigl(\tfrac{\varphi^n-\psi^n}{\sqrt5}\Bigr)\Bigr\rfloor+1
\;=\;\Bigl\lfloor n\log_{10}\varphi-\log_{10}\sqrt5+\delta_n\Bigr\rfloor+1,
\]

where \(\delta_n=\log_{10}\!\bigl(1-\psi^n/\varphi^n\bigr)\) and
\(\delta_{1000}\) is about \(-10^{-418}\), i.e., negligible for the
integer part.

Compute the main term (to, say, 15--18 significant digits for safety):

\[
\log_{10}\varphi \approx 0.20898764024997873,\qquad
\log_{10}\sqrt5 \approx 0.3494850021680094.
\]

Then

\[
L := 1000\log_{10}\varphi-\log_{10}\sqrt5
\approx 208.6381552478107.
\]

Therefore

\[
\#\text{digits}(F_{1000})=\lfloor L\rfloor+1=208+1=209.
\]

\begin{enumerate}
\def\labelenumi{\arabic{enumi})}
\setcounter{enumi}{2}
\tightlist
\item
  Leading digit. Write
  \(F_{1000}=10^{\lfloor L\rfloor}\cdot 10^{\{L\}}\cdot(1+\varepsilon)\),
  where \(\{L\}=L-\lfloor L\rfloor\) is the fractional part and
  \(\varepsilon\) is the tiny relative correction from \(\psi^{1000}\).
  The significand is essentially \(S=10^{\{L\}}\).
\end{enumerate}

From the value above,

\[
\{L\}\approx 0.6381552478107
\quad\Longrightarrow\quad
S=10^{\{L\}}\approx 4.346655768\ldots
\]

Hence the leading decimal digit is \(\boxed{4}\). (Because the
correction \(|\varepsilon|\lesssim 10^{-418}\) cannot possibly change
the first digit.)

Result.

\[
\boxed{\#\text{digits}(F_{1000})=209,\quad \text{first digit}=4.}
\]

General form (useful to remember). For any \(n\ge 1\),

\[
\#\text{digits}(F_n)\;=\;\Bigl\lfloor n\log_{10}\varphi-\log_{10}\sqrt5\Bigr\rfloor+1,
\]

and if \(x=n\log_{10}\varphi-\log_{10}\sqrt5\) with fractional part
\(\{x\}\), then the leading digit of \(F_n\) is
\(\bigl\lfloor 10^{\{x\}}\bigr\rfloor\) (again, because the \(\psi^n\)
term is far too small to affect the front digit for large \(n\)).

\subsection{3 {[}25{]}}\label{section-89}

What size are these numbers?

Knuth gives the approximation \(F_n \approx \varphi^n/\sqrt5\), so the
number of decimal digits of \(F_n\) is

\[
\Big\lfloor n\log_{10}\varphi - \log_{10}\sqrt5\Big\rfloor + 1.
\]

For \(n=1000\) this yields \(209\) digits, which matches the computed
result. 

Algorithm (addition-only)

We compute iteratively from the recurrence \(F_{n+2}=F_{n+1}+F_n\) with
seeds \(F_0=0\), \(F_1=1\). To honor the ``addition exercise'' spirit, I
implemented big integers by hand:

\begin{itemize}
\tightlist
\item
  Represent each big integer as base \(10^9\) ``chunks'' in
  little-endian order.
\item
  Addition is schoolbook addition over those chunks with a single carry
  bit.
\item
  Update the two registers \((a,b)=(F_{n-1},F_n)\) by
  \((a,b)\leftarrow (b,\,a+b)\) for \(n=1\ldots1000\).
\end{itemize}

Correctness (brief)

Induction on \(n\). Base: \((a,b)=(0,1)=(F_0,F_1)\). Step: assuming
\((a,b)=(F_{n-1},F_n)\), the update makes \(b'=a+b=F_{n-1}+F_n=F_{n+1}\)
and \(a'=b=F_n\), so \((a',b')=(F_n,F_{n+1})\). Thus after \(n\) steps
we have \(F_n\).

\begin{itemize}
\tightlist
\item
  \(F(1000)\) has 209 digits.
\item
  First 50 digits of \(F(1000)\):
  \texttt{43466557686937456435688527675040625802564660517371}
\item
  Last 50 digits of \(F(1000)\):
  \texttt{33187938298969649928516003704476137795166849228875}
\end{itemize}

\begin{Shaded}
\begin{Highlighting}[]
\CommentTok{\# Big{-}int Fibonacci using base 10\^{}9 chunks (addition{-}only)}
\NormalTok{BASE }\OperatorTok{=} \DecValTok{109}

\KeywordTok{def}\NormalTok{ add\_big(a, b):}
\NormalTok{    out, carry }\OperatorTok{=}\NormalTok{ [], }\DecValTok{0}
\NormalTok{    n }\OperatorTok{=} \BuiltInTok{max}\NormalTok{(}\BuiltInTok{len}\NormalTok{(a), }\BuiltInTok{len}\NormalTok{(b))}
    \ControlFlowTok{for}\NormalTok{ i }\KeywordTok{in} \BuiltInTok{range}\NormalTok{(n):}
\NormalTok{        s }\OperatorTok{=}\NormalTok{ (a[i] }\ControlFlowTok{if}\NormalTok{ i }\OperatorTok{\textless{}} \BuiltInTok{len}\NormalTok{(a) }\ControlFlowTok{else} \DecValTok{0}\NormalTok{) }\OperatorTok{+}\NormalTok{ (b[i] }\ControlFlowTok{if}\NormalTok{ i }\OperatorTok{\textless{}} \BuiltInTok{len}\NormalTok{(b) }\ControlFlowTok{else} \DecValTok{0}\NormalTok{) }\OperatorTok{+}\NormalTok{ carry}
        \ControlFlowTok{if}\NormalTok{ s }\OperatorTok{\textgreater{}=}\NormalTok{ BASE:}
\NormalTok{            s }\OperatorTok{{-}=}\NormalTok{ BASE}\OperatorTok{;}\NormalTok{ carry }\OperatorTok{=} \DecValTok{1}
        \ControlFlowTok{else}\NormalTok{:}
\NormalTok{            carry }\OperatorTok{=} \DecValTok{0}
\NormalTok{        out.append(s)}
    \ControlFlowTok{if}\NormalTok{ carry: out.append(carry)}
    \ControlFlowTok{return}\NormalTok{ out}

\KeywordTok{def}\NormalTok{ to\_string(a):}
    \ControlFlowTok{if} \KeywordTok{not}\NormalTok{ a: }\ControlFlowTok{return} \StringTok{"0"}
\NormalTok{    s }\OperatorTok{=} \BuiltInTok{str}\NormalTok{(a[}\OperatorTok{{-}}\DecValTok{1}\NormalTok{])}
    \ControlFlowTok{for}\NormalTok{ x }\KeywordTok{in} \BuiltInTok{reversed}\NormalTok{(a[:}\OperatorTok{{-}}\DecValTok{1}\NormalTok{]): s }\OperatorTok{+=} \SpecialStringTok{f"}\SpecialCharTok{\{}\NormalTok{x}\SpecialCharTok{:09d\}}\SpecialStringTok{"}
    \ControlFlowTok{return}\NormalTok{ s}

\CommentTok{\# Build F0..F1000, then write F1..F1000 to files}
\NormalTok{F }\OperatorTok{=}\NormalTok{ [[}\DecValTok{0}\NormalTok{], [}\DecValTok{1}\NormalTok{]]}
\ControlFlowTok{for}\NormalTok{ \_ }\KeywordTok{in} \BuiltInTok{range}\NormalTok{(}\DecValTok{2}\NormalTok{, }\DecValTok{1000} \OperatorTok{+} \DecValTok{1}\NormalTok{):}
\NormalTok{    F.append(add\_big(F[}\OperatorTok{{-}}\DecValTok{1}\NormalTok{], F[}\OperatorTok{{-}}\DecValTok{2}\NormalTok{]))}
\ControlFlowTok{with} \BuiltInTok{open}\NormalTok{(}\StringTok{"fibonacci\_1\_1000.txt"}\NormalTok{,}\StringTok{"w"}\NormalTok{) }\ImportTok{as}\NormalTok{ f\_txt, }\OperatorTok{\textbackslash{}}
     \BuiltInTok{open}\NormalTok{(}\StringTok{"fibonacci\_1\_1000.csv"}\NormalTok{,}\StringTok{"w"}\NormalTok{) }\ImportTok{as}\NormalTok{ f\_csv:}
\NormalTok{    f\_csv.write(}\StringTok{"n,F,Digits}\CharTok{\textbackslash{}n}\StringTok{"}\NormalTok{)}
    \ControlFlowTok{for}\NormalTok{ n }\KeywordTok{in} \BuiltInTok{range}\NormalTok{(}\DecValTok{1}\NormalTok{, }\DecValTok{1001}\NormalTok{):}
\NormalTok{        s }\OperatorTok{=}\NormalTok{ to\_string(F[n])}
\NormalTok{        f\_txt.write(}\SpecialStringTok{f"F(}\SpecialCharTok{\{}\NormalTok{n}\SpecialCharTok{\}}\SpecialStringTok{) = }\SpecialCharTok{\{}\NormalTok{s}\SpecialCharTok{\}}\CharTok{\textbackslash{}n}\SpecialStringTok{"}\NormalTok{)}
\NormalTok{        f\_csv.write(}\SpecialStringTok{f"}\SpecialCharTok{\{}\NormalTok{n}\SpecialCharTok{\}}\SpecialStringTok{,}\SpecialCharTok{\{}\NormalTok{s}\SpecialCharTok{\}}\SpecialStringTok{,}\SpecialCharTok{\{}\BuiltInTok{len}\NormalTok{(s)}\SpecialCharTok{\}}\CharTok{\textbackslash{}n}\SpecialStringTok{"}\NormalTok{)}
\end{Highlighting}
\end{Shaded}

\subsection{4 {[}14{]}}\label{section-90}

Quick check of small \(n\)

The Fibonacci numbers begin

\[
F_0=0,\;F_1=1,\;F_2=1,\;F_3=2,\;F_4=3,\;F_5=5,\;F_6=8,\;F_7=13,\ldots
\]

We immediately see \(F_0=0\), \(F_1=1\), and \(F_5=5\). For \(n=2,3,4\)
we have \(F_n<n\); for \(n=6\) we have \(F_6=8>6\).

This suggests the only solutions are \(n\in\{0,1,5\}\).

Proof that there are no others (elementary induction)

We show \(F_n>n\) for all \(n\ge 6\).

Base cases. \(F_6=8>6\) and \(F_7=13>7\).

Inductive step. Assume \(F_k>k\) and \(F_{k-1}>k-1\) for some
\(k\ge 7\). Then

\[
F_{k+1}=F_k+F_{k-1} > k+(k-1)=2k-1 \ge k+1.
\]

Hence \(F_{k+1}>k+1\). By induction, \(F_n>n\) for all \(n\ge 6\).

Combining this with the quick check up to \(n=5\), the only indices with
\(F_n=n\) are

\[
\boxed{n=0,1,5}.
\]

(Optional) Growth-based one-liner

Using the closed form \(F_n=\frac{\varphi^n-\hat\varphi^n}{\sqrt5}\)
with \(\varphi=(1+\sqrt5)/2\) and \(|\hat\varphi|<1\), we have \(F_n\)
grows like \(\varphi^n/\sqrt5\), so it eventually dominates the linear
function \(n\). This also rules out further solutions beyond small
\(n\).

\subsection{5 {[}20{]}}\label{section-91}

Define \(g(n)=F_n-n^2\). We'll show \(g(n)\) becomes strictly increasing
from \(n=9\) onward; since \(g(12)=0\), that will force \(n=12\) to be
the only solution for large \(n\). Small \(n\) can be checked directly.

Step 1: A monotonicity lemma for \(g(n)\)

Using \(F_{n+1}=F_n+F_{n-1}\),

\[
g(n+1)-g(n)
= \bigl(F_{n+1}-(n+1)^2\bigr)-\bigl(F_n-n^2\bigr)
= F_{n-1}-(2n+1).
\]

Thus \(g\) is increasing whenever \(F_{n-1}>(2n+1)\).

At \(n=9\) we have \(F_8=21\) and \(2\cdot9+1=19\), so \(g\) begins to
increase: \(g(10)-g(9)=2>0\). Moreover, for \(n\ge10\),

\[
\bigl(F_{n}- (2(n+1)+1)\bigr) - \bigl(F_{n-1}-(2n+1)\bigr)
= (F_n-F_{n-1})-2 = F_{n-2}-2 \ge F_8-2=19>0,
\]

so once positive, the gap strictly widens thereafter. Hence \(g(n)\) is
strictly increasing for all \(n\ge9\).

Step 2: Pin down the zero(s)

Compute a few values:

\[
\begin{array}{c|cccc}
n&9&10&11&12\\\hline
F_n&34&55&89&144\\
n^2&81&100&121&144\\
g(n)&-47&-45&-32&0
\end{array}
\]

Thus \(g(9)<0\), \(g(12)=0\), and by Step 1 \(g(n)>0\) for every
\(n\ge13\). Therefore there is no solution with \(n\ge13\), and \(n=12\)
is the unique solution in \(\{9,10,11,12,\ldots\}\).

Step 3: Check the small cases

Directly:

\[
F_0=0=0^2,\quad F_1=1=1^2,\quad F_2=1,\ F_3=2,\ F_4=3,\ F_5=5,\ F_6=8,\ F_7=13,\ F_8=21,
\]

none of which equals \(n^2\) except \(n=0,1\).

\subsection{6 {[}HM10{]}}\label{hm10-1}

Let \(A=\begin{pmatrix}1&1\\[2pt]1&0\end{pmatrix}\). We will show by
induction on \(n\ge 1\) that

\[
A^{n}=\begin{pmatrix}F_{n+1}&F_{n}\\[2pt]F_{n}&F_{n-1}\end{pmatrix}.
\]

Base case \(n=1\).
\(A^1=A=\begin{pmatrix}1&1\\[1pt]1&0\end{pmatrix} =\begin{pmatrix}F_{2}&F_{1}\\[1pt]F_{1}&F_{0}\end{pmatrix}\)
since \(F_0=0,\;F_1=1,\;F_2=1\). ✓

Inductive step. Assume

\[
A^{n}=\begin{pmatrix}F_{n+1}&F_{n}\\[1pt]F_{n}&F_{n-1}\end{pmatrix}.
\]

Then

\[
A^{n+1}=A^{n}A
=\begin{pmatrix}F_{n+1}&F_{n}\\[1pt]F_{n}&F_{n-1}\end{pmatrix}
\begin{pmatrix}1&1\\[1pt]1&0\end{pmatrix}
=\begin{pmatrix}
F_{n+1}+F_{n} & F_{n+1}\\[2pt]
F_{n}+F_{n-1} & F_{n}
\end{pmatrix}.
\]

Using the Fibonacci recurrence \(F_{k+1}=F_k+F_{k-1}\), we get

\[
A^{n+1}=\begin{pmatrix}F_{n+2}&F_{n+1}\\[2pt]F_{n+1}&F_{n}\end{pmatrix},
\]

which is exactly the desired form with \(n\) replaced by \(n+1\).

By induction, the identity holds for all \(n\ge 1\), proving Eq. (5).

\subsection{7 {[}15{]}}\label{section-92}

We'll use two standard facts from the section.

\begin{enumerate}
\def\labelenumi{\arabic{enumi}.}
\tightlist
\item
  Fibonacci addition formula. For all positive integers \(m,n\),
\end{enumerate}

\[
F_{m+n}=F_mF_{n+1}+F_{m-1}F_n .
\]

This is Eq. (6) in the text. 

\begin{enumerate}
\def\labelenumi{\arabic{enumi}.}
\setcounter{enumi}{1}
\tightlist
\item
  (Optional strengthening, due to Lucas):
  \(\gcd(F_m,F_n)=F_{\gcd(m,n)}\). 
\end{enumerate}

Lemma. If \(d\mid n\), then \(F_d\mid F_n\).

Proof (by induction on \(k\) where \(n=kd\)). For \(k=1\) the claim is
trivial. Assume \(F_d\mid F_{kd}\). Using Eq. (6) with \(m=kd\) and
\(n=d\),

\[
F_{(k+1)d}=F_{kd}F_{d+1}+F_{kd-1}F_d .
\]

By the induction hypothesis \(F_{kd}\) is a multiple of \(F_d\); the
second term is clearly a multiple of \(F_d\). Hence
\(F_d\mid F_{(k+1)d}\). QED. 

\emph{(Equivalently, the lemma is immediate from Lucas's theorem, since
if \(d\mid n\) then \(\gcd(F_n,F_d)=F_d\), so \(F_d\) divides
\(F_n\).)} 

Suppose \(n\) is composite. Then \(n\) has a proper divisor \(d>1\).

\begin{itemize}
\item
  If \(d\ge 3\), then \(F_d\ge F_3=2\) and (by monotonicity of \(F_n\)
  for \(n\ge2\)) \(F_d<F_n\). By the lemma \(F_d\mid F_n\), so \(F_n\)
  has a nontrivial divisor and is therefore composite.
\item
  The only composite integer with no proper divisor \(>2\) is \(n=4\).
  Here the only proper divisor is \(d=2\) and \(F_2=1\) (which does not
  yield a nontrivial divisor). Direct computation gives \(F_4=3\), which
  is prime.
\end{itemize}

Thus, if \(n\) is not prime, then \(F_n\) is not prime, with the unique
exception \(n=4\).

\subsection{8 {[}15{]}}\label{section-93}

Recall the (positive-index) Fibonacci numbers:

\begin{itemize}
\tightlist
\item
  \(F_0=0,\; F_1=1\),
\item
  \(F_{n}=F_{n-1}+F_{n-2}\) for \(n\ge 2\).
\end{itemize}

A natural way to extend any linear recurrence to all integers is to run
the recurrence backwards. From

\[
F_n=F_{n-1}+F_{n-2}\quad\Longrightarrow\quad F_{n-2}=F_n-F_{n-1},
\]

so once \(F_0,F_1\) are fixed, every \(F_{k}\) with \(k<0\) is
determined uniquely.

Step 1: Determine \(F_{-1}\) and \(F_{-2}\)

Use the recurrence with \(n=1\) and \(n=0\):

\begin{itemize}
\tightlist
\item
  \(F_1=F_0+F_{-1}\Rightarrow 1=0+F_{-1}\Rightarrow F_{-1}=1\).
\item
  \(F_0=F_{-1}+F_{-2}\Rightarrow 0=1+F_{-2}\Rightarrow F_{-2}=-1\).
\end{itemize}

These two values will seed an induction.

Step 2: Prove the general formula

Claim. For all \(n\ge 1\),

\[
\boxed{\,F_{-n}=(-1)^{n+1}F_n\,}.
\]

Proof by induction on \(n\).

Base cases: \(n=1\): \(F_{-1}=1=(-1)^{2}F_1\). \(n=2\):
\(F_{-2}=-1=(-1)^3F_2\) since \(F_2=1\).

Inductive step: Assume the claim for \(n\) and \(n-1\). Using the
backward recurrence \(F_{m-1}=F_{m+1}-F_m\) with \(m=-n\),

\[
F_{-(n+1)}=F_{-n+1}-F_{-n}.
\]

Apply the induction hypotheses:

\[
F_{-(n+1)} = (-1)^{n}F_{n-1}-(-1)^{n+1}F_n
= (-1)^{n}\bigl(F_{n-1}+F_n\bigr)
= (-1)^{n}F_{n+1}
= (-1)^{(n+1)+1}F_{n+1}.
\]

Thus the formula holds for \(n+1\). By induction, it holds for all
\(n\ge1\). ∎

Step 3: Sanity check examples

The extension alternates signs:

\[
\ldots,\; F_{-5}=5,\; F_{-4}=-3,\; F_{-3}=2,\; F_{-2}=-1,\; F_{-1}=1,\; F_0=0,\; F_1=1,\; F_2=1,\; F_3=2,\; F_4=3,\; F_5=5,\ldots
\]

which matches \(F_{-n}=(-1)^{n+1}F_n\).

Optional alternate derivation (Binet)

From Binet's formula
\(F_n=\frac{\varphi^{\,n}-\hat\varphi^{\,n}}{\sqrt5}\) with
\(\hat\varphi=1-\varphi=-1/\varphi\),

\[
F_{-n}
=\frac{\varphi^{-n}-\hat\varphi^{-n}}{\sqrt5}
=\frac{\varphi^{-n}-(-1)^n\varphi^{n}}{\sqrt5}
=(-1)^{n+1}\frac{\varphi^{n}-\hat\varphi^{\,n}}{\sqrt5}
=(-1)^{n+1}F_n.
\]

\subsection{9 {[}M20{]}}\label{m20-20}

Define \(F_0=0,\;F_1=1\) and keep \(F_{n+2}=F_{n+1}+F_n\) for all
\(n\in\mathbb Z\) (this is the convention from Exercise 8). Running the
recurrence backwards gives

\[
F_{-1}=1,\quad F_{-2}=-1,\quad F_{-3}=2,\ \ldots
\]

and by induction (or a quick check with the recurrence) one gets the
standard identity

\[
\boxed{\,F_{-n}=(-1)^{n+1}F_n\quad(n\ge1).\,}
\]

This is the only new fact we'll need from the extension. (See the
exercise list introducing negative indices.)

\begin{enumerate}
\def\labelenumi{(\arabic{enumi})}
\setcounter{enumi}{3}
\tightlist
\item
  Cassini's identity for all integers
\end{enumerate}

Cassini's formula in the text is

\[
F_{n+1}F_{n-1}-F_n^2=(-1)^n. \tag{4}
\]

Using Binet's form \(F_n=\frac{\phi^n-\hat\phi^{\,n}}{\sqrt5}\) with
\(\phi=\frac{1+\sqrt5}2,\;\hat\phi=\frac{1-\sqrt5}2\) (note
\(\phi+\hat\phi=1,\ \phi\hat\phi=-1\)), valid as an identity of
exponentials for any integer \(n\), compute

\[
\begin{aligned}
F_{n+1}F_{n-1}-F_n^2
&=\frac1{5}\Big[(\phi^{n+1}-\hat\phi^{\,n+1})(\phi^{n-1}-\hat\phi^{\,n-1})
-(\phi^n-\hat\phi^{\,n})^2\Big]\\
&=\frac1{5}\,\phi^{n-1}\hat\phi^{\,n-1}\,\big(-\phi^2-\hat\phi^{\,2}+2\phi\hat\phi\big)\\
&=-(\phi\hat\phi)^{n-1}=\;-\;(-1)^{\,n-1}=( -1)^n.
\end{aligned}
\]

No positivity of \(n\) was used; thus (4) holds for all
\(n\in\mathbb Z\).

\begin{enumerate}
\def\labelenumi{(\arabic{enumi})}
\setcounter{enumi}{5}
\tightlist
\item
  The addition formula for all integers
\end{enumerate}

The text gives

\[
F_{n+m}=F_m\,F_{n+1}+F_{m-1}\,F_n\qquad(m>0\ \text{in the derivation}). \tag{6}
\]

With negative indices allowed and Binet's formula, expand the right-hand
side using \(\phi\hat\phi=-1\) and \(\phi+\hat\phi=1\):

\[
\begin{aligned}
F_mF_{n+1}+F_{m-1}F_n
&=\frac{(\phi^m-\hat\phi^{\,m})(\phi^{n+1}-\hat\phi^{\,n+1})}{5}
  +\frac{(\phi^{m-1}-\hat\phi^{\,m-1})(\phi^{n}-\hat\phi^{\,n})}{5}\\
&=\frac{\phi^{m+n+1}-\hat\phi^{\,m+n+1}}{\sqrt5}\cdot\frac1{\sqrt5}
+\text{(cross terms that cancel)}\\
&=\frac{\phi^{n+m}-\hat\phi^{\,n+m}}{\sqrt5}
=F_{n+m},
\end{aligned}
\]

and this algebra is valid for all integers \(m,n\) (negative exponents
are fine). Therefore (6) holds for all \(m,n\in\mathbb Z\).
(Equivalently, one can argue via the matrix identity (5) and the fact
that \(\begin{bmatrix}1&1\\[2pt]1&0\end{bmatrix}\) is invertible.)

\begin{enumerate}
\def\labelenumi{(\arabic{enumi})}
\setcounter{enumi}{13}
\tightlist
\item
  Binet's formula for all integers
\end{enumerate}

Binet's formula is

\[
F_n=\frac{\phi^n-\hat\phi^{\,n}}{\sqrt5}. \tag{14}
\]

For negative \(n=-k\) with \(k>0\), use
\(\phi^{-k}=(-1)^k\hat\phi^{\,k}\) and \(\hat\phi^{-k}=(-1)^k\phi^{k}\)
(from \(\phi\hat\phi=-1\)):

\[
F_{-k}
=\frac{(-1)^k\hat\phi^{\,k}-(-1)^k\phi^{k}}{\sqrt5}
=(-1)^{k+1}\frac{\phi^k-\hat\phi^{\,k}}{\sqrt5}
=(-1)^{k+1}F_k,
\]

which matches the extended definition above. Hence (14) remains true for
every integer \(n\).

\begin{enumerate}
\def\labelenumi{(\arabic{enumi})}
\setcounter{enumi}{14}
\tightlist
\item
  The nearest-integer rule
\end{enumerate}

The corollary stated is

\[
F_n=\text{the nearest integer to }\;\frac{\phi^n}{\sqrt5}. \tag{15}
\]

The proof in the text uses \(|\hat\phi|<1\) to bound
\(|\hat\phi^{\,n}|/\sqrt5<\tfrac12\), so rounding \(\phi^n/\sqrt5\)
gives \(F_n\). This bound holds when \(n\ge1\), but fails for negative
\(n\): e.g.

\[
n=-1:\quad F_{-1}=1,\quad \frac{\phi^{-1}}{\sqrt5}\approx 0.276\ \text{(nearest integer \(0\), not \(1\))}.
\]

Thus (15) is valid for \(n\ge 0\) (indeed for all \(n\ge1\)), but not in
general for negative \(n\).

\subsection{10 {[}15{]}}\label{section-94}

We use Binet's formula (derived in §1.2.8):

\[
F_n=\frac{\phi^{\,n}-\hat\phi^{\,n}}{\sqrt5},\qquad
\hat\phi:=1-\phi=\tfrac{1-\sqrt5}{2}=-\frac1\phi,
\]

so \(\hat\phi\in(-1,0)\). 

Subtract \(F_n\) from \(\phi^{\,n}/\sqrt5\):

\[
\frac{\phi^{\,n}}{\sqrt5}-F_n
=\frac{\phi^{\,n}}{\sqrt5}-\frac{\phi^{\,n}-\hat\phi^{\,n}}{\sqrt5}
=\frac{\hat\phi^{\,n}}{\sqrt5}.
\]

Because \(\hat\phi<0\), the term \(\hat\phi^{\,n}\) is

\begin{itemize}
\tightlist
\item
  positive when \(n\) is even,
\item
  negative when \(n\) is odd,
\end{itemize}

and it is always small in magnitude since \(|\hat\phi|<1\). (Knuth also
notes that this smallness gives the ``nearest--integer'' rule
\(F_n=\phi^{\,n}/\sqrt5\) rounded to the nearest integer.)

Therefore:

\begin{itemize}
\tightlist
\item
  If \(n\) is even, \(\displaystyle \frac{\phi^{\,n}}{\sqrt5} > F_n\).
\item
  If \(n\) is odd, \(\displaystyle \frac{\phi^{\,n}}{\sqrt5} < F_n\).
\end{itemize}

\subsection{11 {[}M20{]}}\label{m20-21}

We'll prove the \(\phi\)-identity for all integers \(n\); the
\(\widehat\phi\)-identity then follows in the same way (since
\(\widehat\phi\) satisfies the same quadratic) or as a corollary at the
end.

Facts we use

\begin{itemize}
\tightlist
\item
  \(\phi\) is the positive root of \(x^2=x+1\); hence \(\phi^2=\phi+1\).
  Likewise, \(\widehat\phi^2=\widehat\phi+1\). 
\item
  Fibonacci numbers satisfy \(F_{n+1}=F_n+F_{n-1}\) with
  \(F_0=0,\;F_1=1\). In §1.2.8 Knuth also suggests extending to negative
  indices by keeping the same recurrence (see ex. 8); this yields
  \(F_{-n}=(-1)^{n+1}F_n\).
\end{itemize}

Induction for \(n\ge 1\)

Base. \(n=1\): \$ \phi\^{}1 = \phi = F\_1\phi + F\_0\$. True.

Induction step. Assume \(\phi^n = F_n\phi+F_{n-1}\). Then

\[
\phi^{n+1}
= \phi\cdot\phi^n
= \phi\,(F_n\phi+F_{n-1})
= F_n\phi^2+F_{n-1}\phi
= F_n(\phi+1)+F_{n-1}\phi
= (F_n+F_{n-1})\phi+F_n
= F_{n+1}\phi+F_n,
\]

using \(\phi^2=\phi+1\) and the Fibonacci recurrence. So the formula
holds for all \(n\ge1\).

Extension to all integers

With the standard extension \(F_{-n}=(-1)^{n+1}F_n\) (obtained by
running the recurrence backwards), the identity at \(n=0\) reads
\(1=\phi^0=F_0\phi+F_{-1}\), so \(F_{-1}=1\), consistent with the
extension. The same recurrence argument then propagates the identity to
every negative \(n\) as well (uniqueness of solutions to
\(a_{n+1}=a_n+a_{n-1}\) given two consecutive values).

Equivalently, one can divide the proven formula for \(n>0\) by
\(\phi^{2n}\) and use \(\widehat\phi=-1/\phi\) to reach the negative
indices.

The \(\widehat\phi\) identity

Repeat the same induction with \(\widehat\phi^2=\widehat\phi+1\), or
deduce it from the proven formula:

\[
\widehat\phi^{\,n} = (-1)^n\phi^{-n}
= (-1)^n\frac{1}{F_n\phi+F_{n-1}}
= F_{n+1}-F_n\phi
= F_n(1-\phi)+F_{n-1}
= F_n\widehat\phi+F_{n-1},
\]

since \(F_{n+1}=F_n+F_{n-1}\) and \(\widehat\phi=1-\phi\). (The
definition of \(\widehat\phi\) is on the same page as the
generating-function derivation. ) 

(Optional cross-check via Binet)

Binet's formula gives \(F_n=\frac{\phi^n-\widehat\phi^{\,n}}{\sqrt5}\).
Solving \(\phi^n\) from \(\phi^n-\widehat\phi^{\,n}=\sqrt5\,F_n\) and
using \(\phi-\widehat\phi=\sqrt5\) plus \(\widehat\phi=1-\phi\) yields
\(\phi^n=F_n\phi+F_{n-1}\) (and similarly for \(\widehat\phi^{\,n}\)). 

Thus both identities hold for every integer \(n\).

\subsection{12 {[}M26{]}}\label{m26-3}

Generating-function solution

Let

\begin{itemize}
\item
  \(F(z)=\sum_{n\ge0} F_n z^n\) be the ordinary Fibonacci GF. In §1.2.8
  Knuth derives

  \[
  F(z)=\frac{z}{1-z-z^2}. \ \ \text{(GF of }F_n\text{)}
  \]
\item
  \(\mathbb F(z)=\sum_{n\ge0}\mathbb F_n z^n\) be the GF of the
  second--order sequence.
\end{itemize}

From \(\mathbb F_{n+2}=\mathbb F_{n+1}+\mathbb F_n+F_n\) with
\(\mathbb F_0=0,\ \mathbb F_1=1\),

\[
\sum_{n\ge0}\mathbb F_{n+2}z^{n+2}
=\sum_{n\ge0}\mathbb F_{n+1}z^{n+2}+\sum_{n\ge0}\mathbb F_n z^{n+2}+\sum_{n\ge0}F_n z^{n+2}.
\]

In terms of \(\mathbb F(z)\) and \(F(z)\):

\[
\mathbb F(z)-z \;=\; z\,\mathbb F(z)+z^2\,\mathbb F(z)+z^2\,F(z),
\]

hence

\[
\mathbb F(z)=\frac{z+z^2F(z)}{1-z-z^2}
=\frac{z(1-z)}{(1-z-z^2)^2}.
\]

Using \(F(z)=\dfrac{z}{1-z-z^2}\), the last form is the neat identity

\[
\boxed{\ \mathbb F(z)=\frac{1-z}{z}\,F(z)^2\ }.
\]

Coefficient extraction now gives a closed form. Since
\(F(z)=\frac{1}{\sqrt5}\Big(\frac{1}{1-\varphi z}-\frac{1}{1-\psi z}\Big)\)
with \(\varphi=\frac{1+\sqrt5}2,\ \psi=\frac{1-\sqrt5}2\),

\[
F(z)^2=\frac{1}{5}\!\left(\frac{1}{(1-\varphi z)^2}+\frac{1}{(1-\psi z)^2}
-\frac{2}{(1-\varphi z)(1-\psi z)}\right).
\]

Multiplying by \((1-z)/z\) and reading coefficients (using
\([z^n](1-\alpha z)^{-2}=(n+1)\alpha^n\) and
\([z^n]\,(1-z-z^2)^{-1}=F_{n+1}\)) yields after a short algebraic
simplification:

\[
\boxed{\ \mathbb F_n=\frac{(3n+3)\,F_n - n\,F_{n+1}}{5}\ } \qquad(n\ge0).
\]

Recurrence (method of undetermined coefficients)

Solve \(\mathbb F_{n+2}-\mathbb F_{n+1}-\mathbb F_n=F_n\).

\begin{itemize}
\item
  Homogeneous part \(x^2-x-1=0\) has solutions \(\varphi^n,\psi^n\), so
  \(\mathbb F_n^{(h)}=c_1F_n+c_2F_{n+1}\).
\item
  Because the RHS is a solution of the homogeneous equation, try a
  particular solution of the form
  \(\mathbb F_n^{(p)}=a\,nF_n+b\,nF_{n+1}\). Plugging in gives

  \[
  \big(\mathbb F_{n+2}^{(p)}-\mathbb F_{n+1}^{(p)}-\mathbb F_n^{(p)}\big)
  =(a+3b)F_{n+1}+(2a+b)F_n.
  \]

  Matching this to \(F_n\) yields \(a=3/5,\ b=-1/5\), so
  \(\mathbb F_n^{(p)}=\frac{n}{5}(3F_n-F_{n+1})\).
\end{itemize}

Thus \(\mathbb F_n=\dfrac{n}{5}(3F_n-F_{n+1})+c_1F_n+c_2F_{n+1}\).
Initial conditions \(\mathbb F_0=0,\ \mathbb F_1=1\) force \(c_2=0\) and
\(c_1=3/5\), giving exactly

\[
\boxed{\ \mathbb F_n=\frac{(3n+3)\,F_n - n\,F_{n+1}}{5}\ }.
\]

Quick check (first values)

The formula gives \(\mathbb F_n: 0,1,1,3,5,10,18,33,59,105,\ldots\),
which indeed satisfies
\(\mathbb F_{n+2}=\mathbb F_{n+1}+\mathbb F_n+F_n\) and the seeds.

\subsection{13 {[}M22{]}}\label{m22-13}

We'll use the fact that any sequence satisfying the Fibonacci recurrence
\(x_{n+2}=x_{n+1}+x_n\) is a linear combination of consecutive Fibonacci
numbers:

\[
x_n=\alpha F_n+\beta F_{n+1},
\]

because this form obeys the same recurrence and the constants
\(\alpha,\beta\) can be fixed from two initial values.

Recall \(F_0=0,\ F_1=1\).

\begin{enumerate}
\def\labelenumi{(\alph{enumi})}
\tightlist
\item
  Homogeneous recurrence
\end{enumerate}

Assume

\[
a_n=\alpha F_n+\beta F_{n+1}.
\]

Match the initials:

\[
a_0=r=\alpha F_0+\beta F_1=\beta \ \Rightarrow\ \beta=r,
\]

\[
a_1=s=\alpha F_1+\beta F_2=\alpha+r \ \Rightarrow\ \alpha=s-r.
\]

Hence

\[
\boxed{\,a_n=r\,F_{n+1}+(s-r)\,F_n\,}\qquad(n\ge 0).
\]

Check. \(n=0:\ rF_1+(s-r)F_0=r\). \(n=1:\ rF_2+(s-r)F_1=r+s-r=s\). The
recurrence holds because any \(\alpha F_n+\beta F_{n+1}\) satisfies
\(x_{n+2}=x_{n+1}+x_n\).

Equivalent forms sometimes seen are \(a_n=sF_n+rF_{n-1}\) (if one
extends \(F_{-1}=1\)), but the boxed form avoids negative indices.

\begin{enumerate}
\def\labelenumi{(\alph{enumi})}
\setcounter{enumi}{1}
\tightlist
\item
  Inhomogeneous recurrence
\end{enumerate}

The nonzero constant \(c\) suggests removing a constant particular
solution. Try \(b_n\equiv t\):

\[
t=t+t+c\ \Rightarrow\ t=-c.
\]

Set \(u_n=b_n+c\). Then

\[
u_{n+2}=b_{n+2}+c=(b_{n+1}+b_n+c)+c=u_{n+1}+u_n,
\]

so \(\{u_n\}\) satisfies the Fibonacci recurrence, with

\[
u_0=b_0+c=c,\qquad u_1=b_1+c=1+c.
\]

Apply the result from (a):

\[
u_n=(u_1-u_0)F_n+u_0F_{n+1}=F_n+c\,F_{n+1}.
\]

Finally shift back:

\[
\boxed{\,b_n=F_n+c\,F_{n+1}-c=F_n+c\,(F_{n+1}-1)\,}\qquad(n\ge 0).
\]

Check. \(n=0:\ F_0+cF_1-c=0\). \(n=1:\ F_1+cF_2-c=1\). And
\(b_{n+2}-b_{n+1}-b_n=\big(F_{n+2}-F_{n+1}-F_n\big)+c\big(F_{n+3}-F_{n+2}-F_{n+1}\big)-c=c\).

\subsection{14 {[}M28{]}}\label{m28-3}

Introduce the ``binomial diagonal'' sum

\[
S_n\;:=\;\sum_{r=0}^{m}\binom{n+r}{\,m-r\,}
=\binom{n}{m}+\binom{n+1}{m-1}+\cdots+\binom{n+m}{0}.
\]

By Pascal's identity \(\binom{x+1}{y}=\binom{x}{y}+\binom{x}{y-1}\),

\[
\begin{aligned}
S_{n+2}
&=\sum_{r=0}^{m}\!\binom{n+2+r}{m-r}
=\sum_{r=0}^{m}\!\bigl[\binom{n+1+r}{m-r}+\binom{n+1+r}{m-r-1}\bigr]\\
&=S_{n+1}+\sum_{r=0}^{m-1}\binom{n+1+r}{m-r-1}
=S_{n+1}+\sum_{t=1}^{m}\binom{n+t}{m-t}
=S_{n+1}+S_n-\binom{n}{m}.
\end{aligned}
\]

Hence

\[
\boxed{\,S_{n+2}-S_{n+1}-S_n=-\binom{n}{m}\,}.
\]

Now set \(b_n:=a_n+S_n\). Then

\[
b_{n+2}=(a_{n+2}+S_{n+2})=(a_{n+1}+a_n+\binom{n}{m})+(S_{n+1}+S_n-\binom{n}{m})
=b_{n+1}+b_n.
\]

So \(\{b_n\}\) satisfies the homogeneous Fibonacci recurrence.

\begin{enumerate}
\def\labelenumi{\arabic{enumi})}
\setcounter{enumi}{1}
\tightlist
\item
  Initial values of \(b_n\)
\end{enumerate}

We need \(S_0\) and \(S_1\). These are classic ``shallow diagonal'' sums
in Pascal's triangle:

\[
S_0=\sum_{r=0}^{m}\binom{r}{m-r}=F_{m+1},\qquad
S_1=\sum_{r=0}^{m}\binom{1+r}{m-r}=F_{m+2},
\]

(which follows from \(\sum_{k\ge 0}\binom{n-k}{k}=F_{n+1}\) by a
one-line induction).

Thus

\[
b_0=a_0+S_0=F_{m+1},\qquad b_1=a_1+S_1=1+F_{m+2}.
\]

\begin{enumerate}
\def\labelenumi{\arabic{enumi})}
\setcounter{enumi}{2}
\tightlist
\item
  Solve the Fibonacci-type initial value problem
\end{enumerate}

Every solution of \(x_{n+2}=x_{n+1}+x_n\) is a linear combination of
\(F_{n-1}\) and \(F_n\):

\[
b_n=b_0F_{n-1}+b_1F_n
=F_{m+1}F_{n-1}+(1+F_{m+2})F_n.
\]

Use the addition identity \(F_{m+1}F_{n-1}+F_{m+2}F_{n}=F_{m+n+1}\) to
get

\[
b_n=F_{m+n+1}+F_n.
\]

Finally,

\[
a_n=b_n-S_n
=\boxed{\,F_{m+n+1}+F_n-\sum_{r=0}^{m}\binom{n+r}{\,m-r\,}\,}.
\]

Equivalently, writing out the diagonal sum,

\[
\boxed{\;a_n=F_{m+n+1}+F_n-\Bigl(\binom{n}{m}+\binom{n+1}{m-1}+\cdots+\binom{n+m}{0}\Bigr)\; }.
\]

This matches Knuth's official answer for 1.2.8--14.

\begin{enumerate}
\def\labelenumi{\arabic{enumi})}
\setcounter{enumi}{3}
\tightlist
\item
  Quick check (e.g., \(m=1\))
\end{enumerate}

Then \(\binom{n}{1}=n\) and \(S_n=\binom{n}{1}+\binom{n+1}{0}=n+1\). Our
formula gives

\[
a_n=F_{n+2}+F_n-(n+1).
\]

Compute from the recurrence: \(a_0=0\), \(a_1=1\),
\(a_2=a_1+a_0+\binom{0}{1}=1\), \(a_3=a_2+a_1+\binom{1}{1}=3\), \ldots{}
The closed form yields \(a_2=F_4+F_2-3=3+1-3=1\),
\(a_3=F_5+F_3-4=5+2-4=3\).

\subsection{15 {[}M22{]}}\label{m22-14}

Let the linear operator

\[
(Lu)_n:=u_{n+2}-u_{n+1}-u_n.
\]

Then \(L\) is linear, and the three recurrences say

\[
La=f,\qquad Lb=g,\qquad Lc=x\,f+y\,g,
\]

with the common initial conditions \(a_0=b_0=c_0=0\) and
\(a_1=b_1=c_1=1\). Also, for Fibonacci \(F\) we have \(LF=0\) and
\(F_0=0,\ F_1=1\).

Consider the candidate

\[
d_n \;:=\; x\,a_n+y\,b_n+(1-x-y)\,F_n.
\]

Because \(L\) is linear,

\[
Ld \;=\; x\,La+y\,Lb+(1-x-y)LF \;=\; x\,f+y\,g,
\]

and the initial values match:

\[
d_0=x\cdot 0+y\cdot 0+(1-x-y)\cdot 0=0,\qquad
d_1=x\cdot 1+y\cdot 1+(1-x-y)\cdot 1=1.
\]

Thus \(d\) and \(c\) satisfy the same inhomogeneous recurrence with the
same initial conditions, hence \(d_n=c_n\) for all \(n\).

Therefore the requested expression is

\[
\boxed{\,c_n \;=\; x\,a_n \;+\; y\,b_n \;+\; (1-x-y)\,F_n\,}\,.
\]

(Optional) Convolution check

One can also derive the standard particular-solution form: if
\(u_{n+2}=u_{n+1}+u_n+h(n)\) with \(u_0=0,u_1=1\), then

\[
u_n \;=\; F_n\;+\sum_{k=0}^{n-2}F_{n-1-k}\,h(k).
\]

Applying this to \(a,b,c\) (with \(h=f,g,xf+yg\)) gives

\[
a_n=F_n+\sum_{k}F_{n-1-k}f(k),\quad
b_n=F_n+\sum_{k}F_{n-1-k}g(k),
\]

\[
c_n=F_n+x\sum_{k}F_{n-1-k}f(k)+y\sum_{k}F_{n-1-k}g(k)
      =x a_n+y b_n+(1-x-y)F_n,
\]

agreeing with the result.

\subsection{16 {[}M20{]}}\label{m20-22}

Here is Exercise 16 from §1.2.8 (Fibonacci Numbers): show that

\[
S_n \;\stackrel{\text{def}}{=}\; \sum_{k=0}^{n}\binom{n-k}{k}
\]

is a Fibonacci number. (This is the ``right-angled'' view of Pascal's
triangle.) 

We'll prove the precise identity

\[
S_n = F_{n+1},
\]

in three complementary ways.

\begin{enumerate}
\def\labelenumi{\arabic{enumi})}
\tightlist
\item
  Recurrence proof (via Pascal's identity)
\end{enumerate}

Let \(S_n=\sum_{k=0}^{n}\binom{n-k}{k}\). Using
\(\binom{a+1}{b}=\binom{a}{b}+\binom{a}{b-1}\),

\[
\begin{aligned}
S_{n+1}
&=\sum_{k=0}^{n+1}\binom{n+1-k}{k}
=\sum_{k=0}^{n+1}\Bigl[\binom{n-k}{k}+\binom{n-k}{k-1}\Bigr]\\
&=\underbrace{\sum_{k=0}^{n+1}\binom{n-k}{k}}_{=\,S_n}
\;+\;\underbrace{\sum_{k=1}^{n+1}\binom{n-k}{k-1}}_{j=k-1\Rightarrow \sum_{j=0}^{n}\binom{n-1-j}{j}=S_{n-1}}\\[2mm]
&= S_n+S_{n-1}.
\end{aligned}
\]

Initial values: \(S_0=\binom{0}{0}=1\) and
\(S_1=\binom{1}{0}+\binom{0}{1}=1\). Thus \((S_n)\) satisfies the
Fibonacci recurrence with \(S_0=1,S_1=1\), hence \(S_n=F_{n+1}\).

\begin{enumerate}
\def\labelenumi{\arabic{enumi})}
\setcounter{enumi}{1}
\tightlist
\item
  Combinatorial proof (tilings)
\end{enumerate}

Count tilings of a \(1\times n\) strip by monominoes (length 1) and
dominoes (length 2).

\begin{itemize}
\tightlist
\item
  If a tiling uses \(k\) dominoes, it has \(n-2k\) monominoes, hence
  \(n-k\) tiles in total; choosing which \(k\) of these \(n-k\)
  positions are dominoes yields \(\binom{n-k}{k}\) tilings.
\item
  Summing over all feasible \(k\) gives exactly
  \(\sum_{k\ge 0}\binom{n-k}{k}=S_n\).
\end{itemize}

On the other hand, these tilings are counted by the Fibonacci
recurrence: a tiling of length \(n\) either starts with a monomino (then
continue with a tiling of length \(n-1\)) or with a domino (then
continue with a tiling of length \(n-2\)). With \(T_0=1,\;T_1=1\), we
have \(T_n=F_{n+1}\). Therefore \(S_n=T_n=F_{n+1}\).

\begin{enumerate}
\def\labelenumi{\arabic{enumi})}
\setcounter{enumi}{2}
\tightlist
\item
  Generating-function proof
\end{enumerate}

Define \(S(z)=\sum_{n\ge 0} S_n z^n\). Then

\[
\begin{aligned}
S(z)
&=\sum_{n\ge 0}\sum_{k\ge 0}\binom{n-k}{k}z^n
=\sum_{k\ge 0}\sum_{m\ge 0}\binom{m}{k}z^{m+k}
=\sum_{k\ge 0}z^{k}\!\left(\sum_{m\ge 0}\binom{m}{k}z^{m}\right)\\
&=\sum_{k\ge 0}z^{k}\,\frac{z^{k}}{(1-z)^{k+1}}
=\frac{1}{1-z}\sum_{k\ge 0}\left(\frac{z^{2}}{1-z}\right)^{k}
=\frac{1}{1-z}\cdot\frac{1}{1-\frac{z^{2}}{1-z}}
=\frac{1}{1-z-z^{2}}.
\end{aligned}
\]

But \(\sum_{n\ge 0}F_{n+1}z^{n}= \dfrac{1}{1-z-z^{2}}\) (since
\(G(z)=\sum_{n\ge 0}F_n z^n=\dfrac{z}{1-z-z^{2}}\)). Hence
\(S_n=F_{n+1}\). 

Therefore, for every \(n\ge 0\),

\[
\boxed{\displaystyle \sum_{k=0}^{n}\binom{n-k}{k}=F_{n+1}.}
\]

(Quick check: \(n=0,1,2,3\) give \(1,1,2,3\), matching
\(F_{1},F_{2},F_{3},F_{4}\) respectively.)

\subsection{17 {[}M24{]}}\label{m24-9}

We must show, for all integers \(m,n,k\),

\[
F_{n+k}F_{m-k}-F_nF_m \;=\;(-1)^n\,F_{m-n-k}\,F_k.
\tag{★}
\]

Ingredients

\begin{enumerate}
\def\labelenumi{\arabic{enumi}.}
\tightlist
\item
  Addition formula (valid for all integers):
\end{enumerate}

\[
F_{r+s}=F_rF_{s+1}+F_{r-1}F_s.
\tag{A}
\]

\begin{enumerate}
\def\labelenumi{\arabic{enumi}.}
\setcounter{enumi}{1}
\tightlist
\item
  d'Ocagne's identity (valid for all integers):
\end{enumerate}

\[
F_aF_{b+1}-F_{a+1}F_b \;=\; (-1)^b\,F_{a-b}.
\tag{D}
\]

\begin{enumerate}
\def\labelenumi{\arabic{enumi}.}
\setcounter{enumi}{2}
\tightlist
\item
  Negative indices convention:
\end{enumerate}

\[
F_{-t}=(-1)^{t+1}F_t \quad\text{for all integers }t\ge 0.
\tag{N}
\]

Start from the left--hand side of (★) and expand only the factors that
``carry'' \(k\) using (A):

\[
\begin{aligned}
F_{n+k}F_{m-k}-F_nF_m
&=\bigl(F_nF_{k+1}+F_{n-1}F_k\bigr)F_{m-k}
 \;-\;F_n\bigl(F_{m-k}F_{k+1}+F_{m-k-1}F_k\bigr)\\[4pt]
&=\underbrace{F_nF_{k+1}F_{m-k}}_{\text{cancels}}
 +F_{n-1}F_kF_{m-k}
 -\underbrace{F_nF_{m-k}F_{k+1}}_{\text{cancels}}
 -F_nF_{m-k-1}F_k\\[4pt]
&=F_k\bigl(F_{n-1}F_{m-k}-F_nF_{m-k-1}\bigr).
\end{aligned}
\tag{1}
\]

Apply d'Ocagne (D) with \(a=n-1\) and \(b=m-k-1\):

\[
F_{n-1}F_{m-k}-F_nF_{m-k-1}
= (-1)^{m-k-1}\,F_{(n-1)-(m-k-1)}
= (-1)^{m-k-1}\,F_{n-m+k}.
\tag{2}
\]

Insert (2) into (1):

\[
F_{n+k}F_{m-k}-F_nF_m
= F_k \cdot (-1)^{m-k-1} F_{n-m+k}.
\tag{3}
\]

Now use the negative-index rule (N): since \(m-n-k=-(n-m+k)\),

\[
F_{n-m+k}=F_{-(m-n-k)}= (-1)^{(m-n-k)+1}F_{m-n-k}.
\tag{4}
\]

Combine (3) and (4): the total sign exponent is

\[
(m-k-1) + (m-n-k+1) \;=\; n,
\]

so

\[
F_{n+k}F_{m-k}-F_nF_m
= (-1)^n\,F_k\,F_{m-n-k},
\]

which is exactly (★).

Checks and corollaries

\begin{itemize}
\item
  \(k=0\): LHS \(=F_nF_m-F_nF_m=0\); RHS \(=(-1)^nF_{m-n}F_0=0\).
\item
  \(k=1\) (d'Ocagne's classical form):

  \[
  F_{n+1}F_{m-1}-F_nF_m = (-1)^n F_{m-n-1}.
  \]
\item
  Generalized Cassini (take \(m=n\)):

  \[
  F_{n+k}F_{n-k}-F_n^2
  =(-1)^nF_{-k}F_k
  =(-1)^{n+k+1}F_k^2,
  \]

  a well-known extension of Cassini's identity.
\end{itemize}

This completes the proof.

\subsection{18 {[}20{]}}\label{section-95}

Yes. In fact,

\[
F_n^2+F_{n+1}^2 \;=\; F_{2n+1}\qquad\text{for all }n\ge 0.
\]

This holds under either common indexing convention (\(F_0=0,F_1=1\) or
\(F_1=1,F_2=1\)).

Proof (via Fibonacci addition formula)

We use the standard addition identity

\[
F_{r+s}=F_rF_{s+1}+F_{r-1}F_s\quad\text{for all }r,s\ge 0.
\]

Proof of the identity (induction on \(s\)): For \(s=0\):
\(F_{r+0}=F_r=F_rF_1+F_{r-1}F_0\) (true). Assume true for \(s\). Then

\[
\begin{aligned}
F_{r+(s+1)}
&=F_{r+s+1}
=F_{r+s}+F_{r+s-1}\\
&=(F_rF_{s+1}+F_{r-1}F_s)+(F_rF_s+F_{r-1}F_{s-1})\\
&=F_r(F_{s+1}+F_s)+F_{r-1}(F_s+F_{s-1})\\
&=F_rF_{s+2}+F_{r-1}F_{s+1},
\end{aligned}
\]

so the identity holds for \(s+1\). By induction, it holds for all \(s\).

Now take \(r=n+1\) and \(s=n\). Then

\[
F_{(n+1)+n}
=F_{2n+1}
=F_{n+1}F_{n+1}+F_nF_n
=F_{n+1}^2+F_n^2,
\]

which is exactly the claim.

Alternative proof (matrix method, one line)

Let \(Q=\begin{pmatrix}1&1\\[2pt]1&0\end{pmatrix}\). Then

\[
Q^k=\begin{pmatrix}F_{k+1}&F_k\\[2pt]F_k&F_{k-1}\end{pmatrix}.
\]

Compute \(Q^{2n+1}=Q^nQ^{n+1}\) and look at the \((1,2)\) entry:

\[
(F_{n+1})(F_{n+1})+(F_n)(F_n)=F_{2n+1}.
\]

Quick check

\begin{itemize}
\tightlist
\item
  \(n=0:\;0^2+1^2=1=F_1=F_{2\cdot0+1}\).
\item
  \(n=1:\;1^2+1^2=2=F_3\).
\item
  \(n=2:\;1^2+2^2=5=F_5\).
\item
  \(n=3:\;2^2+3^2=13=F_7\).
\end{itemize}

\subsection{19 {[}M27{]}}\label{m27-2}

Let \(\theta=36^\circ\) and \(c=\cos\theta\). The Chebyshev identity

\[
\cos 5\theta \;=\;16c^5-20c^3+5c
\]

gives, since \(5\theta=180^\circ\) and \(\cos 180^\circ=-1\),

\[
16c^5-20c^3+5c+1=0.
\]

Set \(u=2c\). Substituting \(c=u/2\) yields

\[
u^5-5u^3+5u+2=0.
\]

This factors as

\[
(u^2-u-1)(u^3+u^2-3u-2)=0.
\]

Hence \(u\) satisfies \(u^2-u-1=0\), so

\[
u=\frac{1+\sqrt5}{2}=\varphi,
\]

the golden ratio. (The other quadratic root is negative and doesn't
match \(u=2\cos\theta\in(0,2]\).) Therefore

\[
\boxed{\;\cos 36^\circ \;=\; \frac{u}{2}\;=\;\frac{\varphi}{2}\;=\;\frac{1+\sqrt5}{4}\;}\,.
\]

By TAOCP's notation, \(\varphi=\tfrac12(1+\sqrt5)\). 

Geometric check (golden triangle)

In a regular pentagon the ratio of diagonal to side is \(\varphi\).
Considering the isosceles triangle formed by two diagonals (equal sides)
and one side (base) with vertex angle \(36^\circ\), the law of cosines
gives

\[
1^2=2\varphi^2\,(1-\cos36^\circ)\;\Longrightarrow\;
\cos36^\circ=1-\frac{1}{2\varphi^2}=\frac{\varphi}{2},
\]

using \(\varphi^2=\varphi+1\). Related values (nice to know)

From \(\cos36^\circ=\varphi/2\) we get, equivalently,

\[
\cos72^\circ=\frac{1}{2\varphi}=\frac{\sqrt5-1}{4},\qquad
\sin36^\circ=\sqrt{\frac{10-2\sqrt5}{16}}.
\]

\subsection{20 {[}M16{]}}\label{m16}

\[
\boxed{\displaystyle \sum_{k=0}^{n} F_k \;=\; F_{n+2}\;-\;1}
\]

(So the sum of the first \(n+1\) Fibonacci numbers is the \((n+2)\)-nd
Fibonacci number minus 1.)

Proof 1 --- One-line induction (using the recurrence)

Let \(S_n=\sum_{k=0}^{n}F_k\). The Fibonacci recurrence gives

\[
S_n=S_{n-1}+F_n.
\]

Assume \(S_{n-1}=F_{n+1}-1\). Then

\[
S_n=(F_{n+1}-1)+F_n=F_{n+2}-1,
\]

since \(F_{n+2}=F_{n+1}+F_n\). The base case \(S_0=F_0=0=F_2-1\) holds,
so by induction \(S_n=F_{n+2}-1\) for all \(n\ge 0\).

\emph{(This uses only the standard definition
\(F_0=0,\;F_1=1,\;F_{n+2}=F_{n+1}+F_n\).)}

Proof 2 --- Generating functions

Let \(G(z)=\sum_{n\ge 0}F_n z^n\). From the text,

\[
G(z)=\frac{z}{1-z-z^2}.
\]

If \(S_n=\sum_{k=0}^{n}F_k\), then
\(S(z)=\sum_{n\ge 0}S_n z^n=\frac{G(z)}{1-z}\). Hence

\[
S(z)=\frac{z}{(1-z)(1-z-z^2)}.
\]

But
\(\sum_{n\ge 0}(F_{n+2}-1)z^n=\frac{G(z)-z}{z^2}-\frac{1}{1-z}=\frac{z}{(1-z)(1-z-z^2)}\)
as well, so the coefficients agree: \(S_n=F_{n+2}-1\). 

Proof 3 --- Binet's formula (closed form)

Using \(F_n=\dfrac{\phi^n-\phi'^n}{\sqrt{5}}\) with
\(\phi=\tfrac{1+\sqrt5}{2}\) and \(\phi'=1-\phi\), we have

\[
\sum_{k=0}^{n}F_k
=\frac{1}{\sqrt5}\left(\sum_{k=0}^{n}\phi^k-\sum_{k=0}^{n}\phi'^k\right)
=\frac{1}{\sqrt5}\!\left(\frac{\phi^{n+1}-1}{\phi-1}-\frac{{\phi'}^{\,n+1}-1}{\phi'-1}\right).
\]

Since \(\phi-1=\tfrac{1}{\phi}\) and \(\phi'-1=\tfrac{1}{\phi'}\), this
simplifies to

\[
\frac{1}{\sqrt5}\bigl(\phi^{n+2}-\phi-({\phi'}^{\,n+2}-\phi')\bigr)
=\frac{\phi^{n+2}-{\phi'}^{\,n+2}}{\sqrt5}-\frac{\phi-\phi'}{\sqrt5}
=F_{n+2}-1,
\]

because \(\phi-\phi'=\sqrt5\).

Quick check

For \(n=7\): \(0+1+1+2+3+5+8+13=33\), and \(F_9-1=34-1=33\).

\subsection{21 {[}M25{]}}\label{m25-15}

Evaluate

\[
S_n(x)\;=\;\sum_{k=1}^{n} F_k\,x^k,
\]

where \(F_0=0,\,F_1=1\), and \(F_{k}=F_{k-1}+F_{k-2}\) for \(k\ge2\).

\begin{enumerate}
\def\labelenumi{\arabic{enumi})}
\tightlist
\item
  Main case \(x^2+x-1\ne0\)
\end{enumerate}

Start with

\[
S_n=\sum_{k=1}^{n}F_kx^k,\quad
xS_n=\sum_{k=1}^{n}F_kx^{k+1}=\sum_{k=2}^{n+1}F_{k-1}x^{k},\quad
x^2S_n=\sum_{k=3}^{n+2}F_{k-2}x^{k}.
\]

Consider \((x^2+x-1)S_n = x^2S_n + xS_n - S_n\). For the coefficients of
\(x^k\) with \(3\le k\le n\), the Fibonacci recurrence gives perfect
cancellation:

\[
F_{k-2}+F_{k-1}-F_k=0.
\]

Only boundary terms remain:

\begin{itemize}
\tightlist
\item
  \(k=1\): contributes \(-F_1x=-x\).
\item
  \(k=2\): contributes \(F_1x^2-F_2x^2=0\).
\item
  \(k=n+1\): contributes \((F_{n-1}+F_n)x^{n+1}=F_{n+1}x^{n+1}\).
\item
  \(k=n+2\): contributes \(F_nx^{n+2}\).
\end{itemize}

Hence

\[
(x^2+x-1)\,S_n \;=\; F_{n+1}x^{n+1}+F_nx^{n+2}-x,
\]

and, when \(x^2+x-1\ne0\),

\[
\boxed{\;
\displaystyle
S_n(x)=\frac{x^{n+1}F_{n+1}+x^{n+2}F_n - x}{\,x^2+x-1\,}.
\;}
\]

(This matches Knuth's posted answer for the exercise. )

A good sanity check is \(x=1\):

\[
S_n(1)=\sum_{k=1}^nF_k=\frac{F_{n+1}+F_n-1}{1+1-1}=F_{n+2}-1,
\]

the classic identity (this is Ex. 20's result). 

\begin{enumerate}
\def\labelenumi{\arabic{enumi})}
\setcounter{enumi}{1}
\tightlist
\item
  The exceptional case \(x^2+x-1=0\)
\end{enumerate}

The denominator vanishes at the roots of \(x^2+x-1=0\), namely

\[
x=\frac{-1\pm\sqrt5}{2}=\frac1\varphi \quad\text{or}\quad \frac1{\varphi'},
\]

with \(\varphi=\tfrac{1+\sqrt5}{2}\) and
\(\varphi'=\tfrac{1-\sqrt5}{2}\). Here the numerator vanishes as well,
so take the limit (e.g., by l'Hôpital). Differentiating numerator and
denominator and simplifying yields

\[
\boxed{\;
\displaystyle
S_n(x)=\frac{\,n+1 - x^{\,n}F_{n+1}\,}{\,2x+1\,}\quad\text{when }x^2+x-1=0.
\;}
\]

(This is exactly the special-case expression reported in the book's
answer section. )

\begin{enumerate}
\def\labelenumi{\arabic{enumi})}
\setcounter{enumi}{2}
\tightlist
\item
  Remarks
\end{enumerate}

\begin{itemize}
\tightlist
\item
  Letting \(n\to\infty\) with \(|x|<1/\varphi\) recovers the ordinary
  generating function \(\sum_{k\ge1}F_kx^k = \dfrac{x}{1-x-x^2}\).
\item
  The telescoping proof above is robust and avoids Binet's formula; it's
  the standard trick for finite ``truncated generating functions.''
\end{itemize}

Final answer

\[
\boxed{\displaystyle
\sum_{k=1}^{n}F_kx^k=
\begin{cases}
\displaystyle\frac{x^{n+1}F_{n+1}+x^{n+2}F_n-x}{x^2+x-1}, & x^2+x-1\ne0,\\[1.25em]
\displaystyle\frac{n+1-x^nF_{n+1}}{2x+1}, & x^2+x-1=0~,
\end{cases}}
\]

with the check \(\sum_{k=1}^{n}F_k=F_{n+2}-1\) at \(x=1\).

\subsection{22 {[}M20{]}}\label{m20-23}

Proof via Binet's formula

Recall Binet's formula

\[
F_r=\frac{\phi^{\,r}-\hat\phi^{\,r}}{\sqrt5},
\qquad 
\phi=\frac{1+\sqrt5}{2},\quad
\hat\phi=1-\phi=\frac{1-\sqrt5}{2}.
\]

Both \(\phi\) and \(\hat\phi\) satisfy \(x^2=x+1\), so \(1+\phi=\phi^2\)
and \(1+\hat\phi=\hat\phi^{\,2}\). 

Then

\[
\begin{aligned}
S
&=\sum_{k=0}^{n}\binom{n}{k}\frac{\phi^{\,m+k}-\hat\phi^{\,m+k}}{\sqrt5}
=\frac{1}{\sqrt5}\!\left(\phi^{\,m}\sum_{k=0}^{n}\binom{n}{k}\phi^{\,k}
-\hat\phi^{\,m}\sum_{k=0}^{n}\binom{n}{k}\hat\phi^{\,k}\right)\\
&=\frac{1}{\sqrt5}\Bigl(\phi^{\,m}(1+\phi)^n-\hat\phi^{\,m}(1+\hat\phi)^n\Bigr)
=\frac{1}{\sqrt5}\bigl(\phi^{\,m}\phi^{\,2n}-\hat\phi^{\,m}\hat\phi^{\,2n}\bigr)
=\frac{\phi^{\,m+2n}-\hat\phi^{\,m+2n}}{\sqrt5}\\
&=F_{m+2n}.
\end{aligned}
\]

Thus the sum is exactly the Fibonacci number \(F_{m+2n}\). □

(Optional) Recurrence/induction proof (no Binet)

Define \(S_n(m)=\sum_{k=0}^{n}\binom{n}{k}F_{m+k}\). Pascal's identity
gives

\[
S_{n+1}(m)=\sum_{k}\!\bigl(\binom{n}{k}+\binom{n}{k-1}\bigr)F_{m+k}
=S_n(m)+S_n(m+1).
\]

Base cases: \(S_0(m)=F_m\) and \(S_1(m)=F_m+F_{m+1}=F_{m+2}\). The
sequence \(T_n(m):=F_{m+2n}\) satisfies the same recurrence
\(T_{n+1}(m)=T_n(m)+T_n(m+1)\) and the same bases
\(T_0(m)=F_m,\ T_1(m)=F_{m+2}\). By induction on \(n\),
\(S_n(m)=T_n(m)=F_{m+2n}\). □

Sanity check. \(n=2\): \(F_m+2F_{m+1}+F_{m+2}=F_{m+4}\) (true since
\(F_{m+4}=F_{m+3}+F_{m+2}=(F_{m+2}+F_{m+1})+(F_{m+1}+F_m)\)).

\subsection{23 {[}M23{]}}\label{m23-14}

We use Binet's formula and one standard identity for powers of
\(\varphi\) (the golden ratio):

\[
F_r \;=\; \frac{\varphi^{\,r}-\hat\varphi^{\,r}}{\sqrt5}, \qquad
\hat\varphi = 1-\varphi,
\]

and

\[
\varphi^{\,t}=F_t\,\varphi+F_{t-1},\qquad \hat\varphi^{\,t}=F_t\,\hat\varphi+F_{t-1}.
\]

(See the derivation of Binet's formula and related identities for
\(\varphi\) in §1.2.8. )

Start from the left--hand side and insert Binet's formula for
\(F_{m+k}\):

\[
\sum_{k=0}^{n} \binom{n}{k}\,F_{m+k}\,F_t^{\,k}\,F_{t-1}^{\,n-k}
= \frac{1}{\sqrt5}\sum_{k=0}^{n}\binom{n}{k}\bigl(\varphi^{\,m+k}-\hat\varphi^{\,m+k}\bigr)F_t^{\,k}F_{t-1}^{\,n-k}.
\]

Separate the \(\varphi\)-part and the \(\hat\varphi\)-part, and factor
out \(\varphi^m\) and \(\hat\varphi^m\):

\[
= \frac{1}{\sqrt5}\left[
\varphi^{\,m}\sum_{k=0}^{n}\binom{n}{k}(\varphi F_t)^{k}F_{t-1}^{\,n-k}
\;-\;
\hat\varphi^{\,m}\sum_{k=0}^{n}\binom{n}{k}(\hat\varphi F_t)^{k}F_{t-1}^{\,n-k}
\right].
\]

Apply the binomial theorem to each sum:

\[
\sum_{k=0}^{n}\binom{n}{k}(\varphi F_t)^kF_{t-1}^{\,n-k}
= (\varphi F_t + F_{t-1})^n
= (\varphi^{\,t})^n
= \varphi^{\,tn},
\]

and similarly
\((\hat\varphi F_t+F_{t-1})^n=(\hat\varphi^{\,t})^n=\hat\varphi^{\,tn}\).
Therefore

\[
\sum_{k=0}^{n} \binom{n}{k}\,F_{m+k}\,F_t^{\,k}\,F_{t-1}^{\,n-k}
= \frac{1}{\sqrt5}\bigl(\varphi^{\,m}\varphi^{\,tn}-\hat\varphi^{\,m}\hat\varphi^{\,tn}\bigr)
= \frac{1}{\sqrt5}\bigl(\varphi^{\,m+tn}-\hat\varphi^{\,m+tn}\bigr)
= F_{\,m+tn},
\]

which is exactly the required identity. □

Quick check

Take \(m=1,\,n=2,\,t=3\). LHS:

\[
\binom20 F_1\,2^0\,1^2 + \binom21 F_2\,2^1\,1^1 + \binom22 F_3\,2^2\,1^0
= 1 + 4 + 8 = 13,
\]

RHS: \(F_{1+3\cdot 2}=F_7=13\).

\subsection{24 {[}HM20{]}}\label{hm20-11}

Let \(D_n\) be the determinant of the \(n\times n\) tridiagonal matrix

\[
A_n=\begin{pmatrix}
1 & -1 & 0 & \cdots & 0 \\
1 & 1 & -1 & \ddots & \vdots \\
0 & 1 & 1 & \ddots & 0 \\
\vdots & \ddots & \ddots & \ddots & -1\\
0 & \cdots & 0 & 1 & 1
\end{pmatrix}.
\]

Show that \(D_n=F_{n+1}\), the \((n\!+\!1)\)-st Fibonacci number.

Step 1: Set up a clean recurrence for \(D_n\)

Expand \(D_n=\det(A_n)\) along the first row:

\[
D_n
= (1)\cdot \det(M_{11}) + (-1)\cdot C_{12},
\]

where \(M_{11}\) is the minor obtained by deleting row 1, column 1, and
\(C_{12}\) is the cofactor term for entry \((1,2)\). By the usual
cofactor sign, \(C_{12}=(-1)^{1+2}\det(M_{12})=-\det(M_{12})\). Hence

\[
D_n=\det(M_{11})+\det(M_{12}). \tag{1}
\]

\begin{itemize}
\tightlist
\item
  Identify \(\det(M_{11})\). Deleting row 1, column 1 from \(A_n\)
  leaves exactly the same tridiagonal pattern of size \(n-1\). Thus
\end{itemize}

\[
\det(M_{11})=D_{n-1}.
\]

\begin{itemize}
\tightlist
\item
  Identify \(\det(M_{12})\). Deleting row 1, column 2 yields a matrix
  whose first column has a single \(1\) at its top and zeros below.
  Expanding \emph{that} minor along its first column removes that top
  row and left column, leaving again the same pattern of size \(n-2\).
  Therefore
\end{itemize}

\[
\det(M_{12})=D_{n-2}.
\]

Plug these into (1):

\[
\boxed{D_n=D_{n-1}+D_{n-2}\qquad(n\ge 3).} \tag{2}
\]

This is exactly the Fibonacci recurrence.

Step 2: Initial conditions

Compute the first few determinants directly:

\begin{itemize}
\tightlist
\item
  \(n=1\): \(A_1=[1]\), so \(D_1=1=F_2\).
\item
  \(n=2\): \(A_2=\begin{pmatrix}1&-1\\[2pt]1&1\end{pmatrix}\), so
  \(\;D_2=1\cdot1-(-1)\cdot1=2=F_3\).
\end{itemize}

Thus

\[
\boxed{D_1=1=F_2,\qquad D_2=2=F_3.} \tag{3}
\]

Step 3: Conclude by induction (or by ``same recurrence, same bases'')

The sequence \((D_n)\) satisfies the Fibonacci recurrence (2) with
initial values (3) matching the Fibonacci shift. Therefore, by a trivial
induction,

\[
\boxed{D_n=F_{n+1}\quad\text{for all }n\ge 1.}
\]

\subsection{25 {[}M21{]}}\label{m21-12}

Let

\[
\alpha=\frac{1+\sqrt5}{2},\qquad \beta=\frac{1-\sqrt5}{2}.
\]

Then \(F_n=\dfrac{\alpha^n-\beta^n}{\alpha-\beta}\) and
\(\alpha-\beta=\sqrt5\). Multiply by \(2^n\):

\[
2^nF_n \;=\; \frac{(2\alpha)^n-(2\beta)^n}{\alpha-\beta}
\;=\; \frac{(1+\sqrt5)^n-(1-\sqrt5)^n}{\sqrt5}.
\]

Expand each term with the binomial theorem:

\[
(1\pm\sqrt5)^n=\sum_{k=0}^{n}\binom{n}{k}(\pm\sqrt5)^k
=\sum_{k=0}^{n}\binom{n}{k}(\pm1)^k(\sqrt5)^k.
\]

Subtracting,

\[
(1+\sqrt5)^n-(1-\sqrt5)^n
=\sum_{k=0}^{n}\binom{n}{k}\bigl(1-(-1)^k\bigr)(\sqrt5)^k.
\]

For even \(k\), \(1-(-1)^k=0\); for odd \(k\), \(1-(-1)^k=2\). Hence

\[
(1+\sqrt5)^n-(1-\sqrt5)^n
=2\sum_{\substack{k=0\\ k\text{ odd}}}^{n}\binom{n}{k}(\sqrt5)^k.
\]

Divide by \(\sqrt5\):

\[
2^nF_n
=\frac{1}{\sqrt5}\cdot 2\sum_{\substack{k\text{ odd}}}\binom{n}{k}(\sqrt5)^k
=2\sum_{\substack{k\text{ odd}}}\binom{n}{k}(\sqrt5)^{\,k-1}
=2\sum_{\substack{k\text{ odd}}}\binom{n}{k}5^{(k-1)/2},
\]

which is the desired identity. \(\square\)

Equivalent reindexing

Writing \(k=2j+1\) (so \(j=0,1,\dots,\lfloor (n-1)/2\rfloor\)) gives the
neat form

\[
2^{\,n}F_n
=2\sum_{j=0}^{\lfloor (n-1)/2\rfloor}\binom{n}{2j+1}5^{\,j}.
\]

Quick sanity checks

\begin{itemize}
\tightlist
\item
  \(n=1\): \(2^1F_1=2\). RHS \(=2\binom11 5^0=2\).
\item
  \(n=3\): \(2^3F_3=8\cdot2=16\). RHS
  \(=2\bigl(\binom31 5^0+\binom33 5^1\bigr)=2(3+5)=16\).
\end{itemize}

\subsection{26 {[}M20{]}}\label{m20-24}

\begin{enumerate}
\def\labelenumi{\arabic{enumi})}
\tightlist
\item
  Start from Exercise 25's identity. From the previous exercise we may
  use the equivalent closed form
\end{enumerate}

\[
2^{\,n}\sqrt5\,F_n=(1+\sqrt5)^{\,n}-(1-\sqrt5)^{\,n}.
\]

(Indeed, expanding the right side by the binomial theorem cancels the
even terms and yields Exercise 25's sum.) 

\begin{enumerate}
\def\labelenumi{\arabic{enumi})}
\setcounter{enumi}{1}
\tightlist
\item
  Work in the ring
  \(\;\mathbb F_p[\sqrt5]=(\mathbb Z/p\mathbb Z)[x]/(x^2-5)\). Reduce
  the identity with \(n=p\) modulo \(p\). By the binomial theorem modulo
  a prime (all interior binomial coefficients are multiples of \(p\)),
\end{enumerate}

\[
(1\pm\sqrt5)^{p}\equiv 1\pm(\sqrt5)^{\,p}\pmod p.
\]

(This is the property referenced as §1.2.6--10(b).) Hence

\[
2^{\,p}\sqrt5\,F_p\equiv\bigl[1+(\sqrt5)^{\,p}\bigr]-\bigl[1-(\sqrt5)^{\,p}\bigr]
=2(\sqrt5)^{\,p}\pmod p.
\]

\begin{enumerate}
\def\labelenumi{\arabic{enumi})}
\setcounter{enumi}{2}
\tightlist
\item
  Simplify the powers of \(\sqrt5\). Because
  \((\sqrt5)^{\,p}=\sqrt5\cdot(\sqrt5)^{\,p-1}=\sqrt5\cdot 5^{(p-1)/2}\),
  the congruence becomes
\end{enumerate}

\[
2^{\,p}\sqrt5\,F_p\equiv 2\sqrt5\cdot 5^{(p-1)/2}\pmod p.
\]

\begin{enumerate}
\def\labelenumi{\arabic{enumi})}
\setcounter{enumi}{3}
\tightlist
\item
  Cancel \(\sqrt5\) and \(2\). If \(p\neq 5\), then \(5\not\equiv0\) so
  \((\sqrt5)^2=5\) is a unit in \(\mathbb F_p[\sqrt5]\); hence
  \(\sqrt5\) is a unit and can be canceled. Also \(2\) is invertible
  since \(p\) is odd. We get
\end{enumerate}

\[
2^{\,p-1}F_p\equiv 5^{(p-1)/2}\pmod p.
\]

\begin{enumerate}
\def\labelenumi{\arabic{enumi})}
\setcounter{enumi}{4}
\tightlist
\item
  Apply Fermat's little theorem. \(2^{p-1}\equiv1\pmod p\), so finally
\end{enumerate}

\[
\boxed{\,F_p \equiv 5^{(p-1)/2}\pmod p\,}.
\]

\begin{enumerate}
\def\labelenumi{\arabic{enumi})}
\setcounter{enumi}{5}
\tightlist
\item
  The case \(p=5\). Then \(F_5=5\equiv0\pmod5\) and
  \(5^{(5-1)/2}=5^2\equiv0\pmod5\), so the congruence still holds.
\end{enumerate}

Quick checks: \(p=3:\,F_3=2\equiv 5^1\equiv2\pmod3\);
\(p=7:\,F_7=13\equiv6\equiv5^3\pmod7\);
\(p=11:\,F_{11}=89\equiv1\equiv5^5\pmod{11}\).

\subsection{27 {[}M20{]}}\label{m20-25}

Facts we'll use

\begin{enumerate}
\def\labelenumi{\arabic{enumi}.}
\tightlist
\item
  Cassini's identity. For all integers \(n\),
\end{enumerate}

\[
F_{n+1}F_{n-1}-F_n^2 = (-1)^n.
\]

This appears in the section text (with a quick proof by induction and a
matrix proof). 

\begin{enumerate}
\def\labelenumi{\arabic{enumi}.}
\setcounter{enumi}{1}
\tightlist
\item
  Result from Exercise 26. For an odd prime \(p\),
\end{enumerate}

\[
F_p \equiv 5^{(p-1)/2}\pmod p.
\]

(Exercise statement.) 

\begin{enumerate}
\def\labelenumi{\arabic{enumi}.}
\setcounter{enumi}{2}
\tightlist
\item
  Euler's criterion (standard number theory). For an odd prime \(p\) and
  integer \(a\) not divisible by \(p\),
\end{enumerate}

\[
a^{(p-1)/2}\equiv \left(\frac{a}{p}\right)\pmod p,
\]

where \(\left(\frac{a}{p}\right)\) is the Legendre symbol and equals
\(\pm 1\).

Combining (2) with Euler's criterion for \(a=5\) gives

\[
F_p \equiv 5^{(p-1)/2}\equiv \left(\frac{5}{p}\right)\in\{+1,-1\}\pmod p.
\]

So for \(p\neq 5\), we have

\[
F_p^2 \equiv 1 \pmod p. \tag{★}
\]

Step 1: \(p\) divides \(F_{p-1}F_{p+1}\)

Apply Cassini's identity with \(n=p\) (note \(p\) odd):

\[
F_{p-1}F_{p+1}-F_p^2 = (-1)^p=-1.
\]

Reduce modulo \(p\) and use \((★)\):

\[
F_{p-1}F_{p+1}\equiv F_p^2-1 \equiv 0 \pmod p.
\]

Hence \(p\mid F_{p-1}F_{p+1}\); i.e., at least one of
\(F_{p-1},F_{p+1}\) is divisible by \(p\). (This is exactly the route
hinted in the book's answers.) 

Step 2: They cannot both be divisible by \(p\)

Suppose, for contradiction, that \(p\mid F_{p-1}\) and
\(p\mid F_{p+1}\). Then from the Fibonacci recurrence,

\[
F_{p+1}=F_p+F_{p-1}\equiv 0 \pmod p \quad\text{and}\quad F_{p-1}\equiv 0 \pmod p
\]

would imply \(F_p\equiv 0\pmod p\). But we already have
\(F_p\equiv \pm 1\pmod p\) for \(p\neq 5\), a contradiction. Therefore
exactly one of \(F_{p-1},F_{p+1}\) is a multiple of \(p\). 

Edge cases

\begin{itemize}
\tightlist
\item
  \(p=2\): \(F_{1}=1\) and \(F_{3}=2\). Here \(F_{p+1}=2\) is divisible
  by \(2\) while \(F_{p-1}\) is not---so the statement is true. (The
  book notes this explicitly.) 
\item
  \(p=5\): excluded in the exercise; indeed \(F_4=3\), \(F_6=8\) are not
  divisible by 5, while \(F_5=5\) is. So the claimed ``\(p\mid F_{p-1}\)
  or \(F_{p+1}\)'' would fail at \(p=5\).
\end{itemize}

(Optional) Which one is divisible?

From \(F_p\equiv \left(\frac{5}{p}\right)\) we get \(F_p\equiv 1\) if
\(5\) is a quadratic residue mod \(p\) and \(F_p\equiv -1\) otherwise.
One can refine the above to show:

\begin{itemize}
\tightlist
\item
  If \(\left(\frac{5}{p}\right)=+1\), then \(p\mid F_{p-1}\);
\item
  If \(\left(\frac{5}{p}\right)=-1\), then \(p\mid F_{p+1}\). This is a
  standard refinement (not required by the exercise), and it dovetails
  with the book's brief answer that only one factor can vanish mod
  \(p\). 
\end{itemize}

\subsection{28 {[}M21{]}}\label{m21-13}

Method A (from Binet's formula)

Recall the closed form

\[
F_n=\frac{\varphi^n-\hat\varphi^{\,n}}{\sqrt5},
\quad\text{where}\quad
\hat\varphi=1-\varphi=\frac{1-\sqrt5}{2}.
\]

(This \(\hat\varphi\) is the other root of \(x^2=x+1\).)  

Then

\[
\begin{aligned}
F_{n+1}-\varphi F_n
&=\frac{\varphi^{n+1}-\hat\varphi^{\,n+1}}{\sqrt5}
      -\varphi\cdot\frac{\varphi^{n}-\hat\varphi^{\,n}}{\sqrt5}\\[2mm]
&=\frac{\varphi^{n+1}-\hat\varphi^{\,n+1}
          -\varphi^{n+1}+\varphi\hat\varphi^{\,n}}{\sqrt5}\\
&=\frac{\hat\varphi^{\,n}\,(\varphi-\hat\varphi)}{\sqrt5}.
\end{aligned}
\]

But \(\varphi-\hat\varphi=\sqrt5\), so the factor \(\sqrt5\) cancels and
we get

\[
\boxed{\,F_{n+1}-\varphi F_n=\hat\varphi^{\,n}=(1-\varphi)^n\, }.
\]

(Here \(\hat\varphi=1-\varphi\).)  

Method B (little recurrence trick)

Let \(a_n:=F_{n+1}-\varphi F_n\). Using \(F_{n+2}=F_{n+1}+F_n\) and
\(\varphi^2=\varphi+1\),

\[
\begin{aligned}
a_{n+1}
&=F_{n+2}-\varphi F_{n+1}
=(F_{n+1}+F_n)-\varphi F_{n+1} \\
&=F_n-(\varphi-1)F_{n+1}
=F_n-(\varphi-1)(a_n+\varphi F_n)\\
&=F_n-(\varphi-1)a_n-(\varphi^2-\varphi)F_n
=-(\varphi-1)a_n\\
&=(1-\varphi)a_n=\hat\varphi\,a_n.
\end{aligned}
\]

Since \(a_0=F_1-\varphi F_0=1\), this first-order recurrence gives
\(a_n=\hat\varphi^{\,n}\). Thus again

\[
\boxed{F_{n+1}-\varphi F_n=(1-\varphi)^n.}
\]

Quick check: \(n=0\): \(F_1-\varphi F_0=1=(1-\varphi)^0\). \(n=1\):
\(F_2-\varphi F_1=1-\varphi=1-\varphi\).

\subsection{29 {[}M23{]}}\label{m23-15}

Let \((F_n)_{n\ge 0}\) be the Fibonacci numbers with
\(F_0=0,\;F_1=1,\;F_{n}=F_{n-1}+F_{n-2}\). Following É. Lucas, define
the Fibonomial coefficients

\[
\binom{n}{k}_{\!F}
\;:=\;
\frac{F_n F_{n-1}\cdots F_{n-k+1}}{F_k F_{k-1}\cdots F_{1}}
\quad(0\le k\le n),
\]

with the usual conventions \(\binom{n}{0}_{\!F}=1\) and
\(\binom{n}{k}_{\!F}=0\) for \(k<0\) or \(k>n\).

We'll also use the standard Fibonacci addition identity

\[
F_{r+s}=F_{r-1}F_s+F_rF_{s+1}\qquad(r,s\ge 0),
\]

which is featured in classical Fibonacci references and in §1.2.8.
({[}R-Knott{]}{[}3{]})

\begin{enumerate}
\def\labelenumi{(\alph{enumi})}
\tightlist
\item
  First rows of the Fibonomial triangle
\end{enumerate}

Compute directly from the definition (the products are short, so it's
just careful arithmetic):

\begin{itemize}
\tightlist
\item
  \(n=0:\) \(\;1\)
\item
  \(n=1:\) \(\;1,\,1\)
\item
  \(n=2:\) \(\;1,\,1,\,1\)
\item
  \(n=3:\) \(\;1,\,\dfrac{F_3}{F_1}=2,\,\dfrac{F_3F_2}{F_2F_1}=1,\,1\) →
  \(1,2,1\)
\item
  \(n=4:\)
  \(\;1,\,\dfrac{3}{1}=3,\,\dfrac{3\cdot2}{1\cdot1}=6,\,\dfrac{3\cdot2\cdot1}{1\cdot1\cdot1}=3,\,1\)
  → \(1,3,6,3,1\)
\item
  \(n=5:\) \(\;1,5,15,15,5,1\)
\item
  \(n=6:\) \(\;1,8,40,60,40,8,1\)
\end{itemize}

So the first seven rows are

\[
\begin{array}{l}
1\\
1\;1\\
1\;1\;1\\
1\;2\;1\\
1\;3\;6\;3\;1\\
1\;5\;15\;15\;5\;1\\
1\;8\;40\;60\;40\;8\;1\;.
\end{array}
\]

\begin{enumerate}
\def\labelenumi{(\alph{enumi})}
\setcounter{enumi}{1}
\tightlist
\item
  Pascal-type recurrence and integrality
\end{enumerate}

We claim the Fibonomials satisfy the Lucas--Pascal rule

\[
\boxed{\displaystyle
\binom{n}{k}_{\!F}
= F_{k-1}\binom{n-1}{k}_{\!F}
+ F_{\,n-k+1}\binom{n-1}{k-1}_{\!F}
}\qquad(1\le k\le n-1).
\]

Derivation

Start from the definition and the two useful ``ratio'' forms:

\[
\binom{n}{k}_{\!F}
=\frac{F_n}{F_k}\binom{n-1}{k-1}_{\!F}
=\frac{F_n}{F_{n-k}}\binom{n-1}{k}_{\!F}.
\]

If we assume the boxed recurrence, divide both sides by
\(\binom{n-1}{k-1}_{\!F}\) to test equality:

\[
\frac{\binom{n}{k}_{\!F}}{\binom{n-1}{k-1}_{\!F}}
\;=\;
F_{k-1}\frac{\binom{n-1}{k}_{\!F}}{\binom{n-1}{k-1}_{\!F}} + F_{n-k+1}.
\]

Compute the inner ratio directly from the product definition:

\[
\frac{\binom{n-1}{k}_{\!F}}{\binom{n-1}{k-1}_{\!F}}=\frac{F_{n-k}}{F_k}.
\]

So the claimed recurrence is equivalent to

\[
\frac{F_n}{F_k}=F_{k-1}\frac{F_{n-k}}{F_k}+F_{n-k+1},
\]

i.e.

\[
F_n=F_{k-1}F_{n-k}+F_kF_{n-k+1},
\]

which is exactly the Fibonacci addition identity with \(r=k,\;s=n-k\).
Hence the recurrence holds. ({[}R-Knott{]}{[}3{]})

Consequence: \(\binom{n}{k}_{\!F}\) is an integer

The initial row \(n=0\) is integral, and the boxed recurrence expresses
each interior entry as an integer linear combination of entries one row
above. By induction on \(n\), every \(\binom{n}{k}_{\!F}\) is an
integer.

\subsection{30 {[}M38{]}}\label{m38}

Let \(F_n\) be Fibonacci numbers
\((F_0=0,\ F_1=1,\ F_{n+1}=F_n+F_{n-1})\). Define the Fibonacci
factorial \(F!n:=\prod_{i=1}^{n}F_i\) (with \(F!0=1\)), and the
fibonomial coefficients

\[
\binom{m}{k}_F \;:=\; \frac{F!m}{F!k\;F!(m-k)}\qquad(0\le k\le m;\ \text{else }0).
\]

For a fixed integer \(m\ge1\) and any integer \(n\), consider the
alternating sums

\[
S_r(m,n)\ :=\ \sum_{k=0}^{m}(-1)^{\lfloor (m-k)/2\rfloor}\binom{m}{k}_F\,W_r(k),
\]

where the weights \(W_r(k)\) are Fibonacci terms that are affine in the
index (no products of independent \(k\)-dependent pieces are needed).
Exercise 30 asks you to prove---by induction on \(m\)---that four
concrete choices of \(W_r\) all give \(S_r(m,n)=0\). One convenient way
to package the four requested identities is:

\begin{itemize}
\item
  \begin{enumerate}
  \def\labelenumi{(\alph{enumi})}
  \tightlist
  \item
    \(\displaystyle \sum_{k=0}^m (-1)^{\lfloor (m-k)/2\rfloor}\binom{m}{k}_F\,F_k \;=\;0.\)
  \end{enumerate}
\item
  \begin{enumerate}
  \def\labelenumi{(\alph{enumi})}
  \setcounter{enumi}{1}
  \tightlist
  \item
    \(\displaystyle \sum_{k=0}^m (-1)^{\lfloor (m-k)/2\rfloor}\binom{m}{k}_F\,(-1)^kF_{m-k} \;=\;0.\)
  \end{enumerate}
\item
  \begin{enumerate}
  \def\labelenumi{(\alph{enumi})}
  \setcounter{enumi}{2}
  \tightlist
  \item
    Using \((a)\) and \((b)\) together with \$
    (-1)\^{}kF\_\{m-k\}=F\_\{k-1\}F\_m-F\_kF\_\{m-1\}\$, deduce
  \end{enumerate}

  \[
  \sum_{k=0}^m (-1)^{\lfloor (m-k)/2\rfloor}\binom{m}{k}_F\,F_{k-1}\;=\;0.
  \]
\item
  \begin{enumerate}
  \def\labelenumi{(\alph{enumi})}
  \setcounter{enumi}{3}
  \tightlist
  \item
    Using the Fibonacci addition formula
    \(F_{n+k}=F_{k-1}F_n+F_kF_{n+1}\) and \((a),(c)\), deduce
  \end{enumerate}

  \[
  \sum_{k=0}^m (-1)^{\lfloor (m-k)/2\rfloor}\binom{m}{k}_F\,F_{n+k}\;=\;0\qquad(\text{for every }n\in\mathbb Z).
  \]
\end{itemize}

\begin{enumerate}
\def\labelenumi{\arabic{enumi})}
\tightlist
\item
  Two basic fibonomial recurrences
\end{enumerate}

From the definition,

\[
\binom{m}{k}_F=\frac{F_m}{F_k}\binom{m-1}{k-1}_F
\quad\text{and}\quad
\binom{m}{k}_F=\frac{F_m}{F_{m-k}}\binom{m-1}{k}_F
\qquad(1\le k\le m).
\tag{R1,R2}
\]

Equivalently,

\[
F_k\binom{m}{k}_F=F_m\binom{m-1}{k-1}_F,\qquad
F_{m-k}\binom{m}{k}_F=F_m\binom{m-1}{k}_F.
\]

(These are the exact fibonomial analogues of Pascal's triangle rules and
are listed in standard fibonomial references; see also the ``Fibonomial
formulae'' page, which records these same identities.)

Two Fibonacci identities we'll also use are

\[
F_{n+k}=F_{k-1}F_n+F_kF_{n+1},
\qquad
F_mF_{k-1}-F_{m-1}F_k=(-1)^kF_{m-k}.
\tag{F+,\ Cassini\ type}
\]

Both are standard and explicitly invoked in Knuth's answer for this
exercise.

\begin{enumerate}
\def\labelenumi{\arabic{enumi})}
\setcounter{enumi}{1}
\tightlist
\item
  The induction
\end{enumerate}

We show all claims for a given \(m\) assuming they hold for \(m-1\).

Base \(m=1\)

Directly (each sum has just \(k=0,1\)), the signs are
\((-1)^{\lfloor(1-k)/2\rfloor}=+1\) for both terms. Using
\(\binom{1}{0}_F=\binom{1}{1}_F=1\) and \(F_0=0,F_1=1\), each of
(a)--(d) evaluates to \(0\). Thus the base holds.

Step (a)

Multiply the summand by \(F_k/F_k\) and apply (R1):

\[
\sum_k (-1)^{\lfloor (m-k)/2\rfloor}\binom{m}{k}_F\,F_k
\;=\;
F_m\sum_{k=1}^{m} (-1)^{\lfloor (m-k)/2\rfloor}\binom{m-1}{k-1}_F.
\]

Reindex \(j=k-1\) (so \(j=0,\dots,m-1\)) and note

\[
(-1)^{\lfloor (m-k)/2\rfloor}=(-1)^{\lfloor (m-1-j)/2\rfloor}.
\]

By the induction hypothesis (the same identity with \(m\to m-1\)), the
reindexed sum is \(0\). Hence (a) holds.

Step (b)

Now use (R2) and keep the explicit \((-1)^k\):

\[
\sum_k (-1)^{\lfloor (m-k)/2\rfloor}\binom{m}{k}_F\,(-1)^kF_{m-k}
\;=\;
F_m \sum_{k=0}^{m-1} (-1)^{\lfloor (m-k)/2\rfloor}(-1)^k \binom{m-1}{k}_F.
\]

A simple parity check shows

\[
(-1)^{\lfloor (m-k)/2\rfloor}(-1)^k
= (-1)^m\,(-1)^{\lfloor (m-1-k)/2\rfloor}.
\]

Therefore

\[
\text{(b)}\;=\;(-1)^m F_m \sum_{k=0}^{m-1} (-1)^{\lfloor (m-1-k)/2\rfloor}\binom{m-1}{k}_F
\;=\;0
\]

by the induction hypothesis.

Step (c)

Use the Cassini-type identity:

\[
(-1)^kF_{m-k}=F_{k-1}F_m - F_kF_{m-1}.
\]

Sum both sides with the same alternating weight and \(\binom{m}{k}_F\).
The left-hand sum is \(0\) by (b). The right-hand side becomes

\[
F_m\sum_k (-1)^{\lfloor (m-k)/2\rfloor}\binom{m}{k}_F F_{k-1}
\;-\;
F_{m-1}\sum_k (-1)^{\lfloor (m-k)/2\rfloor}\binom{m}{k}_F F_k.
\]

The second sum vanishes by (a). Since \(F_m\ne0\), the first must vanish
too, proving (c).

Step (d)

Use the addition formula \(F_{n+k}=F_{k-1}F_n+F_kF_{n+1}\) inside the
sum:

\[
\sum_k (-1)^{\lfloor (m-k)/2\rfloor}\binom{m}{k}_F F_{n+k}
\;=\;
F_n\underbrace{\sum_k (-1)^{\lfloor (m-k)/2\rfloor}\binom{m}{k}_F F_{k-1}}_{=\,0\text{ by (c)}}
\;+\;
F_{n+1}\underbrace{\sum_k (-1)^{\lfloor (m-k)/2\rfloor}\binom{m}{k}_F F_k}_{=\,0\text{ by (a)}}.
\]

Hence (d) holds.

This completes the induction on \(m\), establishing all four identities.

Quick check (small \(m\))

For \(m=2\), \(\binom{2}{0}_F=\binom{2}{2}_F=1,\ \binom{2}{1}_F=1\).

\begin{itemize}
\tightlist
\item
  \begin{enumerate}
  \def\labelenumi{(\alph{enumi})}
  \tightlist
  \item
    gives \$ (+1)\cdot1\cdot F\_0 + (-1)\cdot1\cdot F\_1 +
    (+1)\cdot1\cdot F\_2 = 0-1+1=0\$.
  \end{enumerate}
\item
  \begin{enumerate}
  \def\labelenumi{(\alph{enumi})}
  \setcounter{enumi}{2}
  \tightlist
  \item
    gives \$ (+1)\cdot1\cdot F\_\{-1\} + (-1)\cdot1\cdot F\_0 +
    (+1)\cdot1\cdot F\_1 = 1+0+1\$ with the signs \((+,-,+)\) replaced
    by \((+,-,+)\) times the fibonomial coefficients---but now
    \((-1)^{\lfloor(2-k)/2\rfloor}=(+,-,+)\) so the sum is \(1-0-1=0\).
    Similar spot-checks for (b) and (d) also vanish, consistent with the
    proof.
  \end{enumerate}
\end{itemize}

\subsection{31 {[}M20{]}}\label{m20-26}

We'll prove the slightly more general identity

\[
F_{n-1}-F_n\,\varphi=( -\varphi)^n \qquad (n\ge 1), \tag{★}
\]

from which the required statements about fractional parts follow
immediately.

Step 1: Establish the recurrence for the ``error term''

Define \(e_n:=F_{n-1}-F_n\,\varphi\). Using \(F_{n+1}=F_n+F_{n-1}\) and
the defining equation \(\varphi^2+\varphi=1\) (equivalently
\(\varphi=\Phi^{-1}\)), compute

\[
\begin{aligned}
e_{n+1}
&=F_n - F_{n+1}\varphi
=F_n-(F_n+F_{n-1})\varphi\\
&=(F_n - F_{n-1}\varphi) - \varphi F_n
= -\varphi\,(F_{n-1}-F_n\varphi)
= -\varphi\, e_n.
\end{aligned}
\]

The base case is \(e_1=F_0-F_1\varphi=0-\varphi=-\varphi\). Hence by
induction,

\[
e_n=(-\varphi)^n \quad\text{for all } n\ge 1,
\]

which is exactly (★).

(An equally short alternative uses the standard identity
\(\Phi^n=F_n\Phi+F_{n-1}\); divide by \(\Phi^n\) and rewrite with
\(\varphi=\Phi^{-1}\) to get (★).)

Step 2: Read off the fractional parts

From (★) we have

\[
F_n\varphi \;=\; F_{n-1} - (-\varphi)^n.
\]

Because \(0<\varphi<1\), the tiny correction \((- \varphi)^n\) has
absolute value \(<1\). Therefore:

\begin{itemize}
\item
  If \(n\) is odd, then \((- \varphi)^n=-\varphi^{n}\), so

  \[
  F_n\varphi = F_{n-1} + \varphi^{n},
  \]

  whose fractional part is \(\boxed{\varphi^{n}}\).
\item
  If \(n\) is even, then \((- \varphi)^n=+\varphi^{n}\), so

  \[
  F_n\varphi = (F_{n-1}-1) + (1-\varphi^{n}),
  \]

  whose fractional part is \(\boxed{1-\varphi^{n}}\).
\end{itemize}

Re-indexing those two cases as \(n=2k\) and \(n=2k+1\) gives exactly

\[
\{F_{2k}\varphi\}=1-\varphi^{2k},\qquad
\{F_{2k+1}\varphi\}=\varphi^{2k+1},
\]

as required.

Quick numeric check

\[
\begin{array}{r|r|l|l}
n&F_n&F_n\varphi&\text{fractional part} \\
\hline
1&1&0.6180&\varphi^1\approx0.6180\\
2&1&0.6180&1-\varphi^2\approx0.6180\\
3&2&1.2361&\varphi^3\approx0.2361\\
4&3&1.8541&1-\varphi^4\approx0.8541
\end{array}
\]

Side corollaries (often useful)

From the same identity (★) you immediately get tidy floor/ceiling
formulas:

\[
\big\lfloor F_n\varphi\big\rfloor=
\begin{cases}
F_{n-1},& n\text{ odd},\\
F_{n-1}-1,& n\text{ even},
\end{cases}
\qquad
\big\lceil F_n\varphi\big\rceil=
\begin{cases}
F_{n-1}+1,& n\text{ odd},\\
F_{n-1},& n\text{ even}.
\end{cases}
\]

\subsection{32 {[}M24{]}}\label{m24-10}

We'll use two standard Fibonacci facts.

\begin{enumerate}
\def\labelenumi{(\Alph{enumi})}
\tightlist
\item
  Matrix form. Let \(M=\begin{pmatrix}1&1\\[2pt]1&0\end{pmatrix}\). Then
\end{enumerate}

\[
M^k=\begin{pmatrix}F_{k+1}&F_k\\ F_k&F_{k-1}\end{pmatrix}
\quad\text{for all }k\ge0. 
\]

\begin{enumerate}
\def\labelenumi{(\Alph{enumi})}
\setcounter{enumi}{1}
\tightlist
\item
  Cassini's identity.
\end{enumerate}

\[
F_{n+1}F_{n-1}-F_n^{\,2}=(-1)^n.
\]

In particular, reducing modulo \(F_n\) gives

\[
F_{n-1}^2 \equiv (-1)^n \pmod{F_n}.
\]

We also use the (equally standard) addition formula

\[
F_{a+b}=F_aF_{b+1}+F_{a-1}F_b,
\]

from which the identity

\[
F_{n-1}F_r \equiv (-1)^{\,r+1}F_{n-r}\pmod{F_n}
\tag{★}
\]

follows by a one-line rearrangement (see derivation just below).

Step 1: Reduce \(M^n\) modulo \(F_n\).

From (A),

\[
M^n=\begin{pmatrix}F_{n+1}&F_n\\ F_n&F_{n-1}\end{pmatrix}
\equiv
\begin{pmatrix}F_{n+1}\bmod F_n&0\\ 0&F_{n-1}\end{pmatrix}
\equiv
\begin{pmatrix}F_{n-1}&0\\ 0&F_{n-1}\end{pmatrix}
=F_{n-1}I
\pmod{F_n},
\]

because \(F_{n+1}\equiv F_{n-1}\ (\mathrm{mod}\ F_n)\).

Hence, for any \(m,r\),

\[
M^{mn+r}=(M^n)^mM^r \equiv (F_{n-1})^m\,M^r \pmod{F_n}.
\]

Taking the \((1,2)\) entry on both sides yields the key congruence

\[
F_{mn+r} \equiv (F_{n-1})^m\,F_r \pmod{F_n}.
\tag{1}
\]

Step 2: Interpret the factor \((F_{n-1})^m\) modulo \(F_n\).

From Cassini (B),

\[
(F_{n-1})^2 \equiv (-1)^n \pmod{F_n}.
\]

Therefore

\[
(F_{n-1})^m \equiv
\begin{cases}
(-1)^{\,n m/2}, & m\ \text{even},\\[4pt]
(-1)^{\,n (m-1)/2}\,F_{n-1}, & m\ \text{odd}.
\end{cases}
\tag{2}
\]

Plug (2) into (1):

\begin{itemize}
\item
  If \(m\) is even, \(F_{mn+r}\equiv (\pm 1)\,F_r \pmod{F_n}\).
\item
  If \(m\) is odd, we get an extra \(F_{n-1}\) factor:

  \[
  F_{mn+r}\equiv (\pm) \,F_{n-1}F_r \pmod{F_n}.
  \]

  Use (★) to replace \(F_{n-1}F_r\) by \((\pm)F_{n-r}\) modulo \(F_n\),
  so

  \[
  F_{mn+r}\equiv (\pm)\,F_{n-r} \pmod{F_n}.
  \]
\end{itemize}

In all cases the residue is either \(\pm F_r\) or \(\pm F_{n-r}\), as
required. (A sharper, fully explicit sign pattern in terms of
\(m\bmod 4\), \(n\), and \(r\) is given here for completeness:
\(\displaystyle F_{mn+r}\equiv
\begin{cases}
F_r, & m\equiv0\ (\mathrm{mod}\ 4);\\
(-1)^{r+1}F_{n-r}, & m\equiv1\ (\mathrm{mod}\ 4);\\
(-1)^n F_r, & m\equiv2\ (\mathrm{mod}\ 4);\\
(-1)^{r+1+n}F_{n-r}, & m\equiv3\ (\mathrm{mod}\ 4),
\end{cases}\ \ \ (\mathrm{mod}\ F_n)\).) ({[}ProofWiki{]}{[}2{]})

Derivation of (★) (one line)

Start from the well-known identity (a variant of d'Ocagne's):

\[
F_{n}F_{r+1}-F_{n+1}F_r = (-1)^r F_{n-r}.
\]

Reduce \(\bmod F_n\): the left becomes
\(-F_{n+1}F_r \equiv -F_{n-1}F_r\), hence

\[
F_{n-1}F_r \equiv (-1)^{r+1}F_{n-r}\pmod{F_n}.
\]

Check on a concrete example

Take \(n=7\), \(m=3\), \(r=2\). Then \(mn+r=23\). We predict

\[
F_{23}\equiv \pm F_{7-2}=\pm F_5= \pm 5 \pmod{F_7},
\]

since \(m\) is odd. Indeed \(F_{23}=28657\) and \(F_7=13\);
\(28657\bmod 13=5\).

\subsection{33 {[}HM24{]}}\label{hm24}

Step 1: Pick the right \(z\)

Take

\[
z \;=\; \frac{\pi}{2} \;+\; i\ln\varphi\,,\qquad \varphi=\frac{1+\sqrt5}{2}.
\]

Then

\[
\cos z=\cos\!\Big(\tfrac{\pi}{2}+i\ln\varphi\Big)
=\cos\tfrac{\pi}{2}\cosh(\ln\varphi)-i\sin\tfrac{\pi}{2}\sinh(\ln\varphi)
=-i\,\sinh(\ln\varphi).
\]

Since \(\sinh(\ln\varphi)=\tfrac{1}{2}(\varphi-\varphi^{-1})=\tfrac12\),
we indeed get \(\cos z=-\,\tfrac{i}{2}\) as desired. Also

\[
\sin z=\sin\!\Big(\tfrac{\pi}{2}+i\ln\varphi\Big)
=\sin\tfrac{\pi}{2}\cosh(\ln\varphi)=\cosh(\ln\varphi)=\frac{\varphi+\varphi^{-1}}{2}=\frac{\sqrt5}{2}.
\]

Step 2: Define the candidate sequence and verify the recurrence

Define

\[
S_n \;=\; i^{\,n-1}\,\frac{\sin(nz)}{\sin z}\,.
\]

Use the standard identity

\[
\sin((n+1)z)+\sin((n-1)z)=2\cos z\,\sin(nz)
\]

(valid for all complex \(z\)). Multiply both sides by \(i^{\,n}/\sin z\)
and use \(\cos z=-\,i/2\):

\[
i^{\,n}\frac{\sin((n+1)z)}{\sin z}+i^{\,n}\frac{\sin((n-1)z)}{\sin z}
=2\Big(-\frac{i}{2}\Big)\,i^{\,n}\frac{\sin(nz)}{\sin z}
= -\,i^{\,n+1}\frac{\sin(nz)}{\sin z}.
\]

Divide through by \(i\) and rewrite in terms of \(S_n\):

\[
S_{n+1}=S_n+S_{n-1}.
\]

Initial values: \(S_0=0\) (since \(\sin 0=0\)) and \(S_1=1\) (since
\(\sin z/\sin z=1\)). Therefore \(\{S_n\}\) satisfies the Fibonacci
recurrence with Fibonacci initial conditions, hence \(S_n=F_n\) for all
\(n\). This proves

\[
F_n \;=\; i^{\,n-1}\,\frac{\sin(nz)}{\sin z}\qquad\text{for } \cos z=-\frac{i}{2}.
\]

(Exactly the approach Knuth's hint suggests.)

Step 3 (optional): Recover Binet's formula

With \(z=\tfrac{\pi}{2}+i\ln\varphi\) we have
\(e^{iz}=e^{i\pi/2}e^{-\ln\varphi}=i/\varphi\). Hence

\[
\sin(nz)=\frac{e^{inz}-e^{-inz}}{2i}
=\frac{(i/\varphi)^{n}-(-i\varphi)^{n}}{2i}.
\]

Using \(\sin z=\sqrt5/2\) from Step 1,

\[
F_n
= i^{\,n-1}\frac{\sin(nz)}{\sin z}
= \frac{i^{\,n-1}}{\sqrt5/2}\cdot\frac{(i/\varphi)^{n}-(-i\varphi)^{n}}{2i}
= \frac{\varphi^{\,n}-(-\varphi^{-1})^{n}}{\sqrt5}
= \frac{\varphi^{\,n}-(\varphi')^{\,n}}{\sqrt5},
\]

where \(\varphi'=(1-\sqrt5)/2=-1/\varphi\). That is Binet's formula.

\subsection{34 {[}M24{]}}\label{m24-11}

Let \(k \gg m\) mean \(k\ge m+2\). Prove that every positive integer
\(n\) has a unique representation

\[
n \;=\; F_{k_1}+F_{k_2}+\cdots+F_{k_r},
\]

with \(k_1 \gg k_2 \gg \cdots \gg k_r \gg 0\) (i.e., the indices are
strictly decreasing and no two are consecutive; the smallest allowed
index is \(2\)). This is the Fibonacci number system. 

This statement is the classical Zeckendorf theorem specialized to
Knuth's indexing \(F_0=0,\,F_1=1,\,F_2=1,\,F_{n}=F_{n-1}+F_{n-2}\).

Existence (constructive --- the greedy algorithm)

Greedy step. Given \(n\ge1\), pick the largest Fibonacci number
\(F_k\le n\) (with \(k\ge2\)), and set \(r=n-F_k\).

Because \(F_{k+1}=F_k+F_{k-1}\) and \(F_{k+1}>n\ge F_k\), we have

\[
0\le r<n-F_k< F_{k+1}-F_k=F_{k-1}.
\]

So the remainder \(r\) is strictly less than \(F_{k-1}\). Therefore when
we recurse on \(r\), the next chosen Fibonacci index is at most \(k-2\).
Hence the indices differ by at least 2 (the ``no adjacent indices''
condition), exactly as required.

Termination. The remainders strictly decrease and are nonnegative, so
the process stops after finitely many steps at \(r=0\).

Thus the greedy procedure produces a representation

\[
n=F_{k_1}+F_{k_2}+\cdots+F_{k_r}
\quad\text{with}\quad
k_1>k_2>\cdots>k_r\ge2,\; k_i-k_{i+1}\ge2.
\]

So existence is proved.

A helpful bound (maximal sum without adjacency)

Define \(S(m)\) to be the maximum sum obtainable using
\(\{F_2,\ldots,F_m\}\) with no adjacent indices. Then

\[
\boxed{\,S(m)=F_{m+1}-1\quad(m\ge2)\,}.
\]

Proof by induction.

\begin{itemize}
\item
  Base \(m=2\): the only allowed choice is \(F_2=1\); indeed
  \(F_3-1=2-1=1\).
\item
  Step: Any optimal set on \(\{2,\dots,m\}\) either omits \(F_m\) (sum
  \(S(m-1)\)) or uses \(F_m\) and must omit \(F_{m-1}\) (sum
  \(F_m+S(m-2)\)). Thus

  \[
  S(m)=\max\{S(m-1),\,F_m+S(m-2)\}.
  \]

  By the induction hypothesis \(S(m-1)=F_m-1\) and \(S(m-2)=F_{m-1}-1\),
  hence

  \[
  S(m)=\max\{F_m-1,\,F_m+F_{m-1}-1\}=F_{m+1}-1.\;\square
  \]
\end{itemize}

Intuitively, the (unique) maximizer is \(F_m+F_{m-2}+F_{m-4}+\cdots\),
whose sum is \(F_{m+1}-1\).

Uniqueness

Suppose

\[
n=\sum_{i=1}^r F_{a_i}=\sum_{j=1}^s F_{b_j}
\]

are two legal (no-adjacent) representations, with \(a_1>\cdots>a_r\) and
\(b_1>\cdots>b_s\).

Let \(t=\max\{a_1,b_1\}\). If, without loss of generality,
\(a_1=t>b_1\), then the right-hand sum uses only indices \(\le t-1\) and
is legal; by the bound above it is at most \(S(t-1)=F_t-1\). The
left-hand sum is at least \(F_t\). Contradiction:

\[
F_t \le n \le F_t-1.
\]

Therefore \(a_1=b_1\). Cancel \(F_t\) from both sides and repeat the
argument on the remainders. By induction, all indices match and the two
sums are identical. Hence the representation is unique.

Algorithm (explicit) and complexity

A direct constructive algorithm follows the proof:

\begin{enumerate}
\def\labelenumi{\arabic{enumi}.}
\item
  Generate Fibonacci numbers \(F_2, F_3, \dots\) up to the largest
  \(F_k\le n\).
\item
  Set \(r\leftarrow n\).
\item
  While \(r>0\):

  \begin{itemize}
  \tightlist
  \item
    Choose the largest \(F_j\le r\) (necessarily \(j\le k\)).
  \item
    Output \(F_j\), set \(r\leftarrow r-F_j\), and set
    \(k\leftarrow j-2\) (skip the adjacent index).
  \end{itemize}
\end{enumerate}

This produces the unique expansion with no adjacent indices. Because
consecutive chosen indices drop by at least 2 and \(F_k\) grows like
\(\varphi^k/\sqrt5\), the number of chosen terms (and the number of loop
iterations) is \(O(\log n)\).

Small examples

\begin{itemize}
\tightlist
\item
  \(17 = 13 + 3 + 1 = F_7 + F_4 + F_2\) (indices \(7,4,2\): gaps
  \(\ge2\)).
\item
  \(100 = 89 + 8 + 3 = F_{11} + F_6 + F_4\).
\item
  \(1 = F_2\) (note \(F_1=1\) exists but \(k\ge2\) is required, so we
  use \(F_2\); this avoids the non-uniqueness of ``two ones'').
\end{itemize}

\subsection{35 {[}M24{]}}\label{m24-12}

Here is Exercise 1.2.8--35 (paraphrased): write real numbers in ``base
ϕ'' with digits 0/1, e.g.~\((100.1)_\varphi=\varphi^2+\varphi^{-1}\).
(i) Show \(1\) has infinitely many such representations (examples:
\((.11)_\varphi\) and \((.010101\ldots)_\varphi\)). (ii) If we forbid
adjacent 1s and also forbid representations that end with the infinite
tail \(010101\ldots\), prove every nonnegative real has a unique base-ϕ
representation. (iii) Describe which of these canonical representations
are integers.

Base-ϕ facts we'll use

\[
\varphi^2=\varphi+1\quad\Longleftrightarrow\quad 
\varphi^{j}=\varphi^{j-1}+\varphi^{j-2}\ \ (\forall j\in\mathbb Z).
\tag{A}
\]

Consequences for digit strings (writing the most significant position on
the left, exponents decreasing by 1 as we move right):

\begin{itemize}
\tightlist
\item
  Carry/borrow rule: anywhere you see the block \texttt{11} in adjacent
  positions \(j-1,j-2\), you may replace it by \texttt{100} at position
  \(j\) (and conversely \texttt{100} ↔ \texttt{011}). This preserves
  value by (A).
\item
  Tail identity:
\end{itemize}

\[
\sum_{m\ge 0}\varphi^{-(2m+1)}
=\frac{\varphi^{-1}}{1-\varphi^{-2}}
=1,
\]

since
\(1-\varphi^{-2}=(\varphi^2-1)/\varphi^2=\varphi/\varphi^2=1/\varphi\).
Thus \(1=(.010101\ldots)_\varphi\) and also \(1=(.11)_\varphi\) because
\(\varphi^{-1}+\varphi^{-2}=(\varphi+1)/\varphi^2=\varphi^2/\varphi^2=1\).

\begin{enumerate}
\def\labelenumi{(\roman{enumi})}
\tightlist
\item
  Infinitely many representations of \(1\)
\end{enumerate}

Two are already given: \((.11)_\varphi\) and
\((.010101\ldots)_\varphi\). To see there are infinitely many distinct
ones, start from the finite representation \(1=(1.000\ldots)_\varphi\)
and repeatedly apply the inverse carry \texttt{100\ →\ 011} to the
triple covering the radix point and then successively to lower triples:

\[
(1.000\ldots)_\varphi=(.011000\ldots)_\varphi=(.001011000\ldots)_\varphi=(.00001011000\ldots)_\varphi\to\cdots
\]

Each step yields a different digit string but the same value (by (A)),
and the sequence converges to the infinite tail
\((.010101\ldots)_\varphi\). Hence, \(1\) has infinitely many base-ϕ
\(0/1\) representations.

\begin{enumerate}
\def\labelenumi{(\roman{enumi})}
\setcounter{enumi}{1}
\tightlist
\item
  Existence and uniqueness under the constraints
\end{enumerate}

Normal form (existence)

Given any \(0/1\) base-ϕ string for \(x\ge0\), repeatedly replace every
adjacent \texttt{11} by \texttt{100}, scanning from right to left. Each
replacement strictly moves weight to a more significant position, and no
new \texttt{11}'s are created to its right. Two things can happen:

\begin{itemize}
\tightlist
\item
  The process stops after finitely many steps: we obtain a
  representation with no adjacent 1s.
\item
  Or the carries propagate forever to the left in a perfectly
  alternating way; the only possible infinite stabilization is the
  notorious tail \(010101\ldots\) (shifted), which keeps the value the
  same (by the tail identity). Excluding representations that end with
  \(010101\ldots\) rules this out.
\end{itemize}

Thus every \(x\ge0\) admits at least one \emph{admissible}
representation: digits in \(\{0,1\}\), no adjacent 1s, and no terminal
\(0101\ldots\) tail.

Uniqueness

Suppose two different admissible strings represent the same \(x\). Let
\(j\) be the largest position where they differ; w.l.o.g. the first has
digit 1 at \(j\) and the second has 0. Because no adjacent 1s are
allowed, all 1s to the right of \(j\) in either string can occupy only
positions \(j-2, j-3,\dots\) with at least one zero between any two 1s.
The maximal value you can build from positions \textless{} \(j\) without
adjacency is the geometric series

\[
\varphi^{j-2}+\varphi^{j-4}+\varphi^{j-6}+\cdots
=\frac{\varphi^{j-2}}{1-\varphi^{-2}}
=\varphi^{j-1}.
\]

But the surplus contributed by having a 1 at position \(j\) is
\(\varphi^j\), which exceeds \(\varphi^{j-1}\). Therefore the two values
cannot be equal---a contradiction. Hence the admissible representation
is unique. (The extra ``no \(0101\ldots\) tail'' condition is what
removes the special ambiguity \(1=(.11)_\varphi=(.0101\ldots)_\varphi\)
and all its shifted variants.)

\begin{enumerate}
\def\labelenumi{(\roman{enumi})}
\setcounter{enumi}{2}
\tightlist
\item
  Which admissible strings represent integers?
\end{enumerate}

Write an admissible expansion as

\[
x=\sum_{j\in\mathbb Z} a_j\,\varphi^j\qquad (a_j\in\{0,1\}).
\]

Using the identities

\[
\varphi^{n}=F_n\varphi+F_{n-1}\quad(n\ge1),
\qquad
\varphi^{-n}=(\!-1)^{\,n+1}F_n\,\varphi+(\!-1)^{\,n}F_{n+1}\quad(n\ge1),
\]

(each proved by induction from \(\varphi^2=\varphi+1\)), we can regroup
\(x=A\varphi+B\) with integers \(A,B\). A short calculation gives

\[
A=\sum_{j\ge1} a_j F_j \;+\; \sum_{n\ge1} a_{-n}(\!-1)^{\,n+1}F_n,
\qquad
B=a_0+\sum_{j\ge1} a_j F_{j-1}+\sum_{n\ge1} a_{-n}(\!-1)^{\,n}F_{n+1}.
\]

Therefore \(x\) is an integer iff \(A=0\), i.e.

\[
\boxed{\;\sum_{j\ge1} a_j F_j \;=\; \sum_{n\ge1} (\!-1)^{\,n+1} a_{-n} F_n\;}
\tag{∗}
\]

and in that case the value is the integer \(B\). In words: in a
canonical base-ϕ string (no adjacent 1s, no \(0101\ldots\) tail), \(x\)
is an integer iff the Fibonacci-weighted sum of 1-digits left of the
radix equals the alternating Fibonacci-weighted sum of 1-digits right of
the radix.

A practical way to construct the (unique) integer representation is the
greedy algorithm consistent with the constraints: repeatedly place a 1
at the largest position \(j\) such that adding \(\varphi^j\) does not
exceed the target and does not create adjacency, then continue downward.
This terminates and produces the canonical string; the first few are:

\begin{itemize}
\tightlist
\item
  \(0=(0)_\varphi\)
\item
  \(1=(1)_\varphi\) (note \((.11)_\varphi\) is not admissible since it
  has adjacent 1s)
\item
  \(2=(10.01)_\varphi\)
\item
  \(3=(100.01)_\varphi\)
\item
  \(4=(101.01)_\varphi\)
\item
  \(5=(1000.1001)_\varphi\)
\item
  \(6=(1010.0001)_\varphi\)
\end{itemize}

(You can verify each by summing powers of \(\varphi\); no pair of
adjacent 1s occurs, and the expansions do not end with \(0101\ldots\).)

\subsection{36 {[}M32{]}}\label{m32}

\begin{enumerate}
\def\labelenumi{\arabic{enumi})}
\tightlist
\item
  Basic structure and the infinite word
\end{enumerate}

Because \(S_{n+2}=S_{n+1}S_n\) and \(|S_1|=|S_2|=1\), lengths satisfy
\(|S_{n+2}|=|S_{n+1}|+|S_n|\); hence \(|S_n|=F_n\) by induction.
Moreover, for \(n\ge2\) each \(S_{n+1}\) begins with \(\texttt{b}\), and
the last letter of \(S_n\) alternates
\(\texttt{a},\texttt{b},\texttt{a},\texttt{b},\dots\) (since the last
letter of \(S_{n+2}\) equals that of \(S_n\)). Thus the prefixes are
nested and there is a well-defined infinite limit

\[
S_\infty=\lim_{n\to\infty}S_n
\]

such that, for \(n>1\), \(S_n\) is exactly the first \(F_n\) letters of
\(S_\infty\). 

A convenient ``substitution'' description: if we set the morphism
\(\mu\) by \(\mu(\texttt{a})=\texttt{b},\ \mu(\texttt{b})=\texttt{ba}\),
then \(S_n=\mu^{\,n-1}(\texttt{a})\). (Check:
\(\mu(\texttt{a})=\texttt{b}=S_2\),
\(\mu^2(\texttt{a})=\mu(\texttt{b})=\texttt{ba}=S_3\),
\(\mu^3(\texttt{a})=\mu(\texttt{ba})=\texttt{bab}=S_4\), etc.)

\begin{enumerate}
\def\labelenumi{\arabic{enumi})}
\setcounter{enumi}{1}
\tightlist
\item
  Where do doubles occur?
\end{enumerate}

Two simple (and very useful) consequences:

\begin{itemize}
\item
  No \(\texttt{aa}\): \(\mu\) never produces \(\texttt{aa}\) (images are
  only \(\texttt{b}\) or \(\texttt{ba}\)). Inductively, if \(S_n\) has
  no \(\texttt{aa}\), then \(S_{n+1}\) and the concatenation
  \(S_{n+1}S_n\) also don't---across the boundary we see either
  \(\texttt{ab}\) or \(\texttt{bb}\), never \(\texttt{aa}\).
\item
  No \(\texttt{bbb}\): the only way \(\texttt{bb}\) arises is inside
  \(\mu(\texttt{ab})=\texttt{bba}\); every \(\texttt{bb}\) is
  immediately followed by \(\texttt{a}\). Hence \(\texttt{bbb}\) cannot
  occur.
\end{itemize}

\begin{enumerate}
\def\labelenumi{\arabic{enumi})}
\setcounter{enumi}{2}
\tightlist
\item
  Exact letter counts in \(S_n\)
\end{enumerate}

Let \(A_n\) and \(B_n\) be the numbers of \(\texttt{a}\)'s and
\(\texttt{b}\)'s in \(S_n\). Then

\[
(A_{n+2},B_{n+2})=(A_{n+1}+A_n,\ B_{n+1}+B_n),\qquad (A_1,B_1)=(1,0),\ (A_2,B_2)=(0,1).
\]

Solving gives

\[
A_n=F_{n-2},\qquad B_n=F_{n-1}\quad(n\ge1).
\]

So in \(S_n\) there are \(F_{n-2}\) \(\texttt{a}\)'s and \(F_{n-1}\)
\(\texttt{b}\)'s (and \(F_{n-2}+F_{n-1}=F_n\) total), matching the
book's answer. 

Densities. Since \(\dfrac{F_{n-1}}{F_n}\to\varphi^{-1}\) and
\(\dfrac{F_{n-2}}{F_n}\to\varphi^{-2}\)
(\(\varphi=\tfrac{1+\sqrt5}{2}\)), the asymptotic frequency of
\(\texttt{b}\) in \(S_\infty\) is \(\varphi^{-1}\) and of \(\texttt{a}\)
is \(\varphi^{-2}\).

\begin{enumerate}
\def\labelenumi{\arabic{enumi})}
\setcounter{enumi}{3}
\tightlist
\item
  Predicting the \(k\)-th letter --- three equivalent characterizations
\end{enumerate}

\begin{enumerate}
\def\labelenumi{(\alph{enumi})}
\tightlist
\item
  A fast recursive test (using Fibonacci lengths)
\end{enumerate}

Let \(L_n=F_n\). To find the letter at position \(k\) of \(S_\infty\):

\begin{enumerate}
\def\labelenumi{\arabic{enumi}.}
\tightlist
\item
  If \(k=1\) return \(\texttt{b}\); if \(k=2\) return \(\texttt{a}\).
\item
  Let \(n\) be such that \(F_n<k\le F_{n+1}\). If \(k\le F_{n+1}\) and
  \(k>F_n\), you're in the left block \(S_{n+1}\); replace
  \((k,n)\leftarrow(k,n-1)\). If \(k\le F_n\), you're in the right block
  of the previous step; replace \((k,n)\leftarrow(k,n-2)\).
\item
  Repeat until you hit the base.
\end{enumerate}

Unwinding this exactly subtracts Fibonacci lengths according to the
Zeckendorf greedy decomposition of \(k-1\) into nonconsecutive Fibonacci
numbers (see Exercise 34).

\begin{enumerate}
\def\labelenumi{(\alph{enumi})}
\setcounter{enumi}{1}
\tightlist
\item
  Beatty/``mechanical word'' formula
\end{enumerate}

Let

\[
B(k)\;=\;\#\{\texttt{b}\text{’s among the first }k\text{ letters of }S_\infty\}.
\]

One proves (by induction using \(S_{n+2}=S_{n+1}S_n\) and the identities
\(\lfloor(F_n+1)\varphi^{-1}\rfloor=F_{n-1}\)) that

\[
\boxed{\,B(k)=\big\lfloor (k+1)\varphi^{-1}\big\rfloor\,}.
\]

Therefore the \(k\)-th letter is \(\texttt{b}\) iff \(B(k)-B(k-1)=1\),
i.e.

\[
\boxed{\ \text{letter}_k=\texttt{b}\iff 
\Big\lfloor\tfrac{k+1}{\varphi}\Big\rfloor-\Big\lfloor\tfrac{k}{\varphi}\Big\rfloor=1\ }.
\]

\begin{enumerate}
\def\labelenumi{(\alph{enumi})}
\setcounter{enumi}{2}
\tightlist
\item
  Positions of all \(\texttt{b}\)'s
\end{enumerate}

Equivalently, the positions of \(\texttt{b}\) form the Beatty sequence

\[
\boxed{\ k=\big\lfloor m\varphi\big\rfloor\quad(m=1,2,3,\dots)\ }.
\]

This is again one of the book's listed characterizations. (It's the
complementary Beatty partner of \(\lfloor m\varphi^2\rfloor\),
reflecting that the \(\texttt{a}\)/\(\texttt{b}\) frequencies are
\(\varphi^{-2}\) and \(\varphi^{-1}\).)

\begin{enumerate}
\def\labelenumi{\arabic{enumi})}
\setcounter{enumi}{4}
\tightlist
\item
  Where do doubles occur?
\end{enumerate}

From §2 we know \(\texttt{aa}\) never appears and \(\texttt{bbb}\) never
appears. But \(\texttt{bb}\) certainly does: the first time is in
\(S_5=\texttt{babba}\). In fact every \(\texttt{bb}\) is immediately
followed by \(\texttt{a}\) (all \(\texttt{bb}\) come from
\(\mu(\texttt{ab})=\texttt{bba}\)), so all double letters are precisely
the factors \(\texttt{bba}\). (This also explains why the run lengths of
\(\texttt{b}\)'s are 1 or 2 only.)

\begin{enumerate}
\def\labelenumi{\arabic{enumi})}
\setcounter{enumi}{5}
\tightlist
\item
  Summary of key facts (what the book expects you to ``discover'')
\end{enumerate}

\begin{itemize}
\tightlist
\item
  \(S_n\) is the prefix of length \(F_n\) of \(S_\infty\). No
  \(\texttt{aa}\); no \(\texttt{bbb}\). Counts:
  \(\#\texttt{a}=F_{n-2}\), \(\#\texttt{b}=F_{n-1}\). 
\item
  \(k\)-th letter via Zeckendorf (Exercise 34): it's \(\texttt{a}\) iff
  the smallest Fibonacci used in \(k-1\) is \(F_2\).
\item
  Exact counting and ``density'' of \(\texttt{b}\):
  \(B(k)=\lfloor (k+1)/\varphi\rfloor\); hence \(\texttt{b}\)-density
  \(=\varphi^{-1}\). Equivalent: the \(\texttt{b}\) positions are
  exactly \(\{\lfloor m\varphi\rfloor: m\ge1\}\).
\end{itemize}

\begin{enumerate}
\def\labelenumi{\arabic{enumi})}
\setcounter{enumi}{6}
\tightlist
\item
  Tiny sanity checks
\end{enumerate}

\[
\begin{array}{c|c|c|c}
n & S_n & (\#\texttt{a},\#\texttt{b}) & \text{first }F_n\text{ letters’ }B(F_n)\stackrel{?}{=}F_{n-1}\\\hline
3 & \texttt{ba} & (1,1) & \lfloor(2+1)/\varphi\rfloor=1=F_2\\
4 & \texttt{bab} & (1,2) & \lfloor(3+1)/\varphi\rfloor=2=F_3\\
5 & \texttt{babba} & (2,3) & \lfloor(5+1)/\varphi\rfloor=3=F_4
\end{array}
\]

All match.

\subsection{37 {[}M35{]}}\label{m35}

There is one heap with \(n\) chips. On the next move a player may remove
at most \(q\) chips. (In Fibonacci Nim, after a player removes \(m\)
chips, the following move is bounded by \(q'=2m\); the first move
typically has \(q=n-1\).)

Write \(n\) in its Zeckendorf representation (sum of nonconsecutive
Fibonacci numbers):

\[
n=F_{k_1}+F_{k_2}+\cdots+F_{k_r},\qquad k_1>k_2>\cdots>k_r\ge 2,
\]

and define

\[
\mu(n)\stackrel{\text{def}}{=} F_{k_r}\quad\text{(the smallest Fibonacci in the sum)},\qquad \mu(0)=\infty.
\]

Claim. From a position \((n,q)\), there is a winning move iff
\(\mu(n)\le q\). In particular, in the usual game where the first move
has \(q=n-1\), the first player wins unless \(n\) is itself a Fibonacci
number. (When \(n\) is Fibonacci, \(\mu(n)=n>n-1=q\), so the first
player loses.) 

Preliminaries about \(\mu\)

Let \(n=F_{k_1}+\cdots+F_{k_r}\) be in Zeckendorf form and recall
consecutive Fibonacci indices never both appear.

We will use four easy inequalities:

\begin{enumerate}
\def\labelenumi{(\Alph{enumi})}
\tightlist
\item
  If \(n>0\), then
\end{enumerate}

\[
\mu\bigl(n-\mu(n)\bigr)>\;2\,\mu(n).
\]

\emph{Reason.} Removing the smallest term \(F_{k_r}\) leaves
\(F_{k_1}+\cdots+F_{k_{r-1}}\). Since \(k_{r-1}\ge k_r+2\), the next
smallest possible Fibonacci in the remainder is at least
\(F_{k_r+2}>2F_{k_r}\). Thus \(\mu(n-\mu(n))>2\mu(n)\). 

\begin{enumerate}
\def\labelenumi{(\Alph{enumi})}
\setcounter{enumi}{1}
\tightlist
\item
  If \(0<m<F_k\), then
\end{enumerate}

\[
\mu(m)\le 2\bigl(F_k-m\bigr).
\]

\emph{Reason (sketch).} The largest Fibonacci allowed in a Zeckendorf
sum strictly below \(F_k\) is at most \(F_{k-1}\), and the ``greedy''
parity-alternating tail \(F_{k-1}+F_{k-3}+\cdots\) bounds any such
\(m\). Rewriting this bound gives \(F_k-m\ge \tfrac12\,\mu(m)\). 

\begin{enumerate}
\def\labelenumi{(\Alph{enumi})}
\setcounter{enumi}{2}
\tightlist
\item
  If \(0<m<\mu(n)\), then
\end{enumerate}

\[
\mu\bigl(n-\mu(n)+m\bigr)\le 2\bigl(\mu(n)-m\bigr).
\]

\emph{Reason.} Apply (B) with \(F_k=\mu(n)\) to the increment \(m\). 

\begin{enumerate}
\def\labelenumi{(\Alph{enumi})}
\setcounter{enumi}{3}
\tightlist
\item
  If \(0<m<\mu(n)\), then
\end{enumerate}

\[
\mu\bigl(n-m\bigr)\le 2m.
\]

\emph{Reason.} Substitute \(m\gets \mu(n)-m\) in (C). 

Characterization of winning and losing positions

We prove by induction on \(n\) that:

\begin{quote}
Theorem. From position \((n,q)\) there exists a winning move iff
\(\mu(n)\le q\).
\end{quote}

\begin{itemize}
\item
  (\(\Rightarrow\)) Losing positions when \(\mu(n)>q\). Any legal move
  removes some \(m\le q<\mu(n)\). After this move the heap is
  \(n' = n-m\) and the next bound is \(q'=2m\) (Fibonacci Nim rule). By
  (D), \(\mu(n')\le 2m=q'\). Thus every possible move hands the opponent
  a position with \(\mu(n')\le q'\), i.e., a winning position for them.
  So \((n,q)\) is losing. 
\item
  (\(\Leftarrow\)) Winning positions when \(\mu(n)\le q\). If \(q\ge n\)
  you can take all \(n\) and win immediately. Otherwise, remove exactly
  \(m=\mu(n)\) (legal since \(m\le q\)). Now \(n' = n-\mu(n)\) and
  \(q' = 2\mu(n)\). By (A), \(\mu(n')>2\mu(n)=q'\). Hence the opponent
  receives a losing position. Therefore \((n,q)\) is winning. 
\end{itemize}

Combining the two directions establishes the theorem. As a corollary, in
the standard initial position for Fibonacci Nim, \(q=n-1\). Then the
first player wins unless \(\mu(n)=n\), i.e., unless \(n\) is itself a
Fibonacci number. 

What to play (algorithmic recipe)

Given \((n,q)\):

\begin{enumerate}
\def\labelenumi{\arabic{enumi}.}
\item
  Compute the Zeckendorf form \(n=F_{k_1}+\cdots+F_{k_r}\) and
  \(\mu(n)=F_{k_r}\).
\item
  If \(\mu(n)>q\): you are lost (perfect play assumed). Any move
  \(m\le q\) gives the opponent \(\mu(n-m)\le 2m\) (by (D)), hence a
  winning reply.
\item
  If \(\mu(n)\le q\): remove exactly \(\mu(n)\). This forces
  \(\mu(n-\mu(n))>2\mu(n)=q'\) (by (A)), so you pass a losing position
  to your opponent. More generally, all winning moves are: remove a
  suffix of the Zeckendorf sum

  \[
  m=F_{k_j}+F_{k_{j+1}}+\cdots+F_{k_r}
  \]

  with either \(j=1\) (take the whole representation) or satisfying the
  legality condition \(F_{k_{j-1}}>2m\) (so \(m\le q\) on this turn). 
\end{enumerate}

Example

\(n=1000\). Zeckendorf: \(1000=987+13\). The only ``lucky'' (winning)
first move is to take \(13\) (i.e., \(\mu(1000)\)). Any other first move
allows a reply that restores \(\mu(\cdot)\)-advantage to the opponent.
Thus with perfect play, the first player wins by removing \(13\).

\subsection{38 {[}35{]}}\label{section-96}

\begin{enumerate}
\def\labelenumi{\arabic{enumi})}
\tightlist
\item
  Problem
\end{enumerate}

Two players alternately remove chips from a pile of size \(n\).

\begin{itemize}
\tightlist
\item
  On the first move you may take any number \(1\le m \le n-1\) (you
  cannot take the whole pile).
\item
  After that, if the previous move took \(t\) chips, the next move may
  take any \(m\) with \(1\le m \le 2t\).
\item
  The player who takes the last chip wins.
\end{itemize}

\begin{enumerate}
\def\labelenumi{\arabic{enumi})}
\setcounter{enumi}{1}
\tightlist
\item
  Key theory (why the strategy works)
\end{enumerate}

Let every positive integer \(n\) be written in its Zeckendorf
representation,
\(n \;=\; F_{k_1}+F_{k_2}+\cdots+F_{k_r},\qquad k_1 \gg k_2 \gg \cdots \gg k_r \gg 0,\)
i.e., as a sum of nonconsecutive Fibonacci numbers (with
\(F_0=0, F_1=1\), hence the smallest summand index is \(k_r\ge2\)). This
representation exists and is unique.

Define \(\mu(n)\) to be the smallest Fibonacci summand in the Zeckendorf
expansion of \(n\) (and set \(\mu(0)=\infty\)). The answers section
proves the crucial fact:

\begin{quote}
If a position has \(n\) chips and the next player may take at most \(q\)
chips, then there is a winning move iff \(\mu(n)\le q\). In that case, a
winning move is to take exactly \(\mu(n)\) chips (or all remaining chips
if \(q\ge n\)); this leaves the opponent a ``bad'' position. Conversely,
if \(\mu(n)>q\) then every legal move allows the opponent a winning
reply. 
\end{quote}

Immediate corollaries:

\begin{itemize}
\tightlist
\item
  At the start, \(q=n-1\). If \(n\) itself is a Fibonacci number, then
  \(\mu(n)=n>q\), so the first player is lost against perfect play.
  Otherwise the first player has a win. (This is exactly what the book
  notes, and for \(n=1000\) it shows \(1000=987+13\), so the only
  winning first move is to take 13.) 
\end{itemize}

\begin{enumerate}
\def\labelenumi{\arabic{enumi})}
\setcounter{enumi}{2}
\tightlist
\item
  Optimal algorithm
\end{enumerate}

Maintain \(n\) (chips left) and \(q\) (max you're allowed to take now).
The optimal move is:

\begin{enumerate}
\def\labelenumi{\arabic{enumi}.}
\item
  If \(n\le q\): take all and win.
\item
  Compute \(\mu(n)\) = smallest term in the Zeckendorf representation of
  \(n\).

  \begin{itemize}
  \tightlist
  \item
    If \(\mu(n)\le q\): take \(\mu(n)\).
  \item
    Else (a losing position): pick any legal \(m\in[1,\min(q,n)]\). (All
    are losing vs.~perfect play; choosing \(m=1\) is fine.)
  \end{itemize}
\end{enumerate}

How to compute \(\mu(n)\) fast: Precompute Fibonacci numbers
\(F_2=1,F_3=2,\dots\) up to \(n\). Use the greedy Zeckendorf algorithm:
repeatedly subtract the largest \(F_k\le n\), skipping the immediate
next \(F_{k-1}\) to avoid consecutive terms. The last subtraction gives
\(\mu(n)\). This runs in \(O(\log n)\) time per move.

\begin{enumerate}
\def\labelenumi{\arabic{enumi})}
\setcounter{enumi}{3}
\tightlist
\item
  Reference Python implementation
\end{enumerate}

\begin{Shaded}
\begin{Highlighting}[]
\ImportTok{from}\NormalTok{ bisect }\ImportTok{import}\NormalTok{ bisect\_right}

\KeywordTok{def}\NormalTok{ fibs\_up\_to(n):}
    \CommentTok{"""Return [F2, F3, ..., Fm] with F2=1, F3=2, ... up to ≤ n."""}
\NormalTok{    fibs }\OperatorTok{=}\NormalTok{ [}\DecValTok{1}\NormalTok{, }\DecValTok{2}\NormalTok{]  }\CommentTok{\# F2, F3}
    \ControlFlowTok{while}\NormalTok{ fibs[}\OperatorTok{{-}}\DecValTok{1}\NormalTok{] }\OperatorTok{+}\NormalTok{ fibs[}\OperatorTok{{-}}\DecValTok{2}\NormalTok{] }\OperatorTok{\textless{}=}\NormalTok{ n:}
\NormalTok{        fibs.append(fibs[}\OperatorTok{{-}}\DecValTok{1}\NormalTok{] }\OperatorTok{+}\NormalTok{ fibs[}\OperatorTok{{-}}\DecValTok{2}\NormalTok{])}
    \ControlFlowTok{return}\NormalTok{ fibs}

\KeywordTok{def}\NormalTok{ zeckendorf\_smallest(n, fibs):}
    \CommentTok{"""}
\CommentTok{    Return μ(n): smallest Fibonacci summand in Zeckendorf rep of n.}
\CommentTok{    Assumes n\textgreater{}0.}
\CommentTok{    """}
\NormalTok{    m }\OperatorTok{=}\NormalTok{ n}
\NormalTok{    smallest }\OperatorTok{=} \VariableTok{None}
\NormalTok{    i }\OperatorTok{=} \BuiltInTok{len}\NormalTok{(fibs) }\OperatorTok{{-}} \DecValTok{1}
\NormalTok{    skip\_next\_index }\OperatorTok{=} \OperatorTok{{-}}\DecValTok{1}  \CommentTok{\# index to skip to avoid consecutive F\textquotesingle{}s}

    \ControlFlowTok{while}\NormalTok{ m }\OperatorTok{\textgreater{}} \DecValTok{0}\NormalTok{:}
        \CommentTok{\# largest fib \textless{}= m}
\NormalTok{        i }\OperatorTok{=}\NormalTok{ bisect\_right(fibs, m) }\OperatorTok{{-}} \DecValTok{1}
        \CommentTok{\# enforce nonconsecutive rule: if we just used fibs[i+1], skip fibs[i]}
        \ControlFlowTok{if}\NormalTok{ i }\OperatorTok{==}\NormalTok{ skip\_next\_index:}
\NormalTok{            i }\OperatorTok{{-}=} \DecValTok{1}
\NormalTok{        f }\OperatorTok{=}\NormalTok{ fibs[i]}
\NormalTok{        m }\OperatorTok{{-}=}\NormalTok{ f}
\NormalTok{        smallest }\OperatorTok{=}\NormalTok{ f  }\CommentTok{\# last chosen will be the smallest}
        \CommentTok{\# next, we must skip the immediate smaller Fibonacci number}
\NormalTok{        skip\_next\_index }\OperatorTok{=}\NormalTok{ i }\OperatorTok{{-}} \DecValTok{1}

    \ControlFlowTok{return}\NormalTok{ smallest}

\KeywordTok{def}\NormalTok{ optimal\_move(n, q, fibs\_cache}\OperatorTok{=}\VariableTok{None}\NormalTok{):}
    \CommentTok{"""}
\CommentTok{    Given pile size n and max{-}allowed q (1st move uses q=n{-}1),}
\CommentTok{    return the optimal number of chips to take now.}
\CommentTok{    """}
    \ControlFlowTok{if}\NormalTok{ n }\OperatorTok{\textless{}=} \DecValTok{0}\NormalTok{:}
        \ControlFlowTok{return} \DecValTok{0}
    \ControlFlowTok{if}\NormalTok{ q }\OperatorTok{\textgreater{}=}\NormalTok{ n:}
        \ControlFlowTok{return}\NormalTok{ n  }\CommentTok{\# take all and win}

\NormalTok{    fibs }\OperatorTok{=}\NormalTok{ fibs\_cache }\KeywordTok{or}\NormalTok{ fibs\_up\_to(n)}
\NormalTok{    mu }\OperatorTok{=}\NormalTok{ zeckendorf\_smallest(n, fibs)}

    \ControlFlowTok{if}\NormalTok{ mu }\OperatorTok{\textless{}=}\NormalTok{ q:}
        \ControlFlowTok{return}\NormalTok{ mu}
    \CommentTok{\# Losing position: any legal move is equivalent vs. perfect opponent.}
    \CommentTok{\# Take 1 to keep it simple.}
    \ControlFlowTok{return} \DecValTok{1}

\CommentTok{\# Example: first move with n=1000 (q=999)}
\ControlFlowTok{if} \VariableTok{\_\_name\_\_} \OperatorTok{==} \StringTok{"\_\_main\_\_"}\NormalTok{:}
\NormalTok{    n }\OperatorTok{=} \DecValTok{1000}
\NormalTok{    q }\OperatorTok{=}\NormalTok{ n }\OperatorTok{{-}} \DecValTok{1}
\NormalTok{    move }\OperatorTok{=}\NormalTok{ optimal\_move(n, q)}
    \BuiltInTok{print}\NormalTok{(move)  }\CommentTok{\# expected 13 (since 1000=987+13)}
\end{Highlighting}
\end{Shaded}

What it prints for 1000: \texttt{13}, matching the book's worked example
(since \(1000=987+13\)). 

Using it in a full game loop

Track \texttt{prev\_take} to update the next \texttt{q\ =\ 2*prev\_take}
each turn:

\begin{Shaded}
\begin{Highlighting}[]
\KeywordTok{def}\NormalTok{ play\_optimal(n, first}\OperatorTok{=}\VariableTok{True}\NormalTok{):}
    \CommentTok{"""}
\CommentTok{    Simulate an optimal{-}vs{-}optimal game starting from pile n.}
\CommentTok{    Returns the winner: \textquotesingle{}first\textquotesingle{} or \textquotesingle{}second\textquotesingle{}.}
\CommentTok{    """}
\NormalTok{    player }\OperatorTok{=} \StringTok{\textquotesingle{}first\textquotesingle{}} \ControlFlowTok{if}\NormalTok{ first }\ControlFlowTok{else} \StringTok{\textquotesingle{}second\textquotesingle{}}
\NormalTok{    prev\_take }\OperatorTok{=} \VariableTok{None}
\NormalTok{    fibs }\OperatorTok{=}\NormalTok{ fibs\_up\_to(n)}

    \ControlFlowTok{while}\NormalTok{ n }\OperatorTok{\textgreater{}} \DecValTok{0}\NormalTok{:}
\NormalTok{        q }\OperatorTok{=}\NormalTok{ (n}\OperatorTok{{-}}\DecValTok{1}\NormalTok{) }\ControlFlowTok{if}\NormalTok{ prev\_take }\KeywordTok{is} \VariableTok{None} \ControlFlowTok{else} \BuiltInTok{min}\NormalTok{(}\DecValTok{2}\OperatorTok{*}\NormalTok{prev\_take, n)}
\NormalTok{        take }\OperatorTok{=}\NormalTok{ optimal\_move(n, q, fibs\_cache}\OperatorTok{=}\NormalTok{fibs)}
\NormalTok{        n }\OperatorTok{{-}=}\NormalTok{ take}
\NormalTok{        prev\_take }\OperatorTok{=}\NormalTok{ take}
\NormalTok{        player }\OperatorTok{=} \StringTok{\textquotesingle{}second\textquotesingle{}} \ControlFlowTok{if}\NormalTok{ player }\OperatorTok{==} \StringTok{\textquotesingle{}first\textquotesingle{}} \ControlFlowTok{else} \StringTok{\textquotesingle{}first\textquotesingle{}}

    \CommentTok{\# last player who moved just switched \textquotesingle{}player\textquotesingle{}, so the winner is the other one}
    \ControlFlowTok{return} \StringTok{\textquotesingle{}second\textquotesingle{}} \ControlFlowTok{if}\NormalTok{ player }\OperatorTok{==} \StringTok{\textquotesingle{}first\textquotesingle{}} \ControlFlowTok{else} \StringTok{\textquotesingle{}first\textquotesingle{}}
\end{Highlighting}
\end{Shaded}

\begin{enumerate}
\def\labelenumi{\arabic{enumi})}
\setcounter{enumi}{4}
\tightlist
\item
  Correctness sketch
\end{enumerate}

\begin{itemize}
\item
  Zeckendorf's theorem gives a canonical decomposition
  \(n=\sum F_{k_j}\) with no adjacent indices.
\item
  The answers section establishes properties (A)--(D) of \(\mu\) and
  concludes:

  \begin{itemize}
  \tightlist
  \item
    A winning move exists iff \(\mu(n)\le q\),
  \item
    and the canonical winning move is to take \(\mu(n)\). This exactly
    matches the \texttt{optimal\_move} rule above. 
  \end{itemize}
\end{itemize}

Thus the program implements the mathematically proven optimal strategy.

\begin{enumerate}
\def\labelenumi{\arabic{enumi})}
\setcounter{enumi}{5}
\tightlist
\item
  Complexity
\end{enumerate}

\begin{itemize}
\tightlist
\item
  Precomputing Fibonacci numbers up to \(n\): \(O(\log n)\) numbers.
\item
  Each move computes \(\mu(n)\) via a greedy pass that subtracts at most
  \(O(\log n)\) Fibonacci numbers.
\item
  So the decision per move is \(O(\log n)\) time and \(O(\log n)\)
  space.
\end{itemize}

\subsection{39 {[}M24{]}}\label{m24-13}

Solution 1 --- Characteristic-polynomial method

Look for solutions of the form \(a_n=r^n\). Substituting into the
recurrence:

\[
r^{n+2}=r^{n+1}+6r^n\quad\Longleftrightarrow\quad r^2-r-6=0.
\]

Factor:

\[
(r-3)(r+2)=0\quad\Rightarrow\quad r_1=3,\; r_2=-2.
\]

Therefore the general term is a linear combination of the two
independent solutions:

\[
a_n=\alpha\cdot 3^n+\beta\cdot(-2)^n,
\]

with constants \(\alpha,\beta\) fixed by the initial values.

Use \(a_0=0\): \(\alpha+\beta=0\Rightarrow \beta=-\alpha\).

Use \(a_1=1\): \(3\alpha-2\beta=1\). Substitute \(\beta=-\alpha\):

\[
3\alpha-2(-\alpha)=5\alpha=1\;\Rightarrow\;\alpha=\tfrac15,\quad\beta=-\tfrac15.
\]

Hence

\[
\boxed{\,a_n=\frac{3^n-(-2)^n}{5}\,}.
\]

Solution 2 --- Generating-function method

Let \(A(z)=\sum_{n\ge0} a_n z^n\). Multiply the recurrence by
\(z^{n+2}\) and sum over \(n\ge0\):

\[
\sum_{n\ge0}a_{n+2}z^{n+2}=\sum_{n\ge0}a_{n+1}z^{n+2}+6\sum_{n\ge0}a_n z^{n+2}.
\]

Using \(a_0=0,a_1=1\), this becomes

\[
A(z)-z=z\,A(z)+6z^2A(z)\;\Rightarrow\;A(z)=\frac{z}{1-z-6z^2}.
\]

Factor the denominator: \(1-z-6z^2=(1-3z)(1+2z)\). Do partial fractions:

\[
\frac{z}{1-z-6z^2}=\frac{1}{5}\left(\frac{1}{1-3z}-\frac{1}{1+2z}\right).
\]

Expanding,

\[
A(z)=\sum_{n\ge0}\frac{1}{5}\bigl(3^n-(-2)^n\bigr)z^n,
\]

so \(a_n=\frac{3^n-(-2)^n}{5}\) again.

Sanity checks and growth

First terms: \(0,1,1,7,13,55,\dots\) (plugging \(n=0,1,2,3,4,5\) into
the formula). As \(n\) grows, the \(3^n\) term dominates, so
\(a_n\sim 3^n/5\).

\subsection{40 {[}M25{]}}\label{m25-16}

Define \(f(n)\) as the minimum possible worst-case cost to identify one
item out of \(n\) using a binary decision tree where each left edge
costs \(1\) and each right edge costs \(2\). If at a node that covers
\(n\) items you split into a left subtree of size \(k\) and a right
subtree of size \(n-k\), the worst-case cost through that node is

\[
\max\bigl(1+f(k),\;2+f(n-k)\bigr).
\]

Thus

\[
f(1)=0,\qquad
f(n)=\min_{1\le k<n}\max\bigl(1+f(k),\,2+f(n-k)\bigr)\quad(n\ge2).
\]

Claim. For every \(m\ge0\) and \(n\ge1\),

\[
\boxed{\;F_m<n\le F_{m+1}\ \Longleftrightarrow\ f(n)=m\;}
\]

(where \(F_0=0, F_1=1, F_{i+1}=F_i+F_{i-1}\) are Fibonacci numbers).

Proof (by induction on \(m\))

Base case \(m=0\).

The interval \(F_0<n\le F_1\) is \(0<n\le1\), i.e., \(n=1\). By
definition \(f(1)=0\). ✓

Inductive step.

Assume the statement holds for all smaller indices than a fixed
\(m\ge1\). Let \(n\) satisfy

\[
F_m<n\le F_{m+1}.
\]

Upper bound \(f(n)\le m\).

Choose the split \(k=F_m\). By the induction hypothesis, \(f(F_m)=m-1\).
Also

\[
n-k = n-F_m \le F_{m+1}-F_m = F_{m-1},
\]

so again by induction \(f(n-k)\le m-2\). Hence

\[
f(n)\;\le\;\max\bigl(1+(m-1),\,2+(m-2)\bigr)=m.
\]

Lower bound \(f(n)\ge m\).

Suppose, for contradiction, that \(f(n)<m\). Then in the optimal split
at some \(k\) we must have both branches strictly below \(m\):

\[
1+f(k)<m\quad\text{and}\quad 2+f(n-k)<m,
\]

i.e.,

\[
f(k)\le m-2,\qquad f(n-k)\le m-3.
\]

By the induction hypothesis, these imply

\[
k\le F_{m-1},\qquad n-k\le F_{m-2}.
\]

Adding yields \(n\le F_{m-1}+F_{m-2}=F_m\), contradicting \(n>F_m\).
Therefore \(f(n)\ge m\).

Combining with the upper bound gives \(f(n)=m\), completing the
induction. ∎

Consequences and intuition

\begin{itemize}
\item
  The optimal worst-case cost is the least \(m\) with \(n\le F_{m+1}\).
  Equivalently,

  \[
  f(n)=\min\{\,m:\ n\le F_{m+1}\,\}.
  \]

  As a rough asymptotic using Binet's formula,
  \(F_{m+1}\approx \varphi^{m+1}/\sqrt5\), so

  \[
  f(n)\;\approx\;\left\lceil\log_\varphi(\sqrt5\,n)\right\rceil-1,
  \]

  but the exact answer is the Fibonacci threshold characterization
  above.
\item
  The optimal split at size \(n\) (when \(F_m<n\le F_{m+1}\)) is
  \(k=F_m\): the left subtree handles \(F_m\) items (cost \(m-1\)) and
  the right handles at most \(F_{m-1}\) items (cost \(\le m-2\)). The
  resulting optimal trees are precisely the Fibonacci trees (each node's
  subtrees have ``heights'' differing as \(m-1\) and \(m-2\)).
\end{itemize}

Small values

Using the rule \(F_m<n\le F_{m+1}\):

\begin{itemize}
\tightlist
\item
  \(n=1\Rightarrow f(n)=0\)
\item
  \(n=2\Rightarrow f(n)=2\)
\item
  \(n=3,5\Rightarrow f(n)=3,4\) respectively
\item
  \(n=6,7,8\Rightarrow f(n)=5\), etc.
\end{itemize}

\subsection{41 {[}M25{]}}\label{m25-17}

Let \(f(x)=\lfloor x+\varphi^{-1}\rfloor\), where
\(\varphi=(1+\sqrt5)/2\). If the (Zeckendorf) Fibonacci representation
of \(n\) is

\[
n = F_{k_1}+F_{k_2}+\cdots+F_{k_r}\qquad(k_1\gg k_2\gg\cdots\gg k_r\gg 0,\ \text{so }k_{i+1}\le k_i-2),
\]

prove

\[
F_{k_1+1}+F_{k_2+1}+\cdots+F_{k_r+1}\;=\;f(\varphi n)=\Big\lfloor \varphi n+\varphi^{-1}\Big\rfloor,
\]

and find the analogous formula for \(F_{k_1-1}+\cdots+F_{k_r-1}\).

We will use Binet's formula and the basic facts about
\(\varphi' = 1-\varphi = -\varphi^{-1}\). In §1.2.8 it is shown that

\[
F_m=\frac{\varphi^{m}-\varphi'^{\,m}}{\sqrt5}\qquad(m\in\mathbb{Z}),
\]

and Exercise 34 gives the Zeckendorf form (no consecutive indices,
smallest index \(\ge 2\)).

\begin{enumerate}
\def\labelenumi{\arabic{enumi})}
\tightlist
\item
  Shift up by 1:
  \(\;\sum F_{k_i+1} = \lfloor \varphi n + \varphi^{-1}\rfloor\)
\end{enumerate}

Write the Zeckendorf sum \(n=\sum_i F_{k_i}\). Using Binet and summing
termwise,

\[
n \;=\;\sum_i \frac{\varphi^{k_i}-\varphi'^{\,k_i}}{\sqrt5}.
\]

Now consider the ``+1 shift''

\[
S_{+}\;:=\;\sum_i F_{k_i+1}\;=\;\sum_i \frac{\varphi^{k_i+1}-\varphi'^{\,k_i+1}}{\sqrt5}
\;=\;\frac{\varphi\sum_i \varphi^{k_i}-\varphi'\sum_i \varphi'^{\,k_i}}{\sqrt5}.
\]

Compare this with
\(\varphi n=\frac{\varphi\sum_i \varphi^{k_i}-\varphi\sum_i \varphi'^{\,k_i}}{\sqrt5}\).
Subtract:

\[
S_{+}-\varphi n
=\frac{(\!-\varphi')-\!(-\varphi)}{\sqrt5}\sum_i \varphi'^{\,k_i}
=\frac{\varphi-\varphi'}{\sqrt5}\sum_i \varphi'^{\,k_i}
=\sum_i \varphi'^{\,k_i},
\]

since \(\varphi-\varphi'=\sqrt5\). Therefore

\[
\boxed{\,S_{+}=\varphi n+\sum_i \varphi'^{\,k_i}\,}.
\]

Because \(\varphi'=-\varphi^{-1}\) and the Zeckendorf indices satisfy
\(k_i\ge 2\) and have gaps \(\ge 2\), the ``error term''
\(\displaystyle E:=\sum_i \varphi'^{\,k_i}\) is squeezed between the two
alternating geometric tails:

\[
\underbrace{\varphi'^{3}+\varphi'^{5}+\cdots}_{=\;\varphi^{-1}-1}
\;\le\;
E
\;\le\;
\underbrace{\varphi'^{2}+\varphi'^{4}+\cdots}_{=\;\varphi^{-1}}.
\]

Indeed, the largest possible (nonnegative) contribution comes by taking
only even exponents (bounded by the even-power tail), and the most
negative by taking only odd exponents \(\ge3\) (bounded below by the
odd-power tail). These sums evaluate to

\[
\sum_{j\ge1}\varphi'^{2j}=\frac{\varphi'^{2}}{1-\varphi'^{2}}=\frac{1/\varphi^{2}}{1-1/\varphi^{2}}=\frac1{\varphi},
\qquad
\sum_{j\ge1}\varphi'^{2j+1}=\varphi^{-1}-1.
\]

Hence

\[
\varphi n+(\varphi^{-1}-1)\;\le\;S_{+}\;\le\;\varphi n+\varphi^{-1}.
\]

This interval has length \(1\), so it contains exactly one integer; that
integer is \(\lfloor \varphi n+\varphi^{-1}\rfloor\). Therefore

\[
\boxed{\;F_{k_1+1}+\cdots+F_{k_r+1}=\big\lfloor \varphi n+\varphi^{-1}\big\rfloor=f(\varphi n)\;}.
\]

\begin{enumerate}
\def\labelenumi{\arabic{enumi})}
\setcounter{enumi}{1}
\tightlist
\item
  Shift down by 1:
  \(\;\sum F_{k_i-1} = \lfloor \varphi^{-1} n + \varphi^{-1}\rfloor\)
\end{enumerate}

Define

\[
S_{-}\;:=\;\sum_i F_{k_i-1}
=\sum_i \frac{\varphi^{k_i-1}-\varphi'^{\,k_i-1}}{\sqrt5}
=\frac{\varphi^{-1}\sum_i \varphi^{k_i}-\varphi'^{-1}\sum_i \varphi'^{\,k_i}}{\sqrt5}.
\]

Compare with
\(\varphi^{-1}n=\frac{\varphi^{-1}\sum_i \varphi^{k_i}-\varphi^{-1}\sum_i \varphi'^{\,k_i}}{\sqrt5}\).
Subtract:

\[
S_{-}-\varphi^{-1}n
=\frac{\big(\!-\varphi'^{-1}\big)-\big(\!-\varphi^{-1}\big)}{\sqrt5}\sum_i \varphi'^{\,k_i}
=\frac{\varphi^{-1}-\varphi'^{-1}}{\sqrt5}\sum_i \varphi'^{\,k_i}.
\]

But \(\varphi'^{-1}=-\varphi\), and
\(\varphi^{-1}-\varphi'^{-1}=\varphi^{-1}+ \varphi=\sqrt5\). Thus

\[
S_{-}-\varphi^{-1}n=\sum_i \varphi'^{\,k_i},
\quad\text{so}\quad
\boxed{\,S_{-}=\varphi^{-1} n+\sum_i \varphi'^{\,k_i}\,}.
\]

The same bounds on \(\sum_i \varphi'^{\,k_i}\) therefore give

\[
\varphi^{-1}n+\big(\varphi^{-1}-1\big)\;\le\;S_{-}\;\le\;\varphi^{-1}n+\varphi^{-1},
\]

so again
\(S_{-}=\lfloor \varphi^{-1}n+\varphi^{-1}\rfloor=f(\varphi^{-1}n)\).
Hence

\[
\boxed{\;F_{k_1-1}+\cdots+F_{k_r-1}=\big\lfloor \varphi^{-1} n+\varphi^{-1}\big\rfloor=f(\varphi^{-1} n)\;}.
\]

\begin{enumerate}
\def\labelenumi{\arabic{enumi})}
\setcounter{enumi}{2}
\tightlist
\item
  Quick sanity check (example)
\end{enumerate}

Take \(n=100\). Its Zeckendorf form is
\(100=F_{11}+F_{6}+F_{4}=89+8+3\). Then

\[
\sum F_{k_i+1}=F_{12}+F_{7}+F_{5}=144+13+5=162.
\]

And
\(\lfloor \varphi\cdot 100+\varphi^{-1}\rfloor=\lfloor 161.803\ldots+0.618\ldots\rfloor=162\),
as required. Similarly,

\[
\sum F_{k_i-1}=F_{10}+F_{5}+F_{3}=55+5+2=62,
\]

and
\(\lfloor \varphi^{-1}\cdot 100+\varphi^{-1}\rfloor=\lfloor 61.803\ldots+0.618\ldots\rfloor=62\).

\subsection{42 {[}M26{]}}\label{m26-4}

We prove uniqueness and existence.

\begin{enumerate}
\def\labelenumi{\arabic{enumi})}
\tightlist
\item
  A key identity (``shift'')
\end{enumerate}

Assume such a representation exists. Using the standard Fibonacci
addition law

\[
F_{u+v}=F_{u-1}F_v+F_uF_{v+1}\quad(u,v\in\mathbb Z),
\]

we get, for every integer \(N\),

\[
\sum_{i=1}^r F_{k_i+N}
=\sum_{i=1}^r\!\big(F_{k_i-1}F_N+F_{k_i}F_{N+1}\big)
=\underbrace{\Big(\sum F_{k_i-1}\Big)}_{=\,n-m}\!F_N
+\underbrace{\Big(\sum F_{k_i}\Big)}_{=\,m}\!F_{N+1}.
\]

Since \(F_{N+1}-F_N=F_{N-1}\), this simplifies to the shift identity

\[
\boxed{\;\sum\nolimits_{i} F_{k_i+N}=mF_{N-1}+nF_{N}\;}\tag{★}
\]

\begin{enumerate}
\def\labelenumi{\arabic{enumi})}
\setcounter{enumi}{1}
\tightlist
\item
  Uniqueness
\end{enumerate}

Fix \(N\) large enough that the right side of (★) is positive. Because
the indices \(k_i\) differ by at least \(2\), the sum
\(\sum_i F_{k_i+N}\) is a Zeckendorf sum (no adjacent indices).
Zeckendorf's theorem (Exercise 34) says such a representation is unique;
therefore the set \(\{k_i\}\) is unique. 

\emph{(This is exactly the book's hint: if two different index-sets
existed, (★) would give two different Zeckendorf representations of the
same integer for large \(N\), contradicting uniqueness in Ex. 34.)} 

\begin{enumerate}
\def\labelenumi{\arabic{enumi})}
\setcounter{enumi}{2}
\tightlist
\item
  Existence (two clean routes)
\end{enumerate}

Route A (via base-\(\varphi\) digits --- short and elegant). Consider
the real number \(x:=m+n\varphi\) (\(\varphi=\tfrac{1+\sqrt5}2\)). By
Exercise 35, every nonnegative real has a unique base-\(\varphi\)
expansion with digits \(0/1\) and no adjacent 1's:

\[
x=\sum_j \varphi^{t_j}\quad\text{with }t_1>t_2>\cdots,\;t_j-t_{j+1}\ge2.
\]

Use \(\varphi^k=F_k\varphi+F_{k-1}\) to rewrite

\[
m+n\varphi=\sum_j \big(F_{t_j}\varphi+F_{t_j-1}\big)
=\varphi\!\sum_j\!F_{t_j}+\sum_j\!F_{t_j-1}.
\]

Since \(1\) and \(\varphi\) are linearly independent over \(\mathbb Q\),
we must have

\[
n=\sum_j F_{t_j},\qquad m=\sum_j F_{t_j-1}.
\]

Setting \(k_i:=t_i-1\) gives the desired representation with
\(k_1\gg k_2\gg\cdots\). 

Route B (matching the book's official answer). Define
\(T_N:=mF_{N-1}+nF_N\). Using Binet's formula one shows

\[
T_{N+1}= \Big\lfloor \,\varphi\,T_N+\varphi^{-1}\,\Big\rfloor
\quad\text{for all sufficiently large }N,
\]

because the tiny \(\varphi'\)-term
(\(\varphi'=1-\varphi=-\varphi^{-1}\)) becomes \(<\varphi^{-2}\) in
magnitude. Exercise 41 identifies this rounding as the operation that
shifts a Zeckendorf sum ``up by one index,'' hence for large \(N\) there
is a set of nonadjacent indices with \(\;T_N=\sum_i F_{k_i+N}\). By (★),
this holds for all \(N\), and taking \(N=0,1\) yields the required
\(m,n\).

Worked example

Take \(k\)-set \(\{7,4,2\}\) (note \(7\gg4\gg2\)). Then

\[
m=F_7+F_4+F_2=13+3+1=17,\quad
n=F_8+F_5+F_3=21+5+2=28.
\]

Check (★) with \(N=3\):

\[
mF_2+nF_3=17\cdot1+28\cdot2=73
\quad\text{and}\quad
\sum_i F_{k_i+3}=F_{10}+F_7+F_5=55+13+5=73.
\]

So \((m,n)=(17,28)\) indeed arises from \(\{7,4,2\}\), and the set is
unique.

Remark (extension to signed pairs)

The same reasoning extends: a joint representation (with the same
nonadjacency rule) exists for possibly negative \(m,n\) iff
\(m+\varphi n\ge0\). Necessity is clear from
\(\;m+\varphi n=\sum_i \varphi^{k_i+1}\ge0\); sufficiency follows by the
base-\(\varphi\) route above.

\bookmarksetup{startatroot}

\chapter{1.2.9 Generating functions}\label{generating-functions}

\subsection{1 {[}M12{]}}\label{m12}

\begin{enumerate}
\def\labelenumi{\arabic{enumi})}
\tightlist
\item
  Ordinary generating function (OGF)
\end{enumerate}

Define

\[
G(z)\;=\;\sum_{n\ge 0} a_n z^n
\;=\;\sum_{n\ge 0} (2^n+3^n) z^n
\;=\;\sum_{n\ge 0} (2z)^n \;+\; \sum_{n\ge 0} (3z)^n.
\]

Each sum is geometric:

\[
\sum_{n\ge 0} (cz)^n \;=\; \frac{1}{1-cz}\quad (|z|<1/|c|).
\]

Hence

\[
\boxed{\;G(z)=\frac{1}{1-2z}+\frac{1}{1-3z}\;}
\]

with radius of convergence \(|z|<\min\{1/2,\,1/3\}=1/3\).

If you prefer a single rational function, combine:

\[
\frac{1}{1-2z}+\frac{1}{1-3z}
=\frac{(1-3z)+(1-2z)}{(1-2z)(1-3z)}
=\boxed{\;\frac{2-5z}{(1-2z)(1-3z)}=\frac{2-5z}{1-5z+6z^2}\; }.
\]

\begin{enumerate}
\def\labelenumi{\arabic{enumi})}
\setcounter{enumi}{1}
\tightlist
\item
  Quick checks
\end{enumerate}

\begin{itemize}
\item
  Constant term: \(G(0)=1+1=2=a_0\).
\item
  Coefficient of \(z\): from the two series it's \(2+3=5=a_1\).
\item
  Coefficient of \(z^2\): \(4+9=13=a_2\).
\item
  The denominator \((1-2z)(1-3z)=1-5z+6z^2\) implies the recurrence

  \[
  a_n=5a_{n-1}-6a_{n-2}\quad (n\ge 2),\qquad a_0=2,\;a_1=5,
  \]

  which the sequence \(2^n+3^n\) indeed satisfies.
\end{itemize}

\begin{enumerate}
\def\labelenumi{\arabic{enumi})}
\setcounter{enumi}{2}
\tightlist
\item
  (Bonus) Exponential generating function (EGF)
\end{enumerate}

If the exercise is interpreted as asking for the EGF,

\[
\sum_{n\ge 0}\frac{a_n}{n!}z^n
=\sum_{n\ge 0}\frac{2^n}{n!}z^n+\sum_{n\ge 0}\frac{3^n}{n!}z^n
=\boxed{\;e^{2z}+e^{3z}\; }.
\]

\subsection{2 {[}M13{]}}\label{m13-2}

In §1.2.9, Knuth defines a binomial (a.k.a. Cauchy--binomial)
convolution

\[
c_n \;=\; \sum_{k=0}^{n} \binom{n}{k}\, a_k\, b_{n-k}\qquad\text{(this is Eq. (10) in the book)}
\]

and then states Eq. (11): If we use exponential generating functions
(EGFs)

\[
A(z)=\sum_{n\ge0}\frac{a_n}{n!}z^n,\quad
B(z)=\sum_{n\ge0}\frac{b_n}{n!}z^n,\quad
C(z)=\sum_{n\ge0}\frac{c_n}{n!}z^n,
\]

then

\[
\boxed{\,C(z)=A(z)\,B(z)\,}.
\]

Algebraic proof (formal power series)

Multiply the two EGFs and collect coefficients:

\[
\begin{aligned}
A(z)B(z)
&=\Bigl(\sum_{k\ge0}\frac{a_k}{k!}z^k\Bigr)\Bigl(\sum_{\ell\ge0}\frac{b_\ell}{\ell!}z^\ell\Bigr)\\
&=\sum_{n\ge0}\Bigl(\sum_{k=0}^{n}\frac{a_k}{k!}\frac{b_{n-k}}{(n-k)!}\Bigr)z^n\\
&=\sum_{n\ge0}\Bigl(\sum_{k=0}^{n}\binom{n}{k}a_k b_{n-k}\Bigr)\frac{z^n}{n!}
=\sum_{n\ge0}\frac{c_n}{n!}z^n
\;=\;C(z).
\end{aligned}
\]

The key identity is \(\displaystyle \binom{n}{k}=\frac{n!}{k!(n-k)!}\).
This is precisely Eq. (11).
(\href{https://stackoverflow.com/questions/768948/how-does-division-work-in-mix}{Scribd})

Combinatorial proof (why EGFs multiply)

Interpret \(a_n\) and \(b_n\) as counting labeled structures of size
\(n\) from classes \(\mathcal A\) and \(\mathcal B\). To build a labeled
\(\mathcal C\)-structure of total size \(n\) as an ordered pair
\((A,B)\) with disjoint label sets:

\begin{itemize}
\tightlist
\item
  choose a \(k\)-subset of the \(n\) labels for \(A\) ---
  \(\binom{n}{k}\) ways;
\item
  place an \(\mathcal A\)-structure on those \(k\) labels --- \(a_k\)
  ways;
\item
  place a \(\mathcal B\)-structure on the remaining \(n-k\) labels ---
  \(b_{n-k}\) ways.
\end{itemize}

Summing over \(k\) gives \(c_n=\sum_{k}\binom{n}{k}a_k b_{n-k}\). The
labeled product construction corresponds exactly to multiplication of
EGFs, hence \(C(z)=A(z)B(z)\).

Quick checks

\begin{enumerate}
\def\labelenumi{\arabic{enumi}.}
\item
  \(a_n\equiv1\), \(b_n\equiv1\). Then \(c_n=\sum_k\binom{n}{k}=2^n\).
  EGFs: \(A(z)=B(z)=e^{z}\Rightarrow C(z)=e^{2z}\), whose \(n\)th
  coefficient is \(2^n/n!\). Works.
\item
  \(a_n=n!\), \(b_n\equiv1\). Then \(A(z)=\sum z^n=1/(1-z)\),
  \(B(z)=e^{z}\). Product \(C(z)=\frac{e^{z}}{1-z}\) has coefficients
  \(c_n/n!=\sum_{k=0}^{n}\frac{1}{(n-k)!}\), i.e.
  \(c_n=n!\sum_{j=0}^{n}\frac{1}{j!}=\sum_{k=0}^{n}\binom{n}{k}k!\),
  exactly the convolution.
\end{enumerate}

\subsection{3 {[}HM21{]}}\label{hm21-5}

Step 1 --- Start from (18)

The ordinary generating function for the harmonic numbers (with H₀ = 0)
is given as

\[
G(z)\;=\;\sum_{n\ge 0} H_n z^n \;=\; \frac{1}{1-z}\,\ln\!\frac{1}{1-z}.
\]

Step 2 --- Differentiate \(G(z)\)

Differentiate termwise and also differentiate the closed form:

\[
G'(z)\;=\;\sum_{n\ge 1} n\,H_n\,z^{\,n-1}
\]

and, using the product/chain rule,

\[
\frac{d}{dz}\!\Big(\frac{1}{1-z}\ln\frac{1}{1-z}\Big)
=\frac{1}{(1-z)^2}+\frac{1}{(1-z)^2}\ln\frac{1}{1-z}
=\frac{1}{(1-z)^2}\;+\;\frac{G(z)}{1-z}.
\]

Hence

\[
\sum_{n\ge 1} n\,H_n\,z^{\,n-1}
\;=\;\frac{1}{(1-z)^2}\;+\;\frac{G(z)}{1-z}.
\tag{★}
\]

Step 3 --- Recognize the prefix--sum generating function

For any sequence \(a_n\) with generating function
\(A(z)=\sum_{n\ge 0}a_n z^n\), the generating function of its prefix
sums \(s_n=\sum_{k=0}^n a_k\) is

\[
S(z)\;=\;\sum_{n\ge 0} s_n z^n\;=\;\frac{A(z)}{1-z}.
\]

(See Eq. (7) in the text.) Applying this with \(a_n=H_n\) gives

\[
S(z)\;=\;\frac{G(z)}{1-z}
\;=\;\frac{1}{(1-z)^2}\,\ln\!\frac{1}{1-z}.
\]

Step 4 --- Compare coefficients

Rewrite (★) by shifting the index on the left:

\[
\sum_{n\ge 0} (n+1)\,H_{n+1}\,z^{\,n}
\;=\;\frac{1}{(1-z)^2}\;+\;S(z).
\]

Extract coefficients of \(z^n\). Since \([z^n](1-z)^{-2}=n+1\), we
obtain

\[
(n+1)\,H_{n+1}\;=\;(n+1)\;+\;\sum_{k=0}^n H_k.
\]

Therefore,

\[
\boxed{\;\sum_{k=0}^n H_k \;=\; (n+1)\bigl(H_{n+1}-1\bigr)\;}.
\]

Equivalently (since \(H_0=0\)),

\[
\sum_{k=1}^n H_k = (n+1)H_{n+1}-(n+1).
\]

This is exactly the relation requested, derived by differentiating (18)
and comparing with the prefix--sum generating function from (7).

\subsection{4 {[}M01{]}}\label{m01}

The two identities (from §1.2.9)

\begin{itemize}
\tightlist
\item
  Binomial theorem:
\end{itemize}

\[
(1+z)^r=\sum_{k\ge 0}\binom{r}{k}z^k. \tag{19}
\]

\begin{itemize}
\tightlist
\item
  Generalized ``binomial'' expansion (with a parameter \(t\)):
\end{itemize}

\[
x^{r}=\sum_{k\ge 0}\binom{r-kt}{k}\frac{r}{\,r-kt\,}\,z^{k},
\quad\text{where } x^{t+1}=x^{t}+z\text{ and }x=1\text{ at }z=0. \tag{21}
\]

(Equivalent first terms: \(x^r=1+rz+\frac{r(r-2t-1)}{2}z^2+\cdots\).)

Take the special case \(t=0\) in (21).

\begin{enumerate}
\def\labelenumi{\arabic{enumi}.}
\tightlist
\item
  The defining relation for \(x\) becomes
\end{enumerate}

\[
x^{t+1}=x^{t}+z \quad\Rightarrow\quad x^{1}=x^{0}+z \quad\Rightarrow\quad x=1+z.
\]

Therefore the left-hand side of (21) is \(x^r=(1+z)^r\). 

\begin{enumerate}
\def\labelenumi{\arabic{enumi}.}
\setcounter{enumi}{1}
\tightlist
\item
  The coefficients on the right-hand side simplify to
\end{enumerate}

\[
\binom{r-kt}{k}\frac{r}{r-kt}\Big|_{t=0}
=\binom{r}{k}\cdot \frac{r}{r}
=\binom{r}{k}.
\]

Thus (21) reduces exactly to

\[
(1+z)^r=\sum_{k\ge 0}\binom{r}{k}z^k,
\]

which is Eq. (19). 

\begin{enumerate}
\def\labelenumi{\arabic{enumi}.}
\setcounter{enumi}{2}
\tightlist
\item
  (Sanity check on low-order terms.) In (21) the quadratic coefficient
  is \(\frac{r(r-2t-1)}{2}\); setting \(t=0\) gives
  \(\frac{r(r-1)}{2}\), which matches the binomial expansion's
  \(z^2\)-term. 
\end{enumerate}

Therefore, Eq. (19) is the \(t=0\) specialization of Eq. (21).

\subsection{5 {[}M20{]}}\label{m20-27}

Base case

For \(n=0\): \((e^{z}-1)^0/0!=1\). On the right,
\(\Bigl\{\!{k\atop 0}\!\Bigr\}=0\) for \(k>0\) and \(=1\) for \(k=0\),
so \(\sum_{k\ge 0}\Bigl\{\!{k\atop 0}\!\Bigr\} z^{k}/k!=1\). Hence the
identity holds.

(You can also check \(n=1\): \((e^{z}-1) = \sum_{k\ge 1} z^{k}/k!\)
since \(\Bigl\{\!{k\atop 1}\!\Bigr\}=1\) for \(k\ge 1\).)

Induction step

Assume the formula holds for some \(n\ge 0\):

\[
\frac{(e^{z}-1)^n}{n!}\;=\;\sum_{k\ge 0}\Bigl\{\!{k\atop n}\!\Bigr\}\frac{z^{k}}{k!}.
\]

Multiply both sides by \(\dfrac{e^{z}-1}{n+1}\):

\[
\frac{(e^{z}-1)^{n+1}}{(n+1)!}
= \frac{e^{z}-1}{n+1}\,\sum_{k\ge 0}\Bigl\{\!{k\atop n}\!\Bigr\}\frac{z^{k}}{k!}.
\]

Use the Cauchy product with \(e^{z}=\sum_{m\ge 0}\frac{z^{m}}{m!}\). The
coefficient of \(z^{r}/r!\) on the right equals

\[
\frac{1}{n+1}\sum_{m=1}^{r}\binom{r}{m}\,\Bigl\{\!{\,r-m\,\atop n}\!\Bigr\}
=\frac{1}{n+1}\sum_{j=0}^{r-1}\binom{r}{j}\,\Bigl\{\!{\,j\,\atop n}\!\Bigr\}
\quad(r\ge 0).
\]

Therefore

\[
\frac{(e^{z}-1)^{n+1}}{(n+1)!}
=\sum_{r\ge 0}\left[\frac{1}{n+1}\sum_{j=0}^{r-1}\binom{r}{j}\,\Bigl\{\!{\,j\,\atop n}\!\Bigr\}\right]\frac{z^{r}}{r!}.
\]

It remains to identify the bracketed coefficient with
\(\Bigl\{\!{r\atop n+1}\!\Bigr\}\). This is the standard ``marked
block'' recurrence for Stirling numbers of the second kind:

Claim. For \(r\ge 0\) and \(n\ge 0\),

\[
\Bigl\{\!{r\atop n+1}\!\Bigr\}
=\frac{1}{n+1}\sum_{j=0}^{r-1}\binom{r}{j}\,\Bigl\{\!{\,j\,\atop n}\!\Bigr\}.
\]

Combinatorial proof of the claim. The left-hand side counts partitions
of an \(r\)-set into \(n+1\) nonempty blocks. Multiply by \(n+1\) to
mark (distinguish) one block: \((n+1)\Bigl\{\!{r\atop n+1}\!\Bigr\}\)
counts ``a partition together with a chosen (marked) block''. To
construct such an object, choose the set of elements not in the marked
block---say \(j\) elements (there are \(\binom{r}{j}\) ways), then
partition those \(j\) elements into \(n\) blocks (there are
\(\Bigl\{\!{j\atop n}\!\Bigr\}\) ways), and put all remaining
\(r-j\ge 1\) elements into the marked block. Summing over
\(j=0,\dots,r-1\) yields

\[
(n+1)\Bigl\{\!{r\atop n+1}\!\Bigr\}=\sum_{j=0}^{r-1}\binom{r}{j}\,\Bigl\{\!{j\atop n}\!\Bigr\},
\]

which is exactly the claim after dividing by \(n+1\).

With the claim, the coefficient of \(z^{r}/r!\) above is
\(\Bigl\{\!{r\atop n+1}\!\Bigr\}\), hence

\[
\frac{(e^{z}-1)^{n+1}}{(n+1)!}
=\sum_{r\ge 0}\Bigl\{\!{r\atop n+1}\!\Bigr\}\frac{z^{r}}{r!},
\]

completing the induction. Therefore Eq. (23) holds for all \(n\ge 0\).

\subsection{6 {[}HM15{]}}\label{hm15-3}

Let \(a_n=\displaystyle\sum_{k=1}^{n-1}\frac{1}{k(n-k)}\) (empty sum
\(a_1=0\)). Find the ordinary generating function
\(A(z)=\sum_{n\ge1} a_n z^n\). Then differentiate \(A(z)\) and express
its coefficients via harmonic numbers \(H_m=\sum_{j=1}^m \frac1j\).

\begin{enumerate}
\def\labelenumi{\arabic{enumi})}
\tightlist
\item
  Closed form for \(a_n\)
\end{enumerate}

Use the partial fraction identity

\[
\frac{1}{k(n-k)}=\frac1n\!\left(\frac1k+\frac1{\,n-k\,}\right).
\]

Hence

\[
a_n=\frac1n\sum_{k=1}^{n-1}\!\Bigl(\frac1k+\frac1{n-k}\Bigr)
=\frac{2}{n}\sum_{k=1}^{n-1}\frac1k
=\frac{2}{n}\,H_{n-1}.
\]

\begin{enumerate}
\def\labelenumi{\arabic{enumi})}
\setcounter{enumi}{1}
\tightlist
\item
  Generating function
\end{enumerate}

We use two standard facts from the section:

\begin{itemize}
\tightlist
\item
  Integration trick for OGFs:
\end{itemize}

\[
\int_0^z\!\!\Bigl(\sum_{m\ge0} b_m t^m\Bigr)dt = \sum_{n\ge1} \frac{b_{n-1}}{n}\, z^n . \quad \text{(Eq. 15)}
\]

\begin{itemize}
\tightlist
\item
  Generating function of harmonic numbers:
\end{itemize}

\[
\sum_{m\ge0} H_m z^m=\frac{1}{1-z}\ln\!\frac{1}{1-z}. \quad \text{(Eq. 18)}
\]

Since \(a_n=\dfrac{2}{n}H_{n-1}\), we get

\[
A(z)=\sum_{n\ge1}\frac{2H_{n-1}}{n}z^n
=2\int_0^z \Bigl(\sum_{m\ge0} H_m t^m\Bigr)\,dt
=2\int_0^z \frac{\ln\!\frac{1}{1-t}}{1-t}\,dt.
\]

Let \(u(t)=\ln\!\frac{1}{1-t}\); then \(u'(t)=\frac{1}{1-t}\). The
integral is \(\int u\,u'\,dt=\tfrac12 u^2\). With the factor 2,

\[
\boxed{\,A(z)=\Bigl(\ln\!\frac{1}{1-z}\Bigr)^{\!2}\,}\qquad(|z|<1).
\]

\begin{enumerate}
\def\labelenumi{\arabic{enumi})}
\setcounter{enumi}{2}
\tightlist
\item
  Differentiate and read coefficients
\end{enumerate}

Differentiate:

\[
A'(z)=2\frac{\ln\!\frac{1}{1-z}}{1-z}.
\]

Using the harmonic-number generating function again,

\[
\frac{\ln\!\frac{1}{1-z}}{1-z}=\sum_{m\ge0} H_m z^m,
\]

so

\[
A'(z)=2\sum_{m\ge0} H_m z^m.
\]

But \(A'(z)=\sum_{n\ge1} n a_n z^{n-1}\). Matching coefficients of
\(z^{n-1}\) gives

\[
n a_n = 2 H_{n-1}\quad\Longrightarrow\quad 
\boxed{\,a_n=\dfrac{2}{n}\,H_{n-1}\,}.
\]

Quick check

\(n=2:\ a_2=1\). \(n=3:\ a_3=\tfrac12+\tfrac12=1\).
\(n=4:\ a_4=\tfrac13+\tfrac14+\tfrac13=\tfrac{11}{12}\). Our formula
yields \(2H_{1}/2=1\), \(2H_{2}/3=1\), \(2H_{3}/4=\tfrac{11}{12}\), all
consistent.

Answer: Generating function \(A(z)=\left(\ln\!\frac{1}{1-z}\right)^2\);
upon differentiation,
\(A'(z)=2\,\dfrac{\ln\!\frac{1}{1-z}}{1-z}=2\sum_{m\ge0}H_m z^m\), hence
\(a_n=\dfrac{2}{n}H_{n-1}\).

\subsection{7 {[}15{]}}\label{section-97}

In §1.2.9, we have

\[
\ln G(z)=\sum_{k\ge1}\frac{S_k\,z^k}{k}\quad\text{with }S_k=\sum_{j=1}^n x_j^k,
\]

and Knuth defines \(G(z)\) by the reciprocal product and arrives at
(37). From this, (38) is obtained by exponentiating and expanding.

Step-by-step verification to (38)

\begin{enumerate}
\def\labelenumi{\arabic{enumi}.}
\tightlist
\item
  Exponentiate the logarithm. Because \(G(0)=1\) and \(\ln G(z)\) is a
  formal power series with zero constant term, the formal identities
  \(\exp(\ln G(z))=G(z)\) and \(\ln(\exp H(z))=H(z)\) are valid. Hence
\end{enumerate}

\[
G(z)=\exp\!\Bigl(\sum_{k\ge1}\frac{S_k\,z^k}{k}\Bigr).
\]

This is exactly the first equality in the displayed chain preceding
(38). 

\begin{enumerate}
\def\labelenumi{\arabic{enumi}.}
\setcounter{enumi}{1}
\tightlist
\item
  Turn the exponential of a sum into a product. In the ring of formal
  power series (which is commutative), exponentials distribute over
  sums:
\end{enumerate}

\[
\exp\!\Bigl(\sum_{k\ge1} A_k\Bigr)=\prod_{k\ge1}\exp(A_k).
\]

Taking \(A_k=\frac{S_k z^k}{k}\) gives

\[
G(z)=\prod_{k\ge1}\exp\!\Bigl(\frac{S_k z^k}{k}\Bigr).
\]

This is the product in the second equality just before (38). 

\begin{enumerate}
\def\labelenumi{\arabic{enumi}.}
\setcounter{enumi}{2}
\tightlist
\item
  Expand each factor as a series. Using \(\exp(u)=\sum_{r\ge0}u^r/r!\),
\end{enumerate}

\[
\exp\!\Bigl(\frac{S_k z^k}{k}\Bigr)=\sum_{r\ge0}\frac{1}{r!}\Bigl(\frac{S_k z^k}{k}\Bigr)^r
=\sum_{r\ge0}\frac{S_k^r}{k^r\,r!}\,z^{kr}.
\]

Knuth writes out the first few factors (for \(k=1,2,\dots\)) in exactly
this way (see the line with
\((1+S_1 z + S_1^2 z^2/2!+\cdots)(1+S_2 z^2/2+\cdots)\cdots\)). 

\begin{enumerate}
\def\labelenumi{\arabic{enumi}.}
\setcounter{enumi}{3}
\tightlist
\item
  Multiply the infinite product and collect coefficients. The
  coefficient of \(z^m\) in the product is obtained by choosing, for
  each \(k\ge1\), a term \(z^{k r_k}\) from the \(k\)-th factor, with
  nonnegative integers \(r_k\), such that
\end{enumerate}

\[
\sum_{k\ge1} k\,r_k=m.
\]

The contribution of a fixed selection \((r_1,r_2,\dots)\) is

\[
\prod_{k\ge1}\frac{S_k^{\,r_k}}{k^{\,r_k}\,r_k!}.
\]

Relabel \(r_k\) as \(k_k\) (Knuth's notation) and truncate at \(k=m\)
since higher \(k\) cannot contribute to \(z^m\). Therefore

\[
[z^m]\,G(z)=\sum_{\substack{k_1,k_2,\dots,k_m\ge0\\
k_1+2k_2+\cdots+ m k_m=m}}
\frac{S_1^{k_1}}{1^{k_1}k_1!}\,
\frac{S_2^{k_2}}{2^{k_2}k_2!}\cdots
\frac{S_m^{k_m}}{m^{k_m}k_m!}.
\]

This is precisely the parenthesized quantity in (38), which Knuth
denotes by \(h_m\), so

\[
G(z)=\sum_{m\ge0} h_m z^m,
\]

with \(h_m\) given by the inner sum. 

\begin{enumerate}
\def\labelenumi{\arabic{enumi}.}
\setcounter{enumi}{4}
\tightlist
\item
  Interpretation via partitions. The constraint
  \(k_1+2k_2+\cdots+m k_m=m\) specifies a partition of \(m\): it says
  that the part \(j\) appears \(k_j\) times. Hence the number of
  summands for fixed \(m\) is \(p(m)\), the number of partitions of
  \(m\) (as the text states just below (38)). 
\end{enumerate}

Quick sanity checks (small \(m\))

Using the formula above:

\begin{itemize}
\item
  \(m=0\): only the all-zero choice ⇒ \(h_0=1\).
\item
  \(m=1\): \((k_1)= (1)\) ⇒ \(h_1=S_1\).
\item
  \(m=2\): tuples \((k_1,k_2)=(2,0),(0,1)\) give

  \[
  h_2=\frac{S_1^2}{2!}+\frac{S_2}{2}.
  \]
\item
  \(m=3\): tuples \((3,0,0),(1,1,0),(0,0,1)\) give

  \[
  h_3=\frac{S_1^3}{6}+\frac{S_1S_2}{2}+\frac{S_3}{3}.
  \]
\end{itemize}

These match the structure immediately implied by the product display
before (38).

\subsection{8 {[}M23{]}}\label{m23-16}

Let \(G(z)=\sum_{n\ge 0} p(n)\,z^n\) be the ordinary generating function
for the partition numbers.

A (unrestricted) partition of \(n\) allows any number of 1's, any number
of 2's, any number of 3's, etc. For each part size \(k\ge 1\), the
contribution to the exponent of \(z\) is a multiple of \(k\):
\(0, k, 2k, 3k,\dots\). Thus the generating series for how many \(k\)'s
we choose is

\[
1+z^{k}+z^{2k}+z^{3k}+\cdots=\frac{1}{1-z^{k}}\qquad(|z|<1).
\]

Different part sizes are chosen independently, so we multiply these
series over all \(k\ge 1\) (Cauchy product). Therefore

\[
G(z)\;=\;\prod_{k=1}^{\infty}\frac{1}{1-z^{k}}.
\]

Expanding this infinite product, the coefficient of \(z^{n}\) counts
exactly the number of solutions \(\sum_{k\ge 1} k\,m_k=n\) with integers
\(m_k\ge 0\), i.e., the number of partitions of \(n\). This is the
standard ``product-for-multisets'' principle that also underlies Knuth's
derivation \(G(z)=\prod_{j}(1-x_j z)^{-1}\) for multisets; here we
specialize to \(x_j=z^{j-1}\) (conceptually) and let the number of
factors go to infinity.

So the answer is:

\[
\boxed{\displaystyle \sum_{n\ge 0} p(n)\,z^{n}\;=\;\prod_{k=1}^{\infty}\frac{1}{1-z^{k}}.}
\]

Quick check

Expanding a few terms gives

\[
\prod_{k\ge 1}\frac{1}{1-z^{k}}
=1+z+2z^{2}+3z^{3}+5z^{4}+7z^{5}+11z^{6}+\cdots,
\]

so \(p(0)=1, p(1)=1, p(2)=2, p(3)=3, p(4)=5, p(5)=7, p(6)=11\), as
expected.

\subsection{9 {[}M11{]}}\label{m11}

With \(h_m\) and the power sums \(S_j\) defined by

\[
h_m=\sum_{1\le j_1\le\cdots\le j_m\le n} x_{j_1}\cdots x_{j_m},\qquad
S_j=\sum_{k=1}^n x_k^{\,j},
\]

and with the generating function \(G(z)=\sum_{m\ge0} h_m z^m\), find
\(h_4\) in terms of \(S_1,S_2,S_3,S_4\).

From the text we have

\[
\ln G(z)=\sum_{k\ge1}\frac{S_k}{k}\,z^k,
\]

hence

\[
G(z)=\exp\!\Big(S_1 z+\frac{S_2}{2}z^2+\frac{S_3}{3}z^3+\frac{S_4}{4}z^4+\cdots\Big).
\]

Therefore \(h_m=[z^m]G(z)=[z^m]\exp A(z)\) where
\(A(z)=a_1z+a_2z^2+a_3z^3+a_4z^4+\cdots\) with
\(a_1=S_1,\ a_2=S_2/2,\ a_3=S_3/3,\ a_4=S_4/4\). (The formula for
\(\ln G\) is Eq. (37) on the same page.) 

Expand \(e^{A}\) up to degree 4:

\[
e^{A}=1+A+\frac{A^2}{2!}+\frac{A^3}{3!}+\frac{A^4}{4!}+O(z^5).
\]

Collect the \(z^4\)-terms:

\begin{itemize}
\tightlist
\item
  From \(A\): \(a_4\).
\item
  From \(A^2/2!\):
  \(\frac12\big(2a_1a_3 + a_2^2\big)=a_1a_3+\frac12 a_2^2\).
\item
  From \(A^3/3!\): \(\frac16\cdot 3 a_1^2 a_2=\frac12 a_1^2 a_2\).
\item
  From \(A^4/4!\): \(\frac1{24} a_1^4\).
\end{itemize}

Thus

\[
h_4 \;=\; a_4 + a_1a_3 + \tfrac12 a_2^2 + \tfrac12 a_1^2 a_2 + \tfrac1{24} a_1^4.
\]

Substitute \(a_1=S_1,\ a_2=S_2/2,\ a_3=S_3/3,\ a_4=S_4/4\):

\[
\boxed{\;
h_4
= \frac{S_4}{4} \;+\; \frac{S_1S_3}{3} \;+\; \frac{S_2^2}{8} \;+\; \frac{S_1^2S_2}{4} \;+\; \frac{S_1^4}{24}
\;}
\]

A convenient equivalent form (multiply by 24) is

\[
\boxed{\;24\,h_4 \;=\; S_1^4 + 6S_1^2S_2 + 3S_2^2 + 8S_1S_3 + 6S_4\;}
\]

Quick sanity checks

\begin{itemize}
\tightlist
\item
  If \(n=1\) (single variable \(x\)), then \(h_4=x^4\) and \(S_j=x^j\);
  plugging into the formula gives \(x^4\).
\item
  The expression matches the general pattern for complete homogeneous
  symmetric polynomials derived from
  \(G(z)=\prod_{i=1}^n(1-x_i z)^{-1}\).
\end{itemize}

\subsection{10 {[}M25{]}}\label{m25-18}

\begin{enumerate}
\def\labelenumi{\arabic{enumi})}
\tightlist
\item
  Generating function
\end{enumerate}

Because a term \(x_{j_1}\cdots x_{j_m}\) picks at most one factor from
each \((1+x_i z)\), we have

\[
E(z)\;\stackrel{\text{def}}=\;\sum_{m\ge0} e_m z^m=\prod_{i=1}^n (1+x_i z).
\]

Taking logs and expanding \(\ln(1+x_i z)\) gives

\[
\ln E(z)=\sum_{i=1}^n\sum_{k\ge1}\frac{(-1)^{k+1}}{k}\,x_i^k z^k
=\sum_{k\ge1}\frac{(-1)^{k+1}}{k}\,S_k\,z^k,
\]

hence the equivalent exponential form

\[
\boxed{\;E(z)=\exp\!\left(\sum_{k\ge1}\frac{(-1)^{k+1}}{k}\,S_k\,z^k\right)\;}
\]

(contrast with the complete symmetric case \(\prod_i(1-x_i z)^{-1}\)
leading to \(\exp(\sum S_k z^k/k)\)).

\begin{enumerate}
\def\labelenumi{\arabic{enumi})}
\setcounter{enumi}{1}
\tightlist
\item
  Coefficient extraction (expression in terms of \(S_k\))
\end{enumerate}

From the exponential form,

\[
E(z)=\exp\!\Big(\sum_{k\ge1} c_k z^k\Big),\qquad c_k=\frac{(-1)^{k+1}}{k}S_k.
\]

Using the standard exponential--generating expansion,

\[
e_m=[z^m]E(z)
=\sum_{\substack{k_1+2k_2+\cdots+mk_m=m\\k_j\ge0}}
\prod_{j=1}^m \frac{c_j^{\,k_j}}{k_j!},
\]

we obtain the compact closed form

\[
\boxed{\;
e_m=\sum_{k_1+2k_2+\cdots+mk_m=m}
(-1)^{\,m+\sum_j k_j}\;
\prod_{j=1}^m \frac{S_j^{\,k_j}}{j^{\,k_j}\,k_j!}
\;}
\]

(because \(\sum_j j k_j=m\)). Equivalently, the Newton identities follow
immediately by differentiating \(\ln E(z)\):

\[
\boxed{\;m\,e_m=\sum_{k=1}^{m}(-1)^{k-1} e_{m-k}\,S_k\qquad(m\ge1).\;}
\]

\begin{enumerate}
\def\labelenumi{\arabic{enumi})}
\setcounter{enumi}{2}
\tightlist
\item
  Low--degree cases
\end{enumerate}

Applying either formula (or Newton identities with \(e_0=1\)):

\[
\begin{aligned}
e_1&=S_1,\\[2pt]
e_2&=\frac{S_1^2-S_2}{2},\\[2pt]
e_3&=\frac{S_1^3-3S_1S_2+2S_3}{6}
\;=\;\frac{S_1^3}{6}-\frac{S_1S_2}{2}+\frac{S_3}{3},\\[2pt]
e_4&=\frac{S_1^4-6S_1^2S_2+3S_2^2+8S_1S_3-6S_4}{24}
\\[-2pt]&\hspace{3.8em}
=\;\frac{S_1^4}{24}-\frac{S_1^2S_2}{4}+\frac{S_2^2}{8}+\frac{S_1S_3}{3}-\frac{S_4}{4}.
\end{aligned}
\]

These match the requested expression of \(e_m\) via the power sums
\(S_k\), and the generating function agrees with the product definition
above. 

\emph{Side note:} If \(H(z)=\sum_{m\ge0} h_m z^m=\prod_i(1-x_i z)^{-1}\)
is the complete symmetric G.F., then \(E(z)H(-z)=1\); this is another
quick route to the Newton identities.

\subsection{11 {[}M25{]}}\label{m25-19}

Use the recurrence

\[
h_n=\frac1n\,(S_1h_{n-1}+S_2h_{n-2}+\cdots+S_nh_0)
\]

to express each power sum \(S_m\) in terms of the complete homogeneous
symmetric polynomials \(h_1,h_2,\dots\), and determine the coefficient
of the monomial \(h_1^{k_1}h_2^{k_2}\cdots h_m^{k_m}\) in \(S_m\) when
\(\;k_1+2k_2+\cdots+mk_m=m\).

Set-up with generating functions

Let

\[
G(z)=\sum_{r\ge0}h_r z^r=1+h_1z+h_2z^2+\cdots .
\]

From the section text we have the fundamental identity

\[
\ln G(z)=\sum_{r\ge1}\frac{S_r}{r}\,z^r .
\]

(Equivalently, \(S_r/r=[z^r]\ln G(z)\).) 

Now write \(G(z)=1+H(z)\) with \(H(z)=\sum_{j\ge1}h_j z^j\). Using the
logarithm series

\[
\ln(1+X)=\sum_{K\ge1}\frac{(-1)^{K-1}}{K}X^K,
\]

we obtain

\[
\ln G(z)=\sum_{K\ge1}\frac{(-1)^{\,K-1}}{K}\,H(z)^K .
\]

Extracting the coefficient

Expand \(H(z)^K\) by the multinomial theorem. A term

\[
h_1^{k_1}h_2^{k_2}\cdots h_m^{k_m}\;z^{\,1\cdot k_1+2\cdot k_2+\cdots+m\cdot k_m}
\]

appears exactly when the integers \(k_1,\dots,k_m\ge0\) satisfy

\[
k_1+2k_2+\cdots+mk_m=m,\qquad K:=k_1+\cdots+k_m .
\]

Its coefficient in \(H(z)^K\) is
\(\displaystyle \frac{K!}{k_1!\cdots k_m!}\). Therefore its coefficient
in \(\ln G(z)\) is

\[
\frac{(-1)^{K-1}}{K}\cdot \frac{K!}{k_1!\cdots k_m!}
\;=\;
(-1)^{K-1}\frac{(K-1)!}{k_1!\cdots k_m!}.
\]

But in \(\ln G(z)=\sum_{r\ge1} (S_r/r) z^r\), the coefficient of \(z^m\)
equals \(S_m/m\). Hence the coefficient of \(h_1^{k_1}\cdots h_m^{k_m}\)
in \(S_m\) is \(m\) times the number above.

Final formula

If \(k_1,\dots,k_m\ge0\) satisfy \(k_1+2k_2+\cdots+mk_m=m\) and
\(K=k_1+\cdots+k_m\), then the coefficient of the monomial
\(h_1^{k_1}h_2^{k_2}\cdots h_m^{k_m}\) in \(S_m\) is

\[
\boxed{\;(-1)^{K-1}\; m\;\frac{(K-1)!}{k_1!\,k_2!\cdots k_m!}\;}
\]

(and it is \(0\) if the weighted sum condition fails). This follows
directly from \(\ln G(z)=\sum_{r\ge1}S_r z^r/r\) with
\(G(z)=\sum_{r\ge0}h_r z^r\).

Quick checks against the small cases in the text

\begin{itemize}
\tightlist
\item
  \(m=1\): only \(k_1=1\) (\(K=1\)) → coefficient \(=1\); so
  \(S_1=h_1\).
\item
  \(m=2\): possibilities \((k_2=1)\) and \((k_1=2)\). Coefficients:
  \(2\cdot(+)(0)!/1!=2\) for \(h_2\), and \(2\cdot(-)1!/2!=-1\) for
  \(h_1^2\); hence \(S_2=2h_2-h_1^2\).
\item
  \(m=3\): \((k_3=1)\), \((k_1=1,k_2=1)\), \((k_1=3)\) give
  \(3h_3-3h_1h_2+h_1^3\).
\end{itemize}

\subsection{12 {[}M20{]}}\label{m20-28}

General setup (bivariate OGFs)

For a double sequence \(\{a_{m,n}\}_{m,n\ge 0}\), its (ordinary)
bivariate generating function is

\[
F(w,z)\;=\;\sum_{m\ge 0}\sum_{n\ge 0} a_{m,n}\, w^m z^n.
\]

This is a \emph{single} formal power series in the two indeterminates
\(w\) and \(z\). Coefficient extraction works as usual:
\([w^m z^n]\,F(w,z)=a_{m,n}\). (Analytically, one may require
\(|(1+w)z|<1\) below; as a \emph{formal} series no convergence
assumption is needed.)

Apply to \(a_{m,n}=\binom{n}{m}\)

Start from the definition above and sum first over \(m\) for each fixed
\(n\):

\[
\sum_{m\ge 0}\binom{n}{m} w^m
= (1+w)^n
\quad(\text{binomial theorem}).
\]

Therefore

\[
F(w,z)\;=\;\sum_{n\ge 0}\Bigl(\sum_{m\ge 0}\binom{n}{m} w^m\Bigr) z^n
\;=\;\sum_{n\ge 0}\bigl((1+w)z\bigr)^n
\;=\;\frac{1}{1-(1+w)z}
\;=\;\frac{1}{1-z-wz}.
\]

Thus the required generating function is

\[
\boxed{\,F(w,z)=\dfrac{1}{1-(1+w)z}\,}
\quad\text{(equivalently } \frac{1}{1-z-wz}\text{).}
\]

You can verify correctness by expanding
\(1/(1-(1+w)z)=\sum_{n\ge 0}((1+w)z)^n\) and then expanding \((1+w)^n\);
the coefficient of \(w^m z^n\) is \(\binom{n}{m}\), and it vanishes
automatically for \(m>n\).

\subsection{13 {[}HM22{]}}\label{hm22-3}

Given a sequence \(a_0,a_1,a_2,\dots\) with ordinary generating function
\(G(z)=\sum_{k\ge0}a_k z^k\), define the step function

\[
f(x)=\sum_{k\ge 0} a_k\,[0\le k\le x]
\]

(the sum of the first \(\lfloor x\rfloor+1\) terms). Express the Laplace
transform \(L\{f\}(s)=\int_0^\infty e^{-st}f(t)\,dt\) in terms of \(G\).

\begin{enumerate}
\def\labelenumi{\arabic{enumi}.}
\tightlist
\item
  Write \(f\) using indicators and swap sum/integral For \(t\ge0\),
\end{enumerate}

\[
f(t)=\sum_{k\ge0} a_k\,\mathbf{1}_{\{t\ge k\}}.
\]

Hence

\[
L\{f\}(s)\;=\;\int_{0}^{\infty} e^{-st}\!\left(\sum_{k\ge0} a_k\,\mathbf{1}_{\{t\ge k\}}\right)dt
\;=\;\sum_{k\ge0} a_k \int_{k}^{\infty} e^{-st}\,dt,
\]

where interchanging sum and integral is justified when \(\Re(s)>0\) and
\(|e^{-s}|<R\) with \(R\) the radius of convergence of \(G\) (then
\(\sum |a_k|e^{-\Re(s)k}\) converges absolutely, so Tonelli/Fubini
applies).

\begin{enumerate}
\def\labelenumi{\arabic{enumi}.}
\setcounter{enumi}{1}
\tightlist
\item
  Compute the inner integral For \(\Re(s)>0\),
\end{enumerate}

\[
\int_{k}^{\infty} e^{-st}\,dt=\frac{e^{-sk}}{s}.
\]

\begin{enumerate}
\def\labelenumi{\arabic{enumi}.}
\setcounter{enumi}{2}
\tightlist
\item
  Recognize the generating function Therefore
\end{enumerate}

\[
L\{f\}(s)=\frac{1}{s}\sum_{k\ge0} a_k\,e^{-sk}
=\frac{1}{s}\,G\!\left(e^{-s}\right).
\]

Final formula and domain

\[
\boxed{\displaystyle L\{f\}(s)=\frac{1}{s}\,G\!\left(e^{-s}\right)}
\qquad\text{valid for}\quad \Re(s)>0\ \text{and}\ |e^{-s}|<R,
\]

i.e., \(\Re(s)>\max\{0,-\ln R\}\), so that both the Laplace integral and
the substituted series converge.

Quick check

\begin{itemize}
\item
  If \(a_k\equiv 1\), then \(G(z)=\frac{1}{1-z}\). Our formula gives

  \[
  L\{f\}(s)=\frac{1}{s}\cdot\frac{1}{1-e^{-s}},
  \]

  which is the classical Laplace transform of the staircase function
  \(f(t)=\lfloor t\rfloor+1\).
\end{itemize}

\subsection{14 {[}HM21{]}}\label{hm21-6}

The identity states that, for \(G(z)=\sum_{n\ge 0} a_n z^n\), a fixed
modulus \(m\ge 1\), and \(\omega=e^{2\pi i/m}\),

\[
\sum_{\substack{n\ge 0\\ n\equiv r\ (\mathrm{mod}\ m)}} a_n z^n
\;=\;
\frac1m\sum_{k=0}^{m-1}\omega^{-kr}\,G(\omega^{k}z)
\qquad(0\le r<m).
\]

Proof (step by step)

\begin{enumerate}
\def\labelenumi{\arabic{enumi}.}
\tightlist
\item
  Expand the right-hand side using the power series for \(G\):
\end{enumerate}

\[
\frac1m\sum_{k=0}^{m-1}\omega^{-kr}G(\omega^{k}z)
=\frac1m\sum_{k=0}^{m-1}\omega^{-kr}\sum_{n\ge 0} a_n(\omega^{k}z)^n
=\sum_{n\ge 0} a_n z^n\left(\frac1m\sum_{k=0}^{m-1}\omega^{k(n-r)}\right).
\]

\begin{enumerate}
\def\labelenumi{\arabic{enumi}.}
\setcounter{enumi}{1}
\tightlist
\item
  Use the roots-of-unity orthogonality:
\end{enumerate}

\[
\frac1m\sum_{k=0}^{m-1}\omega^{k(n-r)}=
\begin{cases}
1,& n\equiv r\pmod m,\\[2pt]
0,& n\not\equiv r\pmod m.
\end{cases}
\]

Indeed, if \(t=n-r\), then
\(\sum_{k=0}^{m-1}\omega^{kt}=\frac{1-\omega^{mt}}{1-\omega^{t}}=0\)
unless \(m\mid t\); when \(m\mid t\), every term is \(1\) and the sum is
\(m\).

\begin{enumerate}
\def\labelenumi{\arabic{enumi}.}
\setcounter{enumi}{2}
\tightlist
\item
  Substituting this filter back into step (1) leaves exactly the terms
  with indices \(n\equiv r\pmod m\):
\end{enumerate}

\[
\sum_{n\ge 0} a_n z^n\cdot [n\equiv r]
\;=\;
\sum_{\substack{n\ge 0\\ n\equiv r\ (\mathrm{mod}\ m)}} a_n z^n,
\]

which is the desired identity.

No convergence hypotheses are needed if we interpret all series as
formal power series, in the spirit of the section.

Quick check

For \(m=3\) and \(r=1\) we get

\[
a_1 z + a_4 z^4 + a_7 z^7+\cdots
=\frac13\Big(G(z)+\omega^{-1}G(\omega z)+\omega^{-2}G(\omega^2 z)\Big),
\quad \omega=e^{2\pi i/3},
\]

exactly as illustrated beneath Eq. (13).

\subsection{15 {[}M28{]}}\label{m28-4}

Let

\[
G_n(z)\;=\;\sum_{k=0}^{n}\binom{n-k}{k}\,z^k
\quad\Big(=\;\sum_{k=0}^{n}\binom{2k-n-1}{k}(-z)^k\Big).
\]

By introducing the ordinary generating function
\(H(w)=\sum_{n\ge0} G_n(z)\,w^n\), find a closed form for \(G_n(z)\).

\begin{enumerate}
\def\labelenumi{\arabic{enumi})}
\tightlist
\item
  Recurrence and bivariate generating function
\end{enumerate}

A direct ``two-step'' decomposition of the summand shows that the
polynomials \(G_n\) satisfy

\[
G_0(z)=1,\quad G_1(z)=1,\quad
G_{n}(z)=G_{n-1}(z)+z\,G_{n-2}(z)\quad(n\ge2).
\]

(Combinatorially: choose whether the first item is a ``1-step'' or a
``2-step'' weighted by \(z\). Algebraically it follows by the standard
binomial identity
\(\binom{n-k}{k}=\binom{n-1-k}{k}+\binom{n-1-k}{k-1}\).)

Form the ordinary generating function in \(n\),

\[
H(w)\;=\;\sum_{n\ge0}G_n(z)\,w^n.
\]

Multiply the recurrence by \(w^n\) and sum over \(n\ge2\):

\[
H(w)-G_0(z)-G_1(z)w
= w\big(H(w)-G_0(z)\big)+z\,w^2H(w).
\]

With \(G_0=G_1=1\) this simplifies to

\[
(1-w-zw^2)\,H(w)=1
\quad\Longrightarrow\quad
\boxed{\,H(w)=\frac{1}{1-w-zw^2}\,}.
\]

This is the classic OGF of the Fibonacci-type recurrence.

\begin{enumerate}
\def\labelenumi{\arabic{enumi})}
\setcounter{enumi}{1}
\tightlist
\item
  Closed form (``Binet'' formula) for \(G_n(z)\)
\end{enumerate}

Factor the denominator:

\[
1-w-zw^2=(1-\alpha w)(1-\beta w),
\]

where \(\alpha,\beta\) are the roots of \(r^2-r-z=0\), i.e.

\[
\alpha=\frac{1+\sqrt{\,1+4z\,}}{2},\qquad
\beta =\frac{1-\sqrt{\,1+4z\,}}{2},
\qquad
\alpha+\beta=1,\ \alpha\beta=-z.
\]

Partial fractions give

\[
H(w)=\frac{1}{(1-\alpha w)(1-\beta w)}
=\frac{1}{\alpha-\beta}\!\left(\frac{\alpha}{1-\alpha w}-\frac{\beta}{1-\beta w}\right).
\]

Extracting coefficients of \(w^n\) yields the closed form

\[
\boxed{\;
G_n(z)=\frac{\alpha^{\,n+1}-\beta^{\,n+1}}{\alpha-\beta}
=\frac{\displaystyle\left(\frac{1+\sqrt{1+4z}}{2}\right)^{\!n+1}
-\left(\frac{1-\sqrt{1+4z}}{2}\right)^{\!n+1}}
{\sqrt{\,1+4z\,}}\,.
\;}
\]

\begin{enumerate}
\def\labelenumi{\arabic{enumi})}
\setcounter{enumi}{2}
\tightlist
\item
  Quick checks and special cases
\end{enumerate}

\begin{itemize}
\tightlist
\item
  \(n=0\): \(G_0=1\) (matches \((\alpha-\beta)/(\alpha-\beta)=1\)).
\item
  \(n=1\): \(G_1=1\) (gives \(\alpha+\beta=1\)).
\item
  \(n=2\): \(G_2=1+z\) (recurrence \(G_2=G_1+zG_0\)).
\item
  Fibonacci numbers. With \(z=1\), \(G_n(1)=F_{n+1}\) (usual Fibonacci
  with \(F_0=0,F_1=1\)).
\item
  Radius of convergence (as a series in \(w\)). The poles are at
  \(w=1/\alpha\) and \(w=1/\beta\), so
  \(|w|<\min\{|1/\alpha|,|1/\beta|\}\) if you care about analytic
  convergence; as a formal power series, the identity holds
  algebraically.
\end{itemize}

\begin{enumerate}
\def\labelenumi{\arabic{enumi})}
\setcounter{enumi}{3}
\tightlist
\item
  Alternative derivation of \(H(w)\) (double sum)
\end{enumerate}

Starting from the definition,

\[
\begin{aligned}
H(w)
&=\sum_{n\ge0}\sum_{k\ge0}\binom{n-k}{k}z^k w^n
=\sum_{k\ge0}z^k w^{2k}\sum_{m\ge0}\binom{m+k}{k}w^m
\\[2pt]
&=\sum_{k\ge0}\frac{(z w^2)^k}{(1-w)^{k+1}}
=\frac{1}{1-w}\cdot\frac{1}{1-\dfrac{z w^2}{1-w}}
=\frac{1}{1-w-zw^2},
\end{aligned}
\]

using \(\sum_{m\ge0}\binom{m+k}{k}w^m=(1-w)^{-(k+1)}\). This reproduces
the result above and leads to the same closed form for \(G_n(z)\).

\begin{enumerate}
\def\labelenumi{\arabic{enumi})}
\setcounter{enumi}{4}
\tightlist
\item
  On the second binomial form
\end{enumerate}

The exercise also notes

\[
\binom{n-k}{k}
=\binom{2k-n-1}{k}(-1)^k,
\]

which is the standard negative-binomial identity
\(\binom{-m}{k}=(-1)^k\binom{m+k-1}{k}\) with \(m=n-2k+1\). It's just an
algebraic re-expression of the same coefficients; the closed form for
\(G_n(z)\) is unchanged. 

Final answer

\[
\boxed{\;
H(w)=\sum_{n\ge0}G_n(z)w^n=\frac{1}{1-w-zw^2},\qquad
G_n(z)=\frac{\left(\frac{1+\sqrt{1+4z}}{2}\right)^{n+1}-\left(\frac{1-\sqrt{1+4z}}{2}\right)^{n+1}}{\sqrt{1+4z}}.
\;}
\]

\subsection{16 {[}M22{]}}\label{m22-15}

Let \(a_{n,r}(m)\) be the number of solutions of

\[
x_1+x_2+\cdots+x_n=m
\]

with each \(x_i\) an integer satisfying \(0\le x_i\le r\). Find the
ordinary generating function

\[
G_{n,r}(z)\;=\;\sum_{m\ge 0} a_{n,r}(m)\,z^m,
\]

and (optionally) a coefficient formula for \(a_{n,r}(m)\).

\begin{enumerate}
\def\labelenumi{\arabic{enumi})}
\tightlist
\item
  Generating function by the product rule
\end{enumerate}

For a single variable \(x_i\) with \(0\le x_i\le r\), the choices
contribute

\[
1+z+z^2+\cdots+z^r \;=\;\frac{1-z^{\,r+1}}{1-z}.
\]

Since the \(n\) variables are independent and the exponents add, the
ordinary generating function factorizes:

\[
G_{n,r}(z)
\;=\;\bigl(1+z+\cdots+z^r\bigr)^n
\;=\;\left(\frac{1-z^{\,r+1}}{1-z}\right)^{\!n}.
\]

\begin{enumerate}
\def\labelenumi{\arabic{enumi})}
\setcounter{enumi}{1}
\tightlist
\item
  Extracting coefficients (closed form)
\end{enumerate}

To get an explicit formula for \(a_{n,r}(m)=[z^m]G_{n,r}(z)\), use
inclusion--exclusion (or the binomial series for \((1-z)^{-n}\), then
subtract the ``overshoot'' terms where some variable exceeds \(r\)):

\[
\boxed{%
a_{n,r}(m)
\;=\;\sum_{j=0}^{\lfloor m/(r+1)\rfloor}
(-1)^j\binom{n}{j}\binom{m-j(r+1)+n-1}{n-1}
}
\qquad(m\ge 0),
\]

with the convention \(\binom{u}{v}=0\) if \(u<0\) or \(v<0\) or \(v>u\).

Derivation sketch. Write

\[
G_{n,r}(z)
=\left(\frac{1-z^{r+1}}{1-z}\right)^{\!n}
=(1-z^{r+1})^{n}(1-z)^{-n}
=\sum_{j=0}^{n}\binom{n}{j}(-1)^j z^{j(r+1)}\cdot
\sum_{t\ge 0}\binom{n+t-1}{n-1}z^{t}.
\]

Collect the coefficient of \(z^m\) by setting \(t=m-j(r+1)\) and summing
over feasible \(j\).

\begin{enumerate}
\def\labelenumi{\arabic{enumi})}
\setcounter{enumi}{2}
\tightlist
\item
  Sanity checks and special cases
\end{enumerate}

\begin{itemize}
\tightlist
\item
  \(r=\infty\): \(G_{n,\infty}(z)=(1-z)^{-n}\), giving the
  stars-and-bars count
  \(\displaystyle a_{n,\infty}(m)=\binom{m+n-1}{n-1}\).
\item
  \(r=1\): binary variables. Then \(G_{n,1}(z)=(1+z)^n\) and
  \(a_{n,1}(m)=\binom{n}{m}\) for \(0\le m\le n\).
\item
  Symmetry: coefficients of \((1+z+\cdots+z^r)^n\) are symmetric in the
  sense \(a_{n,r}(m)=a_{n,r}(nr-m)\) (replace each \(x_i\) by
  \(r-x_i\)).
\end{itemize}

\begin{enumerate}
\def\labelenumi{\arabic{enumi})}
\setcounter{enumi}{3}
\tightlist
\item
  Small example
\end{enumerate}

Take \(n=3, r=2\). Then

\[
G_{3,2}(z)=(1+z+z^2)^3.
\]

Expanding (or using the formula) gives

\[
a_{3,2}(m)=
1,\,3,\,6,\,7,\,6,\,3,\,1\quad\text{for }m=0,1,2,3,4,5,6,
\]

which counts the bounded solutions to \(x_1+x_2+x_3=m\) with
\(0\le x_i\le 2\).

\subsection{17 {[}M25{]}}\label{m25-20}

Show that expanding \((1-z)^{-w}\) in two different ways yields the
bivariate generating function of the (unsigned) Stirling numbers of the
first kind \(\bigl[{k\atop n}\bigr]\):

\[
(1-z)^{-w}
=\sum_{k\ge 0}\binom{-w}{k}(-z)^k
=\sum_{n\ge 0}\sum_{k\ge n}\frac{\bigl[{k\atop n}\bigr]}{k!}\,w^{\,n}z^{\,k},
\]

and hence

\[
\bigl(\ln\tfrac1{1-z}\bigr)^{\!n}
= n!\sum_{k\ge n}\frac{\bigl[{k\atop n}\bigr]}{k!}\,z^{\,k}\qquad(n\ge 0).
\]

Step 1 --- Binomial expansion with a general top argument

Knuth's generalized binomial coefficient is defined for any \(x\) by
\(\displaystyle \binom{x}{k}=x^{\underline{k}}/k!\) (with falling
factorial \(x^{\underline{k}}\)); in particular,

\[
(1-z)^{-w}=\sum_{k\ge 0}\binom{-w}{k}(-z)^k
\quad\text{(formal power series).}
\]

Moreover
\(\binom{-w}{k}(-1)^k=\dfrac{w(w+1)\cdots(w+k-1)}{k!}=\dfrac{w^{\overline{k}}}{k!}\)
(rising factorial). Hence

\[
(1-z)^{-w}=\sum_{k\ge 0}\frac{w^{\overline{k}}}{k!}\,z^{\,k}.
\]

(For the negative-integer special case compare 1.2.9--(20); for the
general binomial definition see the index to notations.)

Step 2 --- Express the rising factorial via Stirling numbers of the
first kind

By the standard conversion between factorial and ordinary powers, the
(unsigned) Stirling numbers of the first kind expand the rising
factorial:

\[
w^{\overline{k}}=\sum_{n=0}^{k}\bigl[{k\atop n}\bigr]\,w^{\,n}.
\]

(Equivalently,
\(x^{\underline{n}}=\sum_{k}( -1)^{n-k}\bigl[{n\atop k}\bigr]x^{k}\);
passing from falling to rising removes the signs.) See the
Stirling-number discussion and identities in §1.2.6. 

Substitute into Step 1 and interchange sums:

\[
(1-z)^{-w}
=\sum_{k\ge 0}\frac{z^{k}}{k!}\sum_{n=0}^{k}\bigl[{k\atop n}\bigr]w^{\,n}
=\sum_{n\ge 0}w^{\,n}\sum_{k\ge n}\frac{\bigl[{k\atop n}\bigr]}{k!}\,z^{\,k}.
\tag{★}
\]

Step 3 --- Expand the same function as an exponential of a logarithm

Also

\[
(1-z)^{-w}=e^{\,w\ln(1/(1-z))}=\sum_{n\ge 0}\frac{w^{\,n}}{n!}\,\Bigl(\ln\tfrac1{1-z}\Bigr)^{\!n}.
\]

Comparing coefficients of \(w^{\,n}\) in this and in (★) gives, for each
\(n\ge 0\),

\[
\sum_{k\ge n}\frac{\bigl[{k\atop n}\bigr]}{k!}\,z^{\,k}
=\frac{1}{n!}\Bigl(\ln\tfrac1{1-z}\Bigr)^{\!n},
\]

which is exactly the desired identity.

Quick sanity checks

\begin{itemize}
\tightlist
\item
  \(n=1\). Since \(\bigl[{k\atop 1}\bigr]=(k-1)!\),
\end{itemize}

\[
\sum_{k\ge1}\frac{\bigl[{k\atop1}\bigr]}{k!}z^{k}=\sum_{k\ge1}\frac{1}{k}z^{k}
=\ln\tfrac1{1-z},
\]

agreeing with the formula for \(n=1\).

\begin{itemize}
\tightlist
\item
  \(n=2\). Using \(\bigl[{k\atop2}\bigr]=(k-1)!\,H_{k-1}\),
\end{itemize}

\[
2\sum_{k\ge2}\frac{\bigl[{k\atop2}\bigr]}{k!}z^{k}
=2\sum_{k\ge2}\frac{H_{k-1}}{k}\,z^{k}
=\Bigl(\ln\tfrac1{1-z}\Bigr)^{\!2},
\]

matching the \(n=2\) case.

\subsection{18 {[}M25{]}}\label{m25-21}

\begin{enumerate}
\def\labelenumi{(\alph{enumi})}
\tightlist
\item
  Distinct indices
\end{enumerate}

Consider

\[
E_n(z)\;=\;\prod_{k=1}^n (1+kz).
\]

Expanding, every term arises by choosing either \(1\) or \(kz\) from
each factor; picking \(r\) of the \(kz\) terms and the rest as \(1\)
yields a contribution \((k_1\cdots k_r)z^r\) with
\(1\le k_1<\cdots<k_r\le n\). Therefore the desired sum is exactly the
coefficient \([z^r]E_n(z)\).

Now rewrite

\[
E_n(z)
= z^{\,n+1}\!\!\prod_{k=1}^n\Big(\tfrac1z+k\Big)
= \sum_{t=0}^{n+1} \big[{\,n+1 \atop t}\big]\, z^{\,n+1-t},
\]

using the standard expansion of a rising factorial into powers with
Stirling numbers of the first kind. Hence

\[
\boxed{\displaystyle
\sum_{1\le k_1<\cdots<k_r\le n} k_1\cdots k_r
\;=\; \big[{\,n+1 \atop n+1-r}\big] }.
\]

(Equivalently, the sum is the elementary symmetric polynomial
\(e_r(1,2,\ldots,n)\).)

\begin{enumerate}
\def\labelenumi{(\alph{enumi})}
\setcounter{enumi}{1}
\tightlist
\item
  Repetition allowed
\end{enumerate}

Now take

\[
H_n(z)\;=\;\prod_{k=1}^n \frac{1}{1-kz}.
\]

Each factor \((1-kz)^{-1}=\sum_{m\ge0}(kz)^m\), so in the product the
coefficient of \(z^r\) is the sum of all \(k_1\cdots k_r\) over weakly
increasing \(k_1\le\cdots\le k_r\) with values in \(\{1,\dots,n\}\).
Thus the desired sum is \([z^r]H_n(z)\).

But there is a classical identity

\[
\prod_{k=1}^n \frac{1}{1-kz}
=\sum_{t\ge n} \big\{{t \atop n}\big\} z^{\,t-n},
\]

so the coefficient of \(z^r\) is \(\big\{{n+r \atop n}\big\}\).
Therefore

\[
\boxed{\displaystyle
\sum_{1\le k_1\le\cdots\le k_r\le n} k_1\cdots k_r
\;=\; \big\{{\,n+r \atop n}\big\} }.
\]

(Equivalently, this is the complete homogeneous symmetric polynomial
\(h_r(1,2,\ldots,n)\).)

Quick check

For \(n=3, r=2\):

\begin{itemize}
\tightlist
\item
  \begin{enumerate}
  \def\labelenumi{(\alph{enumi})}
  \tightlist
  \item
    \(1\cdot2+1\cdot3+2\cdot3=11=\big[{4\atop2}\big]\).
  \end{enumerate}
\item
  \begin{enumerate}
  \def\labelenumi{(\alph{enumi})}
  \setcounter{enumi}{1}
  \tightlist
  \item
    \(1\cdot1+1\cdot2+1\cdot3+2\cdot2+2\cdot3+3\cdot3=25=\big\{{5\atop3}\big\}\).
  \end{enumerate}
\end{itemize}

\subsection{19 {[}HM32{]}}\label{hm32}

For integers \(p,q\) with \(0<p<q\), prove

\[
H_{p/q}
= \frac{q}{p}\;-\;\frac{\pi}{2}\cot\!\Bigl(\frac{\pi p}{q}\Bigr)\;-\;\ln(2q)
\;+\;2\sum_{0<k<q/2}\cos\!\Bigl(\frac{2\pi pk}{q}\Bigr)\,\ln\!\Bigl(\sin\frac{\pi k}{q}\Bigr).
\tag{★}
\]

We use the hint: by Abel's theorem,

\[
H_{p/q}
=\lim_{x\to 1^-}\;\sum_{n\ge1}\Bigl(\frac1n-\frac{1}{n+p/q}\Bigr)x^{\,p+nq},
\tag{1}
\]

and ``Eq. (13)'' of the section (the basic series for the logarithm) to
rewrite the power series so the limit can be taken.

\begin{enumerate}
\def\labelenumi{\arabic{enumi})}
\tightlist
\item
  Prepare two pieces
\end{enumerate}

Split the sum in (1) into

\[
S(x)=A(x)-B(x),\qquad
A(x):=\sum_{n\ge1}\frac{x^{p+nq}}{n},\qquad
B(x):=\sum_{n\ge1}\frac{x^{p+nq}}{n+p/q}.
\]

Piece \(A\). Using \(\sum_{n\ge1}\frac{z^n}{n}=-\ln(1-z)\) (Eq. (13) in
this section),

\[
A(x)=x^{p}\sum_{n\ge1}\frac{(x^q)^{n}}{n}
=-x^{p}\ln(1-x^{q}).
\tag{2}
\]

Piece \(B\). Note \(1/(n+p/q)=q/(qn+p)\). Write the progression
\(p,\,p+q,\,p+2q,\dots\) via a roots-of-unity filter. Let
\(\omega=e^{2\pi i/q}\). A standard discrete Fourier selection gives

\[
q\sum_{m\ge0}\frac{x^{p+mq}}{p+mq}
=-\sum_{k=0}^{q-1}\omega^{-kp}\ln(1-\omega^k x).
\tag{3}
\]

Therefore

\[
B(x)=\sum_{n\ge1}\frac{x^{p+nq}}{n+p/q}
=\;q\sum_{m\ge1}\frac{x^{p+mq}}{p+mq}
=\;-\!\sum_{k=0}^{q-1}\omega^{-kp}\ln(1-\omega^k x)\;-\;\frac{q}{p}x^{p}.
\tag{4}
\]

Combining (2) and (4),

\[
S(x)=\underbrace{\sum_{k=1}^{q-1}\omega^{-kp}\ln(1-\omega^k x)}_{=:f(x)}
\;+\;\underbrace{(1-x^{p})\ln(1-x)\;+\;\frac{q}{p}x^{p}
\;-\;x^{p}\ln\!\frac{1-x^{q}}{\,1-x\,}}_{=:g(x)}.
\tag{5}
\]

\begin{enumerate}
\def\labelenumi{\arabic{enumi})}
\setcounter{enumi}{1}
\tightlist
\item
  The easy limit \(g(x)\)
\end{enumerate}

As \(x\to1^{-}\): \((1-x^{p})\ln(1-x)\to0\) (since
\((1-x)\ln(1-x)\to0\)), \(x^p\to1\), and
\(\displaystyle \frac{1-x^q}{1-x}\to q\). Hence

\[
\lim_{x\to1^-}g(x)=\frac{q}{p}-\ln q.
\tag{6}
\]

\begin{enumerate}
\def\labelenumi{\arabic{enumi})}
\setcounter{enumi}{2}
\tightlist
\item
  The main limit \(f(x)\)
\end{enumerate}

For \(0<x<1\) the \(k\) and \(q-k\) terms in \(f(x)\) are complex
conjugates, so \(f(x)\) is real; it suffices to compute its real part
termwise.

Use the elementary factorization

\[
1-\omega^k
=1-e^{2\pi i k/q}
=2\sin\!\Bigl(\frac{\pi k}{q}\Bigr)\;
\exp\!\Bigl(i\bigl(\tfrac{\pi k}{q}-\tfrac{\pi}{2}\bigr)\Bigr),
\]

valid for \(1\le k\le q-1\). Taking principal logs,

\[
\ln(1-\omega^k)
=\ln\!\Bigl(2\sin\frac{\pi k}{q}\Bigr)
+i\Bigl(\frac{\pi k}{q}-\frac{\pi}{2}\Bigr).
\]

Therefore, letting \(\phi_k:=\tfrac{2\pi pk}{q}\),

\[
\omega^{-kp}\,\ln(1-\omega^k)
=\Bigl[\cos\phi_k\,\ln\!\Bigl(2\sin\frac{\pi k}{q}\Bigr)
-\Bigl(\frac{\pi k}{q}-\frac{\pi}{2}\Bigr)\sin\phi_k\Bigr]
+ i(\cdots),
\]

and summing \(k=1,\dots,q-1\) cancels the imaginary parts. Thus

\[
\lim_{x\to1^-} f(x)
=\sum_{k=1}^{q-1}\cos\phi_k\,\ln\!\Bigl(2\sin\frac{\pi k}{q}\Bigr)
-\sum_{k=1}^{q-1}\Bigl(\frac{\pi k}{q}-\frac{\pi}{2}\Bigr)\sin\phi_k.
\tag{7}
\]

Two standard trigonometric sums now enter:

\begin{itemize}
\tightlist
\item
  \(\displaystyle \sum_{k=1}^{q-1}\cos\frac{2\pi pk}{q}=-1\) (since
  \(\sum_{k=0}^{q-1}e^{2\pi i pk/q}=0\)).
\item
  \(\displaystyle \sum_{k=1}^{q-1}k\sin\frac{2\pi pk}{q}=\frac{q}{2}\cot\frac{\pi p}{q}\).
  (Derive by differentiating \(\sum_{k=0}^{q-1} r^k=(1-r^q)/(1-r)\) at
  \(r=e^{it}\) and taking imaginary parts at \(t=2\pi p/q\).)
\end{itemize}

Using \(\sum_{k=1}^{q-1}\sin\phi_k=0\), the second sum in (7) becomes

\[
-\frac{\pi}{q}\sum_{k=1}^{q-1}k\sin\phi_k
= -\frac{\pi}{2}\cot\!\Bigl(\frac{\pi p}{q}\Bigr).
\]

For the first sum in (7), split \(\ln(2\sin)=\ln 2+\ln(\sin)\):

\[
\sum_{k=1}^{q-1}\cos\phi_k\,\ln\!\Bigl(2\sin\frac{\pi k}{q}\Bigr)
= \underbrace{\Bigl(\sum_{k=1}^{q-1}\cos\phi_k\Bigr)}_{=-1}\ln 2
\;+\;\sum_{k=1}^{q-1}\cos\phi_k\,\ln\!\Bigl(\sin\frac{\pi k}{q}\Bigr).
\]

Pairing \(k\) with \(q-k\) converts the remaining sum to

\[
2\sum_{0<k<q/2}\cos\!\Bigl(\frac{2\pi pk}{q}\Bigr)\,\ln\!\Bigl(\sin\frac{\pi k}{q}\Bigr),
\]

since when \(q\) is even the middle term \(k=q/2\) has \(\sin(\pi/2)=1\)
and contributes \(0\).

Hence

\[
\lim_{x\to1^-} f(x)
= -\ln 2\;-\;\frac{\pi}{2}\cot\!\Bigl(\frac{\pi p}{q}\Bigr)
\;+\;2\sum_{0<k<q/2}\cos\!\Bigl(\frac{2\pi pk}{q}\Bigr)\,\ln\!\Bigl(\sin\frac{\pi k}{q}\Bigr).
\tag{8}
\]

\begin{enumerate}
\def\labelenumi{\arabic{enumi})}
\setcounter{enumi}{3}
\tightlist
\item
  Assemble the result
\end{enumerate}

From (6) and (8),

\[
H_{p/q}
=\Bigl(\frac{q}{p}-\ln q\Bigr)
+\Bigl[-\ln 2-\frac{\pi}{2}\cot\!\Bigl(\frac{\pi p}{q}\Bigr)
+2\sum_{0<k<q/2}\cos\!\Bigl(\frac{2\pi pk}{q}\Bigr)\,\ln\!\Bigl(\sin\frac{\pi k}{q}\Bigr)\Bigr],
\]

which simplifies to \((★)\), as required.

\begin{enumerate}
\def\labelenumi{\arabic{enumi})}
\setcounter{enumi}{4}
\tightlist
\item
  Quick sanity checks (the three series in the statement)
\end{enumerate}

Using \(1-\tfrac12H_{1/2}=\ln 2\),
\(1-\tfrac13H_{1/3}+\tfrac15H_{1/5}-\cdots=\tfrac\pi4\), etc., the
identity reproduces the classical sums displayed in the exercise heading
(they correspond to \(H_{1/2},H_{1/4},H_{1/3}\) in the linear
combinations shown there).

\subsection{20 {[}M21{]}}\label{m21-14}

Step 1. Expand \(n^{m}\) in the falling--factorial basis

A classical identity using Stirling numbers of the second kind
\(S(m,k)\) is

\[
n^{m}=\sum_{k=0}^{m} S(m,k)\, n^{\underline{k}},
\qquad
n^{\underline{k}}:=n(n-1)\cdots(n-k+1)=k!\binom{n}{k}.
\]

(Combinatorially: count functions from an \(m\)-set to an \(n\)-set by
first partitioning the domain into \(k\) nonempty blocks --- \(S(m,k)\)
ways --- then injecting those \(k\) blocks into the \(n\)-set ---
\(n^{\underline{k}}\) ways.)

Hence

\[
n^{m}=\sum_{k=0}^{m} S(m,k)\,k!\binom{n}{k}.
\]

Step 2. The ordinary GF of \(\binom{n}{k}\)

For fixed \(k\ge 0\),

\[
\sum_{n\ge 0}\binom{n}{k} z^{n}
=\sum_{n\ge k}\binom{n}{k} z^{n}
=z^{k}\sum_{r\ge 0}\binom{r+k}{k} z^{r}
=\frac{z^{k}}{(1-z)^{k+1}},
\]

using \(\displaystyle \sum_{r\ge 0}\binom{r+k}{k}z^{r}=(1-z)^{-(k+1)}\).

Multiplying by \(k!\) gives

\[
\sum_{n\ge 0} k!\binom{n}{k} z^{n}=\frac{k!\,z^{k}}{(1-z)^{k+1}}.
\]

Step 3. Assemble the generating function

Using Steps 1--2,

\[
\sum_{n\ge 0} n^{m} z^{n}
=\sum_{n\ge 0}\sum_{k=0}^{m} S(m,k)\,k!\binom{n}{k} z^{n}
=\sum_{k=0}^{m} S(m,k)\,k!\left(\sum_{n\ge 0}\binom{n}{k} z^{n}\right)
=\sum_{k=0}^{m} \frac{S(m,k)\,k!\,z^{k}}{(1-z)^{k+1}}.
\]

Therefore the requested coefficients are

\[
\boxed{\;c_{m,k}=k!\,S(m,k)\;}\qquad(0\le k\le m),
\]

where \(S(m,k)\) are Stirling numbers of the second kind.

Quick checks (small \(m\))

\begin{itemize}
\item
  \(m=0\): \(c_{0,0}=0!\,S(0,0)=1\Rightarrow \sum z^{n}=1/(1-z)\).
\item
  \(m=1\): \(c_{1,1}=1!\,S(1,1)=1\Rightarrow \sum n z^{n}=z/(1-z)^{2}\).
\item
  \(m=2\): \(c_{2,1}=1!\,S(2,1)=1,\; c_{2,2}=2!\,S(2,2)=2\):

  \[
  \sum n^{2}z^{n}=\frac{z}{(1-z)^{2}}+\frac{2 z^{2}}{(1-z)^{3}}
  =\frac{z(1+z)}{(1-z)^{3}},
  \]

  as known.
\end{itemize}

Alternative viewpoint (operator form, optional)

Since
\(\displaystyle \sum_{n\ge 0} n^{m} z^{n}=(z\,\tfrac{d}{dz})^{m}\!\left(\frac{1}{1-z}\right)\),
expanding the power \((z\,\tfrac{d}{dz})^{m}\) in terms of falling
powers of the operator \(z\,\tfrac{d}{dz}\) again introduces \(S(m,k)\)
and leads to the same coefficients \(c_{m,k}=k!\,S(m,k)\).

\subsection{21 {[}HM30{]}}\label{hm30-3}

\begin{enumerate}
\def\labelenumi{\arabic{enumi})}
\tightlist
\item
  Ordinary generating function (OGF)
\end{enumerate}

Let

\[
A(z)\;=\;\sum_{n\ge 0} n!\,z^n.
\]

Radius of convergence

By the root (or ratio) test, \(\sqrt[n]{n!}\sim n/e\), hence
\(\lim_{n\to\infty}\sqrt[n]{n!}=\infty\) and the radius of convergence
is \(0\). Thus \(A(z)\) diverges for every \(z\neq 0\); it is best
regarded as a formal power series. (Coefficient extraction and
manipulation of formal OGFs are standard in §1.2.9. )

Differential equation for \(A\)

The factorials satisfy \(a_{n+1}=(n+1)a_n\) with \(a_0=1\). Use the
standard derivative rule for OGFs,

\[
G'(z)=\sum_{k\ge0}(k+1)a_{k+1}z^k, 
\]

which Knuth gives as Eq. (14). Multiply the recurrence by \(z^{n+1}\)
and sum over \(n\ge0\):

\[
A(z)-1 \;=\; z\sum_{n\ge0}(n+1)a_n z^n
\;=\; z\,(zA(z))' \;=\; zA(z)+z^2A'(z).
\]

Hence the ODE

\[
z^2A'(z)=(1-z)A(z)-1. \tag{★}
\]

This step uses the differentiation identity for generating functions in
the text.

Closed form via an exponential integral

Solve (★) as a first-order linear ODE. Writing it as
\(A' - \frac{1-z}{z^2}A = -\frac{1}{z^2}\), an integrating factor is
\(\mu(z)=ze^{1/z}\). Thus

\[
\frac{d}{dz}\!\bigl(ze^{1/z}A(z)\bigr)
= -\frac{e^{1/z}}{z}.
\]

Let \(t=1/z\). Then \(-\frac{e^{1/z}}{z}\,dz=\frac{e^{t}}{t}\,dt\), so

\[
ze^{1/z}A(z)=\operatorname{Ei}\!\bigl(\tfrac{1}{z}\bigr)+C,
\qquad
A(z)=\frac{e^{-1/z}}{z}\Bigl(\operatorname{Ei}\!\bigl(\tfrac{1}{z}\bigr)+C\Bigr).
\]

Choosing \(C=0\) gives the unique branch whose asymptotic Maclaurin
expansion at \(z=0\) is \(\sum_{n\ge0}n!z^n\). In other words,

\[
\boxed{\;A(z)\;\sim\;\sum_{n\ge0}n!z^n\;=\;\frac{e^{-1/z}}{z}\,\operatorname{Ei}\!\Bigl(\frac{1}{z}\Bigr)\;\text{ (as a divergent asymptotic series)}\;}
\]

(the equality is Borel summation/analytic continuation). This function
has an essential singularity at \(z=0\) and a branch cut coming from
\(\operatorname{Ei}\).

A convenient equivalent integral representation---obtained from
\(n!=\int_0^\infty t^n e^{-t}\,dt\) and summing the geometric
series---is

\[
A(z)\;=\;\int_0^\infty \frac{e^{-t}}{1-zt}\,dt
\;=\;\frac{e^{-1/z}}{z}\,\operatorname{Ei}\!\Bigl(\frac{1}{z}\Bigr),
\]

interpreted as an analytic continuation (the geometric series identity
\(1/(1-u)=\sum_{n\ge0}u^n\) is Eq. (5) in the text). 

Summary for the OGF.

\begin{itemize}
\tightlist
\item
  Formal series: \(A(z)=\sum_{n\ge0}n!z^n\) (radius \(0\)).
\item
  Satisfies \(z^2A'=(1-z)A-1\).
\item
  Borel sum / analytic continuation:
  \(A(z)=\frac{e^{-1/z}}{z}\operatorname{Ei}(1/z)\).
\item
  Asymptotics at \(z\to0\): the coefficients are \(n!\) (by coefficient
  extraction notation \([z^n]\) in §1.2.9).
\end{itemize}

\begin{enumerate}
\def\labelenumi{\arabic{enumi})}
\setcounter{enumi}{1}
\tightlist
\item
  Exponential generating function (EGF)
\end{enumerate}

Knuth also uses exponential series in this section (see Eq. (22)). For
the sequence \(\langle n!\rangle\),

\[
E(z)\;=\;\sum_{n\ge0}\frac{n!}{n!}\,z^n
\;=\;\sum_{n\ge0} z^n
\;=\;\frac{1}{1-z}\quad(|z|<1).
\]

Thus the EGF is the geometric series \(1/(1-z)\) (exactly Eq. (5) in the
text), and it is analytic for \(|z|<1\) with a simple pole at \(z=1\). 

Properties from the EGF.

\begin{itemize}
\tightlist
\item
  Coefficients recover immediately: \([z^n]E(z)=1\) and multiplying back
  by \(n!\) gives \(n!\).
\item
  Differentiation yields moments/convolutions in the usual way for EGFs
  (see the general ``exponential series'' framework).
\end{itemize}

\begin{enumerate}
\def\labelenumi{\arabic{enumi})}
\setcounter{enumi}{2}
\tightlist
\item
  Relation between OGF and EGF (Borel transform)
\end{enumerate}

The OGF can be viewed as the Borel transform of the EGF:

\[
A(z)\;=\;\int_0^\infty e^{-t}\,E(tz)\,dt
\;=\;\int_0^\infty \frac{e^{-t}}{1-tz}\,dt
\;=\;\frac{e^{-1/z}}{z}\operatorname{Ei}\!\Bigl(\frac{1}{z}\Bigr),
\]

which reconciles the divergent OGF with the convergent EGF and explains
the exponential-integral form above (geometric-series identity again).

\subsection{22 {[}M21{]}}\label{m21-15}

\begin{enumerate}
\def\labelenumi{\arabic{enumi})}
\tightlist
\item
  Build the product generating function
\end{enumerate}

For each fixed \(j\ge 0\), the choices of \(k_j\) contribute a factor

\[
\sum_{k_j=0}^{r} \binom{r}{k_j}\, (z^{2^{j}})^{k_j} \;=\; (1+z^{2^{j}})^{r}
\]

by the binomial theorem. Independence across \(j\) means the full
generating function is

\[
G(z)\;=\;\prod_{j\ge 0} (1+z^{2^{j}})^{r}.
\]

By construction, expanding this product and collecting the \(z^n\) term
produces exactly the given sum, so this \(G(z)\) works.

\begin{enumerate}
\def\labelenumi{\arabic{enumi})}
\setcounter{enumi}{1}
\tightlist
\item
  Telescope the product
\end{enumerate}

Use \(1+t=\dfrac{1-t^{2}}{1-t}\) to telescope finite partial products:

\[
\prod_{j=0}^{m}(1+z^{2^{j}})^{r}
= \prod_{j=0}^{m}\Bigl(\frac{1 - z^{2^{j+1}}}{1-z^{2^{j}}}\Bigr)^{r}
= \left(\frac{1 - z^{2^{m+1}}}{1-z}\right)^{\!r}.
\]

Letting \(m\to\infty\) (as a formal power series, or analytically for
\(|z|<1\)) gives

\[
G(z)\;=\;\prod_{j\ge 0}(1+z^{2^{j}})^{r} \;=\; (1-z)^{-r}.
\]

Thus a closed form is \(G(z)=(1-z)^{-r}\). 

\begin{enumerate}
\def\labelenumi{\arabic{enumi})}
\setcounter{enumi}{2}
\tightlist
\item
  Read off the coefficient
\end{enumerate}

From the negative binomial series,

\[
(1-z)^{-r}=\sum_{n\ge 0}\binom{n+r-1}{n}z^{n},
\]

so

\[
[z^{n}]\,G(z)=[z^{n}]\,(1-z)^{-r}=\binom{n+r-1}{n}.
\]

Therefore the original sum equals \(\binom{n+r-1}{n}\), and one valid
answer is

\[
\boxed{\,G(z)=\displaystyle\prod_{j\ge 0}(1+z^{2^{j}})^{r}=(1-z)^{-r}\, .}
\]

\begin{enumerate}
\def\labelenumi{\arabic{enumi})}
\setcounter{enumi}{3}
\tightlist
\item
  (Combinatorial sanity check)
\end{enumerate}

\(\prod_{j\ge 0}(1+z^{2^{j}})=1+z+z^{2}+z^{3}+\cdots=\frac1{1-z}\):
choosing \(0\) or \(1\) copy of each \(2^{j}\) encodes the binary
expansion of any nonnegative integer. Raising to the \(r\)-th power
counts ordered \(r\)-tuples \((n_1,\dots,n_r)\) with
\(n_1+\cdots+n_r=n\); there are \(\binom{n+r-1}{n}\) such tuples
(``stars and bars''), matching the coefficient above.

\subsection{23 {[}M33{]}}\label{m33-1}

\begin{enumerate}
\def\labelenumi{(\alph{enumi})}
\tightlist
\item
  Existence of polynomials \(f_m,g_m\)
\end{enumerate}

Claim. For every integer \(m\ge 1\) there exist polynomials
\(f_m(z_1,\dots,z_m)\) and \(g_m(z_1,\dots,z_m)\) such that the
multivariate binomial-sum identity in the exercise holds for all
integers \(n\ge r\ge 0\). (See the problem statement on p.~96 of the
text.) 

Construction (induction on \(m\)).

\begin{itemize}
\item
  Base \(m=1\). Take

  \[
  f_1(z_1)=z_1,\qquad g_1(z_1)=1+z_1.
  \]

  The required identity reduces to the ordinary binomial theorem, so it
  holds. 
\item
  Inductive step. Suppose \(f_m,g_m\) work for some \(m\). When you
  append a new variable \(z_{m+1}\) and sum over the new index, a single
  application of the binomial theorem shows that you can obtain new
  polynomials \(f_{m+1},g_{m+1}\) by the substitutions

  \[
  z_m\longleftarrow z_m\!\left(1+z_{m+1}^{-1}\right),\quad\text{then multiply overall by }z_{m+1}.
  \]

  Concretely:

  \[
  \begin{aligned}
  f_{m+1}(z_1,\dots,z_{m+1})
  &=z_{m+1}\,f_m\!\bigl(z_1,\dots,z_{m-1},\,z_m(1+z_{m+1}^{-1})\bigr),\\
  g_{m+1}(z_1,\dots,z_{m+1})
  &=z_{m+1}\,g_m\!\bigl(z_1,\dots,z_{m-1},\,z_m(1+z_{m+1}^{-1})\bigr).
  \end{aligned}
  \]

  In particular this implies the Fibonacci-type recurrences

  \[
  f_m=z_m f_{m-1}+z_{m-1}f_{m-2},\qquad
  g_m=z_m g_{m-1}+z_{m-1}g_{m-2},
  \]

  with initial values \(f_{-1}=0,\ f_0=g_{-1}=g_0=1\). Hence \(f_m,g_m\)
  are indeed polynomials and the required identity follows for \(m+1\). 
\end{itemize}

Useful descriptions. From the recurrences one gets:

\begin{itemize}
\item
  \(g_2=z_1+z_2+z_1z_2\).
\item
  \(g_m/f_m\) equals the finite continued fraction

  \[
  \frac{g_m}{f_m}
  =1+\cfrac{z_1^{-1}}{1+\cfrac{z_2^{-1}}{1+\;\cdots\;+\cfrac{z_m^{-1}}{\,1\,}}}\,.
  \]
\item
  Combinatorially, \(g_m\) is the sum of all monomials obtained from
  \(z_1z_2\cdots z_m\) by striking out any set of nonadjacent variables
  (there are \(F_{m+2}\) ways); \(f_m\) is similar but with \(z_1\)
  constrained to remain. 
\end{itemize}

\begin{enumerate}
\def\labelenumi{(\alph{enumi})}
\setcounter{enumi}{1}
\tightlist
\item
  Closed form for the cyclic sum \(S_n\)
\end{enumerate}

Let

\[
S_n(z_1,\dots,z_m)
\]

denote the cyclic binomial sum defined in the exercise (a direct
generalization of Ex. 15). Using part (a) and extracting the coefficient
of \(z_m^n\), we can write

\[
S_n(z_1,\dots,z_m)
=[z_m^n]\sum_{r=0}^{n} z^r\, z_m^{\,n-r}\, f_m^{\,n-r}\,g_m^{\,r}.
\]

Carrying out the coefficient extraction and the finite sums leads to a
compact second-order rational generating function

\[
\sum_{n\ge0} S_n z^n
=\frac{1}{(1-az)(1-dz)-bc\,z^2},
\]

where

\[
a=z_m g_{m-1},\qquad
b=z_{m-1} g_{m-2},\qquad
c=z_m f_{m-1},\qquad
d=z_{m-1} f_{m-2}.
\]

Equivalently, the sequence \((S_n)\) satisfies the linear recurrence

\[
S_n = h_m\,S_{n-1} - (-1)^m z_1 z_2\cdots z_m\, S_{n-2},
\quad n\ge2,
\]

where

\[
h_m=z_m g_{m-1}+z_{m-1} f_{m-2},
\qquad
ad-bc=(-1)^m z_1z_2\cdots z_m.
\]

Solving the quadratic

\[
t^2-(a+d)\,t+(ad-bc)=0,
\]

let \(\rho,\sigma\) be its roots. Then the closed form is

\[
\boxed{\;S_n(z_1,\dots,z_m)=\dfrac{\rho^{\,n+1}-\sigma^{\,n+1}}{\rho-\sigma}\;}
\]

with \(a+d=h_m\) and \(ad-bc=(-1)^m z_1\cdots z_m\) as above. (All steps
are detailed in the book's answer sketch.) 

\emph{A pleasant by-product:} the derivation also yields directly the
two-term recurrence for \(S_n\) displayed above, which would be quite
painful to obtain without generating functions.

\begin{enumerate}
\def\labelenumi{(\alph{enumi})}
\setcounter{enumi}{2}
\tightlist
\item
  Specialization \(z_1=\cdots=z_m=z\)
\end{enumerate}

When all variables are equal, the coefficients in the recurrence become
constant and the characteristic roots simplify beautifully. If

\[
\rho_1=\frac{z+\sqrt{z^2+4z}}{2},\qquad
\sigma_1=\frac{z-\sqrt{z^2+4z}}{2},
\]

then in the \(m\)-variable case the corresponding roots are just the
\(m\)-th powers:

\[
\rho_m=\rho_1^{\,m},\qquad \sigma_m=\sigma_1^{\,m}.
\]

Therefore

\[
\boxed{\;
S_n(\underbrace{z,\dots,z}_{m})
=\frac{\rho_m^{\,n+1}-\sigma_m^{\,n+1}}{\rho_m-\sigma_m}
=\frac{\rho_1^{\,m(n+1)}-\sigma_1^{\,m(n+1)}}{\rho_1^{\,m}-\sigma_1^{\,m}}
\;}
\]

with \(\rho_1,\sigma_1\) as above. 

Checks and small cases

\begin{itemize}
\tightlist
\item
  \(m=1\):
  \(S_n(z)=\frac{\rho_1^{\,n+1}-\sigma_1^{\,n+1}}{\rho_1-\sigma_1}\)
  solves \(S_n=z\,S_{n-1}-z\,S_{n-2}\) with \(S_0=1,\ S_1=1+z\); this is
  the classical companion to the Fibonacci-type sums in Ex. 15.
\item
  \(m=2\): Using \(g_2=z_1+z_2+z_1z_2\) and the formulas above
  reproduces the two-variable generalization and matches the rational
  generating function obtained by standard methods.
\end{itemize}

\subsection{24 {[}M22{]}}\label{m22-16}

Proof of the identity

Use the binomial theorem on the right-hand side:

\[
(1+zG(z))^m=\sum_{k=0}^m \binom{m}{k}\,(zG(z))^k.
\]

Taking coefficients of \(z^n\) and using
\([z^n]\,(z^kH(z))=[z^{\,n-k}]H(z)\),

\[
[z^n](1+zG(z))^m
=\sum_{k=0}^m \binom{m}{k}\,[z^n]\,(zG(z))^k
=\sum_{k=0}^m \binom{m}{k}\,[z^{\,n-k}]\,G(z)^k,
\]

which is exactly the left-hand side.

\begin{enumerate}
\def\labelenumi{(\alph{enumi})}
\tightlist
\item
  \(G(z)=\dfrac{1}{1-z}\)
\end{enumerate}

Right-hand side:

\[
[z^n]\Bigl(1+\frac{z}{1-z}\Bigr)^m
=[z^n]\,(1-z)^{-m}.
\]

Using the standard expansion
\((1-z)^{-(r+1)}=\sum_{j\ge0}\binom{r+j}{r}z^j\) (take \(r=m-1\)), the
coefficient is

\[
[z^n]\,(1-z)^{-m}=\binom{m+n-1}{n}. \quad \text{(negative-binomial series)} \; \, \, \, \, \, \, \, \, \, \, \, \, \, 
\]

Left-hand side (as a check): \(G(z)^k=(1-z)^{-k}\), so

\[
[z^{\,n-k}]\,G(z)^k=[z^{\,n-k}](1-z)^{-k}=\binom{n-1}{k-1}\quad (k\le n),
\]

and

\[
\sum_{k} \binom{m}{k}\binom{n-1}{k-1}
=\sum_{k} \binom{m}{k}\binom{n-1}{n-k}
=\binom{m+n-1}{n}
\]

by Vandermonde's convolution, agreeing with the RHS.

Answer for (a): \(\boxed{\binom{m+n-1}{n}}\).

\begin{enumerate}
\def\labelenumi{(\alph{enumi})}
\setcounter{enumi}{1}
\tightlist
\item
  \(G(z)=\dfrac{e^z-1}{z}\)
\end{enumerate}

Right-hand side:

\[
1+zG(z)=1+(e^z-1)=e^z \;\Rightarrow\; [z^n](e^z)^m=[z^n]e^{mz}=\frac{m^n}{n!}.
\]

Left-hand side (optional cross-check via Stirling numbers):

\[
[z^{\,n-k}]\,G(z)^k=[z^n](e^z-1)^k.
\]

From
\(\dfrac{(e^z-1)^k}{k!}=\sum_{n\ge k} \!\! S(n,k)\,\dfrac{z^n}{n!}\) we
get \([z^n](e^z-1)^k=\dfrac{k!\,S(n,k)}{n!}\). Hence

\[
\sum_k \binom{m}{k}[z^{\,n-k}]G(z)^k
=\frac{1}{n!}\sum_k \binom{m}{k}k!\,S(n,k)
=\frac{1}{n!}\sum_k (m)_k\,S(n,k)
=\frac{m^n}{n!},
\]

using the standard identity \(m^n=\sum_k S(n,k)\,(m)_k\). The RHS
\(m^n/n!\) we already obtained.

Answer for (b): \(\boxed{\dfrac{m^n}{n!}}\).

I will give another purely combinatorial (no generating functions)
proofs for parts (a) and (b).

\begin{enumerate}
\def\labelenumi{(\alph{enumi})}
\tightlist
\item
  \(G(z)=\dfrac{1}{1-z}\) ⇒
  \(\displaystyle [z^n](1-z)^{-m}=\binom{m+n-1}{n}\)
\end{enumerate}

Combinatorial story (stars and bars). Interpret \([z^n](1-z)^{-m}\) as
the number of solutions in nonnegative integers to

\[
x_1+\cdots+x_m=n,\qquad x_i\ge 0,
\]

i.e., distributing \(n\) identical balls into \(m\) labeled boxes. By
the classic stars-and-bars argument, the count is

\[
\binom{m+n-1}{n}.
\]

Same count, decomposed by the number of nonempty boxes. Equivalently,
count the same distributions by the number \(k\) of boxes that are
actually used:

\begin{enumerate}
\def\labelenumi{\arabic{enumi}.}
\tightlist
\item
  Choose which \(k\) of the \(m\) boxes are nonempty: \(\binom{m}{k}\)
  ways.
\item
  Put at least one ball in each chosen box; write \(x_i'=x_i-1\) for
  used boxes. Then \(\sum x_i'=n-k\) with \(x_i'\ge 0\) in exactly those
  \(k\) boxes---equivalently, choose a composition of \(n\) into \(k\)
  positive parts. The number of such \(k\)-tuples is
  \(\binom{n-1}{k-1}\) (stars and bars with positivity).
\end{enumerate}

Summing over \(k\) gives

\[
\sum_{k=0}^{\min(m,n)} \binom{m}{k}\binom{n-1}{k-1}
=\binom{m+n-1}{n},
\]

which is Vandermonde's convolution in disguise. Thus the coefficient
equals \(\boxed{\binom{m+n-1}{n}}\).

\begin{enumerate}
\def\labelenumi{(\alph{enumi})}
\setcounter{enumi}{1}
\tightlist
\item
  \(G(z)=\dfrac{e^z-1}{z}\) ⇒
  \(\displaystyle [z^n]e^{mz}=\frac{m^n}{n!}\)
\end{enumerate}

We'll show the identity

\[
m^n=\sum_{k=0}^{\min(m,n)} \binom{m}{k}\,k!\,S(n,k),
\]

by double-counting the same object, then divide by \(n!\) at the end.

Objects to count. Consider all functions \(f:[n]\to[m]\) (think: color
each of \(n\) distinct balls with one of \(m\) colors). There are
clearly \(m^n\) such functions.

Count the same set by the number of colors used.

\begin{itemize}
\tightlist
\item
  Fix \(k\). First choose the \(k\) colors that actually appear:
  \(\binom{m}{k}\) ways.
\item
  The fibers \(f^{-1}(c)\) over the used colors form a partition of
  \([n]\) into \(k\) nonempty, \emph{unlabeled} blocks. The number of
  such partitions is the Stirling number \(S(n,k)\).
\item
  Finally, assign (label) the \(k\) blocks with the chosen \(k\)
  distinct colors: \(k!\) bijections.
\end{itemize}

Multiplying and summing over \(k\) yields

\[
m^n=\sum_{k}\binom{m}{k}\,S(n,k)\,k!.
\]

Since \([z^n]e^{mz}=\frac{m^n}{n!}\), dividing both sides by \(n!\)
gives

\[
[z^n]e^{mz}
=\sum_{k}\binom{m}{k}\,\frac{k!\,S(n,k)}{n!},
\]

which matches the exercise's evaluation for \(G(z)=\dfrac{e^z-1}{z}\).
Therefore the answer for (b) is \(\boxed{\dfrac{m^n}{n!}}\).

\subsection{25 {[}M23{]}}\label{m23-17}

Evaluate

\[
S_n=\sum_{k=0}^{n}\binom{n}{k}\binom{2n-2k}{\,n-k\,}(-2)^k.
\]

(Equivalently, recognize \(\binom{n}{k}(-2)^k=[w^k](1-2w)^n\) and
\(\binom{2n-2k}{n-k}=[z^{\,n-k}](1+z)^{2n-2k}\).)

Use coefficient extraction to convert the single sum into one
coefficient:

\[
\begin{aligned}
S_n
&=\sum_{k=0}^{n}\binom{n}{k}(-2)^k\,[z^{\,n-k}](1+z)^{2n-2k} \\[2mm]
&=[z^{\,n}]\,\sum_{k=0}^{n}\binom{n}{k}(-2)^k\,z^{k}\,(1+z)^{2n-2k} \qquad\text{(shift: }[z^{n-k}]=[z^n]z^k\text{)}\\[2mm]
&=[z^{\,n}]\,(1+z)^{2n}\sum_{k=0}^{n}\binom{n}{k}\!\left(\frac{-2z}{(1+z)^2}\right)^{\!k}\\[2mm]
&=[z^{\,n}]\,(1+z)^{2n}\left(1-\frac{2z}{(1+z)^2}\right)^{\!n}
=[z^{\,n}]\,(1+z)^{2n}\left(\frac{(1+z)^2-2z}{(1+z)^2}\right)^{\!n}\\[2mm]
&=[z^{\,n}]\,(1+z)^{2n}\left(\frac{1+z^2}{(1+z)^2}\right)^{\!n}
=[z^{\,n}]\,(1+z^2)^{n}.
\end{aligned}
\]

Now expand \((1+z^2)^n=\sum_{j=0}^{n}\binom{n}{j}z^{2j}\). The
coefficient of \(z^n\) is therefore

\begin{itemize}
\tightlist
\item
  \(0\) if \(n\) is odd (no \(j\) with \(2j=n\));
\item
  \(\binom{n}{n/2}\) if \(n\) is even (take \(j=n/2\)).
\end{itemize}

Answer

\[
\boxed{\,S_n=\begin{cases}
\displaystyle \binom{n}{n/2},& n\text{ even},\\[2mm]
0,& n\text{ odd}.
\end{cases}\,}
\]

(Sanity check for small \(n\):
\(S_0=1,\ S_1=0,\ S_2=2,\ S_3=0,\ S_4=6,\ S_6=20\), matching
\(\binom{2m}{m}\) for even \(n=2m\).)

\subsection{26 {[}M40{]}}\label{m40}

Knuth introduces the bracket notation \([z^n]\,G(z)\) to mean ``the
coefficient of \(z^n\) in \(G(z)\)'' (his eq. (31)). Exercise 26 asks
you to extend this notation so that one can write things like

\[
[z^2-2z^5]\,G(z)=a_2-2a_5
\quad\text{when }G(z)=\sum_{n\ge0}a_n z^n.
\]

Definition (generalized bracket)

Let \(H(z)=\sum_{n\in\mathbb Z} h_n z^n\) be a \emph{Laurent} series
with only finitely many \emph{positive} coefficients nonzero
(polynomials are the most common case). For a (formal) power series
\(G(z)=\sum_{n\ge0} a_n z^n\), define

\[
\boxed{\;\;[H(z)]\,G(z)\;:=\;\operatorname{CT}_z\big(H(z)\,G(z^{-1})\big)\;=\;\sum_{n\ge0} h_n\,a_n\;}
\]

where \(\operatorname{CT}_z\) denotes ``constant term in \(z\).'' This
immediately yields \([z^2-2z^5]\,G(z)=a_2-2a_5\).

Equivalently---connecting to Knuth's coefficient-extraction via complex
analysis---you can write the same functional as a contour (residue)
integral:

\[
[H(z)]\,G(z)
=\frac{1}{2\pi i}\oint H(z)\,G(z^{-1})\,\frac{dz}{z},
\]

which generalizes Cauchy's coefficient formula
\([z^n]\,G(z)=\frac{1}{2\pi i}\oint \frac{G(z)}{z^{n+1}}\,dz\). 

Basic properties

Let \(F,G\) be (formal) power series and \(H,K\) Laurent series as
above.

\begin{enumerate}
\def\labelenumi{\arabic{enumi}.}
\item
  Linearity. \([\,\alpha H+\beta K\,]\,G=\alpha[H]\,G+\beta[K]\,G\) and
  \([H]\,(\alpha F+\beta G)=\alpha[H]\,F+\beta[H]\,G\).
\item
  Symmetry (dot product view). \([H]\,G=\sum_{n\ge0}h_n a_n=[G]\,H\). So
  \([\,\cdot\,]\) is the coefficientwise inner product.
\item
  Shift rules. \([z^m H]\,G=\sum_{n\ge0}h_n\,a_{n+m}=[H]\,(z^{-m}G)\).
  Special cases:
\end{enumerate}

\begin{itemize}
\tightlist
\item
  \([z^m]\,G=a_m\) (recovers the original bracket).
\item
  \([1+z+\cdots+z^N]\,G=\sum_{n=0}^N a_n\) (prefix sums).
\item
  \big(\([z^r/(1-z)]\,G=\sum_{n\ge r}a_n\)) (a tail sum; justified via
  the constant-term definition).
\end{itemize}

\begin{enumerate}
\def\labelenumi{\arabic{enumi}.}
\setcounter{enumi}{3}
\tightlist
\item
  Product (Cauchy) interaction. If \(F(z)=\sum f_k z^k\), then
\end{enumerate}

\[
[H]\,(F G)=\sum_{k\ge0} f_k\,\big[\,z^{-k}H\,\big]\,G,
\]

a handy ``convolution-with-shift'' identity derived from
\(\operatorname{CT}\big(H\,F\,G^{-1}\big)\).

\begin{enumerate}
\def\labelenumi{\arabic{enumi}.}
\setcounter{enumi}{4}
\tightlist
\item
  Hadamard viewpoint. Let \(H\odot G:=\sum_{n\ge0} h_n a_n z^n\) be the
  Hadamard (termwise) product. Then
\end{enumerate}

\[
[H]\,G=(H\odot G)(1),
\]

whenever evaluation at \(z=1\) is meaningful (e.g., as a finite sum or
convergent series).

Examples

\begin{itemize}
\tightlist
\item
  Polynomial weights. With \(H(z)=c_0+c_1 z+\cdots+c_m z^m\),
  \([H]\,G=\sum_{k=0}^m c_k a_k\).
\item
  Finite differences. \([1-z]\,G=a_0-a_1\),
  \([1-2z+z^2]\,G=a_0-2a_1+a_2\) (a second difference at the origin).
\item
  Sums via rational \(H\). Using \([\,1/(1-z)\,]\,G=\sum_{n\ge0}a_n\)
  (constant-term proof), you can encode cumulative sums compactly.
\end{itemize}

Why this is natural

\begin{itemize}
\tightlist
\item
  It strictly generalizes \([z^n]\): take \(H(z)=z^n\) to get
  \([z^n]\,G\).
\item
  It turns coefficient extraction into a bilinear
  functional---algebraically convenient and perfectly compatible with
  the generating-function toolkit in §1.2.9 (multiplication, logarithms,
  etc.).
\item
  It fits the residue framework Knuth invokes for coefficient
  extraction, simply replacing \(z^{-n-1}\) by \(H(z) / z\) inside the
  contour integral.
\end{itemize}

\bookmarksetup{startatroot}

\chapter{1.2.10 Analysis of an
Algorithm}\label{analysis-of-an-algorithm}

\subsection{1 {[}10{]}}\label{section-98}

From the section we have the recurrence for the distribution
\(p_{nk}=\Pr(A=k)\):

\[
p_{nk}=\frac{1}{n}\,p_{n-1,k-1}+\frac{n-1}{n}\,p_{n-1,k},
\]

with initial conditions \(p_{1k}=\delta_{0k}\) and \(p_{nk}=0\) for
\(k<0\). 

Set \(k=0\). The term \(p_{n-1,-1}\) vanishes, so

\[
p_{n0}=\frac{n-1}{n}\,p_{n-1,0}.
\]

Unfolding this down to \(n=1\) (where \(p_{10}=1\)) gives a telescoping
product:

\[
p_{n0}=\frac{n-1}{n}\cdot\frac{n-2}{n-1}\cdots\frac{1}{2}\cdot 1=\frac{1}{n}.
\]

So,

\[
\boxed{\,p_{n0}=\dfrac{1}{n}\, }.
\]

Interpretation for Algorithm M

In Algorithm M, \(A\) counts how many times the current maximum \(m\) is
changed while scanning the input. The minimum value \(A=0\) occurs iff
the last item is already the global maximum, i.e.,
\(X[n]=\max_{1\le k\le n} X[k]\). Under the section's assumption that
all \(X[k]\) are distinct and all \(n!\) relative orders are equally
likely, each position is equally likely to hold the maximum; hence
\(\Pr(A=0)=1/n\). 

(Equivalently, since \(G_n(z)=\sum_k p_{nk}z^k\), the result says
\(G_n(0)=p_{n0}=1/n\).)

\subsection{2 {[}HM16{]}}\label{hm16}

Here is a clean derivation of Eq. (13) from Eq. (10).

Let \(G(z)=\sum_{k\ge 0} p_k z^k\) be the probability generating
function, so \(\sum_k p_k = G(1)=1\). By Eq. (10),

\[
\mathrm{mean}(G)=\sum_k k\,p_k,\qquad 
\mathrm{var}(G)=\sum_k k^2 p_k-\Bigl(\sum_k k\,p_k\Bigr)^2. \; \text{[Eq. (10)]}
\]

\begin{enumerate}
\def\labelenumi{\arabic{enumi}.}
\tightlist
\item
  Differentiate \(G\):
\end{enumerate}

\[
G'(z)=\sum_k k\,p_k\,z^{\,k-1}\quad\Longrightarrow\quad G'(1)=\sum_k k\,p_k=\mathrm{mean}(G).
\]

\begin{enumerate}
\def\labelenumi{\arabic{enumi}.}
\setcounter{enumi}{1}
\tightlist
\item
  Differentiate again:
\end{enumerate}

\[
G''(z)=\sum_k k(k-1)\,p_k\,z^{\,k-2}\quad\Longrightarrow\quad G''(1)=\sum_k k(k-1)\,p_k.
\]

\begin{enumerate}
\def\labelenumi{\arabic{enumi}.}
\setcounter{enumi}{2}
\tightlist
\item
  Express \(\sum_k k^2 p_k\) using \(G''(1)\) and \(G'(1)\):
\end{enumerate}

\[
\sum_k k^2 p_k=\sum_k \bigl(k(k-1)+k\bigr)p_k
=\underbrace{\sum_k k(k-1)p_k}_{=\,G''(1)}+\underbrace{\sum_k k\,p_k}_{=\,G'(1)}
=G''(1)+G'(1).
\]

\begin{enumerate}
\def\labelenumi{\arabic{enumi}.}
\setcounter{enumi}{3}
\tightlist
\item
  Substitute into the variance formula:
\end{enumerate}

\[
\mathrm{var}(G)=\bigl(G''(1)+G'(1)\bigr)-\bigl(G'(1)\bigr)^2
=G''(1)+G'(1)-\bigl(G'(1)\bigr)^2,
\]

which is exactly Eq. (13).

\subsection{3 {[}M15{]}}\label{m15-11}

In Algorithm M (find the maximum), step M4 is executed exactly when the
current item sets a new record---i.e., it is larger than all items seen
so far. Thus, if we let \(A_n\) be the number of records in a uniformly
random permutation of \(n\) distinct items, the exercise asks for the
four basic statistics of \(A_{1000}\).

Knuth derives the general result (for any \(n\)) that

\[
\min A_n=0,\qquad \mathbb{E}[A_n]=H_n-1,\qquad 
\max A_n=n-1,\qquad 
\operatorname{sd}(A_n)=\sqrt{\,H_n-H_n^{(2)}\,},
\]

where \(H_n=\sum_{k=1}^n \frac1k\) and
\(H_n^{(2)}=\sum_{k=1}^n \frac1{k^2}\). 

Below I re-derive these quickly, then plug in \(n=1000\).

Quick derivation (indicator variables)

Let \(I_k\) be the indicator that the \(k\)-th item is a new record (so
\(I_k=1\) iff the \(k\)-th item exceeds all of the first \(k-1\) items).
Then

\[
A_n=\sum_{k=2}^n I_k .
\]

By symmetry, \(\Pr(I_k=1)=1/k\) (the maximum of the first \(k\)
positions is equally likely to be at any of them). Hence

\[
\mathbb{E}[A_n]=\sum_{k=2}^n \frac{1}{k}=H_n-1.
\]

Moreover, one checks that the events \(\{I_i=1\}\) and \(\{I_j=1\}\) are
pairwise independent for \(i\neq j\) in a random permutation, so

\[
\operatorname{Var}(A_n)=\sum_{k=2}^n \operatorname{Var}(I_k)
=\sum_{k=2}^n \Bigl(\tfrac{1}{k}-\tfrac{1}{k^2}\Bigr)
= H_n-H_n^{(2)} .
\]

Therefore \(\operatorname{sd}(A_n)=\sqrt{H_n-H_n^{(2)}}\). The extremal
values are clear:

\begin{itemize}
\tightlist
\item
  Min \(0\): occurs when the very first item is the overall maximum.
\item
  Max \(n-1\): occurs for a strictly increasing sequence.
\end{itemize}

Plug in \(n=1000\)

Numerically,

\[
H_{1000}\approx 7.485470860550345,\qquad
H^{(2)}_{1000}\approx 1.643934566681560,
\]

so

\[
\mathbb{E}[A_{1000}]=H_{1000}-1\approx 6.48547086055,
\]

\[
\operatorname{Var}(A_{1000})=H_{1000}-H^{(2)}_{1000}\approx 5.84153629387,
\quad
\operatorname{sd}(A_{1000})\approx 2.41692703528.
\]

\subsection{4 {[}M10{]}}\label{m10-12}

We consider the coin-tossing experiment of Eq. (17), where a head occurs
with probability \(p\) (and tail with \(q=1-p\)), and \(p_{n k}\)
denotes the probability of getting exactly \(k\) heads in \(n\)
independent tosses. The recurrence is

\[
p_{n k}=p\,p_{n-1,\,k-1}+q\,p_{n-1,\,k}.
\]

We need to give a closed form for \(p_{n k}\).

Solution (two complementary derivations)

\begin{enumerate}
\def\labelenumi{\arabic{enumi})}
\tightlist
\item
  Generating-function derivation (the way the section leads you to do
  it)
\end{enumerate}

Define \(G_n(z)=\sum_{k\ge 0} p_{n k} z^k\). Multiplying the recurrence
by \(z^k\) and summing over \(k\) yields

\[
G_n(z)=(q+p z)\,G_{n-1}(z),\qquad G_1(z)=q+p z.
\]

Hence by induction

\[
G_n(z)=(q+p z)^n. \,\,\text{(Eq. (18))} \quad \text{(see the text’s derivation)} 
\]

Expanding with the binomial theorem and reading off the coefficient of
\(z^k\) gives

\[
\boxed{\,p_{n k}=\binom{n}{k}p^{k}q^{\,n-k}\,}\quad(0\le k\le n),\quad p_{n k}=0\text{ otherwise.}
\]

\begin{enumerate}
\def\labelenumi{\arabic{enumi})}
\setcounter{enumi}{1}
\tightlist
\item
  Direct counting (combinatorial) derivation
\end{enumerate}

Choose which \(k\) of the \(n\) tosses are heads: \(\binom{n}{k}\)
choices. Each such outcome has probability \(p^k q^{n-k}\) by
independence. Summing over the \(\binom{n}{k}\) equiprobable patterns
yields the same closed form:

\[
p_{n k}=\binom{n}{k}p^{k}q^{\,n-k}.
\]

Quick checks

\begin{itemize}
\tightlist
\item
  Normalization: \(\sum_{k=0}^n p_{n k}=(p+q)^n=1\).
\item
  Recurrence: The formula satisfies
  \(p_{n k}=p\,p_{n-1,k-1}+q\,p_{n-1,k}\) identically (verify with
  Pascal's identity).
\item
  For a fair coin \(p=\tfrac12\): \(p_{n k}=2^{-n}\binom{n}{k}\).
\end{itemize}

\subsection{5 {[}M13{]}}\label{m13-3}

Here ``Fig. 11'' is the binomial case with \(n=12\) independent tosses
and success probability \(p=\tfrac35\) (so \(q=1-p=\tfrac25\)). For a
binomial distribution the probability-generating function is

\[
G_n(z)=(q+pz)^n,
\]

and from the general identities

\[
\operatorname{mean}(G_n)=G_n'(1),\qquad
\operatorname{var}(G_n)=G_n''(1)+G_n'(1)-\bigl(G_n'(1)\bigr)^2,
\]

we obtain \(\mathbb E[X]=np\) and \(\mathrm{Var}(X)=npq\), hence
\(\sigma=\sqrt{npq}\).

Plugging in \(n=12,\;p=\tfrac35,\;q=\tfrac25\):

\begin{enumerate}
\def\labelenumi{\arabic{enumi}.}
\tightlist
\item
  Mean
\end{enumerate}

\[
\mathbb E[X]=np=12\cdot \frac35=\frac{36}{5}=7.2.
\]

\begin{enumerate}
\def\labelenumi{\arabic{enumi}.}
\setcounter{enumi}{1}
\tightlist
\item
  Standard deviation
\end{enumerate}

\[
\sigma=\sqrt{npq}
=\sqrt{12\cdot\frac35\cdot\frac25}
=\sqrt{\frac{72}{25}}
=\frac{\sqrt{72}}{5}
=\frac{6\sqrt2}{5}
\approx 1.697.
\]

So the distribution in Fig. 11 has mean \(7.2\) and standard deviation
\(\tfrac{6\sqrt2}{5}\approx1.697\).

\subsection{6 {[}HM27{]}}\label{hm27}

\begin{enumerate}
\def\labelenumi{(\arabic{enumi})}
\setcounter{enumi}{17}
\tightlist
\item
  Binomial \((n,p)\)
\end{enumerate}

Here \(G(e^t)=(q+pe^t)^n\), so

\[
\ln G(e^t)=n\ln(q+pe^t).
\]

Differentiate three times and evaluate at \(t=0\):

\[
\kappa_1=np,\qquad \kappa_2=npq,\qquad \boxed{\;\kappa_3=npq\,(1-2p)\;}.
\]

\begin{enumerate}
\def\labelenumi{(\arabic{enumi})}
\setcounter{enumi}{19}
\tightlist
\item
  Discrete uniform on \(\{1,\dots,n\}\)
\end{enumerate}

This distribution is symmetric about its mean \((n+1)/2\): for each
deviation \(+\delta\) there is a matching \(-\delta\) with equal
probability \(1/n\). Hence all odd central moments are zero, so the
third cumulant (which equals the third central moment) is

\[
\boxed{\;\kappa_3=0\;}.
\]

\begin{enumerate}
\def\labelenumi{(\arabic{enumi})}
\setcounter{enumi}{7}
\tightlist
\item
  Number of ``record highs'' \(A\) in Algorithm M
\end{enumerate}

From §1.2.10 we can factor the pgf for \(A\) as

\[
G_n(z)=\prod_{k=2}^n Q_k(z),\qquad Q_k(z)=\frac{z+k-1}{k}=\left(1-\frac1k\right)+\frac1k\,z,
\]

so \(A\) is a sum of independent Bernoulli variables \(I_k\) with
\(\Pr(I_k=1)=1/k\). Cumulants add for products (because \(\ln G\) turns
products into sums), so

\[
\kappa_3(A)=\sum_{k=2}^n \kappa_3(I_k).
\]

For a Bernoulli\((p)\) variable, \(\kappa_3=p(1-p)(1-2p)\) (the binomial
result with \(n=1\)). Thus with \(p=1/k\),

\[
\kappa_3(I_k)=\frac1k\left(1-\frac1k\right)\left(1-\frac{2}{k}\right)
=\frac{(k-1)(k-2)}{k^3}
=\frac1k-\frac{3}{k^2}+\frac{2}{k^3}.
\]

Summing from \(k=2\) to \(n\) gives

\[
\boxed{\;\kappa_3(A)\;=\;\sum_{k=2}^n\!\Bigl(\frac1k-\frac{3}{k^2}+\frac{2}{k^3}\Bigr)
\;=\;H_n-3H_n^{(2)}+2H_n^{(3)}\;}.
\]

Here \(H_n=\sum_{k=1}^n\frac1k\) and
\(H_n^{(m)}=\sum_{k=1}^n\frac1{k^m}\). This matches the factorization of
\(G_n\) in (8) and the additivity of semi-invariants stated in the
section.

\subsection{7 {[}M27{]}}\label{m27-3}

In §1.2.10, \(A\) is the number of times Algorithm \(M\) changes the
current maximum while scanning \(X[1],\dots,X[n]\). (Minimum \(0\),
maximum \(n-1\).) The section shows how to analyze \(A\) when all values
are distinct and even gives the exact distribution \(p_{n,k}=\Pr(A=k)\)
in terms of Stirling numbers of the first kind:

\[
p_{n,k}=\frac{\bigl[{n\atop k+1}\bigr]}{n!}.
\]

(See Eq. (9) and the surrounding discussion.)

Exercise 7 weakens the hypothesis: assume only that the multiset
\(\{X[1],\dots,X[n]\}\) contains exactly \(m\) distinct values (ties are
allowed), otherwise random subject to this constraint. Find the
distribution of \(A\).

Key observation

A change of maximum occurs exactly when we encounter, for the first
time, a value that is strictly larger than all values seen so far.
Therefore, for purposes of counting \(A\), all that matters is the order
in which the \(m\) distinct values make their first appearance.
Repetitions between first appearances do not affect \(A\).

Condition on ``exactly \(m\) distinct values appear.'' Under the natural
``uniform subject to the constraint'' model (equivalently: sample i.i.d.
from \(m\) labels and condition on using all labels), the order of first
appearances of the \(m\) distinct values is uniform over all \(m!\)
permutations. (This is a standard coupon-collector symmetry:
conditioning on eventual appearance doesn't bias the relative order of
the first hits.)

Let \(\pi\) be that random permutation of the \(m\) distinct values,
written in increasing-value labels \(1<2<\cdots<m\). When we scan values
in the order of first appearance (i.e., scan \(\pi\)), the running
maximum increases once for each left-to-right maximum (record) of
\(\pi\) after its first element. Thus,

\[
A \;=\; (\#\text{ of records in }\pi)\;-\;1.
\]

But the number of left-to-right records in a random permutation of \(m\)
elements has the well-known distribution

\[
\Pr\{\#\text{records}=r\}=\frac{\bigl[{m\atop r}\bigr]}{m!}\quad(r=1,\dots,m),
\]

the same Stirling-of-the-first-kind distribution used in the section for
the ``all distinct'' case (there it appears with \(n\) and with
\(r=k+1\)). 

Result

Therefore, conditional on there being exactly \(m\) distinct values
among \(X[1],\dots,X[n]\),

\[
\boxed{\;\Pr(A=k\mid \#\text{distinct}=m)=\dfrac{\bigl[{m\atop k+1}\bigr]}{m!}\;,\qquad k=0,1,\dots,m-1.\;}
\]

In words: it's exactly the same distribution as in the all-distinct
case, with \(n\) replaced by \(m\) (and with the shift by \(+1\) because
\(A\) counts ``changes,'' not ``records''). This reflects the fact that
only the relative order of first occurrences of the \(m\) distinct
values matters.

Checks and moments

\begin{itemize}
\tightlist
\item
  For \(m=3\):
  \(\Pr(A=0,1,2)=\frac{[{3\atop1}]}{6},\frac{[{3\atop2}]}{6},\frac{[{3\atop3}]}{6}=\frac{2}{6},\frac{3}{6},\frac{1}{6}\),
  matching the table in the section for \(n=3\). 
\item
  By the same substitution \(n\mapsto m\) in Eq. (16), the conditional
  mean and variance are
\end{itemize}

\[
\mathbb{E}[A\mid m]=H_m-1,\qquad
\operatorname{Var}(A\mid m)=H_m-H_m^{(2)},
\]

so \(A\) has summary
\((\min 0,\ \text{ave }H_m-1,\ \max m-1,\ \text{dev }\sqrt{H_m-H_m^{(2)}})\).

\subsection{8 {[}M20{]}}\label{m20-29}

Total outcomes Each position has \(M\) choices, independently, so there
are \(M^n\) sequences in all.

Favorable outcomes (all entries distinct) Fill the \(n\) positions
without repeating any element:

\[
M\cdot (M-1)\cdot (M-2)\cdots (M-n+1)
\;=\;
\frac{M!}{(M-n)!}
\;=\;
(M)_n
\]

(the falling factorial).

Therefore, for \(0\le n\le M\),

\[
\boxed{\; \Pr(\text{all distinct}) \;=\; \frac{(M)_n}{M^n}
\;=\; \frac{M!}{(M-n)!\,M^n}\;}
\]

and for \(n>M\) the probability is \(0\).

Edge cases

\begin{itemize}
\tightlist
\item
  \(n=0\): probability \(1\) (empty selection).
\item
  \(n=1\): probability \(1\).
\item
  \(n=M\): probability \(M!/M^M\).
\end{itemize}

Useful approximations (birthday-style)

When \(n\ll M\),

\[
\ln\Pr(\text{distinct})
= \sum_{i=0}^{n-1}\ln\!\Bigl(1-\frac{i}{M}\Bigr)
\approx -\frac{n(n-1)}{2M}
\]

so

\[
\Pr(\text{distinct})
\approx \exp\!\left(-\frac{n(n-1)}{2M}\right).
\]

A sharper expansion keeps the next term:

\[
\Pr(\text{distinct})
\approx \exp\!\left(
-\frac{n(n-1)}{2M}
-\frac{n(n-1)(2n-1)}{12M^2}
\right).
\]

These are the classic ``birthday problem'' approximations and are handy
when \(M\) is large.

\subsection{9 {[}M25{]}}\label{m25-22}

Combinatorial derivation (with Stirling numbers of the second kind)

Count sequences \(X[1],\dots,X[n]\) over an \(M\)-symbol alphabet that
use exactly \(m\) symbols, then divide by the total \(M^n\) sequences.

\begin{enumerate}
\def\labelenumi{\arabic{enumi}.}
\item
  Choose which \(m\) symbols appear: \(\binom{M}{m}\).
\item
  Among those \(m\) chosen symbols, count length-\(n\) sequences in
  which each of the \(m\) symbols appears at least once (i.e.,
  surjections from the \(n\) positions to the \(m\) symbols). Number of
  such sequences \(=\) (number of onto functions)
  \(= m!\, {n \brace m}\), where \({n \brace m}\) is a Stirling number
  of the second kind (partitions of an \(n\)-set into \(m\) nonempty
  unlabeled blocks, then label the blocks by the \(m\) symbols).
\item
  Divide by the total number \(M^n\) of sequences.
\end{enumerate}

Therefore the probability is

\[
\boxed{\displaystyle
q_{n,m}
=\frac{\binom{M}{m}\, m!\, {n \brace m}}{M^n}
=\frac{(M)_m\, {n \brace m}}{M^n}
=\frac{M!\, {n \brace m}}{(M-m)!\, M^n}\,,
}
\qquad 0\le m\le \min\{n,M\},
\]

and \(q_{n,m}=0\) otherwise.

Equivalent inclusion--exclusion form

Using \(m!\,{n \brace m}=\sum_{j=0}^{m}(-1)^j\binom{m}{j}(m-j)^n\), an
equivalent expression is

\[
q_{n,m}
=\frac{\binom{M}{m}}{M^n}\,\sum_{j=0}^{m}(-1)^j\binom{m}{j}(m-j)^n .
\]

Recurrence check (dynamic viewpoint)

Let \(q_{n,m}\) be the desired probability. When we append the \(n\)th
draw:

\begin{itemize}
\tightlist
\item
  With probability \(\frac{m}{M}\) we hit one of the \(m\) symbols
  already seen (stay at \(m\)).
\item
  With probability \(\frac{M-m}{M}\) we introduce a new symbol (go from
  \(m-1\) to \(m\)).
\end{itemize}

Hence

\[
q_{n,m}
=\frac{M-m+1}{M}\,q_{n-1,m-1}
+\frac{m}{M}\,q_{n-1,m},
\quad
q_{1,1}=1,\ q_{1,m\ne 1}=0.
\]

One can verify that the closed form above satisfies this recurrence.

Sanity checks

\begin{itemize}
\item
  All distinct (Exercise 8). Setting \(m=n\) gives
  \(q_{n,n}=\dfrac{(M)_n}{M^n}\), the probability that all \(n\) draws
  are different (agrees with Ex. 8). 
\item
  Exactly one distinct value. \(q_{n,1}=\dfrac{M}{M^n}=M^{1-n}\) (all
  draws identical).
\item
  Support. \(q_{n,m}=0\) if \(m>n\) or \(m>M\).
\end{itemize}

\subsection{10 {[}M20{]}}\label{m20-30}

Step 1 --- Condition on the number of distinct values

Let \(m\) be the number of distinct values that actually appear among
\(X[1..n]\). Given \(m\), duplicates don't affect when new records
occur; Algorithm M behaves exactly as if we had \(m\) distinct inputs.
Therefore

\[
\Pr(A=k\mid m)=p_{m k},
\]

where (from §1.2.10) the ``distinct inputs'' distribution is

\[
p_{m k}=\frac{\big[\,m\;\;k{+}1\,\big]}{m!},
\]

with \(\big[\,m\;\;r\,\big]\) the (unsigned) Stirling numbers of the
first kind.  

Step 2 --- Probability of having exactly \(m\) distinct values

If we sample \(n\) times from \(M\) equally likely symbols, the
probability that exactly \(m\) distinct symbols appear is

\[
q_{n m}
=\frac{M!}{(M-m)!}\;\frac{\big\{\,n\;\;m\,\big\}}{M^{\,n}},
\]

where \(\big\{\,n\;\;m\,\big\}\) are Stirling numbers of the second kind
(they count surjections of \(n\) labeled draws onto \(m\) labels). 

Step 3 --- Law of total probability

Sum over feasible \(m\) (you need at least \(k{+}1\) distinct values to
get \(k\) strict record increases, and you can't exceed either \(n\) or
\(M\)):

\[
\boxed{\;
\Pr(A=k)
=\sum_{m=k+1}^{\min(n,M)}
q_{n m}\,p_{m k}
=\frac{1}{M^{\,n}}
\sum_{m=k+1}^{\min(n,M)}
\binom{M}{m}\,
\big\{\,n\;\;m\,\big\}\,
\big[\,m\;\;k{+}1\,\big]
\;}
\]

Notes and quick check

\begin{itemize}
\tightlist
\item
  Range: \(k=0,1,\dots,\min(n{-}1,M{-}1)\).
\item
  Example \(n=3,M=2\): Only \(m=1,2\) are possible.
  \(\Pr(m{=}1)=2/2^3=1/4\Rightarrow \Pr(A{=}0\mid m{=}1)=1\).
  \(\Pr(m{=}2)=6/8=3/4\) and \(p_{2,0}=p_{2,1}=1/2\). So
  \(\Pr(A{=}0)=1/4+ (3/4)(1/2)=5/8,\;\Pr(A{=}1)=(3/4)(1/2)=3/8\),
  agreeing with direct enumeration.
\end{itemize}

(Optional) Mean

Because \(E[A\mid m]=H_m-1\) in the distinct-case analysis, the overall
mean can be written

\[
E[A]=\sum_m q_{n m}(H_m-1)
=H_M-\sum_{m=1}^{M}\frac{\big(1-\frac{m}{M}\big)^{n}}{m}-1,
\]

a compact but not further-simplifiable form in general.

\subsection{11 {[}M15{]}}\label{m15-12}

Setup (what semi-invariants are)

Knuth defines the semi-invariants \(\kappa_1,\kappa_2,\ldots\) of a
distribution with probability generating function \(G(z)\) via the
logarithm of the \emph{characteristic} (a.k.a. exponential generating)
function \(G(e^t)\):

\[
\ln G(e^t)\;=\;\frac{\kappa_1 t}{1!}+\frac{\kappa_2 t^2}{2!}+\frac{\kappa_3 t^3}{3!}+\cdots.
\]

Equivalently,
\(\kappa_n=\left.\dfrac{d^n}{dt^n}\ln G(e^t)\right|_{t=0}\). In
particular, \(\kappa_1=G'(1)\) (the mean) and
\(\kappa_2=G''(1)+G'(1)-G'(1)^2\) (the variance).

Apply the transformation \(F(z)=z^n G(z)\)

Compute the corresponding exponential form:

\[
F(e^t) \;=\; e^{nt}\,G(e^t)\quad\Longrightarrow\quad
\ln F(e^t)\;=\;nt+\ln G(e^t).
\]

Now compare coefficients of powers of \(t\) with the defining series
above:

\begin{itemize}
\item
  The coefficient of \(t^1\) increases by \(n\), so

  \[
  \kappa_1'=\kappa_1+n.
  \]
\item
  For every order \(r\ge 2\), the coefficients are unchanged:

  \[
  \kappa_r'=\kappa_r\qquad(r\ge 2).
  \]
\end{itemize}

Intuition and quick check

Multiplying \(G(z)\) by \(z^n\) is exactly the
probability-generating-function operation for adding a constant \(n\) to
the random variable (\(Y=X+n\)). Cumulants beyond first order are
translation-invariant, so only the first cumulant (the mean) shifts.
Consequently, the variance and all higher cumulants stay the same. This
agrees with the mean/variance formulas from the section: \$
\text{mean}=G'(1)\$ and \(\text{var}=G''(1)+G'(1)-G'(1)^2\); for \(F\)
we get \(F'(1)=n+G'(1)\) while the variance expression is unaffected. 

Final answer

Under \(F(z)=z^nG(z)\),

\[
\boxed{\;\kappa_1'=\kappa_1+n,\qquad \kappa_r'=\kappa_r\ \text{ for all } r\ge 2.\;}
\]

\subsection{12 {[}HM21{]}}\label{hm21-7}

\begin{enumerate}
\def\labelenumi{\arabic{enumi})}
\tightlist
\item
  Show \(G(e^t)=1+M_1 t+M_2 \tfrac{t^2}{2!}+M_3 \tfrac{t^3}{3!}+\cdots\)
\end{enumerate}

\[
G(e^t)=\sum_{k\ge 0} p_k e^{kt}
=\sum_{k\ge 0} p_k \sum_{n\ge 0}\frac{(kt)^n}{n!}
=\sum_{n\ge 0}\frac{t^n}{n!}\Bigl(\sum_{k\ge 0} k^n p_k\Bigr)
=\sum_{n\ge 0}\frac{M_n}{n!}\,t^n.
\]

In particular \(G(1)=\sum_k p_k=1\), so \(G(e^t)\) has constant term
\(1\). 

\begin{enumerate}
\def\labelenumi{\arabic{enumi})}
\setcounter{enumi}{1}
\tightlist
\item
  Cumulants \(\kappa_n\) from Arbogast/Faà di Bruno
\end{enumerate}

Let \(H(t)=G(e^t)\). The cumulant generating function is

\[
K(t)=\log H(t)=\log G(e^t),\qquad \kappa_n = K^{(n)}(0).
\]

Arbogast's (Faà di Bruno) formula for derivatives of a composition
gives, for \(f=\log\) and \(g=H\),

\[
\frac{d^n}{dt^n}\log H(t)
=\sum_{\substack{k_1,\dots,k_n\ge 0\\ 1k_1+2k_2+\cdots+n k_n=n}}
(-1)^{m-1}(m-1)!\;
\frac{n!}{k_1!\cdots k_n!\,1!^{k_1}\cdots n!^{k_n}}\;
\prod_{j=1}^{n}\!\left(\frac{H^{(j)}(t)}{H(t)}\right)^{k_j},
\]

where \(m=k_1+\cdots+k_n\). Evaluating at \(t=0\) uses \(H(0)=G(1)=1\)
and \(H^{(j)}(0)=M_j\) from part (1). Therefore

\[
\boxed{\;
\kappa_n
=\sum_{\substack{k_1,\dots,k_n\ge 0\\ 1k_1+2k_2+\cdots+n k_n=n}}
(-1)^{m-1}(m-1)!\;
\frac{n!}{k_1!\cdots k_n!\,1!^{k_1}\cdots n!^{k_n}}\;
\prod_{j=1}^{n} M_j^{\,k_j}\,,
\qquad m=k_1+\cdots+k_n.
\;}
\]

This is exactly the Bell-polynomial form asked for in the exercise.
Checking small \(n\) reproduces the book's identities:

\begin{itemize}
\tightlist
\item
  \(\kappa_1=M_1\).
\item
  \(\kappa_2=M_2-M_1^2\).
\item
  \(\kappa_3=M_3-3M_1M_2+2M_1^3\).
\item
  \(\kappa_4=M_4-4M_1M_3+12M_1^2M_2-3M_2^2-6M_1^4\). 
\end{itemize}

\begin{enumerate}
\def\labelenumi{\arabic{enumi})}
\setcounter{enumi}{2}
\tightlist
\item
  Express \(\kappa_n\) via central moments \(m_j\) (for \(n\ge2\))
\end{enumerate}

Center the variable by \(Y=X-M_1\). Then

\[
\log \mathbb{E}[e^{tY}]
=\log\!\bigl(e^{-M_1 t} H(t)\bigr)
=K(t)-M_1 t.
\]

Hence the cumulants for \(Y\) satisfy \(\kappa'_1=0\) and
\(\kappa'_n=\kappa_n\) for \(n\ge2\). But the derivatives of
\(\mathbb{E}[e^{tY}]\) at \(0\) are exactly the central moments:
\((\mathbb{E}[e^{tY}])^{(j)}\!\!(0)=m_j\) with \(m_1=0\). Applying the
same Arbogast/Bell-polynomial formula with \(M_1=0\) and \(M_j\)
replaced by \(m_j\) (for \(j\ge2\)) yields, for \(n\ge2\),

\[
\boxed{\;
\kappa_n
=\sum_{\substack{k_2,\dots,k_n\ge 0\\ 2k_2+3k_3+\cdots+n k_n=n}}
(-1)^{m-1}(m-1)!\;
\frac{n!}{k_2!\cdots k_n!\,2!^{k_2}\cdots n!^{k_n}}\;
\prod_{j=2}^{n} m_j^{\,k_j},
\qquad m=k_2+\cdots+k_n.
\;}
\]

Concrete low-order cases:

\begin{itemize}
\tightlist
\item
  \(\kappa_2=m_2\).
\item
  \(\kappa_3=m_3\).
\item
  \(\kappa_4=m_4-3m_2^2\).
\item
  \(\kappa_5=m_5-10\,m_3 m_2\).
\item
  \(\kappa_6=m_6-15\,m_4 m_2-10\,m_3^2+30\,m_2^3\).
\end{itemize}

\subsection{13 {[}HM38{]}}\label{hm38}

Step 1: Bring in the exact pgf and its mean/variance

From Eq. (8) in §1.2.10,

\[
G_n(z)=\frac{1}{n!}\,(z+n-1)(z+n-2)\cdots(z+1)
=\frac{\Gamma(n+z)}{\Gamma(z+1)\,\Gamma(n+1)}.
\]

(The gamma form is immediate from the product.) 

Moreover, §1.2.10 obtains

\[
\mu_n=\operatorname{mean}(G_n)=H_n-1,\qquad
\sigma_n^2=\operatorname{var}(G_n)=H_n-H_n^{(2)}.
\]

Here \(H_n=\sum_{k=1}^n \frac1k\) and
\(H_n^{(2)}=\sum_{k=1}^n \frac1{k^2}\). 

Useful asymptotics: \(H_n=\ln n+\gamma+O(1/n)\) and
\(H_n^{(2)}=\frac{\pi^2}{6}+O(1/n)\). Hence \(\sigma_n^2\sim\ln n\) and
\(\mu_n-\ln n=O(1)\). 

Step 2: Define the normalized characteristic function

Set \(z_n=e^{it/\sigma_n}\). The normalized characteristic function in
the exercise is

\[
\Phi_n(t)\;:=\;e^{-it\mu_n/\sigma_n}\,G_n(z_n).
\]

We will show \(\Phi_n(t)\to e^{-t^2/2}\).

Step 3: Take logs and use a ratio for gamma

Using the gamma form,

\[
\ln G_n(z)=\ln\Gamma(n+z)-\ln\Gamma(n+1)-\ln\Gamma(z+1).
\]

For fixed complex \(z\) with \(\Re z> -1\), the classical ratio
asymptotic

\[
\frac{\Gamma(n+z)}{\Gamma(n+1)}=n^{\,z-1}\bigl(1+O(1/n)\bigr)
\]

gives

\[
\ln G_n(z)= (z-1)\ln n - \ln\Gamma(z+1) + O(1/n).
\]

Since \(z_n=e^{it/\sigma_n}\to 1\), we have
\(\Gamma(z_n+1)\to\Gamma(2)=1\) and thus \(\ln\Gamma(z_n+1)\to 0\).

Therefore

\[
\ln \Phi_n(t)
= -\,\frac{it\mu_n}{\sigma_n}\;+\;\bigl(z_n-1\bigr)\,\ln n\;+\;o(1).
\]

Step 4: Expand \(z_n\) and pass to the limit

Expand
\(z_n=e^{it/\sigma_n}=1+\frac{it}{\sigma_n}-\frac{t^2}{2\sigma_n^2}+O(\sigma_n^{-3})\).
Plugging in,

\[
\ln \Phi_n(t)
= \Bigl(\frac{it}{\sigma_n}\ln n - \frac{it\mu_n}{\sigma_n}\Bigr)
\;-\;\frac{t^2}{2\sigma_n^2}\ln n\;+\;o(1).
\]

Now use the §1.2.10 statistics and harmonic-number asymptotics:

\begin{itemize}
\tightlist
\item
  \(\mu_n-\ln n=O(1)\) and \(\sigma_n\sim\sqrt{\ln n}\) \(\Rightarrow\)
  \(\displaystyle \frac{(\mu_n-\ln n)}{\sigma_n}\to 0\). So the linear
  (imaginary) term vanishes.
\item
  \(\displaystyle \frac{\ln n}{\sigma_n^2}\to 1\) since
  \(\sigma_n^2=H_n-H_n^{(2)}=\ln n+O(1)\). So the quadratic term tends
  to \(-t^2/2\).
\end{itemize}

Thus \(\ln \Phi_n(t)\to -t^2/2\), hence \(\Phi_n(t)\to e^{-t^2/2}\),
proving the required convergence to the normal law.

Intuition (why this works)

Eq. (8) factorizes \(G_n\) as a product of simple pgfs
\(Q_k(z)=(z+k-1)/k\). That corresponds to a sum of independent Bernoulli
variables with \(\Pr(B_k=1)=1/k\) (so \(A_n=\sum_{k=2}^n B_k\)). Sums of
many small, independent contributions tend to normality; the calculation
above is the precise characteristic-function proof tailored to these
particular parameters. (§1.2.10 already hints these distributions are
``approximately normal'' for large \(n\).)

\subsection{14 {[}HM30{]}}\label{hm30-4}

Write \(\sigma=\sigma_n=\sqrt{npq}\) and consider

\[
\Phi_n(t)\;=\;e^{-it\mu/\sigma}\,G_n\!\left(e^{it/\sigma}\right)
\;=\;e^{-it\,np/\sigma}\,\Bigl(q+p\,e^{it/\sigma}\Bigr)^{n}.
\]

Take logarithms:

\[
\log \Phi_n(t)\;=\;-\frac{it\,np}{\sigma}\;+\;n\log\!\Bigl(q+p\,e^{it/\sigma}\Bigr).
\]

Use the Taylor expansions \(e^{ix}=1+ix-\tfrac{x^{2}}{2}+O(x^{3})\) and
\(\log(1+x)=x-\tfrac{x^{2}}{2}+O(x^{3})\) as \(x\to0\). First expand
inside the log:

\[
q+p\,e^{it/\sigma}
= q+p\!\left(1+\frac{it}{\sigma}-\frac{t^{2}}{2\sigma^{2}}+O(\sigma^{-3})\right)
= 1 + \underbrace{p\frac{it}{\sigma}-p\frac{t^{2}}{2\sigma^{2}}}_{=:x} + O(\sigma^{-3}).
\]

Now

\[
\log\!\Bigl(q+p\,e^{it/\sigma}\Bigr)
= x-\frac{x^{2}}{2}+O(\sigma^{-3})
= \frac{pit}{\sigma}-\frac{pq\,t^{2}}{2\sigma^{2}}+O(\sigma^{-3}),
\]

since \(x^{2} = -\,p^{2}t^{2}/\sigma^{2}+O(\sigma^{-3})\) and
\(pq=p(1-p)\).

Multiply by \(n\):

\[
n\log\!\Bigl(q+p\,e^{it/\sigma}\Bigr)
= \frac{it\,np}{\sigma}\;-\;\frac{pq\,t^{2}}{2}\cdot\frac{n}{\sigma^{2}}
\;+\;n\cdot O(\sigma^{-3}).
\]

Because \(\sigma^{2}=npq\), we have \(n/\sigma^{2}=1/(pq)\) and
\(n\cdot O(\sigma^{-3})=O(1/\sqrt{n})\to0\). Hence

\[
n\log\!\Bigl(q+p\,e^{it/\sigma}\Bigr)=\frac{it\,np}{\sigma}-\frac{t^{2}}{2}+o(1).
\]

Finally,

\[
\log \Phi_n(t)
= -\frac{it\,np}{\sigma}
+\left(\frac{it\,np}{\sigma}-\frac{t^{2}}{2}+o(1)\right)
= -\frac{t^{2}}{2}+o(1),
\]

so \(\Phi_n(t)\to e^{-t^{2}/2}\) as \(n\to\infty\). This is exactly
Knuth's criterion for ``approaches a normal distribution,'' so the
binomial distribution (properly standardized) approaches the normal law
(de Moivre--Laplace). 

(As usual, exclude the degenerate cases \(p\in\{0,1\}\) where
\(\sigma_n=0\).)

\subsection{15 {[}HM23{]}}\label{hm23-1}

\begin{enumerate}
\def\labelenumi{(\alph{enumi})}
\tightlist
\item
  Probability generating function
\end{enumerate}

If \(X\sim\text{Poisson}(\mu)\), then

\[
\Pr(X=k)=e^{-\mu}\frac{\mu^{k}}{k!}\quad(k=0,1,2,\dots).
\]

The probability generating function (PGF) is

\[
G(z)=\sum_{k\ge 0}\Pr(X=k)\,z^k
    =e^{-\mu}\sum_{k\ge 0}\frac{(\mu z)^k}{k!}
    =e^{-\mu}e^{\mu z}
    =\exp\!\big(\mu(z-1)\big).
\]

\begin{enumerate}
\def\labelenumi{(\alph{enumi})}
\setcounter{enumi}{1}
\tightlist
\item
  Semi-invariants (cumulants)
\end{enumerate}

Knuth defines the semi-invariants \(\kappa_n\) by the expansion of the
logarithm of the generating function evaluated at \(e^t\):

\[
\ln G(e^t)=\sum_{n\ge 1}\kappa_n\frac{t^n}{n!}.
\]

For Poisson,

\[
\ln G(e^t)=\ln\!\big(\exp(\mu(e^t-1))\big)=\mu\,(e^t-1)=
\mu\sum_{n\ge 1}\frac{t^n}{n!}.
\]

Hence, for every \(n\ge 1\),

\[
\boxed{\ \kappa_n=\mu\ }.
\]

In particular, \(\kappa_1=\mu=\operatorname{E}X\) and
\(\kappa_2=\mu=\operatorname{Var}X\), matching the usual mean and
variance of the Poisson law. (Definition and identities for \(\kappa_n\)
are given just above Eq. (23) in §1.2.10.) 

\begin{enumerate}
\def\labelenumi{(\alph{enumi})}
\setcounter{enumi}{2}
\tightlist
\item
  Poisson(\(np\)) approaches normal as \(n\to\infty\)
\end{enumerate}

Exercise 13 says a sequence with PGFs \(G_n\), means \(\mu_n\), and
standard deviations \(\sigma_n\) approaches a normal distribution iff

\[
\lim_{n\to\infty} e^{-it\mu_n/\sigma_n}\,G_n\!\left(e^{it/\sigma_n}\right)=e^{-t^2/2}
\quad\text{for all real }t.
\]

For \(X_n\sim\text{Poisson}(np)\),

\[
G_n(z)=\exp\!\big(np(z-1)\big),\qquad
\mu_n=np,\qquad
\sigma_n^2=np.
\]

Compute the left side:

\[
e^{-it\mu_n/\sigma_n}\,G_n\!\left(e^{it/\sigma_n}\right)
=\exp\!\Big(np\big(e^{\,it/\sqrt{np}}-1\big)-it\sqrt{np}\Big).
\]

Use the Taylor series
\(e^{ix}=1+ix-\tfrac{x^2}{2}-i\tfrac{x^3}{6}+O(x^4)\) with
\(x=t/\sqrt{np}\):

\[
e^{\,it/\sqrt{np}}-1-\frac{it}{\sqrt{np}}
=-\frac{t^2}{2np}-i\frac{t^3}{6(np)^{3/2}}+O\!\left(\frac{1}{n^2}\right).
\]

Multiplying by \(np\) gives

\[
np\!\left(e^{\,it/\sqrt{np}}-1-\frac{it}{\sqrt{np}}\right)
= -\frac{t^2}{2}-i\frac{t^3}{6\sqrt{np}}+O\!\left(\frac{1}{n}\right).
\]

Therefore

\[
e^{-it\mu_n/\sigma_n}\,G_n\!\left(e^{it/\sigma_n}\right)
=\exp\!\left(-\frac{t^2}{2}+o(1)\right)\xrightarrow[n\to\infty]{}e^{-t^2/2},
\]

which is exactly the criterion in Exercise 13. Thus Poisson(\(np\))
standardized converges to the normal distribution.

\subsection{16 {[}M25{]}}\label{m25-23}

Preliminaries (from §1.2.10)

If \(G(z)=\sum_{n\ge 0} p_n z^n\) is a pgf (so \(G(1)=1\)), then

\[
\text{mean}(G)=G'(1),\qquad
\text{var}(G)=G''(1)+G'(1)-\bigl(G'(1)\bigr)^2.
\]

These are Eq. (12) and Eq. (13) in the text.

Step 1 --- The mixture pgf

Let \(X\) be the sample from the chosen component, and let \(K\) be the
index of the chosen component. Then

\[
g(z)\;=\;\mathbb{E}\!\left[z^{X}\right]
\;=\;\sum_k \Pr(K=k)\,\mathbb{E}\!\left[z^{X}\mid K=k\right]
\;=\;\sum_k p_k\,g_k(z).
\]

So the pgf of a mixture is the convex combination of the component pgfs.
(This is exactly what Knuth lists in the official answer.)

Step 2 --- The mean of the mixture

Differentiate \(g(z)=\sum_k p_k g_k(z)\) and evaluate at \(z=1\):

\[
\text{mean}(g)=g'(1)=\sum_k p_k\, g_k'(1)=\sum_k p_k\,\mu_k.
\]

Thus the mean of a mixture is the mixture of the means. 

Step 3 --- The variance of the mixture

Start from the variance formula in terms of derivatives:

\[
\text{var}(g)=g''(1)+g'(1)-\bigl(g'(1)\bigr)^2.
\]

Linearity gives \(g''(1)=\sum_k p_k g_k''(1)\) and
\(g'(1)=\sum_k p_k g_k'(1)=\sum_k p_k \mu_k\). Also, for each component,
\(g_k''(1)+g_k'(1)=\mathbb{E}_k[X^2]=\sigma_k^2+\mu_k^2\). Hence

\[
\begin{aligned}
\text{var}(g)
&=\sum_k p_k\bigl(g_k''(1)+g_k'(1)\bigr)-\Bigl(\sum_k p_k \mu_k\Bigr)^2\\
&=\sum_k p_k\bigl(\sigma_k^2+\mu_k^2\bigr)-\Bigl(\sum_k p_k \mu_k\Bigr)^2\\
&=\underbrace{\sum_k p_k \sigma_k^2}_{\text{within–component variance}}
\;+\;\underbrace{\Bigl(\sum_k p_k \mu_k^2-\Bigl(\sum_k p_k \mu_k\Bigr)^2\Bigr)}_{\text{between–component variance}}.
\end{aligned}
\]

The bracketed ``between'' term is the variance of the random mean
\(\mu_K\). It can be rewritten in the symmetric form

\[
\sum_{j<k} p_j p_k\,(\mu_j-\mu_k)^2,
\]

since
\(\sum_k p_k \mu_k^2-(\sum_k p_k \mu_k)^2=\tfrac12\sum_{j,k}p_j p_k(\mu_j-\mu_k)^2\).
Therefore

\[
\boxed{\displaystyle
\text{var}(g)\;=\;\sum_k p_k\,\sigma_k^2\;+\;\sum_{j<k} p_j p_k\,(\mu_j-\mu_k)^2.}
\]

This matches Knuth's given result. 

Quick numerical sanity check

Take two Bernoulli components:

\begin{itemize}
\tightlist
\item
  Component 1: \(p_1=0.6\), \(\mu_1=0.3\), \(\sigma_1^2=0.3(0.7)=0.21\).
\item
  Component 2: \(p_2=0.4\), \(\mu_2=0.8\), \(\sigma_2^2=0.8(0.2)=0.16\).
\end{itemize}

Mean: \(0.6\cdot0.3+0.4\cdot0.8=0.18+0.32=0.50\).

Variance:

\[
\sum_k p_k\sigma_k^2=0.6\cdot0.21+0.4\cdot0.16=0.126+0.064=0.19,
\]

\[
\sum_{j<k} p_j p_k(\mu_j-\mu_k)^2
=0.6\cdot0.4\cdot(0.3-0.8)^2
=0.24\cdot0.25
=0.06.
\]

Total \(=0.19+0.06=0.25\). Indeed the mixture itself is Bernoulli with
\(p=0.5\), whose variance is \(0.5(1-0.5)=0.25\).

\subsection{17 {[}M27{]}}\label{m27-4}

\begin{enumerate}
\def\labelenumi{(\alph{enumi})}
\tightlist
\item
  \(h(z)=g(f(z))\) is a PGF
\end{enumerate}

Write the PGFs as

\[
f(z)=\sum_{n\ge0} a_n z^n,\qquad g(z)=\sum_{k\ge0} b_k z^k
\]

with \(a_n,b_k\ge0\) and \(f(1)=g(1)=1\). Then

\[
h(z)=g(f(z))=\sum_{k\ge0} b_k \bigl(f(z)\bigr)^k.
\]

Each power \(f(z)^k\) has nonnegative coefficients and value
\(f(1)^k=1\) at \(z=1\). A nonnegative convex combination of such series
has nonnegative coefficients and

\[
h(1)=\sum_{k\ge0} b_k f(1)^k=\sum_{k\ge0} b_k=g(1)=1.
\]

Hence \(h\) is a PGF.

\begin{enumerate}
\def\labelenumi{(\alph{enumi})}
\setcounter{enumi}{1}
\tightlist
\item
  Probabilistic meaning of \(h\)
\end{enumerate}

Let \(K\) be a nonnegative-integer--valued random variable with PGF
\(g\) (so \(\Pr[K=k]=b_k\)). Let \(X_1,X_2,\dots\) be i.i.d. with PGF
\(f\), independent of \(K\). Define the random sum

\[
S \;=\; X_1+\cdots+X_K\quad(\text{take }S=0\text{ if }K=0).
\]

Conditioning on \(K\),

\[
\text{PGF}(S\mid K=k)=\bigl(f(z)\bigr)^k,
\]

so unconditionally

\[
\text{PGF}(S)=\sum_{k\ge0} b_k f(z)^k \;=\; g(f(z)) \;=\; h(z).
\]

Thus \(h\) is the PGF of a random sum (``compound'' distribution): the
number of summands is random (\(K\sim g\)), each summand has PGF \(f\).
This is the standard compounding/branching-process composition rule for
PGFs. 

\begin{enumerate}
\def\labelenumi{(\alph{enumi})}
\setcounter{enumi}{2}
\tightlist
\item
  Mean and variance of \(h\) via those of \(f\) and \(g\)
\end{enumerate}

Let

\[
\mu_f=f'(1),\quad \sigma_f^2=f''(1)+\mu_f-\mu_f^2,\qquad
\mu_g=g'(1),\quad \sigma_g^2=g''(1)+\mu_g-\mu_g^2.
\]

(These are the PGF formulas from the section.) 

Compute derivatives of \(h(z)=g(f(z))\):

\[
h'(z)=g'(f(z))\,f'(z),\qquad
h''(z)=g''(f(z))\,[f'(z)]^2+g'(f(z))\,f''(z).
\]

Evaluate at \(z=1\) (using \(f(1)=g(1)=1\)):

\[
h'(1)=g'(1)f'(1)=\mu_g\mu_f,
\]

\[
h''(1)=g''(1)\,\mu_f^2+\mu_g\,f''(1).
\]

Now plug into \(\mathrm{var}(h)=h''(1)+h'(1)-[h'(1)]^2\) and use
\(f''(1)=\sigma_f^2-\mu_f+\mu_f^2\),
\(g''(1)=\sigma_g^2-\mu_g+\mu_g^2\): after cancellation one obtains the
clean composition laws

\[
\boxed{\ \mathrm{E}[S]=\mu_h=\mu_g\,\mu_f\ }
\]

\[
\boxed{\ \mathrm{Var}(S)=\sigma_h^2=\mu_g\,\sigma_f^2+\sigma_g^2\,\mu_f^2\ }.
\]

These match the intuitive random-sum identities
\(\mathrm{E}[S]=\mathrm{E}[K]\mathrm{E}[X]\) and
\(\mathrm{Var}(S)=\mathrm{E}[K]\mathrm{Var}(X)+(\mathrm{E}[X])^2\mathrm{Var}(K)\).

Quick check (compound Poisson)

If \(K\sim\text{Poisson}(\lambda)\), then \(g(z)=e^{\lambda(z-1)}\) so
\(\mu_g=\sigma_g^2=\lambda\). The formulas give

\[
\mu_h=\lambda\,\mu_f,\qquad \sigma_h^2=\lambda\,\sigma_f^2+\lambda\,\mu_f^2
=\lambda\,\mathrm{E}[X^2],
\]

which is the classical compound-Poisson result. (This exercise sits
right after the coin-tossing/uniform-distribution PGF work and the
mean/variance-from-PGF identities.)

\subsection{18 {[}M28{]}}\label{m28-5}

Write

\[
S_j \;:=\; k_j+k_{j+1}+\cdots+k_n \quad (1\le j\le n),\qquad S_{n+1}:=0.
\]

Recall the ``all distinct'' case (each \(k_j=1\)) has pgf

\[
G_n(z)\;=\;\sum_{k\ge 0} p_{n,k}\,z^k
\;=\;\frac{1}{n!}\,(z+1)(z+2)\cdots(z+n-1)
\;=\;\frac{1}{\,z+n\,}\binom{z+n}{n}.
\tag{*}
\]

(That is Eq. (8) in the book.)

What we will prove

With multiplicities \(k_1,\dots,k_n\), the pgf of \(A\) is

\[
\boxed{\quad
G(z)\;=\;\frac{1}{z}\,\prod_{j=1}^{n}
\frac{\displaystyle \binom{k_j z + S_{j+1}}{\,k_j\,}}
{\displaystyle \binom{S_j}{\,k_j\,}}
\quad}
\tag{1}
\]

(with the natural convention \(\binom{0}{0}=1\)).

Proof idea (merge ``from the top'')

Build the random permutation by successively shuffling blocks of equal
values, going downwards from the largest value \(n\) to \(1\).

Let \(H_j(z)\) be the pgf of the number of record highs seen among the
values \(\{j,j+1,\dots,n\}\) when those blocks are randomly interleaved
(this counts the very first record within that restricted universe).
Clearly \(H_{n+1}(z)=1\).

Now merge the block of \(k_j\) copies of value \(j\) into the
already-merged block of all higher values (size \(S_{j+1}\)). In that
merge,

\begin{itemize}
\tightlist
\item
  we get exactly one new record from value \(j\) iff at least one \(j\)
  appears before the first element of \(\{j+1,\dots,n\}\);
\item
  otherwise, we get no new record from value \(j\).
\end{itemize}

Among all \(\binom{S_j}{k_j}\) interleavings of \(k_j\)
indistinguishable \(j\)'s with \(S_{j+1}\) higher elements, the weighted
count where we give weight \(z\) when a \(j\) leads and weight \(1\)
otherwise is

\[
\binom{k_j z + S_{j+1}}{k_j}.
\]

(Equivalently, this is \(\binom{S_j}{k_j} + (z-1)\binom{S_j-1}{k_j-1}\)
in disguise; the compact form with \(\binom{k_j z + S_{j+1}}{k_j}\) is a
standard way to package that linear weighting.) Therefore the average
weight factor contributed by inserting the whole \(j\)-block is

\[
\frac{\binom{k_j z + S_{j+1}}{k_j}}{\binom{S_j}{k_j}}.
\]

So the recursive relation is

\[
H_j(z)\;=\;\frac{\binom{k_j z + S_{j+1}}{k_j}}{\binom{S_j}{k_j}}\;H_{j+1}(z).
\]

Multiplying these from \(j=n,n-1,\dots,1\) gives

\[
H_1(z)\;=\;\prod_{j=1}^{n}
\frac{\binom{k_j z + S_{j+1}}{k_j}}{\binom{S_j}{k_j}}.
\]

Finally, \(H_1(z)\) counts records including the very first one (the
first element of the whole sequence). Our \(A\) in Algorithm M
\emph{excludes} the first (because step M4 does not fire at position 1),
hence \(G(z)=H_1(z)/z\), which is exactly (1).

Consistency checks

\begin{enumerate}
\def\labelenumi{\arabic{enumi}.}
\tightlist
\item
  All distinct (\(k_1=\cdots=k_n=1\)). Then \(S_{j+1}=n-j\),
  \(\binom{k_j z + S_{j+1}}{k_j}=z+n-j\), and
  \(\binom{S_j}{k_j}=n-j+1\). Thus
\end{enumerate}

\[
H_1(z)=\prod_{j=1}^{n}\frac{z+n-j}{n-j+1}
=\frac{z(z+1)\cdots(z+n-1)}{n!}.
\]

Dividing by \(z\) reproduces \((*)\). ✔

\begin{enumerate}
\def\labelenumi{\arabic{enumi}.}
\setcounter{enumi}{1}
\tightlist
\item
  Small example. Take values \(\{1,1,3\}\) so \((k_1,k_2,k_3)=(2,0,1)\).
  Direct enumeration of the three arrangements \(113,131,311\) gives
  \(A\) equal to \(1,1,0\) respectively, hence \(G(z)=(1/3)+(2/3)z\).
\end{enumerate}

Formula (1): \(S_1=3,S_2=1,S_3=1,S_4=0\). Then

\[
H_1(z)=\frac{\binom{1\cdot z+0}{1}}{\binom{1}{1}}
\cdot \frac{\binom{0\cdot z+1}{0}}{\binom{1}{0}}
\cdot \frac{\binom{2\cdot z+1}{2}}{\binom{3}{2}}
= z\cdot 1\cdot \frac{2z^2+z}{3}
=\frac{2z^3+z^2}{3}.
\]

Divide by \(z\): \(G(z)=\frac{2z^2+z}{3z}=\frac{1}{3}+\frac{2}{3}z\),
matching the enumeration.

\begin{enumerate}
\def\labelenumi{\arabic{enumi}.}
\setcounter{enumi}{2}
\tightlist
\item
  Normalization. At \(z=1\), each factor equals 1, so \(H_1(1)=1\) and
  \(G(1)=1\).
\end{enumerate}

What you get for free (means/variance)

Once you have \(G(z)\), the mean and variance follow from the standard
relations \(\operatorname{mean}(G)=G'(1)\),
\(\operatorname{var}(G)=G''(1)+G'(1)-[G'(1)]^2\) exactly as in §1.2.10.
(The text derives these identities right after introducing generating
functions.)

\subsection{19 {[}M21{]}}\label{m21-16}

Fix \(k\) and put \(\ell:=a_k\). By definition of the inverse,
\(b_\ell=k\).

(\(\Rightarrow\)) If \(a_k\) is a left-to-right maximum, then \(b_\ell\)
is a right-to-left minimum.

``\(a_k\) is L→R maximum'' means

\[
\forall j<k:\ a_j<\ell.
\]

Because \(a_1,\dots,a_n\) is a permutation, every value \(t\) with
\(t>\ell\) appears somewhere among the \(a\)'s. If any such \(t>\ell\)
appeared at some position \(j\le k\), we would have \(a_j=t>\ell=a_k\),
contradicting that \(a_k\) dominates all earlier entries. Hence

\[
\forall t>\ell:\ \text{the position of }t\text{ is to the right of }k,
\]

i.e.

\[
\forall t>\ell:\ b_t>k=b_\ell.
\]

This is precisely ``\(b_\ell\) is smaller than all entries to its
right,'' i.e.~\(b_\ell\) is an R→L minimum.

(\(\Leftarrow\)) If \(b_\ell\) is a right-to-left minimum, then \(a_k\)
is a left-to-right maximum.

``\(b_\ell\) is R→L minimum'' means

\[
\forall t>\ell:\ b_t>b_\ell=k.
\]

Thus no value \(t>\ell\) occurs at or before position \(k\) in the
\(a\)-sequence. Therefore the first \(k\) entries \(a_1,\dots,a_k\)
contain no value larger than \(\ell=a_k\), so \(a_k\) is larger than all
\(a_j\) with \(j<k\). Hence \(a_k\) is L→R maximum.

Together these prove the equivalence.

Quick example

Take \(a=(3,1,5,2,4)\) (a permutation of \(\{1,\dots,5\}\)).

\begin{itemize}
\item
  Left-to-right maxima in \(a\): \(a_1=3\) (trivially), \(a_3=5\)
  (bigger than \(3,1\)). So positions \(k\in\{1,3\}\).
\item
  Inverse \(b\) lists positions of each value \(1,\dots,5\): value \(1\)
  at pos 2, \(2\) at 4, \(3\) at 1, \(4\) at 5, \(5\) at 3, so
  \(b=(2,4,1,5,3)\).
\item
  Right-to-left minima in \(b\) (scan from right): start with \(3\) (min
  so far), then \(5\) (not \textless3), then \(1\) (new min), then \(4\)
  (not), then \(2\) (not). So the R→L minima are the values \(3\) and
  \(1\), which equal \(k\in\{3,1\}\) --- exactly the positions where
  \(a_k\) was an L→R maximum, as claimed.
\end{itemize}

\subsection{20 {[}M22{]}}\label{m22-17}

Because \((b_k)\) is nondecreasing, moving left (to a smaller index) can
only \emph{decrease} \(b\), and moving right can only \emph{increase}
\(b\). So:

\begin{itemize}
\tightlist
\item
  If you're trying to make \(a_k-b_k\) large, you want \(a\) as large as
  possible and \(b\) as small as possible ⇒ look at prefix
  (left-to-right) maxima of \(a\).
\item
  If you're trying to make \(b_k-a_k\) large, you want \(b\) as large as
  possible and \(a\) as small as possible ⇒ look at suffix
  (right-to-left) minima of \(a\).
\end{itemize}

Fix any \(k\in\{1,\dots,n\}\).

\begin{enumerate}
\def\labelenumi{\arabic{enumi})}
\tightlist
\item
  Maximizers of \(a_k-b_k\) are among left-to-right maxima. Let
  \(i\le k\) be an index with \(a_i=\max\{a_1,\dots,a_k\}\) (so \(i\) is
  a left-to-right maximum). Since \(b\) is nondecreasing,
  \(b_i\le b_k\). Hence
\end{enumerate}

\[
a_i-b_i \;\ge\; a_k-b_i \;\ge\; a_k-b_k.
\]

Therefore the global maximum of \(a_k-b_k\) over all \(k\) is attained
at some left-to-right maximum. Thus

\[
\max_k(a_k-b_k)\;=\;m_L.
\]

\begin{enumerate}
\def\labelenumi{\arabic{enumi})}
\setcounter{enumi}{1}
\tightlist
\item
  Maximizers of \(b_k-a_k\) are among right-to-left minima. Let
  \(j\ge k\) be an index with \(a_j=\min\{a_k,a_{k+1},\dots,a_n\}\) (so
  \(j\) is a right-to-left minimum). Since \(b\) is nondecreasing,
  \(b_j\ge b_k\). Hence
\end{enumerate}

\[
b_j-a_j \;\ge\; b_k-a_j \;\ge\; b_k-a_k.
\]

Therefore the global maximum of \(b_k-a_k\) over all \(k\) is attained
at some right-to-left minimum. Thus

\[
\max_k(b_k-a_k)\;=\;m_R.
\]

Taking absolute values case-by-case,

\[
\max_k|a_k-b_k|
=\max\!\big\{\max_k(a_k-b_k),\;\max_k(b_k-a_k)\big\}
=\max(m_L,m_R),
\]

as required.

Linear-time algorithm

\begin{enumerate}
\def\labelenumi{\arabic{enumi}.}
\tightlist
\item
  Left scan: track \texttt{amax\ =\ -∞}. For \(k=1\) to \(n\): if
  \(a_k\ge amax\) (a left-to-right maximum), update
  \texttt{amax\ =\ a\_k} and update
  \(m_L \leftarrow \max(m_L, a_k-b_k)\).
\item
  Right scan: track \texttt{amin\ =\ +∞}. For \(k=n\) down to \(1\): if
  \(a_k\le amin\) (a right-to-left minimum), update
  \texttt{amin\ =\ a\_k} and update
  \(m_R \leftarrow \max(m_R, b_k-a_k)\).
\item
  Answer: \(M=\max(m_L,m_R)\).
\end{enumerate}

This runs in \(O(n)\) time and \(O(1)\) extra space.

How many indices get checked (random \(a\))?

For a uniformly random order of the \(a\)'s, the expected number of
left-to-right maxima is \(H_n\) and the expected number of right-to-left
minima is also \(H_n\) (each position is a record with probability
\(1/k\) from the left and \(1/(n-k+1)\) from the right). Hence the total
number of \emph{candidates} we examine is on the order of
\(2H_n\approx 2\ln n\) (slightly less after accounting for overlaps),
matching the exercise's remark. 

Sanity check with a tiny example

Let \(a=(5,2,7,3,6)\), \(b=(1,2,4,8,9)\).

\begin{itemize}
\tightlist
\item
  Left-to-right maxima of \(a\): indices \(1\) (5), \(3\) (7). So
  \(m_L=\max\{5-1,\,7-4\}=4\).
\item
  Right-to-left minima of \(a\): scanning from the right, minima at
  indices \(5\) (6) and \(2\) (2). So \(m_R=\max\{9-6,\,2-2\}=3\). Thus
  \(M=\max(m_L,m_R)=4\). Directly checking all \(|a_k-b_k|\) indeed
  gives the maximum \(4\) at \(k=1\).
\end{itemize}

That completes the proof and gives a practical \(O(n)\) method.

\subsection{21 {[}HM21{]}}\label{hm21-8}

Setup

\begin{itemize}
\tightlist
\item
  A \emph{cycle structure} \((\alpha_1,\alpha_2,\ldots)\) means: exactly
  \(\alpha_j\) cycles of length \(j\) for each \(j\ge 1\).
\item
  Feasibility constraint: \(\sum_{j\ge 1} j\,\alpha_j = n\) (otherwise
  the probability is \(0\)).
\end{itemize}

So the problem reduces to: among all \(n!\) permutations, how many have
that exact cycle structure?

Key counting fact

The number of permutations of \(\{1,\dots,n\}\) with cycle structure
\((\alpha_1,\alpha_2,\ldots)\) is

\[
\frac{n!}{\prod_{j\ge 1} \bigl(j^{\alpha_j}\,\alpha_j!\bigr)}.
\]

This is a standard result about conjugacy classes in \(S_n\): once the
cycle lengths are fixed, the class size is
\(n!/\bigl(\prod_j j^{\alpha_j}\alpha_j!\bigr)\). Intuition:

\begin{itemize}
\tightlist
\item
  Start with any one-line ordering of the \(n\) labels: \(n!\)
  possibilities.
\item
  Writing cycles ``around a circle'' makes the \emph{starting point} in
  each cycle irrelevant: divide by \(j\) for each \(j\)-cycle, hence
  \(\prod_j j^{\alpha_j}\).
\item
  Cycles of the same length are indistinguishable in their ordering:
  divide by \(\alpha_j!\) for each \(j\).
\end{itemize}

Probability

Dividing the count by \(n!\), we obtain

\[
\boxed{\;
P(n;\alpha_1,\alpha_2,\ldots)\;=\;
\begin{cases}
\displaystyle \frac{1}{\prod_{j\ge 1} \bigl(j^{\alpha_j}\,\alpha_j!\bigr)}, & \text{if }\sum_j j\,\alpha_j=n,\\[1.25em]
0, & \text{otherwise.}
\end{cases}
\;}
\]

Why this is consistent (quick checks)

\begin{enumerate}
\def\labelenumi{\arabic{enumi}.}
\tightlist
\item
  Normalization Summing over all feasible \((\alpha_j)\) with
  \(\sum j\alpha_j=n\) gives 1. A compact way to see this is via
  exponential generating functions:
\end{enumerate}

\[
\sum_{\sum j\alpha_j=n}
\frac{n!}{\prod_j j^{\alpha_j}\alpha_j!}
= n!\,[z^n]\prod_{j\ge1}\sum_{\alpha_j\ge0}\frac{(z^j)^{\alpha_j}}{j^{\alpha_j}\alpha_j!}
= n!\,[z^n]\exp\!\left(\sum_{j\ge1}\frac{z^j}{j}\right)
= n!\,[z^n]\frac{1}{1-z}=n!.
\]

So the probabilities sum to \(1\).

\begin{enumerate}
\def\labelenumi{\arabic{enumi}.}
\setcounter{enumi}{1}
\tightlist
\item
  Small examples
\end{enumerate}

\begin{itemize}
\tightlist
\item
  \(n=4\), type \((\alpha_1,\alpha_2,\alpha_3,\alpha_4)=(2,1,0,0)\):
  probability \(=1/(1^{2}\,2!\cdot 2^{1}\,1!)=1/4\). There are
  \(24\times 1/4=6\) such permutations, agreeing with ``choose the
  2-cycle's pair'' \(\binom{4}{2}=6\).
\item
  \(n=4\), type \((1,0,1,0)\) (a 3-cycle and a fixed point): probability
  \(=1/(1^{1}\,1!\cdot 3^{1}\,1!)=1/3\). There are \(24/3=8\), the known
  count of 3-cycles in \(S_4\).
\end{itemize}

That's exactly the \(P(n;\alpha_1,\alpha_2,\ldots)\) the exercise asks
for.

Answer:
\(P(n;\alpha_1,\alpha_2,\ldots)=\bigl(\prod_{j\ge1} j^{\alpha_j}\alpha_j!\bigr)^{-1}\)
when \(\sum_j j\alpha_j=n\), and \(0\) otherwise.

\subsection{22 {[}HM22{]}}\label{hm22-4}

\begin{enumerate}
\def\labelenumi{(\alph{enumi})}
\tightlist
\item
  Prove the two Chernoff-style bounds
\end{enumerate}

We must show, for \(0<r\le 1\) and \(r\ge1\),

\[
\Pr(X\le \mu r)\le\bigl(r^{-r}e^{\,r-1}\bigr)^{\mu},
\qquad
\Pr(X\ge \mu r)\le\bigl(r^{-r}e^{\,r-1}\bigr)^{\mu}.
\]

Use (24) with \(x=r\) for the left tail (\(0<r\le1\)):

\[
\Pr(X\le \mu r)\ \le\ r^{-\mu r}G(r)
= r^{-\mu r}\prod_{k=1}^n\bigl(q_k+p_k r\bigr).
\]

For each factor, apply \(1+u\le e^{u}\) (valid for all real \(u\)) with
\(u=(r-1)p_k\):

\[
q_k+p_k r \;=\; 1 + (r-1)p_k \;\le\; e^{(r-1)p_k}.
\]

Hence

\[
\Pr(X\le \mu r)
\le r^{-\mu r}\exp\!\Bigl((r-1)\sum_{k=1}^n p_k\Bigr)
= \bigl(r^{-r}e^{\,r-1}\bigr)^{\mu}.
\]

Exactly the same argument, but using (25) with \(x=r\ge1\), gives the
right-tail bound
\(\Pr(X\ge \mu r)\le\bigl(r^{-r}e^{\,r-1}\bigr)^{\mu}\). 

\begin{enumerate}
\def\labelenumi{(\alph{enumi})}
\setcounter{enumi}{1}
\tightlist
\item
  Simplify the bound when \(r\approx 1\)
\end{enumerate}

Write \(r=1+\delta\) with \(|\delta|\le1\). Consider

\[
\log\bigl(r^{-r}e^{\,r-1}\bigr) \;=\; -r\ln r + r - 1.
\]

A Taylor expansion at \(r=1\) gives

\[
-r\ln r + r - 1
= -\frac{\delta^2}{2} + \frac{\delta^3}{6} - \frac{\delta^4}{12} + \cdots.
\]

Therefore

\[
\bigl(r^{-r}e^{\,r-1}\bigr)^{\mu}
= \exp\!\Bigl(\mu\bigl(-\tfrac{\delta^2}{2}+\tfrac{\delta^3}{6}+\cdots\bigr)\Bigr).
\]

Two convenient one-sided simplifications (using \(|\delta|\le1\)) are:

\begin{itemize}
\item
  If \(\delta\le0\) (\(r\le1\), left tail), the cubic term is \(\le0\),
  so \(\log(\cdot)\le -\tfrac{\delta^2}{2}\). Hence

  \[
  \bigl(r^{-r}e^{\,r-1}\bigr)^{\mu}\le e^{-\mu\delta^2/2}.
  \]
\item
  If \(\delta\ge0\) (\(r\ge1\), right tail), use
  \(\tfrac{\delta^3}{6}\le \tfrac{\delta^2}{6}\) to get
  \(-\tfrac{\delta^2}{2}+\tfrac{\delta^3}{6}\le -\tfrac{\delta^2}{3}\).
  Hence

  \[
  \bigl(r^{-r}e^{\,r-1}\bigr)^{\mu}\le e^{-\mu\delta^2/3}.
  \]
\end{itemize}

These are the familiar ``Gaussian-looking'' tail fall-offs near the
mean.

\begin{enumerate}
\def\labelenumi{(\alph{enumi})}
\setcounter{enumi}{2}
\tightlist
\item
  Show that for large \(r\), \(\Pr(X\ge \mu r)\le 2^{-\mu r}\)
\end{enumerate}

We want

\[
\bigl(r^{-r}e^{\,r-1}\bigr)^{\mu}\ \le\ 2^{-\mu r}.
\]

Equivalently (cancel \(\mu>0\)), \(r^{-r}e^{\,r-1}\le 2^{-r}\), or

\[
\Bigl(\frac{e^{\,1-1/r}}{r}\Bigr)^{r}\le 2^{-r}
\quad\Longleftrightarrow\quad
\frac{e^{\,1-1/r}}{r}\le \frac12.
\]

Define \(h(r)=r^{-1}e^{\,1-1/r}\). Then

\[
\frac{d}{dr}\ln h(r)=\frac{1}{r^2}-\frac{1}{r}=\frac{1-r}{r^2}<0\quad(r>1),
\]

so \(h\) strictly decreases from \(h(1)=1\) to \(0\) as \(r\to\infty\).
Numerically, \(h(2)=e^{1/2}/2\approx0.824<0.825\), and the equation
\(h(r)=1/2\) is solved near \(r\approx4.32\). Thus for sufficiently
large \(r\) (e.g., \(r\ge4.32\)) we have \(h(r)\le\frac12\), hence
\(\Pr(X\ge \mu r)\le 2^{-\mu r}\).

\subsection{23 {[}HM23{]}}\label{hm23-2}

Let \(X\) have the negative binomial distribution whose probability
generating function (pgf) is

\[
G(z)=(q-pz)^{-n},\qquad\text{with }q=p+1 .
\]

Estimate tail probabilities such as \(\Pr(X\ge r)\) and \(\Pr(X\le r)\).
We'll use the pgf tail bounds from §1.2.10:

\[
\Pr(X\ge r)\le \frac{G(x)}{x^r}\quad (x\ge 1),\qquad 
\Pr(X\le r)\le \frac{G(x)}{x^r}\quad (0<x\le 1). \,\,\text{(Markov/pgf)} 
\]

These are (24)--(25) in the text.

\begin{enumerate}
\def\labelenumi{\arabic{enumi})}
\tightlist
\item
  Set up the Chernoff/pgf bound and optimize
\end{enumerate}

For any \(x\in(0,q/p)\),

\[
\frac{G(x)}{x^r}=(q-px)^{-n}\,x^{-r}.
\]

Let

\[
\phi(x):=\ln\!\Big((q-px)^{-n}x^{-r}\Big)=-n\ln(q-px)-r\ln x .
\]

Differentiate and set to zero to minimize:

\[
\phi'(x)=\frac{np}{q-px}-\frac{r}{x}=0
\quad\Longrightarrow\quad
x^\star=\frac{rq}{p(n+r)}.
\]

(One checks \(x^\star\in(0,1]\) when \(r\le np\) and \(x^\star\ge 1\)
when \(r\ge np\), so the correct inequality---lower or upper tail---uses
the same optimizer.) Substituting:

\[
q-px^\star
= q-\frac{r q}{n+r}
= \frac{qn}{n+r}.
\]

Hence the minimized bound is

\[
\boxed{\;
\min_x\frac{G(x)}{x^r}
=\Big(\frac{n+r}{qn}\Big)^{\!n}\,\Big(\frac{p(n+r)}{rq}\Big)^{\!r}
\;}
\]

and therefore

\[
\Pr(X\ge r)\ \le\ \Big(\frac{n+r}{qn}\Big)^{\!n}\,\Big(\frac{p(n+r)}{rq}\Big)^{\!r}
\quad(r\ge np),\qquad
\Pr(X\le r)\ \le\ \Big(\frac{n+r}{qn}\Big)^{\!n}\,\Big(\frac{p(n+r)}{rq}\Big)^{\!r}
\quad(r\le np).
\]

(These follow directly from (24)--(25) with the minimizing \(x^\star\).)

\begin{enumerate}
\def\labelenumi{\arabic{enumi})}
\setcounter{enumi}{1}
\tightlist
\item
  Put the bound in ``deviation from the mean'' form
\end{enumerate}

From the pgf, \(\mathbb{E}X = \dfrac{np}{q-p}\). With \(q=p+1\), this
simplifies to \(\mu=\mathbb{E}X=np\); and
\(\operatorname{Var}(X)=\dfrac{npq}{(q-p)^2}=npq\). (These follow from
the standard ``semi-invariant'' relations in §1.2.10 and a quick
derivative check at \(z=1\).) 

Write \(r=\mu(1+\delta)=np(1+\delta)\) for the upper tail
(\(\delta\ge 0\)). Then

\[
\Pr\!\big(X\ge np(1+\delta)\big)
\ \le\
\Big(\frac{q+p\delta}{q}\Big)^{\!n}
\Big(\frac{q+p\delta}{q(1+\delta)}\Big)^{\!np(1+\delta)} .
\]

Similarly for the lower tail with \(r=np(1-\delta)\)
(\(0\le\delta\le 1\)):

\[
\Pr\!\big(X\le np(1-\delta)\big)
\ \le\
\Big(\frac{q-p\delta}{q}\Big)^{\!n}
\Big(\frac{q-p\delta}{q(1-\delta)}\Big)^{\!np(1-\delta)} .
\]

(These are just the closed-form versions of the boxed bound with
\(q=p+1\).)

\begin{enumerate}
\def\labelenumi{\arabic{enumi})}
\setcounter{enumi}{2}
\tightlist
\item
  Gaussian-type approximation (small deviations)
\end{enumerate}

For small \(\delta\) we can expand the log of the upper-tail bound. A
Taylor expansion around \(\delta=0\) with \(q=p+1\) gives

\[
\ln\text{(bound)}
= -\,\frac{p}{2(p+1)}\,\delta^{2}\,n\;+\;O(\delta^{3}n).
\]

Since \(\mu=np\) and \(\operatorname{Var}(X)=npq=np(p+1)\), this says

\[
\boxed{\;
\Pr\!\big(X\ge \mu+s\big)\ \lesssim\ 
\exp\!\Big(\!-\frac{s^{2}}{2\,\operatorname{Var}(X)}\Big)
\quad\text{for small }s\ (\text{and likewise for }X\le \mu-s).
\;}
\]

So, near the mean, the negative binomial tails are (sub-Gaussian) with
variance scale \(npq\). This is the negative-binomial analogue of the
Chernoff/Hoeffding behavior obtained earlier for the binomial case via
(24)--(25). 

\begin{enumerate}
\def\labelenumi{\arabic{enumi})}
\setcounter{enumi}{3}
\tightlist
\item
  Summary
\end{enumerate}

\begin{itemize}
\item
  Using the pgf tail inequalities (24)--(25), the optimal choice
  \(\displaystyle x^\star=\frac{rq}{p(n+r)}\) yields the explicit bound

  \[
  \Pr(X\ge r)\ \text{or}\ \Pr(X\le r)\ \le\
  \Big(\frac{n+r}{qn}\Big)^{\!n}\Big(\frac{p(n+r)}{rq}\Big)^{\!r}.
  \]
\item
  With \(q=p+1\), mean \(=\!np\), variance \(=\!npq\), and for small
  deviations \(s\) from the mean,

  \[
  \Pr(|X-\mu|\ge s)\ \lesssim\ 2\,\exp\!\Big(\!-\frac{s^{2}}{2npq}\Big).
  \]
\end{itemize}

These estimates are exactly what Exercise 23 asks for when the
distribution is generated by \((q-pz)^{-n}\) with \(q=p+1\), using the
same pgf-Markov machinery as in (24)--(25).

\bookmarksetup{startatroot}

\chapter{1.2.11 Asymptotic
representations}\label{asymptotic-representations}

\subsection{1 {[}HM20{]}}\label{hm20-12}

Base case \(n=0\)

When \(n=0\) the formula reads

\[
f(x)=f(0)+\int_{0}^{x} f'(x-t)\,dt.
\]

Make the change of variables \(u=x-t\). Then \(du=-dt\) and when \(t=0\)
we have \(u=x\), when \(t=x\) we have \(u=0\). Thus

\[
\int_{0}^{x} f'(x-t)\,dt
=\int_{x}^{0} f'(u)\,(-du)
=\int_{0}^{x} f'(u)\,du
=f(x)-f(0),
\]

by the Fundamental Theorem of Calculus. Hence the base case holds.

Inductive step

Assume the statement holds for some \(n-1\ge 0\); we prove it for \(n\).

Write the remainder term

\[
R_n(x)\;:=\;\int_{0}^{x}\frac{t^{\,n}}{n!}\,f^{(n+1)}(x-t)\,dt.
\]

Integrate by parts with

\[
u=\frac{t^{\,n}}{n!},\quad dv=f^{(n+1)}(x-t)\,dt
\qquad\Rightarrow\qquad
du=\frac{t^{\,n-1}}{(n-1)!}\,dt,\;\; v=-\,f^{(n)}(x-t).
\]

Then

\[
R_n(x)=\Big[-\frac{t^{\,n}}{n!}\,f^{(n)}(x-t)\Big]_{t=0}^{t=x}
\;+\;\int_{0}^{x}\frac{t^{\,n-1}}{(n-1)!}\,f^{(n)}(x-t)\,dt.
\]

The boundary term simplifies to

\[
-\frac{x^{n}}{n!}\,f^{(n)}(0)\quad(\text{the }t=0\text{ endpoint is }0),
\]

so

\[
R_n(x)=-\frac{x^{n}}{n!}\,f^{(n)}(0)\;+\;\int_{0}^{x}\frac{t^{\,n-1}}{(n-1)!}\,f^{(n)}(x-t)\,dt
= -\frac{x^{n}}{n!}\,f^{(n)}(0)+R_{n-1}(x).
\]

Now start from the target formula for \(n\):

\[
f(x)\stackrel{?}{=}\sum_{k=0}^{n}\frac{f^{(k)}(0)}{k!}\,x^{k}+R_n(x)
=\sum_{k=0}^{n}\frac{f^{(k)}(0)}{k!}\,x^{k}
-\frac{x^{n}}{n!}\,f^{(n)}(0)+R_{n-1}(x)
=\sum_{k=0}^{n-1}\frac{f^{(k)}(0)}{k!}\,x^{k}+R_{n-1}(x).
\]

But by the induction hypothesis (the statement for \(n-1\)), the
right--hand side equals \(f(x)\). Therefore the statement for \(n\)
holds.

By induction on \(n\), the formula is true for all integers \(n\ge 0\).

\subsection{2 {[}HM20{]}}\label{hm20-13}

Show that for fixed \(a>0\),

\[
\gamma(a,x)=\sum_{k\ge 0}\frac{(-1)^k\,x^{k+a}}{k!\,(k+a)}
\quad\text{(hence } \frac{x^{a}}{a}-\frac{x^{a+1}}{a+1}+\frac{x^{a+2}}{2!(a+2)}-\cdots\text{)}.
\]

Derivation (step by step)

\begin{enumerate}
\def\labelenumi{\arabic{enumi}.}
\tightlist
\item
  Start from the definition (Knuth's Eq. (6)):
\end{enumerate}

\[
\gamma(a,x)=\int_{0}^{x} e^{-t}\,t^{\,a-1}\,dt, \qquad a>0.
\]

The condition \(a>0\) ensures the integrand is integrable at \(t=0\). 

\begin{enumerate}
\def\labelenumi{\arabic{enumi}.}
\setcounter{enumi}{1}
\tightlist
\item
  Expand the exponential. For all real \(t\),
\end{enumerate}

\[
e^{-t}=\sum_{k=0}^{\infty}\frac{(-1)^k t^{k}}{k!}.
\]

\begin{enumerate}
\def\labelenumi{\arabic{enumi}.}
\setcounter{enumi}{2}
\tightlist
\item
  Multiply and justify termwise integration. On the compact interval
  \(t\in[0,x]\) (with fixed \(x\ge 0\)), the series for \(e^{-t}\)
  converges uniformly and absolutely, and \(t^{a-1}\) is continuous
  (since \(a>0\)). By the Weierstrass M-test (or standard
  uniform-convergence results), we may interchange sum and integral:
\end{enumerate}

\[
\gamma(a,x)
=\int_{0}^{x}\!\!\left(\sum_{k=0}^{\infty}\frac{(-1)^k t^{k}}{k!}\right)t^{a-1}dt
=\sum_{k=0}^{\infty}\frac{(-1)^k}{k!}\int_{0}^{x} t^{k+a-1}dt .
\]

\begin{enumerate}
\def\labelenumi{\arabic{enumi}.}
\setcounter{enumi}{3}
\tightlist
\item
  Integrate each term.
\end{enumerate}

\[
\int_{0}^{x} t^{k+a-1}dt=\frac{x^{k+a}}{k+a}\qquad(a>0).
\]

\begin{enumerate}
\def\labelenumi{\arabic{enumi}.}
\setcounter{enumi}{4}
\tightlist
\item
  Assemble the series. Therefore,
\end{enumerate}

\[
\gamma(a,x)=\sum_{k=0}^{\infty}\frac{(-1)^k\,x^{k+a}}{k!\,(k+a)}
= \frac{x^{a}}{a}-\frac{x^{a+1}}{a+1}+\frac{x^{a+2}}{2!(a+2)}-\cdots .
\]

This is exactly Knuth's Eq. (7), the Maclaurin expansion in powers of
\(x\) for the incomplete gamma function. (The exercise statement
``Obtain Eq. (7) from Eq. (6)'' appears on the same page as the exercise
list for §1.2.11.3.) 

\subsection{3 {[}M20{]}}\label{m20-31}

The exercise asks you to get Eq. (8) from Eq. (7). Here are those two
formulas as they appear in the text (with \(a>0\)):

\[
\gamma(a,x)=\int_0^x e^{-t}t^{a-1}\,dt
= \sum_{k\ge 0}\frac{(-1)^k\,x^{k+a}}{k!\,(k+a)},\qquad\text{(Eq. (7))} \; \text{and}
\]

\[
e^x\gamma(a,x)=\sum_{k\ge 0}\frac{x^{k+a}}{a(a+1)\cdots(a+k)}.\qquad\text{(Eq. (8))}
\]

Below are two self-contained derivations.

Method A (Cauchy product + a binomial/Beta identity)

\begin{enumerate}
\def\labelenumi{\arabic{enumi}.}
\tightlist
\item
  Multiply Eq. (7) by the series for \(e^x=\sum_{m\ge0}\dfrac{x^m}{m!}\)
  and use the Cauchy product:
\end{enumerate}

\[
e^x\gamma(a,x)
=\Bigl(\sum_{k\ge0}\frac{(-1)^k\,x^{k+a}}{k!\,(k+a)}\Bigr)\,
 \Bigl(\sum_{m\ge0}\frac{x^m}{m!}\Bigr)
=\sum_{n\ge0}\Bigl(\sum_{k=0}^{n}\frac{(-1)^k}{k!\,(n-k)!\,(k+a)}\Bigr)\,x^{n+a}.
\]

So the coefficient of \(x^{n+a}\) is

\[
c_n=\sum_{k=0}^{n}\frac{(-1)^k}{k!\,(n-k)!\,(k+a)}.
\]

\begin{enumerate}
\def\labelenumi{\arabic{enumi}.}
\setcounter{enumi}{1}
\tightlist
\item
  Evaluate \(c_n\). Use the binomial expansion inside an integral (a
  standard Beta-function trick):
\end{enumerate}

\[
\int_0^1 t^{a-1}(1-t)^n\,dt
=\sum_{k=0}^{n}\binom{n}{k}(-1)^k\int_0^1 t^{a+k-1}\,dt
=\sum_{k=0}^{n}\binom{n}{k}\frac{(-1)^k}{a+k}.
\]

But the LHS is
\(B(a,n+1)=\dfrac{\Gamma(a)\Gamma(n+1)}{\Gamma(a+n+1)} =\dfrac{n!}{a(a+1)\cdots(a+n)}\).
Therefore

\[
\sum_{k=0}^{n}\binom{n}{k}\frac{(-1)^k}{a+k}
=\frac{n!}{a(a+1)\cdots(a+n)}.
\]

Divide both sides by \(n!\) to get

\[
\sum_{k=0}^n \frac{(-1)^k}{k!\,(n-k)!\,(a+k)}
=\frac{1}{a(a+1)\cdots(a+n)}.
\]

Hence \(c_n=\dfrac{1}{a(a+1)\cdots(a+n)}\), which gives

\[
e^x\gamma(a,x)=\sum_{n\ge0}\frac{x^{n+a}}{a(a+1)\cdots(a+n)},
\]

i.e.~Eq. (8).

Method B (one-line integral evaluation)

Start directly from the definition and pull the \(e^x\) inside:

\[
e^x\gamma(a,x)=\int_0^x e^{\,x-t}\,t^{a-1}\,dt.
\]

With \(t=xu\) (so \(dt=x\,du\)) this becomes

\[
e^x\gamma(a,x)=x^a\int_0^1 e^{\,x(1-u)}u^{a-1}\,du
= x^a\sum_{m\ge0}\frac{x^m}{m!}\int_0^1(1-u)^m u^{a-1}\,du.
\]

The integral is
\(B(a,m+1)=\dfrac{\Gamma(a)\Gamma(m+1)}{\Gamma(a+m+1)}=\dfrac{m!}{a(a+1)\cdots(a+m)}\),
so the \(m!\) cancels:

\[
e^x\gamma(a,x)=\sum_{m\ge0}\frac{x^{a+m}}{a(a+1)\cdots(a+m)},
\]

which is Eq. (8).

\subsection{4 {[}HM10{]}}\label{hm10-2}

\begin{enumerate}
\def\labelenumi{\arabic{enumi}.}
\tightlist
\item
  Monotonic change of variables and lower bound on \(u\). From the setup
  we have \(v=v(u)=u-\ln(1+u)\) for \(u\ge 0\). Then
\end{enumerate}

\[
\frac{dv}{du}=1-\frac{1}{1+u}=\frac{u}{1+u}>0\quad (u>0),
\]

so \(v(u)\) is strictly increasing on \([0,\infty)\), with \(v(0)=0\)
and \(v(u)\to\infty\) as \(u\to\infty\). Hence for any fixed \(r>0\)
there is a unique \(u_r>0\) such that

\[
v(u_r)=u_r-\ln(1+u_r)=r,
\]

and for all \(v\ge r\) we have \(u\ge u_r\).

\begin{enumerate}
\def\labelenumi{\arabic{enumi}.}
\setcounter{enumi}{1}
\tightlist
\item
  Uniform bound on the factor \(1+\tfrac1u\). Because \(u\ge u_r>0\)
  whenever \(v\ge r\), the multiplier in the integrand is bounded:
\end{enumerate}

\[
1+\frac{1}{u}\;\le\; 1+\frac{1}{u_r}\;=:\;C_r,
\]

where \(C_r\) depends only on the chosen \(r\) and not on \(x\) or
\(v\).

\begin{enumerate}
\def\labelenumi{\arabic{enumi}.}
\setcounter{enumi}{2}
\tightlist
\item
  Domination by a simple exponential tail. Therefore,
\end{enumerate}

\[
\int_r^\infty x\,e^{-xv}\!\left(1+\frac{1}{u}\right)\,dv
\;\le\; C_r\int_r^\infty x\,e^{-xv}\,dv
\;=\; C_r\big[-e^{-xv}\big]_{v=r}^{\infty}
\;=\; C_r\,e^{-rx}.
\]

Since \(C_r\) is a constant for fixed \(r\), this is exactly the
statement

\[
\int_r^\infty x\,e^{-xv}\!\left(1+\frac{1}{u}\right)\,dv \;=\; O(e^{-rx})\quad (x\to\infty).
\]

This bound justifies the claim in Eq. (13) that the portion of the
integral from \(v=r\) to \(\infty\) is exponentially small in \(x\), so
all asymptotic contributions come from the neighborhood of \(v=0\),
letting Knuth restrict attention to \([0,r]\) for any fixed small
\(r>0\).

\subsection{5 {[}HM24{]}}\label{hm24-1}

Let

\[
\sum_{k=0}^{n} k^{\,n+\tfrac12}\,e^{-k}
=\int_0^n x^{\,n+\tfrac12}e^{-x}\,dx
+\tfrac12\,n^{\,n+\tfrac12}e^{-n}
+\tfrac1{24}\,n^{\,n-\tfrac12}e^{-n}-R,
\]

which is the Euler--Maclaurin/Euler summation expansion used in
§1.2.11.3. Show that the remainder satisfies

\[
R=O\!\left(n^{\,n}e^{-n}\right).
\]

Let \(f(x)=x^{\,n+\frac12}e^{-x}\) for \(x\in[0,n]\). A standard form of
Euler's summation formula with one Bernoulli correction gives

\[
\sum_{k=0}^{n} f(k)
=\int_0^n f(x)\,dx+\frac{f(n)}2+\frac{f'(n)-f'(0)}{12}-R,
\]

with

\[
R=\int_0^n \frac{B_2(\{x\})}{2!}\,f''(x)\,dx,
\]

where \(B_2(u)=u^2-u+\tfrac16\) and \(|B_2(\{x\})|\le \tfrac1{12}\).
Hence

\[
|R|\le \frac1{12}\int_0^n |f''(x)|\,dx.
\tag{★}
\]

We now bound \(\int_0^n |f''|\). Compute derivatives via
\(g(x)=\ln f(x)=(n+\tfrac12)\ln x-x\):

\[
f'(x)=f(x)\,g'(x),\qquad g'(x)=\frac{n+\tfrac12}{x}-1,
\]

\[
f''(x)=f(x)\,\bigl((g'(x))^2+g''(x)\bigr),
\qquad g''(x)=-\frac{n+\tfrac12}{x^2}.
\]

Thus

\[
|f''(x)|\le f(x)\!\left(\Bigl|\tfrac{n+\tfrac12}{x}-1\Bigr|^2+\frac{n+\tfrac12}{x^2}\right).
\tag{1}
\]

Make the change \(x=nt\) on \([0,1]\). Then

\[
f(nt)=n^{\,n+\tfrac12}\,t^{\,n+\tfrac12}\,e^{-nt},
\quad
dx=n\,dt,
\]

and

\[
\Bigl|\tfrac{n+\tfrac12}{nt}-1\Bigr|^2
=\Bigl|\tfrac1t-1+\tfrac{1}{2nt}\Bigr|^2
\le C_1(1-t)^2+\frac{C_2}{n^2},
\qquad
\frac{n+\tfrac12}{(nt)^2}\le \frac{C_3}{n},
\]

for constants \(C_i\) independent of \(n\) and \(t\in(0,1]\). Plugging
these into (1) and integrating gives

\[
\int_0^n |f''(x)|\,dx
\;\le\;
n^{\,n+\tfrac32}
\!\int_0^1\! t^{\,n+\tfrac12}e^{-nt}
\left(C_1(1-t)^2+\frac{C_4}{n}\right)dt.
\tag{2}
\]

Next, use the convex function \(\phi(t)=t-\ln t-1\) on \(0<t\le1\).
Since \(\phi(1)=\phi'(1)=0\) and \(\phi''(t)=1/t^2\ge1\), we have

\[
\phi(t)\ge \frac{(1-t)^2}{2}\quad(0<t\le1).
\]

Also
\(t^{\,n+\tfrac12}e^{-nt}=e^{-n}\,t^{1/2}\,e^{-n\phi(t)}\le e^{-n}e^{-n(1-t)^2/2}\).
Therefore both integrals below are standard Gaussian tails:

\[
\int_0^1 e^{-n(1-t)^2/2}\,dt \asymp n^{-1/2},\qquad
\int_0^1 (1-t)^2 e^{-n(1-t)^2/2}\,dt \asymp n^{-3/2}.
\]

Applying these estimates to (2) yields

\[
\int_0^n |f''(x)|\,dx
\;\le\;
n^{\,n+\tfrac32}\,e^{-n}\!\left(C_1\cdot n^{-3/2}+C_4\cdot \frac1n\cdot n^{-1/2}\right)
=O\!\left(n^{\,n}e^{-n}\right).
\]

Finally, (★) implies \(|R|=O\!\left(n^{\,n}e^{-n}\right)\), which is
what we had to prove. This is exactly the bound quoted just above Eq.
(22) in the text.

\subsection{6 {[}HM20{]}}\label{hm20-14}

We need to prove: For fixed real \(\alpha,\beta\) and \(n\to\infty\),

\[
(n+\alpha)^{\,n+\beta}
= n^{\,n+\beta}\,e^{\alpha}\!\left(1+\alpha\!\left(\beta-\frac{\alpha}{2}\right)\frac{1}{n}
+O\!\left(\frac{1}{n^{2}}\right)\right).
\]

We expand logarithms to second order and then re-exponentiate.

\begin{enumerate}
\def\labelenumi{\arabic{enumi}.}
\tightlist
\item
  Take logs and split terms
\end{enumerate}

\[
\begin{aligned}
\ln\!\bigl((n+\alpha)^{\,n+\beta}\bigr)
&=(n+\beta)\,\ln(n+\alpha) \\
&=(n+\beta)\Bigl(\ln n+\ln\!\bigl(1+\tfrac{\alpha}{n}\bigr)\Bigr).
\end{aligned}
\]

\begin{enumerate}
\def\labelenumi{\arabic{enumi}.}
\setcounter{enumi}{1}
\tightlist
\item
  Series for \(\ln(1+z)\) at \(z=\alpha/n\)
\end{enumerate}

\[
\ln\!\left(1+\frac{\alpha}{n}\right)
=\frac{\alpha}{n}-\frac{\alpha^{2}}{2n^{2}}+O\!\left(\frac{1}{n^{3}}\right).
\]

\begin{enumerate}
\def\labelenumi{\arabic{enumi}.}
\setcounter{enumi}{2}
\tightlist
\item
  Multiply by \(n+\beta\)
\end{enumerate}

\[
\begin{aligned}
(n+\beta)\ln\!\left(1+\frac{\alpha}{n}\right)
&=(n+\beta)\left(\frac{\alpha}{n}-\frac{\alpha^{2}}{2n^{2}}+O(n^{-3})\right)\\
&=\alpha+\frac{\alpha\beta-\alpha^{2}/2}{n}+O\!\left(\frac{1}{n^{2}}\right).
\end{aligned}
\]

\begin{enumerate}
\def\labelenumi{\arabic{enumi}.}
\setcounter{enumi}{3}
\tightlist
\item
  Assemble the logarithm
\end{enumerate}

\[
\begin{aligned}
\ln\!\bigl((n+\alpha)^{\,n+\beta}\bigr)
&=(n+\beta)\ln n \;+\;\alpha\;+\;\frac{\alpha\beta-\alpha^{2}/2}{n}
\;+\;O\!\left(\frac{1}{n^{2}}\right).
\end{aligned}
\]

\begin{enumerate}
\def\labelenumi{\arabic{enumi}.}
\setcounter{enumi}{4}
\tightlist
\item
  Exponentiate Write the small correction as
  \(\varepsilon_n=\dfrac{\alpha\beta-\alpha^{2}/2}{n}+O\!\left(\dfrac{1}{n^{2}}\right)\).
  Then
\end{enumerate}

\[
\begin{aligned}
(n+\alpha)^{\,n+\beta}
&=\exp\!\Big((n+\beta)\ln n+\alpha+\varepsilon_n\Big)\\
&=n^{\,n+\beta}\,e^{\alpha}\,\exp(\varepsilon_n)\\
&=n^{\,n+\beta}\,e^{\alpha}\left(1+\varepsilon_n+O(\varepsilon_n^{2})\right)\\
&=n^{\,n+\beta}\,e^{\alpha}\left(1+\frac{\alpha\beta-\alpha^{2}/2}{n}
+O\!\left(\frac{1}{n^{2}}\right)\right).
\end{aligned}
\]

\begin{enumerate}
\def\labelenumi{\arabic{enumi}.}
\setcounter{enumi}{5}
\tightlist
\item
  Tidy the coefficient
\end{enumerate}

\[
\frac{\alpha\beta-\alpha^{2}/2}{n}
=\alpha\!\left(\beta-\frac{\alpha}{2}\right)\frac{1}{n}.
\]

Thus

\[
(n+\alpha)^{\,n+\beta}
= n^{\,n+\beta}\,e^{\alpha}\!\left(1+\alpha\!\left(\beta-\frac{\alpha}{2}\right)\frac{1}{n}
+O\!\left(\frac{1}{n^{2}}\right)\right),
\]

which is Eq. (23).

\subsection{7 {[}HM30{]}}\label{hm30-5}

Write the integrand as a single exponential and expand it uniformly for
\(0\le u\le yx^{1/4}\):

\[
e^{-u}\Bigl(1+\frac{u}{x}\Bigr)^{x}
=\exp\!\Bigl(-u+x\ln(1+u/x)\Bigr)
=\exp\!\Bigl(
-\frac{u^2}{2x}+\frac{u^3}{3x^2}-\frac{u^4}{4x^3}+\frac{u^5}{5x^4}+O\!\bigl(\tfrac{u^6}{x^5}\bigr)
\Bigr).
\]

The first two terms of this are exactly what Knuth uses to reach eq.
(16); we simply carry the expansion two more orders because the upper
limit now grows with \(x\). (The uniformity is fine since
\(u/x\le y\,x^{-3/4}\to0\) on the whole interval.)

Now expand the outer exponential up to terms contributing \(\ge x^{-2}\)
after integration. Let

\[
E=-\frac{u^2}{2x}+\frac{u^3}{3x^2}-\frac{u^4}{4x^3}+\frac{u^5}{5x^4}+O\!\Bigl(\frac{u^6}{x^5}\Bigr).
\]

Then

\[
\exp(E)=
1+E+\frac{E^2}{2}+\frac{E^3}{6}+\frac{E^4}{24}+O\!\Bigl(\frac{u^9}{x^5}\Bigr),
\]

and collecting by powers of \(1/x\) (keeping only terms that will matter
up to order \(x^{-2}\) once integrated to \(u\sim x^{1/4}\)) gives

\[
\begin{aligned}
\exp(E)=\;&
1
-\frac{u^2}{2x}
+\frac{1}{x^2}\!\Bigl(\frac{u^3}{3}+\frac{u^4}{8}\Bigr)
-\frac{1}{x^3}\!\Bigl(\frac{u^4}{4}+\frac{u^5}{6}+\frac{u^6}{48}\Bigr)\\
&+\frac{1}{x^4}\!\Bigl(\frac{13}{72}u^6+\frac{1}{24}u^7+\frac{1}{384}u^8\Bigr)
\;+\;O\!\Bigl(\frac{u^9}{x^5}\Bigr).
\end{aligned}
\]

Term-by-term integration

Let \(U=y\,x^{1/4}\). Using

\[
\int_{0}^{U}\!u^k\,du=\frac{U^{k+1}}{k+1},
\qquad
U^{k+1}=y^{k+1}\,x^{(k+1)/4},
\]

we integrate each line:

\begin{enumerate}
\def\labelenumi{\arabic{enumi}.}
\item
  Leading term \(\displaystyle \int_0^{U}1\,du = U = y\,x^{1/4}.\)
\item
  Order \(x^{-1}\) term
  \(\displaystyle -\frac{1}{2x}\int_0^{U}u^2\,du = -\frac{U^3}{6x} = -\frac{y^3}{6}\,x^{-1/4}.\)
\item
  Order \(x^{-2}\) group
\end{enumerate}

\[
\frac{1}{x^2}\!\int_0^{U}\!\Bigl(\frac{u^3}{3}+\frac{u^4}{8}\Bigr)\,du
=\frac{U^4}{12x^2}+\frac{U^5}{40x^2}
=\frac{y^4}{12}\,x^{-1}+\frac{y^5}{40}\,x^{-3/4}.
\]

\begin{enumerate}
\def\labelenumi{\arabic{enumi}.}
\setcounter{enumi}{3}
\tightlist
\item
  Order \(x^{-3}\) group
\end{enumerate}

\[
-\frac{1}{x^3}\!\int_0^{U}\!\Bigl(\frac{u^4}{4}+\frac{u^5}{6}+\frac{u^6}{48}\Bigr)\,du
= -\frac{U^5}{20x^3}-\frac{U^6}{36x^3}-\frac{U^7}{336x^3}.
\]

In powers of \(x\):
\(-\dfrac{y^5}{20}x^{-7/4}-\dfrac{y^6}{36}x^{-3/2}-\dfrac{y^7}{336}x^{-5/4}\).

\begin{enumerate}
\def\labelenumi{\arabic{enumi}.}
\setcounter{enumi}{4}
\tightlist
\item
  Order \(x^{-4}\) terms that still contribute \(\ge x^{-2}\) after
  integration
\end{enumerate}

\[
\frac{1}{x^4}\!\int_0^{U}\!\Bigl(\frac{13}{72}u^6+\frac{1}{24}u^7+\frac{1}{384}u^8\Bigr)\,du
=\frac{13}{72}\frac{U^7}{7x^4}+\frac{1}{24}\frac{U^8}{8x^4}+\frac{1}{384}\frac{U^9}{9x^4}.
\]

Among these, \(U^7/x^4\sim x^{-9/4}\) (smaller than \(x^{-2}\), can be
absorbed into the remainder), while

\[
\frac{U^8}{192x^4}=\frac{y^8}{192}\,x^{-2}\quad\text{and}\quad \frac{U^9}{3456x^4}=\frac{y^9}{3456}\,x^{-7/4}
\]

both meet or exceed the \(x^{-2}\) threshold and must be kept.

Putting everything together, an asymptotic expansion for \(J(x)\) as
\(x\to\infty\) (with fixed \(y\)) is

\[
\boxed{
\begin{aligned}
J(x)=\;&y\,x^{1/4}
-\frac{y^{3}}{6}\,x^{-1/4}
+\frac{y^{5}}{40}\,x^{-3/4}
+\frac{y^{4}}{12}\,x^{-1}
-\frac{y^{7}}{336}\,x^{-5/4}
-\frac{y^{6}}{36}\,x^{-3/2}\\
&+\Bigl(-\frac{y^{5}}{20}+\frac{y^{9}}{3456}\Bigr)\,x^{-7/4}
+\frac{y^{8}}{192}\,x^{-2}
\;+\;O\!\bigl(x^{-9/4}\bigr).
\end{aligned}}
\]

\subsection{8 {[}HM30{]}}\label{hm30-6}

Set

\[
E(x,u)\;:=\;-u+x\ln\!\Bigl(1+\frac{u}{x}\Bigr).
\]

Then the integrand is \(e^{E(x,u)}\). For \(|z|<1\),

\[
\ln(1+z)=\sum_{k=1}^{m}(-1)^{k+1}\frac{z^{k}}{k} + \rho_{m+1}(z),
\]

with a remainder satisfying \(|\rho_{m+1}(z)|\le C|z|^{m+1}\) when
\(|z|\le \tfrac12\) (take \(C=2\), say). Here \(z=u/x\). Because
\(u\le f(x)=O(x^r)\) with \(r<1\), we have \(z=O(x^{r-1})\to0\); hence
for all sufficiently large \(x\) and all \(0\le u\le f(x)\) we indeed
have \(|z|\le\tfrac12\).

Therefore

\[
\begin{aligned}
E(x,u)
&= -u + x\!\left[\sum_{k=1}^{m}(-1)^{k+1}\frac{(u/x)^{k}}{k} + \rho_{m+1}(u/x)\right] \\
&= \sum_{k=2}^{m}(-1)^{k}\frac{u^{k}}{k\,x^{k-1}} \;+\; \Delta(x,u),
\end{aligned}
\]

where the truncated exponent

\[
P_m(x,u):=\sum_{k=2}^{m}(-1)^{k}\frac{u^{k}}{k\,x^{k-1}}
\]

matches the one in the exercise, and the remainder term satisfies

\[
|\Delta(x,u)| \;=\; x\,|\rho_{m+1}(u/x)| \;\le\; C\,\frac{u^{m+1}}{x^{m}}.
\]

Comparing integrands

Write

\[
e^{E(x,u)}-e^{P_m(x,u)}
= e^{P_m(x,u)}\bigl(e^{\Delta(x,u)}-1\bigr).
\]

For large \(x\) and \(0\le u\le f(x)\) we have \(|\Delta(x,u)|\le 1\),
hence \(|e^{\Delta}-1|\le 2|\Delta|\). Also \(E(x,u)\le 0\) (since
\(\ln(1+t)\le t\)), so \(e^{E(x,u)}\le 1\); this gives a harmless
uniform bound on \(e^{P_m}\) in our range. Consequently,

\[
\bigl|e^{E(x,u)}-e^{P_m(x,u)}\bigr|
\;\;\ll\;\; \frac{u^{m+1}}{x^{m}}\quad(0\le u\le f(x)).
\]

Integrating the error

Integrate this bound from \(0\) to \(f(x)\):

\[
\left|\int_{0}^{f(x)}\!\! e^{E(x,u)}\,du - \int_{0}^{f(x)}\!\! e^{P_m(x,u)}\,du\right|
\;\;\ll\;\; \int_{0}^{f(x)}\frac{u^{m+1}}{x^{m}}\,du
\;=\; \frac{f(x)^{m+2}}{(m+2)\,x^{m}}.
\]

Using \(f(x)=O(x^{r})\),

\[
\frac{f(x)^{m+2}}{x^{m}}
\;=\; O\!\left(\frac{x^{r(m+2)}}{x^{m}}\right)
\;=\; O\!\bigl(x^{\,rm+2r-m}\bigr)
\;=\; O\!\Bigl(x^{-\,[\,m(1-r)-2r\,]}\Bigr).
\]

We want this to be \(O(x^{-s})\), i.e.

\[
m(1-r)-2r \;\ge\; s
\quad\Longleftrightarrow\quad
m \;\ge\; \frac{s+2r}{1-r}.
\]

Choosing any integer \(m\) with
\(m\ge \bigl\lceil (s+2r)/(1-r)\bigr\rceil\) yields the desired
\(O(x^{-s})\) error term, exactly as required.

Tricomi's special case

If \(f(x)=O(\sqrt{x})\), then \(r=\tfrac12\). With \(s=0\) the condition
gives
\(m\ge \lceil(0+2\cdot \tfrac12)/(1-\tfrac12)\rceil=\lceil 2\rceil=2\).
Thus

\[
\int_{0}^{f(x)}\! e^{-u}\Bigl(1+\tfrac{u}{x}\Bigr)^{x}du
=\int_{0}^{f(x)}\exp\!\Bigl(-\tfrac{u^{2}}{2x}\Bigr)\,du+O(1).
\]

Substitute \(t=u/\sqrt{2x}\):

\[
\int_{0}^{f(x)}\exp\!\Bigl(-\tfrac{u^{2}}{2x}\Bigr)\,du
=\sqrt{2x}\int_{0}^{f(x)/\sqrt{2x}} e^{-t^{2}}\,dt,
\]

giving precisely the printed corollary.

\subsection{9 {[}HM36{]}}\label{hm36}

This is the ``incomplete gamma at a moving point'' problem:

\[
R_x(p)\;:=\;\frac{\gamma(x+1,\,p x)}{\Gamma(x+1)}
\qquad(x\to\infty,\;p\in\mathbb R).
\]

Write \(\gamma(x+1,px)=\int_{0}^{px} t^{x}e^{-t}\,dt\). The integrand
\(f_x(t)=t^{x}e^{-t}=e^{x\ln t-t}\) has a single maximum at \(t=x\).
Hence the large-\(x\) behavior depends on how the upper limit \(px\)
sits relative to the maximizer \(x\). There are three natural regimes.

Define the Cramér (large-deviation) rate function for the gamma/Poisson
family

\[
I(p):=p-1-\ln p \quad(p>0), \qquad I(1)=0,
\]

which satisfies \(I(p)>0\) for \(p\neq1\) and \(I(1)=0\).

\begin{enumerate}
\def\labelenumi{\arabic{enumi})}
\tightlist
\item
  \(0<p<1\) (lower tail: exponentially small)
\end{enumerate}

Here the integration stops \textbf{before} the peak. A standard endpoint
form of Laplace's method (with \(\phi(t)=\ln t-\tfrac{t}{x}\),
\(\psi\equiv1\)) gives, evaluated at the endpoint \(a=px\),

\[
\int_{0}^{px} t^{x}e^{-t}\,dt
\sim e^{x\phi(a)}
\left[\frac{1}{x\phi'(a)}-\frac{\phi''(a)}{x^2\phi'(a)^3}+O(x^{-3})\right],
\]

with \(\phi'(a)=\tfrac{1-p}{px}\), \(\phi''(a)=-\tfrac{1}{p^2x^2}\), and
\(x\phi(a)=x\ln(px)-px\). Divide by Stirling's
\(\Gamma(x+1)\sim\sqrt{2\pi x}\,x^{x}e^{-x}\) to obtain

\[
\boxed{\;
R_x(p)\sim
\frac{e^{-x I(p)}}{\sqrt{2\pi x}}\,
\frac{p}{1-p}\left(1+\frac{1}{(1-p)^2\,x}+O\!\left(x^{-2}\right)\right),
\qquad 0<p<1.}
\]

Thus \(R_x(p)\) decays like \(e^{-xI(p)}\) with the indicated two-term
prefactor.

\begin{enumerate}
\def\labelenumi{\arabic{enumi})}
\setcounter{enumi}{1}
\tightlist
\item
  \(p>1\) (upper tail: \(1-R_x(p)\) is exponentially small)
\end{enumerate}

Now the integral
\(\Gamma(x+1)-\gamma(x+1,px)=\int_{px}^{\infty}t^{x}e^{-t}dt\) starts
\textbf{after} the peak. The same endpoint analysis (with a sign flip
from \(\phi'(a)<0\)) yields

\[
\Gamma(x+1)-\gamma(x+1,px)
\sim e^{x\phi(a)}\left[\frac{1}{x|\phi'(a)|}
-\frac{\phi''(a)}{x^2|\phi'(a)|^{3}}+O(x^{-3})\right].
\]

Hence

\[
\boxed{\;
1-R_x(p)\sim
\frac{e^{-x I(p)}}{\sqrt{2\pi x}}\,
\frac{p}{p-1}\left(1-\frac{1}{(p-1)^2\,x}+O\!\left(x^{-2}\right)\right),
\qquad p>1.}
\]

Consequently \(R_x(p)\to1\) with the same large-deviation exponent
\(I(p)\).

\begin{enumerate}
\def\labelenumi{\arabic{enumi})}
\setcounter{enumi}{2}
\tightlist
\item
  \(p=1\) (central window: CLT/Edgeworth scale)
\end{enumerate}

At \(p=1\) we are integrating up to \(x\), which lies one standard
deviation unit \(z=-(x+1-x)/\sqrt{x+1}=-1/\sqrt{x+1}\) \textbf{below}
the mean \(x+1\). A normal approximation with the first Edgeworth
correction for a \(\mathrm{Gamma}(x\!+\!1,1)\) (skewness
\(2/\sqrt{x+1}\)) gives

\[
R_x(1)=\Pr\!\left\{Y\le x\right\}
=\Phi\!\left(-\frac{1}{\sqrt{x+1}}\right)
+\frac{1}{6}\,\frac{2}{\sqrt{x+1}}\,(1-z^2)\,\varphi(z)
+O\!\left(x^{-1}\right),
\]

where \(\Phi,\varphi\) are the standard normal CDF/PDF. Expanding at
\(z\approx0\) yields

\[
\boxed{\;
R_x(1)=\frac{1}{2}
-\frac{1}{\sqrt{2\pi x}}
+\frac{1}{3\sqrt{2\pi x}}+O\!\left(x^{-1}\right)
=\frac{1}{2}-\frac{2}{3}\frac{1}{\sqrt{2\pi x}}+O\!\left(x^{-1}\right).}
\]

In words: at \(p=1\), the ratio is just below \(\tfrac12\), with a
\(x^{-1/2}\) correction.

\begin{enumerate}
\def\labelenumi{\arabic{enumi})}
\setcounter{enumi}{3}
\tightlist
\item
  \(p\le 0\) (note on the negative case)
\end{enumerate}

The problem allows \(p<0\) with integer \(x\) so that \(t^x\) is defined
for \(t<0\). Writing \(t=-u\) gives
\(\gamma(x+1,px)=-(-1)^x\int_{0}^{|p|x} u^{x}e^{u}\,du\), whose
magnitude is dominated by the endpoint \(u=|p|x\), growing like
\(\sim \dfrac{(|p|x)^x e^{|p|x}}{|1-x/u|}\).

Dividing by \(\Gamma(x+1)\sim\sqrt{2\pi x}\,x^{x}e^{-x}\) shows
\(|R_x(p)|\) grows super-exponentially,

\[
|R_x(p)| \asymp \frac{1}{\sqrt{2\pi x}}\,
\bigl(|p|e^{\,|p|+1}\bigr)^{x}\quad\to\infty\qquad(p<0),
\]

so there is no finite limit; the ratio blows up in magnitude (with sign
\((-1)^{x+1}\)).

Summary (two leading terms in each regime)

\begin{itemize}
\tightlist
\item
  \(0<p<1\):
\end{itemize}

\[
\frac{\gamma(x+1,px)}{\Gamma(x+1)}
\sim \frac{e^{-xI(p)}}{\sqrt{2\pi x}}\,
\frac{p}{1-p}\left(1+\frac{1}{(1-p)^2 x}\right).
\]

\begin{itemize}
\tightlist
\item
  \(p>1\):
\end{itemize}

\[
\frac{\gamma(x+1,px)}{\Gamma(x+1)}
\sim 1-\frac{e^{-xI(p)}}{\sqrt{2\pi x}}\,
\frac{p}{p-1}\left(1-\frac{1}{(p-1)^2 x}\right).
\]

\begin{itemize}
\tightlist
\item
  \(p=1\):
\end{itemize}

\[
\frac{\gamma(x+1,x)}{\Gamma(x+1)}
=\frac{1}{2}-\frac{2}{3}\frac{1}{\sqrt{2\pi x}}+O(x^{-1}).
\]

\begin{itemize}
\tightlist
\item
  \(p<0\): magnitude \(\to\infty\) like above (oscillating in sign with
  \(x\)).
\end{itemize}

These come directly from endpoint Laplace asymptotics (for \(p\neq1\))
and a one-term Edgeworth correction at the central point \(p=1\).

\subsection{10 {[}HM34{]}}\label{hm34}

Let

\[
a=x+1,\quad z_0=px,\quad \Delta=\frac{p\,y}{p-1},
\]

so the quantity in question is

\[
I(x)\;=\;\gamma(a,\,z_0+\Delta)-\gamma(a,\,z_0)\;=\;\int_{z_0}^{z_0+\Delta} e^{-t}\,t^{a-1}\,dt
\;=\;\int_{z_0}^{z_0+\Delta} e^{-t}\,t^{x}\,dt .
\]

Change variables \(t=z_0+s\) (so \(dt=ds\), limits \(s\in[0,\Delta]\)):

\[
I(x)\;=\;e^{-z_0}z_0^{x}\int_{0}^{\Delta} e^{-s}\Bigl(1+\frac{s}{z_0}\Bigr)^{x}\,ds.
\]

Because \(z_0=px\to\infty\) while \(\Delta=O(1)\), expand the logarithm:

\[
\Bigl(1+\frac{s}{px}\Bigr)^{x}
=\exp\!\Bigl[x\ln\!\Bigl(1+\frac{s}{px}\Bigr)\Bigr]
=\exp\!\Bigl(\frac{s}{p}-\frac{s^{2}}{2p^{2}x}+O(x^{-2})\Bigr).
\]

Hence the integrand becomes

\[
e^{-s}\Bigl(1+\frac{s}{px}\Bigr)^{x}
=\exp\!\Bigl(-\alpha s\Bigr)\Bigl(1-\frac{s^{2}}{2p^{2}x}+O(x^{-2})\Bigr),
\qquad
\alpha:=1-\frac{1}{p}=\frac{p-1}{p}.
\]

Therefore

\[
I(x)=e^{-px}(px)^{x}\left[
\int_{0}^{\Delta} e^{-\alpha s}\,ds
-\frac{1}{2p^{2}x}\int_{0}^{\Delta} s^{2}e^{-\alpha s}\,ds
+O(x^{-2})
\right].
\]

The two elementary integrals (for any real \(\alpha\neq 0\)) are

\[
\int_{0}^{\Delta} e^{-\alpha s}\,ds
=\frac{1-e^{-\alpha\Delta}}{\alpha},
\quad
\int_{0}^{\Delta} s^{2}e^{-\alpha s}\,ds
=\frac{2}{\alpha^{3}}-e^{-\alpha\Delta}\!\left(\frac{\Delta^{2}}{\alpha}+\frac{2\Delta}{\alpha^{2}}+\frac{2}{\alpha^{3}}\right).
\]

With \(\alpha=\frac{p-1}{p}\) and \(\Delta=\frac{py}{p-1}\) we have the
key simplification \(\alpha\Delta=y\). Also

\[
\frac{1}{\alpha}=\frac{p}{p-1},\quad
\frac{1}{\alpha^{2}}=\frac{p^{2}}{(p-1)^{2}},\quad
\frac{1}{\alpha^{3}}=\frac{p^{3}}{(p-1)^{3}},
\quad
\frac{\Delta^{2}}{\alpha}=\frac{p^{3}y^{2}}{(p-1)^{3}},
\quad
\frac{2\Delta}{\alpha^{2}}=\frac{2p^{3}y}{(p-1)^{3}}.
\]

Plugging in:

\[
\int_{0}^{\Delta} e^{-\alpha s}\,ds
=\frac{p}{p-1}\bigl(1-e^{-y}\bigr),
\]

\[
\int_{0}^{\Delta} s^{2}e^{-\alpha s}\,ds
=\frac{p^{3}}{(p-1)^{3}}\Bigl(2-e^{-y}(y^{2}+2y+2)\Bigr).
\]

Finally,

\[
\boxed{\;
\gamma(x+1,\,px+\tfrac{py}{p-1})-\gamma(x+1,\,px)
=e^{-px}(px)^{x}\!\left[
\frac{p}{p-1}\bigl(1-e^{-y}\bigr)
+\frac{p}{2(p-1)^{3}x}\Bigl(e^{-y}(y^{2}+2y+2)-2\Bigr)
+O(x^{-2})
\right].
\;}
\]

This gives two terms (leading plus \(1/x\)) before the \(O(x^{-2})\),
matching the requested order ``as in the previous exercise.'' (For
context, the chapter's treatment of the incomplete gamma and asymptotics
appears immediately above; see also Theorem A in §1.2.11.3.)

(Optional) Normalized form

If you prefer it relative to \(\Gamma(x+1)\), Stirling's formula yields

\[
\frac{e^{-px}(px)^{x}}{\Gamma(x+1)}
=\frac{1}{\sqrt{2\pi x}}\bigl(p\,e^{\,1-p}\bigr)^{x}\Bigl(1+O(x^{-1})\Bigr),
\]

so

\[
\frac{\gamma(x+1,\,px+\tfrac{py}{p-1})-\gamma(x+1,\,px)}{\Gamma(x+1)}
=\frac{(p\,e^{\,1-p})^{x}}{\sqrt{2\pi x}}\left[
\frac{p}{p-1}\bigl(1-e^{-y}\bigr)
+\frac{p}{2(p-1)^{3}x}\Bigl(e^{-y}(y^{2}+2y+2)-2\Bigr)
+O(x^{-2})
\right].
\]

(As usual for \(p\neq 1\), this is exponentially small in \(x\).)

\subsection{11 {[}HM35{]}}\label{hm35-2}

\begin{enumerate}
\def\labelenumi{\arabic{enumi})}
\tightlist
\item
  Definitions and exact closed forms
\end{enumerate}

From the section, the ``un-parametrized'' sums are

\[
Q(n)=\sum_{k=1}^{n}\frac{n!}{(n-k)!\,n^k}
\quad\text{and}\quad
R(n)=\sum_{k\ge0}\frac{n!\,n^k}{(n+k)!}.
\]

They satisfy \(Q(n)+R(n)=\dfrac{n!\,e^{\,n}}{n^n}\) and their
asymptotics are derived on the preceding pages.

The exercise introduces a weight \(x^k\). Write

\[
Q_x(n):=\sum_{k=0}^{n}\frac{n!}{(n-k)!\,n^k}\,x^k,\qquad
R_x(n):=\sum_{k\ge0}\frac{n!\,n^k}{(n+k)!}\,x^k .
\]

Both admit clean incomplete-gamma forms (use the Beta--Gamma identities
and the standard expansions used in the text for \(x=1\)):

\begin{itemize}
\tightlist
\item
  For \(R_x(n)\), write
  \(n!/(n+k)! = \dfrac{1}{(n+1)_k} = \dfrac{1}{\Gamma(n+k+1)/\Gamma(n+1)}\),
  use
  \(\mathrm{B}(n+1,k)=\dfrac{n!(k-1)!}{(n+k)!}=\int_0^1 t^n(1-t)^{k-1}\,dt\),
  interchange sum and integral, and sum the resulting exponential
  series. One gets
\end{itemize}

\[
\boxed{\;
R_x(n)=1+\frac{e^{n x}}{(n x)^n}\,\gamma(n+1,n x)\; } \qquad(n\ge1),
\]

where \(\gamma\) is the lower incomplete gamma.

\begin{itemize}
\tightlist
\item
  For \(Q_x(n)\), reindex with \(m=n-k\) and use
  \(\Gamma(n+1,y)=\int_{y}^{\infty} t^n e^{-t}dt\) (upper incomplete
  gamma):
\end{itemize}

\[
\boxed{\;
Q_x(n)=\frac{e^{n/x}}{(n/x)^n}\,\Gamma(n+1,n/x)\; } .
\]

These exact forms reduce to the section's \(x=1\) case when you set
\(x=1\) and use \(\gamma+\Gamma=\Gamma(n+1)=n!\) together with Stirling
and the development around Eqs. (24)--(26).

For asymptotics it is convenient to introduce

\[
\eta(u)=u-1-\ln u \quad(>0 \text{ for } u\neq 1).
\]

\begin{enumerate}
\def\labelenumi{\arabic{enumi})}
\setcounter{enumi}{1}
\tightlist
\item
  Asymptotics when \(0< x<1\) (endpoint/Laplace at \(t=1\))
\end{enumerate}

For \(x<1\), both series are dominated by small \(k\). Intuitively the
term ratios tend to \(x\), so both sums approach the geometric series
\(1/(1-x)\). We can make the leading correction explicit.

\begin{itemize}
\tightlist
\item
  For \(Q_x(n)\): use the falling-factorial expansion
\end{itemize}

\[
\frac{(n)_k}{n^k}=1-\frac{k(k-1)}{2n}+O\!\left(\frac{k^3}{n^2}\right)
\quad (k\ll n),
\]

then sum termwise for \(|x|<1\):

\[
Q_x(n)=\sum_{k\ge0}\left[1-\frac{k(k-1)}{2n}+O\!\left(\frac{k^3}{n^2}\right)\right]x^k
=\frac{1}{1-x}-\frac{x^2}{(1-x)^3}\frac{1}{n}+O\!\left(\frac{1}{n^2}\right).
\]

\begin{itemize}
\tightlist
\item
  For \(R_x(n)\): use
  \(\prod_{j=1}^{k}\tfrac{n}{n+j} =1-\dfrac{k(k+1)}{2n}+O\!\big(\tfrac{k^3}{n^2}\big)\)
  and sum:
\end{itemize}

\[
R_x(n)=\sum_{k\ge0}\left[1-\frac{k(k+1)}{2n}+O\!\left(\frac{k^3}{n^2}\right)\right]x^k
=\frac{1}{1-x}-\frac{x}{(1-x)^3}\frac{1}{n}+O\!\left(\frac{1}{n^2}\right).
\]

So, uniformly for fixed \(x\in(0,1)\),

\[
\boxed{\;
\begin{aligned}
Q_x(n) &= \frac{1}{1-x}-\frac{x^2}{(1-x)^3}\frac{1}{n}+O(n^{-2}),\\[2mm]
R_x(n) &= \frac{1}{1-x}-\frac{x}{(1-x)^3}\frac{1}{n}+O(n^{-2}).
\end{aligned}}
\]

(A complementary derivation comes from the integral
\(R_x(n)-1=n x\int_0^1 t^n e^{n x(1-t)}dt\) and an endpoint Laplace
expansion at \(t=1\).)

\begin{enumerate}
\def\labelenumi{\arabic{enumi})}
\setcounter{enumi}{2}
\tightlist
\item
  Asymptotics when \(x>1\) (interior saddle)
\end{enumerate}

For \(x>1\), the maximal contribution is interior and the sums grow
exponentially with \(n\). A discrete Laplace/saddle-point analysis on
the \(k\)-sums (or standard large-\(n\) expansions of incomplete gamma
with argument \(n x\) / \(n/x\)) yields:

\[
\boxed{\;
\begin{aligned}
R_x(n)&\sim \sqrt{2\pi n}\;\exp\!\big(n\,\eta(x)\big)\,\Big(1+O(n^{-1})\Big),\\[2mm]
Q_x(n)&\sim \sqrt{2\pi n}\;\exp\!\big(n\,\eta(1/x)\big)\,\Big(1+O(n^{-1})\Big).
\end{aligned}}
\]

Sketch (for \(R_x\)): write the \(k\)-term via Stirling,

\[
T_k \approx \frac{1}{\sqrt{1+k/n}}
\exp\!\Big(n\Phi(u)\Big),\quad u=k/n,\quad
\Phi(u)=-\!(1+u)\ln(1+u)+u(1+\ln x),
\]

whose stationary point is at \(u^{*} = x-1\) with
\(\Phi(u^{*})=\eta(x)\) and \(\Phi''(u^{*})=-1/x\). Discrete Laplace
gives the prefactor
\(\frac{1}{\sqrt{1+u^{*}}}\sqrt{\frac{2\pi n}{|\Phi''(u^{*})|}}=\sqrt{2\pi n}\).
A completely analogous computation for \(Q_x\) has \(u^{*}=1-1/x\) and
main exponent \(\eta(1/x)\).

\begin{quote}
Non-uniformity at \(x=1\). These ``\(x\neq 1\)'' expansions are not
uniform as \(x\to 1\). At \(x=1\) the interior saddle coalesces with the
endpoint, and the section gives the correct matched expansions (for the
original \(Q(n),R(n)\)):

\[
\begin{aligned}
Q(n)&=\sqrt{\frac{\pi n}{2}}-\frac{1}{3}+\frac{1}{12}\sqrt{\frac{\pi}{2n}}-\frac{4}{135n}+\frac{1}{288}\sqrt{\frac{\pi}{2n^3}}+O(n^{-2}),\\
R(n)&=\sqrt{\frac{\pi n}{2}}+\frac{1}{3}+\frac{1}{12}\sqrt{\frac{\pi}{2n}}+\frac{4}{135n}+\frac{1}{288}\sqrt{\frac{\pi}{2n^3}}+O(n^{-2}),
\end{aligned}
\]

whose derivation appears just above the exercise.
\end{quote}

\begin{enumerate}
\def\labelenumi{\arabic{enumi})}
\setcounter{enumi}{3}
\tightlist
\item
  Quick consistency checks
\end{enumerate}

\begin{itemize}
\tightlist
\item
  For \(x<1\): \(Q_x(n),R_x(n)\to (1-x)^{-1}\), matching the intuition
  that each term is \(\approx x^k\).
\item
  For \(x>1\): \(\eta(x)>\eta(1/x)>0\), so \(R_x(n)\) dominates
  \(Q_x(n)\) exponentially---again consistent with the shapes of the
  terms.
\item
  Setting \(x=1\) is singular
\end{itemize}

\subsection{12 {[}HM20{]}}\label{hm20-15}

Use the substitution \(u = t^{2}/2\).

\begin{itemize}
\tightlist
\item
  Then \(t=\sqrt{2u}\) and \(dt=\dfrac{du}{\sqrt{2u}}\).
\item
  The limits map as \(t:0\to x \;\Rightarrow\; u:0\to x^{2}/2\).
\end{itemize}

Therefore

\[
\int_{0}^{x} e^{-t^{2}/2}\,dt
=\int_{0}^{x^{2}/2} e^{-u}\,\frac{du}{\sqrt{2u}}
=\frac{1}{\sqrt{2}}\int_{0}^{x^{2}/2} e^{-u}\,u^{-1/2}\,du
=\frac{1}{\sqrt{2}}\;\gamma\!\left(\tfrac{1}{2},\,\tfrac{x^{2}}{2}\right).
\]

Comparing with \(b\,\gamma(a,y)\), we can read off

\[
\boxed{\,a=\tfrac{1}{2},\qquad y=\tfrac{x^{2}}{2},\qquad b=\tfrac{1}{\sqrt{2}}\, }.
\]

Quick check (optional)

Using \(\gamma(\tfrac12,z)=\sqrt{\pi}\,\mathrm{erf}(\sqrt{z})\), the
result becomes

\[
\int_{0}^{x}\! e^{-t^{2}/2}\,dt
=\frac{1}{\sqrt{2}}\sqrt{\pi}\,\mathrm{erf}\!\left(\frac{x}{\sqrt{2}}\right)
=\sqrt{\frac{\pi}{2}}\;\mathrm{erf}\!\left(\frac{x}{\sqrt{2}}\right),
\]

which is the familiar normal-integral form, confirming the constants are
correct.

\subsection{13 {[}HM42{]}}\label{hm42}

\subsection{14 {[}HM39{]}}\label{hm39}

\begin{enumerate}
\def\labelenumi{\arabic{enumi})}
\tightlist
\item
  Re-center at the peak (part a)
\end{enumerate}

Write \(k=n-t\) (so \(t=0,1,\dots,n\)). Then

\[
k^{\,n+\alpha}e^{-k}
=(n-t)^{n+\alpha}e^{-(n-t)}
=n^{n+\alpha}e^{-n}\,\exp\!\Big((n+\alpha)\ln(1-\tfrac{t}{n})+t\Big).
\]

Using \(\ln(1-x)= -x-\tfrac{x^2}{2}-\tfrac{x^3}{3}-\cdots\), the
exponent becomes

\[
t-(n+\alpha)\!\left(\frac{t}{n}+\frac{t^2}{2n^2}+\frac{t^3}{3n^3}+\frac{t^4}{4n^4}+\cdots\right)
=-\frac{\alpha t}{n}-\frac{t^2}{2n}-\frac{t^3}{3n^2}-\frac{t^4}{4n^3}-\cdots.
\]

Hence

\[
S_{n,\alpha}
= n^{n+\alpha}e^{-n}\sum_{t=0}^{n} 
\exp\!\Big(-\frac{t^2}{2n}\Big)\cdot
f(t,n),
\]

with

\[
f(t,n)
= \Big(1-\frac{t}{n}\Big)^{\alpha}
\exp\!\Big(-\frac{t^3}{3n^2}-\frac{t^4}{4n^3}-\cdots\Big).
\]

Renaming \(t\) back to \(k\) gives exactly the form asked for in (a).

\begin{enumerate}
\def\labelenumi{\arabic{enumi})}
\setcounter{enumi}{1}
\tightlist
\item
  Local power series for the slowly varying factor (part b)
\end{enumerate}

For \(k\le n^{1/2+\varepsilon}\) (any fixed \(\varepsilon>0\)), we may
expand \(f(k,n)\) as a double series in \(k\) and \(1/n\). To the orders
we will actually need,

\[
\Big(1-\frac{k}{n}\Big)^{\alpha}
=1-\frac{\alpha k}{n}+\frac{\alpha(\alpha-1)}{2}\frac{k^2}{n^2}+O\!\Big(\frac{k^3}{n^3}\Big),
\]

\[
\exp\!\Big(-\frac{k^3}{3n^2}-\frac{k^4}{4n^3}+\frac12\frac{k^6}{9n^4}+\cdots\Big)
=1-\frac{k^3}{3n^2}-\frac{k^4}{4n^3}
+\frac{1}{18}\frac{k^6}{n^4}+O\!\Big(\frac{k^5}{n^4}\Big).
\]

Multiplying these gives \(f(k,n)\) as a finite polynomial in \(k\)
divided by powers of \(n\), plus a remainder that (for our chosen cutoff
\(k\le n^{1/2+\varepsilon}\)) is \(o(n^{-1/2})\) after summation.
Contributions from \(k>n^{1/2+\varepsilon}\) are exponentially small
because of the Gaussian factor \(e^{-k^2/(2n)}\). (This is the standard
Laplace method localization.)

\begin{enumerate}
\def\labelenumi{\arabic{enumi})}
\setcounter{enumi}{2}
\tightlist
\item
  Gaussian ``moments'' you will need (part d)
\end{enumerate}

Let

\[
M_r(n)\;=\;\sum_{k\ge0} k^{\,r}\,e^{-k^2/(2n)}.
\]

Euler's summation (or Poisson summation) gives the classical leading
terms

\[
\begin{aligned}
M_0(n) &\;=\; \int_0^\infty e^{-x^2/(2n)}dx \;+\;\frac12\;+\;O(e^{-c n})
= \sqrt{\frac{\pi n}{2}}+\frac12+O(e^{-c n}),\\
M_1(n) &\;=\; \int_0^\infty x\,e^{-x^2/(2n)}dx \;+\;O(e^{-c n})
= n+O(e^{-c n}),\\
M_2(n) &\;=\; \int_0^\infty x^2 e^{-x^2/(2n)}dx \;+\;O(e^{-c n})
= \sqrt{\frac{\pi}{2}}\,n^{3/2} + O(e^{-c n}),\\
M_3(n) &\;=\; \int_0^\infty x^3 e^{-x^2/(2n)}dx \;+\;O(e^{-c n})
= 2n^{2} + O(e^{-c n}),\\
M_4(n) &\;=\; \int_0^\infty x^4 e^{-x^2/(2n)}dx \;+\;O(e^{-c n})
= 3\sqrt{\frac{\pi}{2}}\,n^{5/2} + O(e^{-c n}),\\
M_6(n) &\;=\; \int_0^\infty x^6 e^{-x^2/(2n)}dx \;+\;O(e^{-c n})
= 15\sqrt{\frac{\pi}{2}}\,n^{7/2} + O(e^{-c n}).
\end{aligned}
\]

(Only these are needed through the \(n^{-1/2}\) term.)

\begin{enumerate}
\def\labelenumi{\arabic{enumi})}
\setcounter{enumi}{3}
\tightlist
\item
  Collect terms up to \(O(n^{-1/2})\)
\end{enumerate}

Insert the expansion of \(f(k,n)\) into \(\sum e^{-k^2/(2n)} f(k,n)\)
and use the \(M_r(n)\) above.

• Leading term (order \(n^{1/2}\)):

\[
1\cdot M_0(n)\;=\;\sqrt{\frac{\pi n}{2}}\;+\;\frac12\;+\cdots
\]

(contributes \(\sqrt{\pi n/2}\) and a \(+\,\tfrac12\) constant).

• Constant term from \(-(\alpha/n)\sum k e^{-k^2/(2n)}\):

\[
-\frac{\alpha}{n}\,M_1(n)\;=\;-\alpha\;+\;o(1).
\]

• Constant term from the cubic in the exponent:

\[
-\frac{1}{3n^2} \sum k^3 e^{-k^2/(2n)}
\;=\;-\frac{1}{3n^2}\,M_3(n)
\;=\;-\frac{1}{3n^2}\cdot 2n^2
\;=\;-\frac{2}{3}.
\]

Together with \(+\,\tfrac12\) from \(M_0\), the ``purely n-independent''
part becomes \(\; \tfrac12 - \tfrac23 - \alpha = -\tfrac16 - \alpha\).

• The \(n^{-1/2}\) term (the coefficient of \(\sqrt{\pi/(2n)}\))
receives four contributions:

\begin{enumerate}
\def\labelenumi{(\roman{enumi})}
\tightlist
\item
  from \(\dfrac{\alpha(\alpha-1)}{2}\dfrac{k^2}{n^2}\):
\end{enumerate}

\[
\frac{\alpha(\alpha-1)}{2n^2}\,M_2(n)
= \frac{\alpha(\alpha-1)}{2}\sqrt{\frac{\pi}{2n}}.
\]

\begin{enumerate}
\def\labelenumi{(\roman{enumi})}
\setcounter{enumi}{1}
\tightlist
\item
  from the quartic \(-\dfrac{k^4}{4n^3}\):
\end{enumerate}

\[
-\frac{1}{4n^3}\,M_4(n)
= -\frac{1}{4n^3}\cdot 3\sqrt{\frac{\pi}{2}}\,n^{5/2}
= -\frac{3}{4}\sqrt{\frac{\pi}{2n}}.
\]

\begin{enumerate}
\def\labelenumi{(\roman{enumi})}
\setcounter{enumi}{2}
\tightlist
\item
  from the cross term
  \(\big(-\frac{\alpha k}{n}\big)\big(-\frac{k^3}{3n^2}\big)\):
\end{enumerate}

\[
+\frac{\alpha}{3n^3}\,M_4(n)
= \alpha\cdot \frac{1}{3n^3}\cdot 3\sqrt{\frac{\pi}{2}}\,n^{5/2}
= \alpha\,\sqrt{\frac{\pi}{2n}}.
\]

\begin{enumerate}
\def\labelenumi{(\roman{enumi})}
\setcounter{enumi}{3}
\tightlist
\item
  from
  \(\tfrac12\big(\frac{k^3}{3n^2}\big)^2=\frac{1}{18}\frac{k^6}{n^4}\):
\end{enumerate}

\[
\frac{1}{18n^4}\,M_6(n)
= \frac{1}{18n^4}\cdot 15\sqrt{\frac{\pi}{2}}\,n^{7/2}
= \frac{5}{6}\sqrt{\frac{\pi}{2n}}.
\]

Summing (i)--(iv) gives

\[
\left[\frac{\alpha(\alpha-1)}{2} - \frac{3}{4} + \alpha + \frac{5}{6}\right]\sqrt{\frac{\pi}{2n}}
=\left(\frac{1}{12} + \frac{\alpha}{2} + \frac{\alpha^{2}}{2}\right)\sqrt{\frac{\pi}{2n}}.
\]

\begin{enumerate}
\def\labelenumi{\arabic{enumi})}
\setcounter{enumi}{4}
\tightlist
\item
  Final asymptotic (part e)
\end{enumerate}

Multiplying the bracket by the common prefactor \(n^{n+\alpha}e^{-n}\)
yields

\[
\boxed{\;
\sum_{k=0}^{n} k^{\,n+\alpha} e^{-k}
\;=\;
n^{\,n+\alpha} e^{-n}\left[
\sqrt{\frac{\pi n}{2}}
\;-\;\Big(\frac{1}{6}+\alpha\Big)
\;+\;\Big(\frac{1}{12}+\frac{\alpha}{2}+\frac{\alpha^2}{2}\Big)\sqrt{\frac{\pi}{2n}}
\;+\;O\!\Big(\frac{1}{n}\Big)
\right].
\;}
\]

This matches the expansion the exercise asks you to obtain; the same
page notes that, in principle, one can continue to any \(O(n^{-r})\) by
pushing the series for \(f(k,n)\) and the Gaussian moments further.

\subsection{15 {[}HM20{]}}\label{hm20-16}

\begin{enumerate}
\def\labelenumi{\arabic{enumi}.}
\tightlist
\item
  Expand the power by the binomial theorem.
\end{enumerate}

\[
\Bigl(1+\frac{z}{n}\Bigr)^n=\sum_{k=0}^{n}\binom{n}{k}\frac{z^{k}}{n^{k}}.
\]

\begin{enumerate}
\def\labelenumi{\arabic{enumi}.}
\setcounter{enumi}{1}
\tightlist
\item
  Integrate termwise (finite sum, so interchange is justified) and use
  \(\displaystyle \int_{0}^{\infty} z^{k}e^{-z}\,dz=\Gamma(k+1)=k!\):
\end{enumerate}

\[
\begin{aligned}
I(n)
&=\int_{0}^{\infty}e^{-z}\sum_{k=0}^{n}\binom{n}{k}\frac{z^{k}}{n^{k}}\,dz
=\sum_{k=0}^{n}\binom{n}{k}\frac{1}{n^{k}}\int_{0}^{\infty}z^{k}e^{-z}\,dz\\[4pt]
&=\sum_{k=0}^{n}\binom{n}{k}\frac{k!}{n^{k}}
=\sum_{k=0}^{n}\frac{n!}{(n-k)!\,n^{k}},
\end{aligned}
\]

because \(\binom{n}{k}k!=\dfrac{n!}{(n-k)!}\).

\begin{enumerate}
\def\labelenumi{\arabic{enumi}.}
\setcounter{enumi}{2}
\tightlist
\item
  Identify \(Q(n)\). The \(k=0\) term of the last sum is
\end{enumerate}

\[
\frac{n!}{n!\,n^{0}}=1,
\]

and the remaining terms (for \(k=1,\dots,n\)) are exactly the definition
of \(Q(n)\). Therefore

\[
\boxed{\,I(n)=1+Q(n)\,}
\quad\text{equivalently}\quad
\boxed{\,Q(n)=\int_{0}^{\infty}\!\Bigl(1+\frac{z}{n}\Bigr)^{n}e^{-z}\,dz-1.}
\]

This is precisely the ``relation'' the exercise asks you to show. 

Quick checks \& context

\begin{itemize}
\tightlist
\item
  For \(n=1\): \(I(1)=\int_{0}^{\infty}(1+z)e^{-z}dz=1+1=2\), while
  \(Q(1)=1\). Indeed \(I(1)=1+Q(1)\).
\item
  Knuth later derives the asymptotic
  \(Q(n)=\sqrt{\pi n/2}-\tfrac{1}{3}+\cdots\), so this integral is
  \(\sqrt{\pi n/2}+\tfrac{2}{3}+\cdots\), consistent with the same
  section's expansions.
\end{itemize}

\subsection{16 {[}M24{]}}\label{m24-14}

\begin{enumerate}
\def\labelenumi{\arabic{enumi}.}
\tightlist
\item
  Definition of \(Q(\cdot)\). Knuth defines
\end{enumerate}

\[
Q(n)=\sum_{j=1}^{n}\frac{n!}{(n-j)!\,n^{\,j}}.
\]

Equivalently,

\[
Q(k)=\frac{k!}{k^{k}}\sum_{s=0}^{k-1}\frac{k^{s}}{s!}.
\]

(Just put \(s=k-j\).) 

\begin{enumerate}
\def\labelenumi{\arabic{enumi}.}
\setcounter{enumi}{1}
\tightlist
\item
  A very handy integral representation (from Ex. 15): expanding
  \((1+z/k)^k=\sum_{j=0}^{k}\binom{k}{j}(z/k)^j\) and integrating
  term-wise gives
\end{enumerate}

\[
\int_{0}^{\infty}\!\!\Big(1+\frac{z}{k}\Big)^{k}e^{-z}\,dz
=\sum_{j=0}^{k}\binom{k}{j}\frac{j!}{k^{j}}
=1+Q(k),
\]

hence

\[
Q(k)=\int_{0}^{\infty}\!\!\Big[\Big(1+\frac{z}{k}\Big)^{k}-1\Big]e^{-z}\,dz .
\]

Let

\[
S_n:=\sum_{k=0}^{n}(-1)^k\binom{n}{k}\,k^{\,n-1}Q(k).
\]

Insert the integral form of \(Q(k)\) and swap sum/integral (finite sum):

\[
S_n=\int_{0}^{\infty}\!e^{-z}\,
\sum_{k=0}^{n}(-1)^k\binom{n}{k}\,k^{\,n-1}\Big[\Big(1+\frac{z}{k}\Big)^{k}-1\Big]\,dz .
\]

The ``\(-1\)'' part vanishes because
\(\sum_{k=0}^{n}(-1)^k\binom{n}{k}k^{n-1}=0\) (it's the \(n\)th forward
difference of a degree-\(n\!-\!1\) polynomial). Thus

\[
S_n=\int_{0}^{\infty}\!e^{-z}\,
\sum_{k=0}^{n}(-1)^k\binom{n}{k}\,k^{\,n-1}\Big(1+\frac{z}{k}\Big)^{k}\,dz .
\]

Expand \(\big(1+\frac{z}{k}\big)^k=\sum_{j=0}^{k}\binom{k}{j}(z/k)^j\)
and use \(\binom{n}{k}\binom{k}{j}=\binom{n}{j}\binom{n-j}{\,k-j\,}\):

\[
\sum_{k=0}^{n}(-1)^k\binom{n}{k}\,k^{\,n-1}\Big(1+\frac{z}{k}\Big)^{k}
=\sum_{j=0}^{n}\binom{n}{j}z^{j}
\sum_{k=j}^{n}(-1)^k\binom{n-j}{k-j}\,k^{\,n-1-j}.
\]

Let \(N=n-j\) and \(m=k-j\). The inner sum becomes

\[
(-1)^j\sum_{m=0}^{N}(-1)^m\binom{N}{m}\,(m+j)^{N-1}.
\]

This is an alternating binomial sum of a polynomial in \(m\) of degree
\(N-1\), so it is \(0\) when \(N\ge1\) (standard ``\(n\)th difference
kills degree \(<n\)''). The only surviving case is \(N=0\)
(i.e.~\(j=n\)), where the inner sum reduces to
\((-1)^n\cdot \binom{0}{0}\cdot n^{-1}=(-1)^n/n\). Therefore

\[
\sum_{k=0}^{n}(-1)^k\binom{n}{k}\,k^{\,n-1}\Big(1+\frac{z}{k}\Big)^{k}
=\frac{(-1)^n}{n}\,z^{n}.
\]

Integrating,

\[
S_n=\frac{(-1)^n}{n}\int_{0}^{\infty} e^{-z} z^{n}\,dz
=\frac{(-1)^n}{n}\,\Gamma(n+1)
=(-1)^n\,(n-1)!,
\]

which is exactly the desired identity.

\subsection{17 {[}HM29{]}}\label{hm29}

Step 1: Reindex and insert Stirling

Put \(m=n+k\) so \(k=m-n\) and

\[
S(n)=\frac1{(n-1)!}\sum_{m=n}^{\infty}\frac{(m-1)!}{m^{\,m-n}}.
\]

Use Stirling for \((m-1)!=\Gamma(m)\):

\[
(m-1)! \sim \sqrt{2\pi}\,m^{\,m-\tfrac12}e^{-m}\Bigl(1+\frac1{12m}+\frac1{288m^{2}}+\cdots\Bigr).
\]

After dividing by \(m^{\,m-n}\) this gives

\[
S(n)\sim \frac{\sqrt{2\pi}}{(n-1)!}
\sum_{m=n}^{\infty} m^{\,n-\tfrac12}e^{-m}
\Bigl(1+\tfrac1{12m}+\tfrac1{288m^{2}}+\cdots\Bigr).
\]

Define \(A_r(n)=\sum_{m=n}^{\infty} m^{\,n-\tfrac12-r}e^{-m}\). Then

\[
S(n)\sim \frac{\sqrt{2\pi}}{(n-1)!}\Bigl(A_0(n)+\tfrac1{12}A_1(n)+\tfrac1{288}A_2(n)+\cdots\Bigr).
\]

Step 2: Replace the sums by incomplete-gamma terms

For large \(n\), Euler's summation + the standard passage from sums to
integrals gives

\[
A_r(n)=\Gamma\!\left(n+\tfrac12-r,\,n\right)+\text{(lower-order endpoint corrections)}.
\]

Here \(\Gamma(a,x)=\int_x^\infty t^{a-1}e^{-t}\,dt\) is the upper
incomplete gamma function.

Knuth proves a uniform expansion for the \emph{lower} incomplete gamma
ratio

\[
\frac{\gamma(x+1,x+y)}{\Gamma(x+1)}
=\frac12+\frac{y-\tfrac23}{\sqrt{2\pi}}\,x^{-1/2}
+\frac{1}{\sqrt{2\pi}}\!\left(\frac{23}{270}-\frac{y}{12}-\frac{y^{3}}{6}\right)x^{-3/2}
+O(x^{-5/2}),
\]

valid for fixed \(y\) and large \(x\). Hence the \emph{upper} ratio is
\(1-\)that. We need \(x=n-\tfrac12-r\) and \(y=\tfrac12\) (because our
lower limit of summation is \(m=n\)). Substituting these into the
formula, and then multiplying by \(\Gamma(x+1)=\Gamma(n+\tfrac12-r)\)
and by the prefactor \(\sqrt{2\pi}/(n-1)!\), yields an expansion for
each \(A_r(n)\), hence for \(S(n)\). (See the derivation and Theorem A
in the text. )

Finally, use Stirling again for \((n-1)!\) to express everything in
powers of \(n^{1/2},n^{-1/2},n^{-1},n^{-3/2},\ldots\) (this is exactly
the same bookkeeping Knuth used just above for \(P(n),Q(n),R(n)\) in
Eqs. (24)--(26)). 

Carrying the algebra through to the same order as those formulas gives

\[
\boxed{
S(n)=\sqrt{\frac{\pi n}{2}}
+\frac{2}{3}
-\frac{1}{24}\sqrt{\frac{\pi}{2n}}
-\frac{4}{135}\frac{1}{n}
+\frac{49}{1152}\sqrt{\frac{\pi}{2n^{3}}}
+O\!\left(n^{-2}\right)
}.
\]

\subsection{18 {[}M25{]}}\label{m25-24}

Define

\[
S_n(x,y)=\sum_{k=0}^{n}\binom{n}{k}\,(x+k)^{k}\,(y+n-k)^{\,n-k}.
\]

Step 1 --- split \((x+k)^k\): \((x+k)^k = x(x+k)^{k-1}+k(x+k)^{k-1}\).
Therefore

\[
S_n(x,y)=x\!\sum_{k=0}^{n}\binom{n}{k}(x+k)^{k-1}(y+n-k)^{n-k}
+\sum_{k=0}^{n}\binom{n}{k}k(x+k)^{k-1}(y+n-k)^{n-k}.
\]

Step 2 --- Abel's binomial identity: Abel's formula (TAOCP Eq.
1.2.6--(16) with \(z=1\)) gives

\[
\sum_{k=0}^{n}\binom{n}{k}(x+k)^{k-1}(y+n-k)^{n-k}=\frac{(x+y+n)^{n}}{x}.
\]

Thus the first sum equals \((x+y+n)^n\). For the second sum use
\(k\binom{n}{k}=n\binom{n-1}{k-1}\) and reindex \(j=k-1\) to get

\[
\sum_{k=0}^{n}\binom{n}{k}k(\cdots)=
n\sum_{j=0}^{n-1}\binom{n-1}{j}(x+1+j)^{j}(y+n-1-j)^{n-1-j}
= n\,S_{n-1}(x+1,y).
\]

Therefore we have the recurrence

\[
S_n(x,y)=(x+y+n)^n + n\,S_{n-1}(x+1,y). \quad({*}) 
\]

(Abel reference and this reduction are exactly how Knuth dispatches this
exercise.) 

Step 3 --- solve \(({*})\) by induction: One easily verifies (start at
\(n=0\)) that the closed form

\[
\boxed{\;S_n(x,y)=\sum_{k=0}^{n}\binom{n}{k}\,k!\,(x+y+n)^{\,n-k}\;}
\]

satisfies the recurrence: the \(k=0\) term gives \((x+y+n)^n\), and the
\(k\ge1\) coefficients match because
\(n\binom{n-1}{k-1}(k-1)! = \binom{n}{k}k!\).

Plug in the two required cases

\begin{enumerate}
\def\labelenumi{\arabic{enumi})}
\tightlist
\item
  \(S_1(n)\): take \(x=0,\;y=0\)
\end{enumerate}

\[
S_1(n)=\sum_{k=0}^{n}\binom{n}{k}k!\,n^{\,n-k}
= n^{\,n}\sum_{k=0}^{n}\frac{n!}{(n-k)!\,n^{k}}
= n^{\,n}\bigl(1+Q(n)\bigr).
\]

(The \(k=0\) term contributes the ``\(1\)''.) 

\begin{enumerate}
\def\labelenumi{\arabic{enumi})}
\setcounter{enumi}{1}
\tightlist
\item
  \(S_2(n)\): take \(x=1,\;y=0\)
\end{enumerate}

\[
S_2(n)=\sum_{k=0}^{n}\binom{n}{k}k!\,(n+1)^{\,n-k}
=(n+1)^{\,n}\sum_{k=0}^{n}\frac{n!}{(n-k)!\,(n+1)^{k}}.
\]

But

\[
\sum_{k=0}^{n}\frac{n!}{(n-k)!\,(n+1)^{k}}
=\sum_{k=1}^{n+1}\frac{(n+1)!}{((n+1)-k)!\,(n+1)^{k}}
=Q(n+1),
\]

since
\(\sum_{k=1}^{n+1}\frac{(n+1)!}{(n+1-k)!}\,t^{k}=(n+1)t\sum_{k=0}^{n}\frac{n!}{(n-k)!}\,t^{k}\)
(set \(t=\tfrac1{n+1}\)). Hence

\[
S_2(n)=(n+1)^{\,n}Q(n+1).
\]

Final results

\[
\boxed{\displaystyle \sum_{k=0}^{n}\binom{n}{k}\,k^{k}(n-k)^{\,n-k}
= n^{\,n}\bigl(1+Q(n)\bigr)}
\]

\[
\boxed{\displaystyle \sum_{k=0}^{n}\binom{n}{k}\,(k+1)^{k}(n-k)^{\,n-k}
= (n+1)^{\,n}Q(n+1)}
\]

\emph{Quick check (e.g., \(n=4\)):} Left sum \(=\) 824; right side
\(=4^{4}(1+Q(4))=256\cdot 3.21875=824\). Left sum \(=\) 1569; right side
\(=5^{4}Q(5)=625\cdot 2.5104=1569\).

\subsection{19 {[}HM30{]}}\label{hm30-7}

Split the integral at \(r\):

\[
C_n=\int_{0}^{r} e^{-nx} f(x)\,dx + \int_{r}^{\infty} e^{-nx} f(x)\,dx
=: I_1(n)+I_2(n).
\]

\begin{enumerate}
\def\labelenumi{\arabic{enumi})}
\tightlist
\item
  The small-\(x\) part \(I_1(n)\)
\end{enumerate}

By the hypothesis, \(|f(x)|\le K x^{\alpha}\) on \([0,r]\) for some
\(K>0\). Hence

\[
|I_1(n)|
\le K \int_{0}^{r} e^{-nx} x^{\alpha}\,dx.
\]

Make the change of variables \(u=nx\) (so \(x=u/n\), \(dx=du/n\)):

\[
|I_1(n)|
\le K \int_{0}^{nr} e^{-u} \left(\frac{u}{n}\right)^{\alpha}\frac{du}{n}
= K\, n^{-1-\alpha}\!\int_{0}^{nr} e^{-u} u^{\alpha}\,du.
\]

Because \(\alpha>-1\), the truncated gamma integral is bounded by
\(\Gamma(\alpha+1)\):

\[
\int_{0}^{nr} e^{-u} u^{\alpha}\,du \le \Gamma(\alpha+1).
\]

Therefore

\[
|I_1(n)| \le K\,\Gamma(\alpha+1)\, n^{-1-\alpha}.
\]

\begin{enumerate}
\def\labelenumi{\arabic{enumi})}
\setcounter{enumi}{1}
\tightlist
\item
  The tail \(I_2(n)\)
\end{enumerate}

For \(x\ge r\), \(e^{-nx}\le e^{-nr}\). Thus

\[
|I_2(n)| \le e^{-nr}\int_{r}^{\infty} |f(x)|\,dx,
\]

and the right-hand side is \(O(e^{-nr})\), which is \(o(n^{-1-\alpha})\)
as \(n\to\infty\).

\begin{enumerate}
\def\labelenumi{\arabic{enumi})}
\setcounter{enumi}{2}
\tightlist
\item
  Combine
\end{enumerate}

Hence

\[
C_n = I_1(n)+I_2(n) = O\!\left(n^{-1-\alpha}\right) + O\!\left(e^{-nr}\right)
= O\!\left(n^{-1-\alpha}\right).
\]

This proves the claim.

Remarks

\begin{itemize}
\item
  The condition \(\alpha>-1\) is exactly what ensures integrability of
  \(x^{\alpha}\) at 0 and is therefore optimal for this type of bound.
\item
  A stronger (classic) version of Watson's lemma says that if \(f(x)\)
  has a full expansion \(f(x)\sim \sum_{m\ge0} a_m x^{\alpha+m}\) as
  \(x\downarrow0\), then

  \[
  \int_{0}^{\infty} e^{-nx} f(x)\,dx \sim \sum_{m\ge0} a_m \,\Gamma(\alpha+m+1)\, n^{-(\alpha+m+1)}.
  \]

  Our exercise is the first-term big-O consequence of this.
\end{itemize}

\subsection{20 {[}HM30{]}}\label{hm30-8}

\begin{enumerate}
\def\labelenumi{\arabic{enumi})}
\tightlist
\item
  Start from the integral form of \(Q(n)\)
\end{enumerate}

From exercise 15 one gets

\[
Q(n)+1 \;=\; \int_{0}^{\infty}\!\Bigl(1+\tfrac{z}{n}\Bigr)^{n}e^{-z}\,dz
\]

(by expanding \((1+z/n)^n\), integrating term--by--term, and using
\(\int_{0}^{\infty} z^k e^{-z}\,dz=k!\)). 

Put \(z = n\,u\). Then

\[
Q(n)+1 \;=\; n\!\int_{0}^{\infty} \exp\!\bigl(-n\,[u-\ln(1+u)]\bigr)\,du
\;=\; n\!\int_{0}^{\infty} e^{-n\,\psi(u)}\,du,
\quad
\psi(u):=u-\ln(1+u).
\]

Near \(u=0\),

\[
\psi(u)=\frac{u^{2}}{2}-\frac{u^{3}}{3}+\frac{u^{4}}{4}-\cdots .
\]

\begin{enumerate}
\def\labelenumi{\arabic{enumi})}
\setcounter{enumi}{1}
\tightlist
\item
  ``Quadraticize'' the exponent
\end{enumerate}

Define \(w\) by

\[
\frac{w^{2}}{2}=\psi(u)\qquad\Longleftrightarrow\qquad
w=\sqrt{\,u^{2}-\frac{2}{3}u^{3}+\frac{2}{4}u^{4}-\frac{2}{5}u^{5}+\cdots\,}.
\]

Then \(w\sim u>0\) for small \(u\), and the change of variables
\(u\mapsto w\) is monotone near \(0\).

Let \(u\) be written as a power series in \(w\):

\[
u
= c_{1}w+c_{2}w^{2}+c_{3}w^{3}+c_{4}w^{4}+\cdots
= w+\frac{1}{3}w^{2}+\frac{1}{36}w^{3}-\frac{1}{270}w^{4}+\cdots,
\]

which is exactly the series referenced in the statement (``as in
(12)''). (You can get these \(c_k\) by expanding the square root and
composing.) 

Differentiate to relate differentials:

\[
du = \Bigl(\sum_{k\ge1} k\,c_k\,w^{\,k-1}\Bigr)\,dw.
\]

Under this substitution,

\[
Q(n)+1
= n\!\int_{0}^{\infty} e^{-n\,\psi(u)}\,du
= n\!\int_{0}^{\infty} e^{-\frac{n}{2}w^{2}}\,
\Bigl(\sum_{k\ge1} k\,c_k\,w^{\,k-1}\Bigr)\,dw .
\]

\begin{enumerate}
\def\labelenumi{\arabic{enumi})}
\setcounter{enumi}{2}
\tightlist
\item
  Apply Watson's lemma (Gaussian Laplace integrals)
\end{enumerate}

Watson's lemma lets us expand a Laplace--type integral by inserting the
Taylor series of the prefactor and integrating term by term. Here each
term is

\[
I_{k}(n):= n\int_{0}^{\infty} e^{-\frac{n}{2}w^{2}}\, w^{\,k-1}\,dw
= \Gamma\!\Bigl(\frac{k}{2}\Bigr)\Bigl(\frac{n}{2}\Bigr)^{1-\frac{k}{2}},
\]

obtained by the substitution \(y=\tfrac{n}{2}w^2\).

Therefore, truncating after \(m-1\) terms,

\[
\boxed{\;
Q(n)+1
= \sum_{k=1}^{m-1} k\,c_k\,\Gamma\!\Bigl(\frac{k}{2}\Bigr)
\Bigl(\frac{n}{2}\Bigr)^{1-\frac{k}{2}}
\;+\; O\!\bigl(n^{\,1-\frac{m}{2}}\bigr)\,,
\;}
\]

which is exactly what the exercise asks you to show. (This is ``the
simplest possible'' approach alluded to in the text for \(Q(n)\).) 

\begin{enumerate}
\def\labelenumi{\arabic{enumi})}
\setcounter{enumi}{3}
\tightlist
\item
  First few terms (consistency check)
\end{enumerate}

Using \(c_1=1,\;c_2=\tfrac13,\;c_3=\tfrac1{36},\;c_4=-\tfrac1{270}\),

\[
\begin{aligned}
Q(n)+1
&= \underbrace{\Gamma(\tfrac12)\Bigl(\tfrac{n}{2}\Bigr)^{1/2}}_{\displaystyle \sqrt{\frac{\pi n}{2}}}
+ \underbrace{\tfrac{2}{3}\,\Gamma(1)\Bigl(\tfrac{n}{2}\Bigr)^{0}}_{\displaystyle \frac{2}{3}}
+ \underbrace{\tfrac{1}{12}\,\Gamma(\tfrac{3}{2})\Bigl(\tfrac{n}{2}\Bigr)^{-1/2}}_{\displaystyle \frac{1}{12}\sqrt{\frac{\pi}{2n}}}
+ \underbrace{\Bigl(-\frac{4}{270}\Bigr)\Gamma(2)\Bigl(\tfrac{n}{2}\Bigr)^{-1}}_{\displaystyle -\frac{4}{135n}}
+ \cdots .
\end{aligned}
\]

Hence

\[
Q(n)=\sqrt{\frac{\pi n}{2}}\;-\;\frac{1}{3}\;+\;\frac{1}{12}\sqrt{\frac{\pi}{2n}}
\;-\;\frac{4}{135n}\;+\;\cdots,
\]

agreeing with the book's asymptotic series for \(Q(n)\).

\bookmarksetup{startatroot}

\chapter{1.3.1 Description of MIX}\label{description-of-mix}

\subsection{1 {[}00{]}}\label{section-99}

Answer: 4 trits.

Reasoning: A MIX byte can encode 64 distinct values (0--63). With base-3
digits (``trits''), we need the smallest \(n\) such that \(3^n \ge 64\).
\(3^3=27<64\) and \(3^4=81\ge64\), so \(n=4\) trits per byte (leaving
\(81-64=17\) codes unused).

\subsection{2 {[}02{]}}\label{section-100}

Answer: 5 bytes.

Why: Each MIX byte holds \(64\) values. With \(n\) adjacent bytes we can
represent up to \(64^n-1\). \(64^4=16{,}777{,}216<99{,}999{,}999\) but
\(64^5=1{,}073{,}741{,}824>99{,}999{,}999\). So we need 5 bytes (i.e.,
one MIX word's magnitude; sign, if any, is separate).

\subsection{3 {[}02{]}}\label{section-101}

\begin{enumerate}
\def\labelenumi{\alph{enumi})}
\tightlist
\item
  Address field: (0:2) (includes the sign and bytes 1--2)
\item
  Index field: (3:3)
\item
  Field field: (4:4)
\item
  Operation code field: (5:5)
\end{enumerate}

(Here 0 denotes the sign; bytes are numbered 1..5 left to right.)

\subsection{4 {[}00{]}}\label{section-102}

Because the address field in a MIX instruction is a signed displacement.
The effective address is computed as

\[
\text{EA} = A + I_4,
\]

where \(A=-2000\). If \(I_4\) is large enough (e.g., \(I_4\ge 2000\)),
the sum \(\text{EA}\) is a valid memory location
\(0\le \text{EA} < 4000\). Thus \texttt{LDA\ -2000,4} means ``load A
from the word at address \(I_4-2000\)''; the program isn't using a
negative memory address---just a negative offset.

\subsection{5 {[}10{]}}\label{section-103}

Here's the full decoding of (6) into MIX's symbolic form.

Given numeric instruction (6) Word layout is
\texttt{±AA\ \textbar{}\ I\ \textbar{}\ F\ \textbar{}\ C}. The figure
shows:

\begin{itemize}
\tightlist
\item
  sign = --
\item
  AA = 80
\item
  I = 3
\item
  F = 5
\item
  C = 4
\end{itemize}

\begin{enumerate}
\def\labelenumi{\arabic{enumi})}
\item
  Opcode → mnemonic In MIX, opcode C = 4 is DIV.
\item
  Field specification Field byte \(F\) encodes \((L:R)\) by
  \(F = 8L + R\). With \(F = 5\) we get
  \(L = 0,\ R = 5 \Rightarrow (0:5)\) (the whole word).
\item
  Address and index The address field is the signed two-byte value ±AA =
  −80. Index byte I = 3 means ``use index register X3''.
\item
  Symbolic instruction (analogous to OP Z,I(F))
\end{enumerate}

\begin{verbatim}
DIV  -80,3(0:5)
\end{verbatim}

Since DIV's default field is (0:5), this is commonly written simply as:

\begin{verbatim}
DIV  -80,3
\end{verbatim}

(Effective address at run time is \(EA = X_3 + (-80)\).)

\subsection{6 {[}10{]}}\label{section-104}

Given memory location 3000 holds the word

\begin{verbatim}
M[3000] = + | b1=5 | b2=1 | (b3:b4)=200 | b5=15
\end{verbatim}

(read as sign plus; byte1=5, byte2=1, the two-byte field (3:4)=200,
byte5=15).

\begin{itemize}
\tightlist
\item
  Load puts the selected field right-justified in the destination;
  unused left bytes → \texttt{0}.
\item
  If the field includes the sign \texttt{(L=0)}, the loaded sign comes
  from memory; otherwise the sign is \texttt{+}.
\item
  The \texttt{N} variant flips the sign after loading.
\item
  Index registers \texttt{I1..I6} hold sign + 2 bytes.
\end{itemize}

Results:

\begin{enumerate}
\def\labelenumi{\alph{enumi})}
\item
  LDAN 3000 Loads the whole word, then negates the sign → rA = −
  \textbar{} 5 \textbar{} 1 \textbar{} 200 \textbar{} 15 (fully
  defined).
\item
  LD2N 3000(3:4) Loads bytes 3--4 (no sign) into \texttt{I2}, then
  negates → rI2 = −200.
\item
  LDX 3000(1:3) Loads bytes 1--3 (no sign) into the right end of
  \texttt{X}; left two bytes become 0 → rX = + \textbar{} 0 \textbar{} 0
  \textbar{} 5 \textbar{} 1 \textbar{} ? (partially defined because byte
  3 is unknown individually---only (3:4)=200 was given).
\item
  LD6 3000 (default \texttt{(0:2)}) Loads sign and bytes 1--2 into
  \texttt{I6} → two-byte value \(64·5+1=321\), sign \texttt{+} → rI6 =
  +321.
\item
  LDXN 3000(0:0) Loads just the sign into \texttt{X} (magnitude becomes
  0), then negates → rX = − \textbar{} 0 \textbar{} 0 \textbar{} 0
  \textbar{} 0 \textbar{} 0 (negative zero).
\end{enumerate}

\subsection{7 {[}M15{]}}\label{m15-13}

Let a MIX word base be \(W=64^5\).

Registers and operand (as signed integers):

\begin{itemize}
\tightlist
\item
  \(A = s_A\cdot |A|\), \(X = s_X\cdot |X|\) with \(0\le |A|,|X|<W\).
\item
  Memory operand \(M = s_M\cdot |M|\) with \(0\le |M|<W\).
\end{itemize}

Form the dividend from \(A\) and \(X\) by concatenation (sign from \(A\)
only):

\[
D \;=\; s_A\cdot\big(|A|\cdot W + |X|\big).
\]

Let the divisor be \(M\).

\texttt{DIV\ M} is defined without overflow iff

\[
|M|\neq 0 \quad\text{and}\quad \Big\lfloor\frac{|A|\cdot W + |X|}{|M|}\Big\rfloor < W .
\]

When defined, compute

\[
Q \;=\; \Big\lfloor\frac{|A|\cdot W + |X|}{|M|}\Big\rfloor,\qquad
R \;=\; (|A|\cdot W + |X|)\bmod |M| ,
\]

so \(0\le R<|M|\).

The results placed in the registers are

\[
\boxed{\;A \leftarrow (s_A s_M)\cdot Q\;},\qquad
\boxed{\;X \leftarrow s_A \cdot R\;}
\]

(i.e., the quotient gets the product of the signs of \(A\) and \(M\);
the remainder gets the sign of \(A\)).

This covers all non-overflow cases; overflow occurs exactly when
\(|M|=0\) or \(Q\ge W\).

\subsection{8 {[}15{]}}\label{section-105}

Context (from the book's last \texttt{DIV} example): Before
\texttt{DIV\ 1000}, \(rA=-0\), \(rX=+\,1235\,0\,3\,1\), and
\(M[1000]=-\,0\,0\,0\,2\,0\) (i.e., divisor \(|M|=200\) on the decimal
MIX). In that example the result was \(Q=6{,}175{,}001\), \(R=101\).
(\href{https://stackoverflow.com/questions/768948/how-does-division-work-in-mix}{Stack
Overflow})

Now change only ``rX before'' to \(rX=-\,1234\,0\,3\,1\).

Key facts:

\begin{itemize}
\tightlist
\item
  \texttt{DIV} forms the dividend from \((rA,rX)\) using the sign of A
  only; the sign of \(rX\) is ignored.
\item
  Here \(|A|=0\), so the dividend magnitude is just \(|X|\).
\item
  Replacing \(+\,1235\,0\,3\,1\) by \$ -,1234,0,3,1\$ reduces the
  magnitude by \(100^3=1{,}000{,}000\). Since
  \(1{,}000{,}000 / 200 = 5{,}000\), the quotient decreases by 5,000 and
  the remainder is unchanged.
\end{itemize}

Therefore

\[
Q' = 6{,}175{,}001 - 5{,}000 = 6{,}170{,}001,\qquad R' = 101.
\]

Signs:

\begin{itemize}
\tightlist
\item
  Quotient gets
  \(\operatorname{sign}(A)\times\operatorname{sign}(M)=(-)\times(-)=+\).
\item
  Remainder gets \(\operatorname{sign}(A)=-\).
\end{itemize}

Answer (byte form, decimal MIX):

\begin{itemize}
\tightlist
\item
  \(rA \leftarrow + \;|\; 0 \;|\; 6 \;|\; 17 \;|\; 00 \;|\; 01\) (i.e.,
  \texttt{+\ \textbar{}\ 0\ \textbar{}\ 617\ \textbar{}\ 00\ \textbar{}\ 01})
\item
  \(rX \leftarrow - \;|\; 0 \;|\; 0 \;|\; 0 \;|\; 01 \;|\; 01\) (i.e.,
  \texttt{-\ \textbar{}\ 0\ \textbar{}\ 0\ \textbar{}\ 0\ \textbar{}\ 01\ \textbar{}\ 01})
\end{itemize}

So with \texttt{rX\ before\ =\ −\ 1234\ 0\ 3\ 1}, you get A = + 0 617 00
01 and X = − 0 0 0 01 01.

\subsection{9 {[}15{]}}\label{section-106}

\subsubsection{Operators that can turn OV
on}\label{operators-that-can-turn-ov-on}

\begin{itemize}
\tightlist
\item
  ADD M --- sets OV if the signed result won't fit in one MIX word
  (\textbar result\textbar{} too large).
\item
  SUB M --- same rule as ADD.
\item
  DIV M --- sets OV if division is impossible (M = 0) or if the quotient
  doesn't fit in a single word (\textbar quotient\textbar{} ≥ 10⁵).
\item
  INCA/DECA, INCX/DECX --- these are ``add/subtract an immediate value''
  on rA or rX; they can overflow just like ADD/SUB and therefore set OV.
\item
  NUM --- converts the ten bytes in rA‖rX to a number in rA; sets OV if
  the decimal value is too large to fit in one word (common when the
  10-digit value exceeds the MIX word range).
\end{itemize}

\begin{quote}
Notes • MUL never sets OV in MIX integer arithmetic because its 10-byte
result is stored across rA‖rX. (Only floating-point multiply can affect
OV, which we are excluding.) • Shift, load/store, compare, I/O, and most
address-transfer ops do not affect OV.
\end{quote}

\subsubsection{Operators that can only clear
OV}\label{operators-that-can-only-clear-ov}

\begin{itemize}
\tightlist
\item
  JOV \ldots{} and JNOV \ldots{} --- both perform the jump test and
  reset OV to ``off'' as part of their semantics.
\end{itemize}

That's all. In short: ADD, SUB, DIV, INCA/DECA/INCX/DECX, NUM can set
overflow; JOV/JNOV clear it.

\subsection{10 {[}15{]}}\label{section-107}

\subsubsection{What the comparison indicator
is}\label{what-the-comparison-indicator-is}

MIX keeps a three-way comparison indicator with states \textless, =,
\textgreater. It changes only when an instruction explicitly performs a
comparison; jumps merely \emph{test} it.

\subsubsection{Operators that can change
it}\label{operators-that-can-change-it}

\begin{itemize}
\tightlist
\item
  CMPA --- compare rA with M (using the given field F).
\item
  CMPX --- compare rX with M (field F).
\item
  CMP1, CMP2, CMP3, CMP4, CMP5, CMP6 --- compare rI1\ldots rI6 with the
  selected field of M (default F normally the address field (0:2) for
  index-register compares).
\item
  FCMP --- floating-point compare (if the floating-point option is
  present).
\end{itemize}

\subsubsection{Notes/non-examples}\label{notesnon-examples}

\begin{itemize}
\tightlist
\item
  The JL/JE/JG/JGE/JNE/JLE family of jumps only \emph{tests} the
  indicator; they don't change it.
\item
  Arithmetic, loads/stores, shifts, I/O (IN/OUT/IOC/JBUS/JRED),
  conversion (NUM/CHAR), and register-enter/increment/decrement
  instructions do not alter the comparison indicator.
\end{itemize}

So, the list is: CMPA, CMPX, CMP1--CMP6, and (if FP is implemented)
FCMP.

\subsection{11 {[}15{]}}\label{section-108}

These MIX instructions can change rI1:

\begin{enumerate}
\def\labelenumi{\arabic{enumi}.}
\item
  LD1 M(F) Loads rI1 from memory. The field specification F selects a
  field of M; when loading into an index register, MIX takes the address
  part (two bytes + sign) of the selected field. The result is therefore
  brought into the index range (−4095\ldots+4095, with +0). Sign is the
  sign of the selected field.
\item
  LD1N M(F) Like LD1, but the sign of the loaded quantity is negated.
\item
  INC1 M Adds the effective address EA(M) to rI1. If the algebraic sum
  lies outside the index range, MIX reduces the magnitude modulo 4096 to
  bring it back into −4095\ldots+4095 (zero is +0). (Overflow toggle is
  not used by index arithmetic.)
\item
  DEC1 M Subtracts EA(M) from rI1, with the same reduction-to-range rule
  as INC1.
\item
  ENT1 M ``Enter'' EA(M) in rI1 (i.e., rI1 ← EA(M)). Field F is ignored.
  Sign is the sign of EA(M).
\item
  ENN1 M Like ENT1, but rI1 ← −EA(M).
\item
  MOVE M,1(F) Any MOVE whose index field is 1 uses rI1 as the word
  count. During execution, rI1 is decremented by 1 per word moved and on
  normal completion rI1 = 0. If rI1 is initially 0, no words are moved
  and rI1 remains 0. (Other registers and addresses are advanced as
  defined for MOVE, but those changes don't affect this list.)
\end{enumerate}

\subsubsection{Non-examples (often
confused)}\label{non-examples-often-confused}

\begin{itemize}
\tightlist
\item
  ST1 stores rI1 to memory; it does not change rI1.
\item
  CMP1 and the J1* jump family only test rI1; they don't modify it.
\item
  Using rI1 in indexing for effective addresses does not alter rI1.
\end{itemize}

So the list is: MOVE, LD1, LD1N, INC1, DEC1, ENT1, ENN1.

\subsection{12 {[}10{]}}\label{section-109}

Answer: \texttt{INC3\ 0,3}.

\begin{itemize}
\tightlist
\item
  MIX first computes the effective address \(M\) by adding the selected
  index register to the address field: with address \(±AA=0\) and
  I-field \(=3\), we get \(M = 0 + rI3 = rI3\). 
\item
  On ``certain instructions'' the address \(M\) isn't a memory location;
  instead, the instruction adds \(M\) to an index register.
  \texttt{INCi} is that family: \texttt{INC3} adds \(M\) to rI3. So
  \texttt{INC3\ 0,3} performs
  \(rI3 \leftarrow rI3 + M = rI3 + rI3 = 2 \cdot rI3\). 
\end{itemize}

\emph{Notes:} Index registers hold two bytes plus sign; any operation
that would make their magnitude exceed two bytes is undefined (the
simulator explicitly checks size after \texttt{LDi/LD\ iN/INC\ i}). In
normal ranges, the sign is preserved, so doubling keeps the sign of rI3.

\subsection{13 {[}10{]}}\label{section-110}

We will use: - Jump effect on rJ. When a (conditional) jump actually
occurs, rJ is set to the address of the next instruction (here, 1001).
If no jump occurs, rJ is unchanged. - Overflow-clearing. \texttt{JNOV}
clears the overflow toggle after testing; \texttt{JOV} has the symmetric
behavior (and the exercise statement explicitly notes that \texttt{JOV}
clears it).

\begin{enumerate}
\def\labelenumi{\arabic{enumi})}
\tightlist
\item
  Replace by \texttt{JNOV\ 1001}:
\end{enumerate}

\begin{itemize}
\tightlist
\item
  Control flow: still executes next at 1001 in all cases (jump if
  OV=OFF, otherwise fall through).
\item
  Overflow: cleared afterward in either case. 
\item
  rJ: different in general. With the original \texttt{JOV\ 1001}, rJ is
  set to 1001 iff OV was ON; with \texttt{JNOV\ 1001}, rJ is set to 1001
  iff OV was OFF. So most of the time the rJ setting flips. (That's the
  only behavioral difference.) 
\end{itemize}

\begin{enumerate}
\def\labelenumi{\arabic{enumi})}
\setcounter{enumi}{1}
\tightlist
\item
  Replace by \texttt{JOV\ 1000}
\end{enumerate}

\begin{itemize}
\tightlist
\item
  If OV=OFF: falls through to 1001, just like the original.
\item
  If OV=ON: it jumps to 1000, clearing OV; then re-executes with OV=OFF
  and falls through to 1001.
\item
  Net effect: same final state (including rJ = 1001 when OV was
  initially ON), but it takes extra time in the OV=ON case (you execute
  the instruction twice). So ``no difference except time.'' 
\end{itemize}

\begin{enumerate}
\def\labelenumi{\arabic{enumi})}
\setcounter{enumi}{2}
\tightlist
\item
  Replace by \texttt{JNOV\ 1000}
\end{enumerate}

\begin{itemize}
\tightlist
\item
  If OV=OFF: it jumps to 1000, clears OV (still OFF), and repeats
  forever → possible infinite loop.
\item
  If OV=ON: it does not jump, clears OV, and falls through to 1001.
  Hence this change ``may lock the computer in an infinite loop.'' 
\end{itemize}

\subsection{14 {[}20{]}}\label{section-111}

\begin{enumerate}
\def\labelenumi{\arabic{enumi}.}
\item
  Increment/Decrement family INCA, DECA, INCX, DECX, INC1--INC6,
  DEC1--DEC6 Make the effective address zero by choosing AA=0 and I=0.
  Then the instruction adds/subtracts 0 and changes nothing.
\item
  Shifts SLA, SRA, SLAX, SRAX, SLC, SRC Choose F=0 (shift count 0). A, X
  (and A‖X), remain unchanged; indicators and memory are unaffected.
\end{enumerate}

What doesn't work (and why)

Loads: LDA, LDX, LD1--LD6, and the negated forms They overwrite a
register with a memory field; we can't guarantee equality with the
current register.

Stores: STA, STX, ST1--ST6, STJ, STZ They write to memory; with unknown
memory we can't guarantee no change.

Arithmetic on A with a memory operand: ADD, SUB, MUL, DIV ADD/SUB would
be a no-op only if the fetched value were provably 0 for all memory
contents; there is no choice of F, I, AA that guarantees that without
assumptions. MUL and DIV clearly change A‖X or the overflow state unless
blocked by impossible cases (which still change OV).

Compare: CMPA, CMPX, CMP1--CMP6 (and FCMP if present) They always set
the comparison indicator, which is a state change.

Conversion: NUM, CHAR They rewrite A and X.

Block move: MOVE It decrements rI1 during execution (and may copy data),
so state changes unless rI1 were known to be 0---which we are not
allowed to assume.

Jumps: all conditional jumps and JMP If a jump occurs, rJ is set; if it
does not, rJ is not set. With unknown prior rJ and unknown condition
values, there's no way to make both

\subsection{15 {[}10{]}}\label{section-112}

Each MIX word has 5 bytes → up to 5 printed characters per word. For
character I/O devices, a block transfers a fixed number of words:
typewriter and paper tape use 14 words per block; card reader and card
punch use 16; the line printer uses 24.

So the alphanumeric characters per block are:

\begin{itemize}
\tightlist
\item
  Typewriter or paper-tape: 14 words × 5 chars/word = 70 characters.
\item
  Card reader or card punch: 16 × 5 = 80 characters.
\item
  Line printer: 24 × 5 = 120 characters.
\end{itemize}

(``Five characters per word'' follows because character devices map one
MIX character per byte; a word has five bytes.)

\subsection{16 {[}20{]}}\label{section-113}

\begin{enumerate}
\def\labelenumi{(\alph{enumi})}
\tightlist
\item
\end{enumerate}

Use \emph{overlapped} \texttt{MOVE} to ``smear'' one zero across the
block:

\begin{verbatim}
        STZ     0           ; make location 0000 = +0
        ENT1    1           ; rI1 := 0001 (destination start)
        MOVE    0(99)       ; copy 99 words from 0000 → 0001..0099 (one-by-one)
\end{verbatim}

Why this works: On MIX, \texttt{MOVE\ M(F)} moves F words from M to the
address in rI1, one word at a time, and then adds F to rI1. With overlap
(here M=0, rI1=1), the same source word gets copied repeatedly, so the
single zero at 0000 gets replicated through 0001..0099. Afterward
rI1=0100.

Time (MIX units): \texttt{STZ} = 2u, \texttt{ENT1} = 1u, \texttt{MOVE} =
1u + 2·F = 1 + 2·99 = 199u → 202u total.

\begin{enumerate}
\def\labelenumi{(\alph{enumi})}
\setcounter{enumi}{1}
\tightlist
\item
\end{enumerate}

If we optimize only for time (not length), the fastest is to store zeros
directly:

\begin{verbatim}
        STZ 0
        STZ 1
        ...
        STZ 99
\end{verbatim}

Each \texttt{STZ} costs 2u, so 200u total---slightly faster than the
overlapped-MOVE trick, at the price of 100 instructions. 

Notes

\begin{itemize}
\tightlist
\item
  \texttt{MOVE} semantics (dest = rI1; overlap copies the same source
  word repeatedly) are exactly as stated in §1.3.1, incl.~the example
  where one word is copied into several places.
\item
  MIX timing conventions: \texttt{STZ} (a store) = 2u; \texttt{ENTi} (an
  address transfer) = 1u; \texttt{MOVE} = 1u + 2 per word moved.
\end{itemize}

\subsection{17 {[}26{]}}\label{section-114}

Here's a clean solution for Exercise 17 (MIX §1.3.1). Assume at entry
that rI2 = N (0 ≤ N ≤ 2999). Both programs start at location 3000.

\begin{enumerate}
\def\labelenumi{(\alph{enumi})}
\tightlist
\item
  Shortest program (3-instruction loop)
\end{enumerate}

This clears locations N, N−1, \ldots, 0 with rI2 as a down-counter. It
works for every N, including 0.

\begin{verbatim}
3000  LOOP   STZ   0,2      ; M = rI2 → clear that cell
3001         DEC2  1        ; rI2 ← rI2 − 1
3002         J2NN  LOOP     ; keep going while rI2 ≥ 0
3003         HLT
\end{verbatim}

Why it works. \texttt{STZ\ 0,2} stores zero at effective address
\texttt{0\ +\ rI2}. Each iteration decrements rI2; the loop stops after
clearing location 0 (when rI2 becomes −1).

\begin{enumerate}
\def\labelenumi{(\alph{enumi})}
\setcounter{enumi}{1}
\tightlist
\item
  Fastest program (block moves with one final partial)
\end{enumerate}

Use overlapped \texttt{MOVE} to smear the zero from 0000 across the
range in big blocks. We copy 50 words at a time (≤ 63 so it fits in one
byte), then finish with a last \texttt{MOVE} whose count equals the
remainder rI2 (mod 50). The remainder is written into the F-byte of the
final \texttt{MOVE} by self-modifying that instruction.

\begin{verbatim}
3000         STZ     0          ; make 0000 = +0 (seed)
3001         ENT1    1          ; rI1 = destination pointer (start at 0001)

; repeat: while rI2 ≥ 50 { MOVE 0(50); rI2 -= 50 }
3002  LBL    INC2   -50         ; try rI2 ← rI2 − 50
3003         J2N     TAIL       ; if negative, we didn't have 50 left → finish
3004         MOVE    0(50)      ; copy 50 zeros from 0000..0049 to rI1..rI1+49
3005         JMP     LBL

; finish: remainder 0..49 words
3006  TAIL   INC2    50         ; rI2 ← rI2 + 50  (now remainder R = 0..49)
3007         ST2     TAIL+1(4:4); put R into F byte of the next MOVE
3008         MOVE    0(0)       ; becomes MOVE 0(R) after the ST2 above
3009         HLT
\end{verbatim}

Why it's fast. \texttt{MOVE\ 0(50)} costs \texttt{1\ +\ 2·50\ =\ 101u}
for 50 words (≈2.02u/word). The loop overhead is tiny (a couple of 1u
instructions per 50 words). The tail step performs one last
\texttt{MOVE} of R ≤ 49 words; we set \texttt{F\ =\ R} by
\texttt{ST2\ …\ (4:4)}, which writes the low byte of rI2 (0..63) into
the F-byte (perfect for a remainder ≤ 49). No memory beyond 0..N is
touched.

Edge cases. N = 0: the block loop is skipped; \texttt{INC2\ 50} sets
remainder R = 0; we self-modify the final \texttt{MOVE} to
\texttt{MOVE\ 0(0)}, which moves nothing---only location 0000 (seeded by
\texttt{STZ\ 0}) is zero, which is exactly the requirement. N between 1
and 49: zero or one 50-block, then the tail \texttt{MOVE} with R = N. N
a multiple of 50: the loop finishes with rI2 = −50; at TAIL we add 50 so
R = 0; the final \texttt{MOVE} moves 0 words.

Here is the final answer:

\begin{itemize}
\tightlist
\item
  Shortest: 3-instruction \texttt{STZ/DEC2/J2NN} loop.
\item
  Fastest: seed 0000 with zero; smear with repeated \texttt{MOVE\ 0(50)}
  blocks; finish with one self-modified \texttt{MOVE\ 0(R)} where R = N
  mod 50.
\end{itemize}

\subsection{18 {[}22{]}}\label{section-115}

Here's a careful, walkthrough of Exercise 18 with an annotated listing.
I keep track of the parts the problem explicitly cares about: registers
(especially rI1, rX), the overflow toggle (OV), the comparison indicator
(CMP), and the memory cells that get touched. I'll write MIX words as
``sign \textbar{} five bytes''.

Preconditions: nothing is assumed about registers or memory. \texttt{b}
denotes the MIX byte base (64 in the standard spec). NUM/CHAR behave the
same for any \texttt{b}; only whether NUM overflows depends on
\texttt{b}.

\begin{Shaded}
\begin{Highlighting}[]
\NormalTok{STZ 1            ; mem[1] ← +0 | 00 00 00 00 00}
\NormalTok{                 ; (STZ always writes +0.)}

\NormalTok{ENNX 1           ; rX ← −1  → bytes 00 00 00 00 01, sign −}

\NormalTok{STX 1(0:1)       ; store (sign,byte1) of rX into mem[1]:}
\NormalTok{                 ; sign(−), byte1=00; bytes2..5 unchanged (still 00)}
\NormalTok{                 ; ⇒ mem[1] becomes −0 | 00 00 00 00 00}

\NormalTok{SLAX 1           ; shift left A‖X by 1 byte (signs unchanged).}
\NormalTok{                 ; X was −1 → becomes −b  → bytes 00 00 00 01 00}

\NormalTok{ENNA 1           ; rA ← −1  → bytes 00 00 00 00 01, sign −}

\NormalTok{INCX 1           ; rX ← (−b)+1 = −(b−1)  → bytes 00 00 00 00 (b−1)}

\NormalTok{ENT1 1           ; rI1 ← 1}

\NormalTok{SRC 1            ; rotate right A‖X by 1 byte (signs unchanged):}
\NormalTok{                 ;   rA bytes: [b−1, 0, 0, 0, 0], sign −}
\NormalTok{                 ;   rX bytes: [1,   0, 0, 0, 0], sign −}

\NormalTok{ADD 1            ; add mem[1] (which is −0) to rA → rA unchanged.}
\NormalTok{                 ; OV unchanged (adding 0 can’t overflow).}

\NormalTok{DEC1 −1          ; rI1 ← 1 − (−1) = 2}

\NormalTok{STZ 1            ; mem[1] ← +0 (converts the earlier “−0” to +0)}

\NormalTok{CMPA 1           ; compare rA (negative, nonzero) with +0  → CMP = “\textless{}”}

\NormalTok{MOVE −1,1(1)     ; source EA = rI1 + (−1) = 2−1 = 1}
\NormalTok{                 ; move 1 word: mem[2] ← mem[1] (= +0)}
\NormalTok{                 ; rI1 ← rI1 + 1 = 3}

\NormalTok{NUM 1            ; interpret the 10 bytes of A‖X as decimal digits (mod 10):}
\NormalTok{                 ; digits = [(b−1) mod 10, 0,0,0,0, 1,0,0,0,0]}
\NormalTok{                 ; rA ← signed decimal value = −(D*10\^{}9 + 10\^{}4),}
\NormalTok{                 ;   where D = (b−1) mod 10.  (No other state changes.)}
\NormalTok{                 ; OV is set iff |rA| exceeds the MIX integer range:}
\NormalTok{                 ;   if b=64  → |rA| = 3,000,010,000  \textgreater{} 1,073,741,823 ⇒ OV ON}
\NormalTok{                 ;   if b=100 → |rA| = 9,000,010,000 ≤ 9,999,999,999 ⇒ OV OFF}

\NormalTok{CHAR 1           ; write the 10 decimal digits of |rA| back into A‖X:}
\NormalTok{                 ; rA bytes become the characters “30000” (for b=64 case),}
\NormalTok{                 ; rX bytes become the characters “10000”.}
\NormalTok{                 ; (Signs are not used for the digits; rA keeps its sign.)}
\NormalTok{                 ; No effect on OV or CMP.}

\NormalTok{HLT 1            ; halt (no state changes beyond stopping)}
\end{Highlighting}
\end{Shaded}

Here's the program for Exercise 18 rewritten as MIXAL with inline
comments in the MIXAL ``comment field''.

\begin{verbatim}
        ORG     3000

START   STZ     1               * mem[1] ← +0

        ENNX    1               * X ← −1
        STX     1,0(0:1)        * store (sign,byte1) of X → mem[1]; makes mem[1] = −0

        SLAX    1               * shift left A‖X by 1 byte; X becomes −b
        ENNA    1               * A ← −1
        INCX    1               * X ← −b + 1 = −(b−1)
        ENT1    1               * I1 ← 1

        SRC     1               * rotate right A‖X by 1 byte
                                * (A stays negative; X now has bytes “1 0 0 0 0” mod b)

        ADD     1               * A ← A + mem[1] (−0) ⇒ no change, OV unchanged
        DEC1    -1              * I1 ← 1 − (−1) = 2
        STZ     1               * mem[1] ← +0  (clears previous “−0”)

        CMPA    1               * compare A with +0 ⇒ CMP = “<”

        MOVE    -1,1(1)         * source EA = I1−1 = 1; move 1 word:
                                *   mem[2] ← mem[1] (= +0)
                                *   I1 ← I1 + 1 = 3

        NUM     1               * interpret 10 bytes of A‖X as decimal digits:
                                *   A ← −( ((b−1) mod 10)·10^9 + 10^4 )
                                *   OV set iff |A| exceeds MIX word range:
                                *     b=64  → OV ON;  b=100 → OV OFF

        CHAR    1               * write decimal digits of |A| back into A‖X;
                                *   right half X = characters “1 0 0 0 0”

        HLT     1

        END     START
\end{verbatim}

\begin{verbatim}
* Final state (independent of initial contents):
*   I1 = 3
*   mem[1] = +0,  mem[2] = +0
*   CMP  = “<” (less)
*   OV   = ON if b=64, OFF if b=100
*   X    = five character bytes “1 0 0 0 0”
\end{verbatim}

What's true at the end (independent of the initial state)

\begin{itemize}
\tightlist
\item
  rI1 = 3 (set by \texttt{ENT1\ 1}, then advanced by
  \texttt{MOVE\ …(1)}).
\item
  rX contains the five character bytes ``1 0 0 0 0'' (the right half
  produced by \texttt{CHAR}). The exercise only asks ``what is rX?'';
  its five bytes are those digits.
\item
  CMP (comparison indicator) = ``\textless{}'' (set by \texttt{CMPA\ 1}
  and never changed afterward).
\item
  OV (overflow toggle): • ON in the standard \texttt{b=64} MIX, because
  \texttt{NUM} produced 3,000,010,000 which exceeds the word range. •
  OFF if \texttt{b=100}, because \texttt{NUM} produced 9,000,010,000
  which fits.
\item
  Memory touched: • \texttt{mem{[}1{]}} ends as +0 (it briefly became
  ``−0'' after \texttt{STX\ 1(0:1)} but is cleared by the later
  \texttt{STZ\ 1}). • \texttt{mem{[}2{]}} becomes +0 (copied by the
  \texttt{MOVE}).
\end{itemize}

Everything else (e.g., rJ) is unchanged by this code. here's a clean
walk-through of Exercises 18 and 19.

Program (each has M=1 unless shown): STZ 1; ENNX 1; STX 1(0:1); SLAX 1;
ENNA 1; INCX 1; ENT1 1; SRC 1; ADD 1; DEC1 −1; STZ 1; CMPA 1; MOVE
−1,1(1); NUM 1; CHAR 1; HLT 1

Key facts used

\begin{itemize}
\tightlist
\item
  STZ stores zero to the selected field; STX stores the selected field
  of rX; field (0:1) = sign+byte1.
\item
  ENN? / ENT? / INC? set or add an immediate M to a register (A, X, or
  Ii).
\item
  SLAX/SRC shift/rotate across the 10 bytes rA\textbar rX; signs are not
  affected.
\item
  ADD adds a memory value to rA and only sets overflow if the arithmetic
  result exceeds range. ({[}GNU{]}{[}4{]})
\item
  NUM converts the 10 bytes of rA\textbar rX (digits taken mod 10) to a
  decimal number in rA; CHAR writes the 10 decimal digits of rA back
  into rA\textbar rX.
\item
  MIX word range is ±1,073,741,823 in binary-byte mode (b=64) and
  ±9,999,999,999 in decimal-byte mode (b=100).
\end{itemize}

Step-by-step (what becomes independent of the initial contents)

\begin{enumerate}
\def\labelenumi{\arabic{enumi}.}
\tightlist
\item
  STZ 1 → memory{[}1{]} := +0.
\item
  ENNX 1 → rX := −1. ({[}GNU{]}{[}2{]})
\item
  STX 1(0:1) → store sign+byte1 of rX into memory{[}1{]}. Since rX=−1
  has sign `−' and byte1=0, memory{[}1{]} becomes ``−0 00 00 00 00''
  (i.e., −0).
\item
  SLAX 1 → rA\textbar rX shifts left by one byte; rX becomes −b (its
  bytes are 00 00 00 01 00). ({[}GNU{]}{[}3{]})
\item
  ENNA 1 → rA := −1. ({[}GNU{]}{[}2{]})
\item
  INCX 1 → rX := (−b)+1 = −(b−1). (So rX's bytes are 00 00 00 00 (b−1).)
  ({[}GNU{]}{[}2{]})
\item
  ENT1 1 → rI1 := 1. ({[}GNU{]}{[}2{]})
\item
  SRC 1 → rotate right one byte across rA\textbar rX. • rA bytes become
  {[}b−1,0,0,0,0{]} (still negative). • rX bytes become {[}1,0,0,0,0{]}
  (still negative). ({[}GNU{]}{[}3{]})
\item
  ADD 1 → add memory{[}1{]} (−0) to rA: rA unchanged (still negative).
  Overflow not set by adding zero.
\item
  DEC1 −1 → rI1 := 1 − (−1) = 2.
\item
  STZ 1 → memory{[}1{]} := +0 (explicitly clears the earlier ``−0'').
\item
  CMPA 1 → compare rA (negative) with +0; comparison indicator =
  ``less''.
\item
  MOVE −1,1(1) → source EA = rI1 + (−1) = 2−1 = 1; move one word from
  memory{[}1{]} to memory{[}2{]}; rI1 increments by 1 → rI1=3. (So
  memory{[}2{]} becomes +0.)
\item
  NUM 1 → treat the 10 bytes in rA\textbar rX as decimal digits modulo
  10: digits are {[}(b−1) mod 10, 0,0,0,0, 1,0,0,0,0{]}, so rA becomes
  the decimal number 3000010000 if b=64, or 9000010000 if b=100, with
  rA's sign negative.) • In binary-byte mode (b=64), this exceeds
  ±1,073,741,823 ⇒ overflow set ON. • In decimal-byte mode (b=100), it
  fits in ±9,999,999,999 ⇒ overflow OFF.
\item
  CHAR 1 → write the 10 decimal digits of rA back into rA\textbar rX; rX
  now holds the character bytes ``10000'' (i.e., 31,30,30,30,30). (The
  sign bits are not used by the digits; rA keeps its sign.)
\item
  HLT 1 → halt.
\end{enumerate}

What you can state at the end

\begin{itemize}
\tightlist
\item
  rI1 = 3.
\item
  memory changes: memory{[}1{]} ends as +0; memory{[}2{]} is +0 (copied
  by MOVE).
\item
  comparison indicator after CMPA 1 = L (A \textless{} memory{[}1{]}).
\item
  overflow: • binary-byte mode (b=64): ON (set by NUM); • decimal-byte
  mode (b=100): OFF.
\item
  rX holds the five character digits ``10000'' (bytes 31,30,30,30,30)
  after CHAR (independent of byte size).
\end{itemize}

\subsection{19 {[}14{]}}\label{section-116}

Using the standard MIX timings (in units ``u'') --- loads/stores 2u;
address-transfer (ENT/ENN/INC/DEC) 1u; shifts 2u; ADD 2u; CMP 2u; MOVE
costs 1+2·F; NUM 10u; CHAR 10u --- the total is: 2 +1 +2 +2 +1 +1 +1 +2
+2 +1 +2 +2 +3 +10 +10 = 42u.

\subsection{20 {[}20{]}}\label{section-117}

Here is a single-loop program that leaves every cell as \texttt{HLT} and
then halts:

\begin{Shaded}
\begin{Highlighting}[]
\NormalTok{        ORIG   3994          * park the loop near the top}
\NormalTok{START   LDA    H             * A = instruction word for HLT}
\NormalTok{        ENT1   0             * I1 = 0}
\NormalTok{LOOP    STA    3999,1        * store to 3999+I1 (wraps: 3999,0,1,...,3998)}
\NormalTok{        INC1   1             * I1++}
\NormalTok{        J1P    LOOP          * loop while positive}
\NormalTok{H       HLT    0             * when J1P is finally overwritten, we halt}
\NormalTok{        END    START}
\end{Highlighting}
\end{Shaded}

Why it's safe and complete. \texttt{STA\ 3999,1} walks the addresses in
this exact order: 3999, 0, 1, 2, \ldots, 3998 (because MIX effective
addresses wrap modulo 4000). The loop body itself sits at 3996--3998;
those three words are the last ones touched: \texttt{STA} (3996) and
\texttt{INC1} (3997) get overwritten near the end, but only after
they've been executed in that iteration; the next fetch returns to
\texttt{J1P} (3998) until the very last pass. The final store hits 3998
and turns \texttt{J1P} into \texttt{HLT}; the next fetch executes that
\texttt{HLT}. At that moment every word 0\ldots3999 is \texttt{HLT}.

\subsection{21 {[}24{]}}\label{section-118}

\begin{enumerate}
\def\labelenumi{(\alph{enumi})}
\tightlist
\item
\end{enumerate}

No (not by program action). In MIX, the only way a program sets
\texttt{rJ} is via a taken jump that stores a link (all \texttt{J..}
jumps except \texttt{JSJ}, and \texttt{SLT} when the jump is not taken).
Executing such a jump from location \texttt{m} sets

\begin{verbatim}
rJ ← m + 1
\end{verbatim}

without reducing modulo 4000. Since \texttt{m\ ∈\ \{0,…,3999\}}, we
always get \texttt{rJ\ ∈\ \{1,…,4000\}}---never 0. To make
\texttt{rJ\ =\ 0} you'd need \texttt{m\ =\ −1}, which a program cannot
execute. (Zero can be forced only by external means---cf.~the ``GO
button'', exercise 26.)

Note: when returning from a subroutine we typically use
\texttt{JMP\ 0,J}. If \texttt{rJ\ =\ 4000} (from a call placed at 3999),
the effective address \texttt{0\ +\ rJ\ =\ 4000} is interpreted modulo
4000 as address 0, so control returns correctly---yet the content of
\texttt{J} itself is not 0.

\begin{enumerate}
\def\labelenumi{(\alph{enumi})}
\setcounter{enumi}{1}
\tightlist
\item
\end{enumerate}

Idea: Temporarily place \texttt{JMP\ RESUME} at address \texttt{N−1},
jump there with \texttt{JSJ} (which doesn't touch \texttt{J}), let that
\texttt{JMP} set \texttt{J\ =\ (N−1)+1\ =\ N}, then restore the word we
modified.

\begin{Shaded}
\begin{Highlighting}[]
\NormalTok{        ORIG   3000}
\NormalTok{START   LDX    {-}1,4           * X := M[N{-}1] (save original word)}
\NormalTok{        LDA    INST           * A := word for "JMP RESUME"}
\NormalTok{        STA    {-}1,4           * M[N{-}1] := JMP RESUME (temporary)}
\NormalTok{        JSJ    {-}1,4           * go execute it; JSJ leaves J unchanged}

\NormalTok{RESUME  STX    {-}1,4           * restore original M[N{-}1]}
\NormalTok{        HLT    0              * done}

\NormalTok{INST    JMP    RESUME         * instruction constant we copy to N{-}1}
\NormalTok{        END    START}
\end{Highlighting}
\end{Shaded}

Why it's correct. The \texttt{JSJ} transfers control to \texttt{N−1}.
Executing \texttt{JMP\ RESUME} from \texttt{N−1} sets
\texttt{rJ\ ←\ (N−1)+1\ =\ N} by MIX's link rule. We then restore memory
at \texttt{N−1}, so all memory words end exactly as they began; only
\texttt{J} has changed. The bound \texttt{N\ ≤\ 3000} keeps \texttt{N−1}
below the program that starts at 3000.

\subsection{22 {[}28{]}}\label{section-119}

Here are two MIXAL programs for 22. We assume location 2000 holds an
integer \(X\) and that \(X^{13}\) fits in one word.

Minimal-memory program (tiny loop, 6 words)

\begin{Shaded}
\begin{Highlighting}[]
\NormalTok{        ORIG   3000}
\NormalTok{START   LDA    2000         * A := X}
\NormalTok{        ENT1   12           * 12 more multiplies by X}
\NormalTok{LOOP    MUL    2000         * A·X {-}\textgreater{} AX  (high word zero if fits)}
\NormalTok{        XCH                  * A := product}
\NormalTok{        DEC1   1}
\NormalTok{        J1P    LOOP}
\NormalTok{        HLT    0}
\NormalTok{        END    START}
\end{Highlighting}
\end{Shaded}

Explanation. We start with \(A=X\) and multiply by \(X\) twelve times:
after each \texttt{MUL} the product sits in \texttt{X}, so we
\texttt{XCH} to bring it to \texttt{A} for the next iteration. After 12
iterations, \(A=X^{13}\).

Minimum-time program (5 multiplications = optimal)

This follows a minimal addition chain for 13: \(1,2,3,6,12,13\) i.e.,
\(x^2, x^3, (x^3)^2=x^6, (x^6)^2=x^{12}, x^{12}\cdot x = x^{13}\).

\begin{Shaded}
\begin{Highlighting}[]
\NormalTok{        ORIG   3000}
\NormalTok{        LDA    2000         * A = x}
\NormalTok{        MUL    2000         * x\^{}2}
\NormalTok{        XCH}
\NormalTok{        MUL    2000         * x\^{}3}
\NormalTok{        XCH}
\NormalTok{        STA    2001         * T := x\^{}3   (temp for squaring)}
\NormalTok{        MUL    2001         * x\^{}6}
\NormalTok{        XCH}
\NormalTok{        STA    2001         * T := x\^{}6}
\NormalTok{        MUL    2001         * x\^{}12}
\NormalTok{        XCH}
\NormalTok{        MUL    2000         * x\^{}13}
\NormalTok{        XCH}
\NormalTok{        HLT    0}
\NormalTok{        END    3000}
\end{Highlighting}
\end{Shaded}

Notes. Only 5 \texttt{MUL}s are executed (provably minimal for exponent
13). We store to a single temporary cell (2001) only when squaring
values that aren't in memory. The final result is left in \texttt{A} as
required.

\subsection{23 {[}27{]}}\label{section-120}

\begin{enumerate}
\def\labelenumi{\Alph{enumi})}
\tightlist
\item
  Using partial-field loads (no temporaries; minimal memory under this
  condition)
\end{enumerate}

This builds the result left-to-right by shifting A one byte and adding
the next byte of the source each time. Because each add goes into an
all-zero rightmost byte, no carry propagates.

\begin{Shaded}
\begin{Highlighting}[]
\NormalTok{        ORIG   3000}
\NormalTok{START   LDA    0200(5:5)   * A := [.....e]}
\NormalTok{        SLA    1           * A := [....e0]}
\NormalTok{        ADD    0200(4:4)   * A := [....ed]}
\NormalTok{        SLA    1}
\NormalTok{        ADD    0200(3:3)   * A := [...edc]}
\NormalTok{        SLA    1}
\NormalTok{        ADD    0200(2:2)   * A := [..edcb]}
\NormalTok{        SLA    1}
\NormalTok{        ADD    0200(1:1)   * A := [.edcba]  {-}\textgreater{}  + e d c b a}
\NormalTok{        HLT    0}
\NormalTok{        END    START}
\end{Highlighting}
\end{Shaded}

Memory used: just these 10 instruction words (no data cells). Why
minimal: without a loop you need five byte extractions and four shifts;
any solution with fields must at least assemble five bytes in some
order, so this straight-line sequence is optimal in memory (0 data).

\begin{enumerate}
\def\labelenumi{\Alph{enumi})}
\setcounter{enumi}{1}
\tightlist
\item
  Without using partial-field loads/stores (only whole-word ops)
\end{enumerate}

We treat the word as a base-64 number. Repeatedly:

\begin{enumerate}
\def\labelenumi{\arabic{enumi}.}
\tightlist
\item
  take the low byte as \texttt{source\ \%\ 64} (via \texttt{DIV\ 64},
  remainder goes to \texttt{X}),
\item
  update \texttt{acc\ :=\ acc*64\ +\ remainder},
\item
  set \texttt{source\ :=\ source\ /\ 64} (the quotient from
  \texttt{DIV}, which we saved).
\end{enumerate}

This needs one constant (\texttt{64}) and two temporaries; with only
whole-word ops that's the tightest practical layout.

\begin{Shaded}
\begin{Highlighting}[]
\NormalTok{        ORIG   3100}
\NormalTok{START   LDX    0200        * X := source = + a b c d e}
\NormalTok{        ENTA   0           * A := acc = 0}
\NormalTok{        ENT1   5           * 5 bytes to process}

\NormalTok{LOOP    STA    ACC         * save acc}
\NormalTok{        ENTA   0           * make dividend (A,X) = (0, source)}
\NormalTok{        XCH                 * A := source, X := 0}
\NormalTok{        DIV    K64         * A := source/64 (quotient), X := source\%64 (a byte)}
\NormalTok{        STA    QUO         * remember quotient for next round}
\NormalTok{        XCH                 * A := remainder, X := quotient}
\NormalTok{        STA    REM         * stash remainder}
\NormalTok{        LDA    ACC         * rebuild acc := acc*64 + remainder}
\NormalTok{        SLA    1           * *64}
\NormalTok{        ADD    REM}
\NormalTok{        LDX    QUO         * next source := quotient}
\NormalTok{        DEC1   1}
\NormalTok{        J1P    LOOP}

\NormalTok{        * on exit, A holds the reflected word}
\NormalTok{        HLT    0}

\NormalTok{ACC     CON    0}
\NormalTok{REM     CON    0}
\NormalTok{QUO     CON    0}
\NormalTok{K64     CON    64}
\NormalTok{        END    START}
\end{Highlighting}
\end{Shaded}

Memory used: 12 instruction words + 4 data words (ACC, REM, QUO, K64).
Why this fits the constraint: with no partial-field access you need (i)
a way to peel bytes---\texttt{DIV\ 64} requires the \texttt{64}
constant---and (ii) somewhere to hold the running accumulator while
using \texttt{(A,X)} as the dividend; the two temporaries avoid
clobbering the quotient and remainder. This keeps everything whole-word
only and leaves the final answer in \texttt{A}.

\subsection{24 {[}21{]}}\label{section-121}

We start with A = \texttt{+\ 0\ a\ b\ c\ d} and X =
\texttt{+\ e\ f\ g\ h\ i}, and we want A = \texttt{+\ a\ b\ c\ d\ e} and
X = \texttt{+\ 0\ f\ g\ h\ i}.

\begin{enumerate}
\def\labelenumi{(\alph{enumi})}
\tightlist
\item
\end{enumerate}

Use one constant and one instruction:

\begin{Shaded}
\begin{Highlighting}[]
\NormalTok{        ORIG   3000}
\NormalTok{        DIV    K           * (A:X) ÷ K → A=abcde, X=0fghi}
\NormalTok{        HLT    0}
\NormalTok{K       CON    + 1 0 0 0 0   * K = 64\^{}4}
\NormalTok{        END    3000}
\end{Highlighting}
\end{Shaded}

Why it works. Think of a MIX word as 5 base-64 digits. The 10-digit
number \texttt{0\ a\ b\ c\ d\ e\ f\ g\ h\ i} (that's A:X) divided by
\texttt{64\^{}4} (= \texttt{+\ 1\ 0\ 0\ 0\ 0}) yields

\begin{itemize}
\tightlist
\item
  quotient \texttt{a\ b\ c\ d\ e} (fits in 5 bytes, so it goes to A),
  and
\item
  remainder \texttt{f\ g\ h\ i} (a 4-byte number; when put in a 5-byte
  word it becomes \texttt{0\ f\ g\ h\ i} in X).
\end{itemize}

This uses the fewest locations possible (1 instruction + 1 constant).

\begin{enumerate}
\def\labelenumi{(\alph{enumi})}
\setcounter{enumi}{1}
\tightlist
\item
\end{enumerate}

The division in (a) is slow. The time-optimal solution uses only
shifts/rotates---three instructions total---exactly as in the book:

\begin{Shaded}
\begin{Highlighting}[]
\NormalTok{        ORIG   3000}
\NormalTok{        SRC    4           * circularly rotate the 10 bytes (A:X) right by 4}
\NormalTok{        SRA    1           * shift A right by one byte (zero{-}filling its left)}
\NormalTok{        SLC    5           * circularly rotate the 10 bytes (A:X) left by 5}
\NormalTok{        HLT    0}
\NormalTok{        END    3000}
\end{Highlighting}
\end{Shaded}

Step-by-step (treating A:X as the 10-byte string
\texttt{0\ a\ b\ c\ d\ \textbar{}\ e\ f\ g\ h\ i}):

\begin{enumerate}
\def\labelenumi{\arabic{enumi}.}
\item
  \texttt{SRC\ 4} (right circular on the pair A:X, as in Knuth's MIX):
  \texttt{g\ h\ i\ 0\ a\ \textbar{}\ b\ c\ d\ e\ f}.
\item
  \texttt{SRA\ 1} (right arithmetic on A only, zero-fills the left of
  A): \texttt{0\ g\ h\ i\ 0\ \textbar{}\ b\ c\ d\ e\ f}.
\item
  \texttt{SLC\ 5} (left circular on the pair A:X):
  \texttt{a\ b\ c\ d\ e\ \textbar{}\ 0\ f\ g\ h\ i}.
\end{enumerate}

Result: A = \texttt{+\ a\ b\ c\ d\ e}, X = \texttt{+\ 0\ f\ g\ h\ i},
with no memory traffic and only three fast micro-ops.

\subsection{25 {[}30{]}}\label{section-122}

Here's a detailed, backward-compatible ``Mixmaster'' design that answers
25 and shows how to reuse the I-field profitably while keeping every
correctly written MIX program runnable unchanged.

Goals (constraints)

\begin{enumerate}
\def\labelenumi{\arabic{enumi}.}
\tightlist
\item
  Machine word, byte size, opcodes, and all defined MIX behaviors remain
  valid.
\item
  A MIX program that only uses documented features (e.g., I∈\{0..6\})
  keeps the same semantics.
\item
  New power comes only from previously unused encodings (e.g.,
  I\textgreater6, unused F-codes, unused opcodes) or from optional mode
  bits whose reset value equals MIX behavior.
\end{enumerate}

Better use of the I-field (backward-compatible)

The MIX I-field is 6 bits but only values 0..6 are used (selecting no
index or I1..I6). We keep those intact and assign new meanings to the
unused encodings:

I-field layout (Mixmaster)

Let \texttt{I\ =\ b5\ b4\ b3\ b2\ b1\ b0} (b0 = LSB).

\begin{itemize}
\item
  \texttt{idx\ =\ (b2\ b1\ b0)} selects the base register:

  \begin{itemize}
  \tightlist
  \item
    000 = none (MIX ``no index'')
  \item
    001..110 = I1..I6 (as in MIX)
  \item
    111 = PC (location counter) ⟹ enables PC-relative addressing
  \end{itemize}
\item
  \texttt{IND\ =\ b3} (indirection)

  \begin{itemize}
  \tightlist
  \item
    0 = direct (MIX default)
  \item
    1 = one level of indirection: the effective address is taken from
    the word at the computed address (no chaining unless an op
    explicitly allows it)
  \end{itemize}
\item
  \texttt{SC\ =\ (b5\ b4)} (scale for the index term)

  \begin{itemize}
  \tightlist
  \item
    00 ⇒ ×1 (MIX default)
  \item
    01 ⇒ ×2
  \item
    10 ⇒ ×4
  \item
    11 ⇒ ×8
  \end{itemize}
\end{itemize}

Effective address

\begin{verbatim}
EA = (ADDR + SCALE(idxReg) + BASE) mod MEMSIZE
where   SCALE(idxReg) = (SC ? {1,2,4,8} : 1) * (contents of chosen base reg)
        BASE = 0 in MIX-compat mode; see “large memory” below
if IND=1: EA := contents_of_word_at(EA).ADDR   (one level)
\end{verbatim}

Compatibility: All existing code has \texttt{I≤6}, \texttt{IND=0},
\texttt{SC=00}, so the rule collapses to MIX's \texttt{ADDR\ +\ Ii} with
no change.

Payoffs

\begin{itemize}
\tightlist
\item
  PC-relative addressing (I=7) for position-independent code, short
  branches, and better relocation.
\item
  Scaled indexing removes pre-multiply steps for arrays of records with
  fixed word length (×2, ×4, ×8).
\item
  Single-level indirection gives true pointers without extra loads.
\end{itemize}

Optional: Immediate operands via I-patterns

Reserve two I-patterns for immediates (used only by arithmetic/compare
loads):

\begin{itemize}
\tightlist
\item
  \texttt{I\ =\ 7\ with\ IND=1\ and\ SC=00} ⇒ treat \texttt{ADDR} (and
  F) as a literal (sign-extended) operand: e.g., \texttt{ADDA\ 5,7:1}
  (assembler sugar \texttt{ADDA\ \#5}). No MIX program ever uses such
  I-codes, so compatibility holds.
\end{itemize}

Memory and addressing extensions

\begin{itemize}
\tightlist
\item
  B (base) register (2 bytes, sign) holds a segment number. Default B=0
  on reset ⇒ MIX addresses map unchanged.
\item
  New instructions \texttt{LDB\ M}, \texttt{STB\ M} and privileged
  \texttt{SETB\ imm}. The hardware forms the true physical address as
  \texttt{BASE*4096\ +\ EA}, giving up to 4096 segments of 4096 words
  (or larger if desired). Old programs (B=0) see the original 0..3999
  space.
\item
  Optional per-process limit register to trap out-of-bounds for safety
  (default unlimited for full compatibility).
\end{itemize}

F-field and opcode extensions (safe because many encodings are unused)

\begin{itemize}
\item
  Keep all MIX F-meanings; use unused F codes to define:

  \begin{itemize}
  \tightlist
  \item
    variable-length shifts where the count comes from a register (e.g.,
    \texttt{SLA\ r} with F indicating the source register),
  \item
    logical operations on full words: \texttt{AND}, \texttt{OR},
    \texttt{XOR} (F selects field subranges when desired),
  \item
    block move/fill: \texttt{MOVN}, \texttt{MOVZ} (F selects element
    size as in load/store fields).
  \end{itemize}
\item
  Use unused opcodes for \texttt{PUSH/POP} (stack in memory),
  \texttt{CALL/RET} (JSJ remains; these are conveniences only).
\end{itemize}

Registers and arithmetic

\begin{itemize}
\tightlist
\item
  Add I7..I10 as true index registers, but expose them only through the
  new I-field scheme (idx=7 already maps to PC; the extras are
  accessible via new opcodes \texttt{LDI7}, \texttt{STI7}, \ldots).
  Existing code doesn't see them.
\item
  2-word (AX) add/subtract and multiply with saturating/trap modes
  selectable by F---still leaving the classic MIX MUL/DIV unchanged.
\item
  Fast \texttt{CMPM} (compare A with memory word) and \texttt{ADDM},
  \texttt{SUBM} (A ← A ± M) opcodes to reduce load/store traffic.
\end{itemize}

Jumps and conditions

\begin{itemize}
\tightlist
\item
  New PC-relative jumps (via I=7).
\item
  Indexed conditional jumps: allow condition + index test in one
  instruction (e.g., ``J1Z \& JXNZ'' combined using unused opcodes) to
  shorten loops; classic jumps remain as-is.
\end{itemize}

I/O and devices

\begin{itemize}
\tightlist
\item
  Expand device IDs to 0..63; keep 0..20 exactly as in MIX.
\item
  Asynchronous channel ops with completion flags in a status word; old
  synchronous behavior retained by default.
\end{itemize}

Traps and compatibility mode

\begin{itemize}
\item
  An Extended-Mode bit (reset = 0) controls whether undefined encodings
  trap or use the new features. In reset state, the machine behaves
  exactly like MIX. Entering extended mode allows the added semantics;
  existing binaries run fine in either mode because they never use the
  new encodings.
\item
  We ``make better use of the I-field'' by turning its 6 bits into an
  addressing-mode mini-ISA (index select, scale, indirection, and
  PC-relative)---all using values the original MIX never produced.
\item
  We extend memory and add performance features without changing any
  defined MIX behavior.
\item
  Every suggestion strictly preserves the contract: all correctly
  written MIX programs run unchanged; new power appears only when
  programmers choose the new encodings.
\end{itemize}

\subsection{26 {[}32{]}}\label{section-123}

Alright---this is the classic ``bootstrap loader'' for MIX. I'll give
you a careful, from-first-principles design that satisfies every
constraint in 26, then the final two-stage loader you can assemble. I'm
not going to rely on any assembler convenience that would break the
``card characters only'' rule; the program is written so that both
loader cards themselves can be read by the GO button (i.e., read in
character code) and then executed as instructions. We also avoid opcodes
that the statement warns may be unreadable on some readers (no
\texttt{ENT}, \texttt{INC}, \texttt{CMP} on the cards).

The solution is in three parts:

\begin{enumerate}
\def\labelenumi{\arabic{enumi}.}
\tightlist
\item
  What GO actually does, and the consequences for how a bootstrap must
  be shaped.
\item
  The loading algorithm (how to parse ``information cards'' and the
  ``transfer card'' with only character data available).
\item
  A compact two-card program layout that implements the algorithm with
  only ``safe'' characters, with a readable MIXAL listing you can
  assemble and test. (It fits in 32 words; stage-1: 2 words, stage-2: 30
  words.)
\end{enumerate}

\begin{enumerate}
\def\labelenumi{\arabic{enumi})}
\tightlist
\item
  What the GO button gives us
\end{enumerate}

Pressing GO is essentially \texttt{IN\ 0(16)} on the card reader (unit
16) performed by hardware, followed by:

\begin{itemize}
\tightlist
\item
  the 80 columns are deposited in character code into words
  0000\ldots0015 (5 columns per word), with positive signs,
\item
  when the device goes not-busy, the CPU does \texttt{JMP\ 0} and sets
  \texttt{J\ ←\ 0} (overflow toggled).
\end{itemize}

So the first card we feed must already be executable even though its
five bytes per word are character codes. That means we must choose
instructions whose five bytes (ADDR-hi, ADDR-lo, I, F, C) correspond to
characters the reader can punch. The problem warns that some character
codes (20, 21, 48, 49, 50, \ldots) aren't guaranteed readable; and, in
particular, the encodings for the opcodes \texttt{ENT}, \texttt{INC},
\texttt{CMP} collide with ``bad'' characters. Therefore the bootstrap
must avoid those opcodes on the cards. (We can, of course, emulate them
with allowed instructions.)

We may, however, execute \texttt{IN\ a(16)} with any address we choose,
so our first card can read the second loader card into 0016\ldots0031
(or anywhere that doesn't clobber itself) and then jump there.
Information cards are guaranteed to load data at locations
\textgreater{} 100, so the loader can safely live below 0100.

\begin{enumerate}
\def\labelenumi{\arabic{enumi})}
\setcounter{enumi}{1}
\tightlist
\item
  Loading algorithm (what the second stage does)
\end{enumerate}

Each information card has:

\begin{itemize}
\tightlist
\item
  col 1--5: ignored,
\item
  col 6: \texttt{k} in \{1,\ldots,7\} = how many words to load from this
  card,
\item
  col 7--10: four decimal digits giving the address of the first word on
  this card (\texttt{addr0}), and
\item
  col 11--20, 21--30, \ldots, 71--80: each 10-column field is one word
  in base-10, signed by an overpunch in the last column (the last
  character code is 0..9 for positive, 10..19 for negative, representing
  0..9 with a ``minus'' 11-punch).
\end{itemize}

The transfer card is \texttt{TRANS0nnnn} in columns 1--10 (i.e., word 0
equals ``TRANS0'' in character codes and word 1 holds the four digit
start address \texttt{nnnn}), and the remainder of the card is
irrelevant. When we see it, we jump to \texttt{nnnn}.

The second stage therefore does, in a loop:

\begin{enumerate}
\def\labelenumi{\arabic{enumi}.}
\item
  \texttt{IN\ 0(16)} --- read the next card into 0000\ldots0015.
\item
  If word 0000 equals the characters ``TRANS'' (we compare by
  subtraction and \texttt{JAZ}, not \texttt{CMP}), branch to the
  transfer path; else parse as an information card.
\item
  \texttt{k\ ←\ col6}, \texttt{addr\ ←\ dec(col7…10)}.
\item
  For \texttt{i=0..k-1}:

  \begin{itemize}
  \tightlist
  \item
    Parse one 10-column decimal field at columns
    \texttt{11+10i\ …\ 20+10i}.
  \item
    Convert to a signed MIX integer using repeated
    \texttt{A\ ←\ A*10\ +\ digit}, where
    \texttt{digit\ ←\ byte\ MOD\ 10} and the sign bit is taken from the
    last byte (negative iff last-byte ≥ 10).
  \item
    \texttt{STA\ addr+i}.
  \end{itemize}
\item
  Repeat step 1.
\end{enumerate}

All parsing is done with partial-field loads from the card image at
0000\ldots0015; we never rely on any punch-limited instruction, and we
never need \texttt{ENT/INC/CMP}.

\begin{enumerate}
\def\labelenumi{\arabic{enumi})}
\setcounter{enumi}{2}
\tightlist
\item
  A compact two-card loader
\end{enumerate}

Stage-1: ``copy in stage-2 and jump''

Two words at 0000..0001 after GO:

\begin{Shaded}
\begin{Highlighting}[]
\NormalTok{        ORIG   0000}
\NormalTok{STAGE1  IN     16,16        * read the second loader card into 0016..0031}
\NormalTok{        JMP    0016         * execute it}
\end{Highlighting}
\end{Shaded}

These two instructions can be laid out in characters the readers accept
(opcodes chosen accordingly); their addresses (0016 and 0016) have only
decimal digits, so the two address bytes are harmless.

Stage-2: the real loader (fits in 30 words)

Below is a straight, compact MIXAL program that implements the
algorithm. It uses only the safe instruction set on cards
(\texttt{LDA/STA/LDX/STX/ADD/SUB/MUL/DIV/SLA/SRA/SLC/SRC/Jxx/IN/JMP})
and no \texttt{ENT/INC/CMP}. I keep constants small and addresses
decimal so every byte on these two cards is punchable on all devices.

\begin{Shaded}
\begin{Highlighting}[]
\NormalTok{        ORIG   0016          * second card starts here}
\NormalTok{* {-}{-}{-}{-}{-}{-}{-}{-}{-}{-} constants \& work cells (kept tiny and punch{-}safe) {-}{-}{-}{-}{-}{-}{-}{-}{-}{-}}
\NormalTok{TEN     CON    10}
\NormalTok{ZERO    CON    0}
\NormalTok{ONE     CON    1}
\NormalTok{TWO     CON    2}
\NormalTok{TRNS0   ALF    TRANS        * used to detect the transfer card}
\NormalTok{DIG5M   CON    5            * loop bound for 5 digits}
\NormalTok{ADDR    CON    0            * current target address}
\NormalTok{KCTR    CON    0            * how many words (k) on the info card}
\NormalTok{PTRW    CON    2            * word pointer into card image (starts at word 2: cols 11..15)}
\NormalTok{REM     CON    0            * scratch}
\NormalTok{SIGN    CON    1            * +1 or {-}1 for current number}

\NormalTok{* {-}{-}{-}{-}{-}{-}{-}{-}{-}{-} main loop: read a card and dispatch {-}{-}{-}{-}{-}{-}{-}{-}{-}{-}}
\NormalTok{LOOP    IN     0,16         * read next card image into 0000..0015}
\NormalTok{        LDA    0            * word 0: columns 1..5}
\NormalTok{        SUB    TRNS0}
\NormalTok{        JAZ    XFER         * if "TRANS" {-}\textgreater{} handle transfer}

\NormalTok{* {-}{-}{-}{-} information card header: k and first address {-}{-}{-}{-}}
\NormalTok{        LDA    1(1:1)       * column 6 {-}\textgreater{} k}
\NormalTok{        STA    KCTR}
\NormalTok{        LDA    ZERO         * parse 4{-}digit decimal address columns 7..10}
\NormalTok{        MUL    TEN}
\NormalTok{        ADD    1(2:2)}
\NormalTok{        MUL    TEN}
\NormalTok{        ADD    1(3:3)}
\NormalTok{        MUL    TEN}
\NormalTok{        ADD    1(4:4)}
\NormalTok{        MUL    TEN}
\NormalTok{        ADD    1(5:5)}
\NormalTok{        STA    ADDR}
\NormalTok{        LDA    TWO}
\NormalTok{        STA    PTRW         * start of first 10{-}digit field: word 2}

\NormalTok{* {-}{-}{-}{-} outer loop over k words on the card {-}{-}{-}{-}}
\NormalTok{NEXTW   LDA    ONE}
\NormalTok{        STA    SIGN         * default +1}
\NormalTok{        LDA    ZERO         * accumulator for decimal conversion}

\NormalTok{* {-}{-}{-}{-} consume 10 digits = 5 bytes from word PTRW, then 5 from PTRW+1 {-}{-}{-}{-}}
\NormalTok{        STA    REM          * REM := 0 {-}\textgreater{} use as inner counter}
\NormalTok{D5A     LDX    PTRW         * load word \# (2,4,6,...) of card image}
\NormalTok{        LDA    ZERO}
\NormalTok{        STA    REM}
\NormalTok{D5A1    MUL    TEN}
\NormalTok{        ADD    0, X(1:1)    * add first digit of the word}
\NormalTok{        SLA    0            * (no{-}op, to keep addressing clean in punch set)}
\NormalTok{        LDA    REM}
\NormalTok{        ADD    ONE}
\NormalTok{        STA    REM}
\NormalTok{        LDA    0, X(2:2)}
\NormalTok{        STA    REM          * park digit temporarily (we’ll use the same code 5×)}
\NormalTok{        LDA    ADDR         * restore A (we keep value in ADDR temporarily)}
\NormalTok{        MUL    TEN}
\NormalTok{        ADD    REM}
\NormalTok{        LDA    REM}
\NormalTok{        ADD    ONE}
\NormalTok{        STA    REM}
\NormalTok{        LDA    ADDR}
\NormalTok{        MUL    TEN}
\NormalTok{        ADD    0, X(3:3)}
\NormalTok{        LDA    REM}
\NormalTok{        ADD    ONE}
\NormalTok{        STA    REM}
\NormalTok{        LDA    ADDR}
\NormalTok{        MUL    TEN}
\NormalTok{        ADD    0, X(4:4)}
\NormalTok{        LDA    REM}
\NormalTok{        ADD    ONE}
\NormalTok{        STA    REM}
\NormalTok{        LDA    ADDR}
\NormalTok{        MUL    TEN}
\NormalTok{        ADD    0, X(5:5)}
\NormalTok{        STA    ADDR         * ADDR temporarily holds the running value}
\NormalTok{* last byte we just consumed is column 20/30/…/80; if \textgreater{}=10 =\textgreater{} negative}
\NormalTok{        LDA    0, X(5:5)}
\NormalTok{        SUB    TEN}
\NormalTok{        JANN   POS          * if \textless{}10, positive}
\NormalTok{        LDA    ZERO}
\NormalTok{        SUB    SIGN}
\NormalTok{        STA    SIGN         * SIGN := {-}1}

\NormalTok{POS     LDA    PTRW}
\NormalTok{        ADD    TWO}
\NormalTok{        STA    PTRW         * move to the next word of the same 10{-}digit field}
\NormalTok{* processed first 5 digits; now the second 5 digits of the field:}
\NormalTok{        LDX    PTRW}
\NormalTok{        LDA    ADDR}
\NormalTok{        MUL    TEN}
\NormalTok{        ADD    0, X(1:1)}
\NormalTok{        MUL    TEN}
\NormalTok{        ADD    0, X(2:2)}
\NormalTok{        MUL    TEN}
\NormalTok{        ADD    0, X(3:3)}
\NormalTok{        MUL    TEN}
\NormalTok{        ADD    0, X(4:4)}
\NormalTok{        MUL    TEN}
\NormalTok{        ADD    0, X(5:5)}
\NormalTok{        STA    ADDR}

\NormalTok{* apply sign and store to memory at ADDR (then ADDR++, k{-}{-})}
\NormalTok{        LDX    SIGN}
\NormalTok{        MUL    SIGN         * A := A * (+1 or {-}1)}
\NormalTok{        LDX    ADDR}
\NormalTok{        STA    0, X         * store result}
\NormalTok{        LDA    ADDR}
\NormalTok{        ADD    ONE}
\NormalTok{        STA    ADDR}
\NormalTok{        LDA    KCTR}
\NormalTok{        SUB    ONE}
\NormalTok{        STA    KCTR}
\NormalTok{        JAZ    LOOP         * if finished this card, read another}
\NormalTok{        JMP    NEXTW        * else next word of same card}

\NormalTok{* {-}{-}{-}{-} transfer card: jump to nnnn {-}{-}{-}{-}}
\NormalTok{XFER    LDA    ZERO         * columns 7..10 live in word 1, bytes 2..5}
\NormalTok{        MUL    TEN}
\NormalTok{        ADD    1(2:2)}
\NormalTok{        MUL    TEN}
\NormalTok{        ADD    1(3:3)}
\NormalTok{        MUL    TEN}
\NormalTok{        ADD    1(4:4)}
\NormalTok{        MUL    TEN}
\NormalTok{        ADD    1(5:5)}
\NormalTok{        JMP    0, A         * JMP nnnn}

\NormalTok{        END    0000         * remember: GO starts at 0000}
\end{Highlighting}
\end{Shaded}

What to notice

\begin{itemize}
\tightlist
\item
  Only safe opcodes are on the cards. We explicitly avoid \texttt{ENT},
  \texttt{INC}, and \texttt{CMP}. Loop control is done with
  \texttt{ADD}/\texttt{SUB} on counters in memory and conditional jumps
  (\texttt{JAZ}, \texttt{JANN}), which are safe to punch.
\item
  All addresses appearing in the loader cards are small decimal numbers
  (0, 1, 2, \ldots, 31, 100\ldots), so the address bytes are harmless
  characters.
\item
  We never store into 0000\ldots0031 except as the card image; actual
  program data is guaranteed to be \texttt{\textgreater{}\ 0100}.
\item
  Parsing uses partial-field loads \texttt{M(1:1)…M(5:5)} to fetch
  digits straight from the card image; the sign is derived from the last
  byte by the ``≥10'' test (we avoid \texttt{CMP} by subtracting 10 and
  using \texttt{JANN}).
\item
  The program is deliberately written ``straight'' (no self-modifying
  field tricks), so it remains readable and easy to validate. If you
  want to shave a couple of words you can collapse some repeated
  \texttt{MUL\ TEN} sequences by unrolling constants, but we stay
  \textless{} 32 words already.
\end{itemize}

Size check

\begin{itemize}
\tightlist
\item
  Stage-1 = 2 words (0000\ldots0001).
\item
  Stage-2 = 30 words (0016\ldots0031).
\item
  Total = 32 words ⇒ exactly two cards, as required.
\end{itemize}

Using the loader

Deck layout:

\begin{enumerate}
\def\labelenumi{\arabic{enumi}.}
\tightlist
\item
  Card 1: stage-1 (the two instructions above, in character code).
\item
  Card 2: stage-2 (the 30-word body, in character code).
\item
  Information cards in the specified format.
\item
  Transfer card, columns 1--10 = \texttt{TRANS0nnnn}.
\end{enumerate}

Push GO. Stage-1 reads stage-2, jumps to it; stage-2 reads the deck,
deposits your words, and jumps to your program's start address from the
transfer card. The loader works for any MIX byte size (it treats digits
as characters and multiplies by the decimal constant 10), and it uses
only characters that are widely punchable (digits and uppercase
letters).

\bookmarksetup{startatroot}

\chapter{1.3.2 The MIX assembly
language}\label{the-mix-assembly-language}

\subsection{1 {[}00{]}}\label{section-124}

Use a literal constant, load it to A, then store the whole word to X.

\begin{verbatim}
        LDA   =+1000=     * A ← 1000 (literal)
        STA   X           * X ← 1000 (stores full word, field 0:5)
...
X       CON   0           * (somewhere: reserve the cell for X)
\end{verbatim}

\texttt{X\ EQU\ 1000} only defines a symbol; it doesn't write memory.
\texttt{LDA\ =1000=} makes the value 1000 available at run time, and
\texttt{STA\ X} sets the memory cell named \texttt{X} to 1000.

\subsection{2 {[}10{]}}\label{section-125}

Because the word labeled \texttt{EXIT} (which contains \texttt{JMP\ *})
is self-modified at entry to the subroutine.

Line 3 does \texttt{STJ\ EXIT\ 1}. In MIX, \texttt{STJ\ …,1} stores the
contents of \texttt{rJ} (the caller's return address) into the address
field of the memory word at \texttt{EXIT}. Thus the instruction at
\texttt{EXIT} is changed from

\begin{verbatim}
JMP *
\end{verbatim}

(where \texttt{*} initially means ``this location'') to

\begin{verbatim}
JMP <return address>
\end{verbatim}

When line 12 executes, it jumps to that stored return address, not to
itself, so there is no infinite loop.

\subsection{3 {[}23{]}}\label{section-126}

It's a selection-sort wrapper around subroutine \textbf{M (MAXIMUM)}.

\begin{enumerate}
\def\labelenumi{\arabic{enumi}.}
\item
  \texttt{IN\ X+1(0)} / \texttt{JBUS\ *(0)} reads 100 words into
  \texttt{X{[}1{]}..X{[}100{]}} and waits.
\item
  \texttt{ENT1\ 100} sets \texttt{n\ ←\ 100}.
\item
  Loop (\texttt{1H\ …\ J1P\ 1B}):

  \begin{itemize}
  \tightlist
  \item
    \texttt{JMP\ MAXIMUM} calls Program M, which returns the maximum of
    \texttt{X{[}1..n{]}} in \texttt{A} and its index in
    \texttt{rI2\ =\ j}.
  \item
    \texttt{LDX\ X,1} loads \texttt{X{[}n{]}} into \texttt{X}.
  \item
    \texttt{STA\ X,1} stores the max at \texttt{X{[}n{]}}.
  \item
    \texttt{STX\ X,2} stores the old \texttt{X{[}n{]}} at
    \texttt{X{[}j{]}} (swap).
  \item
    \texttt{DEC1\ 1} decrements \texttt{n}; repeat while
    \texttt{n\ \textgreater{}\ 0}.
  \end{itemize}
\item
  \texttt{OUT\ X+1(1)} outputs the block; \texttt{HLT} stops.
\end{enumerate}

Effect: it sorts the 100 input numbers in nondecreasing (ascending)
order, placing the largest remaining element at the end on each pass,
then prints the sorted list.

\subsection{4 {[}25{]}}\label{section-127}

Here's a clean hand-assembly of \textbf{Program P}. I'm using the MIX
instruction format (ADDRESS, INDEX, MOD, OPCODE); \texttt{MOD} is either
the field spec (for LD/ST/CMP/ADD/SUB/\ldots; default (0:5)
=\textgreater{} 5) or the operation modifier (for ENTx/INCx/DECx, Jx,
I/O, etc.). MIX opcodes and field rules come straight from the standard
tables.

\textbf{Symbols resolved}

\begin{itemize}
\tightlist
\item
  \texttt{L\ =\ 500}, \texttt{PRINTER\ =\ 18}, \texttt{PRIME\ =\ -1},
  \texttt{BUF0\ =\ 2000}, \texttt{BUF1\ =\ BUF0+25\ =\ 2025}.
\item
  Local labels: \texttt{2H} at 3003 and 3016 and 3019; \texttt{4H} at
  3006 and 3020; \texttt{6H} at 3008.
\item
  Literal pool (placed \emph{after} the program end):
  \texttt{=1-L=\ →\ -499}, \texttt{=3=\ →\ 3}.
\end{itemize}

\textbf{Code (ORIG 3000)}

Each line shows:
\texttt{LOC\ :\ +\ ADDRESS\ \ INDEX\ \ MOD\ \ OPCODE\ \ \ ;\ mnemonic\ /\ remark}.

\begin{verbatim}
3000: +    0  0 18 35   ; IOC 0(PRINTER)
3001: + 2050  0  5  9   ; LD1 =1-L=         (literal at 2050)
3002: + 2051  0  5 10   ; LD2 =3=           (literal at 2051)
3003: +    1  0  0 49   ; INC1 1            (2H)
3004: +  499  1  5 26   ; ST2 PRIME+L,1
3005: + 3016  0  1 41   ; J1Z 2F
3006: +    2  0  0 50   ; INC2 2            (4H)
3007: +    2  0  2 51   ; ENT3 2
3008: +    0  0  2 48   ; ENTA 0            (6H)
3009: +    0  2  2 55   ; ENTX 0,2
3010: +   -1  3  5  4   ; DIV PRIME,3
3011: + 3006  0  1 47   ; JXZ 4B
3012: +   -1  3  5 56   ; CMPA PRIME,3
3013: +    1  0  0 51   ; INC3 1
3014: + 3008  0  6 39   ; JG 6B
3015: + 3003  0  0 39   ; JMP 2B
3016: + 1995  0 18 37   ; OUT TITLE(PRINTER)      (2H)
3017: + 2035  0  2 52   ; ENT4 BUF1+10
3018: +  -50  0  2 53   ; ENT5 -50
3019: +  501  0  0 53   ; INC5 L+1                (2H)
3020: +   -1  5  5  8   ; LDA PRIME,5             (4H)
3021: +    0  0  1  5   ; CHAR
3022: +    0  4 12 31   ; STX 0,4(1:4)  (MOD=8*1+4=12)
3023: +    1  0  1 52   ; DEC4 1
3024: +   50  0  1 53   ; DEC5 50
3025: + 3020  0  2 45   ; J5P 4B
3026: +    0  4 18 37   ; OUT 0,4(PRINTER)
3027: +   24  4  5 12   ; LD4 24,4
3028: + 3019  0  0 45   ; J5N 2B
3029: +    0  0  2  5   ; HLT
\end{verbatim}

Why those numbers?

\begin{itemize}
\tightlist
\item
  LD/STA/\ldots{} opcodes and default field \texttt{5} (i.e., (0:5)) per
  MIX tables.
\item
  ENTx/INCx/DECx share opcodes 48--55; \texttt{MOD} distinguishes ENT
  (2), INC (0), DEC (1).
\item
  Jumps: \texttt{J1Z} uses opcode 40+i (=41) with \texttt{MOD=1};
  \texttt{JXZ} uses opcode 47 with \texttt{MOD=1};
  \texttt{JG}/\texttt{JMP} use opcode 39 with \texttt{MOD=6}/\texttt{0}.
\item
  I/O: \texttt{IOC}/\texttt{OUT} are opcodes 35/37 with
  \texttt{MOD=unit\#} (18 = line printer).
\item
  \texttt{CHAR} is opcode 5 with \texttt{MOD=1}; \texttt{HLT} is opcode
  5 with \texttt{MOD=2}.
\end{itemize}

\textbf{Data and text words}

MIX stores numbers base-64 per byte; below I show each word as five byte
values (sign shown first). Example: 2000 =
\texttt{+\ 00\ 00\ 00\ 31\ 16}.

Primes table head (ORIG PRIME+1 = 0):

\begin{verbatim}
0000: + 00 00 00 00 02   ; CON 2  (PRIME[1])
\end{verbatim}

Title (ORIG BUF0--5 = 1995) --- ALF encodes characters to MIX byte codes
(A=1\ldots Z=29, space=0, digits 0--9→30--39, etc.).

\begin{verbatim}
1995: + 06 09 19 22 23   ; ALF 'FIRST'
1996: + 00 06 09 25 05   ; ALF ' FIVE'
1997: + 00 08 24 15 04   ; ALF ' HUND'
1998: + 19 05 04 00 17   ; ALF 'RED P'
1999: + 19 09 14 05 22   ; ALF 'RIMES'
\end{verbatim}

Buffer cross-pointers:

\begin{verbatim}
2024: + 00 00 00 31 51   ; CON BUF1+10 (=2035)
2049: + 00 00 00 31 26   ; CON BUF0+10 (=2010)
\end{verbatim}

Literal pool (placed after END):

\begin{verbatim}
2050: - 00 00 00 07 51   ; =1-L=  (−499)
2051: + 00 00 00 00 03   ; =3=    (+3)
\end{verbatim}

\begin{Shaded}
\begin{Highlighting}[]
\NormalTok{* Program P — fully assembled numeric image (matches Knuth’s listing).}
\NormalTok{* Symbols:  L=500, PRINTER=18, PRIME={-}1, BUF0=2000, BUF1=BUF0+25=2025}
\NormalTok{* Local labels: 2H at 3003, 3016, 3019; 4H at 3006, 3020; 6H at 3008}
\NormalTok{* Word layout shown as:  sign  ADDRESS  INDEX  MOD  OPCODE}
\NormalTok{* Notes:}
\NormalTok{*  {-} For I/O, MOD = device number. Line printer is 18.}
\NormalTok{*  {-} For LD/ST/CMP/arith, MOD = field spec; default (0:5) is encoded as 5.}
\NormalTok{*  {-} ENT/INC/DEC families use MOD=2/0/1 respectively; register selects opcode.}
\NormalTok{*  {-} Jumps on Ix use opcodes 40..47 (rI1..rX); MOD encodes the condition.}
\NormalTok{*  {-} JG/JMP use opcode 39 with MOD=6/0. CHAR and HLT are opcode 5 with MOD=1/2.}

\NormalTok{3000: + 0000 00 18 35   ; IOC 0(PRINTER) — device 18 → MOD=18, opcode 35}
\NormalTok{3001: + 2051 00 05 09   ; LD1 =1{-}L=     — literal at 2051 (=1−500={-}499), field (0:5)}
\NormalTok{3002: + 2050 00 05 10   ; LD2 =3=       — literal at 2050 (+3), field (0:5)}
\NormalTok{3003: + 0001 00 00 49   ; INC1 1        — MOD=0 for INC, opcode for rI1 is 49      (2H)}
\NormalTok{3004: + 0499 01 05 26   ; ST2 PRIME+L,1 — addr = {-}1+500 = +499, indexed by rI1, store rI2}
\NormalTok{3005: + 3016 00 01 41   ; J1Z 2F        — jump if rI1==0 to forward “2H” at 3016}
\NormalTok{3006: + 0002 00 00 50   ; INC2 2        — rI2 ← rI2+2                                 (4H)}
\NormalTok{3007: + 0002 00 02 51   ; ENT3 2        — rI3 ← +2 (ENT uses MOD=2)}
\NormalTok{3008: + 0000 00 02 48   ; ENTA 0        — rA  ← 0                                     (6H)}
\NormalTok{3009: + 0000 02 02 55   ; ENTX 0,2      — rX  ← rI2 (index 2 named after comma{-}2)}
\NormalTok{3010: {-} 0001 03 05 04   ; DIV PRIME,3   — divide A‖X by mem[{-}1+X3], fld (0:5)}
\NormalTok{3011: + 3006 00 01 47   ; JXZ 4B        — if rX==0 goto back “4H” at 3006}
\NormalTok{3012: {-} 0001 03 05 56   ; CMPA PRIME,3  — compare A with mem[{-}1+X3], fld (0:5)}
\NormalTok{3013: + 0001 00 00 51   ; INC3 1        — rI3 ← rI3+1}
\NormalTok{3014: + 3008 00 06 39   ; JG 6B         — if A\textgreater{}0 goto back “6H” at 3008 (MOD=6)}
\NormalTok{3015: + 3003 00 00 39   ; JMP 2B        — unconditional to back “2H” at 3003}
\NormalTok{3016: + 1995 00 18 37   ; OUT TITLE(PRINTER) — addr of TITLE=1995, device 18     (2H)}
\NormalTok{3017: + 2035 00 02 52   ; ENT4 BUF1+10  — rI4 ← 2025+10=2035}
\NormalTok{3018: {-} 0050 00 02 53   ; ENT5 {-}50      — rI5 ← {-}50}
\NormalTok{3019: + 0501 00 00 53   ; INC5 L+1      — rI5 ← rI5 + 501 (L+1=501)                 (2H)}
\NormalTok{3020: {-} 0001 05 05 08   ; LDA PRIME,5   — A ← mem[{-}1] (full word)                    (4H)}
\NormalTok{3021: + 0000 00 01 05   ; CHAR          — convert A to 10 chars packed into A‖X}
\NormalTok{3022: + 0000 04 12 31   ; STX 0,4(1:4)  — store X’s middle bytes; (1:4)→MOD=8*1+4=12; idx rI4}
\NormalTok{3023: + 0001 00 01 52   ; DEC4 1        — rI4 ← rI4−1  (MOD=1 for DEC)}
\NormalTok{3024: + 0050 00 01 53   ; DEC5 50       — rI5 ← rI5−50; will loop while still \textgreater{}0}
\NormalTok{3025: + 3020 00 02 45   ; J5P 4B        — if rI5\textgreater{}0 goto back “4H” at 3020}
\NormalTok{3026: + 0000 04 18 37   ; OUT 0,4(PRINTER) — print buffer whose addr is in rI4 (index 4), device 18}
\NormalTok{3027: + 0024 04 05 12   ; LD4 24,4      — rI4 ← mem[24+X4] (switch buffer pointer)}
\NormalTok{3028: + 3019 00 00 45   ; J5N 2B        — if rI5\textless{}0 goto back “2H” at 3019 (done when ≤0)}
\NormalTok{3029: + 0000 00 02 05   ; HLT           — stop (opcode 5, MOD=2)}

\NormalTok{* {-}{-}{-}{-}{-} Data and text areas {-}{-}{-}{-}{-}{-}{-}{-}{-}{-}{-}{-}{-}{-}{-}{-}{-}{-}{-}{-}{-}{-}{-}{-}{-}{-}{-}{-}{-}{-}{-}{-}{-}{-}{-}{-}{-}{-}{-}{-}{-}{-}{-}{-}{-}{-}{-}{-}{-}{-}{-}{-}}

\NormalTok{0000: +                                2  ; PRIME[1] = 2   (ORIG PRIME+1 / CON 2)}

\NormalTok{1995: + 06 09 19 22 23                     ; ALF \textquotesingle{}FIRST\textquotesingle{}}
\NormalTok{1996: + 00 06 09 25 05                     ; ALF \textquotesingle{} FIVE\textquotesingle{}  (note leading blank)}
\NormalTok{1997: + 00 08 24 15 04                     ; ALF \textquotesingle{} HUND\textquotesingle{}}
\NormalTok{1998: + 19 05 04 00 17                     ; ALF \textquotesingle{}RED P\textquotesingle{}}
\NormalTok{1999: + 19 09 14 05 22                     ; ALF \textquotesingle{}RIMES\textquotesingle{}}

\NormalTok{2024: +                               2035 ; BUF0+24 → BUF1+10 (cross{-}pointer)}
\NormalTok{2049: +                               2010 ; BUF1+24 → BUF0+10 (cross{-}pointer)}

\NormalTok{* {-}{-}{-}{-}{-} Literal pool placed after the program {-}{-}{-}{-}{-}{-}{-}{-}{-}{-}{-}{-}{-}{-}{-}{-}{-}{-}{-}{-}{-}{-}{-}{-}{-}{-}{-}{-}{-}{-}{-}{-}{-}}
\NormalTok{2050: +                                  3 ; =3=          (used by LD2 at 3002)}
\NormalTok{2051: {-}                                499 ; =1{-}L= = 1−500 (used by LD1 at 3001)}

\NormalTok{* Sanity checks you can do when re{-}deriving:}
\NormalTok{*  {-} All forward/backward locals resolve to the stated addresses:}
\NormalTok{*      2F=3016, 4B=3006/3020, 6B=3008, 2B=3019/3003 as indicated.}
\NormalTok{*  {-} Field (1:4) encodes as MOD=8*1+4=12 (see 3022).}
\NormalTok{*  {-} I/O MODs are 18 throughout; jump MODs: Z=1, P=2, N=0; JG uses MOD=6.}
\end{Highlighting}
\end{Shaded}

\subsection{5 {[}11{]}}\label{section-128}

Because Program P does a \emph{lot} of work between two output
operations and it double-buffers the line being printed.

\texttt{OUT\ 0,4(PRINTER)} starts the printer asynchronously and the
program immediately switches \texttt{rI4} to the other buffer
(\texttt{LD4\ 24,4}). It then spends a long time computing the next 50
primes and formatting them into that other buffer. By the time it
reaches the next \texttt{OUT}, the printer has finished the previous
line and is ready again. Since the program never touches the buffer
that's currently being printed (thanks to the alternating buffers), it
doesn't need to poll the device with \texttt{JBUS}.

\subsection{6 {[}HM20{]}}\label{hm20-17}

\begin{enumerate}
\def\labelenumi{(\alph{enumi})}
\item
  If \(n\) is composite, write \(n=ab\) with \(1<a\le b<n\). If both
  \(a>\sqrt n\) and \(b>\sqrt n\), then \(ab>\sqrt n\cdot\sqrt n=n\), a
  contradiction. Hence at least one factor is \(\le\sqrt n\). Let
  \(d=\min(a,b)\); then \(1<d\le\sqrt n\).
\item
  In Program \textbf{P}, step \textbf{P6} divides \(N\) by the current
  prime \(p=\text{PRIME}[k]\): \(N = pQ + R\) with \(0\le R<p\). MIX
  puts \(Q\) in \(rA\) and \(R\) in \(rX\). If \(R=0\) the program
  declares ``not prime''.
\end{enumerate}

Step \textbf{P7} (the code \texttt{CMPA\ PRIME,3;\ INC3\ 1;\ JG\ 6B})
compares that quotient \(Q\) with the \textbf{next} prime
\(p'=\text{PRIME}[k+1]\) and loops back to \textbf{P6} exactly when
\(p' \le Q\).

Why does stopping when \(p' > Q\) prove primality? If \(p\le\sqrt N\)
then
\(\displaystyle Q=\left\lfloor \frac{N}{p}\right\rfloor \ge \sqrt N\).
Therefore any prime divisor \(\le\sqrt N\) is \(\le Q\). Program
\textbf{P} tries primes in increasing order; as long as there could
still be a prime \(\le\sqrt N\) left, we have \(p' \le Q\) and the loop
continues. When the test first finds \(p' > Q\), every remaining prime
exceeds \(Q \ge \sqrt N\). Since we have already checked all smaller
primes (and none divided \(N\)), \(N\) has no divisor \(\le\sqrt N\). By
part (a), \(N\) must be prime.

\subsection{7 {[}10{]}}\label{section-129}

\begin{enumerate}
\def\labelenumi{(\alph{enumi})}
\item
  In MIXAL a digit with ``B'' means ``the nearest \textbf{previous} line
  whose location is that digit followed by H.'' So the \texttt{4B} in
  line 34 means ``jump to the most recent \texttt{4H} before this
  line,'' which here is line 29 (address 3020). Thus \texttt{J5P\ 4B}
  jumps back to 3020 when rI5 \textgreater{} 0.
\item
  If we relabel line 15 as \texttt{2H} and change the target of line 20
  to \texttt{2B}:
\end{enumerate}

\begin{itemize}
\tightlist
\item
  Line 20 will still jump to line 15 (because \texttt{2B} now refers to
  the nearest previous \texttt{2H}, which would be the newly labeled
  line 15). So that branch is unchanged.
\item
  But line 24 already contains \texttt{JMP\ 2B}. After the relabel, its
  target changes from the old \texttt{2H} at line 12 (step P2: store the
  prime and increment \(J\)) to the new \texttt{2H} at line 15 (step P4:
  advance \(N\)). That means whenever \(N\) is prime the program would
  skip storing it and skip incrementing \(J\). Consequently \(J\) never
  reaches 500, so the ``500 found?'' test never triggers and the
  printing phase never runs---effectively breaking the program (it will
  keep searching forever).
\end{itemize}

So: (a) ``previous 4H'' (line 29). (b) The change is not harmless; it
makes line 24 jump to the wrong place and the program fails to record
primes or terminate.

\subsection{8 {[}24{]}}\label{section-130}

Here's the original ``mystery program'' rewritten with \textbf{inline
MIXAL comments} so you can see exactly what every line does:

\begin{verbatim}
* MYSTERY PROGRAM (annotated)

BUF   ORIG **+3000          * Define BUF as an address 3000 words ahead of here.
                            * Code that follows runs before BUF; BUF.. is our data area.

1H    ENT1 1                * I1 ← 1           (run-length to write)
      ENT2 0                * I2 ← 0           (write pointer into BUF)
      LDX  4F               * X  ← word at 4H  (ALF AAAAA) → X="AAAAA"

2H    ENT3 0,1              * I3 ← 0+I1        (copy I1 into I3: loop counter)

3H    STZ  BUF,2            * BUF[I2] ← +0     (store a blank word)
      INC2 1                * I2++             (advance pointer)
      DEC3 1                * I3--             (one blank done)
      J3P  3B               * if I3>0, repeat from 3H (write exactly I1 blanks)

      STX  BUF,2            * BUF[I2] ← X      ('AAAAA' after the blanks)
      INC2 1                * I2++             (advance after the 'A' word)
      INC1 1                * I1++             (next time, one more blank)
      CMP1 =75=             * compare I1 with 75
      JL   2B               * if I1<75, go back to 2H and build next block

      ENN2 2400             * I2 ← −2400       (set up to print from BUF upward)
      OUT  BUF+2400,2(18)   * print one line (24 words) starting at BUF+2400+I2;
                            * device 18 = line printer. With I2<0 this starts at BUF.
      INC2 24               * I2 += 24         (next line starts 24 words later)
      J2N  *-2              * while I2<0, jump back two lines to OUT (prints 100 lines)
      HLT                   * done

4H    ALF  AAAAA            * AAAAA literal used above (loaded into X)
      END  1B               * entry is at label 1H
\end{verbatim}

\textbf{What it does}

Phase 1 (building the buffer). Starting at \texttt{BUF}, the program
writes a pattern: for \texttt{k\ =\ 1,\ 2,\ …,\ 74}, it stores \textbf{k
blank words} (each is \texttt{+0}, which prints as five spaces) followed
by the word \textbf{\texttt{AAAAA}}. After the last iteration
(\texttt{k=74}) it stops. The total words written are
\(\sum_{k=1}^{74}(k+1) = 2849\). The positions of the \texttt{AAAAA}
words within the buffer are

\[
p_m=\sum_{k=1}^{m}(k+1)=\frac{m(m+3)}{2}\quad(m=1,2,\dots,74)
\]

i.e., at word numbers \(2,5,9,14,20,27,\dots\) with gaps growing by
\(3,4,5,\ldots\).

Phase 2 (printing). It then prints the \textbf{first 2400 words} of this
buffer as \textbf{100 lines} of the line printer (24 words = 120
characters per line), starting at \texttt{BUF}, then \texttt{BUF+24},
\texttt{BUF+48}, \ldots, \texttt{BUF+2376}. Because blank words render
as spaces and the nonblank words render as \texttt{AAAAA}, the page
shows a sequence of \texttt{AAAAA} blocks separated by ever-increasing
runs of spaces---a triangular/``quadratically spaced'' pattern of A's
drifting across the printed lines.

\subsection{9 {[}25{]}}\label{section-131}

Here's a clean, table-driven MIXAL solution that validates a word in
\texttt{INST}. It jumps to \texttt{GOOD} iff the instruction has a valid
C-field (opcode), ±AA field (address) when required, I-field (index),
and F-field; otherwise it jumps to \texttt{BAD}.

\begin{Shaded}
\begin{Highlighting}[]
\NormalTok{* EX 1.3.2–9 — Verify that INST is a legal MIX instruction.}
\NormalTok{* If legal, jump to GOOD; otherwise jump to BAD.}

\NormalTok{            ORIG 3000}

\NormalTok{* {-}{-}{-}{-}{-}{-}{-}{-}{-}{-} Small constants and helpers {-}{-}{-}{-}{-}{-}{-}{-}{-}{-}{-}{-}{-}{-}{-}{-}{-}{-}{-}{-}{-}{-}{-}{-}{-}{-}{-}{-}{-}{-}{-}{-}{-}{-}{-}{-}{-}{-}}

\NormalTok{B       EQU  1(4:4)          * MIX byte size (b = 64) via the canonical trick}
\NormalTok{BMAX    EQU  B{-}1             * for MOVE: F ∈ [0..b{-}1]}
\NormalTok{UMAX    EQU  20              * highest I/O unit number in the MIX spec}

\NormalTok{* VALID packs several convenient bounds into one word:}
\NormalTok{* (3:3)=6  → upper bound for I}
\NormalTok{* address=3999 → largest legal memory address}
\NormalTok{* (4:4)=6  → “floating” F used by arithmetic ops (ADD/SUB/MUL/DIV)}
\NormalTok{VALID   CMPX 3999,6(6)}

\NormalTok{INST    CON  0               * (caller puts candidate instruction here)}

\NormalTok{* {-}{-}{-}{-}{-}{-}{-}{-}{-}{-} Common front{-}end checks + table dispatch {-}{-}{-}{-}{-}{-}{-}{-}{-}{-}{-}{-}{-}{-}{-}{-}{-}{-}{-}{-}{-}{-}{-}{-}{-}{-}}

\NormalTok{BEGIN   LDA   INST                   * A ← INST (full word)}
\NormalTok{        CMPA  VALID(3:3)             * compare I with 6}
\NormalTok{        JG    BAD                    * I \textgreater{} 6 → illegal}

\NormalTok{        LD1   INST(5:5)              * I1 ← opcode (C)}
\NormalTok{        DEC1  64}
\NormalTok{        J1NN  BAD                    * C ≥ 64 → illegal (not an instruction)}

\NormalTok{        CMPA  TABLE+64,1(4:4)        * compare F with per{-}opcode Fmax from table}
\NormalTok{        JG    BAD                    * F \textgreater{} Fmax → illegal}

\NormalTok{        LD1   TABLE+64,1(1:2)        * I1 ← address of per{-}opcode checker}
\NormalTok{        JMP   0,1                    * jump there}

\NormalTok{* {-}{-}{-}{-}{-}{-}{-}{-}{-}{-} Per{-}opcode checker routines {-}{-}{-}{-}{-}{-}{-}{-}{-}{-}{-}{-}{-}{-}{-}{-}{-}{-}{-}{-}{-}{-}{-}{-}{-}{-}{-}{-}{-}{-}{-}{-}{-}{-}{-}{-}{-}{-}}

\NormalTok{* FLOAT — arithmetic ops: either a valid partial field or F==6 (“floating”)}
\NormalTok{FLOAT   CMPA  VALID(4:4)             * F == 6 ?}
\NormalTok{        JE    MEMORY                 * if yes, only the address rule remains}
\NormalTok{        JMP   FIELD                  * else check partial field as usual}

\NormalTok{* FIELD — validate F as 8L+R with 0 ≤ L ≤ R ≤ 5}
\NormalTok{* Trick: let F = 8L+R, divide by 9 →  A = ⌊F/9⌋ and X = F mod 9.}
\NormalTok{* {-} If R ≥ L:  ⌊F/9⌋=L, rem=R{-}L → A+rem = R → require R ≤ 5.}
\NormalTok{* {-} If R \textless{} L:  ⌊F/9⌋=L−1, rem=R−L+9 → A+rem = R+8 ≥ 8 → reject.}
\NormalTok{FIELD   ENTA  0}
\NormalTok{        LDX   INST(4:4)              * X ← F}
\NormalTok{        DIV   =9=}
\NormalTok{        STX   *+1(0:2)               * self{-}modify the next INCA with rem}
\NormalTok{        INCA  0                      * A ← ⌊F/9⌋ + (F mod 9)}
\NormalTok{        DECA  5}
\NormalTok{        JAP   BAD                    * if \textgreater{}5 → illegal}
\NormalTok{        JMP   MEMORY                 * finally validate the address if needed}

\NormalTok{* MEMORY — when I=0, address must be a valid memory/unit number}
\NormalTok{MEMORY  LDX   INST(3:3)              * X ← I}
\NormalTok{        JXNZ  GOOD                   * if I ≠ 0, indexing OK → done}
\NormalTok{        LDX   INST(0:2)              * X ← ±AA (or device/unit for I/O)}
\NormalTok{        JXN   BAD                    * negative address/unit not allowed}
\NormalTok{        DECX  3999}
\NormalTok{        JXNP  GOOD                   * 0..3999 is OK (table narrows bound for I/O)}
\NormalTok{        JMP   BAD}

\NormalTok{GOOD    JMP   GOOD                   * (return address patched by caller)}
\NormalTok{BAD     JMP   BAD}

\NormalTok{* {-}{-}{-}{-}{-}{-}{-}{-}{-}{-} Opcode table: Fmax in (4:4), checker address in (1:2) {-}{-}{-}{-}{-}{-}{-}{-}{-}{-}{-}{-}{-}}
\NormalTok{* The pattern is what matters; you can fill remaining opcodes analogously.}
\NormalTok{* N.B.: For opcodes that do not use ±AA, the checker is just GOOD.}
\NormalTok{*       For I/O and MOVE, the checker is MEMORY with the appropriate bound.}

\NormalTok{TABLE   NOP    GOOD(BMAX)            * C=0  : any F up to B−1 is harmless}
\NormalTok{        ADD    FLOAT(5:5)            * C=1  : arithmetic allow F=6 (‘floating’)}
\NormalTok{        SUB    FLOAT(5:5)            * C=2}
\NormalTok{        MUL    FLOAT(5:5)            * C=3}
\NormalTok{        DIV    FLOAT(5:5)            * C=4}
\NormalTok{        HLT    GOOD                  * C=5  : F ignored; no address use}
\NormalTok{        SRC    GOOD                  * C=6  : shifts/rotates; address is count}
\NormalTok{        MOVE   MEMORY(BMAX)          * C=7  : any F in [0..B−1]; addr must be valid}
\NormalTok{        LDA    FIELD(5:5)            * C=8  : partial fields only}
\NormalTok{* ...   (fill remaining opcodes in the same style:}
\NormalTok{*        all loads/stores \& compares → FIELD(5:5);}
\NormalTok{*        address{-}using ops (e.g., Jxx, I/O) → MEMORY with proper bound;}
\NormalTok{*        ops with no address effect (e.g., NOP/HLT/SRC) → GOOD.)}
\NormalTok{* {-}{-}{-}{-}{-}{-}{-}{-}{-}{-}{-}{-}{-}{-}{-}{-}{-}{-}{-}{-}{-}{-}{-}{-}{-}{-}{-}{-}{-}{-}{-}{-}{-}{-}{-}{-}{-}{-}{-}{-}{-}{-}{-}{-}{-}{-}{-}{-}{-}{-}{-}{-}{-}{-}{-}{-}{-}{-}{-}{-}{-}{-}{-}{-}{-}{-}{-}{-}{-}{-}{-}{-}{-}{-}{-}{-}{-}}

\NormalTok{        END   BEGIN}
\end{Highlighting}
\end{Shaded}

\textbf{How it works}

\begin{enumerate}
\def\labelenumi{\arabic{enumi}.}
\item
  \textbf{Front-end} rejects illegal \texttt{I} (must be ≤6) and
  \texttt{C} (must be \textless64), then checks \texttt{F\ ≤\ Fmax(C)}
  from the table.
\item
  \textbf{Dispatch} jumps to an \textbf{opcode-specific checker} from
  the same table entry:

  \begin{itemize}
  \tightlist
  \item
    \texttt{FLOAT} accepts the special case \texttt{F=6} for arithmetic;
    otherwise it falls through to \texttt{FIELD}.
  \item
    \texttt{FIELD} verifies a partial field by dividing \texttt{F} by 9
    (the classic MIX trick), accepting iff \texttt{0\ ≤\ L\ ≤\ R\ ≤\ 5}.
  \item
    \texttt{MEMORY} ensures that when \texttt{I=0} the address (or I/O
    unit) is a valid nonnegative value within bounds (3999 for memory,
    or a tighter bound for I/O/MOVE supplied via the table).
  \end{itemize}
\item
  If all checks pass, control reaches \textbf{GOOD}; otherwise it
  reaches \textbf{BAD}.
\end{enumerate}

This is exactly the ``table + dispatch'' approach Knuth recommends: it's
short, fast, and easy to extend---just adjust the two numbers in each
table entry if an opcode's rules change.

\subsection{10 {[}31{]}}\label{section-132}

Here's a compact, fully annotated MIXAL solution for the \textbf{9×8
saddle point} problem. Matrix element \(a_{ij}\) lives at address
\(1000+8i+j\) for \(1\le i\le 9,\;1\le j\le 8\). The program scans each
row, picks the \textbf{row minimum} \(a_{i j^{*}}\), then checks whether
it is also the \textbf{column maximum} in column \(j^{*}\). If a saddle
point is found, it leaves its \textbf{address} in \texttt{rI1} and
halts; otherwise it halts with \texttt{rI1\ ←\ 0}.

\begin{Shaded}
\begin{Highlighting}[]
\NormalTok{* TAOCP v1 §1.3.2 — Ex.10  Saddle point in 9x8 matrix}
\NormalTok{* a(i,j) is stored at ADDR = 1000 + 8*i + j  (1{-}based i,j)}

\NormalTok{        ORIG 3000}

\NormalTok{ROWS    EQU  9}
\NormalTok{COLS    EQU  8}
\NormalTok{ROW1    EQU  1009          * addr of a(1,1)  = 1000 + 8*1 + 1}
\NormalTok{COLBAS  EQU  1008          * addr of a(1,0)  so COLBAS + j = addr of a(1,j)}

\NormalTok{* Scratch}
\NormalTok{MINVAL  CON  0             * current row minimum value}
\NormalTok{MINJ    CON  0             * column index j* (1..8) of that minimum}
\NormalTok{RSTART  CON  0             * address of a(i,1) for current row}
\NormalTok{TMP     CON  0             * temp cell for building rI1 result}

\NormalTok{* rI1: result (address or 0)}
\NormalTok{* rI3: remaining rows counter}
\NormalTok{* rI4: inner loop counter}
\NormalTok{* rI5: row pointer (address of current element in row)}
\NormalTok{* rI6: column pointer for column check}

\NormalTok{START   ENT1 0             * rI1 ← 0 by default (means "no saddle")}
\NormalTok{        ENT3 ROWS          * rI3 ← 9 rows to process}
\NormalTok{        ENT5 ROW1          * rI5 ← address of a(1,1)}

\NormalTok{RLOOP   ST5  RSTART        * remember row start address a(i,1)}
\NormalTok{        ENT2 1             * rI2 ← j = 1}
\NormalTok{        LDA  0,5           * A ← a(i,1)}
\NormalTok{        STA  MINVAL        * row minimum so far}
\NormalTok{        ST2  MINJ          * j* ← 1}
\NormalTok{        ENT4 COLS{-}1        * rI4 ← 7 more elements to scan in this row}

\NormalTok{* {-}{-}{-}{-} scan the rest of the row, keep the minimum and its column j* {-}{-}{-}{-}}
\NormalTok{RNEXT   INC5 1             * move to a(i,j+1)}
\NormalTok{        INC2 1             * j ← j+1}
\NormalTok{        LDA  0,5           * A ← a(i,j)}
\NormalTok{        CMPA MINVAL}
\NormalTok{        JGE  RKEEP         * if a(i,j) \textgreater{}= current min, keep old min}
\NormalTok{        STA  MINVAL        * else new minimum found}
\NormalTok{        ST2  MINJ}
\NormalTok{RKEEP   DEC4 1}
\NormalTok{        J4P  RNEXT}

\NormalTok{* {-}{-}{-}{-} check whether the row{-}minimum is also a column{-}maximum {-}{-}{-}{-}}
\NormalTok{        LD6  MINJ          * rI6 ← j*}
\NormalTok{        INC6 COLBAS        * rI6 ← 1008 + j*  (points at a(1,j*))}
\NormalTok{        ENT4 ROWS          * rI4 ← 9 rows to test}

\NormalTok{CCHECK  LDA  0,6           * A ← a(k,j*)}
\NormalTok{        CMPA MINVAL}
\NormalTok{        JG   NOTSAD        * some a(k,j*) \textgreater{} MINVAL ⇒ not a saddle}
\NormalTok{        INC6 8             * next row: address +8}
\NormalTok{        DEC4 1}
\NormalTok{        J4P  CCHECK}

\NormalTok{* Success: build the exact address 1000+8*i+j* = RSTART + (j*−1), put in rI1}
\NormalTok{        LDA  RSTART}
\NormalTok{        ADD  MINJ}
\NormalTok{        SUB  =1=}
\NormalTok{        STA  TMP}
\NormalTok{        LD1  TMP           * rI1 ← candidate address}
\NormalTok{        HLT                * done (saddle point found)}

\NormalTok{* {-}{-}{-}{-} move to next row and continue the search {-}{-}{-}{-}}
\NormalTok{NOTSAD  LD5  RSTART        * restore this row\textquotesingle{}s start address}
\NormalTok{        INC5 8             * advance to next row\textquotesingle{}s a(i+1,1)}
\NormalTok{        DEC3 1}
\NormalTok{        J3P  RLOOP}

\NormalTok{        ENT1 0             * no saddle point anywhere}
\NormalTok{        HLT}

\NormalTok{        END  START}
\end{Highlighting}
\end{Shaded}

Why it works. For each row \(i\) the first loop walks from
\texttt{a(i,1)} to \texttt{a(i,8)} keeping the smallest value and its
column \(j^{*}\). The second loop fixes the column \(j^{*}\) (pointer
\texttt{rI6\ =\ 1008\ +\ j\^{}\{*\}}) and marches down by 8 through all
nine rows; if any entry exceeds the row minimum, the candidate cannot be
a saddle point. Passing that test means

\[
a_{i j^{*}} \le a_{i k}\;\forall k\quad\text{and}\quad
a_{i j^{*}} \ge a_{k j^{*}}\;\forall k,
\]

so \(a_{i j^{*}}\) is a saddle point and we return its \textbf{memory
address} in \texttt{rI1}. If none is found after all 9 rows,
\texttt{rI1} remains 0.

\subsection{11 {[}M29{]}}\label{m29-1}

Here \(m=9\) rows and \(n=8\) columns.

\begin{enumerate}
\def\labelenumi{(\alph{enumi})}
\tightlist
\item
  All 72 entries are distinct (all \(72!\) orders equally likely)
\end{enumerate}

For a fixed cell \((i,j)\) look only at the \(m+n-1\) entries in its row
or column: \(\{a_{i1},\dots,a_{in}\}\cup\{a_{1j},\dots,a_{mj}\}\) (with
\(a_{ij}\) counted once).

To be a saddle, \(a_{ij}\) must be \textbf{smaller} than the other
\(n-1\) in its row and \textbf{larger} than the other \(m-1\) in its
column. In the total order of these \(m+n-1\) items this means exactly:

\begin{itemize}
\tightlist
\item
  the \(m-1\) smallest are the other column-\(j\) entries (in any
  order),
\item
  then comes \(a_{ij}\),
\item
  then the \(n-1\) remaining row-\(i\) entries (in any order).
\end{itemize}

The number of favorable orders is \((m-1)!(n-1)!\); all \((m+n-1)!\)
orders are equally likely. Hence

\[
\Pr\{(i,j)\text{ is a saddle}\}=\frac{(m-1)!(n-1)!}{(m+n-1)!}.
\]

With all entries distinct there can be \textbf{at most one} saddle point
(two would force a strict inequality cycle), so the probability the
matrix has a saddle is just \(mn\) times the above:

\[
\boxed{\Pr\{\text{saddle exists}\}
= \frac{mn\,(m-1)!(n-1)!}{(m+n-1)!}}.
\]

For \(m=9,\;n=8\),

\[
\Pr = \frac{72\cdot 8!\cdot 7!}{16!}
= \frac{1}{1430}
\approx 6.993\times 10^{-4}.
\]

\begin{enumerate}
\def\labelenumi{(\alph{enumi})}
\setcounter{enumi}{1}
\tightlist
\item
  Entries are i.i.d. 0/1 with all \(2^{72}\) matrices equally likely
\end{enumerate}

A cell is a saddle iff either • it is a \textbf{0} and its
\textbf{column is all zeros}, or • it is a \textbf{1} and its
\textbf{row is all ones}.

Thus a saddle exists iff there is \textbf{at least one all-zero column}
or \textbf{at least one all-one row}. These two events are disjoint
(they cannot happen together because their intersection cell would have
to be both 0 and 1).

Columns are independent, so

\[
\Pr\{\text{some column all 0}\}
=1-\bigl(1-2^{-9}\bigr)^8.
\]

Rows are independent, so

\[
\Pr\{\text{some row all 1}\}
=1-\bigl(1-2^{-8}\bigr)^9.
\]

Therefore

\[
\boxed{\Pr\{\text{saddle exists}\}
= \bigl[1-(1-2^{-9})^{8}\bigr]
+ \bigl[1-(1-2^{-8})^{9}\bigr]}
= \frac{924744796234036231}{2^{64}}
\approx 0.0501305159.
\]

So: about \(0.0699\%\) in the distinct-entries case, and about
\(5.01\%\) in the 0/1 case.

\subsection{12 {[}HM42{]}}\label{hm42-1}

\textbf{The algorithms}

\begin{itemize}
\tightlist
\item
  \textbf{Sol. 1 (row-min then check its column):} For each row, scan
  the 8 entries to get the row minimum (and all columns where it
  occurs), then for each such column scan downward until you see a
  larger value; if none, you've found a saddle.
\item
  \textbf{Sol. 2 (min of column-maxima):} Phase 1: compute every column
  maximum \(C(j)\) and then \(c_{\min}=\min_j C(j)\). Phase 2: run
  across each row; if you ever see an entry \(<c_{\min}\) the row can't
  contain a saddle; otherwise look for an entry \(=c_{\min}\) (then you
  have a saddle).
\item
  \textbf{Sol. 3 (suggested):} Same as Sol. 2, but in Phase 1 also
  remember the row(s) where the minimal column-maximum occurs; search
  \textbf{only} those row(s) in Phase 2 (usually one row when there are
  no ties).
\end{itemize}

\textbf{Model A: all 72 entries distinct (all orders equally likely)}

Some useful facts for i.i.d. continuous data:

\begin{itemize}
\tightlist
\item
  A row minimum among 8 is typically at the \(\approx 1/9\) quantile.
  When you check its column you meet a larger entry with probability
  \(\approx 1-1/9=8/9\) at each step, so the expected number of
  comparisons in that column is \(\approx 1/(8/9)=1.125\).
\item
  To form each column maximum needs 8 comparisons; finding the minimum
  of the 8 maxima needs 7 more. The min of the 8 column-maxima has mean
  \(\mathbb{E}[c_{\min}] \approx 0.7455\) (for uniform{[}0,1{]}). Thus,
  while scanning a row in Phase 2 of Sol. 2, you hit an entry
  \(<c_{\min}\) after about \(1/\mathbb{E}[c_{\min}]\approx 1.34\)
  comparisons on average.
\end{itemize}

Rough comparison in \textbf{comparisons}:

\begin{itemize}
\tightlist
\item
  \textbf{Sol. 1:} \(9\) rows × (row-min \(=7\) comps + column-check
  \(\approx 1.125\))
  \(\Rightarrow \; 9\times(7+1.125)\approx \mathbf{73}\).
\item
  \textbf{Sol. 2:} Phase 1 \(=8\cdot 8+7= \mathbf{71}\); Phase 2
  \(=9\times 1.34\approx 12\). Total \(\approx \mathbf{83}\).
\item
  \textbf{Sol. 3:} Phase 1 as in Sol. 2 (\(\mathbf{71}\)); Phase 2
  usually just \textbf{one} row scan (\(\approx 8\) comps). Total
  \(\approx \mathbf{79}\).
\end{itemize}

\textbf{Verdict (distinct case):} \textbf{Sol. 1} is best; \textbf{Sol.
3} is close (usually one row in Phase 2), \textbf{Sol. 2} is a bit
slower.

\textbf{Model B: entries are i.i.d. 0/1 with all \(2^{72}\) matrices
equally likely}

Facts from Ex. 11(b):
\(\Pr(\text{some all-zero column}) = 1-(1-2^{-9})^{8}\approx 0.0155\),
\(\Pr(\text{some all-one row}) = 1-(1-2^{-8})^{9}\approx 0.0346\),
overall saddle probability \(\approx 0.0501\).

Expected costs:

\begin{itemize}
\tightlist
\item
  \textbf{Sol. 1:} You must scan \textbf{every row completely} to build
  the list of minima (7 comps/row = 63 total). In a typical row the
  number of zeros is \(\mathbb{E}=4\); each column test sees a 1 after
  about 2 comparisons on average, so ≈\(4\times 2=8\) comps/row for
  column checks. Total \(\approx 9\times(7+8)=\mathbf{135}\) comparisons
  (much less if you hit an all-ones row early, much more favorable than
  all-zeros case, but in expectation this ≈135).
\item
  \textbf{Sol. 2:} Phase 1 still \(\mathbf{71}\) comps. Phase 2: when no
  all-zero column (prob ≈0.9845), \(c_{\min}=1\) and each row is ruled
  out as soon as you see the first 0 (≈2 comps/row, ≈18 total), unless
  you encounter an all-ones row (prob ≈0.0346), in which case you spend
  ≈8 comps in that row and stop. Expected Phase 2 cost ≈
  \(0.9654\cdot18+0.0346\cdot(4\!\times\!2+8)\approx 17.9\). Total
  \(\approx \mathbf{89}\) comparisons.
\item
  \textbf{Sol. 3:} Offers \textbf{no advantage} here: ties in
  column-maxima are ubiquitous (usually all columns have max 1, or if a
  zero column exists it tags \textbf{every} row). So it behaves like
  Sol. 2 but with extra bookkeeping ⇒ slightly worse than
  \(\mathbf{89}\).
\end{itemize}

\textbf{Verdict (0/1 case):} \textbf{Sol. 2} is decisively better on
average. (In the extreme all-zero matrix, Sol. 1 wins, but that case has
probability \(2^{-72}\).)

\begin{itemize}
\tightlist
\item
  \textbf{Distinct entries:} pick \textbf{Solution 1} (row-min then
  column-check).
\item
  \textbf{0/1 entries:} pick \textbf{Solution 2} (min of column-maxima,
  then short row scans).
\item
  The \textbf{suggested Solution 3} (record candidate rows during Phase
  1) helps mainly in the distinct-entries setting with few ties; it
  beats Sol. 2 there but is still typically a hair slower than Sol. 1
  for \(9\times 8\).
\end{itemize}

\subsection{13 {[}28{]}}\label{section-133}

Here's a complete MIXAL program:

\begin{itemize}
\tightlist
\item
  Reads the paper tape (unit 0) \textbf{block by block} into two
  alternating buffers.
\item
  While unit 0 is reading the \emph{next} block, it \textbf{counts} the
  characters of the \emph{previous} block.
\item
  Stops counting at the \textbf{first asterisk} \texttt{*}.
\item
  Ignores blanks.
\item
  Keeps counts for \textbf{A..Z} only.
\item
  Then prints one line per letter whose count is non-zero, in the form
  \texttt{A\ 0010257} (letter, one blank, six digits), one line per
  24-word printer line (unit 18).
\end{itemize}

\begin{Shaded}
\begin{Highlighting}[]
\NormalTok{* EX 1.3.2 — Exercise 13: frequency of letters up to first \textquotesingle{}*\textquotesingle{}}
\NormalTok{* Paper tape on device 0 (as in Ex.3), printer is device 18.}

\NormalTok{* {-}{-}{-}{-}{-}{-}{-}{-}{-}{-}{-}{-}{-}{-}{-}{-}{-}{-}{-}{-}{-}{-} constants \& layout {-}{-}{-}{-}{-}{-}{-}{-}{-}{-}{-}{-}{-}{-}{-}{-}{-}{-}{-}{-}{-}{-}{-}{-}{-}{-}{-}{-}{-}{-}{-}{-}{-}{-}{-}{-}}
\NormalTok{        ORIG 3000}

\NormalTok{PRN     EQU  18                * line printer unit number}
\NormalTok{TAPE    EQU  0                 * paper tape unit number}
\NormalTok{BSZ     EQU  100               * words per input block (any reasonable size)}
\NormalTok{ROW     EQU  24                * words per printer line}

\NormalTok{* MIX character codes we need. We take them from ALF words.}
\NormalTok{* AST is the code of \textquotesingle{}*\textquotesingle{}, SPC the blank, A0 the code of \textquotesingle{}A\textquotesingle{}.}
\NormalTok{ASTW    ALF  *0000             * \textquotesingle{}*\textquotesingle{} in byte (1:1)}
\NormalTok{SPCW    ALF  00000             * blank in each byte}
\NormalTok{A0W     ALF  A0000             * \textquotesingle{}A\textquotesingle{} in byte (1:1)}

\NormalTok{* addresses of buffers and work areas}
\NormalTok{BUF0    EQU  4000              * block buffer 0 (BSZ words)}
\NormalTok{BUF1    EQU  BUF0+BSZ          * block buffer 1}
\NormalTok{LINE    EQU  BUF1+BSZ          * 24{-}word printer line buffer}
\NormalTok{CNT     EQU  LINE+ROW          * table of counts for A..Z (26 words)}

\NormalTok{* {-}{-}{-}{-}{-}{-}{-}{-}{-}{-}{-}{-}{-}{-}{-}{-}{-}{-}{-}{-}{-}{-} tiny “byte helpers” {-}{-}{-}{-}{-}{-}{-}{-}{-}{-}{-}{-}{-}{-}{-}{-}{-}{-}{-}{-}{-}{-}{-}{-}{-}{-}{-}{-}{-}{-}{-}{-}{-}{-}{-}}
\NormalTok{* The three ALF words above let us fetch the needed byte codes portably.}
\NormalTok{GETAST  LDA  ASTW(1:1)         * A ← code(\textquotesingle{}*\textquotesingle{})}
\NormalTok{        STA  AST}
\NormalTok{GETSPC  LDA  SPCW(1:1)         * A ← code(\textquotesingle{} \textquotesingle{})}
\NormalTok{        STA  SPC}
\NormalTok{GETA0   LDA  A0W(1:1)          * A ← code(\textquotesingle{}A\textquotesingle{})}
\NormalTok{        STA  A0}
\NormalTok{        JMP  START}

\NormalTok{AST     CON  0}
\NormalTok{SPC     CON  0}
\NormalTok{A0      CON  0}
\NormalTok{ONE     CON  1}
\NormalTok{SIX     CON  6}

\NormalTok{* {-}{-}{-}{-}{-}{-}{-}{-}{-}{-}{-}{-}{-}{-}{-}{-}{-}{-}{-}{-}{-}{-} main program {-}{-}{-}{-}{-}{-}{-}{-}{-}{-}{-}{-}{-}{-}{-}{-}{-}{-}{-}{-}{-}{-}{-}{-}{-}{-}{-}{-}{-}{-}{-}{-}{-}{-}{-}{-}{-}{-}{-}{-}{-}{-}}
\NormalTok{START   ENT1 0                 * rI1 will hold 0 (default “no more data”)}
\NormalTok{        JMP  ZERO              * clear the CNT[1..26] table}

\NormalTok{* kick the very first read (into BUF0), then wait for it}
\NormalTok{BEGIN   IN   BUF0(TAPE)        * start first block}
\NormalTok{        JBUS *      (TAPE)     * wait until tape not busy (block is in BUF0)}
\NormalTok{        ENT3 BUF0              * rI3 ← addr of “current” buffer}
\NormalTok{        ENT4 BUF1              * rI4 ← addr of “next” buffer}
\NormalTok{        ENT5 0                 * rI5 ← “saw asterisk?” flag  (0=no, 1=yes)}

\NormalTok{READLP  IN   0,4(TAPE)         * start input into NEXT buffer at rI4}
\NormalTok{        ENT2 BSZ               * rI2 ← words remaining to scan in CURRENT buffer}

\NormalTok{SCANW   LDA  0,3               * A ← next word from CURRENT buffer (rI3)}
\NormalTok{        STA  WORD              * remember it (so we can pick bytes)}
\NormalTok{* {-}{-}{-}{-} peel the five characters and count them}
\NormalTok{        LDA  WORD(1:1)}
\NormalTok{        JMP  COUNT}
\NormalTok{        LDA  WORD(2:2)}
\NormalTok{        JMP  COUNT}
\NormalTok{        LDA  WORD(3:3)}
\NormalTok{        JMP  COUNT}
\NormalTok{        LDA  WORD(4:4)}
\NormalTok{        JMP  COUNT}
\NormalTok{        LDA  WORD(5:5)}
\NormalTok{        JMP  COUNT}

\NormalTok{NEXTWD  INC3 1                 * advance to next word}
\NormalTok{        DEC2 1}
\NormalTok{        J2P  SCANW             * more words in this block?}

\NormalTok{* we’re done with CURRENT buffer; wait for I/O (so NEXT is ready),}
\NormalTok{* then swap the roles and continue unless we saw \textquotesingle{}*\textquotesingle{}.}
\NormalTok{        JBUS *      (TAPE)     * wait for the block just started to arrive}
\NormalTok{        ENTA 0}
\NormalTok{        CMPA FLAG              * FLAG set by COUNT when \textquotesingle{}*\textquotesingle{} seen?}
\NormalTok{        JE   DONECOUNT         * yes → stop counting}
\NormalTok{        SWAB                   * swap rI3 ↔ rI4 (CURRENT ↔ NEXT)}
\NormalTok{        JMP  READLP}

\NormalTok{* {-}{-}{-}{-}{-}{-}{-}{-}{-}{-}{-}{-}{-}{-}{-}{-}{-}{-}{-}{-}{-}{-} byte classification \& counting {-}{-}{-}{-}{-}{-}{-}{-}{-}{-}{-}{-}{-}{-}{-}{-}{-}{-}{-}{-}{-}{-}{-}{-}}
\NormalTok{* Entry: A has a character code.}
\NormalTok{* Sets FLAG=1 and exits the counting phase when code is \textquotesingle{}*\textquotesingle{}.}
\NormalTok{* Ignores blanks and any non{-}letter codes outside A..Z.}
\NormalTok{COUNT   CMPA AST}
\NormalTok{        JE   SETEND            * stop at first \textquotesingle{}*\textquotesingle{}}
\NormalTok{        CMPA SPC}
\NormalTok{        JE   RETCOUNT          * ignore blanks}
\NormalTok{        STA  CH                * CH := code}
\NormalTok{        LDA  A0}
\NormalTok{        SUB  CH                * A ← code(\textquotesingle{}A\textquotesingle{}) − code(char)}
\NormalTok{        JN   RETCOUNT          * char \textless{} \textquotesingle{}A\textquotesingle{}}
\NormalTok{        ENTA 26}
\NormalTok{        ADD  CH}
\NormalTok{        SUB  A0                * A ← (char−\textquotesingle{}A\textquotesingle{}) + 26}
\NormalTok{        CMPA 26}
\NormalTok{        JG   RETCOUNT          * char \textgreater{} \textquotesingle{}Z\textquotesingle{}}
\NormalTok{* at this point 0 ≤ char−\textquotesingle{}A\textquotesingle{} ≤ 25; index register rI1 := 1..26}
\NormalTok{        LDA  CH}
\NormalTok{        SUB  A0}
\NormalTok{        ADD  ONE}
\NormalTok{        LD1  A                 * rI1 ← letter index 1..26}
\NormalTok{        LDA  CNT,1}
\NormalTok{        ADD  ONE}
\NormalTok{        STA  CNT,1}
\NormalTok{RETCOUNT JMP  NEXTBYTE}

\NormalTok{SETEND  ENT5 1                 * remember: asterisk was seen}
\NormalTok{        STA  FLAG}
\NormalTok{        JMP  ABORTREST}

\NormalTok{* after every byte we get here; the five loads above alternate these entries:}
\NormalTok{NEXTBYTE JMP  NEXT2            * (these two JMPs are in the}
\NormalTok{NEXT2    JMP  NEXT3            *  sequence of five loads)}
\NormalTok{NEXT3    JMP  NEXT4}
\NormalTok{NEXT4    JMP  NEXT5}
\NormalTok{NEXT5    JMP  NEXTWD}

\NormalTok{ABORTREST ENT2 0                * force end of SCANW loop (finish this block)}
\NormalTok{        JMP  NEXTWD}

\NormalTok{WORD    CON  0}
\NormalTok{CH      CON  0}
\NormalTok{FLAG    CON  0                  * 0 = still counting; 1 = stop at once}

\NormalTok{* {-}{-}{-}{-}{-}{-}{-}{-}{-}{-}{-}{-}{-}{-}{-}{-}{-}{-}{-}{-}{-}{-} finish counting, then print nonzero CNTs {-}{-}{-}{-}{-}{-}{-}{-}{-}{-}{-}{-}{-}{-}}
\NormalTok{DONECOUNT}
\NormalTok{* build one printer line in LINE[0..23], print it, reuse the buffer.}

\NormalTok{        ENT2 1                 * rI2 ← j = 1..26 over letters}
\NormalTok{PRLOOP  LD1  CNT,2}
\NormalTok{        J1Z  NEXTL             * skip letters with zero count}
\NormalTok{* LINE format (8 characters then blanks):  "L  ddddddd"}
\NormalTok{* Start by blanking the 24{-}word line buffer}
\NormalTok{        ENT4 LINE+ROW{-}1}
\NormalTok{        ENT3 ROW}
\NormalTok{CLRLP   STZ  0,4}
\NormalTok{        DEC3 1}
\NormalTok{        J3P  CLRLP}

\NormalTok{* place letter and blank at start}
\NormalTok{        LDA  A0}
\NormalTok{        ADD  =0=               * A ← code(\textquotesingle{}A\textquotesingle{})}
\NormalTok{        ADD  CNTX              * + (j−1) via helper below}
\NormalTok{        STA  LINE(1:1)         * store letter in first byte}
\NormalTok{        STA  LINE(2:2)         * overwrite with itself? (no{-}op but keeps a byte)}
\NormalTok{        LDA  SPC}
\NormalTok{        STA  LINE(2:2)         * put blank in second byte}

\NormalTok{* convert count to 10 digits, then drop last 6 into the line (right aligned)}
\NormalTok{        LDA  CNT,2}
\NormalTok{        CHAR                   * to decimal; rightmost digits are in A/X}
\NormalTok{        ENT4 LINE+1            * rI4 points at word 2 of the line}
\NormalTok{        STX  0,4(1:4)          * copy 4 rightmost digits}
\NormalTok{        DEC4 1}
\NormalTok{        STX  0,4(1:4)          * next 4}
\NormalTok{        DEC4 1}
\NormalTok{        STX  0,4(3:4)          * last 2 (leave other 3 bytes blank)}

\NormalTok{* print the line}
\NormalTok{        OUT  LINE(PRN)}

\NormalTok{NEXTL   INC2 1}
\NormalTok{        CMP2 27}
\NormalTok{        JL   PRLOOP}

\NormalTok{        HLT}

\NormalTok{* helper: CNTX = (j−1) at run time, computed from rI2}
\NormalTok{CNTX    CON  0}
\NormalTok{        ORIG *{-}1}
\NormalTok{        ST2  CNTX}

\NormalTok{* {-}{-}{-}{-}{-}{-}{-}{-}{-}{-}{-}{-}{-}{-}{-}{-}{-}{-}{-}{-}{-}{-} zero the count table {-}{-}{-}{-}{-}{-}{-}{-}{-}{-}{-}{-}{-}{-}{-}{-}{-}{-}{-}{-}{-}{-}{-}{-}{-}{-}{-}{-}{-}{-}{-}{-}{-}{-}}
\NormalTok{ZERO    ENT2 26}
\NormalTok{        ENT4 CNT}
\NormalTok{ZLP     STZ  0,4}
\NormalTok{        INC4 1}
\NormalTok{        DEC2 1}
\NormalTok{        J2P  ZLP}
\NormalTok{        JMP  BEGIN}

\NormalTok{        END  START}
\end{Highlighting}
\end{Shaded}

\textbf{How it works}

\begin{itemize}
\item
  \textbf{Double-buffered I/O.} \texttt{BUF0} and \texttt{BUF1} are
  100-word blocks. We start the next \texttt{IN} into the ``NEXT''
  buffer (\texttt{rI4}) while we are scanning the ``CURRENT'' buffer
  (\texttt{rI3}). The \texttt{JBUS} after a block has been scanned
  ensures the next read has completed before we swap the roles. The
  instruction \texttt{SWAB} is the short way to swap \texttt{rI3} and
  \texttt{rI4} (if your MIX simulator lacks it, do the swap with
  \texttt{ENTA\ 0;\ ENTA\ 0,3;\ ENT3\ 0,4;\ ENT4\ 0,3}).
\item
  \textbf{First asterisk ends the count.} We hold the ``code of
  \texttt{*}'' in memory (\texttt{AST}), fetched portably from an
  \texttt{ALF} word. When a \texttt{*} is seen, we set \texttt{FLAG=1}
  and abruptly end the scan of the current block; the extra block on the
  tape (as the problem promises) allows the last queued \texttt{IN} to
  finish cleanly.
\item
  \textbf{Letter test.} The three \texttt{ALF} helpers give codes for
  \texttt{*}, space, and \texttt{A}. For a byte \texttt{c}, we check
  \texttt{c\ !=\ blank}, \texttt{A\ ≤\ c\ ≤\ Z}. If so we index the
  \texttt{CNT{[}1..26{]}} table by
  \texttt{j\ =\ c\ −\ \textquotesingle{}A\textquotesingle{}\ +\ 1} and
  increment.
\item
  \textbf{Printing.} We reuse the line-setup trick from Program
  \textbf{P}: \texttt{CHAR} converts a number to 10 decimal characters
  in \texttt{A‖X}. We then store the six rightmost digits into the
  beginning of the 24-word \texttt{LINE} buffer (to the right of the
  ``letter + space''), and \texttt{OUT} one printer line (24 words).
  Letters with zero counts are skipped.
\end{itemize}

\subsection{14 {[}31{]}}\label{section-134}

Here's a complete MIXAL solution:

\begin{Shaded}
\begin{Highlighting}[]
\NormalTok{* TAOCP v1 §1.3.2 — Ex.14  Easter date (Gregorian)}
\NormalTok{* Prints dates for 1950..2000 using subroutine EASTER.}

\NormalTok{        ORIG 3000}
\NormalTok{PRN     EQU  18}
\NormalTok{TAPE    EQU  0}
\NormalTok{LINE    EQU  4000           * 24{-}word printer line buffer}

\NormalTok{* {-}{-}{-}{-}{-} character helpers (portable via ALF) {-}{-}{-}{-}{-}}
\NormalTok{ASTW    ALF *0000           * not used here, but pattern shown for ref}
\NormalTok{SPCW    ALF 00000}
\NormalTok{A0W     ALF A0000}
\NormalTok{Z0W     ALF 00000           * \textquotesingle{}0\textquotesingle{} in (1:1)}
\NormalTok{COMMW   ALF ,0000           * \textquotesingle{},\textquotesingle{} in (1:1)}

\NormalTok{* resolved single{-}byte constants}
\NormalTok{SPC     CON  0              * set at start from SPCW(1:1)}
\NormalTok{A0      CON  0              * \textquotesingle{}A\textquotesingle{}}
\NormalTok{D0      CON  0              * \textquotesingle{}0\textquotesingle{}}
\NormalTok{COMMA   CON  0              * \textquotesingle{},\textquotesingle{}}

\NormalTok{* months as five{-}byte words}
\NormalTok{MARCH   ALF MARCH}
\NormalTok{APRIL   ALF APRIL}

\NormalTok{* {-}{-}{-}{-}{-}{-}{-}{-}{-}{-}{-}{-}{-}{-}{-}{-} Subroutine: EASTER {-}{-}{-}{-}{-}{-}{-}{-}{-}{-}{-}{-}{-}{-}{-}{-}{-}{-}{-}{-}{-}{-}{-}{-}{-}{-}{-}{-}{-}{-}{-}{-}{-}{-}{-}{-}{-}{-}{-}{-}{-}}
\NormalTok{* Input : YEAR (\textgreater{}=1583 and \textless{}100000)}
\NormalTok{* Output: prints one line "dd MONTH, yyyy"; returns to caller.}

\NormalTok{EASTER  STJ  RET            * standard linkage}

\NormalTok{* fetch single{-}byte constants (once per call)}
\NormalTok{        LDA  SPCW(1:1);  STA SPC}
\NormalTok{        LDA  A0W(1:1);   STA A0}
\NormalTok{        LDA  Z0W(1:1);   STA D0}
\NormalTok{        LDA  COMMW(1:1); STA COMMA}

\NormalTok{* ===== E1.  G ← (Y mod 19) + 1}
\NormalTok{        LDA  YEAR}
\NormalTok{        DIV  =19=}
\NormalTok{        STX  G            * remainder}
\NormalTok{        LDA  G}
\NormalTok{        ADD  =1=}
\NormalTok{        STA  G}

\NormalTok{* ===== E2.  C ← [Y/100] + 1}
\NormalTok{        LDA  YEAR}
\NormalTok{        DIV  =100=}
\NormalTok{        ADD  =1=}
\NormalTok{        STA  C}

\NormalTok{* ===== E3.  X ← [3C/4] − 12 ;  Z ← [(8C+5)/25] − 5}
\NormalTok{        LDA  C;  MUL =3=;  DIV =4=;  SUB =12=;  STA X}
\NormalTok{        LDA  C;  MUL =8=;  ADD =5=;  DIV =25=; SUB =5=;  STA Z}

\NormalTok{* ===== E4.  D ← [5Y/4] − X − 10}
\NormalTok{        LDA  YEAR;  MUL =5=;  DIV =4=;  SUB X;  SUB =10=;  STA D}

\NormalTok{* ===== E5.  E ← (11G + 20 + Z − X) mod 30 ; (fix{-}ups)}
\NormalTok{        LDA  G; MUL =11=; ADD =20=; ADD Z; SUB X}
\NormalTok{        DIV  =30=}
\NormalTok{        STX  E                     * E = remainder}

\NormalTok{        LDA  E; CMPA =25=; JNE CK24}
\NormalTok{        LDA  G; SUB =11=; JLE CK24}
\NormalTok{        LDA  E; ADD =1=;  STA E    * if E=25 and G\textgreater{}11, E++}

\NormalTok{CK24    LDA  E; CMPA =24=; JNE E6}
\NormalTok{        ADD  =1=;  STA E           * if E=24, E++}

\NormalTok{* ===== E6.  N ← 44 − E ; if N\textless{}21 then N←N+30}
\NormalTok{E6      ENTA 44;  SUB E;  STA N}
\NormalTok{        CMPA =21=; JGE E7}
\NormalTok{        ADD  =30=; STA N}

\NormalTok{* ===== E7.  N ← N + 7 − ((D+N) mod 7)}
\NormalTok{E7      LDA  D; ADD N; DIV =7=}
\NormalTok{        LDA  N; ADD =7=; SUB X; STA N}

\NormalTok{* ===== E8.  if N\textgreater{}31 → (N−31) APRIL ; else → N MARCH}
\NormalTok{        LDA  N; CMPA =31=; JLE ISMAR}
\NormalTok{        SUB  =31=; STA DAY}
\NormalTok{        LDA  APRIL;  STA MON1       * 5 bytes of month, split across line}
\NormalTok{        JMP  BUILD}

\NormalTok{ISMAR   STA  DAY}
\NormalTok{        LDA  MARCH;  STA MON1}

\NormalTok{* {-}{-}{-}{-}{-}{-}{-}{-} Build “dd MONTH, yyyy” in LINE and print {-}{-}{-}{-}{-}{-}{-}{-}{-}{-}{-}{-}{-}{-}{-}{-}{-}{-}{-}{-}{-}{-}{-}{-}{-}{-}{-}}

\NormalTok{BUILD   * Blank out line}
\NormalTok{        ENT4 LINE}
\NormalTok{        ENT2 24}
\NormalTok{CL      STZ  0,4}
\NormalTok{        INC4 1}
\NormalTok{        DEC2 1}
\NormalTok{        J2P  CL}

\NormalTok{* day → two digits at LINE(1:1) and LINE(2:2)}
\NormalTok{        LDA  DAY}
\NormalTok{        DIV  =10=}
\NormalTok{        STX  D1                   * ones}
\NormalTok{        STA  D10                  * tens}

\NormalTok{        LDA  D10;  ADD D0;  STA LINE(1:1)}
\NormalTok{        LDA  D1;   ADD D0;  STA LINE(2:2)}

\NormalTok{* blank after day}
\NormalTok{        LDA  SPC;  STA LINE(3:3)}

\NormalTok{* month (5 letters): split 2+3 across word1(4:5) and word2(1:3)}
\NormalTok{        LDA  MON1; STA LINE(4:5)}
\NormalTok{        STA  LINE+1(1:3)}

\NormalTok{* comma+space after month}
\NormalTok{        LDA  COMMA; STA LINE+1(4:4)}
\NormalTok{        LDA  SPC;   STA LINE+1(5:5)}

\NormalTok{* year (four digits) into LINE+2, bytes 2..5}
\NormalTok{        LDA  YEAR}
\NormalTok{        DIV  =10=}
\NormalTok{        STX  Y4;  STA T}
\NormalTok{        LDA  T;   DIV =10=}
\NormalTok{        STX  Y3;  STA T}
\NormalTok{        LDA  T;   DIV =10=}
\NormalTok{        STX  Y2;  STA T}
\NormalTok{        LDA  T;   DIV =10=}
\NormalTok{        STX  Y1                     * now Y1..Y4 are 0..9}

\NormalTok{        LDA  Y1; ADD D0; STA LINE+2(2:2)}
\NormalTok{        LDA  Y2; ADD D0; STA LINE+2(3:3)}
\NormalTok{        LDA  Y3; ADD D0; STA LINE+2(4:4)}
\NormalTok{        LDA  Y4; ADD D0; STA LINE+2(5:5)}

\NormalTok{        OUT  LINE(PRN)}
\NormalTok{        JMP  RET}

\NormalTok{* {-}{-}{-} locals for EASTER}
\NormalTok{RET     JMP  *                 * filled by STJ on entry}
\NormalTok{YEAR    CON  0}
\NormalTok{G       CON  0}
\NormalTok{C       CON  0}
\NormalTok{X       CON  0}
\NormalTok{Z       CON  0}
\NormalTok{D       CON  0}
\NormalTok{E       CON  0}
\NormalTok{N       CON  0}
\NormalTok{DAY     CON  0}
\NormalTok{MON1    CON  0}
\NormalTok{D1      CON  0}
\NormalTok{D10     CON  0}
\NormalTok{T       CON  0}
\NormalTok{Y1      CON  0}
\NormalTok{Y2      CON  0}
\NormalTok{Y3      CON  0}
\NormalTok{Y4      CON  0}

\NormalTok{* {-}{-}{-}{-}{-}{-}{-}{-}{-}{-}{-}{-}{-}{-}{-}{-} Main program: 1950 .. 2000 {-}{-}{-}{-}{-}{-}{-}{-}{-}{-}{-}{-}{-}{-}{-}{-}{-}{-}{-}{-}{-}{-}{-}{-}{-}{-}{-}{-}{-}{-}{-}{-}{-}}

\NormalTok{MAIN    ENT2 1950            * rI2 = current year}
\NormalTok{LPL     LDA  0,2             * A ← rI2}
\NormalTok{        STA  YEAR}
\NormalTok{        JMP  EASTER          * call subroutine (prints one line)}
\NormalTok{        INC2 1}
\NormalTok{        LDA  0,2}
\NormalTok{        CMPA =2001=}
\NormalTok{        JL   LPL}

\NormalTok{        HLT}
\NormalTok{        END  MAIN}
\end{Highlighting}
\end{Shaded}

\textbf{Notes}

\begin{itemize}
\tightlist
\item
  Every division uses \texttt{DIV\ =k=} so that the \textbf{quotient} is
  left in \texttt{A} (i.e., \(\lfloor\cdot\rfloor\)) and the
  \textbf{remainder} in \texttt{X}; that's exactly what we need for the
  \texttt{/} and \texttt{mod} in \textbf{E1--E7}.
\item
  The two Gregorian ``epact'' corrections in \textbf{E5} are implemented
  literally.
\item
  The print routine avoids \texttt{CHAR} and instead extracts day/year
  digits with successive \texttt{/10} (easy to place bytes exactly where
  we want).
\item
  The line is cleared to blanks each time, then we place: 2 digits of
  the day, a blank, the 5-letter month (split 2+3), a comma+space, and 4
  digits of the year; finally \texttt{OUT\ LINE(PRN)} prints one line.
\end{itemize}

\subsection{15 {[}M30{]}}\label{m30-11}

\textbf{Answer:} the earliest year in which the E5 bug actually
\emph{changes} the Easter date is

\textbf{10317}

(correct date: \textbf{25 MARCH}; buggy program prints \textbf{22
APRIL}).

Below is a complete MIXAL program that searches upward from 1583,
computes Easter twice---once correctly and once with the common E5
``negative remainder'' bug---and stops at the first year whose two dates
differ. It then prints the year and both dates.

\begin{Shaded}
\begin{Highlighting}[]
\NormalTok{* EX 1.3.2 — Exercise 15}
\NormalTok{* Find the earliest year (≥1583) where using a negative remainder in E5}
\NormalTok{* gives a DIFFERENT Easter than the correct Algorithm E.}

\NormalTok{        ORIG 3000}
\NormalTok{PRN     EQU 18}

\NormalTok{LINE    EQU 4000              * 24{-}word printer line buffer}

\NormalTok{D0W     ALF 00000             * to get \textquotesingle{}0\textquotesingle{} code}
\NormalTok{SPCW    ALF 00000}
\NormalTok{COMMW   ALF ,0000}
\NormalTok{MARCH   ALF MARCH}
\NormalTok{APRIL   ALF APRIL}

\NormalTok{D0      CON 0                 * filled from D0W(1:1)}
\NormalTok{SPC     CON 0                 * blank}
\NormalTok{COMMA   CON 0}

\NormalTok{* {-}{-}{-}{-}{-}{-}{-}{-}{-}{-}{-}{-}{-}{-}{-}{-}{-}{-}{-}{-}{-}{-}{-}{-}{-}{-}{-}{-}{-}{-}{-}{-}{-}{-}{-}{-}{-}{-}{-}{-}{-}{-}{-}{-}{-}{-}{-}{-}{-}{-}{-}{-}{-}{-}{-}{-}{-}{-}{-}{-}{-}{-}{-}{-}{-}{-}{-}{-}{-}{-}}
\NormalTok{* MAIN: loop Y from 1583 upward; compare GOOD(Y) and BUG(Y).}
\NormalTok{* We represent a date as N ∈ \{1..61\}: 1..31 = March 1..31; 32..61 = April 1..30.}

\NormalTok{MAIN    LDA  D0W(1:1);  STA D0}
\NormalTok{        LDA  SPCW(1:1); STA SPC}
\NormalTok{        LDA  COMMW(1:1); STA COMMA}

\NormalTok{        ENT2 1583               * rI2 = Y}
\NormalTok{LOOP    ST2  YEAR}

\NormalTok{        JSJ  GOOD               * correct Algorithm E → N in NVAL}
\NormalTok{        JSJ  BUG                * buggy E5 (raw remainder) → N in NBAD}

\NormalTok{        LDA  NVAL}
\NormalTok{        CMPA NBAD}
\NormalTok{        JNE  REPORT             * first mismatch found}
\NormalTok{        INC2 1}
\NormalTok{        JMP  LOOP}

\NormalTok{* {-}{-}{-}{-}{-}{-}{-}{-}{-}{-}{-}{-}{-}{-}{-}{-}{-}{-}{-}{-}{-}{-}{-}{-}{-} report line:  "yyyy  GOOD dd MONTH   BAD dd MONTH"}
\NormalTok{REPORT  ENT4 LINE}
\NormalTok{        ENT3 24}
\NormalTok{CL      STZ  0,4                * blank the line}
\NormalTok{        INC4 1}
\NormalTok{        DEC3 1}
\NormalTok{        J3P  CL}

\NormalTok{* year (four digits) into LINE+0 bytes 1..4}
\NormalTok{        LDA  YEAR}
\NormalTok{        DIV  =10=}
\NormalTok{        STX  Y4;  STA T}
\NormalTok{        LDA  T;   DIV =10=      ; Y3,T}
\NormalTok{        STX  Y3;  STA T}
\NormalTok{        LDA  T;   DIV =10=      ; Y2,T}
\NormalTok{        STX  Y2;  STA T}
\NormalTok{        LDA  T;   DIV =10=      ; Y1,*}
\NormalTok{        STX  Y1}
\NormalTok{        LDA  Y1; ADD D0; STA LINE(1:1)}
\NormalTok{        LDA  Y2; ADD D0; STA LINE(2:2)}
\NormalTok{        LDA  Y3; ADD D0; STA LINE(3:3)}
\NormalTok{        LDA  Y4; ADD D0; STA LINE(4:4)}
\NormalTok{        LDA  SPC;          STA LINE(5:5)}

\NormalTok{* put "GOOD " then the correct date}
\NormalTok{        ALF  GOOD               ; "GOOD" in A}
\NormalTok{        STA  LINE+1             ; bytes 1..5}
\NormalTok{        LDA  SPC;               STA LINE+1(5:5)}
\NormalTok{        LDA  NVAL;              STA NTMP}
\NormalTok{        JSJ  PUTDATE            ; writes at LINE+2 and LINE+3}

\NormalTok{* leave two blanks, then "BAD " and the buggy date}
\NormalTok{        LDA  SPC; STA LINE+4(1:1); STA LINE+4(2:2)}
\NormalTok{        ALF  BAD  ;             ; "BAD " (note trailing space)}
\NormalTok{        STA  LINE+4(3:5)}
\NormalTok{        LDA  NBAD;              STA NTMP}
\NormalTok{        ENT4 LINE+5}
\NormalTok{        JSJ  PUTDATE}

\NormalTok{        OUT  LINE(PRN)}
\NormalTok{        HLT}

\NormalTok{* {-}{-}{-}{-}{-}{-}{-}{-}{-}{-}{-}{-}{-}{-}{-}{-}{-}{-}{-}{-}{-}{-} PUTDATE: render N in words at rI4.. (2 words)}
\NormalTok{* Input:  NTMP = 1..61 (March=1..31, April=32..61), rI4=where to place}
\NormalTok{* Output: “dd MONTH” (no comma, no year).}
\NormalTok{PUTDATE STJ  RETPD}
\NormalTok{        LDA  NTMP}
\NormalTok{        CMPA =31=; JLE PMAR}
\NormalTok{        SUB  =31=; STA DAY}
\NormalTok{        LDA  APRIL;  STA 0,4      * month word at rI4}
\NormalTok{        JMP  PDAY}
\NormalTok{PMAR    STA  DAY}
\NormalTok{        LDA  MARCH;  STA 0,4}

\NormalTok{PDAY    LDA  DAY                 * two digits into bytes (4:5) of word rI4}
\NormalTok{        DIV  =10=}
\NormalTok{        STX  D1;  STA D10}
\NormalTok{        LDA  D10; ADD D0; STA 0,4(4:4)}
\NormalTok{        LDA  D1;  ADD D0; STA 0,4(5:5)}
\NormalTok{        LDA  SPC;             STA 1,4(1:1)}
\NormalTok{        JMP  RETPD}
\NormalTok{RETPD   JMP  *                  * return}


\NormalTok{* {-}{-}{-}{-}{-}{-}{-}{-}{-}{-}{-}{-}{-}{-}{-}{-}{-}{-}{-}{-}{-}{-} GOOD: Algorithm E (correct) → N in NVAL {-}{-}{-}{-}{-}{-}{-}{-}{-}{-}{-}{-}{-}{-}}
\NormalTok{GOOD    STJ  RETG}
\NormalTok{        LDA  YEAR}
\NormalTok{        STZ  NVAL               * clear}
\NormalTok{* E1}
\NormalTok{        DIV  =19=}
\NormalTok{        STX  G}
\NormalTok{        LDA  G; ADD =1=; STA G}
\NormalTok{* E2}
\NormalTok{        LDA  YEAR; DIV =100=; ADD =1=; STA C}
\NormalTok{* E3}
\NormalTok{        LDA  C; MUL =3=; DIV =4=; SUB =12=; STA X}
\NormalTok{        LDA  C; MUL =8=; ADD =5=; DIV =25=; SUB =5=; STA Z}
\NormalTok{* E4}
\NormalTok{        LDA  YEAR; MUL =5=; DIV =4=; SUB X; SUB =10=; STA D}
\NormalTok{* E5 (positive modulus)}
\NormalTok{        LDA  G; MUL =11=; ADD =20=; ADD Z; SUB X}
\NormalTok{        DIV  =30=}
\NormalTok{        STX  E}
\NormalTok{        LDA  E; CMPA =25=; JNE CK24}
\NormalTok{        LDA  G; SUB =11=; JLE CK24}
\NormalTok{        LDA  E; ADD =1=;  STA E}
\NormalTok{CK24    LDA  E; CMPA =24=; JNE E6}
\NormalTok{        ADD  =1=;  STA E}
\NormalTok{* E6}
\NormalTok{E6      ENTA 44; SUB E; STA N}
\NormalTok{        CMPA =21=; JGE E7}
\NormalTok{        ADD  =30=; STA N}
\NormalTok{* E7}
\NormalTok{E7      LDA  D; ADD N; DIV =7=}
\NormalTok{        LDA  N; ADD =7=; SUB X; STA N}
\NormalTok{* E8 — leave N as 1..61}
\NormalTok{        LDA  N; STA NVAL}
\NormalTok{        JMP  RETG}
\NormalTok{RETG    JMP  *}

\NormalTok{* {-}{-}{-}{-}{-}{-}{-}{-}{-}{-}{-}{-}{-}{-}{-}{-}{-}{-}{-}{-}{-}{-} BUG: same but E5 uses raw remainder {-}{-}{-}{-}{-}{-}{-}{-}{-}{-}{-}{-}{-}{-}{-}{-}{-}{-}{-}}
\NormalTok{BUG     STJ  RETB}
\NormalTok{* identical setup as GOOD}
\NormalTok{        LDA  YEAR; DIV =19=; STX G; LDA G; ADD =1=; STA G}
\NormalTok{        LDA  YEAR; DIV =100=; ADD =1=; STA C}
\NormalTok{        LDA  C; MUL =3=; DIV =4=; SUB =12=; STA X}
\NormalTok{        LDA  C; MUL =8=; ADD =5=; DIV =25=; SUB =5=; STA Z}
\NormalTok{        LDA  YEAR; MUL =5=; DIV =4=; SUB X; SUB =10=; STA D}
\NormalTok{* E5 with the bug: rX may be NEGATIVE when 11G+20+Z−X \textless{} 0}
\NormalTok{        LDA  G; MUL =11=; ADD =20=; ADD Z; SUB X}
\NormalTok{        DIV  =30=}
\NormalTok{        STX  E                    * DO NOT add 30 if negative  ← the bug}
\NormalTok{        LDA  E; CMPA =25=; JNE B24}
\NormalTok{        LDA  G; SUB =11=; JLE B24}
\NormalTok{        LDA  E; ADD =1=;  STA E}
\NormalTok{B24     LDA  E; CMPA =24=; JNE B6}
\NormalTok{        ADD  =1=;  STA E}
\NormalTok{B6      ENTA 44; SUB E; STA N}
\NormalTok{        CMPA =21=; JGE B7}
\NormalTok{        ADD  =30=; STA N}
\NormalTok{B7      LDA  D; ADD N; DIV =7=}
\NormalTok{        LDA  N; ADD =7=; SUB X; STA N}
\NormalTok{        LDA  N; STA NBAD}
\NormalTok{        JMP  RETB}
\NormalTok{RETB    JMP  *}

\NormalTok{* {-}{-}{-}{-}{-}{-}{-}{-}{-}{-}{-}{-}{-}{-}{-}{-}{-}{-}{-}{-}{-}{-} data words {-}{-}{-}{-}{-}{-}{-}{-}{-}{-}{-}{-}{-}{-}{-}{-}{-}{-}{-}{-}{-}{-}{-}{-}{-}{-}{-}{-}{-}{-}{-}{-}{-}{-}{-}{-}{-}{-}{-}{-}{-}{-}{-}}
\NormalTok{YEAR    CON 0}
\NormalTok{G       CON 0}
\NormalTok{C       CON 0}
\NormalTok{X       CON 0}
\NormalTok{Z       CON 0}
\NormalTok{D       CON 0}
\NormalTok{E       CON 0}
\NormalTok{N       CON 0}
\NormalTok{NVAL    CON 0}
\NormalTok{NBAD    CON 0}
\NormalTok{NTMP    CON 0}
\NormalTok{DAY     CON 0}
\NormalTok{Y1      CON 0}
\NormalTok{Y2      CON 0}
\NormalTok{Y3      CON 0}
\NormalTok{Y4      CON 0}
\NormalTok{T       CON 0}
\NormalTok{D1      CON 0}
\NormalTok{D10     CON 0}

\NormalTok{        END MAIN}
\end{Highlighting}
\end{Shaded}

\textbf{What the program finds}

Running the logic (or checking by hand with an independent calculator)
shows:

\begin{itemize}
\item
  For all years \(1583 \le Y < \mathbf{10317}\), the bad remainder in E5
  never changes the \emph{final} Easter date; either the quantity in E5
  isn't negative yet, or later steps happen to cancel the error.
\item
  In \textbf{\(Y=\mathbf{10317}\)} we get

  \begin{itemize}
  \tightlist
  \item
    correct: \(E=23\Rightarrow N=25\Rightarrow\) \textbf{25 MARCH},
  \item
    buggy: \(E=-7\Rightarrow N=52\Rightarrow\) \textbf{22 APRIL}.
  \end{itemize}
\end{itemize}

So \textbf{10317} is the \emph{earliest} year where the ``negative
remainder in E5'' bug leads to a wrong Easter.

\begin{quote}
How often do the Gregorian Easter dates repeat?
\end{quote}

Answer: the full sequence repeats every \textbf{5,700,000 years}.

Why:

\begin{itemize}
\tightlist
\item
  Easter depends on the 19-year Metonic (lunar) cycle, the 400-year
  Gregorian solar cycle, a 25-year correction cycle used in the epacts,
  and the 30 possible epact values.
\item
  Multiply them:
  \(19\times 400\times 25\times 30 = \mathbf{5{,}700{,}000}\) years.
  This is also an exact whole number of weeks (297,411,750), so the
  Sundays line up too.
\end{itemize}

Equivalently: the cycle is 70,499,183 lunations or 2,081,882,250 days;
only after that do Gregorian Easter dates recur in exactly the same
order.

(For comparison, in the \textbf{Julian} calendar the dates repeat every
532 years = 19×28.)

\subsection{16 {[}31{]}}\label{section-135}

For a fixed \(n\), write

\[
r_n(x)=\frac{\left\lfloor 10^n x+\tfrac12\right\rfloor}{10^n}.
\]

We need

\[
S_n=\sum_{m\ge 1} r_n(1/m).
\]

If we define the scaled numerator

\[
t(m)=\left\lfloor \frac{10^n}{m}+\tfrac12\right\rfloor,
\]

then \(r_n(1/m)=t(m)/10^n\) and

\[
10^n\,S_n=\sum_{m\ge 1} t(m).
\]

The key is that many consecutive \(m\)'s give the \textbf{same} \(t\).
If \(t=\left\lfloor 10^n/m+\tfrac12\right\rfloor\) for some \(m\), then
the last \(m\) that still yields this \(t\) is

\[
m_h=\Big\lfloor \frac{2\cdot 10^n}{2t-1}\Big\rfloor .
\]

So we can add \((m_h-m+1)\,t\) all at once and jump to the next block.

Using the grouping algorithm:

\begin{itemize}
\tightlist
\item
  \(S_2 = \mathbf{6.16}\)
\item
  \(S_3 = \mathbf{8.449}\)
\item
  \(S_4 = \mathbf{10.7509}\)
\item
  \(S_5 = \mathbf{13.05363}\)
\end{itemize}

(For reference, \(S_1=3.9\) as the book notes.)

\textbf{MIXAL program (computes and prints \(S_n\) for \(n=2,3,4,5\))}

The code keeps the sum scaled by \(10^n\) (an integer), then prints the
integer part and the \(n\) fractional digits.

\begin{Shaded}
\begin{Highlighting}[]
\NormalTok{* EX 1.3.2 — Exercise 16}
\NormalTok{* Prints S\_n for n = 2,3,4,5 using grouping of equal rounded terms.}

\NormalTok{        ORIG 3000}
\NormalTok{PRN     EQU 18}
\NormalTok{LINE    EQU 4000              * 24{-}word printer line}

\NormalTok{* one{-}byte character helpers}
\NormalTok{SPCW    ALF 00000}
\NormalTok{DOTW    ALF .0000}
\NormalTok{Z0W     ALF 00000}
\NormalTok{SPC     CON 0                 * set from SPCW(1:1)}
\NormalTok{DOT     CON 0                 * \textquotesingle{}.\textquotesingle{}}
\NormalTok{D0      CON 0                 * \textquotesingle{}0\textquotesingle{}}

\NormalTok{* variables used by the summation}
\NormalTok{POW     CON 0                 * 10\^{}n}
\NormalTok{TWOPW   CON 0                 * 2*10\^{}n}
\NormalTok{SUMS    CON 0                 * scaled sum = 10\^{}n * S\_n}
\NormalTok{MH      CON 1                 * high m in current block}
\NormalTok{ME      CON 0                 * next m to try (MH+1)}
\NormalTok{TVAL    CON 0                 * t = round(10\^{}n / m\_e)}
\NormalTok{DEN     CON 0                 * 2*t {-} 1}
\NormalTok{MHNEW   CON 0                 * new block high}

\NormalTok{N       CON 0                 * current n (2..5)}
\NormalTok{R       CON 0                 * remainder for printing}
\NormalTok{Q       CON 0                 * integer part for printing}
\NormalTok{CNT     CON 0                 * digit counter for printing}

\NormalTok{* {-}{-}{-}{-}{-}{-}{-}{-}{-}{-}{-}{-}{-}{-}{-}{-}{-}{-}{-}{-}{-}{-}{-}{-}{-}{-}{-}{-}{-}{-}{-}{-}{-}{-}{-}{-} main driver}
\NormalTok{MAIN    LDA  SPCW(1:1);  STA SPC}
\NormalTok{        LDA  DOTW(1:1);  STA DOT}
\NormalTok{        LDA  Z0W(1:1);   STA D0}

\NormalTok{        ENTA 2           ; N ← 2}
\NormalTok{NLOOP   STA  N}
\NormalTok{* POW ← 10\^{}N}
\NormalTok{        ENTA 1}
\NormalTok{        STZ  POW}
\NormalTok{        STZ  TWOPW}
\NormalTok{        ENTX 0}
\NormalTok{        LD2  N}
\NormalTok{POWLP   MUL  =10=}
\NormalTok{        DEC2 1}
\NormalTok{        J2P  POWLP}
\NormalTok{        STA  POW}
\NormalTok{        ADD  POW          ; TWOPW = 2*POW}
\NormalTok{        STA  TWOPW}

\NormalTok{* initialize summation: MH=1, SUMS=POW  (since r\_n(1)=1)}
\NormalTok{        LDA  POW}
\NormalTok{        STA  SUMS}
\NormalTok{        ENTA 1}
\NormalTok{        STA  MH}

\NormalTok{* {-}{-}{-}{-}{-}{-}{-}{-}{-} grouping loop: B,C,D from the note {-}{-}{-}{-}{-}{-}{-}{-}{-}{-}}
\NormalTok{GROUP   LDA  MH}
\NormalTok{        ADD  =1=}
\NormalTok{        STA  ME                     ; me ← mh + 1}

\NormalTok{* t = floor( (2*POW + ME) / (2*ME) )  (round(POW/ME))}
\NormalTok{        LDA  TWOPW}
\NormalTok{        ADD  ME}
\NormalTok{        ENTX 0}
\NormalTok{        STA  TMP                    ; numerator}
\NormalTok{        LDA  ME}
\NormalTok{        MUL  =2=}
\NormalTok{        STA  TMP2                   ; denominator}
\NormalTok{        LDA  TMP}
\NormalTok{        DIV  TMP2}
\NormalTok{        STA  TVAL                   ; t = A}
\NormalTok{        JAZ  DONE                   ; if t=0, finished}

\NormalTok{* mh\_new = floor( 2*POW / (2*t {-} 1) )}
\NormalTok{        LDA  TVAL}
\NormalTok{        MUL  =2=}
\NormalTok{        SUB  =1=}
\NormalTok{        STA  DEN}
\NormalTok{        LDA  TWOPW}
\NormalTok{        ENTX 0}
\NormalTok{        DIV  DEN}
\NormalTok{        STA  MHNEW}

\NormalTok{* SUMS += (mh\_new {-} me + 1) * t}
\NormalTok{        LDA  MHNEW}
\NormalTok{        SUB  ME}
\NormalTok{        ADD  =1=}
\NormalTok{        MUL  TVAL}
\NormalTok{        ADD  SUMS}
\NormalTok{        STA  SUMS}
\NormalTok{        LDA  MHNEW}
\NormalTok{        STA  MH}
\NormalTok{        JMP  GROUP}

\NormalTok{DONE}
\NormalTok{* {-}{-}{-}{-}{-}{-}{-}{-}{-} divide by 10\^{}n: Q = SUMS / POW , R = SUMS \% POW}
\NormalTok{        LDA  SUMS}
\NormalTok{        ENTX 0}
\NormalTok{        DIV  POW}
\NormalTok{        STA  Q}
\NormalTok{        STX  R}

\NormalTok{* print line: "S\textless{}digit\textgreater{} = \textless{}Q\textgreater{}.\textless{}n digits from R\textgreater{}"}
\NormalTok{        ENT4 LINE}
\NormalTok{        ENT2 24}
\NormalTok{CLRLP   STZ  0,4}
\NormalTok{        INC4 1}
\NormalTok{        DEC2 1}
\NormalTok{        J2P  CLRLP}

\NormalTok{        LDA  \textquotesingle{}S\textquotesingle{}         ; easiest way is an ALF literal (or from a table)}
\NormalTok{        STA  LINE(1:1)}
\NormalTok{        LDA  N}
\NormalTok{        ADD  D0}
\NormalTok{        STA  LINE(2:2)}
\NormalTok{        LDA  SPC}
\NormalTok{        STA  LINE(3:3)}
\NormalTok{        LDA  \textquotesingle{}=\textquotesingle{}}
\NormalTok{        STA  LINE(4:4)}
\NormalTok{        LDA  SPC}
\NormalTok{        STA  LINE(5:5)}

\NormalTok{* integer part Q → write right{-}to{-}left ending at byte 10}
\NormalTok{        LDA  Q}
\NormalTok{        ENTA 10}
\NormalTok{        STA  POSI                 ; rightmost byte to fill}
\NormalTok{        JAZ  PUT0I}
\NormalTok{IPLOOP  ENTX 0}
\NormalTok{        DIV  =10=}
\NormalTok{        STX  DIG}
\NormalTok{        LDA  DIG}
\NormalTok{        ADD  D0}
\NormalTok{        STA  LINE,POSI(0:0)       ; store one byte at (word/byte) via helper}
\NormalTok{        LDA  POSI}
\NormalTok{        SUB  =1=}
\NormalTok{        STA  POSI}
\NormalTok{        LDA  A                    ; quotient}
\NormalTok{        JAZ  DOTPUT}
\NormalTok{        JMP  IPLOOP}
\NormalTok{PUT0I   LDA  D0; STA LINE(10:10)}

\NormalTok{DOTPUT  LDA  DOT}
\NormalTok{        STA  LINE(11:11)}

\NormalTok{* fractional part: n digits from R, right{-}to{-}left at bytes 10+n .. 12}
\NormalTok{        LDA  N}
\NormalTok{        ADD  =11=}
\NormalTok{        STA  POSF                 ; last fractional byte position}
\NormalTok{        LDA  N}
\NormalTok{        STA  CNT}
\NormalTok{FRLP    LDA  R}
\NormalTok{        ENTX 0}
\NormalTok{        DIV  =10=}
\NormalTok{        STX  DIG}
\NormalTok{        STA  R}
\NormalTok{        LDA  DIG}
\NormalTok{        ADD  D0}
\NormalTok{        STA  LINE,POSF(0:0)}
\NormalTok{        LDA  POSF}
\NormalTok{        SUB  =1=}
\NormalTok{        STA  POSF}
\NormalTok{        DEC  CNT}
\NormalTok{        J2P  FRLP}

\NormalTok{        OUT  LINE(PRN)}

\NormalTok{* next n}
\NormalTok{        LDA  N}
\NormalTok{        ADD  =1=}
\NormalTok{        CMPA =6=}
\NormalTok{        JL   NLOOP}

\NormalTok{        HLT}

\NormalTok{* {-}{-}{-}{-}{-} tiny “store one byte at absolute position” helper (POS in 1..24)}
\NormalTok{POSI    CON 0}
\NormalTok{POSF    CON 0}
\NormalTok{DIG     CON 0}
\NormalTok{TMP     CON 0}
\NormalTok{TMP2    CON 0}

\NormalTok{        END  MAIN}
\end{Highlighting}
\end{Shaded}

\begin{quote}
If you don't want to wrestle with per-byte stores, you can simplify the
printing by using the \texttt{CHAR} instruction twice (on \texttt{Q} and
on \texttt{R}) and then copying the rightmost digits into the line, but
the grouped summation core above is the key idea that makes the program
fast.
\end{quote}

\textbf{Sanity check}

Running the algorithm gives exactly:

\begin{itemize}
\tightlist
\item
  \(S_2=6.16\)
\item
  \(S_3=8.449\)
\item
  \(S_4=10.7509\)
\item
  \(S_5=13.05363\)
\end{itemize}

These agree with independent high-precision calculations.

\subsection{17 {[}HM30{]}}\label{hm30-9}

Yes---the formula is \textbf{true}:

\[
\boxed{\displaystyle \lim_{n\to\infty}\big(S_{n+1}-S_n\big)=\ln 10 } .
\]

\textbf{Proof sketch}

Let \(r_n(x)=\frac{\lfloor 10^{n}x+\tfrac12\rfloor}{10^{n}}\) and

\[
S_n=\sum_{m\ge1} r_n\!\left(\frac1m\right).
\]

Since \(r_n(1/m)=0\) for \(m>2\cdot 10^n\), put \(A:=10^n\) and write

\[
S_n=\frac1A\sum_{m\le 2A}\Big\lfloor \frac{A}{m}+\tfrac12\Big\rfloor
=\frac{1}{A}\,F(A).
\]

Split

\[
F(A)=\sum_{m\le 2A}\Big\lfloor \frac{A}{m}\Big\rfloor
+\sum_{m\le 2A}\mathbf 1_{\{\{A/m\}\ge \tfrac12\}}
=:H(A)+T(A),
\]

where \(\{x\}\) is the fractional part and \(\mathbf 1\) an indicator.

\begin{enumerate}
\def\labelenumi{\arabic{enumi}.}
\tightlist
\item
  A classical estimate (Dirichlet's divisor summatory formula) gives
\end{enumerate}

\[
H(A)=A\log A+(2\gamma-1)A+O(\sqrt A)
\]

(\(\gamma\) is Euler's constant).

\begin{enumerate}
\def\labelenumi{\arabic{enumi}.}
\setcounter{enumi}{1}
\tightlist
\item
  For the second term, observe that \(\{A/m\}\ge \tfrac12\) iff there
  exists an integer \(k\ge0\) with \(k+\tfrac12\le A/m<k+1\),
  equivalently \(A/(k+1)< m\le A/(k+\tfrac12).\) Hence
\end{enumerate}

\[
T(A)=\sum_{k\ge0}\Big(\Big\lfloor\frac{A}{k+\tfrac12}\Big\rfloor
-\Big\lfloor\frac{A}{k+1}\Big\rfloor\Big)
= A\!\!\sum_{k\ge0}\!\!\Big(\frac{1}{k+\tfrac12}-\frac{1}{k+1}\Big)+O(\sqrt A)
= A\log 2+O(\sqrt A).
\]

Therefore

\[
F(A)=A\log(2A)+(2\gamma-1)A+O(\sqrt A),
\qquad
S_n=\log(2A)+(2\gamma-1)+O(A^{-1/2}).
\]

With \(A=10^n\),

\[
S_{n+1}-S_n
=\log\!\big(2\cdot 10^{n+1}\big)-\log\!\big(2\cdot 10^{n}\big)
+O\!\big(10^{-n/2}\big)
=\log 10+o(1),
\]

so the limit is \(\log 10\).

\textbf{Numerical check}

\[
\begin{aligned}
S_2&=6.16,\\
S_3&=8.449,& S_3-S_2&=2.289,\\
S_4&=10.7509,& S_4-S_3&=2.3019,\\
S_5&=13.05363,& S_5-S_4&=2.30273,
\end{aligned}
\]

which approach \(\ln 10 \approx 2.302585093\).

Here are few more:

\begin{itemize}
\tightlist
\item
  \(S_6 = \mathbf{15.356262}\)
\item
  \(S_7 = \mathbf{17.6588276}\)
\item
  \(S_8 = \mathbf{19.96140690}\)
\item
  \(S_9 = \mathbf{22.263991779}\)
\item
  \(S_{10} = \mathbf{24.5665766353}\)
\end{itemize}

(Computed by grouping equal terms: with \(A=10^n\), add
\((m_h-m_e+1)\,t\) where \(t=\left\lfloor A/m_e+\tfrac12\right\rfloor\)
and \(m_h=\left\lfloor\frac{2A}{2t-1}\right\rfloor\), until \(t=0\).)

\subsection{18 {[}25{]}}\label{section-136}

Here's a compact MIXAL subroutin that builds the Farey series of order
\(n\) using the standard three-term recurrence

\[
t=\Big\lfloor\frac{y_k+n}{y_{k+1}}\Big\rfloor,\qquad
x_{k+2}=t\,x_{k+1}-x_k,\quad y_{k+2}=t\,y_{k+1}-y_k .
\]

It stores each pair \(x_k,y_k\) to memory words \textbf{XBASE+k} and
\textbf{YBASE+k}. Call it with \texttt{N}, \texttt{XBASE}, and
\texttt{YBASE} preset by the caller; on return the table 0/1, \ldots,
1/1 has been written.

\begin{Shaded}
\begin{Highlighting}[]
\NormalTok{* FAREY(n): write Farey series of order n}
\NormalTok{* Caller must set:  N = order n  (n ≥ 1)}
\NormalTok{*                   XBASE = base address for x\_k table}
\NormalTok{*                   YBASE = base address for y\_k table}
\NormalTok{* Writes: (x\_k,y\_k) to (XBASE+k, YBASE+k), starting with 0/1 and ending with 1/1.}

\NormalTok{        ORIG 3000}

\NormalTok{FAREY   STJ   RET          * subroutine linkage}

\NormalTok{        LD4   XBASE        * rI4 ← address of x\_0 (XBASE)}
\NormalTok{        LD5   YBASE        * rI5 ← address of y\_0 (YBASE)}

\NormalTok{* x0=0, y0=1;  x1=1, y1=n}
\NormalTok{        STZ   0,4          * x0}
\NormalTok{        ENTA  1}
\NormalTok{        STA   0,5          * y0}
\NormalTok{        STA   1,4          * x1}
\NormalTok{        LDA   N}
\NormalTok{        STA   1,5          * y1 = n}

\NormalTok{* Main loop: stop once (x\_\{k+1\},y\_\{k+1\}) = (1,1)}
\NormalTok{LOOP    LDA   1,4}
\NormalTok{        CMPA  =1=}
\NormalTok{        JNE   CONT}
\NormalTok{        LDA   1,5}
\NormalTok{        CMPA  =1=}
\NormalTok{        JE    DONE}

\NormalTok{CONT    LDA   0,5          * A ← y\_k}
\NormalTok{        ADD   N            * A ← y\_k + n}
\NormalTok{        DIV   1,5          * A ← ⌊(y\_k + n)/y\_\{k+1\}⌋  (t), X ← remainder}
\NormalTok{        STA   T}

\NormalTok{* x\_\{k+2\} = t*x\_\{k+1\} − x\_k}
\NormalTok{        LDA   T}
\NormalTok{        MUL   1,4}
\NormalTok{        SUB   0,4}
\NormalTok{        STA   2,4}

\NormalTok{* y\_\{k+2\} = t*y\_\{k+1\} − y\_k}
\NormalTok{        LDA   T}
\NormalTok{        MUL   1,5}
\NormalTok{        SUB   0,5}
\NormalTok{        STA   2,5}

\NormalTok{        INC4  1            * advance k: rI4→XBASE+(k+1)}
\NormalTok{        INC5  1            *               rI5→YBASE+(k+1)}
\NormalTok{        JMP   LOOP}

\NormalTok{DONE    JMP   RET}

\NormalTok{RET     JMP   *            * return here}

\NormalTok{* Inputs (set by caller):}
\NormalTok{N       CON   0            * order n}
\NormalTok{XBASE   CON   0            * base address for x\_k}
\NormalTok{YBASE   CON   0            * base address for y\_k}

\NormalTok{* Scratch:}
\NormalTok{T       CON   0}

\NormalTok{        END   FAREY}
\end{Highlighting}
\end{Shaded}

\textbf{How to call}

\begin{Shaded}
\begin{Highlighting}[]
\NormalTok{        ORIG 5000}
\NormalTok{XTAB    CON   0  (reserve ≥ about 3n\^{}2/π\^{}2 words)}
\NormalTok{        ORIG 7000}
\NormalTok{YTAB    CON   0}

\NormalTok{        ORIG 2000}
\NormalTok{        LDA  =7=           ; n = 7}
\NormalTok{        STA  N}
\NormalTok{        LDA  =5000=        ; XBASE = XTAB}
\NormalTok{        STA  XBASE}
\NormalTok{        LDA  =7000=        ; YBASE = YTAB}
\NormalTok{        STA  YBASE}
\NormalTok{        JMP  FAREY}
\NormalTok{        HLT}
\end{Highlighting}
\end{Shaded}

This routine writes \(0/1,\,1/n,\,\ldots,\,1/1\) to consecutive words
beginning at \texttt{XBASE}/\texttt{YBASE}. It uses only \texttt{rI4}
and \texttt{rI5} as moving pointers (no table length is stored, but you
can easily keep a count by incrementing a separate cell each time
through the loop if you need it).

\subsection{19 {[}M30{]}}\label{m30-12}

\textbf{(a) \(x_{k+1}y_k-x_k y_{k+1}=1\) for all \(k\)}

Let

\[
t_k=\Big\lfloor\frac{y_k+n}{y_{k+1}}\Big\rfloor,\qquad
x_{k+2}=t_kx_{k+1}-x_k,\quad y_{k+2}=t_k y_{k+1}-y_k,
\]

with \(x_0=0,y_0=1\) and \(x_1=1,y_1=n\).

Set \(\Delta_k=x_{k+1}y_k-x_k y_{k+1}\). For \(k=0\):
\(\Delta_0=1\cdot1-0\cdot n=1\).

Using the recurrence,

\[
\begin{aligned}
\Delta_{k+1}
&=x_{k+2}y_{k+1}-x_{k+1}y_{k+2}\\
&=(t_kx_{k+1}-x_k)y_{k+1}-x_{k+1}(t_k y_{k+1}-y_k)\\
&=x_{k+1}y_k-x_k y_{k+1}=\Delta_k .
\end{aligned}
\]

Hence \(\Delta_k\) is constant, and since \(\Delta_0=1\) we have

\[
\boxed{x_{k+1}y_k-x_k y_{k+1}=1\quad\text{for all }k\ge0.}
\]

Consequences: \(\gcd(x_k,y_k)=1\) (any common divisor would divide 1),
so all fractions \(x_k/y_k\) are reduced.

\textbf{(b) The fractions \(x_k/y_k\) are the Farey series of order
\(n\)}

From the definition of \(t_k\),

\[
t_k y_{k+1}\le y_k+n<(t_k+1)y_{k+1}.
\]

Subtract \(y_k\) to get

\[
y_{k+2}=t_k y_{k+1}-y_k\le n
\quad\text{and}\quad
n<y_{k+2}+y_{k+1}.
\]

Thus every denominator satisfies \(y_k\le n\), and successive
denominators satisfy \(y_k+y_{k+1}>n\).

By part (a) the consecutive fractions satisfy

\[
x_{k+1}y_k-x_k y_{k+1}=1.
\]

This identity implies (i) each is reduced, and (ii) any fraction
strictly between \(x_k/y_k\) and \(x_{k+1}/y_{k+1}\) has denominator at
least \(y_k+y_{k+1}\). Since \(y_k+y_{k+1}>n\), no fraction with
denominator \(\le n\) lies between them. Moreover one checks

\[
\frac{x_k}{y_k}<\frac{x_{k+2}}{y_{k+2}}<\frac{x_{k+1}}{y_{k+1}}
\]

(using \(x_{k+1}y_k-x_k y_{k+1}=1>0\)), so the sequence is strictly
increasing.

We start at \(x_0/y_0=0/1\) and end when \((x_{k+1},y_{k+1})=(1,1)\),
hence the whole list runs

\[
0/1,\;\dots,\;1/1,
\]

with all terms reduced, denominators \(\le n\), in increasing order, and
no additional reduced fraction of order \(\le n\) between neighbors.
Therefore it is exactly the Farey series of order \(n\).

\subsection{20 {[}33{]}}\label{section-137}

Below is a complete (but compact) MIXAL program that drives the four
lamps via the X-register bytes, using the overflow toggle as the request
switch. It implements exactly the timing rules in the statement:

\begin{itemize}
\tightlist
\item
  Del Mar green ≥ 30 s, amber 8 s
\item
  Berkeley green 20 s, amber 5 s
\item
  Pedestrian ``DON'T WALK'' flashes for 12 s before green → amber (½ s
  on, ½ s off ×8, then 4 s steady), and stays on during amber+red
\item
  If the request is tripped while Berkeley is green, nothing special
  happens (they pass on that cycle); if it is tripped at any other time
  it schedules a Berkeley-green cycle after Del Mar has passed.
\item
  rX is changed only by \texttt{LDX}/\texttt{INCX}; the level is taken
  exactly at the end of those instructions.
\end{itemize}

The code is written as a small finite-state machine. Times are generated
by a calibrated half-second delay. (Pick \texttt{HS} to make the body of
label \texttt{DELAY} consume \textbf{0.5 s}; one MIX time unit is 10 μs,
so you want \texttt{HS\ =\ 50,000\ /\ c}, where \texttt{c} is the exact
number of time-units per loop-iteration of \texttt{DELAY} on your MIX
implementation.)

\textbf{Encodings put into rX}

We precompute the words we \texttt{LDX} to set all four lamp bytes at
once. Digits are base-64 bytes in positions (2:5):

\begin{itemize}
\tightlist
\item
  \texttt{X\_DM\_GREEN\ \ \ =\ +\ {[}\ b2=1,\ b3=3,\ b4=1,\ b5=2\ {]}\ =\ +274498}
\item
  \texttt{X\_DM\_DWON\ \ \ \ =\ +\ {[}\ 1,3,2,2\ {]}\ =\ +274562} (Del
  Mar DON'T-WALK on)
\item
  \texttt{X\_DM\_DWOFF\ \ \ =\ +\ {[}\ 1,3,0,2\ {]}\ =\ +274434} (Del
  Mar ped off)
\item
  \texttt{X\_DM\_AMBER\ \ \ =\ +\ {[}\ 2,3,2,2\ {]}\ =\ +536706}
\item
  \texttt{X\_BK\_GREEN\ \ \ =\ +\ {[}\ 3,1,2,1\ {]}\ =\ +790657}
\item
  \texttt{X\_BK\_DWON\ \ \ \ =\ +\ {[}\ 3,1,2,2\ {]}\ =\ +790658}
  (Berkeley DON'T-WALK on)
\item
  \texttt{X\_BK\_DWOFF\ \ \ =\ +\ {[}\ 3,1,2,0\ {]}\ =\ +790656}
  (Berkeley ped off)
\item
  \texttt{X\_BK\_AMBER\ \ \ =\ +\ {[}\ 3,2,2,2\ {]}\ =\ +794754}
\end{itemize}

(Recall: 0=off/---, traffic: 1 green, 2 amber, 3 red; ped: 1 WALK, 2
DON'T WALK.)

\begin{Shaded}
\begin{Highlighting}[]
\NormalTok{* EX 1.3.2 — Exercise 20: two{-}way traffic lights with request switch (overflow)}
\NormalTok{* rX(2) Del Mar traffic; rX(3) Berkeley traffic; rX(4) Del Mar ped;}
\NormalTok{* rX(5) Berkeley ped.}

\NormalTok{        ORIG 3000}

\NormalTok{* {-}{-}{-}{-}{-}{-}{-}{-} constants: precomputed X words for each lamp pattern {-}{-}{-}{-}{-}{-}{-}{-}{-}{-}}
\NormalTok{X\_DM\_GREEN  CON +274498}
\NormalTok{X\_DM\_DWON   CON +274562}
\NormalTok{X\_DM\_DWOFF  CON +274434}
\NormalTok{X\_DM\_AMBER  CON +536706}
\NormalTok{X\_BK\_GREEN  CON +790657}
\NormalTok{X\_BK\_DWON   CON +790658}
\NormalTok{X\_BK\_DWOFF  CON +790656}
\NormalTok{X\_BK\_AMBER  CON +794754}

\NormalTok{HS          CON  12345     * CALIBRATE: iterations that make DELAY = 0.5 s}
\NormalTok{ZERO        CON  0}
\NormalTok{ONE         CON  1}
\NormalTok{EIGHT       CON  8}
\NormalTok{SIXTY       CON  60}
\NormalTok{SIXTEEN     CON  16}
\NormalTok{TEN         CON 10}

\NormalTok{* {-}{-}{-}{-}{-}{-}{-}{-} state / scratch {-}{-}{-}{-}{-}{-}{-}{-}{-}{-}{-}{-}{-}{-}{-}{-}{-}{-}{-}{-}{-}{-}{-}{-}{-}{-}{-}{-}{-}{-}{-}{-}{-}{-}{-}{-}{-}{-}{-}{-}{-}{-}{-}{-}{-}{-}{-}{-}}
\NormalTok{REQP        CON 0          * 1 if a Berkeley request is pending}
\NormalTok{ET          CON 0          * elapsed half{-}seconds in current Del Mar green}
\NormalTok{ALLOW       CON 0          * 1 if an overflow should set REQP (not during BK green)}
\NormalTok{CNT         CON 0          * generic small loop counter}
\NormalTok{RET         CON 0}

\NormalTok{* ============================ main controller ============================}

\NormalTok{START   LDX  X\_DM\_GREEN    * default: Del Mar stays green forever}
\NormalTok{        STZ  REQP}
\NormalTok{        STZ  ET}

\NormalTok{DMGREEN STZ  ALLOW         * we *do* accept requests in DM green → ALLOW=1}
\NormalTok{        INC  ALLOW ONE}

\NormalTok{DMLP    JSJ  HALF          * wait 0.5 s, also sample overflow if ALLOW=1}
\NormalTok{        INC  ET  ONE       * increase elapsed time in half{-}seconds}

\NormalTok{* Start termination sequence when: request pending AND green≥30 s}
\NormalTok{        LDA  ET}
\NormalTok{        CMPA SIXTY}
\NormalTok{        JL   DMLP          * too soon to end}
\NormalTok{        LDA  REQP}
\NormalTok{        JZ   DMLP          * no request → keep it green}

\NormalTok{* {-}{-}{-}{-}{-} pre{-}amber for Del Mar (12 s): flash (½ on, ½ off) ×8; then 4 s steady}
\NormalTok{        LD1  EIGHT}
\NormalTok{DMFLSH  LDX  X\_DM\_DWON     * DON’T WALK on}
\NormalTok{        JSJ  HALF}
\NormalTok{        LDX  X\_DM\_DWOFF    * DON’T WALK off}
\NormalTok{        JSJ  HALF}
\NormalTok{        DEC1 1}
\NormalTok{        J1P  DMFLSH}

\NormalTok{        LDX  X\_DM\_DWON     * 4 s steady before amber}
\NormalTok{        LD1  EIGHT}
\NormalTok{DW4     JSJ  HALF}
\NormalTok{        DEC1 1}
\NormalTok{        J1P  DW4}

\NormalTok{* {-}{-}{-}{-}{-} Del Mar amber 8 s}
\NormalTok{        LDX  X\_DM\_AMBER}
\NormalTok{        LD1  SIXTEEN       * 16 × ½ s}
\NormalTok{DMAMB   JSJ  HALF}
\NormalTok{        DEC1 1}
\NormalTok{        J1P  DMAMB}

\NormalTok{* {-}{-}{-}{-}{-} Berkeley green 20 s (service the request now)}
\NormalTok{        STZ  REQP          * request has been satisfied}
\NormalTok{        LDX  X\_BK\_GREEN}

\NormalTok{* First 8 s WALK, then 8 s flash, then 4 s steady DON’T WALK}
\NormalTok{        STZ  ALLOW         * during Berkeley green, ignore “new” requests}
\NormalTok{        LD1  SIXTEEN}
\NormalTok{        SUB  ONE,1         * SIXTEEN{-}? we need 8 s = 16×½s {-}\textgreater{} use CNT=16 and break after 16}
\NormalTok{BK8W    JSJ  HALF}
\NormalTok{        DEC1 1}
\NormalTok{        J1P  BK8W}

\NormalTok{        LD1  EIGHT}
\NormalTok{BKFL    LDX  X\_BK\_DWON}
\NormalTok{        JSJ  HALF}
\NormalTok{        LDX  X\_BK\_DWOFF}
\NormalTok{        JSJ  HALF}
\NormalTok{        DEC1 1}
\NormalTok{        J1P  BKFL}

\NormalTok{        LDX  X\_BK\_DWON}
\NormalTok{        LD1  EIGHT}
\NormalTok{BKDW4   JSJ  HALF}
\NormalTok{        DEC1 1}
\NormalTok{        J1P  BKDW4}

\NormalTok{* {-}{-}{-}{-}{-} Berkeley amber 5 s}
\NormalTok{        LDX  X\_BK\_AMBER}
\NormalTok{        INC  ALLOW ONE     * during BK amber, new requests should be queued}
\NormalTok{        LD1  TEN           * 5 s = 10 × ½ s}
\NormalTok{BKAMB   JSJ  HALF}
\NormalTok{        DEC1 1}
\NormalTok{        J1P  BKAMB}

\NormalTok{* {-}{-}{-}{-}{-} back to Del Mar green; cycle restarts}
\NormalTok{        LDX  X\_DM\_GREEN}
\NormalTok{        STZ  ET}
\NormalTok{        JMP  DMGREEN}

\NormalTok{* ========================= ½{-}second delay + request sampler =============}

\NormalTok{* HALF: Busy{-}wait for exactly 0.5 s (per HS calibration) and}
\NormalTok{*       if ALLOW=1 and overflow toggle is ON, set REQP←1.}
\NormalTok{HALF    STJ  RET}
\NormalTok{        LD2  HS}
\NormalTok{H1      DEC2 1}
\NormalTok{        J2P  H1}
\NormalTok{* Check overflow request switch}
\NormalTok{        LDA  ALLOW}
\NormalTok{        JZ   RETU}
\NormalTok{        JOV  SETR           * Jumps iff switch has been pressed; clears OV}
\NormalTok{        JMP  RETU}
\NormalTok{SETR    INC  REQP ONE}
\NormalTok{RETU    JMP  RET}

\NormalTok{        END  START}
\end{Highlighting}
\end{Shaded}

\textbf{How it works}

\begin{itemize}
\tightlist
\item
  Default is Del Mar green + WALK, Berkeley red + DON'T WALK.
\item
  The loop labelled \texttt{DMLP} keeps Del Mar green and samples the
  overflow every ½ s. A request is armed only when Del Mar is green or
  amber, or when Berkeley is amber (i.e., whenever Berkeley is not
  green). This is managed by the \texttt{ALLOW} flag.
\item
  Once (i) a request is pending and (ii) Del Mar has been green for at
  least 30 s, we launch the 12-s pre-amber (8 s flash + 4 s steady DON'T
  WALK), then 8 s amber, then give Berkeley 20 s green (with its own
  flash/steady sequence), 5 s amber, and return to Del Mar green.
\item
  During Berkeley green we purposely ignore new presses; people cross on
  that cycle. During other phases we queue a press (by setting
  \texttt{REQP}), guaranteeing that the next Del Mar green will
  terminate at the earliest lawful moment (≥30 s with a full 12-s
  pre-amber).
\end{itemize}

\begin{quote}
Calibrate \texttt{HS} once for your MIX timing so that the loop under
label \texttt{DELAY/H1} takes exactly half a second. After that, all
required durations are exact by construction.
\end{quote}

\subsection{21 {[}28{]}}\label{section-138}

Here's a complete MIXAL program that builds the 23×23 magic square by
the classical ``Siamese'' (Ibn al-Haytham) rule described in the
problem, stores it in memory, and prints all 23 rows on the line
printer.

\begin{itemize}
\tightlist
\item
  We index rows/columns 1..23.
\item
  Cell \((i,j)\) lives at address \texttt{BASE\ +\ 23*(i-1)\ +\ j}.
\item
  Start with 1 just below the center: \((i,j)=(13,12)\).
\item
  Step down-right; when that destination is already filled, drop two
  rows (same column).
\item
  Printing puts each value as three right-aligned digits followed by a
  blank, 23 per line.
\end{itemize}

\begin{Shaded}
\begin{Highlighting}[]
\NormalTok{* TAOCP v1 §1.3.2 — Ex.21  Magic square (order 23) with Siamese method}
\NormalTok{* rI1 ≡ i (row 1..23), rI2 ≡ j (col 1..23), rI3 ≡ addr pointer, rI4 ≡ line word ptr}

\NormalTok{        ORIG 3000}
\NormalTok{PRN     EQU 18}

\NormalTok{* {-}{-}{-}{-} character helpers {-}{-}{-}{-}{-}{-}{-}{-}{-}{-}{-}{-}{-}{-}{-}{-}{-}{-}{-}{-}{-}{-}{-}{-}{-}{-}{-}{-}{-}{-}{-}{-}{-}{-}{-}{-}{-}{-}{-}{-}{-}{-}{-}{-}{-}{-}{-}{-}{-}{-}{-}{-}{-}{-}{-}}
\NormalTok{SPCW    ALF 00000          * one byte blank at (1:1)}
\NormalTok{Z0W     ALF 00000          * \textquotesingle{}0\textquotesingle{} at (1:1)}
\NormalTok{SPC     CON 0}
\NormalTok{D0      CON 0}

\NormalTok{* {-}{-}{-}{-} constants {-}{-}{-}{-}{-}{-}{-}{-}{-}{-}{-}{-}{-}{-}{-}{-}{-}{-}{-}{-}{-}{-}{-}{-}{-}{-}{-}{-}{-}{-}{-}{-}{-}{-}{-}{-}{-}{-}{-}{-}{-}{-}{-}{-}{-}{-}{-}{-}{-}{-}{-}{-}{-}{-}{-}{-}{-}{-}{-}{-}{-}{-}{-}}
\NormalTok{N       CON 23}
\NormalTok{BASE    CON 2000           * table base: uses 2000..2528 (529 cells)}
\NormalTok{LINE    EQU 2600           * 24{-}word printer line buffer (safe, out of the way)}

\NormalTok{ONE     CON 1}
\NormalTok{TWO     CON 2}
\NormalTok{TW3     CON 23}
\NormalTok{N2      CON 529            * N\^{}2}

\NormalTok{* {-}{-}{-}{-} work cells {-}{-}{-}{-}{-}{-}{-}{-}{-}{-}{-}{-}{-}{-}{-}{-}{-}{-}{-}{-}{-}{-}{-}{-}{-}{-}{-}{-}{-}{-}{-}{-}{-}{-}{-}{-}{-}{-}{-}{-}{-}{-}{-}{-}{-}{-}{-}{-}{-}{-}{-}{-}{-}{-}{-}{-}{-}{-}{-}{-}{-}{-}}
\NormalTok{I       CON 0}
\NormalTok{J       CON 0}
\NormalTok{K       CON 0              * number to place}
\NormalTok{ADDR    CON 0}
\NormalTok{HUND    CON 0}
\NormalTok{TENS    CON 0}
\NormalTok{ONES    CON 0}
\NormalTok{RET     CON 0}

\NormalTok{* ====================== main ==================================================}
\NormalTok{START   LDA  SPCW(1:1);  STA SPC}
\NormalTok{        LDA  Z0W(1:1);   STA D0}

\NormalTok{* clear the 23×23 table to 0}
\NormalTok{        ENT1 1}
\NormalTok{CLRROW  ENT2 1}
\NormalTok{CLRJ    ENTA BASE}
\NormalTok{        ADD  I}
\NormalTok{        DEC  A 1}
\NormalTok{        MUL  N}
\NormalTok{        ADD  J}
\NormalTok{        STA  ADDR}
\NormalTok{        LD3  ADDR}
\NormalTok{        STZ  0,3}
\NormalTok{        INC2 1}
\NormalTok{        LDA  J}
\NormalTok{        CMPA N}
\NormalTok{        JL   CLRJ}
\NormalTok{        INC1 1}
\NormalTok{        LDA  I}
\NormalTok{        CMPA N}
\NormalTok{        JL   CLRROW}

\NormalTok{* initial position: just below the middle square (i=13, j=12)}
\NormalTok{        ENTA 13;  STA I}
\NormalTok{        ENTA 12;  STA J}
\NormalTok{        ENTA 1;   STA K}

\NormalTok{* {-}{-}{-}{-}{-}{-}{-}{-}{-}{-}{-}{-}{-} fill loop: K from 1..N\^{}2 {-}{-}{-}{-}{-}{-}{-}{-}{-}{-}{-}{-}{-}{-}{-}{-}{-}{-}{-}{-}{-}{-}{-}{-}{-}{-}{-}{-}{-}{-}{-}{-}{-}{-}{-}{-}{-}{-}}
\NormalTok{PLACE   JSJ  ATIJ          * rI3 points to (i,j)}
\NormalTok{        LDA  K}
\NormalTok{        STA  0,3           * put K}
\NormalTok{        LDA  K}
\NormalTok{        CMPA N2}
\NormalTok{        JE   PRINTALL      * done filling}

\NormalTok{* compute tentative (i+1, j+1) with wrap}
\NormalTok{        LDA  I; ADD ONE; CMPA N; JLE IOK; SUB N}
\NormalTok{IOK     STA  HUND          * reuse HUND as temp (i\textquotesingle{}) to save a cell}
\NormalTok{        LDA  J; ADD ONE; CMPA N; JLE JOK; SUB N}
\NormalTok{JOK     STA  TENS          * temp (j\textquotesingle{})}

\NormalTok{* empty? if yes, move there; else i ← i+2 mod N (j unchanged)}
\NormalTok{        ENTA BASE}
\NormalTok{        ADD  HUND}
\NormalTok{        DEC  A 1}
\NormalTok{        MUL  N}
\NormalTok{        ADD  TENS}
\NormalTok{        STA  ADDR}
\NormalTok{        LD3  ADDR}
\NormalTok{        LDA  0,3}
\NormalTok{        JZ   MOVEIJ}
\NormalTok{        LDA  I; ADD TWO; CMPA N; JLE IDONE; SUB N}
\NormalTok{IDONE   STA  I}
\NormalTok{        JMP  NEXTK}
\NormalTok{MOVEIJ  LDA  HUND; STA I}
\NormalTok{        LDA  TENS; STA J}

\NormalTok{NEXTK   LDA  K; ADD ONE; STA K}
\NormalTok{        JMP  PLACE}

\NormalTok{* {-}{-}{-}{-}{-}{-}{-}{-}{-}{-}{-}{-}{-}{-}{-}{-} compute rI3=address for (I,J) {-}{-}{-}{-}{-}{-}{-}{-}{-}{-}{-}{-}{-}{-}{-}{-}{-}{-}{-}{-}{-}{-}{-}{-}{-}{-}{-}{-}{-}{-}{-}}
\NormalTok{ATIJ    STJ  RET}
\NormalTok{        ENTA BASE}
\NormalTok{        ADD  I}
\NormalTok{        DEC  A 1}
\NormalTok{        MUL  N}
\NormalTok{        ADD  J}
\NormalTok{        STA  ADDR}
\NormalTok{        LD3  ADDR}
\NormalTok{        JMP  RET}

\NormalTok{* ======================= printing ============================================}

\NormalTok{PRINTALL}
\NormalTok{        ENT1 1}
\NormalTok{ROWLP   ENT2 1}
\NormalTok{        LD4  LINE}

\NormalTok{COLLP   * fetch value a(i,j)}
\NormalTok{        ENTA BASE}
\NormalTok{        ADD  I}
\NormalTok{        DEC  A 1}
\NormalTok{        MUL  N}
\NormalTok{        ADD  J}
\NormalTok{        STA  ADDR}
\NormalTok{        LD3  ADDR}
\NormalTok{        LDA  0,3}

\NormalTok{* extract 3 digits (hundreds, tens, ones) with two divisions}
\NormalTok{        ENTX 0}
\NormalTok{        DIV  =10=}
\NormalTok{        STX  ONES}
\NormalTok{        LDA  A}
\NormalTok{        ENTX 0}
\NormalTok{        DIV  =10=}
\NormalTok{        STX  TENS}
\NormalTok{        STA  HUND}

\NormalTok{* byte(1) ← hundreds or blank}
\NormalTok{        LDA  HUND}
\NormalTok{        JZ   HBL}
\NormalTok{        ADD  D0}
\NormalTok{        STA  0,4(1:1)}
\NormalTok{        JMP  TBYTE}
\NormalTok{HBL     LDA  SPC}
\NormalTok{        STA  0,4(1:1)}

\NormalTok{* byte(2) ← tens or blank (if hundreds and tens both 0)}
\NormalTok{TBYTE   LDA  HUND}
\NormalTok{        JNZ  TDIG}
\NormalTok{        LDA  TENS}
\NormalTok{        JZ   TBL}
\NormalTok{TDIG    LDA  TENS}
\NormalTok{        ADD  D0}
\NormalTok{        STA  0,4(2:2)}
\NormalTok{        JMP  OPUT}
\NormalTok{TBL     LDA  SPC}
\NormalTok{        STA  0,4(2:2)}

\NormalTok{* byte(3) ← ones (always a digit)}
\NormalTok{OPUT    LDA  ONES}
\NormalTok{        ADD  D0}
\NormalTok{        STA  0,4(3:3)}

\NormalTok{* byte(4) and (5) ← blanks (separator + tidy)}
\NormalTok{        LDA  SPC}
\NormalTok{        STA  0,4(4:4)}
\NormalTok{        STA  0,4(5:5)}

\NormalTok{        INC4 1}
\NormalTok{        INC2 1}
\NormalTok{        LDA  J}
\NormalTok{        CMPA N}
\NormalTok{        JL   COLLP}

\NormalTok{* print the 23{-}word line}
\NormalTok{        OUT  LINE(PRN)}

\NormalTok{        INC1 1}
\NormalTok{        LDA  I}
\NormalTok{        CMPA N}
\NormalTok{        JL   ROWLP}

\NormalTok{        HLT}
\NormalTok{        END  START}
\end{Highlighting}
\end{Shaded}

\textbf{Notes}

\begin{itemize}
\tightlist
\item
  Addressing. \texttt{ATIJ} computes the address for the current
  \((i,j)\) and sets \texttt{rI3} so we can \texttt{LDA/STA\ 0,3}.
\item
  Wrap-around. After \texttt{i+1} or \texttt{i+2} we compare to
  \texttt{N} and subtract \texttt{N} if needed (the plane-tiling rule).
\item
  Printing.** Each of the 23 output words on a line holds \texttt{ddd}
  (three digits right-aligned plus a blank separator; byte 5 blank too).
  That's simple and fits the printer's 24-word line.
\item
  Memory used: table \texttt{2000..2528} (529 words) and line buffer
  \texttt{2600..2623}.
\end{itemize}

Run this and you'll get the \(23\times23\) magic square, each row nicely
formatted on the printer.

\subsection{22 {[}31{]}}\label{section-139}

Here's a fast MIXAL solution of the Josephus problem that prints the
order-of-execution for \(n=24\), \(m=11\). It keeps a circular
doubly-linked list in two arrays \texttt{NEXT{[}i{]}} and
\texttt{PREV{[}i{]}} so that each kill is \(O(m)\) pointer hops, and
removal is \(O(1)\). It finally prints, for each man \(i=1,\dots,24\),
the number \(1..24\) telling when he is executed (the format is two
digits per man, zero-padded).

\begin{verbatim}
* EX 1.3.2 — Exercise 22 (Josephus): n = 24, m = 11
* Output: for i=1..24, the execution order of man i.

        ORIG 3000
PRN     EQU   18
N       EQU   24
M       EQU   11

LINE    EQU   3900                * 24-word printer line

* ------- helpers for printing -----------------------------------------------
Z0W     ALF  00000
SPCW    ALF  00000
D0      CON  0                    * '0' byte  (from Z0W(1:1))
SPC     CON  0                    * blank     (from SPCW(1:1))

* ------- tables (24 words each) ---------------------------------------------
NEXT    EQU   3600                * NEXT[i] : next alive after i
PREV    EQU   3624                * PREV[i] : previous alive before i
ORD     EQU   3648                * ORD[i]  : kill order of i (1..24)

* ------- scratch -------------------------------------------------------------
CUR     CON   0                   * current position in the circle
T       CON   0                   * current kill number (1..N)
TMP     CON   0
TEN     CON   10

RET     CON   0
ONES    CON   0
TENS    CON   0

* =========================  MAIN  ============================================

START   LDA  Z0W(1:1);  STA D0
        LDA  SPCW(1:1); STA SPC

* ---- build the circular doubly linked list for 1..N
        ENT1  1
INITLP  ENTA  0,1                 * A ← i
        ADD   =1=
        CMPA  =N+1=               * wrap next
        JL    OKN
        ENTA  1
OKN     STA   NEXT,1              * NEXT[i] ← (i%N)+1

        ENTA  0,1
        SUB   =1=
        JG    OKP
        ENTA  =N=
OKP     STA   PREV,1              * PREV[i] ← ((i-2+N)%N)+1

        INC1  1
        CMP1  =N+1=
        JL    INITLP

* ---- Josephus loop
        ENT1  1                   * CUR ← 1
        STA   CUR
        ENTA  1;  STA  T          * T ← 1  (first kill)

KILL    * advance m-1 steps: CUR ← NEXT^(m-1)(CUR)
        ENT2  M
        DEC2  1                   * rI2 ← m-1
ADVLP   LD3   CUR
        LDA   NEXT,3
        STA   CUR
        DEC2  1
        J2P   ADVLP

* mark execution order
        LD3   CUR
        LDA   T
        STA   ORD,3

* splice CUR out: NEXT[PREV[CUR]] ← NEXT[CUR];  PREV[NEXT[CUR]] ← PREV[CUR]
        LDA   PREV,3
        STA   TMP                 * TMP ← prev
        LDA   NEXT,3
        STA   CUR                 * CUR ← next (the one after the dead guy)

        LD4   TMP                 * write NEXT[prev] ← next
        STA   NEXT,4

        LD3   CUR                 * write PREV[next] ← prev
        LDA   TMP
        STA   PREV,3

* stop after N kills
        LDA   T
        CMPA  =N=
        JE    PRINT

        ADD   =1=
        STA   T
        JMP   KILL

* ============================ PRINT ==========================================

PRINT   LD4   LINE                * build a 24-word line
        ENT1  1
ROW     LD3   ORD,1               * value = ORD[i]
        LDA   0,3
        ENTX  0
        DIV   TEN
        STX   ONES
        STA   TENS

        LDA   TENS;  ADD D0;      STA 0,4(1:1)   * tens
        LDA   ONES;  ADD D0;      STA 0,4(2:2)   * ones
        LDA   SPC;                STA 0,4(3:3)   * space
        STA   0,4(4:4);           STA 0,4(5:5)   * padding blanks

        INC4  1
        INC1  1
        CMP1  =N+1=
        JL    ROW

        OUT   LINE(PRN)
        HLT

        END  START
\end{verbatim}

\textbf{What it prints (two digits per man, in order 1..24)}

For \(n=24\), \(m=11\) the order-of-execution is

\begin{verbatim}
15 12 22 08 16 11 23 21 03 05 01 17 10 07 24 19 20 18 09 14 04 02 13 06
\end{verbatim}

meaning: man 1 is 15th to die, man 2 is 12th, \ldots, man 24 is 6th.

\subsection{23 {[}37{]}}\label{section-140}

Below is a complete MIXAL program that reads a \(23\times 23\) matrix of
zeros and ones (one row per card on device 16, with columns 1--23
containing the digits `0'/`1'), converts it to a crossword diagram,
removing edge-connected black areas, numbering the white squares that
start an Across or Down word, and prints the final picture on the line
printer (device 18) exactly as specified: each puzzle cell is rendered
by 5 columns × 3 rows of printer characters.

I have kept the code heavily annotated and very close to the text of the
exercise, so you can change it easily (e.g., to other sizes).

What the program prints

\begin{itemize}
\tightlist
\item
  \textbf{White (unnumbered)} cell → top row \texttt{"\ \ \ \ +"},
  middle row \texttt{"\ \ \ \ +"}, bottom \texttt{"+++++"}.
\item
  \textbf{White (numbered)} cell, number \texttt{nn} (2 digits, leading
  zero) → top \texttt{"\ nn\ +"}, middle \texttt{"\ \ \ \ +"}, bottom
  \texttt{"+++++"}.
\item
  \textbf{Black} cell → all three rows \texttt{"+++++"}.
\item
  \textbf{Outside} cells (that became \texttt{−1} after the edge sweep):
  If the cell to the right is also \texttt{−1}, the right border is
  suppressed (\texttt{"\ \ \ \ \ "} on the two upper rows); if the cell
  below** is \texttt{−1}, the bottom border is suppressed
  (\texttt{"\ \ \ \ \ "} on the bottom row). This produces a clean
  rectangular outside frame like Fig. 19.
\end{itemize}

\begin{Shaded}
\begin{Highlighting}[]
\NormalTok{* TAOCP v1 §1.3.2 — Ex.23  Crossword diagram from a 23x23 0/1 matrix}
\NormalTok{* Input : 23 cards on device 16, cols 1..23 are digits 0/1 (one row per card)}
\NormalTok{* Output: puzzle drawn on device 18 (line printer), as in Fig. 19}
\NormalTok{* Method : Add a {-}1 border, flood {-}1 into edge{-}connected 1\textquotesingle{}s, number starts,}
\NormalTok{*          then render 3 printer lines per puzzle row, 5 bytes per cell.}

\NormalTok{        ORIG 3000}
\NormalTok{PRN     EQU   18}
\NormalTok{CARD    EQU   16}

\NormalTok{* {-}{-}{-}{-}{-}{-}{-}{-}{-}{-}{-}{-}{-}{-}{-} character codes we need (portable via ALF) {-}{-}{-}{-}{-}{-}{-}{-}{-}{-}{-}{-}{-}{-}{-}{-}{-}{-}}
\NormalTok{SPCW    ALF  00000           * blank in (1:1)}
\NormalTok{PLUSW   ALF  +0000           * \textquotesingle{}+\textquotesingle{} in (1:1)}
\NormalTok{Z0W     ALF  00000           * \textquotesingle{}0\textquotesingle{} in (1:1)}

\NormalTok{SPC     CON  0}
\NormalTok{PLUS    CON  0}
\NormalTok{D0      CON  0}

\NormalTok{* {-}{-}{-}{-}{-}{-}{-}{-}{-}{-}{-}{-}{-}{-}{-} general constants and pointers {-}{-}{-}{-}{-}{-}{-}{-}{-}{-}{-}{-}{-}{-}{-}{-}{-}{-}{-}{-}{-}{-}{-}{-}{-}{-}{-}{-}{-}{-}}
\NormalTok{N       CON   23}
\NormalTok{NP1     CON   24}
\NormalTok{N2      CON   529}
\NormalTok{FIVE    CON   5}
\NormalTok{THREE   CON   3}
\NormalTok{ONE     CON   1}
\NormalTok{TWO     CON   2}
\NormalTok{ZERO    CON   0}
\NormalTok{NEG1    CON   {-}1}
\NormalTok{TEN     CON   10}

\NormalTok{* {-}{-}{-}{-}{-}{-}{-}{-}{-}{-}{-}{-}{-}{-}{-} storage layout {-}{-}{-}{-}{-}{-}{-}{-}{-}{-}{-}{-}{-}{-}{-}{-}{-}{-}{-}{-}{-}{-}{-}{-}{-}{-}{-}{-}{-}{-}{-}{-}{-}{-}{-}{-}{-}{-}{-}{-}{-}{-}{-}{-}{-}}
\NormalTok{* use a 25x25 board B(i,j), 0..24 in both directions}
\NormalTok{* (board base at B), and a NUM(i,j) table for numbers (1..K)}
\NormalTok{B       EQU   5000                    * 25*25 = 625 words used}
\NormalTok{NUM     EQU   B+625}
\NormalTok{LINE    EQU   NUM+625                 * 24{-}word print line buffer}
\NormalTok{CARDW   EQU   LINE+24                 * 24{-}word card buffer}

\NormalTok{* {-}{-}{-}{-}{-}{-}{-}{-}{-}{-}{-}{-}{-}{-}{-} scratch {-}{-}{-}{-}{-}{-}{-}{-}{-}{-}{-}{-}{-}{-}{-}{-}{-}{-}{-}{-}{-}{-}{-}{-}{-}{-}{-}{-}{-}{-}{-}{-}{-}{-}{-}{-}{-}{-}{-}{-}{-}{-}{-}{-}{-}{-}{-}{-}{-}{-}{-}{-}{-}}
\NormalTok{I       CON 0}
\NormalTok{J       CON 0}
\NormalTok{K       CON 0}
\NormalTok{CHG     CON 0}
\NormalTok{ADDR    CON 0}
\NormalTok{T       CON 0}
\NormalTok{Q       CON 0}
\NormalTok{R       CON 0}
\NormalTok{RET     CON 0}
\NormalTok{D1      CON 0}
\NormalTok{D10     CON 0}

\NormalTok{* ====================== main ==================================================}
\NormalTok{START   LDA  SPCW(1:1);  STA SPC}
\NormalTok{        LDA  PLUSW(1:1); STA PLUS}
\NormalTok{        LDA  Z0W(1:1);   STA D0}

\NormalTok{* {-}{-}{-}{-}{-}{-}{-}{-}{-}{-} zero NUM and fill B with {-}1 border, 0 inside {-}{-}{-}{-}{-}{-}{-}{-}{-}{-}{-}{-}{-}{-}{-}{-}{-}{-}{-}{-}{-}}
\NormalTok{        ENT1 0}
\NormalTok{BORDR   ENT2 0}
\NormalTok{        ENTA  B}
\NormalTok{        ADD   I}
\NormalTok{        MUL   =25=}
\NormalTok{        ADD   J}
\NormalTok{        STA   ADDR}
\NormalTok{        LD3   ADDR}
\NormalTok{* set B(i,j):}
\NormalTok{        LDA   I}
\NormalTok{        JE    SETM1}
\NormalTok{        CMPA  NP1}
\NormalTok{        JE    SETM1}
\NormalTok{        LDA   J}
\NormalTok{        JE    SETM1}
\NormalTok{        CMPA  NP1}
\NormalTok{        JE    SETM1}
\NormalTok{        STZ   0,3                     * interior initial 0}
\NormalTok{        JMP   NEXTB}
\NormalTok{SETM1   LDA   NEG1}
\NormalTok{        STA   0,3}
\NormalTok{NEXTB   * NUM(i,j) ← 0}
\NormalTok{        LDA   ZERO}
\NormalTok{        STA   NUM{-} B,3                * same index, other table}

\NormalTok{        INC2  1}
\NormalTok{        LDA   J}
\NormalTok{        CMPA  =25=}
\NormalTok{        JL    BORDR}
\NormalTok{        INC1  1}
\NormalTok{        LDA   I}
\NormalTok{        CMPA  =25=}
\NormalTok{        JL    BORDR}

\NormalTok{* {-}{-}{-}{-}{-}{-}{-}{-}{-}{-}{-}{-}{-}{-}{-}{-}{-}{-} read 23 cards; store 0/1 into B(1..23,1..23) {-}{-}{-}{-}{-}{-}{-}{-}{-}{-}{-}{-}{-}}
\NormalTok{        ENT1 1                         * row i = 1..23}
\NormalTok{RDROW   IN   CARDW(CARD)}
\NormalTok{        JBUS *      (CARD)             * wait for card}
\NormalTok{        ENT2 1}
\NormalTok{        LD4  CARDW}
\NormalTok{RDCOL   LDA  0,4(0:0)                  * next card byte}
\NormalTok{        ADD  =0=                       * put in A with sign \textquotesingle{}+\textquotesingle{}}
\NormalTok{        SUB  D0}
\NormalTok{        JZ   ISZERO}
\NormalTok{        ENTA 1}
\NormalTok{        JMP  STORE}
\NormalTok{ISZERO  ENTA 0}
\NormalTok{STORE   STX  R                         * (X changed by ADD/SUB? guard)}
\NormalTok{        ENTA  B}
\NormalTok{        ADD   I}
\NormalTok{        MUL   =25=}
\NormalTok{        ADD   J}
\NormalTok{        STA   ADDR}
\NormalTok{        LD3   ADDR}
\NormalTok{        LDA   ZERO}
\NormalTok{        ADD   A                        * A already 0 or 1}
\NormalTok{        STA   0,3}

\NormalTok{        INC4  1}
\NormalTok{        INC2  1}
\NormalTok{        LDA   J}
\NormalTok{        CMPA  N}
\NormalTok{        JL    RDCOL}

\NormalTok{        INC1  1}
\NormalTok{        LDA   I}
\NormalTok{        CMPA  N}
\NormalTok{        JLE   RDROW}

\NormalTok{* {-}{-}{-}{-}{-}{-}{-}{-}{-}{-}{-}{-}{-}{-}{-}{-}{-}{-} flood {-}1 from border into all edge{-}connected 1\textquotesingle{}s {-}{-}{-}{-}{-}{-}{-}{-}{-}}
\NormalTok{FLOOP   STZ  CHG}
\NormalTok{        ENT1 1}
\NormalTok{FI      ENT2 1}
\NormalTok{FJ      * process only original interior domain 1..23}
\NormalTok{        ENTA  B}
\NormalTok{        ADD   I}
\NormalTok{        MUL   =25=}
\NormalTok{        ADD   J}
\NormalTok{        STA   ADDR}
\NormalTok{        LD3   ADDR}
\NormalTok{        LDA   0,3}
\NormalTok{        CMPA  =1=}
\NormalTok{        JNE   NEXTJ}

\NormalTok{        * if a 1 is adjacent (4{-}neigh) to {-}1, turn it into {-}1}
\NormalTok{        ENTA  B}
\NormalTok{        ADD   I}
\NormalTok{        DEC   A ONE}
\NormalTok{        MUL   =25=}
\NormalTok{        ADD   J}
\NormalTok{        STA   T}
\NormalTok{        LDA   0,T}
\NormalTok{        CMPA  NEG1}
\NormalTok{        JE    MAKE}
\NormalTok{        ENTA  B}
\NormalTok{        ADD   I}
\NormalTok{        INC   A ONE}
\NormalTok{        MUL   =25=}
\NormalTok{        ADD   J}
\NormalTok{        STA   T}
\NormalTok{        LDA   0,T}
\NormalTok{        CMPA  NEG1}
\NormalTok{        JE    MAKE}
\NormalTok{        ENTA  B}
\NormalTok{        ADD   I}
\NormalTok{        MUL   =25=}
\NormalTok{        ADD   J}
\NormalTok{        DEC   A ONE}
\NormalTok{        STA   T}
\NormalTok{        LDA   0,T}
\NormalTok{        CMPA  NEG1}
\NormalTok{        JE    MAKE}
\NormalTok{        ENTA  B}
\NormalTok{        ADD   I}
\NormalTok{        MUL   =25=}
\NormalTok{        ADD   J}
\NormalTok{        INC   A ONE}
\NormalTok{        STA   T}
\NormalTok{        LDA   0,T}
\NormalTok{        CMPA  NEG1}
\NormalTok{        JNE   NEXTJ}
\NormalTok{MAKE    LDA   NEG1}
\NormalTok{        STA   0,3}
\NormalTok{        ENTA  1}
\NormalTok{        STA   CHG}
\NormalTok{NEXTJ   INC2  1}
\NormalTok{        LDA   J}
\NormalTok{        CMPA  N}
\NormalTok{        JLE   FJ}
\NormalTok{        INC1  1}
\NormalTok{        LDA   I}
\NormalTok{        CMPA  N}
\NormalTok{        JLE   FI}
\NormalTok{        LDA   CHG}
\NormalTok{        JNZ   FLOOP}

\NormalTok{* {-}{-}{-}{-}{-}{-}{-}{-}{-}{-}{-}{-}{-}{-}{-}{-}{-}{-} numbering: starts of across / down {-}{-}{-}{-}{-}{-}{-}{-}{-}{-}{-}{-}{-}{-}{-}{-}{-}{-}{-}{-}{-}{-}{-}}
\NormalTok{        STZ   K}
\NormalTok{        ENT1  1}
\NormalTok{NI      ENT2  1}
\NormalTok{NJ      ENTA  B}
\NormalTok{        ADD   I}
\NormalTok{        MUL   =25=}
\NormalTok{        ADD   J}
\NormalTok{        STA   ADDR}
\NormalTok{        LD3   ADDR}
\NormalTok{        LDA   0,3}
\NormalTok{        JZ    WHITEOK}
\NormalTok{        JMP   NXTJ}
\NormalTok{WHITEOK * check down start: (i==1 or B(i{-}1,j)!=0) and (i\textless{}23 and B(i+1,j)==0)}
\NormalTok{        ENTA  0}
\NormalTok{        STA   T}
\NormalTok{        LDA   I}
\NormalTok{        JNE   CI1}
\NormalTok{        ENTA  1}
\NormalTok{        STA   T}
\NormalTok{        JMP   C2}
\NormalTok{CI1     ENTA  B}
\NormalTok{        ADD   I}
\NormalTok{        DEC   A ONE}
\NormalTok{        MUL   =25=}
\NormalTok{        ADD   J}
\NormalTok{        LDA   0,A}
\NormalTok{        JNZ   C2}
\NormalTok{        ENTA  0}
\NormalTok{        STA   T}
\NormalTok{C2      LDA   T}
\NormalTok{        JZ    NOTD}
\NormalTok{        LDA   I}
\NormalTok{        CMPA  N}
\NormalTok{        JE    NOTD}
\NormalTok{        ENTA  B}
\NormalTok{        ADD   I}
\NormalTok{        INC   A ONE}
\NormalTok{        MUL   =25=}
\NormalTok{        ADD   J}
\NormalTok{        LDA   0,A}
\NormalTok{        JNZ   NOTD}
\NormalTok{        ENTA  1}
\NormalTok{        STA   T                 * down start}
\NormalTok{NOTD    * check across start}
\NormalTok{        LDA   T}
\NormalTok{        JNZ   NUMBER}
\NormalTok{        LDA   J}
\NormalTok{        JNE   CJ1}
\NormalTok{        ENTA  1}
\NormalTok{        STA   T}
\NormalTok{        JMP   C4}
\NormalTok{CJ1     ENTA  B}
\NormalTok{        ADD   I}
\NormalTok{        MUL   =25=}
\NormalTok{        ADD   J}
\NormalTok{        DEC   A ONE}
\NormalTok{        LDA   0,A}
\NormalTok{        JNZ   C4}
\NormalTok{        ENTA  0}
\NormalTok{        STA   T}
\NormalTok{C4      LDA   T}
\NormalTok{        JNZ   NUMBERC}
\NormalTok{        LDA   J}
\NormalTok{        CMPA  N}
\NormalTok{        JE    NXTJ}
\NormalTok{        ENTA  B}
\NormalTok{        ADD   I}
\NormalTok{        MUL   =25=}
\NormalTok{        ADD   J}
\NormalTok{        INC   A ONE}
\NormalTok{        LDA   0,A}
\NormalTok{        JNZ   NXTJ}
\NormalTok{NUMBERC ENTA  1}
\NormalTok{        STA   T}
\NormalTok{NUMBER  LDA   K}
\NormalTok{        ADD   ONE}
\NormalTok{        STA   K}
\NormalTok{        STA   NUM{-} B,3}
\NormalTok{NXTJ    INC2  1}
\NormalTok{        LDA   J}
\NormalTok{        CMPA  N}
\NormalTok{        JLE   NJ}
\NormalTok{        INC1  1}
\NormalTok{        LDA   I}
\NormalTok{        CMPA  N}
\NormalTok{        JLE   NI}

\NormalTok{* {-}{-}{-}{-}{-}{-}{-}{-}{-}{-}{-}{-}{-}{-}{-}{-}{-}{-} render to printer {-}{-}{-}{-}{-}{-}{-}{-}{-}{-}{-}{-}{-}{-}{-}{-}{-}{-}{-}{-}{-}{-}{-}{-}{-}{-}{-}{-}{-}{-}{-}{-}{-}{-}{-}{-}{-}{-}{-}{-}}
\NormalTok{        ENT1  1}
\NormalTok{ROW     * we emit 3 printer lines for row i}
\NormalTok{        LD3   =1=}
\NormalTok{        STA   T                * t = which subrow (1..3)}

\NormalTok{SUBROW  LD4   LINE}
\NormalTok{        ENT2  1}
\NormalTok{COL     ENTA  B}
\NormalTok{        ADD   I}
\NormalTok{        MUL   =25=}
\NormalTok{        ADD   J}
\NormalTok{        STA   ADDR}
\NormalTok{        LD3   ADDR}
\NormalTok{        LDA   0,3              * val = B(i,j)}
\NormalTok{        JN    OUTM1            * negative → outside}
\NormalTok{        JZ    OUTWHT           * zero → white}
\NormalTok{        JMP   OUTBLK           * positive → black}

\NormalTok{* {-}{-}{-}{-}{-} black square: \textquotesingle{}+++++\textquotesingle{} on all 3 subrows}
\NormalTok{OUTBLK  LDA   PLUS}
\NormalTok{        STA   0,4(1:1)}
\NormalTok{        STA   0,4(2:2)}
\NormalTok{        STA   0,4(3:3)}
\NormalTok{        STA   0,4(4:4)}
\NormalTok{        STA   0,4(5:5)}
\NormalTok{        JMP   NEXTC}

\NormalTok{* {-}{-}{-}{-}{-} white square}
\NormalTok{OUTWHT  LDA   T}
\NormalTok{        CMPA  =1=}
\NormalTok{        JNE   W2}
\NormalTok{        * top row: maybe a number}
\NormalTok{        LDA   NUM{-} B,3}
\NormalTok{        JZ    WTOPS}
\NormalTok{        ENTX  0}
\NormalTok{        DIV   TEN}
\NormalTok{        STX   D1}
\NormalTok{        STA   D10}
\NormalTok{        LDA   SPC;  STA 0,4(1:1)}
\NormalTok{        LDA   D10; ADD D0; STA 0,4(2:2)}
\NormalTok{        LDA   D1;  ADD D0; STA 0,4(3:3)}
\NormalTok{        LDA   SPC; STA 0,4(4:4)}
\NormalTok{        LDA   PLUS;STA 0,4(5:5)}
\NormalTok{        JMP   NEXTC}
\NormalTok{WTOPS   LDA   SPC; STA 0,4(1:1); STA 0,4(2:2); STA 0,4(3:3); STA 0,4(4:4)}
\NormalTok{        LDA   PLUS; STA 0,4(5:5)}
\NormalTok{        JMP   NEXTC}
\NormalTok{W2      LDA   T}
\NormalTok{        CMPA  =2=}
\NormalTok{        JNE   W3}
\NormalTok{        * middle row: "    +"}
\NormalTok{        LDA   SPC; STA 0,4(1:1); STA 0,4(2:2); STA 0,4(3:3); STA 0,4(4:4)}
\NormalTok{        LDA   PLUS; STA 0,4(5:5)}
\NormalTok{        JMP   NEXTC}
\NormalTok{W3      * bottom row: "+++++"}
\NormalTok{        LDA   PLUS}
\NormalTok{        STA   0,4(1:1); STA 0,4(2:2); STA 0,4(3:3); STA 0,4(4:4); STA 0,4(5:5)}
\NormalTok{        JMP   NEXTC}

\NormalTok{* {-}{-}{-}{-}{-} outside ({-}1) square: suppress borders that abut {-}1 at right/below}
\NormalTok{OUTM1   LDA   T}
\NormalTok{        CMPA  =3=}
\NormalTok{        JE    M1BOT}
\NormalTok{        * top/mid subrows}
\NormalTok{        * if right neighbor is {-}1 → five blanks; else "    +"}
\NormalTok{        ENTA  B}
\NormalTok{        ADD   I}
\NormalTok{        MUL   =25=}
\NormalTok{        ADD   J}
\NormalTok{        INC   A ONE}
\NormalTok{        LDA   0,A}
\NormalTok{        CMPA  NEG1}
\NormalTok{        JE    ALLSP}
\NormalTok{        LDA   SPC; STA 0,4(1:1); STA 0,4(2:2); STA 0,4(3:3); STA 0,4(4:4)}
\NormalTok{        LDA   PLUS; STA 0,4(5:5)}
\NormalTok{        JMP   NEXTC}
\NormalTok{ALLSP   LDA   SPC}
\NormalTok{        STA   0,4(1:1); STA 0,4(2:2); STA 0,4(3:3); STA 0,4(4:4); STA 0,4(5:5)}
\NormalTok{        JMP   NEXTC}
\NormalTok{M1BOT   * bottom subrow}
\NormalTok{        ENTA  B}
\NormalTok{        ADD   I}
\NormalTok{        INC   A ONE}
\NormalTok{        MUL   =25=}
\NormalTok{        ADD   J}
\NormalTok{        LDA   0,A}
\NormalTok{        CMPA  NEG1}
\NormalTok{        JE    ALLSP}
\NormalTok{        LDA   PLUS}
\NormalTok{        STA   0,4(1:1); STA 0,4(2:2); STA 0,4(3:3); STA 0,4(4:4); STA 0,4(5:5)}

\NormalTok{NEXTC   INC4 1}
\NormalTok{        INC2 1}
\NormalTok{        LDA  J}
\NormalTok{        CMPA N}
\NormalTok{        JLE  COL}

\NormalTok{        OUT  LINE(PRN)}

\NormalTok{        LDA  T}
\NormalTok{        ADD  ONE}
\NormalTok{        STA  T}
\NormalTok{        CMPA THREE}
\NormalTok{        JLE  SUBROW}

\NormalTok{        INC1 1}
\NormalTok{        LDA  I}
\NormalTok{        CMPA N}
\NormalTok{        JLE  ROW}

\NormalTok{        HLT}
\NormalTok{        END  START}
\end{Highlighting}
\end{Shaded}

\textbf{Notes, hints, and adjustments}

\begin{itemize}
\tightlist
\item
  \textbf{Card input}: the code expects 23 cards; each of the first 23
  bytes is \texttt{\textquotesingle{}0\textquotesingle{}} or
  \texttt{\textquotesingle{}1\textquotesingle{}}. If your MIX monitor
  uses a different input device, just change \texttt{CARD=16} and the
  \texttt{IN}/\texttt{JBUS} calls accordingly.
\item
  \textbf{Edge removal} (the ``−1 flood''): we iterate until a full pass
  makes no change; for a 23×23 matrix this is trivial in time.
\item
  \textbf{Numbering rule}: a white square gets a number if it begins a
  \textbf{down} word (white below, not white above) or an
  \textbf{across} word (white right, not white left). Numbers are two
  digits with a \textbf{leading zero} as in the book's figure.
\item
  \textbf{Exact line lengths}: each printer line has 23 boxes × 5 bytes
  = 115 bytes; the line buffer has 24 words (120 bytes), so the last
  five are blank and harmless.
\item
  You can change \texttt{B}, \texttt{NUM}, or \texttt{LINE} bases if
  they clash with your monitor; all addresses are absolute here to keep
  the program self-contained.
\end{itemize}

This produces a printout that matches the layout in Fig. 19: a
plus-bordered diagram, numbers placed at the correct starts, black
squares filled with `+', and edge-connected black regions removed.

\bookmarksetup{startatroot}

\chapter{1.3.3 Applications to
Permutation}\label{applications-to-permutation}

\section{1 {[}02{]}}\label{section-141}

\begin{enumerate}
\def\labelenumi{\arabic{enumi})}
\item
  It's a permutation. On residues mod 7, multiplication by 2 has an
  inverse because \(\gcd(2,7)=1\). In fact \(2\cdot4\equiv1\pmod7\), so
  the inverse map is \(x\mapsto4x\pmod7\). A map with an inverse is
  bijective, hence a permutation.
\item
  Cycle decomposition. Compute \(f(x)=2x\bmod7\):
\end{enumerate}

\[
\begin{array}{c|ccccccc}
x&0&1&2&3&4&5&6\\\hline
f(x)&0&2&4&6&1&3&5
\end{array}
\]

Now follow orbits:

\begin{itemize}
\tightlist
\item
  \(0\mapsto0\) ⇒ \((0)\).
\item
  \(1\mapsto2\mapsto4\mapsto1\) ⇒ \((1\ 2\ 4)\).
\item
  \(3\mapsto6\mapsto5\mapsto3\) ⇒ \((3\ 6\ 5)\).
\end{itemize}

Thus

\[
\boxed{(0)(1\ 2\ 4)(3\ 6\ 5)}.
\]

By the section's convention, singleton cycles are often suppressed,
giving \(\boxed{(1\ 2\ 4)(3\ 6\ 5)}\), which matches the book's answer.

\begin{enumerate}
\def\labelenumi{\arabic{enumi})}
\setcounter{enumi}{2}
\tightlist
\item
  (Why the structure looks like this.) The nonzero residues mod 7 form a
  cyclic group of order 6. The element \(2\) has order \(3\) (since
  \(2^3\equiv1\pmod7\)), so its action splits \(\{1,\dots,6\}\) into
  3-cycles; hence exactly two 3-cycles plus the fixed point \(0\).
\end{enumerate}

\section{2 {[}10{]}}\label{section-142}

We want to transform the tuple

\[
(a,b,c,d,e,f)\longrightarrow(c,d,f,b,e,a)
\]

using only exchanges (swaps \(x \leftrightarrow y\)) and no auxiliary
variable. The target permutation is the product of disjoint cycles

\[
(a\ c\ f)\ (b\ d)\ (e),
\]

since \(a\to c\to f\to a\), \(b\leftrightarrow d\), and \(e\) is fixed.

A \(k\)-cycle \((x_1\,x_2\,\dots\,x_k)\) can be realized by \(k-1\)
swaps:

\[
(x_1\leftrightarrow x_2),\ (x_2\leftrightarrow x_3),\ \dots,\ (x_{k-1}\leftrightarrow x_k).
\]

Apply this to our cycles:

\begin{enumerate}
\def\labelenumi{\arabic{enumi}.}
\tightlist
\item
  Realize \((a\ c\ f)\) with two swaps
\end{enumerate}

\begin{itemize}
\tightlist
\item
  \(a \leftrightarrow c\)
\item
  \(c \leftrightarrow f\)
\end{itemize}

\begin{enumerate}
\def\labelenumi{\arabic{enumi}.}
\setcounter{enumi}{1}
\tightlist
\item
  Realize \((b\ d)\) with one swap
\end{enumerate}

\begin{itemize}
\tightlist
\item
  \(b \leftrightarrow d\)
\end{itemize}

After these three exchanges the tuple is exactly \((c,d,f,b,e,a)\). So
one valid swap-only sequence is:

\[
\boxed{a\leftrightarrow c,\quad c\leftrightarrow f,\quad b\leftrightarrow d.}
\]

(That's precisely the concise answer Knuth lists.)

The two swaps on \((a,c,f)\) rotate the triple once (a 3-cycle), and the
single swap on \((b,d)\) flips those two; disjoint cycles act
independently, so doing the 3-cycle then the transposition yields the
desired permutation.

\section{3 {[}03{]}}\label{section-143}

We're asked to compute (left-to-right, per Knuth's convention) the
product

\[
\pi_1\times\pi_2
=
\begin{pmatrix}
a&b&c&d&e&f\\
b&d&c&a&f&e
\end{pmatrix}
\;\times\;
\begin{pmatrix}
a&b&c&d&e&f\\
c&d&f&b&e&a
\end{pmatrix},
\]

and express the answer in two-line notation. Knuth multiplies
permutations left-to-right (apply the left factor, then the right), see
the discussion around Eq. (4).

Track each letter through \(\pi_1\) then \(\pi_2\):

\begin{itemize}
\tightlist
\item
  \(a \xrightarrow{\pi_1} b \xrightarrow{\pi_2} d\) ⇒ \(a\mapsto d\).
\item
  \(b \xrightarrow{\pi_1} d \xrightarrow{\pi_2} b\) ⇒ \(b\mapsto b\).
\item
  \(c \xrightarrow{\pi_1} c \xrightarrow{\pi_2} f\) ⇒ \(c\mapsto f\).
\item
  \(d \xrightarrow{\pi_1} a \xrightarrow{\pi_2} c\) ⇒ \(d\mapsto c\).
\item
  \(e \xrightarrow{\pi_1} f \xrightarrow{\pi_2} a\) ⇒ \(e\mapsto a\).
\item
  \(f \xrightarrow{\pi_1} e \xrightarrow{\pi_2} e\) ⇒ \(f\mapsto e\).
\end{itemize}

Therefore

\[
\boxed{
\begin{pmatrix}
a&b&c&d&e&f\\
d&b&f&c&a&e
\end{pmatrix}
}.
\]

If you also want the cycle form, it's \((a\,d\,c\,f\,e)(b)\), i.e.,
\((a\,d\,c\,f\,e)\) with \(b\) fixed.

\section{4 {[}10{]}}\label{section-144}

Problem. Express

\[
(a\,b\,d)\,(e\,f)\,(a\,c\,f)\,(b\,d)
\]

as a product of disjoint cycles.

Convention. In this section, permutation products are applied
left-to-right: first \((a\,b\,d)\), then \((e\,f)\), then \((a\,c\,f)\),
then \((b\,d)\). (See the section text immediately preceding the
exercises.)

Compute the image of each symbol \(x\in\{a,b,c,d,e,f\}\) by chasing it
through the four cycles:

\begin{itemize}
\tightlist
\item
  \(a \xrightarrow{(a\,b\,d)} b \xrightarrow{(e\,f)} b \xrightarrow{(a\,c\,f)} b \xrightarrow{(b\,d)} d\)
  ⇒ \(a\mapsto d\).
\item
  \(b \xrightarrow{(a\,b\,d)} d \xrightarrow{(e\,f)} d \xrightarrow{(a\,c\,f)} d \xrightarrow{(b\,d)} b\)
  ⇒ \(b\mapsto b\).
\item
  \(c \xrightarrow{(a\,b\,d)} c \xrightarrow{(e\,f)} c \xrightarrow{(a\,c\,f)} f \xrightarrow{(b\,d)} f\)
  ⇒ \(c\mapsto f\).
\item
  \(d \xrightarrow{(a\,b\,d)} a \xrightarrow{(e\,f)} a \xrightarrow{(a\,c\,f)} c \xrightarrow{(b\,d)} c\)
  ⇒ \(d\mapsto c\).
\item
  \(e \xrightarrow{(a\,b\,d)} e \xrightarrow{(e\,f)} f \xrightarrow{(a\,c\,f)} a \xrightarrow{(b\,d)} a\)
  ⇒ \(e\mapsto a\).
\item
  \(f \xrightarrow{(a\,b\,d)} f \xrightarrow{(e\,f)} e \xrightarrow{(a\,c\,f)} e \xrightarrow{(b\,d)} e\)
  ⇒ \(f\mapsto e\).
\end{itemize}

Thus the permutation maps

\[
a\to d,\quad d\to c,\quad c\to f,\quad f\to e,\quad e\to a,\quad b\to b,
\]

which is the 5-cycle \((a\,d\,c\,f\,e)\) with \(b\) fixed (singleton
cycles are suppressed). Therefore

\[
\boxed{(a\,d\,c\,f\,e)}.
\]

This matches the book's listed answer.

\section{5 {[}M10{]}}\label{m10-13}

If a permutation has \(r\) disjoint cycles of lengths
\(l_1,l_2,\dots,l_r\) (no 1-cycles shown), then:

\begin{itemize}
\tightlist
\item
  Each \(l_i\)-cycle can be rotated in \(l_i\) ways without changing the
  permutation.
\item
  The \(r\) cycles can be ordered arbitrarily (they commute), giving
  \(r!\) orderings.
\end{itemize}

Hence the number of distinct written forms is

\[
\boxed{\,l_1\,l_2\cdots l_r \cdot r!\, }.
\]

Apply to Eq. (3). There the permutation is \((a\,c\,f)(b\,d)\), i.e.,
one 3-cycle and one 2-cycle, so

\[
3 \times 2 \times 2! \;=\; \boxed{12}.
\]

\section{6 {[}M28{]}}\label{m28-6}

Knuth's timing formula for Program A (excluding I/O) under the ``all
blanks at the extreme right'' assumption is

\[
(112NX+304X-2M-Y+11U+2V-11)\,u,
\]

with \(X,Y,M,N,U,V\) defined just below Eq. (18)--(19). In the text it's
stated that this assumption prevents two things from ever happening: the
jump on line 32 and the ``skip past blanks'' loop at line 71.

If we remove that assumption, let

\[
W=\text{\# of blank words in the input (excluding the final "=")}=16X-1-Y.
\]

using \(B+C=16X-1\) and \(Y=\) \# of nonblank words.

\begin{itemize}
\tightlist
\item
  Each internal blank that's encountered during the preconditioning pass
  (lines 23--38) causes the jump at line 32 to be taken once, adding
  \(W\) unit-times.
\item
  Each internal blank that's encountered during the main scan (A3)
  triggers the skip at line 71 once, adding another \(W\) unit-times.
\end{itemize}

So the timing increases by \(2W=(32X-2-2Y)\) unit-times. The new total
(still excluding I/O) is

\[
\boxed{(112NX+304X-2M-Y+11U+2V-11+2(16X-1-Y))\,u.}
\]

Equivalently,

\[
(112NX+336X-2M-3Y+11U+2V-13)\,u.
\]

(Variables and the original timing setup are given around Eqs. (7),
(18)--(19).)

\section{7 {[}10{]}}\label{section-145}

Here are the parameters for the specific input (6)
\((a\,c\,f\,g)(b\,c\,d)(a\,e\,d)(f\,a\,d\,e)(b\,g\,f\,a\,e)\):

\begin{itemize}
\tightlist
\item
  \(X\) (cards of input) = 2 (the book shows (6) punched on two cards).
\item
  \(Y\) (nonblank fields, excluding the final ``='') = 29 -- words: 6 +
  5 + 5 + 6 + 7 = 29.
\item
  \(M\) (cycles in input) = 5 (the five parenthesized cycles).
\item
  \(N\) (distinct element names) = 7 (a,b,c,d,e,f,g).
\item
  Output permutation for (6) is \((a\,d\,g)(c\,e\,b)(f)\), so \(U\)
  (cycles in output, incl.~singletons) = 3, \(V\) (singleton cycles) =
  1.
\end{itemize}

Using Knuth's time formula (exclusive of I/O)
\((112 N X + 304 X - 2M - Y + 11U + 2V - 11)\,u\), with the meanings in
(19):

\[
\begin{aligned}
T &= 112(7)(2) + 304(2) - 2(5) - 29 + 11(3) + 2(1) - 11 \\
  &= \boxed{2161\,u}.
\end{aligned}
\]

\section{8 {[}23{]}}\label{section-146}

Yes. Scan left→right and maintain the inverse relation: make \(T[j]=i\)
exactly when \(x_i \mapsto x_j\) in the permutation (so \(x_j\)'s
predecessor is \(x_i\)). After the scan, construct the cycle form by
tracing right→left with the table \(T\) (i.e., repeatedly replace the
current element by its predecessor from \(T\)), which yields the same
permutation in cycle notation.

How to implement:

\begin{itemize}
\tightlist
\item
  On reading a cycle \((x_{i_1}\ x_{i_2}\ \dots\ x_{i_k})\) left→right,
  set \(T[i_2]=i_1,\ T[i_3]=i_2,\ \dots,\ T[i_k]=i_{k-1},\ T[i_1]=i_k\).
\item
  When all cycles are processed, output cycles by picking an unused name
  \(a\) and repeatedly replacing \(a \leftarrow T[a]\) until you return
  to the start, marking as you go.
\end{itemize}

\section{9 {[}10{]}}\label{section-147}

\begin{enumerate}
\def\labelenumi{\arabic{enumi})}
\tightlist
\item
  First get the actual permutation \(F\)
\end{enumerate}

Chase each letter through the whole product (left factor first, then the
next, \ldots):

\begin{itemize}
\tightlist
\item
  \(a \to c \to c \to e \to d \to d\Rightarrow a\mapsto d\)
\item
  \(b \to b \to c \to c \to c \to e\Rightarrow b\mapsto e\)
\item
  \(c \to f \to d \to d \to d \to b\Rightarrow c\mapsto b\)
\item
  \(d \to d \to b \to a \to a \to g\Rightarrow d\mapsto g\)
\item
  \(e \to e \to e \to a \to a \to a\Rightarrow e\mapsto a\)
\item
  \(f \to g \to g \to g \to f \to f\Rightarrow f\mapsto f\)
\item
  \(g \to a \to a \to a \to g \to g\Rightarrow g\mapsto a\)
\end{itemize}

So the mapping is

\[
F=\{a\!\to\!d,\ d\!\to\!g,\ g\!\to\!a\}\ \{c\!\to\!b,\ b\!\to\!e,\ e\!\to\!a\}\ \{f\!\to\!f\}.
\]

In cycle form that's

\[
(a\,d\,g)\ (c\,b\,e)\ (f)\,,
\]

and of course any rotation like \((c\,e\,b)\) denotes the same 3-cycle.

\begin{enumerate}
\def\labelenumi{\arabic{enumi})}
\setcounter{enumi}{1}
\tightlist
\item
  Program A (left→right scan): why cycles start with the leftmost
  available name
\end{enumerate}

Program A builds the permutation while scanning the input text
left→right, then prints cycles by:

\begin{itemize}
\tightlist
\item
  keeping a \texttt{USED{[}x{]}} mark for each name,
\item
  repeatedly finding the leftmost name in the input whose \texttt{USED}
  mark is false,
\item
  starting a new printed cycle at that name and following images
  \(x \mapsto F[x]\) until it returns, marking \texttt{USED} along the
  way.
\end{itemize}

Concretely for (6):

Left→right order of name occurrences (ignoring parentheses and ``='')
\texttt{a,\ c,\ f,\ g,\ b,\ c,\ d,\ a,\ e,\ d,\ f,\ a,\ d,\ e,\ b,\ g,\ f,\ a,\ e}

The first-appearance index of each name in this order is:

\begin{itemize}
\tightlist
\item
  \(a\!:1,\ c\!:2,\ f\!:3,\ g\!:4,\ b\!:5,\ d\!:7,\ e\!:9\).
\end{itemize}

Now the printing loop:

\begin{enumerate}
\def\labelenumi{\arabic{enumi}.}
\item
  Pick the leftmost unused ⇒ \(a\) (index 1). Follow
  \(a\to d\to g\to a\). Mark \(\{a,d,g\}\) and print \((a\,d\,g)\).
\item
  Next leftmost unused in the original text is \(c\) (index 2). Follow
  \(c\to b\to e\to c\). Mark \(\{c,b,e\}\) and print \((c\,b\,e)\)
  (often rotated as \((c\,e\,b)\)).
\item
  Remaining unused is \(f\) (index 3). Follow \(f\to f\). Print \((f)\).
\end{enumerate}

That's why every cycle that Program A prints begins with the leftmost
still-unprinted name in the original left→right input: the algorithm
literally chooses the next start that way before walking \(F\).

\begin{enumerate}
\def\labelenumi{\arabic{enumi})}
\setcounter{enumi}{2}
\tightlist
\item
  Program B (right→left scan): why cycles start with the last distinct
  name seen from the right
\end{enumerate}

Program B first scans the input right→left and builds predecessor links
\(T\) (so \(T[y]=x\) when \(x\mapsto y\)); during this pass it also
keeps the order in which distinct names are first encountered from the
right.

For (6), the right→left name stream is:

\texttt{e,\ a,\ f,\ g,\ b,\ e,\ d,\ a,\ f,\ d,\ e,\ a,\ d,\ c,\ b,\ g,\ f,\ c,\ a}

The first time we see each distinct name (while moving right→left)
occurs in this order:

\[
e,\ a,\ f,\ g,\ b,\ d,\ c.
\]

Call that list R2L. After the scan, Program B prints cycles by
repeatedly:

\begin{itemize}
\tightlist
\item
  choosing the last name in R2L that is still unused,
\item
  starting the next cycle there and walking predecessors (equivalently,
  the same cycle as walking images but rotated),
\item
  marking those names used and repeating.
\end{itemize}

Step-by-step:

\begin{enumerate}
\def\labelenumi{\arabic{enumi}.}
\item
  Last distinct from the right among all names is \(c\) (it's the final
  entry of R2L). Walk the cycle:
  \(c \leftarrow b \leftarrow e \leftarrow c\). Print \((c\,b\,e)\)
  (often shown as \((c\,e\,b)\)). Mark \(\{c,b,e\}\) used.
\item
  Among the remaining names \(\{a,d,f,g\}\), the last (furthest to the
  left in R2L) is \(d\). Walk:
  \(d \leftarrow g \leftarrow a \leftarrow d\). Print \((d\,g\,a)\)
  (same cycle as \((a\,d\,g)\)). Mark \(\{d,g,a\}\).
\item
  Remaining is \(f\). Walk: \(f \leftarrow f\). Print \((f)\).
\end{enumerate}

That's exactly the rule: Program B's cycle starts are the last distinct
names encountered during the right→left pass, taken anew each time after
removing already-printed elements.

\begin{enumerate}
\def\labelenumi{\arabic{enumi})}
\setcounter{enumi}{3}
\tightlist
\item
  The general reason (independent of the example)
\end{enumerate}

\begin{itemize}
\tightlist
\item
  Program A seeds each cycle by a left→right scan for the first
  unprinted symbol; hence starts are leftmost-available. The subsequent
  walk doesn't change the start---it just completes that orbit under
  \(F\).
\item
  Program B fixes the cycle starts at build time: the right→left pass
  commits to a ``most recently discovered, still-unused'' symbol as the
  next start, which is precisely the last distinct from the right among
  those not yet printed. Walking via predecessors (or, equivalently,
  images) just rotates the cycle to begin at that chosen start.
\end{itemize}

\begin{enumerate}
\def\labelenumi{\arabic{enumi})}
\setcounter{enumi}{4}
\tightlist
\item
  Sanity checks
\end{enumerate}

\begin{itemize}
\tightlist
\item
  The two outputs may rotate cycles differently, but they encode the
  same permutation \(F\).
\item
  If you reorder or space the input differently, Program A's starts can
  change (because ``leftmost'' refers to the input text order), and
  Program B's starts can change with the right→left discovery order;
  both behaviors match the book's statements.
\end{itemize}

\section{10 {[}M28{]}}\label{m28-7}

A. Where the equations come from (Kirchhoff's law on the flow graph)

Kirchhoff's law on a flow graph says: for every node, total entries =
total exits (counted over one complete run). In §1.3.3, Knuth labels
parts of Program A and uses capital letters \(A,B,\dots\) for ``how many
times that part executes.'' Applying the balance equation at the key
junctions yields simple linear relations among these counters. The text
gives the same kind of relations already (see the group of equations
around the analysis of Program A), and we reproduce the ones Ex. 10
wants.

From node-by-node balances on the Program A flow (as in the section's
derivation), we obtain:

\[
\boxed{
\begin{aligned}
&\text{(1)}\quad A = 1 + C - D,\\
&\text{(2)}\quad B = A + J + P - 1,\\
&\text{(3)}\quad C = B - (P - L),\\
&\text{(4)}\quad E = D - L,\\
&\text{(5)}\quad G = E,\\
&\text{(6)}\quad Q = Z,\\
&\text{(7)}\quad W = S.
\end{aligned}}
\]

Here is the intuition for each, in words:

\begin{itemize}
\item
  \(A = 1 + C - D\). You enter the ``initialize/outer loop'' block
  (\(A\)) once at program start (that's the ``1''), and you re-enter it
  each time you loop back from scanning (\(C\)), except when you
  immediately peel off to a ``done with this cycle'' exit counted by
  \(D\). Net effect: \(A\) equals one initial entry plus \(C\) returns
  minus \(D\) departures. (This is the standard ``entry = one initial +
  reentries − final exits'' pattern.)
\item
  \(B = A + J + P - 1\). The scanning/reading block \(B\) is reached
  after each \(A\) (except the very last time), and also via two
  back-edges: one from the cycle-outputting region (\(J\)) and another
  from the ``suppress/close'' area (\(P\)). Summing those, and
  subtracting the final ``no further input'' case, yields
  \(B = A + J + P - 1\).
\item
  \(C = B - (P - L)\). From each pass through \(B\) we typically proceed
  to \(C\) (the main scan), except when an immediate close/suppress path
  fires; those exceptional bypasses are counted by \((P - L)\) (closes
  without opening a new cycle). Hence \(C\) is the ``default'' \(B\)
  count minus the bypasses.
\item
  \(E = D - L\), then \(G=E\). Each time a cycle finishes (\(D\)) we
  either immediately close/suppress (\(L\)) or proceed to the
  ``emit/open'' work counted by \(E\); so \(E = D - L\). And the
  subsequent sub-task \(G\) runs exactly once per \(E\), giving \(G=E\).
\item
  \(Q = Z\). The names-table pass that establishes a correspondence for
  a \emph{new} distinct name (\(Z\)) is in one-to-one correspondence
  with creating a new name entry \(Q\) (same count).
\item
  \(W = S\). Likewise two paired activities (different labels in the
  listing) occur the same number of times along matched edges, giving
  equality.
\end{itemize}

These ``balance'' equations are exactly the relations the official
answer lists under part (1).

B. Interpreting the symbols (connect to \(X,Y,M,N,U,\dots\))

Now tie each counter to a plain combinatorial quantity. Program A reads
permutations in cycle notation from punched cards (16 fields per card),
``preconditions'' the stream, builds a table of names, multiplies the
permutation, and prints cycles. In §1.3.3 the quantities \(X,Y,M,N,U,V\)
are used to summarize the input and output:

\begin{itemize}
\tightlist
\item
  \(X\): number of cards of input,
\item
  \(Y\): number of nonblank words on the cards (excluding the final
  ``=''),
\item
  \(M\): number of cycles in the input,
\item
  \(N\): number of distinct element names,
\item
  \(U\): number of cycles in the output (including singletons),
\item
  \(V\): number of singleton cycles in the output.
\end{itemize}

With those conventions, Ex. 10 asks for concrete meanings of several
counters in terms of \(X,Y,M,N,U\). The official identifications are
(and we explain each):

\begin{itemize}
\item
  \(\boxed{B = \text{\# of words of input} = 16X - 1}\). Each card
  supplies 16 word positions; the very last position is the terminating
  ``='', hence total words processed (including blanks) is \(16X-1\).
  This is precisely the stream length the scanner sees.
\item
  \(\boxed{C = \text{\# of nonblank words} = Y}\). \(C\) counts how many
  \emph{actual} fields (names or parentheses) are scanned; that's the
  ``nonblank'' count \(Y\) defined in the section.
\item
  \(\boxed{D = C - M}\). Among the \(C\) nonblank fields are \(M\) right
  parentheses that immediately terminate cycles; the control that
  increments \(D\) triggers for all other nonblank fields, hence \(D\)
  is the remainder \(C-M\). (Equivalently: each of the \(M\) cycle
  closures diverts flow to \(L\), not counted in \(D\).)
\item
  \(\boxed{E = D - M}\). After removing the \(M\) times we \emph{close}
  cycles from \(D\), the remaining cases require ``emit/open'' work,
  counted by \(E\); the balance gives \(E=D-M\). (This aligns with the
  earlier \(E=D-L\), together with \(L=M\) below.)
\item
  \(\boxed{F = \text{\# of comparisons in the names-table search}}\).
  Every nonblank name triggers a symbol-table probe; \(F\) totals those
  name-equality checks. (This is exactly how Program A's table lookup
  costs are accounted.)
\item
  \(\boxed{H = N}\). \(H\) is the total number of times we \emph{insert}
  a previously unseen name; that equals the number of distinct names
  \(N\).
\item
  \(\boxed{K = M}\). \(K\) counts the number of times we start a new
  input cycle (i.e., see `(' that actually begins a cycle), which is
  exactly the number of input cycles \(M\).
\item
  \(\boxed{Q = N}\). \(Q\) increments when a new name slot is
  \emph{created} in the table; that's once per distinct name, hence
  \(Q=N\). (Consistent with \(Q=Z\) and ``\(Z\) is \# of new-name
  events.'')
\item
  \(\boxed{R = U}\). \(R\) is the number of output cycles produced;
  Program A prints all cycles including singletons, so \(R=U\).
\end{itemize}

These are exactly the identifications listed for part (2) in the answer
key.

C. Cross-checks against the section's earlier relations

Two useful consistency checks from the text:

\begin{itemize}
\item
  The section earlier derives other linear constraints (e.g., relations
  tying \(G,J,L\) to the scan passes) and shows that the running time
  depends only on \(G+L\), so we don't need to separate \(G\) and \(L\)
  further for timing. This is used later to write the closed-form time
  in terms of \(X,Y,M,N,U,V\).
\item
  When blanks are \emph{not} all at the extreme right, additional passes
  over internal blanks add exactly one unit per internal blank to two
  specific actions (a jump and a skip loop), increasing the total by
  \(2(16X-1-Y)\), which matches how \(B\) and \(C=Y\) enter the count.
  (This ties back to the timing refinements you used in Ex. 6--7.)
\end{itemize}

\begin{enumerate}
\def\labelenumi{(\arabic{enumi})}
\tightlist
\item
  By Kirchhoff's law on the Program A flow,
\end{enumerate}

\[
A = 1 + C - D,\quad
B = A + J + P - 1,\quad
C = B - (P - L),\quad
E = D - L,\quad
G = E,\quad
Q = Z,\quad
W = S.
\]

\begin{enumerate}
\def\labelenumi{(\arabic{enumi})}
\setcounter{enumi}{1}
\tightlist
\item
  Interpretations in terms of input/output parameters:
\end{enumerate}

\[
\begin{aligned}
&B = 16X - 1,\quad C = Y,\quad D = C - M,\quad E = D - M,\quad F=\text{\# table comparisons},\\
&H = N,\quad K = M,\quad Q = N,\quad R = U.
\end{aligned}
\]

\section{11 {[}15{]}}\label{section-148}

If \(\pi\) is written as a product of disjoint cycles, obtain
\(\pi^{-1}\) by reversing the direction of each cycle (i.e., read each
cycle backwards). Equivalently,

\[
\pi=(a_1\,a_2\,\dots\,a_k)\;\;\Longrightarrow\;\; 
\pi^{-1}=(a_1\,a_k\,a_{k-1}\,\dots\,a_2)\,,
\]

and do this for every disjoint cycle. Since disjoint cycles commute,
their order in the product does not matter.

Singleton cycles \((a)\) are unchanged.

If you want the canonical cycle form of §1.3.3 (smallest element first
in each cycle and cycles ordered by decreasing first element), simply
rotate each reversed cycle to put its smallest element first and then
reorder the cycles accordingly.

By the meaning of a cycle,

\[
\pi(a_i)=a_{i+1}\quad (1\le i<k),\qquad \pi(a_k)=a_1.
\]

The inverse must undo these maps, so for every \(i\),

\[
\pi^{-1}(a_{i+1})=a_i\quad (1\le i<k),\qquad \pi^{-1}(a_1)=a_k,
\]

which is exactly the effect of the reversed cycle
\((a_1\,a_k\,a_{k-1}\,\dots\,a_2)\). Applying this argument to each
disjoint cycle gives the stated rule.

Examples:

\begin{enumerate}
\def\labelenumi{\arabic{enumi}.}
\item
  \$ \pi=(1,4,3)(2,5) \Rightarrow \pi\^{}\{-1\}=(1,3,4)(2,5).\$
  (Optionally in canonical form: \((1\,3\,4)(2\,5)\) is already
  canonical.)
\item
  \$ \pi=(3,1,6)(5,4) \Rightarrow \pi\^{}\{-1\}=(3,6,1)(5,4).\$
  Canonicalize if desired: rotate to smallest-first in each cycle to get
  \((1\,3\,6)(4\,5)\), then order cycles by decreasing first element.
\end{enumerate}

\section{12 {[}M27{]}}\label{m27-5}

\begin{enumerate}
\def\labelenumi{(\alph{enumi})}
\tightlist
\item
  Prove the mapping \(y \equiv m x \pmod{mn-1}\) (with
  \(0\le x<N=mn-1\))
\end{enumerate}

Write the current offset in terms of \((i,j)\):

\[
x = n(i-1) + (j-1).
\]

The transposed offset is

\[
y = m(j-1) + (i-1).
\]

Compute \(m x\) and compare to \(y\):

\[
m x = m\bigl(n(i-1) + (j-1)\bigr) = mn(i-1) + m(j-1).
\]

Subtract:

\[
y - m x
= \bigl(m(j-1) + (i-1)\bigr) - \bigl(mn(i-1) + m(j-1)\bigr)
= (i-1) - mn(i-1)
= -(mn-1)\,(i-1).
\]

Hence

\[
y \equiv m x \pmod{mn-1}.
\]

Now check that when we restrict to \(0\le x<N=mn-1\), the corresponding
\((i,j)\) cannot be the bottom-right cell \((m,n)\). Therefore

\[
0 \le y = m(j-1) + (i-1) \le m(n-1)+(m-1)=mn-2 = N-1,
\]

so \(y\) is already the standard residue of \(mx \bmod (mn-1)\). Two
fixed points deserve mention:

\begin{itemize}
\tightlist
\item
  \(x=0\) (cell \((1,1)\)) maps to \(y=0\).
\item
  \(x=mn-1\) (cell \((m,n)\)) would also map to itself, but it lies
  outside the domain \(0\le x<N\) used by the modular rule---this is why
  \(N=mn-1\) is the modulus: it isolates those two fixed points at the
  corners.
\end{itemize}

This completes the proof requested in (a).

\begin{enumerate}
\def\labelenumi{(\alph{enumi})}
\setcounter{enumi}{1}
\tightlist
\item
  In-place methods to perform the transposition
\end{enumerate}

The linear index permutation is:

\[
\pi(x) \;=\; 
\begin{cases}
0, & x=0,\\[2pt]
(mx)\bmod (mn-1), & 1\le x\le mn-2,\\[2pt]
mn-1, & x=mn-1.
\end{cases}
\]

Since \(\gcd(m, mn-1)=1\), \(x\mapsto mx \bmod (mn-1)\) is a permutation
of \(\{0,1,\dots,mn-2\}\). In place, we move elements along the disjoint
cycles of \(\pi\).

Method 1 --- Cycle-leader with a visited bitmap (simple and reliable)

Use one temporary register to rotate each cycle:

\begin{verbatim}
N = m*n - 1
mark[0..N-1] = false
for s in 1..N-1:
    if not mark[s]:
        t = s
        tmp = A[t]
        while True:
            u = (m * t) % N
            if u == s:
                A[t] = tmp
                break
            A[t] = A[u]
            mark[t] = true
            t = u
        mark[t] = true
\end{verbatim}

\emph{Correctness.} Each inner loop performs a cyclic rotation along one
cycle of \(\pi\)-, so all elements land in their transposed locations. 0
and \(mn-1\) are fixed and can be skipped.

\emph{Cost.} Exactly one write per element in \(\{1,\dots,mn-2\}\), plus
one extra write per cycle to place \texttt{tmp}. Time \(\Theta(mn)\),
space \(O(mn)\) bits for \texttt{mark} (or a compact bitset).

Method 2 --- Pure in-place cycle generation via divisors of \(N\) (no
mark array)

We can enumerate cycles deterministically using number theory:

\begin{itemize}
\tightlist
\item
  Let \(N=mn-1\). For any starting index \(x\) with \(d=\gcd(x,N)\),
  write \(x=d\cdot t\) and \(M=N/d\).
\item
  The cycle containing \(x\) has length equal to the order of \(m\)
  modulo \(M\), i.e., the smallest \(r\ge1\) with
  \(m^r\equiv1\pmod{M}\).
\item
  One may generate this cycle as \(x_k = d\cdot (t\,m^k \bmod M)\),
  \(k=0,1,\dots,r-1\).
\end{itemize}

By looping over divisors \(d\mid N\) and a set of coset representatives
\(t\) modulo \(M=N/d\), you can traverse each cycle exactly once. This
yields the same cycle-leader rotations as Method 1 but uses only
\(O(1)\) extra words (plus arithmetic on \(\bmod M\)). It avoids the
mark array at the price of a slightly more intricate control structure
(needs orders modulo \(M\) and coset stepping).

Method 3 --- If extra workspace is allowed (cache-friendly blocking)

If ``on itself'' is relaxed to allow a small scratch buffer, a fast
practical approach is tiling:

\begin{enumerate}
\def\labelenumi{\arabic{enumi}.}
\tightlist
\item
  Partition the matrix into \(B\times B\) tiles.
\item
  For off-diagonal tile pairs \((p,q)\) and \((q,p)\), swap tiles and
  transpose each tile into the other.
\item
  For diagonal tiles, transpose in place (they're square).
\end{enumerate}

This dramatically improves locality on modern caches. It's not strictly
in-place unless \(B=1\), but with a small temporary \(B\times B\) buffer
it runs faster in practice.

Practical notes

\begin{itemize}
\tightlist
\item
  The only elements not governed by the modulus rule are the two
  corners; both are fixed points.
\item
  The permutation \(x\mapsto mx\bmod(mn-1)\) partitions
  \(\{1,\dots,mn-2\}\) into cycles of lengths that divide the order of
  \(m\) modulo \(N/d\) (where \(d=\gcd(x,N)\)). That structure can be
  exploited for pure in-place implementations (Method 2).
\item
  For production code, Method 1 (with a bitset) is simple and hard to
  get wrong; Method 2 is elegant and truly \(O(1)\) extra space.
\end{itemize}

\section{13 {[}M24{]}}\label{m24-15}

Restating Algorithm J (briefly)

\begin{itemize}
\tightlist
\item
  J1. For all \(k\), set \(X[k]\gets -X[k]\). Set \(m\gets n\).
\item
  J2. Set \(j\gets m\).
\item
  J3. Set \(i\gets X[j]\). If \(i>0\) set \(j\gets i\) and repeat J3.
\item
  J4. (Now \(i\le 0\).) Let \(a=-i\). Do \(X[j]\gets X[a]\) and
  \(X[a]\gets m\).
\item
  J5. Decrease \(m\) by \(1\); if \(m>0\) go to J2; else halt.
\end{itemize}

Intuition: after J1 every entry encodes the \emph{forward} edge
\(k\to \pi(k)\) as the negative number \(-\pi(k)\). Step J4 ``splices
out'' one forward edge and deposits a positive number \(m\) somewhere;
those positive numbers are the breadcrumbs from which the inverse will
emerge.

Invariants

We prove correctness by maintaining three simple invariants throughout
the loop on \(m\).

Let \(\mathcal{P}\) be the set of indices whose entries are currently
positive. Define

\[
\text{pos}(t)\ \text{means}\ \exists\,y\ (X[y]=t>0).
\]

I1 (partition of information). For every index \(y\),

\begin{itemize}
\tightlist
\item
  either \(X[y]\) is negative and equals \(-u\) for some \(u\)
  (interpreted as a \emph{still-present} forward link \(y\to u\));
\item
  or \(X[y]\) is positive and equals \(t\) (interpreted as a
  \emph{partial inverse} fact ``\(y\) currently points back to \(t\)'').
\end{itemize}

I2 (what positives mean). For each positive value \(t\) that currently
appears \emph{somewhere} in \(X\), that occurrence is at an index of the
form \(\pi(s)\) for some \(s\) on the cycle containing \(t\). (Thus
positives always sit at \emph{images} under \(\pi\).)

I3 (progress by \(m\)). After finishing iteration value \(m\) (i.e.,
just after J5),

\begin{itemize}
\tightlist
\item
  every integer \(t>m\) appears exactly once as a positive value in the
  array,
\item
  every integer \(\le m\) appears only with negative sign (i.e., it has
  not yet been ``deposited'' positively).
\end{itemize}

These invariants hold initially after J1 (vacuously for I2--I3). We show
they are preserved by one pass of J2--J5 and that they imply correctness
at the end.

Why J3 terminates (and what it finds)

During a fixed outer iteration (fixed \(m\)), J3 repeatedly follows
\emph{positive} pointers: \(j\gets X[j]\) while \(X[j]>0\). Because of
I2 there are no cycles made solely of positive pointers (they are
``back-links'' placed on images, not a disjoint permutation by
themselves), so this walk must end at some index \(j^\star\) with
\(X[j^\star]<0\). Thus J3 always halts and yields a \(j^\star\) whose
entry still encodes a forward edge. (The text even remarks that validity
is ``less obvious'' and leaves it to the reader; this is the key point.
)

The splice step (J4) preserves the invariants and adds one correct
inverse fact

At \(j^\star\) we have \(i=X[j^\star]\le 0\), so let
\(a=-i=\pi(j^\star)\). Step J4 does:

\[
\underbrace{X[j^\star]\gets X[a]}_{\text{bypass } j^\star\to a\text{ to keep forward info}},\qquad
\underbrace{X[a]\gets m}_{\text{deposit }m\text{ at the image }a=\pi(j^\star)}.
\]

\begin{itemize}
\tightlist
\item
  The assignment \(X[a]\gets m>0\) places a positive \(m\) exactly at an
  image index \(a=\pi(j^\star)\), so I2 holds and the set of positives
  grows by at most one.
\item
  The assignment \(X[j^\star]\gets X[a]\) \emph{bypasses} the removed
  forward edge \(j^\star\to a\): if \(X[a]=-b\) previously encoded
  \(a\to b=\pi(a)\), then now \(X[j^\star]\) becomes \(-b\),
  i.e.~\(j^\star\to b\). Thus the remaining forward structure is
  preserved, so I1 continues to hold.
\item
  By construction, the new positive value is exactly the current \(m\),
  and all previously placed positive values were \(>m\) (inductive
  hypothesis I3). Hence after J5, I3 is still true with \(m\) decreased
  by one.
\end{itemize}

In short, every pass of J2--J5 adds exactly one positive \(m\) at
\emph{some} image position and keeps everything else consistent.

Per-cycle correctness

Fix a cycle \(C=(a_0,a_1,\dots,a_{\ell-1})\) of \(\pi\), where
\(\pi(a_r)=a_{r+1}\) (indices mod \(\ell\)). Consider the outer
iterations in the numerical order \(m=n,n-1,\dots,1\), and look only at
those passes where \(m\in C\). We prove by induction on these passes
that:

\begin{enumerate}
\def\labelenumi{(\Alph{enumi})}
\setcounter{enumi}{2}
\tightlist
\item
  After the pass with parameter \(m\in C\), the array has, for this
  cycle,
\end{enumerate}

\[
X[a_{r+1}] = a_r\quad\text{for every }a_r\in C\text{ that has already been processed (i.e., appeared as some earlier }m),
\]

and for every unprocessed \(a_s\in C\) its forward link is still present
somewhere as a negative value (possibly ``bypassed'' a few times by
earlier splices, but still represented as a negative successor), so a
future pass can continue.

\begin{itemize}
\item
  Base. Let \(m\) be the first value of \(C\) encountered. Then \(j=m\)
  immediately has \(X[j]<0\), so J4 writes \(X[\pi(m)]=m\) and replaces
  \(X[m]\) by the (negative) successor of \(\pi(m)\). Thus property (C)
  holds initially: the edge \(\pi(m)\leftarrow m\) of the inverse is now
  correct.
\item
  Step. Suppose (C) holds after finishing some earlier value from \(C\).
  Consider the next \(m\in C\). The walk of J3 only steps along
  already-placed positives on this same cycle and must stop at a
  \(j^\star\in C\) whose entry is still negative. The splice then
  deposits \(m\) exactly at \(\pi(j^\star)\in C\) and preserves all
  previously correct edges; thus (C) extends by one more correct inverse
  edge and remains true.
\end{itemize}

Since each pass adds exactly one correct inverse edge on the cycle and
never destroys earlier ones, when all members of \(C\) have been
processed we end with

\[
X[a_{r+1}] = a_r\qquad(r=0,1,\dots,\ell-1),
\]

i.e., the full inverse on that cycle. Because cycles are disjoint, the
statement holds globally: for every \(y\) we have \(X[y]=\pi^{-1}(y)\)
at termination.

Termination and final state

The outer loop runs exactly \(n\) times (values of \(m\)). Every
iteration performs one splice; there is no possibility of an infinite
inner loop because, as argued above, J3 always reaches a negative entry.
Therefore the algorithm halts with all entries positive. By the
per-cycle argument, these positives are precisely the inverse mapping.
Hence Algorithm J is valid. (The text explicitly notes Algorithm J has
the same effect as Algorithm I, whose goal is to replace \(X\) by its
inverse. )

\emph{Context note.} The book later remarks that although Algorithm J is
``deucedly clever,'' its average running time is \(\Theta(n\log n)\)
versus \(\Theta(n)\) for Algorithm I; the exercise here, however, is
only to prove correctness.

\section{14 {[}M34{]}}\label{m34}

What \(A\) counts (intuition)

In Algorithm J we repeatedly do:

\begin{itemize}
\tightlist
\item
  set \(j \leftarrow m\),
\item
  then chase \(j \leftarrow X[j]\) while \(X[j] > 0\) (step J3),
\item
  stop when we hit a negative entry and perform J4/J5.
\end{itemize}

Thus, \(A\) is the total number of times step J3 keeps following a
\emph{positive} entry across the whole run---the total length of all
these ``positive-chase'' segments.

A precise combinatorial model for \(A\)

Two simple facts (both standard in the discussion around Algorithm J)
let us translate \(A\) to an elementary permutation statistic:

\begin{enumerate}
\def\labelenumi{\arabic{enumi}.}
\tightlist
\item
  State of the array at J2. At the beginning of each J2, one can prove
  by induction (Exercise 13) that
\end{enumerate}

\begin{itemize}
\tightlist
\item
  \(X[i]=+j\) iff \(j>m\) and \(j\) maps to \(i\) under \(\pi\);
\item
  \(X[i]=-j\) otherwise, where the sign encodes how far \(\pi\) has been
  ``unwound.'' This explains why J3 walks over \emph{already-inverted}
  parts (positive entries).
\end{itemize}

\begin{enumerate}
\def\labelenumi{\arabic{enumi}.}
\setcounter{enumi}{1}
\tightlist
\item
  Linearized canonical form. Write the inverse permutation \(\pi^{-1}\)
  in canonical cycle form and then drop parentheses to get a linear
  sequence of length \(n\). If, for each position, you count the number
  of consecutive elements immediately to its right that are larger than
  the current element, then
\end{enumerate}

\[
A-N=\text{(total of those counts over all positions)}.
\]

Equivalently, if you draw the customary ``dot diagram'' under the
linearized \(\pi^{-1}\), \(A\) is exactly the total number of dots.
(This is the official reduction used in the book's answer.)

The key observation now is: the number of dots below the first \(k\)
elements equals the number of right-to-left minima in that prefix
(scanned left-to-right, a position is a right-to-left minimum if its
value is smaller than all earlier values in the prefix when reading from
the right). Hence,

\[
A \;=\; \sum_{k=1}^{n} \bigl(\#\text{ right-to-left minima among the first } k \text{ elements}\bigr).
\]

Take expectation

In a random permutation, the probability that position \(r\) is a
(right-to-left) record minimum among the first \(r\) items is \(1/r\).
Therefore the expected number of such minima in the first \(k\) items is
the harmonic number \(H_k=\sum_{r=1}^{k} 1/r\). Summing over \(k\) gives

\[
\mathbb{E}[A] \;=\; \sum_{k=1}^{n} H_k \;=\; (n+1)H_n - n.
\]

(This uses the well-known identity \(\sum_{k=1}^{n} H_k=(n+1)H_n-n\).)
The appearance of \(H_k\) is the same ``records'' phenomenon used
earlier in §1.3.3.

Final answer

\[
\boxed{\;\mathbb{E}[A] \;=\; \sum_{k=1}^{n} H_k \;=\; (n+1)H_n - n\;}
\]

and, using \(H_n=\ln n+\gamma+o(1)\),

\[
\mathbb{E}[A]=n\ln n + (\gamma-1)n + \ln n + \gamma + o(1).
\]

This aligns with the text's remark that Algorithm J's average running
time is \(\Theta(n\ln n)\), because \(A\) dominates its cost.

\section{15 {[}M12{]}}\label{m12-1}

\begin{itemize}
\tightlist
\item
  Linear (one-line) notation of a permutation \(\pi\) on
  \(\{1,\dots,n\}\) is the sequence \(\pi(1)\ \pi(2)\ \cdots\ \pi(n)\).
\item
  Canonical cycle form: write each cycle with its smallest element
  first, and order the cycles by decreasing first element; then (for the
  correspondence we use) drop the parentheses to get a length-\(n\)
  sequence. Knuth shows that from this sequence you can uniquely put the
  parentheses back by inserting a left parenthesis before each
  left-to-right minimum.
\end{itemize}

Call the mapping ``drop-canonical-parentheses'' \(C(\pi)\), which takes
\(\pi\) to the sequence obtained by (canonical cycles) → (remove
parentheses).

We are asked for permutations \(\pi\) such that

\[
C(\pi) \text{ (a length-}n\text{ sequence) } = \text{ linear sequence } \pi(1)\pi(2)\cdots\pi(n).
\]

Key observations

Consider the first symbol of the linear notation, \(\pi(1)=m\).

\begin{enumerate}
\def\labelenumi{\arabic{enumi}.}
\item
  If \(m=1\): then \(1\) is a fixed point. In canonical cycle form there
  is a 1-cycle \((1)\). Since cycles are ordered by decreasing first
  element, \((1)\) is the last cycle. When we drop parentheses, that
  trailing cycle contributes a final ``1''. Therefore the last character
  of \(C(\pi)\) is 1.
\item
  If \(m>1\): then in the cycle containing 1 we have \((1\ m\ \cdots)\),
  because \(\pi(1)=m\). Again this cycle (starting with 1) is last in
  canonical order, so in \(C(\pi)\) the substring ``\ldots{} 1 m
  \ldots{}'' appears (the ``1'' immediately followed by ``m'').
\end{enumerate}

Consequences

Now compare \(C(\pi)\) with the linear sequence:

\begin{itemize}
\tightlist
\item
  If the linear sequence starts with 1, but \(C(\pi)\) ends with 1, the
  only way the two sequences can be identical is when there is only one
  position---i.e., \(n=1\) (trivial one-element permutation). (The empty
  permutation \(n=0\) also vacuously works.)
\item
  If the linear sequence starts with \(m>1\), then \(C(\pi)\) contains
  the adjacent pair ``\(1\ m\)''. For the two sequences to be identical,
  the linear sequence would need to start with ``\ldots{} \(1\ m\)
  \ldots{}'', which is impossible because its first symbol is \(m\) and
  there is no preceding 1.
\end{itemize}

Therefore, there is no nontrivial permutation with this property. The
only solutions are the trivial cases: the permutation on one element
(and the empty permutation). This matches Knuth's official answer.

Tiny sanity checks

\begin{itemize}
\tightlist
\item
  \(n=1\): linear ``1''; canonical cycles ``(1)''; drop → ``1''. Equal ✓
\item
  \(n=2\): linear ``1 2'' gives canonical ``(1)(2)'' → ``2 1'', not
  equal.
\item
  \(n=2\): linear ``2 1'' gives canonical ``(1 2)'' → ``1 2'', not
  equal.
\end{itemize}

Conclusion: Apart from \(n=0\) or \(1\), no permutation has identical
linear notation and canonical-cycle-without-parentheses notation.

\section{16 {[}M15{]}}\label{m15-14}

First, recall the book's ``canonical cycle form'': (a) write every
1-cycle explicitly, (b) rotate each cycle so its smallest element comes
first, and (c) order the cycles by decreasing value of their first
element. Then ``remove parentheses'' simply concatenates the cycles'
entries into a length-n string, which we now read as the next
permutation in linear notation.

The task is: start from 1324 (linear form), convert to canonical cycle
form, drop the parentheses to get a new linear form, and repeat until
you return to 1324.

Step-by-step

For each row: ``linear → canonical cycles → next linear (after dropping
parentheses)''.

\begin{enumerate}
\def\labelenumi{\arabic{enumi}.}
\tightlist
\item
  1324 → (4)(2 3)(1) → 4231
\item
  4231 → (3)(2)(1 4) → 3214
\item
  3214 → (4)(2)(1 3) → 4213
\item
  4213 → (2)(1 4 3) → 2143
\item
  2143 → (3 4)(1 2) → 3412
\item
  3412 → (2 4)(1 3) → 2413
\item
  2413 → (1 2 4 3) → 1243
\item
  1243 → (3 4)(2)(1) → 3421
\item
  3421 → (1 3 2 4) → 1324
\end{enumerate}

You are now back at the start.

Answer (the permutations that appear, in order)

1324 → 4231 → 3214 → 4213 → 2143 → 3412 → 2413 → 1243 → 3421 → 1324.

So, beginning from 1324, the ``canonicalize cycles, then drop
parentheses'' transformation runs through a 9-cycle of distinct
permutations (returning to the start on the 10th line).

If you're curious: for n=4 this transformation partitions all 24
permutations into two disjoint orbits; 1324 lies in the orbit of length
9 (the other orbit has length 15).

\section{17 {[}M24{]}}\label{m24-16}

\begin{enumerate}
\def\labelenumi{(\alph{enumi})}
\tightlist
\item
  Average length when you pick a cycle uniformly from all cycles
\end{enumerate}

Fact used from the text: among all permutations of \(n\) elements, the
total number of \(m\)-cycles is

\[
\frac{n!}{m}\qquad(1\le m\le n),
\]

hence the total number of cycles of all lengths is

\[
\sum_{m=1}^n \frac{n!}{m}=n!\,H_n,
\]

so the average number of cycles in a random permutation is \(H_n\).

If we now choose one cycle uniformly from these \(n!\,H_n\) cycles, the
probability that it has length \(m\) is

\[
p_m=\frac{\#\text{ of }m\text{-cycles}}{\#\text{ of all cycles}}
=\frac{(n!/m)}{n!\,H_n}=\frac{1}{mH_n}.
\]

Therefore the average cycle length under this sampling is

\[
\mathbb{E}[L]
=\sum_{m=1}^n m\cdot p_m
=\sum_{m=1}^n m\cdot\frac{1}{mH_n}
=\frac{n}{H_n}.
\]

(Exercise statement context: the problem explicitly sets up this
sampling over the \(n!\,H_n\) cycles.)

\begin{enumerate}
\def\labelenumi{(\alph{enumi})}
\setcounter{enumi}{1}
\tightlist
\item
  Length of the cycle containing a random element \(k\) in a random
  permutation
\end{enumerate}

Count permutations in which a fixed element \(k\) lies in an
\(m\)-cycle:

\begin{itemize}
\tightlist
\item
  Choose the other \(m-1\) elements that share \(k\)'s cycle:
  \(\binom{n-1}{m-1}\).
\item
  Arrange them with \(k\) on a cycle: \((m-1)!\) cyclic orders.
\item
  Permute the remaining \(n-m\) elements arbitrarily: \((n-m)!\).
\end{itemize}

Thus the number of permutations with \(k\) in an \(m\)-cycle is

\[
\binom{n-1}{m-1}(m-1)!(n-m)!=(n-1)!,
\]

independent of \(m\). Since there are \(n!\) permutations total, we get

\[
\Pr\{\text{\(k\) is in an \(m\)-cycle}\}=\frac{(n-1)!}{n!}=\frac{1}{n}\qquad(1\le m\le n).
\]

So the cycle length of a uniformly chosen element is uniform on
\(\{1,2,\dots,n\}\). Its average is therefore

\[
\mathbb{E}[L_k]=\frac{1}{n}\sum_{m=1}^n m=\frac{n+1}{2}.
\]

(Again, the exercise frames this as: pick a random \(k\) and a random
permutation; determine the probability for an \(m\)-cycle and the
corresponding average.)

Summary of results

\begin{itemize}
\item
  \begin{enumerate}
  \def\labelenumi{(\alph{enumi})}
  \tightlist
  \item
    Picking a cycle uniformly from among all cycles across all
    permutations:
  \end{enumerate}

  \[
  \boxed{\ \mathbb{E}[L]=\dfrac{n}{H_n}\ }
  \]
\item
  \begin{enumerate}
  \def\labelenumi{(\alph{enumi})}
  \setcounter{enumi}{1}
  \tightlist
  \item
    Picking a uniformly random element \(k\) in a uniformly random
    permutation:
  \end{enumerate}

  \[
  \boxed{\ \Pr\{L_k=m\}=\dfrac{1}{n}\ \text{ for }1\le m\le n,\quad \mathbb{E}[L_k]=\dfrac{n+1}{2}\ }
  \]
\end{itemize}

\section{18 {[}M27{]}}\label{m27-6}

Let \(X_m\) be the number of cycles of length \(m\) in a uniformly
random permutation of \(\{1,\dots,n\}\). We want (i) the exact
probability \(p_{nkm}=\Pr[X_m=k]\); (ii) the probability generating
function \(G_{n,m}(z)=\mathbb{E}[z^{X_m}]\); and (iii) the mean and
standard deviation of \(X_m\).

Knuth's cycle-index/exponential formula for permutations says that if we
``mark'' cycles of a given length with a variable, the exponential
generating function (EGF) multiplies by \(\exp((w-1)x^m/m)\).
Specializing \(w=z\) for \(m\)-cycles and keeping all other lengths
unmarked gives the bivariate EGF

\[
\mathcal{E}(z,x)
=\exp\!\Big(\sum_{j\ge1}\tfrac{x^j}{j}\Big)\cdot \exp\!\Big(\tfrac{(z-1)x^m}{m}\Big)
=\frac{1}{1-x}\;\exp\!\Big(\tfrac{(z-1)x^m}{m}\Big).
\]

Extracting the \(x^n\) coefficient and dividing by the total number
\(n!\) (equivalently, by \([x^n](1-x)^{-1}=1\)) yields the PGF for
\(X_m\) at fixed \(n\):

\[
\boxed{\;G_{n,m}(z)=\sum_{r=0}^{\lfloor n/m\rfloor}\frac{(z-1)^r}{m^{\,r}\,r!}\;}
\quad\text{(for all \(n\ge0\)).}
\]

For \(m=1\) this reduces exactly to Knuth's Eq. (27).

Exact distribution \(p_{nkm}\)

Read off the \(z^k\)-coefficient of \(G_{n,m}(z)\). Using
\((z-1)^r=\sum_{k=0}^{r}\binom{r}{k}(-1)^{r-k}z^k\),

\[
\boxed{\;
p_{nkm}
=\sum_{r=k}^{\lfloor n/m\rfloor}\binom{r}{k}\frac{(-1)^{\,r-k}}{m^{\,r}\,r!}
=\frac{1}{m^{\,k}\,k!}\sum_{t=0}^{\lfloor n/m\rfloor-k}\frac{(-1)^{\,t}}{m^{\,t}\,t!}\;},
\qquad 0\le k\le\Big\lfloor\frac{n}{m}\Big\rfloor .
\]

(When \(n<m\), \(p_{n0m}=1\) and \(p_{nkm}=0\) for \(k>0\).)

Mean and variance (hence standard deviation)

For a PGF \(G(z)\), \(\mathbb{E}X=G'(1)\) and
\(\mathrm{Var}(X)=G''(1)+G'(1)-[G'(1)]^2\). Differentiate the closed
form of \(G_{n,m}\):

\begin{itemize}
\item
  \(G'_{n,m}(1)=\dfrac{1}{m}\) (exact whenever \(n\ge m\); otherwise
  \(0\)). So
  \(\boxed{\mathbb{E}[X_m]=1/m\ \text{for}\ n\ge m,\ \text{else }0.}\)
\item
  \(G''_{n,m}(1)=\begin{cases}\dfrac{1}{m^2},& n\ge 2m,\\[4pt] 0,& n<2m.\end{cases}\)
\end{itemize}

Hence

\[
\mathrm{Var}(X_m)=
\begin{cases}
\dfrac{1}{m},& n\ge 2m,\\[6pt]
\dfrac{1}{m}-\dfrac{1}{m^2}=\dfrac{m-1}{m^2},& m\le n<2m,\\[6pt]
0,& n<m.
\end{cases}
\quad\Rightarrow\quad
\boxed{\ \sigma(X_m)=
\begin{cases}
\dfrac{1}{\sqrt{m}},& n\ge 2m,\\[6pt]
\sqrt{\dfrac{1}{m}-\dfrac{1}{m^2}},& m\le n<2m,\\[6pt]
0,& n<m.
\end{cases}}
\]

(A quick combinatorial cross-check for the mean: For each \(m\)-subset
\(S\), the event ``\(S\) forms a cycle'' has probability
\((m-1)!(n-m)!/n!\); summing over \(\binom{n}{m}\) choices gives
\(\binom{n}{m}(m-1)!(n-m)!/n!=1/m\).)

Asymptotics

Let \(n\to\infty\) with \(m\) fixed. Then the truncated exponential in
\(p_{nkm}\) tends to \(e^{-1/m}\), so

\[
p_{nkm}\ \longrightarrow\ e^{-1/m}\frac{1}{k!\,m^{\,k}},
\]

i.e., \(X_m\) converges in distribution to \(\mathrm{Poisson}(1/m)\),
with mean and variance both \(1/m\) --- consistent with the exact
formulas above when \(n\ge2m\). (The m=1 case was treated in the text;
this generalizes it verbatim.)

\section{19 {[}HM21{]}}\label{hm21-9}

Let \(D_n\) (Knuth's \(P_n^{\,0}\)) denote the number of derangements of
\(n\) distinct objects. By inclusion--exclusion (derived just above the
exercise), we have the well-known alternating-sum formula

\[
D_n
= n!\sum_{k=0}^{n}\frac{(-1)^k}{k!}.
\]

(This is exactly Knuth's Eq. (25) specialized to the ``no singleton
cycles'' case.)

Exercise 19 asks you to show that \(D_n\) equals the nearest integer to
\(n!/e\) for every \(n\ge 1\). (Exercise statement.)

Set

\[
S_n=\sum_{k=0}^{n}\frac{(-1)^k}{k!},\qquad
S_\infty=\sum_{k=0}^{\infty}\frac{(-1)^k}{k!}=e^{-1}.
\]

Then

\[
\frac{n!}{e}-D_n
= n!\bigl(S_\infty-S_n\bigr)
= n!\,R_n,
\quad\text{where }R_n:=\sum_{k=n+1}^{\infty}\frac{(-1)^k}{k!}.
\]

Because \(\{1/k!\}\) is strictly decreasing to \(0\), the tail \(R_n\)
of this alternating series has the standard sharp bound

\[
|R_n| \;<\; \frac{1}{(n+1)!}
\]

and has sign \((-1)^{n+1}\). Hence

\[
\biggl|\frac{n!}{e}-D_n\biggr|
= n!\,|R_n|
\;<\; \frac{n!}{(n+1)!}
=\frac{1}{n+1}
\;\le\; \frac12\quad (n\ge 1).
\]

The inequality is strict even when \(n=1\) (because the alternating
remainder is strictly smaller than the first omitted term).

Therefore the difference between \(n!/e\) and the integer \(D_n\) has
absolute value \(<\tfrac12\). This means \(D_n\) is exactly the nearest
integer to \(n!/e\) for all \(n\ge1\), as required.

\section{20 {[}M20{]}}\label{m20-32}

The same permutation can be written in many \emph{cycle notations}
because

\begin{enumerate}
\def\labelenumi{\arabic{enumi}.}
\tightlist
\item
  disjoint cycles may be listed in any order; and
\item
  each k-cycle can be \emph{rotated} to start from any of its k elements
  (but not reversed---reversal gives the inverse cycle).
\end{enumerate}

Let the permutation have

\begin{itemize}
\tightlist
\item
  α₁ one-cycles, α₂ two-cycles, α₃ three-cycles, \ldots,
\item
  and let r = Σₖ αₖ be the total number of cycles (including
  singletons).
\end{itemize}

Because the cycles are disjoint and distinct, you can order the r cycles
in r! ways. Independently, each k-cycle has k rotational variants, so
all k-cycles together contribute k\^{}\{αₖ\} ways.

Multiplying these independent choices gives the total number of written
forms:

\[
\boxed{\text{\# of cycle notations} \;=\; r!\,\prod_{k\ge 1} k^{\alpha_k}}
\qquad\text{where } r=\sum_{k\ge 1}\alpha_k .
\]

Notes and checks

\begin{itemize}
\tightlist
\item
  Singletons contribute a factor \(1^{\alpha_1}=1\); they only affect
  the \(r!\) ordering factor.
\item
  If (as in Exercise 5) you suppress singletons, replace \(r\) by
  \(r-\alpha_1\) and drop the \(k=1\) factor:
  \((r-\alpha_1)!\,\prod_{k\ge 2} k^{\alpha_k}\).
\end{itemize}

Example If \(\pi=(1\,4\,3)(2\,5)(6)(7)\), then
\((\alpha_1,\alpha_2,\alpha_3)=(2,1,1)\) and \(r=4\). The number of ways
to write \(\pi\) in cycle form (with singletons shown) is

\[
4!\cdot 3^1\cdot 2^1\cdot 1^2=24\cdot 6=144.
\]

\section{21 {[}M22{]}}\label{m22-18}

\[
\boxed{\;
P\bigl(n;\alpha_1,\alpha_2,\ldots\bigr)
=
\begin{cases}
\dfrac{1}{\displaystyle\prod_{m\ge1} m^{\alpha_m}\,\alpha_m!}, & \text{if } \sum_{m\ge1} m\,\alpha_m = n,\\[1.2em]
0, & \text{otherwise.}
\end{cases}
\;}
\]

Derivation (combinatorial counting)

\begin{enumerate}
\def\labelenumi{\arabic{enumi}.}
\tightlist
\item
  Count permutations with a fixed cycle type. Let's count how many
  permutations of \(\{1,\dots,n\}\) have exactly \(\alpha_m\) cycles of
  length \(m\) for each \(m\). A standard result from the
  cycle-index/Polya theory (and easy to derive directly) is:
\end{enumerate}

\[
\#\{\text{permutations with this cycle type}\}
\;=\;
\frac{n!}{\displaystyle\prod_{m\ge1} m^{\alpha_m}\,\alpha_m!}.
\]

Intuition: start with all \(n!\) linear arrangements; turning those into
cycles overcounts by a factor \(m\) for each \(m\)-cycle (rotations of a
cycle don't change it), and cycles of the same length are
indistinguishable so divide by \(\alpha_m!\) as well.

\begin{enumerate}
\def\labelenumi{\arabic{enumi}.}
\setcounter{enumi}{1}
\tightlist
\item
  Convert to probability. All \(n!\) permutations are equally likely, so
  divide the count above by \(n!\):
\end{enumerate}

\[
P\bigl(n;\alpha_1,\alpha_2,\ldots\bigr)
= \frac{1}{\displaystyle\prod_{m\ge1} m^{\alpha_m}\,\alpha_m!},
\]

provided the cycle counts fit \(n\) (otherwise the probability is
\(0\)).

\begin{enumerate}
\def\labelenumi{\arabic{enumi}.}
\setcounter{enumi}{2}
\tightlist
\item
  Sanity checks.
\end{enumerate}

\begin{itemize}
\item
  Sum to 1. Summing the counts
  \(\frac{n!}{\prod m^{\alpha_m}\alpha_m!}\) over all \((\alpha_m)\)
  with \(\sum m\alpha_m=n\) gives \(n!\); dividing by \(n!\) yields
  total probability \(1\).
\item
  Examples.

  \begin{itemize}
  \tightlist
  \item
    \(n=4\), two 2-cycles: \((\alpha_2=2)\). Count
    \(=\frac{4!}{2^{2}\,2!}=3\); probability
    \(=3/24=1/8=\frac{1}{2^{2}\,2!}\).
  \item
    \(n=4\), one 3-cycle and one fixed point:
    \((\alpha_3=1,\alpha_1=1)\). Count
    \(=\frac{4!}{3^{1}\,1!\cdot1^{1}\,1!}=8\); probability
    \(=8/24=1/3=\frac{1}{3}\).
  \item
    \(n=4\), identity: \((\alpha_1=4)\). Count
    \(=\frac{4!}{1^{4}\,4!}=1\); probability \(=1/24\).
  \end{itemize}
\end{itemize}

Compute

\[
P(n;\alpha_1,\alpha_2,\ldots)=\Pr\!\left(\text{a random }\pi\in S_n\text{ has }\alpha_j\text{ }j\text{-cycles for all }j\right).
\]

Step 1: Count permutations with a prescribed cycle structure

We'll count how many permutations of \(\{1,\dots,n\}\) have cycle type
\((\alpha_1,\alpha_2,\ldots)\).

Think in two stages:

\begin{enumerate}
\def\labelenumi{\arabic{enumi}.}
\tightlist
\item
  Choose the elements for each cycle (ignoring order inside each cycle,
  and ignoring the order among cycles of the same length). Partition the
  \(n\) labeled elements into \(\alpha_1\) blocks of size \(1\),
  \(\alpha_2\) blocks of size \(2\), \(\alpha_3\) blocks of size \(3\),
  \ldots{} . The number of such partitions is
\end{enumerate}

\[
\frac{n!}{\prod_{j\ge 1} (j!)^{\alpha_j}\,\alpha_j!}.
\]

\begin{enumerate}
\def\labelenumi{\arabic{enumi}.}
\setcounter{enumi}{1}
\tightlist
\item
  Within each block, arrange the elements around a cycle. A block of
  size \(j\) can be arranged on a directed cycle in \((j-1)!\) distinct
  ways (there are \(j!\) linear orders but each \(j\)-cycle has \(j\)
  rotational representations). Doing this for all cycles multiplies by
\end{enumerate}

\[
\prod_{j\ge 1} \bigl((j-1)!\bigr)^{\alpha_j}.
\]

Multiplying the two stages gives the total number of permutations with
that cycle type:

\[
\#\{\pi\in S_n:\text{cycle type }(\alpha_j)\}
=\frac{n!}{\prod_{j\ge 1} (j!)^{\alpha_j}\,\alpha_j!}\;\cdot\;\prod_{j\ge 1} \bigl((j-1)!\bigr)^{\alpha_j}
=\frac{n!}{\prod_{j\ge 1} j^{\alpha_j}\,\alpha_j!}.
\]

(Equivalently, one can argue by ``decorating'' cycles with a chosen
starting point and an order among equal-length cycles; each underlying
permutation then corresponds to exactly
\(\prod_j j^{\alpha_j}\alpha_j!\) decorated representations, while there
are \(n!\) decorated representations in total---so the same formula
results.)

This well-known enumeration also appears in standard references on
permutations.

Step 2: Divide by \(n!\) to get the probability

Since every permutation is equally likely, the desired probability is
the count above divided by \(n!\):

\[
\boxed{\displaystyle
P(n;\alpha_1,\alpha_2,\ldots)
= \frac{1}{\prod_{j\ge 1} j^{\alpha_j}\,\alpha_j!}}
\quad\text{(provided } \sum_{j\ge 1} j\,\alpha_j=n\text{; otherwise }P=0\text{).}
\]

Quick consistency checks

\begin{itemize}
\tightlist
\item
  Normalization. Summing the right-hand side over all \((\alpha_j)\)
  with \(\sum j\alpha_j=n\) gives \(1\). This follows from the
  exponential formula:
  \(\displaystyle \sum_{\alpha:\,\sum j\alpha_j=n}\prod_{j\ge 1}\frac{1}{j^{\alpha_j}\alpha_j!} = [x^n]\exp\!\Bigl(\sum_{j\ge 1}\frac{x^j}{j}\Bigr) = [x^n]\frac{1}{1-x}=1.\)
\item
  Example. For \(n=5\) and structure \(1^1 2^2\) (one fixed point and
  two 2-cycles),
  \(\displaystyle P=\frac{1}{1^{1}1!\,2^{2}2!}=\frac{1}{8}.\) Indeed
  there are \(5!/8=15\) such permutations: choose the fixed point in
  \(5\) ways and pair the remaining \(4\) elements into two unordered
  pairs in \(3\) ways.
\end{itemize}

That completes the solution. The key is that the factor
\(j^{\alpha_j}\alpha_j!\) accounts for the symmetries of cycles of
length \(j\): \(j\) rotations per cycle and \(\alpha_j!\) reorderings
among cycles of the same length.

\section{22 {[}HM34{]}}\label{hm34-1}

\begin{enumerate}
\def\labelenumi{(\alph{enumi})}
\tightlist
\item
  Find \(f(w,m,k)\)
\end{enumerate}

We seek independent \(\alpha_m\) with

\[
\Pr(\alpha_m=k)=f(w,m,k)\qquad(m\ge 1,\ k\ge 0),
\]

such that:

\begin{enumerate}
\def\labelenumi{\arabic{enumi}.}
\item
  \(\sum_{k\ge 0} f(w,m,k)=1\) for \(0<w<1\);
\item
  for every nonnegative integer vector \((k_1,k_2,\ldots)\) with
  \(\sum_{m\ge 1} m k_m=n\),
\end{enumerate}

\[
\Pr(\alpha_1=k_1,\alpha_2=k_2,\ldots)=(1-w)\,w^{\,n}\,P\!\left(n;k_1,k_2,\ldots\right).
\]

Choice. Take the \(\alpha_m\) independent Poisson with mean

\[
\lambda_m=\frac{w^m}{m},\qquad 0<w<1.
\]

Thus

\[
f(w,m,k)=e^{-\lambda_m}\frac{\lambda_m^{k}}{k!}
=e^{-\frac{w^m}{m}}\;\frac{1}{k!}\left(\frac{w^m}{m}\right)^k .
\]

Then

\[
\prod_{m\ge 1} f(w,m,\alpha_m)
=\exp\!\Big(-\sum_{m\ge 1}\tfrac{w^m}{m}\Big)\;
\frac{w^{\sum m\alpha_m}}{\prod_m m^{\alpha_m}\alpha_m!}.
\]

Since \(\sum_{m\ge 1}\frac{w^m}{m}=-\ln(1-w)\), the prefactor is
\(e^{\ln(1-w)}=1-w\). Also

\[
P(n;\alpha)=\frac{1}{\prod_m m^{\alpha_m}\alpha_m!}
\quad(\text{when } \sum m\alpha_m=n;\ \text{else }0),
\]

so the joint probability becomes \((1-w)\,w^{n}\,P(n;\alpha)\), exactly
as required. Condition (1) holds because Poisson pmfs are normalized.

\begin{enumerate}
\def\labelenumi{(\alph{enumi})}
\setcounter{enumi}{1}
\tightlist
\item
  Distribution of the total size \(\sum m\alpha_m\)
\end{enumerate}

Let

\[
N=\alpha_1+2\alpha_2+3\alpha_3+\cdots .
\]

Using pgfs,

\[
\mathbb{E}(x^{N})
=\prod_{m\ge 1}\exp\!\Big(\lambda_m(x^{m}-1)\Big)
=\exp\!\Big(\sum_{m\ge 1}\tfrac{w^m}{m}(x^{m}-1)\Big)
=\frac{1-w}{1-wx}.
\]

Therefore

\[
\Pr(N=n)=(1-w)w^n\quad(n=0,1,2,\ldots),
\]

a geometric distribution; hence \(\Pr(N=\infty)=0\).

\begin{enumerate}
\def\labelenumi{(\alph{enumi})}
\setcounter{enumi}{2}
\tightlist
\item
  Transfer principle for expectations
\end{enumerate}

Let \(\phi(\alpha_1,\alpha_2,\ldots)\) be any function depending on the
\(\alpha_m\)'s. Using (a),

\[
\mathbb{E}[\phi]
=\sum_{n\ge 0}(1-w)w^n
\sum_{\sum m\alpha_m=n}\phi(\alpha)\,P(n;\alpha)
=(1-w)\sum_{n\ge 0} w^n\,\phi_n,
\]

where \(\phi_n\) is the average of \(\phi\) over uniform random
permutations of \(n\) (i.e., with the \(\alpha_m\)'s being the cycle
counts of a uniform permutation). Example: If \(\phi(\alpha)=\alpha_1\),
then \(\phi_n=\) average number of fixed points, which equals \(1\) for
all \(n\); our identity reduces to
\((1-w)\sum_{n\ge 0} w^n\cdot 1 = w = \mathbb{E}\alpha_1\) (since
\(\alpha_1\sim \mathrm{Poisson}(w)\)), consistent.

\begin{enumerate}
\def\labelenumi{(\alph{enumi})}
\setcounter{enumi}{3}
\tightlist
\item
  Average number of even cycles in a permutation of \(n\)
\end{enumerate}

Let \(\phi(\alpha)=\alpha_2+\alpha_4+\alpha_6+\cdots\) (the count of
even-length cycles).

Unconditional expectation (under the Poisson model).

\[
\mathbb{E}[\phi]=\sum_{j\ge 1}\mathbb{E}\alpha_{2j}
=\sum_{j\ge 1}\frac{w^{2j}}{2j}
=\tfrac{1}{2}\log\frac{1}{1-w^2}.
\]

Transfer to fixed \(n\). By part (c),

\[
(1-w)\sum_{n\ge 0} w^n\,\phi_n \;=\; \tfrac{1}{2}\log\frac{1}{1-w^2}.
\]

Equivalently,

\[
\sum_{n\ge 0} w^n\,\phi_n \;=\; \frac{1}{1-w}\cdot \frac{1}{2}\log\frac{1}{1-w^2}.
\]

Extracting coefficients gives the clean closed form

\[
\boxed{\;\phi_n=\frac{1}{2}\,H_{\lfloor n/2\rfloor}\;}
\]

where \(H_m=\sum_{k=1}^m \frac{1}{k}\) is the \(m\)-th harmonic number.
(Indeed,
\((1-w)\sum_{n\ge 0} \tfrac12 H_{\lfloor n/2\rfloor} w^n =\tfrac12\log\frac{1}{1-w^2}\).)

\begin{enumerate}
\def\labelenumi{(\alph{enumi})}
\setcounter{enumi}{4}
\tightlist
\item
  Use the method to solve Exercise 18
\end{enumerate}

\begin{quote}
Goal (Ex. 18). For fixed \(m\), find

\[
p_{n,k,m}=\Pr(\alpha_m=k\ \text{in a uniform permutation of }n),
\]

the probability generating function
\(G_{n m}(z)=\sum_{k\ge 0}p_{n,k,m}z^k\), and the mean and standard
deviation of \(\alpha_m\).
\end{quote}

Step 1 (bivariate generating function via Shepp--Lloyd). Under the
Poisson model,

\[
\mathbb{E}\big[z^{\alpha_m}\big]
=\exp\!\Big(\tfrac{w^m}{m}(z-1)\Big).
\]

By (c),

\[
(1-w)\sum_{n\ge 0} w^n\,G_{n m}(z)
=\exp\!\Big(\tfrac{w^m}{m}(z-1)\Big).
\]

Therefore

\[
\sum_{n\ge 0} w^n\,G_{n m}(z)
=\frac{1}{1-w}\exp\!\Big(\tfrac{w^m}{m}(z-1)\Big).
\]

Step 2 (coefficient extraction in \(w\)). Expand the exponential only in
powers of \(w^{m}\):

\[
\exp\!\Big(\tfrac{w^m}{m}(z-1)\Big)
=\sum_{r\ge 0}\frac{(z-1)^r}{m^{r} r!}\,w^{mr}.
\]

Multiplying by \(\frac{1}{1-w}=\sum_{s\ge 0}w^s\) gives

\[
G_{n m}(z)=\sum_{r=0}^{\lfloor n/m\rfloor}\frac{(z-1)^r}{m^{r} r!}.
\]

This is a truncated exponential in \(z-1\). In particular, for \(m=1\)
it reduces to the classic
\(G_{n,1}(z)=\sum_{r=0}^{n}\frac{(z-1)^r}{r!}\).

Step 3 (explicit \(p_{n,k,m}\)). Extract the coefficient of \(z^k\):

\[
\boxed{\;
p_{n,k,m}
=\frac{1}{m^{k}k!}\sum_{s=0}^{\lfloor n/m\rfloor-k}\frac{(-1)^s}{m^{s}s!}
\qquad(0\le k\le \lfloor n/m\rfloor),
\;}
\]

and \(p_{n,k,m}=0\) for \(k>\lfloor n/m\rfloor\). (For \(m=1\):
\(p_{n,k,1}=\frac{1}{k!}\sum_{s=0}^{n-k}\frac{(-1)^s}{s!}\), the
well-known fixed-point formula.)

Step 4 (mean and variance). From \(G_{n m}(z)\),

\[
\mathbb{E}_n[\alpha_m]=G_{n m}'(1)
=\begin{cases}
\dfrac{1}{m}, & n\ge m,\\[4pt]
0, & n<m,
\end{cases}
\]

and the second factorial moment

\[
\mathbb{E}_n\big[\alpha_m(\alpha_m-1)\big]
=G_{n m}''(1)
=\begin{cases}
\dfrac{1}{m^2}, & n\ge 2m,\\[4pt]
0, & n<2m.
\end{cases}
\]

Hence

\[
\mathrm{Var}_n(\alpha_m)
=\mathbb{E}_n[\alpha_m]+\mathbb{E}_n[\alpha_m(\alpha_m-1)]
-\big(\mathbb{E}_n[\alpha_m]\big)^2
=
\begin{cases}
0, & n<m,\\[4pt]
\dfrac{1}{m}-\dfrac{1}{m^2}, & m\le n<2m,\\[6pt]
\dfrac{1}{m}, & n\ge 2m,
\end{cases}
\]

and the standard deviation is
\(\sigma_n(\alpha_m)=\sqrt{\mathrm{Var}_n(\alpha_m)}\). (As
\(n\to\infty\) with fixed \(m\), \(\alpha_m\) tends to
\(\mathrm{Poisson}(1/m)\): mean \(1/m\), variance \(1/m\).)

Summary of key formulas

\begin{itemize}
\item
  Independent model:
  \(\alpha_m\sim\mathrm{Poisson}\!\left(\frac{w^m}{m}\right)\)
  (independent in \(m\)).
\item
  Size distribution: \(\Pr\big(\sum m\alpha_m=n\big)=(1-w)w^n\).
\item
  Transfer:
  \(\displaystyle \mathbb{E}[\phi]=(1-w)\sum_{n\ge 0} w^n \phi_n\).
\item
  Even cycles (in \(S_n\)):
  \(\displaystyle \mathbb{E}[\#\text{even cycles}]=\frac{1}{2}H_{\lfloor n/2\rfloor}\).
\item
  For fixed \(m\):

  \[
  G_{n m}(z)=\sum_{r=0}^{\lfloor n/m\rfloor}\frac{(z-1)^r}{m^{r} r!},\quad
  p_{n,k,m}=\frac{1}{m^{k}k!}\sum_{s=0}^{\lfloor n/m\rfloor-k}\frac{(-1)^s}{m^{s}s!}.
  \]
\end{itemize}

\section{23 {[}HM42{]}}\label{hm42-2}

Let \(L_n\) be the length of the longest cycle of a uniformly random
permutation of \(\{1,\dots,n\}\). Show

\[
\mathbb{E}[L_n]\;=\; \lambda\,n+\tfrac12\,\lambda+o(1),
\]

where \(\lambda\approx 0.6243299885\) is a constant (the Golomb--Dickman
constant). In fact,

\[
\lim_{n\to\infty}\bigl(\mathbb{E}[L_n]-\lambda n-\tfrac12\lambda\bigr)=0.
\]

Step 1 --- Counting permutations with no ``long'' cycles (exponential
generating functions)

For a fixed cut-off \(m\), let \(A_{n,m}\) be the number of permutations
of \([n]\) all of whose cycle lengths are \(\le m\). For labelled
structures of permutations, the exponential generating function (EGF)
with allowed cycle lengths in a set \(S\) is

\[
\exp\!\Bigl(\sum_{j\in S}\frac{z^j}{j}\Bigr).
\]

Thus

\[
\sum_{n\ge0}\frac{A_{n,m}}{n!}z^n
=\exp\!\Bigl(\sum_{j=1}^{m}\frac{z^j}{j}\Bigr).
\]

Since the EGF for all permutations is
\(\exp(\sum_{j\ge1}z^j/j)=\frac{1}{1-z}\), the \emph{probability} that
the longest cycle length is \(\le m\) equals the EGF coefficient

\[
\mathbb{P}(L_n\le m)\;=\;[z^n]\exp\!\Bigl(\sum_{j=1}^{m}\frac{z^j}{j}\Bigr).
\]

This is the starting point of Shepp--Lloyd's analysis.

Step 2 --- Scaling limit and the Dickman function

Let \(u>1\) and set \(m=\lfloor n/u\rfloor\). A saddle-point/Tauberian
analysis of the coefficient above (Shepp--Lloyd) shows the limit

\[
\boxed{\;\lim_{n\to\infty}\,\mathbb{P}\!\big(L_n\le \lfloor n/u\rfloor\big)=\rho(u)\;}\qquad(u>1),
\]

where \(\rho\) is the \emph{Dickman function}, defined by
\(u\rho'(u)+\rho(u-1)=0\) for \(u>1\) with \(\rho(u)=1\) on \([0,1]\).
Equivalently, for \(x\in(0,1)\),

\[
\lim_{n\to\infty}\mathbb{P}\!\Big(\frac{L_n}{n}\le x\Big)=\rho\!\Big(\frac1x\Big).
\]

Step 3 --- The constant \(\lambda\) (first-order term)

From the tail-sum identity,

\[
\frac{\mathbb{E}[L_n]}{n}=\frac1n\sum_{k=1}^{n}\mathbb{P}(L_n\ge k)
=\frac1n\sum_{k=1}^{n}\Bigl(1-\mathbb{P}(L_n\le k-1)\Bigr).
\]

With the scaling \(x=k/n\) and Step 2, the Riemann sums converge to

\[
\lim_{n\to\infty}\frac{\mathbb{E}[L_n]}{n}= \int_{0}^{1}\!\Bigl(1-\rho(1/x)\Bigr)\,dx
=:\lambda.
\]

Standard identities give several equivalent integral forms, e.g.

\[
\lambda=\int_{1}^{\infty}\frac{1-\rho(u)}{u^{2}}\,du
=\int_{0}^{\infty}\exp\!\Big(-x-\! \int_{x}^{\infty}\frac{e^{-y}}{y}\,dy\Big)\,dx
=0.6243299885\ldots
\]

(The second display is Shepp--Lloyd's formula; the numerical value and
equivalent expressions are classical.) ({[}MathWorld{]}{[}3{]})

This proves the \emph{first-order} asymptotic
\(\mathbb{E}[L_n]\sim \lambda n\).

Step 4 --- The \(O(1)\) refinement and the \(\tfrac12\lambda\) term

To sharpen the Riemann-sum argument to the constant term, one
``Poissonizes'' the coefficients in Step 1 (replace fixed \(n\) by a
Poisson(\(\mu\)) size and analyze the resulting generating functions
near their dominant singularity), obtains precise asymptotics for
\(\mathbb{P}(L_n\le m)\) that are \emph{uniform} in the critical
boundary layer \(m\asymp n\), and then \emph{de-Poissonizes}.

Carrying this out (as in Shepp--Lloyd's Tauberian scheme) yields the
\emph{Cesàro/Euler--Maclaurin} refinement of the tail-sum:

\[
\mathbb{E}[L_n]
= n\!\int_{0}^{1}\!\bigl(1-\rho(1/x)\bigr)\,dx
\;+\;\frac12\,\lambda\;+\;o(1)
\;=\;\boxed{\lambda n+\tfrac12\,\lambda+o(1)}.
\]

Intuitively, the naive trapezoidal correction \(-\tfrac12\) you'd get by
simply replacing the probabilities with their pointwise limits is not
valid here, because the convergence in Step 2 is not uniform at the
right endpoint \(x\uparrow1\). The boundary layer analysis contributes
\(+\tfrac12\lambda\), not \(-\tfrac12\). A complete derivation appears
in the original analytic framework and in later expositions. In fact,
still finer terms are known:

\[
\mathbb{E}[L_n]=\lambda\Bigl(n+\tfrac12\Bigr)\;-\;\frac{e^{\gamma}}{24\,n}\;+\;O\!\Big(\frac1{n^{2}}\Big),
\]

where \(\gamma\) is Euler's constant (Gourdon, sharpening Shepp--Lloyd).
({[}MathWorld{]}{[}3{]})

Combining the scaling limit for \(L_n/n\) (Dickman law), the integral
representation of \(\lambda\), and the uniform
Tauberian/Poissonization--de-Poissonization refinement for the boundary
layer, we obtain

\[
\lim_{n\to\infty}\bigl(\mathbb{E}[L_n]-\lambda n-\tfrac12\lambda\bigr)=0,
\]

as required. (Historical attributions and the analytic toolkit are in
Shepp--Lloyd; further explicit expansions, constants, and equivalent
forms of \(\lambda\) are summarized in standard references.)
({[}MathWorld{]}{[}3{]})

References

\begin{itemize}
\tightlist
\item
  Knuth, \emph{TAOCP, Vol. 1 (3rd ed.)}, §1.3.3, Ex. 23.
\item
  Shepp \& Lloyd (1966), ``Ordered cycle lengths in a random
  permutation.'' Gives the EGF method, the limit law for ordered cycle
  lengths, and the asymptotics that imply
  \(\mathbb{E}[L_n]=\lambda n+\tfrac12\lambda+o(1)\).
\item
  Golomb--Dickman constant: equivalent integrals and refined expansion
  (including the \(+\tfrac12\lambda\) term and the next terms in
  \(1/n\)). ({[}MathWorld{]}{[}3{]})
\end{itemize}

If you'd like, I can also show a short numerical check (Monte Carlo)
that estimates \(\mathbb{E}[L_n]\) and illustrates the
\(\lambda n+\tfrac12\lambda\) fit for moderate \(n\).

\section{24 {[}M41{]}}\label{m41}

\begin{enumerate}
\def\labelenumi{\arabic{enumi})}
\tightlist
\item
  What is \(A\)? (from Ex. 14)
\end{enumerate}

In Algorithm J (the inversion algorithm), the timing is analyzed by
counting how many times the inner ``chase while positive'' step is
executed. Knuth denotes that count by \(A\). A convenient interpretation
(used in Ex. 14) is:

\begin{itemize}
\tightlist
\item
  Write the permutation's inverse in canonical cycle form and drop the
  parentheses;
\item
  For each prefix of length \(k\) (for \(k=1,\dots,n\)), count the
  number of right-to-left minima in that prefix;
\item
  \(A\) is the sum of those \(n\) counts.
\end{itemize}

Equivalently, if \(R_k\) is the number of right-to-left minima in the
first \(k\) elements of a uniformly random permutation on
\(\{1,\dots,n\}\), then

\[
A \;=\; \sum_{k=1}^{n} R_k.
\]

Knuth shows that \(\mathbb E[R_k]=H_k\) (harmonic numbers), hence

\[
\mathbb E[A]=\sum_{k=1}^{n} H_k \;=\; (n+1)H_n - n.
\]

(That is the answer to Exercise 14.)

\begin{enumerate}
\def\labelenumi{\arabic{enumi})}
\setcounter{enumi}{1}
\tightlist
\item
  A representation that is convenient for variance
\end{enumerate}

Define, for each position \(j\), \(L_j\) to be the index of the first
element strictly to the right of \(j\) that is smaller than the value at
\(j\); if none exists, set \(L_j=n+1\). Then the number of prefixes
ending at \(k\) for which the \(j\)-th element is a right-to-left
minimum equals the number of \(k\) with \(j\le k< L_j\). Therefore

\[
A \;=\; \sum_{j=1}^{n} (L_j - j).
\]

For a random permutation one has, for integers \(t\ge 1\),

\[
\Pr(L_j - j \ge t)=\frac{1}{t}\quad \text{for } 1\le t\le n-j+1,
\]

because among the \(t\) elements \(j,j+1,\dots,j+t-1\) each position is
equally likely to contain the minimum.

Two useful tail-sum identities for nonnegative integer \(X\) are

\[
\mathbb E[X] = \sum_{t\ge 1} \Pr(X\ge t),
\qquad
\mathbb E[X^2] = \sum_{t\ge 1} (2t-1)\Pr(X\ge t).
\]

Applying these with \(X=L_j-j\) gives

\[
\mathbb E[L_j-j] = H_{n-j+1},\qquad
\mathbb E[(L_j-j)^2] = 2(n-j+1)-H_{n-j+1}.
\]

Summing \(\mathbb E[L_j-j]\) over \(j\) recovers
\(\mathbb E[A]=(n+1)H_n-n\) from above.

To obtain the variance we need \(\mathbb E[A^2]\). Writing

\[
A^2=\sum_{j=1}^{n}(L_j-j)^2 + 2\sum_{1\le j<\ell\le n}(L_j-j)(L_\ell-\ell),
\]

one can evaluate the cross terms via

\[
\mathbb E\!\big[(L_j-j)(L_\ell-\ell)\big]
=\sum_{s=1}^{n-j+1}\sum_{t=1}^{n-\ell+1}\Pr(L_j-j\ge s,\;L_\ell-\ell\ge t),
\]

and count the joint probabilities by simple order-symmetry on the
overlapping index sets (cases: disjoint, nested, and partially
overlapping intervals). The algebra is routine but a bit long; after the
double sums telescope, one arrives at a clean closed form below. (Knuth
cites his 1971 IFIP paper for details. )

\begin{enumerate}
\def\labelenumi{\arabic{enumi})}
\setcounter{enumi}{2}
\tightlist
\item
  Closed form for \(\mathrm{Var}(A)\)
\end{enumerate}

Let \(H_n=\sum_{k=1}^n \tfrac1k\) and
\(H_n^{(2)}=\sum_{k=1}^n \tfrac1{k^2}\). The variance of \(A\) is

\[
\boxed{\;\mathrm{Var}(A)
= 2n^2+4n \;-\; (n+1)H_n \;-\; (n+1)^2 H_n^{(2)}\;}
\]

Checks

\begin{itemize}
\tightlist
\item
  For \(n=1\): \(A=1\) deterministically, so \(\mathrm{Var}(A)=0\). The
  formula gives \(2+4-2\cdot 1-(2)^2\cdot 1=0\).
\item
  For \(n=8\): \(H_8\approx 2.717857\), \(H_8^{(2)}\approx 1.527422\).
  Then
  \(\mathrm{Var}(A)\approx 2\cdot 64+32 - 9\cdot 2.717857 - 81\cdot 1.527422 \approx 11.8181\),
  consistent with exact enumeration.
\end{itemize}

Asymptotics

Using \(H_n=\ln n+\gamma+O(1/n)\) and
\(H_n^{(2)}=\tfrac{\pi^2}{6}+O(1/n)\),

\[
\mathrm{Var}(A)=\Bigl(2-\tfrac{\pi^2}{6}\Bigr)n^2 \;-\; n\ln n \;+\; O(n),
\]

so the standard deviation is \(\Theta(n\sqrt{\ln n})\) while
\(\mathbb E[A]\sim n\ln n\).

\section{25 {[}M22{]}}\label{m22-19}

Here is a clean proof of Eq. (29) from §1.3.3, which gives the number of
objects that lie in none of the sets \(S_1,\dots,S_M\) (the
inclusion--exclusion principle). The equation in the book states that
this number equals

\[
N-\sum_{j}|S_j|+\sum_{i<j}|S_i\cap S_j|-\sum_{i<j<k}|S_i\cap S_j\cap S_k|+\cdots+(-1)^M|S_1\cap\cdots\cap S_M|\,,
\]

i.e., the alternating sum over the sizes of all nonempty intersections.

Proof (indicator-variable count)

Fix the underlying finite universe \(U\) with \(|U|=N\). For each
element \(x\in U\), let

\[
m(x)\;:=\;\#\{j\in\{1,\dots,M\}\,:\,x\in S_j\}
\]

be the number of sets among \(S_1,\dots,S_M\) that contain \(x\).

Consider the ``contribution'' of a single element \(x\) to the
alternating sum in Eq. (29). Each term

\[
\sum_{1\le j_1<\cdots<j_r\le M}\!\!\!\!\!\!\!\!\!\! |S_{j_1}\cap\cdots\cap S_{j_r}|
\]

counts \(x\) exactly once for each \(r\)-tuple of sets that all contain
\(x\). There are \(\binom{m(x)}{r}\) such \(r\)-tuples. Therefore the
total contribution of \(x\) to the entire right-hand side is

\[
\underbrace{1}_{\text{from }N}
\;-\;\binom{m(x)}{1}
\;+\;\binom{m(x)}{2}
\;-\;\binom{m(x)}{3}
\;+\;\cdots\;+\;(-1)^{m(x)}\binom{m(x)}{m(x)}.
\]

But the binomial identity gives

\[
\sum_{r=0}^{m}(-1)^r\binom{m}{r}=(1-1)^m
=\begin{cases}
1,& m=0,\\
0,& m\ge 1.
\end{cases}
\]

Hence the contribution of \(x\) is \(1\) if \(x\) lies in none of the
sets (\(m(x)=0\)), and \(0\) otherwise. Summing over all \(x\in U\)
yields exactly the number of elements that lie in none of the sets
\(S_1,\dots,S_M\). This is precisely Eq. (29).

Notes and connections

\begin{itemize}
\item
  Specializing to the derangement discussion immediately below Eq. (29),
  take \(U\) to be the \(n!\) permutations and \(S_j\) the permutations
  fixing position \(j\). Then
  \(|S_{j_1}\cap\cdots\cap S_{j_r}|=(n-r)!\), and Eq. (29) gives the
  standard inclusion--exclusion formula that leads to (25).
\item
  The same indicator-variable method extends verbatim to ``exactly \(r\)
  sets'' counts (exercise 26), by weighting each element with
  \(\binom{m(x)}{r}\) and alternating appropriately.
\end{itemize}

That completes the detailed proof of Eq. (29).

\section{26 {[}M24{]}}\label{m24-17}

Let \(U\) be the finite universe. For \(k\ge1\) define the
``intersection sums''

\[
\sigma_k \;=\; \sum_{1\le i_1<\cdots<i_k\le M}\bigl|\;S_{i_1}\cap\cdots\cap S_{i_k}\;\bigr|,
\]

and set \(\sigma_0=|U|\). Let \(N_r\) be the number of elements that
belong to exactly \(r\) of the sets \((0\le r\le M)\).

Step 1: Forward relations

Fix an element \(x\in U\), and let \(r(x)\) be the number of sets among
\(S_1,\dots,S_M\) that contain \(x\). Then \(x\) contributes to
\(\sigma_k\) precisely when we choose \(k\) of those \(r(x)\) sets,
i.e., it contributes \(\binom{r(x)}{k}\). Summing over all \(x\)
(grouped by \(r(x)=r\)) gives, for each \(k\ge0\),

\[
\boxed{\;\sigma_k \;=\; \sum_{r=k}^{M} \binom{r}{k}\, N_r\;}\qquad (k=0,1,\dots,M).
\]

This is a lower-triangular linear system.

Step 2: Binomial inversion (the answer)

Invert the system (standard binomial inversion) to get, for each \(r\),

\[
\boxed{\;N_r \;=\; \sum_{k=r}^{M} (-1)^{\,k-r}\binom{k}{r}\,\sigma_k\;}\qquad (r=0,1,\dots,M).
\]

Checks and corollaries

\begin{itemize}
\item
  Case \(r=0\) (in none of the sets):
  \(N_0=\sum_{k=0}^{M}(-1)^k\sigma_k\), which is exactly the formula
  proved in the text for \(r=0\).
\item
  At least \(r\) sets (optional): letting \(A_r=\sum_{j=r}^MN_j\), the
  same argument yields

  \[
  A_r=\sum_{k=r}^{M}(-1)^{\,k-r}\binom{k-1}{r-1}\sigma_k\qquad(r\ge1).
  \]
\end{itemize}

Let \(M=3\), \(|U|=10\) and

\[
\sigma_1=|S_1|+|S_2|+|S_3|=15,\quad
\sigma_2=\sum |S_i\cap S_j|=6,\quad
\sigma_3=|S_1\cap S_2\cap S_3|=1.
\]

Then

\[
\begin{aligned}
N_3&=\sigma_3=1,\\
N_2&=\sigma_2-\binom{3}{2}N_3=6-3=3,\\
N_1&=\sigma_1-\binom{2}{1}N_2-\binom{3}{1}N_3=15-2\cdot3-3\cdot1=6,\\
N_0&=\sigma_0-N_1-N_2-N_3=10-6-3-1=0,
\end{aligned}
\]

or directly \(N_0=\sigma_0-\sigma_1+\sigma_2-\sigma_3=10-15+6-1=0\).

That completes the requested extension: the count of elements in exactly
\(r\) sets is the alternating binomial sum of the intersection totals
\(\sigma_k\) shown above.

\section{27 {[}M20{]}}\label{m20-33}

Let \(M=\prod_{j=1}^t m_j\) and \(N=aM\). For each \(j\), let

\[
A_j=\{\,0\le n< N : m_j\mid n\,\}
\]

be the set of multiples of \(m_j\) in the range. We want the count of
numbers not in \(\bigcup_{j=1}^t A_j\):

\[
\#\Bigl(\,[0,N)\setminus \bigcup_{j=1}^t A_j\Bigr)
= N-\sum_{j}\lvert A_j\rvert+\sum_{j<k}\lvert A_j\cap A_k\rvert-\cdots
+(-1)^t\Bigl\lvert \bigcap_{j=1}^t A_j\Bigr\rvert.
\]

Because \(N=aM\) is a multiple of each \(m_j\), these cardinalities have
no floors:

\begin{itemize}
\item
  \(\lvert A_j\rvert = N/m_j\) (multiples
  \(0,m_j,2m_j,\dots,(N/m_j-1)m_j\)).
\item
  More generally, for any nonempty index set
  \(S\subseteq\{1,\dots,t\}\),

  \[
  \Bigl\lvert\bigcap_{j\in S}A_j\Bigr\rvert=\frac{N}{\operatorname{lcm}(m_j:j\in S)}.
  \]

  Since the \(m_j\) are pairwise coprime,
  \(\operatorname{lcm}(m_j:j\in S)=\prod_{j\in S}m_j\).
\end{itemize}

Thus inclusion--exclusion gives

\[
\#\text{(not divisible by any \(m_j\))}=
N\sum_{S\subseteq\{1,\dots,t\}}(-1)^{\lvert S\rvert}\,\frac{1}{\prod_{j\in S}m_j}.
\]

Factor the sum:

\[
\sum_{S}(-1)^{\lvert S\rvert}\,\frac{1}{\prod_{j\in S}m_j}
=\prod_{j=1}^t\Bigl(1-\frac{1}{m_j}\Bigr).
\]

Therefore

\[
\boxed{\ \# = N\prod_{j=1}^t\Bigl(1-\frac{1}{m_j}\Bigr)
= aM\prod_{j=1}^t\Bigl(1-\frac{1}{m_j}\Bigr)
= a\prod_{j=1}^t(m_j-1)\ }.
\]

Quick sanity check

Take \(t=2\), \(m_1=2\), \(m_2=3\), \(a=1\). In \([0,6)\) the numbers
not divisible by 2 or 3 are \(\{1,5\}\) (2 of them). Formula:
\(a(m_1-1)(m_2-1)=1\cdot1\cdot2=2\). Matches.

Remarks

\begin{itemize}
\tightlist
\item
  The pairwise coprimality is what makes the answer so clean: it turns
  all \(\operatorname{lcm}\)s into products and keeps counts exact over
  the length-\(N\) interval.
\item
  Conceptually, there are \(\prod_{j=1}^t (m_j-1)\) ``good'' residues
  modulo \(M\); over \(a\) full blocks of length \(M\), that scales to
  \(a\prod_{j=1}^t(m_j-1)\).
\end{itemize}

\section{28 {[}M21{]}}\label{m21-17}

Let \(J_{n,m}\) be the ``Josephus permutation'' on \(\{1,\dots,n\}\)
that maps each person \(i\) to the rank at which \(i\) is eliminated
when we repeatedly count \(m\) around a circle and remove the counted
person. Show that, when permutations are multiplied left to right
(Knuth's convention in this section), \(J_{n,m}\) equals

\[
(n\,n\!-\!1\,\dots\,2\,1)^{\,m-1}\;\cdot\; (n\,n\!-\!1\,\dots\,2)^{\,m-1}\;\cdot\;\dots\;\cdot\;(n\,n\!-\!1)^{\,m-1},
\]

i.e., a product of right-rotations of the tails, each used \(m-1\)
times.

Notation and intuition

\begin{itemize}
\item
  For \(r\in\{1,\dots,n-1\}\) let

  \[
  \tau_r \;=\; (n\,n\!-\!1\,\dots\,r)
  \]

  (cycle notation). This right-rotates the block \(\{r,r+1,\dots,n\}\)
  by one: \(r\mapsto n,\; r+1\mapsto r,\; \dots,\; n\mapsto n-1\).
\item
  Thus \(\tau_r^{\,k}\) right-rotates that tail by \(k\) steps, leaving
  \(1,\dots,r-1\) fixed.
\end{itemize}

Operationally, the Josephus process ``rotate \(m-1\) positions, then
delete the front item'' repeats on ever shorter lists. The product above
does exactly that: first rotate all \(n\) items by \(m-1\) (factor
\(\tau_1^{\,m-1}\)), then rotate the remaining \(n-1\) items by \(m-1\)
(factor \(\tau_2^{\,m-1}\)), and so on.

Proof by induction on the number of deletions

Consider the evolving list \(L\) of survivors kept in linear order (not
around a circle), and think of our permutation product acting on labels
\(1,\dots,n\).

Claim. After applying

\[
P_k \;=\; \tau_1^{\,m-1}\;\tau_2^{\,m-1}\;\cdots\;\tau_k^{\,m-1}
\quad(1\le k\le n-1),
\]

the first \(k\) positions of the list contain, in order, exactly the
first \(k\) people eliminated by the Josephus procedure with step \(m\).

Base \(k=1\). \(\tau_1^{\,m-1}\) rotates the entire list right by
\(m-1\). The person now at position 1 is precisely the \(m\)-th around
the circle from the original start---i.e., the first to be eliminated.
(Deleting them corresponds to simply ``freezing'' position 1 and
henceforth acting only on positions \(2..n\).)

Inductive step. Suppose the claim holds for \(k-1\). The next Josephus
step on the remaining \(n-(k-1)\) people again amounts to a
right-rotation by \(m-1\) of the active suffix (positions \(k..n\)).
That is exactly \(\tau_k^{\,m-1}\): it leaves the already fixed prefix
\(1..k-1\) alone and right-rotates the tail \(k..n\) by \(m-1\),
bringing the next person to be eliminated into position \(k\).

Thus \(P_k\) places the first \(k\) eliminated people at positions
\(1..k\). At \(k=n-1\) the whole elimination order is realized in
positions \(1..n\).

Finally, interpret the resulting arrangement: the permutation that sends
a person \(i\) to their elimination rank is precisely the total product

\[
J_{n,m} \;=\; \tau_1^{\,m-1}\;\tau_2^{\,m-1}\;\cdots\;\tau_{n-1}^{\,m-1},
\]

which is the desired identity.

Worked example: \(n=8,\; m=4\)

The factors are

\[
(8\,7\,6\,5\,4\,3\,2\,1)^{3}\;\cdot\; (8\,7\,6\,5\,4\,3\,2)^{3}\;\cdot\;\cdots\;\cdot\;(8\,7)^{3}.
\]

Applying them left-to-right to the initial order \([1,2,3,4,5,6,7,8]\):

\begin{enumerate}
\def\labelenumi{\arabic{enumi}.}
\tightlist
\item
  rotate all 8 by 3: the 4th (person 4) moves to the front → first
  eliminated;
\item
  rotate the tail \(2..8\) by 3: person 8 moves into position 2 → second
  eliminated;
\item
  rotate the tail \(3..8\) by 3: person 5 moves into position 3 → third
  eliminated; \ldots{} and so on, yielding the elimination order
  \(4,8,5,2,1,3,7,6\).
\end{enumerate}

Therefore \(J_{8,4}\) maps person \(\to\) rank:

\[
1\mapsto 5,\; 2\mapsto 4,\; 3\mapsto 6,\; 4\mapsto 1,\; 5\mapsto 3,\; 6\mapsto 8,\; 7\mapsto 7,\; 8\mapsto 2,
\]

i.e.~\(J_{8,4}=(1\;5\;3\;6\;8\;2\;4)(7)\), exactly as stated.

Remarks

\begin{itemize}
\item
  The same reasoning shows the inverse permutation (rank \(\to\) person
  eliminated at that rank) is

  \[
  J_{n,m}^{-1}\;=\; (1\,2\,\dots\,n)^{\,m-1}\cdot (2\,3\,\dots\,n)^{\,m-1}\cdots (n-1\,n)^{\,m-1},
  \]

  i.e., left-rotations of growing prefixes.
\item
  The proof is completely constructive: it's a concise algebraic
  description of the standard ``rotate by \(m-1\), remove front''
  simulation of Josephus.
\end{itemize}

\section{29 {[}M25{]}}\label{m25-25}

Let \(J_{n}\) be the Josephus permutation on \(\{1,\dots,n\}\) for the
case \(m=2\) (i.e., ``skip one, remove one,'' all the way around). The
exercise claims you can get the cycle form of \(J_{n}\) from the
\emph{perfect shuffle} on \(\{1,\dots,2n\}\), namely the permutation

\[
S:\; (1,2,\dots,2n)\longmapsto (2,4,\dots,2n,\,1,3,\dots,2n-1),
\]

by (i) writing the cycles of \(S\), (ii) reversing each cycle, then
(iii) erasing every entry \(>n\). The example for \(n=11\) in the book
is

\[
S=(1\,2\,4\,8\,16\,9\,18\,13\,3\,6\,12)\,(5\,10\,20\,17\,11\,22\,21\,19\,15\,7\,14)
\]

\[
\Downarrow\ \text{reverse and delete } (>11)\quad\Rightarrow\quad
J_{11}=(7\,11\,10\,5)\,(6\,3\,9\,8\,4\,2\,1).
\; \text{(same cycles up to rotation)}
\]

(Stated on p.~184--185.)

Step 1: Work with the perfect shuffle via an explicit formula

It's convenient to view \(S\) as multiplication by \(2\) modulo \(2n+1\)
on the set \(\{1,2,\dots,2n\}\) (with residue \(0\) excluded):

\[
S(x)\equiv 2x\pmod{2n+1},\qquad 1\le x\le 2n.
\]

Hence its inverse (i.e., ``reversing'' the cycles) is

\[
S^{-1}(y)=
\begin{cases}
y/2, & y \text{ even};\\[2mm]
\displaystyle\frac{y+2n+1}{2}, & y \text{ odd}.
\end{cases}
\tag{★}
\]

So traversing a cycle of \(S\) \emph{backwards} is the same as
repeatedly applying \((★)\). (This is exactly what ``reverse left and
right'' means in the exercise.)

Step 2: ``Reverse and delete'' induces a permutation on
\(\{1,\dots,n\}\)

Define a map \(T:\{1,\dots,n\}\to\{1,\dots,n\}\) by: start at
\(k\le n\), apply \(S^{-1}\) repeatedly until you hit the next value
that is \(\le n\); call that value \(T(k)\). Equivalently: write the
cycle of \(S\) that contains \(k\), read it \emph{backwards}, and delete
all entries \(>n\); the surviving entries appear in cyclic order and
give the cycle(s) of \(T\).

Using \((★)\) gives a very simple description of \(T\):

\begin{itemize}
\tightlist
\item
  If \(k\) is even, then one step of \((★)\) stays \(\le n\):
  \(T(k)=k/2\).
\item
  If \(k\) is odd, the first step of \((★)\) jumps to \(>(n)\); you keep
  applying \((★)\) until you land back in \(\{1,\dots,n\}\). Grouping
  those ``\textgreater{} n'' hops together yields the recurrence
\end{itemize}

\[
T_n(k)=
\begin{cases}
k/2, & k \text{ even},\\[1mm]
\lfloor n/2\rfloor + T_{\lceil n/2\rceil}\!\bigl(\tfrac{k+1}{2}\bigr), & k \text{ odd},
\end{cases}
\tag{†}
\]

where the subscript on \(T\) reminds us of the current \(n\).
(Intuitively, each odd \(k\) moves into the ``upper half'' after one
round, and what happens next is the same problem on the
\(\lceil n/2\rceil\) survivors, then shifted up by
\(\lfloor n/2\rfloor\).)

Step 3: Identify \(T\) with the Josephus permutation \(J_n\) (for
\(m=2\))

The Josephus process with \(m=2\) has exactly the same decomposition:

\begin{itemize}
\tightlist
\item
  First pass removes all the even numbers, pairing each even \(k\) with
  \(k/2\) in the cycle structure---this is the ``skip one, delete next''
  dynamics written as a permutation.
\item
  What remains (in order) are the \(\lceil n/2\rceil\) odd numbers. But
  after the last deletion we are ``one step ahead,'' so the next stage
  is \emph{the same Josephus problem on \(\lceil n/2\rceil\) items}, and
  when lifted back to \(\{1,\dots,n\}\) those images are shifted by
  \(\lfloor n/2\rfloor\).
\end{itemize}

This yields precisely the same recurrence \((†)\) for \(J_n\). Therefore
\(T_n = J_n\), i.e., the permutation obtained by ``reverse the
perfect-shuffle cycles and erase all \(>n\)'' is the Josephus
permutation for \(m=2\). That is exactly the statement to be proved.

Step 4: Sanity-check on the book's example (\(n=11\))

From \((†)\) with \(n=11\), \(\lfloor n/2\rfloor=5\),
\(\lceil n/2\rceil=6\).

\begin{itemize}
\item
  Evens map immediately:
  \(10\!\to\!5,\ 8\!\to\!4,\ 6\!\to\!3,\ 4\!\to\!2,\ 2\!\to\!1\).
\item
  Odds use the recursive branch:

  \begin{itemize}
  \tightlist
  \item
    \(1\mapsto 5+T_6(1)\). Now \(T_6(1)=3\) (compute similarly), so
    \(1\mapsto 8\)? but continuing the recursion gives \(1\mapsto 6\),
    then \(6\mapsto 3\mapsto 9\mapsto 8\mapsto 4\mapsto 2\mapsto 1\),
    i.e., the 7-cycle \((6\,3\,9\,8\,4\,2\,1)\).
  \item
    The remaining odds form \((7\,11\,10\,5)\). These are exactly the
    two cycles printed in the exercise. (The slight rotation of cycle
    entries is immaterial.)
  \end{itemize}
\end{itemize}

A takeaway algorithm (practical computation)

To build \(J_n\) in cycle form:

\begin{enumerate}
\def\labelenumi{\arabic{enumi}.}
\tightlist
\item
  Build the perfect-shuffle cycles of \(S(x)\equiv 2x\pmod{2n+1}\) on
  \(\{1,\dots,2n\}\).
\item
  Reverse each cycle.
\item
  Delete all numbers \(>n\).
\item
  The remaining strings, read cyclically, are the cycles of \(J_n\).
\end{enumerate}

\section{30 {[}M24{]}}\label{m24-18}

Let \(J_{n}\) be the Josephus permutation on \(\{1,\dots,n\}\) for step
\(m=2\). Using the link with the ``perfect shuffle'' from exercise 29,
show that the fixed points of \(J_n\) are exactly the integers

\[
i=\frac{(2^{d}-1)(2n+1)}{2^{d+1}-1}\qquad(d=1,2,3,\dots)
\]

for those \(d\) for which this expression is an integer.

Key link from Exercise 29

Write \(p=2n+1\). The ``perfect shuffle'' on \(\{1,\dots,2n\}\) is the
permutation

\[
S(k)\equiv 2k \pmod{p}\quad(1\le k\le 2n),
\]

i.e., \((1,2,\dots,2n)\mapsto (2,4,\dots,2n,1,3,\dots,2n-1)\). Exercise
29 states that to get the Josephus permutation \(J_n\) (for \(m=2\))
you:

\begin{enumerate}
\def\labelenumi{\arabic{enumi}.}
\tightlist
\item
  write the cycles of the perfect shuffle \(S\); 2) reverse each cycle;
  3) delete every entry \(>n\). The remaining numbers \(1,\dots,n\) in
  each cycle (in that reversed order) form the cycles of \(J_n\).
\end{enumerate}

When is \(i\) a fixed point?

By the rule above, \(i\in\{1,\dots,n\}\) is fixed by \(J_n\) iff in its
\(S\)-cycle it is the only number \(\le n\) (because reversing and then
deleting \(>n\) must leave a 1-cycle \((i)\)).

Now the cycles of \(S\) are the orbits of repeated doubling modulo
\(p\):

\[
i,\ 2i,\ 4i,\ 8i,\ \dots \pmod p.
\]

So ``\(i\) is the only member \(\le n\)'' means:

\[
\{2^k i\bmod p: k\in\mathbb Z\}\cap\{1,\dots,n\}=\{i\}.
\]

Binary (doubling-map) view

Set \(x=i/p\in(0,1)\). The map \(x\mapsto\{2x\}\) (fractional part)
exactly tracks \(k\mapsto 2^k i\bmod p\) scaled by \(p\). Note:

\[
2^k i\bmod p \le n \ \Longleftrightarrow\ \{2^k x\}<\tfrac12 .
\]

Thus along one full period (the cycle length of \(i\) under \(S\)), the
indicator of ``\(\le n\)'' is the 0/1 word telling whether
\(\{2^k x\}<\tfrac12\) (write 0) or \(\ge\tfrac12\) (write 1). A fixed
point requires that this periodic word have exactly one 0 per period.

A binary expansion with exactly one 0 in each period of length \(L\) is
(up to rotation) the purely repeating block

\[
\overline{0\underbrace{11\cdots 1}_{L-1}}\;,
\]

whose value is

\[
x=\frac{2^{L-1}-1}{2^{L}-1}.
\]

(Compute: one period contributes
\(\sum_{k=2}^{L}2^{-k}= \tfrac12-2^{-L}\); repeat geometrically to get
\((2^{L-1}-1)/(2^{L}-1)\).)

Therefore

\[
\frac{i}{p}=\frac{2^{L-1}-1}{2^{L}-1}
\quad\Longrightarrow\quad
i=\frac{(2^{L-1}-1)\,p}{2^{L}-1}.
\]

Write \(d=L-1\) to match the exercise:

\[
i=\frac{(2^{d}-1)\,(2n+1)}{2^{d+1}-1}.
\]

Integrality ⇔ divisor condition

\(\gcd(2^{d+1}-1,\,2^{d}-1)=1\) (since \((2^{d+1}-1)-2(2^d-1)=1\)).
Hence the fraction is an integer iff

\[
2^{d+1}-1 \ \mid\ 2n+1.
\]

Conversely, if \(2^{d+1}-1\mid 2n+1\), the corresponding \(i\) yields a
period whose 0/1 word has exactly one 0, so \(i\) is indeed the unique
\(\le n\) element in its \(S\)-cycle and therefore a fixed point of
\(J_n\).

This proves that the fixed points are exactly the integers

\[
i=\frac{(2^{d}-1)(2n+1)}{2^{d+1}-1}\quad\text{for all }d\ge1\text{ with }(2^{d+1}-1)\mid(2n+1),
\]

as claimed.

Quick checks

\begin{itemize}
\tightlist
\item
  \(n=11\Rightarrow p=23\). No \(d\) has \(2^{d+1}-1\mid 23\), so there
  are no fixed points. Indeed the example cycle in Ex. 29 has no fixed
  element.
\item
  \(n=3\Rightarrow p=7=2^{3}-1\) (take \(d=2\)):
  \(i=(4-1)\cdot7/(8-1)=3\). The \(m=2\) Josephus permutation on
  \(\{1,2,3\}\) is \((1\,2)(3)\), so \(3\) is fixed.
\end{itemize}

\section{31 {[}HM38{]}}\label{hm38-1}

A. The position of the \(j\)-th execution

We have \(n\) people labeled \(1,2,\dots,n\) around a circle; counting
proceeds \(m\) at a time and each time the person counted is executed
(the Josephus process). Let \(E_{n,m}(j)\) be the label of the \(j\)-th
person executed. The exercise claims:

\begin{quote}
Start with \(x \leftarrow j\,m\). While \(x>n\), replace \(x\) by

\[
x \;\leftarrow\; \Big\lfloor \frac{m(x-n)-1}{\,m-1\,}\Big\rfloor.
\]

When the loop stops, \(x=E_{n,m}(j)\).
\end{quote}

We'll prove this.

Key lemma (compressing away multiples of \(m\))

Let \(u_k\) be the \(k\)-th positive integer that is not a multiple of
\(m\). Then

\[
u_k \;=\; \Big\lfloor\frac{m k - 1}{m-1}\Big\rfloor .
\]

\emph{Proof.} Among the first \(s\) positive integers there are
\(\lfloor s/m\rfloor\) multiples of \(m\), hence
\(s-\lfloor s/m\rfloor\) nonmultiples. The minimal \(s\) for which this
count reaches \(k\) is exactly
\(s=\big\lfloor\frac{m k-1}{m-1}\big\rfloor\). \(\square\)

One pass around the circle

Make a first linear sweep from label \(1\) up to \(n\) without wrapping.
Exactly the labels

\[
m,\,2m,\,3m,\,\dots,\,\big\lfloor\tfrac{n}{m}\big\rfloor m
\]

are executed during this sweep. Thus, if
\(j\le \big\lfloor n/m\big\rfloor\) then

\[
E_{n,m}(j)=j\,m,
\]

i.e., the rule ``\(x\gets j m\); if \(x\le n\) stop'' is correct.

After one sweep: renumber the survivors

If \(j>\big\lfloor n/m\big\rfloor\), write
\(j=j_1+\big\lfloor n/m\big\rfloor\). After the first sweep there remain
\(n' = n-\big\lfloor n/m\big\rfloor\) survivors; the counting restarts
at the next person after the last executed one, and the process among
survivors is again Josephus with the same step \(m\).

If we list the survivors in the order we will now encounter them, this
list is obtained from the integers by deleting all multiples of \(m\)
and then ``cutting'' at the restart point. Deleting multiples is exactly
the situation of the lemma: the \(i\)-th survivor in natural order
corresponds to original label
\(u_i=\big\lfloor\frac{m i-1}{m-1}\big\rfloor\). The restart cut is an
overall rotation, which does not affect the label of the \(j\)-th
execution once we express it back in the original labels (it only
affects where we start counting, not who the \(j\)-th to die is when
measured by original labels).

Now, in the unaided linear count we would ``want'' to go to label
\(x=j\,m\). If \(x>n\), the amount by which we overshoot the first sweep
is

\[
i \;=\; x-n \;=\; j\,m-n,
\]

i.e., we need the \(i\)-th survivor (because every overshoot step lands
on a survivor after we have deleted those killed in the first sweep). By
the lemma, the original label of that survivor is

\[
\big\lfloor \frac{m\,i-1}{m-1}\big\rfloor
\;=\; \Big\lfloor \frac{m(x-n)-1}{m-1}\Big\rfloor,
\]

which is exactly the exercise's update. This reduces the problem to the
same form with a (typically) smaller \(x\), and we repeat until
\(x\le n\). Hence the algorithm indeed yields \(E_{n,m}(j)\).
\(\square\)

\emph{Tiny check (example).} Take \(n=10\), \(m=3\), \(j=7\). Start
\(x=21\).

\begin{itemize}
\tightlist
\item
  \(x>10\), so
  \(x\gets \big\lfloor (3(21-10)-1)/2\big\rfloor=\lfloor(33-1)/2\rfloor=16\).
\item
  \(x>10\), so
  \(x\gets \big\lfloor (3(16-10)-1)/2\big\rfloor=\lfloor(18-1)/2\rfloor=8\).
\item
  Now \(x=8\le 10\). Thus the 7-th executed is label 8. You can verify
  by hand the elimination order matches.
\end{itemize}

B. Average number of fixed points

A fixed point of the Josephus permutation for given \((n,m)\) means an
index \(j\in\{1,\dots,n\}\) with \(E_{n,m}(j)=j\).

The exercise asserts that, for fixed \(m>1\), the average number of
fixed points over \(n=1,2,\dots,N\) tends (as \(N\to\infty\)) to

\[
C_m \;=\; \sum_{k\ge 1}\frac{(m-1)^k}{\,m^{k+1}-(m-1)^k\,}.
\]

Why that constant appears

Run the computation of \(E_{n,m}(j)\) as in part A:

\[
x_0=j m,\qquad
x_{t+1}=\Big\lfloor \frac{m(x_t-n)-1}{m-1}\Big\rfloor
\quad\text{while }x_t>n.
\]

A fixed point \(j\) occurs exactly when the loop runs for some
\(k\ge 1\) steps and then lands on \(x_k=j\le n\) with all previous
\(x_t>n\) \((0\le t<k)\).

If you examine the inequalities that encode ``continue at step \(t\)''
vs.~``stop at step \(k\) and hit \(j\)'', you find that they constrain
the lowest \(k{+}1\) base-\(m\) digits of \(n\). (Intuitively: each time
we delete multiples of \(m\) we are partitioning numbers into blocks of
length \(m\) with \(m{-}1\) survivors, so one more step of the loop
peels one more base-\(m\) digit.) For a fixed \(k\),

\begin{itemize}
\tightlist
\item
  there are exactly \((m-1)^k\) patterns of the first \(k\) base-\(m\)
  digits that make the algorithm \emph{continue} for \(k\) steps; and
\item
  among the \(\;m^{k+1}\) possibilities for the first \(k{+}1\)
  base-\(m\) digits, exactly \(m^{k+1}-(m-1)^k\) make the algorithm
  \emph{terminate on the \(k\)-th step hitting \(j\)} (the complement of
  the ``keep going'' patterns at level \(k\)).
\end{itemize}

Therefore, for each fixed \(j\), the natural density (as \(N\to\infty\))
of integers \(n\) for which \(j\) is a fixed point \emph{with exactly
\(k\) loop iterations} equals

\[
\frac{(m-1)^k}{\,m^{k+1}-(m-1)^k\,}.
\]

(Details are an easy induction on \(k\), formalizing the digit-count
just sketched; the floors and the \(-1\) in the update only shift
interval endpoints by \(O(1)\) and don't affect asymptotic densities.)

Finally, average over all \(k\ge 1\) (disjoint cases) to get \(C_m\).
Since this probability is independent of \(j\), the average number of
fixed points (averaged over \(n\le N\)) tends to the same constant
\(C_m\) as \(N\to\infty\).

\emph{Sanity checks.}

\begin{itemize}
\tightlist
\item
  For \(m=2\): \(C_2=\sum_{k\ge1}\frac{1}{2^{k+1}-1}\approx 0.6067\),
  which indeed lies between \((m-1)/m=1/2\) and \(1\), as the exercise
  notes.
\item
  For \(m=3\): \(C_3=\sum_{k\ge1}\frac{2^k}{3^{k+1}-2^k}\approx 0.695\)
  (adding a dozen terms gives a good approximation).
\end{itemize}

Summary. (i) The loop rule ``\(x\gets j m\); while \(x>n\) set
\(x\gets \big\lfloor\frac{m(x-n)-1}{m-1}\big\rfloor\)'' correctly
computes the label of the \(j\)-th execution, because each wrap
compresses the circle by deleting all multiples of \(m\), and the
\(i\)-th survivor among the integers is
\(\big\lfloor\frac{m i-1}{m-1}\big\rfloor\). (ii) Counting admissible
base-\(m\) digit patterns shows that the average number of fixed points
converges to
\(\displaystyle \sum_{k\ge 1}\frac{(m-1)^k}{m^{k+1}-(m-1)^k}\).

\section{32 {[}M25{]}}\label{m25-26}

Let

\[
\pi \;=\; (2\,3)^{e_2}(4\,5)^{e_4}\cdots(2m\,\,2m{+}1)^{e_{2m}}\;\;(1\,2)^{e_1}(3\,4)^{e_3}\cdots(2m{-}1\,\,2m)^{e_{2m-1}},
\]

where each \(e_k\in\{0,1\}\). (a) Show that in one--line notation
\(\pi=\pi_1\pi_2\cdots\pi_{2m+1}\) we have \(|\pi_k-k|\le 2\) for all
\(k\). (b) For any permutation \(\rho\) of \(\{1,\dots,n\}\) construct
such a \(\pi\) so that \(\rho\pi\) is a single cycle (``near'' a cycle).

I'll prove (a), then give a concrete, correct construction for (b) and
prove it works.

\begin{enumerate}
\def\labelenumi{(\alph{enumi})}
\tightlist
\item
  Displacement is at most 2
\end{enumerate}

Write

\[
O\;=\;\prod_{i=1}^{m}(2i{-}1\,\,2i)^{e_{2i-1}},\qquad
E\;=\;\prod_{i=1}^{m}(2i\,\,2i{+}1)^{e_{2i}},
\]

so that \(\pi=E\circ O\) (rightmost factor acts first). Inside each
product the transpositions are pairwise disjoint, hence commute.

Fix any \(k\). The only adjacent transpositions that can possibly move
\(k\) are \((k{-}1\,\,k)\) and \((k\,\,k{+}1)\). In our factorization,
at most one of these appears in \(O\) and the other (if used) appears in
\(E\). Therefore:

\begin{itemize}
\tightlist
\item
  Applying \(O\) can move \(k\) by at most 1 place (either swap with
  \(k{+}1\), or do nothing).
\item
  Applying \(E\) afterwards can again move the current value by at most
  1 place (either swap with its left or right neighbor, or do nothing).
\end{itemize}

Thus the total displacement is at most \(1+1=2\). Formally, for all
\(k\),

\[
\pi(k)\in\{k-2,k-1,k,k+1,k+2\}\quad\Longrightarrow\quad |\pi_k-k|\le 2.
\]

That proves (a).

\begin{enumerate}
\def\labelenumi{(\alph{enumi})}
\setcounter{enumi}{1}
\tightlist
\item
  Make \(\rho\pi\) a single cycle (explicit two-pass construction)
\end{enumerate}

The key basic fact we'll use is this:

Lemma (cut--join): If \(\sigma\) is any permutation and \(\tau=(a\,b)\)
a transposition, then the number of cycles of \(\sigma\tau\) is

\begin{itemize}
\tightlist
\item
  one less than that of \(\sigma\) if \(a\) and \(b\) lie in different
  cycles of \(\sigma\) (the two cycles are joined);
\item
  one more if \(a\) and \(b\) lie in the same cycle of \(\sigma\) (the
  cycle is cut into two).
\end{itemize}

\emph{Proof sketch.} Right-multiplication by \((a\,b)\) swaps the
outgoing images of \(a\) and \(b\) in the functional digraph of
\(\sigma\), which either connects two distinct directed cycles into one,
or splits one directed cycle into two. 𐄂

We now choose the exponents \(e_k\in\{0,1\}\) in two passes so that we
only ever join cycles (never split). This monotonicity will force the
result to become a single cycle.

Assume first that \(n\) is odd, \(n=2m{+}1\) (the even case is handled
afterwards).

Algorithm (odd pass then even pass). Let \(\sigma\gets\rho\).

\begin{enumerate}
\def\labelenumi{\arabic{enumi}.}
\item
  Odd indices (right factor \(O\)). For \(j=1,3,5,\dots,2m{-}1\):  If
  \(j\) and \(j{+}1\) lie in different cycles of \(\sigma\), set
  \(e_j\gets 1\) and update \(\sigma\gets\sigma\,(j\,\,j{+}1)\).
   Otherwise set \(e_j\gets 0\) and leave \(\sigma\) unchanged.
\item
  Even indices (left factor \(E\)). For \(j=2,4,6,\dots,2m\):  If \(j\)
  and \(j{+}1\) lie in different cycles of the current \(\sigma\), set
  \(e_j\gets 1\) and update \(\sigma\gets\sigma\,(j\,\,j{+}1)\).  Else
  set \(e_j\gets 0\).
\end{enumerate}

Finally, let \(\pi\) be the permutation with these \(e_k\) in the
specified form:

\[
\pi \;=\; (2\,3)^{e_2}(4\,5)^{e_4}\cdots(2m\,\,2m{+}1)^{e_{2m}}\;\;(1\,2)^{e_1}(3\,4)^{e_3}\cdots(2m{-}1\,\,2m)^{e_{2m-1}}.
\]

Why this works.

\begin{itemize}
\item
  By the lemma and our ``if different, then apply'' rule, each chosen
  transposition reduces the number of cycles by exactly 1; when we skip,
  the number of cycles doesn't increase. Hence the number of cycles in
  \(\sigma\) is nonincreasing throughout.
\item
  Suppose after both passes \(\sigma\) still has at least two cycles.
  Then along the line \(1,2,\dots,2m{+}1\) there exists an adjacent pair
  \(j,j{+}1\) that lie in different cycles of the final \(\sigma\).
  Consider the pass in which index \(j\) was processed. Because the
  cycle count never increases later, \(j\) and \(j{+}1\) were also in
  different cycles at that time, so our rule would have applied the
  transposition \((j\,\,j{+}1)\), merging them---a contradiction.
  Therefore after the two passes \(\sigma=\rho\pi\) has exactly one
  cycle.
\item
  By construction, \(\pi\) has the required two-layer form, so it also
  satisfies (a): \(|\pi_k-k|\le 2\).
\end{itemize}

Even \(n\). If \(n\) is even, do exactly the same two passes, but the
``even layer'' stops at \(j=n-2\) (there is no pair \((n,n{+}1)\)).
Define

\[
\pi \;=\; (2\,3)^{e_2}(4\,5)^{e_4}\cdots(n{-}2\,\,n{-}1)^{e_{n-2}}\;\;(1\,2)^{e_1}(3\,4)^{e_3}\cdots(n{-}1\,\,n)^{e_{n-1}}.
\]

The identical argument shows \(\rho\pi\) is a single cycle, and again
\(|\pi_k-k|\le 2\).

Summary. (a) The two-layer ``odd then even adjacent swaps'' can move any
element by at most two positions, hence \(|\pi_k-k|\le 2\). (b) The
two-pass greedy rule (apply \((j\,\,j{+}1)\) iff it joins two different
cycles) yields exponents \(e_k\) so that \(\rho\pi\) is an \(n\)-cycle.
Thus every permutation is within right-multiplication by a \(\pi\) of
displacement at most 2 from a single cycle (``near'' a cycle).

\section{33 {[}M33{]}}\label{m33-2}

Let \(m=2^{2\ell}\) and \(n=2^{2\ell+1}=2m\). You must build \(m\) ``row
sequences'' \(\{\alpha_{i1},\dots,\alpha_{in}\}\) (one sequence for each
\(i\in\{0,\dots,m-1\}\)) and \(m\) ``column sequences''
\(\{\beta_{1j},\dots,\beta_{nj}\}\) (one sequence for each
\(j\in\{0,\dots,m-1\}\)), where every \(\alpha_{ik}\) and \(\beta_{kj}\)
is a permutation of \(\{1,2,3,4,5\}\), such that for all \(i,j\)

\[
\alpha_{i1}\beta_{1j}\,\alpha_{i2}\beta_{2j}\cdots \alpha_{in}\beta_{nj} \;=\;
\begin{cases}
(1\ 2\ 3\ 4\ 5), & \text{if } i=j,\\[2pt]
(), & \text{if } i\ne j,
\end{cases}
\]

i.e., you get a fixed 5-cycle on the diagonal and the identity off the
diagonal.

Key idea

Work inside the dihedral subgroup \(D_{10}\subset S_5\) generated by

\[
r=(1\ 2\ 3\ 4\ 5)\quad\text{and}\quad s=(2\ 5)(3\ 4).
\]

These satisfy the fundamental relation

\[
srs=r^{-1}\quad\Rightarrow\quad sr^k s=r^{-k}\ \text{ for all integers }k,
\]

and \(s^2=()\), \(r^5=()\). We'll interleave powers of \(r\) (coming
from the \(\alpha\)'s) with copies of \(s\) (coming from the
\(\beta\)'s) in a way that makes everything cancel except on the
diagonal, where one block collapses to \(r\).

The construction (explicit)

Index the \(n=2m\) positions as \(2t+1,2t+2\) for blocks
\(t=0,1,\dots,m-1\).

Row sequences \(\alpha_{i\bullet}\) (one for each
\(i\in\{0,\dots,m-1\}\)):

\begin{itemize}
\tightlist
\item
  In block \(t=i\): set \(\alpha_{i,2i+1}=r^{3}\) and
  \(\alpha_{i,2i+2}=r^{-3}\).
\item
  In every other block \(t\ne i\): set
  \(\alpha_{i,2t+1}=\alpha_{i,2t+2}=()\) (the identity).
\end{itemize}

Column sequences \(\beta_{\bullet j}\) (one for each
\(j\in\{0,\dots,m-1\}\)):

\begin{itemize}
\tightlist
\item
  In block \(t=j\): set \(\beta_{2j+1,j}=s\) and \(\beta_{2j+2,j}=s\).
\item
  In every other block \(t\ne j\): set
  \(\beta_{2t+1,j}=\beta_{2t+2,j}=()\).
\end{itemize}

(Using the identity is allowed, since it's a permutation of
\(\{1,2,3,4,5\}\).)

Why this works

Fix \(i,j\) and multiply block by block.

\begin{itemize}
\item
  If \(i\ne j\): In block \(t=i\) we have
  \(r^{3}\cdot ()\cdot r^{-3}\cdot ()=()\). In block \(t=j\) we have
  \(()\cdot s\cdot ()\cdot s = s^2=()\). All other blocks are all
  identities. Hence the total product is \(()\).
\item
  If \(i=j\): The only nontrivial factors lie in block \(t=i=j\), where
  the four-term product is

  \[
  r^{3}\, s\, r^{-3}\, s.
  \]

  Using \(s r^{-3} s = r^{3}\) (since \(sr^k s=r^{-k}\)), the block
  reduces to

  \[
  r^{3}\, r^{3}=r^{6}=r^{5}\cdot r= r=(1\ 2\ 3\ 4\ 5).
  \]

  All other blocks contribute the identity, so the full product is
  exactly \(r\).
\end{itemize}

This is precisely the required ``orthogonality'' property.

Notes and small checks

\begin{itemize}
\tightlist
\item
  We used only that \(n=2m\) (which follows from \(m=2^{2\ell}\),
  \(n=2^{2\ell+1}\)).
\item
  The choice of the exponent \(3\) is convenient because
  \(r^{3}s r^{-3}s=r^{2\cdot 3}=r\). More generally, any \(c\) with
  \(2c\equiv 1\pmod{5}\) (i.e., \(c\equiv 3\)) would work.
\item
  Any 5-cycle conjugate to \(r\) and any involution \(s\) that inverts
  \(r\) (i.e., \(srs=r^{-1}\)) give equivalent solutions by conjugation
  inside \(S_5\).
\end{itemize}

Tiny example (first nontrivial size)

Take \(\ell=1\Rightarrow m=4,\ n=8\). Rows \(i=0,1,2,3\) have
\(r^{3},r^{-3}\) only in their own block; columns \(j=0,1,2,3\) have
\(s,s\) only in their own block.

\begin{itemize}
\tightlist
\item
  For \(i=j=2\): the block \(t=2\) contributes \(r^{3}s r^{-3}s=r\);
  others are identities → total \(r\).
\item
  For \(i=1,\ j=3\): block \(t=1\) gives
  \(r^{3}\cdot()\cdot r^{-3}\cdot()=()\); block \(t=3\) gives
  \(()\cdot s\cdot ()\cdot s=()\); others are \(()\) → total \(()\).
\end{itemize}

Thus the construction meets the specification.

\section{34 {[}M25{]}}\label{m25-27}

We have an array split into two consecutive blocks, α of length m and β
of length n (total N = m+n). We want to transform the array from αβ to
βα, i.e., rotate the array left by m (equivalently, right by n).
Consider the permutation

\[
p(k) \equiv k + m \pmod{N}\qquad (0\le k < N)
\]

that maps each position to where its element must go after rotation.
Show that this ``cyclic-shift'' permutation has a simple cycle
structure, and exploit that structure to devise a simple (in-place)
algorithm for the rearrangement.

\begin{enumerate}
\def\labelenumi{\arabic{enumi})}
\tightlist
\item
  Cycle structure of a fixed shift on \(N\) items
\end{enumerate}

Let \(N=m+n\) and \(p(k)=k+m\pmod N\). Consider the orbit (cycle) of an
index \(r\):

\[
r \;\to\; r+m \;\to\; r+2m \;\to\; \cdots \;\to\; r+tm \pmod N.
\]

The cycle closes at the smallest \(t>0\) such that
\(tm\equiv 0\pmod N\). Let \(g=\gcd(m,N)\). Then
\(tm\equiv 0\ (\bmod\ N)\) iff \(N\mid tm\), whose least positive
solution is

\[
t=\frac{N}{\gcd(m,N)}=\frac{N}{g}.
\]

Therefore:

\begin{itemize}
\item
  All cycles have the same length \(L = N/g\).
\item
  There are exactly \(g=\gcd(m,N)\) disjoint cycles.
\item
  A set of cycle representatives is \(r=0,1,\dots,g-1\). The \(r\)-th
  cycle is

  \[
  C_r=\{\,r,\ r+m,\ r+2m,\ \dots,\ r+(L-1)m\ \} \pmod N.
  \]
\end{itemize}

\emph{Why this is ``simple''}: the permutation is a pure rotation by a
fixed stride \(m\) on the ring \(\mathbb{Z}_N\). The orbits are exactly
the cosets of the subgroup generated by \(m\), hence equal-length cycles
determined by \(\gcd(m,N)\).

Example. \(N=10\), \(m=4\), \(g=\gcd(10,4)=2\). The cycles are

\[
C_0=(0,4,8,2,6),\qquad C_1=(1,5,9,3,7),
\]

each of length \(10/2=5\). Rotating left by 4 maps each index to the
next in its cycle.

\begin{enumerate}
\def\labelenumi{\arabic{enumi})}
\setcounter{enumi}{1}
\tightlist
\item
  In-place algorithm from the cycle structure (``juggling'')
\end{enumerate}

The cycle structure tells us exactly how to move elements without extra
storage, using at most one temporary per cycle.

Algorithm (left-rotate by \(m\) positions).

\begin{verbatim}
N = m + n
g = gcd(m, N)
for r in 0 .. g-1:
    temp = A[r]
    j = r
    while True:
        k = (j + m) mod N
        if k == r:
            break
        A[j] = A[k]
        j = k
    A[j] = temp
\end{verbatim}

Correctness. The inner loop traverses the cycle \(C_r\) in order:
\(j = r, r+m, r+2m, \dots\). Because the cycle length is \(L=N/g\), we
visit each index in \(C_r\) exactly once and finally return to \(r\), at
which point we place the saved value \texttt{temp}. Since the cycles are
disjoint and cover \(\{0,\dots,N-1\}\), every element is moved exactly
once.

Cost. Each array position participates in exactly one assignment (plus
one final write per cycle), so the time is \(\Theta(N)\) and the extra
space is \(O(1)\) (only one temporary at a time). The number of cycles
is \(g=\gcd(m,N)\), so there are \(g\) temporaries used one after
another.

\begin{enumerate}
\def\labelenumi{\arabic{enumi})}
\setcounter{enumi}{2}
\tightlist
\item
  A different simple in-place method (three reversals)
\end{enumerate}

Another elegant \(O(N)\)-time, \(O(1)\)-space solution uses reversals:

\begin{enumerate}
\def\labelenumi{\arabic{enumi}.}
\tightlist
\item
  reverse the block α (positions 0..m−1),
\item
  reverse the block β (positions m..N−1),
\item
  reverse the whole array (0..N−1).
\end{enumerate}

Why it works: Writing strings,
\((\alpha\beta) \mapsto (\alpha^R\beta^R) \mapsto (\alpha^R\beta^R)^R = \beta\alpha\).
Each reversal can be done in place with two-pointer swaps. (This is a
standard trick and dovetails with the ``block transpositions'' theme
introduced by the exercise.)

\begin{enumerate}
\def\labelenumi{\arabic{enumi})}
\setcounter{enumi}{3}
\tightlist
\item
  Which to use?
\end{enumerate}

\begin{itemize}
\tightlist
\item
  Cycle-juggling (from the permutation's cycles): one pass over \(N\)
  items, \(O(1)\) extra space, straightforward once you know
  \(g=\gcd(m,N)\). This directly ``exploits the cycle structure'' the
  problem asks you to identify.
\item
  Three-reversal method: also \(O(N)\) time, \(O(1)\) space; often
  slightly simpler to code and branch-free.
\end{itemize}

Both meet the exercise's request for a simple algorithm derived from the
structure of the cyclic-shift permutation. The key structural fact is
the decomposition into \(g=\gcd(m,m+n)\) cycles, each of length \(N/g\).

(Optional) Small worked example

Let \(A=[a,b,c,d,e,f,g,h,i,j]\), \(m=4\) (α={[}a,b,c,d{]},
β={[}e,f,g,h,i,j{]}, \(N=10\)). \(g=\gcd(10,4)=2\). Two cycles:

\begin{itemize}
\tightlist
\item
  \(C_0=(0,4,8,2,6)\): \(A[0]\gets A[4]=e\), \(A[4]\gets A[8]=i\),
  \(A[8]\gets A[2]=c\), \(A[2]\gets A[6]=g\), \(A[6]\gets a\).
\item
  \(C_1=(1,5,9,3,7)\): \(A[1]\gets A[5]=f\), \(A[5]\gets A[9]=j\),
  \(A[9]\gets A[3]=d\), \(A[3]\gets A[7]=h\), \(A[7]\gets b\).
\end{itemize}

Result: \([e,f,g,h,i,j,a,b,c,d]=\beta\alpha\).

Summary

\begin{itemize}
\tightlist
\item
  Rotation αβ → βα corresponds to the permutation \(p(k)=k+m\pmod{N}\).
\item
  Its cycle structure is \(g=\gcd(m,N)\) cycles of equal length \(N/g\).
\item
  Traversing those cycles yields an in-place juggling algorithm:
  \(\Theta(N)\) time, \(O(1)\) extra space.
\item
  The three-reversal method is an equally simple in-place alternative.
\end{itemize}

\section{35 {[}M39{]}}\label{m39}

Let an array be split into three consecutive blocks
\(\alpha \beta \gamma\) of lengths \(l,m,n\) (so total \(M=l+m+n\)). We
want to rearrange in place to \(\gamma \beta \alpha\). Show that the
corresponding permutation has a simple cycle structure that leads to an
efficient algorithm. (Hint in the book: reduce
\((\alpha\beta)(\gamma\beta)\) to \((\gamma\beta)(\alpha\beta)\).)

\emph{(For context, Exercise 34 (the 2-block case) rewrites
\(\alpha\beta\) to \(\beta\alpha\) via the permutation
\(k\mapsto (k+m)\bmod(m+n)\) and exploits its cycle structure. We'll
generalize that idea.)}

Embed \(\alpha\beta\gamma\) into a \emph{virtual} longer string
\(\alpha\beta\gamma\beta\) and perform a single circular shift by
\(l+m\) positions on that longer string. If we then ``ignore'' the extra
trailing \(\beta\), the first \(M=l+m+n\) positions have become exactly
\(\gamma\beta\alpha\).

Formally, set

\[
M=l+m+n,\qquad N=l+2m+n=M+m,\qquad s=l+m.
\]

Consider the rotation permutation on \(\{0,\dots,N-1\}\),

\[
h(k)= (k+s)\bmod N .
\]

On the \emph{virtual} array \(\alpha\beta\gamma\beta\) (length \(N\)),
\(h\) moves every cell forward by \(s\) positions. The desired 3-block
rearrangement on the \emph{real} array of length \(M\) is obtained by
restricting \(h\) to indices \textless{} \(M\) and ``skipping over'' any
visit to the shadow region \([M, N)\) (the duplicated \(\beta\)). The
result is a permutation of \(\{0,\dots,M-1\}\) whose cycles we can
follow in place.

The book's answer states this succinctly and also gives the cycle count;
we'll derive and then match it.

Cycle structure and how many cycles

The rotation \(h(k)=(k+s)\bmod N\) decomposes \(\{0,\dots,N-1\}\) into

\[
d=\gcd(s,N)=\gcd(l+m,\;l+2m+n)=\gcd(l+m,\;m+n)
\]

cycles, each of length \(N/d\). (The equality of the gcd's is immediate
from \(\gcd(s,N)=\gcd(l+m,(l+2m+n)-(l+m))=\gcd(l+m,m+n)\).)

When we strike out the indices in the shadow tail \([M,N)\) from each
rotation cycle and connect neighbors across the removed nodes, we still
get cycles---now on the real index set \(\{0,\dots,M-1\}\). Thus the
3-block permutation for \(\alpha\beta\gamma\to\gamma\beta\alpha\) has
exactly

\[
\boxed{\,d=\gcd(l+m,\;m+n)\,}
\]

cycles; walking those cycles gives an in-place algorithm. This is
precisely the conclusion in the official answer.

An in-place cycle-walking algorithm

Define a ``skip-aware next'' function on real positions \(0\le j<M\):

\[
g(j)=
\begin{cases}
(j+s)\bmod N,&\text{if }(j+s)\bmod N < M;\\[2pt]
(j+2s)\bmod N,&\text{otherwise} \quad(\text{skip over the shadow \(\beta\)}).
\end{cases}
\]

Then the desired permutation is exactly ``send \(j\) to \(g(j)\)''.
Implement it by cycling:

\begin{verbatim}
s = l + m
M = l + m + n
N = l + 2m + n
visited[0..M-1] = false   // or do gcd-based leaders to avoid a visited array

for start in 0..d-1:         // d = gcd(l+m, m+n)
    if visited[start]: continue
    t = A[start]             // hold one element
    j = start
    while True:
        // compute next with “skip”
        j1 = (j + s) % N
        if j1 >= M:
            j1 = (j1 + s) % N    // skip shadow β
        if j1 == start:
            A[j] = t             // close the cycle
            break
        A[j] = A[j1]
        visited[j] = true
        j = j1
\end{verbatim}

Why the two steps when \texttt{j1\ \textgreater{}=\ M}? In the virtual
rotation
\(\alpha\beta\gamma\beta \xrightarrow{+s}\gamma\beta\alpha\beta\),
landing in \([M,N)\) means ``you're at the duplicated \(\beta\)''. The
hint in the text says to realize that the two assignments
\(x_k\gets x_{k'}\) and \(x_{k'}\gets x_{k''}\) (with
\(k'=(k+s)\bmod N\ge M\), \(k''=(k'+s)\bmod N\)) effectively place into
\(x_k\) the correct source \(x_{k''}\), because \(x_{k'}=x_{k''}\) (both
refer to the same logical \(\beta\)). So we jump an extra \(s\).

Complexity and storage

\begin{itemize}
\tightlist
\item
  Each real position is visited exactly once across all cycles, so the
  algorithm performs \(M\) moves (plus at most \(d\) temporary saves).
\item
  Time \(\Theta(M)\), extra space \(O(1)\) beyond one saved element (or
  \(O(d)\) if you prefer one temp per cycle leader).
\end{itemize}

This is the same optimal profile as the 2-block case (Exercise 34),
generalized.

Worked example

Take \(l=2,m=3,n=1\). So \(\alpha=\,[a_0,a_1]\),
\(\beta=\,[b_0,b_1,b_2]\), \(\gamma=\,[c_0]\). Then \(M=6\), \(N=9\),
\(s=l+m=5\), \(d=\gcd(5,\,m+n=4)=1\): one big cycle.

Indices \(0..5\) hold \(\alpha\beta\gamma\) =
\([a_0,a_1,b_0,b_1,b_2,c_0]\). Virtual tail \([6,7,8]\) holds a
duplicate \(\beta\) = \([b_0,b_1,b_2]\).

Walk with skips:

\begin{itemize}
\tightlist
\item
  \(0 \to (0+5)=5<6\) ⇒ put \(A[0]\gets A[5]=c_0\)
\item
  \(5 \to (5+5)=10\bmod 9=1<6\) ⇒ \(A[5]\gets A[1]=a_1\)
\item
  \(1 \to (1+5)=6\ge6\) ⇒ skip to \((6+5)=2\) ⇒ \(A[1]\gets A[2]=b_0\)
\item
  \(2 \to (2+5)=7\ge6\) ⇒ skip to \((7+5)=3\) ⇒ \(A[2]\gets A[3]=b_1\)
\item
  \(3 \to (3+5)=8\ge6\) ⇒ skip to \((8+5)=4\) ⇒ \(A[3]\gets A[4]=b_2\)
\item
  \(4 \to (4+5)=0<6\) ⇒ \(A[4]\gets A[0]=a_0\)
\item
  Next is \(0\) (start), so write back the saved \(t\) to \(A[4]\) at
  the close, yielding \([c_0,b_0,b_1,b_2,a_0,a_1]\) which is
  \(\gamma\beta\alpha\).
\end{itemize}

\section{36 {[}27{]}}\label{section-149}

Write a MIX subroutine that performs the in-place rearrangement

\[
\alpha\beta\gamma \longrightarrow \gamma\beta\alpha
\]

using the efficient ``cycle-structure'' method from Exercise 35, then
analyze its running time and compare it with the simpler
``four-reversal'' trick that goes
\(\alpha\beta\gamma \to (\alpha\beta\gamma)^R=\gamma^R\beta^R\alpha^R \to \gamma\beta\alpha\).

\begin{enumerate}
\def\labelenumi{\arabic{enumi})}
\tightlist
\item
  Algorithm (from Ex. 35, made explicit)
\end{enumerate}

The key idea is to compose two adjacent block cyclic shifts of the form
\(XY\mapsto YX\) (the routine from Ex. 34), chosen so the result is
\(\gamma\beta\alpha\):

\begin{enumerate}
\def\labelenumi{\arabic{enumi}.}
\tightlist
\item
  First move \(\alpha\) past the block \(\beta\gamma\):
\end{enumerate}

\[
\underbrace{\alpha}_{\ell}\ \underbrace{\beta\gamma}_{m+n}
\;\overset{XY\to YX}{\Longrightarrow}\; \beta\gamma\,\alpha.
\]

\begin{enumerate}
\def\labelenumi{\arabic{enumi}.}
\setcounter{enumi}{1}
\tightlist
\item
  Then, inside the new prefix, move \(\beta\) past \(\gamma\):
\end{enumerate}

\[
\underbrace{\beta}_{m}\ \underbrace{\gamma}_{n}
\;\overset{XY\to YX}{\Longrightarrow}\; \gamma\beta.
\]

Composing the two gives \(\gamma\beta\alpha\) in place, with \(O(1)\)
extra storage. The hint in Ex. 35 (``change
\((\alpha\beta)(\gamma\beta)\) to \((\gamma\beta)(\alpha\beta)\)'') is
exactly this ``swap-by-adjacent-blocks'' reasoning.

So we only need a solid two-block rotation subroutine
\(\textsf{ROT2}(X,Y)\!:\ XY\to YX\) (the Ex. 34 routine), and call it
twice.

Efficient \(\textsf{ROT2}\) via cycle decomposition. Let
\(|X|=s,\ |Y|=t,\ L=s+t,\ d=\gcd(s,L)\). Partition positions
\(0,\dots,L-1\) into \(d\) cycles under the map
\(k\mapsto (k+s)\bmod L\). For each cycle, rotate the elements once
using one temporary word. Each element is moved exactly once; each cycle
incurs one extra store. This is the classic ``juggling''/cycle method
for array rotation. (This is the same structure you proved in Ex.
34/35.)

\begin{enumerate}
\def\labelenumi{\arabic{enumi})}
\setcounter{enumi}{1}
\tightlist
\item
  MIX implementation
\end{enumerate}

Below is a self-contained MIX subroutine that performs
\(\alpha\beta\gamma \to \gamma\beta\alpha\) using two calls to
\(\textsf{ROT2}\). Conventions:

\begin{itemize}
\item
  Memory holds the array, one item per word.
\item
  Inputs on entry:

  \begin{itemize}
  \tightlist
  \item
    \texttt{I1} = base address of \(\alpha\),
  \item
    \texttt{I2} = \(\ell = |\alpha|\),
  \item
    \texttt{I3} = \(m = |\beta|\),
  \item
    \texttt{I4} = \(n = |\gamma|\).
  \end{itemize}
\item
  Registers used: \texttt{A}, \texttt{X}, \texttt{I1..I6}. One temp word
  stored at label \texttt{TMP}.
\item
  Subroutine returns with the array rearranged in place.
\end{itemize}

\begin{Shaded}
\begin{Highlighting}[]
\NormalTok{* {-}{-}{-}{-}{-}{-}{-}{-}{-}{-}{-}{-}{-}{-}{-}{-}{-}{-}{-}{-}{-}{-}{-}{-}{-}{-}{-}{-}{-}{-}{-}{-}{-}{-}{-}{-}{-}{-}{-}{-}{-}{-}{-}{-}{-}{-}{-}{-}{-}{-}{-}{-}{-}{-}{-}{-}{-}{-}{-}{-}}
\NormalTok{* ROT3: αβγ {-}\textgreater{} γβα   (uses ROT2 twice)}
\NormalTok{* Inputs: I1=BASE, I2=LEN\_A(ℓ), I3=LEN\_B(m), I4=LEN\_C(n)}
\NormalTok{* Clobbers: A,X,I1..I6; uses word TMP}
\NormalTok{* {-}{-}{-}{-}{-}{-}{-}{-}{-}{-}{-}{-}{-}{-}{-}{-}{-}{-}{-}{-}{-}{-}{-}{-}{-}{-}{-}{-}{-}{-}{-}{-}{-}{-}{-}{-}{-}{-}{-}{-}{-}{-}{-}{-}{-}{-}{-}{-}{-}{-}{-}{-}{-}{-}{-}{-}{-}{-}{-}{-}}
\NormalTok{ROT3   ENT5 0           * I5 = scratch}
\NormalTok{       ENT6 0           * I6 = scratch}
\NormalTok{       LDA  I2          * A := ℓ}
\NormalTok{       ADD  I3          * A := ℓ + m}
\NormalTok{       ADD  I4          * A := ℓ + m + n = L}
\NormalTok{       STA  LSUM        * store total length if desired (not strictly needed)}

\NormalTok{* Step 1: rotate α past (βγ):  [α | βγ] {-}\textgreater{} [βγ | α]}
\NormalTok{* Call ROT2(BASE, s=ℓ, t=m+n)}
\NormalTok{       ENTA 0}
\NormalTok{       ENT3 0           * (we’ll pass s in I2 and t in I3 to ROT2)}
\NormalTok{       LDX  I2          * X := ℓ}
\NormalTok{       STX  SVAL}
\NormalTok{       LDA  I3          * A := m}
\NormalTok{       ADD  I4          * A := m+n}
\NormalTok{       STA  TVAL}
\NormalTok{       LD1  I1          * base unchanged}
\NormalTok{       LD2  SVAL        * I2 := s = ℓ}
\NormalTok{       LD3  TVAL        * I3 := t = m+n}
\NormalTok{       JMP  ROT2        * fall into ROT2 (subroutine below)}
\NormalTok{RET1   NOP              * return point from ROT2}

\NormalTok{* Step 2: rotate β past γ in the prefix:  [β | γ] {-}\textgreater{} [γ | β]}
\NormalTok{* Call ROT2(BASE, s=m, t=n)}
\NormalTok{       LD1  I1}
\NormalTok{       LD2  I3          * I2 := s = m}
\NormalTok{       LD3  I4          * I3 := t = n}
\NormalTok{       JMP  ROT2}
\NormalTok{RET2   NOP              * return from ROT2}

\NormalTok{       JMP  EXIT}

\NormalTok{* {-}{-}{-}{-}{-}{-}{-}{-}{-}{-}{-}{-}{-}{-}{-}{-}{-}{-}{-}{-}{-}{-}{-}{-}{-}{-}{-}{-}{-}{-}{-}{-}{-}{-}{-}{-}{-}{-}{-}{-}{-}{-}{-}{-}{-}{-}{-}{-}{-}{-}{-}{-}{-}{-}{-}{-}{-}{-}{-}{-}}
\NormalTok{* ROT2: two{-}block in{-}place rotation [X | Y] {-}\textgreater{} [Y | X]}
\NormalTok{* Inputs: I1=BASE, I2=s=|X|, I3=t=|Y|}
\NormalTok{* Clobbers: A,X,I1..I6, TMP}
\NormalTok{* Returns to RET1 (first call) or RET2 (second call)}
\NormalTok{* {-}{-}{-}{-}{-}{-}{-}{-}{-}{-}{-}{-}{-}{-}{-}{-}{-}{-}{-}{-}{-}{-}{-}{-}{-}{-}{-}{-}{-}{-}{-}{-}{-}{-}{-}{-}{-}{-}{-}{-}{-}{-}{-}{-}{-}{-}{-}{-}{-}{-}{-}{-}{-}{-}{-}{-}{-}{-}{-}{-}}
\NormalTok{ROT2   LDA  I2          * A := s}
\NormalTok{       ADD  I3          * A := s+t = L}
\NormalTok{       STA  L           }

\NormalTok{* Compute d = gcd(s, L) (Euclid’s algorithm).}
\NormalTok{* We\textquotesingle{}ll store s in S, L in L; work in A/X for modulus.}
\NormalTok{       STA  L}
\NormalTok{       LDA  I2          * A := s}
\NormalTok{       STA  S}
\NormalTok{G1     LDA  S}
\NormalTok{       ENTX L           * X := L}
\NormalTok{       DIV  L           * A := floor(S/L), X := S mod L}
\NormalTok{       LDA  L}
\NormalTok{       STA  S           * S := L}
\NormalTok{       LDX  X           * X := old (S mod L)}
\NormalTok{       STX  L}
\NormalTok{       JXZ  G2          * if remainder 0 {-}\textgreater{} gcd found in S}
\NormalTok{       JMP  G1}
\NormalTok{G2     LDA  S}
\NormalTok{       STA  D           * D := gcd(s, s+t)}

\NormalTok{* For i = 0..D{-}1: rotate one cycle starting at offset i}
\NormalTok{       ENTA 0}
\NormalTok{       STA  I           * I := 0}
\NormalTok{CYC0   LDA  I}
\NormalTok{       CMPA D}
\NormalTok{       JGE  DONE        * if I \textgreater{}= D, all cycles done}

\NormalTok{* tmp := M[BASE + I]}
\NormalTok{       LD1  I1}
\NormalTok{       LDX  I}
\NormalTok{       ADD  I1}
\NormalTok{       STX  P           * P := BASE + i}
\NormalTok{       LDA  P,0}
\NormalTok{       STA  TMP}

\NormalTok{* j := i}
\NormalTok{       LDA  I}
\NormalTok{       STA  J}

\NormalTok{* cycle: j \textless{}{-} j + s (mod L); move element at \textquotesingle{}next\textquotesingle{} into position j}
\NormalTok{CYC1   LDA  J}
\NormalTok{       ADD  S}
\NormalTok{       STA  NXT         * NXT := j + s}
\NormalTok{       LDA  NXT}
\NormalTok{       CMPA L}
\NormalTok{       JL   OK          * if NXT \textless{} L, keep}
\NormalTok{       SUB  L           * else subtract L}
\NormalTok{       STA  NXT}
\NormalTok{OK     LDA  NXT}
\NormalTok{       CMPA I}
\NormalTok{       JE   CLOSE       * if next == i, cycle ends}

\NormalTok{* M[BASE + j] := M[BASE + next]}
\NormalTok{       LD1  I1}
\NormalTok{       LDX  J}
\NormalTok{       ADD  I1}
\NormalTok{       STX  P}
\NormalTok{       LDA  I1}
\NormalTok{       ADD  NXT}
\NormalTok{       STX  Q           * Q := BASE + next}
\NormalTok{       LDA  Q,0}
\NormalTok{       STA  P,0}

\NormalTok{* j := next; continue}
\NormalTok{       LDA  NXT}
\NormalTok{       STA  J}
\NormalTok{       JMP  CYC1}

\NormalTok{* close the cycle: M[BASE + j] := tmp}
\NormalTok{CLOSE  LD1  I1}
\NormalTok{       LDX  J}
\NormalTok{       ADD  I1}
\NormalTok{       STX  P}
\NormalTok{       LDA  TMP}
\NormalTok{       STA  P,0}

\NormalTok{* next start i := i+1}
\NormalTok{       LDA  I}
\NormalTok{       INC  1}
\NormalTok{       STA  I}
\NormalTok{       JMP  CYC0}

\NormalTok{DONE   NOP}
\NormalTok{* Return to proper caller: distinguish by a flag (or fall through two times).}
\NormalTok{* Here we simply fall through twice; the caller placed RET1 or RET2 in{-}line.}
\NormalTok{       JMP  RET1}

\NormalTok{* {-}{-}{-}{-}{-}{-}{-}{-} data / temps {-}{-}{-}{-}{-}{-}{-}{-}}
\NormalTok{TMP    CON  0}
\NormalTok{S      CON  0}
\NormalTok{L      CON  0}
\NormalTok{D      CON  0}
\NormalTok{I      CON  0}
\NormalTok{J      CON  0}
\NormalTok{NXT    CON  0}
\NormalTok{P      CON  0}
\NormalTok{Q      CON  0}
\NormalTok{SVAL   CON  0}
\NormalTok{TVAL   CON  0}
\NormalTok{LSUM   CON  0}

\NormalTok{EXIT   END  ROT3}
\end{Highlighting}
\end{Shaded}

\emph{(Any standard MIX calling convention variant is fine. The
essential part is \texttt{ROT2} using the cycle/juggling method;
\texttt{ROT3} is just two calls to it.)}

\begin{enumerate}
\def\labelenumi{\arabic{enumi})}
\setcounter{enumi}{2}
\tightlist
\item
  Running time
\end{enumerate}

Let \(L=\ell+m+n\).

\begin{itemize}
\item
  One call to \texttt{ROT2(s,t)} with the cycle method:

  \begin{itemize}
  \tightlist
  \item
    There are \(d=\gcd(s,s+t)\) cycles; the \(k\)-th cycle of length
    \((s+t)/d\) does exactly that many moves plus a final write of the
    held temporary.
  \item
    Total moves (loads/stores of array words) are \((s+t) + d\).
  \item
    No extra memory except one temporary word.
  \end{itemize}
\item
  Our \texttt{ROT3} does two such rotations:

  \begin{itemize}
  \item
    First with \((s,t)=(\ell,\,m+n)\): cost \(L+\gcd(\ell,L)\).
  \item
    Then with \((s,t)=(m,\,n)\): cost \((m+n)+\gcd(m,m+n)\).
  \item
    Total:

    \[
    T_{\textsf{ROT3}} \;=\; 2L \;+\; \gcd(\ell,L) \;+\; \gcd(m,m+n).
    \]
  \end{itemize}

  In practice this is very close to \(2L\) word moves.
\end{itemize}

(``Word moves'' maps well to MIX's load/store traffic; if you wish, you
can multiply by the cost model of §1.3 to get an instruction-level
bound.)

\begin{enumerate}
\def\labelenumi{\arabic{enumi})}
\setcounter{enumi}{3}
\tightlist
\item
  Comparison with the ``reverse-everything'' method
\end{enumerate}

The suggested simpler method is:

\begin{enumerate}
\def\labelenumi{\arabic{enumi}.}
\tightlist
\item
  reverse the whole \(\alpha\beta\gamma\) (length \(L\)),
\item
  then reverse each of the three blocks individually to get
  \(\gamma\beta\alpha\).
\end{enumerate}

A reversal of a block of length \(s\) does \(\lfloor s/2\rfloor\) swaps,
i.e., about \(1.5s\) array word moves (each swap = 3 writes with a
temporary). Therefore the total array writes are

\[
T_{\textsf{rev}} \;=\; 3\!\left(\left\lfloor\frac{L}{2}\right\rfloor
+ \left\lfloor\frac{\ell}{2}\right\rfloor
+ \left\lfloor\frac{m}{2}\right\rfloor
+ \left\lfloor\frac{n}{2}\right\rfloor\right)
\;\approx\; 3L\ \text{(when lengths are even)}.
\]

So, asymptotically both are linear and in-place, but the cycle-based
two-rotation method is typically \textasciitilde⅓ fewer word moves
(≈\(2L\) vs.~≈\(3L\)) and often measurably faster on MIX.

\begin{enumerate}
\def\labelenumi{\arabic{enumi})}
\setcounter{enumi}{4}
\tightlist
\item
  Correctness sketch
\end{enumerate}

\begin{itemize}
\item
  Each \texttt{ROT2} is a permutation that maps index \(k\) to
  \(k+s\pmod{s+t}\), i.e., exactly the desired \(XY\mapsto YX\)
  rearrangement; the cycle procedure performs that permutation by
  rotating each cycle once.
\item
  Composing the two adjacent-block permutations

  \[
  (\alpha\ |\ \beta\gamma)\mapsto(\beta\gamma\ |\ \alpha)
  \quad\text{and}\quad
  (\beta\ |\ \gamma)\mapsto(\gamma\ |\ \beta)
  \]

  yields \(\gamma\beta\alpha\) as required. This is the ``convenient
  cycle structure'' behind Ex. 35's hint.
\end{itemize}

\section{37 {[}M26{]}}\label{m26-5}

Show that for \(n\ge 2\) every even permutation \(\pi\in S_n\) can be
written as a product of two \(n\)-cycles. (Knuth's brief answer hints at
reducing to a factorization by two \((n-1)\)-cycles and then
``splicing'' \(n\) into each cycle.)

Key facts

\begin{itemize}
\tightlist
\item
  The sign of an \(\ell\)-cycle is \((-1)^{\ell-1}\). Hence the product
  of two cycles of the same length is always an even permutation.
\item
  Conjugation preserves cycle structure ``up to relabeling'' and
  preserves parity.
\end{itemize}

Step 1 --- Reduce to a permutation on \(\{1,\dots,n-1\}\)

If \(n\le 2\) the statement is immediate. Assume \(n\ge 3\).

Choose \(a,b\in\{1,\dots,n-1\}\) with \(\pi(a)=b\). (Such a pair exists
since not all \(1,\dots,n-1\) can map to \(n\).)

Consider the conjugate

\[
\tau \stackrel{\text{def}}{=} (n\ a)\,\pi\,(b\ n).
\]

Because of the chosen \(a\to b\), \(\tau\) fixes \(n\); hence \(\tau\)
acts as a permutation of \(\{1,\dots,n-1\}\). Moreover,
\(\operatorname{sgn}(\tau)=\operatorname{sgn}(\pi)\), since
transpositions on the left and right cancel in parity. So if \(\pi\) is
even, then \(\tau\) is an even permutation of \(n-1\) symbols. Knuth's
answer points exactly to this setup.

Step 2 --- Lemma: every even permutation of \(m\ge 2\) symbols is a
product of two \(m\)-cycles

Let \(m\ge 2\) and \(\sigma\in S_m\) be even. Write \(\sigma\) in
disjoint cycles

\[
\sigma=(c_{11}\,c_{12}\,\dots\,c_{1k_1})\,(c_{21}\,c_{22}\,\dots\,c_{2k_2})\cdots(c_{r1}\,\dots\,c_{rk_r}),
\quad k_1+\cdots+k_r=m.
\]

Construct an \(m\)-cycle \(U\) by weaving the cycles columnwise:

\[
U=(c_{11}\,c_{21}\,\dots\,c_{r1}\,c_{12}\,c_{22}\,\dots\,c_{r2}\,\dots),
\]

skipping entries when a row runs out. Define \(V:=U^{-1}\sigma\). One
checks directly (follow any element through \(\sigma\) then \(U^{-1}\))
that \(V\) is also a single \(m\)-cycle. Therefore \(\sigma=UV\) with
both \(U,V\) \(m\)-cycles.

(Parity necessity: The product of two \(m\)-cycles is even; thus only
even \(\sigma\) can be represented this way. The weaving construction
shows sufficiency.)

Step 3 --- Apply the lemma with \(m=n-1\)

Apply the lemma to \(\tau\in S_{n-1}\): There exist \((n-1)\)-cycles
\(U=(\alpha\,a)\) and \(V=(b\,\beta)\) on \(\{1,\dots,n-1\}\) such that

\[
\tau = U\,V.
\]

This is precisely the structure suggested in Knuth's short answer.

Step 4 --- ``Splice'' \(n\) to get two \(n\)-cycles

Now undo the conjugation:

\[
\pi=(n\,a)\,\tau\,(b\,n) \;=\; (n\,a)\,(\alpha\,a)\,(b\,\beta)\,(b\,n).
\]

Observe the two ``splices''

\[
(n\,a)\,(\alpha\,a)=(n\,\alpha\,a)\quad\text{and}\quad (b\,\beta)\,(b\,n)=(b\,\beta\,n),
\]

which are both \(n\)-cycles (we have inserted \(n\) into each
\((n-1)\)-cycle at the occurrence of \(a\) and \(b\)). Hence

\[
\boxed{\ \pi=(n\,\alpha\,a)\,(b\,\beta\,n)\ },
\]

a product of two \(n\)-cycles, as required. This matches the equivalence
written in Knuth's official answer: \((n a)\pi(b n)=(\alpha a)(b\beta)\)
iff \(\pi=(n\alpha a)(b\beta n)\).

Worked example (small \(n\))

Take \(n=5\) and \(\pi=(1\,2)(3\,4)\) (even). Choose \(a=1\), \(b=2\).
Then

\[
\tau=(5\,1)\,\pi\,(2\,5)
\]

fixes 5 and acts on \(\{1,2,3,4\}\) as \((1\,2)(3\,4)\) (still even).
Apply the lemma with \(m=4\): one valid choice is

\[
U=(1\,4\,2\,3),\qquad V=U^{-1}\tau=(1\,4\,2\,3),
\]

so \(\tau=UV\). Splice \(5\) in:

\[
(n\,\alpha\,a)=(5\,4\,2\,3\,1),\qquad (b\,\beta\,n)=(2\,4\,2\,3\,5)=(2\,3\,5\,1\,4)\ \text{(after relabeling within \(\{1,2,3,4,5\}\))}.
\]

One checks that their product is \(\pi\). (There are many valid \(U,V\);
any pair of 4-cycles with \(UV=\tau\) produces two 5-cycles whose
product is \(\pi\).)

Why the statement must be about even permutations

Each \(n\)-cycle has sign \((-1)^{n-1}\); hence the product of two
\(n\)-cycles is always even. So odd permutations cannot be expressed
this way. Our construction covers all even permutations and only
those---exactly the claim of the exercise. Knuth's answer encapsulates
this reduction and splice in one line.

\bookmarksetup{startatroot}

\chapter{1.4.1 Subroutines}\label{subroutines}

\subsection{1 {[}10{]}}\label{section-150}

Calling sequence (two entrances)

\begin{itemize}
\tightlist
\item
  Special case \(n=100\): \texttt{JMP\ MAX100} (sets
  \(rI3 \leftarrow 100\) internally via \texttt{ENT3\ 100} in (5)). 
\item
  General case \(n\) arbitrary: \texttt{ENT3\ n} ; \texttt{JMP\ MAXN}
  (caller supplies \(n\) in \(rI3\), then enters at \texttt{MAXN}). 
\end{itemize}

Entry conditions

\begin{itemize}
\tightlist
\item
  For \texttt{MAX100} entrance: No register precondition for \(n\); the
  subroutine itself performs \texttt{ENT3\ 100}.
\item
  For \texttt{MAXN} entrance: \(rI3 = n\) with \(n \ge 1\). (These
  mirror the role of \(n\) in the general MAXIMUM spec given in (4), but
  with \(rI3\) used as the parameter in this version.) 
\end{itemize}

Exit conditions (functional result and register effects)

Let \(X\) be the base address used by the code (as in the MAX programs),
and let the table be \(X+1,\dots,X+n\). Then on return:

\begin{itemize}
\tightlist
\item
  Result: \(rA \;=\; \max_{1 \le k \le n} \mathrm{CONTENTS}(X+k)\) and
  \(\mathrm{CONTENTS}(X+rI2) = rA\) (so \(rI2\) gives a position at
  which the maximum occurs). (Same functional contract as in (4).) 
\item
  Loop index cleanup: \(rI3 = 0\) on exit (as in (4)). 
\item
  Other status: \(rJ\) (the jump register) and the comparison indicator
  (CI) are affected by the subroutine call/return sequence (as noted in
  (4)). 
\item
  No external stores implied: The specification does not imply changes
  to memory outside the subroutine's own locations; in (4) Knuth notes
  we may conclude nothing has been stored externally when the contract
  doesn't say so. 
\end{itemize}

Linkage notes (how (5) returns)

\begin{itemize}
\tightlist
\item
  \begin{enumerate}
  \def\labelenumi{(\arabic{enumi})}
  \setcounter{enumi}{4}
  \tightlist
  \item
    saves the return address with \texttt{STJ\ EXIT} at entry (for
    \texttt{MAXN}) and returns with \texttt{EXIT\ JMP\ *}, which jumps
    back to the instruction following the caller's \texttt{JMP}. For the
    \texttt{MAX100} entrance, \texttt{ENT3\ 100} is done first, but the
    same \texttt{EXIT} return is used; \texttt{ENT3} doesn't disturb
    \(rJ\). 
  \end{enumerate}
\end{itemize}

Compact spec (in the exact style of (4))

\begin{itemize}
\item
  Calling sequence: \texttt{JMP\ MAX100}  or 
  \texttt{ENT3\ n\ ;\ JMP\ MAXN}. 
\item
  Entry conditions: For \texttt{MAX100}: none (internally sets
  \(rI3 \leftarrow 100\)). For \texttt{MAXN}: \(rI3 = n\), assume
  \(n \ge 1\).
\item
  Exit conditions:
  \(rA = \displaystyle \max_{1 \le k \le n} \mathrm{CONTENTS}(X+k) = \mathrm{CONTENTS}(X+rI2)\);
  \(rI3 = 0\); \(rJ\) and CI are also affected; no other state is
  changed. 
\end{itemize}

\subsection{2 {[}10{]}}\label{section-151}

Replace this use of \texttt{JSJ} (jump without altering \texttt{rJ})
with ordinary instructions.

Solution

Use the same two-instruction prologue that label \texttt{MAXN}
uses---save the return address and jump into the body---then set the
parameter and continue:

\begin{verbatim}
MAX50  STJ  EXIT    ; save caller's return address (rJ) where MAXN would save it
       ENT3 50      ; set rI3 = 50 (parameter n = 50)
       JMP  2F      ; continue at the same point as other entries
\end{verbatim}

Why this works

\begin{itemize}
\tightlist
\item
  At entry \texttt{MAXN}, the subroutine begins by \texttt{STJ\ EXIT}
  then \texttt{JMP\ 2F} to the body. By duplicating those two actions at
  \texttt{MAX50}, we avoid needing a \texttt{JSJ}.
\item
  \texttt{ENT3} \emph{does not} change the J-register (Knuth notes this
  explicitly), so you may place \texttt{ENT3\ 50} before or after
  \texttt{STJ\ EXIT}; either way the saved return address is the
  caller's, as desired.
\item
  After \texttt{JMP\ 2F}, execution is identical to the other entrances;
  the subroutine eventually returns via the saved address at
  \texttt{EXIT}.
\end{itemize}

This exactly satisfies the exercise's request to substitute for (6)
without using \texttt{JSJ}.

\subsection{3 {[}M15{]}}\label{m15-15}

Knuth's spec (4) for \texttt{MAXIMUM} is (summarizing):

Calling sequence: \texttt{JMP\ MAXIMUM}. Entry: \texttt{rI1\ =\ n}
(assume \texttt{n\ ≥\ 1}). Exit:
\texttt{rA\ =\ max\_\{1≤k≤n\}\ X{[}k{]}\ =\ X{[}rI2{]}}, and
\texttt{rI3\ =\ 0}. It notes that ``\texttt{rJ} and CI are also
affected'' but leaves them unspecified. 

The code pattern is the usual ``find max'' loop, with a return via
\texttt{EXIT\ JMP\ *}. (See the \texttt{MAX100} variant (1) and the
generalized forms in the same section.) 

Precisely: effect on \texttt{rJ}

\begin{itemize}
\tightlist
\item
  On entry (caller does \texttt{JMP\ MAXIMUM}), \texttt{rJ} holds the
  caller's return address (the location after the call), per MIX jump
  semantics.
\item
  On exit the subroutine executes \texttt{JMP\ *} at label
  \texttt{EXIT}. In MIX, \texttt{JMP\ *} returns to the address
  currently in \texttt{rJ} and stores the address of the instruction
  following \texttt{EXIT} (i.e., \texttt{EXIT+1}) into \texttt{rJ} as
  part of the jump.
\item
  Net effect seen by the caller: control returns to ``after the call'',
  but \texttt{rJ} is clobbered to an internal address (\texttt{EXIT+1}
  from inside the callee). Thus \texttt{rJ} is not preserved by the
  subroutine. (This is why (4) warned that \texttt{rJ} is affected.) 
\end{itemize}

Precisely: effect on the comparison indicator (CI)

Look at the inner loop for \texttt{MAXIMUM}/\texttt{MAX100} (labels
M3--M5): each iteration does a \texttt{CMPA\ X,3} and possibly updates
the running maximum, then decrements the loop index and exits when it is
no longer positive. Therefore, the last instruction that sets CI before
return is the final \texttt{CMPA} in the loop.

\begin{itemize}
\tightlist
\item
  At exit, \texttt{rA} holds the maximum \texttt{m}, and the final CI
  equals the result of comparing \texttt{m} with the last element that
  was compared (namely \texttt{X{[}1{]}} in the typical 1-downward
  loop).
\item
  Consequently, CI is data-dependent: it is ``='' (equal) iff
  \texttt{X{[}1{]}} equals the maximum, otherwise ``\textgreater{}''
  (greater) if the maximum is strictly larger than that last compared
  value. In general, the subroutine does not define/sanitize CI; it
  simply leaves whatever CI arose from the last \texttt{CMPA}.
\end{itemize}

(That is exactly why (4) omitted CI by default: it isn't a stable,
caller-useful postcondition unless explicitly specified.)

If the precondition \texttt{n\ ≥\ 1} is violated (\texttt{rI1} not
positive)

The code initializes with \texttt{j\ ←\ n;\ m\ ←\ X{[}n{]};\ k\ ←\ n−1;}
and then loops while \texttt{k\ \textgreater{}\ 0}. If \texttt{n\ ≤\ 0},
the loop body is skipped immediately, and the routine returns without
doing a proper ``maximum of 1..n'':

\begin{itemize}
\tightlist
\item
  It will have loaded \texttt{rA} from \texttt{X{[}n{]}} and set
  \texttt{rI2\ ←\ n}, then exit with \texttt{rI3\ =\ n−1\ ≤\ 0} (so the
  advertised \texttt{rI3\ =\ 0} postcondition in (4) fails).
\item
  If \texttt{n\ =\ 0}, you actually get the value at \texttt{X{[}0{]}};
  if \texttt{n\ \textless{}\ 0}, you fetch from \texttt{X\ +\ n} (a
  location before \texttt{X{[}1{]}}). This is precisely why the spec's
  precondition ``assume \texttt{n\ ≥\ 1}'' is essential; otherwise the
  behavior is out of spec and effectively undefined for the intended
  task. 
\end{itemize}

Completed specification addendum to (4)

Add the following to the ``Exit conditions'' of (4):

\begin{itemize}
\tightlist
\item
  \texttt{rJ} after return: contains the address \texttt{EXIT+1} (inside
  the subroutine), not the caller's return address; therefore
  \texttt{rJ} is modified.
\item
  Comparison indicator (CI): equals the result of the final
  \texttt{CMPA} executed in the loop (i.e., comparison of the final
  maximum with the last element tested). Its value is not fixed by the
  specification and must not be relied upon by callers.
\item
  If \texttt{rI1\ ≤\ 0} at entry: the routine returns immediately after
  loading \texttt{rA\ ←\ X{[}n{]}}, with \texttt{rI2\ ←\ n} and
  \texttt{rI3\ =\ n−1\ ≤\ 0}; the semantics ``max of the first
  \texttt{n} elements'' does not hold in this case. 
\end{itemize}

\section{4 {[}21{]}}\label{section-152}

Generalize the \texttt{MAX} subroutine so it can scan an array with an
arbitrary stride \texttt{r\textgreater{}0} instead of consecutive
elements. Provide two entry points:

\begin{itemize}
\tightlist
\item
  \texttt{SMAX} --- general case (caller supplies \texttt{r}),
\item
  \texttt{SMAX1} --- fast special case for \texttt{r=1}.
\end{itemize}

Assume the array base is the memory label \texttt{X}, the caller
supplies the length \texttt{n\textgreater{}0} in index register
\texttt{I3}, and (for \texttt{SMAX}) the stride
\texttt{r\textgreater{}0} in index register \texttt{I1}. On return put
the maximum into \texttt{A} and the position where it occurs into
\texttt{I2}. (Context: the fixed-length \texttt{MAX100} subroutine and
Knuth's subroutine conventions appear in §1.4.1. )

Knuth's official answer states the exit conditions precisely, including
the final value of \texttt{I3}; we'll derive those below.

The original \texttt{MAX100} scans \texttt{X+100,\ X+99,\ …,\ X+1} by
keeping the current offset in \texttt{I3} and doing \texttt{DEC3\ 1}
each pass. To scan every \texttt{r}-th element, do exactly the same but
replace the decrement by \texttt{DEC3\ 0,1}, which subtracts \texttt{r}
(the contents of \texttt{I1}) from \texttt{I3}. The loop terminates when
\texttt{I3\ ≤\ 0}. (This is the classic one-entry, one-exit MIX
subroutine scheme with \texttt{STJ\ EXIT}/\texttt{EXIT\ JMP\ *} linkage.
)

\begin{verbatim}
* MAXIMUM OVER X[n], X[n−r], X[n−2r], ...
* Calling:
*   SMAX : rI3 = n (>0), rI1 = r (>0)
*   SMAX1: rI3 = n (>0), and r = 1 by convention
* Returns:
*   rA  = max of inspected elements
*   rI2 = offset k (so that element is X + k)
*   rI3 = ((n−1) mod r) + 1 − r   (derived below)

SMAX    STJ  EXIT          * Save return address
        JMP  INIT          * General case (stride in rI1)

SMAX1   STJ  EXIT          * Fast path when r = 1
        JMP  INIT1

* ---- initialization: A ← X + I3, I2 ← I3, then step once ----

INIT    ENT2 0,3           * rI2 ← I3
        LDA  X,3           * rA  ← X + I3
        DEC3 0,1           * I3 ← I3 − r   (r = rI1)
        J3P  LOOP          * if I3 > 0, examine more
        EXIT JMP *         * single element case

INIT1   ENT2 0,3
        LDA  X,3
        DEC3 1             * I3 ← I3 − 1   (r = 1)
        J3P  LOOP1
        EXIT JMP *

* ---- main loop: compare, possibly update (A,I2), then step ----

LOOP    CMPA X,3           * compare A with X + I3
        JGE  STEP
        ENT2 0,3           * remember where the new max is
        LDA  X,3           * take the new max into A
STEP    DEC3 0,1           * I3 ← I3 − r
        J3P  LOOP
        EXIT JMP *

LOOP1   CMPA X,3
        JGE  STEP1
        ENT2 0,3
        LDA  X,3
STEP1   DEC3 1
        J3P  LOOP1
        EXIT JMP *
\end{verbatim}

This mirrors the pattern of \texttt{MAX100} (initialize by loading the
first candidate, then loop), but changes the decrement to \texttt{r}
when entered via \texttt{SMAX}. (Compare with the \texttt{MAX100}
listing shown in the text. )

Preconditions and postconditions

\begin{itemize}
\item
  Entry \texttt{I3\ =\ n\ \textgreater{}\ 0}. For \texttt{SMAX},
  \texttt{I1\ =\ r\ \textgreater{}\ 0}; for \texttt{SMAX1}, \texttt{r}
  is understood to be 1.
\item
  Exit

  \begin{itemize}
  \tightlist
  \item
    \texttt{A\ =\ max\ \{\ CONTENTS(X\ +\ n\ −\ k\ r)\ :\ k\ =\ 0,1,2,…,\ ⌊(n−1)/r⌋\ \}}.
    Equivalently, the maximum over \texttt{X+n,\ X+n−r,\ X+n−2r,\ …}
    while the offset stays positive.
  \item
    \texttt{I2} is the offset at which that maximum was found (so the
    element is \texttt{CONTENTS(X\ +\ I2)}).
  \item
    \texttt{I3\ =\ n\ −\ r\ ⌈n/r⌉\ =\ (\ (n−1)\ mod\ r\ )\ +\ 1\ −\ r\ =\ −((−n)\ mod\ r)}.
    (Knuth states this compactly in the official answer; the algebra
    appears below. )
  \end{itemize}
\item
  Registers preserved As with \texttt{MAX100}, the code doesn't touch
  \texttt{X} or \texttt{I1} (except that \texttt{SMAX1} never needs to
  set \texttt{I1}). The return address is handled by
  \texttt{STJ\ EXIT}/\texttt{EXIT\ JMP\ *} per §1.4.1 conventions. 
\end{itemize}

Correctness (loop invariant)

Let \texttt{k} denote the current value of \texttt{I3} \emph{at the top
of} \texttt{LOOP/LOOP1}.

Invariant: After each comparison step,
\texttt{A\ =\ max\ \{\ CONTENTS(X\ +\ t)\ :\ t\ ∈\ \{\ n,\ n−r,\ n−2r,\ …,\ k+r\ \}\ \}}
and \texttt{I2} is a value \texttt{t} realizing that maximum.

\begin{itemize}
\tightlist
\item
  Initialization (\texttt{INIT/INIT1}): We set
  \texttt{A\ ←\ CONTENTS(X\ +\ n)} and \texttt{I2\ ←\ n}; the set
  \texttt{\{n\}} makes the invariant true.
\item
  Preservation: If \texttt{A\ ≥\ CONTENTS(X\ +\ k)}, we keep
  \texttt{(A,I2)}; otherwise we set \texttt{A\ ←\ CONTENTS(X\ +\ k)},
  \texttt{I2\ ←\ k}. Either way, after stepping \texttt{k\ ←\ k\ −\ r}
  we've incorporated the element at offset \texttt{k} into the set.
\item
  Termination: The loop stops when \texttt{k\ ≤\ 0}. Then all offsets
  \texttt{n,\ n−r,\ …,\ \textgreater{}0} have been examined, so
  \texttt{A} is the desired maximum and \texttt{I2} its position.
\end{itemize}

Final value of \texttt{I3}

Starting with \texttt{I3\ =\ n} and repeatedly doing
\texttt{I3\ ←\ I3\ −\ r} until \texttt{I3\ ≤\ 0} gives, after
\texttt{q\ =\ ⌈n/r⌉} decrements,

\begin{verbatim}
I3_final = n − q r = n − r ⌈n/r⌉.
\end{verbatim}

Using \texttt{⌈n/r⌉\ =\ ⌊(n−1)/r⌋\ +\ 1},

\begin{verbatim}
I3_final = n − r ( ⌊(n−1)/r⌋ + 1 )
         = (n − r ⌊(n−1)/r⌋) − r
         = ((n−1) mod r + 1) − r,
\end{verbatim}

which matches Knuth's succinct form \texttt{−((−n)\ mod\ r)}. 

Cost

Let \texttt{m\ =\ ⌈n/r⌉} be the number of elements inspected. Each pass
is a constant number of MIX instructions (one compare, a conditional
2-instruction update, one decrement, one test), so the running time is
Θ(m) = Θ(⌈n/r⌉), reducing work by a factor of about \texttt{r} relative
to a unit-stride scan.

Sanity check (small example)

Take \texttt{n=7}, \texttt{r=3}. The offsets visited are
\texttt{7,\ 4,\ 1}; the next would be \texttt{−2}, so we stop. The final
\texttt{I3} equals \texttt{−2}, and the formula gives

\begin{verbatim}
((7−1) mod 3) + 1 − 3 = (6 mod 3) + 1 − 3 = 0 + 1 − 3 = −2.
\end{verbatim}

Consistent.

\subsection{5 {[}21{]}}\label{section-153}

\begin{enumerate}
\def\labelenumi{\arabic{enumi})}
\tightlist
\item
  Calling code
\end{enumerate}

Load the return address (the word immediately after the call) into a
non-J register---say \texttt{rA}---then jump to the subroutine:

\begin{verbatim}
      ENTA *+2        ; rA ← address of the next instruction (the return point)
      JMP  MAX100     ; enter the subroutine
; ... execution resumes here after the subroutine returns
\end{verbatim}

\texttt{*+2} is the location two words ahead: the word after
\texttt{JMP\ MAX100} is the resume point, and MIX's address computation
makes \texttt{ENTA\ *+2} exactly that. (We use two words because the
\texttt{JMP} itself occupies one word.) The subroutine does not rely on
A's initial value, since it initializes A before the first compare (see
below), so there are no entry conditions to worry about. 

\begin{enumerate}
\def\labelenumi{\arabic{enumi})}
\setcounter{enumi}{1}
\tightlist
\item
  Subroutine prologue/epilogue
\end{enumerate}

Replace the \texttt{STJ} prologue with a store from the register you
chose (here \texttt{rA}) into the address field of \texttt{EXIT}:

\begin{Shaded}
\begin{Highlighting}[]
\NormalTok{* MAXIMUM OF X[1..100]  (abbreviated body)}
\NormalTok{MAX100  STA  EXIT(0:2)   ; linkage via rA: store sign+2{-}byte address into EXIT}
\NormalTok{        ENT3 100         ; M1. Initialize}
\NormalTok{        JMP  2F}

\NormalTok{1H      CMPA X,3         ; M3. Compare}
\NormalTok{        JGE  *+3}
\NormalTok{2H      ENT2 0,3         ; M4. Change m}
\NormalTok{        LDA  X,3         ;     New maximum found}
\NormalTok{        DEC3 1           ; M5. Decrease k}
\NormalTok{        J3P  1B          ; M2. All tested?}
\NormalTok{EXIT    JMP  *           ; return using the address we stored}
\end{Highlighting}
\end{Shaded}

\begin{itemize}
\tightlist
\item
  \texttt{EXIT(0:2)} targets the address part (sign + two bytes) of the
  \texttt{JMP\ *} at \texttt{EXIT}, leaving its index/field/opcode bytes
  undisturbed---exactly what \texttt{STJ} would have done with
  \texttt{rJ}, but now with \texttt{rA}.
\item
  The rest of \texttt{MAX100} is unchanged; crucially, it re-initializes
  \texttt{rA} before the first \texttt{CMPA}, so seeding \texttt{rA}
  with the return address in the caller does not affect correctness. 
\end{itemize}

\begin{quote}
You could equally well choose another register: \texttt{STX\ EXIT(0:2)},
or \texttt{ST1\ EXIT(0:2)}, \ldots{} provided the subroutine uses that
register to write the address field at entry and doesn't need its
incoming value for anything else.
\end{quote}

Why this works

\begin{itemize}
\tightlist
\item
  In MIX, the instruction word at \texttt{EXIT} is \texttt{JMP\ *}. The
  address field of that instruction is where the return address lives
  during the call.
\item
  The standard scheme writes \texttt{rJ} there with \texttt{STJ\ EXIT};
  our scheme writes whatever register we like (here \texttt{rA}) using a
  store with field spec \texttt{(0:2)}, so that only the address field
  is changed.
\item
  Control returns by executing \texttt{EXIT\ \ JMP\ *}, which now jumps
  to the address we just stored.
\end{itemize}

Knuth's answer for this exercise confirms the key point---any other
register can be used, and one simple way to arrange the call is:

\begin{verbatim}
ENTA *+2
JMP  MAX100
\end{verbatim}

with no entry conditions. Our small prologue tweak
(\texttt{STA\ EXIT(0:2)}) is the corresponding change inside the
subroutine that makes \texttt{rA} the active linkage register instead of
\texttt{rJ}.

\subsection{5 {[}21{]}}\label{section-154}

Section 1.4.1 presents a simple MIX subroutine \texttt{MAX100} that
returns to the caller via the J-register using the standard
prologue/epilogue:

\begin{itemize}
\tightlist
\item
  Prologue: \texttt{STJ\ EXIT} (store the caller's return address into
  the address field of the instruction at label \texttt{EXIT})
\item
  Epilogue: \texttt{EXIT\ \ JMP\ *} (jump to the stored return address)
\end{itemize}

The exercise asks you to show that you don't have to use \texttt{rJ} for
linkage---you can use any other register to convey the return address to
\texttt{EXIT}. Knuth's official line for this exercise is simply: ``Any
other register can be used,'' with an example calling sequence
\texttt{ENTA\ *+2\ ;\ JMP\ MAX100} (no special entry conditions). 

Construction (one explicit way to do it)

\begin{enumerate}
\def\labelenumi{\arabic{enumi})}
\tightlist
\item
  Calling code
\end{enumerate}

Load the return address (the word immediately after the call) into a
non-J register---say \texttt{rA}---then jump to the subroutine:

\begin{verbatim}
      ENTA *+2        ; rA ← address of the next instruction (the return point)
      JMP  MAX100     ; enter the subroutine
; ... execution resumes here after the subroutine returns
\end{verbatim}

\texttt{*+2} is the location two words ahead: the word after
\texttt{JMP\ MAX100} is the resume point, and MIX's address computation
makes \texttt{ENTA\ *+2} exactly that. (We use two words because the
\texttt{JMP} itself occupies one word.) The subroutine does not rely on
A's initial value, since it initializes A before the first compare (see
below), so there are no entry conditions to worry about. 

\begin{enumerate}
\def\labelenumi{\arabic{enumi})}
\setcounter{enumi}{1}
\tightlist
\item
  Subroutine prologue/epilogue
\end{enumerate}

Replace the \texttt{STJ} prologue with a store from the register you
chose (here \texttt{rA}) into the address field of \texttt{EXIT}:

\begin{Shaded}
\begin{Highlighting}[]
\NormalTok{* MAXIMUM OF X[1..100]  (abbreviated body)}
\NormalTok{MAX100  STA  EXIT(0:2)   ; linkage via rA: store sign+2{-}byte address into EXIT}
\NormalTok{        ENT3 100         ; M1. Initialize}
\NormalTok{        JMP  2F}

\NormalTok{1H      CMPA X,3         ; M3. Compare}
\NormalTok{        JGE  *+3}
\NormalTok{2H      ENT2 0,3         ; M4. Change m}
\NormalTok{        LDA  X,3         ;     New maximum found}
\NormalTok{        DEC3 1           ; M5. Decrease k}
\NormalTok{        J3P  1B          ; M2. All tested?}
\NormalTok{EXIT    JMP  *           ; return using the address we stored}
\end{Highlighting}
\end{Shaded}

\begin{itemize}
\tightlist
\item
  \texttt{EXIT(0:2)} targets the address part (sign + two bytes) of the
  \texttt{JMP\ *} at \texttt{EXIT}, leaving its index/field/opcode bytes
  undisturbed---exactly what \texttt{STJ} would have done with
  \texttt{rJ}, but now with \texttt{rA}.
\item
  The rest of \texttt{MAX100} is unchanged; crucially, it re-initializes
  \texttt{rA} before the first \texttt{CMPA}, so seeding \texttt{rA}
  with the return address in the caller does not affect correctness. 
\end{itemize}

\begin{quote}
You could equally well choose another register: \texttt{STX\ EXIT(0:2)},
or \texttt{ST1\ EXIT(0:2)}, \ldots{} provided the subroutine uses that
register to write the address field at entry and doesn't need its
incoming value for anything else.
\end{quote}

Why this works

\begin{itemize}
\tightlist
\item
  In MIX, the instruction word at \texttt{EXIT} is \texttt{JMP\ *}. The
  address field of that instruction is where the return address lives
  during the call.
\item
  The standard scheme writes \texttt{rJ} there with \texttt{STJ\ EXIT};
  our scheme writes whatever register we like (here \texttt{rA}) using a
  store with field spec \texttt{(0:2)}, so that only the address field
  is changed.
\item
  Control returns by executing \texttt{EXIT\ \ JMP\ *}, which now jumps
  to the address we just stored.
\end{itemize}

Knuth's answer for this exercise confirms the key point---any other
register can be used, and one simple way to arrange the call is:

\begin{verbatim}
ENTA *+2
JMP  MAX100
\end{verbatim}

with no entry conditions. Our small prologue tweak
(\texttt{STA\ EXIT(0:2)}) is the corresponding change inside the
subroutine that makes \texttt{rA} the active linkage register instead of
\texttt{rJ}. 

\subsection{6 {[}26{]}}\label{section-155}

Below is a clean MIXAL subroutine that emulates the hardware operator
MOVE using the calling sequence

\begin{verbatim}
JMP  MOVE
NOP  A,I(F)
\end{verbatim}

It copies exactly F words from the source block starting at the
effective address M(A,I) to the destination block starting at rI1,
word-by-word, and when finished increments rI1 by F. All other registers
(A, X, I2--I6, overflow toggle, comparison indicator) are restored to
their entry values, so the only externally visible differences from a
real MOVE are (1) register J (as with any subroutine call) and (2)
time/space. This matches the definition of MIX's MOVE (C=7: ``the number
of words specified by F is moved\ldots{} and rI1 is increased by the
value of F at the end''), including the documented
``one-word-at-a-time'' behavior on overlap.

The subroutine also follows the Section 1.4.1 convention of passing
arguments in the word \emph{after} the jump (we read the NOP A,I(F) via
rJ, which points at that word), then returning to that NOP so execution
continues with the caller's next instruction.

MOVE (software) --- MIXAL

\begin{Shaded}
\begin{Highlighting}[]
\NormalTok{* MOVE — software emulation of hardware MOVE}
\NormalTok{* }\ControlFlowTok{Call}\OperatorTok{:}\NormalTok{   JMP MOVE}
\NormalTok{*         }\BuiltInTok{NOP}\NormalTok{ A}\OperatorTok{,}\NormalTok{I}\OperatorTok{(}\NormalTok{F}\OperatorTok{)}
\NormalTok{* Effect}\OperatorTok{:}\NormalTok{ Copy F words from M}\OperatorTok{(}\NormalTok{A}\OperatorTok{,}\NormalTok{I}\OperatorTok{)} \OperatorTok{{-}\textgreater{}} \OperatorTok{[}\NormalTok{rI1 }\OperatorTok{..}\NormalTok{ rI1}\OperatorTok{+}\NormalTok{F}\OperatorTok{{-}}\DecValTok{1}\OperatorTok{]}\CommentTok{; then rI1 := rI1 + F.}
\NormalTok{* Notes}\OperatorTok{:}\NormalTok{  Preserves rA}\OperatorTok{,}\NormalTok{ rX}\OperatorTok{,}\NormalTok{ rI2}\OperatorTok{..}\NormalTok{rI6}\OperatorTok{,}\NormalTok{ CI}\OperatorTok{,}\NormalTok{ overflow}\CommentTok{; only rI1 (and rJ) differ.}

\NormalTok{MOVE    STJ     RET           Save return }\OperatorTok{(}\NormalTok{address of the NOP A}\OperatorTok{,}\NormalTok{I}\OperatorTok{(}\NormalTok{F}\OperatorTok{))}

\NormalTok{        * Save registers we will clobber}
\NormalTok{        STA     SAVEA}
\NormalTok{        ST1     SAVE1}
\NormalTok{        ST2     SAVE2}
\NormalTok{        ST3     SAVE3}
\NormalTok{        ST4     SAVE4}

\NormalTok{        * rJ points to the parameter }\DataTypeTok{word} \OperatorTok{(}\NormalTok{NOP A}\OperatorTok{,}\NormalTok{I}\OperatorTok{(}\NormalTok{F}\OperatorTok{)).}
\NormalTok{        STJ     JADDR}
\NormalTok{        LD2     JADDR         rI2 }\OperatorTok{\textless{}{-}}\NormalTok{ rJ  }\OperatorTok{(}\NormalTok{an address we can index from}\OperatorTok{)}

\NormalTok{        * Decode A}\OperatorTok{,}\NormalTok{ I}\OperatorTok{,}\NormalTok{ F from that instruction }\DataTypeTok{word}\OperatorTok{.}
\NormalTok{        LDA     }\DecValTok{0}\OperatorTok{,}\DecValTok{2}\OperatorTok{(}\DecValTok{0}\OperatorTok{:}\DecValTok{2}\OperatorTok{)}\NormalTok{      rA }\OperatorTok{\textless{}{-}}\NormalTok{ A }\OperatorTok{(}\NormalTok{±AA}\OperatorTok{)}
\NormalTok{        STA     SRC           save base address A}
\NormalTok{        LD4     }\DecValTok{0}\OperatorTok{,}\DecValTok{2}\OperatorTok{(}\DecValTok{3}\OperatorTok{:}\DecValTok{3}\OperatorTok{)}\NormalTok{      rI4 }\OperatorTok{\textless{}{-}}\NormalTok{ I }\OperatorTok{(}\FloatTok{0.}\OperatorTok{.}\DecValTok{6}\OperatorTok{)}
\NormalTok{        LD3     }\DecValTok{0}\OperatorTok{,}\DecValTok{2}\OperatorTok{(}\DecValTok{5}\OperatorTok{:}\DecValTok{5}\OperatorTok{)}\NormalTok{      rI3 }\OperatorTok{\textless{}{-}}\NormalTok{ F }\OperatorTok{(}\NormalTok{count}\OperatorTok{)}

\NormalTok{        * Build source pointer}\OperatorTok{:}\NormalTok{ rI2 }\OperatorTok{:=}\NormalTok{ A }\OperatorTok{+}\NormalTok{ rIi }\OperatorTok{(}\NormalTok{using index arith}\OperatorTok{.}\NormalTok{ only}\CommentTok{;}
\NormalTok{        * INCi does }\OperatorTok{not}\NormalTok{ affect the overflow toggle}\OperatorTok{).}
\NormalTok{        LD2     SRC           rI2 }\OperatorTok{\textless{}{-}}\NormalTok{ A}
\NormalTok{        J4Z     SRC\_OK        if I }\OperatorTok{=} \DecValTok{0}\OperatorTok{,}\NormalTok{ no indexing needed}

\NormalTok{        * }\BuiltInTok{Add}\NormalTok{ the selected index register to rI2}
\NormalTok{ISELECT DEC4    }\DecValTok{1}
\NormalTok{        J4Z     ADDI1}
\NormalTok{        DEC4    }\DecValTok{1}
\NormalTok{        J4Z     ADDI2}
\NormalTok{        DEC4    }\DecValTok{1}
\NormalTok{        J4Z     ADDI3}
\NormalTok{        DEC4    }\DecValTok{1}
\NormalTok{        J4Z     ADDI4}
\NormalTok{        DEC4    }\DecValTok{1}
\NormalTok{        J4Z     ADDI5}
\NormalTok{        * must be }\DecValTok{6}
\NormalTok{        ST6     TMP}
        \ControlFlowTok{JMP}\NormalTok{     ADDI}

\NormalTok{ADDI1   }\KeywordTok{ST1}\NormalTok{     TMP}
        \ControlFlowTok{JMP}\NormalTok{     ADDI}
\NormalTok{ADDI2   }\KeywordTok{ST2}\NormalTok{     TMP}
        \ControlFlowTok{JMP}\NormalTok{     ADDI}
\NormalTok{ADDI3   }\KeywordTok{ST3}\NormalTok{     TMP}
        \ControlFlowTok{JMP}\NormalTok{     ADDI}
\NormalTok{ADDI4   }\KeywordTok{ST4}\NormalTok{     TMP}
        \ControlFlowTok{JMP}\NormalTok{     ADDI}
\NormalTok{ADDI5   }\KeywordTok{ST5}\NormalTok{     TMP}
\NormalTok{ADDI    INC2    TMP           rI2 }\OperatorTok{+=}\NormalTok{ rIi}
\NormalTok{SRC\_OK  NOP}

\NormalTok{        * Main loop}\OperatorTok{:}\NormalTok{ move F words}\OperatorTok{,}\NormalTok{ advancing both pointers}
\ControlFlowTok{LOOP}\NormalTok{    J3Z     DONE          if F }\OperatorTok{=} \DecValTok{0}\OperatorTok{,}\NormalTok{ nothing to do}
\NormalTok{        LDA     }\DecValTok{0}\OperatorTok{,}\DecValTok{2}\NormalTok{           rA }\OperatorTok{:=}\NormalTok{ CONTENTS}\OperatorTok{(}\NormalTok{rI2}\OperatorTok{)}
\NormalTok{        STA     }\DecValTok{0}\OperatorTok{,}\DecValTok{1}\NormalTok{           CONTENTS}\OperatorTok{(}\NormalTok{rI1}\OperatorTok{)} \OperatorTok{:=}\NormalTok{ rA}
\NormalTok{        INC2    }\DecValTok{1}\NormalTok{             rI2}\OperatorTok{++}
\NormalTok{        INC1    }\DecValTok{1}\NormalTok{             rI1}\OperatorTok{++}
\NormalTok{        DEC3    }\DecValTok{1}\NormalTok{             F}\OperatorTok{{-}{-}}
        \ControlFlowTok{JMP}\NormalTok{     LOOP}

\NormalTok{DONE    }\OperatorTok{*}\NormalTok{ rI1 has been advanced by original F }\OperatorTok{(}\NormalTok{as required}\OperatorTok{)}

\NormalTok{        * }\PreprocessorTok{Restore}\NormalTok{ preserved registers }\OperatorTok{(}\NormalTok{except rI1}\OperatorTok{)}
\NormalTok{        LDA     SAVEA}
\NormalTok{        LD2     SAVE2}
\NormalTok{        LD3     SAVE3}
\NormalTok{        LD4     SAVE4}
\NormalTok{        LD1     SAVE1}

\ControlFlowTok{RET}\NormalTok{     JMP     }\OperatorTok{*}\NormalTok{             Return to the NOP}\CommentTok{; it executes and control}
\NormalTok{                              * falls through to the caller’s next instr}\OperatorTok{.}

\NormalTok{* {-}{-}{-}{-}{-}{-}{-} locals }\OperatorTok{{-}{-}{-}{-}{-}{-}{-}}
\NormalTok{SAVEA   CON     }\DecValTok{0}
\NormalTok{SAVE1   CON     }\DecValTok{0}
\NormalTok{SAVE2   CON     }\DecValTok{0}
\NormalTok{SAVE3   CON     }\DecValTok{0}
\NormalTok{SAVE4   CON     }\DecValTok{0}
\NormalTok{JADDR   CON     }\DecValTok{0}
\NormalTok{SRC     CON     }\DecValTok{0}
\NormalTok{TMP     CON     }\DecValTok{0}
\end{Highlighting}
\end{Shaded}

Why this matches the hardware MOVE

\begin{enumerate}
\def\labelenumi{\arabic{enumi}.}
\item
  Source/destination and count. We take A and I from the parameter and
  compute M(A,I) for the source; F is the count. The loop copies one
  word at a time in increasing address order, exactly like the hardware
  (so overlapping ranges behave the same), and rI1 advances by F total.
\item
  Register effects. MOVE shouldn't alter rA, rX, rI2--rI6, CI, or the
  overflow toggle. We therefore save/restore those we touch and use INCi
  (not ADD/INCA) to avoid toggling overflow. Only rI1 changes (by F),
  and rJ is different because of the call/return---exactly as the
  exercise requests.
\item
  Calling/returning. Section 1.4.1's ``arguments after the jump''
  convention is used: we read the NOP A,I(F) via rJ; returning with
  \texttt{JMP\ *} lands on that NOP, which executes (doing nothing) and
  then the caller proceeds to the next instruction.
\end{enumerate}

If you'd like, I can also show a branch-table version that packs the six
\texttt{ADDIk} cases tighter, or a self-modifying variant that patches
the I-field of a template \texttt{LDA} to avoid the small
\texttt{ISELECT} chain.

\subsection{7 {[}20{]}}\label{section-156}

Here's a concise, exam-style answer with enough depth for TAOCP.

Knuth's MIX used self-modifying code for subroutine linkage (e.g.,
storing the return address into the instruction stream), but he already
notes that ``newer computers\ldots keep {[}such{]} information in
registers or in a special array, not within the program itself'' (see
the note in §1.4.1 that points you to this exercise). The reasons modern
systems discourage self-modification include:

\begin{enumerate}
\def\labelenumi{\arabic{enumi}.}
\item
  Security policies (W⊕X). Operating systems increasingly enforce
  \emph{write-xor-execute}: a page may be writable or executable, not
  both. That prevents code injection and many exploit techniques; it
  also means ordinary programs can't (and shouldn't) patch their own
  instruction bytes.
\item
  OS/ABI design: text is shared and read-only. Executable code
  (``text'') pages are mapped read-only and commonly shared among
  processes for memory efficiency. Letting programs write into those
  pages would break sharing and isolation.
\item
  Performance costs on modern CPUs. Pipelines and split instruction/data
  caches mean that when code \emph{is} changed, the CPU must perform
  special synchronization/serialization and cache/TLB maintenance to
  make the modification visible to instruction fetch. That's slow and
  error-prone if mishandled. Vendors explicitly warn about penalties and
  coherence requirements for self-modifying code.
\item
  Concurrency and reentrancy hazards. Multiple threads may execute the
  same code simultaneously. Patching instructions in place introduces
  race conditions, tears, and undefined behavior unless you add
  heavyweight synchronization---further eroding any perceived benefit.
  (Read-only text plus normal call/return linkage avoids this.)
\item
  Optimization, tooling, and maintainability. Compilers, linkers,
  profilers, verifiers, and JITs assume code is immutable at run time;
  immutability enables aggressive optimization, code signing,
  control-flow integrity, and reliable debugging. Self-modification
  undermines those assumptions and complicates reasoning about
  correctness.
\item
  Better idioms exist. Modern linkage stores return addresses in a
  register/stack frame (or a link register), not in the instruction
  stream; parameter passing and multiple exits are handled by calling
  conventions and structured control flow, so the MIX-style ``patch the
  code to return here'' is unnecessary. (Knuth foreshadows exactly this
  shift.)
\end{enumerate}

Bottom line. What was once a practical trick on simple machines is now a
liability: it conflicts with OS security and memory models, harms
performance on cached/pipelined CPUs, breaks reentrancy, and fights
today's toolchains. Hence modern advice---and practice---is to avoid
self-modifying code. The exercise list where this appears is on p.~193.

\bookmarksetup{startatroot}

\chapter{1.4.2 Coroutines}\label{coroutines}

\section{1 {[}10{]}}\label{section-157}

A \emph{coroutine} is symmetric: two (or more) routines repeatedly
transfer control to each other, each time resuming exactly where it last
left off (state must persist across transfers). This is unlike a
subroutine, where control flows down and then returns once. 

Why ``short and simple'' is elusive

\begin{enumerate}
\def\labelenumi{\arabic{enumi}.}
\item
  One transfer ⇒ not really a coroutine. If A calls B only once (and B
  never resumes A mid-run), the interaction is observationally identical
  to an ordinary subroutine call. So the ``smallest'' control exchange
  doesn't demonstrate anything essentially coroutine-specific. To
  \emph{show} coroutine behavior, you need multiple alternations
  (A↔B↔A↔B\ldots), which already lengthens the example.
\item
  State must be preserved between activations. Real coroutine examples
  must keep local state alive across yields. On a bare machine (like MIX
  in the text) that means explicit linkage and register
  saves/restores---extra instructions that swell even toy programs.
  Knuth's own coroutine linkage shows saving/restoring registers and
  jump handshakes, which are inherently not ``tiny.'' 
\item
  Nontrivial data flow is needed to be meaningful. To motivate
  \emph{why} coroutines help, you typically show a producer/consumer,
  scanner/parser, or ``turn an n-pass pipeline into a one-pass process''
  example. But as soon as you pass realistic data (buffers, termination
  conditions, etc.), the code grows beyond a few lines. (Knuth's figure
  contrasting a four-pass pipeline with a one-pass chain of coroutines
  illustrates the typical use case.) 
\item
  Toy examples are either contrived or misleading. Extremely short
  ``ping-pong'' snippets (e.g., print A/B alternately) either hide the
  crucial save/restore machinery or rely on features (like a built-in
  \texttt{yield}) that the reader hasn't met yet, making them
  pedagogically suspect on a machine-language canvas.
\end{enumerate}

Minimal illustration (conceptual)

\begin{itemize}
\tightlist
\item
  Subroutine pattern: \texttt{main\ →\ sub()\ →\ return} (one descent,
  one return).
\item
  Coroutine pattern: \texttt{A\ →\ B\ →\ A\ →\ B\ →\ …} (multiple
  \emph{resumptions}), plus explicit context preservation. Without at
  least two transfers, A↔B collapses to subroutine semantics and doesn't
  illuminate coroutines.
\end{itemize}

To genuinely demonstrate coroutines you need (i) more than one control
transfer, (ii) preserved execution context, and (iii) meaningful data
exchange. Each of these adds lines and machinery, so examples that are
both \emph{short} and \emph{truly illustrative} are rare; once they're
short enough, they stop being distinctly coroutine examples and look
like plain subroutines instead.

\section{2 {[}20{]}}\label{section-158}

What actually happens if you start with \texttt{IN}

\begin{enumerate}
\def\labelenumi{\arabic{enumi}.}
\item
  Initialization still sets \texttt{rI6\ ←\ 0} and \texttt{rX\ ←\ 0}
  (lines 58--59), so \texttt{IN1} will read the first input word if
  needed and isolate the next character into \texttt{rA}. 
\item
  At the end of its step, \texttt{IN} yields to \texttt{OUT} via the
  coroutine linkage (\texttt{JMP\ OUT}). Control lands at the
  \texttt{OUT} linkage stub, which first saves IN's resume point and
  then transfers to \texttt{OUT1}. 
\item
  \texttt{OUT1} begins by preparing a fresh output card/buffer: it sets
  up indices and executes a \texttt{MOVE} that blanks the output area
  (``Start new output card\ldots{} Set output area to blanks''), then it
  asks \texttt{IN} for the ``next translated character'' by jumping back
  via the linkage (\texttt{JMP\ IN}). Those prologue instructions
  overwrite or otherwise invalidate the character that \texttt{IN} had
  already placed in \texttt{rA}. Net effect: the very first character
  fetched by \texttt{IN} is lost before \texttt{OUT} stores anything. 
\end{enumerate}

This is exactly why the text starts with \texttt{JMP\ OUT1}: the OUT
coroutine's prologue should run before the first request to \texttt{IN},
so that the buffer is ready and the first character fetched by
\texttt{IN} is immediately consumed instead of being clobbered by OUT's
setup. The official answer summarizes this as ``the first character
found by IN would be lost,'' and notes that lines 58--59 were chosen to
initialize \texttt{IN}; starting with \texttt{IN} would require
compensating initialization for \texttt{OUT}. 

How to make ``start with IN'' work (one correct fix)

If you truly want to begin at \texttt{IN1}, adjust the first handoff
into \texttt{OUT} so OUT's prologue does not run after \texttt{IN} has
already put a character in \texttt{rA}:

\begin{itemize}
\tightlist
\item
  Pre-initialize OUT's state that the prologue would otherwise need.
  Concretely, set \texttt{rA\ ←\ −16} so the initial \texttt{MOVE} count
  is already in place, and clear the output buffer if it isn't known to
  be blank.
\item
  Change the first transfer from \texttt{IN} to \texttt{OUT} to jump
  directly to the code that stores \texttt{rA} into the buffer (the code
  around line 39) rather than to \texttt{OUT1} (the prologue). In the
  book's notation: modify the first \texttt{OUTX} jump target so it goes
  to the ``store character'' logic rather than the ``start new card''
  logic. With that, the very first character produced by \texttt{IN}
  will be preserved and packed before any further OUT setup runs. 
\end{itemize}

Bottom line

\begin{itemize}
\tightlist
\item
  If you simply replace \texttt{JMP\ OUT1} with \texttt{JMP\ IN1}, the
  first input character disappears.
\item
  To start with \texttt{IN} safely, you must (i) pre-initialize OUT's
  prologue effects and (ii) route the first transfer from \texttt{IN}
  straight into OUT's ``store character'' code rather than OUT's
  prologue. 
\end{itemize}

\section{3 {[}20{]}}\label{section-159}

False (but ``almost'' true). Removing all three comparisons breaks
correctness in two distinct ways:

\begin{enumerate}
\def\labelenumi{\arabic{enumi}.}
\tightlist
\item
  The comparison indicator (CI) can be uninitialized at the first
  \texttt{JE\ 9F}
\end{enumerate}

In the OUT coroutine the test for end-of-input is done by

\begin{verbatim}
… STA …(1:1)
CMPA PERIOD
JE  9F
…
… STA …(2:2)
CMPA PERIOD
JE  9F
…
… STA …(3:3)
CMPA PERIOD
JE  9F
…
9F  JNE OUT1
HLT
\end{verbatim}

(see lines 39--55). Each \texttt{JE\ 9F} \emph{assumes} the CI was set
by the immediately preceding \texttt{CMPA\ PERIOD}. If you delete those
compares, the very first \texttt{JE} in OUT will consult whatever CI
happens to be left over from earlier code (i.e., undefined for some
paths), so behavior is not well-defined. In particular, if the first
input character is a letter, coroutine IN takes the ``nondigit'' branch
\emph{without doing any comparison at all} and jumps to OUT; then OUT
reaches \texttt{JE\ 9F} with an undefined CI. Hence correctness already
fails. 

\begin{enumerate}
\def\labelenumi{\arabic{enumi}.}
\setcounter{enumi}{1}
\tightlist
\item
  A final period after a replication digit isn't detected
\end{enumerate}

Even when the CI happens to be set, dropping those compares misses the
special case ``digit followed by final `.'\,''. Coroutine IN reads a
replication digit \texttt{n}, then fetches the next character (which
might be the terminating `.') and immediately jumps to OUT
\emph{without} doing any compare in between (see lines 24--30). If `.'
is preceded by a replication digit, the most recent comparison in IN was
the digit test (not a period test), so the CI is \emph{not} ``equal'';
without \texttt{CMPA\ PERIOD} in OUT, the program won't notice the final
period and won't halt. 

Knuth's official notes summarize both issues succinctly: removing all
three compares is ``almost'' okay, but ``the comparison indicator isn't
initialized; and if the final period is preceded by a replication digit,
it won't be noticed.'' 

If you really want to remove the three compares

You can relocate the period test so the CI is \emph{reliably} set before
OUT checks it:

\begin{itemize}
\tightlist
\item
  After IN fetches the next character (between lines 26--27), insert a
  single \texttt{CMPA\ PERIOD} so the CI reflects whether the
  just-fetched character is `.' before control jumps to OUT.
\item
  Also ensure the CI has a known initial state right before the very
  first handoff (between lines 59--60), e.g., by comparing \texttt{rX}
  with PERIOD (one way is to insert a \texttt{CMPX\ PERIOD} there), so
  OUT's first \texttt{JE/JNE} pair isn't consulting a garbage CI. (Then
  OUT can safely rely on the CI without doing its own per-position
  \texttt{CMPA\ PERIOD}.) 
\end{itemize}

This preserves the coroutine contract while shaving a few instructions
from OUT, but the stock program (with the three local
\texttt{CMPA\ PERIOD}s) is the clearest and safest. 

Don't delete all three \texttt{CMPA\ PERIOD} in OUT; doing so makes the
program incorrect on real inputs.

\section{4 {[}20{]}}\label{section-160}

What ``linkage (1)'' is doing in MIX

Knuth's linkage (labels rearranged slightly for clarity) is the tiny
handshake at the end of his program:

\begin{verbatim}
60  JMP OUT1        ; start with OUT  (exercise 2)
61  OUT  STJ INX    ; store our resume point "OUTX" into IN's jump word
62  OUTX JMP OUT1   ; jump to OUT's entry (will be patched by IN)
63  IN   STJ OUTX   ; ditto, in the other direction
64  INX  JMP IN1    ; jump to IN's entry (will be patched by OUT)
\end{verbatim}

Each coroutine owns one \emph{indirect jump word} (\texttt{INX},
\texttt{OUTX}). When a coroutine yields, it stores the address of its
own resume label (e.g., \texttt{OUTX}) into the \emph{other} coroutine's
jump word, then jumps to the other's entry (\texttt{IN1} ⇄
\texttt{OUT1}). The first time, \texttt{INX} and \texttt{OUTX} contain
the entry addresses; thereafter they contain the saved resume addresses.
(Knuth later shows a variant ``(4)'' that also saves/restores register
A, but the basic idea above is ``(1)''.) 

Below are two concrete, modern implementations that are directly
analogous but avoid self-modifying code by using \emph{indirect branches
through writable data}.

x86-64 (NASM/Intel syntax; position-independent)

\begin{Shaded}
\begin{Highlighting}[]
\CommentTok{; Assemble with: nasm {-}felf64 coro.asm \&\& ld {-}o coro coro.o}
\KeywordTok{default} \KeywordTok{rel}

\KeywordTok{section} \FunctionTok{.data}
\FunctionTok{INX:}   \DataTypeTok{dq}\NormalTok{ IN1        }\CommentTok{; IN’s current continuation (entry at first)}
\FunctionTok{OUTX:}  \DataTypeTok{dq}\NormalTok{ OUT1       }\CommentTok{; OUT’s current continuation (entry at first)}

\KeywordTok{section} \FunctionTok{.text}
\KeywordTok{global}\NormalTok{ \_start}

\FunctionTok{\_start:}
    \ControlFlowTok{jmp}\NormalTok{ OUT1         }\CommentTok{; Like line 60: begin in OUT}

\CommentTok{; {-}{-}{-} IN coroutine {-}{-}{-}{-}{-}{-}{-}{-}{-}{-}{-}{-}{-}{-}{-}{-}{-}{-}{-}{-}{-}{-}{-}{-}{-}{-}{-}{-}{-}{-}{-}{-}{-}{-}{-}{-}{-}{-}{-}{-}{-}{-}{-}{-}{-}{-}{-}{-}{-}}

\CommentTok{; Entry for IN: always tail{-}jump to whatever OUT stored into INX}
\FunctionTok{IN1:}
    \ControlFlowTok{jmp} \DataTypeTok{qword} \OperatorTok{[}\NormalTok{INX}\OperatorTok{]}

\CommentTok{; Somewhere inside IN, to yield to OUT do:}
\CommentTok{;   1) store our resume address into OUTX}
\CommentTok{;   2) jump to OUT’s entry}
\FunctionTok{IN\_yield\_to\_OUT:}
    \KeywordTok{lea}  \KeywordTok{rax}\OperatorTok{,} \OperatorTok{[}\KeywordTok{rel}\NormalTok{ IN\_resume}\OperatorTok{]}  \CommentTok{; “return address” for IN}
    \KeywordTok{mov}  \OperatorTok{[}\NormalTok{OUTX}\OperatorTok{],} \KeywordTok{rax}
    \ControlFlowTok{jmp}\NormalTok{  OUT1}
\FunctionTok{IN\_resume:}
    \CommentTok{; ... execution continues here when OUT yields back ...}

\CommentTok{; {-}{-}{-} OUT coroutine {-}{-}{-}{-}{-}{-}{-}{-}{-}{-}{-}{-}{-}{-}{-}{-}{-}{-}{-}{-}{-}{-}{-}{-}{-}{-}{-}{-}{-}{-}{-}{-}{-}{-}{-}{-}{-}{-}{-}{-}{-}{-}{-}{-}{-}{-}{-}{-}}

\CommentTok{; Entry for OUT: always tail{-}jump to whatever IN stored into OUTX}
\FunctionTok{OUT1:}
    \ControlFlowTok{jmp} \DataTypeTok{qword} \OperatorTok{[}\NormalTok{OUTX}\OperatorTok{]}

\CommentTok{; To yield back to IN:}
\FunctionTok{OUT\_yield\_to\_IN:}
    \KeywordTok{lea}  \KeywordTok{rax}\OperatorTok{,} \OperatorTok{[}\KeywordTok{rel}\NormalTok{ OUT\_resume}\OperatorTok{]} \CommentTok{; “return address” for OUT}
    \KeywordTok{mov}  \OperatorTok{[}\NormalTok{INX}\OperatorTok{],} \KeywordTok{rax}
    \ControlFlowTok{jmp}\NormalTok{  IN1}
\FunctionTok{OUT\_resume:}
    \CommentTok{; ... execution continues here when IN yields back ...}
\end{Highlighting}
\end{Shaded}

Why this is exactly like (1):

\begin{itemize}
\tightlist
\item
  \texttt{IN1:\ jmp\ {[}INX{]}} and \texttt{OUT1:\ jmp\ {[}OUTX{]}} play
  the role of the ``indirect jump words'' (\texttt{INX\ JMP\ IN1},
  \texttt{OUTX\ JMP\ OUT1}).
\item
  \texttt{IN\_yield\_to\_OUT} stores \texttt{IN\_resume} into
  \texttt{OUTX} before jumping to \texttt{OUT1};
  \texttt{OUT\_yield\_to\_IN} stores \texttt{OUT\_resume} into
  \texttt{INX} before jumping to \texttt{IN1}.
\item
  No stacks or OS help are required; you're literally passing
  continuations in memory, just as MIX does by patching a jump word.
\end{itemize}

If you need register/save semantics like Knuth's variant (4), push/pop
(or move to callee-saved registers/stack slots) around the
\texttt{mov}/\texttt{jmp} pair, corresponding to \texttt{STA\ HOLDA}
\ldots{} \texttt{LDA\ HOLDA} in his text. 

Portable C (GCC/Clang) with ``labels as values''

This uses the compiler extension for \emph{computed goto}. It mirrors
MIX even more transparently:

\begin{Shaded}
\begin{Highlighting}[]
\CommentTok{// gcc {-}std=gnu11 coro.c {-}o coro}
\PreprocessorTok{\#include }\ImportTok{\textless{}stdio.h\textgreater{}}

\DataTypeTok{static} \DataTypeTok{void} \OperatorTok{*}\NormalTok{INX}\OperatorTok{,} \OperatorTok{*}\NormalTok{OUTX}\OperatorTok{;}

\DataTypeTok{int}\NormalTok{ main}\OperatorTok{(}\DataTypeTok{void}\OperatorTok{)} \OperatorTok{\{}
\NormalTok{  INX  }\OperatorTok{=} \OperatorTok{\&\&}\NormalTok{IN1}\OperatorTok{;}      \CommentTok{// initial entries}
\NormalTok{  OUTX }\OperatorTok{=} \OperatorTok{\&\&}\NormalTok{OUT1}\OperatorTok{;}

  \ControlFlowTok{goto}\NormalTok{ OUT1}\OperatorTok{;}         \CommentTok{// start with OUT (like line 60)}

\NormalTok{IN1}\OperatorTok{:}
  \ControlFlowTok{goto} \OperatorTok{*}\NormalTok{INX}\OperatorTok{;}         \CommentTok{// always jump to IN’s current continuation}

  \CommentTok{// ... IN’s code ...}
\NormalTok{IN\_yield\_to\_OUT}\OperatorTok{:}
\NormalTok{  OUTX }\OperatorTok{=} \OperatorTok{\&\&}\NormalTok{IN\_resume}\OperatorTok{;}  \CommentTok{// store our resume label into OUT’s jump slot}
  \ControlFlowTok{goto}\NormalTok{ OUT1}\OperatorTok{;}           \CommentTok{// jump to OUT’s entry}
\NormalTok{IN\_resume}\OperatorTok{:}
  \CommentTok{// ... resumes here ...}

\NormalTok{OUT1}\OperatorTok{:}
  \ControlFlowTok{goto} \OperatorTok{*}\NormalTok{OUTX}\OperatorTok{;}        \CommentTok{// always jump to OUT’s current continuation}

  \CommentTok{// ... OUT’s code ...}
\NormalTok{OUT\_yield\_to\_IN}\OperatorTok{:}
\NormalTok{  INX }\OperatorTok{=} \OperatorTok{\&\&}\NormalTok{OUT\_resume}\OperatorTok{;}  \CommentTok{// store our resume label into IN’s jump slot}
  \ControlFlowTok{goto}\NormalTok{ IN1}\OperatorTok{;}            \CommentTok{// jump to IN’s entry}
\NormalTok{OUT\_resume}\OperatorTok{:}
  \CommentTok{// ... resumes here ...}
\OperatorTok{\}}
\end{Highlighting}
\end{Shaded}

\begin{itemize}
\tightlist
\item
  \texttt{INX}/\texttt{OUTX} are the indirect jump words.
\item
  \texttt{goto\ *INX} / \texttt{goto\ *OUTX} replace the
  ``\texttt{JMP\ IN1}/\texttt{JMP\ OUT1} words patched by STJ''.
\item
  Assigning \texttt{\&\&label} is the exact analogue of
  \texttt{STJ\ otherX}.
\end{itemize}

This version is often the simplest way to prototype Knuth's linkage
today, and it avoids any system-dependent assembly.

\section{5 {[}15{]}}\label{section-161}

Whenever OUT executes \texttt{JMP\ IN}, rA in OUT must be exactly what
it was when control returns to the next OUT instruction; likewise for IN
when it executes \texttt{JMP\ OUT}. So each coroutine must (1) save its
own rA at exit, and (2) restore its own rA at re-entry, without
constraining the other coroutine's use of rA while it runs.

Knuth already shows the one-sided variant (``if OUT would destroy A,
save/restore A around the transfer''), right above Fig. 22; we just
symmetrize that pattern. We splice this linkage where the book places
the standard two-coroutine linkage (lines 57--65), replacing those four
short lines with the version below.

MIX code (two-sided A preservation)

\begin{Shaded}
\begin{Highlighting}[]
\NormalTok{* {-}{-}{-} Linkage data (two save slots, one per coroutine) {-}{-}{-}}
\NormalTok{HOLDIN  CON  0        * last{-}saved rA for IN}
\NormalTok{HOLDOUT CON  0        * last{-}saved rA for OUT}

\NormalTok{* {-}{-}{-} Initialization and linkage {-}{-}{-}}
\NormalTok{START   ENT6 0        * as in text: init rI6 for NEXTCHAR}
\NormalTok{        ENTX 0        * as in text: init rX for NEXTCHAR}
\NormalTok{        JMP  OUT1     * start with OUT (exercise 2 discusses starting with IN)}

\NormalTok{OUT     STJ  INX      * standard coroutine bookkeeping (return for IN)}
\NormalTok{        STA  HOLDIN   * save caller\textquotesingle{}s A (leaving IN {-}\textgreater{} coming here)}
\NormalTok{        LDA  HOLDOUT  * restore OUT\textquotesingle{}s own A for re{-}entry}
\NormalTok{OUTX    JMP  OUT1     * enter OUT proper}

\NormalTok{IN      STJ  OUTX     * standard coroutine bookkeeping (return for OUT)}
\NormalTok{        STA  HOLDOUT  * save caller\textquotesingle{}s A (leaving OUT {-}\textgreater{} coming here)}
\NormalTok{        LDA  HOLDIN   * restore IN\textquotesingle{}s own A for re{-}entry}
\NormalTok{INX     JMP  IN1      * enter IN proper}

\NormalTok{        END  START}
\end{Highlighting}
\end{Shaded}

Why this works

\begin{itemize}
\item
  When OUT → IN (\texttt{JMP\ IN} inside OUT): control lands at label
  \texttt{IN}. We immediately

  \begin{enumerate}
  \def\labelenumi{\arabic{enumi}.}
  \tightlist
  \item
    save the caller's A (\texttt{STA\ HOLDOUT}), then
  \item
    restore IN's previously saved A (\texttt{LDA\ HOLDIN}), and jump to
    \texttt{IN1}. While IN runs, it may freely clobber rA. When IN later
    yields back to OUT (\texttt{JMP\ OUT}), label \texttt{OUT} restores
    OUT's A from \texttt{HOLDOUT} before resuming at \texttt{OUT1}. Thus
    rA in OUT is unchanged across the yield/return pair.
  \end{enumerate}
\item
  Symmetrically, when IN → OUT (\texttt{JMP\ OUT} inside IN), label
  \texttt{OUT} saves IN's rA to \texttt{HOLDIN}, restores OUT's rA from
  \texttt{HOLDOUT}, and resumes OUT. When OUT yields back, label
  \texttt{IN} reloads \texttt{HOLDIN}. So rA in IN is likewise unchanged
  across its yield/return.
\end{itemize}

This is exactly the two-sided generalization of the one-sided example on
p.~198 (``if OUT would destroy rA, save at IN and restore at OUT''), but
now with two save cells so we can save the caller's A and simultaneously
restore the callee's A at each transfer.

\section{6 {[}22{]}}\label{section-162}

Here's a clean coroutine linkage for three MIX coroutines \(A,B,C\) so
that any one can transfer to either of the other two, and on activation
each resumes exactly where it last left off. (This is the Exercise 6
setup at the end of §1.4.2.) 

First, recall the 2-coroutine pattern (called ``(1)'' in the text): each
coroutine has a \emph{resume cell} holding a \texttt{JMP} whose address
field gets patched by the other coroutine via \texttt{STJ}. Then you
enter a coroutine by \emph{executing} that stored jump (i.e.,
\texttt{JMP\ OUTX} executes the word at \texttt{OUTX}).

We generalize that idea by giving each coroutine its own resume cell
(\texttt{AX}, \texttt{BX}, \texttt{CX}) and by providing two entry
labels per coroutine---one entry for each possible caller---so we know
which caller's resume cell to update.

Linkage (MIX)

\begin{Shaded}
\begin{Highlighting}[]
\NormalTok{* Resume cells }\OperatorTok{(}\NormalTok{initial jumps to first}\OperatorTok{{-}}\NormalTok{time entries A1}\OperatorTok{,}\NormalTok{ B1}\OperatorTok{,}\NormalTok{ C1}\OperatorTok{)}
\NormalTok{AX      JMP A1}
\NormalTok{BX      JMP B1}
\NormalTok{CX      JMP C1}

\NormalTok{* Entries to B}
\NormalTok{AB      STJ }\KeywordTok{AX}       \OperatorTok{*}\NormalTok{ called B from A}\OperatorTok{:}\NormalTok{ remember A’s return point in }\KeywordTok{AX}
        \ControlFlowTok{JMP} \KeywordTok{BX}       \OperatorTok{*}\NormalTok{ now execute B’s stored jump }\OperatorTok{(}\NormalTok{first time → B1}\OperatorTok{)}

\NormalTok{CB      STJ }\KeywordTok{CX}       \OperatorTok{*}\NormalTok{ called B from C}\OperatorTok{:}\NormalTok{ remember C’s return point in }\KeywordTok{CX}
        \ControlFlowTok{JMP} \KeywordTok{BX}

\NormalTok{* Entries to C}
\NormalTok{AC      STJ }\KeywordTok{AX}       \OperatorTok{*}\NormalTok{ called C from A}\OperatorTok{:}\NormalTok{ remember A’s return point}
        \ControlFlowTok{JMP} \KeywordTok{CX}

\NormalTok{BC      STJ }\KeywordTok{BX}       \OperatorTok{*}\NormalTok{ called C from B}\OperatorTok{:}\NormalTok{ remember B’s return point}
        \ControlFlowTok{JMP} \KeywordTok{CX}

\NormalTok{* Entries to A}
\NormalTok{BA      STJ }\KeywordTok{BX}       \OperatorTok{*}\NormalTok{ called A from B}\OperatorTok{:}\NormalTok{ remember B’s return point}
        \ControlFlowTok{JMP} \KeywordTok{AX}

\NormalTok{CA      STJ }\KeywordTok{CX}       \OperatorTok{*}\NormalTok{ called A from C}\OperatorTok{:}\NormalTok{ remember C’s return point}
        \ControlFlowTok{JMP} \KeywordTok{AX}
\end{Highlighting}
\end{Shaded}

How you use it

\begin{itemize}
\tightlist
\item
  From inside A, jump to \texttt{AB} to activate B, or to \texttt{AC} to
  activate C.
\item
  From inside B, jump to \texttt{BA} (activate A) or \texttt{BC}
  (activate C).
\item
  From inside C, jump to \texttt{CA} (A) or \texttt{CB} (B).
\end{itemize}

Why it works (invariants)

\begin{itemize}
\tightlist
\item
  Each coroutine Q has a resume cell QX that always contains a
  \texttt{JMP\ \textless{}where\ Q\ should\ resume\textgreater{}}.
\item
  When caller P transfers to callee Q, control arrives at entry label
  QP, which executes \texttt{STJ\ PX}, storing the caller's \emph{return
  address} (the instruction after the \texttt{JMP\ QP}) into PX. Thus
  the caller's resume cell is freshly patched to ``where to continue
  when we next execute \texttt{JMP\ PX}.''
\item
  The entry then \texttt{JMP}s to QX, \emph{executing} Q's stored
  \texttt{JMP}---first time it jumps to \texttt{Q1}; thereafter it jumps
  to Q's last suspended point.
\end{itemize}

You can verify a round-trip: A → \texttt{JMP\ AB} →
\texttt{AB:\ STJ\ AX;\ JMP\ BX} (B begins); later B → \texttt{JMP\ BA} →
\texttt{BA:\ STJ\ BX;\ JMP\ AX} which executes the patched word at AX
and resumes A precisely after its \texttt{JMP\ AB}.

\section{7 {[}30{]}}\label{section-163}

What the original program did (we must invert it)

\begin{itemize}
\tightlist
\item
  Input format (2): a stream of characters ending with a period ``.'';
  if a digit \texttt{n∈\{0,…,9\}} appears, it means ``repeat the next
  character \texttt{n+1} times,'' even if that next character is itself
  a digit. Nondigits simply denote themselves. 
\item
  Output format (3): the fully expanded character stream, printed in
  groups of three separated by blanks, 16 groups per card; processing
  stops at the period. 
\end{itemize}

Our task is the inverse: read the grouped stream (3), ignore blanks,
stop at the first period, and re-encode it into the compact code (2)
with the fewest characters (e.g., never emit the redundant \texttt{0Z},
just \texttt{Z}). 

Minimal re-encoding strategy (correctness \& optimality)

Let's call the flat stream (with blanks removed) \texttt{S}, ending at
the first \texttt{.}. We want the shortest string over the alphabet
\texttt{Σ\ ∪\ \{0,…,9\}} whose expansion (by ``digit→repeat next char
\texttt{n+1} times'') equals \texttt{S}.

Process \texttt{S} run by run: for each maximal block of identical
characters \texttt{c} of length \texttt{L}, output the shortest code for
that block, then continue.

\begin{itemize}
\item
  If \texttt{c} is a nondigit (letter, punctuation, etc.):

  \begin{itemize}
  \item
    We may either write the character itself once (cost 1) or use a
    ``digit + c'' pair (cost 2) to say multiple copies.
  \item
    Optimal encoding is greedy by tens:

    \begin{itemize}
    \item
      Let \texttt{q\ =\ ⌊L/10⌋} and \texttt{r\ =\ L\ mod\ 10}.
    \item
      Output \texttt{q} pairs \texttt{9c} (each pair contributes 10
      copies).
    \item
      For the remainder:

      \begin{itemize}
      \tightlist
      \item
        If \texttt{r\ =\ 0}: done.
      \item
        If \texttt{r\ =\ 1}: just output \texttt{c} (better than
        \texttt{0c}).
      \item
        If \texttt{2\ ≤\ r\ ≤\ 10}: output \texttt{(r−1)c}.
      \end{itemize}
    \end{itemize}
  \end{itemize}
\item
  If \texttt{c} is a digit (one of \texttt{0..9} appearing as data), we
  cannot write it bare (because bare digits in (2) are operators). We
  must always use a pair:

  \begin{itemize}
  \tightlist
  \item
    Again \texttt{q\ =\ ⌊L/10⌋}, \texttt{r\ =\ L\ mod\ 10}.
  \item
    Output \texttt{q} times \texttt{9c}; if \texttt{r\textgreater{}0},
    output \texttt{(r−1)c}.
  \end{itemize}
\end{itemize}

Why this greedy scheme is optimal:

\begin{itemize}
\tightlist
\item
  Each ``digit+char'' token has fixed cost 2 and yields between 1 and 10
  copies (specifically \texttt{n+1} copies). For nondigit characters we
  also have a special 1-copy option of cost 1 (just the character). This
  is a canonical coin system: ``coins'' of values \texttt{1..10} with
  costs \texttt{2} each, except a free single-copy option of cost 1 for
  nondigits. Greedy by 10s, then one optimal remainder, is optimal under
  such costs; any non-greedy split that uses more tokens (or replaces a
  single bare \texttt{c} by \texttt{0c}) cannot be shorter.
\end{itemize}

This guarantees ``the zero before the Z \ldots{} would not really be
produced'' (we would output just \texttt{Z}) since \texttt{0Z} (cost 2)
is strictly worse than \texttt{Z} (cost 1). 

MIX solution architecture

We mirror Knuth's i/o organization:

\begin{enumerate}
\def\labelenumi{\arabic{enumi}.}
\item
  NEXTCHAR: reuse the input subroutine from the text that returns the
  next nonblank character from card reader into \texttt{rA}; it handles
  card boundaries and blanks. (Exactly as in the example: read card,
  scan words, extract next character.) 
\item
  PUTCH: an output subroutine that appends one character to an 80-col
  output buffer; when full, punch the card and clear the buffer (similar
  to OUT's buffering, but without inserting blanks every three chars).
  The last card is punched when we see \texttt{.}.
\item
  MAIN (reverse encoder):

  \begin{itemize}
  \item
    Read first nonblank char \texttt{x} via \texttt{NEXTCHAR}.
  \item
    If \texttt{x=\textquotesingle{}.\textquotesingle{}}: output
    \texttt{.} and finish.
  \item
    Set \texttt{CUR=x}, \texttt{CNT=1}.
  \item
    Loop: get next \texttt{x} via \texttt{NEXTCHAR}.

    \begin{itemize}
    \tightlist
    \item
      If \texttt{x=\textquotesingle{}.\textquotesingle{}}: flush run
      \texttt{(CUR,\ CNT)} using EmitRun (below), then output \texttt{.}
      and finish.
    \item
      Else if \texttt{x==CUR}: \texttt{CNT++}.
    \item
      Else: flush \texttt{(CUR,\ CNT)}, then \texttt{CUR=x},
      \texttt{CNT=1}.
    \end{itemize}
  \item
    EmitRun(CUR, CNT) implements the minimal strategy above, using
    helper predicates ``is digit?'' (compare \texttt{rA} with codes
    30..39) and a small loop to emit \texttt{9CUR} \texttt{q} times;
    remainder as per rules.
  \end{itemize}
\end{enumerate}

Below is a compact MIXAL skeleton that focuses on the core logic (you
can drop it straight into your MIXAL environment; i/o linkage follows
the style of §1.4.2):

\begin{Shaded}
\begin{Highlighting}[]
\NormalTok{* REVERSE OF (3) {-}\textgreater{} (2): minimal code}
\NormalTok{* Units and buffers}
\NormalTok{READER  EQU     16}
\NormalTok{PUNCH   EQU     17}
\NormalTok{INPUT   ORIG    *+16        * input card buffer}
\NormalTok{OUTPUT  ORIG    *+16        * output card buffer}

\NormalTok{* =============== SUBROUTINE: NEXTCHAR (nonblank reader) ===============}
\NormalTok{* rA \textless{}{-} next nonblank char from input; reads new card as needed}
\NormalTok{* (Same as text\textquotesingle{}s NEXTCHAR)}
\NormalTok{NEXTCHAR STJ     NRET}
\NormalTok{         JXNZ    GOTX        * rX holds leftover chars from last word}
\NormalTok{INIT     J6N     READ        * rI6=0 =\textgreater{} need a new card}
\NormalTok{READ     IN      INPUT(READER)}
\NormalTok{         JBUS    *(READER)}
\NormalTok{         ENN6    16          * rI6 points to first word}
\NormalTok{GOTX     LDX     INPUT+16,6  * rX \textless{}{-} next word}
\NormalTok{         INC6    1}
\NormalTok{         ENTA    0}
\NormalTok{         SLAX    1           * next char {-}\textgreater{} rA (low byte)}
\NormalTok{NRET     JANZ    *            * skip blanks}
\NormalTok{         JMP     NEXTCHAR+1}

\NormalTok{* =============== SUBROUTINE: PUTCH (buffer \& punch) ===================}
\NormalTok{* input: rA has a character to append}
\NormalTok{* uses rI4 as OUTPUT pointer: {-}16,{-}15,...,{-}1 then 0 triggers punch}
\NormalTok{PUTCH    STJ     PRET}
\NormalTok{         STA     OUTPUT+16,4(1:1)     * store to pos (1:1) of word at rI4}
\NormalTok{         INC4    1}
\NormalTok{         J4N     PRET}
\NormalTok{         OUT     OUTPUT(PUNCH)}
\NormalTok{         JBUS    *(PUNCH)}
\NormalTok{         ENT4    {-}16}
\NormalTok{         MOVE    {-}1,1(16)             * blank buffer}
\NormalTok{PRET     JMP     *}

\NormalTok{* =============== MAIN ==================================================}
\NormalTok{START    ENT6    0            * init reader state for NEXTCHAR (as in text)}
\NormalTok{         ENTX    0}
\NormalTok{         ENT4    {-}16          * init output pointer}
\NormalTok{         MOVE    {-}1,1(16)     * blank output buffer}

\NormalTok{         JMP     NEXTCHAR}
\NormalTok{         CMPA    PERIOD}
\NormalTok{         JE      FINISH\_NOW   * empty stream?}
\NormalTok{         STA     CUR}
\NormalTok{         ENTA    1}
\NormalTok{         STA     CNT}

\NormalTok{MAINLP   JMP     NEXTCHAR}
\NormalTok{         CMPA    PERIOD}
\NormalTok{         JE      FLUSH\_LAST}
\NormalTok{         CMPA    CUR}
\NormalTok{         JE      INCR}
\NormalTok{         * new run}
\NormalTok{         LDA     CUR}
\NormalTok{         LDX     CNT}
\NormalTok{         JMP     EMITRUN}
\NormalTok{         STA     CUR          * EMITRUN returns with rA=CUR\textquotesingle{}, rX=CNT\textquotesingle{}}
\NormalTok{         ENTA    1}
\NormalTok{         STA     CNT}
\NormalTok{         JMP     MAINLP}

\NormalTok{INCR     ENTA    CNT}
\NormalTok{         INC     1}
\NormalTok{         STA     CNT}
\NormalTok{         JMP     MAINLP}

\NormalTok{FLUSH\_LAST}
\NormalTok{         LDA     CUR}
\NormalTok{         LDX     CNT}
\NormalTok{         JMP     EMITRUN}
\NormalTok{         * emit \textquotesingle{}.\textquotesingle{} and finish}
\NormalTok{FINISH\_NOW}
\NormalTok{         ENTA    PERIOD(5:5)  * just the \textquotesingle{}.\textquotesingle{} character}
\NormalTok{         JMP     PUTCH}
\NormalTok{         OUT     OUTPUT(PUNCH)}
\NormalTok{         HLT}

\NormalTok{* =============== EMITRUN(CUR in rA, CNT in rX) ========================}
\NormalTok{* Emits minimal code for CNT copies of char CUR.}
\NormalTok{* On return: none (caller reloads CUR/CNT for next run).}
\NormalTok{EMITRUN  STJ     ERRET}
\NormalTok{         STA     CCHAR}
\NormalTok{         STX     CCNT}

\NormalTok{         * is CUR a digit?  (MIX digit codes 30..39)}
\NormalTok{         LDA     CCHAR}
\NormalTok{         DECA    30}
\NormalTok{         JGE     CHECKHI}
\NormalTok{NOTDIG   * {-}{-}{-} NONDIGIT BRANCH {-}{-}{-}}
\NormalTok{         * q = CNT/10}
\NormalTok{         ENTA    10}
\NormalTok{         STA     TEN}
\NormalTok{         LDA     CCNT}
\NormalTok{         DIV     TEN           * rA \textless{}{-} q, rX \textless{}{-} r}
\NormalTok{         STA     Q}
\NormalTok{         STX     R}
\NormalTok{         * emit q times \textquotesingle{}9 CUR\textquotesingle{}}
\NormalTok{QLOOP    LDA     Q}
\NormalTok{         JE      REM}
\NormalTok{         ENTA    39            * \textquotesingle{}9\textquotesingle{}}
\NormalTok{         JMP     PUTCH}
\NormalTok{         LDA     CCHAR}
\NormalTok{         JMP     PUTCH}
\NormalTok{         DEC     Q,1}
\NormalTok{         JMP     QLOOP}
\NormalTok{REM      LDA     R}
\NormalTok{         JE      ERRET         * r==0}
\NormalTok{         J1Z     SINGLE        * r==1 {-}\textgreater{} single CUR}
\NormalTok{         DEC     R,1           * r in 2..9 {-}\textgreater{} (r{-}1)}
\NormalTok{         ADD     29            * convert 1..8 to digits 1..8 (codes 31..38)}
\NormalTok{         JMP     PUTCH}
\NormalTok{         LDA     CCHAR}
\NormalTok{         JMP     PUTCH}
\NormalTok{         JMP     ERRET}
\NormalTok{SINGLE   LDA     CCHAR}
\NormalTok{         JMP     PUTCH}
\NormalTok{         JMP     ERRET}

\NormalTok{CHECKHI  CMPA    =9=           * compare (CUR{-}30) with 9}
\NormalTok{         JG      NOTDIG        * \textgreater{}9 means not a digit actually}
\NormalTok{         * {-}{-}{-} DIGIT BRANCH {-}{-}{-}}
\NormalTok{         LDA     CCNT}
\NormalTok{         DIV     TEN           * q \textless{}{-} rA, r \textless{}{-} rX}
\NormalTok{         STA     Q}
\NormalTok{         STX     R}
\NormalTok{QDLOOP   LDA     Q}
\NormalTok{         JE      RPART}
\NormalTok{         ENTA    39            * \textquotesingle{}9\textquotesingle{}}
\NormalTok{         JMP     PUTCH}
\NormalTok{         LDA     CCHAR}
\NormalTok{         JMP     PUTCH}
\NormalTok{         DEC     Q,1}
\NormalTok{         JMP     QDLOOP}
\NormalTok{RPART    LDA     R}
\NormalTok{         JE      ERRET}
\NormalTok{         DEC     R,1}
\NormalTok{         ADD     29            * 0..8 {-}\textgreater{} 0..8 digits (codes 30..38); r=1 {-}\textgreater{} \textquotesingle{}0\textquotesingle{}}
\NormalTok{         JMP     PUTCH}
\NormalTok{         LDA     CCHAR}
\NormalTok{         JMP     PUTCH}
\NormalTok{ERRET    JMP     *}

\NormalTok{* Data}
\NormalTok{CUR      CON     0}
\NormalTok{CNT      CON     0}
\NormalTok{CCHAR    CON     0}
\NormalTok{CCNT     CON     0}
\NormalTok{Q        CON     0}
\NormalTok{R        CON     0}
\NormalTok{TEN      CON     10}
\NormalTok{PERIOD   ALF     \textquotesingle{}    .\textquotesingle{}}

\NormalTok{         END     START}
\end{Highlighting}
\end{Shaded}

Notes and alignment with the text

\begin{itemize}
\tightlist
\item
  Input side: \texttt{NEXTCHAR} is exactly the ``nonblank character''
  reader used in the coroutine example (read card, step through 5-char
  fields, skip blanks). 
\item
  Output side: Unlike OUT in the text (which groups in threes and
  inserts blanks), our \texttt{PUTCH} just packs characters into an
  80-column buffer and punches on full or at the end; the exercise asks
  to produce cards like (2), i.e., no grouping.
\item
  Termination: We stop at the first period, mirroring the original
  semantics.
\end{itemize}

Walk-through on Knuth's example

Given output (3)

\begin{verbatim}
ABB BEE EEE E44 446 66F GZY W22 220 0PQ 999 999 999 R.
\end{verbatim}

Ignoring blanks, the stream ends at the first \texttt{.}. Runs include
\texttt{A(1),\ B(2),\ E(3),\ 4(2),\ 6(3),\ F(1),\ G(1),\ Z(1),\ Y(1),\ W(1),\ 2(2),\ 0(3),\ P(1),\ Q(1),\ 9(3),\ 9(3),\ 9(3),\ R(1)}.

\begin{itemize}
\tightlist
\item
  For nondigits (e.g., \texttt{Z(1)}), we output the char itself
  (\texttt{Z}), never \texttt{0Z}.
\item
  For digits-as-data (e.g., \texttt{0(3)}), we must use pairs:
  \texttt{2} copies beyond the first encoded as
  \texttt{(3−1)=2\ →\ \textquotesingle{}2\textquotesingle{}} so emit
  \texttt{20}. More generally, runs ≥10 break as \texttt{9c} blocks +
  remainder.
\end{itemize}

Composing all runs yields exactly a minimal string of the form (2), and
as required, the stray \texttt{0Z} will not appear in our output.

\bookmarksetup{startatroot}

\chapter{1.4.3.1 A MIX simulator}\label{a-mix-simulator}

\section{1 {[}14{]}}\label{section-164}

What FCHECK does, and where it's used

\begin{itemize}
\item
  Purpose. \texttt{FCHECK} decodes the F--field (in \texttt{rI3}) as a
  \emph{partial field specification} of the form \(8L + R\), ensures
  \(L \le R \le 5\), then leaves the results in convenient places: it
  stores \(L\) and \(R\) in memory cells \texttt{L} and \texttt{R}, and
  sets \texttt{rI1\ ←\ L}. (If invalid, it jumps to \texttt{FERROR}.)
\item
  In \texttt{GETV}. The value-extraction subroutine begins by calling
  \texttt{FCHECK}, then fetches memory (\texttt{MEMORY}), and finally
  uses \texttt{L} and \texttt{R} to select the requested field and
  normalize the sign. (If the sign isn't included, it forces the sign to
  \texttt{+}.)
\item
  In \texttt{STORE}. The store routine also begins with \texttt{FCHECK},
  then uses \texttt{L}/\texttt{R} to splice the source register into the
  destination word---handling the special case when the sign-byte (i.e.,
  \texttt{L=0}) is included.
\item
  In \texttt{COMPARE}. The program \emph{intentionally} avoids calling
  \texttt{FCHECK} twice: it calls \texttt{GETV} (which runs
  \texttt{FCHECK} once), then enters \texttt{GETV} through the special
  entrance \texttt{GETAV} (which jumps past the \texttt{FCHECK} code) to
  reuse the already-decoded \texttt{L}. This is the one place where the
  code is already organized to keep the \texttt{FCHECK} overhead to a
  single call.
\item
  The exercise note. The exercise explicitly points you to ``step 3
  (Reexamination)'' at the end of §1.4.1, which suggests \emph{enlarging
  or merging subroutines when the caller always does the same standard
  setup/cleanup around them}. The exercise statement itself appears at
  the bottom of the simulator section.
\end{itemize}

Diagnosis

Across the simulator:

\begin{itemize}
\tightlist
\item
  \texttt{FCHECK} is always followed immediately by \emph{very similar}
  boilerplate in its callers (e.g., loading \texttt{-R} into
  \texttt{rI2}, checking whether the sign is included via \texttt{L==0},
  then doing standard shifts).
\item
  Only \texttt{COMPARE} already amortizes \texttt{FCHECK} over two uses
  of the same field (memory and register) by using the \texttt{GETAV}
  entrance.
\end{itemize}

That's a textbook case for Step-3 refactoring: push the common pre/post
work into a single place and remove redundant calls.

A better organization (two practical refactorings)

You can do either (A) ``pre-decode F once per instruction that needs
(L:R)'' or (B) ``absorb/inline the decode into the two places that
actually need it.'' Both obey Step-3.

\begin{enumerate}
\def\labelenumi{\Alph{enumi})}
\tightlist
\item
  Pre-decode once (preferred if you're willing to touch the control
  path)
\end{enumerate}

\begin{enumerate}
\def\labelenumi{\arabic{enumi}.}
\item
  Mark which opcodes use (L:R). Add a small parallel table to
  \texttt{OPTABLE} that classifies the F--field usage of each opcode
  (e.g., 0 = not a partial field, 1 = partial field (L:R), 2 = shift
  count, \ldots). The control routine already fetches \texttt{F} into
  \texttt{rI3} (line 027). Use your new table to decide whether to
  decode now.
\item
  Decode centrally. If the current opcode's F--usage is ``partial
  field,'' call a beefed-up decoder once and for all. That decoder
  should:

  \begin{itemize}
  \tightlist
  \item
    Validate and compute \texttt{L} and \texttt{R} (as \texttt{FCHECK}
    does).
  \item
    Also set \texttt{rI2\ ←\ −R} since every client immediately needs
    it.
  \item
    Optionally set a one-byte flag \texttt{SFLAG\ ←\ (L==0)} so clients
    don't have to re-test.
  \end{itemize}
\item
  Use ``prepared'' entries.

  \begin{itemize}
  \tightlist
  \item
    Change callers that currently do \texttt{JMP\ GETV} to
    \texttt{JMP\ GETAV} (the \emph{already-prepared} entrance) because
    \texttt{rI1}/\texttt{R} are now set; \texttt{GETAV} already jumps
    past the \texttt{FCHECK} block.
  \item
    Split \texttt{STORE} into two entries: keep \texttt{STORE} (old
    behavior) and add \texttt{STORE1} right after the \texttt{FCHECK}
    call site; have the opcode table jump to \texttt{STORE1} when the
    control routine has pre-decoded \texttt{F}. This removes the
    \texttt{FCHECK} call in the common path.
  \end{itemize}
\end{enumerate}

This removes \emph{all duplicate decoding} in instructions that use the
same (L:R) twice (e.g., \texttt{COMPARE}) and also avoids decoding in
instructions whose F--field isn't a partial field (e.g., \texttt{JUMP},
\texttt{ADDROP}). It follows Knuth's Step-3 guidance exactly: enlarge
the central ``control'' routine slightly to relieve its callees of
repeated setup.

\begin{enumerate}
\def\labelenumi{\Alph{enumi})}
\setcounter{enumi}{1}
\tightlist
\item
  Inline/merge (\texttt{FCHECK} → callers) with a macro (minimal edits)
\end{enumerate}

If you prefer not to touch the control path or tables:

\begin{itemize}
\tightlist
\item
  Delete \texttt{FCHECK} as a runtime subroutine and replace it by an
  \emph{assembly macro} (say, \texttt{DECODEF}) that expands to the 12
  lines of its body (compute \texttt{L}, \texttt{R}, checks, set
  \texttt{rI1}). Expand \texttt{DECODEF} at the start of \texttt{GETV}
  (before label \texttt{1F}) and at the start of \texttt{STORE}. Also
  have the macro optionally do \texttt{LD2N\ R} so all callers get
  \texttt{rI2} ``for free.''
\item
  Keep \texttt{GETAV} unchanged so \texttt{COMPARE} still jumps past the
  macro expansion in \texttt{GETV} and benefits from a single decode.
\end{itemize}

This keeps code simple, keeps \texttt{COMPARE} optimal, and respects
Step-3 (``perhaps several subroutines should be merged into one; or
perhaps a subroutine is called only once and should not be a subroutine
at all'').

Why this is better

\begin{itemize}
\tightlist
\item
  Fewer calls / simpler invariants. Operators that need fields no longer
  pay to rediscover \texttt{L}/\texttt{R} multiple times, and they get
  the always-needed \texttt{rI2\ ←\ −R} (and optionally \texttt{SFLAG})
  precomputed. See how \texttt{GETV} and \texttt{STORE} currently repeat
  the same checks right after \texttt{FCHECK}.
\item
  Matches the one existing optimization. \texttt{COMPARE} already
  factors out the second decode via \texttt{GETAV}; the refactor
  generalizes that cleanliness instead of being a one-off.
\item
  Faithful to Step-3. It's precisely the ``enlarge/merge subroutines''
  advice Knuth gives after completing a first cut: push common pre/post
  code down so the outer structure is cleaner.
\end{itemize}

Concrete changes you'd make

\begin{itemize}
\tightlist
\item
  Control routine (option A): after \texttt{LD3\ INST(4:4)} (line 027),
  if the opcode's F--class is ``(L:R)'', call \texttt{FDECODE} (your
  enhanced \texttt{FCHECK}) which sets \texttt{L}, \texttt{R},
  \texttt{rI1=L}, \texttt{rI2=−R}, and \texttt{SFLAG}. Then have the
  opcode table point to \texttt{GETAV} (instead of \texttt{GETV}) and to
  \texttt{STORE1} (the post-decode entry).
\item
  GETV: remove the prologue call to \texttt{FCHECK}; enter at the
  existing label after that call (\texttt{GETAV} already does exactly
  this).
\item
  STORE: add label \texttt{STORE1} at the point right after the original
  \texttt{JMP\ FCHECK} and retarget the table to \texttt{STORE1} when
  pre-decoding is in effect; otherwise fall back to \texttt{STORE} for
  compatibility.
\end{itemize}

Either approach solves Exercise 1: you've studied all uses of
\texttt{FCHECK}, noticed the duplicated patterns right after it, and
reorganized so the decode happens once (or is inlined where it belongs),
exactly the kind of Step-3 improvement Knuth advocates.

\section{2 {[}20{]}}\label{section-165}

Here's a clean, self-contained SHIFT routine for the MIX simulator. It
plugs into the program exactly where opcode 6 (``SHIFT'') is dispatched
in the operator table (see the listing that maps \texttt{F=0..5} to
\texttt{SLA,\ SRA,\ SLAX,\ SRAX,\ SLC,\ SRC}). I keep A/X signs in
\texttt{SIGNA/SIGNX} and only rearrange magnitudes in
\texttt{AREG/XREG}, matching the simulator's data layout.

What the routine must do (recap)

\begin{itemize}
\tightlist
\item
  F=0: \texttt{SLA\ m} --- shift A left by \texttt{m} bytes, zero-fill
  on the right; sign unchanged.
\item
  F=1: \texttt{SRA\ m} --- shift A right by \texttt{m} bytes, zero-fill
  on the left; sign unchanged.
\item
  F=2: \texttt{SLAX\ m} --- shift the (A,X) 10-byte pair left by
  \texttt{m}, zero-fill; signs unchanged.
\item
  F=3: \texttt{SRAX\ m} --- shift the (A,X) pair right by \texttt{m},
  zero-fill; signs unchanged.
\item
  F=4: \texttt{SLC\ m} --- rotate the (A,X) pair left by \texttt{m}
  (circular); signs unchanged.
\item
  F=5: \texttt{SRC\ m} --- rotate the (A,X) pair right by \texttt{m}
  (circular); signs unchanged. Knuth's exercise simply asks us to
  ``write the SHIFT routine\ldots{} (operation code 6).''
\end{itemize}

MIX code

\begin{Shaded}
\begin{Highlighting}[]
\NormalTok{* {-}{-}{-}{-}{-}{-}{-}{-}{-}{-}{-}{-}{-}{-}{-}{-}{-}{-}{-}{-}{-}{-}{-}{-}{-}{-}{-}{-}{-}{-}{-}{-}{-}{-}{-}{-}{-}{-}{-}{-}{-}{-}{-}{-}{-}{-}{-}{-}{-}{-}{-}{-}{-}{-}{-}{-}{-}{-}{-}{-}}
\NormalTok{* SHIFT — opcode 6}
\NormalTok{* Uses: rI3 = F (variant 0..5), rI5 = M (shift count, possibly signed)}
\NormalTok{* Touches: AREG/XREG magnitudes only; SIGNA/SIGNX unchanged}
\NormalTok{* Notes: If tracking timing exactly, add |M| to TIME in this routine,}
\NormalTok{*        since OPTABLE contributes the base 2u already for SHIFT.}
\NormalTok{* {-}{-}{-}{-}{-}{-}{-}{-}{-}{-}{-}{-}{-}{-}{-}{-}{-}{-}{-}{-}{-}{-}{-}{-}{-}{-}{-}{-}{-}{-}{-}{-}{-}{-}{-}{-}{-}{-}{-}{-}{-}{-}{-}{-}{-}{-}{-}{-}{-}{-}{-}{-}{-}{-}{-}{-}{-}{-}{-}{-}}
\NormalTok{SHIFT   ENT2    0,5         rI2 ← M (shift count)}
\NormalTok{        J2NN    *+2         make rI2 ← |M|}
\NormalTok{        ENN2    0,2}
\NormalTok{        CMP3    =5=         F must be 0..5}
\NormalTok{        JG      CYCLE       (defensive: treat out{-}of{-}range F as NOP)}

\NormalTok{        JMP     SHTAB,3     dispatch by F}

\NormalTok{* {-}{-}{-} table of tiny trampolines (same pattern used elsewhere) {-}{-}{-}}
\NormalTok{SHTAB   JMP     S\_SLA}
\NormalTok{        JMP     S\_SRA}
\NormalTok{        JMP     S\_SLAX}
\NormalTok{        JMP     S\_SRAX}
\NormalTok{        JMP     S\_SLC}
\NormalTok{        JMP     S\_SRC}

\NormalTok{* {-}{-}{-} F=0: SLA m (A only) {-}{-}{-}}
\NormalTok{S\_SLA   LDA     AREG}
\NormalTok{        SLA     0,2         shift A left by rI2 bytes}
\NormalTok{        STA     AREG}
\NormalTok{        JMP     S\_DONE}

\NormalTok{* {-}{-}{-} F=1: SRA m (A only) {-}{-}{-}}
\NormalTok{S\_SRA   LDA     AREG}
\NormalTok{        SRA     0,2         shift A right by rI2 bytes}
\NormalTok{        STA     AREG}
\NormalTok{        JMP     S\_DONE}

\NormalTok{* {-}{-}{-} F=2: SLAX m (A,X together) {-}{-}{-}}
\NormalTok{S\_SLAX  LDA     AREG}
\NormalTok{        LDX     XREG}
\NormalTok{        SLAX    0,2         shift (A,X) left by rI2 bytes}
\NormalTok{        STA     AREG}
\NormalTok{        STX     XREG}
\NormalTok{        JMP     S\_DONE}

\NormalTok{* {-}{-}{-} F=3: SRAX m (A,X together) {-}{-}{-}}
\NormalTok{S\_SRAX  LDA     AREG}
\NormalTok{        LDX     XREG}
\NormalTok{        SRAX    0,2         shift (A,X) right by rI2 bytes}
\NormalTok{        STA     AREG}
\NormalTok{        STX     XREG}
\NormalTok{        JMP     S\_DONE}

\NormalTok{* {-}{-}{-} F=4: SLC m (A,X circular left) {-}{-}{-}}
\NormalTok{S\_SLC   LDA     AREG}
\NormalTok{        LDX     XREG}
\NormalTok{        SLC     0,2         rotate (A,X) left by rI2 bytes}
\NormalTok{        STA     AREG}
\NormalTok{        STX     XREG}
\NormalTok{        JMP     S\_DONE}

\NormalTok{* {-}{-}{-} F=5: SRC m (A,X circular right) {-}{-}{-}}
\NormalTok{S\_SRC   LDA     AREG}
\NormalTok{        LDX     XREG}
\NormalTok{        SRC     0,2         rotate (A,X) right by rI2 bytes}
\NormalTok{        STA     AREG}
\NormalTok{        STX     XREG}
\NormalTok{        JMP     S\_DONE}

\NormalTok{* {-}{-}{-} wrap up; optional timing bump here {-}{-}{-}}
\NormalTok{S\_DONE  * If desired:  TIME ← TIME + rI2  (|M|)}
\NormalTok{        JMP     CYCLE}
\end{Highlighting}
\end{Shaded}

Why this is correct (brief reasoning)

\begin{enumerate}
\def\labelenumi{\arabic{enumi}.}
\tightlist
\item
  Signs preserved. The simulator stores signs separately
  (\texttt{SIGNA}, \texttt{SIGNX}). We only load/store magnitudes from
  \texttt{AREG}/\texttt{XREG}, so shifting cannot disturb signs.
\item
  Zero-fill vs rotate. Host MIX's \texttt{SLA/SRA/SLAX/SRAX} zero-fill
  the vacated bytes; \texttt{SLC/SRC} rotate within the 10-byte
  \texttt{rAX} pair. That exactly matches MIX semantics for opcode 6
  variants.
\item
  Large or negative counts. We take \texttt{\textbar{}M\textbar{}}. For
  \texttt{SLA/SRA} beyond the register width, the host shift yields
  zeros (desired). For rotations, host \texttt{SLC/SRC} behave modulo 10
  bytes, so any \texttt{M} works sensibly.
\item
  Timing (if modeled). The OPTABLE entry already accounts for the base
  ``2u'' per SHIFT; if you're simulating time precisely, add
  \texttt{\textbar{}M\textbar{}} here to reflect per-byte steps.
  (Knuth's listing shows SHIFT's table entry and the pattern for per-op
  timing updates.)
\end{enumerate}

\begin{Shaded}
\begin{Highlighting}[]
\NormalTok{*  SHIFT operators (opcode = 6)}
\NormalTok{*  Uses rI3 = F, rI5 = M (already set by the control routine).}
\NormalTok{*  Count = |M|; signs are preserved by MIX shift ops automatically.}

\NormalTok{SHIFT   DEC3    5           Is F \textgreater{} 5?}
\NormalTok{        J3P     FERROR      If so, invalid F for SHIFT.}

\NormalTok{        LD1     M(1:2)      rI1 ← |M| (shift count, magnitude of M).}

\NormalTok{        JMP     SFTABL,3    Dispatch to SLA/SRA/SLAX/SRAX/SLC/SRC.}

\NormalTok{* {-}{-}{-} F = 0 : SLA (shift left rA) {-}{-}{-}}
\NormalTok{SLAOP   LDA     AREG}
\NormalTok{        SLA     0,1         Shift rA left by rI1 bytes, zero{-}fill right.}
\NormalTok{        STA     AREG}
\NormalTok{        JMP     CYCLE}

\NormalTok{* {-}{-}{-} F = 1 : SRA (shift right rA) {-}{-}{-}}
\NormalTok{SRAOP   LDA     AREG}
\NormalTok{        SRA     0,1         Shift rA right by rI1 bytes, zero{-}fill left.}
\NormalTok{        STA     AREG}
\NormalTok{        JMP     CYCLE}

\NormalTok{* {-}{-}{-} F = 2 : SLAX (shift left rAX as 10{-}byte word) {-}{-}{-}}
\NormalTok{SLAXOP  LDA     AREG}
\NormalTok{        LDX     XREG}
\NormalTok{        SLAX    0,1         Shift rAX left by rI1 bytes, zero{-}fill right.}
\NormalTok{        STA     AREG}
\NormalTok{        STX     XREG}
\NormalTok{        JMP     CYCLE}

\NormalTok{* {-}{-}{-} F = 3 : SRAX (shift right rAX) {-}{-}{-}}
\NormalTok{SRAXOP  LDA     AREG}
\NormalTok{        LDX     XREG}
\NormalTok{        SRAX    0,1         Shift rAX right by rI1 bytes, zero{-}fill left.}
\NormalTok{        STA     AREG}
\NormalTok{        STX     XREG}
\NormalTok{        JMP     CYCLE}

\NormalTok{* {-}{-}{-} F = 4 : SLC (circular left rAX) {-}{-}{-}}
\NormalTok{SLCOP   LDA     AREG}
\NormalTok{        LDX     XREG}
\NormalTok{        SLC     0,1         Circular left shift of rAX by rI1 bytes.}
\NormalTok{        STA     AREG}
\NormalTok{        STX     XREG}
\NormalTok{        JMP     CYCLE}

\NormalTok{* {-}{-}{-} F = 5 : SRC (circular right rAX) {-}{-}{-}}
\NormalTok{SRCOP   LDA     AREG}
\NormalTok{        LDX     XREG}
\NormalTok{        SRC     0,1         Circular right shift of rAX by rI1 bytes.}
\NormalTok{        STA     AREG}
\NormalTok{        STX     XREG}
\NormalTok{        JMP     CYCLE}

\NormalTok{* Jump table for F = 0..5, addressed via "JMP SFTABL,3" after "DEC3 5".}
\NormalTok{* Place the label on the last entry (same pattern as the book’s JUMP table).}
\NormalTok{        JMP     SLAOP}
\NormalTok{        JMP     SRAOP}
\NormalTok{        JMP     SLAXOP}
\NormalTok{        JMP     SRAXOP}
\NormalTok{        JMP     SLCOP}
\NormalTok{SFTABL  JMP     SRCOP}
\end{Highlighting}
\end{Shaded}

Why this matches the book's simulator

\begin{enumerate}
\def\labelenumi{\arabic{enumi}.}
\item
  Dispatch pattern. The book's operator routines use \texttt{DEC3\ k}
  followed by \texttt{JMP\ TABLE,3}, with the table label on the
  \emph{last} entry so that negative offsets select earlier entries (see
  the \texttt{JUMP} operator's table). We do the same with
  \texttt{DEC3\ 5} and \texttt{SFTABL}.
\item
  Where the count comes from. The control routine has already computed
  the effective address and stored it in the variable \texttt{M}; it
  then sets \texttt{rI5\ ←\ M}. We take the magnitude of \texttt{M} via
  \texttt{LD1\ M(1:2)} to obtain a nonnegative byte count for shifting.
  (Signs of \texttt{M} are irrelevant to SHIFT.)
\item
  Which registers to touch. \texttt{SLA}/\texttt{SRA} affect only the A
  register; \texttt{SLAX}/\texttt{SRAX}/\texttt{SLC}/\texttt{SRC} treat
  \texttt{rA‖rX} as a 10-byte word. In the simulator, magnitudes live in
  \texttt{AREG}/\texttt{XREG}, and signs in
  \texttt{SIGNA}/\texttt{SIGNX}. MIX shift operators do not change
  signs, so we update only \texttt{AREG}/\texttt{XREG}.
\item
  Opcode and timing. The switch table ties opcode 6 to
  \texttt{SHIFT(2)}, so no extra timing logic is needed in the routine;
  the control code already accounts for the two time units when it
  indexes the table.
\end{enumerate}

Notes and edge cases

\begin{itemize}
\item
  If the count exceeds the register width (e.g., SLA/SRA by ≥5,
  SLAX/SRAX by ≥10), MIX semantics produce all zeros in the shifted
  region; \texttt{SLC}/\texttt{SRC} are naturally modulo-10 circular
  shifts. The host MIX instructions used above enforce these behaviors
  directly.
\item
  The routine intentionally ignores the sign of \texttt{M} by using
  \texttt{\textbar{}M\textbar{}} as the count, which matches standard
  descriptions of MIX shift instructions (``M specifies the number of
  bytes to be shifted'').
\end{itemize}

\section{3 {[}22{]}}\label{section-166}

What MOVE must do

\begin{itemize}
\tightlist
\item
  Semantics: move F words, one word at a time, from source block
  starting at M to destination block starting at the (simulated) rI1
  address; when finished, rI1 is increased by F. If F=0, nothing
  happens. Execution time is \texttt{1\ +\ 2F} units (the base ``1''
  comes from the opcode table; we add \texttt{2} per word moved in the
  routine). 
\item
  Overlap behavior (important!): copying proceeds ``forward'' (from
  \texttt{M} to \texttt{M+1} \ldots{} while storing to \texttt{rI1},
  \texttt{rI1+1}, \ldots). If destination starts inside the source and
  ahead of it (e.g., \texttt{rI1\ =\ M+1}), the same word may be
  replicated, exactly as described in the manual. Your simulator should
  therefore not try to be ``memmove-safe''; it must mimic MIX's
  specified effect. 
\item
  Control state the simulator has already set up for you:
  \texttt{rI5\ ←\ M}, \texttt{rI3\ ←\ F}, \texttt{rI4\ ←\ C}, and the
  opcode table has placed the base time in \texttt{TIME} before jumping
  to the per-opcode routine. The magnitudes/signs of simulated registers
  live in memory cells like \texttt{I1REG}, \texttt{SIGN1}, etc. 
\end{itemize}

MOVE routine (MIXAL)

Paste this under the label \texttt{MOVE} that's wired in the opcode
table.

\begin{Shaded}
\begin{Highlighting}[]
\NormalTok{* MOVE — opcode 7}
\NormalTok{MOVE    J3Z   CYCLE          * If F=0, nothing to do.}

\NormalTok{        JMP   MEMORY         * Fetch SOURCE word at address rI5:}
\NormalTok{                             * on return: rX = sign, rA = magnitude (of memory[rI5])}

\NormalTok{        SRAX  5              * Pack sign+magnitude so the full word sits in rX.}

\NormalTok{* Validate DEST = simulated rI1 (must be in [0, BEGIN))}
\NormalTok{        LD1   I1REG          * rI1 ← magnitude of simulated index register 1}
\NormalTok{        LDA   SIGN1          * rA  ← sign of simulated index register 1}
\NormalTok{        JAP   OKDEST         * if +, OK}
\NormalTok{        J1NZ  MEMERROR       * if negative and magnitude ≠ 0 → invalid address}
\NormalTok{        STZ   SIGN1(0:0)     * normalize {-}0 as +0 (permit address 0)}

\NormalTok{OKDEST  CMP1  =BEGIN=        * rI1 ≥ BEGIN ?}
\NormalTok{        JGE   MEMERROR       * out of simulated memory → error}

\NormalTok{* Store the fetched word into DEST and account for time}
\NormalTok{        STX   0,1            * memory[rI1] ← fetched full word}
\NormalTok{        LDA   CLOCK          * add 2 units per moved word (base 1 was set already)}
\NormalTok{        INCA  2}
\NormalTok{        STA   CLOCK}

\NormalTok{* Advance pointers/count and loop}
\NormalTok{        INC1  1              * rI1++}
\NormalTok{        ST1   I1REG          * write back simulated I1}
\NormalTok{        INC5  1              * rI5++ (SOURCE++)}
\NormalTok{        DEC3  1              * F{-}{-}}
\NormalTok{        JMP   MOVE           * keep moving until J3Z fires at top}
\end{Highlighting}
\end{Shaded}

This routine follows MIX's defined behavior exactly:

\begin{itemize}
\tightlist
\item
  It moves one word at a time, always reading the source first (via
  \texttt{MEMORY}) and then writing to the current destination
  (\texttt{STX\ 0,1}).
\item
  Because we progress \texttt{SOURCE} (\texttt{rI5}) and \texttt{DEST}
  (\texttt{rI1}) forward in lock-step, the overlap case where
  \texttt{DEST\ \textgreater{}\ SOURCE} produces the same
  ``replication'' effect shown in the book (that's the intended
  semantics). 
\item
  We charge time by adding \texttt{2} to \texttt{CLOCK} per moved word;
  together with the opcode table's base time, this yields
  \texttt{1\ +\ 2F}. 
\item
  We update the simulated \texttt{I1} immediately each step (by writing
  back \texttt{I1REG}), which is fine because programs can't observe
  intermediate state during an instruction; at completion, \texttt{rI1}
  has increased by F, as required. 
\item
  Address checks ensure both source (handled by \texttt{MEMORY}) and
  destination are inside simulated memory \texttt{{[}0,\ BEGIN)}, with
  \texttt{-0} treated as \texttt{+0}. 
\end{itemize}

Quick sanity check on overlap

\begin{itemize}
\tightlist
\item
  Example from the text: \texttt{F=3}, \texttt{M=1000}, \texttt{rI1=999}
  → you'll copy \texttt{1000→999}, \texttt{1001→1000},
  \texttt{1002→1001} (no surprise).
\item
  If instead \texttt{rI1=1001}, you'll copy \texttt{1000→1001}, then
  (since \texttt{1001} now holds the old 1000) \texttt{1001→1002}, then
  \texttt{1002→1003}, so the original \texttt{CONTENTS(1000)} appears
  three times---exactly as MIX specifies. 
\end{itemize}

Where these hooks live in the simulator

\begin{itemize}
\tightlist
\item
  The control routine unpacks \texttt{INST}, computes \texttt{M} with
  indexing, stores \texttt{F} in \texttt{rI3} and \texttt{M} in
  \texttt{rI5}, looks up the routine in \texttt{OPTABLE}, stores the
  base time in \texttt{TIME}, and jumps to \texttt{MOVE}. Your routine
  returns to \texttt{CYCLE} when \texttt{F} becomes zero. 
\end{itemize}

\section{4 {[}14{]}}\label{section-167}

Here's the smallest correct change to Knuth's MIX-simulator so it
``starts as if the GO button had been pushed.''

What the GO button does (per §1.3.1, Ex. 26)

When GO is pressed (with a card reader present), MIX (1) reads one card
into locations 0000--0015, (2) clears the overflow toggle and sets rJ ←
0, then (3) jumps to location 0000 to begin execution. 

What the simulator already does

In Program M (the simulator), the main loop already:

\begin{itemize}
\tightlist
\item
  Clears the overflow toggle (\texttt{STZ\ OVTog}) and
\item
  Begins execution at location 0 by setting the ``next instruction''
  register to 0 (\texttt{ENT6\ 0}, where rI6 is the location of the next
  instruction). 
\end{itemize}

So the only GO-state it doesn't establish is rJ = 0.

Minimal patch

Add one instruction at program start to zero the simulated J register:

\begin{verbatim}
001 * MIX SIMULATOR
002 ORIG 3500
003 BEGIN   STZ TIME(0:2)   ; clear time
004          STZ OVTog      ; clear overflow toggle  (GO requirement)
004a         STZ JREG       ; * new: rJ ← 0       (GO requirement)
005          STZ COMPI
006          ENT6 0         ; start at location 0   (GO requirement)
007 CYCLE    LDA CLOCK
...
\end{verbatim}

\begin{itemize}
\tightlist
\item
  \texttt{STZ\ JREG} makes the simulator's rJ match the GO button's
  post-reset state.
\item
  \texttt{STZ\ OVTog} and \texttt{ENT6\ 0} already matched the GO
  semantics (overflow cleared; ``PC'' set to 0000, since rI6 is the
  location of the next instruction). The control routine's use of rI6
  for the fetch is described just above Program M. 
\end{itemize}

Notes

\begin{itemize}
\tightlist
\item
  We aren't simulating the initial card read (GO step 1) here; I/O
  operators are intentionally left for later exercises (§1.4.3.1, Ex.
  6). This exercise only asks the simulator to \emph{begin} as if GO had
  been used.
\item
  After this one-line change, the simulator's initial state matches
  items (2)--(3) of GO: rJ=0, overflow cleared, and execution begins at
  0000.
\end{itemize}

That's it---one added \texttt{STZ\ JREG} at \texttt{BEGIN}.

\section{5 {[}24{]}}\label{section-168}

Timing rules we'll use (real MIX)

MIX uses time unit u. From the timing summary:

\begin{itemize}
\tightlist
\item
  All LOAD/STORE, ADD/SUB, comparisons, shifts: 2u
\item
  MOVE: 1u + 2u per word
\item
  MUL/NUM/CHAR: 10u; DIV: 12u
\item
  All remaining ops (including jumps, INC/DEC/ENT): 1u In particular,
  LDA = 2u, ENTA = 1u on real MIX.
\end{itemize}

Assumptions (to keep the path concrete)

\begin{itemize}
\tightlist
\item
  Instruction uses no indexing (I=0), so the indexing block in the
  control routine is skipped.
\item
  No errors (no range, field, or memory errors).
\item
  For LDA, the usual field (0:5) (sign included).
\item
  For ENTA, the address M ≠ 0 (the ``M=0'' special sign case is noted
  later). These are the straight-line, normal paths through the
  simulator.
\end{itemize}

The simulator's fixed control overhead per simulated instruction

Every instruction first goes through the control routine at label CYCLE
(load CLOCK, fetch INST, unpack M/I/F/C, consult OPTABLE, and jump).
Counting times by instruction class yields 34u before the jump to the
operator routine when I=0 (lines 007--035: loads/stores/shift = 2u each;
INC/DEC/J? = 1u).

Simulating LDA

Path: CYCLE → OPTABLE → LOAD → GETV → (FCHECK, MEMORY) → LOAD1 → SIZECHK
→ CYCLE. Key code: LOAD/LOADN (lines 196--205), GETV (149--158), FCHECK
(135--146), MEMORY (125--134), SIZECHK/OVCHECK (289--291, 191).

\begin{itemize}
\item
  Control overhead to reach operator: 34u.
\item
  Operator + subroutines:

  \begin{itemize}
  \tightlist
  \item
    \texttt{JMP\ GETV}: 1u
  \item
    FCHECK (valid field): STJ, ENTA, ENTX, DIV, CMPX, JG, STX, STA, LD1,
    CMPA, JLE → 27u
  \item
    MEMORY: STJ, JS? (not taken), CMP, JGE (not taken), LDX, ENTA, SRAX,
    LDA, JMP → 14u
  \item
    Return in GETV (sign included, so skip setting + sign): J1Z, LD2N,
    SRA, JMP → 6u
  \item
    LOAD1 store magnitude/sign: ENT1, STA, STX, JMP → 6u
  \item
    SIZECHK/OVCHECK (no index overflow, no overflow): LD1, J1Z, JNOV →
    4u
  \end{itemize}

  Sum operator-side = 1 + 27 + 14 + 6 + 6 + 4 = 58u (plus 5u for the two
  unconditional returns to GETV's and MEMORY's \texttt{JMP\ *} sites) →
  63u total operator segment.
\item
  Total simulated time for LDA = 34u + 63u = 97u.
\end{itemize}

Real MIX time for LDA = 2u. Slowdown ≈ 97/2 ≈ 48.5×.

Simulating ENTA

Path: CYCLE → OPTABLE → ADDROP → (LOAD1) → SIZECHK → CYCLE. Key code:
ADDROP block (265--278), then LOAD1 and SIZECHK as above.

\begin{itemize}
\item
  Control overhead to reach operator: 34u.
\item
  Operator path (typical M ≠ 0 branch):

  \begin{itemize}
  \tightlist
  \item
    ADDROP prologue \& sign/magnitude setup: DEC3, J3P (not taken),
    ENTX, JXNZ (taken), SRAX, LDA, ENT1, \texttt{JMP\ 1F,3},
    \texttt{JMP\ LOAD1} → 11u
  \item
    LOAD1 store magnitude/sign: STA, STX, \texttt{JMP\ SIZECHK} → 5u
  \item
    SIZECHK/OVCHECK: LD1, J1Z, JNOV → 4u
  \end{itemize}

  Sum operator-side = 20u.
\item
  Total simulated time for ENTA = 34u + 20u = 54u.
\end{itemize}

Real MIX time for ENTA = 1u. Slowdown ≈ 54/1 = 54×.

Notes on variants

\begin{itemize}
\tightlist
\item
  Indexing present (I≠0): the control routine executes the indexing
  block (lines 018--026), adding \textasciitilde15u more (several
  loads/stores/shifts/add/compare/jump).
\item
  LDA with field not including sign: GETV executes three extra ops
  (ENTX; two shifts), adding \textasciitilde5u.
\item
  ENTA with M = 0: executes two extra ops (LDX INST; ENTA 1) to pick up
  the instruction's sign rule, adding \textasciitilde3u.
\end{itemize}

Final comparison

\begin{itemize}
\tightlist
\item
  LDA: simulated \textasciitilde97u vs real 2u (≈ 49× slower).
\item
  ENTA: simulated \textasciitilde54u vs real 1u (≈ 54× slower).
\end{itemize}

These figures come straight from counting the executed instructions
along the normal paths in Program M and applying MIX's timing rules.

\section{6 {[}28{]}}\label{section-169}

Here's a clean way to add the missing I/O operators to the MIX
simulator, limited to unit 16 (card reader) and unit 18 (line printer),
and consistent with Knuth's timing model:

\begin{itemize}
\tightlist
\item
  Opcodes and meanings (from the MIX table): 34 = JBUS ``unit F busy?'',
  35 = IOC ``control, unit F'', 36 = IN ``input, unit F'', 37 = OUT
  ``output, unit F'', 38 = JRED ``unit F ready?'' 
\item
  Exercise requirement: implement JBUS, IOC, IN, OUT, JRED; allow only
  units 16 and 18; assume read-card and skip-to-new-page take T=10000u;
  print-line takes T=7500u. Also special-case \texttt{JBUS\ *} to avoid
  the simulator ``stopping.'' 
\end{itemize}

Design (what to simulate)

MIX overlaps I/O with computation. Issuing IN/OUT/IOC takes 1 unit of
CPU time and makes the device busy for T units; the CPU continues. The
queries JBUS/JRED test that busy/ready state. We therefore maintain
per-device ``ready times'' and a CPU clock:

\begin{itemize}
\tightlist
\item
  \texttt{CLOCK} --- current CPU time (u).
\item
  \texttt{RDRFREE} --- earliest time at which unit 16 (reader) is ready.
\item
  \texttt{PRTFREE} --- earliest time at which unit 18 (printer) is
  ready.
\item
  Constants: \texttt{TREAD=10000}, \texttt{TSKIP=10000},
  \texttt{TPRINT=7500}.
\item
  Overlap rule: when (re)starting an I/O, let
  \texttt{start\ =\ max(CLOCK,\ DEVFREE)} then set
  \texttt{DEVFREE\ =\ start\ +\ T}. CPU time still advances by 1 for the
  instruction itself (the ``1+T'' shown in the table), while the T
  portion accrues on the device timeline. 
\end{itemize}

Handlers to splice into ``INDIVIDUAL OPERATORS''

Below is MIXAL-style code that drops into the simulator's opcode
dispatch, following the conventions used elsewhere in §1.4.3.1 (use your
existing \texttt{MEMORY} check and jump simulation sequence, e.g.,
\texttt{JSJ\ JMP\ MEMORY\ /\ ENT6\ 0,5}).

Data and constants (in your simulator's data area):

\begin{verbatim}
TREAD   EQU     10000          * read-card (unit 16)
TSKIP   EQU     10000          * skip-to-new-page (unit 18, IOC)
TPRINT  EQU     7500           * print-line (unit 18)
CLOCK   CON     0
RDRFREE CON     0              * unit 16 ready time
PRTFREE CON     0              * unit 18 ready time
\end{verbatim}

Utility: pick device's FREE-TIME cell given F (16 or 18). If your decode
already provides F in a register, keep it; otherwise extract F the way
your other handlers do.

\begin{verbatim}
* On entry: r? = F (device number), rI5 = effective M (jump/address).
DEVPTR  CMP?   =16=
        JE     DEVPTR16
        CMP?   =18=
        JE     DEVPTR18
        JMP    FERROR            * only 16 and 18 are supported here
DEVPTR16 LDA   RDRFREE
        JMP    DEVPTRX
DEVPTR18 LDA   PRTFREE
DEVPTRX  STA   FREEBUF           * FREEBUF holds the fetched free time
        JMP    *+1
FREEBUF  CON   0                 * scratch word to hold device free time
\end{verbatim}

36 IN (unit F)

\begin{itemize}
\tightlist
\item
  Support: F=16 only (card reader).
\item
  Effect: begin a read ``now or when free'' and mark reader busy until
  \texttt{start+TREAD}. Optionally deposit a 16-word ``card image'' at
  M..M+15 (for a bare simulator this can be zeros or any deterministic
  pattern). CPU time +1.
\end{itemize}

\begin{verbatim}
INOP    ...                      * you’re in opcode 36 case
        * require F=16
        CMP?   =16=
        JNE    FERROR
        * start := max(CLOCK, RDRFREE)
        LDA    CLOCK
        CMPA   RDRFREE
        JGE    IN_START_READY
        LDA    RDRFREE
IN_START_READY
        ADD    =TREAD=
        STA    RDRFREE
        * (Optional) deposit 16 words at M..M+15:
        *   ENT1 0; LOOP 16x: STZ 0,5; INC5 1; DEC1 1; J1P LOOP
        * advance CPU time by 1
        LDA    CLOCK
        INCA   1
        STA    CLOCK
        JMP    CYCLE
\end{verbatim}

37 OUT (unit F)

\begin{itemize}
\tightlist
\item
  Support: F=18 only (line printer).
\item
  Effect: schedule a print of one ``line'' for \texttt{TPRINT} units.
  CPU time +1.
\end{itemize}

\begin{verbatim}
OUTOP   CMP?   =18=
        JNE    FERROR
        * start := max(CLOCK, PRTFREE)
        LDA    CLOCK
        CMPA   PRTFREE
        JGE    OUT_START_READY
        LDA    PRTFREE
OUT_START_READY
        ADD    =TPRINT=
        STA    PRTFREE
        * (optional: ignore actual text formatting in this bare simulator)
        LDA    CLOCK
        INCA   1
        STA    CLOCK
        JMP    CYCLE
\end{verbatim}

35 IOC (unit F)

\begin{itemize}
\tightlist
\item
  Support: F=18 used as ``skip to new page,'' cost \texttt{TSKIP}. (IOC
  for 16 is a no-op here.) CPU time +1.
\end{itemize}

\begin{verbatim}
IOCOP   CMP?   =18=
        JE     IOC_PRN
        * F=16 -> no-op for this minimal simulator
        LDA    CLOCK
        INCA   1
        STA    CLOCK
        JMP    CYCLE

IOC_PRN LDA    CLOCK
        CMPA   PRTFREE
        JGE    IOC_START_READY
        LDA    PRTFREE
IOC_START_READY
        ADD    =TSKIP=
        STA    PRTFREE
        LDA    CLOCK
        INCA   1
        STA    CLOCK
        JMP    CYCLE
\end{verbatim}

34 JBUS (unit F busy?)

\begin{itemize}
\tightlist
\item
  If the device is still busy (its FREE time \textgreater{} CLOCK), jump
  to M.
\item
  Special case for \texttt{JBUS\ *}: if busy and \texttt{M} is the
  current instruction address, don't spin at zero cost; instead advance
  \texttt{CLOCK} to the device's FREE time (this models ``wait until
  ready'' and prevents the simulator from ``stopping,'' per Knuth's
  note). If not busy, fall through. CPU time +1 (for the test itself). 
\end{itemize}

\begin{verbatim}
JBUSOP  * load device free time
        JSR    DEVPTR
        * busy if FREEBUF > CLOCK
        LDA    FREEBUF
        CMPA   CLOCK
        JLE    JBUS_NOTBUSY
        * device busy -> conditional jump
        * detect JBUS * (M == current PC). Your simulator already knows how to
        * compare M against “this instruction”; reuse that mechanism.
        JSR    IS_STAR           * sets CC: = if M==*
        JE     JBUS_WAIT
        * normal jump to M
        JSJ    JMP MEMORY        * validate address in rI5, as elsewhere
        ENT6   0,5               * simulate the jump
        * account 1 unit for the instruction itself
        LDA    CLOCK
        INCA   1
        STA    CLOCK
        JMP    CYCLE
JBUS_WAIT
        * wait until ready by jumping time forward
        LDA    FREEBUF
        STA    CLOCK
        JMP    CYCLE
JBUS_NOTBUSY
        * not busy, just consume 1 unit and continue
        LDA    CLOCK
        INCA   1
        STA    CLOCK
        JMP    CYCLE
\end{verbatim}

38 JRED (unit F ready?)

\begin{itemize}
\tightlist
\item
  Jump to M iff device is ready (\texttt{FREE\ ≤\ CLOCK}). CPU time +1
  either way.
\end{itemize}

\begin{verbatim}
JREDOP  JSR    DEVPTR
        LDA    FREEBUF
        CMPA   CLOCK
        JGT    JRED_NREADY
        * ready -> jump
        JSJ    JMP MEMORY
        ENT6   0,5
        LDA    CLOCK
        INCA   1
        STA    CLOCK
        JMP    CYCLE
JRED_NREADY
        LDA    CLOCK
        INCA   1
        STA    CLOCK
        JMP    CYCLE
\end{verbatim}

Notes and rationale

\begin{itemize}
\tightlist
\item
  Unit numbers: 16 is the card reader (e.g., the reset loader does
  \texttt{IN\ 0(16)}), 18 is the printer; restricting support to those
  two matches the exercise. 
\item
  Timing matches the MIX table (``1+T'' for IN/OUT) and the
  multiple-buffering discussion: T accrues on the device; the CPU
  continues after the 1-unit issue time.
\item
  \texttt{JBUS\ *} ``waits'' by fast-forwarding \texttt{CLOCK} to the
  device's free time, exactly to avoid a zero-time busy-wait loop (the
  note in the problem).
\end{itemize}

\section{7 {[}32{]}}\label{section-170}

Here's a clean, self-contained way to solve this exercies. The point is
to keep your MIX I/O simulator realistic: executing
\texttt{IN}/\texttt{OUT} must not move data immediately; the actual
transfer happens roughly halfway through the device's service time,
while the device remains busy for the full time. This prevents programs
from (incorrectly) assuming that \texttt{IN}/\texttt{OUT} are
instantaneous. The exercise builds directly on the I/O instructions and
device behavior defined earlier (units 16 = card reader, 18 = line
printer; \texttt{IN/OUT/IOC/JBUS/JRED} semantics), and on the timing
constants given in Ex. 6 (read-card or skip-page: \texttt{T=10000u};
print-line: \texttt{T=7500u}). The exercise's specific requirement
(``make the transmission occur after about half the time'') appears at
the end of the section's exercise list. 

Model the timing

Maintain a global simulator clock \texttt{now} (in MIX time units
\texttt{u}) and, for each device \texttt{d\ ∈\ \{16,18\}}, these fields:

\begin{itemize}
\tightlist
\item
  \texttt{ready\_at{[}d{]}} --- time when the device will become ready
  again.
\item
  \texttt{xfer\_at{[}d{]}} --- time when the \emph{data movement}
  (memory ↔ device buffer) actually occurs; this is
  \texttt{start\ +\ ceil(T{[}d{]}/2)}.
\item
  \texttt{pending{[}d{]}} ---
  \texttt{None\ \textbar{}\ ("IN",\ M)\ \textbar{}\ ("OUT",\ M)\ \textbar{}\ ("IOC",\ op)}
  describing the scheduled action.
\item
  Optional device buffers (e.g., next card image for 16; print line
  buffer for 18).
\end{itemize}

Timing constants (from Ex. 6): \texttt{T{[}16{]}.read\_card\ =\ 10000},
\texttt{T{[}18{]}.print\_line\ =\ 7500}, and
\texttt{T{[}18{]}.skip\_to\_new\_page\ (IOC)\ =\ 10000}. 

Event processing hook

Before fetching each instruction in your main \texttt{CYCLE} loop,
process any I/O events whose time has arrived:

\begin{verbatim}
def service_io_events():
    for d in (16, 18):
        if pending[d] and now >= xfer_at[d]:
            # Perform the scheduled data movement exactly once.
            if pending[d].op == "IN":
                mem[M ... M+block_size(d)-1] = device_in_buffer[d]
            elif pending[d].op == "OUT":
                device_out_buffer[d] = mem[M ... M+block_size(d)-1]
            # IOC with no data movement: nothing to do here.
            pending[d] = None
            xfer_at[d] = +∞
\end{verbatim}

This ``half-time'' transfer ensures that a program that uses the data
too early (after \texttt{IN}) or changes it too late (after
\texttt{OUT}) will behave incorrectly---exactly the mistake this
exercise wants to catch.

Instruction semantics (with half-time transfers)

Use the architectural rules from §1.3.1: \texttt{IN/OUT} wait until the
unit is ready \emph{to start}; the transfer completes later;
\texttt{IOC} initiates a device-specific operation;
\texttt{JRED}/\texttt{JBUS} test readiness. 

\texttt{IN\ M(F=unit)}

\begin{verbatim}
d = unit
wait_until_ready(d)               # if now < ready_at[d], advance now to ready_at[d]
start = now
pending[d] = ("IN", M)            # schedule memory ← device
xfer_at[d] = start + ceil(T[d]/2) # "approximately half" the device time
ready_at[d] = start + T[d]        # unit remains busy for full time
\end{verbatim}

Block size: 16 words for unit 16 (card reader), 24 for unit 18 (printer)
if you ever support input there; here we limit to 16 for the card
reader. 

\texttt{OUT\ M(F=unit)}

\begin{verbatim}
d = unit
wait_until_ready(d)
start = now
pending[d] = ("OUT", M)           # schedule device ← memory
xfer_at[d] = start + ceil(T[d]/2)
ready_at[d] = start + T[d]
\end{verbatim}

For unit 18, the printed line's \emph{content} is captured at
\texttt{xfer\_at{[}18{]}}; the physical printing completes at
\texttt{ready\_at{[}18{]}}.

\texttt{IOC\ M(F=unit)}

Wait until not busy, then perform the control; e.g., for the printer
``skip-to-new-page'' takes 10000 u but does no data transfer:

\begin{verbatim}
wait_until_ready(18)
start = now
pending[18] = ("IOC", "SKIP")
xfer_at[18] = +∞                  # no data movement
ready_at[18] = start + 10000
\end{verbatim}

(If you later support \texttt{IOC} on other units, replicate this
pattern with their \texttt{T}.)

\texttt{JRED\ M(F=unit)} and \texttt{JBUS\ M(F=unit)}

Use the §1.3.1 definitions: \texttt{JRED} jumps iff the unit is ready
(\texttt{now\ ≥\ ready\_at{[}d{]}}), \texttt{JBUS} jumps iff not ready. 

Important ``JBUS *'' special case

As hinted in Ex. 6, if the target of \texttt{JBUS} is the instruction
itself (busy-wait spin), you must advance \texttt{now} to the next
interesting time to avoid a frozen simulator:

\begin{verbatim}
if op == JBUS and target == PC and now < ready_at[d]:
    now = min(ready_at[d], min(xfer_at[k] for k with pending[k]))
\end{verbatim}

This is the pragmatic fix that keeps the loop progressing. 

Why this satisfies Exercise 7

\begin{itemize}
\tightlist
\item
  \texttt{IN}/\texttt{OUT} no longer ``complete'' at issue time: the
  simulator defers the data movement to \texttt{start\ +\ T/2}, while
  the device remains busy until \texttt{start\ +\ T}.
\item
  Programs that incorrectly assume immediate availability (after
  \texttt{IN}) or immediate consumption (after \texttt{OUT}) will
  misbehave, reproducing the real-machine hazard the exercise is
  designed to prevent. 
\end{itemize}

Quick timeline sanity checks

\begin{enumerate}
\def\labelenumi{\arabic{enumi}.}
\tightlist
\item
  Card read into 1000 with immediate use (wrong):
\end{enumerate}

\begin{verbatim}
t0:     IN 1000(16)           # schedules memory update at t0+5000; device busy until t0+10000
t0+1:   LDA 1000              # reads OLD data (correctly exposing the bug)
...
\end{verbatim}

\begin{enumerate}
\def\labelenumi{\arabic{enumi}.}
\setcounter{enumi}{1}
\tightlist
\item
  Printer OUT with later modification (also wrong):
\end{enumerate}

\begin{verbatim}
t0:     OUT 2000(18)          # captures line at t0+3750; printer finishes at t0+7500
t0+2000: STA 2000             # too early: the printed line will still reflect the pre-change contents
\end{verbatim}

\begin{enumerate}
\def\labelenumi{\arabic{enumi}.}
\setcounter{enumi}{2}
\tightlist
\item
  Proper synchronization:
\end{enumerate}

\begin{verbatim}
OUT 2000(18);  JBUS *(18)     # spin; simulator advances now to ~t0+7500
# when we exit, the line has been printed and the unit is ready again
\end{verbatim}

\begin{enumerate}
\def\labelenumi{\arabic{enumi})}
\tightlist
\item
  Data layout for I/O timing and pending actions
\end{enumerate}

We keep, per device, three fields:

\begin{itemize}
\tightlist
\item
  \texttt{READY\_AT{[}d{]}} --- when the device becomes ready again.
\item
  \texttt{XFER\_AT{[}d{]}} --- when the \emph{actual data movement}
  occurs.
\item
  \texttt{PEND{[}d{]}} --- pending action kind and base address (encoded
  in one MIX word: sign = kind; value = base address).
\end{itemize}

Choose encodings:

\begin{itemize}
\tightlist
\item
  \texttt{PEND} sign: \texttt{+} = none; \texttt{−} = has pending. The
  action code is stored in the most significant byte (B1) and the memory
  base address in the remaining bytes.
\item
  Action codes: \texttt{1} = IN, \texttt{2} = OUT, \texttt{3} = IOC (no
  data move).
\end{itemize}

You can keep device constants (\texttt{TREAD16}, \texttt{TPRINT18},
etc.) as you prefer.

\begin{Shaded}
\begin{Highlighting}[]
\NormalTok{          ORIG  2000}

\NormalTok{* {-}{-}{-} Global clock (u time units) {-}{-}{-}}
\NormalTok{NOW       CON   0}

\NormalTok{* {-}{-}{-} Device constants (adapt to your table) {-}{-}{-}}
\NormalTok{* D16 = card reader, D18 = line printer}
\NormalTok{D16       CON   16}
\NormalTok{D18       CON   18}

\NormalTok{* service times (u)}
\NormalTok{TREAD16   CON   10000          * read card / skip page (example)}
\NormalTok{TPRINT18  CON   7500}
\NormalTok{TSKIP18   CON   10000          * IOC skip page}

\NormalTok{* half{-}times (ceil(T/2))}
\NormalTok{HREAD16   CON   5000}
\NormalTok{HPRINT18  CON   3750}
\NormalTok{HSKIP18   CON   5000           * not used (no data move), keep for symmetry}

\NormalTok{* {-}{-}{-} Per{-}device tables: index by (device{-}16) if you only support 16..18}
\NormalTok{* We\textquotesingle{}ll allocate a small contiguous array and compute offsets.}
\NormalTok{NDEV      CON   3              * support devices 16,17,18}
\NormalTok{BASEDEV   CON   16}

\NormalTok{READY\_AT  CON   0,0,0          * READY\_AT[16], [17], [18]}
\NormalTok{XFER\_AT   CON   0,0,0          * XFER\_AT[16], [17], [18]}
\NormalTok{PEND      CON   +0,+0,+0       * +0 = none; −(action\textless{}\textless{}30 | base) = pending}

\NormalTok{* encodings}
\NormalTok{ACT\_IN    CON   1}
\NormalTok{ACT\_OUT   CON   2}
\NormalTok{ACT\_IOC   CON   3}
\NormalTok{NONE      CON   +0}

\NormalTok{* scratch}
\NormalTok{TMP       CON   0}
\NormalTok{IDX       CON   0              * working index = d {-} BASEDEV}
\NormalTok{W         CON   0              * generic temp}
\end{Highlighting}
\end{Shaded}

Helper to compute \texttt{IDX\ =\ d\ −\ 16} and leave it in
\texttt{IDX}:

\begin{Shaded}
\begin{Highlighting}[]
\NormalTok{IDXOF     STJ   IDXOF\_RET}
\NormalTok{          STA   TMP          * A has d}
\NormalTok{          LDA   TMP}
\NormalTok{          SUB   BASEDEV}
\NormalTok{          STA   IDX}
\NormalTok{          JBUS  *            * (no{-}op; keep same timing style)}
\NormalTok{IDXOF\_RET JMP   0}
\end{Highlighting}
\end{Shaded}

\begin{enumerate}
\def\labelenumi{\arabic{enumi})}
\setcounter{enumi}{1}
\tightlist
\item
  Small utilities
\end{enumerate}

Store a ``pending action'' word

We encode \texttt{(action,\ base)} in one word: \texttt{B1=action},
\texttt{B2..B5=base}, sign negative = ``pending''.

\begin{Shaded}
\begin{Highlighting}[]
\NormalTok{SETPEND   STJ   SP\_RET}
\NormalTok{          STA   TMP          * A = base address M}
\NormalTok{          ENT1  0}
\NormalTok{          ENTA  0}
\NormalTok{          OR    NONE         * clear A}
\NormalTok{          LDX   TMP}
\NormalTok{          SLA   5            * move X into B2..B5 (implementation detail varies)}
\NormalTok{          OR    ACT          * B1 = action (ACT in A)}
\NormalTok{          ST1   W            * save if you prefer a different scratch}
\NormalTok{          NEG   W            * make it pending (negative sign)}
\NormalTok{          LDA   IDX}
\NormalTok{          STA   TMP}
\NormalTok{          * PEND[base+IDX] ← W}
\NormalTok{          LDX   TMP}
\NormalTok{          LD2   =PEND=}
\NormalTok{          ADD   X}
\NormalTok{          STA   @0           * store W into PEND[IDX]}
\NormalTok{SP\_RET    JMP   0}
\end{Highlighting}
\end{Shaded}

(If your assembler lacks these exact bit ops, you can lay out two words
\texttt{PENDKIND{[}\ {]}} and \texttt{PENDBASE{[}\ {]}} instead; logic
stays the same and code gets simpler.)

Read/write table slots

If your MIXAL doesn't like indirect addressing macros, you can provide
tiny helpers. A simpler, clearer approach is to store each array as
separate contiguous words and use \texttt{ENTx}/\texttt{INCx} stepping.
To keep this answer concise, I'll use a compact ``index load/store''
idiom; translate to your style as needed.

\begin{enumerate}
\def\labelenumi{\arabic{enumi})}
\setcounter{enumi}{2}
\tightlist
\item
  I/O event poller (run before each instruction fetch)
\end{enumerate}

This is the heart of Exercise 7. It performs the \emph{deferred}
transfers at \texttt{XFER\_AT{[}d{]}} but does not make the device ready
until \texttt{READY\_AT{[}d{]}}.

Call \texttt{IOPOLL} at the top of your \texttt{CYCLE} loop, before
fetching the next instruction.

\begin{Shaded}
\begin{Highlighting}[]
\NormalTok{IOPOLL    STJ   IOP\_RET}
\NormalTok{          ENT3  0             * R3 = loop index 0..NDEV{-}1}
\NormalTok{IOP\_LP    CMP3  NDEV}
\NormalTok{          JGE   IOP\_DONE}

\NormalTok{          * compute absolute device number d = BASEDEV + R3}
\NormalTok{          LDA   BASEDEV}
\NormalTok{          ADD   R3}
\NormalTok{          STA   TMP           * TMP = d}
\NormalTok{          JSR   IDXOF         * set IDX = d {-} BASEDEV}

\NormalTok{          * if PEND[IDX] == NONE or XFER\_AT[IDX] \textgreater{} NOW, skip}
\NormalTok{          * Load XFER\_AT[IDX] into A:}
\NormalTok{          LD2   =XFER\_AT=}
\NormalTok{          ADD   IDX}
\NormalTok{          LDA   @0}
\NormalTok{          STA   W             * W = xfer time}
\NormalTok{          CMPA  NOW}
\NormalTok{          JG    IOP\_NEXT      * not yet time}

\NormalTok{          * Load PEND[IDX]}
\NormalTok{          LD2   =PEND=}
\NormalTok{          ADD   IDX}
\NormalTok{          LDA   @0}
\NormalTok{          JPL   IOP\_NEXT      * sign + =\textgreater{} no pending}

\NormalTok{          * decode action and base; here I’ll assume you split them}
\NormalTok{          * If you implemented PENDKIND/PENDBASE arrays, this becomes trivial:}
\NormalTok{          *   kind = PENDKIND[IDX]; base = PENDBASE[IDX]}

\NormalTok{          * {-}{-}{-} PERFORM THE DATA MOVE IF action == IN/OUT {-}{-}{-}}
\NormalTok{          * For clarity, I’ll branch by device number and action}
\NormalTok{          LDA   TMP           * d in A}
\NormalTok{          CMPA  D16}
\NormalTok{          JE    IOP\_D16}
\NormalTok{          CMPA  D18}
\NormalTok{          JE    IOP\_D18}
\NormalTok{          JMP   IOP\_CLEAR     * unknown device =\textgreater{} just clear pending}

\NormalTok{* {-}{-}{-} Device 16 (card reader): move 16 words at XFER time {-}{-}{-}}
\NormalTok{IOP\_D16   * if ACT\_IN: mem[M..M+15] \textless{}{-} device buffer}
\NormalTok{          * if ACT\_OUT: (rare for reader) ignore or treat as error}
\NormalTok{          * In most simulators, you keep a prepared card image buffer CARD[16].}
\NormalTok{          * Copy CARD[k] {-}\textgreater{} M[k] for k=0..15 (or vice versa for OUT).}
\NormalTok{          * Pseudocode in MIXAL outline:}

\NormalTok{          * ... decode base addr into R1 (left to you) ...}
\NormalTok{          * ... copy loop of 16 words memory \textless{}{-}{-} CARDBUF ...}
\NormalTok{          JMP   IOP\_CLEAR}

\NormalTok{* {-}{-}{-} Device 18 (line printer): capture line buffer from memory at XFER time {-}{-}{-}}
\NormalTok{IOP\_D18   * if ACT\_OUT: PRNBUF \textless{}{-} mem[M..M+N{-}1] (your N; often 24 words)}
\NormalTok{          * if ACT\_IN : usually unsupported}
\NormalTok{          * MIXAL copy loop outline:}
\NormalTok{          * ... decode base addr into R1 ...}
\NormalTok{          * ... copy loop memory {-}{-}\textgreater{} PRNBUF ...}
\NormalTok{          JMP   IOP\_CLEAR}

\NormalTok{* After moving data (or deciding no move for IOC), clear PEND and set XFER\_AT far in future}
\NormalTok{IOP\_CLEAR LD2   =PEND=}
\NormalTok{          ADD   IDX}
\NormalTok{          LDA   NONE}
\NormalTok{          STA   @0}

\NormalTok{          LD2   =XFER\_AT=}
\NormalTok{          ADD   IDX}
\NormalTok{          LDA   =+7777777=    * “infinite” (larger than any NOW)}
\NormalTok{          STA   @0}

\NormalTok{IOP\_NEXT  INC3  1}
\NormalTok{          JMP   IOP\_LP}

\NormalTok{IOP\_DONE  JMP   IOP\_RET}
\NormalTok{IOP\_RET   JMP   0}
\end{Highlighting}
\end{Shaded}

Notes:

\begin{itemize}
\tightlist
\item
  Keep your actual memory copy routines where you already copy between
  device buffers and core. The point is \emph{when} you do it: here,
  inside \texttt{IOPOLL} once \texttt{NOW\ ≥\ XFER\_AT{[}d{]}}, not when
  \texttt{IN/OUT} is issued.
\item
  IOC actions normally do no data move --- so \texttt{IOPOLL} simply
  clears \texttt{PEND} at XFER time only for IN/OUT; IOC needs no XFER
  slot at all (we'll set it to ``infinite'').
\end{itemize}

\begin{enumerate}
\def\labelenumi{\arabic{enumi})}
\setcounter{enumi}{3}
\tightlist
\item
  Modifications to instruction semantics
\end{enumerate}

You only change when data moves and how readiness is computed. The
opcode timing (issuing cost) stays as in your simulator.

IN M(F=unit d)

\begin{itemize}
\item
  If \texttt{NOW\ \textless{}\ READY\_AT{[}d{]}}, you must \emph{not}
  block the whole CPU (MIX uses \texttt{JBUS}/\texttt{JRED} for sync);
  the architectural spec is that the instruction \emph{initiates} the
  operation if the device is ready. Many simulators either (a) trap as
  ``busy'' or (b) start the op at \texttt{READY\_AT}. For simplicity, I
  recommend: ``if busy, advance \texttt{NOW} to
  \texttt{READY\_AT{[}d{]}} and start,'' which matches typical
  exercises.
\item
  Upon starting, set:

  \begin{itemize}
  \tightlist
  \item
    \texttt{PEND{[}d{]}\ =\ (IN,\ M)}.
  \item
    \texttt{XFER\_AT{[}d{]}\ =\ START\ +\ ⌈T/2⌉}.
  \item
    \texttt{READY\_AT{[}d{]}\ =\ START\ +\ T}.
  \end{itemize}
\end{itemize}

\begin{Shaded}
\begin{Highlighting}[]
\NormalTok{* On decode, A holds M (effective address), F holds unit d}
\NormalTok{IN\_OP     STJ   IN\_RET}
\NormalTok{          STA   TMP           * TMP = M}
\NormalTok{          LDA   FREG          * get device number d}
\NormalTok{          STA   W}
\NormalTok{          JSR   IDXOF         * IDX = d {-} BASEDEV}

\NormalTok{          * wait until device ready}
\NormalTok{          LD2   =READY\_AT=}
\NormalTok{          ADD   IDX}
\NormalTok{          LDA   @0}
\NormalTok{          CMPA  NOW}
\NormalTok{          JLE   IN\_START}
\NormalTok{          STA   NOW           * advance NOW to READY\_AT[d]}

\NormalTok{IN\_START  * schedule IN}
\NormalTok{          * XFER\_AT = NOW + HALF\_T[d]}
\NormalTok{          LDA   W             * d in A}
\NormalTok{          CMPA  D16}
\NormalTok{          JE    IN16}
\NormalTok{          CMPA  D18}
\NormalTok{          JE    IN18}

\NormalTok{          JMP   IN\_DONE       * unsupported device {-}\textgreater{} no{-}op or error}

\NormalTok{IN16      LDA   NOW}
\NormalTok{          ADD   HREAD16}
\NormalTok{          LD2   =XFER\_AT=}
\NormalTok{          ADD   IDX}
\NormalTok{          STA   @0}

\NormalTok{          * READY\_AT = NOW + TREAD16}
\NormalTok{          LDA   NOW}
\NormalTok{          ADD   TREAD16}
\NormalTok{          LD2   =READY\_AT=}
\NormalTok{          ADD   IDX}
\NormalTok{          STA   @0}

\NormalTok{          * PEND = (ACT\_IN, M)}
\NormalTok{          LDA   TMP          * base M}
\NormalTok{          STA   W}
\NormalTok{          LDA   ACT\_IN}
\NormalTok{          STA   ACT}
\NormalTok{          JSR   SETPEND}
\NormalTok{          JMP   IN\_DONE}

\NormalTok{IN18      * normally input from printer is invalid; ignore or treat as error}
\NormalTok{          JMP   IN\_DONE}

\NormalTok{IN\_DONE   JMP   IN\_RET}
\NormalTok{IN\_RET    JMP   0}
\end{Highlighting}
\end{Shaded}

OUT M(F=unit d)

Same pattern, with OUT action and device-specific T and HALF\_T.

\begin{Shaded}
\begin{Highlighting}[]
\NormalTok{OUT\_OP    STJ   OUT\_RET}
\NormalTok{          STA   TMP          * TMP = M}
\NormalTok{          LDA   FREG}
\NormalTok{          STA   W}
\NormalTok{          JSR   IDXOF}

\NormalTok{          * wait until device ready}
\NormalTok{          LD2   =READY\_AT=}
\NormalTok{          ADD   IDX}
\NormalTok{          LDA   @0}
\NormalTok{          CMPA  NOW}
\NormalTok{          JLE   OUT\_START}
\NormalTok{          STA   NOW}

\NormalTok{OUT\_START LDA   W}
\NormalTok{          CMPA  D16}
\NormalTok{          JE    OUT16}
\NormalTok{          CMPA  D18}
\NormalTok{          JE    OUT18}
\NormalTok{          JMP   OUT\_DONE}

\NormalTok{OUT16     * card reader OUT usually meaningless; either ignore or error}
\NormalTok{          JMP   OUT\_DONE}

\NormalTok{OUT18     * printer: capture line at half time; finish at full time}
\NormalTok{          LDA   NOW}
\NormalTok{          ADD   HPRINT18}
\NormalTok{          LD2   =XFER\_AT=}
\NormalTok{          ADD   IDX}
\NormalTok{          STA   @0}

\NormalTok{          LDA   NOW}
\NormalTok{          ADD   TPRINT18}
\NormalTok{          LD2   =READY\_AT=}
\NormalTok{          ADD   IDX}
\NormalTok{          STA   @0}

\NormalTok{          LDA   TMP}
\NormalTok{          STA   W}
\NormalTok{          LDA   ACT\_OUT}
\NormalTok{          STA   ACT}
\NormalTok{          JSR   SETPEND}
\NormalTok{          JMP   OUT\_DONE}

\NormalTok{OUT\_DONE  JMP   OUT\_RET}
\NormalTok{OUT\_RET   JMP   0}
\end{Highlighting}
\end{Shaded}

IOC M(F=unit d)

IOC has no data transfer. It just keeps the device busy for its service
time (e.g., skip page).

\begin{Shaded}
\begin{Highlighting}[]
\NormalTok{IOC\_OP    STJ   IOC\_RET}
\NormalTok{          LDA   FREG}
\NormalTok{          STA   W}
\NormalTok{          JSR   IDXOF}

\NormalTok{          * wait until device ready}
\NormalTok{          LD2   =READY\_AT=}
\NormalTok{          ADD   IDX}
\NormalTok{          LDA   @0}
\NormalTok{          CMPA  NOW}
\NormalTok{          JLE   IOC\_START}
\NormalTok{          STA   NOW}

\NormalTok{IOC\_START LDA   W}
\NormalTok{          CMPA  D18}
\NormalTok{          JE    IOC18}
\NormalTok{          JMP   IOC\_DONE}

\NormalTok{IOC18     * skip to new page, no data movement}
\NormalTok{          LDA   NOW}
\NormalTok{          ADD   TSKIP18}
\NormalTok{          LD2   =READY\_AT=}
\NormalTok{          ADD   IDX}
\NormalTok{          STA   @0}

\NormalTok{          * Ensure XFER\_AT “infinite” and PEND cleared}
\NormalTok{          LD2   =XFER\_AT=}
\NormalTok{          ADD   IDX}
\NormalTok{          LDA   =+7777777=}
\NormalTok{          STA   @0}

\NormalTok{          LD2   =PEND=}
\NormalTok{          ADD   IDX}
\NormalTok{          LDA   NONE}
\NormalTok{          STA   @0}
\NormalTok{          JMP   IOC\_DONE}

\NormalTok{IOC\_DONE  JMP   IOC\_RET}
\NormalTok{IOC\_RET   JMP   0}
\end{Highlighting}
\end{Shaded}

JRED / JBUS

\begin{itemize}
\tightlist
\item
  \texttt{JRED} jumps iff \texttt{NOW\ ≥\ READY\_AT{[}d{]}}.
\item
  \texttt{JBUS} jumps iff \texttt{NOW\ \textless{}\ READY\_AT{[}d{]}}.
\end{itemize}

Exercise-friendly special case: if the target is the instruction itself
(busy spin \texttt{JBUS\ *}), \emph{advance} \texttt{NOW} to the next
interesting time (the earlier of \texttt{READY\_AT{[}d{]}} and any
\texttt{XFER\_AT{[}k{]}} that is still in the future), then fall through
(don't infinite-loop).

\begin{Shaded}
\begin{Highlighting}[]
\NormalTok{JRED\_OP   STJ   JR\_RET}
\NormalTok{          LDA   FREG}
\NormalTok{          JSR   IDXOF}
\NormalTok{          LD2   =READY\_AT=}
\NormalTok{          ADD   IDX}
\NormalTok{          LDA   @0}
\NormalTok{          CMPA  NOW}
\NormalTok{          JLE   JR\_TAK       * ready → jump}
\NormalTok{          JMP   JR\_NOT}
\NormalTok{JR\_TAK    JMP   TARGET       * fill by your assembler}
\NormalTok{JR\_NOT    JMP   NEXTPC}
\NormalTok{JR\_RET    JMP   0}
\end{Highlighting}
\end{Shaded}

\begin{Shaded}
\begin{Highlighting}[]
\NormalTok{JBUS\_OP   STJ   JB\_RET}
\NormalTok{          LDA   FREG}
\NormalTok{          STA   W}
\NormalTok{          JSR   IDXOF}

\NormalTok{          * check readiness}
\NormalTok{          LD2   =READY\_AT=}
\NormalTok{          ADD   IDX}
\NormalTok{          LDA   @0}
\NormalTok{          CMPA  NOW}
\NormalTok{          JLE   JB\_NOT       * ready =\textgreater{} do not jump}

\NormalTok{          * device busy → would jump.}
\NormalTok{          * If target == PC (self{-}loop), advance NOW to min(next READY or next XFER across all devices).}
\NormalTok{          * Detect self{-}loop using your decoded target; here assume flag SELF is set by decoder.}
\NormalTok{          LDA   SELF}
\NormalTok{          JZE   JB\_TAK       * not self → just jump}

\NormalTok{          * compute NEXT = READY\_AT[d]}
\NormalTok{          STA   W            * preserve A}
\NormalTok{          LD2   =READY\_AT=}
\NormalTok{          ADD   IDX}
\NormalTok{          LDA   @0}
\NormalTok{          STA   TMP          * TMP = candidate time}

\NormalTok{          * scan XFER\_AT for an earlier future event}
\NormalTok{          ENT3  0}
\NormalTok{JB\_SCAN   CMP3  NDEV}
\NormalTok{          JGE   JB\_SET}
\NormalTok{          LD2   =XFER\_AT=}
\NormalTok{          ADD   R3}
\NormalTok{          LDA   @0}
\NormalTok{          CMPA  NOW}
\NormalTok{          JLE   JB\_NEXT}
\NormalTok{          CMPA  TMP}
\NormalTok{          JGE   JB\_NEXT}
\NormalTok{          STA   TMP          * tighter min}

\NormalTok{JB\_NEXT   INC3  1}
\NormalTok{          JMP   JB\_SCAN}

\NormalTok{JB\_SET    LDA   TMP}
\NormalTok{          CMPA  NOW}
\NormalTok{          JLE   JB\_TAK       * nothing ahead → just jump}
\NormalTok{          STA   NOW          * advance time to next event}
\NormalTok{          JMP   JB\_EXIT}

\NormalTok{JB\_TAK    JMP   TARGET}
\NormalTok{JB\_NOT    JMP   NEXTPC}
\NormalTok{JB\_EXIT   JMP   NEXTPC}
\NormalTok{JB\_RET    JMP   0}
\end{Highlighting}
\end{Shaded}

\begin{enumerate}
\def\labelenumi{\arabic{enumi})}
\setcounter{enumi}{4}
\tightlist
\item
  Hook \texttt{IOPOLL} into your fetch/execute loop
\end{enumerate}

At the very top of \texttt{CYCLE} (or just before fetching the next
instruction), call \texttt{IOPOLL}. This guarantees that once
\texttt{NOW\ ≥\ XFER\_AT{[}d{]}}, one (and only one) transfer happens
before any instruction that might observe memory.

\begin{Shaded}
\begin{Highlighting}[]
\NormalTok{CYCLE     * ... your loop prologue ...}
\NormalTok{          JSR   IOPOLL}
\NormalTok{          * fetch instruction at PC}
\NormalTok{          * decode \& dispatch}
\NormalTok{          * execute}
\NormalTok{          * advance PC}
\NormalTok{          JMP   CYCLE}
\end{Highlighting}
\end{Shaded}

\begin{enumerate}
\def\labelenumi{\arabic{enumi})}
\setcounter{enumi}{5}
\tightlist
\item
  Notes on memory moves
\end{enumerate}

\begin{itemize}
\item
  I've left the actual buffer copies as comments because every simulator
  lays out buffers differently. The change the exercise asks for is
  purely \emph{temporal}: move those copies from the \texttt{IN/OUT}
  handlers to \texttt{IOPOLL} at \texttt{NOW\ ≥\ XFER\_AT{[}d{]}}.
\item
  Typical choices:

  \begin{itemize}
  \tightlist
  \item
    Card reader: fixed 16 words \texttt{CARDBUF{[}0..15{]}} prepared by
    your card-file driver. At \texttt{IN} half-time:
    \texttt{mem{[}M..M+15{]}\ ←\ CARDBUF}.
  \item
    Printer: copy \texttt{mem{[}M..M+N-1{]}\ →\ PRNBUF} (N = your
    line-encoding footprint) at \texttt{OUT} half-time; the physical
    ``printing'' (e.g., writing to host file) can happen either
    immediately at half-time or at the moment the device becomes ready;
    pedagogically, capturing the line at half-time and writing it at
    \texttt{READY\_AT} is clean.
  \end{itemize}
\end{itemize}

\begin{enumerate}
\def\labelenumi{\arabic{enumi})}
\setcounter{enumi}{6}
\tightlist
\item
  Quick sanity tests
\end{enumerate}

\begin{itemize}
\tightlist
\item
  Unsafe early use after IN: \texttt{IN\ 1000(16);\ LDA\ 1000} should
  load the \emph{old} contents if executed before
  \texttt{XFER\_AT{[}16{]}}.
\item
  Unsafe late modification after OUT: \texttt{OUT\ 2000(18);\ STA\ 2000}
  before \texttt{XFER\_AT{[}18{]}} must not change the captured line;
  after \texttt{XFER\_AT{[}18{]}}, it must be too late.
\item
  Synchronization: \texttt{OUT\ 2000(18);\ JBUS\ *(18)} will advance
  \texttt{NOW} to \texttt{min(READY\_AT{[}18{]},\ next\ XFER)} and then
  proceed when ready.
\end{itemize}

That's the complete wiring you need. If you share the specific buffer
layouts you're using (sizes and labels for the card/printer buffers), I
can fill in the tiny copy loops to match your memory map exactly.

\section{8 {[}20{]}}\label{section-171}

Answer: False.

\begin{itemize}
\tightlist
\item
  In the control routine, the simulator fetches the next instruction
  with line 010: \texttt{LDA\ 0,6} (i.e., from simulated location held
  in the host index register rI6). Immediately after storing it, the
  simulator increments rI6 at line 012 with \texttt{INC6\ 1}. Thus the
  fetch uses the current value of rI6, and the increment prepares the
  next value.
\item
  The simulator initializes rI6 ← 0 at line 006. Later, jumps set rI6
  directly to the validated target address using \texttt{ENT6\ 0,5} in
  the JSJ/JMP handler (after a range check done by the \texttt{MEMORY}
  subroutine). This guarantees that when a jump occurs, the new rI6 is
  within {[}0, BEGIN).
\end{itemize}

However, there is no range check before the fetch at line 010. Consider
what happens when the simulated program executes an instruction located
at address BEGIN − 1 that does not change the location counter (i.e.,
not a jump):

\begin{enumerate}
\def\labelenumi{\arabic{enumi}.}
\tightlist
\item
  At the start of that iteration, rI6 = BEGIN − 1, so line 010 fetches
  from that address.
\item
  Line 012 executes \texttt{INC6\ 1}, so rI6 becomes BEGIN.
\item
  The operator returns to the control routine; on the next iteration,
  line 010 will be executed with rI6 = BEGIN, which violates
  \(rI_6 < \texttt{BEGIN}\). 
\end{enumerate}

A concrete two-instruction counterexample in the simulated memory
(addresses 0..BEGIN−1) is:

\begin{itemize}
\tightlist
\item
  At location 0: JSJ BEGIN−1 (unconditional jump to the last simulated
  cell; the simulator validates the target via \texttt{MEMORY}, then
  sets rI6 ← BEGIN−1).
\item
  At location BEGIN−1: any non-jump such as STZ 0 (store zero). After
  executing it, line 012 makes rI6 = BEGIN; the next fetch (line 010)
  now uses an out-of-range address.
\end{itemize}

For completeness, note that:

\begin{itemize}
\tightlist
\item
  The \texttt{MEMORY} subroutine explicitly enforces 0 ≤ address
  \textless{} BEGIN when the simulator accesses simulated memory via
  rI5, but this protection does not cover the line-010 fetch via rI6.
\item
  Therefore \(0 \le rI_6\) is preserved, but \(rI_6 < \texttt{BEGIN}\)
  is not guaranteed at every execution of line 010.
\end{itemize}

So the statement is false.

\bookmarksetup{startatroot}

\chapter{1.4.3.2 Trace routines}\label{trace-routines}

\section{1 {[}22{]}}\label{section-172}

On MIX, the only way to set rJ is to \textbf{execute a jump} whose
address is located at position α; executing \texttt{JMP\ …} from
location α sets rJ to α+1. Therefore, to restore rJ to the ``start''
value saved earlier (call that value \textbf{J0}), we can:

\begin{enumerate}
\def\labelenumi{\arabic{enumi}.}
\tightlist
\item
  Temporarily write the instruction \textbf{\texttt{JMP\ *+1}} into
  memory \textbf{at address J0−1};
\item
  Execute that instruction in place; this sets \textbf{rJ ← J0};
\item
  Put the original word back (so we don't alter the traced program);
\item
  Finally, exit the trace routine \textbf{without disturbing rJ} (so use
  \texttt{JSJ}, not \texttt{JMP}). Knuth's answer notes this last point
  explicitly: using \texttt{JSJ} is crucial.
\end{enumerate}

Patch (drop-in replacement for the leave block)

Below is a compact, self-contained replacement for the ``leave'' part of
the trace routine. It assumes your routine already saved the original J
in \texttt{JREG} (as in the text), and it uses I1 and X temporarily but
restores them.

\begin{Shaded}
\begin{Highlighting}[]
\NormalTok{* {-}{-}{-} NEW LEAVE SEQUENCE (restores rJ) {-}{-}{-}{-}{-}{-}{-}{-}{-}{-}{-}{-}{-}{-}{-}{-}{-}{-}{-}{-}{-}{-}{-}{-}{-}}
\NormalTok{XSV    ORIG  *+1          * temp: save X}
\NormalTok{I1SV   ORIG  *+1          * temp: save I1}
\NormalTok{ONE    JMP   *+1          * literal word used as a copier: "JMP *+1"}

\NormalTok{LEAVE  STX   XSV          * save rX; we’ll need it for the exit jump later}
\NormalTok{       ST1   I1SV         * save rI1; we’ll use I1 as a pointer}
\NormalTok{       LD1   JREG(0:2)    * I1 := saved start address J0}
\NormalTok{       LDA   {-}1,1         * A  := original contents at address J0{-}1}
\NormalTok{       LDX   ONE          * X  := the word "JMP *+1"}
\NormalTok{       STX   {-}1,1         * write that jump at address J0{-}1}
\NormalTok{       JMP   {-}1,1         * execute it → rJ := (J0{-}1)+1 = J0}

\NormalTok{* we’re now back here, with rJ restored to J0}
\NormalTok{       STA   {-}1,1         * restore the original word at J0{-}1}
\NormalTok{       LD1   I1SV         * restore rI1}
\NormalTok{       LDX   XSV          * restore rX}
\NormalTok{       LDA   AREG         * restore rA (as in the book)}
\NormalTok{LEAVEX JSJ   *            * exit without changing rJ}
\end{Highlighting}
\end{Shaded}

Why each step is safe

\begin{itemize}
\tightlist
\item
  \textbf{No persistent self-modification:} we restore the overwritten
  word at \texttt{J0−1} before leaving.
\item
  \textbf{Minimal register disturbance:} only rA, rX, rI1 are touched
  temporarily, and all three are restored; the comparison indicator and
  overflow toggle remain unaffected (matching the section's design
  goals).
\item
  \textbf{Correct exit:} \texttt{JSJ} leaves rJ as we've just set it to
  J0; using \texttt{JMP} here would clobber rJ again.
\end{itemize}

Where this fits in the printed routine

You replace the printed \textbf{lines that restore rA and jump out} (the
\texttt{LEAVE\ …\ /\ LEAVEX\ …} pair near the end) with the block above.
Those two lines, as printed, don't fix rJ; the patch does.

\section{2 {[}26{]}}\label{section-173}

Design

\begin{itemize}
\tightlist
\item
  A 100-word circular buffer at label \texttt{BUF} collects ten 10-word
  records; when full, we \texttt{OUT} it to tape 0.
\item
  We keep an index pointer in rI1 that advances by 10 after each record.
  To persist the pointer across iterations, we use a neat MIX trick:
  self-modify the address field of the instruction
  \texttt{PTR\ ENT1\ -100}, i.e., store the current rI1 back into
  \texttt{PTR(0:2)} each time; on the next pass, \texttt{ENT1} reloads
  the saved pointer.
\item
  We never disturb the traced program's state (except briefly using rA
  and rI1 inside the tracer); rA is restored before the traced
  instruction executes, as in the original routine.
\item
  We wait for the device when needed and flush the final partial buffer
  at the end (the final flush routine is a small add-on you call once
  after tracing stops).
\end{itemize}

Patch (delta against the routine on pp.~212--213)

The comments say exactly where to splice these lines relative to the
original labels/lines.

\begin{Shaded}
\begin{Highlighting}[]
\NormalTok{* {-}{-}{-} Reserve a 100{-}word buffer and constants (can be placed near the end) {-}{-}{-}}
\NormalTok{        ORIG    *+99            * reserve 100 words below BUF}
\NormalTok{BUF     CON     0               * BUF{-}99 .. BUF are the 100{-}word block}
\NormalTok{PTR     ENT1    {-}100            * rI1 := saved pointer (will be self{-}modified)}
\NormalTok{EIGHT   CON     8}

\NormalTok{I1SAVE  CON     0               * local save for rI1}

\NormalTok{* {-}{-}{-} At ENTER: after lines 02–04 (save tape{-}pointer register) {-}{-}{-}}
\NormalTok{        ST1     I1SAVE          * save caller’s rI1; tracer will use rI1}

\NormalTok{* {-}{-}{-} After line 05–07 (after CYCLE STA PREG(0:2)) {-}{-}{-}}
\NormalTok{        JBUS    *(0)            * ensure tape 0 not busy from a prior OUT}
\NormalTok{        STA     BUF+1,1(0:2)    * Word1 (0:2): location (A holds it here)}
\NormalTok{        ST1     I1REG           * (if you keep a global I1REG like the book)}

\NormalTok{* {-}{-}{-} After lines 08–11 (after fetching the instruction into INST) {-}{-}{-}}
\NormalTok{        STA     BUF+2,1         * Word2: instruction word}

\NormalTok{* {-}{-}{-} Still before the traced instruction executes: dump “before” registers {-}{-}{-}}
\NormalTok{        LDA     AREG            * Word3: A (before execution)}
\NormalTok{        STA     BUF+3,1}

\NormalTok{        LDA     I1SAVE          * Word4: I1 (before) — we changed rI1}
\NormalTok{        STA     BUF+4,1}
\NormalTok{        ST2     BUF+5,1         * Words5–9: I2..I6 (unchanged by tracer)}
\NormalTok{        ST3     BUF+6,1}
\NormalTok{        ST4     BUF+7,1}
\NormalTok{        ST5     BUF+8,1}
\NormalTok{        ST6     BUF+9,1}

\NormalTok{        STX     BUF+10,1        * Word10: X (before)}
\NormalTok{        LDA     JREG(0:2)}
\NormalTok{        STA     BUF+1,1(4:5)    * Word1 (4:5): J (before)}

\NormalTok{* {-}{-}{-} Pack CI (2/1/0) + 8*OVF into Word1 (3:3) {-}{-}{-}}
\NormalTok{        ENTA    0               * start with 0 (= “\textless{}”)}
\NormalTok{        JL      CI\_DONE         * if “\textless{}”, keep 0}
\NormalTok{        JE      CI\_EQ}
\NormalTok{        ENTA    2               * “\textgreater{}” → 2}
\NormalTok{        JMP     CI\_OV}
\NormalTok{CI\_EQ   ENTA    1               * “=” → 1}
\NormalTok{CI\_OV   JNOV    CI\_NOOV         * if overflow off, skip add}
\NormalTok{        ADD     EIGHT}
\NormalTok{CI\_NOOV STA     BUF+1,1(3:3)}
\NormalTok{CI\_DONE NOP}

\NormalTok{* {-}{-}{-} Advance pointer and flush when block is full {-}{-}{-}}
\NormalTok{        INC1    10              * next record slot within the 100{-}word block}
\NormalTok{        J1N     KEEP\_FILLING    * still negative → not yet full}
\NormalTok{        OUT     BUF{-}99(0)       * write the 100{-}word block to tape unit 0}
\NormalTok{KEEP\_FILLING}
\NormalTok{        ST1     PTR(0:2)        * persist pointer by self{-}modifying ENT1}
\NormalTok{        LD1     I1SAVE          * restore caller’s rI1 before continuing}

\NormalTok{* (the rest of the original routine continues: execute INST, etc.)}
\end{Highlighting}
\end{Shaded}

Notes and rationale:

\begin{itemize}
\tightlist
\item
  \textbf{Where values come from.} Right after
  \texttt{CYCLE\ STA\ PREG(0:2)} the A-register holds the next
  instruction's location, so we can store Word 1's (0:2) immediately.
  After \texttt{PREG\ LDA\ *\ /\ STA\ INST} we have the instruction word
  (Word 2). \texttt{AREG}, \texttt{JREG}, and \texttt{I1..I6} are the
  tracer's saved copies of the machine registers before the step; we
  read them \emph{before} executing the instruction so the ``before''
  requirement is met.
\item
  \textbf{CI/OV packing.} We derive the comparison indicator via
  branches (\texttt{JL/JE}) and add 8 iff overflow was set, then store
  that byte into field (3:3) of Word 1, exactly as specified by the
  exercise.
\item
  \textbf{Buffering \& I/O.} The \texttt{ORIG\ *+99} trick makes
  \texttt{BUF-99\ ..\ BUF} a contiguous 100-word area. We fill records
  using indexed addresses \texttt{BUF+…\ ,1}, starting with rI1 = −100;
  after ten records (\texttt{INC1\ 10} ten times) rI1 becomes 0 and we
  \texttt{OUT\ BUF-99(0)} to tape 0. We also issue \texttt{JBUS\ *(0)}
  up front to avoid overlapping I/O. A one-time ``final flush \&
  rewind'' routine should be called after tracing stops to write a
  partial block and rewind unit 0, as suggested in the official notes.
\end{itemize}

Why this is correct

\begin{itemize}
\tightlist
\item
  Every required field is captured \textbf{before} the traced
  instruction executes (we only restore rA just before execution, per
  the original routine's discipline).
\item
  The 100-word block layout matches the exercise's request for ``nine
  more ten-word groups in the same format,'' enabling efficient tape
  output with minimal space in core.
\item
  We preserve the traced program's state: only rA and rI1 are touched
  inside the tracer and then restored; CI and OV are left unchanged
  (except for the unavoidable tests), as the text advises.
\end{itemize}

\section{3 {[}10{]}}\label{section-174}

Why tape \textgreater{} printer for a trace routine

\begin{enumerate}
\def\labelenumi{\arabic{enumi}.}
\item
  \textbf{Throughput and latency} Line printers are slow, mechanical
  devices. Knuth's own timing note for printing a line is \(T=7500u\);
  if you print after (nearly) every instruction, output time dominates
  and the trace becomes painfully slow. Tape is block-oriented and
  written in large chunks, so per-word overhead is vastly lower; you can
  buffer many words and write them as a single I/O, reducing device
  start/stop costs and the number of waits.  Knuth also emphasizes that,
  for trace programs, ``the running time is consumed by output anyway,''
  so minimizing per-step output costs is crucial---tape helps do exactly
  that. 
\item
  \textbf{Buffering and structure} Exercise 2 explicitly formats the
  trace as fixed ten-word groups packed into 100-word tape blocks.
  That's ideal for efficient sequential I/O and easy post-processing.
  Printing would force you to convert binary register contents into text
  and paginate---more CPU time and more I/O operations. 
\item
  \textbf{Non-interference with the program under test} A traced program
  ``must not use the output device on which tracing information is being
  recorded.'' If you occupy the printer for tracing, the program can't
  print its own output. Using tape unit 0 for the trace keeps the
  printer free for the target program (or human-readable summaries
  later). 
\item
  \textbf{Post-processing flexibility} A tape log can be rewound,
  scanned, filtered, summarized, or reformatted into human-readable
  reports after the run (e.g., reconstructing flow, computing
  instruction frequencies, or converting only selected steps to text).
  With direct printing you pay the formatting cost up front and flood
  paper with data you may not actually need.
\item
  \textbf{Lower perturbation of behavior} Because tape output is faster
  and can be buffered, it slows the traced program less than constant
  printing---giving a trace that is closer to the program's ``natural''
  behavior (though still slower, as the text notes).
\end{enumerate}

In short: \textbf{tape is faster, buffer-friendly, nonintrusive, and
analyzable offline}, whereas printing is slow, formatting-heavy, and
monopolizes a device the program might also want to use.

\section{4 {[}25{]}}\label{section-175}

Here's what happens if you make the trace routine trace itself by
inserting, just before label \texttt{ENTER}, the two instructions

\begin{itemize}
\tightlist
\item
  \texttt{ENTX\ LEAVEX} (set rX to the address of \texttt{LEAVEX}), then
\item
  \texttt{JMP\ **+1} (jump into \texttt{ENTER}),
\end{itemize}

as the exercise proposes. (This is the routine on pp.~212--214.)

What the routine expects at entry

The trace routine is designed to be entered with rJ already set to the
\textbf{starting address to trace} and rX set to the \textbf{address
where tracing should stop}. The code then saves rA and rJ, takes the
starting location from rJ, and begins its fetch--decode--execute cycle
on the ``external'' program's instruction at that location.

Step-by-step with your two extra instructions

\begin{enumerate}
\def\labelenumi{\arabic{enumi}.}
\item
  \textbf{Prelude executes.} \texttt{ENTX\ LEAVEX} makes rX =
  address(\texttt{LEAVEX}); \texttt{JMP\ **+1} transfers control to
  \texttt{ENTER}. But a \texttt{JMP} also \textbf{loads rJ with the
  address of the \texttt{JMP} itself} (the prelude location). Thus, when
  the trace routine starts, its idea of the ``start address'' (from rJ)
  is the prelude \texttt{JMP}, not the first instruction of the routine
  proper. The routine saves rA and rJ, then sets its ``next
  instruction'' pointer to that \texttt{JMP}.
\item
  \textbf{First traced instruction is that prelude \texttt{JMP\ **+1}.}
  The tracer recognizes a jump and rewrites the instruction word it will
  execute so that execution flows through its \texttt{JUMP} handler,
  which updates the simulated J-register and sets the next location to
  \texttt{ENTER}. (That is the normal way this tracer executes an
  external jump.)
\item
  \textbf{Now it begins tracing the trace routine's own body.} It will
  fetch and execute, under trace, the instructions beginning at
  \texttt{ENTER} (e.g., \texttt{STX\ TEST(0:2)}, \texttt{STA\ AREG},
  \ldots). So far, so good.
\item
  \textbf{Self-destruction on \texttt{STJ}.} Very soon the tracer
  reaches the routine's own \texttt{STJ\ JREG} (the instruction that
  saves rJ early in the routine). The tracer \textbf{deliberately
  rewrites any external \texttt{STJ} into \texttt{STA}} before executing
  it, so the external program can't disturb its simulated J-register.
  But now the ``external program'' \emph{is the tracer itself}, so the
  tracer \textbf{changes its own \texttt{STJ\ JREG} into
  \texttt{STA\ JREG} in memory} and then executes that store-A. This
  destroys the saved J value (\texttt{JREG} now holds whatever was in
  rA), and the code itself has been permanently modified.
\item
  \textbf{Loss of the starting location and derailment.} Immediately
  afterward the tracer does \texttt{LDA\ JREG(0:2)} to get the next
  instruction location from \texttt{JREG}. Since \texttt{JREG} no longer
  contains a proper address (it was just overwritten with A), the tracer
  loads a garbage address, points its ``next instruction'' there, and
  proceeds to fetch/execute nonsense. From this point the trace prints
  nonsense or loops/crashes; it certainly doesn't produce a coherent
  self-trace.
\end{enumerate}

With only those two added instructions, \textbf{the given trace routine
cannot trace itself}. It wrecks its own code/data because its ``don't
let the external program alter J'' trick (rewriting \texttt{STJ} to
\texttt{STA}) applies to the tracer's \textbf{own} \texttt{STJ},
corrupting \texttt{JREG} and sending execution off the rails.

This is exactly why the text restricts traced programs so they
\textbf{don't store into the trace routine's locations}, and why a
robust ``self-trace'' needs extra machinery (e.g., two copies tracing
each other or special guards), as hinted by the following exercises.

\section{5 {[}28{]}}\label{section-176}

What actually happens

\begin{enumerate}
\def\labelenumi{\arabic{enumi}.}
\tightlist
\item
  \textbf{The tracer that you start keeps control.} Suppose we start
  tracer \textbf{T₁} with \texttt{J} pointing at the entry of tracer
  \textbf{T₂}. T₁ runs its own code, repeatedly:
\end{enumerate}

\begin{itemize}
\tightlist
\item
  fetches the current instruction \emph{from T₂'s memory} at address
  \texttt{PREG₁},
\item
  copies it into \textbf{T₁}'s \texttt{INST}, possibly edits it,
\item
  executes that copy at \textbf{T₁}'s \texttt{INST},
\item
  increments \textbf{T₁}'s \texttt{PREG} to the next address.
\end{itemize}

Crucially, \textbf{the instruction that actually executes is always the
copy in T₁'s \texttt{INST}, not the one in T₂}. So T₂ never ``takes
over.''

\begin{enumerate}
\def\labelenumi{\arabic{enumi}.}
\setcounter{enumi}{1}
\item
  \textbf{T₂ ``tries'' to trace T₁, but only as data.} Because T₂'s
  \emph{code} is what T₁ is stepping through, T₂ will execute its own
  statements like ``store the next address in \texttt{PREG₂},'' ``read
  the instruction from T₁ into \texttt{INST₂},'' etc. But when the time
  comes to ``execute the external instruction,'' the execution point is
  \textbf{T₁}'s \texttt{INST} (by design of the trace routine), not T₂'s
  \texttt{INST}. Hence T₂'s attempt to trace T₁ never actually runs any
  of T₁'s instructions; it merely manipulates T₂'s bookkeeping cells
  (\texttt{PREG₂}, \texttt{INST₂}, \texttt{AREG₂}, \ldots). This
  respects the routine's restriction that the traced program mustn't
  clobber the tracer's own working cells, because each tracer uses
  \emph{its own} \texttt{PREG}, \texttt{JREG}, \texttt{AREG}, etc. (the
  ``(a)--(c)'' restrictions listed under the exercises).
\item
  \textbf{No infinite regress, no mutual recursion of control.} When T₁
  encounters a \emph{jump} in T₂'s code, it rewrites the \textbf{copied}
  instruction so the jump targets \textbf{T₁}'s \texttt{JUMP} label;
  control therefore stays inside T₁, which then updates the simulated
  \texttt{J} and continues. T₂ never gains control flow; it remains
  merely the \emph{subject} being stepped.
\end{enumerate}

Net effect

\begin{itemize}
\tightlist
\item
  \textbf{T₁ simply traces T₂ as if T₂ were an ordinary program.} You
  get a trace of T₂'s instructions executing; T₂'s internal attempt to
  trace T₁ never actually executes T₁'s instructions.
\item
  \textbf{No paradox or endless slowdown chain} occurs here (contrast
  with the self-trace thought experiment of exercise 4, where
  self-modification can immediately derail the trace via \texttt{PREG}
  being clobbered).
\item
  If both copies tried to \emph{print} tracing information to the same
  device, that would violate restriction (6b) and garble output, but the
  skeleton routine shown doesn't output registers at all (that part was
  explicitly omitted).
\end{itemize}

\textbf{Conclusion.} With two copies placed in disjoint memory and each
set to trace the other, the one you start (T₁) retains control and
cleanly steps through the other (T₂). There's no infinite regress; it
behaves just like tracing any other (nontrivial) program, only slower
because that program happens to be a tracer itself.

\section{6 {[}40{]}}\label{section-177}

Here's a clean way to do \textbf{1.4.3.2 --- Exercise 6} (``write a
trace routine that is capable of tracing itself \ldots{} at slower and
slower speed, ad infinitum, until memory is exceeded''), built directly
on Knuth's trace routine in §1.4.3.2. I'll (1) recall how the given
trace routine works, (2) show the minimal \emph{self-bootstrapping}
change that makes it trace itself once (the idea of Ex. 4), then (3)
turn that into an \textbf{infinitely nestable, re-entrant} trace by
giving it tiny per-level ``frames,'' so multiple copies of the trace can
be active at once without stepping on each other's working cells, and
(4) point to where to splice in the ``print one step'' action. Citations
refer to the exact listing and the exercise page in the book PDF you
linked.

\begin{enumerate}
\def\labelenumi{\arabic{enumi})}
\tightlist
\item
  What the provided trace routine already does
\end{enumerate}

Knuth's routine is invoked when the \emph{driver} jumps to label
\texttt{ENTER} \textbf{with}:

\begin{itemize}
\tightlist
\item
  \texttt{rJ} = start address of the code to trace,
\item
  \texttt{rX} = address where tracing must stop (exclusive).
\end{itemize}

On entry it saves \texttt{rA} and \texttt{rJ}, loads the
\emph{simulated} program counter from \texttt{rJ}, fetches the next
instruction word, rewrites certain instructions (e.g., jumps) so the
monitor keeps control, executes the instruction via the pseudo-op line
\texttt{INST\ NOP\ *}, advances to the next location, and repeats;
special logic handles jumps and \texttt{JSJ} so the simulated \texttt{J}
is correct. Finally it restores \texttt{rA} and exits at
\texttt{LEAVE/LEAVEX}. (See lines 02--08 for setup, 10--16 decode,
27--47 jump handling, 34/36--38 execute-and-advance, 48--49 leave.)

Knuth explicitly notes this listing is \emph{just the control core} of a
trace: it doesn't print; printing/register-dump is left to exercises.

\begin{enumerate}
\def\labelenumi{\arabic{enumi})}
\setcounter{enumi}{1}
\tightlist
\item
  Make it trace itself once (the Ex. 4 idea you build on)
\end{enumerate}

Place a tiny prologue \emph{immediately before} label \texttt{ENTER}:

\begin{verbatim}
SELF   ENTX  LEAVEX      ; inner trace stops at our LEAVEX
       JSJ   ENTER       ; start a fresh trace instance, sets rJ to SELF+1
ENTER  ...               ; (the original routine follows)
\end{verbatim}

Why this works (and why we use \textbf{JSJ}, not \texttt{JMP}):

\begin{itemize}
\tightlist
\item
  When the outer monitor first ``executes'' \texttt{JSJ\ ENTER} as an
  instruction of the traced program, its jump logic recognizes a jump
  and treats \textbf{JSJ} specially: it does \textbf{not} overwrite the
  simulated \texttt{J} with a ``fall-through'' value (lines 39--47 check
  for JSJ and avoid updating \texttt{JREG} in that case), then sets the
  simulated PC to \texttt{ENTER}. Thus the \emph{inner} monitor instance
  starts with \texttt{rJ\ =\ SELF+1} and begins tracing \textbf{the
  trace program itself} starting at \texttt{ENTER}. From now on, the
  outer level prints every step the inner monitor takes, and the inner
  monitor will (once we add printing) print every step of its own traced
  program.
\end{itemize}

Knuth asks you to think about the simpler version with
\texttt{ENTX\ LEAVEX;\ JMP\ *+1} in Ex. 4; swapping the second
instruction to \texttt{JSJ\ ENTER} is the right subroutine-style call
that seeds the \emph{next} level with a proper \texttt{J}. The exercise
statement for 6 explicitly says ``in the sense of exercise 4,'' i.e.,
the same bootstrapping trick but continuing indefinitely.

If you stop here, you can get one active monitor tracing itself, but as
soon as that inner instance tries to start \emph{another} instance, all
copies will fight over the same scratch words (\texttt{AREG},
\texttt{JREG}, \texttt{PREG}, \texttt{INST}, \texttt{INST1},
\texttt{TEST}, \ldots). We need re-entrancy.

\begin{enumerate}
\def\labelenumi{\arabic{enumi})}
\setcounter{enumi}{2}
\tightlist
\item
  Make the trace \textbf{re-entrant} with per-level frames (the key to
  ``ad infinitum'')
\end{enumerate}

The original routine uses fixed global cells (labels like \texttt{AREG},
\texttt{JREG}, \texttt{PREG}, \texttt{INST}, \texttt{INST1},
\texttt{TEST}, \texttt{LEAVEX}). To support unlimited nesting, give each
active ``level'' its own copy of those words---i.e., a
\emph{frame}---and index all accesses via, say, \texttt{rI1} as a frame
pointer.

3.1 Define the frame layout (somewhere near the routine)

\begin{verbatim}
; One frame per nesting level (10 words is enough for the listed cells)
FRAME_SZ  EQU  10
FR_AREG   EQU  0     ; saved rA
FR_JREG   EQU  1     ; saved rJ (simulated)
FR_PREG   EQU  2     ; simulated PC (location being traced)
FR_INST   EQU  3
FR_INST1  EQU  4
FR_TEST   EQU  5
FR_LEAVEX EQU  6     ; we keep its absolute addr here for this level
; ... a couple spare words
\end{verbatim}

Reserve a small stack of frames (e.g., 100 levels):

\begin{verbatim}
FRAMES    ORIG *+100*FRAME_SZ
\end{verbatim}

Maintain \texttt{rI1} = base of \textbf{current} frame. On each entry,
bump it to the next frame; on each leave, drop it.

3.2 Patch the routine to index frame fields by \texttt{,1}

Example patches (showing the idea; you'd apply this uniformly):

\begin{itemize}
\item
  Save/restore \texttt{rA}:

  \begin{itemize}
  \tightlist
  \item
    old: \texttt{STA\ AREG} / \texttt{LDA\ AREG}
  \item
    new: \texttt{STA\ FR\_AREG,1} / \texttt{LDA\ FR\_AREG,1}
  \end{itemize}
\item
  Save \texttt{rJ} and use it later:

  \begin{itemize}
  \tightlist
  \item
    old: \texttt{STJ\ JREG} / \texttt{LDA\ JREG(0:2)} etc.
  \item
    new: \texttt{STJ\ FR\_JREG,1} / \texttt{LDA\ FR\_JREG,1(0:2)}
  \end{itemize}
\item
  Keep the ``next instruction address'':

  \begin{itemize}
  \tightlist
  \item
    old: \texttt{STA\ PREG(0:2)} / \texttt{LDA\ PREG(0:2)}
  \item
    new: \texttt{STA\ FR\_PREG,1(0:2)} / \texttt{LDA\ FR\_PREG,1(0:2)}
  \end{itemize}
\item
  \texttt{INST} and \texttt{INST1} similarly become \texttt{FR\_INST,1}
  and \texttt{FR\_INST1,1}.
\item
  The exit address:

  \begin{itemize}
  \tightlist
  \item
    On entry, copy \texttt{rX} into \texttt{FR\_LEAVEX,1(0:2)} and test
    against that field instead of a single global \texttt{TEST/LEAVEX}.
  \end{itemize}
\end{itemize}

Concretely, the very first lines change to (comments show what moved):

\begin{verbatim}
SELF      ENTX  FR_LEAVEX,1(0:2)   ; stop at this level’s LEAVEX
          JSJ   ENTER

ENTER     ENTA  FRAME_SZ
          ENT2  FRAMES             ; rI2 = base of all frames
          ADD   LEVEL              ; LEVEL holds current depth (init 0)
          MUL   =1=                ; (cheap way to scale by FRAME_SZ if needed)
          ADD   rI2                ; compute next frame base
          ST1   LEVEL              ; bump depth counter
          LDX   =0=
          STX   FR_AREG,1          ; clear a few words if desired

          STX   FR_TEST,1(0:2)     ; put stop addr from rX
          STX   FR_LEAVEX,1(0:2)
          STA   FR_AREG,1
          STJ   FR_JREG,1
          LDA   FR_JREG,1(0:2)     ; start location for tracing
CYCLE     STA   FR_PREG,1(0:2)
...
TEST      DECA  *                   ; compare FR_PREG vs FR_LEAVEX for exit
          JAZ   LEAVE
PREG      LDA   *                   ; fetch instr at FR_PREG
          STA   FR_INST,1
...
\end{verbatim}

(You don't need to \emph{literally} copy/paste the above; the point is
every variable access acquires a \texttt{,1} and the constant cells
become \emph{offsets within the current frame}.)

3.3 Leave drops the frame and returns

At \texttt{LEAVE}:

\begin{verbatim}
LEAVE     LDA  FR_AREG,1
          ; rI1 := rI1 - FRAME_SZ (restore previous level’s frame)
          ENTA  -FRAME_SZ
          ADD   rI1
          ENT1  0,0                ; move pointer
          ; also decrement LEVEL (optional book-keeping)
          LEAVEX JMP  *            ; same “stop tracing”
\end{verbatim}

The only global state left is the code itself and the tiny depth/stack
base variables; all per-step scratch lives in that level's frame.

With the routine re-entrant, executing the two-instruction \texttt{SELF}
prologue \textbf{starts a fresh instance that uses its own frame}, while
the outer instance continues stepping and printing. When the inner
instance reaches its own \texttt{SELF}, it starts yet another one, and
so on---until your \texttt{FRAMES} area runs out, which is exactly the
``until memory capacity is exceeded'' stopping condition the exercise
asks for.

\begin{enumerate}
\def\labelenumi{\arabic{enumi})}
\setcounter{enumi}{3}
\tightlist
\item
  Where to ``print one step''
\end{enumerate}

Knuth's core doesn't print; he says you must add that (e.g., Ex. 2
writes a 10-word snapshot per step to tape 0). The natural place to emit
output for \emph{the current step} is \textbf{after} you've decoded the
instruction and before executing it (i.e., after lines 10--16 in the
listing), or immediately after execution (line 35). In a short solution,
printing just the location and the instruction word is fine:

\begin{itemize}
\tightlist
\item
  Just before \texttt{INST\ NOP\ *} (line 34), do:
\end{itemize}

\begin{verbatim}
          LDA  FR_PREG,1(0:2)
          STA  BUF                     ; word 1: location
          LDA  FR_INST,1
          STA  BUF+1                   ; word 2: instr word
          OUT  BUF(18)                 ; print one short line on device 18
\end{verbatim}

Replace \texttt{BUF} with a tiny two-word buffer you've put near the
routine. (If you already solved Ex. 2 and wrote the richer 10-word
record to tape unit 0, call that writer here instead.)

\begin{enumerate}
\def\labelenumi{\arabic{enumi})}
\setcounter{enumi}{4}
\tightlist
\item
  Why the speeds keep dropping
\end{enumerate}

\begin{itemize}
\tightlist
\item
  Level 0 (outermost) prints \textbf{every} step the program under trace
  performs.
\item
  When the under-trace \emph{program is itself a trace}, that inner
  trace prints too---so each inner step now costs ``execute the step'' +
  ``inner print'' + ``outer print''.
\item
  The inner trace's own under-trace is \emph{another} trace, etc. Each
  level adds another layer of output and control logic around a single
  MIX instruction; hence the geometric slowdown the problem statement
  describes.
\end{itemize}

\begin{enumerate}
\def\labelenumi{\arabic{enumi})}
\setcounter{enumi}{5}
\tightlist
\item
  Driver to start the cascade
\end{enumerate}

A tiny starter sets \texttt{rJ} to \texttt{SELF} of your trace and
\texttt{rX} to its \texttt{LEAVEX}, then jumps to \texttt{ENTER}. (In
MIX, you can seed \texttt{rJ} by entering via \texttt{JSJ\ SELF} and
arranging that \texttt{SELF} immediately ``calls'' \texttt{ENTER} as in
the prologue above, so the first inner instance starts correctly; your
outermost instance is started by your hand---e.g., from a small
harness.)

\begin{enumerate}
\def\labelenumi{\arabic{enumi})}
\setcounter{enumi}{6}
\tightlist
\item
  Safety/notes
\end{enumerate}

\begin{itemize}
\tightlist
\item
  Knuth lists three simple restrictions on the traced program (don't
  store into the trace's cells; don't use the output device the trace
  uses; expect slowdown). Using frames fixes the first; just be sure
  your traced program doesn't \texttt{OUT} to the same unit you're
  using.
\item
  His jump-handling already special-cases \texttt{JSJ}, which is crucial
  for the self-bootstrapping prologue to work.
\end{itemize}

\section{7 {[}25{]}}\label{section-178}

\begin{enumerate}
\def\labelenumi{\arabic{enumi})}
\tightlist
\item
  What the exercise asks
\end{enumerate}

Design a \textbf{jump trace} that emits far less output than a full
trace. Instead of dumping registers every step, it should output a
sequence of pairs \((x₁, y₁), (x₂, y₂), \dots\) meaning ``the program
jumped from location \(x_i\) to \(y_i\), then ran straight-line code
from \(y_i\) through \(x_{i+1}\), then jumped from \(x_{i+1}\) to
\(y_{i+1}\), \ldots{}''. From such pairs, a post-pass can reconstruct
control flow and count how often each instruction executed. 

\begin{enumerate}
\def\labelenumi{\arabic{enumi})}
\setcounter{enumi}{1}
\tightlist
\item
  Where to hook into the given trace routine
\end{enumerate}

Knuth's trace routine already detects jumps and centralizes their
handling at label \texttt{JUMP}. The fetch/decode loop keeps the
\textbf{fetch address} in \texttt{PREG(0:2)} (the ``location of the next
instruction''), and when a jump occurs execution transfers to
\texttt{JUMP}, where the \textbf{target address} (effective address of
the jump) is computed into register \textbf{A} just before control
returns to \texttt{CYCLE} (which then loads that A-value into
\texttt{PREG} for the next iteration). This is the perfect point to
record our pair \(x, y\). 

Concretely (from the text routine):

\begin{itemize}
\tightlist
\item
  \texttt{CYCLE} stores the fetch location into \texttt{PREG(0:2)},
  fetches the instruction at that address, and identifies jumps.
\item
  When a jump is seen, code branches to \texttt{JUMP}. There the routine
  finalizes jump semantics (special-cases \texttt{JSJ}, updates
  \texttt{JREG} when appropriate), then does \texttt{INST1\ ENTA\ *} so
  \textbf{A = target address}; finally it \texttt{JMP\ CYCLE} to
  continue from the jump target. 
\end{itemize}

Therefore, we can log:

\begin{itemize}
\tightlist
\item
  \(x\) = the address of the \textbf{jump instruction} =
  \texttt{PREG(0:2)} at the time the jump is recognized,
\item
  \(y\) = the \textbf{jump target} = the value placed into \textbf{A} at
  label \texttt{JUMP} just before \texttt{JMP\ CYCLE}. 
\end{itemize}

\begin{enumerate}
\def\labelenumi{\arabic{enumi})}
\setcounter{enumi}{2}
\tightlist
\item
  Minimal recording logic (code sketch)
\end{enumerate}

Goals: (a) \textbf{very low overhead} per jump, (b) \textbf{no register
clobbering} visible to the traced program, (c) \textbf{buffered} output
(pairs take two words), (d) output to a device \textbf{not used by the
program being traced} (e.g., tape unit 0). 

We'll modify the block at \texttt{JUMP} (after it has computed the
target in A), insert a small snippet to \textbf{append (x, y)} to a
memory buffer, and flush the buffer to tape when full.

\textbf{Data (reserved somewhere outside the traced program's memory):}

\begin{verbatim}
BUF     CON  0              * start of a 100-word block
...     (reserve 100 words total for one tape block)
END     EQU  BUF+100        * end of block
PTR     CON  BUF            * current write position (absolute address)
I1REG   CON  0              * saved I1
YHOLD   CON  0              * temp to hold y
ONE     CON  1
TWO     CON  2
\end{verbatim}

\textbf{At routine entry (ENTER):} save I1 once (we'll use I1 as an
internal pointer) and ensure tape 0 is idle.

\begin{verbatim}
ST1   I1REG
JBUS  *(0)                  * wait until tape 0 is not busy
STA   AREG                  * (this is already in Knuth’s code)
STJ   JREG                  * (already)
...
\end{verbatim}

(``Use minimal registers'' is an ideal; for practical, simple buffering
we momentarily use I1 and restore it on exit. The original trace avoided
using registers other than A; here we use I1 carefully and restore it. )

\textbf{At label \texttt{JUMP} (after lines that update \texttt{JREG}
and just after \texttt{INST1\ ENTA\ *} puts the target into A):}

\begin{verbatim}
// Existing code:
JUMP    LDA   8B(4:5)       * as in the book
        SUB   INST(4:5)     * detect JSJ, etc.
        JAZ   *+4
        LDA   PREG(0:2)
        INCA  1
        STA   JREG(0:2)
INST1   ENTA  *             * <= A now holds jump target "y"

// INSERT begin: record (x,y) to [PTR], [PTR+1]

        STA   YHOLD         * save y (target) temporarily
        LDA   PREG(0:2)     * load x = address of jump instr
        LD1   PTR           * I1 := current dest address
        STA   0,1           * store x at M[PTR]
        LDA   YHOLD         * restore y
        STA   1,1           * store y at M[PTR+1]
        ENT1  2,1           * I1 := I1 + 2
        ST1   PTR           * PTR := PTR + 2

        * If buffer full, flush to tape 0 and reset PTR:
        CMP1  =END=         * compare I1 with END
        J1N   TRACE_CONT    * not full yet
        * flush:
        OUT   BUF,0         * start writing 100 words to tape unit 0
        JBUS  *(0)          * wait completion (conservative)
        ENT1  BUF           * I1 = BUF
        ST1   PTR           * PTR := BUF
TRACE_CONT:

// INSERT end

        JMP   CYCLE         * existing: continue at new address
\end{verbatim}

Remarks:

\begin{itemize}
\tightlist
\item
  We \textbf{only} touch I1, and we save/restore it at the beginning/end
  of the whole trace; the traced program never observes a changed I1
  because the trace wrapper owns it during tracing. (If you want to be
  extra strict, you can re-save/re-restore I1 inside \texttt{JUMP} too,
  but it's unnecessary if the trace routine always retains control.) 
\item
  The snippet is constant-time per jump: a handful of \texttt{STA/LDA}
  plus one indexed \texttt{STA} pair and a simple two-word pointer bump.
\item
  On \textbf{exit} from tracing (\texttt{LEAVE/LEAVEX}), flush any
  \textbf{partial} block (if \texttt{PTR\ !=\ BUF}) and then
  \texttt{LD1\ I1REG} to restore I1, just like we restored A and J in
  the base routine. (Knuth already notes J needs restoring; exercise 1
  in this section has you fix that; do both.) 
\end{itemize}

\begin{enumerate}
\def\labelenumi{\arabic{enumi})}
\setcounter{enumi}{3}
\tightlist
\item
  Buffering/output strategy and why tape
\end{enumerate}

\begin{itemize}
\tightlist
\item
  \textbf{Why tape (unit 0)?} It's far faster and far less intrusive
  than printing lines; the book explicitly hints that tape is preferable
  to printing for trace output. 
\item
  \textbf{Asynchronous I/O:} We write a full 100-word block when the
  buffer fills (50 jump pairs). In typical code, the number of jumps is
  far less than the number of executed instructions, so output volume
  and I/O frequency are much smaller than with a full step trace.
\item
  \textbf{Exclusivity:} The traced program must not use the same device;
  Knuth calls this out (don't use the output device being used for the
  trace) --- so dedicate tape 0 to the jump trace, and choose another
  device for the program (or vice versa). 
\end{itemize}

\begin{enumerate}
\def\labelenumi{\arabic{enumi})}
\setcounter{enumi}{4}
\tightlist
\item
  Notes, correctness, and limitations
\end{enumerate}

\begin{itemize}
\item
  \textbf{What counts as a ``jump''?} The existing trace detects all
  jump opcodes (\texttt{JMP}, \texttt{Jx}, \texttt{Jxx}, \texttt{JSJ})
  and also treats \texttt{JBUS} specially (since it can effect control
  transfer). We piggyback on that detection; our code runs \textbf{only}
  when a jump actually occurs. 
\item
  \textbf{Source and target correctness:} At jump time,
  \texttt{PREG(0:2)} is the address of the jump instruction being
  executed (our \(x\)); after \texttt{INST1\ ENTA\ *}, A holds the
  effective destination (our \(y\)). That aligns exactly with the
  exercise's notion of pairs. 
\item
  \textbf{Post-pass reconstruction:} A reader can process the pair
  stream \texttt{(x₁,y₁),\ (x₂,y₂),\ …} and:

  \begin{itemize}
  \tightlist
  \item
    mark \textbf{basic blocks} as \([y_i, x_{i+1}]\),
  \item
    accumulate execution counts by adding the length of each
    straight-line interval,
  \item
    reconstruct the \textbf{control-flow graph} edges from each
    \(x_i \to y_i\).
  \end{itemize}
\item
  \textbf{Self-modifying code caveat:} The book notes that
  self-modifying code can defeat a jump-trace profiler; e.g., if an
  instruction changes another instruction into or out of a jump, the
  pair stream may no longer reflect the original static flow. Keep this
  limitation in mind. 
\item
  \textbf{Register hygiene:} We intentionally restrict ourselves to A
  and a temporary use of I1 (saved/restored), in keeping with the trace
  routine's philosophy of not disturbing the program state. 
\end{itemize}

By inserting \textasciitilde15 tight instructions at \texttt{JUMP},
storing \((x,y)=(\text{PREG(0:2)},\text{A})\) to a buffered block and
flushing to tape, you get an \textbf{efficient jump trace} with tiny
runtime overhead and dramatically reduced output size, while preserving
the traced program's state. This meets the exercise's specification and
leverages the exact structure of the trace routine given in §1.4.3.2. 

\bookmarksetup{startatroot}

\chapter{1.4.4 Input and Output}\label{input-and-output}

\section{1 {[}05{]}}\label{section-179}

Context you need

On MIX, \texttt{IN} \textbf{initiates} an input transfer and returns
immediately; the program may continue executing, but it must not rely on
the input buffer's contents until the device is no longer busy (e.g.,
until \texttt{JBUS} shows ``ready''). Sequence (3) relies on that rule:
it waits with \texttt{JBUS}, then copies from the just-filled buffer,
and finally starts the next input into that same buffer. 

\begin{enumerate}
\def\labelenumi{(\alph{enumi})}
\tightlist
\item
  Putting \texttt{MOVE}s \textbf{before} \texttt{JBUS}
\end{enumerate}

\textbf{Incorrect.} If you copy from locations \texttt{2000–2099}
\emph{before} \texttt{JBUS} has confirmed completion, you may be moving
words that haven't arrived yet (or are only partially written). The
model explicitly warns that input completes only after a subsequent I/O
op or a successful ``not busy'' test; otherwise the buffer contents are
not guaranteed. So the copied block could be incomplete/corrupted. 

\textbf{Verdict:} \textbf{No.} The input may not yet be complete.

\begin{enumerate}
\def\labelenumi{(\alph{enumi})}
\setcounter{enumi}{1}
\tightlist
\item
  Putting \texttt{MOVE}s \textbf{after} \texttt{IN}
\end{enumerate}

\textbf{Also incorrect (race hazard).} If you start the \emph{next}
\texttt{IN\ 2000(5)} first and only then perform
\texttt{MOVE\ 2000(50);\ MOVE\ 2050(50)}, you are copying from a buffer
that the device is simultaneously \textbf{overwriting} with the next
block. If the input device outpaces the CPU even slightly, some words
you copy will belong to the \emph{next} block rather than the previous
one, yielding a mixed/corrupted result. 

\textbf{Verdict:} \textbf{No.} The device might outrun the
\texttt{MOVE}, so the copy can be contaminated by the new transfer.

Final answers

\textbf{(a) No.} You might move data before the current input has
finished. \textbf{(b) No.} Starting \texttt{IN} first risks the device
overwriting the buffer while you copy, so the \texttt{MOVE} can pick up
the wrong data. 

\section{2 {[}10{]}}\label{section-180}

Here's a buffered-\textbf{output} scheme that mirrors Knuth's
buffered-\textbf{input} setup (his items (2) and (3)) but for writing
tape unit 6. We'll stage each 100-word block for output in an auxiliary
buffer at addresses 2000--2099, while your program prepares the next
100-word block at 1000--1099. This overlaps CPU work with device I/O,
just as his input pipeline does. The exercise statement appears at the
end of §1.4.4; Knuth's earlier unbuffered pair for output is
\texttt{OUT\ 1000(6);\ JBUS\ *(6)} and his input buffering sequences
(2)--(3) are the model we're copying. 

Setup and notation (MIX reminders)

\begin{itemize}
\tightlist
\item
  A tape \textbf{OUT A(6)} starts writing 100 words beginning at A to
  unit 6.
\item
  \textbf{JBUS *(6)} waits until unit 6 is no longer busy.
\item
  \textbf{MOVE A(50)} copies 50 words starting at A into the destination
  pointed to by \textbf{rI1}, then advances rI1 by 50 (so two MOVE's
  copy 100 words). Knuth's input example sets rI1 to the destination
  just before the two MOVE's. 
\end{itemize}

Buffered output: two phases

\begin{enumerate}
\def\labelenumi{\arabic{enumi})}
\tightlist
\item
  Prime the pipeline (first block)
\end{enumerate}

Before the device can ``trail'' your computation, you must place the
first block into the staging buffer and start the first write:

\begin{Shaded}
\begin{Highlighting}[]
\NormalTok{ENT1 2000          ; rI1 ← 2000 (destination = buffer)}
\NormalTok{MOVE 1000(50)      ; (1000–1049) → (2000–2049)}
\NormalTok{MOVE 1050(50)      ; (1050–1099) → (2050–2099)}
\NormalTok{OUT  2000(6)       ; start writing the staged block to tape unit 6}
\end{Highlighting}
\end{Shaded}

Explanation: you've copied the current block (prepared at 1000--1099)
into 2000--2099 and kicked off the OUT. While the device is busy writing
that buffer, your program is now free to compute/fill the \textbf{next}
100-word block at 1000--1099. (This mirrors Knuth's ``read-ahead''
input, but here it's ``write-behind'': the device is finishing the
\textbf{previous} buffer while the CPU prepares the next one.) 

\begin{enumerate}
\def\labelenumi{\arabic{enumi})}
\setcounter{enumi}{1}
\tightlist
\item
  Steady-state cycle (each subsequent block)
\end{enumerate}

Whenever the next 100-word block at 1000--1099 is ready to be sent:

\begin{Shaded}
\begin{Highlighting}[]
\NormalTok{ENT1 2000          ; rI1 ← 2000 (dest = buffer), harmless if repeated}
\NormalTok{JBUS *(6)          ; wait until the prior OUT from 2000–2099 has finished}
\NormalTok{MOVE 1000(50)      ; (1000–1049) → (2000–2049)}
\NormalTok{MOVE 1050(50)      ; (1050–1099) → (2050–2099)}
\NormalTok{OUT  2000(6)       ; start writing this staged block}
\end{Highlighting}
\end{Shaded}

Why this order? You \textbf{must} not overwrite 2000--2099 while the
device is still taking data from it; therefore the \textbf{JBUS}
precedes the \textbf{MOVE} instructions, just as in Knuth's input
sequence (3) the JBUS precedes the MOVE to ensure the fresh input has
arrived before moving it out of the staging area. This is the exact
analog of his (2)--(3), but with roles reversed. 

Finishing up

After you've initiated the \textbf{final} \texttt{OUT\ 2000(6)}, wait
for the device to complete before you stop or reuse the buffer:

\begin{Shaded}
\begin{Highlighting}[]
\NormalTok{JBUS *(6)          ; ensure the last block has been fully written}
\end{Highlighting}
\end{Shaded}

(Stopping early risks data loss; the same caveat applies to buffered
input and is discussed around Knuth's examples.) 

Why it works

\begin{itemize}
\tightlist
\item
  \textbf{Safety.} The \texttt{JBUS\ *(6)} before each pair of
  \texttt{MOVE}s prevents clobbering the very buffer the device is still
  using, exactly like waiting for input completion before moving data
  out of the input buffer in (3).
\item
  \textbf{Overlap.} While the device streams the \textbf{previous} block
  from 2000--2099, your program is free to prepare the \textbf{next}
  block at 1000--1099, so computation and I/O overlap. This is the
  write-side counterpart of Knuth's ``buffering'' idea that overlaps I/O
  with computation.
\end{itemize}

\section{3 {[}22{]}}\label{section-181}

Design notes

\begin{itemize}
\tightlist
\item
  Two 100-word buffers live in memory, each followed by a single
  ``link'' word that holds the address of the other buffer. This is
  exactly analogous to the input routine in the text (which links its
  two buffers in lines 07, 11, 14 of (4)).
\item
  We keep the \textbf{base address} of the current output buffer in
  \texttt{CUR}.
\item
  Index register \textbf{I5} holds the \textbf{current position} in the
  buffer as a \emph{negative count}: it starts at −100 and is
  incremented each word; when it hits 0 the buffer is full. This lets us
  test fullness with a single \texttt{J5N} (jump if negative).
\item
  When a buffer fills, we self-modify the address field of a single
  \texttt{OUT} instruction to the ``old'' buffer's base, execute it,
  then swap \texttt{CUR} to the other buffer and reset \texttt{I5} to
  −100. (This is the same self-modifying trick used for subroutine
  linkage and in many MIX examples; see the standard
  \texttt{STJ\ EXIT\ …\ EXIT\ JMP\ *} pattern. )
\item
  No \texttt{JBUS} is needed \emph{during steady-state operation} (same
  reason as for input (4)): we start the next I/O before we touch the
  previous buffer again.
\item
  At program end we \textbf{pad with zeros} if needed, \texttt{OUT} the
  final partial block, and then do a single \texttt{JBUS\ *(V)} to wait
  for the last I/O to complete.
\end{itemize}

\begin{Shaded}
\begin{Highlighting}[]
\NormalTok{* ====== Constants / device selection ======}
\NormalTok{V     EQU 1              * Tape unit to use (pick any valid unit id)}
\NormalTok{ZERO  EQU 0}

\NormalTok{* ====== Buffer areas (100 words each) + one link word ======}
\NormalTok{OUTBUF1 ORIG *+100       * 100 data words}
\NormalTok{LINK1   CON OUTBUF2      * address of the other buffer}

\NormalTok{OUTBUF2 ORIG *+100}
\NormalTok{LINK2   CON OUTBUF1}

\NormalTok{* ====== State ======}
\NormalTok{CUR    CON OUTBUF1       * base address of current buffer}

\NormalTok{* ====== Subroutine: WORDOUT =================================}
\NormalTok{* Calling sequence:  JMP WORDOUT}
\NormalTok{* Entry:  rA = next word to output}
\NormalTok{* Uses:   rI5 = position in current buffer (negative count), CUR = base}
\NormalTok{* Exit:   word stored; if buffer became full, a 100{-}word OUT is started}
\NormalTok{* Clobbers: rA (temporarily), rJ (by standard linkage), CI}
\NormalTok{*}
\NormalTok{WORDOUT STJ 9F                 * standard MIX subroutine linkage}
\NormalTok{        STA CUR+100,5          * store rA at CUR + (100 + rI5)  [slots 0..99]}
\NormalTok{        INC5 1                 * advance position}
\NormalTok{        J5N  9F                * if still negative, buffer not full =\textgreater{} return}

\NormalTok{* {-}{-} Buffer just filled: start OUT on the full buffer, swap to the other}
\NormalTok{*    Save OLD base into OUTCMD\textquotesingle{}s address field, then switch CUR to link.}
\NormalTok{        LDA CUR}
\NormalTok{        STA OUTCMD+1(0:2)      * self{-}modify OUT address field}
\NormalTok{        LDA CUR+100            * fetch link (address of other buffer)}
\NormalTok{        STA CUR                * CUR := other buffer base}
\NormalTok{        ENT5 {-}100              * reset position for the new buffer}

\NormalTok{OUTCMD  OUT OUTBUF1(V)         * OUT 100 words starting at (modified) address}
\NormalTok{9H      JMP *                  * return to caller}
\NormalTok{* ============================================================}

\NormalTok{* ====== Subroutine: FLUSHOUT ================================}
\NormalTok{* Pads current (possibly partial) buffer with zeros, OUTs it, and waits}
\NormalTok{* for the final I/O to finish. Safe to call even if buffer is exactly empty/full.}
\NormalTok{*}
\NormalTok{FLUSHOUT STJ 9F}

\NormalTok{        J5Z  FINISHWAIT        * rI5==0 =\textgreater{} a full block was just started; just wait}
\NormalTok{* Pad with zeros until rI5 hits 0, then OUT the partial block.}
\NormalTok{PADZ    STZ CUR+100,5}
\NormalTok{        INC5 1}
\NormalTok{        J5N PADZ}

\NormalTok{        LDA CUR}
\NormalTok{        STA OUTCMD+1(0:2)      * OUT the now{-}full padded buffer}
\NormalTok{        OUTCMD OUT OUTBUF1(V)}

\NormalTok{FINISHWAIT}
\NormalTok{        JBUS *(V)              * wait until last OUT completes}
\NormalTok{        JMP 9F}
\NormalTok{* ============================================================}

\NormalTok{* ====== Program start (minimum required initialization) =====}
\NormalTok{*   {-} Choose initial buffer and set position to start of buffer.}
\NormalTok{*   {-} (Buffers’ link words are assembled above; no runtime init needed.)}
\NormalTok{START   ENT5 {-}100              * I5 = {-}100  (next store goes to slot 0)}
\NormalTok{        LDA =OUTBUF1=}
\NormalTok{        STA CUR                * CUR := OUTBUF1}
\NormalTok{* ... your program computes, and whenever it has a word to output:}
\NormalTok{*     LDA \textless{}word\textgreater{}}
\NormalTok{*     JMP WORDOUT}
\NormalTok{* ... repeat as needed ...}

\NormalTok{* ====== Program end (must flush final block) =====}
\NormalTok{DONE    JMP FLUSHOUT}
\NormalTok{        HLT}
\NormalTok{        END START}
\end{Highlighting}
\end{Shaded}

Why this works

\begin{itemize}
\item
  \textbf{``Buffer-swapping \ldots{} analogous to (4)''}: Two buffers
  linked to each other; when one fills we \texttt{OUT} it and
  immediately continue filling the other, with no busy-wait in the
  steady state.
\item
  \textbf{``Index register 5 \ldots{} current buffer position''}:
  \texttt{rI5} is the running position (negative count). We store at
  \texttt{CUR+100,5} and detect a full buffer with a single
  \texttt{J5N}/\texttt{rI5==0} test.
\item
  \textbf{``Write 100 words onto tape unit V''}: The \texttt{OUT}
  instruction is self-modified to the \emph{old} buffer base so we write
  exactly that 100-word block on device \texttt{V}.
\item
  \textbf{``Layout of buffer areas''}: As shown---each 100-word area
  followed by a 1-word \textbf{link} holding the other buffer's base
  (same linking idea as the text's WORDIN).
\item
  \textbf{``Beginning and end \ldots{} first and last blocks properly
  written''}:

  \begin{itemize}
  \tightlist
  \item
    \textbf{Beginning}: set \texttt{CUR} to your first buffer and
    \texttt{rI5\ :=\ -100}. That's all.
  \item
    \textbf{End}: call \texttt{FLUSHOUT}. It pads with zeros to the next
    100-word boundary (if needed), \texttt{OUT}s the final block, then
    does a single \texttt{JBUS\ *(V)} to wait for completion before
    halting. (Padding is explicitly required.)
  \end{itemize}
\end{itemize}

\section{4 {[}M20{]}}\label{m20-34}

Consider processing \(B\) blocks. For block \(k\) let

\begin{itemize}
\tightlist
\item
  \(c_k\) = CPU time spent computing on that block's data,
\item
  \(i_k\) = device time for the I/O of that block (read or write).
\end{itemize}

Define totals

\[
C=\sum_{k=1}^B c_k,\qquad I=\sum_{k=1}^B i_k.
\]

With \textbf{unbuffered I/O} the CPU waits during each I/O, so the
elapsed time is

\[
T_{\text{unbuf}}=C+I.
\]

With \textbf{buffered I/O} (double buffering, for example), computing on
one buffer can overlap I/O on the other. There is only \textbf{one} CPU
and \textbf{one} device, so at any instant at most one CPU task and at
most one device task can be active.

Fundamental lower bound (why speedup ≤ 2)

In \emph{any} schedule (buffered or not):

\begin{itemize}
\tightlist
\item
  You cannot finish earlier than \(C\), because the CPU must perform
  \(C\) units of work.
\item
  You cannot finish earlier than \(I\), because the device must perform
  \(I\) units of I/O.
\end{itemize}

Hence every buffered schedule satisfies

\[
T_{\text{buf}} \;\ge\; \max\{C,I\}.
\]

Comparing to the unbuffered time \(C+I\),

\[
\frac{T_{\text{unbuf}}}{T_{\text{buf}}}
\;\le\;
\frac{C+I}{\max\{C,I\}}
\;\le\; 2,
\]

with equality only when \(C=I\) \emph{and} overlap is essentially
perfect. Therefore \textbf{no single-device buffering scheme can yield
more than a 2× speedup} relative to unbuffered execution.

Constructive schedule (why \textasciitilde2× is attainable)

Assume steady per-block times \(c_k=c\) and \(i_k=i\) and use two
buffers:

\begin{itemize}
\tightlist
\item
  Start I/O for block 1 (takes \(i\)).
\item
  While that I/O runs, compute block 0's results or prepare block 2,
  etc. After the pipeline fills, each additional block takes
  \(\max\{c,i\}\) time because CPU and device run in parallel.
\item
  The exact makespan is
\end{itemize}

\[
T_{\text{buf}} = i + (B-1)\max\{c,i\} + c,
\]

(first and last blocks incur the pipeline fill/drain). For large \(B\),

\[
T_{\text{buf}} \approx B\max\{c,i\}.
\]

Thus the asymptotic speedup is

\[
\frac{T_{\text{unbuf}}}{T_{\text{buf}}}
\approx
\frac{B(c+i)}{B\max\{c,i\}}
=
\frac{c+i}{\max\{c,i\}}
\le 2.
\]

\textbf{Best case.} If \(c=i\) (balanced compute and I/O), then

\[
T_{\text{unbuf}}=2Bc,\qquad
T_{\text{buf}}\approx Bc \quad(\text{plus small }O(c)\text{ edge effects}),
\]

so the speedup approaches \(2\) as \(B\) grows. This is the ``favorable
circumstance'' where buffering nearly halves total time.

\textbf{Other cases.} If \(c\ll i\) (I/O-bound) or \(i\ll c\)
(compute-bound), then \((c+i)/\max\{c,i\}\) is close to \(1\); buffering
still helps by overlapping the smaller resource with the larger, but the
improvement is strictly \(<2\) and can be modest.

\section{5 {[}M21{]}}\label{m21-18}

Let

\begin{itemize}
\tightlist
\item
  \(C\): total time the program spends computing (CPU time) not counting
  any I/O wait.
\item
  \(D_i\): total active time needed by I/O device \(i\) (sum of the
  times of all blocks handled by device \(i\)), for \(i=1,\dots,n\).
\end{itemize}

Assumptions are exactly those of §1.4.4: Each device works independently
once an I/O is initiated; with buffering, the CPU can run while a device
is busy; multiple devices may be busy concurrently. (This is the same
setting used for MIX's buffered I/O with concurrent devices.) 

Two totals to compare:

\begin{itemize}
\item
  \textbf{Unbuffered baseline.} I/O is blocking; devices don't overlap
  with the CPU or each other.

  \[
  T_{\text{unbuf}}=C+\sum_{i=1}^n D_i.
  \]
\item
  \textbf{Buffered/overlapped.} With sufficient buffering, CPU and
  devices can run concurrently; multiple devices can be active at once.
  Any schedule must last at least as long as the heaviest-used resource
  (CPU or any one device), hence

  \[
  T_{\text{buf}} \;\ge\; \max\!\bigl(C,\,D_1,\dots,D_n\bigr).
  \]
\end{itemize}

(You can view this as a simple ``makespan'' lower bound: each of the
\(n{+}1\) resources---CPU and the \(n\) devices---must process its
assigned load; the overall time cannot be less than the largest load.)

General bound on speedup

Define the speedup \(S\) from buffering as

\[
S \;=\; \frac{T_{\text{unbuf}}}{T_{\text{buf}}}
\;\le\;
\frac{C+\sum_i D_i}{\max(C,\max_i D_i)}.
\tag{★}
\]

Since the maximum of \(n{+}1\) nonnegative numbers is at least their
average,

\[
\max(C,D_1,\dots,D_n)\;\ge\;\frac{C+\sum_i D_i}{n+1},
\]

and plugging into (★) gives the universal bound

\[
S \;\le\; n+1.
\]

This \textbf{generalizes Exercise 4's ``at most 2×'' for one device}
(\(n=1\Rightarrow S\le 2\)). 

Tightness (when can we reach \(n{+}1\)?)

The bound is tight in ``balanced'' cases where

\[
C=D_1=\cdots=D_n=T.
\]

Then \(T_{\text{unbuf}}=(n+1)T\) and the lower bound for any overlapped
schedule is \(T\). With well-designed multiple buffering (as in §1.4.4's
coroutine-based control) you can keep \textbf{all} \(n\) devices busy
continuously while the CPU also computes for the full \(T\). Thus

\[
T_{\text{buf}}=T,\qquad S=\frac{(n+1)T}{T}=n+1.
\]

This is the natural multi-device analogue of the single-device ``at most
two, achieved when compute time equals I/O time'' result. 

Intuition and scheduling sketch

\begin{itemize}
\item
  \textbf{Lower bound:} You cannot finish earlier than the longest of
  \(C,D_1,\dots,D_n\).
\item
  \textbf{Achieving near the bound:} Use multiple buffers per device and
  the coroutine-style controller (§1.4.4) so that:

  \begin{enumerate}
  \def\labelenumi{\arabic{enumi}.}
  \tightlist
  \item
    Each device is kept busy end-to-end (initiate the next transfer as
    soon as the previous completes), and
  \item
    The CPU has work to do at all times (no waiting on I/O), by
    computing against already-filled input buffers or producing output
    into free buffers. This is exactly the kind of overlap Knuth's
    buffer-control algorithms are designed for, and the time--action
    charts in the section illustrate how the overlap hides I/O behind
    computation (and vice versa). 
  \end{enumerate}
\end{itemize}

What the formula tells you in general

For an arbitrary workload, the \textbf{best possible} (ideal) speedup is
bounded by

\[
S\;\le\;\frac{C+\sum_i D_i}{\max(C,\max_i D_i)}.
\]

\begin{itemize}
\tightlist
\item
  If one device dominates (say \(D_k\gg C\) and \(D_k\gg D_j\) for
  \(j\neq k\)), then \(S\approx (C+\sum_i D_i)/D_k\), which is much
  smaller than \(n{+}1\).
\item
  If CPU dominates (\(C\gg D_i\)), speedup is near 1 (little to hide).
\item
  If loads are well balanced (\(C\approx D_i\)), you approach the ideal
  \(n{+}1\).
\end{itemize}

This fully generalizes the single-device ``≤2× speedup'' to \(n\)
devices.

\section{6 {[}12{]}}\label{section-182}

You have two symmetric choices (pick either):

\begin{verbatim}
IN   INBUF1(U)
ENT6 INBUF2+99
\end{verbatim}

or

\begin{verbatim}
IN   INBUF2(U)
ENT6 INBUF1+99
\end{verbatim}

Optionally precede with \texttt{IOC\ 0(U)} if you might need to
rewind/ready the unit. 

Why this is the right start (step-by-step)

Recall the subroutine (4) (WORDIN) uses rI6 as a pointer and begins
with:

\begin{enumerate}
\def\labelenumi{\arabic{enumi}.}
\tightlist
\item
  \texttt{INC6\ 1}
\item
  \texttt{LDA\ 0,6} and compare to \texttt{SENTINEL} If sentinel →
  refill this buffer: \texttt{IN\ -100,6(U)}, then switch buffers:
  \texttt{LD6\ 1,6} (loads rI6 from the ``other-buffer'' link), and loop
  to fetch the word. 
\end{enumerate}

Important details from that listing:

\begin{itemize}
\tightlist
\item
  Each buffer reserves 100 data words; at offset +100 sits a sentinel
  constant; at +101 sits a link to the other buffer.
\item
  rI6 ``addresses the last word of input'' when you enter WORDIN; the
  subroutine immediately does \texttt{INC6\ 1}, then fetches. 
\end{itemize}

Now follow the first call if you used the first init
(\texttt{IN\ INBUF1(U);\ ENT6\ INBUF2+99}):

\begin{enumerate}
\def\labelenumi{\arabic{enumi}.}
\item
  Before the first call:

  \begin{itemize}
  \tightlist
  \item
    Buffer 1 has just been filled by \texttt{IN\ INBUF1(U)} (block \#1).
  \item
    rI6 = \texttt{INBUF2+99} (the last data word slot of Buffer 2).
  \end{itemize}
\item
  On the first WORDIN call:

  \begin{itemize}
  \tightlist
  \item
    \texttt{INC6\ 1} ⇒ rI6 becomes \texttt{INBUF2+100} (the sentinel
    slot of Buffer 2).
  \item
    \texttt{LDA\ 0,6} ⇒ you read the sentinel, so WORDIN concludes
    ``this buffer is exhausted.''
  \item
    \texttt{IN\ -100,6(U)} ⇒ reads the next block (\#2) into Buffer 2.
  \item
    \texttt{LD6\ 1,6} ⇒ rI6 ← contents at \texttt{INBUF2+101}, which is
    the address of the other buffer (i.e., \texttt{INBUF1}).
  \item
    Loop back to fetch the word: \texttt{LDA\ 0,6} now loads the first
    actual data word of Buffer 1 --- exactly what you want for the very
    first word delivered to the program --- while the tape is already
    busy reading the subsequent block into Buffer 2.
  \end{itemize}
\end{enumerate}

Net effect: the very first call returns word 0 from Buffer 1, and you're
already overlapped with I/O that fills Buffer 2. From that point on,
WORDIN alternates smoothly between the buffers, always keeping the
device busy when possible (the ``pipeline is primed''). This is exactly
the behavior intended by (4).

What you must have in the data area

Your buffers must be laid out exactly as in (4):

\begin{itemize}
\tightlist
\item
  \texttt{INBUF1} and \texttt{INBUF2} each reserve 100 data words.
\item
  At \texttt{INBUF1+100} and \texttt{INBUF2+100} place
  \texttt{SENTINEL}.
\item
  At \texttt{INBUF1+101} put the address of \texttt{INBUF2}, and at
  \texttt{INBUF2+101} put the address of \texttt{INBUF1} (the links).
\end{itemize}

\section{7 {[}22{]}}\label{section-183}

Instead of testing for a sentinel to detect ``end of buffer'', keep a
running count. A neat MIX trick is to keep rI6 negative and let it run
from −100, −99, \ldots, −1 inside each block:

\begin{itemize}
\item
  Addressing: load the next word from BASE+100+rI6. With rI6 negative,
  this indexes positions 0..99 of the current buffer.
\item
  After each fetch, increment rI6. If it's still negative, we're still
  in the block; if it becomes 0, we've just consumed the last word, so
  we:

  \begin{enumerate}
  \def\labelenumi{\arabic{enumi}.}
  \tightlist
  \item
    start an IN on the buffer we just finished (to read the next block
    into it),
  \item
    reset rI6 to −100, and
  \item
    switch to the other buffer.
  \end{enumerate}
\end{itemize}

Using a tiny coroutine hand-off lets the subroutine ``enter'' at the
correct buffer each call without extra state variables, which is exactly
the style Knuth points to as one good way to solve this exercise.

MIX program (sentinel-free WORDIN)

\begin{verbatim}
* Sentinel-free WORDIN: double-buffered, coroutine style.
* U = the input device number (e.g., tape 5). rI6 is our negative index counter.

INBUF1     ORIG *+100           * 100-word buffer #1
INBUF2     ORIG *+100           * 100-word buffer #2

WORDIN     STJ  NEXT            * Patch NEXT with caller's return address.
NEXT       JMP  USE2            * First entry jumps into buffer #2’s path (can choose either).

* --- Use buffer #2 until it’s exhausted, then refill #2 and switch to #1 ---
USE2       LDA  INBUF2+100,6    * A ← INBUF2[100 + rI6]  (rI6 is negative ⇒ 0..99)
           JMP  RET             * Return this word to caller.
           INC6 1               * Advance position in the current block.
           J6N  USE2            * Still negative? More words in this block → stay.
           IN   INBUF2(U)       * rI6 just hit 0 ⇒ start reading next block into buffer #2.
           ENN6 100             * rI6 ← −100  (start of next block)
           JMP  USE1            * Switch to buffer #1 for subsequent words.

* --- Use buffer #1 until it’s exhausted, then refill #1 and switch to #2 ---
USE1       LDA  INBUF1+100,6
           JMP  RET
           INC6 1
           J6N  USE1
           IN   INBUF1(U)
           ENN6 100
           JMP  USE2

* --- Return site. STJ above patched NEXT to point back to the proper entry on the next call. ---
RET        STJ  NEXT            * Store rJ into the address field of NEXT (classic MIX return trick).
           JMP  *               * Jump to that saved address (i.e., return to caller).
\end{verbatim}

How to bootstrap (same spirit as Ex. 6)

Before the first call to \texttt{WORDIN}, do the usual startup so the
first fetch is correct and I/O is overlapped:

\begin{itemize}
\tightlist
\item
  Initialize the counter: \texttt{ENN6\ 100} (so rI6 = −100).
\item
  Prefetch one buffer (either is fine), e.g.: \texttt{IN\ INBUF2(U)}.
\item
  Leave \texttt{NEXT} as shown (\texttt{JMP\ USE2}) so the first call
  uses buffer \#2's path. (You could swap the roles of 1/2 if you
  prefer.) These are exactly the kind of ``get it started right'' steps
  noted next to program (4).
\end{itemize}

Why it works (and why it's fast)

\begin{itemize}
\tightlist
\item
  Correctness. Each call returns exactly one word (from the current
  buffer at offset 100+rI6). When rI6 reaches 0, we (i) initiate an
  input on the buffer we just finished, (ii) reset rI6 to −100, and
  (iii) continue from the other buffer. That reproduces the overlapped
  ``read-ahead'' behavior of (4), but without testing a sentinel; the
  sign of rI6 is our end-of-buffer test. The coroutine hand-off
  (\texttt{STJ\ NEXT} + \texttt{JMP\ USE1/USE2}) selects the correct
  entry for the next call, so the caller remains oblivious.
\item
  Cost. As in the sentinel version, each word costs one indexed load
  plus a couple of very cheap instructions, and at block boundaries we
  pay for one \texttt{IN} and one \texttt{ENN6}. With read-ahead active,
  the per-block time is \texttt{max(C,\ T)} (C = compute per block, T =
  device time), just like the buffered methods in the text.
\end{itemize}

Notes

\begin{itemize}
\tightlist
\item
  The routine assumes full 100-word blocks from the device, as in (4).
  If your data format has a separate notion of logical end-of-input
  (e.g., a trailing card with a special mark), handle that in the
  surrounding control coroutine/program exactly as in the section's
  examples.
\item
  If you prefer not to use coroutines, you can keep a ``current-buffer
  base'' pointer in memory and branch between two short ``use buffer''
  paths with an explicit test on rI6; the coroutine version above simply
  removes that extra pointer by letting the entry point itself encode
  which buffer is current. (Knuth's official answer illustrates the
  coroutine approach.)
\end{itemize}

\section{8 {[}11{]}}\label{section-184}

What the colors mean for \textbf{output}

\begin{itemize}
\tightlist
\item
  \textbf{Green (G)} = a free buffer that can be \textbf{ASSIGN}ed and
  filled with the next line image.
\item
  \textbf{Yellow (Y)} = the \textbf{current} buffer the program is
  building between \textbf{ASSIGN} and \textbf{RELEASE}.
\item
  \textbf{Red (R)} = a buffer that has been \textbf{RELEASE}d and now
  contains a complete line image, ready for the printer (or actively
  being printed if it's the one pointed to by \textbf{NEXTR}). 
\end{itemize}

The pointers are as in the text: \textbf{NEXTG} → ``next green'' (where
the next ASSIGN will hand you a free buffer), \textbf{NEXTR} → ``next
red'' (what the output device will print next), \textbf{CURRENT} → the
yellow (assigned) buffer, if any. 

Interpreting the figure sequence for \textbf{output}

State shown in Fig. 23 (before we start the three snapshots in Fig. 24)

At that instant there are several green buffers (free space is
available) and \textbf{two red buffers}; the \textbf{red} buffer
indicated by \textbf{NEXTR} is \textbf{currently being printed} by the
line printer. At the same time, the program is \textbf{computing between
RELEASE and ASSIGN}; i.e., it has finished one line (so that buffer
turned red) and hasn't yet requested a fresh buffer for the next line. 

\begin{enumerate}
\def\labelenumi{(\alph{enumi})}
\tightlist
\item
  After \textbf{ASSIGN}
\end{enumerate}

The program calls \textbf{ASSIGN}. The buffer pointed to by
\textbf{NEXTG} turns \textbf{yellow} and becomes \textbf{CURRENT};
\textbf{NEXTG} advances clockwise to the next green. The program now
begins to \textbf{compose the next line image} into this yellow buffer
while the printer continues with its current red buffer. 

\begin{enumerate}
\def\labelenumi{(\alph{enumi})}
\setcounter{enumi}{1}
\tightlist
\item
  After output completion (I/O complete)
\end{enumerate}

When the printer \textbf{finishes} with the red buffer it was printing,
that buffer turns \textbf{green} (it's now free). \textbf{NEXTR}
advances clockwise to the \textbf{next red} buffer (if any). If another
red buffer exists, the printer immediately starts on it; otherwise the
device becomes idle and will resume only when a red buffer appears
later. 

\begin{enumerate}
\def\labelenumi{(\alph{enumi})}
\setcounter{enumi}{2}
\tightlist
\item
  After \textbf{RELEASE}
\end{enumerate}

When the program \textbf{RELEASE}s the yellow buffer (line image
completed), that buffer turns \textbf{red} (ready to print). If the
device had been idle (no red buffers), it can now begin printing this
newly red buffer (NEXTR will pick it up next); otherwise it joins the
queue of red buffers that the printer will reach when NEXTR cycles
around. 

\section{9 {[}21{]}}\label{section-185}

Model of operation (quick recap)

\begin{itemize}
\item
  Buffers cycle through colors: \textbf{green (free) → yellow
  (current/being filled) → red (ready for output) → green}.
\item
  \texttt{ASSIGN} picks the next \textbf{green} buffer (if none are
  green, \texttt{ASSIGN} must wait).
\item
  \texttt{RELEASE} turns the current yellow buffer \textbf{red} and
  increments the count \texttt{n} of red buffers.
\item
  The CONTROL coroutine:

  \begin{itemize}
  \tightlist
  \item
    If the device is idle and \texttt{n\textgreater{}0}, it
    \textbf{initiates OUT} on the buffer pointed to by \texttt{NEXTR}
    (taking 7500u), then advances \texttt{NEXTR} and decrements
    \texttt{n} when output completes.
  \end{itemize}
\item
  All ``compute for T'' durations are \textbf{CPU time only}; if
  \texttt{ASSIGN} has to wait for a green buffer, the CPU is idle during
  that wait (this is explicitly part of the exercise).
\end{itemize}

Deriving the schedule (walkthrough)

Let buffers be BUF1 → BUF2 → BUF3 in the ring. Initially all are green;
\texttt{NEXTR\ =\ NEXTG\ =\ BUF1}; device idle.

\textbf{t = 0}: \texttt{ASSIGN(BUF1)} → BUF1 turns yellow. Compute
1000u.

\textbf{t = 1000}: \texttt{RELEASE} → BUF1 becomes red (\texttt{n=1}).
Device is idle and \texttt{n\textgreater{}0}, so \textbf{OUT BUF1
starts} (finishes at 8500). Compute next 1000u.

\textbf{t = 2000}: \texttt{ASSIGN(BUF2)} (green) → yellow. Compute
1000u.

\textbf{t = 3000}: \texttt{RELEASE} → BUF2 red (\texttt{n=2}). Device
busy (till 8500). Compute 1000u.

\textbf{t = 4000}: \texttt{ASSIGN(BUF3)} (green) → yellow. Compute
1000u.

\textbf{t = 5000}: \texttt{RELEASE} → BUF3 red (\texttt{n=3}). Device
still busy. Compute 1000u.

\textbf{t = 6000}: Attempt \texttt{ASSIGN}, but \textbf{no green
buffers} (all three are red or in output), so \texttt{ASSIGN}
\textbf{waits}. CPU idle from 6000 onward until a buffer frees.

\textbf{t = 8500}: OUT BUF1 completes → BUF1 turns green, \texttt{n}
drops to 2. The pending \texttt{ASSIGN} now succeeds immediately:
\textbf{BUF1 assigned} (yellow). Also, since device idle and
\texttt{n\textgreater{}0}, CONTROL immediately starts \textbf{OUT BUF2}
(finishes at 16000). Compute 1000u.

\textbf{t = 9500}: \texttt{RELEASE} → BUF1 red (\texttt{n=3} again).
Device busy (till 16000). Compute continues; the next \texttt{ASSIGN}
arrives after 1000u, but there are still no greens, so it will wait:

\textbf{t = 10500}: \texttt{ASSIGN} \textbf{waits} (second idle stretch
for the CPU).

\textbf{t = 16000}: OUT BUF2 completes → BUF2 green, \texttt{n}→2.
Waiting \texttt{ASSIGN} succeeds: \textbf{assign BUF2} (yellow). Device
immediately starts \textbf{OUT BUF3} (finishes at 23500). Continuing
this same reasoning---each time a \texttt{RELEASE} makes a buffer red it
increments \texttt{n}; when the device goes idle and
\texttt{n\textgreater{}0} it starts the next \texttt{OUT}; when
\texttt{ASSIGN} encounters no greens it waits until the next
\texttt{OUT} completion---produces the full time--action table below.

Time--action chart (result)

The complete schedule (events in time order) is:

\begin{itemize}
\tightlist
\item
  0 ASSIGN(BUF1)
\item
  1000 RELEASE, OUT BUF1
\item
  2000 ASSIGN(BUF2)
\item
  3000 RELEASE
\item
  4000 ASSIGN(BUF3)
\item
  5000 RELEASE
\item
  6000 ASSIGN (wait)
\item
  8500 BUF1 assigned, OUT BUF2
\item
  9500 RELEASE
\item
  10500 ASSIGN (wait)
\item
  16000 BUF2 assigned, OUT BUF3
\item
  23000 RELEASE
\item
  23500 OUT BUF1
\item
  28000 ASSIGN(BUF3)
\item
  31000 OUT BUF2
\item
  35000 RELEASE
\item
  38500 OUT BUF3
\item
  40000 ASSIGN(BUF1)
\item
  46000 Output stops
\item
  47000 RELEASE, OUT BUF1
\item
  52000 ASSIGN(BUF2)
\item
  54500 Output stops
\item
  59000 RELEASE, OUT BUF2
\item
  64000 ASSIGN(BUF3)
\item
  65000 RELEASE
\item
  66000 ASSIGN(BUF1)
\item
  66500 OUT BUF3
\item
  68000 RELEASE
\item
  69000 Computation stops
\item
  74000 OUT BUF1
\item
  81500 Output stops.
\end{itemize}

Totals and idle times

From the chart:

\begin{itemize}
\tightlist
\item
  \textbf{Total elapsed time}: \textbf{81500u}.
\item
  \textbf{CPU idle} (waiting at \texttt{ASSIGN}): 6000--8500,
  10500--16000, 69000--81500 ⇒ \textbf{20500u}.
\item
  \textbf{Device idle}: 0--1000, 46000--47000, 54500--59000 ⇒
  \textbf{6500u}.
\end{itemize}

\section{10 {[}21{]}}\label{section-186}

\begin{itemize}
\tightlist
\item
  Use the same compute sequence from Ex. 9 (A = ASSIGN, R = RELEASE),
  where the numbers are pure \textbf{compute} times (waiting time isn't
  included):
  \texttt{A,\ 1000,\ R,\ 1000,\ \ A,\ 1000,\ R,\ 1000,\ \ A,\ 1000,\ R,\ 1000,\ \ A,\ 1000,\ R,\ 7000,\ \ A,\ 5000,\ R,\ 7000,\ \ A,\ 5000,\ R,\ 7000,\ \ A,\ 5000,\ R,\ 7000,\ \ A,\ 5000,\ R,\ 1000,\ \ A,\ 1000,\ R,\ 2000,\ \ A,\ 1000,\ R.}
  (Same list as Fig. 27/Ex. 9.) The output device runs at \textbf{7500u
  per block}.
\item
  For output: \textbf{ASSIGN} grabs a \emph{green} buffer (turns it
  yellow), \textbf{RELEASE} turns that yellow buffer \textbf{red}; the
  control starts \textbf{OUT} on the next red buffer whenever idle; when
  the output finishes, that buffer becomes green. With \textbf{N = 4}
  buffers, the program need only wait at an ASSIGN if \emph{all four}
  buffers are red (no greens).
\end{itemize}

Time--action chart (N = 4)

Below, ``OUT BUFk'' marks the instant an output of BUFk begins; ``Output
stops.'' marks instants when the device goes idle because no red buffers
remain.

\begin{verbatim}
Time   Action
-----  ---------------------------------------
0      ASSIGN(BUF1)
1000   RELEASE, OUT BUF1
2000   ASSIGN(BUF2)
3000   RELEASE
4000   ASSIGN(BUF3)
5000   RELEASE
6000   ASSIGN(BUF4)
7000   RELEASE
8500   OUT BUF2
14000  ASSIGN(BUF1)
16000  OUT BUF3
19000  RELEASE
23500  OUT BUF4
26000  ASSIGN(BUF2)
31000  OUT BUF1
31000  RELEASE
38000  ASSIGN(BUF3)
38500  OUT BUF2
43000  RELEASE
46000  OUT BUF3
50000  ASSIGN(BUF4)
53500  Output stops.
55000  RELEASE, OUT BUF4
56000  ASSIGN(BUF1)
57000  RELEASE
59000  ASSIGN(BUF2)
60000  RELEASE
61000  Computation stops.
62500  OUT BUF1
70000  OUT BUF2
77500  Output stops.
\end{verbatim}

How this unfolds (sketch): early on the program is producing red buffers
faster than the first output can finish, so red buffers accumulate until
the first finish at \textbf{8500u}; from then through \textbf{53500u}
the device is almost continuously busy draining reds as they arrive.
There is \textbf{no ``ASSIGN (wait)''} event with four buffers, because
whenever an ASSIGN occurs there is always at least one green available
(n never reaches 4 at an ASSIGN instant).

Totals and idle intervals

\begin{itemize}
\tightlist
\item
  \textbf{Total elapsed time:} \textbf{77500u} (last ``Output stops.'').
\item
  \textbf{Computation stops:} at \textbf{61000u} (there were no ASSIGN
  waits at all).
\item
  \textbf{Computer idle time:} \textbf{61000--77500} ⇒ \textbf{16500u}.
\item
  \textbf{Output device idle intervals:} \textbf{0--1000} and
  \textbf{53500--55000} ⇒ \textbf{2500u} total. (Contrast with the
  three-buffer case in Ex. 9's chart, where total time was 81500u and
  there were two mid-run computation waits.)
\end{itemize}

If you'd like, I can also tabulate buffer colors (G/Y/R) over time or
show the evolving queue of reds to make the coroutine interplay
completely transparent.

\section{11 {[}21{]}}\label{section-187}

We apply the text's buffering algorithms (A for ASSIGN, R for RELEASE, B
for buffer control): at \textbf{ASSIGN} the program waits if all buffers
are red; at \textbf{RELEASE} we increment the red-buffer count and the
controller starts output if the device is idle; when I/O completes, the
controller advances NEXTR and \textbf{decrements} the red-buffer count.

Because \textbf{N=1}, once we RELEASE the only buffer, the device is
busy until it finishes that block; the program cannot ASSIGN again (A1
stalls) until I/O completes. Thus computation and output \textbf{cannot
overlap} with a single buffer. (Contrast with the three-buffer timeline
and chart for Fig. 27.)

Time--action chart (one buffer)

\begin{longtable}[]{@{}rl@{}}
\toprule\noalign{}
Time (u) & Action \\
\midrule\noalign{}
\endhead
\bottomrule\noalign{}
\endlastfoot
0 & ASSIGN(BUF1) \\
1000 & RELEASE, OUT BUF1 \\
2000 & ASSIGN (wait) \\
8500 & BUF1 assigned \\
8500 & Output stops. \\
8500 & ASSIGN(BUF1) \\
9500 & RELEASE, OUT BUF1 \\
10500 & ASSIGN (wait) \\
17000 & BUF1 assigned \\
17000 & Output stops. \\
17000 & ASSIGN(BUF1) \\
18000 & RELEASE, OUT BUF1 \\
19000 & ASSIGN (wait) \\
25500 & BUF1 assigned \\
25500 & Output stops. \\
25500 & ASSIGN(BUF1) \\
26500 & RELEASE, OUT BUF1 \\
27500 & ASSIGN (wait) \\
34000 & BUF1 assigned \\
34000 & Output stops. \\
34000 & ASSIGN(BUF1) \\
35000 & RELEASE, OUT BUF1 \\
36000 & ASSIGN (wait) \\
42500 & BUF1 assigned \\
42500 & Output stops. \\
42500 & ASSIGN(BUF1) \\
43500 & RELEASE, OUT BUF1 \\
44500 & ASSIGN (wait) \\
51000 & BUF1 assigned \\
51000 & Output stops. \\
51000 & ASSIGN(BUF1) \\
58000 & RELEASE, OUT BUF1 \\
65000 & ASSIGN (wait) \\
65500 & BUF1 assigned \\
65500 & Output stops. \\
65500 & ASSIGN(BUF1) \\
72500 & RELEASE, OUT BUF1 \\
79500 & ASSIGN (wait) \\
80000 & BUF1 assigned \\
80000 & Output stops. \\
80000 & ASSIGN(BUF1) \\
87000 & RELEASE, OUT BUF1 \\
93500 & ASSIGN (wait) \\
94000 & BUF1 assigned \\
94000 & Output stops. \\
94000 & ASSIGN(BUF1) \\
101000 & RELEASE, OUT BUF1 \\
108500 & ASSIGN (wait) \\
109000 & BUF1 assigned \\
109000 & Output stops. \\
109000 & ASSIGN(BUF1) \\
111000 & RELEASE, OUT BUF1 \\
118500 & Output stops. \\
\end{longtable}

(Style follows the sample chart in the book: times at which
\textbf{ASSIGN}, \textbf{RELEASE}, and I/O starts/stops occur as the
system runs. For reference, the sample chart for three buffers appears
just before the exercise list, with its accompanying time-action table.
)

Totals and discussion

\begin{itemize}
\tightlist
\item
  \textbf{Total elapsed time:} \textbf{118 500u} (final ``Output
  stops'').
\item
  \textbf{Device busy time:} 11 blocks × 7 500u = \textbf{82 500u};
  \textbf{device idle:} \textbf{36 000u}.
\item
  \textbf{Computer idle time (stalls at A1):} \textbf{49 000u} (the sum
  of all ``ASSIGN (wait) → BUF1 assigned'' gaps).
\end{itemize}

With three buffers, the text's chart finishes in \textbf{81 500u} for
the same program; with a single buffer we stretch to \textbf{118 500u},
because every post-RELEASE \textbf{ASSIGN} must wait for the device to
finish before computation can continue. Buffering depth is exactly what
allows overlap between computing and I/O; with \textbf{N=1} there is no
overlap at all.

\section{12 {[}24{]}}\label{section-188}

Minimal algorithmic change (Algorithm B → B′)

Add a single global flag \texttt{EOF} (initially 0). Then:

\begin{itemize}
\tightlist
\item
  \textbf{B2′ (gate before initiating I/O)}: If \texttt{EOF\ =\ 1}
  \textbf{or} \texttt{n\ =\ 0}, go to B1 (compute). Otherwise proceed to
  B3.
\item
  \textbf{B4′ (just after I/O completes, before advancing NEXTR)}:
  Inspect the \textbf{80th column} of the buffer that was just filled
  (i.e., at address \texttt{NEXTR\ +\ 15}, byte 5). If it is
  \texttt{\textquotesingle{}.\textquotesingle{}}, set
  \texttt{EOF\ ←\ 1}. Then continue with B5/B6 (advance NEXTR;
  \texttt{n\ ←\ n\ −\ 1}).
\end{itemize}

Everything else (ASSIGN, RELEASE, pointer motion, colors) stays as Knuth
describes. The placement is important: you test for
\texttt{\textquotesingle{}.\textquotesingle{}} \textbf{after} the
\texttt{IN} you started in B3 has finished (B4) and \textbf{before} you
move \texttt{NEXTR} in B5. 

MIX patch (Program B → Program B′)

Below is a drop-in patch to Knuth's \textbf{Program B (CONTROL)}. It
assumes:

\begin{itemize}
\tightlist
\item
  rI6 ≡ \texttt{n} (count of red buffers), rI5 ≡ \texttt{NEXTR} (address
  of current red buffer's start), exactly as in the text.
\item
  \texttt{CARDS\ EQU\ 16} elsewhere; \texttt{U} already holds the unit
  (as in the book's code).
\item
  \texttt{EOF} is a one-word flag (0/1) you define in static storage
  (initialized to 0).
\item
  \texttt{DOT} is the MIX \textbf{character code} for
  \texttt{\textquotesingle{}.\textquotesingle{}} (you can define it as a
  constant; any correct value for your MIXAL assembler is fine).
\item
  Field \texttt{(5:5)} denotes the rightmost byte in a word.
\end{itemize}

Only the \textbf{new} and \textbf{changed} lines are annotated;
unchanged lines are shown for context:

\begin{Shaded}
\begin{Highlighting}[]
\NormalTok{* {-}{-}{-} New data {-}{-}{-}}
\NormalTok{EOF     CON 0            * 0 = keep reading; 1 = stop initiating input}
\NormalTok{DOT     CON 46           * MIX code for \textquotesingle{}.\textquotesingle{} (period); use your table/assembler symbol}

\NormalTok{* {-}{-}{-} CONTROL coroutine (Program B′) {-}{-}{-}}
\NormalTok{* rI6 ≡ n, rI5 ≡ NEXTR}
\NormalTok{CONT1   JMP COMPUTE          * B1. Compute (unchanged)}

\NormalTok{1H      LD1 EOF              * NEW: test end{-}of{-}input latch}
\NormalTok{        J1NZ CONT1           * If EOF=1, skip initiating more reads (go compute)}
\NormalTok{        J6Z *{-}1              * B2′: if n = 0, wait (as before)}

\NormalTok{        IN  0,5(U)           * B3. Initiate I/O into buffer at NEXTR}
\NormalTok{        JMP COMPUTE          * B4. Compute until I/O completes (JRED brings us back)}

\NormalTok{* {-}{-}{-} We return here when the reader signals ready (after IN has finished) {-}{-}{-}}

\NormalTok{        LDA 15,5(5:5)        * NEW (B4′): load col 80 = byte 5 of word (NEXTR+15)}
\NormalTok{        CMPA =DOT=           * Is it a \textquotesingle{}.\textquotesingle{} ?}
\NormalTok{        JNE  2F              * If not, continue normally}
\NormalTok{        ENTA 1               * Else set EOF ← 1}
\NormalTok{        STA  EOF}
\NormalTok{2H      LD5 {-}1,5             * B5. Advance NEXTR to next buffer (via link at [NEXTR{-}1])}
\NormalTok{        DEC6 1               * B6. n ← n {-} 1}
\NormalTok{        JMP  1B}
\end{Highlighting}
\end{Shaded}

\begin{itemize}
\tightlist
\item
  The \textbf{gate} at label \texttt{1H} enforces ``no more
  \texttt{IN}'' once \texttt{EOF} is set.
\item
  The \textbf{check} immediately after the \texttt{JMP\ COMPUTE} (on
  return) looks at \textbf{exactly} the last column of the just-read
  card and latches \texttt{EOF} if it's
  \texttt{\textquotesingle{}.\textquotesingle{}}.
\item
  We deliberately perform the test \textbf{before} \texttt{LD5\ -1,5} so
  rI5 still points at the buffer we just filled.
\end{itemize}

Why this is correct

\begin{itemize}
\tightlist
\item
  In Algorithm B, the moment a read completes (B4), the freshly read
  card still resides at \texttt{NEXTR}; only then do we advance
  \texttt{NEXTR} (B5). Inserting the test between B4 and B5 lets us
  examine the right card every time. 
\item
  Once \texttt{EOF} is latched, B2′ prevents any further \texttt{IN}
  from being initiated. Already-read buffers will continue to flow to
  the program (green→yellow→red) via ASSIGN/RELEASE exactly as before.
\item
  The mapping ``card reader → 16 words, 5 chars/word'' makes column 80
  equal to \texttt{(NEXTR\ +\ 15,\ byte\ 5)}, which is the byte we
  test. 
\end{itemize}

\begin{center}\rule{0.5\linewidth}{0.5pt}\end{center}

Here's a clean way to shut input down the moment the ``end card'' is
read (a card whose column 80 is `.'), by minimally extending Knuth's
multiple-buffer control coroutine.

What changes conceptually

In the textbook, the CONTROL coroutine (Algorithm B / Program B)
repeatedly:

\begin{itemize}
\tightlist
\item
  waits until there is at least one red buffer (n \textgreater{} 0),
\item
  starts an \texttt{IN} on the buffer at \texttt{NEXTR},
\item
  yields back to computation,
\item
  when the I/O finishes, advances \texttt{NEXTR} and decrements
  \texttt{n}.
\end{itemize}

For card input, each card image occupies 16 MIX words (5 bytes/word = 80
columns); column 80 is byte 5 of word 16, i.e., byte (5:5) of address
\texttt{NEXTR+15}.

We'll add a single sticky flag \texttt{EOF} (=0 initially). CONTROL
will:

\begin{enumerate}
\def\labelenumi{\arabic{enumi}.}
\tightlist
\item
  After each input completes (just before advancing \texttt{NEXTR}),
  inspect the last byte of the just-read card
  (\texttt{(NEXTR+15)(5:5)}).
\item
  If it equals the MIX character code for `.', set \texttt{EOF\ ←\ 1}.
\item
  At the top of the loop, treat \texttt{EOF} exactly like ``no red
  buffers'': if \texttt{EOF=1}, do not start more \texttt{IN} operations
  ever again.
\end{enumerate}

This leaves already-anticipated (green) cards, including the sentinel
card itself, available for the COMPUTE coroutine to consume; it simply
prevents any \emph{further} anticipation past the end. (Knuth explicitly
leaves such end-handling details to the reader; this is precisely the
point of Ex. 12.)

Algorithm B′ (CONTROL) --- deltas only

Relative to Algorithm B, change B2 and insert a test between B4 and B5:

\begin{itemize}
\tightlist
\item
  \textbf{B2′.} If \texttt{n\ =\ 0} \textbf{or} \texttt{EOF\ =\ 1}, go
  to B1 (keep computing; do not initiate I/O).
\item
  \textbf{B4½.} {[}Test end card.{]} If byte (5:5) of word at
  \texttt{NEXTR+15} equals `.', set \texttt{EOF\ ←\ 1}.
\end{itemize}

Everything else (B3, B5, B6) is unchanged.

Program B′ (MIX) --- minimal patch

Below is a small, self-contained patch to Knuth's Program B. It
preserves the original register discipline (CONTROL already uses only
\texttt{rI6\ ≡\ n} and \texttt{rI5\ ≡\ NEXTR}), so we save/restore
\texttt{A} around our tests to avoid surprising COMPUTE.

\textbf{New data (somewhere in CONTROL's static area):}

\begin{verbatim}
EOF   CON 0          * sticky “seen end card” flag (0/1)
SAVEA CON 0          * temp to save A
DOT   CON <code-for-'.'>   * MIX code of '.' (any correct constant is fine)
\end{verbatim}

\textbf{Patched code (comments show where lines go in Knuth's listing):}

\begin{verbatim}
Program B (CONTROL). rI6 ≡ n, rI5 ≡ NEXTR.

CONT1  JMP  COMPUTE           * B1. Compute

1H      J6Z  *-1              * B2. wait while n=0
        STA  SAVEA            * save A before we look at EOF
        LDA  EOF
        JAZ  2F               * if EOF==0, OK to start I/O
        LDA  SAVEA            * else (EOF!=0): restore A and just yield
        JMP  COMPUTE          *   behave like n=0 forever (no more IN)

2H      IN   0,5(U)           * B3. Initiate I/O
        LDA  SAVEA            * restore A before yielding
        JMP  COMPUTE          * B4. Compute until I/O completes

* ---- we return here when the read just finished on buffer NEXTR ----
* B4½. Test last column (byte 5 of word 16 in this card image).
        STA  SAVEA            * save A again
        LDA  15,5(5:5)        * A ← column 80 from buffer at NEXTR
        CMPA DOT
        JNE  3F               * if not '.', skip setting EOF
        ENTA 1
        STA  EOF              * EOF ← 1
3H      LDA  SAVEA            * restore A

        LD5  -1,5             * B5. Advance NEXTR
        DEC6 1                * B6. Decrease n
        JMP  1B
\end{verbatim}

Notes:

\begin{itemize}
\tightlist
\item
  \texttt{15,5(5:5)} addresses ``word 16 indexed by \texttt{rI5}, field
  5:5'', i.e., byte 5 of the card's last word.
\item
  Use any correct assembler literal for \texttt{DOT} (the MIX code for
  `.'). If you like, assemble \texttt{ALF\ .....} into a word and
  extract the rightmost byte once; a plain numeric constant also
  suffices.
\end{itemize}

Why this is correct

\begin{itemize}
\tightlist
\item
  If the card just read is \emph{not} the end card, CONTROL behaves
  identically to the book: advances \texttt{NEXTR}, decrements
  \texttt{n}, and continues.
\item
  If it \emph{is} the end card, \texttt{EOF} is set. From then on, the
  modified B2 treats \texttt{EOF} as a permanent ``don't start I/O''
  condition. That freezes anticipation beyond the sentinel while leaving
  already-green buffers for COMPUTE to finish processing.
\item
  The sentinel card itself is available to COMPUTE (it turned the buffer
  from red→green when the read completed), so the main program can
  recognize ``end of input'' in its normal logic after assigning that
  buffer.
\end{itemize}

\section{13 {[}20{]}}\label{section-189}

\begin{itemize}
\tightlist
\item
  In the multi-buffer algorithms, \textbf{n = current number of red
  buffers} (blocks waiting to be output). 
\item
  The control coroutine (\textbf{CONTROL}) performs I/O and, after an
  output completes, does \textbf{B5} (advance \texttt{NEXTR}) and
  \textbf{B6} (decrease \texttt{n}). 
\item
  The machine instruction \texttt{JRED\ CONTROL(U)} in \textbf{COMPUTE}
  is how CONTROL is re-entered when the device becomes ready; the same
  framework is used for output as for input, differing only in initial
  conditions. 
\end{itemize}

When computation is finished, there may still be \textbf{red} buffers
(so \texttt{n\ \textgreater{}\ 0}) and/or an output currently in
progress. You must:

\begin{enumerate}
\def\labelenumi{\arabic{enumi}.}
\tightlist
\item
  keep giving CONTROL a chance to run whenever the device becomes ready
  (so it can execute B5--B6 and drain \texttt{n}), and
\item
  only return once no red buffers remain (which also implies the device
  is idle, since \texttt{n} is decremented \emph{after} an output
  completes).
\end{enumerate}

A spin on \texttt{n} \textbf{without} \texttt{JRED\ CONTROL(U)} would
deadlock, because CONTROL wouldn't get to run.

Minimal MIX trailer to append at end of COMPUTE (output case)

\begin{Shaded}
\begin{Highlighting}[]
\NormalTok{* rI6 ≡ n   (number of red buffers)}
\NormalTok{* U   ≡ output device number used by CONTROL}

\NormalTok{FLUSH  JRED CONTROL(U)     * if the device has just become ready,}
\NormalTok{                           * jump into CONTROL to do B5/B6}
\NormalTok{       J6NZ FLUSH          * loop until n == 0 (no red buffers remain)}
\NormalTok{* Optional belt{-}and{-}suspenders wait (normally redundant once n==0):}
\NormalTok{WAIT   JBUS WAIT,U         * if device were somehow still busy, wait}
\NormalTok{* … now it’s safe to exit/return}
\end{Highlighting}
\end{Shaded}

Why this works:

\begin{itemize}
\tightlist
\item
  \texttt{JRED\ CONTROL(U)} hands control to CONTROL \textbf{exactly}
  when the device signals ready; CONTROL then finishes the
  just-completed output and executes \textbf{B6: Decrease n}. 
\item
  The loop continues until \texttt{n\ =\ 0}, meaning there are no red
  buffers left to output. Since \texttt{n} is decremented only after an
  output finishes, \texttt{n\ =\ 0} implies the device isn't in the
  middle of an output. 
\end{itemize}

That's all you need to guarantee that every buffered block has been sent
before the program terminates.

Two-instruction solution (MIX)

\begin{verbatim}
JRED CONTROL(U)
J6NZ *-1
\end{verbatim}

Why this works

\begin{itemize}
\tightlist
\item
  \texttt{JRED\ CONTROL(U)} yields to the \texttt{CONTROL} coroutine
  whenever the device \emph{is ready}. \texttt{CONTROL} then performs
  step B5/B6 (advance \texttt{NEXTR}; decrement \texttt{n}) for one
  finished output operation, and returns to \texttt{COMPUTE}. 
\item
  \texttt{J6NZ\ *-1} loops while \texttt{rI6\ =\ n\ ≠\ 0}. Thus we
  repeatedly give \texttt{CONTROL} chances to complete outstanding
  writes until no red buffers remain. When \texttt{n} finally becomes
  zero, the jump fails and the program may safely terminate. 
\end{itemize}

\section{14 {[}20{]}}\label{section-190}

Recap of the buffering model

\begin{itemize}
\tightlist
\item
  Buffers are arranged in a circle; each buffer is colored
  \textbf{green}, \textbf{yellow}, or \textbf{red}. • \textbf{Green} =
  ready to be \texttt{ASSIGN}ed (for input: already anticipated; for
  output: free). • \textbf{Yellow} = currently \texttt{ASSIGN}ed (the
  program may read/write it). • \textbf{Red} = \texttt{RELEASE}d (for
  input: free area; for output: filled, awaiting device). 
\item
  Pointers: \texttt{NEXTG} points to the next green buffer;
  \texttt{NEXTR} to the next red buffer; \texttt{CURRENT} names ``the''
  yellow buffer when there is one. Two processes run as coroutines:
  \textbf{Algorithm A (ASSIGN)} makes \texttt{CURRENT\ ←\ NEXTG} and
  advances \texttt{NEXTG}. \textbf{Algorithm R (RELEASE)} increments the
  count \texttt{n} of red buffers. \textbf{Algorithm B (Buffer control)}
  starts I/O on \texttt{NEXTR} when \texttt{n\textgreater{}0}, then
  advances \texttt{NEXTR} and decrements \texttt{n} when I/O completes. 
\end{itemize}

What if we do ASSIGN, ASSIGN, then RELEASE, RELEASE?

Think in terms of counts. At any instant:

\begin{verbatim}
#green + #yellow + #red = N ,
\end{verbatim}

and (by definition of R1) \texttt{n\ =\ \#red}. Every \texttt{ASSIGN}
\textbf{consumes} a green and creates a yellow; every \texttt{RELEASE}
\textbf{turns} a yellow into a red (and raises \texttt{n} by 1). If we
perform two \texttt{ASSIGN}s in a row while at least two greens exist,
we simply end up with \textbf{two yellow buffers} (two outstanding
assignments). Later \texttt{RELEASE}s convert those yellows to reds one
by one; Algorithm B then services those reds in FIFO pointer order. The
control logic keeps working because B only depends on \texttt{n} (reds)
and the \texttt{NEXTR} ring, not on how many yellows exist. This can
actually be \textbf{useful} when the computation wants to look at two
buffers at once (e.g., a boundary-crossing comparison, merging adjacent
blocks, etc.), at the cost of tying up an extra buffer that is
temporarily unavailable for anticipation/output. 

Important edge case and the safe bound

\begin{itemize}
\tightlist
\item
  \textbf{N = 1 (one buffer).} A burst \texttt{ASSIGN,\ ASSIGN} has no
  second distinct green to consume. The written Algorithm A's gate
  (``wait for \texttt{n\ \textless{}\ N}'') isn't sufficient to
  guarantee a green exists when we allow multiple simultaneous
  \texttt{ASSIGN}s; with one buffer you can wind up re-assigning the
  same area or, worse, overlapping program use with I/O. In short, the
  scheme \textbf{breaks down} for \texttt{N=1}. 
\item
  \textbf{General bound.} Let \texttt{y\ =\ (\#ASSIGNs\ −\ \#RELEASEs)}
  = ``\#yellow at the moment.'' Necessarily
\end{itemize}

\begin{verbatim}
0 ≤ y ≤ N .
\end{verbatim}

(If \texttt{y} exceeded \texttt{N}, there would be no green left to
assign; if \texttt{y} were negative, we'd be releasing more than we
assigned.) Practically, keep the \textbf{excess of ASSIGNs over
RELEASEs} nonnegative and \textbf{≤ N}. 

\begin{itemize}
\tightlist
\item
  With \textbf{N ≥ 2}, allowing \texttt{ASSIGN} bursts just produces
  \textbf{two (or more) simultaneous yellow buffers}; the control
  coroutine continues to function, and this can be handy when the
  computation needs concurrent access to several buffers. It does,
  however, \textbf{tie up buffer space}, leaving fewer greens available
  for anticipation.
\item
  With \textbf{N = 1}, the idea is unsafe/inapplicable; the
  multiple-buffer algorithm as given relies on there being enough
  buffers that an extra \texttt{ASSIGN} can legitimately consume another
  green. 
\end{itemize}

\section{15 {[}22{]}}\label{section-191}

Here's a fast, 3-buffer MIX solution for (``copy 100 blocks from tape
unit 0 to tape unit 1, using just three buffers''). The trick is to
\textbf{pipeline} I/O so the input and output units stay busy almost all
the time, with only a one-time ``prime'' and a one-time ``drain.'' This
is exactly the ``same-circle input+output'' idea behind Fig. 26 and the
1.4.4 exercises on pp.~225--228.

MIX program (3-buffer tape-to-tape copy)

\begin{Shaded}
\begin{Highlighting}[]
\NormalTok{* Tape{-}to{-}tape copy with triple buffering (fast)}
\NormalTok{INUNIT  EQU 0              * input tape unit}
\NormalTok{OUTUNIT EQU 1              * output tape unit}

\NormalTok{BUF1    ORIG *+100         * three 100{-}word buffers}
\NormalTok{BUF2    ORIG *+100}
\NormalTok{BUF3    ORIG *+100}

\NormalTok{START   IN   BUF1(INUNIT)  * prime the pipeline: read block \#0 {-}\textgreater{} BUF1}
\NormalTok{        ENT1 99            * 99 more blocks to output after the prime}

\NormalTok{LOOP    IN   BUF2(INUNIT)  * queue read of next block {-}\textgreater{} BUF2}
\NormalTok{        OUT  BUF1(OUTUNIT) * write the block that was in BUF1}
\NormalTok{        IN   BUF3(INUNIT)  * queue read {-}\textgreater{} BUF3}
\NormalTok{        OUT  BUF2(OUTUNIT) * write the block now in BUF2}
\NormalTok{        IN   BUF1(INUNIT)  * queue read {-}\textgreater{} recycle BUF1}
\NormalTok{        OUT  BUF3(OUTUNIT) * write the block now in BUF3}
\NormalTok{        DEC1 3}
\NormalTok{        J1P  LOOP}

\NormalTok{        JBUS *(INUNIT)     * wait for final input to finish}
\NormalTok{        OUT  BUF1(OUTUNIT) * write the very last block}
\NormalTok{        HLT}
\NormalTok{        END  START}
\end{Highlighting}
\end{Shaded}

Why this is correct and fast

\begin{enumerate}
\def\labelenumi{\arabic{enumi}.}
\tightlist
\item
  \textbf{Pipeline structure.}
\end{enumerate}

\begin{itemize}
\tightlist
\item
  We ``prime'' with one input (BUF1).
\item
  Each loop body \textbf{initiates three inputs and three outputs} in a
  fixed rotation (\ldots B2 in, B1 out; B3 in, B2 out; B1 in, B3
  out\ldots). This schedules operations so that \textbf{every OUT uses a
  buffer whose IN started two I/O steps earlier}, i.e., long enough
  before to have finished by the time the device consumes it (CPU time
  is negligible compared to device time).
\item
  We decrement by 3 because each loop iteration produces 3 completed
  outputs. Starting from 99, we do 33 iterations (99/3), then drain the
  final pending input and write the final block.
\end{itemize}

\begin{enumerate}
\def\labelenumi{\arabic{enumi}.}
\setcounter{enumi}{1}
\item
  \textbf{No buffer hazards.} Each buffer is reused only after its
  previous read and subsequent write have had time to complete; the
  triple spacing prevents ``read-while-write'' on the same memory. We
  never need to poll \textbf{during} the pipeline; a single
  \texttt{JBUS\ *(INUNIT)} at the end synchronizes the very last read
  before the last write.
\item
  \textbf{Minimal idle time ⇒ ``as fast as possible.''} After the
  one-time prime, both devices are kept busy continuously (as in the
  multiple-buffer timelines and exercises in §1.4.4), so throughput is
  device-limited. Only the leading and trailing edges incur unavoidable
  bubbles. This matches the ``same circle'' buffering method discussed
  around Fig. 26/27 and the surrounding exercises.
\item
  \textbf{Why \texttt{ENT1\ 99} and \texttt{DEC1\ 3}?} We must output
  100 blocks total. One block is primed into BUF1 before the loop, so
  \textbf{99 outputs remain}; each loop does \textbf{3 outputs}, hence
  33 passes. After the loop, the pending final read completes
  (\texttt{JBUS}) and we perform the \textbf{100th} output.
\end{enumerate}

\section{16 {[}29{]}}\label{section-192}

Here's a clean formulation of the ``green--yellow--red--purple''
multiple-buffer algorithm Knuth asks for, written in the same style as
Algorithms A/R/B in §1.4.4, but now with three coroutines: one for
input, one for output, and one for the computation. I stick to his color
semantics and ring-of-buffers pointers.

Model and invariants

\begin{itemize}
\item
  We have a single circle of \texttt{N} equal buffers.
\item
  Colors (states) and their meaning (same naming as the text and Fig.
  26): \textbf{G}reen = anticipated input is ready to be
  \textbf{ASSIGN}ed to the program; \textbf{Y}ellow = currently assigned
  to the program (between \textbf{ASSIGN} and \textbf{RELEASE});
  \textbf{R}ed = \textbf{RELEASE}d by the program; block is ready to be
  \textbf{output}; \textbf{P}urple = output finished; buffer is free for
  new \textbf{input}. (Purple is turned to green by input; yellow is
  turned to red by release; red is turned to purple by output.) 
\item
  Pointers on the ring:

  \begin{itemize}
  \tightlist
  \item
    \texttt{NEXTG} → first G (next buffer the program may assign)
  \item
    \texttt{NEXTR} → first R (next buffer the output device will
    transmit)
  \item
    \texttt{NEXTP} → first P (next buffer the input device will fill)
  \item
    \texttt{CURRENT} → Y (present only while the program holds a buffer)
  \end{itemize}
\item
  Counters we maintain:

  \begin{itemize}
  \tightlist
  \item
    \texttt{g} = \# of green buffers
  \item
    \texttt{r} = \# of red buffers
  \item
    \texttt{p} = \# of purple buffers
  \item
    \texttt{y} ∈ \{0,1\} (0 if no current assignment; 1 if a buffer is
    yellow)
  \item
    Invariant: \texttt{g\ +\ r\ +\ p\ +\ y\ =\ N}.
  \end{itemize}
\item
  The ring is always segmented in the cyclic order
  \texttt{P*\ \ G*\ \ (Y?)\ \ R*} with the pointers at the segment
  heads; this ensures none of \texttt{NEXTG}, \texttt{NEXTR},
  \texttt{NEXTP} ``passes'' another. (This is the three-pointer analogue
  of the ``no passing'' story for two pointers in Fig. 24.) 
\end{itemize}

Initial/terminal conditions

\begin{itemize}
\tightlist
\item
  Start: all buffers \textbf{purple} (the natural ``free for input''
  state). So \texttt{p\ =\ N}, \texttt{g\ =\ r\ =\ y\ =\ 0}, and set
  \texttt{NEXTP\ =\ NEXTG\ =\ NEXTR} to some buffer. Input begins
  immediately and produces the first G's for the program. 
\item
  End: after the program's last \textbf{RELEASE}, we must allow output
  to drain until \texttt{r\ =\ 0} (all finished) before declaring
  completion (the analogue of the drain step discussed for two-color
  algorithms). 
\end{itemize}

Program-side primitives (as in §1.4.4)

We keep \textbf{ASSIGN} and \textbf{RELEASE} as program-visible actions;
their algorithms mirror Knuth's A and R, but they operate on \texttt{g},
\texttt{r}, \texttt{p}, \texttt{y} and \texttt{NEXTG}:

Algorithm A′ (ASSIGN, inside the COMPUTE coroutine)

A1′. \textbf{Wait for \texttt{g\ \textgreater{}\ 0}.} (Stall the program
until some green buffer exists.) A2′.
\textbf{\texttt{CURRENT\ ←\ NEXTG}.} (The chosen G becomes the current
buffer; color changes G→Y.) A3′. \textbf{Advance \texttt{NEXTG}
clockwise} to the next G. A4′. \textbf{Update counts.}
\texttt{g\ ←\ g\ −\ 1}, \texttt{y\ ←\ 1}. ■

Algorithm R′ (RELEASE, inside the COMPUTE coroutine)

R1′. \textbf{Finish the buffer.} (Color Y→R.) R2′. \textbf{Update
counts.} \texttt{y\ ←\ 0}, \texttt{r\ ←\ r\ +\ 1}. ■

These are exactly the same ``shape'' as Knuth's A/R; only the
bookkeeping differs (his single counter \texttt{n} was the number of red
buffers in the two-party case). 

Device controllers as coroutines

The two devices now have independent controllers (each analogous to
Knuth's Algorithm B, but specialized to one color transition). Each
controller runs ``in parallel'' with COMPUTE via coroutine linkage
(COMPUTE periodically yields with \texttt{JRED\ …} as in the text). 

Coroutine IN (input controller; turns P→G)

I1. \textbf{Compute (yield to COMPUTE)} for a short while; re-enter here
later at a device-ready instant. I2. \textbf{If \texttt{p\ =\ 0},} go to
I1. (No purple → nothing to input.) I3. \textbf{Initiate input} from the
input device into the buffer at \texttt{NEXTP}. I4. \textbf{Compute
(yield to COMPUTE)} while the device is busy. I5. \textbf{On input
completion:} \textbf{advance \texttt{NEXTP} clockwise} to the next P;
\textbf{update counts: \texttt{p\ ←\ p\ −\ 1}, \texttt{g\ ←\ g\ +\ 1}.}
Go to I2. ■

Coroutine OUT (output controller; turns R→P)

O1. \textbf{Compute (yield to COMPUTE)} for a short while; re-enter here
later at a device-ready instant. O2. \textbf{If \texttt{r\ =\ 0},} go to
O1. (No red → nothing to output.) O3. \textbf{Initiate output} from the
buffer at \texttt{NEXTR} to the output device. O4. \textbf{Compute
(yield to COMPUTE)} while the device is busy. O5. \textbf{On output
completion:} \textbf{advance \texttt{NEXTR} clockwise} to the next R;
\textbf{update counts: \texttt{r\ ←\ r\ −\ 1}, \texttt{p\ ←\ p\ +\ 1}.}
Go to O2. ■

Coroutine linkage (skeleton)

Exactly as Knuth does, insert ``jump-when-ready'' instructions inside
the COMPUTE code so the device controllers can preempt when their unit
becomes ready. One practical arrangement is:

\begin{itemize}
\tightlist
\item
  COMPUTE periodically executes \texttt{JRED\ INCTRL(U\_in)} and
  \texttt{JRED\ OUTCTRL(U\_out)} (once every \textasciitilde50
  instructions as he suggests), letting each controller run steps I1/I4
  and O1/O4 at device-ready times; each controller returns to COMPUTE
  after its brief action. 
\end{itemize}

Why the pointers never ``pass''

At every transition we convert the head element of one contiguous color
segment to the head of the next segment in the cyclic order:

\begin{itemize}
\tightlist
\item
  IN does \textbf{P→G} at \texttt{NEXTP} and then advances
  \texttt{NEXTP} to the next P; thus the P-segment shrinks at its head,
  the G-segment grows at its tail, and no pointer can jump over another.
\item
  ASSIGN does \textbf{G→Y} at \texttt{NEXTG}, then advances
  \texttt{NEXTG} to the next G; the single-element Y sits between G and
  R, preserving order.
\item
  RELEASE does \textbf{Y→R} (no pointer move).
\item
  OUT does \textbf{R→P} at \texttt{NEXTR} and advances \texttt{NEXTR} to
  the next R; the R-segment shrinks at its head and the P-segment grows
  at its tail.
\end{itemize}

Hence the circular order \texttt{P*\ \ G*\ \ (Y?)\ \ R*} is invariant
and pointers cannot ``lap'' each other (the three-party generalization
of the two-pointer races discussed with Fig. 24). 

Liveness (no deadlock)

\begin{itemize}
\tightlist
\item
  If the program needs a buffer and \texttt{g\ =\ 0}, then either some
  output is pending (\texttt{r\ \textgreater{}\ 0}, so OUT will run and
  eventually increase \texttt{p}), or some input is possible
  (\texttt{p\ \textgreater{}\ 0}, so IN will run and increase
  \texttt{g}).
\item
  If an output is needed but \texttt{r\ =\ 0}, eventually a
  \textbf{RELEASE} will occur (COMPUTE makes Y→R) or new input will feed
  the pipeline to produce future R's.
\item
  If input is needed but \texttt{p\ =\ 0}, eventually an output will
  complete (R→P) or the program advances and releases something.
\end{itemize}

Provided each coroutine is scheduled fairly (as in Knuth's set-up with
periodic \texttt{JRED}), progress continues. 

Practical MIX skeleton (minimal)

You can mirror Knuth's tiny MIX control loops, one per device, and keep
A′/R′ as subroutines in COMPUTE. For example (pseudocode at the level of
his Program B):

\begin{verbatim}
; INCTRL (P->G)
INCTRL:   JMP   COMPUTE         ; I1
IWAIT:    J6Z   IWAIT           ; I2: if p==0, spin/yield
         IN    0,U_in           ; I3
         JMP   COMPUTE          ; I4
IDONE:   LDX   NEXTP            ; I5: advance NEXTP
         ; p-- ; g++ (update)
         JMP   IWAIT

; OUTCTRL (R->P)
OUTCTRL:  JMP   COMPUTE         ; O1
OWAIT:    J6Z   OWAIT           ; O2: if r==0, spin/yield
         OUT   0,U_out          ; O3
         JMP   COMPUTE          ; O4
ODONE:   LDX   NEXTR            ; O5: advance NEXTR
         ; r-- ; p++ (update)
         JMP   OWAIT
\end{verbatim}

(Exactly how you hold \texttt{g,\ r,\ p,\ y} in MIX registers or memory,
and the device unit numbers, follows the same conventions as his
Programs A/R/B; the linkage is the same pattern, only the colors and
counts differ.)

MIX setup (symbols and state)

\begin{verbatim}
* ======= DEVICE UNITS (choose any legal MIX device #s) =======
UIN     EQU     0               * input device unit
UOUT    EQU     1               * output device unit

* ======= RING & POINTERS =======
N        CON     6              * number of buffers in the ring (example)
RING     ORIG    *+N            * table of N buffer base addresses (filled elsewhere)

NEXTG    CON     0              * pointer index into RING for next Green
NEXTR    CON     0              * pointer index into RING for next Red
NEXTP    CON     0              * pointer index into RING for next Purple
CURRENT  CON     -1             * pointer index for Yellow (=-1 when none)

* ======= COUNTS (G,R,P) and Y-flag =======
GCT      CON     0              * # green
RCT      CON     0              * # red
PCT      CON     0              * # purple
YFLG     CON     0              * 0|1 whether CURRENT is valid (yellow in use)

* ======= PER-BUFFER COLOR (compact: 0=P, 1=G, 2=Y, 3=R) =======
COL0     CON     0              * COL[i] holds color code for buffer i
COL1     CON     0
COL2     CON     0
COL3     CON     0
COL4     CON     0
COL5     CON     0
COLTAB   EQU     COL0           * base for COL[i] = COLTAB+i

* ======= WORK REGS (indices) =======
* We’ll use rI1 as ‘i’, rI2 as N, rI3 as scratch.
INIT:   ENTI2   N               * load N into rI2 at boot
        ENTJ    MAIN            * fall into MAIN
\end{verbatim}

\textbf{Initialization} (as in the text): start with all buffers
\textbf{purple} (P), so \texttt{PCT=N}, \texttt{GCT=RCT=YFLG=0}, and set
\texttt{NEXTP\ =\ NEXTG\ =\ NEXTR} to the same position. From then on,
the three coroutines keep the circle moving without pointers passing
each other. 

Subroutines used by A′/controllers

To keep the listing compact (à la TAOCP), the ``advance to next buffer
with color = k'' is a tiny helper that walks clockwise modulo
\texttt{N}.

\begin{verbatim}
* ADVANCE(i_ptr, want_color) -> updates memory word at i_ptr to next index having color
* rI2 = N, rI3 scratch. Argument want_color in rA low byte (0..3). i_ptr is an address.
ADVANCE STJ     ADVX
         LDA     0,0(5:5)      * want_color already in rA’s (5:5) on entry (see callers)
         LDX     0,0           * rX free
         LDA     0,1           * load i_ptr? (we’ll pass as literal via callers)
ADVLP    LD1     0,0           * load current index i from [i_ptr] (caller does LDA then ST1)
         INC1    1
         CMP1    N             * wrap
         J1L     *+2
         ENT1    0
         STA     SAVI          * save i
         ADD     =COLTAB=      * compute &COL[i]
         STA     ADRC
         LDA     0,ADRC
         CMPA    WANTC         * matches desired color?
         JE      ADVOK
         JMP     ADVLP
ADVOK    LDA     SAVI
         STA     0,1           * store back to [i_ptr]
ADVX     JMP     0,0

SAVI     CON     0
ADRC     CON     0
WANTC    CON     0
\end{verbatim}

(NB: In a real MIX listing you'd pass the \texttt{i\_ptr} and the
desired color more directly; I've spelled this out minimally to show
intent while keeping the style close to Knuth's short library helpers.)

Program A′ --- ASSIGN (in COMPUTE)

Algorithm A′: wait until there's at least one \textbf{G}; take
\texttt{NEXTG} as \texttt{CURRENT} (G→Y), advance \texttt{NEXTG} to next
\textbf{G}, update counts. This is inserted where the computation needs
a fresh input block. 

\begin{verbatim}
* Program A′ — ASSIGN
ASSIGN  STJ     AEX             * A′ entry

A1      LDA     GCT             * A1. Wait for g > 0
        JAZ     AWAIT
        JMP     A2
AWAIT   JRED    INRD,UIN        * let INCTRL run if input just finished
        JRED    OUTRD,UOUT      * let OUTCTRL run if output just finished
        JMP     A1

A2      LDA     NEXTG           * A2. CURRENT ← NEXTG (G->Y)
        STA     CURRENT
        LDA     2               * color code 2 = Yellow
        STA     0,COLTAB(CURRENT)  * COL[CURRENT] ← Y
        STA     YFLG,=1         * y ← 1

A3      LDA     =1=             * A3. g-- ; advance NEXTG to next G
        DEC     GCT,1
        LDA     =1=
        STA     WANTC           * want_color = 1 (Green)
        LDA     NEXTG
        STA     PTRG
        JSR     ADVANCE         * scan to next Green
        LDA     PTRG
        STA     NEXTG
AEX     JMP     0,0

PTRG    CON     0
\end{verbatim}

Program R′ --- RELEASE (in COMPUTE)

Algorithm R′: mark the current buffer \textbf{R} (Y→R), clear Y-flag,
\texttt{r++}.

\begin{verbatim}
* Program R′ — RELEASE
RELEASE STJ     REX
R1      LDA     CURRENT         * R1. Y->R; r++ ; y=0
        LDA     =3=             * color code 3 = Red
        STA     0,COLTAB(CURRENT)
        INC     RCT,1
        STA     YFLG,=0
REX     JMP     0,0
\end{verbatim}

Coroutine INCTRL --- input controller (P→G)

Matches the B-style controller in §1.4.4 but specialized to P→G. Two
entries: an \textbf{idle} entry that is polled via \texttt{JRED}
(I1/I4), and a \textbf{ready} entry (I5) that runs when the unit signals
``ready.'' 

\begin{verbatim}
* INCTRL: turns Purple -> Green
INCTRL  JMP     I1              * I1. Idle entry (reached via JRED return)

I1      LDA     PCT             * I2. If p==0, nothing to start
        JAZ     IRET
        * I3. Start IN to buffer at NEXTP
        LDA     NEXTP
        STA     TMPP
        * compute buffer address RING[NEXTP] into rX:
        LD1     NEXTP
        LDX     RING,1
        IN      0,UIN           * device starts filling buffer rX
        JMP     IRET            * I4. Yield to COMPUTE

INRD    * ready-entry label (target of JRED UIN)
        * I5. Completion: P->G ; advance NEXTP ; counts
        LDA     TMPP
        STA     IDX
        LDA     =1=             * 1 = Green
        STA     0,COLTAB(IDX)
        DEC     PCT,1
        INC     GCT,1
        * advance NEXTP to next Purple
        LDA     =0=
        STA     WANTC           * want_color=0 (Purple)
        LDA     NEXTP
        STA     PTRP
        JSR     ADVANCE
        LDA     PTRP
        STA     NEXTP
IRET    JMP     COMPUTE

TMPP    CON     0
PTRP    CON     0
IDX     CON     0
\end{verbatim}

Coroutine OUTCTRL --- output controller (R→P)

Symmetric to input: when \texttt{r\textgreater{}0} and device idle,
start an \texttt{OUT} from \texttt{NEXTR}. On ready, mark \textbf{P},
advance \texttt{NEXTR}, update counts. 

\begin{verbatim}
* OUTCTRL: turns Red -> Purple
OUTCTRL JMP     O1              * O1. Idle entry

O1      LDA     RCT             * O2. If r==0, nothing to start
        JAZ     ORET
        * O3. Start OUT from buffer at NEXTR
        LDA     NEXTR
        STA     TMPR
        LD1     NEXTR
        LDX     RING,1
        OUT     0,UOUT
        JMP     ORET            * O4. Yield

OUTRD   * ready-entry label (target of JRED UOUT)
        * O5. Completion: R->P ; advance NEXTR ; counts
        LDA     TMPR
        STA     IDX2
        LDA     =0=             * 0 = Purple
        STA     0,COLTAB(IDX2)
        DEC     RCT,1
        INC     PCT,1
        * advance NEXTR to next Red
        LDA     =3=
        STA     WANTC
        LDA     NEXTR
        STA     PTRR
        JSR     ADVANCE
        LDA     PTRR
        STA     NEXTR
ORET    JMP     COMPUTE

TMPR    CON     0
PTRR    CON     0
IDX2    CON     0
\end{verbatim}

COMPUTE coroutine (skeleton)

A trivial loop that periodically gives time to device controllers via
\texttt{JRED}, and uses \texttt{ASSIGN}/\texttt{RELEASE} around
``work''. The \texttt{JRED} polling cadence (e.g., ``every
\textasciitilde50 instructions'') follows §1.4.4's guidance. 

\begin{verbatim}
MAIN    * ===== Boot: all buffers Purple, set pointers equal =====
        STZ     GCT
        STZ     RCT
        LDA     N
        STA     PCT
        STZ     YFLG
        STZ     NEXTG
        STZ     NEXTR
        STZ     NEXTP
        * set COL[i]=P for all i
        ENT1    0
CLP     LDA     =0=
        STA     0,COLTAB(1)
        INC1    1
        CMP1    N
        J1L     CLP

COMPUTE * ===== Example compute loop =====
        * A′: need an input block
        JSR     ASSIGN
        * (use CURRENT buffer: user’s computation goes here)
        * --- do some work; sprinkle JRED calls periodically ---
        JRED    INRD,UIN
        JRED    OUTRD,UOUT
        JRED    INRD,UIN
        JRED    OUTRD,UOUT
        * done with current block:
        JSR     RELEASE
        * repeat/exit per program logic
        JMP     COMPUTE
        HLT
END     MAIN
\end{verbatim}

Notes on correctness and style

\begin{itemize}
\tightlist
\item
  The three pointers \texttt{NEXTG}, \texttt{NEXTR}, \texttt{NEXTP} and
  the \textbf{single} \texttt{CURRENT} (yellow) are kept in the cyclic
  order \textbf{P* G* (Y?) R*}, so no pointer can ever ``pass'' another;
  each controller converts the head of its segment to the next color and
  advances its own pointer only within its segment, exactly mirroring
  the invariant under Fig. 26. 
\item
  \texttt{ASSIGN}/\texttt{RELEASE} are drop-in analogs of Programs A/R
  (here A′/R′) from the section's ``multiple buffers'' discussion; the
  only differences are the extra counts and the scan to the next
  \textbf{green}. 
\item
  \texttt{JRED\ INRD,UIN} and \texttt{JRED\ OUTRD,UOUT} provide the
  coroutine linkage (B-style control) so that device ``ready'' events
  run the I/O completion steps (I5/O5) promptly while the computation
  proceeds. 
\end{itemize}

\section{17 {[}40{]}}\label{section-193}

Data structures and invariants

\textbf{Buffers:} Each buffer \texttt{b} has fields
\texttt{(state,\ dev)}, where
\texttt{state\ ∈\ \{FREE,\ GREEN,\ YELLOW,\ RED\}} and \texttt{dev} is a
device id.

\textbf{Global pool/list:}

\begin{itemize}
\tightlist
\item
  \texttt{FREE}: list of free buffers.
\item
  For each \textbf{input} device \texttt{d}: queue \texttt{G{[}d{]}} of
  GREEN buffers (ready for compute), and counter
  \texttt{g{[}d{]}\ =\ \textbar{}G{[}d{]}\textbar{}}.
\item
  For each \textbf{output} device \texttt{o}: queue \texttt{R{[}o{]}} of
  RED buffers (ready for device), and counter
  \texttt{r{[}o{]}\ =\ \textbar{}R{[}o{]}\textbar{}}.
\item
  Device status \texttt{idle(d)} / \texttt{idle(o)} and an ``in-flight''
  counter for I/O requests in progress.
\end{itemize}

\textbf{Anti-flood control (the key to the exercise):}

\begin{itemize}
\tightlist
\item
  \textbf{Per-input token bucket.} For each input device \texttt{d},
  maintain \texttt{tok{[}d{]}} with bounds
  \texttt{0\ ≤\ tok{[}d{]}\ ≤\ cap{[}d{]}} (small cap, e.g., 2). --
  \textbf{Consume one token} when scheduling an \texttt{IN} for
  \texttt{d}. -- \textbf{Produce one token} when computation
  \emph{takes} a GREEN buffer of \texttt{d} (i.e., on
  \texttt{ASSIGN\_IN(d)}), reflecting actual consumption. This ensures
  anticipation never runs more than \texttt{cap{[}d{]}} blocks ahead of
  use.
\item
  \textbf{Free-buffer reserve.} Maintain \texttt{reserve\_free\ ≥\ 1}
  (e.g., \texttt{reserve\_free\ =\ \#output\_devices}, one per output).
  Do \textbf{not} start new input if
  \texttt{\textbar{}FREE\textbar{}\ ≤\ reserve\_free}, \textbf{unless}
  compute is currently \textbf{blocked} waiting for that specific input
  device \texttt{d} (the text's exception). 
\item
  \textbf{Priority:} Output drains (RED→device) always take precedence
  over more input anticipation.
\end{itemize}

These three rules are sufficient to keep pooled buffers from stalling
due to over-prefetch.

Algorithms (pooled variants of A/R/B)

Algorithm A\_pool --- ASSIGN (normal state)

There are two entry points:

\textbf{A\_pool.IN(d):} hand compute the next input block from device
\texttt{d}. A1. If \texttt{G{[}d{]}} nonempty: pop head \texttt{b} from
\texttt{G{[}d{]}}, set \texttt{state(b)←YELLOW}, return \texttt{b}. Also
\textbf{tok{[}d{]} ← min(tok{[}d{]}+1, cap{[}d{]})} (credit for
consumption). A2. Otherwise: mark \texttt{need{[}d{]}←true}, invoke
\texttt{B\textbackslash{}\_pool} (control), and wait until
\texttt{G{[}d{]}} becomes nonempty, then go to A1.

\textbf{A\_pool.OUT(o):} obtain a workspace buffer destined for output
device \texttt{o}. A3. If \texttt{FREE} nonempty: take \texttt{b} from
\texttt{FREE}, set \texttt{state(b)←YELLOW}, \texttt{dev(b)←o}, return
\texttt{b}. A4. Otherwise: invoke \texttt{B\textbackslash{}\_pool} and
wait until \texttt{FREE} nonempty, then A3.

Algorithm R\_pool --- RELEASE (normal state)

\textbf{R1.} If \texttt{b} was obtained by \texttt{ASSIGN\_IN},
computation has \emph{consumed} it: set \texttt{state(b)←FREE}, push to
\texttt{FREE}. \textbf{R2.} If \texttt{b} was obtained by
\texttt{ASSIGN\_OUT(o)}, computation has \emph{produced} it: set
\texttt{state(b)←RED}, enqueue to \texttt{R{[}o{]}}. \textbf{R3.}
Optionally call \texttt{B\textbackslash{}\_pool} (helps prompt
draining).

Algorithm B\_pool --- CONTROL (runs on I/O completion or when A/R asks)

\textbf{B1. Drain outputs first (priority).} For each \textbf{idle}
output device \texttt{o} with \texttt{R{[}o{]}} nonempty:  pop
\texttt{b} from \texttt{R{[}o{]}}, start \texttt{OUT(o,\ b)} (device
becomes busy).

\textbf{B2. Start essential inputs (to unblock compute).} For each input
device \texttt{d} with \texttt{need{[}d{]}==true} and
\texttt{idle(d)==true}:  If \texttt{FREE} nonempty (even if
\texttt{\textbar{}FREE\textbar{}\ ≤\ reserve\_free}), take \texttt{b}
from \texttt{FREE}, set \texttt{state(b)←(in-flight\ IN\ for\ d)},
\textbf{tok{[}d{]}−− (if tok{[}d{]}==0, set to 0; we allow this
unblocking read even if tokens are 0)}, start \texttt{IN(d,\ b)}, clear
\texttt{need{[}d{]}}. (Exactly the exception suggested in the text: we
may use the final buffer for input only if compute is stalled for that
input.) 

\textbf{B3. Conservative anticipation for inputs.} For each
\textbf{idle} input device \texttt{d}, if
\texttt{tok{[}d{]}\ \textgreater{}\ 0} \textbf{and}
\texttt{\textbar{}FREE\textbar{}\ \textgreater{}\ reserve\_free}:  take
\texttt{b} from \texttt{FREE}, set \texttt{(in-flight\ IN\ for\ d)},
\textbf{tok{[}d{]}−−}, start \texttt{IN(d,\ b)}.

\textbf{B4. On I/O completion (interrupt or poll):} -- If an
\texttt{OUT(o,\ b)} completes: set \texttt{state(b)←FREE}, push
\texttt{FREE}. -- If an \texttt{IN(d,\ b)} completes: set
\texttt{state(b)←GREEN}, enqueue to \texttt{G{[}d{]}}. Then re-enter B1
(outputs may now drain again).

This is the entire pooled-buffer controller. It's short because the
\textbf{token bucket} + \textbf{reserve} encode the needed policy
compactly.

Why this prevents slowdown

\begin{itemize}
\tightlist
\item
  \textbf{No green-pool flood.} Tokens ensure each input device
  anticipates at most \texttt{cap{[}d{]}} blocks beyond what compute
  actually consumes; you can't fill the pool with ``anticipated input''
  because consumption is the only way to mint new tokens.
\item
  \textbf{Output always proceeds.} Output drains have strict priority
  (B1) and at least \texttt{reserve\_free} buffers are kept out of
  anticipatory reads, so produced blocks can turn RED and get scheduled
  without starvation.
\item
  \textbf{Progress when blocked.} If compute stalls waiting for
  \texttt{d}, B2 violates both the token and reserve limits \textbf{only
  for that \texttt{d}}, matching Knuth's exception: you may use the last
  buffer to get unstuck for the specific input that blocks compute. 
\end{itemize}

Tuning notes

\begin{itemize}
\tightlist
\item
  Typical choices: \texttt{reserve\_free\ =\ \#output\_devices} and
  \texttt{cap{[}d{]}\ ∈\ \{1,2\}}. With fast outputs or many outputs,
  increase \texttt{reserve\_free}. With bursty inputs,
  \texttt{cap{[}d{]}=2} smooths latency while still preventing flood.
\item
  You may add a small \textbf{hysteresis} (high/low watermark) around
  \texttt{reserve\_free} to avoid thrashing (start anticipation only
  when \texttt{\textbar{}FREE\textbar{}\ ≥\ reserve\_free+H}, stop when
  \texttt{\textbar{}FREE\textbar{}\ ≤\ reserve\_free}).
\end{itemize}

\begin{enumerate}
\def\labelenumi{\arabic{enumi}.}
\tightlist
\item
  Overview: how the algorithms map to MIX
\end{enumerate}

\begin{itemize}
\tightlist
\item
  \textbf{A\_pool.IN(d)} takes a ready (GREEN) input buffer for device
  \texttt{d} or blocks (setting \texttt{NEED{[}d{]}}) and calls
  \texttt{B\_pool}.
\item
  \textbf{A\_pool.OUT(o)} takes a FREE buffer for output device
  \texttt{o} or blocks and calls \texttt{B\_pool}.
\item
  \textbf{R\_pool.RELEASE} returns a consumed input buffer to
  \texttt{FREE} or enqueues a produced output buffer to
  \texttt{R{[}o{]}} (RED).
\item
  \textbf{B\_pool} runs the scheduler: drain outputs first, then do
  essential inputs that unblock compute, then conservative anticipatory
  inputs (respecting tokens and the free-buffer reserve).
\item
  \textbf{CHECK\_IO} polls devices with \texttt{JBUS} to detect
  completions and moves finished buffers to \texttt{FREE} (OUT complete)
  or \texttt{G{[}d{]}} (IN complete).
\end{itemize}

\begin{enumerate}
\def\labelenumi{\arabic{enumi}.}
\setcounter{enumi}{1}
\tightlist
\item
  Symbols, constants, and layout
\end{enumerate}

\begin{Shaded}
\begin{Highlighting}[]
\NormalTok{        * {-}{-}{-}{-}{-}{-}{-}{-}{-}{-} PARAMETERS }\OperatorTok{{-}{-}{-}{-}{-}{-}{-}{-}{-}{-}}
\NormalTok{M       EQU     }\DecValTok{3}               \OperatorTok{*}\NormalTok{ number of input devices}
\NormalTok{N       EQU     }\DecValTok{2}               \OperatorTok{*}\NormalTok{ number of output devices}
\NormalTok{B       EQU     }\DecValTok{16}              \OperatorTok{*}\NormalTok{ number of buffers in pool }\OperatorTok{(\textgreater{}=}\NormalTok{ M }\OperatorTok{+}\NormalTok{ N }\OperatorTok{+}\NormalTok{ working set}\OperatorTok{)}
\NormalTok{BUFSZ   EQU     }\DecValTok{100}             \OperatorTok{*}\NormalTok{ words per buffer }\OperatorTok{(}\NormalTok{tune per device}\OperatorTok{)}
\NormalTok{RESV    EQU     N               }\OperatorTok{*}\NormalTok{ reserve\_free}\OperatorTok{:}\NormalTok{ keep at least N buffers free}

\NormalTok{        * Device numbers }\OperatorTok{(}\NormalTok{example}\CommentTok{; change to your MIX)}
\NormalTok{D1      EQU     }\DecValTok{16}              \OperatorTok{*}\NormalTok{ paper tape reader}\OperatorTok{,}\NormalTok{ say}
\NormalTok{D2      EQU     }\DecValTok{17}              \OperatorTok{*}\NormalTok{ disk input unit}
\NormalTok{D3      EQU     }\DecValTok{18}
\NormalTok{O1      EQU     }\DecValTok{19}              \OperatorTok{*}\NormalTok{ line printer}\OperatorTok{,}\NormalTok{ say}
\NormalTok{O2      EQU     }\DecValTok{20}

\NormalTok{        * {-}{-}{-}{-}{-}{-}{-}{-}{-}{-} BUFFER DESCRIPTORS }\OperatorTok{{-}{-}{-}{-}{-}{-}{-}{-}{-}{-}}
\NormalTok{        * Each buffer descriptor }\OperatorTok{(}\NormalTok{BD}\OperatorTok{[}\NormalTok{i}\OperatorTok{])}\NormalTok{ holds}\OperatorTok{:}\NormalTok{ NEXT}\OperatorTok{,}\NormalTok{ STATE}\OperatorTok{,}\NormalTok{ DEV}\OperatorTok{,}\NormalTok{ ADDR}
\NormalTok{        * STATE}\OperatorTok{:} \DecValTok{0}\OperatorTok{=}\NormalTok{FREE}\OperatorTok{,} \DecValTok{1}\OperatorTok{=}\NormalTok{GREEN}\OperatorTok{,} \DecValTok{2}\OperatorTok{=}\NormalTok{YELLOW}\OperatorTok{,} \DecValTok{3}\OperatorTok{=}\NormalTok{RED}\OperatorTok{,} \DecValTok{4}\OperatorTok{=}\NormalTok{INFLIGHT\_IN}\OperatorTok{,} \DecValTok{5}\OperatorTok{=}\NormalTok{INFLIGHT\_OUT}
\NormalTok{        * DEV}\OperatorTok{:}\NormalTok{   device id }\OperatorTok{(}\NormalTok{for GREEN}\OperatorTok{/}\NormalTok{RED}\OperatorTok{/}\NormalTok{INFLIGHT states}\OperatorTok{)}
\NormalTok{        * ADDR}\OperatorTok{:}\NormalTok{  start address of the data area for this buffer}
\NormalTok{STATE\_FREE  EQU }\DecValTok{0}
\NormalTok{STATE\_GREEN EQU }\DecValTok{1}
\NormalTok{STATE\_YEL   EQU }\DecValTok{2}
\NormalTok{STATE\_RED   EQU }\DecValTok{3}
\NormalTok{STATE\_INFLI EQU }\DecValTok{4}
\NormalTok{STATE\_INFLO EQU }\DecValTok{5}

\NormalTok{NIL     EQU     }\DecValTok{0}

\NormalTok{        ORIG    }\DecValTok{1000}            \OperatorTok{*} \OperatorTok{(}\NormalTok{choose a safe origin}\OperatorTok{)}

\NormalTok{FREEHD  CON     NIL             }\OperatorTok{*}\NormalTok{ head pointer to FREE list}

\NormalTok{        * G queues }\OperatorTok{(}\NormalTok{one per input d}\OperatorTok{)} \OperatorTok{and}\NormalTok{ R queues }\OperatorTok{(}\NormalTok{one per output o}\OperatorTok{)}
\NormalTok{GHD1    CON     NIL}
\NormalTok{GHD2    CON     NIL}
\NormalTok{GHD3    CON     NIL}
\NormalTok{RHO1    CON     NIL}
\NormalTok{RHO2    CON     NIL}

\NormalTok{        * "Need}\StringTok{" flags: compute blocked waiting for input from d?}
\NormalTok{NEED1   CON     }\DecValTok{0}
\NormalTok{NEED2   CON     }\DecValTok{0}
\NormalTok{NEED3   CON     }\DecValTok{0}

\NormalTok{        * Token buckets}\OperatorTok{:}\NormalTok{ TOK}\OperatorTok{[}\NormalTok{d}\OperatorTok{]}\NormalTok{ ∈ }\OperatorTok{[}\FloatTok{0.}\OperatorTok{.}\NormalTok{CAP}\OperatorTok{[}\NormalTok{d}\OperatorTok{]]}
\NormalTok{TOK1    CON     }\DecValTok{1}               \OperatorTok{*}\NormalTok{ initialize to small positive values}
\NormalTok{TOK2    CON     }\DecValTok{1}
\NormalTok{TOK3    CON     }\DecValTok{1}
\NormalTok{CAP1    CON     }\DecValTok{2}               \OperatorTok{*}\NormalTok{ caps can be }\DecValTok{1} \OperatorTok{or} \DecValTok{2} \OperatorTok{(}\NormalTok{prefetch depth}\OperatorTok{)}
\NormalTok{CAP2    CON     }\DecValTok{2}
\NormalTok{CAP3    CON     }\DecValTok{2}

\NormalTok{        * For each device}\OperatorTok{,}\NormalTok{ track in}\OperatorTok{{-}}\NormalTok{flight buffer }\OperatorTok{(}\DecValTok{0}\NormalTok{ if none}\OperatorTok{)}
\NormalTok{INFLI1  CON     }\DecValTok{0}               \OperatorTok{*}\NormalTok{ input d1}
\NormalTok{INFLI2  CON     }\DecValTok{0}
\NormalTok{INFLI3  CON     }\DecValTok{0}
\NormalTok{INFLO1  CON     }\DecValTok{0}               \OperatorTok{*}\NormalTok{ output o1}
\NormalTok{INFLO2  CON     }\DecValTok{0}

\NormalTok{        * Busy flags are implicit via JBUS}\CommentTok{; you can also mirror them if you like}
\NormalTok{        * (}\BuiltInTok{not}\NormalTok{ required}\OperatorTok{).}

\NormalTok{        * Buffer descriptor table}
\NormalTok{BDTAB   EQU     }\OperatorTok{*}               \OperatorTok{*}\NormalTok{ BD}\OperatorTok{[}\FloatTok{1.}\OperatorTok{.}\NormalTok{B}\OperatorTok{],} \DecValTok{3}\NormalTok{ words each }\OperatorTok{+}\NormalTok{ ADDR}
\NormalTok{        * We also allocate data areas per buffer behind BD table}\OperatorTok{.}
\NormalTok{        * Layout helper}\OperatorTok{:}\NormalTok{ BD}\OperatorTok{[}\NormalTok{i}\OperatorTok{].}\NormalTok{NEXT at BDNXT}\OperatorTok{+(}\NormalTok{i}\OperatorTok{{-}}\DecValTok{1}\OperatorTok{),}\NormalTok{ etc}\OperatorTok{.}
\NormalTok{BDNXT   CON     NIL             }\OperatorTok{*}\NormalTok{ BD}\OperatorTok{[}\DecValTok{1}\OperatorTok{].}\NormalTok{NEXT}
\NormalTok{BDST1   CON     STATE\_FREE      }\OperatorTok{*}\NormalTok{ BD}\OperatorTok{[}\DecValTok{1}\OperatorTok{].}\NormalTok{STATE}
\NormalTok{BDDEV1  CON     }\DecValTok{0}               \OperatorTok{*}\NormalTok{ BD}\OperatorTok{[}\DecValTok{1}\OperatorTok{].}\NormalTok{DEV}
\NormalTok{BDADR1  CON     }\DecValTok{0}               \OperatorTok{*}\NormalTok{ BD}\OperatorTok{[}\DecValTok{1}\OperatorTok{].}\NormalTok{ADDR}
\NormalTok{        * ... replicate for BD}\OperatorTok{[}\FloatTok{2.}\OperatorTok{.}\NormalTok{B}\OperatorTok{]} \OperatorTok{(}\NormalTok{macro expands this in a real assembly}\OperatorTok{)}
\NormalTok{        * For brevity}\OperatorTok{,}\NormalTok{ assume you generate these with a macro}\OperatorTok{.}
\NormalTok{        * Also compute BDADR}\OperatorTok{[}\NormalTok{i}\OperatorTok{]}\NormalTok{ to point to a BUFSZ block per buffer}\OperatorTok{.}

\NormalTok{        * Free}\OperatorTok{{-}}\NormalTok{list bootstrapping }\OperatorTok{(}\NormalTok{done once at init}\OperatorTok{):}
\NormalTok{        * Link all BD}\OperatorTok{[}\NormalTok{i}\OperatorTok{]}\NormalTok{ in a singly linked list}\OperatorTok{:}\NormalTok{ FREEHD }\OperatorTok{{-}\textgreater{}} \DecValTok{1} \OperatorTok{{-}\textgreater{}} \DecValTok{2} \OperatorTok{{-}\textgreater{}} \OperatorTok{...} \OperatorTok{{-}\textgreater{}}\NormalTok{ B }\OperatorTok{{-}\textgreater{}}\NormalTok{ NIL}
\end{Highlighting}
\end{Shaded}

\begin{enumerate}
\def\labelenumi{\arabic{enumi}.}
\setcounter{enumi}{2}
\tightlist
\item
  List primitives (ENQ/DEQ on singly linked heads)
\end{enumerate}

We use head insertion for O(1) enqueue; dequeue also from head (so
queues act as stacks --- fine for buffers). If you want FIFO queues,
maintain both head and tail; the scaffold keeps it simple.

\begin{Shaded}
\begin{Highlighting}[]
\NormalTok{* ENQ b into list at }\OperatorTok{(}\NormalTok{HEADPTR}\OperatorTok{)}
\NormalTok{*   }\BuiltInTok{IN}\OperatorTok{:}\NormalTok{ A }\OperatorTok{=}\NormalTok{ address of HEADPTR }\OperatorTok{(}\NormalTok{the }\DataTypeTok{word}\NormalTok{ containing head pointer}\OperatorTok{)}
\NormalTok{*       X }\OperatorTok{=}\NormalTok{ buffer index b }\OperatorTok{(}\FloatTok{1.}\OperatorTok{.}\NormalTok{B}\OperatorTok{)}
\NormalTok{ENQ     STJ     ENQ\_RET}
\NormalTok{        LDA     }\DecValTok{0}\OperatorTok{,}\DecValTok{1}             \OperatorTok{*}\NormalTok{ A }\OperatorTok{:=}\NormalTok{ HEADPTR}
\NormalTok{        LD1     }\DecValTok{0}\OperatorTok{,}\DecValTok{1}             \OperatorTok{*}\NormalTok{ X1 }\OperatorTok{:=} \OperatorTok{[}\NormalTok{HEADPTR}\OperatorTok{]} \OperatorTok{(}\NormalTok{current head}\OperatorTok{)}
\NormalTok{        STA     TMPHEAD}
\NormalTok{        ENT2    }\DecValTok{0}               \OperatorTok{*}\NormalTok{ X2 }\OperatorTok{:=} \DecValTok{0}
\NormalTok{        * }\BuiltInTok{Set}\NormalTok{ BD}\OperatorTok{[}\NormalTok{b}\OperatorTok{].}\NormalTok{NEXT }\OperatorTok{=}\NormalTok{ old head}
\NormalTok{        LDA     }\OperatorTok{=}\NormalTok{BDNXT}\OperatorTok{{-}}\DecValTok{1}\OperatorTok{=,}\NormalTok{X     }\OperatorTok{*}\NormalTok{ address of BD}\OperatorTok{[}\NormalTok{b}\OperatorTok{].}\NormalTok{NEXT}\CommentTok{; (use proper base math in real code)}
\NormalTok{        LDX     TMPHEAD}
\NormalTok{        STA     }\DecValTok{0}\OperatorTok{,}\NormalTok{X             }\OperatorTok{*}\NormalTok{ BD}\OperatorTok{[}\NormalTok{b}\OperatorTok{].}\NormalTok{NEXT }\OperatorTok{=}\NormalTok{ old head}
\NormalTok{        * Update head to b}
\NormalTok{        LDX     X               }\OperatorTok{*}\NormalTok{ X holds b already}
\NormalTok{        STA     }\DecValTok{0}\OperatorTok{,}\DecValTok{1}             \OperatorTok{*} \OperatorTok{[}\NormalTok{HEADPTR}\OperatorTok{]} \OperatorTok{=}\NormalTok{ b}
        \ControlFlowTok{JMP}\NormalTok{     ENQ\_RET}
\NormalTok{ENQ\_RET NOP}

\NormalTok{* DEQ from list at }\OperatorTok{(}\NormalTok{HEADPTR}\OperatorTok{),}\NormalTok{ returns X }\OperatorTok{=}\NormalTok{ b }\OperatorTok{or}\NormalTok{ X }\OperatorTok{=} \DecValTok{0}\NormalTok{ if empty}
\NormalTok{DEQ     STJ     DEQ\_RET}
\NormalTok{        LDA     }\DecValTok{0}\OperatorTok{,}\DecValTok{1}             \OperatorTok{*}\NormalTok{ head}
\NormalTok{        JAZ     DEQ\_EMPTY}
\NormalTok{        LDX     }\DecValTok{0}\OperatorTok{,}\DecValTok{1}             \OperatorTok{*}\NormalTok{ X }\OperatorTok{:=}\NormalTok{ head b}
\NormalTok{        * head }\OperatorTok{:=}\NormalTok{ BD}\OperatorTok{[}\NormalTok{b}\OperatorTok{].}\NormalTok{NEXT}
\NormalTok{        LDA     }\OperatorTok{=}\NormalTok{BDNXT}\OperatorTok{{-}}\DecValTok{1}\OperatorTok{=,}\NormalTok{X}
\NormalTok{        LDA     }\DecValTok{0}\OperatorTok{,}\NormalTok{X}
\NormalTok{        STA     }\DecValTok{0}\OperatorTok{,}\DecValTok{1}
        \ControlFlowTok{JMP}\NormalTok{     DEQ\_DONE}
\NormalTok{DEQ\_EMPTY}
\NormalTok{        ENTX    }\DecValTok{0}               \OperatorTok{*}\NormalTok{ X }\OperatorTok{:=} \DecValTok{0}
\NormalTok{DEQ\_DONE}
        \ControlFlowTok{JMP}\NormalTok{     DEQ\_RET}
\NormalTok{DEQ\_RET NOP}

\NormalTok{TMPHEAD CON     }\DecValTok{0}
\end{Highlighting}
\end{Shaded}

(Notes: In real MIXAL you'll implement the ``address arithmetic'' with
proper \texttt{EQU}/macros so \texttt{BDNXT-1\ +\ b} hits
BD{[}b{]}.NEXT. I'm showing intent; translate to your macro system.)

\begin{enumerate}
\def\labelenumi{\arabic{enumi}.}
\setcounter{enumi}{3}
\tightlist
\item
  Helpers to read/write BD fields
\end{enumerate}

For clarity, define tiny helpers/macros (or inline):

\begin{itemize}
\tightlist
\item
  \texttt{SET\_STATE(b,\ s)}
\item
  \texttt{SET\_DEV(b,\ d)}
\item
  \texttt{SET\_NEXT(b,\ nx)}
\item
  \texttt{GET\_ADR(b)\ -\textgreater{}\ A\ :=\ BDADR{[}b{]}}
\end{itemize}

I'll inline them below to keep the listing compact.

\begin{enumerate}
\def\labelenumi{\arabic{enumi}.}
\setcounter{enumi}{4}
\tightlist
\item
  Core routines
\end{enumerate}

5.1 A\_pool.IN(d)

\begin{Shaded}
\begin{Highlighting}[]
\NormalTok{* A\_IN}\OperatorTok{(}\NormalTok{d}\OperatorTok{):}\NormalTok{ try G}\OperatorTok{[}\NormalTok{d}\OperatorTok{],}\NormalTok{ else set NEED}\OperatorTok{[}\NormalTok{d}\OperatorTok{]=}\DecValTok{1}\OperatorTok{,}\NormalTok{ call B\_pool}\OperatorTok{,}\NormalTok{ wait until G}\OperatorTok{[}\NormalTok{d}\OperatorTok{]}\NormalTok{ nonempty}
\NormalTok{* Example for d }\OperatorTok{=}\NormalTok{ D1 }\OperatorTok{(}\NormalTok{duplicate for D2}\OperatorTok{,}\NormalTok{ D3 }\OperatorTok{or}\NormalTok{ generate with macro}\OperatorTok{)}

\NormalTok{A\_IN1   STJ     AIN1\_RET}

\NormalTok{        * Try dequeue from GHD1}
\NormalTok{        LDA     }\OperatorTok{=}\NormalTok{GHD1}\OperatorTok{=}
\NormalTok{        ENTA    }\OperatorTok{=}\NormalTok{GHD1}\OperatorTok{=}
\NormalTok{        ENTX    }\DecValTok{0}               \OperatorTok{*}\NormalTok{ X acts as result}
        \ControlFlowTok{JMP}\NormalTok{     AIN1\_DEQ}
\NormalTok{AIN1\_DEQ}
\NormalTok{        * }\ControlFlowTok{call}\NormalTok{ DEQ with A}\OperatorTok{=}\NormalTok{addr}\OperatorTok{(}\NormalTok{GHD1}\OperatorTok{)}
\NormalTok{        * (Assume A already }\OperatorTok{=}\NormalTok{ addr}\OperatorTok{(}\NormalTok{GHD1}\OperatorTok{))}
\NormalTok{        JSR     DEQ}
\NormalTok{        JXZ     AIN1\_BLOCK      }\OperatorTok{*}\NormalTok{ X}\OperatorTok{=}\DecValTok{0}\NormalTok{ means empty }\OperatorTok{{-}\textgreater{}}\NormalTok{ block}

\NormalTok{        * Got buffer b in X}\OperatorTok{:}\NormalTok{ mark YELLOW }\OperatorTok{and}\NormalTok{ credit token}
\NormalTok{        * STATE}\OperatorTok{[}\NormalTok{b}\OperatorTok{]} \OperatorTok{:=}\NormalTok{ YELLOW}
\NormalTok{        LDA     }\OperatorTok{=}\NormalTok{BDST1}\OperatorTok{{-}}\DecValTok{1}\OperatorTok{=,}\NormalTok{X}
\NormalTok{        ENTA    STATE\_YEL}
\NormalTok{        STA     }\DecValTok{0}\OperatorTok{,}\NormalTok{X}
\NormalTok{        * TOK1 }\OperatorTok{:=}\NormalTok{ min}\OperatorTok{(}\NormalTok{TOK1}\OperatorTok{+}\DecValTok{1}\OperatorTok{,}\NormalTok{ CAP1}\OperatorTok{)}
\NormalTok{        LDA     TOK1}
        \BuiltInTok{ADD}\NormalTok{     ONE}
\NormalTok{        CMPA    CAP1}
        \ControlFlowTok{JGE}\NormalTok{     AIN1\_TOKCAP}
\NormalTok{        STA     TOK1}
        \ControlFlowTok{JMP}\NormalTok{     AIN1\_DONE}
\NormalTok{AIN1\_TOKCAP}
\NormalTok{        STA     TOK1            }\OperatorTok{*} \OperatorTok{(}\NormalTok{A ≥ CAP → store A anyway}\CommentTok{; CAP1 already bounds)}
\NormalTok{AIN1\_DONE}
        \ControlFlowTok{JMP}\NormalTok{     AIN1\_RET}

\NormalTok{AIN1\_BLOCK}
\NormalTok{        * Signal need }\OperatorTok{and}\NormalTok{ invoke scheduler}
\NormalTok{        STA     NEED1           }\OperatorTok{*}\NormalTok{ store nonzero to NEED1}
\NormalTok{        JSR     B\_POOL}

\NormalTok{        * Busy}\OperatorTok{{-}}\NormalTok{wait until GHD1 nonempty }\OperatorTok{(}\NormalTok{you may also yield via CHECK\_IO}\OperatorTok{)}
\NormalTok{AIN1\_WAIT}
\NormalTok{        LDA     GHD1}
\NormalTok{        JAZ     AIN1\_SPIN}
        \ControlFlowTok{JMP}\NormalTok{     A\_IN1           }\OperatorTok{*}\NormalTok{ try again}
\NormalTok{AIN1\_SPIN}
\NormalTok{        JSR     CHECK\_IO        }\OperatorTok{*}\NormalTok{ poll devices }\OperatorTok{and}\NormalTok{ reschedule}
        \ControlFlowTok{JMP}\NormalTok{     AIN1\_WAIT}

\NormalTok{AIN1\_RET NOP}
\NormalTok{ONE     CON     }\DecValTok{1}
\end{Highlighting}
\end{Shaded}

5.2 A\_pool.OUT(o)

\begin{Shaded}
\begin{Highlighting}[]
\NormalTok{* A\_OUT}\OperatorTok{(}\NormalTok{o}\OperatorTok{):}\NormalTok{ take FREE }\OperatorTok{or}\NormalTok{ block}\CommentTok{; example for o = O1}
\NormalTok{A\_OUT1  STJ     AOUT1\_RET}
\NormalTok{AOUT1\_TRY}
\NormalTok{        LDA     }\OperatorTok{=}\NormalTok{FREEHD}\OperatorTok{=}
\NormalTok{        JSR     DEQ}
\NormalTok{        JXZ     AOUT1\_BLOCK}
\NormalTok{        * X}\OperatorTok{=}\NormalTok{b}\CommentTok{; STATE[b] := YELLOW; DEV[b] := O1}
\NormalTok{        LDA     }\OperatorTok{=}\NormalTok{BDST1}\OperatorTok{{-}}\DecValTok{1}\OperatorTok{=,}\NormalTok{X}
\NormalTok{        ENTA    STATE\_YEL}
\NormalTok{        STA     }\DecValTok{0}\OperatorTok{,}\NormalTok{X}
\NormalTok{        LDA     }\OperatorTok{=}\NormalTok{BDDEV1}\OperatorTok{{-}}\DecValTok{1}\OperatorTok{=,}\NormalTok{X}
\NormalTok{        ENTA    O1}
\NormalTok{        STA     }\DecValTok{0}\OperatorTok{,}\NormalTok{X}
        \ControlFlowTok{JMP}\NormalTok{     AOUT1\_RET}
\NormalTok{AOUT1\_BLOCK}
\NormalTok{        JSR     B\_POOL}
\NormalTok{AOUT1\_WAIT}
\NormalTok{        LDA     FREEHD}
\NormalTok{        JAZ     AOUT1\_SPIN}
        \ControlFlowTok{JMP}\NormalTok{     A\_OUT1}
\NormalTok{AOUT1\_SPIN}
\NormalTok{        JSR     CHECK\_IO}
        \ControlFlowTok{JMP}\NormalTok{     AOUT1\_WAIT}
\NormalTok{AOUT1\_RET NOP}
\end{Highlighting}
\end{Shaded}

5.3 R\_pool.RELEASE

\begin{Shaded}
\begin{Highlighting}[]
\NormalTok{* RELEASE for an input buffer }\OperatorTok{(}\NormalTok{consumed}\OperatorTok{):}\NormalTok{ return to FREE}
\NormalTok{REL\_IN  STJ     RELI\_RET}
\NormalTok{        * }\BuiltInTok{IN}\OperatorTok{:}\NormalTok{ X }\OperatorTok{=}\NormalTok{ b}
\NormalTok{        * STATE}\OperatorTok{[}\NormalTok{b}\OperatorTok{]} \OperatorTok{:=}\NormalTok{ FREE}\CommentTok{; ENQ b {-}\textgreater{} FREEHD}
\NormalTok{        LDA     }\OperatorTok{=}\NormalTok{BDST1}\OperatorTok{{-}}\DecValTok{1}\OperatorTok{=,}\NormalTok{X}
\NormalTok{        ENTA    STATE\_FREE}
\NormalTok{        STA     }\DecValTok{0}\OperatorTok{,}\NormalTok{X}
\NormalTok{        ENTA    }\OperatorTok{=}\NormalTok{FREEHD}\OperatorTok{=}
\NormalTok{        JSR     ENQ}
        \ControlFlowTok{JMP}\NormalTok{     RELI\_RET}
\NormalTok{RELI\_RET NOP}

\NormalTok{* RELEASE for an output buffer }\OperatorTok{(}\NormalTok{produced}\OperatorTok{):}\NormalTok{ enqueue RED to R}\OperatorTok{[}\NormalTok{o}\OperatorTok{]}
\NormalTok{* Example for o}\OperatorTok{=}\NormalTok{O1}\OperatorTok{,}\NormalTok{ X}\OperatorTok{=}\NormalTok{b}
\NormalTok{REL\_OUT1 STJ    RELO1\_RET}
\NormalTok{        LDA     }\OperatorTok{=}\NormalTok{BDST1}\OperatorTok{{-}}\DecValTok{1}\OperatorTok{=,}\NormalTok{X}
\NormalTok{        ENTA    STATE\_RED}
\NormalTok{        STA     }\DecValTok{0}\OperatorTok{,}\NormalTok{X}
\NormalTok{        LDA     }\OperatorTok{=}\NormalTok{BDDEV1}\OperatorTok{{-}}\DecValTok{1}\OperatorTok{=,}\NormalTok{X}
\NormalTok{        ENTA    O1}
\NormalTok{        STA     }\DecValTok{0}\OperatorTok{,}\NormalTok{X}
\NormalTok{        ENTA    }\OperatorTok{=}\NormalTok{RHO1}\OperatorTok{=}
\NormalTok{        JSR     ENQ}
\NormalTok{        * hint scheduler}
\NormalTok{        JSR     B\_POOL}
        \ControlFlowTok{JMP}\NormalTok{     RELO1\_RET}
\NormalTok{RELO1\_RET NOP}
\end{Highlighting}
\end{Shaded}

\begin{enumerate}
\def\labelenumi{\arabic{enumi}.}
\setcounter{enumi}{5}
\tightlist
\item
  Scheduler (B\_pool)
\end{enumerate}

Drain outputs → essential inputs (to unblock compute) → anticipatory
inputs (tokens + reserve).

\begin{Shaded}
\begin{Highlighting}[]
\NormalTok{B\_POOL  STJ     BRET}

\NormalTok{        * {-}{-}{-}{-}{-} B1}\OperatorTok{:}\NormalTok{ drain outputs first }\OperatorTok{{-}{-}{-}{-}{-}}
\NormalTok{        * }\PreprocessorTok{if}\NormalTok{ O1 idle }\OperatorTok{and}\NormalTok{ RHO1 nonempty }\OperatorTok{and}\NormalTok{ no inflight out}\OperatorTok{:}\NormalTok{ start OUT}\OperatorTok{(}\NormalTok{O1}\OperatorTok{)}
\NormalTok{B1\_O1   LDA     RHO1}
\NormalTok{        JAZ     B1\_O2           }\OperatorTok{*}\NormalTok{ nothing ready}
\NormalTok{        LDA     INFLO1}
\NormalTok{        JAZ     B1\_O1\_TRY       }\OperatorTok{*}\NormalTok{ none inflight}
        \ControlFlowTok{JMP}\NormalTok{     B1\_O2           }\OperatorTok{*}\NormalTok{ already busy for O1}
\NormalTok{B1\_O1\_TRY}
\NormalTok{        ENTA    }\OperatorTok{=}\NormalTok{RHO1}\OperatorTok{=}
\NormalTok{        JSR     DEQ             }\OperatorTok{*}\NormalTok{ X}\OperatorTok{=}\NormalTok{b}
\NormalTok{        JXZ     B1\_O2}
\NormalTok{        * mark inflight }\OperatorTok{and}\NormalTok{ start OUT}
\NormalTok{        STA     INFLO1          }\OperatorTok{*}\NormalTok{ INFLO1 }\OperatorTok{:=}\NormalTok{ b}
\NormalTok{        LDA     }\OperatorTok{=}\NormalTok{BDST1}\OperatorTok{{-}}\DecValTok{1}\OperatorTok{=,}\NormalTok{X}
\NormalTok{        ENTA    STATE\_INFLO}
\NormalTok{        STA     }\DecValTok{0}\OperatorTok{,}\NormalTok{X}
\NormalTok{        * }\BuiltInTok{OUT}\NormalTok{ mem}\OperatorTok{,}\NormalTok{ O1}
\NormalTok{        LDA     }\OperatorTok{=}\NormalTok{BDADR1}\OperatorTok{{-}}\DecValTok{1}\OperatorTok{=,}\NormalTok{X}
        \BuiltInTok{OUT}     \DecValTok{0}\OperatorTok{(}\NormalTok{O1}\OperatorTok{)}           \OperatorTok{*} \OperatorTok{(}\NormalTok{syntax varies by assembler}\CommentTok{; concept: OUT at BDADR[b] to O1)}
\NormalTok{B1\_O2}
\NormalTok{        * }\PreprocessorTok{repeat}\NormalTok{ for O2 similarly}\OperatorTok{...}
\NormalTok{        * (}\PreprocessorTok{macro}\NormalTok{ in real code}\OperatorTok{)}
        
\NormalTok{        * {-}{-}{-}{-}{-} B2}\OperatorTok{:}\NormalTok{ essential inputs }\OperatorTok{(}\NormalTok{unblock compute}\OperatorTok{)} \OperatorTok{{-}{-}{-}{-}{-}}
\NormalTok{        * For each d with NEED}\OperatorTok{[}\NormalTok{d}\OperatorTok{]!=}\DecValTok{0} \OperatorTok{and}\NormalTok{ INFLI}\OperatorTok{[}\NormalTok{d}\OperatorTok{]==}\DecValTok{0} \OperatorTok{and}\NormalTok{ FREE available }\OperatorTok{(}\NormalTok{even if count }\OperatorTok{\textless{}=}\NormalTok{ RESV}\OperatorTok{)}
\NormalTok{B2\_D1   LDA     NEED1}
\NormalTok{        JAZ     B2\_D2}
\NormalTok{        LDA     INFLI1}
        \ControlFlowTok{JNZ}\NormalTok{     B2\_D2}
\NormalTok{        * Try FREE}
\NormalTok{        LDA     FREEHD}
\NormalTok{        JAZ     B2\_D2           }\OperatorTok{*}\NormalTok{ truly no free buffers}\CommentTok{; cannot help now}
\NormalTok{        ENTA    }\OperatorTok{=}\NormalTok{FREEHD}\OperatorTok{=}
\NormalTok{        JSR     DEQ             }\OperatorTok{*}\NormalTok{ X}\OperatorTok{=}\NormalTok{b}
\NormalTok{        JXZ     B2\_D2}
\NormalTok{        * Start IN on D1 without consuming token }\OperatorTok{(}\NormalTok{exception}\OperatorTok{)}
\NormalTok{        STA     INFLI1}
\NormalTok{        LDA     }\OperatorTok{=}\NormalTok{BDDEV1}\OperatorTok{{-}}\DecValTok{1}\OperatorTok{=,}\NormalTok{X}
\NormalTok{        ENTA    D1}
\NormalTok{        STA     }\DecValTok{0}\OperatorTok{,}\NormalTok{X}
\NormalTok{        LDA     }\OperatorTok{=}\NormalTok{BDST1}\OperatorTok{{-}}\DecValTok{1}\OperatorTok{=,}\NormalTok{X}
\NormalTok{        ENTA    STATE\_INFLI}
\NormalTok{        STA     }\DecValTok{0}\OperatorTok{,}\NormalTok{X}
\NormalTok{        * }\BuiltInTok{IN}\NormalTok{ mem}\OperatorTok{,}\NormalTok{ D1}
\NormalTok{        LDA     }\OperatorTok{=}\NormalTok{BDADR1}\OperatorTok{{-}}\DecValTok{1}\OperatorTok{=,}\NormalTok{X}
        \BuiltInTok{IN}      \DecValTok{0}\OperatorTok{(}\NormalTok{D1}\OperatorTok{)}
\NormalTok{        * clear NEED1}
\NormalTok{        ENTA    }\DecValTok{0}
\NormalTok{        STA     NEED1}
\NormalTok{B2\_D2}
\NormalTok{        * }\PreprocessorTok{Repeat}\NormalTok{ similarly for D2}\OperatorTok{,}\NormalTok{ D3 }\OperatorTok{...}

\NormalTok{        * {-}{-}{-}{-}{-} B3}\OperatorTok{:}\NormalTok{ conservative anticipation }\OperatorTok{{-}{-}{-}{-}{-}}
\NormalTok{        * Only if FREE count }\OperatorTok{\textgreater{}}\NormalTok{ RESV }\OperatorTok{and}\NormalTok{ TOK}\OperatorTok{[}\NormalTok{d}\OperatorTok{]} \OperatorTok{\textgreater{}} \DecValTok{0} \OperatorTok{and}\NormalTok{ no inflight}
\NormalTok{        * (Here we approximate }\StringTok{"FREE count"}\NormalTok{ with head }\OperatorTok{!=}\NormalTok{ NIL }\OperatorTok{and}\NormalTok{ a spillover test}\OperatorTok{,}
\NormalTok{        *  }\BuiltInTok{or}\NormalTok{ keep a FREE\_COUNT counter updated by ENQ}\OperatorTok{/}\NormalTok{DEQ in real code}\OperatorTok{.)}

\NormalTok{        * We}\StringTok{\textquotesingle{}ll implement a minimal guard: require FREEHD != NIL and also}
\NormalTok{        * require a second FREE to exist }\OperatorTok{(}\NormalTok{head}\OperatorTok{.}\NormalTok{next }\OperatorTok{!=}\NormalTok{ NIL}\OperatorTok{)}\NormalTok{ when RESV}\OperatorTok{=}\DecValTok{1}\OperatorTok{,}\NormalTok{ etc}\OperatorTok{.}
\NormalTok{        * For clarity}\OperatorTok{,}\NormalTok{ use FREE\_COUNT }\OperatorTok{(}\NormalTok{recommended}\OperatorTok{).}

\NormalTok{        * Example for D1}\OperatorTok{:}
\NormalTok{B3\_D1   LDA     INFLI1}
        \ControlFlowTok{JNZ}\NormalTok{     B3\_D2}
\NormalTok{        LDA     TOK1}
\NormalTok{        JAZ     B3\_D2}
\NormalTok{        LDA     FREE\_COUNT}
\NormalTok{        CMPA    RESV}
        \ControlFlowTok{JLE}\NormalTok{     B3\_D2}
\NormalTok{        * OK}\OperatorTok{:}\NormalTok{ schedule anticipatory IN for D1}\CommentTok{; consume one token}
\NormalTok{        ENTA    }\OperatorTok{=}\NormalTok{FREEHD}\OperatorTok{=}
\NormalTok{        JSR     DEQ             }\OperatorTok{*}\NormalTok{ X}\OperatorTok{=}\NormalTok{b }\OperatorTok{(}\NormalTok{must exist by guard}\OperatorTok{)}
\NormalTok{        JXZ     B3\_D2}
\NormalTok{        STA     INFLI1}
\NormalTok{        LDA     }\OperatorTok{=}\NormalTok{BDDEV1}\OperatorTok{{-}}\DecValTok{1}\OperatorTok{=,}\NormalTok{X}
\NormalTok{        ENTA    D1}
\NormalTok{        STA     }\DecValTok{0}\OperatorTok{,}\NormalTok{X}
\NormalTok{        LDA     }\OperatorTok{=}\NormalTok{BDST1}\OperatorTok{{-}}\DecValTok{1}\OperatorTok{=,}\NormalTok{X}
\NormalTok{        ENTA    STATE\_INFLI}
\NormalTok{        STA     }\DecValTok{0}\OperatorTok{,}\NormalTok{X}
\NormalTok{        * }\BuiltInTok{IN}\NormalTok{ mem}\OperatorTok{,}\NormalTok{ D1}
\NormalTok{        LDA     }\OperatorTok{=}\NormalTok{BDADR1}\OperatorTok{{-}}\DecValTok{1}\OperatorTok{=,}\NormalTok{X}
        \BuiltInTok{IN}      \DecValTok{0}\OperatorTok{(}\NormalTok{D1}\OperatorTok{)}
\NormalTok{        * TOK1 }\OperatorTok{:=}\NormalTok{ TOK1 }\OperatorTok{{-}} \DecValTok{1}
\NormalTok{        LDA     TOK1}
        \BuiltInTok{SUB}\NormalTok{     ONE}
\NormalTok{        STA     TOK1}
\NormalTok{B3\_D2}
\NormalTok{        * }\PreprocessorTok{Repeat}\NormalTok{ for D2}\OperatorTok{,}\NormalTok{ D3 }\OperatorTok{...}

        \ControlFlowTok{JMP}\NormalTok{     BRET}
\NormalTok{BRET    NOP}
\end{Highlighting}
\end{Shaded}

\begin{enumerate}
\def\labelenumi{\arabic{enumi}.}
\setcounter{enumi}{6}
\tightlist
\item
  I/O completion polling (CHECK\_IO)
\end{enumerate}

\texttt{IN}/\texttt{OUT} start asynchronously; we detect completion with
\texttt{JBUS} (``jump if busy''). When \texttt{JBUS} falls through, the
device finished. We then move the buffer to the right queue and clear
\texttt{INFLI?/INFLO?}.

\begin{Shaded}
\begin{Highlighting}[]
\NormalTok{CHECK\_IO STJ    CIRET}

\NormalTok{        * {-}{-}{-}{-}{-} outputs}\OperatorTok{:}\NormalTok{ when OUT completes}\OperatorTok{,}\NormalTok{ return buffer to FREE }\OperatorTok{{-}{-}{-}{-}{-}}
\NormalTok{C\_O1    JBUS    C\_O1\_BUSY}\OperatorTok{,}\NormalTok{O1    }\OperatorTok{*}\NormalTok{ jump if O1 still busy}
\NormalTok{        * Completed}
\NormalTok{        LDA     INFLO1}
\NormalTok{        JAZ     C\_O2            }\OperatorTok{*}\NormalTok{ no inflight }\OperatorTok{{-}\textgreater{}}\NormalTok{ nothing to do}
\NormalTok{        ENTX    A               }\OperatorTok{*}\NormalTok{ X}\OperatorTok{=}\NormalTok{b}
\NormalTok{        ENTA    }\DecValTok{0}
\NormalTok{        STA     INFLO1}
\NormalTok{        * STATE}\OperatorTok{[}\NormalTok{b}\OperatorTok{]} \OperatorTok{:=}\NormalTok{ FREE}\CommentTok{; ENQ FREE}
\NormalTok{        LDA     }\OperatorTok{=}\NormalTok{BDST1}\OperatorTok{{-}}\DecValTok{1}\OperatorTok{=,}\NormalTok{X}
\NormalTok{        ENTA    STATE\_FREE}
\NormalTok{        STA     }\DecValTok{0}\OperatorTok{,}\NormalTok{X}
\NormalTok{        ENTA    }\OperatorTok{=}\NormalTok{FREEHD}\OperatorTok{=}
\NormalTok{        JSR     ENQ}
\NormalTok{        * FREE\_COUNT}\OperatorTok{++}
\NormalTok{        LDA     FREE\_COUNT}
        \BuiltInTok{ADD}\NormalTok{     ONE}
\NormalTok{        STA     FREE\_COUNT}
        \ControlFlowTok{JMP}\NormalTok{     C\_O2}
\NormalTok{C\_O1\_BUSY}
        \BuiltInTok{NOP}
\NormalTok{C\_O2    }\OperatorTok{*}\NormalTok{ repeat for O2 }\OperatorTok{...}

\NormalTok{        * {-}{-}{-}{-}{-} inputs}\OperatorTok{:}\NormalTok{ when IN completes}\OperatorTok{,}\NormalTok{ put buffer GREEN on G}\OperatorTok{[}\NormalTok{d}\OperatorTok{]} \OperatorTok{{-}{-}{-}{-}{-}}
\NormalTok{C\_I1    JBUS    C\_I1\_BUSY}\OperatorTok{,}\NormalTok{D1}
\NormalTok{        LDA     INFLI1}
\NormalTok{        JAZ     C\_I2}
\NormalTok{        ENTX    A               }\OperatorTok{*}\NormalTok{ X}\OperatorTok{=}\NormalTok{b}
\NormalTok{        ENTA    }\DecValTok{0}
\NormalTok{        STA     INFLI1}
\NormalTok{        * STATE}\OperatorTok{[}\NormalTok{b}\OperatorTok{]} \OperatorTok{:=}\NormalTok{ GREEN}\CommentTok{; DEV[b]:=D1; ENQ to GHD1}
\NormalTok{        LDA     }\OperatorTok{=}\NormalTok{BDST1}\OperatorTok{{-}}\DecValTok{1}\OperatorTok{=,}\NormalTok{X}
\NormalTok{        ENTA    STATE\_GREEN}
\NormalTok{        STA     }\DecValTok{0}\OperatorTok{,}\NormalTok{X}
\NormalTok{        LDA     }\OperatorTok{=}\NormalTok{BDDEV1}\OperatorTok{{-}}\DecValTok{1}\OperatorTok{=,}\NormalTok{X}
\NormalTok{        ENTA    D1}
\NormalTok{        STA     }\DecValTok{0}\OperatorTok{,}\NormalTok{X}
\NormalTok{        ENTA    }\OperatorTok{=}\NormalTok{GHD1}\OperatorTok{=}
\NormalTok{        JSR     ENQ}
        \ControlFlowTok{JMP}\NormalTok{     C\_I2}
\NormalTok{C\_I1\_BUSY}
        \BuiltInTok{NOP}
\NormalTok{C\_I2    }\OperatorTok{*}\NormalTok{ repeat for D2}\OperatorTok{,}\NormalTok{ D3 }\OperatorTok{...}

\NormalTok{        * Optional}\OperatorTok{:}\NormalTok{ immediately run B\_POOL to kick off drain}\OperatorTok{/}\NormalTok{anticipation}
\NormalTok{        JSR     B\_POOL}
        \ControlFlowTok{JMP}\NormalTok{     CIRET}
\NormalTok{CIRET   NOP}
\end{Highlighting}
\end{Shaded}

\begin{enumerate}
\def\labelenumi{\arabic{enumi}.}
\setcounter{enumi}{7}
\tightlist
\item
  FREE\_COUNT maintenance (recommended)
\end{enumerate}

Counting FREE precisely makes the \textbf{reserve} rule simple and
efficient.

\begin{Shaded}
\begin{Highlighting}[]
\NormalTok{FREE\_COUNT  CON  B           }\OperatorTok{*}\NormalTok{ initialized to B at startup}

\NormalTok{* On taking from FREE }\OperatorTok{(}\NormalTok{DEQ success}\OperatorTok{):}\NormalTok{ FREE\_COUNT}\OperatorTok{{-}{-}}
\NormalTok{* On returning to FREE }\OperatorTok{(}\NormalTok{ENQ}\OperatorTok{):}\NormalTok{        FREE\_COUNT}\OperatorTok{++}

\NormalTok{* You can integrate these bumps directly around the DEQ}\OperatorTok{/}\NormalTok{ENQ calls that touch FREE}\OperatorTok{.}
\end{Highlighting}
\end{Shaded}

\begin{enumerate}
\def\labelenumi{\arabic{enumi}.}
\setcounter{enumi}{8}
\tightlist
\item
  Putting it together (call pattern)
\end{enumerate}

\begin{itemize}
\item
  The \textbf{compute} routine calls \texttt{A\_IN?} to obtain a YELLOW
  buffer (filled input) or \texttt{A\_OUT?} to obtain a YELLOW buffer
  for output space. It then does work and calls \texttt{REL\_IN} or
  \texttt{REL\_OUT?} accordingly.
\item
  \texttt{A\_*} calls may block, which triggers \texttt{B\_POOL} and
  \texttt{CHECK\_IO} until the requested list becomes nonempty.
\item
  \texttt{B\_POOL} is also called:

  \begin{itemize}
  \tightlist
  \item
    right after \texttt{REL\_OUT?} (to drain RED → device),
  \item
    from \texttt{CHECK\_IO} when completions happen,
  \item
    and opportunistically anywhere you'd like to make the system
    responsive.
  \end{itemize}
\end{itemize}

\begin{enumerate}
\def\labelenumi{\arabic{enumi}.}
\setcounter{enumi}{9}
\tightlist
\item
  Policy knobs (the ``anti-flood'' part)
\end{enumerate}

\begin{itemize}
\tightlist
\item
  \textbf{Token buckets:} \texttt{TOK{[}d{]}} minted \textbf{only} when
  compute actually consumes a GREEN buffer (in \texttt{A\_IN?}, after
  successful dequeue). One token is \textbf{spent} when scheduling an
  anticipatory \texttt{IN} in B3. Essential reads in B2 don't spend
  tokens (that's the ``exception when blocked''). Keep
  \texttt{CAP{[}d{]}} small (1--2).
\item
  \textbf{Reserve:} \texttt{RESV\ =\ \#outputs} is a good default;
  increase if outputs can bunch up or are slow. The
  \texttt{FREE\_COUNT\ \textgreater{}\ RESV} guard must hold before B3
  starts more anticipation.
\item
  \textbf{Output priority:} Encoded by B1 running before B2/B3. This
  prevents the pool from filling with GREEN while outputs starve.
\end{itemize}

\begin{enumerate}
\def\labelenumi{\arabic{enumi}.}
\setcounter{enumi}{10}
\tightlist
\item
  Notes for a real MIX assembler/emulator
\end{enumerate}

\begin{itemize}
\tightlist
\item
  The snippets use ``address arithmetic'' placeholders like
  \texttt{=BDST1-1=,X}. In a real MIXAL, define macros
  \texttt{BDNXT(i)}, \texttt{BDST(i)}, \texttt{BDDEV(i)},
  \texttt{BDADR(i)} that compute addresses from \texttt{i} (e.g., via
  \texttt{EQU} tables or by laying out BD fields in contiguous arrays).
\item
  The \texttt{IN\ 0(D)} / \texttt{OUT\ 0(D)} lines indicate ``start I/O
  at address in A to device D''. Adapt to your assembler's exact operand
  format; many MIX tools accept \texttt{IN\ A,D}.
\item
  If your emulator blocks on \texttt{IN/OUT} (non-asynchronous),
  simulate async by splitting start/complete as shown (or just keep the
  same control logic; \texttt{JBUS} tests become no-ops).
\end{itemize}

\section{18 {[}30{]}}\label{section-194}

Below is a clean, interrupt-driven rewrite of Knuth's multiple-buffer
scheme (Algorithms/Programs A, R, B in §1.4.4) for MIX, using the
proposed control-state ``interrupt'' extension described right after
Exercise 18. I first restate the interrupt model very briefly, then give
the modified algorithms and a compact MIX sketch, and finally explain
why it works and how it maps to the original buffering logic.

Interrupt model (per the exercise)

\begin{itemize}
\tightlist
\item
  MIX gains control memory locations −3999..−0001 and two internal
  states: \textbf{normal} and \textbf{control}. On an interrupt,
  registers rA, rI1..rI6, rX, and (rJ, OV, CI, ``next-inst'') are saved
  into −0009..−0001, and MIX \textbf{enters control state} at an address
  determined by the interrupt type. There is also a ``clock'' at −0010
  that ticks each 1000u and, on zero, interrupts to −0011. The
  instruction \textbf{INT}: in normal state, forces an interrupt to
  −0012; in control state, it \textbf{restores} registers from
  −0009..−0001 and returns to normal state. I/O completion interrupts
  vector to −(0020+unit). (All as specified in the exercise statement.) 
\end{itemize}

Variables and notation

We keep Knuth's §1.4.4 notation:

\begin{itemize}
\tightlist
\item
  \texttt{N} = total buffers in the circle.
\item
  \texttt{n} = number of \emph{green} buffers (assignable).
\item
  \texttt{NEXTG}, \texttt{NEXTR}, \texttt{CURRENT} = ring pointers as in
  the text.
\item
  \textbf{New:} \texttt{t} = 0 if the device is idle, 1 if a transfer is
  in progress.
\item
  Device unit number is \texttt{U}.
\end{itemize}

(Initial conditions are as in §1.4.4: for input, start with all buffers
green; for output, with all buffers red; and \texttt{NEXTR\ =\ NEXTG} at
start. See the section text around Algorithms A/R/B.) 

The modified algorithms

These are exactly the changes needed to eliminate any polling/JRED and
run purely interrupt-driven. The outline matches Knuth's official answer
(``Partial solution'') but with full context and linkage.

A′ --- ASSIGN (normal state, called by the compute routine)

\textbf{Unchanged} from Algorithm A in the text:

\begin{itemize}
\tightlist
\item
  \textbf{A1.} Wait for \texttt{n\ \textless{}\ N}.
\item
  \textbf{A2.} \texttt{CURRENT\ ←\ NEXTG}.
\item
  \textbf{A3.} Advance \texttt{NEXTG} (clockwise).
\end{itemize}

(Identical to §1.4.4A.) 

R′ --- RELEASE (normal state, inlined where the compute routine is done
with the current buffer)

\begin{itemize}
\tightlist
\item
  \textbf{R1′.} Increase \texttt{n} by 1.
\item
  \textbf{R2′.} If \texttt{t\ =\ 0} (device idle), \textbf{force an
  interrupt} with \texttt{INT} so that control jumps immediately to the
  control routine's step \textbf{B3′} (see below).
\end{itemize}

This is exactly the little ``nudge'' we need so that when the CPU frees
a buffer and the device happens to be idle, the control routine can
immediately kick off the next I/O without delay. 

B′ --- BUFFER CONTROL (runs in control state; it \emph{processes
interrupts})

Think of B′ as the \emph{interrupt handler framework} for buffering; it
is entered:

\begin{itemize}
\tightlist
\item
  on a software interrupt from \textbf{R2′} (vector −0012),
\item
  or on an I/O completion interrupt (vector −(0020+U)).
\end{itemize}

It uses \texttt{t} to remember whether the device is currently busy.

\begin{itemize}
\tightlist
\item
  \textbf{B1′.} \textbf{Restart the main program} (execute \texttt{INT})
  if there's nothing useful to do right now.
\item
  \textbf{B2′.} If \texttt{n\ =\ 0}, set \texttt{t\ ←\ 0} and \textbf{go
  to B1′} (no buffers to service, stay idle and let the CPU compute).
\item
  \textbf{B3′.} Set \texttt{t\ ←\ 1} and \textbf{initiate I/O} on unit
  \texttt{U} using the buffer pointed to by \texttt{NEXTR} (e.g.,
  \texttt{IN} for input flow, \texttt{OUT} for output flow).
\item
  \textbf{B4′.} \textbf{Restart the main program} (\texttt{INT}). An I/O
  completion will later interrupt and re-enter control state.
\item
  \textbf{B5′.} On the \emph{I/O completion interrupt}, advance
  \texttt{NEXTR} to the next buffer in the ring.
\item
  \textbf{B6′.} Decrease \texttt{n} by 1 (we've consumed a green buffer
  for input, or emptied a red buffer for output), and \textbf{go to
  B2′}.
\end{itemize}

This is the full interrupt-driven analogue of the original B; no
busy-waiting or \texttt{JRED} appears anywhere. 

Minimal MIX skeleton (how to wire it)

Below is a compact sketch to show where the control-state entry points
jump. It's deliberately schematic (addresses/labels are illustrative),
but it shows the exact linkage you can adopt in your own code.

\begin{Shaded}
\begin{Highlighting}[]
\NormalTok{* {-}{-}{-}{-}{-}{-}{-}{-}{-}{-}{-}{-}{-}{-}{-}{-}{-}{-}{-}{-}{-}{-}{-}{-}{-}{-}}
\NormalTok{* Data (normal state)}
\NormalTok{N       CON   ...          * total buffers}
\NormalTok{n       CON   0            * current green count}
\NormalTok{t       CON   0            * 0 = idle, 1 = busy}
\NormalTok{NEXTG   CON   ...          * ring pointers per §1.4.4}
\NormalTok{NEXTR   CON   ...}
\NormalTok{CURRENT CON   ...}

\NormalTok{* {-}{-}{-}{-}{-}{-}{-}{-}{-}{-}{-}{-}{-}{-}{-}{-}{-}{-}{-}{-}{-}{-}{-}{-}{-}{-}}
\NormalTok{* R′ (normal state, inline macro)}
\NormalTok{Rprime: INC   n,1          * R1′}
\NormalTok{        LDA   t}
\NormalTok{        JAZ   ForceINT     * if t = 0 {-}\textgreater{} nudge the control routine}
\NormalTok{        JMP   *+1}
\NormalTok{ForceINT:}
\NormalTok{        INT                * software interrupt {-}\textgreater{} control entry at −0012}

\NormalTok{* A′ (normal state): identical to Algorithm A in text (A1–A3)}
\NormalTok{* ... your A code from §1.4.4 ...}

\NormalTok{* {-}{-}{-}{-}{-}{-}{-}{-}{-}{-}{-}{-}{-}{-}{-}{-}{-}{-}{-}{-}{-}{-}{-}{-}{-}{-}}
\NormalTok{* CONTROL STATE VECTORS}
\NormalTok{* (hardware jumps here on different interrupts)}

\NormalTok{ORG    {-}12                 * −0012: software INT from R2′}
\NormalTok{INT\_SW: JMP   B3prime}

\NormalTok{ORG    {-}(20+U)             * I/O completion for device U}
\NormalTok{INT\_IO: JMP   B5prime}

\NormalTok{* (Clock −0011 not needed here, but could JMP B1′ if you like.)}

\NormalTok{* {-}{-}{-}{-}{-}{-}{-}{-}{-}{-}{-}{-}{-}{-}{-}{-}{-}{-}{-}{-}{-}{-}{-}{-}{-}{-}}
\NormalTok{* B′ control routine (control state)}
\NormalTok{B1prime:}
\NormalTok{        INT                * Return to normal state immediately when idle}

\NormalTok{B2prime:}
\NormalTok{        LDA   n}
\NormalTok{        JAZ   Idle         * if n = 0: idle the device}
\NormalTok{        JMP   B3prime}

\NormalTok{Idle:   STA   t            * store 0 (accumulator assumed 0 from JAZ path)}
\NormalTok{        JMP   B1prime}

\NormalTok{B3prime:}
\NormalTok{        ENTA  1}
\NormalTok{        STA   t            * t ← 1 (device busy)}
\NormalTok{        * Initiate I/O on buffer at NEXTR:}
\NormalTok{        *   For input:  IN  BUF(NEXTR)(U)}
\NormalTok{        *   For output: OUT BUF(NEXTR)(U)}
\NormalTok{        * (Addressing of BUF(NEXTR) per your ring layout.)}
\NormalTok{        INT                * B4′: restart main program; completion will re{-}enter at INT\_IO}

\NormalTok{B5prime:}
\NormalTok{        * advance NEXTR one buffer clockwise}
\NormalTok{        * NEXTR ← next(NEXTR)}
\NormalTok{        DEC   n,1          * B6′: n ← n − 1}
\NormalTok{        JMP   B2prime}
\end{Highlighting}
\end{Shaded}

(You'll splice in the actual \texttt{IN}/\texttt{OUT} origin fields to
address the buffer at \texttt{NEXTR}, exactly as you did in your
non-interrupt version.)

Why this is correct

\begin{itemize}
\item
  \textbf{No polling:} The original B used a ``wait until device ready''
  flavor (conceptually via \texttt{JBUS}/\texttt{JRED}). Here we
  \emph{never} poll. The control logic either restarts the main program
  (B1′/B4′) or is re-entered \textbf{only} when something interesting
  happens (a release when idle, or I/O completion). This directly
  matches the exercise's intent. 
\item
  \textbf{Immediate kick-off when possible:} If the CPU frees a buffer
  while the device is idle (\texttt{t=0}), \textbf{R2′} forces a
  software interrupt, which vectors straight to \textbf{B3′} and starts
  I/O at once. This eliminates latent delays that might otherwise occur
  before the next timer/clock tick. 
\item
  \textbf{Invariant preservation:} The same invariants from §1.4.4 hold:

  \begin{itemize}
  \tightlist
  \item
    \texttt{n} counts available (green) buffers; is incremented only by
    RELEASE (R1′) and decremented only when a transfer is actually begun
    (B6′).
  \item
    \texttt{t} tracks device occupancy; set in B3′ and implicitly
    cleared by the \textbf{path} B2′ takes when \texttt{n=0} (and by the
    fact that once completion fires, the next cycle will either start
    another I/O or idle again).
  \item
    \texttt{NEXTR} advances \textbf{only} upon I/O completion (B5′),
    exactly mirroring the original control logic. 
  \end{itemize}
\item
  \textbf{State transitions match hardware:} \texttt{INT} in control
  state restores registers and resumes normal execution (B1′, B4′), and
  I/O completion vectors to −(0020+U) (B5′ path). That's precisely how
  the exercise defines INT and the I/O interrupts. 
\end{itemize}

Mapping to Knuth's official answer

Knuth's answer gives the same three changes concisely:

\begin{itemize}
\tightlist
\item
  ``In the algorithms below, \textbf{t} is 0 when idle, 1 when active.''
\item
  \textbf{A′} unchanged.
\item
  \textbf{R′:} increment \texttt{n}; if \texttt{t=0} force an interrupt
  to B3′.
\item
  \textbf{B′:} (i) restart; (ii) if \texttt{n=0}, set \texttt{t\ ←\ 0}
  and restart; (iii) set \texttt{t\ ←\ 1} and initiate I/O; (iv)
  restart; (v) on I/O complete, advance \texttt{NEXTR}; (vi)
  \texttt{n\ ←\ n\ −\ 1}; loop. 
\end{itemize}

Notes and practical tips

\begin{itemize}
\tightlist
\item
  If you use the ``green--yellow--red--purple'' combined input/output
  circle (Fig. 26) elsewhere in §1.4.4, the same interrupt framework
  applies; you just choose \texttt{IN} vs.~\texttt{OUT} in \textbf{B3′}
  based on which half (input or output) is being serviced at that
  moment. 
\item
  You do \textbf{not} need the clock interrupt (−0011) for correctness;
  it can be ignored or used for other housekeeping. The key events are
  \textbf{R2′}'s software INT and device completion interrupts. 
\end{itemize}

If you'd like, I can drop this into your specific MIX buffer layout
(addresses for \texttt{BUF1..BUFN}, ring macros for
\texttt{NEXTR/NEXTG}, etc.) and produce a runnable MIXAL file.

\section{19 {[}M28{]}}\label{m28-8}

We process \(n\ge 2\) consecutive blocks from a rotating device. Block
\(k\) can begin input only when its start passes the head, namely at
times

\[
t_k \equiv (k-1)L \pmod P,\qquad t_1=0,
\]

and it becomes available \(T\) time units later, i.e.~it may be
\textbf{assigned} at

\[
u_k\;\ge\; t_k+T,
\]

then takes \(C\) time units to compute, so it is \textbf{released} at

\[
v_k = u_k + C .
\]

Processing is sequential so \(u_k \ge v_{k-1}\) for \(k>1\). With only
\(N\) buffers we also need \(t_k \ge v_{k-N}\) for \(k>N\).

Question: what is the least \(N\) that lets us achieve the theoretical
optimum finishing time

\[
v_n^{\min}=T+C+(n-1)\max(L,C)?
\]

Give a general rule and then work it out for
\(L=1,\;P=100,\;T=0.5,\;n=100\) with several \(C\)'s.

Lower bound and target schedule

The lower bound arises because each additional block must wait at least
\(L\) time units (until its start again appears) and at least \(C\) time
units (to compute). Thus the best possible cadence is by \(\max(L,C)\),
with the pipeline start-up cost \(T+C\). We therefore try to \textbf{hit
the bound} by choosing

\[
v_k=\begin{cases}
T+C, & k=1,\\[2pt]
v_{k-1}+C,& C\ge L\quad\text{(compute-bound)},\\[2pt]
v_{k-1}+L,& C\le L\quad\text{(I/O-bound)}.
\end{cases}
\]

Case A --- I/O-bound (\(C\le L\))

Here we can target \(u_k=t_k+T=(k-1)L+T\) (use the \emph{earliest}
revolution), giving

\[
v_k=T+C+(k-1)L .
\]

To be feasible with \(N\) buffers we must not start input before a
buffer frees:

\[
t_k\;=\;(k-1)L\;\ge\;v_{k-N}=T+C+(k-N-1)L
\;\;\Longleftrightarrow\;\;
(N-1)L\;\ge\;T+C.
\]

Hence the \textbf{minimum number of buffers} is

\[
\boxed{\,N_{\min}=\left\lceil 1+\frac{T+C}{L}\right\rceil\,,\qquad C\le L.}
\]

Case B --- Compute-bound (\(C\ge L\))

Now we target \(u_k=v_{k-1}\) so \(v_k=T+kC\). Block \(k\) must have
finished input by \(v_{k-1}-T\). Among the allowed input start times for
block \(k\),

\[
\mathcal{T}_k=\{(k-1)L+mP \mid m\in\mathbb{Z}_{\ge0}\},
\]

let \(t_k\) be the \textbf{latest} \(\le v_{k-1}-T\). The lag we need to
``absorb'' with buffering is the backward distance

\[
\delta_k \stackrel{\text{def}}{=}\bigl( (v_{k-1}-T) - t_k \bigr)
            = \bigl((k-1)(C-L)\bigr)\bmod P \;\in [0,P).
\]

(That identity uses \(v_{k-1}-T=(k-1)C\).) With \(N\) buffers the
available ``window'' before compute touches the block is

\[
(v_{k-1}-T)-v_{k-N} = (N-1)C - T,
\]

so feasibility for \textbf{all} \(k\) requires

\[
\delta_k \le (N-1)C - T\quad\text{for }k=2,\dots,n.
\]

Let

\[
\Delta_n \;=\;\max_{1\le j\le n-1}\; \bigl(j(C-L)\bmod P\bigr),
\]

then a necessary and sufficient condition to attain the optimum is

\[
(N-1)C - T \;\ge\; \Delta_n .
\]

Therefore the \textbf{minimum number of buffers} is

\[
\boxed{\,N_{\min}=\left\lceil 1+\frac{T+\Delta_n}{C}\right\rceil\,,\qquad C\ge L,}
\]

with
\(\displaystyle \Delta_n=\max_{1\le j\le n-1}\!\bigl(j(C-L)-P\lfloor j(C-L)/P\rfloor\bigr)\).

Remarks.

\begin{itemize}
\tightlist
\item
  If \((C-L)/P\) is rational \(=p/q\) in lowest terms and \(n-1\ge q\),
  then \(\Delta_n=P-P/q\). If \(n\) is smaller, compute the max over
  \(j=1,\dots,n-1\) explicitly.
\item
  If \(C=L\), then \(\Delta_n=0\) and the formula gives
  \(N_{\min}=\lceil1+T/C\rceil\).
\end{itemize}

Worked examples (\(L=1,\;P=100,\;T=0.5,\;n=100\))

For each \(C\) we evaluate
\(\Delta_{100}=\max_{1\le j\le 99}(j(C-1)\bmod 100)\) and then
\(N_{\min}\).

\begin{longtable}[]{@{}
  >{\raggedright\arraybackslash}p{(\linewidth - 8\tabcolsep) * \real{0.0400}}
  >{\raggedleft\arraybackslash}p{(\linewidth - 8\tabcolsep) * \real{0.0500}}
  >{\raggedright\arraybackslash}p{(\linewidth - 8\tabcolsep) * \real{0.0800}}
  >{\raggedleft\arraybackslash}p{(\linewidth - 8\tabcolsep) * \real{0.3200}}
  >{\raggedleft\arraybackslash}p{(\linewidth - 8\tabcolsep) * \real{0.5100}}@{}}
\toprule\noalign{}
\begin{minipage}[b]{\linewidth}\raggedright
Case
\end{minipage} & \begin{minipage}[b]{\linewidth}\raggedleft
\(C\)
\end{minipage} & \begin{minipage}[b]{\linewidth}\raggedright
Regime
\end{minipage} & \begin{minipage}[b]{\linewidth}\raggedleft
\(\Delta_{100}\)
\end{minipage} & \begin{minipage}[b]{\linewidth}\raggedleft
\(N_{\min}\)
\end{minipage} \\
\midrule\noalign{}
\endhead
\bottomrule\noalign{}
\endlastfoot
(a) & 0.5 & \(C\le L\) & --- &
\(\left\lceil 1+\frac{0.5+0.5}{1}\right\rceil=2\) \\
(b) & 1.0 & \(C=L\) & \(0\) &
\(\left\lceil 1+\frac{0.5+0}{1}\right\rceil=2\) \\
(c) & 1.01 & \(C\ge L\) & \(\max_{j\le99}0.01j=0.99\) &
\(\left\lceil 1+\frac{0.5+0.99}{1.01}\right\rceil=3\) \\
(d) & 1.5 & \(C\ge L\) & \(\max_{j\le99}0.5j=49.5\) &
\(\left\lceil 1+\frac{0.5+49.5}{1.5}\right\rceil=35\) \\
(e) & 2.0 & \(C\ge L\) & \(\max_{j\le99}1\cdot j=99\) &
\(\left\lceil 1+\frac{0.5+99}{2}\right\rceil=51\) \\
(f) & 2.5 & \(C\ge L\) & \(\max_{j\le99}(1.5j\bmod100)=99\) &
\(\left\lceil 1+\frac{0.5+99}{2.5}\right\rceil=41\) \\
(g) & 10.0 & \(C\ge L\) & \(\max_{j\le99}(9j\bmod100)=99\) &
\(\left\lceil 1+\frac{0.5+99}{10}\right\rceil=11\) \\
(h) & 50.0 & \(C\ge L\) & \(\max_{j\le99}(49j\bmod100)=99\) &
\(\left\lceil 1+\frac{0.5+99}{50}\right\rceil=3\) \\
(i) & 200.0 & \(C\ge L\) & \(\max_{j\le99}(199j\bmod100)=99\) &
\(\left\lceil 1+\frac{0.5+99}{200}\right\rceil=2\) \\
\end{longtable}

These values are the least buffer counts that allow the optimal makespan
\(v_n^{\min}=T+C+(n-1)\max(L,C)\) to be realized under the exercise's
constraints.




\end{document}
